\IfFileExists{stacks-project-book.cls}{%
\documentclass{stacks-project-book}
}{%
\documentclass{amsbook}
}

\usepackage{verbatim}
\newenvironment{reference}{\comment}{\endcomment}
\newenvironment{slogan}{\comment}{\endcomment}
\newenvironment{history}{\comment}{\endcomment}

\usepackage[all]{xy}

\xyoption{2cell}
\UseAllTwocells

\usepackage{multicol}

\usepackage{lmodern}
\usepackage[T1]{fontenc}
\usepackage[french]{babel}
\AtBeginDocument{\shorthandoff{?}}


\usepackage{hyperref}


\theoremstyle{plain}
\newtheorem{theorem}[subsection]{Théorème}
\newtheorem{proposition}[subsection]{Proposition}
\newtheorem{lemma}[subsection]{Lemme}

\theoremstyle{definition}
\newtheorem{definition}[subsection]{Définition}
\newtheorem{example}[subsection]{Exemple}
\newtheorem{exercise}[subsection]{Exercice}
\newtheorem{situation}[subsection]{Situation}

\theoremstyle{remark}
\newtheorem{remark}[subsection]{Remarque}
\newtheorem{remarks}[subsection]{Remarques}

\renewcommand{\proofname}{Démonstration}

\numberwithin{equation}{subsection}

\def\lim{\mathop{\mathrm{lim}}\nolimits}
\def\colim{\mathop{\mathrm{colim}}\nolimits}
\def\Spec{\mathop{\mathrm{Spec}}}
\def\Hom{\mathop{\mathrm{Hom}}\nolimits}
\def\Ext{\mathop{\mathrm{Ext}}\nolimits}
\def\SheafHom{\mathop{\mathcal{H}\!\mathit{om}}\nolimits}
\def\SheafExt{\mathop{\mathcal{E}\!\mathit{xt}}\nolimits}
\def\Sch{\mathit{Sch}}
\def\Mor{\mathop{\mathrm{Mor}}\nolimits}
\def\Ob{\mathop{\mathrm{Ob}}\nolimits}
\def\Sh{\mathop{\mathit{Sh}}\nolimits}
\def\NL{\mathop{N\!L}\nolimits}
\def\CH{\mathop{\mathrm{CH}}\nolimits}
\def\proetale{{pro\text{-}\acute{e}tale}}
\def\etale{{\acute{e}tale}}
\def\QCoh{\mathit{QCoh}}
\def\Ker{\mathop{\mathrm{Ker}}}
\def\Im{\mathop{\mathrm{Im}}}
\def\Coker{\mathop{\mathrm{Coker}}}
\def\Coim{\mathop{\mathrm{Coim}}}

\DeclareMathSymbol{\boxtimes}{\mathbin}{AMSa}{"02}

\def\QCohstack{\mathcal{QC}\!\mathit{oh}}
\def\Cohstack{\mathcal{C}\!\mathit{oh}}
\def\Spacesstack{\mathcal{S}\!\mathit{paces}}
\def\Quotfunctor{\mathrm{Quot}}
\def\Hilbfunctor{\mathrm{Hilb}}
\def\Curvesstack{\mathcal{C}\!\mathit{urves}}
\def\Polarizedstack{\mathcal{P}\!\mathit{olarized}}
\def\Complexesstack{\mathcal{C}\!\mathit{omplexes}}
\def\Pic{\mathop{\mathrm{Pic}}\nolimits}
\def\Picardstack{\mathcal{P}\!\mathit{ic}}
\def\Picardfunctor{\mathrm{Pic}}
\def\Deformationcategory{\mathcal{D}\!\mathit{ef}}
\begin{document}
\title{Le projet Stacks --- édition française, partie 11}
\maketitle
\tableofcontents

% OK, start here.
%

\chapter{Théorie des déformations}



\phantomsection
\label{defos-section-phantom}


\section{Introduction}
\label{defos-section-introduction}

\noindent
Le but de ce chapitre est de donner une introduction (relativement) élémentaire à
la théorie des déformations des modules, des morphismes, etc. Nous y traitons
les résultats que l'on peut démontrer au moyen du complexe cotangent naïf. Dans
le chapitre consacré au complexe cotangent, nous étendrons un peu ces résultats.
Le lecteur averti pourra consulter le traité d'Illusie sur ce
sujet ; voir \cite{cotangent}.





\section{Déformations d'anneaux et complexe cotangent naïf}
\label{defos-section-deformations}

\noindent
Dans cette section, nous utilisons le complexe cotangent naïf pour faire un peu
de théorie des déformations. Partons d'un homomorphisme surjectif d'anneaux $A' \to A$
dont le noyau est un idéal $I$ de carré nul. Supposons en outre
donnés un homomorphisme d'anneaux $A \to B$, un $B$-module $N$ et une application $A$-linéaire
$c : I \to N$. Nous cherchons dans cette section à déterminer
le point d'interrogation qui complète le diagramme suivant :
\begin{equation}
\label{defos-equation-to-solve}
\vcenter{
\xymatrix{
0 \ar[r] & N \ar[r] & {?} \ar[r] & B \ar[r] & 0 \\
0 \ar[r] & I \ar[u]^c \ar[r] & A' \ar[u] \ar[r] & A \ar[u] \ar[r] & 0
}
}
\end{equation}
et à mesurer en outre l'unicité de la solution (si elle existe). Plus précisément,
nous cherchons une surjection de $A'$-algèbres $B' \to B$ dont le noyau soit
un idéal de carré nul identifié à
$N$, telle que $A' \to B'$ induise l'application donnée $c$.
Nous dirons que $B'$ est une {\it solution} de (\ref{defos-equation-to-solve}).

\begin{lemma}
\label{defos-lemma-huge-diagram}
Étant donné un diagramme commutatif
$$
\xymatrix{
& 0 \ar[r] & N_2 \ar[r] & B'_2 \ar[r] & B_2 \ar[r] & 0 \\
& 0 \ar[r]|\hole & I_2 \ar[u]_{c_2} \ar[r] &
A'_2 \ar[u] \ar[r]|\hole & A_2 \ar[u] \ar[r] & 0 \\
0 \ar[r] & N_1 \ar[ruu] \ar[r] & B'_1 \ar[r] & B_1 \ar[ruu] \ar[r] & 0 \\
0 \ar[r] & I_1 \ar[ruu]|\hole \ar[u]^{c_1} \ar[r] &
A'_1 \ar[ruu]|\hole \ar[u] \ar[r] & A_1 \ar[ruu]|\hole \ar[u] \ar[r] & 0
}
$$
dont les faces avant et arrière sont des solutions de (\ref{defos-equation-to-solve}), on a :
\begin{enumerate}
\item Il existe un élément canonique de
$\Ext^1_{B_1}(\NL_{B_1/A_1}, N_2)$
dont l'annulation est une condition nécessaire et suffisante pour qu'il existe
un homomorphisme d'anneaux $B'_1 \to B'_2$ complétant le diagramme.
\item S'il existe un homomorphisme $B'_1 \to B'_2$ complétant le diagramme,
l'ensemble de tous ces homomorphismes est un espace principal homogène sous
$\Hom_{B_1}(\Omega_{B_1/A_1}, N_2)$.
\end{enumerate}
\end{lemma}

\begin{proof}
Soit $E = B_1$ considéré comme un ensemble.
Considérons la surjection $A_1[E] \to B_1$, de noyau $J$, utilisée
pour définir le complexe cotangent naïf par la formule
$$
\NL_{B_1/A_1} = (J/J^2 \to \Omega_{A_1[E]/A_1} \otimes_{A_1[E]} B_1)
$$
dans
Algèbre, section \ref{algebra-section-netherlander}.
Puisque $\Omega_{A_1[E]/A_1} \otimes B_1$ est un
$B_1$-module libre, on a
$$
\Ext^1_{B_1}(\NL_{B_1/A_1}, N_2) =
\frac{\Hom_{B_1}(J/J^2, N_2)}
{\Hom_{B_1}(\Omega_{A_1[E]/A_1} \otimes B_1, N_2)}
$$
Nous allons construire une obstruction dans le module de droite.
Soit $J' = \Ker(A'_1[E] \to B_1)$. Remarquons qu'il existe une surjection
$J' \to J$ dont le noyau est $I_1A'_1[E]$.
Pour tout $e \in E$, notons $x_e \in A_1[E]$ la variable correspondante.
Choisissons un relèvement $y_e \in B'_1$ de l'image de $x_e$ dans $B_1$ et
un relèvement $z_e \in B'_2$ de l'image de $x_e$ dans $B_2$.
Ces choix déterminent des homomorphismes de $A'_1$-algèbres
$$
A'_1[E] \to B'_1 \quad\text{et}\quad A'_1[E] \to B'_2
$$
Le premier donne une application $J' \to N_1$, $f' \mapsto f'(y_e)$,
et le second une application $J' \to N_2$, $f' \mapsto f'(z_e)$.
Un calcul montre que ces applications annulent $(J')^2$.
Comme le carré gauche du diagramme (où interviennent $c_1$ et $c_2$)
est commutatif, ces applications coïncident sur $I_1A'_1[E]$ comme applications à valeurs dans $N_2$.
Remarquons que $B'_1$ est la somme amalgamée de $J' \to A'_1[E]$ et $J' \to N_1$.
Ainsi, si les applications $J' \to N_1 \to N_2$ et $J' \to N_2$ coïncident, on
obtient un homomorphisme $B'_1 \to B'_2$ complétant le diagramme.
On définit donc l'obstruction comme la classe de l'application
$$
J/J^2 \to N_2,\quad f \mapsto f'(z_e) - \nu(f'(y_e))
$$
où $\nu : N_1 \to N_2$ est l'application donnée et où $f' \in J'$
est un relèvement de $f$. Les remarques précédentes montrent que cette application est bien définie.
Nous sommes libres
de remplacer nos choix de $z_e$ par $z_e + \delta_{2, e}$
et de $y_e$ par $y_e + \delta_{1, e}$, pour certains $\delta_{i, e} \in N_i$.
Cela remplace l'application précédente par
$$
f \mapsto f'(z_e + \delta_{2, e}) - \nu(f'(y_e + \delta_{1, e})) =
f'(z_e) - \nu(f'(z_e)) +
\sum (\delta_{2, e} - \nu(\delta_{1, e}))\frac{\partial f}{\partial x_e}
$$
Cela signifie exactement que l'on modifie l'application $J/J^2 \to N_2$
par la composée $J/J^2 \to \Omega_{A_1[E]/A_1} \otimes B_1 \to N_2$,
où la seconde application envoie $\text{d}x_e$ sur
$\delta_{2, e} - \nu(\delta_{1, e})$. Notre obstruction est donc bien définie
et elle est nulle si et seulement s'il existe un relèvement.

\medskip\noindent
L'assertion (2) résulte de l'observation suivante : étant donnés deux homomorphismes
$\varphi, \psi : B'_1 \to B'_2$ complétant le diagramme,
$\varphi - \psi$ se factorise par une application $D : B_1 \to N_2$ qui
est une $A_1$-dérivation :
\begin{align*}
D(fg) & = \varphi(f'g') - \psi(f'g') \\
& =
\varphi(f')\varphi(g') - \psi(f')\psi(g') \\
& =
(\varphi(f') - \psi(f'))\varphi(g') + \psi(f')(\varphi(g') - \psi(g')) \\
& =
gD(f) + fD(g)
\end{align*}
Ainsi, $D$ correspond à une unique application $B_1$-linéaire
$\Omega_{B_1/A_1} \to N_2$. Réciproquement, une telle application linéaire
donne une dérivation $D$ et, étant donné un homomorphisme d'anneaux $\psi : B'_1 \to B'_2$
complétant le diagramme,
l'application $\psi + D$ est un autre homomorphisme d'anneaux qui le complète.
\end{proof}

\begin{lemma}
\label{defos-lemma-choices}
S'il existe une solution de (\ref{defos-equation-to-solve}), l'ensemble des
classes d'isomorphisme de solutions est un espace principal homogène sous
$\Ext^1_B(\NL_{B/A}, N)$.
\end{lemma}

\begin{proof}
Remarquons d'emblée que, pour deux solutions $B'_1$ et $B'_2$
de (\ref{defos-equation-to-solve}), le lemme \ref{defos-lemma-huge-diagram} fournit un
élément d'obstruction $o(B'_1, B'_2) \in \Ext^1_B(\NL_{B/A}, N)$
à l'existence d'un homomorphisme $B'_1 \to B'_2$. Cet élément est clairement
l'obstruction à l'existence d'un isomorphisme et distingue donc
les classes d'isomorphisme. Pour achever la démonstration, il suffit alors de
montrer qu'étant donnés une solution $B'$ et un élément
$\xi \in \Ext^1_B(\NL_{B/A}, N)$,
on peut trouver une seconde solution $B'_\xi$ telle que
$o(B', B'_\xi) = \xi$.

\medskip\noindent
Soit $E = B$ considéré comme un ensemble. Considérons la surjection $A[E] \to B$ de noyau
$J$ utilisée pour définir le complexe cotangent naïf par la formule
$$
\NL_{B/A} = (J/J^2 \to \Omega_{A[E]/A} \otimes_{A[E]} B)
$$
dans Algèbre, section \ref{algebra-section-netherlander}.
Puisque $\Omega_{A[E]/A} \otimes B$ est un $B$-module libre, on a
$$
\Ext^1_B(\NL_{B/A}, N) =
\frac{\Hom_B(J/J^2, N)}
{\Hom_B(\Omega_{A[E]/A} \otimes B, N)}
$$
On peut donc représenter $\xi$ par la classe d'un morphisme $\delta : J/J^2 \to N$.

\medskip\noindent
Pour tout $e \in E$, notons $x_e \in A[E]$ la variable correspondante.
Choisissons un relèvement $y_e \in B'$ de l'image de $x_e$ dans $B$.
Ces choix déterminent un homomorphisme de $A'$-algèbres $\varphi : A'[E] \to B'$.
Soit $J' = \Ker(A'[E] \to B)$. Remarquons que $\varphi$ induit une application
$\varphi|_{J'} : J' \to N$ et que $B'$ est la somme amalgamée, comme dans le
diagramme suivant :
$$
\xymatrix{
0 \ar[r] & N \ar[r] & B' \ar[r] & B \ar[r] & 0 \\
0 \ar[r] & J' \ar[u]^{\varphi|_{J'}} \ar[r] & A'[E] \ar[u] \ar[r] &
B \ar[u]_{=} \ar[r] & 0
}
$$
Soit $\psi : J' \to N$ la somme de l'application $\varphi|_{J'}$ et de la
composée
$$
J' \to J'/(J')^2 \to J/J^2 \xrightarrow{\delta} N.
$$
La somme amalgamée suivant $\psi$ est alors une autre extension d'anneaux $B'_\xi$
qui s'insère dans un diagramme comme ci-dessus. Un calcul montre que
$o(B', B'_\xi) = \xi$, comme voulu.
\end{proof}

\begin{lemma}
\label{defos-lemma-extensions-of-algebras}
Soit $A$ un anneau. Soit $B$ une $A$-algèbre. Soit $N$ un $B$-module.
L'ensemble des classes d'isomorphisme d'extensions de $A$-algèbres
$$
0 \to N \to B' \to B \to 0
$$
où $N$ est un idéal de carré nul est canoniquement en bijection avec
$\Ext^1_B(\NL_{B/A}, N)$.
\end{lemma}

\begin{proof}
Pour le démontrer, appliquons les résultats précédents au cas où
(\ref{defos-equation-to-solve}) est donné par le diagramme
$$
\xymatrix{
0 \ar[r] & N \ar[r] & {?} \ar[r] & B \ar[r] & 0 \\
0 \ar[r] & 0 \ar[u] \ar[r] & A \ar[u] \ar[r]^{\text{id}} & A \ar[u] \ar[r] & 0
}
$$
Le lemme résulte donc du lemme \ref{defos-lemma-choices}
et du fait qu'il existe une solution, à savoir $N \oplus B$.
(Voir la remarque ci-dessous pour une construction directe de la bijection.)
\end{proof}

\begin{remark}
\label{defos-remark-extensions-of-algebras}
Soient $A \to B$ et $N$ comme dans le lemme \ref{defos-lemma-extensions-of-algebras}.
Soit $\alpha : P \to B$ une présentation de $B$ sur $A$ ; voir
Algèbre, section \ref{algebra-section-netherlander}.
Pour $J = \Ker(\alpha)$, le complexe cotangent naïf $\NL(\alpha)$
associé à $\alpha$ est le complexe $J/J^2 \to \Omega_{P/A} \otimes_P B$.
On a
$$
\Ext^1_B(\NL(\alpha), N) =
\Coker\left(\Hom_B(\Omega_{P/A} \otimes_P B, N) \to \Hom_B(J/J^2, N)\right)
$$
car $\Omega_{P/A}$ est un module libre.
Considérons une extension $0 \to N \to B' \to B \to 0$ comme dans le lemme.
Puisque $P$ est une algèbre polynomiale sur $A$, on peut
relever $\alpha$ en un homomorphisme de $A$-algèbres $\alpha' : P' \to B'$.
Alors $\alpha'|_J : J \to N$ se factorise sous la forme $J \to J/J^2 \to N$
car $N$ est de carré nul dans $B'$. Le lemme envoie notre extension
sur la classe de cette application $J/J^2 \to N$ dans le conoyau affiché.
\end{remark}












\begin{lemma}
\label{defos-lemma-extensions-of-algebras-functorial}
Soient des homomorphismes d'anneaux $A \to B \to C$, un $B$-module $M$, un $C$-module $N$,
une application $B$-linéaire $c : M \to N$, et des extensions de
$A$-algèbres à noyaux de carré nul
\begin{enumerate}
\item[(a)] $0 \to M \to B' \to B \to 0$ correspondant à
$\xi \in \Ext^1_B(\NL_{B/A}, M)$, et
\item[(b)] $0 \to N \to C' \to C \to 0$ correspondant à
$\zeta \in \Ext^1_C(\NL_{C/A}, N)$.
\end{enumerate}
Voir le lemme \ref{defos-lemma-extensions-of-algebras}.
Il existe alors un homomorphisme de $A$-algèbres $B' \to C'$ compatible avec
$B \to C$ et $c$ si et seulement si $\xi$ et $\zeta$
s'envoient sur le même élément de
$\Ext^1_B(\NL_{B/A}, N)$.
\end{lemma}

\begin{proof}
L'énoncé a un sens car on dispose des applications
$$
\Ext^1_B(\NL_{B/A}, M) \to \Ext^1_B(\NL_{B/A}, N)
$$
induites par l'application $M \to N$, ainsi que de
$$
\Ext^1_C(\NL_{C/A}, N) \to \Ext^1_B(\NL_{C/A}, N) \to \Ext^1_B(\NL_{B/A}, N)
$$
où la première flèche utilise le foncteur de restriction $D(C) \to D(B)$
et la seconde l'application canonique de complexes
$\NL_{B/A} \to \NL_{C/A}$. L'énoncé du lemme se déduit du
lemme \ref{defos-lemma-huge-diagram} appliqué au diagramme
$$
\xymatrix{
& 0 \ar[r] & N \ar[r] & C' \ar[r] & C \ar[r] & 0 \\
& 0 \ar[r]|\hole & 0 \ar[u] \ar[r] &
A \ar[u] \ar[r]|\hole & A \ar[u] \ar[r] & 0 \\
0 \ar[r] & M \ar[ruu] \ar[r] & B' \ar[r] & B \ar[ruu] \ar[r] & 0 \\
0 \ar[r] & 0 \ar[ruu]|\hole \ar[u] \ar[r] &
A \ar[ruu]|\hole \ar[u] \ar[r] & A \ar[ruu]|\hole \ar[u] \ar[r] & 0
}
$$
et d'une compatibilité entre les constructions des démonstrations
des lemmes \ref{defos-lemma-extensions-of-algebras} et \ref{defos-lemma-huge-diagram},
dont nous omettons l'énoncé et la démonstration. (Voir la remarque ci-dessous pour un argument direct.)
\end{proof}

\begin{remark}
\label{defos-remark-extensions-of-algebras-functorial}
Soient $A \to B \to C$, $M$, $N$, $c : M \to N$,
$0 \to M \to B' \to B \to 0$, $\xi \in \Ext^1_B(\NL_{B/A}, M)$,
$0 \to N \to C' \to C \to 0$ et $\zeta \in \Ext^1_C(\NL_{C/A}, N)$ comme dans le
lemme \ref{defos-lemma-extensions-of-algebras-functorial}.
En prenant la somme amalgamée le long de $c : M \to N$, on construit une extension
$$
\xymatrix{
0 \ar[r] & N \ar[r] & B'_1 \ar[r] & B \ar[r] & 0 \\
0 \ar[r] & M \ar[u]^c \ar[r] & B' \ar[u] \ar[r] & B \ar[u] \ar[r] & 0
}
$$
en posant $B'_1 = (N \times B')/M$, où $M$ est plongé
antidiagonalement. En prenant le produit fibré le long de $B \to C$, on construit une extension
$$
\xymatrix{
0 \ar[r] & N \ar[r] & C' \ar[r] & C \ar[r] & 0 \\
0 \ar[r] & N \ar[u] \ar[r] & B'_2 \ar[u] \ar[r] & B \ar[u] \ar[r] & 0
}
$$
en posant $B'_2 = C' \times_C B$ (produit fibré d'anneaux). Une chasse au diagramme élémentaire
montre qu'il existe un homomorphisme de $A$-algèbres $B' \to C'$
compatible avec $B \to C$ et $c$ si et seulement si $B'_1$ est isomorphe
à $B'_2$ comme extension de $A$-algèbres de $B$ par $N$. Pour vérifier
le lemme \ref{defos-lemma-extensions-of-algebras-functorial},
il suffit donc de montrer que $B'_1$ correspond, par la bijection du
lemme \ref{defos-lemma-extensions-of-algebras},
à l'image de $\xi$ par l'application
$\Ext^1_B(\NL_{B/A}, M) \to \Ext^1_B(\NL_{B/A}, N)$,
et que $B'_2$ correspond à l'image de $\zeta$ par l'application
$\Ext^1_C(\NL_{C/A}, N) \to \Ext^1_B(\NL_{B/A}, N)$.
La première de ces assertions résulte immédiatement de la construction
de la classe dans la remarque \ref{defos-remark-extensions-of-algebras}.
Pour la seconde, choisissons un diagramme commutatif
$$
\xymatrix{
Q \ar[r]_\beta & C \\
P \ar[u]^\varphi \ar[r]^\alpha & B \ar[u]
}
$$
de $A$-algèbres, tel que $\alpha$ soit une présentation de $B$ sur $A$
et $\beta$ une présentation de $C$ sur $A$. Voir la
remarque \ref{defos-remark-extensions-of-algebras} et les références qui y sont données.
Posons $J = \Ker(\alpha)$ et $K = \Ker(\beta)$. L'application $\varphi$
induit une application de complexes $\NL(\alpha) \to \NL(\beta)$
et, en particulier, $\bar\varphi : J/J^2 \to K/K^2$.
Choisissons un homomorphisme de $A$-algèbres $\beta' : Q \to C'$
qui relève $\beta$. Alors
$\alpha' = (\beta' \circ \varphi, \alpha) : P \to B'_2 = C' \times_C B$
est un relèvement de $\alpha$. Pour ces choix, la composée de l'application
$K/K^2 \to N$ induite par $\beta'$ et de l'application $\bar\varphi : J/J^2 \to K/K^2$
est la restriction de $\alpha'$ à $J/J^2$.
En explicitant les constructions de nos classes dans la
remarque \ref{defos-remark-extensions-of-algebras},
on voit bien que
$B'_2$ correspond à l'image de $\zeta$ par l'application
$\Ext^1_C(\NL_{C/A}, N) \to \Ext^1_B(\NL_{B/A}, N)$.
\end{remark}

\begin{lemma}
\label{defos-lemma-parametrize-solutions}
Soient $0 \to I \to A' \to A \to 0$, $A \to B$ et $c : I \to N$ comme dans
(\ref{defos-equation-to-solve}). Notons $\xi \in \Ext^1_A(\NL_{A/A'}, I)$
l'élément correspondant à l'extension $A'$ de $A$ par $I$ par le
lemme \ref{defos-lemma-extensions-of-algebras}. L'ensemble des classes d'isomorphisme
de solutions est canoniquement en bijection avec la fibre de
$$
\Ext^1_B(\NL_{B/A'}, N) \to \Ext^1_A(\NL_{A/A'}, N)
$$
au-dessus de l'image de $\xi$.
\end{lemma}

\begin{proof}
D'après le lemme \ref{defos-lemma-extensions-of-algebras} appliqué à $A' \to B$ et au
$B$-module $N$, les éléments $\zeta$ de $\Ext^1_B(\NL_{B/A'}, N)$
paramètrent les extensions $0 \to N \to B' \to B \to 0$ de $A'$-algèbres.
D'après le lemme \ref{defos-lemma-extensions-of-algebras-functorial} appliqué
à $A' \to A \to B$ et $c : I \to N$, il existe un homomorphisme de $A'$-algèbres
$A' \to B'$ compatible avec $c$ et $A \to B$ si et seulement si
$\zeta$ s'envoie sur $\xi$. Cela revient bien sûr à dire que $B'$ est une
solution de (\ref{defos-equation-to-solve}).
\end{proof}

\begin{remark}
\label{defos-remark-parametrize-solutions}
Remarquons que, dans la situation du lemme \ref{defos-lemma-parametrize-solutions},
on a
$$
\Ext^1_A(\NL_{A/A'}, N) =
\Ext^1_B(\NL_{A/A'} \otimes_A^\mathbf{L} B, N) =
\Ext^1_B(\NL_{A/A'} \otimes_A B, N)
$$
La première égalité résulte de
Compléments d'algèbre, lemme \ref{more-algebra-lemma-tensor-hom-adjoint}, et
la seconde de
Compléments d'algèbre, lemme \ref{more-algebra-lemma-tensor-NL}.
On dispose d'applications de complexes
$$
\NL_{A/A'} \otimes_A B \to \NL_{B/A'} \to \NL_{B/A}
$$
qui forment presque un triangle distingué ; voir
Algèbre, lemme \ref{algebra-lemma-exact-sequence-NL}.
Si elles formaient un triangle distingué, on conclurait
que l'image de $\xi$ dans $\Ext^2_B(\NL_{B/A}, N)$
serait l'obstruction à l'existence d'une solution de
(\ref{defos-equation-to-solve}).
\end{remark}

\noindent
Si notre homomorphisme d'anneaux $A \to B$ est d'intersection complète locale, il existe
une solution. Il s'agit d'un résultat de relèvement ; remarquons que,
pour les homomorphismes syntomiques, nous avons démontré un résultat de relèvement assez fort dans
Lissage des morphismes d'anneaux, Proposition \ref{smoothing-proposition-lift-smooth}.

\begin{lemma}
\label{defos-lemma-existence-lci}
Si $A \to B$ est un homomorphisme d'anneaux d'intersection complète locale, alors
il existe une solution de (\ref{defos-equation-to-solve}).
\end{lemma}

\begin{proof}[Première démonstration]
Écrivons $B = A[x_1, \ldots, x_n]/J$. D'après Compléments d'algèbre, définition
\ref{more-algebra-definition-local-complete-intersection},
l'idéal $J$ est Koszul-régulier. Il s'ensuit que $J$ est $H_1$-régulier et
quasi-régulier ; voir Compléments d'algèbre, section \ref{more-algebra-section-ideals}.
Soit $J' \subset A'[x_1, \ldots, x_n]$
l'image réciproque de $J$. Notons $I[x_1, \ldots, x_n]$ le
noyau de $A'[x_1, \ldots, x_n] \to A[x_1, \ldots, x_n]$.
D'après Compléments d'algèbre, lemme
\ref{more-algebra-lemma-conormal-sequence-H1-regular-ideal}, on a
$I[x_1, \ldots, x_n] \cap (J')^2 = J'I[x_1, \ldots, x_n] =
JI[x_1, \ldots, x_n]$. On obtient donc une suite exacte courte
$$
0 \to I \otimes_A B \to J'/(J')^2 \to J/J^2 \to 0
$$
Puisque $J/J^2$ est projectif (Compléments d'algèbre, lemme
\ref{more-algebra-lemma-quasi-regular-ideal-finite-projective}),
on peut choisir un scindage de cette suite
$$
J'/(J')^2 = I \otimes_A B \oplus J/J^2 
$$
Soit $(J')^2 \subset J'' \subset J'$ l'ensemble des éléments qui s'envoient dans le
second facteur de la décomposition précédente. Alors
$$
0 \to I \otimes_A B \to A'[x_1, \ldots, x_n]/J'' \to B \to 0
$$
est une solution de (\ref{defos-equation-to-solve}) avec $N = I \otimes_A B$.
Le cas général s'obtient en prenant la somme amalgamée le long de l'application donnée
$I \otimes_A B \to N$.
\end{proof}

\begin{proof}[Seconde démonstration]
Veuillez lire la remarque \ref{defos-remark-parametrize-solutions}
avant cette démonstration. D'après
Compléments d'algèbre, lemme \ref{more-algebra-lemma-transitive-lci-at-end},
les applications $\NL_{A/A'} \otimes_A B \to \NL_{B/A'} \to \NL_{B/A}$
forment bien un triangle distingué dans $D(B)$.
Il suffit donc de montrer que $\Ext^2_B(\NL_{B/A}, N)$ s'annule.
D'après Compléments d'algèbre, lemme \ref{more-algebra-lemma-lci-NL},
le complexe $\NL_{B/A}$ est parfait, d'amplitude de Tor contenue dans
$[-1, 0]$. Il s'ensuit que notre $\Ext^2$ s'annule, par exemple
d'après Compléments d'algèbre, lemme \ref{more-algebra-lemma-splitting-unique}, partie (1).
\end{proof}






\section{Épaississements d'espaces annelés}
\label{defos-section-thickenings-spaces}

\noindent
Dans les quelques sections qui suivent, nous utiliserons les notions suivantes :
\begin{enumerate}
\item Un faisceau d'idéaux $\mathcal{I} \subset \mathcal{O}_{X'}$ sur
un espace annelé $(X', \mathcal{O}_{X'})$ est {\it localement nilpotent}
si toute section locale de $\mathcal{I}$ est localement nilpotente.
Comparer avec Algèbre, Item \ref{algebra-item-ideal-locally-nilpotent}.
\item Un {\it épaississement} d'espaces annelés est un morphisme
$i : (X, \mathcal{O}_X) \to (X', \mathcal{O}_{X'})$ d'espaces annelés
tel que
\begin{enumerate}
\item $i$ induise un homéomorphisme $X \to X'$ ;
\item l'application $i^\sharp : \mathcal{O}_{X'} \to i_*\mathcal{O}_X$
soit surjective ; et
\item le noyau de $i^\sharp$ soit un faisceau d'idéaux localement nilpotent.
\end{enumerate}
\item Un {\it épaississement du premier ordre} d'espaces annelés est un épaississement
$i : (X, \mathcal{O}_X) \to (X', \mathcal{O}_{X'})$ d'espaces annelés
tel que $\Ker(i^\sharp)$ soit de carré nul.
\item Il est clair comment définir les {\it morphismes d'épaississements},
les {\it morphismes d'épaississements au-dessus d'un espace annelé de base}, etc.
\end{enumerate}
Si $i : (X, \mathcal{O}_X) \to (X', \mathcal{O}_{X'})$ est un épaississement
d'espaces annelés, nous identifions les espaces topologiques sous-jacents
et considérons $\mathcal{O}_X$, $\mathcal{O}_{X'}$ et
$\mathcal{I} = \Ker(i^\sharp)$ comme des faisceaux sur $X = X'$. On obtient
une suite exacte courte
$$
0 \to \mathcal{I} \to \mathcal{O}_{X'} \to \mathcal{O}_X \to 0
$$
de $\mathcal{O}_{X'}$-modules. D'après
Modules, lemme \ref{modules-lemma-i-star-equivalence},
la catégorie des $\mathcal{O}_X$-modules est équivalente à la catégorie
des $\mathcal{O}_{X'}$-modules annulés par $\mathcal{I}$. En particulier,
si $i$ est un épaississement du premier ordre,
$\mathcal{I}$ est un $\mathcal{O}_X$-module.

\begin{situation}
\label{defos-situation-morphism-thickenings}
Un morphisme d'épaississements $(f, f')$ est donné par un diagramme commutatif
\begin{equation}
\label{defos-equation-morphism-thickenings}
\vcenter{
\xymatrix{
(X, \mathcal{O}_X) \ar[r]_i \ar[d]_f & (X', \mathcal{O}_{X'}) \ar[d]^{f'} \\
(S, \mathcal{O}_S) \ar[r]^t & (S', \mathcal{O}_{S'})
}
}
\end{equation}
d'espaces annelés, dont les flèches horizontales sont des épaississements. Dans cette
situation, on pose
$\mathcal{I} = \Ker(i^\sharp) \subset \mathcal{O}_{X'}$ et
$\mathcal{J} = \Ker(t^\sharp) \subset \mathcal{O}_{S'}$.
Puisque $f = f'$ sur les espaces topologiques sous-jacents, nous identifierons
les foncteurs images réciproques (topologiques) $f^{-1}$ et $(f')^{-1}$.
Remarquons que $(f')^\sharp : f^{-1}\mathcal{O}_{S'} \to \mathcal{O}_{X'}$
induit en particulier une application $f^{-1}\mathcal{J} \to \mathcal{I}$
et donc une application de $\mathcal{O}_{X'}$-modules
$$
(f')^*\mathcal{J} \longrightarrow \mathcal{I}
$$
Si $i$ et $t$ sont des épaississements du premier ordre, alors
$(f')^*\mathcal{J} = f^*\mathcal{J}$ et l'application précédente devient une
application $f^*\mathcal{J} \to \mathcal{I}$.
\end{situation}

\begin{definition}
\label{defos-definition-strict-morphism-thickenings}
Dans la situation \ref{defos-situation-morphism-thickenings}, on dit que $(f, f')$ est un
{\it morphisme strict d'épaississements}
si l'application $(f')^*\mathcal{J} \longrightarrow \mathcal{I}$ est surjective.
\end{definition}

\noindent
Le lemme suivant montre en particulier qu'un morphisme
$(f, f') : (X \subset X') \to (S \subset S')$
d'épaississements de schémas est strict si et seulement si $X = S \times_{S'} X'$.

\begin{lemma}
\label{defos-lemma-strict-morphism-thickenings}
Dans la situation \ref{defos-situation-morphism-thickenings}, le morphisme $(f, f')$
est un morphisme strict d'épaississements si et seulement si
(\ref{defos-equation-morphism-thickenings}) est cartésien dans la catégorie
des espaces annelés.
\end{lemma}

\begin{proof}
Omis.
\end{proof}




\section{Modules sur les épaississements du premier ordre d'espaces annelés}
\label{defos-section-modules-thickenings}

\noindent
Dans cette section, nous présentons quelques préliminaires à la théorie des déformations
des modules. Soit $i : (X, \mathcal{O}_X) \to (X', \mathcal{O}_{X'})$
un épaississement du premier ordre d'espaces annelés. Nous utiliserons librement la notation
introduite dans la section \ref{defos-section-thickenings-spaces} ; en particulier, nous
identifierons les espaces topologiques sous-jacents.
Dans cette section, nous considérons des suites exactes courtes
\begin{equation}
\label{defos-equation-extension}
0 \to \mathcal{K} \to \mathcal{F}' \to \mathcal{F} \to 0
\end{equation}
de $\mathcal{O}_{X'}$-modules, où $\mathcal{F}$, $\mathcal{K}$ sont des
$\mathcal{O}_X$-modules et $\mathcal{F}'$ un $\mathcal{O}_{X'}$-module.
Dans cette situation, on dispose d'une application canonique de $\mathcal{O}_X$-modules
$$
c_{\mathcal{F}'} :
\mathcal{I} \otimes_{\mathcal{O}_X} \mathcal{F}
\longrightarrow
\mathcal{K}
$$
où $\mathcal{I} = \Ker(i^\sharp)$.
En effet, pour des sections locales $f$ de $\mathcal{I}$ et $s$
de $\mathcal{F}$, on pose $c_{\mathcal{F}'}(f \otimes s) = fs'$,
où $s'$ est une section locale de $\mathcal{F}'$ relevant $s$.

\begin{lemma}
\label{defos-lemma-inf-map}
Soit $i : (X, \mathcal{O}_X) \to (X', \mathcal{O}_{X'})$
un épaississement du premier ordre d'espaces annelés. Supposons données des
extensions
$$
0 \to \mathcal{K} \to \mathcal{F}' \to \mathcal{F} \to 0
\quad\text{et}\quad
0 \to \mathcal{L} \to \mathcal{G}' \to \mathcal{G} \to 0
$$
comme dans (\ref{defos-equation-extension}),
ainsi que des applications $\varphi : \mathcal{F} \to \mathcal{G}$ et
$\psi : \mathcal{K} \to \mathcal{L}$.
\begin{enumerate}
\item S'il existe une application de $\mathcal{O}_{X'}$-modules
$\varphi' : \mathcal{F}' \to \mathcal{G}'$ compatible avec $\varphi$
et $\psi$, alors le diagramme
$$
\xymatrix{
\mathcal{I} \otimes_{\mathcal{O}_X} \mathcal{F}
\ar[r]_-{c_{\mathcal{F}'}} \ar[d]_{1 \otimes \varphi} &
\mathcal{K} \ar[d]^\psi \\
\mathcal{I} \otimes_{\mathcal{O}_X} \mathcal{G}
\ar[r]^-{c_{\mathcal{G}'}} &
\mathcal{L}
}
$$
est commutatif.
\item L'ensemble des applications de $\mathcal{O}_{X'}$-modules
$\varphi' : \mathcal{F}' \to \mathcal{G}'$ compatibles avec $\varphi$
et $\psi$ est, s'il est non vide, un espace principal homogène sous
$\Hom_{\mathcal{O}_X}(\mathcal{F}, \mathcal{L})$.
\end{enumerate}
\end{lemma}

\begin{proof}
L'assertion (1) résulte immédiatement de la description des applications.
Pour (2), si $\varphi'$ et $\varphi''$ sont deux applications
$\mathcal{F}' \to \mathcal{G}'$ compatibles avec $\varphi$
et $\psi$, alors $\varphi' - \varphi''$ se factorise sous la forme
$$
\mathcal{F}' \to \mathcal{F} \to \mathcal{L} \to \mathcal{G}'
$$
L'application du milieu provient d'un unique élément de
$\Hom_{\mathcal{O}_X}(\mathcal{F}, \mathcal{L})$, d'après
Modules, lemme \ref{modules-lemma-i-star-equivalence}.
Réciproquement, étant donné un élément $\alpha$ de ce groupe, on peut ajouter la
composée (affichée ci-dessus avec $\alpha$ au milieu)
à $\varphi'$. Nous omettons certains détails.
\end{proof}

\begin{lemma}
\label{defos-lemma-inf-obs-map}
Soit $i : (X, \mathcal{O}_X) \to (X', \mathcal{O}_{X'})$
un épaississement du premier ordre d'espaces annelés. Supposons données des
extensions
$$
0 \to \mathcal{K} \to \mathcal{F}' \to \mathcal{F} \to 0
\quad\text{et}\quad
0 \to \mathcal{L} \to \mathcal{G}' \to \mathcal{G} \to 0
$$
comme dans (\ref{defos-equation-extension}),
ainsi que des applications $\varphi : \mathcal{F} \to \mathcal{G}$ et
$\psi : \mathcal{K} \to \mathcal{L}$. Supposons le diagramme
$$
\xymatrix{
\mathcal{I} \otimes_{\mathcal{O}_X} \mathcal{F}
\ar[r]_-{c_{\mathcal{F}'}} \ar[d]_{1 \otimes \varphi} &
\mathcal{K} \ar[d]^\psi \\
\mathcal{I} \otimes_{\mathcal{O}_X} \mathcal{G}
\ar[r]^-{c_{\mathcal{G}'}} &
\mathcal{L}
}
$$
commutatif. Il existe alors un élément
$$
o(\varphi, \psi) \in
\Ext^1_{\mathcal{O}_X}(\mathcal{F}, \mathcal{L})
$$
dont l'annulation est une condition nécessaire et suffisante pour qu'il existe
une application $\varphi' : \mathcal{F}' \to \mathcal{G}'$ compatible avec
$\varphi$ et $\psi$.
\end{lemma}

\begin{proof}
On peut construire explicitement une extension
$$
0 \to \mathcal{L} \to \mathcal{H} \to \mathcal{F} \to 0
$$
en prenant pour $\mathcal{H}$ la cohomologie au milieu du complexe
$$
\mathcal{K}
\xrightarrow{1, - \psi}
\mathcal{F}' \oplus \mathcal{G}' \xrightarrow{\varphi, 1}
\mathcal{G}
$$
(avec les notations évidentes). Un calcul sur les sections locales,
utilisant l'hypothèse de commutativité du diagramme du lemme,
montre que $\mathcal{H}$ est annulé par $\mathcal{I}$. Par conséquent,
$\mathcal{H}$ définit une classe dans
$$
\Ext^1_{\mathcal{O}_X}(\mathcal{F}, \mathcal{L})
\subset
\Ext^1_{\mathcal{O}_{X'}}(\mathcal{F}, \mathcal{L})
$$
Enfin, la classe de $\mathcal{H}$ est la différence entre la somme amalgamée
de l'extension $\mathcal{F}'$ suivant $\psi$ et l'image réciproque
de l'extension $\mathcal{G}'$ suivant $\varphi$ (calculs omis).
L'annulation de la classe de $\mathcal{H}$ équivaut donc à
l'existence d'un diagramme commutatif
$$
\xymatrix{
0 \ar[r] &
\mathcal{K} \ar[r] \ar[d]_{\psi} &
\mathcal{F}' \ar[r] \ar[d]_{\varphi'} &
\mathcal{F} \ar[r] \ar[d]_\varphi & 0\\
0 \ar[r] &
\mathcal{L} \ar[r] &
\mathcal{G}' \ar[r] &
\mathcal{G} \ar[r] & 0
}
$$
comme voulu.
\end{proof}

\begin{lemma}
\label{defos-lemma-inf-ext}
Soit $i : (X, \mathcal{O}_X) \to (X', \mathcal{O}_{X'})$ un épaississement
du premier ordre d'espaces annelés.
Supposons donnés des $\mathcal{O}_X$-modules $\mathcal{F}$, $\mathcal{K}$
et une application $\mathcal{O}_X$-linéaire
$c : \mathcal{I} \otimes_{\mathcal{O}_X} \mathcal{F} \to \mathcal{K}$.
S'il existe une suite (\ref{defos-equation-extension}) telle que
$c_{\mathcal{F}'} = c$, alors l'ensemble des classes d'isomorphisme de ces
extensions est un espace principal homogène sous
$\Ext^1_{\mathcal{O}_X}(\mathcal{F}, \mathcal{K})$.
\end{lemma}

\begin{proof}
Supposons données des extensions
$$
0 \to \mathcal{K} \to \mathcal{F}'_1 \to \mathcal{F} \to 0
\quad\text{et}\quad
0 \to \mathcal{K} \to \mathcal{F}'_2 \to \mathcal{F} \to 0
$$
telles que $c_{\mathcal{F}'_1} = c_{\mathcal{F}'_2} = c$. Leur différence
(dans le groupe des extensions ; voir
Homology, section \ref{homology-section-extensions})
est une extension
$$
0 \to \mathcal{K} \to \mathcal{E} \to \mathcal{F} \to 0
$$
où $\mathcal{E}$ est annulé par $\mathcal{I}$ (on omet le calcul local).
La suite est donc une extension de $\mathcal{O}_X$-modules ;
voir Modules, lemme \ref{modules-lemma-i-star-equivalence}.
Réciproquement, étant donnée une telle extension $\mathcal{E}$, on peut ajouter l'extension
$\mathcal{E}$ à l'extension de $\mathcal{O}_{X'}$-modules $\mathcal{F}'$ sans
modifier l'application $c_{\mathcal{F}'}$. Nous omettons certains détails.
\end{proof}

\begin{lemma}
\label{defos-lemma-inf-obs-ext}
Soit $i : (X, \mathcal{O}_X) \to (X', \mathcal{O}_{X'})$
un épaississement du premier ordre d'espaces annelés. Supposons donnés des
$\mathcal{O}_X$-modules $\mathcal{F}$, $\mathcal{K}$
et une application $\mathcal{O}_X$-linéaire
$c : \mathcal{I} \otimes_{\mathcal{O}_X} \mathcal{F} \to \mathcal{K}$.
Il existe alors un élément
$$
o(\mathcal{F}, \mathcal{K}, c) \in
\Ext^2_{\mathcal{O}_X}(\mathcal{F}, \mathcal{K})
$$
dont l'annulation est une condition nécessaire et suffisante pour qu'il existe
une suite (\ref{defos-equation-extension}) telle que $c_{\mathcal{F}'} = c$.
\end{lemma}

\begin{proof}
Montrons d'abord que, si $\mathcal{K}$ est un $\mathcal{O}_X$-module injectif,
il existe bien une suite (\ref{defos-equation-extension}) telle que
$c_{\mathcal{F}'} = c$. Pour cela, choisissons un
$\mathcal{O}_{X'}$-module plat $\mathcal{H}'$ et une surjection
$\mathcal{H}' \to \mathcal{F}$
(Modules, lemme \ref{modules-lemma-module-quotient-flat}).
Soit $\mathcal{J} \subset \mathcal{H}'$ le noyau. Comme $\mathcal{H}'$
est plat, on a
$$
\mathcal{I} \otimes_{\mathcal{O}_{X'}} \mathcal{H}' =
\mathcal{I}\mathcal{H}'
\subset \mathcal{J} \subset \mathcal{H}'
$$
Remarquons que l'application
$$
\mathcal{I}\mathcal{H}' =
\mathcal{I} \otimes_{\mathcal{O}_{X'}} \mathcal{H}'
\longrightarrow
\mathcal{I} \otimes_{\mathcal{O}_{X'}} \mathcal{F} =
\mathcal{I} \otimes_{\mathcal{O}_X} \mathcal{F}
$$
annule $\mathcal{I}\mathcal{J}$. En effet, si $f$ est une section locale
de $\mathcal{I}$ et $s$ une section locale de $\mathcal{H}$, alors
$fs$ s'envoie sur $f \otimes \overline{s}$, où $\overline{s}$ est
l'image de $s$ dans $\mathcal{F}$. On obtient ainsi
$$
\xymatrix{
\mathcal{I}\mathcal{H}'/\mathcal{I}\mathcal{J}
\ar@{^{(}->}[r] \ar[d] &
\mathcal{J}/\mathcal{I}\mathcal{J} \ar@{..>}[d]_\gamma \\
\mathcal{I} \otimes_{\mathcal{O}_X} \mathcal{F} \ar[r]^-c &
\mathcal{K}
}
$$
un diagramme de $\mathcal{O}_X$-modules. Si $\mathcal{K}$ est injectif
comme $\mathcal{O}_X$-module, on obtient la flèche pointillée.
Notons $\gamma' : \mathcal{J} \to \mathcal{K}$ la composée
de $\gamma$ avec $\mathcal{J} \to \mathcal{J}/\mathcal{I}\mathcal{J}$.
Un calcul local montre que la somme amalgamée
$$
\xymatrix{
0 \ar[r] &
\mathcal{J} \ar[r] \ar[d]_{\gamma'} &
\mathcal{H}' \ar[r] \ar[d] &
\mathcal{F} \ar[r] \ar@{=}[d] &
0 \\
0 \ar[r] &
\mathcal{K} \ar[r] &
\mathcal{F}' \ar[r] &
\mathcal{F} \ar[r] &
0
}
$$
est une solution au problème posé par le lemme.

\medskip\noindent
Cas général. Choisissons un plongement $\mathcal{K} \subset \mathcal{K}'$,
où $\mathcal{K}'$ est un $\mathcal{O}_X$-module injectif. Soit $\mathcal{Q}$
le quotient, de sorte que l'on ait une suite exacte
$$
0 \to \mathcal{K} \to \mathcal{K}' \to \mathcal{Q} \to 0
$$
Notons
$c' : \mathcal{I} \otimes_{\mathcal{O}_X} \mathcal{F} \to \mathcal{K}'$
la composée. Le paragraphe précédent fournit une suite
$$
0 \to \mathcal{K}' \to \mathcal{E}' \to \mathcal{F} \to 0
$$
comme dans (\ref{defos-equation-extension}), telle que $c_{\mathcal{E}'} = c'$.
Remarquons que la composée de $c'$ avec l'application $\mathcal{K}' \to \mathcal{Q}$
est nulle ; la somme amalgamée de $\mathcal{E}'$ suivant
$\mathcal{K}' \to \mathcal{Q}$ est donc une extension
$$
0 \to \mathcal{Q} \to \mathcal{D}' \to \mathcal{F} \to 0
$$
comme dans (\ref{defos-equation-extension}), telle que $c_{\mathcal{D}'} = 0$.
Cela signifie exactement que $\mathcal{D}'$ est annulé par
$\mathcal{I}$ ; autrement dit, $\mathcal{D}'$ est une extension
de $\mathcal{O}_X$-modules, c'est-à-dire qu'il définit un élément
$$
o(\mathcal{F}, \mathcal{K}, c) \in
\Ext^1_{\mathcal{O}_X}(\mathcal{F}, \mathcal{Q}) =
\Ext^2_{\mathcal{O}_X}(\mathcal{F}, \mathcal{K})
$$
(l'égalité résulte de la suite exacte longue de cohomologie associée
à la suite exacte précédente et de l'annulation des groupes Ext supérieurs
à valeurs dans le module injectif $\mathcal{K}'$). Si
$o(\mathcal{F}, \mathcal{K}, c) = 0$, on peut choisir un scindage
$s : \mathcal{F} \to \mathcal{D}'$ et poser
$$
\mathcal{F}' = \Ker(\mathcal{E}' \to \mathcal{D}'/s(\mathcal{F}))
$$
ce qui donne le diagramme suivant
$$
\xymatrix{
0 \ar[r] &
\mathcal{K} \ar[r] \ar[d] &
\mathcal{F}' \ar[r] \ar[d] &
\mathcal{F} \ar[r] \ar@{=}[d] &
0 \\
0 \ar[r] &
\mathcal{K}' \ar[r] &
\mathcal{E}' \ar[r] &
\mathcal{F} \ar[r] & 0
}
$$
à lignes exactes, qui montre que $c_{\mathcal{F}'} = c$. Réciproquement, si
$\mathcal{F}'$ existe, alors la somme amalgamée de $\mathcal{F}'$ suivant l'application
$\mathcal{K} \to \mathcal{K}'$ est isomorphe à $\mathcal{E}'$, d'après le
lemme \ref{defos-lemma-inf-ext} et l'annulation des groupes Ext supérieurs
à valeurs dans le module injectif $\mathcal{K}'$. On obtient ainsi un diagramme
comme ci-dessus, ce qui implique que $\mathcal{D}'$ est scindée comme extension, c'est-à-dire que
la classe $o(\mathcal{F}, \mathcal{K}, c)$ est nulle.
\end{proof}

\begin{remark}
\label{defos-remark-trivial-thickening}
Soit $(X, \mathcal{O}_X)$ un espace annelé. Un épaississement du premier ordre
$i : (X, \mathcal{O}_X) \to (X', \mathcal{O}_{X'})$ est dit
{\it trivial} s'il existe un morphisme d'espaces annelés
$\pi : (X', \mathcal{O}_{X'}) \to (X, \mathcal{O}_X)$ qui soit un
inverse à gauche de $i$. Le choix d'un tel morphisme
$\pi$ est appelé une {\it trivialisation} de l'épaississement du premier ordre.
Étant donné $\pi$, on obtient un scindage
\begin{equation}
\label{defos-equation-splitting}
\mathcal{O}_{X'} = \mathcal{O}_X \oplus \mathcal{I}
\end{equation}
comme faisceaux d'algèbres sur $X$, en utilisant $\pi^\sharp$ pour scinder la surjection
$\mathcal{O}_{X'} \to \mathcal{O}_X$. Réciproquement, un tel scindage détermine
un morphisme $\pi$. La catégorie des épaississements du premier ordre trivialisés de
$(X, \mathcal{O}_X)$ est équivalente à la catégorie des
$\mathcal{O}_X$-modules.
\end{remark}

\begin{remark}
\label{defos-remark-trivial-extension}
Soit $i : (X, \mathcal{O}_X) \to (X', \mathcal{O}_{X'})$
un épaississement trivial du premier ordre d'espaces annelés,
et soit $\pi : (X', \mathcal{O}_{X'}) \to (X, \mathcal{O}_X)$
une trivialisation. Pour tout triplet
$(\mathcal{F}, \mathcal{K}, c)$ formé d'une paire de
$\mathcal{O}_X$-modules et d'une application
$c : \mathcal{I} \otimes_{\mathcal{O}_X} \mathcal{F} \to \mathcal{K}$,
on peut poser
$$
\mathcal{F}'_{c, triv} = \mathcal{F} \oplus \mathcal{K}
$$
et utiliser le scindage (\ref{defos-equation-splitting}) associé à $\pi$
ainsi que l'application $c$ pour définir la structure de $\mathcal{O}_{X'}$-module
et obtenir une extension (\ref{defos-equation-extension}). On appellera
$\mathcal{F}'_{c, triv}$ l'{\it extension triviale} de $\mathcal{F}$
par $\mathcal{K}$ correspondant
à $c$ et à la trivialisation $\pi$. Pour toute extension
$\mathcal{F}'$ comme dans (\ref{defos-equation-extension}), on peut utiliser
$\pi^\sharp : \mathcal{O}_X \to \mathcal{O}_{X'}$ pour considérer $\mathcal{F}'$
comme une extension de $\mathcal{O}_X$-modules, et donc comme une classe $\xi_{\mathcal{F}'}$
dans $\Ext^1_{\mathcal{O}_X}(\mathcal{F}, \mathcal{K})$.
Le lemme \ref{defos-lemma-inf-ext} assure que
$\mathcal{F}' \mapsto \xi_{\mathcal{F}'}$
induit une bijection
$$
\left\{
\begin{matrix}
\text{isomorphism classes of extensions}\\
\mathcal{F}'\text{ comme dans (\ref{defos-equation-extension}) avec }c = c_{\mathcal{F}'}
\end{matrix}
\right\}
\longrightarrow
\Ext^1_{\mathcal{O}_X}(\mathcal{F}, \mathcal{K})
$$
De plus, l'extension triviale $\mathcal{F}'_{c, triv}$ s'envoie sur la classe nulle.
\end{remark}

\begin{remark}
\label{defos-remark-extension-functorial}
Soit $(X, \mathcal{O}_X)$ un espace annelé. Soient
$(X, \mathcal{O}_X) \to (X'_i, \mathcal{O}_{X'_i})$, $i = 1, 2$,
des épaississements du premier ordre, de faisceaux d'idéaux $\mathcal{I}_i$.
Soit $h : (X'_1, \mathcal{O}_{X'_1}) \to (X'_2, \mathcal{O}_{X'_2})$
un morphisme d'épaississements du premier ordre de $(X, \mathcal{O}_X)$.
On a le diagramme
$$
\xymatrix{
& (X, \mathcal{O}_X) \ar[ld] \ar[rd] & \\
(X'_1, \mathcal{O}_{X'_1}) \ar[rr]^h & & 
(X'_2, \mathcal{O}_{X'_2})
}
$$
Remarquons que $h^\sharp : \mathcal{O}_{X'_2} \to \mathcal{O}_{X'_1}$
induit en particulier une application de $\mathcal{O}_X$-modules
$\mathcal{I}_2 \to \mathcal{I}_1$.
Soit $\mathcal{F}$ un
$\mathcal{O}_X$-module. Soit $(\mathcal{K}_i, c_i)$, pour $i = 1, 2$, une paire
formée d'un $\mathcal{O}_X$-module $\mathcal{K}_i$ et d'une application
$c_i : \mathcal{I}_i \otimes_{\mathcal{O}_X} \mathcal{F} \to
\mathcal{K}_i$. Supposons en outre donnée une application
de $\mathcal{O}_X$-modules $\mathcal{K}_2 \to \mathcal{K}_1$
telle que le diagramme
$$
\xymatrix{
\mathcal{I}_2 \otimes_{\mathcal{O}_X} \mathcal{F}
\ar[r]_-{c_2} \ar[d] &
\mathcal{K}_2 \ar[d] \\
\mathcal{I}_1 \otimes_{\mathcal{O}_X} \mathcal{F}
\ar[r]^-{c_1} &
\mathcal{K}_1
}
$$
soit commutatif. On dispose alors d'une fonctorialité canonique
$$
\left\{
\begin{matrix}
\mathcal{F}'_2\text{ comme dans (\ref{defos-equation-extension}) avec }\\
c_2 = c_{\mathcal{F}'_2}\text{ et }\mathcal{K} = \mathcal{K}_2
\end{matrix}
\right\}
\longrightarrow
\left\{
\begin{matrix}
\mathcal{F}'_1\text{ comme dans (\ref{defos-equation-extension}) avec }\\
c_1 = c_{\mathcal{F}'_1}\text{ et }\mathcal{K} = \mathcal{K}_1
\end{matrix}
\right\}
$$
En effet, en considérant tous les faisceaux $\mathcal{O}_X$, $\mathcal{O}_{X'_i}$,
$\mathcal{F}$, $\mathcal{K}_i$, etc., comme des faisceaux sur $X$, on associe
à $\mathcal{F}'_2$ le faisceau $\mathcal{F}'_1$ défini comme la
somme amalgamée, c'est-à-dire s'insérant dans le diagramme d'extensions suivant :
$$
\xymatrix{
0 \ar[r] &
\mathcal{K}_2 \ar[r] \ar[d] &
\mathcal{F}'_2 \ar[r] \ar[d] &
\mathcal{F} \ar@{=}[d] \ar[r] & 0 \\
0 \ar[r] &
\mathcal{K}_1 \ar[r] &
\mathcal{F}'_1 \ar[r] &
\mathcal{F} \ar[r] & 0
}
$$
Nous omettons la construction de la structure de $\mathcal{O}_{X'_1}$-module
sur la somme amalgamée (elle utilise la commutativité du diagramme
faisant intervenir $c_1$ et $c_2$).
\end{remark}

\begin{remark}
\label{defos-remark-trivial-extension-functorial}
Soient $(X, \mathcal{O}_X)$, $(X, \mathcal{O}_X) \to (X'_i, \mathcal{O}_{X'_i})$,
$\mathcal{I}_i$, et
$h : (X'_1, \mathcal{O}_{X'_1}) \to (X'_2, \mathcal{O}_{X'_2})$
comme dans la remarque \ref{defos-remark-extension-functorial}. Supposons données des
trivialisations $\pi_i : X'_i \to X$ telles que
$\pi_1 = h \circ \pi_2$. Autrement dit, supposons que $h$ soit un morphisme
d'épaississements du premier ordre trivialisés de $(X, \mathcal{O}_X)$. Soit
$(\mathcal{K}_i, c_i)$, pour $i = 1, 2$, une paire formée d'un
$\mathcal{O}_X$-module $\mathcal{K}_i$ et d'une application
$c_i : \mathcal{I}_i \otimes_{\mathcal{O}_X} \mathcal{F} \to
\mathcal{K}_i$. Supposons en outre donnée une application
de $\mathcal{O}_X$-modules $\mathcal{K}_2 \to \mathcal{K}_1$
telle que le diagramme
$$
\xymatrix{
\mathcal{I}_2 \otimes_{\mathcal{O}_X} \mathcal{F}
\ar[r]_-{c_2} \ar[d] &
\mathcal{K}_2 \ar[d] \\
\mathcal{I}_1 \otimes_{\mathcal{O}_X} \mathcal{F}
\ar[r]^-{c_1} &
\mathcal{K}_1
}
$$
soit commutatif. Dans cette situation, la construction de la
remarque \ref{defos-remark-trivial-extension} induit
un diagramme commutatif
$$
\xymatrix{
\{\mathcal{F}'_2\text{ comme dans (\ref{defos-equation-extension}) avec }
c_2 = c_{\mathcal{F}'_2}\text{ et }\mathcal{K} = \mathcal{K}_2\}
\ar[d] \ar[rr] & &
\Ext^1_{\mathcal{O}_X}(\mathcal{F}, \mathcal{K}_2) \ar[d] \\
\{\mathcal{F}'_1\text{ comme dans (\ref{defos-equation-extension}) avec }
c_1 = c_{\mathcal{F}'_1}\text{ et }\mathcal{K} = \mathcal{K}_1\}
\ar[rr] & &
\Ext^1_{\mathcal{O}_X}(\mathcal{F}, \mathcal{K}_1)
}
$$
où l'application verticale de droite est donnée par la fonctorialité de $\Ext$
et l'application $\mathcal{K}_2 \to \mathcal{K}_1$, tandis que l'application verticale de gauche
est celle de la remarque \ref{defos-remark-extension-functorial}.
\end{remark}

\begin{remark}
\label{defos-remark-short-exact-sequence-thickenings}
Soit $(X, \mathcal{O}_X)$ un espace annelé. On dit qu'une suite de morphismes
d'épaississements du premier ordre
$$
(X'_1, \mathcal{O}_{X'_1}) \to
(X'_2, \mathcal{O}_{X'_2}) \to
(X'_3, \mathcal{O}_{X'_3})
$$
de $(X, \mathcal{O}_X)$ est un {\it complexe}
si les applications correspondantes entre
les faisceaux d'idéaux $\mathcal{I}_i$
forment un complexe de $\mathcal{O}_X$-modules
$\mathcal{I}_3 \to \mathcal{I}_2 \to \mathcal{I}_1$
(c'est-à-dire si la composée est nulle). Dans ce cas, la composée
$(X'_1, \mathcal{O}_{X'_1}) \to (X_3', \mathcal{O}_{X'_3})$ se factorise par
$(X, \mathcal{O}_X) \to (X'_3, \mathcal{O}_{X'_3})$ ; autrement dit,
l'épaississement du premier ordre $(X'_1, \mathcal{O}_{X'_1})$ de
$(X, \mathcal{O}_X)$ est trivial et muni d'une
trivialisation canonique
$\pi : (X'_1, \mathcal{O}_{X'_1}) \to (X, \mathcal{O}_X)$.

\medskip\noindent
On dit qu'une suite de morphismes d'épaississements du premier ordre
$$
(X'_1, \mathcal{O}_{X'_1}) \to
(X'_2, \mathcal{O}_{X'_2}) \to
(X'_3, \mathcal{O}_{X'_3})
$$
de $(X, \mathcal{O}_X)$ est une {\it suite exacte courte} si les
applications correspondantes entre les faisceaux d'idéaux forment une suite exacte courte
$$
0 \to \mathcal{I}_3 \to \mathcal{I}_2 \to \mathcal{I}_1 \to 0
$$
de $\mathcal{O}_X$-modules.
\end{remark}

\begin{remark}
\label{defos-remark-complex-thickenings-and-ses-modules}
Soit $(X, \mathcal{O}_X)$ un espace annelé. Soit $\mathcal{F}$ un
$\mathcal{O}_X$-module. Soit
$$
(X'_1, \mathcal{O}_{X'_1}) \to
(X'_2, \mathcal{O}_{X'_2}) \to
(X'_3, \mathcal{O}_{X'_3})
$$
un complexe d'épaississements du premier ordre de $(X, \mathcal{O}_X)$, voir
Remarque \ref{defos-remark-short-exact-sequence-thickenings}.
Soient $(\mathcal{K}_i, c_i)$, $i = 1, 2, 3$, des couples formés d'un
$\mathcal{O}_X$-module $\mathcal{K}_i$ et d'une application
$c_i : \mathcal{I}_i \otimes_{\mathcal{O}_X} \mathcal{F} \to
\mathcal{K}_i$. Supposons donnée une suite exacte courte
de $\mathcal{O}_X$-modules
$$
0 \to \mathcal{K}_3 \to \mathcal{K}_2 \to \mathcal{K}_1 \to 0
$$
telle que
$$
\vcenter{
\xymatrix{
\mathcal{I}_2 \otimes_{\mathcal{O}_X} \mathcal{F}
\ar[r]_-{c_2} \ar[d] &
\mathcal{K}_2 \ar[d] \\
\mathcal{I}_1 \otimes_{\mathcal{O}_X} \mathcal{F}
\ar[r]^-{c_1} &
\mathcal{K}_1
}
}
\quad\text{et}\quad
\vcenter{
\xymatrix{
\mathcal{I}_3 \otimes_{\mathcal{O}_X} \mathcal{F}
\ar[r]_-{c_3} \ar[d] &
\mathcal{K}_3 \ar[d] \\
\mathcal{I}_2 \otimes_{\mathcal{O}_X} \mathcal{F}
\ar[r]^-{c_2} &
\mathcal{K}_2
}
}
$$
soient commutatifs. Supposons enfin donnée une extension
$$
0 \to \mathcal{K}_2 \to \mathcal{F}'_2 \to \mathcal{F} \to 0
$$
comme dans (\ref{defos-equation-extension}), avec $\mathcal{K} = \mathcal{K}_2$,
de $\mathcal{O}_{X'_2}$-modules telle que $c_{\mathcal{F}'_2} = c_2$.
Dans cette situation, nous pouvons appliquer la fonctorialité de la
Remarque \ref{defos-remark-extension-functorial} pour obtenir une extension
$\mathcal{F}'_1$ sur $X'_1$ (nous décrirons ci-dessous $\mathcal{F}'_1$
dans ce cas particulier). Par la
Remarque \ref{defos-remark-trivial-extension},
en utilisant le scindage canonique
$\pi : (X'_1, \mathcal{O}_{X'_1}) \to (X, \mathcal{O}_X)$ de la
Remarque \ref{defos-remark-short-exact-sequence-thickenings},
nous obtenons
$\xi_{\mathcal{F}'_1} \in
\Ext^1_{\mathcal{O}_X}(\mathcal{F}, \mathcal{K}_1)$.
Enfin, nous avons l'obstruction
$$
o(\mathcal{F}, \mathcal{K}_3, c_3) \in
\Ext^2_{\mathcal{O}_X}(\mathcal{F}, \mathcal{K}_3)
$$
du Lemme \ref{defos-lemma-inf-obs-ext}.
Dans cette situation, nous {\bf affirmons} que l'application canonique
$$
\partial :
\Ext^1_{\mathcal{O}_X}(\mathcal{F}, \mathcal{K}_1)
\longrightarrow
\Ext^2_{\mathcal{O}_X}(\mathcal{F}, \mathcal{K}_3)
$$
provenant de la suite exacte courte
$0 \to \mathcal{K}_3 \to \mathcal{K}_2 \to \mathcal{K}_1 \to 0$
envoie $\xi_{\mathcal{F}'_1}$
sur la classe d'obstruction $o(\mathcal{F}, \mathcal{K}_3, c_3)$.

\medskip\noindent
Pour démontrer cette affirmation, choisissons un plongement
$j : \mathcal{K}_3 \to \mathcal{K}$, où $\mathcal{K}$ est un
$\mathcal{O}_X$-module injectif.
Nous pouvons relever $j$ en une application $j' : \mathcal{K}_2 \to \mathcal{K}$.
Posons $\mathcal{E}'_2 = j'_*\mathcal{F}'_2$, égal à la somme amalgamée
de $\mathcal{F}'_2$ par $j'$, de sorte que $c_{\mathcal{E}'_2} = j' \circ c_2$.
Voici le diagramme :
$$
\xymatrix{
0 \ar[r] &
\mathcal{K}_2 \ar[r] \ar[d]_{j'} &
\mathcal{F}'_2 \ar[r] \ar[d] &
\mathcal{F} \ar[r] \ar[d] & 0 \\
0 \ar[r] &
\mathcal{K} \ar[r] &
\mathcal{E}'_2 \ar[r] &
\mathcal{F} \ar[r] & 0
}
$$
Posons $\mathcal{E}'_3 = \mathcal{E}'_2$, mais considérons-le comme un
$\mathcal{O}_{X'_3}$-module au moyen de $\mathcal{O}_{X'_3} \to \mathcal{O}_{X'_2}$.
Alors $c_{\mathcal{E}'_3} = j \circ c_3$.
La démonstration du Lemme \ref{defos-lemma-inf-obs-ext} construit
$o(\mathcal{F}, \mathcal{K}_3, c_3)$
comme le cobord de la classe de l'extension de $\mathcal{O}_X$-modules
$$
0 \to
\mathcal{K}/\mathcal{K}_3 \to
\mathcal{E}'_3/\mathcal{K}_3 \to
\mathcal{F} \to 0
$$
D'autre part, remarquons que
$\mathcal{F}'_1 = \mathcal{F}'_2/\mathcal{K}_3$ ; ainsi, la classe
$\xi_{\mathcal{F}'_1}$ est la classe de l'extension
$$
0 \to \mathcal{K}_2/\mathcal{K}_3 \to \mathcal{F}'_2/\mathcal{K}_3
\to \mathcal{F} \to 0
$$
considérée comme une suite de $\mathcal{O}_X$-modules au moyen de
$\pi^\sharp$, où
$\pi : (X'_1, \mathcal{O}_{X'_1}) \to (X, \mathcal{O}_X)$
est le scindage canonique.
Ainsi, l'affirmation résulte finalement de l'existence du diagramme
commutatif
$$
\xymatrix{
0 \ar[r] &
\mathcal{K}_2/\mathcal{K}_3 \ar[r] \ar[d] &
\mathcal{F}'_2/\mathcal{K}_3 \ar[r] \ar[d] &
\mathcal{F} \ar[r] \ar[d] & 0 \\
0 \ar[r] &
\mathcal{K}/\mathcal{K}_3 \ar[r] &
\mathcal{E}'_3/\mathcal{K}_3 \ar[r] &
\mathcal{F} \ar[r] & 0
}
$$
qui est $\mathcal{O}_X$-linéaire (pour les structures de
$\mathcal{O}_X$-modules données ci-dessus).
\end{remark}








\section{Déformations infinitésimales de modules sur des espaces annelés}
\sectionmark{Déformations de modules sur des espaces annelés}
\label{defos-section-deformation-modules}

\noindent
Soit $i : (X, \mathcal{O}_X) \to (X', \mathcal{O}_{X'})$ un épaississement
du premier ordre d'espaces annelés. Nous employons librement les notations introduites à la
section \ref{defos-section-thickenings-spaces}.
Soit $\mathcal{F}'$ un $\mathcal{O}_{X'}$-module
et posons $\mathcal{F} = i^*\mathcal{F}'$.
Dans cette situation, nous avons une suite exacte courte
$$
0 \to \mathcal{I}\mathcal{F}' \to \mathcal{F}' \to \mathcal{F} \to 0
$$
de $\mathcal{O}_{X'}$-modules. Puisque $\mathcal{I}^2 = 0$, la structure de
$\mathcal{O}_{X'}$-module sur $\mathcal{I}\mathcal{F}'$
provient d'une unique structure de $\mathcal{O}_X$-module.
Ainsi, la suite ci-dessus est une extension comme dans (\ref{defos-equation-extension}).
Dans le cas particulier où $\mathcal{F}' = \mathcal{O}_{X'}$, nous avons
$i^*\mathcal{O}_{X'} = \mathcal{O}_X$ et
$\mathcal{I}\mathcal{O}_{X'} = \mathcal{I}$, et nous retrouvons la
suite des faisceaux structuraux
$$
0 \to \mathcal{I} \to \mathcal{O}_{X'} \to \mathcal{O}_X \to 0
$$

\begin{lemma}
\label{defos-lemma-inf-map-special}
Soit $i : (X, \mathcal{O}_X) \to (X', \mathcal{O}_{X'})$
un épaississement du premier ordre d'espaces annelés.
Soient $\mathcal{F}'$, $\mathcal{G}'$ des $\mathcal{O}_{X'}$-modules.
Posons $\mathcal{F} = i^*\mathcal{F}'$ et $\mathcal{G} = i^*\mathcal{G}'$.
Soit $\varphi : \mathcal{F} \to \mathcal{G}$ une application
$\mathcal{O}_X$-linéaire.
L'ensemble des relèvements de $\varphi$ en une application
$\mathcal{O}_{X'}$-linéaire
$\varphi' : \mathcal{F}' \to \mathcal{G}'$ est, s'il est non vide, un espace
principal homogène sous
$\Hom_{\mathcal{O}_X}(\mathcal{F}, \mathcal{I}\mathcal{G}')$.
\end{lemma}

\begin{proof}
C'est un cas particulier du Lemme \ref{defos-lemma-inf-map}, mais nous en donnons
aussi une démonstration directe. Nous avons des suites exactes courtes de modules
$$
0 \to \mathcal{I} \to \mathcal{O}_{X'} \to \mathcal{O}_X \to 0
\quad\text{et}\quad
0 \to \mathcal{I}\mathcal{G}' \to \mathcal{G}' \to \mathcal{G} \to 0
$$
et de même pour $\mathcal{F}'$.
Comme $\mathcal{I}$ est de carré nul, les structures de
$\mathcal{O}_{X'}$-modules sur $\mathcal{I}$ et
$\mathcal{I}\mathcal{G}'$ proviennent d'uniques structures de
$\mathcal{O}_X$-modules. Il s'ensuit que
$$
\Hom_{\mathcal{O}_{X'}}(\mathcal{F}', \mathcal{I}\mathcal{G}') =
\Hom_{\mathcal{O}_X}(\mathcal{F}, \mathcal{I}\mathcal{G}')
\quad\text{et}\quad
\Hom_{\mathcal{O}_{X'}}(\mathcal{F}', \mathcal{G}) =
\Hom_{\mathcal{O}_X}(\mathcal{F}, \mathcal{G})
$$
Le lemme résulte alors de la suite exacte
$$
0 \to \Hom_{\mathcal{O}_{X'}}(\mathcal{F}', \mathcal{I}\mathcal{G}') \to
\Hom_{\mathcal{O}_{X'}}(\mathcal{F}', \mathcal{G}') \to
\Hom_{\mathcal{O}_{X'}}(\mathcal{F}', \mathcal{G})
$$
voir Homology, lemme \ref{homology-lemma-check-exactness}.
\end{proof}

\begin{lemma}
\label{defos-lemma-deform-module}
Soit $(f, f')$ un morphisme d'épaississements du premier ordre d'espaces annelés
comme dans la Situation \ref{defos-situation-morphism-thickenings}.
Soit $\mathcal{F}'$ un $\mathcal{O}_{X'}$-module
et posons $\mathcal{F} = i^*\mathcal{F}'$.
Supposons que $\mathcal{F}$ soit plat sur $S$
et que $(f, f')$ soit un morphisme strict d'épaississements
(Définition \ref{defos-definition-strict-morphism-thickenings}).
Alors les assertions suivantes sont équivalentes :
\begin{enumerate}
\item $\mathcal{F}'$ est plat sur $S'$ ;
\item l'application canonique
$f^*\mathcal{J} \otimes_{\mathcal{O}_X} \mathcal{F} \to
\mathcal{I}\mathcal{F}'$
est un isomorphisme.
\end{enumerate}
De plus, dans ce cas, les applications
$$
f^*\mathcal{J} \otimes_{\mathcal{O}_X} \mathcal{F} \to
\mathcal{I} \otimes_{\mathcal{O}_X} \mathcal{F} \to
\mathcal{I}\mathcal{F}'
$$
sont des isomorphismes.
\end{lemma}

\begin{proof}
L'application $f^*\mathcal{J} \to \mathcal{I}$ est surjective,
car $(f, f')$ est un morphisme strict d'épaississements.
La dernière assertion résulte donc de (2).

\medskip\noindent
Démontrons l'équivalence de (1) et (2). Nous pouvons vérifier ces conditions
sur les germes. Soit $x \in X \subset X'$
un point d'image $s = f(x) \in S \subset S'$.
Posons $A' = \mathcal{O}_{S', s}$, $B' = \mathcal{O}_{X', x}$,
$A = \mathcal{O}_{S, s}$ et $B = \mathcal{O}_{X, x}$.
Alors $A = A'/J$ et $B = B'/I$ pour certains idéaux de carré nul.
Puisque $(f, f')$ est un morphisme strict d'épaississements, nous avons
$I = JB'$.
Soient $M' = \mathcal{F}'_x$ et $M = \mathcal{F}_x$.
Alors $M'$ est un $B'$-module et $M$ est un $B$-module.
Comme $\mathcal{F} = i^*\mathcal{F}'$, le noyau de la
surjection $M' \to M$ est $IM' = JM'$. Nous avons donc une suite exacte
courte
$$
0 \to JM' \to M' \to M \to 0
$$
En utilisant
Faisceaux, lemme \ref{sheaves-lemma-stalk-pullback-modules}
et
Modules, lemme \ref{modules-lemma-stalk-tensor-product}
pour identifier les germes des images réciproques et des produits tensoriels, nous voyons
que le germe en $x$ de l'application canonique de la lemme est l'application
$$
(J \otimes_A B) \otimes_B M = J \otimes_A M = J \otimes_{A'} M'
\longrightarrow JM'
$$
L'hypothèse que $\mathcal{F}$ est plat sur $S$ signifie que
$M$ est un $A$-module plat.

\medskip\noindent
Supposons (1). La platitude implique
$\text{Tor}_1^{A'}(M', A) = 0$ d'après
Algèbre, lemme \ref{algebra-lemma-characterize-flat}.
Cela signifie que $J \otimes_{A'} M' \to M'$ est injective d'après
Algèbre, remarque \ref{algebra-remark-Tor-ring-mod-ideal}.
Ainsi, $J \otimes_A M \to JM'$ est un isomorphisme.

\medskip\noindent
Supposons (2). Alors $J \otimes_{A'} M' \to M'$ est injective. Donc
$\text{Tor}_1^{A'}(M', A) = 0$ d'après
Algèbre, remarque \ref{algebra-remark-Tor-ring-mod-ideal}.
Ainsi, $M'$ est plat sur $A'$ d'après
Algèbre, lemme \ref{algebra-lemma-what-does-it-mean}.
\end{proof}

\begin{lemma}
\label{defos-lemma-inf-map-rel}
Soit $(f, f')$ un morphisme d'épaississements du premier ordre comme dans la
Situation \ref{defos-situation-morphism-thickenings}.
Soient $\mathcal{F}'$, $\mathcal{G}'$ des $\mathcal{O}_{X'}$-modules et posons
$\mathcal{F} = i^*\mathcal{F}'$ et $\mathcal{G} = i^*\mathcal{G}'$.
Soit $\varphi : \mathcal{F} \to \mathcal{G}$ une application
$\mathcal{O}_X$-linéaire.
Supposons que $\mathcal{G}'$ soit plat sur $S'$ et que
$(f, f')$ soit un morphisme strict d'épaississements.
L'ensemble des relèvements de $\varphi$ en une application
$\mathcal{O}_{X'}$-linéaire
$\varphi' : \mathcal{F}' \to \mathcal{G}'$ est, s'il est non vide, un espace
principal homogène sous
$$
\Hom_{\mathcal{O}_X}(\mathcal{F},
\mathcal{G} \otimes_{\mathcal{O}_X} f^*\mathcal{J})
$$
\end{lemma}

\begin{proof}
Il suffit de combiner les Lemmes \ref{defos-lemma-inf-map-special} et
\ref{defos-lemma-deform-module}.
\end{proof}

\begin{lemma}
\label{defos-lemma-inf-obs-map-special}
Soit $i : (X, \mathcal{O}_X) \to (X', \mathcal{O}_{X'})$
un épaississement du premier ordre d'espaces annelés.
Soient $\mathcal{F}'$, $\mathcal{G}'$ des $\mathcal{O}_{X'}$-modules et posons
$\mathcal{F} = i^*\mathcal{F}'$ et $\mathcal{G} = i^*\mathcal{G}'$.
Soit $\varphi : \mathcal{F} \to \mathcal{G}$ une application
$\mathcal{O}_X$-linéaire.
Il existe un élément
$$
o(\varphi) \in
\Ext^1_{\mathcal{O}_X}(Li^*\mathcal{F}',
\mathcal{I}\mathcal{G}')
$$
dont l'annulation est une condition nécessaire et suffisante pour
l'existence d'un relèvement de $\varphi$ en une application
$\mathcal{O}_{X'}$-linéaire
$\varphi' : \mathcal{F}' \to \mathcal{G}'$.
\end{lemma}

\begin{proof}
Il ressort de la démonstration du Lemme \ref{defos-lemma-inf-map-special} que
l'annulation du cobord de $\varphi$ par l'application
$$
\Hom_{\mathcal{O}_X}(\mathcal{F}, \mathcal{G}) =
\Hom_{\mathcal{O}_{X'}}(\mathcal{F}', \mathcal{G}) \longrightarrow
\Ext^1_{\mathcal{O}_{X'}}(\mathcal{F}', \mathcal{I}\mathcal{G}')
$$
est une condition nécessaire et suffisante pour l'existence d'un relèvement. Nous
concluons puisque
$$
\Ext^1_{\mathcal{O}_{X'}}(\mathcal{F}', \mathcal{I}\mathcal{G}') =
\Ext^1_{\mathcal{O}_X}(Li^*\mathcal{F}', \mathcal{I}\mathcal{G}')
$$
par l'adjonction de $i_* = Ri_*$ et $Li^*$ dans la catégorie dérivée
(Cohomologie, lemme \ref{cohomology-lemma-adjoint}).
\end{proof}

\begin{lemma}
\label{defos-lemma-inf-obs-map-rel}
Soit $(f, f')$ un morphisme d'épaississements du
premier ordre comme dans la Situation \ref{defos-situation-morphism-thickenings}.
Soient $\mathcal{F}'$, $\mathcal{G}'$ des $\mathcal{O}_{X'}$-modules et posons
$\mathcal{F} = i^*\mathcal{F}'$ et $\mathcal{G} = i^*\mathcal{G}'$.
Soit $\varphi : \mathcal{F} \to \mathcal{G}$ une application
$\mathcal{O}_X$-linéaire.
Supposons que $\mathcal{F}'$ et $\mathcal{G}'$ soient plats sur $S'$ et
que $(f, f')$ soit un morphisme strict d'épaississements. Il existe un élément
$$
o(\varphi) \in  \Ext^1_{\mathcal{O}_X}(\mathcal{F},
\mathcal{G} \otimes_{\mathcal{O}_X} f^*\mathcal{J})
$$
dont l'annulation est une condition nécessaire et suffisante pour
l'existence d'un relèvement de $\varphi$ en une application
$\mathcal{O}_{X'}$-linéaire
$\varphi' : \mathcal{F}' \to \mathcal{G}'$.
\end{lemma}

\begin{proof}[Première démonstration]
Cela résulte du Lemme \ref{defos-lemma-inf-obs-map-special},
car nous affirmons que, sous les hypothèses de la lemme, nous avons
$$
\Ext^1_{\mathcal{O}_X}(Li^*\mathcal{F}',
\mathcal{I}\mathcal{G}') =
\Ext^1_{\mathcal{O}_X}(\mathcal{F},
\mathcal{G} \otimes_{\mathcal{O}_X} f^*\mathcal{J})
$$
En effet, nous avons
$\mathcal{I}\mathcal{G}' =
\mathcal{G} \otimes_{\mathcal{O}_X} f^*\mathcal{J}$
d'après le Lemme \ref{defos-lemma-deform-module}.
D'autre part, remarquons que
$$
H^{-1}(Li^*\mathcal{F}') =
\text{Tor}_1^{\mathcal{O}_{X'}}(\mathcal{F}', \mathcal{O}_X)
$$
(nous omettons le calcul local). En utilisant la suite exacte courte
$$
0 \to \mathcal{I} \to \mathcal{O}_{X'} \to \mathcal{O}_X \to 0
$$
nous voyons que ce $\text{Tor}_1$ est calculé par le noyau de l'application
$\mathcal{I} \otimes_{\mathcal{O}_X} \mathcal{F} \to \mathcal{I}\mathcal{F}'$,
qui est nul d'après la dernière assertion du Lemme \ref{defos-lemma-deform-module}.
Ainsi, $\tau_{\geq -1}Li^*\mathcal{F}' = \mathcal{F}$.
D'autre part, nous avons
$$
\Ext^1_{\mathcal{O}_X}(Li^*\mathcal{F}',
\mathcal{I}\mathcal{G}') =
\Ext^1_{\mathcal{O}_X}(\tau_{\geq -1}Li^*\mathcal{F}',
\mathcal{I}\mathcal{G}')
$$
par le dual de
Catégories dérivées, lemme \ref{derived-lemma-negative-vanishing}.
\end{proof}

\begin{proof}[Seconde démonstration]
Nous pouvons appliquer le Lemme \ref{defos-lemma-inf-obs-map} comme suit. Remarquons que
$\mathcal{K} = \mathcal{I} \otimes_{\mathcal{O}_X} \mathcal{F}$ et
$\mathcal{L} = \mathcal{I} \otimes_{\mathcal{O}_X} \mathcal{G}$
d'après le Lemme \ref{defos-lemma-deform-module}, que
$c_{\mathcal{F}'} = 1 \otimes 1$ et $c_{\mathcal{G}'} = 1 \otimes 1$,
et qu'en prenant $\psi = 1 \otimes \varphi$, le diagramme de la lemme
est commutatif. Ainsi, $o(\varphi) = o(\varphi, 1 \otimes \varphi)$
convient.
\end{proof}

\begin{lemma}
\label{defos-lemma-inf-ext-rel}
Soit $(f, f')$ un morphisme d'épaississements du premier ordre comme dans la
Situation \ref{defos-situation-morphism-thickenings}.
Soit $\mathcal{F}$ un $\mathcal{O}_X$-module.
Supposons que $(f, f')$ soit un morphisme strict d'épaississements et que
$\mathcal{F}$ soit plat sur $S$. S'il existe un couple
$(\mathcal{F}', \alpha)$ formé d'un
$\mathcal{O}_{X'}$-module $\mathcal{F}'$ plat sur $S'$ et d'un isomorphisme
$\alpha : i^*\mathcal{F}' \to \mathcal{F}$, alors l'ensemble des
classes d'isomorphisme de tels couples est un espace principal homogène
sous
$\Ext^1_{\mathcal{O}_X}(
\mathcal{F}, \mathcal{I} \otimes_{\mathcal{O}_X} \mathcal{F})$.
\end{lemma}

\begin{proof}
Si nous supposons qu'il existe un tel module, alors l'application canonique
$$
f^*\mathcal{J} \otimes_{\mathcal{O}_X} \mathcal{F} \to
\mathcal{I} \otimes_{\mathcal{O}_X} \mathcal{F}
$$
est un isomorphisme d'après le Lemme \ref{defos-lemma-deform-module}. Appliquons le
Lemme \ref{defos-lemma-inf-ext} avec $\mathcal{K} = 
\mathcal{I} \otimes_{\mathcal{O}_X} \mathcal{F}$
et $c = 1$. D'après le Lemme \ref{defos-lemma-deform-module}, les extensions
$\mathcal{F}'$ correspondantes sont toutes plates sur $S'$.
\end{proof}

\begin{lemma}
\label{defos-lemma-inf-obs-ext-rel}
Soit $(f, f')$ un morphisme d'épaississements du premier ordre comme dans la
Situation \ref{defos-situation-morphism-thickenings}.
Soit $\mathcal{F}$ un $\mathcal{O}_X$-module. Supposons que
$(f, f')$ soit un morphisme strict d'épaississements
et que $\mathcal{F}$ soit plat sur $S$. Il existe un
$\mathcal{O}_{X'}$-module $\mathcal{F}'$ plat sur $S'$ tel que
$i^*\mathcal{F}' \cong \mathcal{F}$, si et seulement si
\begin{enumerate}
\item l'application canonique $
f^*\mathcal{J} \otimes_{\mathcal{O}_X} \mathcal{F} \to
\mathcal{I} \otimes_{\mathcal{O}_X} \mathcal{F}$
est un isomorphisme ;
\item la classe
$o(\mathcal{F}, \mathcal{I} \otimes_{\mathcal{O}_X} \mathcal{F}, 1)
\in \Ext^2_{\mathcal{O}_X}(
\mathcal{F}, \mathcal{I} \otimes_{\mathcal{O}_X} \mathcal{F})$
du Lemme \ref{defos-lemma-inf-obs-ext} est nulle.
\end{enumerate}
\end{lemma}

\begin{proof}
Cela résulte immédiatement de la caractérisation des
$\mathcal{O}_{X'}$-modules plats sur $S'$ donnée par la
Lemme \ref{defos-lemma-deform-module} et de la
Lemme \ref{defos-lemma-inf-obs-ext}.
\end{proof}






\section{Application aux modules plats sur les épaississements plats d'espaces annelés}
\label{defos-section-flat}

\noindent
Considérons un diagramme commutatif
$$
\xymatrix{
(X, \mathcal{O}_X) \ar[r]_i \ar[d]_f & (X', \mathcal{O}_{X'}) \ar[d]^{f'} \\
(S, \mathcal{O}_S) \ar[r]^t & (S', \mathcal{O}_{S'})
}
$$
d'espaces annelés dont les flèches horizontales sont des épaississements du premier ordre,
comme dans la Situation \ref{defos-situation-morphism-thickenings}. Posons
$\mathcal{I} = \Ker(i^\sharp) \subset \mathcal{O}_{X'}$ et
$\mathcal{J} = \Ker(t^\sharp) \subset \mathcal{O}_{S'}$.
Soit $\mathcal{F}$ un $\mathcal{O}_X$-module. Supposons que
\begin{enumerate}
\item $(f, f')$ soit un morphisme strict d'épaississements ;
\item $f'$ soit plat ;
\item $\mathcal{F}$ soit plat sur $S$.
\end{enumerate}
Remarquons que (1) $+$ (2) impliquent que $\mathcal{I} = f^*\mathcal{J}$
(appliquer le Lemme \ref{defos-lemma-deform-module} à $\mathcal{O}_{X'}$).
La théorie de la section précédente prend une forme particulièrement simple
sous ces hypothèses. Nous résumons dans la lemme suivante les résultats déjà obtenus.

\begin{lemma}
\label{defos-lemma-flat}
Dans la situation ci-dessus, on a les assertions suivantes.
\begin{enumerate}
\item Il existe un $\mathcal{O}_{X'}$-module $\mathcal{F}'$ plat sur
$S'$ tel que $i^*\mathcal{F}' \cong \mathcal{F}$, si et seulement si
la classe
$o(\mathcal{F}, f^*\mathcal{J} \otimes_{\mathcal{O}_X} \mathcal{F}, 1)
\in \Ext^2_{\mathcal{O}_X}(
\mathcal{F}, f^*\mathcal{J} \otimes_{\mathcal{O}_X} \mathcal{F})$
du Lemme \ref{defos-lemma-inf-obs-ext} est nulle.
\item S'il existe un tel module, alors l'ensemble des classes d'isomorphisme
des relèvements est un espace principal homogène sous
$\Ext^1_{\mathcal{O}_X}(
\mathcal{F}, f^*\mathcal{J} \otimes_{\mathcal{O}_X} \mathcal{F})$.
\item Étant donné un relèvement $\mathcal{F}'$, l'ensemble des automorphismes de
$\mathcal{F}'$ dont l'image réciproque est $\text{id}_\mathcal{F}$ est canoniquement
isomorphe à $\Ext^0_{\mathcal{O}_X}(
\mathcal{F}, f^*\mathcal{J} \otimes_{\mathcal{O}_X} \mathcal{F})$.
\end{enumerate}
\end{lemma}

\begin{proof}
L'assertion (1) résulte du Lemme \ref{defos-lemma-inf-obs-ext-rel},
car nous avons vu ci-dessus que $\mathcal{I} = f^*\mathcal{J}$.
L'assertion (2) résulte du Lemme \ref{defos-lemma-inf-ext-rel}.
L'assertion (3) résulte du Lemme \ref{defos-lemma-inf-map-rel}.
\end{proof}

\begin{situation}
\label{defos-situation-ses-flat-thickenings}
Soit $f : (X, \mathcal{O}_X) \to (S, \mathcal{O}_S)$ un morphisme
d'espaces annelés. Considérons un diagramme commutatif
$$
\xymatrix{
(X'_1, \mathcal{O}'_1) \ar[r]_h \ar[d]_{f'_1} &
(X'_2, \mathcal{O}'_2) \ar[r] \ar[d]_{f'_2} &
(X'_3, \mathcal{O}'_3) \ar[d]_{f'_3} \\
(S'_1, \mathcal{O}_{S'_1}) \ar[r] &
(S'_2, \mathcal{O}_{S'_2}) \ar[r] &
(S'_3, \mathcal{O}_{S'_3})
}
$$
où (a) la ligne supérieure est une suite exacte courte d'épaississements du premier ordre
de $X$, (b) la ligne inférieure est une suite exacte courte d'épaississements du premier ordre
de $S$, (c) chaque $f'_i$ induit $f$ par restriction, (d) chaque couple
$(f, f_i')$ est un morphisme strict d'épaississements et (e) chaque $f'_i$ est plat.
Enfin, soit $\mathcal{F}'_2$ un $\mathcal{O}'_2$-module plat sur
$S'_2$, et posons $\mathcal{F} = \mathcal{F}'_2|_X$. Soit $\pi : X'_1 \to X$
le scindage canonique
(Remarque \ref{defos-remark-short-exact-sequence-thickenings}).
\end{situation}

\begin{lemma}
\label{defos-lemma-verify-iv}
Dans la Situation \ref{defos-situation-ses-flat-thickenings}, les modules
$\pi^*\mathcal{F}$ et $h^*\mathcal{F}'_2$ sont des $\mathcal{O}'_1$-modules
plats sur $S'_1$ et dont la restriction est $\mathcal{F}$ sur $X$.
Leur différence (Lemme \ref{defos-lemma-flat}) est un élément
$\theta$ de $\Ext^1_{\mathcal{O}_X}(
\mathcal{F}, f^*\mathcal{J}_1 \otimes_{\mathcal{O}_X} \mathcal{F})$
dont le cobord dans
$\Ext^2_{\mathcal{O}_X}(
\mathcal{F}, f^*\mathcal{J}_3 \otimes_{\mathcal{O}_X} \mathcal{F})$
est égal à l'obstruction (Lemme \ref{defos-lemma-flat})
au relèvement de $\mathcal{F}$ en un $\mathcal{O}'_3$-module plat sur $S'_3$.
\end{lemma}

\begin{proof}
Remarquons que $\pi^*\mathcal{F}$ et $h^*\mathcal{F}'_2$
ont tous deux pour restriction $\mathcal{F}$ à $X$ et que les noyaux de
$\pi^*\mathcal{F} \to \mathcal{F}$ et $h^*\mathcal{F}'_2 \to \mathcal{F}$
sont donnés par $f^*\mathcal{J}_1 \otimes_{\mathcal{O}_X} \mathcal{F}$.
Ils sont donc plats d'après le Lemme \ref{defos-lemma-deform-module}.
Prendre le cobord a un sens, car la suite de modules
$$
0 \to f^*\mathcal{J}_3 \otimes_{\mathcal{O}_X} \mathcal{F} \to
f^*\mathcal{J}_2 \otimes_{\mathcal{O}_X} \mathcal{F} \to
f^*\mathcal{J}_1 \otimes_{\mathcal{O}_X} \mathcal{F} \to 0
$$
est exacte courte par les hypothèses de la
Situation \ref{defos-situation-ses-flat-thickenings}
et le fait que $\mathcal{F}$ est plat sur $S$.
L'assertion sur la classe d'obstruction est une traduction directe du résultat de la
Remarque \ref{defos-remark-complex-thickenings-and-ses-modules}
dans cette situation particulière.
\end{proof}





\section{Déformations d'espaces annelés et complexe cotangent naïf}
\label{defos-section-deformations-ringed-spaces}

\noindent
Dans cette section, nous utilisons le complexe cotangent naïf pour faire un peu
de théorie des déformations. Nous partons d'un épaississement du premier ordre
$t : (S, \mathcal{O}_S) \to (S', \mathcal{O}_{S'})$ d'espaces annelés.
Nous notons $\mathcal{J} = \Ker(t^\sharp)$ et identifions les espaces
topologiques sous-jacents à $S$ et $S'$.
Nous nous donnons en outre un morphisme d'espaces annelés
$f : (X, \mathcal{O}_X) \to (S, \mathcal{O}_S)$, un $\mathcal{O}_X$-module
$\mathcal{G}$ et une $f$-application $c : \mathcal{J} \to \mathcal{G}$
de faisceaux de modules (Faisceaux, Définition
\ref{sheaves-definition-f-map} et section
\ref{sheaves-section-ringed-spaces-functoriality-modules}).
Dans cette section, nous cherchons à savoir si l'on peut trouver
le point d'interrogation qui complète le diagramme suivant
\begin{equation}
\label{defos-equation-to-solve-ringed-spaces}
\vcenter{
\xymatrix{
0 \ar[r] & \mathcal{G} \ar[r] & {?} \ar[r] & \mathcal{O}_X \ar[r] & 0 \\
0 \ar[r] & \mathcal{J} \ar[u]^c \ar[r] & \mathcal{O}_{S'} \ar[u] \ar[r] &
\mathcal{O}_S \ar[u] \ar[r] & 0
}
}
\end{equation}
(où les flèches verticales sont des $f$-applications), et aussi dans quelle
mesure la solution est unique (lorsqu'elle existe). Plus précisément,
nous cherchons un épaississement du premier ordre
$i : (X, \mathcal{O}_X) \to (X', \mathcal{O}_{X'})$
et un morphisme d'épaississements $(f, f')$ comme dans
(\ref{defos-equation-morphism-thickenings}), où $\Ker(i^\sharp)$ est identifié
à $\mathcal{G}$, tels que $(f')^\sharp$ induise l'application donnée $c$.
Nous dirons que $X'$ est une {\it solution} de
(\ref{defos-equation-to-solve-ringed-spaces}).

\begin{lemma}
\label{defos-lemma-huge-diagram-ringed-spaces}
Supposons donné un diagramme commutatif de morphismes d'espaces annelés
\begin{equation}
\label{defos-equation-huge-1}
\vcenter{
\xymatrix{
& (X_2, \mathcal{O}_{X_2}) \ar[r]_{i_2} \ar[d]_{f_2} \ar[ddl]_g &
(X'_2, \mathcal{O}_{X'_2}) \ar[d]^{f'_2} \\
& (S_2, \mathcal{O}_{S_2}) \ar[r]^{t_2} \ar[ddl]|\hole &
(S'_2, \mathcal{O}_{S'_2}) \ar[ddl] \\
(X_1, \mathcal{O}_{X_1}) \ar[r]_{i_1} \ar[d]_{f_1} &
(X'_1, \mathcal{O}_{X'_1}) \ar[d]^{f'_1} \\
(S_1, \mathcal{O}_{S_1}) \ar[r]^{t_1} &
(S'_1, \mathcal{O}_{S'_1})
}
}
\end{equation}
dont les flèches horizontales sont des épaississements du premier ordre. Posons
$\mathcal{G}_j = \Ker(i_j^\sharp)$ et supposons donnée
une $g$-application $\nu : \mathcal{G}_1 \to \mathcal{G}_2$ de modules
donnant lieu au diagramme commutatif
\begin{equation}
\label{defos-equation-huge-2}
\vcenter{
\xymatrix{
& 0 \ar[r] & \mathcal{G}_2 \ar[r] &
\mathcal{O}_{X'_2} \ar[r] &
\mathcal{O}_{X_2} \ar[r] & 0 \\
& 0 \ar[r]|\hole &
\mathcal{J}_2 \ar[u]_{c_2} \ar[r] &
\mathcal{O}_{S'_2} \ar[u] \ar[r]|\hole &
\mathcal{O}_{S_2} \ar[u] \ar[r] & 0 \\
0 \ar[r] & \mathcal{G}_1 \ar[ruu] \ar[r] &
\mathcal{O}_{X'_1} \ar[r] &
\mathcal{O}_{X_1} \ar[ruu] \ar[r] & 0 \\
0 \ar[r] & \mathcal{J}_1 \ar[ruu]|\hole \ar[u]^{c_1} \ar[r] &
\mathcal{O}_{S'_1} \ar[ruu]|\hole \ar[u] \ar[r] &
\mathcal{O}_{S_1} \ar[ruu]|\hole \ar[u] \ar[r] & 0
}
}
\end{equation}
dont les faces avant et arrière sont des solutions de
(\ref{defos-equation-to-solve-ringed-spaces}).
\begin{enumerate}
\item Il existe un élément canonique de
$\Ext^1_{\mathcal{O}_{X_2}}(Lg^*\NL_{X_1/S_1}, \mathcal{G}_2)$
dont l'annulation est une condition nécessaire et suffisante pour qu'il existe
un morphisme d'espaces annelés $X'_2 \to X'_1$ complétant
(\ref{defos-equation-huge-1}) de façon compatible à $\nu$.
\item S'il existe un morphisme $X'_2 \to X'_1$ complétant
(\ref{defos-equation-huge-1}) de façon compatible à $\nu$, l'ensemble de tous ces
morphismes est un espace principal homogène sous
$$
\Hom_{\mathcal{O}_{X_1}}(\Omega_{X_1/S_1}, g_*\mathcal{G}_2) =
\Hom_{\mathcal{O}_{X_2}}(g^*\Omega_{X_1/S_1}, \mathcal{G}_2) =
\Ext^0_{\mathcal{O}_{X_2}}(Lg^*\NL_{X_1/S_1}, \mathcal{G}_2).
$$
\end{enumerate}
\end{lemma}

\begin{proof}
Le complexe cotangent naïf $\NL_{X_1/S_1}$ est défini dans Modules,
Définition \ref{modules-definition-cotangent-complex-morphism-ringed-topoi}.
Les égalités de la dernière assertion de la lemme résultent de ce que
$g^*$ est adjoint à $g_*$, de l'égalité
$H^0(\NL_{X_1/S_1}) = \Omega_{X_1/S_1}$ (par construction du complexe
cotangent naïf) et du fait que $Lg^*$ est le foncteur dérivé à gauche de
$g^*$. Nous travaillerons donc, dans la suite de la démonstration, avec les
groupes
$\Ext^k_{\mathcal{O}_{X_2}}(Lg^*\NL_{X_1/S_1}, \mathcal{G}_2)$,
$k = 0, 1$. Montrons d'abord que l'on peut se ramener au cas où les espaces
topologiques sous-jacents à tous les espaces annelés de la lemme sont les mêmes.

\medskip\noindent
Pour cela, remarquons que $g^{-1}\NL_{X_1/S_1}$ est égal au complexe
cotangent naïf de l'homomorphisme de faisceaux d'anneaux
$g^{-1}f_1^{-1}\mathcal{O}_{S_1} \to g^{-1}\mathcal{O}_{X_1}$ ; voir
Modules, Lemme \ref{modules-lemma-pullback-NL}.
De plus, le terme de degré $0$ de $\NL_{X_1/S_1}$ est un
$\mathcal{O}_{X_1}$-module plat ; l'application canonique
$$
Lg^*\NL_{X_1/S_1}
\longrightarrow
g^{-1}\NL_{X_1/S_1} \otimes_{g^{-1}\mathcal{O}_{X_1}} \mathcal{O}_{X_2}
$$
induit donc un isomorphisme sur les faisceaux de cohomologie en degrés $0$ et
$-1$. On peut ainsi remplacer les groupes Ext de la lemme par
$$
\Ext^k_{g^{-1}\mathcal{O}_{X_1}}(g^{-1}\NL_{X_1/S_1}, \mathcal{G}_2) =
\Ext^k_{g^{-1}\mathcal{O}_{X_1}}(
\NL_{g^{-1}\mathcal{O}_{X_1}/g^{-1}f_1^{-1}\mathcal{O}_{S_1}}, \mathcal{G}_2)
$$
L'ensemble des morphismes d'espaces annelés $X'_2 \to X'_1$ complétant
(\ref{defos-equation-huge-1}) de façon compatible à $\nu$ est en bijection avec
l'ensemble des homomorphismes de $g^{-1}f_1^{-1}\mathcal{O}_{S'_1}$-algèbres
$g^{-1}\mathcal{O}_{X'_1} \to \mathcal{O}_{X'_2}$ compatibles avec
$f^\sharp$ et $\nu$. Nous voyons ainsi que nous pouvons supposer donné un
diagramme (\ref{defos-equation-huge-2}) de faisceaux sur $X$ et chercher un
homomorphisme de faisceaux d'anneaux
$\mathcal{O}_{X'_1} \to \mathcal{O}_{X'_2}$ qui le complète.

\medskip\noindent
Dans la suite de la démonstration, nous supposons que tous les espaces
topologiques sous-jacents sont identiques, c'est-à-dire que nous disposons
d'un diagramme (\ref{defos-equation-huge-2}) de faisceaux sur un espace $X$ et
cherchons les homomorphismes de faisceaux d'anneaux
$\mathcal{O}_{X'_1} \to \mathcal{O}_{X'_2}$ qui le complètent.
Comme groupes Ext, nous utiliserons
$\Ext^k_{\mathcal{O}_{X_1}}(
\NL_{\mathcal{O}_{X_1}/\mathcal{O}_{S_1}}, \mathcal{G}_2)$, $k = 0, 1$.

\medskip\noindent
Étape 1. Construction de la classe d'obstruction. Considérons le faisceau
d'ensembles
$$
\mathcal{E} = \mathcal{O}_{X'_1} \times_{\mathcal{O}_{X_2}} \mathcal{O}_{X'_2}
$$
Il est muni d'une application surjective
$\alpha : \mathcal{E} \to \mathcal{O}_{X_1}$ ; nous pouvons donc employer
$\NL(\alpha)$ à la place de
$\NL_{\mathcal{O}_{X_1}/\mathcal{O}_{S_1}}$ ; voir Modules, Lemme
\ref{modules-lemma-NL-up-to-qis}.
Posons
$$
\mathcal{I}' =
\Ker(\mathcal{O}_{S'_1}[\mathcal{E}] \to \mathcal{O}_{X_1})
\quad\text{et}\quad
\mathcal{I} =
\Ker(\mathcal{O}_{S_1}[\mathcal{E}] \to \mathcal{O}_{X_1})
$$
Il existe une surjection $\mathcal{I}' \to \mathcal{I}$ dont le noyau est
$\mathcal{J}_1\mathcal{O}_{S'_1}[\mathcal{E}]$.
Nous obtenons deux homomorphismes de $\mathcal{O}_{S'_1}$-algèbres
$$
a : \mathcal{O}_{S'_1}[\mathcal{E}] \to \mathcal{O}_{X'_1}
\quad\text{et}\quad
b : \mathcal{O}_{S'_1}[\mathcal{E}] \to \mathcal{O}_{X'_2}
$$
qui induisent des applications
$a|_{\mathcal{I}'} : \mathcal{I}' \to \mathcal{G}_1$ et
$b|_{\mathcal{I}'} : \mathcal{I}' \to \mathcal{G}_2$.
Les deux applications $a$ et $b$ annulent $(\mathcal{I}')^2$.
De plus, $a$ et $b$ coïncident sur
$\mathcal{J}_1\mathcal{O}_{S'_1}[\mathcal{E}]$ en tant qu'applications à
valeurs dans $\mathcal{G}_2$, car le carré de gauche de
(\ref{defos-equation-huge-2}) est commutatif. Par conséquent, la différence
$b|_{\mathcal{I}'} - \nu \circ a|_{\mathcal{I}'}$
induit une application $\mathcal{O}_{X_1}$-linéaire bien définie
$$
\xi : \mathcal{I}/\mathcal{I}^2 \longrightarrow \mathcal{G}_2
$$
qui envoie la classe d'une section locale $f$ de $\mathcal{I}$ sur
$\nu(a(f')) - b(f')$, où $f'$ est un relèvement de $f$ en une section locale
de $\mathcal{I}'$. Nous notons
$[\xi] \in \Ext^1_{\mathcal{O}_{X_1}}(\NL(\alpha), \mathcal{G}_2)$
son image (voir ci-dessous).

\medskip\noindent
Étape 2. L'annulation de $[\xi]$ est nécessaire. Écrivons
$\Omega = \Omega_{\mathcal{O}_{S_1}[\mathcal{E}]/\mathcal{O}_{S_1}}
\otimes_{\mathcal{O}_{S_1}[\mathcal{E}]} \mathcal{O}_{X_1}$.
Remarquons que $\NL(\alpha) = (\mathcal{I}/\mathcal{I}^2 \to \Omega)$
s'insère dans un triangle distingué
$$
\Omega[0] \to
\NL(\alpha) \to
\mathcal{I}/\mathcal{I}^2[1] \to
\Omega[1]
$$
On voit donc que $[\xi]$ est nul si et seulement si $\xi$ est un composé
$\mathcal{I}/\mathcal{I}^2 \to \Omega \to \mathcal{G}_2$
pour une certaine application $\Omega \to \mathcal{G}_2$. Supposons qu'il
existe un homomorphisme de faisceaux d'anneaux
$\varphi : \mathcal{O}_{X'_1} \to \mathcal{O}_{X'_2}$ complétant
(\ref{defos-equation-huge-2}). Considérons alors l'application
$\mathcal{O}_{S'_1}[\mathcal{E}] \to \mathcal{G}_2$,
$f' \mapsto b(f') - \varphi(a(f'))$. Un calcul montre qu'elle annule
$\mathcal{J}_1\mathcal{O}_{S'_1}[\mathcal{E}]$ et induit une dérivation
$\mathcal{O}_{S_1}[\mathcal{E}] \to \mathcal{G}_2$.
L'application linéaire $\Omega \to \mathcal{G}_2$ ainsi obtenue atteste que
$[\xi] = 0$ dans ce cas.

\medskip\noindent
Étape 3. L'annulation de $[\xi]$ est suffisante. Soit
$\theta : \Omega \to \mathcal{G}_2$ une application
$\mathcal{O}_{X_1}$-linéaire telle que $\xi$ soit égal à
$\theta \circ (\mathcal{I}/\mathcal{I}^2 \to \Omega)$.
Un calcul montre alors que
$$
b + \theta \circ d : \mathcal{O}_{S'_1}[\mathcal{E}] \to \mathcal{O}_{X'_2}
$$
restreint à $\mathcal{I}'$ coïncide avec
$\nu \circ a : \mathcal{I}' \to \mathcal{G}_2$. Comme
$\mathcal{O}_{X'_1}$ est la somme amalgamée de
$\mathcal{I}' \to \mathcal{O}_{S'_1}[\mathcal{E}]$ et
$\mathcal{I}' \to \mathcal{G}_1$, les applications
$b + \theta \circ d$ et $a$ définissent une application
$\mathcal{O}_{X'_1} \to \mathcal{O}_{X'_2}$ complétant
(\ref{defos-equation-huge-2}).

\medskip\noindent
Démonstration de (2) dans le cas particulier ci-dessus. Omise. Indication :
c'est exactement la même que la démonstration de (2) de la
Lemme \ref{defos-lemma-huge-diagram}.
\end{proof}

\begin{lemma}
\label{defos-lemma-NL-represent-ext-class}
Soit $X$ un espace topologique. Soit $\mathcal{A} \to \mathcal{B}$ un
homomorphisme de faisceaux d'anneaux. Soit $\mathcal{G}$ un
$\mathcal{B}$-module. Soit
$\xi \in \Ext^1_\mathcal{B}(\NL_{\mathcal{B}/\mathcal{A}}, \mathcal{G})$.
Il existe une application de faisceaux d'ensembles
$\alpha : \mathcal{E} \to \mathcal{B}$ telle que
$\xi \in \Ext^1_\mathcal{B}(\NL(\alpha), \mathcal{G})$
soit la classe d'une application
$\mathcal{I}/\mathcal{I}^2 \to \mathcal{G}$
(voir la démonstration pour les notations).
\end{lemma}

\begin{proof}
Rappelons que, si $\alpha : \mathcal{E} \to \mathcal{B}$ est telle que
$\mathcal{A}[\mathcal{E}] \to \mathcal{B}$ soit surjective, de noyau
$\mathcal{I}$, le complexe
$\NL(\alpha) = (\mathcal{I}/\mathcal{I}^2 \to 
\Omega_{\mathcal{A}[\mathcal{E}]/\mathcal{A}}
\otimes_{\mathcal{A}[\mathcal{E}]} \mathcal{B})$ est canoniquement
isomorphe à $\NL_{\mathcal{B}/\mathcal{A}}$ ; voir Modules, Lemme
\ref{modules-lemma-NL-up-to-qis}.
Remarquons de plus que
$\Omega = \Omega_{\mathcal{A}[\mathcal{E}]/\mathcal{A}}
\otimes_{\mathcal{A}[\mathcal{E}]} \mathcal{B}$ est le faisceau associé
au préfaisceau
$U \mapsto \bigoplus_{e \in \mathcal{E}(U)} \mathcal{B}(U)$.
Autrement dit, $\Omega$ est le $\mathcal{B}$-module libre sur le faisceau
d'ensembles $\mathcal{E}$ et il existe en particulier une application
canonique $\mathcal{E} \to \Omega$.

\medskip\noindent
Cela étant, choisissons un $\mathcal{E}$ (par exemple
$\mathcal{E} = \mathcal{B}$ comme dans la définition du complexe cotangent
naïf). L'obstruction à l'écriture de $\xi$ comme classe d'une application
$\mathcal{I}/\mathcal{I}^2 \to \mathcal{G}$ est un élément de
$\Ext^1_\mathcal{B}(\Omega, \mathcal{G})$. Supposons qu'il soit représenté
par l'extension
$0 \to \mathcal{G} \to \mathcal{H} \to \Omega \to 0$
de $\mathcal{B}$-modules. Considérons le faisceau d'ensembles
$\mathcal{E}' = \mathcal{E} \times_\Omega \mathcal{H}$,
muni de l'application induite $\alpha' : \mathcal{E}' \to \mathcal{B}$.
Posons $\mathcal{I}' = \Ker(\mathcal{A}[\mathcal{E}'] \to \mathcal{B})$
et $\Omega' = \Omega_{\mathcal{A}[\mathcal{E}']/\mathcal{A}}
\otimes_{\mathcal{A}[\mathcal{E}']} \mathcal{B}$.
L'image réciproque de $\xi$ par le quasi-isomorphisme
$\NL(\alpha') \to \NL(\alpha)$ s'envoie sur zéro dans
$\Ext^1_\mathcal{B}(\Omega', \mathcal{G})$, car l'image réciproque de
l'extension $\mathcal{H}$ par l'application $\Omega' \to \Omega$ est
scindée : en effet, $\Omega'$ est le $\mathcal{B}$-module libre sur le
faisceau d'ensembles $\mathcal{E}'$ et, par construction, il existe un
diagramme commutatif
$$
\xymatrix{
\mathcal{E}' \ar[r] \ar[d] & \mathcal{E} \ar[d] \\
\mathcal{H} \ar[r] & \Omega
}
$$
Cela termine la démonstration.
\end{proof}

\begin{lemma}
\label{defos-lemma-choices-ringed-spaces}
S'il existe une solution de (\ref{defos-equation-to-solve-ringed-spaces}),
l'ensemble des classes d'isomorphisme de solutions est un espace principal
homogène sous $\Ext^1_{\mathcal{O}_X}(\NL_{X/S}, \mathcal{G})$.
\end{lemma}

\begin{proof}
Remarquons tout de suite qu'étant données deux solutions $X'_1$ et $X'_2$
de (\ref{defos-equation-to-solve-ringed-spaces}), la
Lemme \ref{defos-lemma-huge-diagram-ringed-spaces} fournit un élément d'obstruction
$o(X'_1, X'_2) \in \Ext^1_{\mathcal{O}_X}(\NL_{X/S}, \mathcal{G})$
à l'existence d'une application $X'_1 \to X'_2$. Manifestement, cet élément
est l'obstruction à l'existence d'un isomorphisme, et sépare donc les classes
d'isomorphisme. Pour achever la démonstration, il suffit par conséquent de
montrer qu'étant donnés une solution $X'$ et un élément
$\xi \in \Ext^1_{\mathcal{O}_X}(\NL_{X/S}, \mathcal{G})$,
on peut trouver une seconde solution $X'_\xi$ telle que
$o(X', X'_\xi) = \xi$.

\medskip\noindent
Choisissons $\alpha : \mathcal{E} \to \mathcal{O}_X$ comme dans la
Lemme \ref{defos-lemma-NL-represent-ext-class} pour la classe $\xi$.
Considérons la surjection
$f^{-1}\mathcal{O}_S[\mathcal{E}] \to \mathcal{O}_X$
de noyau $\mathcal{I}$ et le complexe cotangent naïf correspondant
$\NL(\alpha) = (\mathcal{I}/\mathcal{I}^2 \to
\Omega_{f^{-1}\mathcal{O}_S[\mathcal{E}]/f^{-1}\mathcal{O}_S}
\otimes_{f^{-1}\mathcal{O}_S[\mathcal{E}]} \mathcal{O}_X)$.
D'après la lemme, $\xi$ est la classe d'un morphisme
$\delta : \mathcal{I}/\mathcal{I}^2 \to \mathcal{G}$.
Après avoir remplacé $\mathcal{E}$ par
$\mathcal{E} \times_{\mathcal{O}_X} \mathcal{O}_{X'}$, nous pouvons
également supposer que $\alpha$ se factorise par une application
$\alpha' : \mathcal{E} \to \mathcal{O}_{X'}$.

\medskip\noindent
Ces choix déterminent un homomorphisme de
$f^{-1}\mathcal{O}_{S'}$-algèbres
$\varphi : f^{-1}\mathcal{O}_{S'}[\mathcal{E}] \to \mathcal{O}_{X'}$.
Posons
$\mathcal{I}' = \Ker(f^{-1}\mathcal{O}_{S'}[\mathcal{E}] \to \mathcal{O}_X)$.
Remarquons que $\varphi$ induit une application
$\varphi|_{\mathcal{I}'} : \mathcal{I}' \to \mathcal{G}$
et que $\mathcal{O}_{X'}$ est la somme amalgamée, comme dans le diagramme
suivant
$$
\xymatrix{
0 \ar[r] & \mathcal{G} \ar[r] & \mathcal{O}_{X'} \ar[r] &
\mathcal{O}_X \ar[r] & 0 \\
0 \ar[r] & \mathcal{I}' \ar[u]^{\varphi|_{\mathcal{I}'}} \ar[r] &
f^{-1}\mathcal{O}_{S'}[\mathcal{E}] \ar[u] \ar[r] &
\mathcal{O}_X \ar[u]_{=} \ar[r] & 0
}
$$
Soit $\psi : \mathcal{I}' \to \mathcal{G}$ la somme de l'application
$\varphi|_{\mathcal{I}'}$ et du composé
$$
\mathcal{I}' \to \mathcal{I}'/(\mathcal{I}')^2 \to
\mathcal{I}/\mathcal{I}^2 \xrightarrow{\delta} \mathcal{G}.
$$
La somme amalgamée suivant $\psi$ est alors une autre extension d'anneaux
$\mathcal{O}_{X'_\xi}$ qui s'insère dans un diagramme comme ci-dessus.
Un calcul (omis) montre que $o(X', X'_\xi) = \xi$, comme voulu.
\end{proof}

\begin{lemma}
\label{defos-lemma-extensions-of-relative-ringed-spaces}
Soit $f : (X, \mathcal{O}_X) \to (S, \mathcal{O}_S)$ un morphisme
d'espaces annelés. Soit $\mathcal{G}$ un $\mathcal{O}_X$-module.
L'ensemble des classes d'isomorphisme des extensions de
$f^{-1}\mathcal{O}_S$-algèbres
$$
0 \to \mathcal{G} \to \mathcal{O}_{X'} \to \mathcal{O}_X \to 0
$$
où $\mathcal{G}$ est un idéal de carré nul\footnote{Autrement dit,
l'ensemble des classes d'isomorphisme des épaississements du premier ordre
$i : X \to X'$ au-dessus de $S$, munis d'un isomorphisme
$\mathcal{G} \to \Ker(i^\sharp)$ de $\mathcal{O}_X$-modules.}
est canoniquement en bijection avec
$\Ext^1_{\mathcal{O}_X}(\NL_{X/S}, \mathcal{G})$.
\end{lemma}

\begin{proof}
Pour le démontrer, appliquons les résultats précédents au cas où
(\ref{defos-equation-to-solve-ringed-spaces}) est donné par le diagramme
$$
\xymatrix{
0 \ar[r] &
\mathcal{G} \ar[r] &
{?} \ar[r] &
\mathcal{O}_X \ar[r] & 0 \\
0 \ar[r] &
0 \ar[u] \ar[r] &
\mathcal{O}_S \ar[u] \ar[r]^{\text{id}} &
\mathcal{O}_S \ar[u] \ar[r] & 0
}
$$
Le lemme résulte donc du Lemme \ref{defos-lemma-choices-ringed-spaces}
et de l'existence d'une solution, à savoir
$\mathcal{G} \oplus \mathcal{O}_X$.
(Voir la remarque ci-dessous pour une construction directe de la bijection.)
\end{proof}

\begin{remark}
\label{defos-remark-extensions-of-relative-ringed-spaces}
Soient $f : (X, \mathcal{O}_X) \to (S, \mathcal{O}_S)$ et
$\mathcal{G}$ comme dans la
Lemme \ref{defos-lemma-extensions-of-relative-ringed-spaces}.
Considérons une extension
$0 \to \mathcal{G} \to \mathcal{O}_{X'} \to \mathcal{O}_X \to 0$
comme dans la lemme. On peut choisir un faisceau d'ensembles $\mathcal{E}$
et un diagramme commutatif
$$
\xymatrix{
\mathcal{E} \ar[d]_{\alpha'} \ar[rd]^\alpha \\
\mathcal{O}_{X'} \ar[r] & \mathcal{O}_X
}
$$
tels que $f^{-1}\mathcal{O}_S[\mathcal{E}] \to \mathcal{O}_X$
soit surjective, de noyau $\mathcal{J}$.
(On peut par exemple prendre n'importe quel faisceau d'ensembles se
surjectant sur $\mathcal{O}_{X'}$.) Alors
$$
\NL_{X/S} \cong \NL(\alpha) = 
\left(
\mathcal{J}/\mathcal{J}^2
\longrightarrow
\Omega_{f^{-1}\mathcal{O}_S[\mathcal{E}]/f^{-1}\mathcal{O}_S}
\otimes_{f^{-1}\mathcal{O}_S[\mathcal{E}]} \mathcal{O}_X\right)
$$
Voir Modules, section \ref{modules-section-netherlander}, et en particulier
Lemme \ref{modules-lemma-NL-up-to-qis}. Bien entendu, $\alpha'$ détermine une
application $f^{-1}\mathcal{O}_S[\mathcal{E}] \to \mathcal{O}_{X'}$,
qui détermine à son tour une application
$$
\mathcal{J}/\mathcal{J}^2 \longrightarrow \mathcal{G}
$$
laquelle détermine à son tour l'élément de
$\Ext^1_{\mathcal{O}_X}(\NL(\alpha), \mathcal{G}) =
\Ext^1_{\mathcal{O}_X}(\NL_{X/S}, \mathcal{G})$
correspondant à $\mathcal{O}_{X'}$ par la bijection de la lemme.
\end{remark}

\begin{lemma}
\label{defos-lemma-extensions-of-relative-ringed-spaces-functorial}
Soient $f : (X, \mathcal{O}_X) \to (S, \mathcal{O}_S)$ et
$g : (Y, \mathcal{O}_Y) \to (X, \mathcal{O}_X)$ des morphismes
d'espaces annelés. Soit $\mathcal{F}$ un $\mathcal{O}_X$-module.
Soit $\mathcal{G}$ un $\mathcal{O}_Y$-module. Soit
$c : \mathcal{F} \to \mathcal{G}$ une $g$-application. Enfin, considérons
\begin{enumerate}
\item[(a)] $0 \to \mathcal{F} \to \mathcal{O}_{X'} \to \mathcal{O}_X \to 0$,
une extension de $f^{-1}\mathcal{O}_S$-algèbres,
correspondant à
$\xi \in \Ext^1_{\mathcal{O}_X}(\NL_{X/S}, \mathcal{F})$, et
\item[(b)] $0 \to \mathcal{G} \to \mathcal{O}_{Y'} \to \mathcal{O}_Y \to 0$,
une extension de $g^{-1}f^{-1}\mathcal{O}_S$-algèbres,
correspondant à
$\zeta \in \Ext^1_{\mathcal{O}_Y}(\NL_{Y/S}, \mathcal{G})$.
\end{enumerate}
Voir le Lemme \ref{defos-lemma-extensions-of-relative-ringed-spaces}.
Il existe alors un $S$-morphisme $g' : Y' \to X'$ compatible avec $g$ et
$c$ si et seulement si $\xi$ et $\zeta$ ont la même image dans
$\Ext^1_{\mathcal{O}_Y}(Lg^*\NL_{X/S}, \mathcal{G})$.
\end{lemma}

\begin{proof}
L'énoncé a bien un sens, car nous disposons des applications
$$
\Ext^1_{\mathcal{O}_X}(\NL_{X/S}, \mathcal{F}) \to
\Ext^1_{\mathcal{O}_Y}(Lg^*\NL_{X/S}, Lg^*\mathcal{F}) \to
\Ext^1_{\mathcal{O}_Y}(Lg^*\NL_{X/S}, \mathcal{G})
$$
obtenues à l'aide de l'application
$Lg^*\mathcal{F} \to g^*\mathcal{F} \xrightarrow{c} \mathcal{G}$,
ainsi que de l'application
$$
\Ext^1_{\mathcal{O}_Y}(\NL_{Y/S}, \mathcal{G}) \to
\Ext^1_{\mathcal{O}_Y}(Lg^*\NL_{X/S}, \mathcal{G})
$$
obtenue à l'aide de $Lg^*\NL_{X/S} \to \NL_{Y/S}$.
L'énoncé de la lemme se déduit de la
Lemme \ref{defos-lemma-huge-diagram-ringed-spaces} appliquée au diagramme
$$
\xymatrix{
& 0 \ar[r] &
\mathcal{G} \ar[r] &
\mathcal{O}_{Y'} \ar[r] &
\mathcal{O}_Y \ar[r] & 0 \\
& 0 \ar[r]|\hole & 0 \ar[u] \ar[r] &
\mathcal{O}_S \ar[u] \ar[r]|\hole &
\mathcal{O}_S \ar[u] \ar[r] & 0 \\
0 \ar[r] &
\mathcal{F} \ar[ruu] \ar[r] &
\mathcal{O}_{X'} \ar[r] &
\mathcal{O}_X \ar[ruu] \ar[r] & 0 \\
0 \ar[r] & 0 \ar[ruu]|\hole \ar[u] \ar[r] &
\mathcal{O}_S \ar[ruu]|\hole \ar[u] \ar[r] &
\mathcal{O}_S \ar[ruu]|\hole \ar[u] \ar[r] & 0
}
$$
et d'une compatibilité entre les constructions des démonstrations des
Lemmes \ref{defos-lemma-extensions-of-relative-ringed-spaces} et
\ref{defos-lemma-huge-diagram-ringed-spaces}, dont nous omettons l'énoncé et la
démonstration. (Voir la remarque ci-dessous pour un argument direct.)
\end{proof}

\begin{remark}
\label{defos-remark-extensions-of-relative-ringed-spaces-functorial}
Soient $f : (X, \mathcal{O}_X) \to (S, \mathcal{O}_S)$,
$g : (Y, \mathcal{O}_Y) \to (X, \mathcal{O}_X)$,
$\mathcal{F}$,
$\mathcal{G}$,
$c : \mathcal{F} \to \mathcal{G}$,
$0 \to \mathcal{F} \to \mathcal{O}_{X'} \to \mathcal{O}_X \to 0$,
$\xi \in \Ext^1_{\mathcal{O}_X}(\NL_{X/S}, \mathcal{F})$,
$0 \to \mathcal{G} \to \mathcal{O}_{Y'} \to \mathcal{O}_Y \to 0$
et $\zeta \in \Ext^1_{\mathcal{O}_Y}(\NL_{Y/S}, \mathcal{G})$
comme dans la
Lemme \ref{defos-lemma-extensions-of-relative-ringed-spaces-functorial}.
En prenant la somme amalgamée suivant
$c : g^{-1}\mathcal{F} \to \mathcal{G}$, on construit une extension
$$
\xymatrix{
0 \ar[r] &
\mathcal{G} \ar[r] &
\mathcal{O}'_1 \ar[r] &
g^{-1}\mathcal{O}_X \ar[r] & 0 \\
0 \ar[r] &
g^{-1}\mathcal{F} \ar[u]^c \ar[r] &
g^{-1}\mathcal{O}_{X'} \ar[u] \ar[r] &
g^{-1}\mathcal{O}_X \ar@{=}[u] \ar[r] & 0
}
$$
En prenant l'image réciproque suivant
$g^\sharp : g^{-1}\mathcal{O}_X \to \mathcal{O}_Y$,
on construit une extension
$$
\xymatrix{
0 \ar[r] &
\mathcal{G} \ar[r] &
\mathcal{O}_{Y'} \ar[r] &
\mathcal{O}_Y \ar[r] & 0 \\
0 \ar[r] &
\mathcal{G} \ar@{=}[u] \ar[r] &
\mathcal{O}'_2 \ar[u] \ar[r] &
g^{-1}\mathcal{O}_X \ar[u] \ar[r] & 0
}
$$
Une chasse au diagramme montre qu'il existe un $S$-morphisme
$Y' \to X'$ compatible avec $g$ et $c$ si et seulement si
$\mathcal{O}'_1$ est isomorphe à $\mathcal{O}'_2$ comme extension de
$g^{-1}f^{-1}\mathcal{O}_S$-algèbres de $g^{-1}\mathcal{O}_X$ par
$\mathcal{G}$. D'après la
Lemme \ref{defos-lemma-extensions-of-relative-ringed-spaces}, ces extensions sont
classifiées par le membre de gauche de
$$
\Ext^1_{g^{-1}\mathcal{O}_X}(
\NL_{g^{-1}\mathcal{O}_X/g^{-1}f^{-1}\mathcal{O}_S}, \mathcal{G}) =
\Ext^1_{\mathcal{O}_Y}(Lg^*\NL_{X/S}, \mathcal{G})
$$
Ici, l'égalité provient de l'adjonction tensor-hom et des égalités
$$
\NL_{g^{-1}\mathcal{O}_X/g^{-1}f^{-1}\mathcal{O}_S} = g^{-1}\NL_{X/S}
\quad\text{et}\quad
Lg^*\NL_{X/S} =
g^{-1}\NL_{X/S} \otimes_{g^{-1}\mathcal{O}_X}^\mathbf{L} \mathcal{O}_Y
$$
Pour la première, voir Modules, Lemme \ref{modules-lemma-pullback-NL} ;
la seconde résulte de la définition de l'image réciproque dérivée.
Ainsi, pour voir que la
Lemme \ref{defos-lemma-extensions-of-relative-ringed-spaces-functorial} est vraie,
il suffit de montrer que $\mathcal{O}'_1$ correspond à l'image de $\xi$ et
que $\mathcal{O}'_2$ correspond à l'image de $\zeta$.
La correspondance entre $\xi$ et $\mathcal{O}'_1$ résulte immédiatement de
la construction de la classe $\xi$ dans la
Remarque \ref{defos-remark-extensions-of-relative-ringed-spaces}.
Pour la correspondance entre $\zeta$ et $\mathcal{O}'_2$, choisissons d'abord
un diagramme commutatif
$$
\xymatrix{
\mathcal{E} \ar[d]_{\beta'} \ar[rd]^\beta \\
\mathcal{O}_{Y'} \ar[r] & \mathcal{O}_Y
}
$$
tel que $g^{-1}f^{-1}\mathcal{O}_S[\mathcal{E}] \to \mathcal{O}_Y$
soit surjective, de noyau $\mathcal{K}$. Choisissons ensuite un diagramme
commutatif
$$
\xymatrix{
\mathcal{E} \ar[d]_{\beta'} &
\mathcal{E}' \ar[l]^\varphi \ar[d]_{\alpha'} \ar[rd]^\alpha \\
\mathcal{O}_{Y'} &
\mathcal{O}'_2 \ar[l] \ar[r] &
g^{-1}\mathcal{O}_X
}
$$
tel que
$g^{-1}f^{-1}\mathcal{O}_S[\mathcal{E}'] \to g^{-1}\mathcal{O}_X$
soit surjective, de noyau $\mathcal{J}$. (Il suffit par exemple de prendre
$\mathcal{E}' = \mathcal{E} \amalg \mathcal{O}'_2$ comme faisceau
d'ensembles.) L'application $\varphi$ induit une application de complexes
$\NL(\alpha) \to \NL(\beta)$
(notations de Modules, section \ref{modules-section-netherlander})
et en particulier
$\bar\varphi : \mathcal{J}/\mathcal{J}^2 \to \mathcal{K}/\mathcal{K}^2$.
On a alors $\NL(\alpha) \cong \NL_{Y/S}$ et
$\NL(\beta) \cong \NL_{g^{-1}\mathcal{O}_X/g^{-1}f^{-1}\mathcal{O}_S}$,
et l'application de complexes $\NL(\alpha) \to \NL(\beta)$ représente
l'application $Lg^*\NL_{X/S} \to \NL_{Y/S}$ employée dans l'énoncé de la
Lemme \ref{defos-lemma-extensions-of-relative-ringed-spaces-functorial}
(voir la première partie de sa démonstration). Or $\zeta$ correspond à la
classe de l'application $\mathcal{K}/\mathcal{K}^2 \to \mathcal{G}$
induite par $\beta'$ ; voir la
Remarque \ref{defos-remark-extensions-of-relative-ringed-spaces}.
De même, l'extension $\mathcal{O}'_2$ correspond à l'application
$\mathcal{J}/\mathcal{J}^2 \to \mathcal{G}$ induite par $\alpha'$.
Le diagramme commutatif ci-dessus montre que cette application est le composé
de l'application $\mathcal{K}/\mathcal{K}^2 \to \mathcal{G}$ induite par
$\beta'$ et de l'application
$\bar\varphi : \mathcal{J}/\mathcal{J}^2 \to \mathcal{K}/\mathcal{K}^2$.
C'est la compatibilité recherchée.
\end{remark}

\begin{lemma}
\label{defos-lemma-parametrize-solutions-ringed-spaces}
Soient $t : (S, \mathcal{O}_S) \to (S', \mathcal{O}_{S'})$,
$\mathcal{J} = \Ker(t^\sharp)$,
$f : (X, \mathcal{O}_X) \to (S, \mathcal{O}_S)$, $\mathcal{G}$ et
$c : \mathcal{J} \to \mathcal{G}$ comme dans
(\ref{defos-equation-to-solve-ringed-spaces}).
Notons $\xi \in \Ext^1_{\mathcal{O}_S}(\NL_{S/S'}, \mathcal{J})$
l'élément correspondant, par la
Lemme \ref{defos-lemma-extensions-of-relative-ringed-spaces}, à l'extension
$\mathcal{O}_{S'}$ de $\mathcal{O}_S$ par $\mathcal{J}$.
L'ensemble des classes d'isomorphisme de solutions est canoniquement en
bijection avec la fibre de
$$
\Ext^1_{\mathcal{O}_X}(\NL_{X/S'}, \mathcal{G}) \to
\Ext^1_{\mathcal{O}_X}(Lf^*\NL_{S/S'}, \mathcal{G})
$$
au-dessus de l'image de $\xi$.
\end{lemma}

\begin{proof}
D'après le Lemme \ref{defos-lemma-extensions-of-relative-ringed-spaces}, appliqué
à $X \to S'$ et au $\mathcal{O}_X$-module $\mathcal{G}$, les éléments
$\zeta$ de
$\Ext^1_{\mathcal{O}_X}(\NL_{X/S'}, \mathcal{G})$
paramètrent les extensions
$0 \to \mathcal{G} \to \mathcal{O}_{X'} \to \mathcal{O}_X \to 0$
de $f^{-1}\mathcal{O}_{S'}$-algèbres.
D'après la
Lemme \ref{defos-lemma-extensions-of-relative-ringed-spaces-functorial}, appliquée
à $X \to S \to S'$ et à $c : \mathcal{J} \to \mathcal{G}$, il existe un
$S'$-morphisme $X' \to S'$ compatible avec $c$ et
$f : X \to S$ si et seulement si $\zeta$ s'envoie sur $\xi$.
Bien entendu, cela revient à dire que $\mathcal{O}_{X'}$ est une solution de
(\ref{defos-equation-to-solve-ringed-spaces}).
\end{proof}

\begin{remark}
\label{defos-remark-parametrize-solutions-ringed-spaces}
Dans la situation de la
Lemme \ref{defos-lemma-parametrize-solutions-ringed-spaces}, nous avons des
applications de complexes
$$
Lf^*\NL_{S'/S} \to \NL_{X/S'} \to \NL_{X/S}
$$
Ces applications sont proches de former un triangle distingué ; voir
Modules, Lemme \ref{modules-lemma-exact-sequence-NL-ringed-topoi}.
Si elles formaient un triangle distingué, on en conclurait que l'image de
$\xi$ dans $\Ext^2_{\mathcal{O}_X}(\NL_{X/S}, \mathcal{G})$
serait l'obstruction à l'existence d'une solution de
(\ref{defos-equation-to-solve-ringed-spaces}).
\end{remark}

\section{Déformations de schémas}
\label{defos-section-deformations-schemes}

\noindent
Dans cette section, nous explicitons la signification des résultats de la
section \ref{defos-section-deformations-ringed-spaces} pour les déformations de
schémas.

\begin{lemma}
\label{defos-lemma-deform}
Soit $S \subset S'$ un épaississement du premier ordre de schémas.
Soit $f : X \to S$ un morphisme plat de schémas.
S'il existe un morphisme plat $f' : X' \to S'$ de schémas et un isomorphisme
$a : X \to X' \times_{S'} S$ au-dessus de $S$, alors
\begin{enumerate}
\item l'ensemble des classes d'isomorphisme de couples
$(f' : X' \to S', a)$ est un ensemble principal homogène sous
$\Ext^1_{\mathcal{O}_X}(\NL_{X/S}, f^*\mathcal{C}_{S/S'})$, et
\item l'ensemble des automorphismes $\varphi : X' \to X'$ au-dessus de $S'$
qui se réduisent à l'identité sur $X' \times_{S'} S$ est
$\Ext^0_{\mathcal{O}_X}(\NL_{X/S}, f^*\mathcal{C}_{S/S'})$.
\end{enumerate}
\end{lemma}

\begin{proof}
Observons d'abord que les épaississements de schémas tels qu'ils sont définis
dans Morphismes, section \ref{more-morphisms-section-thickenings}, ne sont
autres que les morphismes de schémas qui sont des épaississements au sens de la
section \ref{defos-section-thickenings-spaces}.
Nous pouvons considérer $X$ comme un sous-schéma fermé de $X'$ de sorte que
$(f, f') : (X \subset X') \to (S \subset S')$ soit un morphisme
d'épaississements du premier ordre. Il résulte alors de Morphismes,
Lemme \ref{more-morphisms-lemma-deform} (ou du
Lemme \ref{defos-lemma-deform-module}, plus général) que le faisceau d'idéaux de $X$
dans $X'$ est égal à $f^*\mathcal{C}_{S/S'}$.
Nous avons donc un diagramme commutatif
$$
\xymatrix{
0 \ar[r] & f^*\mathcal{C}_{S/S'} \ar[r] &
\mathcal{O}_{X'} \ar[r] &
\mathcal{O}_X \ar[r] & 0 \\
0 \ar[r] & \mathcal{C}_{S/S'} \ar[u] \ar[r] &
\mathcal{O}_{S'} \ar[u] \ar[r] &
\mathcal{O}_S \ar[u] \ar[r] & 0
}
$$
où les flèches verticales sont des $f$-applications ; comparer avec
(\ref{defos-equation-to-solve-ringed-spaces}).
La partie (1) découle donc du Lemme \ref{defos-lemma-choices-ringed-spaces}, et la
partie (2), de la partie (2) du
Lemme \ref{defos-lemma-huge-diagram-ringed-spaces}.
(Notons que $\NL_{X/S}$, tel qu'il est défini pour un morphisme de schémas dans
Morphismes, section \ref{more-morphisms-section-netherlander}, coïncide avec
$\NL_{X/S}$ tel qu'il est employé dans la
section \ref{defos-section-deformations-ringed-spaces}.)
\end{proof}

\section{Épaississements de topos annelés}
\label{defos-section-thickenings-ringed-topoi}

\noindent
Cette section est l'analogue de la section \ref{defos-section-thickenings-spaces}
pour les topos annelés.
Dans les quelques sections qui suivent, nous utiliserons les notions suivantes :
\begin{enumerate}
\item Un faisceau d'idéaux $\mathcal{I} \subset \mathcal{O}'$ sur un topos
annelé $(\Sh(\mathcal{D}), \mathcal{O}')$ est {\it localement nilpotent} si
toute section locale de $\mathcal{I}$ est localement nilpotente.
\item Un {\it épaississement} de topos annelés est un morphisme
$i : (\Sh(\mathcal{C}), \mathcal{O}) \to (\Sh(\mathcal{D}), \mathcal{O}')$
de topos annelés tel que
\begin{enumerate}
\item $i_*$ soit une équivalence $\Sh(\mathcal{C}) \to \Sh(\mathcal{D})$,
\item l'application $i^\sharp : \mathcal{O}' \to i_*\mathcal{O}$ soit
surjective, et
\item le noyau de $i^\sharp$ soit un faisceau d'idéaux localement nilpotent.
\end{enumerate}
\item Un {\it épaississement du premier ordre} de topos annelés est un
épaississement
$i : (\Sh(\mathcal{C}), \mathcal{O}) \to (\Sh(\mathcal{D}), \mathcal{O}')$
de topos annelés tel que $\Ker(i^\sharp)$ soit de carré nul.
\item On voit clairement comment définir les
{\it morphismes d'épaississements de topos annelés}, les
{\it morphismes d'épaississements de topos annelés au-dessus d'un topos annelé
de base}, etc.
\end{enumerate}
Si
$i : (\Sh(\mathcal{C}), \mathcal{O}) \to (\Sh(\mathcal{D}), \mathcal{O}')$
est un épaississement de topos annelés, nous identifions les topos sous-jacents
et considérons $\mathcal{O}$, $\mathcal{O}'$ et
$\mathcal{I} = \Ker(i^\sharp)$ comme des faisceaux sur $\mathcal{C}$.
Nous obtenons une suite exacte courte
$$
0 \to \mathcal{I} \to \mathcal{O}' \to \mathcal{O} \to 0
$$
de $\mathcal{O}'$-modules. D'après Modules sur les sites,
Lemme \ref{sites-modules-lemma-i-star-equivalence}, la catégorie des
$\mathcal{O}$-modules est équivalente à la catégorie des
$\mathcal{O}'$-modules annulés par $\mathcal{I}$. En particulier, si $i$ est un
épaississement du premier ordre, alors $\mathcal{I}$ est un
$\mathcal{O}$-module.

\begin{situation}
\label{defos-situation-morphism-thickenings-ringed-topoi}
Un morphisme d'épaississements de topos annelés $(f, f')$ est donné par un
diagramme commutatif
\begin{equation}
\label{defos-equation-morphism-thickenings-ringed-topoi}
\vcenter{
\xymatrix{
(\Sh(\mathcal{C}), \mathcal{O}) \ar[r]_i \ar[d]_f &
(\Sh(\mathcal{D}), \mathcal{O}') \ar[d]^{f'} \\
(\Sh(\mathcal{B}), \mathcal{O}_\mathcal{B}) \ar[r]^t &
(\Sh(\mathcal{B}'), \mathcal{O}_{\mathcal{B}'})
}
}
\end{equation}
de topos annelés dont les flèches horizontales sont des épaississements. Dans
cette situation, nous posons
$\mathcal{I} = \Ker(i^\sharp) \subset \mathcal{O}'$ et
$\mathcal{J} = \Ker(t^\sharp) \subset \mathcal{O}_{\mathcal{B}'}$.
Comme $f = f'$ sur les topos sous-jacents, nous identifierons les foncteurs
image réciproque $f^{-1}$ et $(f')^{-1}$.
Observons que
$(f')^\sharp : f^{-1}\mathcal{O}_{\mathcal{B}'} \to \mathcal{O}'$
induit en particulier une application $f^{-1}\mathcal{J} \to \mathcal{I}$,
et donc une application de $\mathcal{O}'$-modules
$$
(f')^*\mathcal{J} \longrightarrow \mathcal{I}
$$
Si $i$ et $t$ sont des épaississements du premier ordre, alors
$(f')^*\mathcal{J} = f^*\mathcal{J}$ et l'application ci-dessus devient une
application $f^*\mathcal{J} \to \mathcal{I}$.
\end{situation}

\begin{definition}
\label{defos-definition-strict-morphism-thickenings-ringed-topoi}
Dans la Situation \ref{defos-situation-morphism-thickenings-ringed-topoi}, nous
disons que $(f, f')$ est un {\it morphisme strict d'épaississements} si
l'application $(f')^*\mathcal{J} \longrightarrow \mathcal{I}$ est surjective.
\end{definition}

\section{Modules sur les épaississements du premier ordre de topos annelés}
\label{defos-section-modules-thickenings-ringed-topoi}

\noindent
Dans cette section, nous présentons quelques préliminaires à la théorie des
déformations des modules. Soit
$i : (\Sh(\mathcal{C}, \mathcal{O}) \to (\Sh(\mathcal{D}), \mathcal{O}')$
un épaississement du premier ordre de topos annelés. Nous utiliserons librement
les notations introduites dans la
section \ref{defos-section-thickenings-ringed-topoi} ; en particulier, nous
identifierons les topos sous-jacents.
Dans cette section, nous considérons des suites exactes courtes
\begin{equation}
\label{defos-equation-extension-ringed-topoi}
0 \to \mathcal{K} \to \mathcal{F}' \to \mathcal{F} \to 0
\end{equation}
de $\mathcal{O}'$-modules, où $\mathcal{F}$ et $\mathcal{K}$ sont des
$\mathcal{O}$-modules, et $\mathcal{F}'$ est un $\mathcal{O}'$-module.
Dans cette situation, nous avons une application canonique de
$\mathcal{O}$-modules
$$
c_{\mathcal{F}'} :
\mathcal{I} \otimes_\mathcal{O} \mathcal{F}
\longrightarrow
\mathcal{K}
$$
où $\mathcal{I} = \Ker(i^\sharp)$.
En effet, étant données des sections locales $f$ de $\mathcal{I}$ et $s$ de
$\mathcal{F}$, nous posons $c_{\mathcal{F}'}(f \otimes s) = fs'$, où $s'$ est
une section locale de $\mathcal{F}'$ relevant $s$.

\begin{lemma}
\label{defos-lemma-inf-map-ringed-topoi}
Soit $i : (\Sh(\mathcal{C}), \mathcal{O}) \to (\Sh(\mathcal{D}), \mathcal{O}')$
un épaississement du premier ordre de topos annelés. Supposons données des
extensions
$$
0 \to \mathcal{K} \to \mathcal{F}' \to \mathcal{F} \to 0
\quad\text{et}\quad
0 \to \mathcal{L} \to \mathcal{G}' \to \mathcal{G} \to 0
$$
comme dans (\ref{defos-equation-extension-ringed-topoi}), ainsi que des applications
$\varphi : \mathcal{F} \to \mathcal{G}$ et
$\psi : \mathcal{K} \to \mathcal{L}$.
\begin{enumerate}
\item S'il existe une application de $\mathcal{O}'$-modules
$\varphi' : \mathcal{F}' \to \mathcal{G}'$ compatible avec $\varphi$ et
$\psi$, alors le diagramme
$$
\xymatrix{
\mathcal{I} \otimes_\mathcal{O} \mathcal{F}
\ar[r]_-{c_{\mathcal{F}'}} \ar[d]_{1 \otimes \varphi} &
\mathcal{K} \ar[d]^\psi \\
\mathcal{I} \otimes_\mathcal{O} \mathcal{G}
\ar[r]^-{c_{\mathcal{G}'}} &
\mathcal{L}
}
$$
est commutatif.
\item L'ensemble des applications de $\mathcal{O}'$-modules
$\varphi' : \mathcal{F}' \to \mathcal{G}'$ compatibles avec $\varphi$ et
$\psi$ est, s'il est non vide, un espace principal homogène sous
$\Hom_\mathcal{O}(\mathcal{F}, \mathcal{L})$.
\end{enumerate}
\end{lemma}

\begin{proof}
La partie (1) résulte immédiatement de la description des applications.
Pour (2), si $\varphi'$ et $\varphi''$ sont deux applications
$\mathcal{F}' \to \mathcal{G}'$ compatibles avec $\varphi$ et $\psi$, alors
$\varphi' - \varphi''$ se factorise sous la forme
$$
\mathcal{F}' \to \mathcal{F} \to \mathcal{L} \to \mathcal{G}'
$$
L'application du milieu provient d'un unique élément de
$\Hom_\mathcal{O}(\mathcal{F}, \mathcal{L})$, d'après Modules sur les sites,
Lemme \ref{sites-modules-lemma-i-star-equivalence}.
Réciproquement, étant donné un élément $\alpha$ de ce groupe, nous pouvons
ajouter la composée affichée ci-dessus, avec $\alpha$ au milieu, à $\varphi'$.
Nous omettons certains détails.
\end{proof}

\begin{lemma}
\label{defos-lemma-inf-obs-map-ringed-topoi}
Soit $i : (\Sh(\mathcal{C}), \mathcal{O}) \to (\Sh(\mathcal{D}), \mathcal{O}')$
un épaississement du premier ordre de topos annelés. Supposons données des
extensions
$$
0 \to \mathcal{K} \to \mathcal{F}' \to \mathcal{F} \to 0
\quad\text{et}\quad
0 \to \mathcal{L} \to \mathcal{G}' \to \mathcal{G} \to 0
$$
comme dans (\ref{defos-equation-extension-ringed-topoi}), ainsi que des applications
$\varphi : \mathcal{F} \to \mathcal{G}$ et
$\psi : \mathcal{K} \to \mathcal{L}$. Supposons le diagramme
$$
\xymatrix{
\mathcal{I} \otimes_\mathcal{O} \mathcal{F}
\ar[r]_-{c_{\mathcal{F}'}} \ar[d]_{1 \otimes \varphi} &
\mathcal{K} \ar[d]^\psi \\
\mathcal{I} \otimes_\mathcal{O} \mathcal{G}
\ar[r]^-{c_{\mathcal{G}'}} &
\mathcal{L}
}
$$
commutatif. Il existe alors un élément
$$
o(\varphi, \psi) \in
\Ext^1_\mathcal{O}(\mathcal{F}, \mathcal{L})
$$
dont l'annulation est une condition nécessaire et suffisante pour qu'il existe
une application $\varphi' : \mathcal{F}' \to \mathcal{G}'$ compatible avec
$\varphi$ et $\psi$.
\end{lemma}

\begin{proof}
Nous pouvons construire explicitement une extension
$$
0 \to \mathcal{L} \to \mathcal{H} \to \mathcal{F} \to 0
$$
en prenant pour $\mathcal{H}$ la cohomologie au milieu du complexe
$$
\mathcal{K}
\xrightarrow{1, - \psi}
\mathcal{F}' \oplus \mathcal{G}' \xrightarrow{\varphi, 1}
\mathcal{G}
$$
(avec les notations évidentes). Un calcul sur les sections locales utilisant
l'hypothèse de commutativité du diagramme du lemme montre que $\mathcal{H}$ est
annulé par $\mathcal{I}$. Par conséquent, $\mathcal{H}$ définit une classe dans
$$
\Ext^1_\mathcal{O}(\mathcal{F}, \mathcal{L})
\subset
\Ext^1_{\mathcal{O}'}(\mathcal{F}, \mathcal{L})
$$
Enfin, la classe de $\mathcal{H}$ est la différence entre la somme amalgamée de
l'extension $\mathcal{F}'$ suivant $\psi$ et l'image réciproque de l'extension
$\mathcal{G}'$ suivant $\varphi$ (calculs omis).
L'annulation de la classe de $\mathcal{H}$ équivaut donc à l'existence d'un
diagramme commutatif
$$
\xymatrix{
0 \ar[r] &
\mathcal{K} \ar[r] \ar[d]_{\psi} &
\mathcal{F}' \ar[r] \ar[d]_{\varphi'} &
\mathcal{F} \ar[r] \ar[d]_\varphi & 0\\
0 \ar[r] &
\mathcal{L} \ar[r] &
\mathcal{G}' \ar[r] &
\mathcal{G} \ar[r] & 0
}
$$
comme voulu.
\end{proof}

\begin{lemma}
\label{defos-lemma-inf-ext-ringed-topoi}
Soit $i : (\Sh(\mathcal{C}), \mathcal{O}) \to (\Sh(\mathcal{D}), \mathcal{O}')$
un épaississement du premier ordre de topos annelés. Supposons donnés des
$\mathcal{O}$-modules $\mathcal{F}$, $\mathcal{K}$ et une application
$\mathcal{O}$-linéaire
$c : \mathcal{I} \otimes_\mathcal{O} \mathcal{F} \to \mathcal{K}$.
S'il existe une suite (\ref{defos-equation-extension-ringed-topoi}) telle que
$c_{\mathcal{F}'} = c$, alors l'ensemble des classes d'isomorphisme de ces
extensions est un espace principal homogène sous
$\Ext^1_\mathcal{O}(\mathcal{F}, \mathcal{K})$.
\end{lemma}

\begin{proof}
Supposons données des extensions
$$
0 \to \mathcal{K} \to \mathcal{F}'_1 \to \mathcal{F} \to 0
\quad\text{et}\quad
0 \to \mathcal{K} \to \mathcal{F}'_2 \to \mathcal{F} \to 0
$$
telles que $c_{\mathcal{F}'_1} = c_{\mathcal{F}'_2} = c$. Leur différence
(dans le groupe des extensions ; voir Homologie,
section \ref{homology-section-extensions}) est une extension
$$
0 \to \mathcal{K} \to \mathcal{E} \to \mathcal{F} \to 0
$$
où $\mathcal{E}$ est annulé par $\mathcal{I}$ (on omet le calcul local).
La suite est donc une extension de $\mathcal{O}$-modules ; voir Modules sur les
sites, Lemme \ref{sites-modules-lemma-i-star-equivalence}.
Réciproquement, étant donnée une telle extension $\mathcal{E}$, nous pouvons
ajouter l'extension $\mathcal{E}$ à l'extension de $\mathcal{O}'$-modules
$\mathcal{F}'$ sans modifier l'application $c_{\mathcal{F}'}$. Nous omettons
certains détails.
\end{proof}

\begin{lemma}
\label{defos-lemma-inf-obs-ext-ringed-topoi}
Soit $i : (\Sh(\mathcal{C}), \mathcal{O}) \to (\Sh(\mathcal{D}), \mathcal{O}')$
un épaississement du premier ordre de topos annelés. Supposons donnés des
$\mathcal{O}$-modules $\mathcal{F}$, $\mathcal{K}$ et une application
$\mathcal{O}$-linéaire
$c : \mathcal{I} \otimes_\mathcal{O} \mathcal{F} \to \mathcal{K}$.
Il existe alors un élément
$$
o(\mathcal{F}, \mathcal{K}, c) \in
\Ext^2_\mathcal{O}(\mathcal{F}, \mathcal{K})
$$
dont l'annulation est une condition nécessaire et suffisante pour qu'il existe
une suite (\ref{defos-equation-extension-ringed-topoi}) telle que
$c_{\mathcal{F}'} = c$.
\end{lemma}

\begin{proof}
Montrons d'abord que, si $\mathcal{K}$ est un $\mathcal{O}$-module injectif,
il existe bien une suite (\ref{defos-equation-extension-ringed-topoi}) telle que
$c_{\mathcal{F}'} = c$. Pour cela, choisissons un $\mathcal{O}'$-module plat
$\mathcal{H}'$ et une surjection $\mathcal{H}' \to \mathcal{F}$
(Modules sur les sites,
Lemme \ref{sites-modules-lemma-module-quotient-flat}).
Soit $\mathcal{J} \subset \mathcal{H}'$ le noyau. Comme $\mathcal{H}'$ est
plat, nous avons
$$
\mathcal{I} \otimes_{\mathcal{O}'} \mathcal{H}' =
\mathcal{I}\mathcal{H}'
\subset \mathcal{J} \subset \mathcal{H}'
$$
Observons que l'application
$$
\mathcal{I}\mathcal{H}' =
\mathcal{I} \otimes_{\mathcal{O}'} \mathcal{H}'
\longrightarrow
\mathcal{I} \otimes_{\mathcal{O}'} \mathcal{F} =
\mathcal{I} \otimes_\mathcal{O} \mathcal{F}
$$
annule $\mathcal{I}\mathcal{J}$. En effet, si $f$ est une section locale de
$\mathcal{I}$ et $s$ une section locale de $\mathcal{H}$, alors $fs$ s'envoie
sur $f \otimes \overline{s}$, où $\overline{s}$ est l'image de $s$ dans
$\mathcal{F}$. Nous obtenons ainsi
$$
\xymatrix{
\mathcal{I}\mathcal{H}'/\mathcal{I}\mathcal{J}
\ar@{^{(}->}[r] \ar[d] &
\mathcal{J}/\mathcal{I}\mathcal{J} \ar@{..>}[d]_\gamma \\
\mathcal{I} \otimes_\mathcal{O} \mathcal{F} \ar[r]^-c &
\mathcal{K}
}
$$
un diagramme de $\mathcal{O}$-modules. Si $\mathcal{K}$ est injectif comme
$\mathcal{O}$-module, nous obtenons la flèche pointillée.
Notons $\gamma' : \mathcal{J} \to \mathcal{K}$ la composée de $\gamma$ avec
$\mathcal{J} \to \mathcal{J}/\mathcal{I}\mathcal{J}$.
Un calcul local montre que la somme amalgamée
$$
\xymatrix{
0 \ar[r] &
\mathcal{J} \ar[r] \ar[d]_{\gamma'} &
\mathcal{H}' \ar[r] \ar[d] &
\mathcal{F} \ar[r] \ar@{=}[d] &
0 \\
0 \ar[r] &
\mathcal{K} \ar[r] &
\mathcal{F}' \ar[r] &
\mathcal{F} \ar[r] &
0
}
$$
est une solution au problème posé par le lemme.

\medskip\noindent
Cas général. Choisissons un plongement $\mathcal{K} \subset \mathcal{K}'$,
où $\mathcal{K}'$ est un $\mathcal{O}$-module injectif. Soit $\mathcal{Q}$ le
quotient, de sorte que nous ayons une suite exacte
$$
0 \to \mathcal{K} \to \mathcal{K}' \to \mathcal{Q} \to 0
$$
Notons
$c' : \mathcal{I} \otimes_\mathcal{O} \mathcal{F} \to \mathcal{K}'$
la composée. Le paragraphe précédent fournit une suite
$$
0 \to \mathcal{K}' \to \mathcal{E}' \to \mathcal{F} \to 0
$$
comme dans (\ref{defos-equation-extension-ringed-topoi}), telle que
$c_{\mathcal{E}'} = c'$.
Remarquons que la composée de $c'$ avec l'application
$\mathcal{K}' \to \mathcal{Q}$ est nulle ; la somme amalgamée de
$\mathcal{E}'$ suivant $\mathcal{K}' \to \mathcal{Q}$ est donc une extension
$$
0 \to \mathcal{Q} \to \mathcal{D}' \to \mathcal{F} \to 0
$$
comme dans (\ref{defos-equation-extension-ringed-topoi}), telle que
$c_{\mathcal{D}'} = 0$.
Cela signifie exactement que $\mathcal{D}'$ est annulé par $\mathcal{I}$ ;
autrement dit, $\mathcal{D}'$ est une extension de $\mathcal{O}$-modules,
c'est-à-dire qu'il définit un élément
$$
o(\mathcal{F}, \mathcal{K}, c) \in
\Ext^1_\mathcal{O}(\mathcal{F}, \mathcal{Q}) =
\Ext^2_\mathcal{O}(\mathcal{F}, \mathcal{K})
$$
(l'égalité résulte de la suite exacte longue de cohomologie associée à la suite
exacte précédente et de l'annulation des groupes Ext supérieurs à valeurs dans
le module injectif $\mathcal{K}'$). Si
$o(\mathcal{F}, \mathcal{K}, c) = 0$, nous pouvons choisir un scindage
$s : \mathcal{F} \to \mathcal{D}'$ et poser
$$
\mathcal{F}' = \Ker(\mathcal{E}' \to \mathcal{D}'/s(\mathcal{F}))
$$
de sorte que nous obtenions le diagramme suivant
$$
\xymatrix{
0 \ar[r] &
\mathcal{K} \ar[r] \ar[d] &
\mathcal{F}' \ar[r] \ar[d] &
\mathcal{F} \ar[r] \ar@{=}[d] &
0 \\
0 \ar[r] &
\mathcal{K}' \ar[r] &
\mathcal{E}' \ar[r] &
\mathcal{F} \ar[r] & 0
}
$$
à lignes exactes, qui montre que $c_{\mathcal{F}'} = c$. Réciproquement, si
$\mathcal{F}'$ existe, alors la somme amalgamée de $\mathcal{F}'$ suivant
l'application $\mathcal{K} \to \mathcal{K}'$ est isomorphe à $\mathcal{E}'$,
d'après le Lemme \ref{defos-lemma-inf-ext-ringed-topoi} et l'annulation des groupes
Ext supérieurs à valeurs dans le module injectif $\mathcal{K}'$.
Nous obtenons ainsi un diagramme comme ci-dessus, ce qui implique que
$\mathcal{D}'$ est scindée comme extension, c'est-à-dire que la classe
$o(\mathcal{F}, \mathcal{K}, c)$ est nulle.
\end{proof}

\begin{remark}
\label{defos-remark-trivial-thickening-ringed-topoi}
Soit $(\Sh(\mathcal{C}), \mathcal{O})$ un topos annelé. Un épaississement du
premier ordre $i : (\Sh(\mathcal{C}), \mathcal{O}) \to
(\Sh(\mathcal{D}), \mathcal{O}')$ est dit {\it trivial} s'il existe un
morphisme de topos annelés
$\pi : (\Sh(\mathcal{D}), \mathcal{O}') \to (\Sh(\mathcal{C}), \mathcal{O})$
qui soit un inverse à gauche de $i$. Le choix d'un tel morphisme $\pi$ est
appelé une {\it trivialisation} de l'épaississement du premier ordre.
Étant donné $\pi$, nous obtenons un scindage
\begin{equation}
\label{defos-equation-splitting-ringed-topoi}
\mathcal{O}' = \mathcal{O} \oplus \mathcal{I}
\end{equation}
comme faisceaux d'algèbres sur $\mathcal{C}$, en utilisant $\pi^\sharp$ pour
scinder la surjection $\mathcal{O}' \to \mathcal{O}$.
Réciproquement, un tel scindage détermine un morphisme $\pi$. La catégorie des
épaississements du premier ordre trivialisés de
$(\Sh(\mathcal{C}), \mathcal{O})$ est équivalente à la catégorie des
$\mathcal{O}$-modules.
\end{remark}

\begin{remark}
\label{defos-remark-trivial-extension-ringed-topoi}
Soit $i : (\Sh(\mathcal{C}), \mathcal{O}) \to (\Sh(\mathcal{D}), \mathcal{O}')$
un épaississement trivial du premier ordre de topos annelés, et soit
$\pi : (\Sh(\mathcal{D}), \mathcal{O}') \to
(\Sh(\mathcal{C}), \mathcal{O})$ une trivialisation. Pour tout triplet
$(\mathcal{F}, \mathcal{K}, c)$ formé d'une paire de $\mathcal{O}$-modules et
d'une application
$c : \mathcal{I} \otimes_\mathcal{O} \mathcal{F} \to \mathcal{K}$,
nous pouvons poser
$$
\mathcal{F}'_{c, triv} = \mathcal{F} \oplus \mathcal{K}
$$
et utiliser le scindage (\ref{defos-equation-splitting-ringed-topoi}) associé à
$\pi$, ainsi que l'application $c$, pour définir la structure de
$\mathcal{O}'$-module et obtenir une extension
(\ref{defos-equation-extension-ringed-topoi}).
Nous appellerons $\mathcal{F}'_{c, triv}$ l'{\it extension triviale} de
$\mathcal{F}$ par $\mathcal{K}$ correspondant à $c$ et à la trivialisation
$\pi$. Pour toute extension $\mathcal{F}'$ comme dans
(\ref{defos-equation-extension-ringed-topoi}), nous pouvons utiliser
$\pi^\sharp : \mathcal{O} \to \mathcal{O}'$ pour considérer $\mathcal{F}'$
comme une extension de $\mathcal{O}$-modules, et donc comme une classe
$\xi_{\mathcal{F}'}$ dans
$\Ext^1_\mathcal{O}(\mathcal{F}, \mathcal{K})$.
Le Lemme \ref{defos-lemma-inf-ext-ringed-topoi} assure que
$\mathcal{F}' \mapsto \xi_{\mathcal{F}'}$ induit une bijection
$$
\left\{
\begin{matrix}
\text{isomorphism classes of extensions}\\
\mathcal{F}'\text{ comme dans (\ref{defos-equation-extension-ringed-topoi}) avec }
c = c_{\mathcal{F}'}
\end{matrix}
\right\}
\longrightarrow
\Ext^1_\mathcal{O}(\mathcal{F}, \mathcal{K})
$$
De plus, l'extension triviale $\mathcal{F}'_{c, triv}$ s'envoie sur la classe
nulle.
\end{remark}

\begin{remark}
\label{defos-remark-extension-functorial-ringed-topoi}
Soit $(\Sh(\mathcal{C}), \mathcal{O})$ un topos annelé. Soient
$(\Sh(\mathcal{C}), \mathcal{O}) \to (\Sh(\mathcal{D}_i), \mathcal{O}'_i)$,
$i = 1, 2$, des épaississements du
premier ordre de faisceaux d'idéaux $\mathcal{I}_i$.
Soit $h : (\Sh(\mathcal{D}_1), \mathcal{O}'_1) \to
(\Sh(\mathcal{D}_2), \mathcal{O}'_2)$ un morphisme d'épaississements du
premier ordre de $(\Sh(\mathcal{C}), \mathcal{O})$.
Voici le diagramme
$$
\xymatrix{
& (\Sh(\mathcal{C}), \mathcal{O}) \ar[ld] \ar[rd] & \\
(\Sh(\mathcal{D}_1), \mathcal{O}'_1) \ar[rr]^h & & 
(\Sh(\mathcal{D}_2), \mathcal{O}'_2)
}
$$
Remarquons que $h^\sharp : \mathcal{O}'_2 \to \mathcal{O}'_1$ induit en
particulier une application de $\mathcal{O}$-modules
$\mathcal{I}_2 \to \mathcal{I}_1$.
Soit $\mathcal{F}$ un $\mathcal{O}$-module.
Soit $(\mathcal{K}_i, c_i)$, $i = 1, 2$, une paire formée d'un
$\mathcal{O}$-module $\mathcal{K}_i$ et d'une application
$c_i : \mathcal{I}_i \otimes_\mathcal{O} \mathcal{F} \to
\mathcal{K}_i$. Supposons en outre donnée une application de
$\mathcal{O}$-modules $\mathcal{K}_2 \to \mathcal{K}_1$ telle que le diagramme
$$
\xymatrix{
\mathcal{I}_2 \otimes_\mathcal{O} \mathcal{F}
\ar[r]_-{c_2} \ar[d] &
\mathcal{K}_2 \ar[d] \\
\mathcal{I}_1 \otimes_\mathcal{O} \mathcal{F}
\ar[r]^-{c_1} &
\mathcal{K}_1
}
$$
soit commutatif. Nous disposons alors d'une fonctorialité canonique
$$
\left\{
\begin{matrix}
\mathcal{F}'_2\text{ comme dans (\ref{defos-equation-extension-ringed-topoi}) avec }\\
c_2 = c_{\mathcal{F}'_2}\text{ et }\mathcal{K} = \mathcal{K}_2
\end{matrix}
\right\}
\longrightarrow
\left\{
\begin{matrix}
\mathcal{F}'_1\text{ comme dans (\ref{defos-equation-extension-ringed-topoi}) avec }\\
c_1 = c_{\mathcal{F}'_1}\text{ et }\mathcal{K} = \mathcal{K}_1
\end{matrix}
\right\}
$$
En effet, en considérant tous les faisceaux $\mathcal{O}$,
$\mathcal{O}'_i$, $\mathcal{F}$, $\mathcal{K}_i$, etc., comme des faisceaux
sur $\mathcal{C}$, nous associons à $\mathcal{F}'_2$ le faisceau
$\mathcal{F}'_1$ défini comme la somme amalgamée, c'est-à-dire s'insérant dans
le diagramme d'extensions suivant
$$
\xymatrix{
0 \ar[r] &
\mathcal{K}_2 \ar[r] \ar[d] &
\mathcal{F}'_2 \ar[r] \ar[d] &
\mathcal{F} \ar@{=}[d] \ar[r] & 0 \\
0 \ar[r] &
\mathcal{K}_1 \ar[r] &
\mathcal{F}'_1 \ar[r] &
\mathcal{F} \ar[r] & 0
}
$$
Nous omettons la construction de la structure de $\mathcal{O}'_1$-module sur
la somme amalgamée (elle utilise la commutativité du diagramme faisant
intervenir $c_1$ et $c_2$).
\end{remark}

\begin{remark}
\label{defos-remark-trivial-extension-functorial-ringed-topoi}
Soient $(\Sh(\mathcal{C}), \mathcal{O})$,
$(\Sh(\mathcal{C}), \mathcal{O}) \to (\Sh(\mathcal{D}_i), \mathcal{O}'_i)$,
$\mathcal{I}_i$, et
$h : (\Sh(\mathcal{D}_1), \mathcal{O}'_1) \to
(\Sh(\mathcal{D}_2), \mathcal{O}'_2)$ comme dans la
Remarque \ref{defos-remark-extension-functorial-ringed-topoi}.
Supposons données des trivialisations
$\pi_i : (\Sh(\mathcal{D}_i), \mathcal{O}'_i) \to
(\Sh(\mathcal{C}), \mathcal{O})$ telles que
$\pi_1 = h \circ \pi_2$. Autrement dit, supposons que $h$ soit un morphisme
d'épaississements du premier ordre trivialisés de
$(\Sh(\mathcal{C}), \mathcal{O})$.
Soit $(\mathcal{K}_i, c_i)$, $i = 1, 2$, une paire formée d'un
$\mathcal{O}$-module $\mathcal{K}_i$ et d'une application
$c_i : \mathcal{I}_i \otimes_\mathcal{O} \mathcal{F} \to
\mathcal{K}_i$. Supposons en outre donnée une application de
$\mathcal{O}$-modules $\mathcal{K}_2 \to \mathcal{K}_1$ telle que le diagramme
$$
\xymatrix{
\mathcal{I}_2 \otimes_\mathcal{O} \mathcal{F}
\ar[r]_-{c_2} \ar[d] &
\mathcal{K}_2 \ar[d] \\
\mathcal{I}_1 \otimes_\mathcal{O} \mathcal{F}
\ar[r]^-{c_1} &
\mathcal{K}_1
}
$$
soit commutatif. Dans cette situation, la construction de la
Remarque \ref{defos-remark-trivial-extension-ringed-topoi} induit un diagramme
commutatif
$$
\xymatrix{
\{\mathcal{F}'_2\text{ comme dans (\ref{defos-equation-extension-ringed-topoi}) avec }
c_2 = c_{\mathcal{F}'_2}\text{ et }\mathcal{K} = \mathcal{K}_2\}
\ar[d] \ar[rr] & &
\Ext^1_\mathcal{O}(\mathcal{F}, \mathcal{K}_2) \ar[d] \\
\{\mathcal{F}'_1\text{ comme dans (\ref{defos-equation-extension-ringed-topoi}) avec }
c_1 = c_{\mathcal{F}'_1}\text{ et }\mathcal{K} = \mathcal{K}_1\}
\ar[rr] & &
\Ext^1_\mathcal{O}(\mathcal{F}, \mathcal{K}_1)
}
$$
où l'application verticale de droite est donnée par la fonctorialité de $\Ext$
et l'application $\mathcal{K}_2 \to \mathcal{K}_1$, tandis que l'application
verticale de gauche est celle de la
Remarque \ref{defos-remark-extension-functorial-ringed-topoi}.
\end{remark}

\begin{remark}
\label{defos-remark-obstruction-extension-functorial-ringed-topoi}
Soient $(\Sh(\mathcal{C}), \mathcal{O})$,
$(\Sh(\mathcal{C}), \mathcal{O}) \to (\Sh(\mathcal{D}_i), \mathcal{O}'_i)$,
$\mathcal{I}_i$, et
$h : (\Sh(\mathcal{D}_1), \mathcal{O}'_1) \to
(\Sh(\mathcal{D}_2), \mathcal{O}'_2)$ comme dans la
Remarque \ref{defos-remark-extension-functorial-ringed-topoi}.
Observons que $h^\sharp : \mathcal{O}'_2 \to \mathcal{O}'_1$ induit en
particulier une application de $\mathcal{O}$-modules
$\mathcal{I}_2 \to \mathcal{I}_1$.
Soit $\mathcal{F}$ un $\mathcal{O}$-module.
Soit $(\mathcal{K}_i, c_i)$, $i = 1, 2$, une paire formée d'un
$\mathcal{O}$-module $\mathcal{K}_i$ et d'une application
$c_i : \mathcal{I}_i \otimes_\mathcal{O} \mathcal{F} \to
\mathcal{K}_i$. Supposons en outre donnée une application de
$\mathcal{O}$-modules $\mathcal{K}_2 \to \mathcal{K}_1$ telle que le diagramme
$$
\xymatrix{
\mathcal{I}_2 \otimes_\mathcal{O} \mathcal{F}
\ar[r]_-{c_2} \ar[d] &
\mathcal{K}_2 \ar[d] \\
\mathcal{I}_1 \otimes_\mathcal{O} \mathcal{F}
\ar[r]^-{c_1} &
\mathcal{K}_1
}
$$
soit commutatif. Nous {\bf affirmons} alors que l'application
$$
\Ext^2_\mathcal{O}(\mathcal{F}, \mathcal{K}_2)
\longrightarrow
\Ext^2_\mathcal{O}(\mathcal{F}, \mathcal{K}_1)
$$
envoie $o(\mathcal{F}, \mathcal{K}_2, c_2)$ sur
$o(\mathcal{F}, \mathcal{K}_1, c_1)$.

\medskip\noindent
Pour démontrer cette affirmation, choisissons un plongement
$j_2 : \mathcal{K}_2 \to \mathcal{K}_2'$, où $\mathcal{K}_2'$ est un
$\mathcal{O}$-module injectif.
Comme dans la démonstration du
Lemme \ref{defos-lemma-inf-obs-ext-ringed-topoi}, nous pouvons choisir une extension
de $\mathcal{O}_2$-modules
$$
0 \to \mathcal{K}_2' \to \mathcal{E}_2 \to \mathcal{F} \to 0
$$
telle que $c_{\mathcal{E}_2} = j_2 \circ c_2$.
La démonstration du Lemme \ref{defos-lemma-inf-obs-ext-ringed-topoi} construit
$o(\mathcal{F}, \mathcal{K}_2, c_2)$ comme la classe de l'extension de Yoneda
(au sens de Catégories dérivées, section \ref{derived-section-ext}) associée à
la suite exacte de $\mathcal{O}$-modules
$$
0 \to
\mathcal{K}_2 \to \mathcal{K}_2' \to
\mathcal{E}_2/\mathcal{K}_2 \to
\mathcal{F} \to 0
$$
Soit $\mathcal{K}_1'$ le conoyau de
$\mathcal{K}_2 \to \mathcal{K}_1 \oplus \mathcal{K}_2'$.
Il existe une injection $j_1 : \mathcal{K}_1 \to \mathcal{K}_1'$ et une
application $\mathcal{K}_2' \to \mathcal{K}_1'$ formant un carré commutatif.
Nous formons la somme amalgamée
$$
\xymatrix{
0 \ar[r] &
\mathcal{K}_2' \ar[r] \ar[d] &
\mathcal{E}_2 \ar[r] \ar[d] &
\mathcal{F} \ar[r] \ar[d] & 0 \\
0 \ar[r] &
\mathcal{K}_1' \ar[r] &
\mathcal{E}_1 \ar[r] &
\mathcal{F} \ar[r] & 0
}
$$
Il existe une structure canonique de $\mathcal{O}_1$-module sur
$\mathcal{E}_1$, et, pour cette structure, nous avons
$c_{\mathcal{E}_1} = j_1 \circ c_1$ (cela utilise la commutativité du diagramme
ci-dessus faisant intervenir $c_1$ et $c_2$).
Le procédé du Lemme \ref{defos-lemma-inf-obs-ext-ringed-topoi} nous dit que
$o(\mathcal{F}, \mathcal{K}_1, c_1)$ est la classe de l'extension de Yoneda
associée à la suite exacte de $\mathcal{O}$-modules
$$
0 \to
\mathcal{K}_1 \to
\mathcal{K}_1' \to
\mathcal{E}_1/\mathcal{K}_1 \to
\mathcal{F} \to 0
$$
Puisque nous avons des applications de suites exactes
$$
\xymatrix{
0 \ar[r] &
\mathcal{K}_2 \ar[d] \ar[r] &
\mathcal{K}_2' \ar[d] \ar[r] &
\mathcal{E}_2/\mathcal{K}_2 \ar[r] \ar[d] &
\mathcal{F} \ar[r] \ar@{=}[d] &
0 \\
0 \ar[r] &
\mathcal{K}_2 \ar[r] &
\mathcal{K}_2' \ar[r] &
\mathcal{E}_2/\mathcal{K}_2 \ar[r] &
\mathcal{F} \ar[r] &
0
}
$$
nous concluons que l'affirmation est vraie.
\end{remark}

\begin{remark}
\label{defos-remark-short-exact-sequence-thickenings-ringed-topoi}
Soit $(\Sh(\mathcal{C}), \mathcal{O})$ un topos annelé.
Nous disons qu'une suite de morphismes d'épaississements du premier ordre
$$
(\Sh(\mathcal{D}_1), \mathcal{O}'_1) \to
(\Sh(\mathcal{D}_2), \mathcal{O}'_2) \to
(\Sh(\mathcal{D}_3), \mathcal{O}'_3)
$$
de $(\Sh(\mathcal{C}), \mathcal{O})$ est un {\it complexe} si les applications
correspondantes entre les faisceaux d'idéaux $\mathcal{I}_i$ forment un
complexe de $\mathcal{O}$-modules
$\mathcal{I}_3 \to \mathcal{I}_2 \to \mathcal{I}_1$ (c'est-à-dire si la
composée est nulle). Dans ce cas, la composée
$(\Sh(\mathcal{D}_1), \mathcal{O}'_1) \to
(\Sh(\mathcal{D}_3), \mathcal{O}'_3)$ se factorise par
$(\Sh(\mathcal{C}), \mathcal{O}) \to
(\Sh(\mathcal{D}_3), \mathcal{O}'_3)$ ; autrement dit, l'épaississement du
premier ordre $(\Sh(\mathcal{D}_1), \mathcal{O}'_1)$ de
$(\Sh(\mathcal{C}), \mathcal{O})$ est trivial et muni d'une trivialisation
canonique
$\pi : (\Sh(\mathcal{D}_1), \mathcal{O}'_1) \to
(\Sh(\mathcal{C}), \mathcal{O})$.

\medskip\noindent
Nous disons qu'une suite de morphismes d'épaississements du premier ordre
$$
(\Sh(\mathcal{D}_1), \mathcal{O}'_1) \to
(\Sh(\mathcal{D}_2), \mathcal{O}'_2) \to
(\Sh(\mathcal{D}_3), \mathcal{O}'_3)
$$
de $(\Sh(\mathcal{C}), \mathcal{O})$ est une {\it suite exacte courte} si les
applications correspondantes entre les faisceaux d'idéaux forment une suite
exacte courte
$$
0 \to \mathcal{I}_3 \to \mathcal{I}_2 \to \mathcal{I}_1 \to 0
$$
de $\mathcal{O}$-modules.
\end{remark}

\begin{remark}
\label{defos-remark-complex-thickenings-and-ses-modules-ringed-topoi}
Soit $(\Sh(\mathcal{C}), \mathcal{O})$ un topos annelé. Soit $\mathcal{F}$ un
$\mathcal{O}$-module. Soit
$$
(\Sh(\mathcal{D}_1), \mathcal{O}'_1) \to
(\Sh(\mathcal{D}_2), \mathcal{O}'_2) \to
(\Sh(\mathcal{D}_3), \mathcal{O}'_3)
$$
un complexe d'épaississements du premier ordre de
$(\Sh(\mathcal{C}), \mathcal{O})$ ; voir la
Remarque \ref{defos-remark-short-exact-sequence-thickenings-ringed-topoi}.
Soient $(\mathcal{K}_i, c_i)$, $i = 1, 2, 3$, des couples formés d'un
$\mathcal{O}$-module $\mathcal{K}_i$ et d'une application
$c_i : \mathcal{I}_i \otimes_\mathcal{O} \mathcal{F} \to
\mathcal{K}_i$. Supposons donnée une suite exacte courte de
$\mathcal{O}$-modules
$$
0 \to \mathcal{K}_3 \to \mathcal{K}_2 \to \mathcal{K}_1 \to 0
$$
telle que
$$
\vcenter{
\xymatrix{
\mathcal{I}_2 \otimes_\mathcal{O} \mathcal{F}
\ar[r]_-{c_2} \ar[d] &
\mathcal{K}_2 \ar[d] \\
\mathcal{I}_1 \otimes_\mathcal{O} \mathcal{F}
\ar[r]^-{c_1} &
\mathcal{K}_1
}
}
\quad\text{et}\quad
\vcenter{
\xymatrix{
\mathcal{I}_3 \otimes_\mathcal{O} \mathcal{F}
\ar[r]_-{c_3} \ar[d] &
\mathcal{K}_3 \ar[d] \\
\mathcal{I}_2 \otimes_\mathcal{O} \mathcal{F}
\ar[r]^-{c_2} &
\mathcal{K}_2
}
}
$$
soient commutatifs. Supposons enfin donnée une extension
$$
0 \to \mathcal{K}_2 \to \mathcal{F}'_2 \to \mathcal{F} \to 0
$$
comme dans (\ref{defos-equation-extension-ringed-topoi}), avec
$\mathcal{K} = \mathcal{K}_2$, de $\mathcal{O}'_2$-modules telle que
$c_{\mathcal{F}'_2} = c_2$.
Dans cette situation, nous pouvons appliquer la fonctorialité de la
Remarque \ref{defos-remark-extension-functorial-ringed-topoi} pour obtenir une
extension $\mathcal{F}'_1$ de $\mathcal{O}'_1$-modules (nous décrirons
ci-dessous $\mathcal{F}'_1$ dans ce cas particulier). D'après la
Remarque \ref{defos-remark-trivial-extension-ringed-topoi}, en utilisant le scindage
canonique
$\pi : (\Sh(\mathcal{D}_1), \mathcal{O}'_1) \to
(\Sh(\mathcal{C}), \mathcal{O})$ de la
Remarque \ref{defos-remark-short-exact-sequence-thickenings-ringed-topoi}, nous
obtenons
$\xi_{\mathcal{F}'_1} \in
\Ext^1_\mathcal{O}(\mathcal{F}, \mathcal{K}_1)$.
Enfin, nous avons l'obstruction
$$
o(\mathcal{F}, \mathcal{K}_3, c_3) \in
\Ext^2_\mathcal{O}(\mathcal{F}, \mathcal{K}_3)
$$
du Lemme \ref{defos-lemma-inf-obs-ext-ringed-topoi}.
Dans cette situation, nous {\bf affirmons} que l'application canonique
$$
\partial :
\Ext^1_\mathcal{O}(\mathcal{F}, \mathcal{K}_1)
\longrightarrow
\Ext^2_\mathcal{O}(\mathcal{F}, \mathcal{K}_3)
$$
provenant de la suite exacte courte
$0 \to \mathcal{K}_3 \to \mathcal{K}_2 \to \mathcal{K}_1 \to 0$ envoie
$\xi_{\mathcal{F}'_1}$ sur la classe d'obstruction
$o(\mathcal{F}, \mathcal{K}_3, c_3)$.

\medskip\noindent
Pour démontrer cette affirmation, choisissons un plongement
$j : \mathcal{K}_3 \to \mathcal{K}$, où $\mathcal{K}$ est un
$\mathcal{O}$-module injectif.
Nous pouvons relever $j$ en une application
$j' : \mathcal{K}_2 \to \mathcal{K}$.
Posons $\mathcal{E}'_2 = j'_*\mathcal{F}'_2$, égal à la somme amalgamée de
$\mathcal{F}'_2$ par $j'$, de sorte que
$c_{\mathcal{E}'_2} = j' \circ c_2$. Voici le diagramme
$$
\xymatrix{
0 \ar[r] &
\mathcal{K}_2 \ar[r] \ar[d]_{j'} &
\mathcal{F}'_2 \ar[r] \ar[d] &
\mathcal{F} \ar[r] \ar[d] & 0 \\
0 \ar[r] &
\mathcal{K} \ar[r] &
\mathcal{E}'_2 \ar[r] &
\mathcal{F} \ar[r] & 0
}
$$
Posons $\mathcal{E}'_3 = \mathcal{E}'_2$, mais considérons-le comme un
$\mathcal{O}'_3$-module au moyen de $\mathcal{O}'_3 \to \mathcal{O}'_2$.
Alors $c_{\mathcal{E}'_3} = j \circ c_3$.
La démonstration du Lemme \ref{defos-lemma-inf-obs-ext-ringed-topoi} construit
$o(\mathcal{F}, \mathcal{K}_3, c_3)$ comme le cobord de la classe de
l'extension de $\mathcal{O}$-modules
$$
0 \to
\mathcal{K}/\mathcal{K}_3 \to
\mathcal{E}'_3/\mathcal{K}_3 \to
\mathcal{F} \to 0
$$
D'autre part, remarquons que
$\mathcal{F}'_1 = \mathcal{F}'_2/\mathcal{K}_3$ ; ainsi, la classe
$\xi_{\mathcal{F}'_1}$ est la classe de l'extension
$$
0 \to \mathcal{K}_2/\mathcal{K}_3 \to \mathcal{F}'_2/\mathcal{K}_3
\to \mathcal{F} \to 0
$$
considérée comme une suite de $\mathcal{O}$-modules au moyen de $\pi^\sharp$,
où
$\pi : (\Sh(\mathcal{D}_1), \mathcal{O}'_1) \to
(\Sh(\mathcal{C}), \mathcal{O})$ est le scindage canonique.
Ainsi, l'affirmation résulte finalement de l'existence du diagramme commutatif
$$
\xymatrix{
0 \ar[r] &
\mathcal{K}_2/\mathcal{K}_3 \ar[r] \ar[d] &
\mathcal{F}'_2/\mathcal{K}_3 \ar[r] \ar[d] &
\mathcal{F} \ar[r] \ar[d] & 0 \\
0 \ar[r] &
\mathcal{K}/\mathcal{K}_3 \ar[r] &
\mathcal{E}'_3/\mathcal{K}_3 \ar[r] &
\mathcal{F} \ar[r] & 0
}
$$
qui est $\mathcal{O}$-linéaire (pour les structures de $\mathcal{O}$-modules
données ci-dessus).
\end{remark}

\section{Déformations infinitésimales de modules sur des topos annelés}
\sectionmark{Déformations de modules sur des topos annelés}
\label{defos-section-deformation-modules-ringed-topoi}

\noindent
Soit $i : (\Sh(\mathcal{C}), \mathcal{O}) \to (\Sh(\mathcal{D}), \mathcal{O}')$
un épaississement du premier ordre de topos annelés. Nous employons librement
les notations introduites dans la
section \ref{defos-section-thickenings-ringed-topoi}.
Soit $\mathcal{F}'$ un $\mathcal{O}'$-module et posons
$\mathcal{F} = i^*\mathcal{F}'$.
Dans cette situation, nous avons une suite exacte courte
$$
0 \to \mathcal{I}\mathcal{F}' \to \mathcal{F}' \to \mathcal{F} \to 0
$$
de $\mathcal{O}'$-modules. Puisque $\mathcal{I}^2 = 0$, la structure de
$\mathcal{O}'$-module sur $\mathcal{I}\mathcal{F}'$ provient d'une unique
structure de $\mathcal{O}$-module.
Ainsi, la suite ci-dessus est une extension comme dans
(\ref{defos-equation-extension-ringed-topoi}).
Dans le cas particulier où $\mathcal{F}' = \mathcal{O}'$, nous avons
$i^*\mathcal{O}' = \mathcal{O}$ et
$\mathcal{I}\mathcal{O}' = \mathcal{I}$, et nous retrouvons la suite des
faisceaux structuraux
$$
0 \to \mathcal{I} \to \mathcal{O}' \to \mathcal{O} \to 0
$$

\begin{lemma}
\label{defos-lemma-inf-map-special-ringed-topoi}
Soit $i : (\Sh(\mathcal{C}), \mathcal{O}) \to (\Sh(\mathcal{D}), \mathcal{O}')$
un épaississement du premier ordre de topos annelés.
Soient $\mathcal{F}'$, $\mathcal{G}'$ des $\mathcal{O}'$-modules.
Posons $\mathcal{F} = i^*\mathcal{F}'$ et
$\mathcal{G} = i^*\mathcal{G}'$.
Soit $\varphi : \mathcal{F} \to \mathcal{G}$ une application
$\mathcal{O}$-linéaire.
L'ensemble des relèvements de $\varphi$ en une application
$\mathcal{O}'$-linéaire $\varphi' : \mathcal{F}' \to \mathcal{G}'$ est, s'il
est non vide, un espace principal homogène sous
$\Hom_\mathcal{O}(\mathcal{F}, \mathcal{I}\mathcal{G}')$.
\end{lemma}

\begin{proof}
C'est un cas particulier du Lemme \ref{defos-lemma-inf-map-ringed-topoi}, mais nous
en donnons aussi une démonstration directe. Nous avons des suites exactes
courtes de modules
$$
0 \to \mathcal{I} \to \mathcal{O}' \to \mathcal{O} \to 0
\quad\text{et}\quad
0 \to \mathcal{I}\mathcal{G}' \to \mathcal{G}' \to \mathcal{G} \to 0
$$
et de même pour $\mathcal{F}'$.
Comme $\mathcal{I}$ est de carré nul, les structures de $\mathcal{O}'$-modules
sur $\mathcal{I}$ et $\mathcal{I}\mathcal{G}'$ proviennent d'uniques structures
de $\mathcal{O}$-modules. Il s'ensuit que
$$
\Hom_{\mathcal{O}'}(\mathcal{F}', \mathcal{I}\mathcal{G}') =
\Hom_\mathcal{O}(\mathcal{F}, \mathcal{I}\mathcal{G}')
\quad\text{et}\quad
\Hom_{\mathcal{O}'}(\mathcal{F}', \mathcal{G}) =
\Hom_\mathcal{O}(\mathcal{F}, \mathcal{G})
$$
Le lemme résulte alors de la suite exacte
$$
0 \to \Hom_{\mathcal{O}'}(\mathcal{F}', \mathcal{I}\mathcal{G}') \to
\Hom_{\mathcal{O}'}(\mathcal{F}', \mathcal{G}') \to
\Hom_{\mathcal{O}'}(\mathcal{F}', \mathcal{G})
$$
voir Homologie, Lemme \ref{homology-lemma-check-exactness}.
\end{proof}

\begin{lemma}
\label{defos-lemma-deform-module-ringed-topoi}
Soit $(f, f')$ un morphisme d'épaississements du premier ordre de topos annelés
comme dans la
Situation \ref{defos-situation-morphism-thickenings-ringed-topoi}.
Soit $\mathcal{F}'$ un $\mathcal{O}'$-module et posons
$\mathcal{F} = i^*\mathcal{F}'$.
Supposons que $\mathcal{F}$ soit plat sur $\mathcal{O}_\mathcal{B}$ et que
$(f, f')$ soit un morphisme strict d'épaississements
(Définition \ref{defos-definition-strict-morphism-thickenings-ringed-topoi}).
Alors les assertions suivantes sont équivalentes :
\begin{enumerate}
\item $\mathcal{F}'$ est plat sur $\mathcal{O}_{\mathcal{B}'}$, et
\item l'application canonique
$f^*\mathcal{J} \otimes_\mathcal{O} \mathcal{F} \to
\mathcal{I}\mathcal{F}'$
est un isomorphisme.
\end{enumerate}
De plus, dans ce cas, les applications
$$
f^*\mathcal{J} \otimes_\mathcal{O} \mathcal{F} \to
\mathcal{I} \otimes_\mathcal{O} \mathcal{F} \to
\mathcal{I}\mathcal{F}'
$$
sont des isomorphismes.
\end{lemma}

\begin{proof}
L'application $f^*\mathcal{J} \to \mathcal{I}$ est surjective, car
$(f, f')$ est un morphisme strict d'épaississements.
La dernière assertion résulte donc de (2).

\medskip\noindent
Démontrons l'équivalence de (1) et (2). Par définition, la platitude sur
$\mathcal{O}_\mathcal{B}$ signifie la platitude sur
$f^{-1}\mathcal{O}_\mathcal{B}$. De même pour la platitude sur
$f^{-1}\mathcal{O}_{\mathcal{B}'}$.
Remarquons que la stricte compatibilité de $(f, f')$ et l'hypothèse
$\mathcal{F} = i^*\mathcal{F}'$ impliquent que
$$
\mathcal{F} = \mathcal{F}'/(f^{-1}\mathcal{J})\mathcal{F}'
$$
comme faisceaux sur $\mathcal{C}$. En outre, observons que
$f^*\mathcal{J} \otimes_\mathcal{O} \mathcal{F} =
f^{-1}\mathcal{J} \otimes_{f^{-1}\mathcal{O}_\mathcal{B}} \mathcal{F}$.
L'équivalence de (1) et (2) résulte donc de Modules sur les sites,
Lemme \ref{sites-modules-lemma-flat-over-thickening}.
\end{proof}

\begin{lemma}
\label{defos-lemma-deform-fp-module-ringed-topoi}
Soit $(f, f')$ un morphisme d'épaississements du premier ordre de topos annelés
comme dans la
Situation \ref{defos-situation-morphism-thickenings-ringed-topoi}.
Soit $\mathcal{F}'$ un $\mathcal{O}'$-module et posons
$\mathcal{F} = i^*\mathcal{F}'$.
Supposons que $\mathcal{F}'$ soit plat sur
$\mathcal{O}_{\mathcal{B}'}$ et que $(f, f')$ soit un morphisme strict
d'épaississements. Alors les assertions suivantes sont équivalentes :
\begin{enumerate}
\item $\mathcal{F}'$ est un $\mathcal{O}'$-module de présentation finie, et
\item $\mathcal{F}$ est un $\mathcal{O}$-module de présentation finie.
\end{enumerate}
\end{lemma}

\begin{proof}
L'implication (1) $\Rightarrow$ (2) résulte de Modules sur les sites,
Lemme \ref{sites-modules-lemma-local-pullback}.
Pour la réciproque, supposons $\mathcal{F}$ de présentation finie.
Nous pouvons et allons supposer que $\mathcal{C} = \mathcal{C}'$.
D'après le Lemme \ref{defos-lemma-deform-module-ringed-topoi}, nous avons une suite
exacte courte
$$
0 \to \mathcal{I} \otimes_{\mathcal{O}_X} \mathcal{F} \to
\mathcal{F}' \to \mathcal{F} \to 0
$$
Soit $U$ un objet de $\mathcal{C}$ tel que $\mathcal{F}|_U$ admette une
présentation
$$
\mathcal{O}_U^{\oplus m} \to \mathcal{O}_U^{\oplus n} \to \mathcal{F}|_U \to 0
$$
Quitte à remplacer $U$ par les membres d'un recouvrement, nous pouvons
supposer que l'application $\mathcal{O}_U^{\oplus n} \to \mathcal{F}|_U$ se
relève en une application
$(\mathcal{O}'_U)^{\oplus n} \to \mathcal{F}'|_U$. L'application induite
$\mathcal{I}^{\oplus n} \to \mathcal{I} \otimes \mathcal{F}$ est surjective
par exactitude à droite de $\otimes$. Ainsi, quitte à remplacer encore $U$ par
les membres d'un recouvrement, nous pouvons trouver un relèvement
$(\mathcal{O}'|_U)^{\oplus m} \to (\mathcal{O}'|_U)^{\oplus n}$ de
l'application donnée
$\mathcal{O}_U^{\oplus m} \to \mathcal{O}_U^{\oplus n}$ tel que
$$
(\mathcal{O}'_U)^{\oplus m} \to (\mathcal{O}'_U)^{\oplus n} \to
\mathcal{F}'|_U \to 0
$$
soit un complexe. Une nouvelle application de l'exactitude à droite de
$\otimes$ montre que ce complexe est exact.
\end{proof}

\begin{lemma}
\label{defos-lemma-inf-map-rel-ringed-topoi}
Soit $(f, f')$ un morphisme d'épaississements du premier ordre comme dans la
Situation \ref{defos-situation-morphism-thickenings-ringed-topoi}.
Soient $\mathcal{F}'$, $\mathcal{G}'$ des $\mathcal{O}'$-modules et posons
$\mathcal{F} = i^*\mathcal{F}'$ et $\mathcal{G} = i^*\mathcal{G}'$.
Soit $\varphi : \mathcal{F} \to \mathcal{G}$ une application
$\mathcal{O}$-linéaire.
Supposons que $\mathcal{G}'$ soit plat sur
$\mathcal{O}_{\mathcal{B}'}$ et que $(f, f')$ soit un morphisme strict
d'épaississements.
L'ensemble des relèvements de $\varphi$ en une application
$\mathcal{O}'$-linéaire $\varphi' : \mathcal{F}' \to \mathcal{G}'$ est, s'il
est non vide, un espace principal homogène sous
$$
\Hom_\mathcal{O}(\mathcal{F},
\mathcal{G} \otimes_\mathcal{O} f^*\mathcal{J})
$$
\end{lemma}

\begin{proof}
Combiner les Lemmes \ref{defos-lemma-inf-map-special-ringed-topoi} et
\ref{defos-lemma-deform-module-ringed-topoi}.
\end{proof}

\begin{lemma}
\label{defos-lemma-inf-obs-map-special-ringed-topoi}
Soit $i : (\Sh(\mathcal{C}), \mathcal{O}) \to (\Sh(\mathcal{D}), \mathcal{O}')$
un épaississement du premier ordre de topos annelés.
Soient $\mathcal{F}'$, $\mathcal{G}'$ des $\mathcal{O}'$-modules et posons
$\mathcal{F} = i^*\mathcal{F}'$ et $\mathcal{G} = i^*\mathcal{G}'$.
Soit $\varphi : \mathcal{F} \to \mathcal{G}$ une application
$\mathcal{O}$-linéaire.
Il existe un élément
$$
o(\varphi) \in
\Ext^1_\mathcal{O}(Li^*\mathcal{F}', \mathcal{I}\mathcal{G}')
$$
dont l'annulation est une condition nécessaire et suffisante pour l'existence
d'un relèvement de $\varphi$ en une application $\mathcal{O}'$-linéaire
$\varphi' : \mathcal{F}' \to \mathcal{G}'$.
\end{lemma}

\begin{proof}
Il ressort de la démonstration du
Lemme \ref{defos-lemma-inf-map-special-ringed-topoi} que l'annulation du cobord de
$\varphi$ par l'application
$$
\Hom_\mathcal{O}(\mathcal{F}, \mathcal{G}) =
\Hom_{\mathcal{O}'}(\mathcal{F}', \mathcal{G}) \longrightarrow
\Ext^1_{\mathcal{O}'}(\mathcal{F}', \mathcal{I}\mathcal{G}')
$$
est une condition nécessaire et suffisante pour l'existence d'un relèvement.
Nous concluons puisque
$$
\Ext^1_{\mathcal{O}'}(\mathcal{F}', \mathcal{I}\mathcal{G}') =
\Ext^1_\mathcal{O}(Li^*\mathcal{F}', \mathcal{I}\mathcal{G}')
$$
par l'adjonction de $i_* = Ri_*$ et $Li^*$ dans la catégorie dérivée
(Cohomologie sur les sites, Lemme \ref{sites-cohomology-lemma-adjoint}).
\end{proof}

\begin{lemma}
\label{defos-lemma-inf-obs-map-rel-ringed-topoi}
Soit $(f, f')$ un morphisme d'épaississements du premier ordre comme dans la
Situation \ref{defos-situation-morphism-thickenings-ringed-topoi}.
Soient $\mathcal{F}'$, $\mathcal{G}'$ des $\mathcal{O}'$-modules et posons
$\mathcal{F} = i^*\mathcal{F}'$ et $\mathcal{G} = i^*\mathcal{G}'$.
Soit $\varphi : \mathcal{F} \to \mathcal{G}$ une application
$\mathcal{O}$-linéaire.
Supposons que $\mathcal{F}'$ et $\mathcal{G}'$ soient plats sur
$\mathcal{O}_{\mathcal{B}'}$ et que $(f, f')$ soit un morphisme strict
d'épaississements. Il existe un élément
$$
o(\varphi) \in
\Ext^1_\mathcal{O}(\mathcal{F},
\mathcal{G} \otimes_\mathcal{O} f^*\mathcal{J})
$$
dont l'annulation est une condition nécessaire et suffisante pour l'existence
d'un relèvement de $\varphi$ en une application $\mathcal{O}'$-linéaire
$\varphi' : \mathcal{F}' \to \mathcal{G}'$.
\end{lemma}

\begin{proof}[Première démonstration]
Cela résulte du Lemme \ref{defos-lemma-inf-obs-map-special-ringed-topoi}, car nous
affirmons que, sous les hypothèses du lemme, nous avons
$$
\Ext^1_\mathcal{O}(Li^*\mathcal{F}', \mathcal{I}\mathcal{G}') =
\Ext^1_\mathcal{O}(\mathcal{F},
\mathcal{G} \otimes_\mathcal{O} f^*\mathcal{J})
$$
En effet, nous avons
$\mathcal{I}\mathcal{G}' =
\mathcal{G} \otimes_\mathcal{O} f^*\mathcal{J}$ d'après le
Lemme \ref{defos-lemma-deform-module-ringed-topoi}.
D'autre part, observons que
$$
H^{-1}(Li^*\mathcal{F}') =
\text{Tor}_1^{\mathcal{O}'}(\mathcal{F}', \mathcal{O})
$$
(nous omettons le calcul local). En utilisant la suite exacte courte
$$
0 \to \mathcal{I} \to \mathcal{O}' \to \mathcal{O} \to 0
$$
nous voyons que ce $\text{Tor}_1$ est calculé par le noyau de l'application
$\mathcal{I} \otimes_\mathcal{O} \mathcal{F} \to \mathcal{I}\mathcal{F}'$,
qui est nul d'après la dernière assertion du
Lemme \ref{defos-lemma-deform-module-ringed-topoi}.
Ainsi, $\tau_{\geq -1}Li^*\mathcal{F}' = \mathcal{F}$.
D'autre part, nous avons
$$
\Ext^1_\mathcal{O}(Li^*\mathcal{F}',
\mathcal{I}\mathcal{G}') =
\Ext^1_\mathcal{O}(\tau_{\geq -1}Li^*\mathcal{F}',
\mathcal{I}\mathcal{G}')
$$
par le dual de Catégories dérivées,
Lemme \ref{derived-lemma-negative-vanishing}.
\end{proof}

\begin{proof}[Seconde démonstration]
Nous pouvons appliquer le Lemme \ref{defos-lemma-inf-obs-map-ringed-topoi} comme
suit. Remarquons que
$\mathcal{K} = \mathcal{I} \otimes_\mathcal{O} \mathcal{F}$ et
$\mathcal{L} = \mathcal{I} \otimes_\mathcal{O} \mathcal{G}$ d'après le
Lemme \ref{defos-lemma-deform-module-ringed-topoi}, que
$c_{\mathcal{F}'} = 1 \otimes 1$ et $c_{\mathcal{G}'} = 1 \otimes 1$, et
qu'en prenant $\psi = 1 \otimes \varphi$, le diagramme du lemme est commutatif.
Ainsi, $o(\varphi) = o(\varphi, 1 \otimes \varphi)$ convient.
\end{proof}

\begin{lemma}
\label{defos-lemma-inf-ext-rel-ringed-topoi}
Soit $(f, f')$ un morphisme d'épaississements du premier ordre comme dans la
Situation \ref{defos-situation-morphism-thickenings-ringed-topoi}.
Soit $\mathcal{F}$ un $\mathcal{O}$-module.
Supposons que $(f, f')$ soit un morphisme strict d'épaississements et que
$\mathcal{F}$ soit plat sur $\mathcal{O}_\mathcal{B}$. S'il existe un couple
$(\mathcal{F}', \alpha)$ formé d'un $\mathcal{O}'$-module $\mathcal{F}'$ plat
sur $\mathcal{O}_{\mathcal{B}'}$ et d'un isomorphisme
$\alpha : i^*\mathcal{F}' \to \mathcal{F}$, alors l'ensemble des classes
d'isomorphisme de tels couples est un espace principal homogène sous
$\Ext^1_\mathcal{O}(
\mathcal{F}, \mathcal{I} \otimes_\mathcal{O} \mathcal{F})$.
\end{lemma}

\begin{proof}
Si nous supposons qu'il existe un tel module, alors l'application canonique
$$
f^*\mathcal{J} \otimes_\mathcal{O} \mathcal{F} \to
\mathcal{I} \otimes_\mathcal{O} \mathcal{F}
$$
est un isomorphisme d'après le
Lemme \ref{defos-lemma-deform-module-ringed-topoi}. Appliquons le
Lemme \ref{defos-lemma-inf-ext-ringed-topoi} avec $\mathcal{K} = 
\mathcal{I} \otimes_\mathcal{O} \mathcal{F}$ et $c = 1$.
D'après le Lemme \ref{defos-lemma-deform-module-ringed-topoi}, les extensions
$\mathcal{F}'$ correspondantes sont toutes plates sur
$\mathcal{O}_{\mathcal{B}'}$.
\end{proof}

\begin{lemma}
\label{defos-lemma-inf-obs-ext-rel-ringed-topoi}
Soit $(f, f')$ un morphisme d'épaississements du premier ordre comme dans la
Situation \ref{defos-situation-morphism-thickenings-ringed-topoi}.
Soit $\mathcal{F}$ un $\mathcal{O}$-module. Supposons que $(f, f')$ soit un
morphisme strict d'épaississements et que $\mathcal{F}$ soit plat sur
$\mathcal{O}_\mathcal{B}$. Il existe un $\mathcal{O}'$-module $\mathcal{F}'$
plat sur $\mathcal{O}_{\mathcal{B}'}$ tel que
$i^*\mathcal{F}' \cong \mathcal{F}$ si et seulement si
\begin{enumerate}
\item l'application canonique
$f^*\mathcal{J} \otimes_\mathcal{O} \mathcal{F} \to
\mathcal{I} \otimes_\mathcal{O} \mathcal{F}$
est un isomorphisme, et
\item la classe
$o(\mathcal{F}, \mathcal{I} \otimes_\mathcal{O} \mathcal{F}, 1)
\in \Ext^2_\mathcal{O}(
\mathcal{F}, \mathcal{I} \otimes_\mathcal{O} \mathcal{F})$
du Lemme \ref{defos-lemma-inf-obs-ext-ringed-topoi} est nulle.
\end{enumerate}
\end{lemma}

\begin{proof}
Cela résulte immédiatement de la caractérisation des $\mathcal{O}'$-modules
plats sur $\mathcal{O}_{\mathcal{B}'}$ donnée par le
Lemme \ref{defos-lemma-deform-module-ringed-topoi} et du
Lemme \ref{defos-lemma-inf-obs-ext-ringed-topoi}.
\end{proof}

\section{Application aux modules plats sur les épaississements plats de topos annelés}
\label{defos-section-flat-ringed-topoi}

\noindent
Considérons un diagramme commutatif
$$
\xymatrix{
(\Sh(\mathcal{C}), \mathcal{O}) \ar[r]_i \ar[d]_f &
(\Sh(\mathcal{D}), \mathcal{O}') \ar[d]^{f'} \\
(\Sh(\mathcal{B}), \mathcal{O}_\mathcal{B}) \ar[r]^t &
(\Sh(\mathcal{B}'), \mathcal{O}_{\mathcal{B}'})
}
$$
de topos annelés dont les flèches horizontales sont des épaississements du
premier ordre, comme dans la
Situation \ref{defos-situation-morphism-thickenings-ringed-topoi}. Posons
$\mathcal{I} = \Ker(i^\sharp) \subset \mathcal{O}'$ et
$\mathcal{J} = \Ker(t^\sharp) \subset \mathcal{O}_{\mathcal{B}'}$.
Soit $\mathcal{F}$ un $\mathcal{O}$-module. Supposons que
\begin{enumerate}
\item $(f, f')$ soit un morphisme strict d'épaississements,
\item $f'$ soit plat, et
\item $\mathcal{F}$ soit plat sur $\mathcal{O}_\mathcal{B}$.
\end{enumerate}
Remarquons que (1) $+$ (2) impliquent que
$\mathcal{I} = f^*\mathcal{J}$ (appliquer le
Lemme \ref{defos-lemma-deform-module-ringed-topoi} à $\mathcal{O}'$).
La théorie de la section précédente prend une forme particulièrement simple
sous ces hypothèses. Nous résumons les résultats déjà obtenus dans le lemme
suivant.

\begin{lemma}
\label{defos-lemma-flat-ringed-topoi}
Dans la situation ci-dessus, les assertions suivantes valent.
\begin{enumerate}
\item Il existe un $\mathcal{O}'$-module $\mathcal{F}'$ plat sur
$\mathcal{O}_{\mathcal{B}'}$ tel que
$i^*\mathcal{F}' \cong \mathcal{F}$ si et seulement si la classe
$o(\mathcal{F}, f^*\mathcal{J} \otimes_\mathcal{O} \mathcal{F}, 1)
\in \Ext^2_\mathcal{O}(
\mathcal{F}, f^*\mathcal{J} \otimes_\mathcal{O} \mathcal{F})$
du Lemme \ref{defos-lemma-inf-obs-ext-ringed-topoi} est nulle.
\item S'il existe un tel module, alors l'ensemble des classes d'isomorphisme
des relèvements est un espace principal homogène sous
$\Ext^1_\mathcal{O}(
\mathcal{F}, f^*\mathcal{J} \otimes_\mathcal{O} \mathcal{F})$.
\item Étant donné un relèvement $\mathcal{F}'$, l'ensemble des automorphismes de
$\mathcal{F}'$ dont l'image inverse est $\text{id}_\mathcal{F}$ est
canoniquement isomorphe à
$\Ext^0_\mathcal{O}(
\mathcal{F}, f^*\mathcal{J} \otimes_\mathcal{O} \mathcal{F})$.
\end{enumerate}
\end{lemma}

\begin{proof}
La partie (1) résulte du
Lemme \ref{defos-lemma-inf-obs-ext-rel-ringed-topoi}, car nous avons vu ci-dessus que
$\mathcal{I} = f^*\mathcal{J}$.
La partie (2) résulte du Lemme \ref{defos-lemma-inf-ext-rel-ringed-topoi}.
La partie (3) résulte du Lemme \ref{defos-lemma-inf-map-rel-ringed-topoi}.
\end{proof}

\begin{situation}
\label{defos-situation-morphism-flat-thickenings-ringed-topoi}
Soit $f : (\Sh(\mathcal{C}), \mathcal{O}) \to
(\Sh(\mathcal{B}), \mathcal{O}_\mathcal{B})$ un morphisme de topos annelés.
Considérons un diagramme commutatif
$$
\xymatrix{
(\Sh(\mathcal{C}'_1), \mathcal{O}'_1) \ar[r]_h \ar[d]_{f'_1} &
(\Sh(\mathcal{C}'_2), \mathcal{O}'_2) \ar[d]_{f'_2} \\
(\Sh(\mathcal{B}'_1), \mathcal{O}_{\mathcal{B}'_1}) \ar[r] &
(\Sh(\mathcal{B}'_2), \mathcal{O}_{\mathcal{B}'_2})
}
$$
où $h$ est un morphisme d'épaississements du premier ordre de
$(\Sh(\mathcal{C}), \mathcal{O})$, la flèche horizontale inférieure est un
morphisme d'épaississements du premier ordre de
$(\Sh(\mathcal{B}), \mathcal{O}_\mathcal{B})$, chaque $f'_i$ se restreint à
$f$, les deux couples $(f, f_i')$ sont des morphismes stricts
d'épaississements, et les deux $f'_i$ sont plats. Enfin, soit $\mathcal{F}$ un
$\mathcal{O}$-module plat sur $\mathcal{O}_\mathcal{B}$.
\end{situation}

\begin{lemma}
\label{defos-lemma-functorial-ringed-topoi}
Dans la Situation \ref{defos-situation-morphism-flat-thickenings-ringed-topoi}, la
classe d'obstruction
$o(\mathcal{F}, f^*\mathcal{J}_2 \otimes_\mathcal{O} \mathcal{F}, 1)$
s'envoie sur la classe d'obstruction
$o(\mathcal{F}, f^*\mathcal{J}_1 \otimes_\mathcal{O} \mathcal{F}, 1)$
par l'application canonique
$$
\Ext^2_\mathcal{O}(
\mathcal{F}, f^*\mathcal{J}_2 \otimes_\mathcal{O} \mathcal{F})
\to \Ext^2_\mathcal{O}(
\mathcal{F}, f^*\mathcal{J}_1 \otimes_\mathcal{O} \mathcal{F})
$$
\end{lemma}

\begin{proof}
Cela résulte de la
Remarque \ref{defos-remark-obstruction-extension-functorial-ringed-topoi}.
\end{proof}

\begin{situation}
\label{defos-situation-ses-flat-thickenings-ringed-topoi}
Soit $f : (\Sh(\mathcal{C}), \mathcal{O}) \to
(\Sh(\mathcal{B}), \mathcal{O}_\mathcal{B})$ un morphisme de topos annelés.
Considérons un diagramme commutatif
$$
\xymatrix{
(\Sh(\mathcal{C}'_1), \mathcal{O}'_1) \ar[r]_h \ar[d]_{f'_1} &
(\Sh(\mathcal{C}'_2), \mathcal{O}'_2) \ar[r] \ar[d]_{f'_2} &
(\Sh(\mathcal{C}'_3), \mathcal{O}'_3) \ar[d]_{f'_3} \\
(\Sh(\mathcal{B}'_1), \mathcal{O}_{\mathcal{B}'_1}) \ar[r] &
(\Sh(\mathcal{B}'_2), \mathcal{O}_{\mathcal{B}'_2}) \ar[r] &
(\Sh(\mathcal{B}'_3), \mathcal{O}_{\mathcal{B}'_3})
}
$$
où (a) la ligne supérieure est une suite exacte courte d'épaississements du
premier ordre de $(\Sh(\mathcal{C}), \mathcal{O})$, (b) la ligne inférieure
est une suite exacte courte d'épaississements du premier ordre de
$(\Sh(\mathcal{B}), \mathcal{O}_\mathcal{B})$, (c) chaque $f'_i$ se restreint
à $f$, (d) chaque couple $(f, f_i')$ est un morphisme strict
d'épaississements et (e) chaque $f'_i$ est plat. Enfin, soit
$\mathcal{F}'_2$ un $\mathcal{O}'_2$-module plat sur
$\mathcal{O}_{\mathcal{B}'_2}$ et posons
$\mathcal{F} = \mathcal{F}'_2 \otimes \mathcal{O}$. Soit
$\pi : (\Sh(\mathcal{C}'_1), \mathcal{O}'_1) \to
(\Sh(\mathcal{C}), \mathcal{O})$ le scindage canonique
(Remarque \ref{defos-remark-short-exact-sequence-thickenings-ringed-topoi}).
\end{situation}

\begin{lemma}
\label{defos-lemma-verify-iv-ringed-topoi}
Dans la Situation \ref{defos-situation-ses-flat-thickenings-ringed-topoi}, les
modules $\pi^*\mathcal{F}$ et $h^*\mathcal{F}'_2$ sont des
$\mathcal{O}'_1$-modules plats sur $\mathcal{O}_{\mathcal{B}'_1}$ qui se
restreignent à $\mathcal{F}$ sur $(\Sh(\mathcal{C}), \mathcal{O})$.
Leur différence (Lemme \ref{defos-lemma-flat-ringed-topoi}) est un élément $\theta$
de
$\Ext^1_\mathcal{O}(\mathcal{F},
f^*\mathcal{J}_1 \otimes_\mathcal{O} \mathcal{F})$
dont le cobord dans
$\Ext^2_\mathcal{O}(\mathcal{F},
f^*\mathcal{J}_3 \otimes_\mathcal{O} \mathcal{F})$
est égal à l'obstruction (Lemme \ref{defos-lemma-flat-ringed-topoi}) au relèvement de
$\mathcal{F}$ en un $\mathcal{O}'_3$-module plat sur
$\mathcal{O}_{\mathcal{B}'_3}$.
\end{lemma}

\begin{proof}
Remarquons que $\pi^*\mathcal{F}$ et $h^*\mathcal{F}'_2$ se restreignent tous
deux à $\mathcal{F}$ sur $(\Sh(\mathcal{C}), \mathcal{O})$, et que les noyaux
de $\pi^*\mathcal{F} \to \mathcal{F}$ et
$h^*\mathcal{F}'_2 \to \mathcal{F}$ sont donnés par
$f^*\mathcal{J}_1 \otimes_\mathcal{O} \mathcal{F}$.
Ils sont donc plats d'après le
Lemme \ref{defos-lemma-deform-module-ringed-topoi}.
Prendre le cobord a un sens, car la suite de modules
$$
0 \to f^*\mathcal{J}_3 \otimes_\mathcal{O} \mathcal{F} \to
f^*\mathcal{J}_2 \otimes_\mathcal{O} \mathcal{F} \to
f^*\mathcal{J}_1 \otimes_\mathcal{O} \mathcal{F} \to 0
$$
est exacte courte par les hypothèses de la
Situation \ref{defos-situation-ses-flat-thickenings-ringed-topoi} et le fait que
$\mathcal{F}$ est plat sur $\mathcal{O}_\mathcal{B}$.
L'assertion sur la classe d'obstruction est une traduction directe du résultat
de la Remarque \ref{defos-remark-complex-thickenings-and-ses-modules-ringed-topoi}
dans cette situation particulière.
\end{proof}

\section{Déformations de topos annelés et complexe cotangent naïf}
\label{defos-section-deformations-ringed-topoi}

\noindent
Dans cette section, nous utilisons le complexe cotangent naïf pour faire un peu
de théorie des déformations. Nous partons d'un épaississement du premier ordre
$t : (\Sh(\mathcal{B}), \mathcal{O}_\mathcal{B}) \to
(\Sh(\mathcal{B}'), \mathcal{O}_{\mathcal{B}'})$ de topos annelés.
Nous notons $\mathcal{J} = \Ker(t^\sharp)$ et identifions les topos sous-jacents
de $\mathcal{B}$ et $\mathcal{B}'$.
Supposons en outre donnés un morphisme de topos annelés
$f : (\Sh(\mathcal{C}), \mathcal{O}) \to
(\Sh(\mathcal{B}), \mathcal{O}_\mathcal{B})$, un $\mathcal{O}$-module
$\mathcal{G}$ et une application $f^{-1}\mathcal{J} \to \mathcal{G}$ de
faisceaux de $f^{-1}\mathcal{O}_\mathcal{B}$-modules.
Dans cette section, nous nous demandons si nous pouvons trouver le point
d'interrogation qui complète le diagramme suivant
\begin{equation}
\label{defos-equation-to-solve-ringed-topoi}
\vcenter{
\xymatrix{
0 \ar[r] & \mathcal{G} \ar[r] & {?} \ar[r] & \mathcal{O} \ar[r] & 0 \\
0 \ar[r] & f^{-1}\mathcal{J} \ar[u]^c \ar[r] &
f^{-1}\mathcal{O}_{\mathcal{B}'} \ar[u] \ar[r] &
f^{-1}\mathcal{O}_\mathcal{B} \ar[u] \ar[r] & 0
}
}
\end{equation}
et, de plus, à quel point la solution est unique (si elle existe). Plus
précisément, nous cherchons un épaississement du premier ordre
$i : (\Sh(\mathcal{C}), \mathcal{O}) \to (\Sh(\mathcal{C}'), \mathcal{O}')$
et un morphisme d'épaississements $(f, f')$ comme dans
(\ref{defos-equation-morphism-thickenings-ringed-topoi}), où
$\Ker(i^\sharp)$ est identifié à $\mathcal{G}$, tels que $(f')^\sharp$ induise
l'application donnée $c$.
Nous dirons que $(\Sh(\mathcal{C}'), \mathcal{O}')$ est une {\it solution} de
(\ref{defos-equation-to-solve-ringed-topoi}).


\begin{lemma}
\label{defos-lemma-huge-diagram-ringed-topoi}
Supposons donné un diagramme commutatif de morphismes de topos annelés
\begin{equation}
\label{defos-equation-huge-1-ringed-topoi}
\vcenter{
\xymatrix{
& (\Sh(\mathcal{C}_2), \mathcal{O}_2) \ar[r]_{i_2} \ar[d]_{f_2} \ar[ddl]_g &
(\Sh(\mathcal{C}'_2), \mathcal{O}'_2) \ar[d]^{f'_2} \\
&
(\Sh(\mathcal{B}_2), \mathcal{O}_{\mathcal{B}_2}) \ar[r]^{t_2} \ar[ddl]|\hole &
(\Sh(\mathcal{B}'_2), \mathcal{O}_{\mathcal{B}'_2}) \ar[ddl] \\
(\Sh(\mathcal{C}_1), \mathcal{O}_1) \ar[r]_{i_1} \ar[d]_{f_1} &
(\Sh(\mathcal{C}'_1), \mathcal{O}'_1) \ar[d]^{f'_1} \\
(\Sh(\mathcal{B}_1), \mathcal{O}_{\mathcal{B}_1}) \ar[r]^{t_1} &
(\Sh(\mathcal{B}'_1), \mathcal{O}_{\mathcal{B}'_1})
}
}
\end{equation}
dont les flèches horizontales sont des épaississements du premier ordre.
Posons $\mathcal{G}_j = \Ker(i_j^\sharp)$ et supposons donnée une application
de $g^{-1}\mathcal{O}_1$-modules
$\nu : g^{-1}\mathcal{G}_1 \to \mathcal{G}_2$ qui donne lieu au diagramme
commutatif
\begin{equation}
\label{defos-equation-huge-2-ringed-topoi}
\vcenter{
\xymatrix{
& 0 \ar[r] & \mathcal{G}_2 \ar[r] &
\mathcal{O}'_2 \ar[r] &
\mathcal{O}_2 \ar[r] & 0 \\
& 0 \ar[r]|\hole &
f_2^{-1}\mathcal{J}_2 \ar[u]_{c_2} \ar[r] &
f_2^{-1}\mathcal{O}_{\mathcal{B}'_2} \ar[u] \ar[r]|\hole &
f_2^{-1}\mathcal{O}_{\mathcal{B}_2} \ar[u] \ar[r] & 0 \\
0 \ar[r] &
\mathcal{G}_1 \ar[ruu] \ar[r] &
\mathcal{O}'_1 \ar[r] &
\mathcal{O}_1 \ar[ruu] \ar[r] & 0 \\
0 \ar[r] &
f_1^{-1}\mathcal{J}_1 \ar[ruu]|\hole \ar[u]^{c_1} \ar[r] &
f_1^{-1}\mathcal{O}_{\mathcal{B}'_1} \ar[ruu]|\hole \ar[u] \ar[r] &
f_1^{-1}\mathcal{O}_{\mathcal{B}_1} \ar[ruu]|\hole \ar[u] \ar[r] & 0
}
}
\end{equation}
dont les faces avant et arrière sont des solutions de
(\ref{defos-equation-to-solve-ringed-topoi}).
(Les flèches vers le nord-nord-ouest sont des applications sur
$\mathcal{C}_2$ après application de $g^{-1}$ à la source.)
\begin{enumerate}
\item Il existe un élément canonique de
$\Ext^1_{\mathcal{O}_2}(
Lg^*\NL_{\mathcal{O}_1/\mathcal{O}_{\mathcal{B}_1}}, \mathcal{G}_2)$
dont l'annulation est une condition nécessaire et suffisante pour qu'il existe
un morphisme de topos annelés
$(\Sh(\mathcal{C}'_2), \mathcal{O}'_2) \to
(\Sh(\mathcal{C}'_1), \mathcal{O}'_1)$ s'insérant dans
(\ref{defos-equation-huge-1-ringed-topoi}) de façon compatible avec $\nu$.
\item S'il existe un morphisme
$(\Sh(\mathcal{C}'_2), \mathcal{O}'_2) \to
(\Sh(\mathcal{C}'_1), \mathcal{O}'_1)$ s'insérant dans
(\ref{defos-equation-huge-1-ringed-topoi}) de façon compatible avec $\nu$, l'ensemble
de tous ces morphismes est un espace principal homogène sous
$$
\Hom_{\mathcal{O}_1}(
\Omega_{\mathcal{O}_1/\mathcal{O}_{\mathcal{B}_1}}, g_*\mathcal{G}_2) =
\Hom_{\mathcal{O}_2}(
g^*\Omega_{\mathcal{O}_1/\mathcal{O}_{\mathcal{B}_1}}, \mathcal{G}_2) =
\Ext^0_{\mathcal{O}_2}(
Lg^*\NL_{\mathcal{O}_1/\mathcal{O}_{\mathcal{B}_1}}, \mathcal{G}_2).
$$
\end{enumerate}
\end{lemma}

\begin{proof}
La démonstration de ce lemme est identique à celle du
Lemme \ref{defos-lemma-huge-diagram-ringed-spaces}.
Nous conseillons vivement au lecteur de lire cette dernière plutôt que celle
qui suit. Nous identifierons les topos sous-jacents de tous les
épaississements considérés (nous avons déjà employé cette convention dans
l'énoncé). Les égalités de la dernière assertion du lemme résultent
immédiatement des définitions. Nous travaillerons donc avec les groupes
$\Ext^k_{\mathcal{O}_2}(
Lg^*\NL_{\mathcal{O}_1/\mathcal{O}_{\mathcal{B}_1}}, \mathcal{G}_2)$,
$k = 0, 1$ dans la suite de la démonstration. Montrons d'abord que nous pouvons
nous ramener au cas où le topos sous-jacent de tous les topos annelés du lemme
est le même.

\medskip\noindent
Pour cela, observons que
$g^{-1}\NL_{\mathcal{O}_1/\mathcal{O}_{\mathcal{B}_1}}$ est égal au complexe
cotangent naïf de l'homomorphisme de faisceaux d'anneaux
$g^{-1}f_1^{-1}\mathcal{O}_{\mathcal{B}_1} \to g^{-1}\mathcal{O}_1$ ; voir
Modules sur les sites,
Lemme \ref{sites-modules-lemma-pullback-differentials}.
De plus, le terme de degré $0$ de
$\NL_{\mathcal{O}_1/\mathcal{O}_{\mathcal{B}_1}}$ est un
$\mathcal{O}_1$-module plat ; l'application canonique
$$
Lg^*\NL_{\mathcal{O}_1/\mathcal{O}_{\mathcal{B}_1}}
\longrightarrow
g^{-1}\NL_{\mathcal{O}_1/\mathcal{O}_{\mathcal{B}_1}}
\otimes_{g^{-1}\mathcal{O}_1} \mathcal{O}_2
$$
induit donc un isomorphisme sur les faisceaux de cohomologie en degrés $0$ et
$-1$. Nous pouvons ainsi remplacer les groupes Ext du lemme par
$$
\Ext^k_{g^{-1}\mathcal{O}_1}(
g^{-1}\NL_{\mathcal{O}_1/\mathcal{O}_{\mathcal{B}_1}}, \mathcal{G}_2) =
\Ext^k_{g^{-1}\mathcal{O}_1}(
\NL_{g^{-1}\mathcal{O}_1/g^{-1}f_1^{-1}\mathcal{O}_{\mathcal{B}_1}},
\mathcal{G}_2)
$$
L'ensemble des morphismes de topos annelés
$(\Sh(\mathcal{C}'_2), \mathcal{O}'_2) \to
(\Sh(\mathcal{C}'_1), \mathcal{O}'_1)$ s'insérant dans
(\ref{defos-equation-huge-1-ringed-topoi}) de façon compatible avec $\nu$ est en
bijection avec l'ensemble des homomorphismes de
$g^{-1}f_1^{-1}\mathcal{O}_{\mathcal{B}'_1}$-algèbres
$g^{-1}\mathcal{O}'_1 \to \mathcal{O}'_2$ compatibles avec $f^\sharp$ et
$\nu$. Nous voyons ainsi que nous pouvons supposer donné un diagramme
(\ref{defos-equation-huge-2-ringed-topoi}) de faisceaux sur un site $\mathcal{C}$
(avec $f_1 = f_2 = \text{id}$ sur les topos sous-jacents) et chercher un
homomorphisme de faisceaux d'anneaux $\mathcal{O}'_1 \to \mathcal{O}'_2$ qui
s'y insère.

\medskip\noindent
Dans la suite de la démonstration du lemme, nous supposons que tous les topos
sous-jacents sont les mêmes, c'est-à-dire que nous avons un diagramme
(\ref{defos-equation-huge-2-ringed-topoi}) de faisceaux sur un site $\mathcal{C}$
(avec $f_1 = f_2 = \text{id}$ sur les topos sous-jacents) et que nous cherchons
des homomorphismes de faisceaux d'anneaux
$\mathcal{O}'_1 \to \mathcal{O}'_2$ s'y insérant.
Comme groupes Ext, nous utiliserons
$\Ext^k_{\mathcal{O}_1}(
\NL_{\mathcal{O}_1/\mathcal{O}_{\mathcal{B}_1}}, \mathcal{G}_2)$,
$k = 0, 1$.

\medskip\noindent
Étape 1. Construction de la classe d'obstruction. Considérons le faisceau
d'ensembles
$$
\mathcal{E} = \mathcal{O}'_1 \times_{\mathcal{O}_2} \mathcal{O}'_2
$$
Il est muni d'une application surjective
$\alpha : \mathcal{E} \to \mathcal{O}_1$ ; nous pouvons donc employer
$\NL(\alpha)$ au lieu de
$\NL_{\mathcal{O}_1/\mathcal{O}_{\mathcal{B}_1}}$ ; voir Modules sur les sites,
Lemme \ref{sites-modules-lemma-NL-up-to-qis}.
Posons
$$
\mathcal{I}' =
\Ker(\mathcal{O}_{\mathcal{B}'_1}[\mathcal{E}] \to \mathcal{O}_1)
\quad\text{et}\quad
\mathcal{I} =
\Ker(\mathcal{O}_{\mathcal{B}_1}[\mathcal{E}] \to \mathcal{O}_1)
$$
Il existe une surjection $\mathcal{I}' \to \mathcal{I}$ dont le noyau est
$\mathcal{J}_1\mathcal{O}_{\mathcal{B}'_1}[\mathcal{E}]$.
Nous obtenons deux homomorphismes de $\mathcal{O}_{\mathcal{B}'_2}$-algèbres
$$
a : \mathcal{O}_{\mathcal{B}'_1}[\mathcal{E}] \to \mathcal{O}'_1
\quad\text{et}\quad
b : \mathcal{O}_{\mathcal{B}'_1}[\mathcal{E}] \to \mathcal{O}'_2
$$
qui induisent des applications
$a|_{\mathcal{I}'} : \mathcal{I}' \to \mathcal{G}_1$ et
$b|_{\mathcal{I}'} : \mathcal{I}' \to \mathcal{G}_2$. Les deux applications
$a$ et $b$ annulent $(\mathcal{I}')^2$. De plus, $a$ et $b$ coïncident sur
$\mathcal{J}_1\mathcal{O}_{\mathcal{B}'_1}[\mathcal{E}]$ comme applications
à valeurs dans $\mathcal{G}_2$, car le carré de gauche de
(\ref{defos-equation-huge-2-ringed-topoi}) est commutatif. La différence
$b|_{\mathcal{I}'} - \nu \circ a|_{\mathcal{I}'}$ induit donc une application
$\mathcal{O}_1$-linéaire bien définie
$$
\xi : \mathcal{I}/\mathcal{I}^2 \longrightarrow \mathcal{G}_2
$$
qui envoie la classe d'une section locale $f$ de $\mathcal{I}$ sur
$a(f') - \nu(b(f'))$, où $f'$ est un relèvement de $f$ en une section locale
de $\mathcal{I}'$. Nous notons
$[\xi] \in \Ext^1_{\mathcal{O}_1}(\NL(\alpha), \mathcal{G}_2)$ son image
(voir ci-dessous).

\medskip\noindent
Étape 2. L'annulation de $[\xi]$ est nécessaire. Écrivons
$\Omega =
\Omega_{\mathcal{O}_{\mathcal{B}_1}[\mathcal{E}]/\mathcal{O}_{\mathcal{B}_1}}
\otimes_{\mathcal{O}_{\mathcal{B}_1}[\mathcal{E}]} \mathcal{O}_1$.
Observons que $\NL(\alpha) = (\mathcal{I}/\mathcal{I}^2 \to \Omega)$ s'insère
dans un triangle distingué
$$
\Omega[0] \to
\NL(\alpha) \to
\mathcal{I}/\mathcal{I}^2[1] \to
\Omega[1]
$$
Nous voyons donc que $[\xi]$ est nul si et seulement si $\xi$ est une composée
$\mathcal{I}/\mathcal{I}^2 \to \Omega \to \mathcal{G}_2$ pour une certaine
application $\Omega \to \mathcal{G}_2$.
Supposons qu'il existe un homomorphisme de faisceaux d'anneaux
$\varphi : \mathcal{O}'_1 \to \mathcal{O}'_2$ s'insérant dans
(\ref{defos-equation-huge-2-ringed-topoi}). Considérons alors l'application
$\mathcal{O}'_1[\mathcal{E}] \to \mathcal{G}_2$,
$f' \mapsto b(f') - \varphi(a(f'))$. Un calcul montre qu'elle annule
$\mathcal{J}_1\mathcal{O}_{\mathcal{B}'_1}[\mathcal{E}]$ et induit une
dérivation $\mathcal{O}_{\mathcal{B}_1}[\mathcal{E}] \to \mathcal{G}_2$.
L'application linéaire $\Omega \to \mathcal{G}_2$ ainsi obtenue témoigne de
l'égalité $[\xi] = 0$ dans ce cas.

\medskip\noindent
Étape 3. L'annulation de $[\xi]$ est suffisante. Soit
$\theta : \Omega \to \mathcal{G}_2$ une application $\mathcal{O}_1$-linéaire
telle que $\xi$ soit égale à
$\theta \circ (\mathcal{I}/\mathcal{I}^2 \to \Omega)$.
Un calcul montre alors que
$$
b + \theta \circ d :
\mathcal{O}_{\mathcal{B}'_1}[\mathcal{E}]
\longrightarrow
\mathcal{O}'_2
$$
annule $\mathcal{I}'$ et définit donc une application
$\mathcal{O}'_1 \to \mathcal{O}'_2$ s'insérant dans
(\ref{defos-equation-huge-2-ringed-topoi}).

\medskip\noindent
Démonstration de (2) dans le cas particulier ci-dessus. Omise. Indication :
c'est exactement la même démonstration que celle de la partie (2) du
Lemme \ref{defos-lemma-huge-diagram}.
\end{proof}

\begin{lemma}
\label{defos-lemma-NL-represent-ext-class-ringed-topoi}
Soit $\mathcal{C}$ un site. Soit $\mathcal{A} \to \mathcal{B}$ un
homomorphisme de faisceaux d'anneaux sur $\mathcal{C}$.
Soit $\mathcal{G}$ un $\mathcal{B}$-module.
Soit
$\xi \in \Ext^1_\mathcal{B}(\NL_{\mathcal{B}/\mathcal{A}}, \mathcal{G})$.
Il existe une application de faisceaux d'ensembles
$\alpha : \mathcal{E} \to \mathcal{B}$ telle que
$\xi \in \Ext^1_\mathcal{B}(\NL(\alpha), \mathcal{G})$ soit la classe d'une
application $\mathcal{I}/\mathcal{I}^2 \to \mathcal{G}$ (voir les notations
dans la démonstration).
\end{lemma}

\begin{proof}
Rappelons qu'étant donnée $\alpha : \mathcal{E} \to \mathcal{B}$ telle que
$\mathcal{A}[\mathcal{E}] \to \mathcal{B}$ soit surjective de noyau
$\mathcal{I}$, le complexe
$\NL(\alpha) = (\mathcal{I}/\mathcal{I}^2 \to 
\Omega_{\mathcal{A}[\mathcal{E}]/\mathcal{A}}
\otimes_{\mathcal{A}[\mathcal{E}]} \mathcal{B})$ est canoniquement isomorphe
à $\NL_{\mathcal{B}/\mathcal{A}}$ ; voir Modules sur les sites,
Lemme \ref{sites-modules-lemma-NL-up-to-qis}.
Observons en outre que
$\Omega = \Omega_{\mathcal{A}[\mathcal{E}]/\mathcal{A}}
\otimes_{\mathcal{A}[\mathcal{E}]} \mathcal{B}$ est le faisceau associé au
préfaisceau
$U \mapsto \bigoplus_{e \in \mathcal{E}(U)} \mathcal{B}(U)$.
Autrement dit, $\Omega$ est le $\mathcal{B}$-module libre sur le faisceau
d'ensembles $\mathcal{E}$ ; il existe en particulier une application
canonique $\mathcal{E} \to \Omega$.

\medskip\noindent
Ceci étant dit, choisissons un $\mathcal{E}$ (par exemple
$\mathcal{E} = \mathcal{B}$ comme dans la définition du complexe cotangent
naïf). L'obstruction à écrire $\xi$ comme la classe d'une application
$\mathcal{I}/\mathcal{I}^2 \to \mathcal{G}$ est un élément de
$\Ext^1_\mathcal{B}(\Omega, \mathcal{G})$. Supposons qu'il soit représenté par
l'extension $0 \to \mathcal{G} \to \mathcal{H} \to \Omega \to 0$ de
$\mathcal{B}$-modules. Considérons le faisceau d'ensembles
$\mathcal{E}' = \mathcal{E} \times_\Omega \mathcal{H}$, muni d'une application
induite $\alpha' : \mathcal{E}' \to \mathcal{B}$.
Posons $\mathcal{I}' = \Ker(\mathcal{A}[\mathcal{E}'] \to \mathcal{B})$ et
$\Omega' = \Omega_{\mathcal{A}[\mathcal{E}']/\mathcal{A}}
\otimes_{\mathcal{A}[\mathcal{E}']} \mathcal{B}$.
L'image inverse de $\xi$ par le quasi-isomorphisme
$\NL(\alpha') \to \NL(\alpha)$ s'envoie sur zéro dans
$\Ext^1_\mathcal{B}(\Omega', \mathcal{G})$, car l'image inverse de l'extension
$\mathcal{H}$ par l'application $\Omega' \to \Omega$ est scindée : en effet,
$\Omega'$ est le $\mathcal{B}$-module libre sur le faisceau d'ensembles
$\mathcal{E}'$ et, par construction, il existe un diagramme commutatif
$$
\xymatrix{
\mathcal{E}' \ar[r] \ar[d] & \mathcal{E} \ar[d] \\
\mathcal{H} \ar[r] & \Omega
}
$$
Cela achève la démonstration.
\end{proof}

\begin{lemma}
\label{defos-lemma-choices-ringed-topoi}
S'il existe une solution de (\ref{defos-equation-to-solve-ringed-topoi}), alors
l'ensemble des classes d'isomorphisme de solutions est un espace principal
homogène sous
$\Ext^1_\mathcal{O}(
\NL_{\mathcal{O}/\mathcal{O}_\mathcal{B}}, \mathcal{G})$.
\end{lemma}

\begin{proof}
Observons d'emblée qu'étant données deux solutions $\mathcal{O}'_1$ et
$\mathcal{O}'_2$ de (\ref{defos-equation-to-solve-ringed-topoi}), le
Lemme \ref{defos-lemma-huge-diagram-ringed-topoi} fournit un élément d'obstruction
$o(\mathcal{O}'_1, \mathcal{O}'_2) \in \Ext^1_\mathcal{O}(
\NL_{\mathcal{O}/\mathcal{O}_\mathcal{B}}, \mathcal{G})$ à l'existence d'une
application $\mathcal{O}'_1 \to \mathcal{O}'_2$.
Cet élément est manifestement l'obstruction à l'existence d'un isomorphisme ;
il sépare donc les classes d'isomorphisme. Pour achever la démonstration, il
suffit par conséquent de montrer qu'étant donnés une solution $\mathcal{O}'$ et
un élément
$\xi \in \Ext^1_\mathcal{O}(
\NL_{\mathcal{O}/\mathcal{O}_\mathcal{B}}, \mathcal{G})$, nous pouvons trouver
une seconde solution $\mathcal{O}'_\xi$ telle que
$o(\mathcal{O}', \mathcal{O}'_\xi) = \xi$.

\medskip\noindent
Choisissons $\alpha : \mathcal{E} \to \mathcal{O}$ comme dans le
Lemme \ref{defos-lemma-NL-represent-ext-class-ringed-topoi} pour la classe $\xi$.
Considérons la surjection
$f^{-1}\mathcal{O}_\mathcal{B}[\mathcal{E}] \to \mathcal{O}$ de noyau
$\mathcal{I}$ et le complexe cotangent naïf correspondant
$\NL(\alpha) = (\mathcal{I}/\mathcal{I}^2 \to
\Omega_{f^{-1}\mathcal{O}_\mathcal{B}[\mathcal{E}]/
f^{-1}\mathcal{O}_\mathcal{B}}
\otimes_{f^{-1}\mathcal{O}_\mathcal{B}[\mathcal{E}]} \mathcal{O})$.
D'après le lemme, $\xi$ est la classe d'un morphisme
$\delta : \mathcal{I}/\mathcal{I}^2 \to \mathcal{G}$.
Après avoir remplacé $\mathcal{E}$ par
$\mathcal{E} \times_\mathcal{O} \mathcal{O}'$, nous pouvons aussi supposer que
$\alpha$ se factorise par une application
$\alpha' : \mathcal{E} \to \mathcal{O}'$.

\medskip\noindent
Ces choix déterminent un homomorphisme de
$f^{-1}\mathcal{O}_{\mathcal{B}'}$-algèbres
$\varphi : \mathcal{O}_{\mathcal{B}'}[\mathcal{E}] \to \mathcal{O}'$.
Posons $\mathcal{I}' = \Ker(\varphi)$.
Observons que $\varphi$ induit une application
$\varphi|_{\mathcal{I}'} : \mathcal{I}' \to \mathcal{G}$ et que
$\mathcal{O}'$ est la somme amalgamée, comme dans le diagramme suivant
$$
\xymatrix{
0 \ar[r] & \mathcal{G} \ar[r] & \mathcal{O}' \ar[r] &
\mathcal{O} \ar[r] & 0 \\
0 \ar[r] & \mathcal{I}' \ar[u]^{\varphi|_{\mathcal{I}'}} \ar[r] &
f^{-1}\mathcal{O}_{\mathcal{B}'}[\mathcal{E}] \ar[u] \ar[r] &
\mathcal{O} \ar[u]_{=} \ar[r] & 0
}
$$
Soit $\psi : \mathcal{I}' \to \mathcal{G}$ la somme de l'application
$\varphi|_{\mathcal{I}'}$ et de la composée
$$
\mathcal{I}' \to \mathcal{I}'/(\mathcal{I}')^2 \to
\mathcal{I}/\mathcal{I}^2 \xrightarrow{\delta} \mathcal{G}.
$$
La somme amalgamée suivant $\psi$ est alors une autre extension d'anneaux
$\mathcal{O}'_\xi$ s'insérant dans un diagramme comme ci-dessus.
Un calcul (omis) montre que $o(\mathcal{O}', \mathcal{O}'_\xi) = \xi$, comme
voulu.
\end{proof}

\begin{lemma}
\label{defos-lemma-extensions-of-relative-ringed-topoi}
Soit $f : (\Sh(\mathcal{C}), \mathcal{O}) \to
(\Sh(\mathcal{B}), \mathcal{O}_\mathcal{B})$ un morphisme de topos annelés.
Soit $\mathcal{G}$ un $\mathcal{O}$-module.
L'ensemble des classes d'isomorphisme d'extensions de
$f^{-1}\mathcal{O}_\mathcal{B}$-algèbres
$$
0 \to \mathcal{G} \to \mathcal{O}' \to \mathcal{O} \to 0
$$
où $\mathcal{G}$ est un idéal de carré nul\footnote{Autrement dit, l'ensemble
des classes d'isomorphisme d'épaississements du premier ordre
$i : (\Sh(\mathcal{C}), \mathcal{O}) \to (\Sh(\mathcal{C}), \mathcal{O}')$
au-dessus de $(\Sh(\mathcal{B}), \mathcal{O}_\mathcal{B})$, munis d'un
isomorphisme $\mathcal{G} \to \Ker(i^\sharp)$ de $\mathcal{O}$-modules.}
est canoniquement en bijection avec
$\Ext^1_\mathcal{O}(\NL_{\mathcal{O}/\mathcal{O}_\mathcal{B}}, \mathcal{G})$.
\end{lemma}

\begin{proof}
Pour le démontrer, nous appliquons les résultats précédents au cas où
(\ref{defos-equation-to-solve-ringed-topoi}) est donné par le diagramme
$$
\xymatrix{
0 \ar[r] &
\mathcal{G} \ar[r] &
{?} \ar[r] &
\mathcal{O} \ar[r] & 0 \\
0 \ar[r] &
0 \ar[u] \ar[r] &
f^{-1}\mathcal{O}_\mathcal{B} \ar[u] \ar[r]^{\text{id}} &
f^{-1}\mathcal{O}_\mathcal{B} \ar[u] \ar[r] & 0
}
$$
Notre lemme résulte donc du Lemme \ref{defos-lemma-choices-ringed-topoi} et du fait
qu'il existe une solution, à savoir $\mathcal{G} \oplus \mathcal{O}$.
(Voir la remarque ci-dessous pour une construction directe de la bijection.)
\end{proof}

\begin{remark}
\label{defos-remark-extensions-of-relative-ringed-topoi}
Soient $f : (\Sh(\mathcal{C}), \mathcal{O}) \to
(\mathcal{B}, \mathcal{O}_\mathcal{B})$ et $\mathcal{G}$ comme dans le
Lemme \ref{defos-lemma-extensions-of-relative-ringed-topoi}.
Considérons une extension
$0 \to \mathcal{G} \to \mathcal{O}' \to \mathcal{O} \to 0$ comme dans le
lemme. Nous pouvons choisir un faisceau d'ensembles $\mathcal{E}$ et un
diagramme commutatif
$$
\xymatrix{
\mathcal{E} \ar[d]_{\alpha'} \ar[rd]^\alpha \\
\mathcal{O}' \ar[r] & \mathcal{O}
}
$$
tels que $f^{-1}\mathcal{O}_\mathcal{B}[\mathcal{E}] \to \mathcal{O}$ soit
surjective de noyau $\mathcal{J}$.
(On peut par exemple prendre n'importe quel faisceau d'ensembles se surjectant
sur $\mathcal{O}'$.) Alors
$$
\NL_{\mathcal{O}/\mathcal{O}_\mathcal{B}} \cong \NL(\alpha) = 
\left(
\mathcal{J}/\mathcal{J}^2
\longrightarrow
\Omega_{f^{-1}\mathcal{O}_\mathcal{B}[\mathcal{E}]/
f^{-1}\mathcal{O}_\mathcal{B}}
\otimes_{f^{-1}\mathcal{O}_\mathcal{B}[\mathcal{E}]} \mathcal{O}\right)
$$
Voir Modules sur les sites,
section \ref{sites-modules-section-netherlander}, et en particulier le
Lemme \ref{sites-modules-lemma-NL-up-to-qis}.
Bien entendu, $\alpha'$ détermine une application
$f^{-1}\mathcal{O}_\mathcal{B}[\mathcal{E}] \to \mathcal{O}'$, qui détermine
à son tour une application
$$
\mathcal{J}/\mathcal{J}^2 \longrightarrow \mathcal{G}
$$
qui détermine à son tour l'élément de
$\Ext^1_\mathcal{O}(\NL(\alpha), \mathcal{G}) =
\Ext^1_\mathcal{O}(\NL_{\mathcal{O}/\mathcal{O}_\mathcal{B}}, \mathcal{G})$
correspondant à $\mathcal{O}'$ par la bijection du lemme.
\end{remark}

\begin{lemma}
\label{defos-lemma-extensions-of-relative-ringed-topoi-functorial}
Soient $f : (\Sh(\mathcal{C}), \mathcal{O}_\mathcal{C}) \to
(\Sh(\mathcal{B}), \mathcal{O}_\mathcal{B})$ et
$g : (\Sh(\mathcal{D}), \mathcal{O}_\mathcal{D}) \to
(\Sh(\mathcal{C}), \mathcal{O}_\mathcal{C})$ des morphismes de topos annelés.
Soit $\mathcal{F}$ un $\mathcal{O}_\mathcal{C}$-module.
Soit $\mathcal{G}$ un $\mathcal{O}_\mathcal{D}$-module. Soit
$c : g^*\mathcal{F} \to \mathcal{G}$ une application
$\mathcal{O}_\mathcal{D}$-linéaire. Considérons enfin
\begin{enumerate}
\item[(a)]
$0 \to \mathcal{F} \to \mathcal{O}_{\mathcal{C}'} \to
\mathcal{O}_\mathcal{C} \to 0$, une extension de
$f^{-1}\mathcal{O}_\mathcal{B}$-algèbres correspondant à
$\xi \in \Ext^1_{\mathcal{O}_\mathcal{C}}(
\NL_{\mathcal{O}_\mathcal{C}/\mathcal{O}_\mathcal{B}}, \mathcal{F})$, et
\item[(b)]
$0 \to \mathcal{G} \to \mathcal{O}_{\mathcal{D}'} \to
\mathcal{O}_\mathcal{D} \to 0$, une extension de
$g^{-1}f^{-1}\mathcal{O}_\mathcal{B}$-algèbres correspondant à
$\zeta \in \Ext^1_{\mathcal{O}_\mathcal{D}}(
\NL_{\mathcal{O}_\mathcal{D}/\mathcal{O}_\mathcal{B}}, \mathcal{G})$.
\end{enumerate}
Voir le Lemme \ref{defos-lemma-extensions-of-relative-ringed-topoi}.
Il existe alors un morphisme
$$
g' :
(\Sh(\mathcal{D}), \mathcal{O}_{\mathcal{D}'})
\longrightarrow
(\Sh(\mathcal{C}), \mathcal{O}_{\mathcal{C}'})
$$
de topos annelés au-dessus de
$(\Sh(\mathcal{B}), \mathcal{O}_\mathcal{B})$, compatible avec $g$ et $c$, si
et seulement si $\xi$ et $\zeta$ s'envoient sur le même élément de
$\Ext^1_{\mathcal{O}_\mathcal{D}}(
Lg^*\NL_{\mathcal{O}_\mathcal{C}/\mathcal{O}_\mathcal{B}}, \mathcal{G})$.
\end{lemma}

\begin{proof}
L'énoncé a un sens, car nous disposons des applications
$$
\Ext^1_{\mathcal{O}_\mathcal{C}}(
\NL_{\mathcal{O}_\mathcal{C}/\mathcal{O}_\mathcal{B}}, \mathcal{F}) \to
\Ext^1_{\mathcal{O}_\mathcal{D}}(
Lg^*\NL_{\mathcal{O}_\mathcal{C}/\mathcal{O}_\mathcal{B}}, Lg^*\mathcal{F}) \to
\Ext^1_{\mathcal{O}_\mathcal{D}}
(Lg^*\NL_{\mathcal{O}_\mathcal{C}/\mathcal{O}_\mathcal{B}}, \mathcal{G})
$$
au moyen de l'application
$Lg^*\mathcal{F} \to g^*\mathcal{F} \xrightarrow{c} \mathcal{G}$, et
$$
\Ext^1_{\mathcal{O}_Y}(
\NL_{\mathcal{O}_\mathcal{D}/\mathcal{O}_\mathcal{B}}, \mathcal{G}) \to
\Ext^1_{\mathcal{O}_Y}(
Lg^*\NL_{\mathcal{O}_\mathcal{C}/\mathcal{O}_\mathcal{B}}, \mathcal{G})
$$
au moyen de l'application
$Lg^*\NL_{\mathcal{O}_\mathcal{C}/\mathcal{O}_\mathcal{B}} \to
\NL_{\mathcal{O}_\mathcal{D}/\mathcal{O}_\mathcal{B}}$.
L'énoncé du lemme se déduit du
Lemme \ref{defos-lemma-huge-diagram-ringed-topoi}, appliqué au diagramme
$$
\xymatrix{
& 0 \ar[r] &
\mathcal{G} \ar[r] &
\mathcal{O}_{\mathcal{D}'} \ar[r] &
\mathcal{O}_\mathcal{D} \ar[r] & 0 \\
& 0 \ar[r]|\hole & 0 \ar[u] \ar[r] &
g^{-1}f^{-1}\mathcal{O}_\mathcal{B} \ar[u] \ar[r]|\hole &
g^{-1}f^{-1}\mathcal{O}_\mathcal{B} \ar[u] \ar[r] & 0 \\
0 \ar[r] &
\mathcal{F} \ar[ruu] \ar[r] &
\mathcal{O}_{\mathcal{C}'} \ar[r] &
\mathcal{O}_\mathcal{C} \ar[ruu] \ar[r] & 0 \\
0 \ar[r] & 0 \ar[ruu]|\hole \ar[u] \ar[r] &
f^{-1}\mathcal{O}_\mathcal{B} \ar[ruu]|\hole \ar[u] \ar[r] &
f^{-1}\mathcal{O}_\mathcal{B} \ar[ruu]|\hole \ar[u] \ar[r] & 0
}
$$
ainsi que d'une compatibilité entre les constructions des démonstrations des
Lemmes \ref{defos-lemma-extensions-of-relative-ringed-topoi} et
\ref{defos-lemma-huge-diagram-ringed-topoi}, dont nous omettons l'énoncé et la
démonstration. (Voir la remarque ci-dessous pour un argument direct.)
\end{proof}

\begin{remark}
\label{defos-remark-extensions-of-relative-ringed-topoi-functorial}
Soient $f : (\Sh(\mathcal{C}), \mathcal{O}_\mathcal{C}) \to
(\Sh(\mathcal{B}), \mathcal{O}_\mathcal{B})$,
$g : (\Sh(\mathcal{D}), \mathcal{O}_\mathcal{D}) \to
(\Sh(\mathcal{C}), \mathcal{O}_\mathcal{C})$,
$\mathcal{F}$,
$\mathcal{G}$,
$c : g^*\mathcal{F} \to \mathcal{G}$,
$0 \to \mathcal{F} \to \mathcal{O}_{\mathcal{C}'} \to
\mathcal{O}_\mathcal{C} \to 0$,
$\xi \in \Ext^1_{\mathcal{O}_\mathcal{C}}(
\NL_{\mathcal{O}_\mathcal{C}/\mathcal{O}_\mathcal{B}}, \mathcal{F})$,
$0 \to \mathcal{G} \to \mathcal{O}_{\mathcal{D}'} \to
\mathcal{O}_\mathcal{D} \to 0$, et
$\zeta \in \Ext^1_{\mathcal{O}_\mathcal{D}}(
\NL_{\mathcal{O}_\mathcal{D}/\mathcal{O}_\mathcal{B}}, \mathcal{G})$
comme dans le
Lemme \ref{defos-lemma-extensions-of-relative-ringed-topoi-functorial}.
En prenant la somme amalgamée suivant
$c : g^{-1}\mathcal{F} \to \mathcal{G}$, nous pouvons construire une extension
$$
\xymatrix{
0 \ar[r] &
\mathcal{G} \ar[r] &
\mathcal{O}'_1 \ar[r] &
g^{-1}\mathcal{O}_\mathcal{C} \ar[r] & 0 \\
0 \ar[r] &
g^{-1}\mathcal{F} \ar[u]^c \ar[r] &
g^{-1}\mathcal{O}_{\mathcal{C}'} \ar[u] \ar[r] &
g^{-1}\mathcal{O}_\mathcal{C} \ar@{=}[u] \ar[r] & 0
}
$$
En prenant l'image réciproque suivant
$g^\sharp : g^{-1}\mathcal{O}_\mathcal{C} \to \mathcal{O}_\mathcal{D}$, nous
pouvons construire une extension
$$
\xymatrix{
0 \ar[r] &
\mathcal{G} \ar[r] &
\mathcal{O}_{\mathcal{D}'} \ar[r] &
\mathcal{O}_\mathcal{D} \ar[r] & 0 \\
0 \ar[r] &
\mathcal{G} \ar@{=}[u] \ar[r] &
\mathcal{O}'_2 \ar[u] \ar[r] &
g^{-1}\mathcal{O}_\mathcal{C} \ar[u] \ar[r] & 0
}
$$
Une chasse au diagramme montre qu'il existe un morphisme
$g' : (\Sh(\mathcal{D}), \mathcal{O}_{\mathcal{D}'}) \to
(\Sh(\mathcal{C}), \mathcal{O}_{\mathcal{C}'})$ au-dessus de
$(\Sh(\mathcal{B}), \mathcal{O}_\mathcal{B})$, compatible avec $g$ et $c$, si
et seulement si $\mathcal{O}'_1$ est isomorphe à $\mathcal{O}'_2$ comme
extension de $g^{-1}f^{-1}\mathcal{O}_\mathcal{B}$-algèbres de
$g^{-1}\mathcal{O}_\mathcal{C}$ par $\mathcal{G}$.
D'après le Lemme \ref{defos-lemma-extensions-of-relative-ringed-topoi}, ces extensions
sont classifiées par le membre de gauche de
$$
\Ext^1_{g^{-1}\mathcal{O}_\mathcal{C}}(
\NL_{g^{-1}\mathcal{O}_\mathcal{C}/g^{-1}f^{-1}\mathcal{O}_\mathcal{B}},
\mathcal{G}) =
\Ext^1_{\mathcal{O}_\mathcal{D}}(
Lg^*\NL_{\mathcal{O}_\mathcal{C}/\mathcal{O}_\mathcal{B}}, \mathcal{G})
$$
Ici, l'égalité résulte de l'adjonction entre produit tensoriel et Hom, et des
égalités
$$
\NL_{g^{-1}\mathcal{O}_\mathcal{C}/g^{-1}f^{-1}\mathcal{O}_\mathcal{B}} =
g^{-1}\NL_{\mathcal{O}_\mathcal{C}/\mathcal{O}_\mathcal{B}}
\quad\text{et}\quad
Lg^*\NL_{\mathcal{O}_\mathcal{C}/\mathcal{O}_\mathcal{B}} =
g^{-1}\NL_{\mathcal{O}_\mathcal{C}/\mathcal{O}_\mathcal{B}}
\otimes_{g^{-1}\mathcal{O}_X}^\mathbf{L} \mathcal{O}_Y
$$
Pour la première, voir Modules sur les sites,
Lemme \ref{sites-modules-lemma-pullback-NL} ; la seconde résulte de la
définition de l'image réciproque dérivée.
Ainsi, pour voir que le
Lemme \ref{defos-lemma-extensions-of-relative-ringed-topoi-functorial} est vrai, il
suffit de montrer que $\mathcal{O}'_1$ correspond à l'image de $\xi$ et que
$\mathcal{O}'_2$ correspond à l'image de $\zeta$.
La correspondance entre $\xi$ et $\mathcal{O}'_1$ résulte immédiatement de la
construction de la classe $\xi$ dans la
Remarque \ref{defos-remark-extensions-of-relative-ringed-topoi}.
Pour la correspondance entre $\zeta$ et $\mathcal{O}'_2$, choisissons d'abord
un diagramme commutatif
$$
\xymatrix{
\mathcal{E} \ar[d]_{\beta'} \ar[rd]^\beta \\
\mathcal{O}_{\mathcal{D}'} \ar[r] & \mathcal{O}_\mathcal{D}
}
$$
tel que $g^{-1}f^{-1}\mathcal{O}_\mathcal{B}[\mathcal{E}] \to
\mathcal{O}_\mathcal{D}$ soit surjective de noyau $\mathcal{K}$. Choisissons
ensuite un diagramme commutatif
$$
\xymatrix{
\mathcal{E} \ar[d]_{\beta'} &
\mathcal{E}' \ar[l]^\varphi \ar[d]_{\alpha'} \ar[rd]^\alpha \\
\mathcal{O}_{\mathcal{D}'} &
\mathcal{O}'_2 \ar[l] \ar[r] &
g^{-1}\mathcal{O}_\mathcal{C}
}
$$
tel que $g^{-1}f^{-1}\mathcal{O}_\mathcal{B}[\mathcal{E}'] \to
g^{-1}\mathcal{O}_\mathcal{C}$ soit surjective de noyau $\mathcal{J}$.
(On peut par exemple prendre simplement
$\mathcal{E}' = \mathcal{E} \amalg \mathcal{O}'_2$ comme faisceau
d'ensembles.)
L'application $\varphi$ induit une application de complexes
$\NL(\alpha) \to \NL(\beta)$ (notations de Modules,
section \ref{modules-section-netherlander}), et en particulier
$\bar\varphi : \mathcal{J}/\mathcal{J}^2 \to \mathcal{K}/\mathcal{K}^2$.
Alors $\NL(\alpha) \cong \NL_{\mathcal{O}_\mathcal{D}/\mathcal{O}_\mathcal{B}}$
et
$\NL(\beta) \cong
\NL_{g^{-1}\mathcal{O}_\mathcal{C}/g^{-1}f^{-1}\mathcal{O}_\mathcal{B}}$, et
l'application de complexes $\NL(\alpha) \to \NL(\beta)$ représente
l'application
$Lg^*\NL_{\mathcal{O}_\mathcal{C}/\mathcal{O}_\mathcal{B}} \to
\NL_{\mathcal{O}_\mathcal{D}/\mathcal{O}_\mathcal{B}}$ employée dans l'énoncé
du Lemme \ref{defos-lemma-extensions-of-relative-ringed-topoi-functorial} (voir la
première partie de sa démonstration). Or $\zeta$ correspond à la classe de
l'application $\mathcal{K}/\mathcal{K}^2 \to \mathcal{G}$ induite par
$\beta'$ ; voir la Remarque \ref{defos-remark-extensions-of-relative-ringed-topoi}.
De même, l'extension $\mathcal{O}'_2$ correspond à l'application
$\mathcal{J}/\mathcal{J}^2 \to \mathcal{G}$ induite par $\alpha'$.
Le diagramme commutatif ci-dessus montre que cette application est la composée
de l'application $\mathcal{K}/\mathcal{K}^2 \to \mathcal{G}$ induite par
$\beta'$ avec l'application
$\bar\varphi : \mathcal{J}/\mathcal{J}^2 \to \mathcal{K}/\mathcal{K}^2$.
C'est la compatibilité recherchée.
\end{remark}

\begin{lemma}
\label{defos-lemma-parametrize-solutions-ringed-topoi}
Soient $t : (\Sh(\mathcal{B}), \mathcal{O}_\mathcal{B}) \to
(\Sh(\mathcal{B}'), \mathcal{O}_{\mathcal{B}'})$,
$\mathcal{J} = \Ker(t^\sharp)$,
$f : (\Sh(\mathcal{C}), \mathcal{O}) \to
(\Sh(\mathcal{B}), \mathcal{O}_\mathcal{B})$, $\mathcal{G}$ et
$c : \mathcal{J} \to \mathcal{G}$ comme dans
(\ref{defos-equation-to-solve-ringed-topoi}).
Notons $\xi \in \Ext^1_{\mathcal{O}_\mathcal{B}}(
\NL_{\mathcal{O}_\mathcal{B}/\mathcal{O}_{\mathcal{B}'}}, \mathcal{J})$
l'élément correspondant à l'extension $\mathcal{O}_{\mathcal{B}'}$ de
$\mathcal{O}_\mathcal{B}$ par $\mathcal{J}$ au moyen du
Lemme \ref{defos-lemma-extensions-of-relative-ringed-topoi}.
L'ensemble des classes d'isomorphisme de solutions est canoniquement en
bijection avec la fibre de
$$
\Ext^1_\mathcal{O}(\NL_{\mathcal{O}/\mathcal{O}_{\mathcal{B}'}},
\mathcal{G})\to
\Ext^1_\mathcal{O}(
Lf^*\NL_{\mathcal{O}_\mathcal{B}/\mathcal{O}_{\mathcal{B}'}}, \mathcal{G})
$$
au-dessus de l'image de $\xi$.
\end{lemma}

\begin{proof}
D'après le Lemme \ref{defos-lemma-extensions-of-relative-ringed-topoi}, appliqué à
$t \circ f : (\Sh(\mathcal{C}), \mathcal{O}) \to
(\Sh(\mathcal{B}'), \mathcal{O}_{\mathcal{B}'})$ et au $\mathcal{O}$-module
$\mathcal{G}$, les éléments $\zeta$ de
$\Ext^1_\mathcal{O}(\NL_{\mathcal{O}/\mathcal{O}_{\mathcal{B}'}},
\mathcal{G})$ paramètrent les extensions
$0 \to \mathcal{G} \to \mathcal{O}' \to \mathcal{O} \to 0$ de
$f^{-1}\mathcal{O}_{\mathcal{B}'}$-algèbres. D'après le
Lemme \ref{defos-lemma-extensions-of-relative-ringed-topoi-functorial}, appliqué à
$$
(\Sh(\mathcal{C}), \mathcal{O}) \xrightarrow{f}
(\Sh(\mathcal{B}), \mathcal{O}_\mathcal{B}) \xrightarrow{t}
(\Sh(\mathcal{B}'), \mathcal{O}_{\mathcal{B}'})
$$
et à $c : \mathcal{J} \to \mathcal{G}$, il existe un morphisme
$$
f' : 
(\Sh(\mathcal{C}), \mathcal{O}')
\longrightarrow
(\Sh(\mathcal{B}'), \mathcal{O}_{\mathcal{B}'})
$$
au-dessus de $(\Sh(\mathcal{B}'), \mathcal{O}_{\mathcal{B}'})$, compatible
avec $c$ et $f$, si et seulement si $\zeta$ s'envoie sur $\xi$.
Bien entendu, cela revient à dire que $\mathcal{O}'$ est une solution de
(\ref{defos-equation-to-solve-ringed-topoi}).
\end{proof}

\section{Déformations d'espaces algébriques}
\label{defos-section-deformations-spaces}

\noindent
Dans cette section, nous explicitons la signification des résultats de la
section \ref{defos-section-deformations-ringed-topoi} pour les déformations d'espaces
algébriques.

\begin{lemma}
\label{defos-lemma-match-thickenings}
Soit $S$ un schéma. Soit $i : Z \to Z'$ un morphisme d'espaces algébriques
au-dessus de $S$. Les assertions suivantes sont équivalentes :
\begin{enumerate}
\item $i$ est un épaississement d'espaces algébriques tel qu'il est défini dans
Compléments sur les morphismes d'espaces, section
\ref{spaces-more-morphisms-section-thickenings}, et
\item le morphisme associé
$i_{small} : (\Sh(Z_\etale), \mathcal{O}_Z) \to
(\Sh(Z'_\etale), \mathcal{O}_{Z'})$
de topos annelés (Propriétés des espaces, Lemme
\ref{spaces-properties-lemma-morphism-ringed-topoi}) est un épaississement au
sens de la section \ref{defos-section-thickenings-ringed-topoi}.
\end{enumerate}
\end{lemma}

\begin{proof}
Soulignons qu'il ne s'agit pas d'une trivialité.

\medskip\noindent
Supposons (1). D'après Compléments sur les morphismes d'espaces,
Lemme \ref{spaces-more-morphisms-lemma-thickening-equivalence}, le morphisme
$i$ induit une équivalence des petits sites étales, et donc en particulier des
topos. Bien entendu, $i^\sharp$ est surjective à noyau localement nilpotent par
définition des épaississements.

\medskip\noindent
Supposons (2). (Cette implication est moins importante et relève davantage de
la curiosité.) Pour tout morphisme étale $Y' \to Z'$, nous voyons que
$Y = Z \times_{Z'} Y'$ a le même topos étale que $Y'$. En particulier, $Y'$ est
quasi-compact si et seulement si $Y$ est quasi-compact, car la quasi-compacité
est une notion toposique (Sites, Lemme \ref{sites-lemma-quasi-compact}).
Ceci étant dit, nous voyons que $Y'$ est quasi-compact et quasi-séparé si et
seulement si $Y$ est quasi-compact et quasi-séparé (car on peut caractériser la
quasi-séparation de $Y'$ en demandant que, pour tous espaces algébriques
quasi-compacts $Y'_1, Y'_2$ étales sur $Y'$, le produit
$Y'_1 \times_{Y'} Y'_2$ soit quasi-compact).
Prenons $Y'$ affine. Alors l'espace algébrique $Y$ est quasi-compact et
quasi-séparé. Pour tout $\mathcal{O}_Y$-module quasi-cohérent $\mathcal{F}$,
nous avons
$H^q(Y, \mathcal{F}) = H^q(Y', (Y \to Y')_*\mathcal{F})$, car les topos étales
sont les mêmes. Ensuite,
$H^q(Y', (Y \to Y')_*\mathcal{F}) = 0$, car l'image directe est
quasi-cohérente (Morphismes d'espaces,
Lemme \ref{spaces-morphisms-lemma-pushforward}) et $Y$ est affine.
Il s'ensuit que $Y'$ est affine d'après Cohomologie des espaces, Proposition
\ref{spaces-cohomology-proposition-vanishing-affine} (il existe sûrement une
démonstration de cette implication qui évite cette référence).
Ainsi, $i$ est un morphisme affine. Dans le cas affine, les conditions de la
section \ref{defos-section-thickenings-ringed-topoi} entraînent aisément que $i$ est
un épaississement d'espaces algébriques.
\end{proof}

\begin{lemma}
\label{defos-lemma-deform-spaces}
Soit $S$ un schéma.
Soit $Y \subset Y'$ un épaississement du premier ordre d'espaces algébriques
au-dessus de $S$.
Soit $f : X \to Y$ un morphisme plat d'espaces algébriques au-dessus de $S$.
S'il existe un morphisme plat $f' : X' \to Y'$ d'espaces algébriques au-dessus
de $S$ et un isomorphisme $a : X \to X' \times_{Y'} Y$ au-dessus de $Y$, alors
\begin{enumerate}
\item l'ensemble des classes d'isomorphisme de couples
$(f' : X' \to Y', a)$ est un espace principal homogène sous
$\Ext^1_{\mathcal{O}_X}(\NL_{X/Y}, f^*\mathcal{C}_{Y/Y'})$, et
\item l'ensemble des automorphismes $\varphi : X' \to X'$ au-dessus de $Y'$
qui se réduisent à l'identité sur $X' \times_{Y'} Y$ est
$\Ext^0_{\mathcal{O}_X}(\NL_{X/Y}, f^*\mathcal{C}_{Y/Y'})$.
\end{enumerate}
\end{lemma}

\begin{proof}
Nous appliquerons les résultats sur les déformations de topos annelés aux
petits topos étales des espaces algébriques du lemme.
Nous pouvons considérer $X$ comme un sous-espace fermé de $X'$ de sorte que
$(f, f') : (X \subset X') \to (Y \subset Y')$ soit un morphisme
d'épaississements du premier ordre.
D'après le Lemme \ref{defos-lemma-match-thickenings}, cela se traduit par un
morphisme d'épaississements de topos annelés.
Il résulte alors de Compléments sur les morphismes d'espaces,
Lemme \ref{spaces-more-morphisms-lemma-deform} (ou du
Lemme \ref{defos-lemma-deform-module-ringed-topoi}, plus général) que le faisceau
d'idéaux de $X$ dans $X'$ est égal à $f^*\mathcal{C}_{Y'/Y}$, et ceci équivaut
en fait à la platitude de $X'$ sur $Y'$.
Nous avons donc un diagramme commutatif
$$
\xymatrix{
0 \ar[r] & f^*\mathcal{C}_{Y/Y'} \ar[r] &
\mathcal{O}_{X'} \ar[r] &
\mathcal{O}_X \ar[r] & 0 \\
0 \ar[r] &
f_{small}^{-1}\mathcal{C}_{Y/Y'} \ar[u] \ar[r] &
f_{small}^{-1}\mathcal{O}_{Y'} \ar[u] \ar[r] &
f_{small}^{-1}\mathcal{O}_Y \ar[u] \ar[r] & 0
}
$$
Comparer avec (\ref{defos-equation-to-solve-ringed-topoi}).
Observons que les automorphismes $\varphi$ de (2) donnent des automorphismes
$\varphi^\sharp : \mathcal{O}_{X'} \to \mathcal{O}_{X'}$ qui s'insèrent dans
le diagramme ci-dessus. Réciproquement, un automorphisme
$\alpha : \mathcal{O}_{X'} \to \mathcal{O}_{X'}$ s'insérant dans le diagramme
de faisceaux ci-dessus est égal à $\varphi^\sharp$ pour un certain
automorphisme $\varphi$ de (2), d'après Compléments sur les morphismes
d'espaces, Lemme
\ref{spaces-more-morphisms-lemma-first-order-thickening-maps}.
Enfin, d'après Compléments sur les morphismes d'espaces,
Lemme \ref{spaces-more-morphisms-lemma-first-order-thickening}, si nous trouvons
un autre faisceau d'anneaux $\mathcal{A}$ sur $X_\etale$ s'insérant dans le
diagramme
$$
\xymatrix{
0 \ar[r] & f^*\mathcal{C}_{Y/Y'} \ar[r] &
\mathcal{A} \ar[r] &
\mathcal{O}_X \ar[r] & 0 \\
0 \ar[r] &
f_{small}^{-1}\mathcal{C}_{Y/Y'} \ar[u] \ar[r] &
f_{small}^{-1}\mathcal{O}_{Y'} \ar[u] \ar[r] &
f_{small}^{-1}\mathcal{O}_Y \ar[u] \ar[r] & 0
}
$$
alors il existe un épaississement du premier ordre $X \subset X''$ tel que
$\mathcal{O}_{X''} = \mathcal{A}$ et, en appliquant encore Compléments sur les
morphismes d'espaces, Lemme
\ref{spaces-more-morphisms-lemma-first-order-thickening-maps}, nous obtenons un
morphisme $(f, f'') : (X \subset X'') \to (Y \subset Y')$ ayant toutes les
propriétés voulues.
La partie (1) résulte donc du Lemme \ref{defos-lemma-choices-ringed-topoi}, et la
partie (2), de la partie (2) du
Lemme \ref{defos-lemma-huge-diagram-ringed-topoi}.
(Notons que $\NL_{X/Y}$, tel qu'il est défini pour un morphisme d'espaces
algébriques dans Compléments sur les morphismes d'espaces, section
\ref{spaces-more-morphisms-section-netherlander}, coïncide avec $\NL_{X/Y}$ tel
qu'il est employé dans la section \ref{defos-section-deformations-ringed-topoi}.)
\end{proof}

\noindent
Soit $S$ un schéma. Soit $f : X \to B$ un morphisme d'espaces algébriques
au-dessus de $S$.
Soit $\mathcal{F} \to \mathcal{G}$ un homomorphisme de
$\mathcal{O}_X$-modules (non nécessairement quasi cohérents).
Considérons le foncteur
$$
F :
\left\{
\begin{matrix}
\text{extensions of }f^{-1}\mathcal{O}_B\text{ algebras}\\
0 \to \mathcal{F} \to \mathcal{O}' \to \mathcal{O}_X \to 0\\
\text{où }\mathcal{F}\text{ est un idéal de carré nul}
\end{matrix}
\right\}
\longrightarrow
\left\{
\begin{matrix}
\text{extensions of }f^{-1}\mathcal{O}_B\text{ algebras}\\
0 \to \mathcal{G} \to \mathcal{O}' \to \mathcal{O}_X \to 0\\
\text{où }\mathcal{G}\text{ est un idéal de carré nul}
\end{matrix}
\right\}
$$
défini par somme amalgamée.

\begin{lemma}
\label{defos-lemma-thickening-space-quasi-coherent}
Dans la situation ci-dessus, supposons que $X$ soit quasi-compact et
quasi-séparé, et que $DQ_X(\mathcal{F}) \to DQ_X(\mathcal{G})$
(Catégories dérivées des espaces, section
\ref{spaces-perfect-section-better-coherator}) soit un isomorphisme.
Alors le foncteur $F$ est une équivalence de catégories.
\end{lemma}

\begin{proof}
Rappelons que $\NL_{X/B}$ est un objet de $D_\QCoh(\mathcal{O}_X)$ ; voir
Compléments sur les morphismes d'espaces, Lemme
\ref{spaces-more-morphisms-lemma-netherlander-quasi-coherent}.
Notre hypothèse implique donc que les applications
$$
\Ext^i_X(\NL_{X/B}, \mathcal{F}) \longrightarrow
\Ext^i_X(\NL_{X/B}, \mathcal{G})
$$
sont des isomorphismes pour tout $i$. Le
Lemme \ref{defos-lemma-huge-diagram-ringed-topoi} implique que notre foncteur est
pleinement fidèle. D'autre part, le foncteur est essentiellement surjectif par
le Lemme \ref{defos-lemma-choices-ringed-topoi}, car nous avons les solutions
$\mathcal{O}_X \oplus \mathcal{F}$ et $\mathcal{O}_X \oplus \mathcal{G}$ dans
les deux catégories.
\end{proof}

\noindent
Soit $S$ un schéma. Soit $B \subset B'$ un épaississement du premier ordre
d'espaces algébriques au-dessus de $S$, de faisceau d'idéaux $\mathcal{J}$,
que nous considérons soit comme un $\mathcal{O}_B$-module quasi-cohérent, soit
comme un faisceau quasi-cohérent d'idéaux sur $B'$ ; voir Compléments sur les
morphismes d'espaces, section
\ref{spaces-more-morphisms-section-thickenings}.
Soit $f : X \to B$ un morphisme d'espaces algébriques au-dessus de $S$.
Soit $\mathcal{F} \to \mathcal{G}$ un homomorphisme de
$\mathcal{O}_X$-modules (non nécessairement quasi cohérents).
Soit $c : f^{-1}\mathcal{J} \to \mathcal{F}$ une application de
$f^{-1}\mathcal{O}_B$-modules, et notons
$c' : f^{-1}\mathcal{J} \to \mathcal{G}$ la composée.
Considérons le foncteur
$$
FT :
\{\text{solutions de }(\ref{defos-equation-to-solve-ringed-topoi})
\text{ pour }\mathcal{F}\text{ et }c\}
\longrightarrow
\{\text{solutions de }(\ref{defos-equation-to-solve-ringed-topoi})
\text{ pour }\mathcal{G}\text{ et }c'\}
$$
défini par somme amalgamée.

\begin{lemma}
\label{defos-lemma-thickening-over-thickening-space-quasi-coherent}
Dans la situation ci-dessus, supposons que $X$ soit quasi-compact et
quasi-séparé, et que $DQ_X(\mathcal{F}) \to DQ_X(\mathcal{G})$
(Catégories dérivées des espaces, section
\ref{spaces-perfect-section-better-coherator}) soit un isomorphisme.
Alors le foncteur $FT$ est une équivalence de catégories.
\end{lemma}

\begin{proof}
Une solution de (\ref{defos-equation-to-solve-ringed-topoi}) pour $\mathcal{F}$ donne
en particulier une extension de $f^{-1}\mathcal{O}_{B'}$-algèbres
$$
0 \to \mathcal{F} \to \mathcal{O}' \to \mathcal{O}_X \to 0
$$
où $\mathcal{F}$ est un idéal de carré nul. De même pour $\mathcal{G}$.
De plus, une telle extension fournit une application
$c_{\mathcal{O}'} : f^{-1}\mathcal{J} \to \mathcal{F}$.
Nous considérons donc la sous-catégorie pleine de ces extensions de
$f^{-1}\mathcal{O}_{B'}$-algèbres telles que $c = c_{\mathcal{O}'}$.
Manifestement, si $\mathcal{O}'' = F(\mathcal{O}')$, où $F$ est l'équivalence
du Lemme \ref{defos-lemma-thickening-space-quasi-coherent} (appliqué cette fois à
$X \to B'$), alors $c_{\mathcal{O}''}$ est la composée de
$c_{\mathcal{O}'}$ et de l'application $\mathcal{F} \to \mathcal{G}$.
Cela démontre le lemme.
\end{proof}

\section{Déformations de complexes}
\label{defos-section-deformations-complexes}

\noindent
Cette section sert de préparation à la suivante.
Nous utiliserons autant que possible les résultats des chapitres consacrés à
l'algèbre commutative.

\begin{lemma}
\label{defos-lemma-canonical-class-algebra}
Soit $R' \to R$ une surjection d'anneaux dont le noyau est un idéal $I$ de
carré nul. Pour tout $K \in D^-(R)$, il existe une application canonique
$$
\omega(K) : K \longrightarrow K \otimes_R^\mathbf{L} I[2]
$$
dans $D(R)$ possédant les propriétés suivantes :
\begin{enumerate}
\item $\omega(K) = 0$ si et seulement s'il existe
$K' \in D(R')$ tel que $K' \otimes_{R'}^\mathbf{L} R = K$,
\item étant donnée une application $K \to L$ dans $D^-(R)$, le diagramme
$$
\xymatrix{
K \ar[d] \ar[rr]_-{\omega(K)} & &
K \otimes^\mathbf{L}_R I[2] \ar[d] \\
L \ar[rr]^-{\omega(L)} & &
L \otimes^\mathbf{L}_R I[2]
}
$$
est commutatif, et
\item la formation de $\omega(K)$ est compatible aux homomorphismes d'anneaux
$R' \to S'$ (voir la démonstration pour un énoncé précis).
\end{enumerate}
\end{lemma}

\begin{proof}
Choisissons un complexe $K^\bullet$ de $R$-modules libres, borné
supérieurement, représentant $K$. Nous pouvons alors choisir des
$R'$-modules libres $(K')^n$ relevant $K^n$.
Nous pouvons choisir des applications de $R'$-modules
$(d')^n_K : (K')^n \to (K')^{n + 1}$ relevant les différentielles
$d^n_K : K^n \to K^{n + 1}$ de $K^\bullet$.
Bien que les composées
$$
(d')^{n + 1}_K \circ (d')^n_K : (K')^n \to (K')^{n + 2}
$$
ne soient pas nécessairement nulles, elles se factorisent sous la forme
$$
(K')^n \to K^n \xrightarrow{\omega^n_K}
K^{n + 2} \otimes_R I = I(K')^{n + 2} \to (K')^{n + 2}
$$
car $d^{n + 1} \circ d^n = 0$.
Un calcul montre que $\omega^n_K$ définit une application de complexes.
Cette application de complexes définit $\omega(K)$.

\medskip\noindent
Montrons que cette construction est compatible avec une application de
complexes $\alpha^\bullet : K^\bullet \to L^\bullet$ de $R$-modules libres
bornés supérieurement et avec des choix de relèvements
$(K')^n, (L')^n, (d')^n_K, (d')^n_L$.
Choisissons pour cela $(\alpha')^n : (K')^n \to (L')^n$ relevant les
composantes $\alpha^n : K^n \to L^n$. Comme ci-dessus, nous obtenons une
factorisation
$$
(K')^n \to K^n \xrightarrow{h^n}
L^{n + 1} \otimes_R I = I(L')^{n + 1} \to (L')^{n + 2}
$$
de $(d')^n_L \circ (\alpha')^n - (\alpha')^{n + 1} \circ (d')_K^n$.
Un calcul agréable montre alors que
$$
\omega^n_L \circ \alpha^n =
(d_L^{n + 1} \otimes \text{id}_I) \circ h^n + h^{n + 1} \circ d_K^n +
(\alpha^{n + 2} \otimes \text{id}_I) \circ \omega^n_K
$$
Cela démontre la commutativité du diagramme de la partie (2) du lemme dans ce
cas particulier. En appliquant ceci à deux choix différents de complexes
libres bornés supérieurement représentant $K$, nous voyons que $\omega(K)$ est
bien définie. Bien entendu, (2) vaut alors aussi en général.

\medskip\noindent
Si $K$ se relève en $K'$ dans $D^-(R')$, nous pouvons représenter $K'$ par un
complexe de $R'$-modules libres borné supérieurement et voyons immédiatement
que $\omega(K) = 0$.
Réciproquement, revenons à nos choix $K^\bullet$,
$(K')^n$, $(d')^n_K$. Si $\omega(K) = 0$, nous pouvons trouver
$g^n : K^n \to K^{n + 1} \otimes_R I$ tel que
$$
\omega^n = (d_K^{n + 1} \otimes \text{id}_I) \circ g^n +
g^{n + 1} \circ d_K^n
$$
En prenant comme différentielles
$(d')^n_K - g^n : (K')^n \to (K')^{n + 1}$, nous obtenons donc un complexe de
$R'$-modules libres relevant $K^\bullet$. Cela démontre (1).

\medskip\noindent
Enfin, la partie (3) signifie ce qui suit. Soit $R' \to S'$ un homomorphisme
d'anneaux. Posons $S = S' \otimes_{R'} R$ et notons
$J = IS' \subset S'$ le noyau de carré nul de $S' \to S$. Étant donné
$K \in D^-(R)$, l'assertion est que nous obtenons un diagramme commutatif
$$
\xymatrix{
K \otimes_R^\mathbf{L} S \ar[d] \ar[rr]_-{\omega(K) \otimes \text{id}} & &
(K \otimes^\mathbf{L}_R I[2]) \otimes_R^\mathbf{L} S \ar[d] \\
K \otimes_R^\mathbf{L} S \ar[rr]^-{\omega(K \otimes_R^\mathbf{L} S)} & &
(K \otimes_R^\mathbf{L} S) \otimes^\mathbf{L}_S J[2]
}
$$
Ici, la flèche verticale de droite provient de
$$
(K \otimes^\mathbf{L}_R I[2]) \otimes_R^\mathbf{L} S =
(K \otimes_R^\mathbf{L} S) \otimes_S^\mathbf{L}
(I \otimes_R^\mathbf{L} S)[2] \longrightarrow
(K \otimes_R^\mathbf{L} S) \otimes_S^\mathbf{L} J[2]
$$
Choisissons $K^\bullet$, $(K')^n$ et $(d')^n_K$ comme ci-dessus.
Nous pouvons alors employer $K^\bullet \otimes_R S$,
$(K')^n \otimes_{R'} S'$ et $(d')^n_K \otimes \text{id}_{S'}$ pour la
construction de $\omega(K \otimes_R^\mathbf{L} S)$.
Avec ces choix, la commutativité se vérifie immédiatement au niveau des
applications de complexes.
\end{proof}

\section{Déformations de complexes sur des topos annelés}
\label{defos-section-thickenings-complexes}

\noindent
Ce matériau est tiré de \cite{lieblich-complexes}.

\medskip\noindent
Les résultats de cette section valent dans le cadre d'un épaississement du
premier ordre de topos annelés tel qu'il est défini dans la
section \ref{defos-section-thickenings-ringed-topoi}.
Cependant, afin de simplifier les notations, nous supposerons que les sites
sous-jacents $\mathcal{C}$ et $\mathcal{D}$ sont les mêmes.
De plus, l'homomorphisme surjectif $\mathcal{O}' \to \mathcal{O}$ de faisceaux
d'anneaux sera noté $\mathcal{O} \to \mathcal{O}_0$, conformément à l'usage
sans doute plus courant dans la littérature.

\begin{lemma}
\label{defos-lemma-lift-complex}
Soit $\mathcal{C}$ un site. Soit $\mathcal{O} \to \mathcal{O}_0$ une
surjection de faisceaux d'anneaux. Supposons données les informations
suivantes :
\begin{enumerate}
\item des $\mathcal{O}$-modules plats $\mathcal{G}^n$,
\item des applications de $\mathcal{O}$-modules
$\mathcal{G}^n \to \mathcal{G}^{n + 1}$,
\item un complexe $\mathcal{K}_0^\bullet$ de $\mathcal{O}_0$-modules,
\item des applications de $\mathcal{O}$-modules
$\mathcal{G}^n \to \mathcal{K}_0^n$
\end{enumerate}
telles que
\begin{enumerate}
\item[(a)] $H^n(\mathcal{K}_0^\bullet) = 0$ pour $n \gg 0$,
\item[(b)] $\mathcal{G}^n = 0$ pour $n \gg 0$,
\item[(c)] en posant
$\mathcal{G}^n_0 = \mathcal{G}^n \otimes_\mathcal{O} \mathcal{O}_0$, les
applications induites déterminent un complexe $\mathcal{G}_0^\bullet$ et une
application de complexes
$\mathcal{G}_0^\bullet \to \mathcal{K}_0^\bullet$.
\end{enumerate}
Il existe alors
\begin{enumerate}
\item[(\romannumeral1)]
des $\mathcal{O}$-modules plats $\mathcal{F}^n$,
\item[(\romannumeral2)]
des applications de $\mathcal{O}$-modules
$\mathcal{F}^n \to \mathcal{F}^{n + 1}$,
\item[(\romannumeral3)]
des applications de $\mathcal{O}$-modules
$\mathcal{F}^n \to \mathcal{K}_0^n$,
\item[(\romannumeral4)]
des applications de $\mathcal{O}$-modules
$\mathcal{G}^n \to \mathcal{F}^n$,
\end{enumerate}
tels que $\mathcal{F}^n = 0$ pour $n \gg 0$, que les diagrammes
$$
\xymatrix{
\mathcal{G}^n \ar[r] \ar[d] & \mathcal{G}^{n + 1} \ar[d] \\
\mathcal{F}^n \ar[r] & \mathcal{F}^{n + 1}
}
$$
soient commutatifs pour tout $n$, que la composée
$\mathcal{G}^n \to \mathcal{F}^n \to \mathcal{K}_0^n$ soit l'application
donnée $\mathcal{G}^n \to \mathcal{K}_0^n$, et qu'en posant
$\mathcal{F}^n_0 = \mathcal{F}^n \otimes_\mathcal{O} \mathcal{O}_0$, nous
obtenions un complexe $\mathcal{F}_0^\bullet$ et une application de complexes
$\mathcal{F}_0^\bullet \to \mathcal{K}_0^\bullet$ qui soit un
quasi-isomorphisme.
\end{lemma}

\begin{proof}
Nous allons démontrer par récurrence descendante sur $e$ que nous pouvons
trouver $\mathcal{F}^n$,
$\mathcal{G}^n \to \mathcal{F}^n$ et
$\mathcal{F}^n \to \mathcal{F}^{n + 1}$ pour $n \geq e$, s'insérant dans un
diagramme commutatif
$$
\xymatrix{
\ldots \ar[r] &
\mathcal{G}^{e - 1} \ar[r] \ar@/_2pc/[dd] &
\mathcal{G}^e \ar[d] \ar[r] \ar@/_2pc/[dd] &
\mathcal{G}^{e + 1} \ar[d] \ar[r] \ar@/_2pc/[dd]|\hole &
\ldots \\
& &
\mathcal{F}^e \ar[d] \ar[r] &
\mathcal{F}^{e + 1} \ar[d] \ar[r] & \ldots \\
\ldots  \ar[r] &
\mathcal{K}_0^{e - 1} \ar[r] &
\mathcal{K}_0^e \ar[r] &
\mathcal{K}_0^{e + 1} \ar[r] & \ldots
}
$$
tel que $\mathcal{F}_0^\bullet$ soit un complexe et que l'application induite
$\mathcal{F}_0^\bullet \to \mathcal{K}_0^\bullet$ induise un isomorphisme sur
$H^n$ pour $n > e$ et une surjection pour $n = e$. Pour $e \gg 0$, ceci est
vrai, car les hypothèses (a) et (b) permettent de prendre
$\mathcal{F}^n = 0$ pour $n \geq e$.

\medskip\noindent
Étape de récurrence. Nous devons construire $\mathcal{F}^{e - 1}$ et les
applications $\mathcal{G}^{e - 1} \to \mathcal{F}^{e - 1}$,
$\mathcal{F}^{e - 1} \to \mathcal{F}^e$ et
$\mathcal{F}^{e - 1} \to \mathcal{K}_0^{e - 1}$.
Nous choisirons $\mathcal{F}^{e - 1} = A \oplus B \oplus C$ comme somme
directe de trois morceaux.

\medskip\noindent
Pour le premier, nous prenons $A = \mathcal{G}^{e - 1}$ et choisissons pour
$\mathcal{G}^{e - 1} \to \mathcal{F}^{e - 1}$ l'inclusion du premier facteur.
Les applications $A \to \mathcal{K}^{e - 1}_0$ et
$A \to \mathcal{F}^e$ seront les applications évidentes.

\medskip\noindent
Pour choisir $B$, considérons la surjection (par hypothèse de récurrence)
$$
\gamma :
\Ker(\mathcal{F}^e_0 \to \mathcal{F}^{e + 1}_0)
\longrightarrow
\Ker(\mathcal{K}^e_0 \to \mathcal{K}^{e + 1}_0)/
\Im(\mathcal{K}^{e - 1}_0 \to \mathcal{K}^e_0)
$$
Nous pouvons choisir un ensemble $I$, pour chaque $i \in I$ un objet $U_i$ de
$\mathcal{C}$, et des sections
$s_i \in \mathcal{F}^e(U_i)$,
$t_i \in \mathcal{K}^{e - 1}_0(U_i)$ telles que
\begin{enumerate}
\item $s_i$ s'envoie sur une section de $\Ker(\gamma) \subset
\Ker(\mathcal{F}^e_0 \to \mathcal{F}^{e + 1}_0)$,
\item $s_i$ et $t_i$ s'envoient sur la même section de
$\mathcal{K}^e_0$,
\item les sections $s_i$ engendrent $\Ker(\gamma)$ comme
$\mathcal{O}_0$-module.
\end{enumerate}
Nous omettons la justification complète ; on utilise le fait que
$\mathcal{F}^e \to \mathcal{F}^e_0$ est une application surjective de
faisceaux d'ensembles. Nous posons alors
$$
B = \bigoplus\nolimits_{i \in I} j_{U_i!}\mathcal{O}_{U_i}
$$
et définissons les applications $B \to \mathcal{F}^e$ et
$B \to \mathcal{K}_0^{e - 1}$ en utilisant $s_i$ et $t_i$ pour déterminer où
envoyer le facteur $j_{U_i!}\mathcal{O}_{U_i}$.

\medskip\noindent
Avec $\mathcal{F}^{e - 1} = A \oplus B$ et les applications ci-dessus, nous
obtenons un diagramme du type précédent pour $e - 1$, tel que
$\mathcal{F}_0^\bullet \to \mathcal{K}_0^\bullet$ induise un isomorphisme sur
$H^n$ pour $n \geq e$.
Pour que l'application soit surjective sur $H^{e - 1}$, nous choisissons le
facteur $C$ comme suit.
Choisissons un ensemble $J$, pour chaque $j \in J$ un objet $U_j$ de
$\mathcal{C}$ et une section $t_j$ de
$\Ker(\mathcal{K}^{e - 1}_0 \to \mathcal{K}^e_0)$ au-dessus de $U_j$, de sorte
que ces sections engendrent ce noyau sur $\mathcal{O}_0$. Posons alors
$$
C = \bigoplus\nolimits_{j \in J} j_{U_j!}\mathcal{O}_{U_j}
$$
et prenons l'application nulle $C \to \mathcal{F}^e$ ainsi que l'application
$C \to \mathcal{K}_0^{e - 1}$ définie en utilisant $s_j$ pour déterminer où
envoyer le facteur $j_{U_j!}\mathcal{O}_{U_j}$. Ceci achève l'étape de
récurrence en prenant $\mathcal{F}^{e - 1} = A \oplus B \oplus C$ et les
applications indiquées.
\end{proof}

\begin{lemma}
\label{defos-lemma-canonical-class}
Soit $\mathcal{C}$ un site. Soit $\mathcal{O} \to \mathcal{O}_0$ une
surjection de faisceaux d'anneaux dont le noyau est un faisceau d'idéaux
$\mathcal{I}$ de carré nul. Pour tout objet $K_0$ de
$D^-(\mathcal{O}_0)$, il existe une application canonique
$$
\omega(K_0) :
K_0 \longrightarrow
K_0 \otimes_{\mathcal{O}_0}^\mathbf{L} \mathcal{I}[2]
$$
dans $D(\mathcal{O}_0)$ telle que, pour toute application
$K_0 \to L_0$ dans $D^-(\mathcal{O}_0)$, le diagramme
$$
\xymatrix{
K_0 \ar[d] \ar[rr]_-{\omega(K_0)} & &
(K_0 \otimes^\mathbf{L}_{\mathcal{O}_0} \mathcal{I})[2] \ar[d] \\
L_0 \ar[rr]^-{\omega(L_0)} & &
(L_0 \otimes^\mathbf{L}_{\mathcal{O}_0} \mathcal{I})[2]
}
$$
soit commutatif.
\end{lemma}

\begin{proof}
Représentons $K_0$ par un complexe quelconque
$\mathcal{K}_0^\bullet$ de $\mathcal{O}_0$-modules.
Appliquons le Lemme \ref{defos-lemma-lift-complex} avec $\mathcal{G}^n = 0$ pour tout
$n$. Notons $d : \mathcal{F}^n \to \mathcal{F}^{n + 1}$ les applications
produites par le lemme. Nous voyons alors que
$d \circ d : \mathcal{F}^n \to \mathcal{F}^{n + 2}$ est nul modulo
$\mathcal{I}$. Comme $\mathcal{F}^n$ est plat, nous voyons que
$\mathcal{I}\mathcal{F}^n =
\mathcal{F}^n \otimes_{\mathcal{O}} \mathcal{I} =
\mathcal{F}^n_0 \otimes_{\mathcal{O}_0} \mathcal{I}$.
Nous obtenons donc une application canonique de complexes
$$
d \circ d : \mathcal{F}_0^\bullet
\longrightarrow
(\mathcal{F}_0^\bullet \otimes_{\mathcal{O}_0} \mathcal{I})[2]
$$
Puisque $\mathcal{F}_0^\bullet$ est un complexe borné supérieurement de
$\mathcal{O}_0$-modules plats, il est K-plat et peut servir à calculer le
produit tensoriel dérivé. De plus, l'application de complexes
$\mathcal{F}_0^\bullet \to \mathcal{K}_0^\bullet$ est un quasi-isomorphisme par
construction. La source et le but de l'application que nous venons de
construire représentent donc $K_0$ et
$K_0 \otimes_{\mathcal{O}_0}^\mathbf{L} \mathcal{I}[2]$ ; nous obtenons ainsi
notre application $\omega(K_0)$.

\medskip\noindent
Montrons que ce procédé est compatible avec les applications de complexes.
Soit $\mathcal{L}_0^\bullet$ un complexe représentant un autre objet de
$D^-(\mathcal{O}_0)$, et supposons que
$$
\mathcal{K}_0^\bullet \longrightarrow \mathcal{L}_0^\bullet
$$
soit une application de complexes. Appliquons le
Lemme \ref{defos-lemma-lift-complex} au complexe $\mathcal{L}_0^\bullet$, aux modules
plats $\mathcal{F}^n$, aux applications
$\mathcal{F}^n \to \mathcal{F}^{n + 1}$ et aux composées
$\mathcal{F}^n \to \mathcal{K}_0^n \to \mathcal{L}_0^n$ (nous nous excusons
de l'inversion des lettres employées).
Nous obtenons des modules plats $\mathcal{G}^n$, des applications
$\mathcal{F}^n \to \mathcal{G}^n$, des applications
$\mathcal{G}^n \to \mathcal{G}^{n + 1}$ et des applications
$\mathcal{G}^n \to \mathcal{L}_0^n$ possédant toutes les propriétés du lemme.
Il est alors clair que
$$
\xymatrix{
\mathcal{F}_0^\bullet \ar[d] \ar[r] &
(\mathcal{F}_0^\bullet \otimes_{\mathcal{O}_0} \mathcal{I})[2] \ar[d] \\
\mathcal{G}_0^\bullet \ar[r] &
(\mathcal{G}_0^\bullet \otimes_{\mathcal{O}_0} \mathcal{I})[2]
}
$$
est un diagramme commutatif de complexes.

\medskip\noindent
Pour voir que $\omega(K_0)$ est bien définie, supposons que nous disposions de
deux complexes $\mathcal{K}_0^\bullet$ et
$(\mathcal{K}'_0)^\bullet$ de $\mathcal{O}_0$-modules représentant $K_0$, et de
deux systèmes
$(\mathcal{F}^n, d : \mathcal{F}^n \to \mathcal{F}^{n + 1},
\mathcal{F}^n \to \mathcal{K}_0^n)$
et
$((\mathcal{F}')^n, d : (\mathcal{F}')^n \to (\mathcal{F}')^{n + 1},
(\mathcal{F}')^n \to \mathcal{K}_0^n)$
comme ci-dessus. Nous pouvons choisir un complexe
$(\mathcal{K}''_0)^\bullet$ et des quasi-isomorphismes
$\mathcal{K}_0^\bullet  \to (\mathcal{K}''_0)^\bullet$
et
$(\mathcal{K}'_0)^\bullet  \to (\mathcal{K}''_0)^\bullet$
réalisant le fait que les deux complexes représentent $K_0$ dans la catégorie
dérivée. Appliquons ensuite le résultat du paragraphe précédent à
$$
(\mathcal{K}_0)^\bullet \oplus (\mathcal{K}'_0)^\bullet
\longrightarrow
(\mathcal{K}''_0)^\bullet
$$
Nous obtenons ainsi un diagramme commutatif
$$
\xymatrix{
\mathcal{F}_0^\bullet \oplus (\mathcal{F}'_0)^\bullet
\ar[d] \ar[r] &
(\mathcal{F}_0^\bullet \otimes_{\mathcal{O}_0} \mathcal{I})[2] \oplus
((\mathcal{F}'_0)^\bullet \otimes_{\mathcal{O}_0} \mathcal{I})[2] \ar[d] \\
\mathcal{G}_0^\bullet \ar[r] &
(\mathcal{G}_0^\bullet \otimes_{\mathcal{O}_0} \mathcal{I})[2]
}
$$
Puisque les flèches verticales donnent des quasi-isomorphismes sur les facteurs,
nous en déduisons la commutativité voulue dans $D(\mathcal{O}_0)$.

\medskip\noindent
La bonne définition étant établie, l'assertion de compatibilité aux
applications découle du résultat du deuxième paragraphe.
\end{proof}

\begin{lemma}
\label{defos-lemma-induced-map}
Soit $(\mathcal{C}, \mathcal{O})$ un site annelé.
Soit $\alpha : K \to L$ une application dans $D^-(\mathcal{O})$.
Soit $\mathcal{F}$ un faisceau de $\mathcal{O}$-modules.
Soit $n \in \mathbf{Z}$.
\begin{enumerate}
\item Si $H^i(\alpha)$ est un isomorphisme pour $i \geq n$, alors
$H^i(\alpha \otimes_\mathcal{O}^\mathbf{L} \text{id}_\mathcal{F})$ est un
isomorphisme pour $i \geq n$.
\item Si $H^i(\alpha)$ est un isomorphisme pour $i > n$ et surjectif pour
$i = n$, alors
$H^i(\alpha \otimes_\mathcal{O}^\mathbf{L} \text{id}_\mathcal{F})$ est un
isomorphisme pour $i > n$ et surjectif pour $i = n$.
\end{enumerate}
\end{lemma}

\begin{proof}
Choisissons un triangle distingué
$$
K \to L \to C \to K[1]
$$
Dans le cas (2), nous voyons que $H^i(C) = 0$ pour $i \geq n$.
Par conséquent,
$H^i(C \otimes_\mathcal{O}^\mathbf{L} \mathcal{F}) = 0$ pour $i \geq n$
d'après le dual de Catégories dérivées,
Lemme \ref{derived-lemma-negative-vanishing}.
Cela montre à son tour que
$H^i(\alpha \otimes_\mathcal{O}^\mathbf{L} \text{id}_\mathcal{F})$ est un
isomorphisme pour $i > n$ et surjectif pour $i = n$.
Dans le cas (1), nous voyons en outre que
$H^{n - 1}(L) \to H^{n - 1}(C)$ est surjectif. En considérant le diagramme
$$
\xymatrix{
H^{n - 1}(L) \otimes_\mathcal{O} \mathcal{F} \ar[r] \ar[d] &
H^{n - 1}(C) \otimes_\mathcal{O} \mathcal{F} \ar@{=}[d] \\
H^{n - 1}(L \otimes_\mathcal{O}^\mathbf{L} \mathcal{F}) \ar[r] &
H^{n - 1}(C \otimes_\mathcal{O}^\mathbf{L} \mathcal{F})
}
$$
nous concluons que la flèche horizontale inférieure est surjective.
Combiné à ce qui précède, cela implique que
$H^n(\alpha \otimes_\mathcal{O}^\mathbf{L} \text{id}_\mathcal{F})$ est un
isomorphisme.
\end{proof}

\begin{lemma}
\label{defos-lemma-canonical-class-obstruction}
Soit $\mathcal{C}$ un site. Soit $\mathcal{O} \to \mathcal{O}_0$ une
surjection de faisceaux d'anneaux dont le noyau est un faisceau d'idéaux
$\mathcal{I}$ de carré nul. Pour tout objet $K_0$ de
$D^-(\mathcal{O}_0)$, les assertions suivantes sont équivalentes :
\begin{enumerate}
\item la classe
$\omega(K_0) \in
\Ext^2_{\mathcal{O}_0}(K_0, K_0 \otimes_{\mathcal{O}_0} \mathcal{I})$
construite dans le Lemme \ref{defos-lemma-canonical-class} est nulle,
\item il existe $K \in D^-(\mathcal{O})$ tel que
$K \otimes_\mathcal{O}^\mathbf{L} \mathcal{O}_0 = K_0$
dans $D(\mathcal{O}_0)$.
\end{enumerate}
\end{lemma}

\begin{proof}
Soit $K$ comme dans (2). Nous pouvons représenter $K$ par un complexe
$\mathcal{F}^\bullet$ de $\mathcal{O}$-modules plats, borné supérieurement.
Alors $\mathcal{F}_0^\bullet =
\mathcal{F}^\bullet \otimes_{\mathcal{O}} \mathcal{O}_0$ représente $K_0$ dans
$D(\mathcal{O}_0)$.
Puisque $d_{\mathcal{F}^\bullet} \circ d_{\mathcal{F}^\bullet} = 0$, car
$\mathcal{F}^\bullet$ est un complexe, la construction même de
$\omega(K_0)$ montre que cette classe est nulle.

\medskip\noindent
Supposons (1). Soient $\mathcal{F}^n$ et
$d : \mathcal{F}^n \to \mathcal{F}^{n + 1}$ comme dans la construction de
$\omega(K_0)$. La nullité de $\omega(K_0)$ implique que l'application
$$
\omega = d \circ d : \mathcal{F}_0^\bullet
\longrightarrow
(\mathcal{F}_0^\bullet \otimes_{\mathcal{O}_0} \mathcal{I})[2]
$$
est nulle dans $D(\mathcal{O}_0)$. Par définition de la catégorie dérivée
comme localisation de la catégorie homotopique des complexes de
$\mathcal{O}_0$-modules, il existe un quasi-isomorphisme
$\alpha : \mathcal{G}_0^\bullet \to \mathcal{F}_0^\bullet$ et des applications
de $\mathcal{O}_0$-modules
$h^n : \mathcal{G}_0^n \to
\mathcal{F}_0^{n + 1} \otimes_\mathcal{O} \mathcal{I}$ telles que
$$
\omega \circ \alpha =
d_{\mathcal{F}_0^\bullet \otimes \mathcal{I}} \circ h +
h \circ d_{\mathcal{G}_0^\bullet}
$$
Posons
$$
\mathcal{H}^n = \mathcal{F}^n \times_{\mathcal{F}^n_0} \mathcal{G}_0^n
$$
et définissons
$$
d' : \mathcal{H}^n \longrightarrow \mathcal{H}^{n + 1},\quad
(f^n, g_0^n) \longmapsto (d(f^n) - h^n(g_0^n), d(g_0^n))
$$
avec les notations évidentes, en utilisant
$\mathcal{F}_0^{n + 1} \otimes_{\mathcal{O}_0} \mathcal{I} =
\mathcal{F}^{n + 1} \otimes_\mathcal{O} \mathcal{I} =
\mathcal{I}\mathcal{F}^{n + 1} \subset \mathcal{F}^{n + 1}$.
On vérifie alors que $d' \circ d' = 0$ grâce au choix de $h^n$
et à la définition de $\omega$.
Ainsi, $\mathcal{H}^\bullet$ définit un objet de $D(\mathcal{O})$.
D'autre part, il existe une suite exacte courte de complexes de
$\mathcal{O}$-modules
$$
0 \to \mathcal{F}_0^\bullet \otimes_{\mathcal{O}_0} \mathcal{I} \to
\mathcal{H}^\bullet \to \mathcal{G}_0^\bullet \to 0
$$
Il reste à montrer que
$\mathcal{H}^\bullet \otimes_\mathcal{O}^\mathbf{L} \mathcal{O}_0$ est
isomorphe à $K_0$.
Choisissons un quasi-isomorphisme
$\mathcal{E}^\bullet \to \mathcal{H}^\bullet$, où
$\mathcal{E}^\bullet$ est un complexe borné supérieurement de
$\mathcal{O}$-modules plats. Nous obtenons un diagramme commutatif
$$
\xymatrix{
0 \ar[r] &
\mathcal{E}^\bullet \otimes_\mathcal{O} \mathcal{I} \ar[d]^\beta \ar[r] &
\mathcal{E}^\bullet \ar[d]^\gamma \ar[r] &
\mathcal{E}_0^\bullet \ar[d]^\delta \ar[r] &
0 \\
0 \ar[r] &
\mathcal{F}_0^\bullet \otimes_{\mathcal{O}_0} \mathcal{I} \ar[r] &
\mathcal{H}^\bullet \ar[r] &
\mathcal{G}_0^\bullet \ar[r] &
0
}
$$
Nous affirmons que $\delta$ est un quasi-isomorphisme. Comme
$H^i(\delta)$ est un isomorphisme pour $i \gg 0$, nous pouvons procéder par
récurrence descendante sur $n$, en supposant que $H^i(\delta)$ est un
isomorphisme pour $i \geq n$.
Observons que
$\mathcal{E}^\bullet \otimes_\mathcal{O} \mathcal{I}$ représente
$\mathcal{E}_0^\bullet \otimes_{\mathcal{O}_0}^\mathbf{L} \mathcal{I}$, que
$\mathcal{F}_0^\bullet \otimes_{\mathcal{O}_0} \mathcal{I}$ représente
$\mathcal{G}_0^\bullet \otimes_{\mathcal{O}_0}^\mathbf{L} \mathcal{I}$, et que
$\beta = \delta \otimes_{\mathcal{O}_0}^\mathbf{L} \text{id}_\mathcal{I}$
comme applications dans $D(\mathcal{O}_0)$. Cela est vrai parce que
$\beta =
(\alpha \otimes \text{id}_\mathcal{I})
\circ
(\delta \otimes \text{id}_\mathcal{I})$.
Supposons que $H^i(\delta)$ soit un isomorphisme en degrés $\geq n$.
Il en est alors de même pour $\beta$, d'après ce que nous venons de dire et le
Lemme \ref{defos-lemma-induced-map}.
Nous pouvons ensuite considérer le diagramme
$$
\xymatrix{
H^{n - 1}(\mathcal{E}^\bullet \otimes_\mathcal{O} \mathcal{I})
\ar[r] \ar[d]^{H^{n - 1}(\beta)} &
H^{n - 1}(\mathcal{E}^\bullet) \ar[r] \ar[d] &
H^{n - 1}(\mathcal{E}_0^\bullet) \ar[r] \ar[d]^{H^{n - 1}(\delta)} &
H^n(\mathcal{E}^\bullet \otimes_\mathcal{O} \mathcal{I})
\ar[r] \ar[d]^{H^n(\beta)} &
H^n(\mathcal{E}^\bullet) \ar[d] \\
H^{n - 1}(\mathcal{F}_0^\bullet \otimes_\mathcal{O} \mathcal{I}) \ar[r] &
H^{n - 1}(\mathcal{H}^\bullet) \ar[r] &
H^{n - 1}(\mathcal{G}_0^\bullet) \ar[r] &
H^n(\mathcal{F}_0^\bullet \otimes_\mathcal{O} \mathcal{I}) \ar[r] &
H^n(\mathcal{H}^\bullet)
}
$$
D'après Homologie, Lemme \ref{homology-lemma-four-lemma}, nous voyons que
$H^{n - 1}(\delta)$ est surjectif.
Cela implique à son tour que $H^{n - 1}(\beta)$ est surjectif, d'après le
Lemme \ref{defos-lemma-induced-map}.
En appliquant de nouveau Homologie, Lemme \ref{homology-lemma-four-lemma}, nous
voyons que $H^{n - 1}(\delta)$ est un isomorphisme.
L'affirmation résulte par récurrence ; ainsi, $\delta$ est un
quasi-isomorphisme, comme nous voulions le montrer.
\end{proof}

\begin{lemma}
\label{defos-lemma-lift-map-complexes}
Soit $\mathcal{C}$ un site. Soit $\mathcal{O} \to \mathcal{O}_0$ une
surjection de faisceaux d'anneaux. Supposons données les informations
suivantes :
\begin{enumerate}
\item un complexe de $\mathcal{O}$-modules $\mathcal{F}^\bullet$,
\item un complexe $\mathcal{K}_0^\bullet$ de $\mathcal{O}_0$-modules,
\item un quasi-isomorphisme $\mathcal{K}_0^\bullet \to
\mathcal{F}^\bullet \otimes_\mathcal{O} \mathcal{O}_0$,
\end{enumerate}
Il existe alors un quasi-isomorphisme
$\mathcal{G}^\bullet \to \mathcal{F}^\bullet$ tel que l'application de
complexes
$\mathcal{G}^\bullet  \otimes_\mathcal{O} \mathcal{O}_0 \to
\mathcal{F}^\bullet \otimes_\mathcal{O} \mathcal{O}_0$ se factorise par
$\mathcal{K}_0^\bullet$ dans la catégorie homotopique des complexes de
$\mathcal{O}_0$-modules.
\end{lemma}

\begin{proof}
Posons $\mathcal{F}_0^\bullet =
\mathcal{F}^\bullet \otimes_\mathcal{O} \mathcal{O}_0$.
D'après Catégories dérivées, Lemme \ref{derived-lemma-make-surjective}, il
existe une factorisation
$$
\mathcal{K}_0^\bullet \to \mathcal{L}_0^\bullet \to \mathcal{F}_0^\bullet
$$
de l'application donnée, telle que la première flèche admette un inverse à
homotopie près et que la seconde soit surjective et scindée terme à terme.
Nous pouvons donc supposer que
$\mathcal{K}_0^\bullet \to \mathcal{F}_0^\bullet$ est surjective terme à
terme. Dans ce cas, nous prenons
$$
\mathcal{G}^n = \mathcal{F}^n \times_{\mathcal{F}^n_0} \mathcal{K}_0^n
$$
et tout est clair.
\end{proof}

\begin{lemma}
\label{defos-lemma-inf-obs-map-defo-complex}
Soit $\mathcal{C}$ un site. Soit $\mathcal{O} \to \mathcal{O}_0$ une
surjection de faisceaux d'anneaux dont le noyau est un faisceau d'idéaux
$\mathcal{I}$ de carré nul. Soient $K, L \in D^-(\mathcal{O})$.
Posons $K_0 = K \otimes_\mathcal{O}^\mathbf{L} \mathcal{O}_0$ et
$L_0 = L \otimes_\mathcal{O}^\mathbf{L} \mathcal{O}_0$ dans
$D^-(\mathcal{O}_0)$. Étant donnée
$\alpha_0 : K_0 \to L_0$ dans $D(\mathcal{O}_0)$, il existe un élément
canonique
$$
o(\alpha_0) \in \Ext^1_{\mathcal{O}_0}(K_0,
L_0 \otimes_{\mathcal{O}_0}^\mathbf{L} \mathcal{I})
$$
dont l'annulation est une condition nécessaire et suffisante pour qu'il existe
une application $\alpha : K \to L$ dans $D(\mathcal{O})$ telle que
$\alpha_0 = \alpha \otimes_\mathcal{O}^\mathbf{L} \text{id}$.
\end{lemma}

\begin{proof}
Trouver une application $\alpha : K \to L$ relevant $\alpha_0$ revient à
trouver $\alpha : K \to L$ telle que la composée
$K \xrightarrow{\alpha} L \to L_0$ soit égale à la composée
$K \to K_0 \xrightarrow{\alpha_0} L_0$.
La suite exacte courte
$0 \to \mathcal{I} \to \mathcal{O} \to \mathcal{O}_0 \to 0$ donne lieu à un
triangle distingué canonique
$$
L \otimes_\mathcal{O}^\mathbf{L} \mathcal{I} \to
L \to
L_0 \to
(L \otimes_\mathcal{O}^\mathbf{L} \mathcal{I})[1]
$$
dans $D(\mathcal{O})$.
D'après Catégories dérivées,
Lemme \ref{derived-lemma-representable-homological}, la composée
$$
K \to K_0 \xrightarrow{\alpha_0} L_0 \to
(L \otimes_\mathcal{O}^\mathbf{L} \mathcal{I})[1]
$$
est nulle si et seulement si nous pouvons trouver
$\alpha : K \to L$ relevant $\alpha_0$. Cette composée est un élément de
$$
\Hom_{D(\mathcal{O})}(K, (L \otimes_\mathcal{O}^\mathbf{L} \mathcal{I})[1]) =
\Hom_{D(\mathcal{O}_0)}(K_0,
(L \otimes_\mathcal{O}^\mathbf{L} \mathcal{I})[1]) =
\Ext^1_{\mathcal{O}_0}(K_0,
L_0 \otimes_{\mathcal{O}_0}^\mathbf{L} \mathcal{I})
$$
par adjonction.
\end{proof}

\begin{lemma}
\label{defos-lemma-first-order-defos-complex}
Soit $\mathcal{C}$ un site. Soit $\mathcal{O} \to \mathcal{O}_0$ une
surjection de faisceaux d'anneaux dont le noyau est un faisceau d'idéaux
$\mathcal{I}$ de carré nul. Soit $K_0 \in D^-(\mathcal{O})$.
Un relèvement de $K_0$ est un couple $(K, \alpha_0)$ formé d'un objet $K$ de
$D^-(\mathcal{O})$ et d'un isomorphisme
$\alpha_0 : K \otimes_\mathcal{O}^\mathbf{L} \mathcal{O}_0 \to K_0$ dans
$D(\mathcal{O}_0)$.
\begin{enumerate}
\item Étant donné un relèvement $(K, \alpha)$, le groupe des automorphismes du
couple est canoniquement le conoyau d'une application
$$
\Ext^{-1}_{\mathcal{O}_0}(K_0, K_0)
\longrightarrow
\Hom_{\mathcal{O}_0}(K_0, K_0 \otimes_{\mathcal{O}_0}^\mathbf{L} \mathcal{I})
$$
\item S'il existe un relèvement, alors l'ensemble des classes d'isomorphisme de
relèvements est un espace principal homogène sous
$\Ext^1_{\mathcal{O}_0}(K_0,
K_0 \otimes_{\mathcal{O}_0}^\mathbf{L} \mathcal{I})$.
\end{enumerate}
\end{lemma}

\begin{proof}
Un automorphisme de $(K, \alpha)$ est une application
$\varphi : K \to K$ dans $D(\mathcal{O})$ telle que
$\varphi \otimes_\mathcal{O} \text{id}_{\mathcal{O}_0} = \text{id}$.
Cela revient à dire que
$$
K \xrightarrow{\varphi - \text{id}} K \to
K \otimes_\mathcal{O}^\mathbf{L} \mathcal{O}_0
$$
est nulle. Nous en concluons que le groupe des automorphismes est le conoyau
d'une application
$$
\Hom_\mathcal{O}(K, K_0[-1])
\longrightarrow
\Hom_\mathcal{O}(K, K_0 \otimes_{\mathcal{O}_0}^\mathbf{L} \mathcal{I})
$$
en vertu du triangle distingué
$$
K \otimes_\mathcal{O}^\mathbf{L} \mathcal{I} \to
K \to
K \otimes_\mathcal{O}^\mathbf{L} \mathcal{O}_0 \to
(K \otimes_\mathcal{O}^\mathbf{L} \mathcal{I})[1]
$$
dans $D(\mathcal{O})$ et de Catégories dérivées,
Lemme \ref{derived-lemma-representable-homological}.
Pour traduire ceci dans les groupes du lemme, nous utilisons l'adjonction entre
le foncteur de restriction $D(\mathcal{O}_0) \to D(\mathcal{O})$ et
$- \otimes_\mathcal{O} \mathcal{O}_0 : D(\mathcal{O}) \to D(\mathcal{O}_0)$.
Cela démontre (1).

\medskip\noindent
Démonstration de (2).
Supposons que $K_0 = K \otimes_\mathcal{O}^\mathbf{L} \mathcal{O}_0$ dans
$D(\mathcal{O})$. D'après le
Lemme \ref{defos-lemma-inf-obs-map-defo-complex}, l'application qui associe à un
relèvement $(K', \alpha_0)$ l'obstruction $o(\alpha_0)$ au relèvement de
$\alpha_0$ définit une injection canonique de l'ensemble des classes
d'isomorphisme de couples dans
$\Ext^1_{\mathcal{O}_0}(K_0,
K_0 \otimes_{\mathcal{O}_0}^\mathbf{L} \mathcal{I})$.
Pour achever la démonstration, montrons qu'elle est surjective.
Choisissons
$\xi : K_0 \to (K_0 \otimes_{\mathcal{O}_0}^\mathbf{L} \mathcal{I})[1]$
dans le groupe $\Ext^1$ du lemme.
Choisissons un complexe $\mathcal{F}^\bullet$ de $\mathcal{O}$-modules plats,
borné supérieurement, représentant $K$.
L'application $\xi$ peut être représentée sous la forme $t \circ s^{-1}$, où
$s : \mathcal{K}_0^\bullet \to \mathcal{F}_0^\bullet$ est un
quasi-isomorphisme et
$t : \mathcal{K}_0^\bullet \to
\mathcal{F}_0^\bullet \otimes_{\mathcal{O}_0} \mathcal{I}[1]$ est une
application de complexes.
D'après le Lemme \ref{defos-lemma-lift-map-complexes}, nous pouvons supposer qu'il
existe un quasi-isomorphisme
$\mathcal{G}^\bullet \to \mathcal{F}^\bullet$ de complexes de
$\mathcal{O}$-modules tel que
$\mathcal{G}_0^\bullet \to \mathcal{F}_0^\bullet$ se factorise par $s$ à
homotopie près.
Nous pouvons et allons remplacer $\mathcal{G}^\bullet$ par un complexe borné
supérieurement de $\mathcal{O}$-modules plats (en choisissant un
quasi-isomorphisme d'un tel complexe vers $\mathcal{G}^\bullet$, puis en
effectuant le remplacement).
Nous voyons alors que $\xi$ est représenté par une application de complexes
$t : \mathcal{G}_0^\bullet \to
\mathcal{F}_0^\bullet \otimes_{\mathcal{O}_0} \mathcal{I}[1]$ et par le
quasi-isomorphisme
$\mathcal{G}_0^\bullet \to \mathcal{F}_0^\bullet$.
Posons
$$
\mathcal{H}^n = \mathcal{F}^n \times_{\mathcal{F}_0^n} \mathcal{G}_0^n
$$
avec les différentielles
$$
\mathcal{H}^n \to \mathcal{H}^{n + 1},\quad
(f^n, g_0^n) \mapsto (d(f^n) + t(g_0^n), d(g_0^n))
$$
Cela a un sens, car
$\mathcal{F}_0^{n + 1} \otimes_{\mathcal{O}_0} \mathcal{I} =
\mathcal{F}^{n + 1} \otimes_\mathcal{O} \mathcal{I} =
\mathcal{I}\mathcal{F}^{n + 1} \subset \mathcal{F}^{n + 1}$.
Nous omettons le calcul montrant que $\mathcal{H}^\bullet$ est un complexe de
$\mathcal{O}$-modules. Par construction, il existe une suite exacte courte
$$
0 \to \mathcal{F}_0^\bullet \otimes_{\mathcal{O}_0} \mathcal{I} \to
\mathcal{H}^\bullet \to \mathcal{G}_0^\bullet \to 0
$$
de complexes de $\mathcal{O}$-modules.
Exactement comme dans la démonstration du
Lemme \ref{defos-lemma-canonical-class-obstruction}, on montre que cette suite induit
un isomorphisme
$\alpha_0 :
\mathcal{H}^\bullet \otimes_\mathcal{O}^\mathbf{L} \mathcal{O}_0 \to
\mathcal{G}_0^\bullet$ dans $D(\mathcal{O}_0)$.
Autrement dit, nous avons construit un couple
$(\mathcal{H}^\bullet, \alpha_0)$.
Nous omettons la vérification de $o(\alpha_0) = \xi$ ; indication :
$o(\alpha_0)$ peut être calculé explicitement dans ce cas, car nous avons des
applications $\mathcal{H}^n \to \mathcal{F}^n$ (non compatibles aux
différentielles) relevant les composantes de $\alpha_0$.
Cela achève la démonstration.
\end{proof}

















% OK, start here.
%

\chapter{Le complexe cotangent}



\phantomsection
\label{cotangent-section-phantom}


\section{Introduction}
\label{cotangent-section-introduction}

\noindent
Le but de ce chapitre est de construire le complexe cotangent d'un
homomorphisme d'anneaux, d'un morphisme de schémas et d'un morphisme d'espaces
algébriques. On pourra consulter les notes \cite{quillenhomology}, l'article
\cite{quillencohomology}, ainsi que les livres \cite{Andre} et
\cite{cotangent}.




\section{Conseils au lecteur}
\label{cotangent-section-advice-reader}

\noindent
En rédigeant ce chapitre, nous avons cherché à réduire au minimum l'emploi
des techniques simpliciales. Nous considérons le choix d'une
{\it résolution} $P_\bullet$ d'un anneau $B$ sur un anneau $A$ comme un outil
pour calculer l'{\it homologie} des faisceaux abéliens sur la catégorie
$\mathcal{C}_{B/A}$ ; voir la remarque \ref{cotangent-remark-resolution}. Ce rôle est
analogue à celui d'un ``bon recouvrement'' dans le calcul de la cohomologie au
moyen du complexe de {\v C}ech. Pour quelques notions d'homologie sur les
catégories, voir Cohomologie sur les sites, section
\ref{sites-cohomology-section-homology}. Le foncteur dérivé gauche $L\pi_!$
joue pour l'homologie le rôle que $R\Gamma(\mathcal{C}_{B/A}, -)$ joue pour
la cohomologie. La catégorie $\mathcal{C}_{B/A}$, étudiée dans la section
\ref{cotangent-section-compute-L-pi-shriek}, est l'opposée de la catégorie des
factorisations $A \to P \to B$, où $P$ est une algèbre polynomiale sur $A$.
Cette catégorie est munie d'homomorphismes de faisceaux d'anneaux
$$
\underline{A} \longrightarrow \mathcal{O} \longrightarrow \underline{B}
$$
où, sur l'objet $U =  (P \to B)$, on a $\mathcal{O}(U) = P$.
On obtient alors le complexe cotangent de $B$ sur $A$ sous la forme
$$
L_{B/A} =
L\pi_!(\Omega_{\mathcal{O}/\underline{A}} \otimes_\mathcal{O} \underline{B})
$$
voir le lemme \ref{cotangent-lemma-compute-cotangent-complex}. Nous nous sommes efforcés
d'adopter systématiquement ce point de vue pour démontrer les propriétés
élémentaires des complexes cotangents d'homomorphismes d'anneaux. Tous les
résultats peuvent notamment se démontrer sans recourir à l'existence de
résolutions standard, bien que nous ne l'ayons pas fait. La théorie est tout à
fait satisfaisante, à ceci près que la démonstration du triangle fondamental
(proposition \ref{cotangent-proposition-triangle}) utilise peut-être un peu plus de
théorie des foncteurs dérivés gauches. Afin d'offrir une autre voie au
lecteur, nous donnons dans les remarques \ref{cotangent-remark-triangle} et
\ref{cotangent-remark-explicit-map} une esquisse assez complète d'une approche fondée
sur des propriétés simples des résolutions standard.

\medskip\noindent
Pour traiter le complexe cotangent des morphismes de topos annelés, des
morphismes de schémas, des morphismes d'espaces algébriques, etc., notre
méthode consiste à déduire autant que possible du cas des ``simples
homomorphismes d'anneaux'' étudié ci-dessus.





\section{Le complexe cotangent d'un homomorphisme d'anneaux}
\label{cotangent-section-cotangent-ring-map}

\noindent
Soit $A$ un anneau. Notons $\textit{Alg}_A$ la catégorie des $A$-algèbres.
Considérons le couple de foncteurs adjoints $(U, V)$ où
$V : \textit{Alg}_A \to \textit{Ensembles}$ est le foncteur d'oubli et où
$U : \textit{Ensembles} \to \textit{Alg}_A$ associe à un ensemble $E$ l'algèbre
polynomiale $A[E]$ sur $E$ au-dessus de $A$. Soit $X_\bullet$ l'objet
simplicial de $\text{Fun}(\textit{Alg}_A, \textit{Alg}_A)$ construit dans
Objets simpliciaux, section \ref{simplicial-section-standard}.

\medskip\noindent
Considérons une $A$-algèbre $B$. Notons $P_\bullet = X_\bullet(B)$ la
$A$-algèbre simpliciale ainsi obtenue. Rappelons que $P_0 = A[B]$,
$P_1 = A[A[B]]$, et ainsi de suite. En particulier, chaque terme $P_n$ est
une $A$-algèbre polynomiale. Rappelons aussi qu'il existe une augmentation
$$
\epsilon : P_\bullet \longrightarrow B
$$
où l'on considère $B$ comme une $A$-algèbre simpliciale constante.

\begin{definition}
\label{cotangent-definition-standard-resolution}
Soit $A \to B$ un homomorphisme d'anneaux. La {\it résolution standard de
$B$ sur $A$} est l'augmentation $\epsilon : P_\bullet \to B$ dont les termes
sont
$$
P_0 = A[B],\quad P_1 = A[A[B]],\quad \ldots
$$
et dont les applications sont construites dans Objets simpliciaux, exemple
\ref{simplicial-example-polynomial-algebra-maps}.
\end{definition}

\noindent
Nous verrons que la résolution standard permet, dans certaines situations,
de calculer les foncteurs dérivés gauches.

\begin{definition}
\label{cotangent-definition-cotangent-complex-ring-map}
Le {\it complexe cotangent} $L_{B/A}$ d'un homomorphisme d'anneaux $A \to B$
est le complexe de $B$-modules associé au $B$-module simplicial
$$
\Omega_{P_\bullet/A} \otimes_{P_\bullet, \epsilon} B
$$
où $\epsilon : P_\bullet \to B$ est la résolution standard de $B$ sur $A$.
\end{definition}

\noindent
Dans Objets simpliciaux, section \ref{simplicial-section-complexes}, nous
associons un complexe de chaînes à un module simplicial ; ici, toutefois, nous
travaillons avec des complexes de cochaînes. Ainsi, le terme $L_{B/A}^{-n}$
en degré $-n$ est le $B$-module
$\Omega_{P_n/A} \otimes_{P_n, \epsilon_n} B$, et $L_{B/A}^m = 0$ pour
$m > 0$.

\begin{remark}
\label{cotangent-remark-variant-cotangent-complex}
Soit $A \to B$ un homomorphisme d'anneaux. Soit $\mathcal{A}$ la catégorie
des flèches $\psi : C \to B$ de $A$-algèbres, et soit $\mathcal{S}$ la
catégorie des applications $E \to B$, où $E$ est un ensemble. On dispose de
foncteurs adjoints $V : \mathcal{A} \to \mathcal{S}$ (le foncteur d'oubli) et
$U : \mathcal{S} \to \mathcal{A}$, qui envoie $E \to B$ sur
$A[E] \to B$. Soit $X_\bullet$ l'objet simplicial de
$\text{Fun}(\mathcal{A}, \mathcal{A})$ construit dans Objets simpliciaux,
section \ref{simplicial-section-standard}. Le diagramme
$$
\xymatrix{
\mathcal{A} \ar[d] \ar[r] & \mathcal{S} \ar@<1ex>[l] \ar[d] \\
\textit{Alg}_A \ar[r] & \textit{Ensembles} \ar@<1ex>[l]
}
$$
est commutatif. Il en résulte que $X_\bullet(\text{id}_B : B \to B)$ est égal
à la résolution standard de $B$ sur $A$.
\end{remark}

\begin{lemma}
\label{cotangent-lemma-colimit-cotangent-complex}
Soit $A_i \to B_i$ un système d'homomorphismes d'anneaux indexé par un
ensemble ordonné filtrant $I$. Alors
$\colim L_{B_i/A_i} = L_{\colim B_i/\colim A_i}$.
\end{lemma}

\begin{proof}
Cela est vrai parce que le foncteur d'oubli
$V : A\textit{-Alg} \to \textit{Ensembles}$ et son adjoint
$U : \textit{Ensembles} \to A\textit{-Alg}$ commutent aux colimites filtrantes.
Il en va de même du foncteur $B/A \mapsto \Omega_{B/A}$
(Algèbre, lemme \ref{algebra-lemma-colimit-differentials}).
\end{proof}





\section{Résolutions simpliciales et foncteur dérivé gauche}
\label{cotangent-section-compute-L-pi-shriek}

\noindent
Soit $A \to B$ un homomorphisme d'anneaux. Considérons la catégorie dont les
objets sont les homomorphismes de $A$-algèbres $\alpha : P \to B$, où $P$ est
une algèbre polynomiale sur $A$ (en un certain ensemble de
variables\footnote{Il suffit de considérer des ensembles dont le cardinal est
au plus celui de $B$.}), et dont les morphismes
$s : (\alpha : P \to B) \to (\alpha' : P' \to B)$ sont les homomorphismes de
$A$-algèbres $s : P \to P'$ tels que $\alpha' \circ s = \alpha$.
Notons $\mathcal{C} = \mathcal{C}_{B/A}$ la catégorie {\bf opposée} de cette
catégorie. Ce passage à l'opposée vient de ce que nous voulons considérer les
objets $(P, \alpha)$ comme correspondant au diagramme de schémas affines
$$
\xymatrix{
\Spec(B) \ar[d] \ar[r] & \Spec(P) \ar[ld] \\
\Spec(A)
}
$$
Munissons $\mathcal{C}$ de la topologie chaotique
(Sites, exemple \ref{sites-example-indiscrete}), c'est-à-dire munissons
$\mathcal{C}$ de la structure de site dans laquelle les recouvrements sont
donnés par les identités, de
sorte que tout préfaisceau est un faisceau. Munissons en outre $\mathcal{C}$
de deux faisceaux d'anneaux. Le premier est le faisceau $\mathcal{O}$ qui
associe $P$ à l'objet $(P, \alpha)$. Le second est le faisceau constant $B$,
que nous noterons $\underline{B}$. On obtient le diagramme suivant de
morphismes de topos annelés
\begin{equation}
\label{cotangent-equation-pi}
\vcenter{
\xymatrix{
(\Sh(\mathcal{C}), \underline{B}) \ar[r]_i \ar[d]_\pi &
(\Sh(\mathcal{C}), \mathcal{O}) \\
(\Sh(*), B)
}
}
\end{equation}
Le morphisme $i$ est l'identité sur les topos sous-jacents, et
$i^\sharp : \mathcal{O} \to \underline{B}$ est l'homomorphisme évident.
L'application $\pi$ est celle de Cohomologie sur les sites, exemple
\ref{sites-cohomology-example-category-to-point}. Les foncteurs dérivés
suivants joueront un rôle important :
$
Li^* : D(\mathcal{O}) \longrightarrow D(\underline{B})
$
adjoint à gauche de $Ri_* = i_* : D(\underline{B}) \to D(\mathcal{O})$, et
$
L\pi_! : D(\underline{B}) \longrightarrow D(B)
$
adjoint à gauche de $\pi^* = \pi^{-1} : D(B) \to D(\underline{B})$.

\begin{lemma}
\label{cotangent-lemma-identify-pi-shriek}
Avec les notations précédentes, soit $P_\bullet$ une $A$-algèbre simpliciale
munie d'une augmentation $\epsilon : P_\bullet \to B$.
Supposons que chaque $P_n$ soit une algèbre polynomiale sur $A$ et que
$\epsilon$ soit une fibration de Kan triviale sur les ensembles simpliciaux
sous-jacents. Alors
$$
L\pi_!(\mathcal{F}) = \mathcal{F}(P_\bullet, \epsilon)
$$
dans $D(\textit{Ab})$, resp.\ $D(B)$, fonctoriellement en $\mathcal{F}$ dans
$\textit{Ab}(\mathcal{C})$, resp.\ $\textit{Mod}(\underline{B})$.
\end{lemma}

\begin{proof}
Nous utiliserons pour cela le critère de Cohomologie sur les sites, lemme
\ref{sites-cohomology-lemma-compute-by-cosimplicial-resolution}.
Étant donné un objet $U = (Q, \beta)$ de $\mathcal{C}$, il faut montrer que
$$
S_\bullet = \Mor_\mathcal{C}((P_\bullet, \epsilon), (Q, \beta))
$$
est homotopiquement équivalent à un singleton.
Écrivons $Q = A[E]$ pour un certain ensemble $E$ (ce qui est possible par
notre choix de la catégorie $\mathcal{C}$). On voit que
$$
S_\bullet = \Mor_{\textit{Ensembles}}((E, \beta|_E), (P_\bullet, \epsilon))
$$
Soit $*$ l'ensemble simplicial constant de valeur un singleton. Pour $b \in B$,
soit $F_{b, \bullet}$ l'ensemble simplicial défini par le diagramme cartésien
$$
\xymatrix{
F_{b, \bullet} \ar[r] \ar[d] & P_\bullet \ar[d]_\epsilon \\
{*} \ar[r]^b & B
}
$$
Avec cette notation, $S_\bullet = \prod_{e \in E} F_{\beta(e), \bullet}$.
Comme $\epsilon$ est par hypothèse une fibration de Kan triviale,
$F_{b, \bullet} \to *$ est une fibration de Kan triviale
(Objets simpliciaux, lemme
\ref{simplicial-lemma-trivial-kan-base-change}). Ainsi,
$S_\bullet \to *$ est une fibration de Kan triviale (Objets simpliciaux,
lemme \ref{simplicial-lemma-product-trivial-kan}). Par conséquent,
$S_\bullet$ est homotopiquement équivalent à $*$ (Objets simpliciaux, lemme
\ref{simplicial-lemma-trivial-kan-homotopy}).
\end{proof}

\noindent
En particulier, on peut employer la résolution standard de $B$ sur $A$ pour
calculer le foncteur dérivé gauche correspondant.

\begin{lemma}
\label{cotangent-lemma-pi-shriek-standard}
Soit $A \to B$ un homomorphisme d'anneaux. Soit
$\epsilon : P_\bullet \to B$ la résolution standard de $B$ sur $A$.
Soit $\pi$ comme dans (\ref{cotangent-equation-pi}). Alors
$$
L\pi_!(\mathcal{F}) = \mathcal{F}(P_\bullet, \epsilon)
$$
dans $D(\textit{Ab})$, resp.\ $D(B)$, fonctoriellement en $\mathcal{F}$ dans
$\textit{Ab}(\mathcal{C})$, resp.\ $\textit{Mod}(\underline{B})$.
\end{lemma}

\begin{proof}[Première démonstration]
Appliquons le lemme \ref{cotangent-lemma-identify-pi-shriek}.
Puisque les termes $P_n$ sont des algèbres polynomiales, la première hypothèse
de ce lemme est satisfaite. Démontrons la seconde. D'après Objets simpliciaux,
lemme \ref{simplicial-lemma-standard-simplicial-homotopy}, l'application
$\epsilon$ est une équivalence d'homotopie des ensembles simpliciaux
sous-jacents. D'après Objets simpliciaux, lemme
\ref{simplicial-lemma-homotopy-equivalence}, il en résulte que $\epsilon$
induit un quasi-isomorphisme des complexes de groupes abéliens associés.
D'après Objets simpliciaux, lemme
\ref{simplicial-lemma-qis-simplicial-abelian-groups}, cela implique que
$\epsilon$ est une fibration de Kan triviale des ensembles simpliciaux
sous-jacents.
\end{proof}

\begin{proof}[Seconde démonstration]
Nous utiliserons le critère de Cohomologie sur les sites, lemme
\ref{sites-cohomology-lemma-compute-by-cosimplicial-resolution}.
Soit $U = (Q, \beta)$ un objet de $\mathcal{C}$.
Il nous faut montrer que
$$
S_\bullet = \Mor_\mathcal{C}((P_\bullet, \epsilon), (Q, \beta))
$$
est homotopiquement équivalent à un singleton. Écrivons $Q = A[E]$ pour un
certain ensemble $E$ (ce qui est possible par notre choix de la catégorie
$\mathcal{C}$). Avec les notations de la remarque
\ref{cotangent-remark-variant-cotangent-complex}, on voit que
$$
S_\bullet = \Mor_\mathcal{S}((E \to B), i(P_\bullet \to B))
$$
D'après Objets simpliciaux, lemme
\ref{simplicial-lemma-standard-simplicial-homotopy}, l'application
$i(P_\bullet \to B) \to i(B \to B)$ est une équivalence d'homotopie dans
$\mathcal{S}$. Par conséquent, $S_\bullet$ est homotopiquement équivalent à
$$
\Mor_\mathcal{S}((E \to B), (B \to B)) = \{*\}
$$
comme voulu.
\end{proof}

\begin{lemma}
\label{cotangent-lemma-compute-cotangent-complex}
Soit $A \to B$ un homomorphisme d'anneaux. Soient $\pi$ et $i$ comme dans
(\ref{cotangent-equation-pi}). Il existe un isomorphisme canonique
$$
L_{B/A} = L\pi_!(Li^*\Omega_{\mathcal{O}/A}) =
L\pi_!(i^*\Omega_{\mathcal{O}/A}) =
L\pi_!(\Omega_{\mathcal{O}/A} \otimes_\mathcal{O} \underline{B})
$$
dans $D(B)$.
\end{lemma}

\begin{proof}
Pour un objet $\alpha : P \to B$ de la catégorie $\mathcal{C}$, le module
$\Omega_{P/A}$ est un $P$-module libre. Ainsi,
$\Omega_{\mathcal{O}/A}$ est un $\mathcal{O}$-module plat. Par conséquent,
$Li^*\Omega_{\mathcal{O}/A} = i^*\Omega_{\mathcal{O}/A}$ est le faisceau de
$\underline{B}$-modules qui associe à $\alpha : P \to A$ le $B$-module
$\Omega_{P/A} \otimes_{P, \alpha} B$. D'après le lemme
\ref{cotangent-lemma-pi-shriek-standard}, le membre de droite se calcule en évaluant ce
faisceau sur la résolution standard ; c'est précisément notre définition du
membre de gauche (définition \ref{cotangent-definition-cotangent-complex-ring-map}).
\end{proof}

\begin{lemma}
\label{cotangent-lemma-pi-lower-shriek-constant-sheaf}
Si $A \to B$ est un homomorphisme d'anneaux, alors
$L\pi_!(\pi^{-1}M) = M$, où $\pi$ est comme dans (\ref{cotangent-equation-pi}).
\end{lemma}

\begin{proof}
Cela résulte du lemme \ref{cotangent-lemma-identify-pi-shriek}, qui affirme que
$L\pi_!(\pi^{-1}M)$ se calcule par $(\pi^{-1}M)(P_\bullet, \epsilon)$, objet
simplicial constant de valeur $M$.
\end{proof}

\begin{lemma}
\label{cotangent-lemma-identify-H0}
Si $A \to B$ est un homomorphisme d'anneaux, alors
$H^0(L_{B/A}) = \Omega_{B/A}$.
\end{lemma}

\begin{proof}
Nous allons le démontrer par un calcul direct, en utilisant l'identification
du lemme \ref{cotangent-lemma-compute-cotangent-complex}. Il existe manifestement un
homomorphisme de $\Omega_{\mathcal{O}/A} \otimes \underline{B}$ vers le
faisceau constant de valeur $\Omega_{B/A}$. Il induit donc un homomorphisme
$$
H^0(L_{B/A}) = H^0(L\pi_!(\Omega_{\mathcal{O}/A} \otimes \underline{B}))
= \pi_!(\Omega_{\mathcal{O}/A} \otimes \underline{B}) \to \Omega_{B/A}
$$
En choisissant un objet $P \to B$ de $\mathcal{C}_{B/A}$ tel que $P \to B$
soit surjectif, on voit que cet homomorphisme est surjectif (d'après Algèbre,
lemme \ref{algebra-lemma-differential-surjective}). Pour montrer son
injectivité, supposons que $P \to B$ soit un objet de $\mathcal{C}_{B/A}$ et
que $\xi \in \Omega_{P/A} \otimes_P B$ soit un élément dont l'image dans
$\Omega_{B/A}$ est nulle. Choisissons d'abord une factorisation
$P \to P' \to B$ telle que $P' \to B$ soit surjectif et que $P'$ soit une
algèbre polynomiale sur $A$. Nous pouvons remplacer $P$ par $P'$. Si
$B = P/I$, alors le noyau de
$\Omega_{P/A} \otimes_P B \to \Omega_{B/A}$ est l'image de $I/I^2$ (Algèbre,
lemme \ref{algebra-lemma-differential-seq}). Disons que $\xi$ est l'image de
$f \in I$. Considérons alors les deux homomorphismes
$a, b : P' = P[x] \to P$ : le premier envoie $x$ sur $0$ et le second envoie
$x$ sur $f$ (dans les deux cas, $P[x] \to B$ envoie $x$ sur zéro). On voit que
$\xi$ et $0$ sont les images de $\text{d}x \otimes 1$ dans
$\Omega_{P'/A} \otimes_{P'} B$. Ainsi, $\xi$ et $0$ ont la même image dans la
colimite (voir Cohomologie sur les sites, exemple
\ref{sites-cohomology-example-category-to-point})
$\pi_!(\Omega_{\mathcal{O}/A} \otimes \underline{B})$, comme voulu.
\end{proof}

\begin{lemma}
\label{cotangent-lemma-pi-lower-shriek-polynomial-algebra}
Si $B$ est une algèbre polynomiale sur l'anneau $A$, alors, pour $\pi$ comme
dans (\ref{cotangent-equation-pi}), le foncteur $\pi_!$ est exact et
$\pi_!\mathcal{F} = \mathcal{F}(B \to B)$.
\end{lemma}

\begin{proof}
Cela résulte du lemme \ref{cotangent-lemma-identify-pi-shriek}, qui affirme que
l'algèbre simpliciale constante de valeur $B$ peut servir à calculer $L\pi_!$.
\end{proof}

\begin{lemma}
\label{cotangent-lemma-cotangent-complex-polynomial-algebra}
Si $B$ est une algèbre polynomiale sur l'anneau $A$, alors $L_{B/A}$ est
quasi-isomorphe à $\Omega_{B/A}[0]$.
\end{lemma}

\begin{proof}
Cela résulte immédiatement des lemmes
\ref{cotangent-lemma-compute-cotangent-complex} et
\ref{cotangent-lemma-pi-lower-shriek-polynomial-algebra}.
\end{proof}





\section{Construction d'une résolution}
\label{cotangent-section-polynomial}

\noindent
Dans le cas noethérien de type fini, on peut construire une ``petite''
résolution simpliciale des homomorphismes d'anneaux de type fini.

\begin{lemma}
\label{cotangent-lemma-polynomial}
Soit $A$ un anneau noethérien. Soit $A \to B$ un homomorphisme d'anneaux de
type fini. Soit $\mathcal{A}$ la catégorie des homomorphismes de $A$-algèbres
$C \to B$. Soit $n \geq 0$, et soit $P_\bullet$ un objet simplicial de
$\mathcal{A}$ tel que :
\begin{enumerate}
\item $P_\bullet \to B$ est une fibration de Kan triviale d'ensembles
simpliciaux ;
\item $P_k$ est de type fini sur $A$ pour $k \leq n$ ;
\item $P_\bullet = \text{cosk}_n \text{sk}_n P_\bullet$ comme objets
simpliciaux de $\mathcal{A}$.
\end{enumerate}
Alors $P_{n + 1}$ est une $A$-algèbre de type fini.
\end{lemma}

\begin{proof}
Bien que la démonstration que nous allons donner soit directe, elle est un peu
laborieuse. Pour en éclairer l'idée, expliquons d'abord ce qui se passe pour
les petites valeurs de $n$, avant de traiter le cas général. Si, par exemple,
$n = 0$, alors (3) signifie que $P_1 = P_0 \times_B P_0$. Comme
l'homomorphisme d'anneaux $P_0 \to B$ est surjectif, cette algèbre est de type
fini sur $A$ d'après Compléments d'algèbre, lemme
\ref{more-algebra-lemma-fibre-product-finite-type}.

\medskip\noindent
Si $n = 1$, alors (3) signifie que
$$
P_2 = \{(f_0, f_1, f_2) \in P_1^3 \mid
d_0f_0 = d_0f_1,\ d_1f_0 = d_0f_2,\ d_1f_1 = d_1f_2 \}
$$
où les égalités ont lieu dans $P_0$. Remarquons que le triplet
$$
(d_0f_0, d_1f_0, d_1f_1) = (d_0f_1, d_0f_2, d_1f_2)
$$
est un élément du produit fibré $P_0 \times_B P_0 \times_B P_0$ au-dessus de
$B$, car les applications $d_i : P_1 \to P_0$ sont des morphismes au-dessus
de $B$. On obtient donc une application
$$
\psi : P_2 \longrightarrow P_0 \times_B P_0 \times_B P_0
$$
La fibre de $\psi$ au-dessus d'un élément
$(g_0, g_1, g_2) \in P_0 \times_B P_0 \times_B P_0$
est l'ensemble des triplets $(f_0, f_1, f_2)$ de $1$-simplexes tels que
$(d_0, d_1)(f_0) = (g_0, g_1)$, $(d_0, d_1)(f_1) = (g_0, g_2)$ et
$(d_0, d_1)(f_2) = (g_1, g_2)$. Puisque $P_\bullet \to B$ est une fibration
de Kan triviale, l'application $(d_0, d_1) : P_1 \to P_0 \times_B P_0$ est
surjective. Ainsi, $P_2$ s'insère dans le diagramme cartésien
$$
\xymatrix{
P_2 \ar[d] \ar[r] & P_1^3 \ar[d] \\
P_0 \times_B P_0 \times_B P_0 \ar[r] & (P_0 \times_B P_0)^3
}
$$
On conclut par Compléments d'algèbre, lemme
\ref{more-algebra-lemma-formal-consequence}. Le cas général est analogue, mais
demande un peu plus de notation.

\medskip\noindent
Traitons le cas $n > 1$. D'après Objets simpliciaux, lemme
\ref{simplicial-lemma-cosk-above-object}, la condition
$P_\bullet = \text{cosk}_n \text{sk}_n P_\bullet$ implique que la même égalité
vaut dans la catégorie des $A$-algèbres simpliciales, puis dans celle des
ensembles (car le foncteur d'oubli des $A$-algèbres vers les ensembles commute
aux limites). Ainsi,
$$
P_{n + 1} =
\Mor(\Delta[n + 1], P_\bullet) =
\Mor(\text{sk}_n \Delta[n + 1], \text{sk}_n P_\bullet)
$$
d'après Objets simpliciaux, lemme \ref{simplicial-lemma-simplex-map}, et
l'équation (\ref{simplicial-equation-cosk}). Nous allons démontrer par
récurrence sur $1 \leq k < m \leq n + 1$ que l'anneau
$$
Q_{k, m} = \Mor(\text{sk}_k \Delta[m], \text{sk}_k P_\bullet)
$$
est de type fini sur $A$. Le cas $k = 1$, $1 < m \leq n + 1$, est entièrement
analogue à la discussion ci-dessus du cas $n = 1$. Plus précisément, on a un
diagramme cartésien
$$
\xymatrix{
Q_{1, m} \ar[d] \ar[r] & P_1^N \ar[d] \\
P_0 \times_B \ldots \times_B P_0 \ar[r] & (P_0 \times_B P_0)^N
}
$$
où $N = {m + 1 \choose 2}$. On conclut comme précédemment.

\medskip\noindent
Soit $1 \leq k_0 \leq n$, et supposons que $Q_{k, m}$ soit de type fini sur
$A$ pour tous $1 \leq k \leq k_0$ et $k < m \leq n + 1$.
Pour $k_0 + 1 < m \leq n + 1$, nous affirmons qu'il existe un carré cartésien
$$
\xymatrix{
Q_{k_0 + 1, m} \ar[d] \ar[r] & P_{k_0 + 1}^N \ar[d] \\
Q_{k_0, m} \ar[r] & Q_{k_0, k_0 + 1}^N
}
$$
où $N$ est le nombre de $(k_0 + 1)$-simplexes non dégénérés de $\Delta[m]$.
En effet, considérons un élément de $Q_{k_0 + 1, m}$ comme une application $f$
du $(k_0 + 1)$-squelette de $\Delta[m]$ vers $P_\bullet$. On peut restreindre
$f$ au $k_0$-squelette, ce qui donne l'application verticale gauche du
diagramme. Sa
restriction à chaque $(k_0 + 1)$-simplexe non dégénéré donne la flèche
horizontale supérieure. De plus, se donner un tel $f$ revient à se donner sa
restriction au $k_0$-squelette et à chaque $(k_0 + 1)$-face non dégénérée,
pourvu que ces restrictions coïncident sur leurs intersections ; c'est
exactement le contenu du diagramme. Enfin, le fait que
$P_\bullet \to B$ soit une fibration de Kan triviale implique que
l'application
$$
P_{k_0} \to Q_{k_0, k_0 + 1} = \Mor(\partial \Delta[k_0 + 1], P_\bullet)
$$
est surjective, car toute application $\partial \Delta[k_0 + 1] \to B$ se
prolonge en une application $\Delta[k_0 + 1] \to B$ pour $k_0 \geq 1$ (nous
omettons un petit argument sur les ensembles simpliciaux constants). Par
l'hypothèse de récurrence, les anneaux $Q_{k_0, m}$,
$Q_{k_0, k_0 + 1}$ sont des $A$-algèbres de type fini ; il en va donc de même
de $Q_{k_0 + 1, m}$, encore une fois d'après Compléments d'algèbre, lemme
\ref{more-algebra-lemma-formal-consequence}.
\end{proof}

\begin{proposition}
\label{cotangent-proposition-polynomial}
Soit $A$ un anneau noethérien. Soit $A \to B$ un homomorphisme d'anneaux de
type fini. Il existe une $A$-algèbre simpliciale $P_\bullet$ munie d'une
augmentation $\epsilon : P_\bullet \to B$ telle que chaque $P_n$ soit une
algèbre polynomiale de type fini sur $A$ et que $\epsilon$ soit une fibration
de Kan triviale d'ensembles simpliciaux.
\end{proposition}

\begin{proof}
Soit $\mathcal{A}$ la catégorie des homomorphismes de $A$-algèbres $C \to B$.
Dans cette démonstration, les objets simpliciaux ainsi que les foncteurs
squelette et cosquelette seront pris dans cette catégorie.

\medskip\noindent
Choisissons une algèbre polynomiale $P_0$ de type fini sur $A$ et une
surjection $P_0 \to B$. En première approximation, prenons
$P_\bullet = \text{cosk}_0(P_0)$. Autrement dit, $P_\bullet$ est la
$A$-algèbre simpliciale dont les termes sont
$P_n = P_0 \times_A \ldots \times_A P_0$. (Dans le dernier paragraphe de la
démonstration, cet objet simplicial sera noté $P^0_\bullet$.) D'après Objets
simpliciaux, lemme \ref{simplicial-lemma-cosk-minus-one-equivalence},
l'application $P_\bullet \to B$ est une fibration de Kan triviale d'ensembles
simpliciaux. Remarquons aussi que
$P_\bullet = \text{cosk}_0 \text{sk}_0 P_\bullet$.

\medskip\noindent
Supposons que, pour un certain $n \geq 0$, nous ayons construit $P_\bullet$
(qui sera noté $P^n_\bullet$ dans le dernier paragraphe de la démonstration)
de telle sorte que :
\begin{enumerate}
\item[(a)] $P_\bullet \to B$ est une fibration de Kan triviale d'ensembles
simpliciaux ;
\item[(b)] $P_k$ est une algèbre polynomiale de type fini pour
$0 \leq k \leq n$ ;
\item[(c)] $P_\bullet = \text{cosk}_n \text{sk}_n P_\bullet$
\end{enumerate}
D'après le lemme \ref{cotangent-lemma-polynomial}, il existe une algèbre polynomiale de
type fini $Q$ sur $A$ et une surjection $Q \to P_{n + 1}$. Comme $P_n$ est une
algèbre polynomiale, les homomorphismes de $A$-algèbres
$s_i : P_n \to P_{n + 1}$ se relèvent en des homomorphismes
$s'_i : P_n \to Q$. Définissons $d'_j : Q \to P_n$ comme la composée de
$Q \to P_{n + 1}$ et de $d_j : P_{n + 1} \to P_n$. On obtient un objet
simplicial tronqué $P'_\bullet$ de $\mathcal{A}$ en posant $P'_k = P_k$ pour
$k \leq n$ et $P'_{n + 1} = Q$, en prenant les morphismes $d'_i = d_i$ et
$s'_i = s_i$ aux degrés $k \leq n - 1$, puis les morphismes $d'_j$ et $s'_i$
au degré $n$. Prolongeons-le en un objet simplicial complet $P'_\bullet$ de
$\mathcal{A}$ au moyen de $\text{cosk}_{n + 1}$. Par fonctorialité des
foncteurs cosquelette, il existe un morphisme d'objets simpliciaux
$P'_\bullet \to P_\bullet$ qui prolonge le morphisme donné d'objets simpliciaux
$(n + 1)$-tronqués. (Ce morphisme sera noté
$P^{n + 1}_\bullet \to P^n_\bullet$ dans le dernier paragraphe de la
démonstration.)

\medskip\noindent
Notons que les conditions (b) et (c) sont satisfaites par $P'_\bullet$ après
remplacement de $n$ par $n + 1$. Nous affirmons que l'application
$P'_\bullet \to P_\bullet$ satisfait aux hypothèses (1), (2), (3) et (4) de
Objets simpliciaux, lemme \ref{simplicial-lemma-section}, avec $n + 1$ à la
place de $n$. Les conditions (1) et (2) sont vraies par construction. D'après
Objets simpliciaux, lemme \ref{simplicial-lemma-cosk-above-object}, on a
$P_\bullet = \text{cosk}_{n + 1}\text{sk}_{n + 1}P_\bullet$
et
$P'_\bullet = \text{cosk}_{n + 1}\text{sk}_{n + 1}P'_\bullet$
non seulement dans $\mathcal{A}$, mais aussi dans la catégorie des
$A$-algèbres, puis dans celle des ensembles (car le foncteur d'oubli des
$A$-algèbres vers les ensembles commute à toutes les limites). Cela démontre
(3) et (4). Le lemme s'applique donc, et
$P'_\bullet \to P_\bullet$ est une fibration de Kan triviale. D'après Objets
simpliciaux, lemme \ref{simplicial-lemma-trivial-kan-composition}, on en déduit
que $P'_\bullet \to B$ est une fibration de Kan triviale ; la condition (a)
est donc également satisfaite.

\medskip\noindent
Pour achever la démonstration, prenons la limite projective
$P_\bullet = \lim P^n_\bullet$ de la suite d'algèbres simpliciales
$$
\ldots \to P^2_\bullet \to P^1_\bullet \to P^0_\bullet
$$
construite ci-dessus. L'application $P_\bullet \to B$ est une fibration de Kan
triviale d'après Objets simpliciaux, lemme
\ref{simplicial-lemma-limit-trivial-kan}. Or la construction ci-dessus se
stabilise, en chaque degré, en une algèbre polynomiale de type fini fixée,
comme voulu.
\end{proof}

\begin{lemma}
\label{cotangent-lemma-pi-shriek-finite}
Soit $A$ un anneau noethérien. Soit $A \to B$ un homomorphisme d'anneaux de
type fini. Soient $\pi$, $\underline{B}$ comme dans (\ref{cotangent-equation-pi}).
Si $\mathcal{F}$ est un $\underline{B}$-module tel que
$\mathcal{F}(P, \alpha)$ soit un $B$-module fini pour tout
$\alpha : P = A[x_1, \ldots, x_n] \to B$, alors les modules de cohomologie de
$L\pi_!(\mathcal{F})$ sont des $B$-modules finis.
\end{lemma}

\begin{proof}
D'après le lemme \ref{cotangent-lemma-identify-pi-shriek} et la proposition
\ref{cotangent-proposition-polynomial}, on peut calculer $L\pi_!(\mathcal{F})$ au moyen
d'un complexe construit à partir des valeurs de $\mathcal{F}$ sur les algèbres
polynomiales de type fini.
\end{proof}

\begin{lemma}
\label{cotangent-lemma-cotangent-finite}
Soit $A$ un anneau noethérien. Soit $A \to B$ un homomorphisme d'anneaux de
type fini. Alors $H^n(L_{B/A})$ est un $B$-module fini pour tout
$n \in \mathbf{Z}$.
\end{lemma}

\begin{proof}
Appliquer les lemmes \ref{cotangent-lemma-compute-cotangent-complex} et
\ref{cotangent-lemma-pi-shriek-finite}.
\end{proof}

\begin{remark}[Resolutions]
\label{cotangent-remark-resolution}
Soit $A \to B$ un homomorphisme d'anneaux quelconque. Appelons une
$A$-algèbre simpliciale augmentée $\epsilon : P_\bullet \to B$ une
{\it résolution de $B$ sur $A$} si chaque $P_n$ est une algèbre polynomiale
et si $\epsilon$ est une fibration de Kan triviale d'ensembles simpliciaux.
Si $P_\bullet \to B$ est une augmentation d'une $A$-algèbre simpliciale telle
que chaque $P_n$ soit une algèbre polynomiale se surjectant sur $B$, alors les
conditions suivantes sont équivalentes :
\begin{enumerate}
\item $\epsilon : P_\bullet \to B$ est une résolution de $B$ sur $A$ ;
\item $\epsilon : P_\bullet \to B$ est un quasi-isomorphisme sur les
complexes associés ;
\item $\epsilon : P_\bullet \to B$ induit une équivalence d'homotopie
d'ensembles simpliciaux.
\end{enumerate}
Pour le voir, utiliser Objets simpliciaux, lemmes
\ref{simplicial-lemma-trivial-kan-homotopy},
\ref{simplicial-lemma-homotopy-equivalence} et
\ref{simplicial-lemma-qis-simplicial-abelian-groups}.
Une résolution $P_\bullet$ de $B$ sur $A$ fournit un objet cosimplicial
$U_\bullet$ de $\mathcal{C}_{B/A}$ comme dans Cohomologie sur les sites, lemme
\ref{sites-cohomology-lemma-compute-by-cosimplicial-resolution}
et il s'ensuit que
$$
L\pi_!\mathcal{F} = \mathcal{F}(P_\bullet)
$$
fonctoriellement en $\mathcal{F}$ ; voir le lemme
\ref{cotangent-lemma-identify-pi-shriek}. La partie formelle de la démonstration de la
proposition \ref{cotangent-proposition-polynomial} montre l'existence de résolutions.
Nous avons aussi vu, dans la première démonstration du lemme
\ref{cotangent-lemma-pi-shriek-standard}, que la résolution standard de $B$ sur $A$ est
une résolution (de sorte que cette terminologie ne crée pas de conflit).
Cependant, l'argument de la démonstration de la proposition
\ref{cotangent-proposition-polynomial} établit l'existence de résolutions sans faire
appel aux calculs simpliciaux d'Objets simpliciaux, section
\ref{simplicial-section-standard}. De plus, pour {\it tout} choix de
résolution, on dispose d'un isomorphisme canonique
$$
L_{B/A} = \Omega_{P_\bullet/A} \otimes_{P_\bullet, \epsilon} B
$$
dans $D(B)$ d'après le lemme \ref{cotangent-lemma-compute-cotangent-complex}. La liberté
de choisir une résolution arbitraire peut être fort utile.
\end{remark}

\begin{lemma}
\label{cotangent-lemma-O-homology-B-homology}
Soit $A \to B$ un homomorphisme d'anneaux. Soient $\pi$, $\mathcal{O}$,
$\underline{B}$ comme dans (\ref{cotangent-equation-pi}). Pour tout $\mathcal{O}$-module
$\mathcal{F}$, on a
$$
L\pi_!(\mathcal{F}) = L\pi_!(Li^*\mathcal{F}) =
L\pi_!(\mathcal{F} \otimes_\mathcal{O}^\mathbf{L} \underline{B})
$$
dans $D(\textit{Ab})$.
\end{lemma}

\begin{proof}
Il suffit de vérifier que les hypothèses de Cohomologie sur les sites, lemme
\ref{sites-cohomology-lemma-O-homology-qis}
sont satisfaites pour $\mathcal{O} \to \underline{B}$ sur
$\mathcal{C}_{B/A}$. Nous utiliserons sans plus le signaler les résultats de
la remarque \ref{cotangent-remark-resolution}. Choisissons une résolution $P_\bullet$
de $B$ sur $A$ afin d'obtenir un objet cosimplicial convenable $U_\bullet$ de
$\mathcal{C}_{B/A}$. Puisque $P_\bullet \to B$ induit un quasi-isomorphisme
sur les complexes associés de groupes abéliens, on voit que
$L\pi_!\mathcal{O} = B$. D'autre part, $L\pi_!\underline{B}$ est calculé par
$\underline{B}(U_\bullet) = B$. Cela vérifie la seconde hypothèse de
Cohomologie sur les sites, lemme
\ref{sites-cohomology-lemma-O-homology-qis}
et achève la démonstration.
\end{proof}

\begin{lemma}
\label{cotangent-lemma-apply-O-B-comparison}
Soit $A \to B$ un homomorphisme d'anneaux. Soient $\pi$, $\mathcal{O}$,
$\underline{B}$ comme dans (\ref{cotangent-equation-pi}). On a
$$
L\pi_!(\mathcal{O}) = L\pi_!(\underline{B}) = B
\quad\text{et}\quad
L_{B/A} = L\pi_!(\Omega_{\mathcal{O}/A} \otimes_\mathcal{O} \underline{B}) =
L\pi_!(\Omega_{\mathcal{O}/A})
$$
dans $D(\textit{Ab})$.
\end{lemma}

\begin{proof}
C'est une application immédiate du lemme
\ref{cotangent-lemma-O-homology-B-homology} (et la première égalité à droite est le
lemme \ref{cotangent-lemma-compute-cotangent-complex}).
\end{proof}

\noindent
Voici un cas particulier du triangle fondamental dont la démonstration est
facile.

\begin{lemma}
\label{cotangent-lemma-special-case-triangle}
Soient $A \to B \to C$ des homomorphismes d'anneaux. Si $B$ est une algèbre
polynomiale sur $A$, alors il existe un triangle distingué 
$L_{B/A} \otimes_B^\mathbf{L} C \to L_{C/A} \to L_{C/B} \to
L_{B/A} \otimes_B^\mathbf{L} C[1]$ dans $D(C)$.
\end{lemma}

\begin{proof}
Nous utiliserons sans plus le signaler les observations de la remarque
\ref{cotangent-remark-resolution}. Choisissons une résolution
$\epsilon : P_\bullet \to C$ de $C$ sur $B$ (par exemple la résolution
standard). Puisque $B$ est une algèbre polynomiale sur $A$, on voit que
$P_\bullet$ est aussi une résolution de $C$ sur $A$. Ainsi, $L_{C/A}$ est
calculé par
$\Omega_{P_\bullet/A} \otimes_{P_\bullet, \epsilon} C$
et $L_{C/B}$ est calculé par
$\Omega_{P_\bullet/B} \otimes_{P_\bullet, \epsilon} C$.
Comme, pour chaque $n$, on a la suite exacte courte
$0 \to \Omega_{B/A} \otimes_B P_n \to \Omega_{P_n/A} \to \Omega_{P_n/B} \to 0$
(Algèbre, lemme \ref{algebra-lemma-ses-formally-smooth}) et comme
$L_{B/A} = \Omega_{B/A}[0]$ (lemme
\ref{cotangent-lemma-cotangent-complex-polynomial-algebra}), on obtient le résultat.
\end{proof}

\begin{example}
\label{cotangent-example-resolution-length-two}
Soit $A \to B$ un homomorphisme d'anneaux. Dans cet exemple, nous allons 
construire une résolution ``explicite'' $P_\bullet$ de $B$ sur $A$ de longueur
$2$. Pour ce faire, nous suivons la procédure de la démonstration de la
proposition \ref{cotangent-proposition-polynomial} ; voir aussi la discussion de la
remarque \ref{cotangent-remark-resolution}.

\medskip\noindent
Choisissons une surjection $P_0 = A[u_i] \to B$, où les $u_i$ forment un
ensemble de variables. Choisissons des générateurs $f_t \in P_0$, $t \in T$,
de l'idéal $\Ker(P_0 \to B)$. Prenons $P_1 = A[u_i, x_t]$, avec pour
applications de face $d_0$ et $d_1$ les uniques homomorphismes de $A$-algèbres
tels que $d_j(u_i) = u_i$, $d_0(x_t) = 0$ et $d_1(x_t) = f_t$.
L'application $s_0 : P_0 \to P_1$ est l'unique homomorphisme de $A$-algèbres
tel que $s_0(u_i) = u_i$. Il est clair que
$$
P_1 \xrightarrow{d_0 - d_1} P_0 \to B \to 0
$$
est exacte ; en particulier, l'application
$(d_0, d_1) : P_1 \to P_0 \times_B P_0$ est surjective. Ainsi, si
$P_\bullet$ désigne l'algèbre simpliciale $1$-tronquée sur $A$ donnée par $P_0$,
$P_1$, $d_0$, $d_1$ et $s_0$, alors l'augmentation
$\text{cosk}_1(P_\bullet) \to B$ est une fibration de Kan triviale. L'étape
suivante de la procédure de la démonstration de la proposition
\ref{cotangent-proposition-polynomial} consiste à choisir une algèbre polynomiale $P_2$
et une surjection
$$
P_2 \longrightarrow \text{cosk}_1(P_\bullet)_2
$$
Rappelons que
$$
\text{cosk}_1(P_\bullet)_2 =
\{(g_0, g_1, g_2) \in P_1^3 \mid d_0(g_0) = d_0(g_1),
d_1(g_0) = d_0(g_2), d_1(g_1) = d_1(g_2)\}
$$
En considérant $g_i \in P_1$ comme un polynôme en les $x_t$, les conditions
sont
$$
g_0(0) = g_1(0),\quad
g_0(f_t) = g_2(0),\quad
g_1(f_t) = g_2(f_t)
$$
Ainsi, $\text{cosk}_1(P_\bullet)_2$ contient les éléments
$y_t = (x_t, x_t, f_t)$ et $z_t = (0, x_t, x_t)$. Tout élément $G$ de
$\text{cosk}_1(P_\bullet)_2$ est de la forme $G = H + (0, 0, g)$, où $H$
appartient à l'image de
$A[u_i, y_t, z_t] \to \text{cosk}_1(P_\bullet)_2$. Ici, $g \in P_1$ est un
polynôme de terme constant nul tel que $g(f_t) = 0$ dans $P_0$. Observons que
\begin{enumerate}
\item $g = x_t x_{t'} - f_t x_{t'}$ ;
\item $g = \sum r_t x_t$ avec $r_t \in P_0$ si $\sum r_t f_t = 0$ dans $P_0$,
\end{enumerate}
sont des éléments de $P_1$ de la forme voulue. Posons
$$
Rel = \Ker(\bigoplus\nolimits_{t \in T} P_0 \longrightarrow P_0),\quad
(r_t) \longmapsto \sum r_tf_t
$$
Prenons $P_2 = A[u_i, y_t, z_t, v_r, w_{t, t'}]$, où
$r = (r_t) \in Rel$, muni de l'application
$$
P_2 \longrightarrow \text{cosk}_1(P_\bullet)_2
$$
définie par $y_t  \mapsto (x_t, x_t, f_t)$,
$z_t \mapsto (0, x_t, x_t)$,
$v_r \mapsto (0, 0, \sum r_t x_t)$ et
$w_{t, t'} \mapsto (0, 0, x_t x_{t'} - f_t x_{t'})$. Un calcul (omis) montre
que cette application est surjective. Notre choix de l'application affichée
ci-dessus détermine les applications $d_0, d_1, d_2 : P_2 \to P_1$.
Enfin, la procédure de la démonstration de la proposition
\ref{cotangent-proposition-polynomial} nous dit de choisir les applications
$s_0, s_1 : P_1 \to P_2$ relevant les deux applications
$P_1 \to \text{cosk}_1(P_\bullet)_2$. Il est clair que l'on peut prendre pour
$s_i$ les uniques homomorphismes de $A$-algèbres déterminés par
$s_0(x_t) = y_t$ et $s_1(x_t) = z_t$.
\end{example}








\section{Fonctorialité}
\label{cotangent-section-functoriality}

\noindent
Dans cette section, nous considérons un carré commutatif
\begin{equation}
\label{cotangent-equation-commutative-square}
\vcenter{
\xymatrix{
B \ar[r] & B' \\
A \ar[u] \ar[r] & A' \ar[u]
}
}
\end{equation}
d'homomorphismes d'anneaux. Nous affirmons qu'il existe une application
canonique $B$-linéaire de complexes
$$
L_{B/A} \longrightarrow L_{B'/A'}
$$
associée à ce diagramme. En effet, si $P_\bullet \to B$ est la résolution
standard de $B$ sur $A$ et si $P'_\bullet \to B'$ est la résolution standard
de $B'$ sur $A'$, alors il existe une application canonique
$P_\bullet \to P'_\bullet$
de $A$-algèbres simpliciales compatible avec les augmentations
$P_\bullet \to B$ et $P'_\bullet \to B'$. Cela se voit au moyen de la
construction des résolutions standard dans Objets simpliciaux, section
\ref{simplicial-section-standard}, mais, dans le cas particulier qui nous
occupe, il suffit probablement de dire simplement que les applications
$$
P_0 = A[B] \longrightarrow A'[B'] = P'_0,\quad
P_1 = A[A[B]] \longrightarrow A'[A'[B']] = P'_1,
$$
et ainsi de suite sont données par les homomorphismes prescrits $A \to A'$ et
$B \to B'$. L'application souhaitée $L_{B/A} \to L_{B'/A'}$ provient alors
des applications associées $\Omega_{P_n/A} \to \Omega_{P'_n/A'}$.

\medskip\noindent
On peut décrire autrement l'application de fonctorialité comme suit. Soient
$\mathcal{C} = \mathcal{C}_{B/A}$ et $\mathcal{C}' = \mathcal{C}_{B'/A}'$
les catégories considérées dans la section
\ref{cotangent-section-compute-L-pi-shriek}. Il existe un foncteur
$$
u : \mathcal{C} \longrightarrow \mathcal{C}',\quad
(P, \alpha) \longmapsto (P \otimes_A A', c \circ (\alpha \otimes 1))
$$
où $c : B \otimes_A A' \to B'$ est l'application évidente. Comme il est
expliqué dans Cohomologie sur les sites, exemple
\ref{sites-cohomology-example-morphism-categories}
on obtient un morphisme de topos $g : \Sh(\mathcal{C}) \to \Sh(\mathcal{C}')$
et un diagramme commutatif d'applications de topos annelés
\begin{equation}
\label{cotangent-equation-double-square}
\vcenter{
\xymatrix{
(\Sh(\mathcal{C}'), \underline{B}) \ar[d]_{\pi'} &
(\Sh(\mathcal{C}'), \underline{B'}) \ar[d]_{\pi'} \ar[l]^h &
(\Sh(\mathcal{C}), \underline{B'}) \ar[d]_\pi \ar[l]^g \\
(\Sh(*), B) &
(\Sh(*), B') \ar[l]_f &
(\Sh(*), B') \ar[l]
}
}
\end{equation}
Ici, $h$ est l'identité sur les topos sous-jacents et est donné, sur les
faisceaux d'anneaux, par l'homomorphisme d'anneaux $B \to B'$. 
D'après Cohomologie sur les sites, remarque
\ref{sites-cohomology-remark-morphism-fibred-categories}
étant donnés $\mathcal{F}$ sur $\mathcal{C}$ et $\mathcal{F}'$ sur
$\mathcal{C}'$, ainsi qu'une transformation
$t : \mathcal{F} \to g^{-1}\mathcal{F}'$, on obtient une application canonique
$L\pi_!(\mathcal{F}) \to L\pi'_!(\mathcal{F}')$. Appliquons ceci aux faisceaux
$$
\mathcal{F} : (P, \alpha) \mapsto \Omega_{P/A} \otimes_P B,\quad
\mathcal{F}' : (P', \alpha') \mapsto \Omega_{P'/A'} \otimes_{P'} B',
$$
et à la transformation $t$ donnée par les applications canoniques
$$
\Omega_{P/A} \otimes_P B \longrightarrow
\Omega_{P \otimes_A A'/A'} \otimes_{P \otimes_A A'} B'
$$
pour obtenir une application canonique
$$
L\pi_!(\Omega_{\mathcal{O}/A} \otimes_\mathcal{O} \underline{B})
\longrightarrow
L\pi'_!(\Omega_{\mathcal{O}'/A'} \otimes_{\mathcal{O}'} \underline{B'})
$$
D'après le lemme \ref{cotangent-lemma-compute-cotangent-complex}, cela fournit
$L_{B/A} \to L_{B'/A'}$. Nous omettons de vérifier que cette application
coïncide avec celle définie ci-dessus à l'aide de résolutions simpliciales.

\begin{lemma}
\label{cotangent-lemma-flat-base-change}
Supposons que (\ref{cotangent-equation-commutative-square}) induise un quasi-isomorphisme
$B \otimes_A^\mathbf{L} A' = B'$. Alors, avec les notations de
(\ref{cotangent-equation-double-square}) et pour
$\mathcal{F}' \in \textit{Ab}(\mathcal{C}')$, on a
$L\pi_!(g^{-1}\mathcal{F}') = L\pi'_!(\mathcal{F}')$.
\end{lemma}

\begin{proof}
Nous utiliserons sans plus le signaler les résultats de la remarque
\ref{cotangent-remark-resolution}. Appliquons Cohomologie sur les sites, lemme
\ref{sites-cohomology-lemma-get-it-now}. Soit $P_\bullet \to B$ une
résolution. Il suffit de montrer que
$u(P_\bullet) = P_\bullet \otimes_A A' \to B'$ est un quasi-isomorphisme. Le
complexe de $A$-modules $s(P_\bullet)$ associé à $P_\bullet$ (considéré comme
un $A$-module simplicial) est une résolution par des $A$-modules libres de
$B$. En effet, $P_n$ est un $A$-module libre et
$s(P_\bullet) \to B$ est un quasi-isomorphisme. Ainsi,
$B \otimes_A^\mathbf{L} A'$ est calculé par
$s(P_\bullet) \otimes_A A' = s(P_\bullet \otimes_A A')$. L'hypothèse du
lemme signifie donc que $\epsilon' : P_\bullet \otimes_A A' \to B'$ est un
quasi-isomorphisme.
\end{proof}

\noindent
Le lemme suivant s'applique en particulier lorsque $A \to A'$ est plat et que
$B' = B \otimes_A A'$ (changement de base plat).

\begin{lemma}
\label{cotangent-lemma-flat-base-change-cotangent-complex}
Si (\ref{cotangent-equation-commutative-square}) induit un quasi-isomorphisme
$B \otimes_A^\mathbf{L} A' = B'$, alors l'application de fonctorialité induit
un isomorphisme
$$
L_{B/A} \otimes_B^\mathbf{L} B' \longrightarrow L_{B'/A'}
$$
\end{lemma}

\begin{proof}
Nous utiliserons les notations introduites dans l'équation
(\ref{cotangent-equation-double-square}). On a
$$
L_{B/A} \otimes_B^\mathbf{L} B' =
L\pi_!(\Omega_{\mathcal{O}/A} \otimes_\mathcal{O} \underline{B})
\otimes_B^\mathbf{L} B' =
L\pi_!(Lh^*(\Omega_{\mathcal{O}/A} \otimes_\mathcal{O} \underline{B}))
$$
où la première égalité résulte du lemme
\ref{cotangent-lemma-compute-cotangent-complex} et la seconde de Cohomologie sur les
sites, lemme \ref{sites-cohomology-lemma-change-of-rings}. Comme
$\Omega_{\mathcal{O}/A}$ est un $\mathcal{O}$-module plat, on voit que
$\Omega_{\mathcal{O}/A} \otimes_\mathcal{O} \underline{B}$ est un
$\underline{B}$-module plat. Ainsi,
$Lh^*(\Omega_{\mathcal{O}/A} \otimes_\mathcal{O} \underline{B}) =
\Omega_{\mathcal{O}/A} \otimes_\mathcal{O} \underline{B'}$
qui est égal à
$g^{-1}(\Omega_{\mathcal{O}'/A'} \otimes_{\mathcal{O}'} \underline{B'})$
par inspection. On conclut par le lemme \ref{cotangent-lemma-flat-base-change} et par le
fait que $L_{B'/A'}$ est calculé par
$L\pi'_!(\Omega_{\mathcal{O}'/A'} \otimes_{\mathcal{O}'} \underline{B'})$.
\end{proof}

\begin{remark}
\label{cotangent-remark-homotopy-triangle}
Supposons donné un carré (\ref{cotangent-equation-commutative-square}) tel qu'il existe
une flèche $\kappa : B \to A'$ qui rende le diagramme commutatif :
$$
\xymatrix{
B \ar[r]_\beta \ar[rd]_\kappa & B' \\
A \ar[u] \ar[r]^\alpha & A' \ar[u]
}
$$
Dans ce cas, nous affirmons que l'application de fonctorialité
$P_\bullet \to P'_\bullet$ est homotope à la composée
$P_\bullet \to B \to A' \to P'_\bullet$. En effet, au moyen de $\kappa$,
l'application de fonctorialité se factorise sous la forme
$$
P_\bullet \to P_{A'/A', \bullet} \to P'_\bullet
$$
où $P_{A'/A', \bullet}$ est la résolution standard de $A'$ sur $A'$. Comme
$A'$ est l'algèbre polynomiale sur l'ensemble vide sur $A'$, on déduit
d'Objets simpliciaux, lemme
\ref{simplicial-lemma-standard-simplicial-homotopy}, que l'augmentation
$\epsilon_{A'/A'} : P_{A'/A', \bullet} \to A'$ est une équivalence d'homotopie
d'anneaux simpliciaux. Observons que l'application inverse à homotopie près
$c : A' \to P_{A'/A', \bullet}$ construite dans la démonstration de ce lemme
n'est autre que le morphisme structural ; on obtient donc le résultat voulu,
car les deux composées
$$
\xymatrix{
P_\bullet \ar[r] &
P_{A'/A', \bullet} \ar@<1ex>[rr]^{\text{id}}
\ar@<-1ex>[rr]_{c \circ \epsilon_{A'/A'}} & &
P_{A'/A', \bullet} \ar[r] &
P'_\bullet
}
$$
sont les deux applications considérées ci-dessus, et elles sont homotopes
(Objets simpliciaux, remarque
\ref{simplicial-remark-homotopy-pre-post-compose}). Puisque la seconde
application $P_\bullet \to P'_\bullet$ induit l'application nulle
$\Omega_{P_\bullet/A} \to \Omega_{P'_\bullet/A'}$, on conclut que, dans ce
cas, l'application de fonctorialité $L_{B/A} \to L_{B'/A'}$ est homotope à
zéro.
\end{remark}

\begin{lemma}
\label{cotangent-lemma-cotangent-complex-product}
Soient $A \to B$ et $A \to C$ des homomorphismes d'anneaux. Alors
l'application $L_{B \times C/A} \to L_{B/A} \oplus L_{C/A}$ est un
isomorphisme dans $D(B \times C)$.
\end{lemma}

\begin{proof}
Bien que ce lemme puisse se déduire du triangle fondamental, nous allons en
donner maintenant une démonstration directe et élémentaire. Factorisons
l'homomorphisme d'anneaux $A \to B \times C$ sous la forme
$A \to A[x] \to B \times C$, où $x \mapsto (1, 0)$. D'après le lemme
\ref{cotangent-lemma-special-case-triangle}, on a un triangle distingué
$$
L_{A[x]/A} \otimes_{A[x]}^\mathbf{L} (B \times C) \to L_{B \times C/A} \to
L_{B \times C/A[x]} \to L_{A[x]/A} \otimes_{A[x]}^\mathbf{L} (B \times C)[1]
$$
dans $D(B \times C)$. De même, on a les triangles distingués
$$
\begin{matrix}
L_{A[x]/A} \otimes_{A[x]}^\mathbf{L} B \to L_{B/A} \to L_{B/A[x]}
\to L_{A[x]/A} \otimes_{A[x]}^\mathbf{L} B[1] \\
L_{A[x]/A} \otimes_{A[x]}^\mathbf{L} C \to L_{C/A} \to L_{C/A[x]}
\to L_{A[x]/A} \otimes_{A[x]}^\mathbf{L} C[1]
\end{matrix}
$$
Il suffit donc de démontrer le résultat pour $B \times C$ sur $A[x]$.
Notons que $A[x] \to A[x, x^{-1}]$ est plat, que
$(B \times C) \otimes_{A[x]} A[x, x^{-1}] = B \otimes_{A[x]} A[x, x^{-1}]$,
et que $C \otimes_{A[x]} A[x, x^{-1}] = 0$. Par changement de base (lemme
\ref{cotangent-lemma-flat-base-change-cotangent-complex}), l'application
$L_{B \times C/A[x]} \to L_{B/A[x]} \oplus L_{C/A[x]}$ devient un
isomorphisme après inversion de $x$. De la même manière, on montre que
l'application devient un isomorphisme après inversion de $x - 1$. Cela
démontre le lemme.
\end{proof}




\section{Le triangle fondamental}
\label{cotangent-section-triangle}

\noindent
Dans cette section, nous considérons une suite d'homomorphismes d'anneaux
$A \to B \to C$. Notre but est de montrer que ce triangle donne naissance
à un triangle distingué
\begin{equation}
\label{cotangent-equation-triangle}
L_{B/A} \otimes_B^\mathbf{L} C \to L_{C/A} \to L_{C/B} \to
L_{B/A} \otimes_B^\mathbf{L} C[1]
\end{equation}
dans $D(C)$. Cela sera démontré dans la proposition
\ref{cotangent-proposition-triangle}. Pour une autre approche, voir la remarque
\ref{cotangent-remark-triangle}.

\medskip\noindent
Considérons la catégorie $\mathcal{C}_{C/B/A}$, qui est la catégorie
{\bf opposée} de la catégorie dont les objets sont les
$(P \to B, Q \to C)$ tels que
\begin{enumerate}
\item $P$ soit une algèbre polynomiale sur $A$,
\item $P \to B$ soit un homomorphisme de $A$-algèbres,
\item $Q$ soit une algèbre polynomiale sur $P$, et
\item $Q \to C$ soit un homomorphisme de $P$-algèbres.
\end{enumerate}
Nous prenons la catégorie opposée parce que nous voulons voir
$(P \to B, Q \to C)$ comme correspondant au diagramme commutatif
$$
\xymatrix{
\Spec(C) \ar[d] \ar[r] & \Spec(Q) \ar[d] \\
\Spec(B) \ar[d] \ar[r] & \Spec(P) \ar[dl] \\
\Spec(A)
}
$$
Soient $\mathcal{C}_{B/A}$, $\mathcal{C}_{C/A}$ et $\mathcal{C}_{C/B}$
les catégories considérées dans la section
\ref{cotangent-section-compute-L-pi-shriek}. Il existe des foncteurs
$$
\begin{matrix}
u_1 : \mathcal{C}_{C/B/A} \to \mathcal{C}_{B/A}, &
(P \to B, Q \to C) \mapsto (P \to B) \\
u_2 : \mathcal{C}_{C/B/A} \to \mathcal{C}_{C/A}, &
(P \to B, Q \to C) \mapsto (Q \to C) \\
u_3 : \mathcal{C}_{C/B/A} \to \mathcal{C}_{C/B}, &
(P \to B, Q \to C) \mapsto (Q \otimes_P B \to C)
\end{matrix}
$$
Ces foncteurs induisent des morphismes de topos correspondants $g_i$.
Posons $\mathcal{O}_i = g_i^{-1}\mathcal{O}$ ; nous obtenons ainsi des
morphismes de topos annelés
\begin{equation}
\label{cotangent-equation-three-maps}
\begin{matrix}
g_1 : (\Sh(\mathcal{C}_{C/B/A}), \mathcal{O}_1)
\longrightarrow (\Sh(\mathcal{C}_{B/A}), \mathcal{O}) \\
g_2 : (\Sh(\mathcal{C}_{C/B/A}), \mathcal{O}_2)
\longrightarrow (\Sh(\mathcal{C}_{C/A}), \mathcal{O}) \\
g_3 : (\Sh(\mathcal{C}_{C/B/A}), \mathcal{O}_3)
\longrightarrow (\Sh(\mathcal{C}_{C/B}), \mathcal{O})
\end{matrix}
\end{equation}
Notons
$\pi : \Sh(\mathcal{C}_{C/B/A}) \to \Sh(*)$,
$\pi_1 : \Sh(\mathcal{C}_{B/A}) \to \Sh(*)$,
$\pi_2 : \Sh(\mathcal{C}_{C/A}) \to \Sh(*)$, et
$\pi_3 : \Sh(\mathcal{C}_{C/B}) \to \Sh(*)$,
de sorte que $\pi = \pi_i \circ g_i$. Nous obtiendrons notre triangle
distingué à partir de l'identification du complexe cotangent dans le lemme
\ref{cotangent-lemma-compute-cotangent-complex} et des lemmes suivants.

\begin{lemma}
\label{cotangent-lemma-triangle-ses}
Avec les notations de (\ref{cotangent-equation-three-maps}), posons
$$
\begin{matrix}
\Omega_1 = \Omega_{\mathcal{O}/A} \otimes_\mathcal{O} \underline{B}
\text{ sur }\mathcal{C}_{B/A} \\
\Omega_2 = \Omega_{\mathcal{O}/A} \otimes_\mathcal{O} \underline{C}
\text{ sur }\mathcal{C}_{C/A} \\
\Omega_3 = \Omega_{\mathcal{O}/B} \otimes_\mathcal{O} \underline{C}
\text{ sur }\mathcal{C}_{C/B}
\end{matrix}
$$
Alors nous avons une suite exacte courte canonique de faisceaux de
$\underline{C}$-modules
$$
0 \to g_1^{-1}\Omega_1 \otimes_{\underline{B}} \underline{C} \to
g_2^{-1}\Omega_2 \to
g_3^{-1}\Omega_3 \to 0
$$
sur $\mathcal{C}_{C/B/A}$.
\end{lemma}

\begin{proof}
Rappelons que $g_i^{-1}$ s'obtient simplement par précomposition avec $u_i$.
Pour un objet $U = (P \to B, Q \to C)$, nous avons une suite exacte courte
scindée
$$
0 \to \Omega_{P/A} \otimes Q \to \Omega_{Q/A} \to \Omega_{Q/P} \to 0
$$
par exemple d'après Algèbre, lemme
\ref{algebra-lemma-ses-formally-smooth}. En tensorisant par $C$ sur $Q$,
nous obtenons une suite exacte courte
$$
0 \to \Omega_{P/A} \otimes C \to \Omega_{Q/A} \otimes C \to
\Omega_{Q/P} \otimes C \to 0
$$
Nous avons $\Omega_{P/A} \otimes C = \Omega_{P/A} \otimes B \otimes C$ ;
c'est donc la valeur de
$g_1^{-1}\Omega_1 \otimes_{\underline{B}} \underline{C}$
en $U$. Le module $\Omega_{Q/A} \otimes C$ est la valeur de
$g_2^{-1}\Omega_2$ en $U$. Nous avons
$\Omega_{Q/P} \otimes C = \Omega_{Q \otimes_P B/B} \otimes C$
d'après Algèbre, lemme \ref{algebra-lemma-differentials-base-change} ;
c'est donc la valeur de $g_3^{-1}\Omega_3$ en $U$. La suite exacte courte
du lemme s'obtient ainsi en associant à $U$ la dernière suite exacte courte
affichée.
\end{proof}

\begin{lemma}
\label{cotangent-lemma-polynomial-on-top}
Avec les notations de (\ref{cotangent-equation-three-maps}), supposons que $C$ soit
une algèbre polynomiale sur $B$. Alors
$L\pi_!(g_3^{-1}\mathcal{F}) = L\pi_{3, !}\mathcal{F} = \pi_{3, !}\mathcal{F}$
pour tout faisceau abélien $\mathcal{F}$ sur $\mathcal{C}_{C/B}$.
\end{lemma}

\begin{proof}
Écrivons $C = B[E]$ pour un certain ensemble $E$. Choisissons une résolution
$P_\bullet \to B$ de $B$ sur $A$. Pour tout $n$, considérons l'objet
$U_n = (P_n \to B, P_n[E] \to C)$ de $\mathcal{C}_{C/B/A}$. Alors
$U_\bullet$ est un objet cosimplicial de $\mathcal{C}_{C/B/A}$. Notons que
$u_3(U_\bullet)$ est l'objet cosimplicial constant de
$\mathcal{C}_{C/B}$ de valeur $(C \to C)$. Nous allons démontrer que l'objet
$U_\bullet$ de $\mathcal{C}_{C/B/A}$ satisfait aux hypothèses de
Cohomologie sur les sites, lemme
\ref{sites-cohomology-lemma-compute-by-cosimplicial-resolution}.
Cela implique le lemme, car on voit ainsi que $L\pi_!(g_3^{-1}\mathcal{F})$
est calculé par le groupe abélien simplicial constant
$\mathcal{F}(C \to C)$, qui est la valeur de
$L\pi_{3, !}\mathcal{F} = \pi_{3, !}\mathcal{F}$ d'après le lemme
\ref{cotangent-lemma-pi-lower-shriek-polynomial-algebra}.

\medskip\noindent
Soit $U = (\beta : P \to B, \gamma : Q \to C)$ un objet de
$\mathcal{C}_{C/B/A}$. Nous pouvons écrire $P = A[S]$ et
$Q = A[S \amalg T]$, par définition de notre catégorie
$\mathcal{C}_{C/B/A}$. Nous devons montrer que
$$
\Mor_{\mathcal{C}_{C/B/A}}(U_\bullet, U)
$$
est homotopiquement équivalent au singleton simplicial $*$. Observons que
cet ensemble simplicial est le produit
$$
\prod\nolimits_{s \in S} F_s \times \prod\nolimits_{t \in T} F'_t
$$
où $F_s$ est l'ensemble simplicial correspondant à
$U_s = (A[\{s\}] \to B, A[\{s\}] \to C)$
et $F'_t$ est l'ensemble simplicial correspondant à
$U_t = (A \to B, A[\{t\}] \to C)$. En effet, l'objet $U$ est le produit
$\prod U_s \times \prod U_t$ dans $\mathcal{C}_{C/B/A}$. Il suffit que
chaque $F_s$ et chaque $F'_t$ soit homotopiquement équivalent à $*$ ; voir
Simplicial, lemme \ref{simplicial-lemma-products-homotopy}. Le cas de $F_s$
résulte de ce que $P_\bullet \to B$ est une fibration de Kan triviale
(puisqu'il s'agit d'une résolution) et que $F_s$ est la fibre de cette
application au-dessus de $\beta(s)$. (Utiliser Simplicial, lemmes
\ref{simplicial-lemma-trivial-kan-base-change} et
\ref{simplicial-lemma-trivial-kan-homotopy}).
Le cas de $F'_t$ est plus intéressant. Il s'agit ici d'affirmer que la
fibre de
$$
P_\bullet[E] \longrightarrow C = B[E]
$$
au-dessus de $\gamma(t) \in C$ est homotopiquement équivalente à un point.
En fait, nous allons montrer que cette application est une fibration de Kan
triviale. En effet, $P_\bullet \to B$ est une fibration de Kan triviale.
Pour tout anneau $R$, nous avons
$$
R[E] =
\colim_{\Sigma \subset \text{Map}(E, \mathbf{Z}_{\geq 0})\text{ finite}}
\prod\nolimits_{I \in \Sigma} R
$$
(colimite filtrante). Ainsi, l'application affichée d'ensembles simpliciaux
est une colimite filtrante de fibrations de Kan triviales, donc une fibration
de Kan triviale d'après Simplicial, lemme
\ref{simplicial-lemma-filtered-colimit-trivial-kan}.
\end{proof}

\begin{lemma}
\label{cotangent-lemma-triangle-compute-lower-shriek}
Avec les notations de (\ref{cotangent-equation-three-maps}), nous avons
$Lg_{i, !} \circ g_i^{-1} = \text{id}$ for $i = 1, 2, 3$
et par conséquent aussi $L\pi_! \circ g_i^{-1} = L\pi_{i, !}$ pour
$i = 1, 2, 3$.
\end{lemma}

\begin{proof}
Démonstration pour $i = 1$. Nous affirmons que le foncteur
$\mathcal{C}_{C/B/A}$ est une catégorie fibrée sur $\mathcal{C}_{B/A}$.
En effet, supposons donnés $(P \to B, Q \to C)$ et un morphisme
$(P' \to B) \to (P \to B)$ de $\mathcal{C}_{B/A}$. Rappelons que cela
signifie que nous avons un homomorphisme de $A$-algèbres $P \to P'$
compatible aux applications vers $B$. Posons alors $Q' = Q \otimes_P P'$,
muni de l'application induite vers $C$ ; le morphisme
$$
(P' \to B, Q' \to C) \longrightarrow (P \to B, Q \to C)
$$
dans $\mathcal{C}_{C/B/A}$ (noter à nouveau le renversement des flèches)
est fortement cartésien dans $\mathcal{C}_{C/B/A}$ au-dessus de
$\mathcal{C}_{B/A}$. Observons en outre que la catégorie fibre de $u_1$
au-dessus de $P \to B$ est la catégorie $\mathcal{C}_{C/P}$. Soit
$\mathcal{F}$ un faisceau abélien sur $\mathcal{C}_{B/A}$. Puisque nous
avons une catégorie fibrée, nous pouvons appliquer Cohomologie sur les sites,
lemme \ref{sites-cohomology-lemma-compute-left-derived-pi-shriek}. Ainsi,
$L_ng_{1, !}g_1^{-1}\mathcal{F}$ est le (pré)faisceau qui associe à
$U \in \Ob(\mathcal{C}_{B/A})$ la $n$-ième homologie de
$g_1^{-1}\mathcal{F}$ restreint à la catégorie fibre au-dessus de $U$.
Comme ces restrictions sont constantes, le résultat voulu découle du lemme
\ref{cotangent-lemma-pi-lower-shriek-constant-sheaf}, grâce aux identifications des
catégories fibres données ci-dessus.

\medskip\noindent
Le cas $i = 2$. Nous affirmons que $\mathcal{C}_{C/B/A}$ est une catégorie
fibrée sur $\mathcal{C}_{C/A}$. En effet, supposons donnés
$(P \to B, Q \to C)$ et un morphisme $(Q' \to C) \to (Q \to C)$ de
$\mathcal{C}_{C/A}$. Rappelons que cela signifie que nous avons un
homomorphisme de $B$-algèbres $Q \to Q'$ compatible aux applications vers
$C$. Alors
$$
(P \to B, Q' \to C) \longrightarrow (P \to B, Q \to C)
$$
est fortement cartésien dans $\mathcal{C}_{C/B/A}$ au-dessus de
$\mathcal{C}_{C/A}$. Notons que la catégorie fibre de $u_2$ au-dessus de
$Q \to C$ possède un objet final (attention au renversement des flèches),
à savoir $(A \to B, Q \to C)$. Soit $\mathcal{F}$ un faisceau abélien sur
$\mathcal{C}_{C/A}$. Puisque nous avons une catégorie fibrée, nous pouvons
appliquer Cohomologie sur les sites, lemme
\ref{sites-cohomology-lemma-compute-left-derived-pi-shriek}. Ainsi,
$L_ng_{2, !}g_2^{-1}\mathcal{F}$ est le (pré)faisceau qui associe à
$U \in \Ob(\mathcal{C}_{C/A})$ la $n$-ième homologie de
$g_1^{-1}\mathcal{F}$ restreint à la catégorie fibre au-dessus de $U$.
Comme ces restrictions sont constantes, le résultat voulu découle de
Cohomologie sur les sites, lemme
\ref{sites-cohomology-lemma-initial-final}, puisque les catégories fibres
possèdent toutes des objets finals.

\medskip\noindent
Le cas $i = 3$. Dans ce cas, nous appliquerons Cohomologie sur les sites,
lemme
\ref{sites-cohomology-lemma-compute-left-derived-g-shriek}
à $u = u_3 : \mathcal{C}_{C/B/A} \to \mathcal{C}_{C/B}$ et à
$\mathcal{F}' = g_3^{-1}\mathcal{F}$, où $\mathcal{F}$ est un faisceau
abélien sur $\mathcal{C}_{C/B}$. Supposons que
$U = (\overline{Q} \to C)$ soit un objet de $\mathcal{C}_{C/B}$. Alors
$\mathcal{I}_U = \mathcal{C}_{\overline{Q}/B/A}$ (attention encore au
renversement des flèches). Le faisceau $\mathcal{F}'_U$ est donné par la
règle $(P \to B, Q \to \overline{Q}) \mapsto \mathcal{F}(Q \otimes_P B \to C)$.
Autrement dit, ce faisceau est l'image inverse d'un faisceau sur
$\mathcal{C}_{\overline{Q}/C}$ par le morphisme
$\Sh(\mathcal{C}_{\overline{Q}/B/A}) \to \Sh(\mathcal{C}_{\overline{Q}/B})$.
Le lemme \ref{cotangent-lemma-polynomial-on-top} montre donc que
$H_n(\mathcal{I}_U, \mathcal{F}'_U) = 0$ for $n > 0$
et est égal à $\mathcal{F}(\overline{Q} \to C)$ pour $n = 0$. Le lemme de
Cohomologie sur les sites déjà cité,
\ref{sites-cohomology-lemma-compute-left-derived-g-shriek}
implique que $Lg_{3, !}(g_3^{-1}\mathcal{F}) = \mathcal{F}$, ce qui achève
la démonstration.
\end{proof}

\begin{proposition}
\label{cotangent-proposition-triangle}
Soient $A \to B \to C$ des homomorphismes d'anneaux. Il existe un triangle
distingué canonique
$$
L_{B/A} \otimes_B^\mathbf{L} C \to L_{C/A} \to L_{C/B} \to
L_{B/A} \otimes_B^\mathbf{L} C[1]
$$
dans $D(C)$.
\end{proposition}

\begin{proof}
Considérons la suite exacte courte de faisceaux du lemme
\ref{cotangent-lemma-triangle-ses} et appliquons-lui le foncteur dérivé $L\pi_!$ ;
nous obtenons un triangle distingué
$$
L\pi_!(g_1^{-1}\Omega_1 \otimes_{\underline{B}} \underline{C}) \to
L\pi_!(g_2^{-1}\Omega_2) \to
L\pi_!(g_3^{-1}\Omega_3) \to
L\pi_!(g_1^{-1}\Omega_1 \otimes_{\underline{B}} \underline{C})[1]
$$
dans $D(C)$. Les lemmes \ref{cotangent-lemma-triangle-compute-lower-shriek} et
\ref{cotangent-lemma-compute-cotangent-complex}
montrent que les deuxième et troisième termes s'identifient respectivement
à $L_{C/A}$ et $L_{C/B}$, tandis que le premier est égal à
$$
L\pi_{1, !}(\Omega_1 \otimes_{\underline{B}} \underline{C}) =
L\pi_{1, !}(\Omega_1) \otimes_B^\mathbf{L} C =
L_{B/A} \otimes_B^\mathbf{L} C
$$
La première égalité résulte de Cohomologie sur les sites, lemme
\ref{sites-cohomology-lemma-change-of-rings}
(et de la platitude de $\Omega_1$ comme faisceau de modules sur
$\underline{B}$), et la seconde du lemme
\ref{cotangent-lemma-compute-cotangent-complex}.
\end{proof}

\begin{remark}
\label{cotangent-remark-triangle}
Esquissons une autre démonstration, peut-être plus simple, de l'existence du
triangle fondamental. Soient $A \to B \to C$ des homomorphismes d'anneaux et
supposons que $B \to C$ soit injectif. Soit $P_\bullet \to B$ la résolution
standard de $B$ sur $A$, et soit $Q_\bullet \to C$ la résolution standard de
$C$ sur $B$. Représentons-les ainsi :
$$
\xymatrix{
P_\bullet : &
A[A[A[B]]] \ar[d]
\ar@<2ex>[r]
\ar@<0ex>[r]
\ar@<-2ex>[r]
&
A[A[B]] \ar[d]
\ar@<1ex>[r]
\ar@<-1ex>[r]
\ar@<1ex>[l]
\ar@<-1ex>[l]
&
A[B] \ar[d] \ar@<0ex>[l] \ar[r] &
B \\
Q_\bullet : &
A[A[A[C]]]
\ar@<2ex>[r]
\ar@<0ex>[r]
\ar@<-2ex>[r]
&
A[A[C]]
\ar@<1ex>[r]
\ar@<-1ex>[r]
\ar@<1ex>[l]
\ar@<-1ex>[l]
&
A[C] \ar@<0ex>[l] \ar[r] &
C
}
$$
Puisque $B \to C$ est injectif, observons que l'anneau $Q_n$ est une algèbre
polynomiale sur $P_n$ pour tout $n$. Nous obtenons donc un objet cosimplicial
de $\mathcal{C}_{C/B/A}$ (attention au renversement des flèches). Posons
maintenant $\overline{Q}_\bullet = Q_\bullet \otimes_{P_\bullet} B$.
Le point clef de la démonstration de la proposition
\ref{cotangent-proposition-triangle} consiste à montrer que $\overline{Q}_\bullet$ est
une résolution de $C$ sur $B$. Cela résulte de Cohomologie sur les sites,
lemme
\ref{sites-cohomology-lemma-O-homology-qis}
appliqué à $\mathcal{C} = \Delta$, $\mathcal{O} = P_\bullet$,
$\mathcal{O}' = B$ et $\mathcal{F} = Q_\bullet$ (on utilise ici que $Q_n$
est plat sur $P_n$ ; pour relier les modules simpliciaux aux faisceaux, voir
Cohomologie sur les sites, remarque
\ref{sites-cohomology-remark-simplicial-modules}). Ce fait clef implique que
le triangle distingué de la proposition \ref{cotangent-proposition-triangle} est celui
qui est associé à la suite exacte courte de $C$-modules simpliciaux
$$
0 \to
\Omega_{P_\bullet/A} \otimes_{P_\bullet} C \to
\Omega_{Q_\bullet/A} \otimes_{Q_\bullet} C \to
\Omega_{\overline{Q}_\bullet/B} \otimes_{\overline{Q}_\bullet} C \to 0
$$
qui se déduit des suites exactes courtes
$0 \to \Omega_{P_n/A} \otimes_{P_n} Q_n \to \Omega_{Q_n/A} \to
\Omega_{Q_n/P_n} \to 0$ of
Algèbre, lemme \ref{algebra-lemma-ses-formally-smooth}.
En effet, d'après la remarque \ref{cotangent-remark-resolution} et le fait clef, le
complexe du membre de droite représente $L_{C/B}$ dans $D(C)$.

\medskip\noindent
Si $B \to C$ n'est pas injectif, nous pouvons appliquer ce qui précède pour
obtenir un triangle fondamental pour $A \to B \to B \times C$. Puisque
$L_{B \times C/B} \to L_{B/B} \oplus L_{C/B}$ et
$L_{B \times C/A} \to L_{B/A} \oplus L_{C/A}$
constituent des quasi-isomorphismes dans $D(B \times C)$
(lemme \ref{cotangent-lemma-cotangent-complex-product})
cela induit le triangle distingué voulu dans $D(C)$ en tensorisant par
l'homomorphisme plat d'anneaux $B \times C \to C$.
\end{remark}

\begin{remark}
\label{cotangent-remark-explicit-map}
Soient $A \to B \to C$ des homomorphismes d'anneaux, avec $B \to C$
injectif. Rappelons les notations $P_\bullet$, $Q_\bullet$ et
$\overline{Q}_\bullet$ de la remarque \ref{cotangent-remark-triangle}. Soit
$R_\bullet$ la résolution standard de $C$ sur $B$. Dans cette remarque,
nous expliquons comment obtenir l'identification canonique de
$\Omega_{\overline{Q}_\bullet/B} \otimes_{\overline{Q}_\bullet} C$ à
$L_{C/B} = \Omega_{R_\bullet/B} \otimes_{R_\bullet} C$. Soit
$S_\bullet \to B$ la résolution standard de $B$ sur $B$. Notons que
l'application de fonctorialité $S_\bullet \to R_\bullet$ fait de $R_n$ une
algèbre polynomiale sur $S_n$, puisque $B \to C$ est injectif.
Par exemple, en degré $0$, nous avons l'application $B[B] \to B[C]$ ; en
degré $1$, l'application $B[B[B]] \to B[B[C]]$ ; et ainsi de suite. Par
conséquent,
$\overline{R}_\bullet = R_\bullet \otimes_{S_\bullet} B$
est elle aussi une algèbre polynomiale simpliciale sur $B$, et il résulte
(comme dans la remarque \ref{cotangent-remark-triangle}) de Cohomologie sur les sites,
lemme
\ref{sites-cohomology-lemma-O-homology-qis}
que $\overline{R}_\bullet \to C$ est une résolution. Comme nous avons un
diagramme commutatif
$$
\xymatrix{
Q_\bullet \ar[r] & R_\bullet \\
P_\bullet \ar[u] \ar[r] & S_\bullet \ar[u] \ar[r] & B
}
$$
nous obtenons une application canonique
$\overline{Q}_\bullet = Q_\bullet \otimes_{P_\bullet} B \to
\overline{R}_\bullet$. Les applications
$$
L_{C/B} = \Omega_{R_\bullet/B} \otimes_{R_\bullet} C
\longrightarrow
\Omega_{\overline{R}_\bullet/B} \otimes_{\overline{R}_\bullet} C
\longleftarrow
\Omega_{\overline{Q}_\bullet/B} \otimes_{\overline{Q}_\bullet} C
$$
sont alors des quasi-isomorphismes (remarque \ref{cotangent-remark-resolution}), et
composer l'une avec l'inverse de l'autre donne l'identification voulue.
\end{remark}





\section{Localisation et homomorphismes d'anneaux \'etales}
\label{cotangent-section-localization}

\noindent
Dans cette section, nous étudions ce qui se passe lorsque nous localisons nos
anneaux. Soient $A \to A' \to B$ des homomorphismes d'anneaux tels que
$B = B \otimes_A^\mathbf{L} A'$. C'est le cas, par exemple, si
$A' = S^{-1}A$ est le localisé de $A$ par une partie multiplicative
$S \subset A$. Dans ce cas, pour un faisceau abélien $\mathcal{F}'$ sur
$\mathcal{C}_{B/A'}$, l'homologie de $g^{-1}\mathcal{F}'$ sur
$\mathcal{C}_{B/A}$ coïncide avec celle de $\mathcal{F}'$ sur
$\mathcal{C}_{B/A'}$ ; voir le lemme \ref{cotangent-lemma-flat-base-change} pour un
énoncé précis.

\begin{lemma}
\label{cotangent-lemma-localize-at-bottom}
Soient $A \to A' \to B$ des homomorphismes d'anneaux tels que
$B = B \otimes_A^\mathbf{L} A'$. Alors $L_{B/A} = L_{B/A'}$ dans $D(B)$.
\end{lemma}

\begin{proof}
D'après la discussion précédente (c'est-à-dire en utilisant le lemme
\ref{cotangent-lemma-flat-base-change}) et le lemme
\ref{cotangent-lemma-compute-cotangent-complex}, il nous faut montrer que le faisceau
défini sur $\mathcal{C}_{B/A}$ par la règle
$(P \to B) \mapsto \Omega_{P/A} \otimes_P B$ est l'image inverse du faisceau
défini par la règle $(P \to B) \mapsto \Omega_{P/A'} \otimes_P B$.
Le foncteur image inverse $g^{-1}$ s'obtient en précomposant avec le foncteur
$u : \mathcal{C}_{B/A} \to \mathcal{C}_{B/A'}$,
$(P \to B) \mapsto (P \otimes_A A' \to B)$. Il nous faut donc montrer que
$$
\Omega_{P/A} \otimes_P B =
\Omega_{P \otimes_A A'/A'} \otimes_{(P \otimes_A A')} B
$$
D'après Algèbre, lemme \ref{algebra-lemma-differentials-base-change}, le
membre de droite est égal à
$$
(\Omega_{P/A} \otimes_A A') \otimes_{(P \otimes_A A')} B
$$
Comme $P$ est une algèbre polynomiale sur $A$, le module $\Omega_{P/A}$ est
libre et l'égalité est évidente.
\end{proof}

\begin{lemma}
\label{cotangent-lemma-derived-diagonal}
Soit $A \to B$ un homomorphisme d'anneaux tel que
$B = B \otimes_A^\mathbf{L} B$. Alors $L_{B/A} = 0$ dans $D(B)$.
\end{lemma}

\begin{proof}
Cela résulte de $L_{B/A} = L_{B/B} = 0$, d'après les lemmes
\ref{cotangent-lemma-localize-at-bottom} et
\ref{cotangent-lemma-cotangent-complex-polynomial-algebra}.
\end{proof}

\begin{lemma}
\label{cotangent-lemma-bootstrap}
Soit $A \to B$ un homomorphisme d'anneaux tel que
$\text{Tor}^A_i(B, B) = 0$ pour $i > 0$ et tel que
$L_{B/B \otimes_A B} = 0$. Alors $L_{B/A} = 0$ dans $D(B)$.
\end{lemma}

\begin{proof}
D'après le lemme \ref{cotangent-lemma-flat-base-change-cotangent-complex}, on a
$L_{B/A} \otimes_B^\mathbf{L} (B \otimes_A B) = L_{B \otimes_A B/B}$.
Utilisons maintenant le triangle distingué (\ref{cotangent-equation-triangle})
$$
L_{B \otimes_A B/B} \otimes^\mathbf{L}_{(B \otimes_A B)} B \to
L_{B/B} \to L_{B/B \otimes_A B} \to
L_{B \otimes_A B/B} \otimes^\mathbf{L}_{(B \otimes_A B)} B[1]
$$
associé aux homomorphismes d'anneaux $B \to B \otimes_A B \to B$, ainsi que
l'annulation de $L_{B/B}$ (lemme
\ref{cotangent-lemma-cotangent-complex-polynomial-algebra}) et celle, supposée, de
$L_{B/B \otimes_A B}$, pour obtenir
$$
0 =
L_{B \otimes_A B/B} \otimes^\mathbf{L}_{(B \otimes_A B)} B =
L_{B/A} \otimes_B^\mathbf{L} (B \otimes_A B)
\otimes^\mathbf{L}_{(B \otimes_A B)} B = L_{B/A}
$$
comme voulu.
\end{proof}

\begin{lemma}
\label{cotangent-lemma-when-zero}
Le complexe cotangent $L_{B/A}$ est nul dans chacun des cas suivants :
\begin{enumerate}
\item $A \to B$ et $B \otimes_A B \to B$ sont plats, c'est-à-dire que
$A \to B$ est faiblement \'etale (Compléments sur l'algèbre, définition
\ref{more-algebra-definition-weakly-etale}),
\item $A \to B$ est un épimorphisme plat d'anneaux,
\item $B = S^{-1}A$ pour une partie multiplicative $S \subset A$,
\item $A \to B$ est non ramifié et plat,
\item $A \to B$ est \'etale,
\item $A \to B$ est une colimite filtrante d'homomorphismes d'anneaux dont le
complexe cotangent s'annule,
\item $B$ est un hensélisé d'un anneau local de $A$,
\item $B$ est un hensélisé strict d'un anneau local de $A$, et
\item ajouter ici d'autres cas.
\end{enumerate}
\end{lemma}

\begin{proof}
Dans le cas (1), nous pouvons appliquer le lemme
\ref{cotangent-lemma-derived-diagonal} à l'homomorphisme d'anneaux plat surjectif
$B \otimes_A B \to B$ pour conclure que $L_{B/B \otimes_A B} = 0$, puis le
lemme \ref{cotangent-lemma-bootstrap} permet de conclure. Les cas (2) -- (5) sont tous
des cas particuliers de (1). L'assertion (6) résulte du lemme
\ref{cotangent-lemma-colimit-cotangent-complex}. Les assertions (7) et (8) résultent du
fait que les hensélisés (stricts) sont des colimites filtrantes d'extensions
d'anneaux \'etales de $A$ ; voir Algèbre, lemmes
\ref{algebra-lemma-henselization-different} et
\ref{algebra-lemma-strict-henselization-different}.
\end{proof}

\begin{lemma}
\label{cotangent-lemma-localize-on-top}
Soient $A \to B \to C$ des homomorphismes d'anneaux tels que $L_{C/B} = 0$.
Alors $L_{C/A} = L_{B/A} \otimes_B^\mathbf{L} C$.
\end{lemma}

\begin{proof}
C'est une conséquence immédiate du triangle distingué
(\ref{cotangent-equation-triangle}).
\end{proof}

\begin{lemma}
\label{cotangent-lemma-localize}
Soit $A \to B$ un homomorphisme d'anneaux, et soient $S \subset A$ et
$T \subset B$ des parties multiplicatives telles que $S$ s'envoie dans $T$.
Alors $L_{T^{-1}B/S^{-1}A} = L_{B/A} \otimes_B T^{-1}B$ dans $D(T^{-1}B)$.
\end{lemma}

\begin{proof}
D'après le lemme \ref{cotangent-lemma-localize-on-top},
$L_{T^{-1}B/A} = L_{B/A} \otimes_B T^{-1}B$
et, d'après le lemme \ref{cotangent-lemma-localize-at-bottom},
$L_{T^{-1}B/A} = L_{T^{-1}B/S^{-1}A}$.
\end{proof}

\begin{lemma}
\label{cotangent-lemma-cotangent-complex-henselization}
Soit $A \to B$ un homomorphisme local d'anneaux locaux. Notons
$A^h \to B^h$, resp.\ $A^{sh} \to B^{sh}$, les homomorphismes induits entre
hensélisés, resp.\ hensélisés stricts. Alors
$$
L_{B^h/A^h} = L_{B^h/A} = L_{B/A} \otimes_B^\mathbf{L} B^h
\quad\text{resp.}\quad
L_{B^{sh}/A^{sh}} = L_{B^{sh}/A} = L_{B/A} \otimes_B^\mathbf{L} B^{sh}
$$
dans $D(B^h)$, resp.\ $D(B^{sh})$.
\end{lemma}

\begin{proof}
Les complexes $L_{A^h/A}$, $L_{A^{sh}/A}$, $L_{B^h/B}$ et
$L_{B^{sh}/B}$ sont tous nuls d'après le lemme \ref{cotangent-lemma-when-zero}. En
appliquant le triangle distingué fondamental (\ref{cotangent-equation-triangle}) à
$A \to B \to B^h$, nous obtenons
$L_{B^h/A} = L_{B/A} \otimes_B^\mathbf{L} B^h$.
Le triangle fondamental appliqué à $A \to A^h \to B^h$ donne
$L_{B^h/A^h} = L_{B^h/A}$. Le raisonnement est le même pour les hensélisés
stricts.
\end{proof}




\section{Homomorphismes d'anneaux lisses}
\label{cotangent-section-smooth}

\noindent
Soit $C \to B$ une surjection d'anneaux de noyau $I$. Disons qu'un tel
homomorphisme d'anneaux est « faiblement quasi régulier » si $I/I^2$ est un
$B$-module plat et si $\text{Tor}_*^C(B, B)$ est l'algèbre extérieure de
$I/I^2$. Pour généraliser aux « homomorphismes d'anneaux lisses » ce qui est
fait dans le lemme \ref{cotangent-lemma-when-zero} pour les « homomorphismes d'anneaux
\'etales », il convient de considérer les homomorphismes plats $A \to B$ tels
que l'homomorphisme de multiplication $B \otimes_A B \to B$ soit faiblement
quasi régulier. Pour le moment, nous nous en tenons aux homomorphismes
d'anneaux lisses.

\begin{lemma}
\label{cotangent-lemma-when-projective}
Si $A \to B$ est un homomorphisme d'anneaux lisse, alors
$L_{B/A} = \Omega_{B/A}[0]$.
\end{lemma}

\begin{proof}
L'identification en degré cohomologique $0$ est donnée par le lemme
\ref{cotangent-lemma-identify-H0}. Il suffit donc de montrer que les autres groupes de
cohomologie sont nuls. Il suffit de le faire localement sur $\Spec(B)$, puisque
$L_{B_g/A} = (L_{B/A})_g$ pour $g \in B$ d'après le lemme
\ref{cotangent-lemma-localize-on-top}. Nous pouvons ainsi supposer que $A \to B$ est
lisse standard (Algèbre, lemme \ref{algebra-lemma-smooth-syntomic}),
c'est-à-dire que $A \to B$ se factorise sous la forme
$A \to A[x_1, \ldots, x_n] \to B$, où
$A[x_1, \ldots, x_n] \to B$ est \'etale. Dans ce cas, les lemmes
\ref{cotangent-lemma-when-zero} et \ref{cotangent-lemma-localize-on-top} montrent que
$L_{B/A} = L_{A[x_1, \ldots, x_n]/A} \otimes B$
d'où la conclusion par le lemme
\ref{cotangent-lemma-cotangent-complex-polynomial-algebra}.
\end{proof}





\section{Caractéristique positive}
\label{cotangent-section-positive-characteristic}

\noindent
Dans cette section, nous fixons un nombre premier $p$. Si $A$ est un anneau
tel que $p = 0$ dans $A$, alors $F_A : A \to A$ désigne l'endomorphisme de
Frobenius $a \mapsto a^p$.

\begin{lemma}
\label{cotangent-lemma-frobenius-homotopy}
Soit $A \to B$ un homomorphisme d'anneaux tel que $p = 0$ dans $A$. Soit
$P_\bullet$ la résolution standard de $B$ sur $A$. L'application
$P_\bullet \to P_\bullet$ induite par le diagramme
$$
\xymatrix{
B \ar[r]_{F_B} & B \\
A \ar[u] \ar[r]^{F_A} & A \ar[u]
}
$$
considéré dans la section \ref{cotangent-section-functoriality} est homotope à
l'endomorphisme de Frobenius $P_\bullet \to P_\bullet$ donné par Frobenius sur
chaque $P_n$.
\end{lemma}

\begin{proof}
Soit $\mathcal{A}$ la catégorie des homomorphismes de $\mathbf{F}_p$-algèbres
$A \to B$. Soit $\mathcal{S}$ la catégorie des couples $(A, E)$, où $A$ est
une $\mathbf{F}_p$-algèbre et $E$ un ensemble. Considérons les foncteurs
adjoints
$$
V : \mathcal{A} \to \mathcal{S}, \quad (A \to B) \mapsto (A, B)
$$
et
$$
U : \mathcal{S} \to \mathcal{A}, \quad (A, E) \mapsto (A \to A[E])
$$
Soit $X$ l'objet simplicial de la catégorie des foncteurs de $\mathcal{A}$
dans $\mathcal{A}$ construit dans Simplicial, section
\ref{simplicial-section-standard}. Il est clair que
$P_\bullet = X(A \to B)$ car, si l'on fixe $A$, alors.

\medskip\noindent
Posons $Y = U \circ V$. Rappelons que $X$ est construit à partir de $Y$ et de
certaines applications, et que ses termes sont
$X_n = Y \circ \ldots \circ Y$, avec $n + 1$ facteurs ; la construction est
donnée dans Simplicial, exemple \ref{simplicial-example-godement}, et l'on
trouvera les détails dans la démonstration de Simplicial, lemme
\ref{simplicial-lemma-standard-simplicial}.

\medskip\noindent
Soit $f : \text{id}_\mathcal{A} \to \text{id}_\mathcal{A}$ l'endomorphisme de
Frobenius du foncteur identité. Autrement dit, posons
$f_{A \to B} = (F_A, F_B) : (A \to B) \to (A \to B)$. Nos deux applications
sur $X(A \to B)$ sont alors données par les transformations naturelles
$f \star 1_X$ et $1_X \star f$. Nous omettons les détails. On conclut donc par
Simplicial, lemme \ref{simplicial-lemma-godement-before-after}.
\end{proof}

\begin{lemma}
\label{cotangent-lemma-frobenius-acts-as-zero}
Soit $p$ un nombre premier. Soit $A \to B$ un homomorphisme d'anneaux et
supposons que $p = 0$ dans $A$. L'application $L_{B/A} \to L_{B/A}$ de la
section \ref{cotangent-section-functoriality} induite par les applications de Frobenius
$F_A$ et $F_B$ est homotope à zéro.
\end{lemma}

\begin{proof}
Soit $P_\bullet$ la résolution standard de $B$ sur $A$. D'après le lemme
\ref{cotangent-lemma-frobenius-homotopy}, l'application $P_\bullet \to P_\bullet$ induite
par $F_A$ et $F_B$ est homotope à l'application
$F_{P_\bullet} : P_\bullet \to P_\bullet$ donnée par Frobenius sur chaque
terme. On obtient donc le résultat voulu, puisque $F_{P_\bullet}$ induit
manifestement l'application nulle $\Omega_{P_n/A} \to \Omega_{P_n/A}$ (la
dérivée d'une puissance $p$-ième étant nulle).
\end{proof}

\begin{lemma}
\label{cotangent-lemma-perfect-zero}
Soit $p$ un nombre premier. Soit $A \to B$ un homomorphisme d'anneaux et
supposons que $p = 0$ dans $A$. Si $A$ et $B$ sont parfaits, alors $L_{B/A}$
est nul dans $D(B)$.
\end{lemma}

\begin{proof}
L'application $(F_A, F_B) : (A \to B) \to (A \to B)$ est un isomorphisme et
induit donc un isomorphisme sur $L_{B/A}$ ; d'autre part, elle induit zéro sur
$L_{B/A}$ d'après le lemme \ref{cotangent-lemma-frobenius-acts-as-zero}.
\end{proof}





\section{Comparaison avec le complexe cotangent naïf}
\label{cotangent-section-surjections}

\noindent
Le complexe cotangent naïf a été introduit dans Algèbre, section
\ref{algebra-section-netherlander}.

\begin{remark}
\label{cotangent-remark-make-map}
Soit $A \to B$ un homomorphisme d'anneaux. Travaillons sur
$\mathcal{C}_{B/A}$ comme dans la section
\ref{cotangent-section-compute-L-pi-shriek}, et notons
$\mathcal{J} \subset \mathcal{O}$ le noyau de
$\mathcal{O} \to \underline{B}$. Remarquons que
$L\pi_!(\mathcal{J}) = 0$ d'après le lemme
\ref{cotangent-lemma-apply-O-B-comparison}. Posons
$\Omega =  \Omega_{\mathcal{O}/A} \otimes_\mathcal{O} \underline{B}$
de sorte que
$L_{B/A} = L\pi_!(\Omega)$ d'après le lemme
\ref{cotangent-lemma-compute-cotangent-complex}. Il s'ensuit que
$L\pi_!(\mathcal{J} \to \Omega) = L\pi_!(\Omega) = L_{B/A}$. Ainsi, pour tout
objet $U = (P \to B)$ de $\mathcal{C}_{B/A}$, nous obtenons une application
\begin{equation}
\label{cotangent-equation-comparison-map-A}
(J \to \Omega_{P/A} \otimes_P B) \longrightarrow L_{B/A}
\end{equation}
où $J = \Ker(P \to B)$ dans $D(A)$ ; voir Cohomologie sur les sites, remarque
\ref{sites-cohomology-remark-map-evaluation-to-derived}. Poursuivant de cette
façon, remarquons que
$L\pi_!(\mathcal{J} \otimes_\mathcal{O}^\mathbf{L} \underline{B}) =
L\pi_!(\mathcal{J}) = 0$ by
lemme \ref{cotangent-lemma-O-homology-B-homology}.
Comme $\text{Tor}_0^\mathcal{O}(\mathcal{J}, \underline{B}) =
\mathcal{J}/\mathcal{J}^2$, la suite spectrale
$$
H_p(\mathcal{C}_{B/A}, \text{Tor}_q^\mathcal{O}(\mathcal{J}, \underline{B}))
\Rightarrow 
H_{p + q}(\mathcal{C}_{B/A},
\mathcal{J} \otimes_\mathcal{O}^\mathbf{L} \underline{B}) = 0
$$
(duale de Catégories dérivées, lemme
\ref{derived-lemma-two-ss-complex-functor}) implique que
$H_0(\mathcal{C}_{B/A}, \mathcal{J}/\mathcal{J}^2) = 0$
et $H_1(\mathcal{C}_{B/A}, \mathcal{J}/\mathcal{J}^2) = 0$.
Il s'ensuit que le complexe de $\underline{B}$-modules
$\mathcal{J}/\mathcal{J}^2 \to \Omega$ vérifie
$\tau_{\geq -1}L\pi_!(\mathcal{J}/\mathcal{J}^2 \to \Omega) =
\tau_{\geq -1}L_{B/A}$.
Ainsi, pour tout objet $U = (P \to B)$ de $\mathcal{C}_{B/A}$, nous obtenons
une application
\begin{equation}
\label{cotangent-equation-comparison-map}
(J/J^2 \to \Omega_{P/A} \otimes_P B) \longrightarrow \tau_{\geq -1}L_{B/A}
\end{equation}
dans $D(B)$ ; voir Cohomologie sur les sites, remarque
\ref{sites-cohomology-remark-map-evaluation-to-derived}.
\end{remark}

\noindent
Le premier cas est celui d'une surjection d'anneaux.

\begin{lemma}
\label{cotangent-lemma-surjection}
\begin{slogan}
La cohomologie du complexe cotangent d'un homomorphisme d'anneaux surjectif est
nulle en degré zéro ; en degré $-1$, elle est le noyau modulo son carré.
\end{slogan}
Soit $A \to B$ un homomorphisme d'anneaux surjectif de noyau $I$. Alors
$H^0(L_{B/A}) = 0$ et $H^{-1}(L_{B/A}) = I/I^2$. Cet isomorphisme provient de
l'application (\ref{cotangent-equation-comparison-map}) pour l'objet $(A \to B)$ de
$\mathcal{C}_{B/A}$.
\end{lemma}

\begin{proof}
Nous allons montrer ci-dessous (en utilisant la surjectivité de $A \to B$)
qu'il existe une suite exacte courte
$$
0 \to \pi^{-1}(I/I^2) \to \mathcal{J}/\mathcal{J}^2 \to \Omega \to 0
$$
de faisceaux sur $\mathcal{C}_{B/A}$. En appliquant $L\pi_!$, puis la suite
exacte longue d'homologie associée, et en utilisant l'annulation de
$H_1(\mathcal{C}_{B/A}, \mathcal{J}/\mathcal{J}^2)$ et de
$H_0(\mathcal{C}_{B/A}, \mathcal{J}/\mathcal{J}^2)$ établie dans la remarque
\ref{cotangent-remark-make-map}, nous obtenons le résultat voulu grâce au lemme
\ref{cotangent-lemma-pi-lower-shriek-constant-sheaf}.

\medskip\noindent
Il reste à vérifier l'énoncé local mentionné ci-dessus. Pour tout objet
$U = (P \to B)$ de $\mathcal{C}_{B/A}$, nous pouvons choisir un isomorphisme
$P = A[E]$ tel que l'homomorphisme $P \to B$ envoie chaque $e \in E$ sur zéro.
Alors
$J = \mathcal{J}(U) \subset P = \mathcal{O}(U)$
est égal à $J = IP + (e; e \in E)$. La valeur en $U$ de la suite exacte courte
de faisceaux ci-dessus est la suite
$$
0 \to I/I^2 \to J/J^2 \to \Omega_{P/A} \otimes_P B \to 0
$$
Nous omettons la vérification (indication : le seul point délicat est l'égalité
$IP \cap J^2 = IJ$, qui résulte par exemple de Compléments sur l'algèbre,
lemme \ref{more-algebra-lemma-conormal-sequence-H1-regular}).
\end{proof}

\begin{lemma}
\label{cotangent-lemma-relation-with-naive-cotangent-complex}
Soit $A \to B$ un homomorphisme d'anneaux. Alors
$\tau_{\geq -1}L_{B/A}$ est canoniquement quasi-isomorphe au complexe
cotangent naïf.
\end{lemma}

\begin{proof}
Considérons $P = A[B] \to B$, de noyau $I$. Le complexe cotangent naïf
$\NL_{B/A}$ de $B$ sur $A$ est le complexe
$I/I^2 \to \Omega_{P/A} \otimes_P B$ ; voir Algèbre, définition
\ref{algebra-definition-naive-cotangent-complex}. Remarquons que nous avons
déjà construit dans (\ref{cotangent-equation-comparison-map}) une application canonique
$$
c : \NL_{B/A} \longrightarrow \tau_{\geq -1}L_{B/A}
$$
Considérons le triangle distingué (\ref{cotangent-equation-triangle})
$$
L_{P/A} \otimes_P^\mathbf{L} B \to L_{B/A} \to L_{B/P} \to 
(L_{P/A} \otimes_P^\mathbf{L} B)[1]
$$
associé aux homomorphismes d'anneaux $A \to A[B] \to B$. Nous savons que
$L_{P/A} = \Omega_{P/A}[0] = \NL_{P/A}$ in $D(P)$
(lemme \ref{cotangent-lemma-cotangent-complex-polynomial-algebra} et Algèbre, lemme
\ref{algebra-lemma-NL-polynomial-algebra}) et que
$\tau_{\geq -1}L_{B/P} = I/I^2[1] = \NL_{B/P}$ in $D(B)$
(lemme \ref{cotangent-lemma-surjection} et Algèbre, lemme
\ref{algebra-lemma-NL-surjection}). Pour montrer que $c$ est un
quasi-isomorphisme, il suffit, d'après Algèbre, lemme
\ref{algebra-lemma-exact-sequence-NL}, et la suite exacte longue de cohomologie
associée au triangle distingué, de montrer que les applications
$L_{P/A} \to L_{B/A} \to L_{B/P}$ sont compatibles sur les groupes de
cohomologie avec les applications correspondantes
$\NL_{P/A} \to \NL_{B/A} \to \NL_{B/P}$
du complexe cotangent naïf. Nous omettons la vérification.
\end{proof}

\begin{remark}
\label{cotangent-remark-explicit-comparison-map}
Nous pouvons expliciter comme suit l'application de comparaison du lemme
\ref{cotangent-lemma-relation-with-naive-cotangent-complex}. Soit $P_\bullet$ la
résolution standard de $B$ sur $A$. Posons $I = \Ker(A[B] \to B)$. Rappelons
que $P_0 = A[B]$. L'application du lemme est donnée par le diagramme
commutatif
$$
\xymatrix{
L_{B/A} \ar[d] & \ldots \ar[r] &
\Omega_{P_2/A} \otimes_{P_2} B
\ar[r] \ar[d] &
\Omega_{P_1/A} \otimes_{P_1} B
\ar[r] \ar[d] &
\Omega_{P_0/A} \otimes_{P_0} B
\ar[d] \\
\NL_{B/A} & \ldots \ar[r] &
0 \ar[r] & 
I/I^2 \ar[r] &
\Omega_{P_0/A} \otimes_{P_0} B
}
$$
Nous construisons la flèche descendante de but $I/I^2$ en envoyant
$\text{d}f \otimes b$ sur la classe de $(d_0(f) - d_1(f))b$ dans $I/I^2$.
Ici, $d_i : P_1 \to P_0$, $i = 0, 1$, sont les deux applications de face de
la structure simpliciale. Cela a un sens puisque $d_0 - d_1$ envoie $P_1$
dans $I = \Ker(P_0 \to B)$. Nous omettons de vérifier que cette règle est bien
définie. Notre application est compatible avec la différentielle
$\Omega_{P_1/A} \otimes_{P_1} B \to \Omega_{P_0/A} \otimes_{P_0} B$
car celle-ci envoie $\text{d}f \otimes b$ sur
$\text{d}(d_0(f) - d_1(f)) \otimes b$. En outre, la différentielle
$\Omega_{P_2/A} \otimes_{P_2} B \to \Omega_{P_1/A} \otimes_{P_1} B$
envoie $\text{d}f \otimes b$ sur
$\text{d}(d_0(f) - d_1(f) + d_2(f)) \otimes b$, qui est annulé par notre
flèche descendante. Nous obtenons ainsi une application de complexes. Nous
omettons de vérifier qu'elle coïncide avec l'application du lemme
\ref{cotangent-lemma-relation-with-naive-cotangent-complex}.
\end{remark}

\begin{remark}
\label{cotangent-remark-surjection}
Reprenons les notations de la remarque \ref{cotangent-remark-make-map}. Les arguments
qui y sont donnés montrent que la différentielle
$$
H_2(\mathcal{C}_{B/A}, \mathcal{J}/\mathcal{J}^2)
\longrightarrow
H_0(\mathcal{C}_{B/A}, \text{Tor}_1^\mathcal{O}(\mathcal{J}, \underline{B}))
$$
de la suite spectrale est un isomorphisme. Notons $\mathcal{C}'_{B/A}$ la
sous-catégorie pleine de $\mathcal{C}_{B/A}$ formée des homomorphismes
surjectifs $P \to B$. L'identification du complexe cotangent avec le complexe
cotangent naïf (lemme \ref{cotangent-lemma-relation-with-naive-cotangent-complex}) montre
que nous avons une suite exacte de faisceaux
$$
0 \to \underline{H_1(L_{B/A})} \to
\mathcal{J}/\mathcal{J}^2 \xrightarrow{\text{d}} \Omega \to
\underline{H_2(L_{B/A})} \to 0
$$
sur $\mathcal{C}'_{B/A}$. Il s'ensuit que $\Ker(d)$ et $\Coker(d)$ sur la
catégorie tout entière $\mathcal{C}_{B/A}$ ont leurs groupes d'homologie
supérieure nuls, puisque ceux-ci sont calculés par les groupes d'homologie de
groupes abéliens simpliciaux constants d'après le lemme
\ref{cotangent-lemma-identify-pi-shriek}. Nous en concluons que
$$
H_n(\mathcal{C}_{B/A}, \mathcal{J}/\mathcal{J}^2) \to H_n(L_{B/A})
$$
est un isomorphisme pour tout $n \geq 2$. Avec la remarque précédente, ceci
donne la formule
$H_2(L_{B/A}) =
H_0(\mathcal{C}_{B/A}, \text{Tor}_1^\mathcal{O}(\mathcal{J}, \underline{B}))$.
\end{remark}





\section{Une suite spectrale de Quillen}
\label{cotangent-section-spectral-sequence}

\noindent
Dans cette section, nous étudions une suite spectrale qui relie le produit
tensoriel dérivé au complexe cotangent.

\begin{lemma}
\label{cotangent-lemma-vanishing-symmetric-powers}
Adoptons les notations et les hypothèses de Cohomologie sur les sites, exemple
\ref{sites-cohomology-example-category-to-point}. Supposons que $\mathcal{C}$
possède un objet cosimplicial comme dans Cohomologie sur les sites, lemme
\ref{sites-cohomology-lemma-compute-by-cosimplicial-resolution}. Soit
$\mathcal{F}$ un $\underline{B}$-module plat tel que
$H_0(\mathcal{C}, \mathcal{F}) = 0$. Alors
$H_l(\mathcal{C}, \text{Sym}_{\underline{B}}^k(\mathcal{F})) = 0$ pour
$l < k$.
\end{lemma}

\begin{proof}
Nous omettons l'indice ${}_{\underline{B}}$ dans les produits tensoriels, les
puissances extérieures et les puissances symétriques. Démontrons le lemme par
récurrence sur $k$. Les cas $k = 0, 1$ résultent des hypothèses. Si $k > 1$,
considérons le complexe exact
$$
\ldots \to
\wedge^2\mathcal{F} \otimes \text{Sym}^{k - 2}\mathcal{F} \to
\mathcal{F} \otimes \text{Sym}^{k - 1}\mathcal{F} \to
\text{Sym}^k\mathcal{F} \to 0
$$
dont les différentielles sont celles du complexe de Koszul. Considéré comme
une résolution de $\text{Sym}^k\mathcal{F}$, ce complexe donne une suite
spectrale du premier quadrant
$$
E_1^{p, q} =
H_p(\mathcal{C},
\wedge^{q + 1}\mathcal{F} \otimes \text{Sym}^{k - q - 1}\mathcal{F})
\Rightarrow
H_{p + q}(\mathcal{C}, \text{Sym}^k(\mathcal{F}))
$$
D'après Cohomologie sur les sites, lemme
\ref{sites-cohomology-lemma-eilenberg-zilber}, nous avons
$$
L\pi_!(\wedge^{q + 1}\mathcal{F} \otimes \text{Sym}^{k - q - 1}\mathcal{F}) =
L\pi_!(\wedge^{q + 1}\mathcal{F}) \otimes_B^\mathbf{L}
L\pi_!(\text{Sym}^{k - q - 1}\mathcal{F}))
$$
Il résulte de la construction des produits tensoriels dérivés que l'hypothèse
de récurrence, jointe à l'annulation de
$H_0(\mathcal{C}, \wedge^{q + 1}(\mathcal{F})) = 0$, donnera le résultat voulu.
Cette annulation est vraie parce que $\wedge^{q + 1}(\mathcal{F})$ est un
quotient de $\mathcal{F}^{\otimes q + 1}$ et que
$H_0(\mathcal{C}, \mathcal{F}^{\otimes q + 1})$ est un quotient de
$H_0(\mathcal{C}, \mathcal{F})^{\otimes q + 1}$, qui est nul.
\end{proof}

\begin{remark}
\label{cotangent-remark-first-homology-symmetric-power}
Dans la situation du lemme \ref{cotangent-lemma-vanishing-symmetric-powers}, on peut
montrer que
$H_k(\mathcal{C}, \text{Sym}^k(\mathcal{F})) =
\wedge^k_B(H_1(\mathcal{C}, \mathcal{F}))$.
En effet, on peut déduire de la démonstration que
$H_k(\mathcal{C}, \text{Sym}^k(\mathcal{F}))$ est le module des coinvariants
sous $S_k$ de
$$
H^{-k}(L\pi_!(\mathcal{F}) \otimes_B^\mathbf{L}
L\pi_!(\mathcal{F}) \otimes_B^\mathbf{L}
\ldots \otimes_B^\mathbf{L} L\pi_!(\mathcal{F})) =
H_1(\mathcal{C}, \mathcal{F})^{\otimes k}
$$
Notre assertion est donc que cette action est l'action usuelle de $S_k$ sur le
produit tensoriel, multipliée par le caractère signature. Pour le démontrer,
il faut examiner les conventions de signes dans la définition du complexe
total associé à un multicomplexe. Nous omettons la vérification.
\end{remark}

\begin{lemma}
\label{cotangent-lemma-map-tors-zero}
Soit $A$ un anneau. Soit $P = A[E]$ un anneau de polynômes. Posons
$I = (e; e \in E) \subset P$. Les applications
$\text{Tor}_i^P(A, I^{n + 1}) \to \text{Tor}_i^P(A, I^n)$
sont nulles pour tous $i$ et $n$.
\end{lemma}

\begin{proof}
Notons $x_e \in P$ la variable correspondant à $e \in E$. Le complexe de
Koszul $K_\bullet$ associé aux $x_e$ fournit une résolution libre de $A$ sur
$P$. Ici, $K_i$ a pour base les produits extérieurs
$e_1 \wedge \ldots \wedge e_i$, $e_1, \ldots, e_i \in E$, et $d(e) = x_e$.
Ainsi, $K_\bullet \otimes_P I^n = I^nK_\bullet$ calcule
$\text{Tor}_i^P(A, I^n)$. Remarquons que tout est gradué, avec
$\deg(x_e) = 1$, $\deg(e) = 1$ et $\deg(a) = 0$ pour $a \in A$. Soit
$\xi \in I^{n + 1}K_i$ un cocycle homogène de degré $m$. On a
$m \geq i + 1 + n$. Il existe alors $\eta \in K_{i + 1}$ tel que
$\xi = \text{d}\eta$, puisque $K_\bullet$ est exact en degrés $ > 0$. (Le cas
$i = 0$ est laissé au lecteur.) Or $\deg(\eta) = m \geq i + 1 + n$. En
écrivant $\eta$ dans la base, on voit donc que ses coordonnées appartiennent à
$I^n$. Ainsi, $\xi$ s'envoie sur zéro dans l'homologie de $I^nK_\bullet$,
comme voulu.
\end{proof}

\begin{theorem}[Suite spectrale de Quillen]
\label{cotangent-theorem-quillen-spectral-sequence}
Soit $A \to B$ un homomorphisme d'anneaux surjectif. Considérons le faisceau
de $\underline{B}$-modules
$\Omega = \Omega_{\mathcal{O}/A} \otimes_\mathcal{O} \underline{B}$ sur
$\mathcal{C}_{B/A}$ ; voir la section \ref{cotangent-section-compute-L-pi-shriek}. Il
existe alors une suite spectrale dont la page $E_1$ est
$$
E_1^{p, q} =
H_{- p - q}(\mathcal{C}_{B/A}, \text{Sym}^p_{\underline{B}}(\Omega))
\Rightarrow \text{Tor}^A_{- p - q}(B, B)
$$
et dans laquelle $d_r$ est de bidegré $(r, -r + 1)$. De plus,
$H_i(\mathcal{C}_{B/A}, \text{Sym}^k_{\underline{B}}(\Omega)) = 0$ pour
$i < k$.
\end{theorem}

\begin{proof}
Soit $I \subset A$ le noyau de $A \to B$. Soit
$\mathcal{J} \subset \mathcal{O}$ le noyau de
$\mathcal{O} \to \underline{B}$. Alors
$I\mathcal{O} \subset \mathcal{J}$. Posons
$\mathcal{K} = \mathcal{J}/I\mathcal{O}$ et
$\overline{\mathcal{O}} = \mathcal{O}/I\mathcal{O}$.

\medskip\noindent
Pour tout objet $U = (P \to B)$ de $\mathcal{C}_{B/A}$, nous pouvons choisir
un isomorphisme $P = A[E]$ tel que l'homomorphisme $P \to B$ envoie chaque
$e \in E$ sur zéro. Alors
$J = \mathcal{J}(U) \subset P = \mathcal{O}(U)$
est égal à $J = IP + (e; e \in E)$. De plus,
$\overline{\mathcal{O}}(U) = B[E]$ et $K = \mathcal{K}(U) = (e; e \in E)$
est l'idéal engendré par les variables dans l'anneau de polynômes $B[E]$. En
particulier, il est clair que
$$
K/K^2 \xrightarrow{\text{d}} \Omega_{P/A} \otimes_P B
$$
est une bijection. Autrement dit, $\Omega = \mathcal{K}/\mathcal{K}^2$ et
$\text{Sym}_B^k(\Omega) = \mathcal{K}^k/\mathcal{K}^{k + 1}$. Remarquons que
$\pi_!(\Omega) = \Omega_{B/A} = 0$ (lemme \ref{cotangent-lemma-identify-H0}), puisque
$A \to B$ est surjectif (Algèbre, lemme
\ref{algebra-lemma-trivial-differential-surjective}). D'après le lemme
\ref{cotangent-lemma-vanishing-symmetric-powers}, nous en concluons que
$$
H_i(\mathcal{C}_{B/A}, \mathcal{K}^k/\mathcal{K}^{k + 1}) =
H_i(\mathcal{C}_{B/A}, \text{Sym}^k_{\underline{B}}(\Omega)) = 0
$$
pour $i < k$. Cela démontre la dernière assertion du théorème.

\medskip\noindent
Pour aborder le théorème, remarquons que
$$
B \otimes_A^\mathbf{L} B = L\pi_!(\mathcal{O}) \otimes_A^\mathbf{L} B =
L\pi_!(\mathcal{O} \otimes_{\underline{A}}^\mathbf{L} \underline{B}) =
L\pi_!(\overline{\mathcal{O}})
$$
La première égalité résulte du lemme \ref{cotangent-lemma-apply-O-B-comparison}, la
deuxième de Cohomologie sur les sites, lemme
\ref{sites-cohomology-lemma-change-of-rings}, et la troisième de la platitude
de $\mathcal{O}$ sur $\underline{A}$. Le faisceau
$\overline{\mathcal{O}}$ possède une filtration
$$
\ldots \subset
\mathcal{K}^3 \subset
\mathcal{K}^2 \subset
\mathcal{K} \subset
\overline{\mathcal{O}}
$$
Celle-ci induit une filtration $F$ sur un complexe $C$ représentant
$L\pi_!(\overline{\mathcal{O}})$, où $F^pC$ représente
$L\pi_!(\mathcal{K}^p)$ (nous omettons la construction de $C$ et de $F$).
Considérons la suite spectrale de Homologie, section
\ref{homology-section-filtered-complex}, associée à $(C, F)$. Sa page $E_1$
est
$$
E_1^{p, q} = H_{- p - q}(\mathcal{C}_{B/A}, \mathcal{K}^p/\mathcal{K}^{p + 1})
\quad\Rightarrow\quad
H_{- p - q}(\mathcal{C}_{B/A}, \overline{\mathcal{O}}) = 
\text{Tor}_{- p - q}^A(B, B)
$$
et ses différentielles sont
$E_r^{p, q} \to E_r^{p + r, q - r + 1}$. Pour montrer la convergence, nous
allons établir que, pour tout $k$, il existe un $c$ tel que
$H_i(\mathcal{C}_{B/A}, \mathcal{K}^n) = 0$
pour $i < k$ et $n > c$\footnote{A posteriori, on pourra en déduire
l'annulation « correcte » $H_i(\mathcal{C}_{B/A}, \mathcal{K}^n) = 0$ pour
$i < n$.}.

\medskip\noindent
Étant donné $k \geq 0$, posons $c = k^2$. Nous affirmons que
$$
H_i(\mathcal{C}_{B/A}, \mathcal{K}^{n + c}) \to
H_i(\mathcal{C}_{B/A}, \mathcal{K}^n)
$$
est nulle pour $i < k$ et tout $n \geq 0$. Remarquons que
$\mathcal{K}^n/\mathcal{K}^{n + c}$ possède une filtration finie dont les
quotients successifs $\mathcal{K}^m/\mathcal{K}^{m + 1}$,
$n \leq m < n + c$, vérifient
$H_i(\mathcal{C}_{B/A}, \mathcal{K}^m/\mathcal{K}^{m + 1}) = 0$ pour $i < n$
(voir ci-dessus). L'assertion implique donc
$H_i(\mathcal{C}_{B/A}, \mathcal{K}^{n + c}) = 0$ pour $i < k$ et tout
$n \geq k$, ce qu'il fallait montrer.

\medskip\noindent
Démonstration de l'assertion. Rappelons que, pour tout $\mathcal{O}$-module
$\mathcal{F}$, l'application
$\mathcal{F} \to \mathcal{F} \otimes_\mathcal{O}^\mathbf{L} B$ induit un
isomorphisme après application de $L\pi_!$ ; voir le lemme
\ref{cotangent-lemma-O-homology-B-homology}. Considérons l'application
$$
\mathcal{K}^{n + k} \otimes_\mathcal{O}^\mathbf{L} B
\longrightarrow
\mathcal{K}^n \otimes_\mathcal{O}^\mathbf{L} B
$$
Nous affirmons que cette application induit l'application nulle sur les
faisceaux de cohomologie en degrés $0, -1, \ldots, - k + 1$. Si cette seconde
assertion est vraie, la composée de $k$ applications
$$
\mathcal{K}^{n + c} \otimes_\mathcal{O}^\mathbf{L} B
\longrightarrow
\mathcal{K}^n \otimes_\mathcal{O}^\mathbf{L} B
$$
se factorise par
$\tau_{\leq -k}\mathcal{K}^n \otimes_\mathcal{O}^\mathbf{L} B$ et induit donc
zéro sur $H_i(\mathcal{C}_{B/A}, -) = L_i\pi_!( - )$ pour $i < k$ ; voir
Catégories dérivées, lemme
\ref{derived-lemma-trick-vanishing-composition}. D'après la remarque
précédente, il en va de même de
$H_i(\mathcal{C}_{B/A}, \mathcal{K}^{n + c}) \to
H_i(\mathcal{C}_{B/A}, \mathcal{K}^n)$
ce qui démontre la première assertion.

\medskip\noindent
Démonstration de la seconde assertion. L'énoncé est local ; nous pouvons donc
travailler au-dessus d'un objet $U = (P \to B)$ comme ci-dessus. Il faut
montrer que les applications
$$
\text{Tor}_i^P(B, K^{n + k}) \to \text{Tor}_i^P(B, K^n)
$$
sont nulles pour $i < k$. Il existe une suite spectrale
$$
\text{Tor}_a^P(P/IP, \text{Tor}_b^{P/IP}(B, K^n))
\Rightarrow
\text{Tor}_{a + b}^P(B, K^n),
$$
voir Compléments sur l'algèbre, exemple
\ref{more-algebra-example-tor-change-rings}. Il suffit donc de montrer que les
applications
$$
\text{Tor}_i^{P/IP}(B, K^{n + 1}) \to \text{Tor}_i^{P/IP}(B, K^n)
$$
sont nulles pour tout $i$. C'est le lemme \ref{cotangent-lemma-map-tors-zero}.
\end{proof}

\begin{remark}
\label{cotangent-remark-elucidate-ss}
Dans la situation du théorème \ref{cotangent-theorem-quillen-spectral-sequence}, posons
$I = \Ker(A \to B)$. Alors
$H^{-1}(L_{B/A}) = H_1(\mathcal{C}_{B/A}, \Omega) = I/I^2$, d'après le
lemme \ref{cotangent-lemma-surjection}.
Ainsi,
$H_k(\mathcal{C}_{B/A}, \text{Sym}^k(\Omega)) = \wedge^k_B(I/I^2)$ d'après la
remarque \ref{cotangent-remark-first-homology-symmetric-power}. La page $E_1$ se présente
donc sous la forme
$$
\begin{matrix}
B \\
0 \\
0 & I/I^2 \\
0 & H^{-2}(L_{B/A}) \\
0 & H^{-3}(L_{B/A}) & \wedge^2(I/I^2) \\
0 & H^{-4}(L_{B/A}) & H_3(\mathcal{C}_{B/A}, \text{Sym}^2(\Omega)) \\
0 & H^{-5}(L_{B/A}) & H_4(\mathcal{C}_{B/A}, \text{Sym}^2(\Omega)) &
\wedge^3(I/I^2)
\end{matrix}
$$
avec une différentielle horizontale. Nous obtenons ainsi des morphismes de bord
$\text{Tor}_i^A(B, B) \to H^{-i}(L_{B/A})$, $i > 0$, et
$\wedge^i_B(I/I^2) \to \text{Tor}_i^A(B, B)$. Enfin,
$\text{Tor}_1^A(B, B) = I/I^2$ et il existe une suite exacte à cinq termes
$$
\text{Tor}_3^A(B, B) \to H^{-3}(L_{B/A}) \to \wedge^2_B(I/I^2) \to
\text{Tor}_2^A(B, B) \to H^{-2}(L_{B/A}) \to 0
$$
formée des termes de bas degré.
\end{remark}

\begin{remark}
\label{cotangent-remark-elucidate-degree-two}
Soit $A \to B$ un homomorphisme d'anneaux. Soit $P_\bullet$ une résolution de
$B$ sur $A$ (remarque \ref{cotangent-remark-resolution}). Posons
$J_n = \Ker(P_n \to B)$. Remarquons que
$$
\text{Tor}_2^{P_n}(B, B) = 
\text{Tor}_1^{P_n}(J_n, B) =
\Ker(J_n \otimes_{P_n} J_n \to J_n^2).
$$
Ainsi, $H_2(L_{B/A})$ est canoniquement égal à
$$
\Coker(\text{Tor}_2^{P_1}(B, B) \to \text{Tor}_2^{P_0}(B, B))
$$
d'après la remarque \ref{cotangent-remark-surjection}. Pour expliciter ceci, choisissons
$P_2$, $P_1$, $P_0$ comme dans l'exemple
\ref{cotangent-example-resolution-length-two}. Nous affirmons que
$$
\text{Tor}_2^{P_1}(B, B) =
\wedge^2(\bigoplus\nolimits_{t \in T} B)\ \oplus
\ \bigoplus\nolimits_{t \in T} J_0\ \oplus
\ \text{Tor}_2^{P_0}(B, B)
$$
Plus précisément, les éléments de base $x_t \wedge x_{t'}$ du premier facteur
direct correspondent à l'élément
$x_t \otimes x_{t'} - x_{t'} \otimes x_t$ de $J_1 \otimes_{P_1} J_1$. Pour
$f \in J_0$, l'élément $x_t \otimes f$ du deuxième facteur direct correspond
à l'élément $x_t \otimes s_0(f) - s_0(f) \otimes x_t$ de
$J_1 \otimes_{P_1} J_1$. Enfin, l'application
$\text{Tor}_2^{P_0}(B, B) \to \text{Tor}_2^{P_1}(B, B)$ est donnée par $s_0$.
L'application
$d_0 - d_1 : \text{Tor}_2^{P_1}(B, B) \to \text{Tor}_2^{P_0}(B, B)$
est nulle sur le dernier facteur direct, envoie $x_t \otimes f$ sur
$f \otimes f_t - f_t \otimes f$, et $x_t \wedge x_{t'}$ sur
$f_t \otimes f_{t'} - f_{t'} \otimes f_t$. En définitive, nous en concluons
qu'il existe une suite exacte
$$
\wedge^2_B(J_0/J_0^2) \to \text{Tor}_2^{P_0}(B, B) \to H^{-2}(L_{B/A}) \to 0
$$
Nous obtenons ainsi une démonstration directe d'une conséquence de la suite
spectrale de Quillen examinée dans la remarque \ref{cotangent-remark-elucidate-ss}.
\end{remark}






\section{Comparaison avec Lichtenbaum-Schlessinger}
\label{cotangent-section-compare-higher}

\noindent
Soit $A \to B$ un homomorphisme d'anneaux. L'article
\cite{Lichtenbaum-Schlessinger} donne une détermination assez explicite de
$\tau_{\geq -2}L_{B/A}$, souvent utilisée pour calculer les espaces de
déformations verselles de singularités. En voici la construction. Choisissons
une algèbre polynomiale $P$ sur $A$ et une surjection $P \to B$ de noyau $I$.
Choisissons des générateurs $f_t$, $t \in T$, de $I$, ce qui induit une
surjection $F = \bigoplus_{t \in T} P \to I$, où $F$ est un $P$-module libre.
Soit $Rel \subset F$ le noyau de $F \to I$ ; autrement dit, $Rel$ est
l'ensemble des relations entre les $f_t$. Soit $TrivRel \subset Rel$ le
sous-module des relations triviales, c'est-à-dire le sous-module de $Rel$
engendré par les éléments
$(\ldots, f_{t'}, 0, \ldots, 0, -f_t, 0, \ldots)$. Considérons le complexe de
$B$-modules
\begin{equation}
\label{cotangent-equation-lichtenbaum-schlessinger}
Rel/TrivRel \longrightarrow
F \otimes_P B \longrightarrow
\Omega_{P/A} \otimes_P B
\end{equation}
où le dernier terme est placé en degré $0$. La première application est
l'application évidente et la seconde envoie l'élément de base correspondant à
$t \in T$ sur $\text{d}f_t \otimes 1$.

\begin{definition}
\label{cotangent-definition-biderivation}
Soit $A \to B$ un homomorphisme d'anneaux. Soit $M$ un $(B, B)$-bimodule sur
$A$. Une {\it $A$-bidérivation} est une application $A$-linéaire
$\lambda : B \to M$ telle que $\lambda(xy) = x\lambda(y) + \lambda(x)y$.
\end{definition}

\noindent
Pour une algèbre polynomiale, les bidérivations sont faciles à décrire.

\begin{lemma}
\label{cotangent-lemma-polynomial-ring-unique}
Soit $P = A[S]$ un anneau de polynômes sur $A$. Soit $M$ un
$(P, P)$-bimodule sur $A$. Étant donnés des $m_s \in M$, pour $s \in S$, il
existe une unique $A$-bidérivation $\lambda : P \to M$ qui envoie $s$ sur
$m_s$ pour $s \in S$.
\end{lemma}

\begin{proof}
Posons
$$
\lambda(s_1 \ldots s_t) =
\sum s_1 \ldots s_{i - 1} m_{s_i} s_{i + 1} \ldots s_t
$$
dans $M$. Le prolongement par $A$-linéarité est une bidérivation.
\end{proof}

\noindent
Voici l'énoncé de comparaison. Le lecteur pourra également consulter
\cite[page 206, Proposition 12]{Andre-Homologie} ou l'article \cite{Doncel},
qui prolonge le complexe (\ref{cotangent-equation-lichtenbaum-schlessinger}) d'un terme
et étend la comparaison à $\tau_{\geq -3}$.

\begin{lemma}
\label{cotangent-lemma-compare-higher}
Dans la situation précédente, notons $L$ le complexe
(\ref{cotangent-equation-lichtenbaum-schlessinger}). Il existe dans $D(B)$ une
application canonique $L_{B/A} \to L$ qui induit un isomorphisme
$\tau_{\geq -2}L_{B/A} \to L$ dans $D(B)$.
\end{lemma}

\begin{proof}
Soit $P_\bullet \to B$ une résolution de $B$ sur $A$ (remarque
\ref{cotangent-remark-resolution}). Nous identifierons $L_{B/A}$ à
$\Omega_{P_\bullet/A} \otimes B$. Pour construire l'application, effectuons
quelques choix.

\medskip\noindent
Choisissons un homomorphisme de $A$-algèbres $\psi : P_0 \to P$ compatible
avec les homomorphismes donnés $P_0 \to B$ et $P \to B$.

\medskip\noindent
Écrivons $P_1 = A[S]$ pour un certain ensemble $S$. Pour $s \in S$, nous
pouvons écrire
$$
\psi(d_0(s) - d_1(s)) = \sum p_{s, t} f_t
$$
pour certains $p_{s, t} \in P$. Considérons
$F = \bigoplus_{t \in T} P$ comme un $(P_1, P_1)$-bimodule au moyen des
homomorphismes $(\psi \circ d_0, \psi \circ d_1)$. D'après le lemme
\ref{cotangent-lemma-polynomial-ring-unique}, il existe une unique $A$-bidérivation
$\lambda : P_1 \to F$ envoyant $s$ sur le vecteur de coordonnées $p_{s, t}$.
Par construction, la composée
$$
P_1 \longrightarrow F \longrightarrow P
$$
envoie $f \in P_1$ sur $\psi(d_0(f) - d_1(f))$, car l'application
$f \mapsto \psi(d_0(f) - d_1(f))$ est une $A$-bidérivation qui coïncide avec la
composée sur les générateurs.

\medskip\noindent
Pour $g \in P_2$, nous affirmons que
$\lambda(d_0(g) - d_1(g) + d_2(g))$ appartient à $Rel$. En effet, d'après la
dernière observation du paragraphe précédent, l'image de
$\lambda(d_0(g) - d_1(g) + d_2(g))$ dans $P$ est
$$
\psi((d_0 - d_1)(d_0(g) - d_1(g) + d_2(g)))
$$
qui est nulle d'après Simplicial, section
\ref{simplicial-section-complexes}.

\medskip\noindent
Le choix de $\psi$ détermine une application
$$
\text{d}\psi \otimes 1 :
\Omega_{P_0/A} \otimes B
\longrightarrow
\Omega_{P/A} \otimes B
$$
La composée de $\lambda$ avec l'application $F \to F \otimes B$ est une
$A$-dérivation usuelle, puisque les deux structures de $P_1$-module sur
$F \otimes B$ coïncident. Ainsi, $\lambda$ détermine une application
$$
\overline{\lambda} :
\Omega_{P_1/A} \otimes B
\longrightarrow
F \otimes B
$$
Enfin, nous obtenons une application $B$-linéaire
$$
q :
\Omega_{P_2/A} \otimes B
\longrightarrow
Rel/TrivRel
$$
en envoyant $\text{d}g$ sur la classe de
$\lambda(d_0(g) - d_1(g) + d_2(g))$ dans le quotient.

\medskip\noindent
Le diagramme
$$
\xymatrix{
\Omega_{P_3/A} \otimes B \ar[r] \ar[d] &
\Omega_{P_2/A} \otimes B \ar[r] \ar[d]_q &
\Omega_{P_1/A} \otimes B \ar[r] \ar[d]_{\overline{\lambda}} &
\Omega_{P_0/A} \otimes B \ar[d]_{\text{d}\psi \otimes 1} \\
0 \ar[r] &
Rel/TrivRel \ar[r] &
F \otimes B \ar[r] &
\Omega_{P/A} \otimes B
}
$$
est commutatif (calcul omis), et nous obtenons l'application du lemme. La
remarque \ref{cotangent-remark-explicit-comparison-map} et le lemme
\ref{cotangent-lemma-relation-with-naive-cotangent-complex} montrent que cette
application induit des isomorphismes $H_1(L_{B/A}) \to H_1(L)$ et
$H_0(L_{B/A}) \to H_0(L)$.

\medskip\noindent
Il reste à voir que notre application $L_{B/A} \to L$ induit un isomorphisme
$H_2(L_{B/A}) \to H_2(L)$. Choisissons une résolution de $B$ sur $A$ telle que
$P_0 = P = A[u_i]$, puis $P_1$ et $P_2$ comme dans l'exemple
\ref{cotangent-example-resolution-length-two}. Dans la remarque
\ref{cotangent-remark-elucidate-degree-two}, nous avons construit une suite exacte
$$
\wedge^2_B(J_0/J_0^2) \to \text{Tor}_2^{P_0}(B, B) \to H^{-2}(L_{B/A}) \to 0
$$
où $P_0 = P$ et $J_0 = \Ker(P \to B) = I$. En calculant le groupe de Tor à
l'aide des suites exactes courtes $0 \to I \to P \to B \to 0$ et
$0 \to Rel \to F \to I \to 0$, nous trouvons que
$\text{Tor}_2^P(B, B) = \Ker(Rel \otimes B \to F \otimes B)$.
Sous cette identification, l'image de l'application
$\wedge^2_B(I/I^2) \to \text{Tor}_2^P(B, B)$ est exactement l'image de
$TrivRel \otimes B$. On voit donc que $H_2(L_{B/A}) \cong H_2(L)$.

\medskip\noindent
Enfin, il faut vérifier que notre application $L_{B/A} \to L$ induit bien cet
isomorphisme. Nous utiliserons sans autre mention les notations et les résultats
de l'exemple \ref{cotangent-example-resolution-length-two} et des remarques
\ref{cotangent-remark-elucidate-degree-two} et \ref{cotangent-remark-surjection}. Choisissons un
élément $\xi$ de
$\text{Tor}_2^{P_0}(B, B) = \Ker(I \otimes_P I \to I^2)$.
Écrivons $\xi = \sum h_{t', t}f_{t'} \otimes f_t$ pour certains
$h_{t', t} \in P$. En suivant les suites exactes précédentes, on trouve que
$\xi$ correspond à l'image dans $Rel \otimes B$ de l'élément
$r \in Rel \subset F = \bigoplus_{t \in T} P$ dont la coordonnée d'indice $t$
est $r_t = \sum_{t' \in T} h_{t', t}f_{t'}$. D'autre part, $\xi$ correspond à
l'élément de $H_2(L_{B/A}) = H_2(\Omega)$ qui est l'image, par
$\text{d} : H_2(\mathcal{J}/\mathcal{J}^2) \to H_2(\Omega)$, du bord de $\xi$
pour la $2$-extension
$$
0 \to
\text{Tor}_2^\mathcal{O}(\underline{B}, \underline{B})
\to 
\mathcal{J} \otimes_\mathcal{O} \mathcal{J} \to \mathcal{J}
\to
\mathcal{J}/\mathcal{J}^2 \to 0
$$
Calculons les transgressions successives de notre élément. Tout d'abord,
$$
\xi = (d_0 - d_1)(- \sum s_0(h_{t', t} f_{t'}) \otimes x_t)
$$
puis
$$
\sum s_0(h_{t', t} f_{t'}) x_t = d_0(v_r) - d_1(v_r) + d_2(v_r)
$$
par notre choix des variables $v$ dans l'exemple
\ref{cotangent-example-resolution-length-two}. Nous pouvons choisir l'application
$\lambda$ ci-dessus de telle sorte que
$\lambda(u_i) = 0$ et $\lambda(x_t) = - e_t$, où $e_t \in F$ désigne le
vecteur de base correspondant à $t \in T$. La
construction de l'application $q$ ci-dessus envoie donc $\text{d}v_r$ sur
$$
\lambda(\sum s_0(h_{t', t} f_{t'}) x_t) =
\sum\nolimits_t \left(\sum\nolimits_{t'} h_{t', t}f_{t'}\right) e_t
$$
qui coïncide avec l'image de $\xi$ dans $Rel \otimes B$ (les deux signes moins
apparus ci-dessus se compensent). Cette coïncidence achève la démonstration.
\end{proof}

\begin{remark}[Fonctorialité du complexe de Lichtenbaum-Schlessinger]
\label{cotangent-remark-functoriality-lichtenbaum-schlessinger}
Considérons un carré commutatif
$$
\xymatrix{
A' \ar[r] & B' \\
A \ar[u] \ar[r] & B \ar[u]
}
$$
d'homomorphismes d'anneaux. Choisissons une factorisation
$$
\xymatrix{
A' \ar[r] & P' \ar[r] & B' \\
A \ar[u] \ar[r] & P \ar[u] \ar[r] & B \ar[u]
}
$$
où $P$ est une algèbre polynomiale sur $A$ et $P'$ une algèbre polynomiale sur
$A'$. Choisissons des générateurs $f_t$, $t \in T$, de $\Ker(P \to B)$. Pour
$t \in T$, notons $f'_t$ l'image de $f_t$ dans $P'$. Choisissons des
$f'_s \in P'$ de sorte que les éléments $f'_t$, pour
$t \in T' = T \amalg S$, engendrent le noyau de $P' \to B'$. Posons
$F = \bigoplus_{t \in T} P$ et $F' = \bigoplus_{t' \in T'} P'$. Soient
$Rel = \Ker(F \to P)$ et $Rel' = \Ker(F' \to P')$, où les applications sont
données sur les coordonnées par la multiplication par $f_t$, resp.\ $f'_t$.
Enfin, soit $TrivRel$, resp.\ $TrivRel'$, le sous-module de $Rel$, resp.\
$Rel'$, engendré par les éléments
$(\ldots, f_{t'}, 0, \ldots, 0, -f_t, 0, \ldots)$
pour $t, t' \in T$, resp.\ $T'$. Ces choix donnent un diagramme commutatif
canonique
$$
\xymatrix{
L' : &
Rel'/TrivRel' \ar[r] &
F' \otimes_{P'} B' \ar[r] &
\Omega_{P'/A'} \otimes_{P'} B' \\
L : \ar[u] &
Rel/TrivRel \ar[r] \ar[u] &
F \otimes_P B \ar[r] \ar[u] &
\Omega_{P/A} \otimes_P B \ar[u]
}
$$
De plus, en suivant les choix effectués dans la démonstration du lemme
\ref{cotangent-lemma-compare-higher}, le lecteur constate que l'on obtient un diagramme
commutatif
$$
\xymatrix{
L_{B'/A'} \ar[r] & L' \\
L_{B/A} \ar[r] \ar[u] & L \ar[u]
}
$$
\end{remark}





\section{Le complexe cotangent d'une intersection complète locale}
\label{cotangent-section-lci}

\noindent
Si $A \to B$ est un homomorphisme d'intersection complète locale, alors
$L_{B/A}$ est un complexe parfait. Le lemme suivant est la clef de la
démonstration.

\begin{lemma}
\label{cotangent-lemma-special-case}
Soit $A = \mathbf{Z}[x_1, \ldots, x_n] \to B = \mathbf{Z}$ l'homomorphisme
d'anneaux qui envoie $x_i$ sur $0$ pour $i = 1, \ldots, n$. Soit
$I = (x_1, \ldots, x_n) \subset A$. Alors $L_{B/A}$ est quasi-isomorphe à
$I/I^2[1]$.
\end{lemma}

\begin{proof}
Il y a plusieurs façons de le démontrer. On peut, par exemple, construire
explicitement une résolution de $B$ sur $A$ et effectuer le calcul. Nous
utiliserons (\ref{cotangent-equation-triangle}). Considérons en effet le triangle
distingué
$$
L_{\mathbf{Z}[x_1, \ldots, x_n]/\mathbf{Z}}
\otimes_{\mathbf{Z}[x_1, \ldots, x_n]} \mathbf{Z} \to
L_{\mathbf{Z}/\mathbf{Z}} \to
L_{\mathbf{Z}/\mathbf{Z}[x_1, \ldots, x_n]}\to
L_{\mathbf{Z}[x_1, \ldots, x_n]/\mathbf{Z}}
\otimes_{\mathbf{Z}[x_1, \ldots, x_n]} \mathbf{Z}[1]
$$
Le complexe $L_{\mathbf{Z}[x_1, \ldots, x_n]/\mathbf{Z}}$ est quasi-isomorphe à
$\Omega_{\mathbf{Z}[x_1, \ldots, x_n]/\mathbf{Z}}$ d'après le lemme
\ref{cotangent-lemma-cotangent-complex-polynomial-algebra}. Le complexe
$L_{\mathbf{Z}/\mathbf{Z}}$ est nul dans $D(\mathbf{Z})$ d'après le lemme
\ref{cotangent-lemma-when-zero}. On voit ainsi que $L_{B/A}$ n'a qu'un seul groupe de
cohomologie non nul, qui est celui décrit dans l'énoncé d'après le lemme
\ref{cotangent-lemma-surjection}.
\end{proof}

\begin{lemma}
\label{cotangent-lemma-mod-regular-sequence}
Soit $A \to B$ un homomorphisme d'anneaux surjectif dont le noyau $I$ est
engendré par une suite régulière au sens de Koszul (par exemple une suite
régulière). Alors $L_{B/A}$ est quasi-isomorphe à $I/I^2[1]$.
\end{lemma}

\begin{proof}
Soit $f_1, \ldots, f_r \in I$ une suite régulière au sens de Koszul qui
engendre $I$. Considérons l'homomorphisme d'anneaux
$\mathbf{Z}[x_1, \ldots, x_r] \to A$ qui envoie $x_i$ sur $f_i$. Comme
$x_1, \ldots, x_r$ est une suite régulière dans
$\mathbf{Z}[x_1, \ldots, x_r]$, le complexe de Koszul associé à
$x_1, \ldots, x_r$ est une résolution libre de
$\mathbf{Z} = \mathbf{Z}[x_1, \ldots, x_r]/(x_1, \ldots, x_r)$
sur $\mathbf{Z}[x_1, \ldots, x_r]$ (voir Compléments sur l'algèbre, lemme
\ref{more-algebra-lemma-regular-koszul-regular}). Ainsi, l'hypothèse que
$f_1, \ldots, f_r$ est régulière au sens de Koszul signifie exactement que
$B = A \otimes_{\mathbf{Z}[x_1, \ldots, x_r]}^\mathbf{L} \mathbf{Z}$.
Par conséquent,
$L_{B/A} = L_{\mathbf{Z}/\mathbf{Z}[x_1, \ldots, x_r]}
\otimes_\mathbf{Z}^\mathbf{L} B$ d'après les lemmes
\ref{cotangent-lemma-flat-base-change-cotangent-complex} et \ref{cotangent-lemma-special-case}.
\end{proof}

\begin{lemma}
\label{cotangent-lemma-mod-Koszul-regular-ideal}
Soit $A \to B$ un homomorphisme d'anneaux surjectif dont le noyau $I$ est un
idéal de Koszul. Alors $L_{B/A}$ est quasi-isomorphe à $I/I^2[1]$.
\end{lemma}

\begin{proof}
Localement sur $\Spec(A)$, l'idéal $I$ est engendré par une suite régulière au
sens de Koszul ; voir Compléments sur l'algèbre, définition
\ref{more-algebra-definition-regular-ideal}. Le résultat découle donc du lemme
\ref{cotangent-lemma-flat-base-change-cotangent-complex}.
\end{proof}

\begin{proposition}
\label{cotangent-proposition-cotangent-complex-local-complete-intersection}
Soit $A \to B$ un homomorphisme d'intersection complète locale. Alors
$L_{B/A}$ est un complexe parfait d'amplitude de Tor contenue dans $[-1, 0]$.
\end{proposition}

\begin{proof}
Choisissons une surjection $P = A[x_1, \ldots, x_n] \to B$ de noyau $J$.
D'après le lemme \ref{cotangent-lemma-relation-with-naive-cotangent-complex}, le
complexe $J/J^2 \to \bigoplus B\text{d}x_i$ est quasi-isomorphe à
$\tau_{\geq -1}L_{B/A}$. Remarquons que $J/J^2$ est projectif de type fini
(Compléments sur l'algèbre, lemme
\ref{more-algebra-lemma-quasi-regular-ideal-finite-projective}) ; ainsi,
$\tau_{\geq -1}L_{B/A}$ est un complexe parfait d'amplitude de Tor contenue
dans $[-1, 0]$. Il suffit donc de montrer que $H^i(L_{B/A}) = 0$ pour
$i \not \in [-1, 0]$. Cela résulte de (\ref{cotangent-equation-triangle})
$$
L_{P/A} \otimes_P^\mathbf{L} B \to L_{B/A} \to L_{B/P} \to
L_{P/A} \otimes_P^\mathbf{L} B[1]
$$
et du lemme \ref{cotangent-lemma-mod-Koszul-regular-ideal}, qui montrent que
$H^i(L_{B/P})$ est nul sauf si $i \in \{-1, 0\}$. (Pour le terme de gauche,
nous utilisons également le lemme
\ref{cotangent-lemma-cotangent-complex-polynomial-algebra}.)
\end{proof}








\section{Produits tensoriels et complexe cotangent}
\label{cotangent-section-tensor-product}

\noindent
Soit $R$ un anneau et soient $A$, $B$ des $R$-algèbres. Dans cette section,
nous étudions $L_{A \otimes_R B/R}$. L'essentiel des renseignements recherchés
est contenu dans le diagramme suivant
\begin{equation}
\label{cotangent-equation-tensor-product}
\vcenter{
\xymatrix{
L_{A/R} \otimes_A^\mathbf{L} (A \otimes_R B) \ar[r] &
L_{A \otimes_R B/B} \ar[r] &
E \\
L_{A/R} \otimes_A^\mathbf{L} (A \otimes_R B) \ar[r] \ar@{=}[u] &
L_{A \otimes_R B/R} \ar[r] \ar[u] &
L_{A \otimes_R B/A} \ar[u] \\
 &
L_{B/R} \otimes_B^\mathbf{L} (A \otimes_R B) \ar[u] \ar@{=}[r] &
L_{B/R} \otimes_B^\mathbf{L} (A \otimes_R B) \ar[u]
}
}
\end{equation}
Explication : la ligne médiane est le triangle fondamental
(\ref{cotangent-equation-triangle}) associé aux homomorphismes d'anneaux
$R \to A \to A \otimes_R B$. La colonne médiane est le triangle fondamental
(\ref{cotangent-equation-triangle}) associé aux homomorphismes d'anneaux
$R \to B \to A \otimes_R B$. Ensuite, $E$ est un objet de
$D(A \otimes_R B)$ qui « complète » le coin supérieur droit, c'est-à-dire qui
fait de la ligne supérieure et de la colonne droite des triangles distingués.
Un tel $E$ existe d'après Catégories dérivées, proposition
\ref{derived-proposition-9}, appliquée au carré inférieur gauche (en plaçant
$0$ à la position manquante). Plus explicitement, nous pourrions par exemple
définir $E$ comme le cône (Catégories dérivées, définition
\ref{derived-definition-cone}) de l'application de complexes
$$
L_{A/R} \otimes_A^\mathbf{L} (A \otimes_R B) \oplus
L_{B/R} \otimes_B^\mathbf{L} (A \otimes_R B)
\longrightarrow
L_{A \otimes_R B/R}
$$
et obtenir les deux applications de but $E$ en appliquant TR3. Dans le cas
Tor-indépendant, l'objet $E$ est nul.

\begin{lemma}
\label{cotangent-lemma-tensor-product-tor-independent}
Si $A$ et $B$ sont des $R$-algèbres Tor-indépendantes, alors l'objet $E$ de
(\ref{cotangent-equation-tensor-product}) est nul. Dans ce cas, nous avons
$$
L_{A \otimes_R B/R} =
L_{A/R} \otimes_A^\mathbf{L} (A \otimes_R B) \oplus
L_{B/R} \otimes_B^\mathbf{L} (A \otimes_R B)
$$
qui est représenté par le complexe
$L_{A/R} \otimes_R B \oplus L_{B/R} \otimes_R A $
de $A \otimes_R B$-modules.
\end{lemma}

\begin{proof}
Les deux premières assertions résultent immédiatement du lemme
\ref{cotangent-lemma-flat-base-change-cotangent-complex}. La dernière résulte du fait que
$L_{A/R}$ est un complexe de $A$-modules libres ; ainsi,
$L_{A/R} \otimes_A^\mathbf{L} (A \otimes_R B)$ est représenté par
$L_{A/R} \otimes_A (A \otimes_R B) = L_{A/R} \otimes_R B$
\end{proof}

\noindent
En général, nous pouvons dire ce qui suit de l'objet $E$.

\begin{lemma}
\label{cotangent-lemma-tensor-product}
Soit $R$ un anneau et soient $A$, $B$ des $R$-algèbres. L'objet $E$ de
(\ref{cotangent-equation-tensor-product}) vérifie
$$
H^i(E) =
\left\{
\begin{matrix}
0 & \text{si} & i \geq -1 \\
\text{Tor}_1^R(A, B) & \text{si} & i = -2
\end{matrix}
\right.
$$
\end{lemma}

\begin{proof}
Nous utilisons la description de $E$ comme le cône de
$L_{B/R} \otimes_B^\mathbf{L} (A \otimes_R B) \to L_{A \otimes_R B/A}$.
D'après le lemme \ref{cotangent-lemma-compare-higher}, les troncations canoniques
$\tau_{\geq -2}L_{B/R}$ et $\tau_{\geq -2}L_{A \otimes_R B/A}$ sont calculées
par le complexe de Lichtenbaum-Schlessinger
(\ref{cotangent-equation-lichtenbaum-schlessinger}). Ces isomorphismes sont compatibles
à la fonctorialité (remarque
\ref{cotangent-remark-functoriality-lichtenbaum-schlessinger}). Nous travaillons donc
avec les complexes de Lichtenbaum-Schlessinger dans cette démonstration.

\medskip\noindent
Choisissons une algèbre polynomiale $P$ sur $R$ et une surjection $P \to B$.
Choisissons des générateurs $f_t \in P$, $t \in T$, du noyau de cette
surjection. Soit $Rel \subset F = \bigoplus_{t \in T} P$ le noyau de
l'application $F \to P$ qui envoie le vecteur de base correspondant à $t$ sur
$f_t$. Posons $P_A = A \otimes_R P$ et
$F_A = A \otimes_R F = P_A \otimes_P F$. Soit $Rel_A$ le noyau de
l'application $F_A \to P_A$. La suite exacte
$$
0 \to Rel \to F \to P \to B \to 0
$$
et les suites exactes courtes usuelles pour Tor donnent une suite exacte
$$
A \otimes_R Rel \to Rel_A \to \text{Tor}_1^R(A, B) \to 0
$$
Remarquons que $P_A \to A \otimes_R B$ est une surjection dont le noyau est
engendré par les éléments $1 \otimes f_t$ de $P_A$. Notons
$TrivRel_A \subset Rel_A$ le sous-$P_A$-module engendré par les éléments
$(\ldots, 1 \otimes f_{t'}, 0, \ldots,
0, - 1 \otimes f_t \otimes 1, 0, \ldots)$.
Comme $TrivRel \otimes_R A \to TrivRel_A$ est surjectif, nous obtenons une
suite exacte canonique
$$
A \otimes_R (Rel/TrivRel) \to Rel_A/TrivRel_A \to \text{Tor}_1^R(A, B) \to 0
$$
L'application entre les complexes de Lichtenbaum-Schlessinger est donnée par
le diagramme
$$
\xymatrix{
Rel_A/TrivRel_A \ar[r] &
F_A \otimes_{P_A} (A \otimes_R B) \ar[r] &
\Omega_{P_A/A \otimes_R B} \otimes_{P_A} (A \otimes_R B) \\
Rel/TrivRel \ar[r] \ar[u]_{-2} &
F \otimes_P B \ar[r] \ar[u]_{-1} &
\Omega_{P/A} \otimes_P B \ar[u]_0
}
$$
Remarquons que les applications verticales $-1$ et $0$ induisent un
isomorphisme après application du foncteur
$A \otimes_R - = P_A \otimes_P -$ à la source, et que l'application verticale
$-2$ donne exactement l'application dont le conoyau est le module de Tor
recherché, comme nous l'avons vu ci-dessus.
\end{proof}












\section{Déformations d'homomorphismes d'anneaux et complexe cotangent}
\label{cotangent-section-deformations}

\noindent
Cette section prolonge Théorie des déformations, section
\ref{defos-section-deformations}, que nous invitons vivement le lecteur à lire
d'abord. Partons d'un homomorphisme d'anneaux surjectif $A' \to A$ dont le
noyau est un idéal $I$ de carré nul. Supposons de plus donnés un homomorphisme
d'anneaux $A \to B$, un $B$-module $N$ et une application de $A$-modules
$c : I \to N$. Nous cherchons ici si l'on peut remplacer le point
d'interrogation dans le diagramme suivant
\begin{equation}
\label{cotangent-equation-to-solve}
\vcenter{
\xymatrix{
0 \ar[r] & N \ar[r] & {?} \ar[r] & B \ar[r] & 0 \\
0 \ar[r] & I \ar[u]^c \ar[r] & A' \ar[u] \ar[r] & A \ar[u] \ar[r] & 0
}
}
\end{equation}
et, si une solution existe, à quel point elle est unique. Plus précisément,
nous cherchons une surjection de $A'$-algèbres $B' \to B$ dont le noyau est
un idéal de carré nul identifié à $N$, et telle que $A' \to B'$ induise
l'application donnée $c$. Nous dirons que $B'$ est une {\it solution} de
(\ref{cotangent-equation-to-solve}).

\begin{lemma}
\label{cotangent-lemma-find-obstruction}
Dans la situation précédente :
\begin{enumerate}
\item il existe un élément canonique $\xi \in \Ext^2_B(L_{B/A}, N)$ dont
l'annulation est une condition nécessaire et suffisante pour l'existence d'une
solution de (\ref{cotangent-equation-to-solve}) ;
\item s'il existe une solution, l'ensemble des classes d'isomorphisme de
solutions est un espace principal homogène sous $\Ext^1_B(L_{B/A}, N)$ ;
\item étant donnée une solution $B'$, l'ensemble des automorphismes de $B'$
qui s'insèrent dans (\ref{cotangent-equation-to-solve}) est canoniquement isomorphe à
$\Ext^0_B(L_{B/A}, N)$.
\end{enumerate}
\end{lemma}

\begin{proof}
Au moyen des identifications $\NL_{B/A} = \tau_{\geq -1}L_{B/A}$ (lemme
\ref{cotangent-lemma-relation-with-naive-cotangent-complex}) et
$H^0(L_{B/A}) = \Omega_{B/A}$ (lemme \ref{cotangent-lemma-identify-H0}), nous avons déjà
établi (2) et (3) dans Théorie des déformations, lemmes
\ref{defos-lemma-huge-diagram} et \ref{defos-lemma-choices}.

\medskip\noindent
Démonstration de (1). En gros, l'assertion découle de la discussion de Théorie
des déformations, remarque \ref{defos-remark-parametrize-solutions}, en
remplaçant le complexe cotangent naïf par le complexe cotangent entier. Voici
une explication plus détaillée. D'après Théorie des déformations, lemme
\ref{defos-lemma-parametrize-solutions} et remarque
\ref{defos-remark-parametrize-solutions}, il existe un élément
$$
\xi' \in
\Ext^1_A(\NL_{A/A'}, N) =
\Ext^1_B(\NL_{A/A'} \otimes_A^\mathbf{L} B, N) =
\Ext^1_B(L_{A/A'} \otimes_A^\mathbf{L} B, N)
$$
(pour les égalités, voir Théorie des déformations, remarque
\ref{defos-remark-parametrize-solutions}, et utiliser
$\NL_{A'/A} = \tau_{\geq -1} L_{A'/A}$) tel qu'une solution existe si et
seulement si cet élément appartient à l'image de l'application
$$
\Ext^1_B(\NL_{B/A'}, N) = \Ext^1_B(L_{B/A'}, N)
\longrightarrow
\Ext^1_B(L_{A/A'} \otimes_A^\mathbf{L} B, N)
$$
Le triangle distingué (\ref{cotangent-equation-triangle}) associé à
$A' \to A \to B$ donne naissance à une suite exacte longue
$$
\ldots \to
\Ext^1_B(L_{B/A'}, N) \to
\Ext^1_B(L_{A/A'} \otimes_A^\mathbf{L} B, N) \to
\Ext^2_B(L_{B/A}, N) \to \ldots
$$
Il suffit donc de prendre pour $\xi$ l'image de $\xi'$.
\end{proof}





\section{La classe d'Atiyah d'un module}
\label{cotangent-section-atiyah}

\noindent
Soit $A \to B$ un homomorphisme d'anneaux. Soit $M$ un $B$-module. Soit
$P \to B$ un objet de $\mathcal{C}_{B/A}$ (section
\ref{cotangent-section-compute-L-pi-shriek}). Considérons l'extension des parties
principales
$$
0 \to \Omega_{P/A} \otimes_P M \to P^1_{P/A}(M) \to M \to 0
$$
voir Algèbre, lemme \ref{algebra-lemma-sequence-of-principal-parts}. Cette
suite est fonctorielle en $P$ d'après Algèbre, remarque
\ref{algebra-remark-functoriality-principal-parts}. Nous obtenons ainsi une
suite exacte courte de faisceaux de $\mathcal{O}$-modules
$$
0 \to \Omega_{\mathcal{O}/\underline{A}} \otimes_\mathcal{O} \underline{M} \to
P^1_{\mathcal{O}/\underline{A}}(M) \to \underline{M} \to 0
$$
sur $\mathcal{C}_{B/A}$. Nous avons
$L\pi_!(\Omega_{\mathcal{O}/\underline{A}} \otimes_\mathcal{O} \underline{M})
= L_{B/A} \otimes_B M = L_{B/A} \otimes_B^\mathbf{L} M$
d'après le lemme \ref{cotangent-lemma-pi-shriek-standard} et la platitude des termes de
$L_{B/A}$. De plus, $L\pi_!(\underline{M}) = M$ d'après le lemme
\ref{cotangent-lemma-pi-lower-shriek-constant-sheaf}. On obtient donc un triangle
distingué
\begin{equation}
\label{cotangent-equation-atiyah}
L_{B/A} \otimes_B^\mathbf{L} M \to
L\pi_!\left(P^1_{\mathcal{O}/\underline{A}}(M)\right) \to M
\to L_{B/A} \otimes_B^\mathbf{L} M [1]
\end{equation}
dans $D(B)$. Nous utilisons ici Cohomologie sur les sites, remarque
\ref{sites-cohomology-remark-O-homology-B-homology-general}, pour obtenir un
triangle distingué dans $D(B)$, et pas seulement dans $D(A)$.

\begin{definition}
\label{cotangent-definition-atiyah-class}
Soit $A \to B$ un homomorphisme d'anneaux. Soit $M$ un $B$-module.
L'application $M \to L_{B/A} \otimes_B^\mathbf{L} M[1]$ de
(\ref{cotangent-equation-atiyah}) est appelée la {\it classe d'Atiyah} de $M$.
\end{definition}




\section{Le complexe cotangent}
\label{cotangent-section-cotangent-complex}

\noindent
Dans cette section, nous étudions le complexe cotangent d'un homomorphisme de
faisceaux d'anneaux sur un site. Dans les sections suivantes, nous
spécialiserons cette construction afin d'obtenir le complexe cotangent d'un
morphisme de topos annelés, d'espaces annelés, de schémas, d'espaces
algébriques, etc.

\medskip\noindent
Soit $\mathcal{C}$ un site et notons $\Sh(\mathcal{C})$ le topos associé. Soit
$\mathcal{A}$ un faisceau d'anneaux sur $\mathcal{C}$. Notons
$\mathcal{A}\textit{-Alg}$ la catégorie des $\mathcal{A}$-algèbres.
Considérons le couple de foncteurs adjoints $(U, V)$, où
$V : \mathcal{A}\textit{-Alg} \to \Sh(\mathcal{C})$ est le foncteur d'oubli et
où $U : \Sh(\mathcal{C}) \to \mathcal{A}\textit{-Alg}$ associe à un faisceau
d'ensembles $\mathcal{E}$ l'algèbre polynomiale
$\mathcal{A}[\mathcal{E}]$ en $\mathcal{E}$ sur $\mathcal{A}$. Soit
$X_\bullet$ l'objet simplicial de
$\text{Fun}(\mathcal{A}\textit{-Alg}, \mathcal{A}\textit{-Alg})$ construit dans
Simplicial, section \ref{simplicial-section-standard}.

\medskip\noindent
Supposons maintenant que $\mathcal{A} \to \mathcal{B}$ soit un homomorphisme
de faisceaux d'anneaux. Alors $\mathcal{B}$ est un objet de la catégorie
$\mathcal{A}\textit{-Alg}$. Notons
$\mathcal{P}_\bullet = X_\bullet(\mathcal{B})$ la $\mathcal{A}$-algèbre
simpliciale ainsi obtenue. Rappelons que
$\mathcal{P}_0 = \mathcal{A}[\mathcal{B}]$,
$\mathcal{P}_1 = \mathcal{A}[\mathcal{A}[\mathcal{B}]]$, et ainsi de suite.
Rappelons également qu'il existe une augmentation
$$
\epsilon : \mathcal{P}_\bullet \longrightarrow \mathcal{B}
$$
où $\mathcal{B}$ est regardée comme une $\mathcal{A}$-algèbre simpliciale
constante.

\begin{definition}
\label{cotangent-definition-standard-resolution-sheaves-rings}
Soit $\mathcal{C}$ un site. Soit $\mathcal{A} \to \mathcal{B}$ un
homomorphisme de faisceaux d'anneaux sur $\mathcal{C}$. La {\it résolution
standard de $\mathcal{B}$ sur $\mathcal{A}$} est l'augmentation
$\epsilon : \mathcal{P}_\bullet \to \mathcal{B}$
dont les termes sont
$$
\mathcal{P}_0 = \mathcal{A}[\mathcal{B}],\quad
\mathcal{P}_1 = \mathcal{A}[\mathcal{A}[\mathcal{B}]],\quad \ldots
$$
et dont les applications sont celles construites ci-dessus.
\end{definition}

\noindent
Cette définition étant posée, le complexe cotangent d'un homomorphisme de
faisceaux d'anneaux se définit comme suit. Nous utiliserons le module des
différentielles défini dans Modules sur les sites, section
\ref{sites-modules-section-differentials}.

\begin{definition}
\label{cotangent-definition-cotangent-complex-morphism-sheaves-rings}
Soit $\mathcal{C}$ un site. Soit $\mathcal{A} \to \mathcal{B}$ un
homomorphisme de faisceaux d'anneaux sur $\mathcal{C}$. Le {\it complexe
cotangent} $L_{\mathcal{B}/\mathcal{A}}$ est le complexe de
$\mathcal{B}$-modules associé au module simplicial
$$
\Omega_{\mathcal{P}_\bullet/\mathcal{A}}
\otimes_{\mathcal{P}_\bullet, \epsilon} \mathcal{B}
$$
où $\epsilon : \mathcal{P}_\bullet \to \mathcal{B}$ est la résolution standard
de $\mathcal{B}$ sur $\mathcal{A}$. Nous considérons généralement
$L_{\mathcal{B}/\mathcal{A}}$ comme un objet de $D(\mathcal{B})$.
\end{definition}

\noindent
Ces constructions possèdent une fonctorialité analogue à celle étudiée dans la
section \ref{cotangent-section-functoriality}. Plus précisément, étant donné un diagramme
commutatif
\begin{equation}
\label{cotangent-equation-commutative-square-sheaves}
\vcenter{
\xymatrix{
\mathcal{B} \ar[r] & \mathcal{B}' \\
\mathcal{A} \ar[u] \ar[r] & \mathcal{A}' \ar[u]
}
}
\end{equation}
de faisceaux d'anneaux sur $\mathcal{C}$, il existe une application canonique
$\mathcal{B}$-linéaire de complexes
$$
L_{\mathcal{B}/\mathcal{A}} \longrightarrow L_{\mathcal{B}'/\mathcal{A}'}
$$
construite de la manière suivante. Si $\mathcal{P}_\bullet \to \mathcal{B}$ est
la résolution standard de $\mathcal{B}$ sur $\mathcal{A}$ et si
$\mathcal{P}'_\bullet \to \mathcal{B}'$ est la résolution standard de
$\mathcal{B}'$ sur $\mathcal{A}'$, il existe une application canonique
$\mathcal{P}_\bullet \to \mathcal{P}'_\bullet$ de $\mathcal{A}$-algèbres
simpliciales, compatible avec les augmentations
$\mathcal{P}_\bullet \to \mathcal{B}$ et
$\mathcal{P}'_\bullet \to \mathcal{B}'$. Les applications
$$
\mathcal{P}_0 = \mathcal{A}[\mathcal{B}]
\longrightarrow
\mathcal{A}'[\mathcal{B}'] = \mathcal{P}'_0,
\quad
\mathcal{P}_1 = \mathcal{A}[\mathcal{A}[\mathcal{B}]]
\longrightarrow
\mathcal{A}'[\mathcal{A}'[\mathcal{B}']] = \mathcal{P}'_1
$$
et ainsi de suite sont données par les homomorphismes
$\mathcal{A} \to \mathcal{A}'$ et $\mathcal{B} \to \mathcal{B}'$. L'application
recherchée $L_{\mathcal{B}/\mathcal{A}} \to L_{\mathcal{B}'/\mathcal{A}'}$
provient alors des applications associées sur les faisceaux de
différentielles.

\begin{lemma}
\label{cotangent-lemma-pullback-cotangent-morphism-topoi}
Soit $f : \Sh(\mathcal{D}) \to \Sh(\mathcal{C})$ un morphisme de topos. Soit
$\mathcal{A} \to \mathcal{B}$ un homomorphisme de faisceaux d'anneaux sur
$\mathcal{C}$. Alors
$f^{-1}L_{\mathcal{B}/\mathcal{A}} = L_{f^{-1}\mathcal{B}/f^{-1}\mathcal{A}}$.
\end{lemma}

\begin{proof}
Le diagramme
$$
\xymatrix{
\mathcal{A}\textit{-Alg} \ar[d]_{f^{-1}} \ar[r] &
\Sh(\mathcal{C}) \ar@<1ex>[l] \ar[d]^{f^{-1}} \\
f^{-1}\mathcal{A}\textit{-Alg} \ar[r] & \Sh(\mathcal{D}) \ar@<1ex>[l]
}
$$
est commutatif.
\end{proof}

\begin{lemma}
\label{cotangent-lemma-compute-L-morphism-sheaves-rings}
Soit $\mathcal{C}$ un site. Soit $\mathcal{A} \to \mathcal{B}$ un
homomorphisme de faisceaux d'anneaux sur $\mathcal{C}$. Alors
$H^i(L_{\mathcal{B}/\mathcal{A}})$ est le faisceau associé au préfaisceau
$U \mapsto H^i(L_{\mathcal{B}(U)/\mathcal{A}(U)})$.
\end{lemma}

\begin{proof}
Soit $\mathcal{C}'$ le site obtenu en munissant $\mathcal{C}$ de la topologie
chaotique (les préfaisceaux sont alors des faisceaux). Il existe un morphisme
de topos $f : \Sh(\mathcal{C}) \to \Sh(\mathcal{C}')$, où $f_*$ est
l'inclusion des faisceaux dans les préfaisceaux et $f^{-1}$ la faisceautisation.
D'après le lemme \ref{cotangent-lemma-pullback-cotangent-morphism-topoi}, il suffit de
démontrer le résultat pour $\mathcal{C}'$, c'est-à-dire lorsque
$\mathcal{C}$ est muni de la topologie chaotique.

\medskip\noindent
Si $\mathcal{C}$ est muni de la topologie chaotique, alors
$L_{\mathcal{B}/\mathcal{A}}(U)$ est égal à
$L_{\mathcal{B}(U)/\mathcal{A}(U)}$, car le diagramme
$$
\xymatrix{
\mathcal{A}\textit{-Alg} \ar[d]_{\text{sections sur }U} \ar[r] &
\Sh(\mathcal{C}) \ar@<1ex>[l] \ar[d]^{\text{sections sur }U} \\
\mathcal{A}(U)\textit{-Alg} \ar[r] & \textit{Ensembles} \ar@<1ex>[l]
}
$$
est commutatif.
\end{proof}

\begin{remark}
\label{cotangent-remark-map-sections-over-U}
Il ressort de la démonstration du lemme
\ref{cotangent-lemma-compute-L-morphism-sheaves-rings} que, pour tout
$U \in \Ob(\mathcal{C})$, il existe une application canonique
$L_{\mathcal{B}(U)/\mathcal{A}(U)} \to L_{\mathcal{B}/\mathcal{A}}(U)$
de complexes de $\mathcal{B}(U)$-modules. De plus, ces applications sont
compatibles aux applications de restriction, et le complexe
$L_{\mathcal{B}/\mathcal{A}}$ est la faisceautisation de la règle
$U \mapsto L_{\mathcal{B}(U)/\mathcal{A}(U)}$.
\end{remark}

\begin{lemma}
\label{cotangent-lemma-H0-L-morphism-sheaves-rings}
Soit $\mathcal{C}$ un site. Soit $\mathcal{A} \to \mathcal{B}$ un
homomorphisme de faisceaux d'anneaux sur $\mathcal{C}$. Alors
$H^0(L_{\mathcal{B}/\mathcal{A}}) = \Omega_{\mathcal{B}/\mathcal{A}}$.
\end{lemma}

\begin{proof}
Cela résulte des lemmes \ref{cotangent-lemma-compute-L-morphism-sheaves-rings} et
\ref{cotangent-lemma-identify-H0}, ainsi que de Modules sur les sites, lemme
\ref{sites-modules-lemma-differentials-sheafify}.
\end{proof}

\begin{lemma}
\label{cotangent-lemma-compute-L-product-sheaves-rings}
Soit $\mathcal{C}$ un site. Soient $\mathcal{A} \to \mathcal{B}$ et
$\mathcal{A} \to \mathcal{B}'$ des homomorphismes de faisceaux d'anneaux sur
$\mathcal{C}$. Alors
$$
L_{\mathcal{B} \times \mathcal{B}'/\mathcal{A}}
\longrightarrow
L_{\mathcal{B}/\mathcal{A}} \oplus L_{\mathcal{B}'/\mathcal{A}}
$$
est un isomorphisme dans $D(\mathcal{B} \times \mathcal{B}')$.
\end{lemma}

\begin{proof}
D'après le lemme \ref{cotangent-lemma-compute-L-morphism-sheaves-rings}, il suffit de le
démontrer pour des homomorphismes d'anneaux. Dans le cas des anneaux, il s'agit
du lemme \ref{cotangent-lemma-cotangent-complex-product}.
\end{proof}

\noindent
Le triangle fondamental pour le complexe cotangent de faisceaux d'anneaux est
une conséquence immédiate du résultat pour les homomorphismes d'anneaux.

\begin{lemma}
\label{cotangent-lemma-triangle-sheaves-rings}
Soit $\mathcal{D}$ un site. Soient
$\mathcal{A} \to \mathcal{B} \to \mathcal{C}$ des homomorphismes de faisceaux
d'anneaux sur $\mathcal{D}$. Il existe un triangle distingué canonique
$$
L_{\mathcal{B}/\mathcal{A}} \otimes_\mathcal{B}^\mathbf{L} \mathcal{C}
\to L_{\mathcal{C}/\mathcal{A}} \to L_{\mathcal{C}/\mathcal{B}} \to
L_{\mathcal{B}/\mathcal{A}} \otimes_\mathcal{B}^\mathbf{L} \mathcal{C}[1]
$$
dans $D(\mathcal{C})$.
\end{lemma}

\begin{proof}
Nous utiliserons la méthode décrite dans les remarques \ref{cotangent-remark-triangle} et
\ref{cotangent-remark-explicit-map} pour construire le triangle, et ferons librement
usage des résultats qui y sont mentionnés. Comme dans ces remarques, nous
construisons d'abord le triangle lorsque
$\mathcal{B} \to \mathcal{C}$ est un homomorphisme injectif de faisceaux
d'anneaux. Dans ce cas, posons :
\begin{enumerate}
\item $\mathcal{P}_\bullet$ est la résolution standard de $\mathcal{B}$ sur
$\mathcal{A}$,
\item $\mathcal{Q}_\bullet$ est la résolution standard de $\mathcal{C}$ sur
$\mathcal{A}$,
\item $\mathcal{R}_\bullet$ est la résolution standard de $\mathcal{C}$ sur
$\mathcal{B}$,
\item $\mathcal{S}_\bullet$ est la résolution standard de $\mathcal{B}$ sur
$\mathcal{B}$,
\item $\overline{\mathcal{Q}}_\bullet =
\mathcal{Q}_\bullet \otimes_{\mathcal{P}_\bullet} \mathcal{B}$, et
\item $\overline{\mathcal{R}}_\bullet =
\mathcal{R}_\bullet \otimes_{\mathcal{S}_\bullet} \mathcal{B}$.
\end{enumerate}
Le triangle distingué est celui qui est associé à la suite exacte courte de
$\mathcal{C}$-modules simpliciaux
$$
0 \to
\Omega_{\mathcal{P}_\bullet/\mathcal{A}}
\otimes_{\mathcal{P}_\bullet} \mathcal{C} \to
\Omega_{\mathcal{Q}_\bullet/\mathcal{A}}
\otimes_{\mathcal{Q}_\bullet} \mathcal{C} \to
\Omega_{\overline{\mathcal{Q}}_\bullet/\mathcal{B}}
\otimes_{\overline{\mathcal{Q}}_\bullet} \mathcal{C} \to 0
$$
Les deux premiers termes sont égaux aux deux premiers termes du triangle de
l'énoncé. L'identification du dernier terme avec
$L_{\mathcal{C}/\mathcal{B}}$ utilise les quasi-isomorphismes de complexes
$$
L_{\mathcal{C}/\mathcal{B}} =
\Omega_{\mathcal{R}_\bullet/\mathcal{B}}
\otimes_{\mathcal{R}_\bullet} \mathcal{C}
\longrightarrow
\Omega_{\overline{\mathcal{R}}_\bullet/\mathcal{B}}
\otimes_{\overline{\mathcal{R}}_\bullet} \mathcal{C}
\longleftarrow
\Omega_{\overline{\mathcal{Q}}_\bullet/\mathcal{B}}
\otimes_{\overline{\mathcal{Q}}_\bullet} \mathcal{C}
$$
Toutes les constructions utilisées ci-dessus peuvent d'abord être effectuées
au niveau des préfaisceaux, puis faisceautisées. Ainsi, pour démontrer que des
suites sont exactes ou que des applications sont des quasi-isomorphismes, il
suffit de démontrer l'énoncé correspondant pour les homomorphismes d'anneaux
$\mathcal{A}(U) \to \mathcal{B}(U) \to \mathcal{C}(U)$
pour lesquels le résultat est connu. Ceci achève la démonstration lorsque
$\mathcal{B} \to \mathcal{C}$ est injectif.

\medskip\noindent
Dans le cas général, nous nous ramenons au cas où
$\mathcal{B} \to \mathcal{C}$ est injectif en remplaçant, si nécessaire,
$\mathcal{C}$ par $\mathcal{B} \times \mathcal{C}$. Cela est possible grâce à
l'argument de la remarque \ref{cotangent-remark-triangle} et au lemme
\ref{cotangent-lemma-compute-L-product-sheaves-rings}.
\end{proof}

\begin{lemma}
\label{cotangent-lemma-stalk-cotangent-complex}
Soit $\mathcal{C}$ un site. Soit $\mathcal{A} \to \mathcal{B}$ un
homomorphisme de faisceaux d'anneaux sur $\mathcal{C}$. Si $p$ est un point de
$\mathcal{C}$, alors
$(L_{\mathcal{B}/\mathcal{A}})_p = L_{\mathcal{B}_p/\mathcal{A}_p}$.
\end{lemma}

\begin{proof}
C'est un cas particulier du lemme
\ref{cotangent-lemma-pullback-cotangent-morphism-topoi}.
\end{proof}

\noindent
Pour la construction du complexe cotangent naïf et ses propriétés, nous
renvoyons à Modules sur les sites, section
\ref{sites-modules-section-netherlander}.

\begin{lemma}
\label{cotangent-lemma-compare-cotangent-complex-with-naive}
Soit $\mathcal{C}$ un site. Soit $\mathcal{A} \to \mathcal{B}$ un
homomorphisme de faisceaux d'anneaux sur $\mathcal{C}$. Il existe une
application canonique
$L_{\mathcal{B}/\mathcal{A}} \to \NL_{\mathcal{B}/\mathcal{A}}$ qui identifie
le complexe cotangent naïf à la troncation
$\tau_{\geq -1}L_{\mathcal{B}/\mathcal{A}}$.
\end{lemma}

\begin{proof}
Soit $\mathcal{P}_\bullet$ la résolution standard de $\mathcal{B}$ sur
$\mathcal{A}$. Posons
$\mathcal{I} = \Ker(\mathcal{A}[\mathcal{B}] \to \mathcal{B})$. Rappelons que
$\mathcal{P}_0 = \mathcal{A}[\mathcal{B}]$. L'application du lemme est donnée
par le diagramme commutatif
$$
\xymatrix{
L_{\mathcal{B}/\mathcal{A}} \ar[d] & \ldots \ar[r] &
\Omega_{\mathcal{P}_2/\mathcal{A}} \otimes_{\mathcal{P}_2} \mathcal{B}
\ar[r] \ar[d] &
\Omega_{\mathcal{P}_1/\mathcal{A}} \otimes_{\mathcal{P}_1} \mathcal{B}
\ar[r] \ar[d] &
\Omega_{\mathcal{P}_0/\mathcal{A}} \otimes_{\mathcal{P}_0} \mathcal{B}
\ar[d] \\
\NL_{\mathcal{B}/\mathcal{A}} & \ldots \ar[r] &
0 \ar[r] & 
\mathcal{I}/\mathcal{I}^2 \ar[r] &
\Omega_{\mathcal{P}_0/\mathcal{A}} \otimes_{\mathcal{P}_0} \mathcal{B}
}
$$
Nous construisons la flèche descendante de but
$\mathcal{I}/\mathcal{I}^2$ en envoyant une section locale
$\text{d}f \otimes b$ sur la classe de $(d_0(f) - d_1(f))b$ dans
$\mathcal{I}/\mathcal{I}^2$. Ici, $d_i : \mathcal{P}_1 \to \mathcal{P}_0$,
$i = 0, 1$, sont les deux applications de face de la structure simpliciale.
Cela a un sens puisque $d_0 - d_1$ envoie $\mathcal{P}_1$ dans
$\mathcal{I} = \Ker(\mathcal{P}_0 \to \mathcal{B})$. Nous omettons de vérifier
que cette règle est bien définie. Notre application est compatible avec la
différentielle
$\Omega_{\mathcal{P}_1/\mathcal{A}} \otimes_{\mathcal{P}_1} \mathcal{B}
\to \Omega_{\mathcal{P}_0/\mathcal{A}} \otimes_{\mathcal{P}_0} \mathcal{B}$
car celle-ci envoie une section locale $\text{d}f \otimes b$ sur
$\text{d}(d_0(f) - d_1(f)) \otimes b$. En outre, la différentielle
$\Omega_{\mathcal{P}_2/\mathcal{A}} \otimes_{\mathcal{P}_2} \mathcal{B}
\to \Omega_{\mathcal{P}_1/\mathcal{A}} \otimes_{\mathcal{P}_1} \mathcal{B}$
envoie une section locale $\text{d}f \otimes b$ sur
$\text{d}(d_0(f) - d_1(f) + d_2(f)) \otimes b$, qui est annulé par notre
flèche descendante. Nous obtenons ainsi une application de complexes.

\medskip\noindent
Pour voir que notre application induit un isomorphisme sur les faisceaux de
cohomologie $H^0$ et $H^{-1}$, raisonnons comme suit. Soit $\mathcal{C}'$ le
site ayant même catégorie sous-jacente que $\mathcal{C}$, mais muni de la
topologie chaotique. Soit $f : \Sh(\mathcal{C}) \to \Sh(\mathcal{C}')$ le
morphisme de topos dont le foncteur image inverse est la faisceautisation.
Considérons l'homomorphisme donné $\mathcal{A}' \to \mathcal{B}'$ comme un
homomorphisme de faisceaux d'anneaux sur $\mathcal{C}'$. La construction
précédente donne sur $\mathcal{C}'$ une application
$L_{\mathcal{B}'/\mathcal{A}'} \to \NL_{\mathcal{B}'/\mathcal{A}'}$ dont la
valeur sur tout objet $U$ de $\mathcal{C}'$ est simplement l'application
$$
L_{\mathcal{B}(U)/\mathcal{A}(U)} \to \NL_{\mathcal{B}(U)/\mathcal{A}(U)}
$$
de la remarque \ref{cotangent-remark-explicit-comparison-map}, qui induit un
isomorphisme sur $H^0$ et $H^{-1}$. Puisque
$f^{-1}L_{\mathcal{B}'/\mathcal{A}'} = L_{\mathcal{B}/\mathcal{A}}$
(lemme \ref{cotangent-lemma-pullback-cotangent-morphism-topoi}) et
$f^{-1}\NL_{\mathcal{B}'/\mathcal{A}'} = \NL_{\mathcal{B}/\mathcal{A}}$
(Modules sur les sites, lemme \ref{sites-modules-lemma-pullback-NL}), le lemme
est démontré.
\end{proof}











\section{La classe d'Atiyah d'un faisceau de modules}
\label{cotangent-section-atiyah-general}

\noindent
Soit $\mathcal{C}$ un site. Soit $\mathcal{A} \to \mathcal{B}$ un
homomorphisme de faisceaux d'anneaux. Soit $\mathcal{F}$ un faisceau de
$\mathcal{B}$-modules. Soit $\mathcal{P}_\bullet \to \mathcal{B}$ la
résolution standard de $\mathcal{B}$ sur $\mathcal{A}$ (section
\ref{cotangent-section-cotangent-complex}). Pour tout $n \geq 0$, considérons l'extension
des parties principales
\begin{equation}
\label{cotangent-equation-atiyah-extension}
0 \to
\Omega_{\mathcal{P}_n/\mathcal{A}} \otimes_{\mathcal{P}_n} \mathcal{F} \to
\mathcal{P}^1_{\mathcal{P}_n/\mathcal{A}}(\mathcal{F}) \to
\mathcal{F} \to 0
\end{equation}
voir Modules sur les sites, lemme
\ref{sites-modules-lemma-sequence-of-principal-parts}. La fonctorialité de
cette construction (Modules sur les sites, remarque
\ref{sites-modules-remark-functoriality-principal-parts}) montre que
(\ref{cotangent-equation-atiyah-extension}) est la partie de degré $n$ d'une suite
exacte courte de $\mathcal{P}_\bullet$-modules simpliciaux (Cohomologie sur les
sites, section \ref{sites-cohomology-section-simplicial-modules}). En utilisant
le foncteur $L\pi_! : D(\mathcal{P}_\bullet) \to D(\mathcal{B})$ de Cohomologie
sur les sites, remarque \ref{sites-cohomology-remark-homology-augmentation}
(nous utilisons ici le fait que $\mathcal{P}_\bullet \to \mathcal{A}$ est une
résolution), nous obtenons un triangle distingué
\begin{equation}
\label{cotangent-equation-atiyah-general}
L_{\mathcal{B}/\mathcal{A}} \otimes_\mathcal{B}^\mathbf{L} \mathcal{F} \to
L\pi_!\left(\mathcal{P}^1_{\mathcal{P}_\bullet/\mathcal{A}}(\mathcal{F})\right)
\to \mathcal{F} \to
L_{\mathcal{B}/\mathcal{A}} \otimes_\mathcal{B}^\mathbf{L} \mathcal{F} [1]
\end{equation}
dans $D(\mathcal{B})$.

\begin{definition}
\label{cotangent-definition-atiyah-class-general}
Soit $\mathcal{C}$ un site. Soit $\mathcal{A} \to \mathcal{B}$ un
homomorphisme de faisceaux d'anneaux. Soit $\mathcal{F}$ un faisceau de
$\mathcal{B}$-modules. L'application $\mathcal{F} \to
L_{\mathcal{B}/\mathcal{A}} \otimes_\mathcal{B}^\mathbf{L} \mathcal{F}[1]$
de (\ref{cotangent-equation-atiyah-general}) est appelée la {\it classe d'Atiyah} de
$\mathcal{F}$.
\end{definition}









\section{Le complexe cotangent d'un morphisme d'espaces annelés}
\label{cotangent-section-cotangent-morphism-ringed-spaces}

\noindent
Le complexe cotangent d'un morphisme d'espaces annelés se définit au moyen du
complexe cotangent défini ci-dessus.

\begin{definition}
\label{cotangent-definition-cotangent-complex-morphism-ringed-spaces}
Soit $f : (X, \mathcal{O}_X) \to (S, \mathcal{O}_S)$ un morphisme d'espaces
annelés. Le {\it complexe cotangent} $L_f$ de $f$ est
$L_f = L_{\mathcal{O}_X/f^{-1}\mathcal{O}_S}$.
Nous utiliserons également la notation
$L_f = L_{X/S} = L_{\mathcal{O}_X/\mathcal{O}_S}$.
\end{definition}

\noindent
Plus précisément, nous considérons le complexe cotangent (définition
\ref{cotangent-definition-cotangent-complex-morphism-sheaves-rings}) de l'homomorphisme
$f^\sharp : f^{-1}\mathcal{O}_S \to \mathcal{O}_X$ de faisceaux d'anneaux sur
le site associé à l'espace topologique $X$ (Sites, exemple
\ref{sites-example-site-topological}).

\begin{lemma}
\label{cotangent-lemma-H0-L-morphism-ringed-spaces}
Soit $f : (X, \mathcal{O}_X) \to (S, \mathcal{O}_S)$ un morphisme d'espaces
annelés. Alors $H^0(L_{X/S}) = \Omega_{X/S}$.
\end{lemma}

\begin{proof}
C'est un cas particulier du lemme \ref{cotangent-lemma-H0-L-morphism-sheaves-rings}.
\end{proof}

\begin{lemma}
\label{cotangent-lemma-triangle-ringed-spaces}
Soient $f : X \to Y$ et $g : Y \to Z$ des morphismes d'espaces annelés. Il
existe alors un triangle distingué canonique
$$
Lf^* L_{Y/Z} \to L_{X/Z} \to L_{X/Y} \to Lf^*L_{Y/Z}[1]
$$
dans $D(\mathcal{O}_X)$.
\end{lemma}

\begin{proof}
Posons $h = g \circ f$, de sorte que
$h^{-1}\mathcal{O}_Z = f^{-1}g^{-1}\mathcal{O}_Z$. D'après le lemme
\ref{cotangent-lemma-pullback-cotangent-morphism-topoi}, nous avons
$f^{-1}L_{Y/Z} = L_{f^{-1}\mathcal{O}_Y/h^{-1}\mathcal{O}_Z}$
et celui-ci est un complexe de $f^{-1}\mathcal{O}_Y$-modules plats. Le triangle
distingué précédent est donc un cas du triangle distingué du lemme
\ref{cotangent-lemma-triangle-sheaves-rings}, avec
$\mathcal{A} = h^{-1}\mathcal{O}_Z$, $\mathcal{B} = f^{-1}\mathcal{O}_Y$ et
$\mathcal{C} = \mathcal{O}_X$.
\end{proof}

\begin{lemma}
\label{cotangent-lemma-compare-cotangent-complex-with-naive-ringed-spaces}
Soit $f : (X, \mathcal{O}_X) \to (Y, \mathcal{O}_Y)$ un morphisme d'espaces
annelés. Il existe une application canonique $L_{X/Y} \to \NL_{X/Y}$ qui
identifie le complexe cotangent naïf à la troncation
$\tau_{\geq -1}L_{X/Y}$.
\end{lemma}

\begin{proof}
C'est un cas particulier du lemme
\ref{cotangent-lemma-compare-cotangent-complex-with-naive}.
\end{proof}






\section{Déformations d'espaces annelés et complexe cotangent}
\label{cotangent-section-deformations-ringed-spaces}

\noindent
Cette section prolonge Théorie des déformations, section
\ref{defos-section-deformations-ringed-spaces}, que nous invitons vivement le
lecteur à lire d'abord. Rappelons brièvement la situation. Nous avons un
épaississement du premier ordre
$t : (S, \mathcal{O}_S) \to (S', \mathcal{O}_{S'})$ d'espaces annelés, avec
$\mathcal{J} = \Ker(t^\sharp)$, un morphisme d'espaces annelés
$f : (X, \mathcal{O}_X) \to (S, \mathcal{O}_S)$, un $\mathcal{O}_X$-module
$\mathcal{G}$ et un $f$-morphisme $c : \mathcal{J} \to \mathcal{G}$ de
faisceaux de modules. Nous cherchons si l'on peut remplacer le point
d'interrogation dans le diagramme suivant
\begin{equation}
\label{cotangent-equation-to-solve-ringed-spaces}
\vcenter{
\xymatrix{
0 \ar[r] & \mathcal{G} \ar[r] & {?} \ar[r] & \mathcal{O}_X \ar[r] & 0 \\
0 \ar[r] & \mathcal{J} \ar[u]^c \ar[r] & \mathcal{O}_{S'} \ar[u] \ar[r] &
\mathcal{O}_S \ar[u] \ar[r] & 0
}
}
\end{equation}
et, si une solution existe, à quel point elle est unique. Plus précisément,
nous cherchons un épaississement du premier ordre
$i : (X, \mathcal{O}_X) \to (X', \mathcal{O}_{X'})$
et un morphisme d'épaississements $(f, f')$ comme dans Théorie des
déformations, équation (\ref{defos-equation-morphism-thickenings}), où
$\Ker(i^\sharp)$ est identifié à $\mathcal{G}$, tels que $(f')^\sharp$ induise
l'application donnée $c$. Nous dirons que $X'$ est une {\it solution} de
(\ref{cotangent-equation-to-solve-ringed-spaces}).

\begin{lemma}
\label{cotangent-lemma-find-obstruction-ringed-spaces}
Dans la situation précédente :
\begin{enumerate}
\item il existe un élément canonique
$\xi \in \Ext^2_{\mathcal{O}_X}(L_{X/S}, \mathcal{G})$
dont l'annulation est une condition nécessaire et suffisante pour l'existence
d'une solution de (\ref{cotangent-equation-to-solve-ringed-spaces}) ;
\item s'il existe une solution, l'ensemble des classes d'isomorphisme de
solutions est un espace principal homogène sous
$\Ext^1_{\mathcal{O}_X}(L_{X/S}, \mathcal{G})$ ;
\item étant donnée une solution $X'$, l'ensemble des automorphismes de $X'$
qui s'insèrent dans (\ref{cotangent-equation-to-solve-ringed-spaces}) est canoniquement
isomorphe à $\Ext^0_{\mathcal{O}_X}(L_{X/S}, \mathcal{G})$.
\end{enumerate}
\end{lemma}

\begin{proof}
Au moyen des identifications $\NL_{X/S} = \tau_{\geq -1}L_{X/S}$ (lemme
\ref{cotangent-lemma-compare-cotangent-complex-with-naive-ringed-spaces}) et
$H^0(L_{X/S}) = \Omega_{X/S}$
(lemme \ref{cotangent-lemma-H0-L-morphism-ringed-spaces}), nous avons déjà établi (2) et
(3) dans Théorie des déformations, lemmes
\ref{defos-lemma-huge-diagram-ringed-spaces} et
\ref{defos-lemma-choices-ringed-spaces}.

\medskip\noindent
Démonstration de (1). En gros, l'assertion découle de la discussion de Théorie
des déformations, remarque
\ref{defos-remark-parametrize-solutions-ringed-spaces}, en remplaçant le
complexe cotangent naïf par le complexe cotangent entier. Voici une explication
plus détaillée. D'après Théorie des déformations, lemme
\ref{defos-lemma-parametrize-solutions-ringed-spaces}, il existe un élément
$$
\xi' \in
\Ext^1_{\mathcal{O}_X}(Lf^*\NL_{S/S'}, \mathcal{G}) =
\Ext^1_{\mathcal{O}_X}(Lf^*L_{S/S'}, \mathcal{G})
$$
tel qu'une solution existe si et seulement si cet élément appartient à l'image
de l'application
$$
\Ext^1_{\mathcal{O}_X}(NL_{X/S'}, \mathcal{G}) =
\Ext^1_{\mathcal{O}_X}(L_{X/S'}, \mathcal{G})
\longrightarrow
\Ext^1_{\mathcal{O}_X}(Lf^*L_{S/S'}, \mathcal{G})
$$
Le triangle distingué du lemme \ref{cotangent-lemma-triangle-ringed-spaces}, appliqué à
$X \to S \to S'$, donne naissance à une suite exacte longue
$$
\ldots \to
\Ext^1_{\mathcal{O}_X}(L_{X/S'}, \mathcal{G}) \to
\Ext^1_{\mathcal{O}_X}(Lf^*L_{S/S'}, \mathcal{G}) \to
\Ext^2_{\mathcal{O}_X}(L_{X/S}, \mathcal{G}) \to \ldots
$$
Il suffit donc de prendre pour $\xi$ l'image de $\xi'$.
\end{proof}










\section{Le complexe cotangent d'un morphisme de topos annelés}
\label{cotangent-section-cotangent-morphism-ringed-topoi}

\noindent
Le complexe cotangent d'un morphisme de topos annelés se définit au moyen du
complexe cotangent défini ci-dessus.

\begin{definition}
\label{cotangent-definition-cotangent-complex-morphism-ringed-topoi}
Soit $(f, f^\sharp) : (\Sh(\mathcal{C}), \mathcal{O}_\mathcal{C}) \to
(\Sh(\mathcal{D}), \mathcal{O}_\mathcal{D})$ un morphisme de topos annelés.
Le {\it complexe cotangent} $L_f$ de $f$ est
$L_f = L_{\mathcal{O}_\mathcal{C}/f^{-1}\mathcal{O}_\mathcal{D}}$.
Nous écrivons parfois
$L_f = L_{\mathcal{O}_\mathcal{C}/\mathcal{O}_\mathcal{D}}$.
\end{definition}

\noindent
Cette définition s'applique dans de nombreuses situations, mais elle ne donne
pas toujours l'objet attendu. Par exemple, si $f : X \to Y$ est un morphisme
de schémas, alors $f$ induit un morphisme de grands sites étales
$f_{big} : (\Sch/X)_\etale \to (\Sch/Y)_\etale$
qui est un morphisme de topos annelés (Descente, remarque
\ref{descent-remark-change-topologies-ringed}). Toutefois,
$L_{f_{big}} = 0$ puisque $(f_{big})^\sharp$ est un isomorphisme. En revanche,
si nous prenons $L_f$ en regardant $f$ comme un morphisme entre les topos
annelés de Zariski sous-jacents, alors $L_f$ coïncide bien avec le complexe
cotangent $L_{X/Y}$ (défini ci-dessous), dont le faisceau de cohomologie de
degré zéro est $\Omega_{X/Y}$.

\begin{lemma}
\label{cotangent-lemma-H0-L-morphism-ringed-topoi}
Soit $f : (\Sh(\mathcal{C}), \mathcal{O}) \to
(\Sh(\mathcal{B}), \mathcal{O}_\mathcal{B})$ un morphisme de topos annelés.
Alors $H^0(L_f) = \Omega_f$.
\end{lemma}

\begin{proof}
C'est un cas particulier du lemme \ref{cotangent-lemma-H0-L-morphism-sheaves-rings}.
\end{proof}

\begin{lemma}
\label{cotangent-lemma-triangle-ringed-topoi}
Soient $f : (\Sh(\mathcal{C}_1), \mathcal{O}_1) \to
(\Sh(\mathcal{C}_2), \mathcal{O}_2)$ et
$g : (\Sh(\mathcal{C}_2), \mathcal{O}_2) \to
(\Sh(\mathcal{C}_3), \mathcal{O}_3)$ des morphismes de topos annelés. Il
existe alors un triangle distingué canonique
$$
Lf^* L_g \to L_{g \circ f} \to L_f \to Lf^*L_g[1]
$$
dans $D(\mathcal{O}_1)$.
\end{lemma}

\begin{proof}
Posons $h = g \circ f$, de sorte que
$h^{-1}\mathcal{O}_3 = f^{-1}g^{-1}\mathcal{O}_3$. D'après le lemme
\ref{cotangent-lemma-pullback-cotangent-morphism-topoi}, nous avons
$f^{-1}L_g = L_{f^{-1}\mathcal{O}_2/h^{-1}\mathcal{O}_3}$
et celui-ci est un complexe de $f^{-1}\mathcal{O}_2$-modules plats. Le triangle
distingué précédent est donc un cas du triangle distingué du lemme
\ref{cotangent-lemma-triangle-sheaves-rings}, avec
$\mathcal{A} = h^{-1}\mathcal{O}_3$, $\mathcal{B} = f^{-1}\mathcal{O}_2$ et
$\mathcal{C} = \mathcal{O}_1$.
\end{proof}

\begin{lemma}
\label{cotangent-lemma-compare-cotangent-complex-with-naive-ringed-topoi}
Soit $f : (\Sh(\mathcal{C}), \mathcal{O}) \to
(\Sh(\mathcal{B}), \mathcal{O}_\mathcal{B})$ un morphisme de topos annelés. Il
existe une application canonique $L_f \to \NL_f$ qui identifie le complexe
cotangent naïf à la troncation $\tau_{\geq -1}L_f$.
\end{lemma}

\begin{proof}
C'est un cas particulier du lemme
\ref{cotangent-lemma-compare-cotangent-complex-with-naive}.
\end{proof}








\section{Déformations de topos annelés et complexe cotangent}
\label{cotangent-section-deformations-ringed-topoi}

\noindent
Cette section prolonge Théorie des déformations, section
\ref{defos-section-deformations-ringed-topoi}, que nous invitons vivement le
lecteur à lire d'abord. Rappelons brièvement la situation. Nous avons un
épaississement du premier ordre
$t : (\Sh(\mathcal{B}), \mathcal{O}_\mathcal{B}) \to
(\Sh(\mathcal{B}'), \mathcal{O}_{\mathcal{B}'})$ de topos annelés, avec
$\mathcal{J} = \Ker(t^\sharp)$, un morphisme de topos annelés
$f : (\Sh(\mathcal{C}), \mathcal{O}) \to
(\Sh(\mathcal{B}), \mathcal{O}_\mathcal{B})$, an $\mathcal{O}$-module
$\mathcal{G}$ et une application $f^{-1}\mathcal{J} \to \mathcal{G}$ de
faisceaux de $f^{-1}\mathcal{O}_\mathcal{B}$-modules. Nous cherchons si l'on
peut remplacer le point d'interrogation dans le diagramme suivant
\begin{equation}
\label{cotangent-equation-to-solve-ringed-topoi}
\vcenter{
\xymatrix{
0 \ar[r] & \mathcal{G} \ar[r] & {?} \ar[r] & \mathcal{O} \ar[r] & 0 \\
0 \ar[r] & f^{-1}\mathcal{J} \ar[u]^c \ar[r] &
f^{-1}\mathcal{O}_{\mathcal{B}'} \ar[u] \ar[r] &
f^{-1}\mathcal{O}_\mathcal{B} \ar[u] \ar[r] & 0
}
}
\end{equation}
et, si une solution existe, à quel point elle est unique. Plus précisément,
nous cherchons un épaississement du premier ordre
$i : (\Sh(\mathcal{C}), \mathcal{O}) \to (\Sh(\mathcal{C}'), \mathcal{O}')$
et un morphisme d'épaississements $(f, f')$ comme dans Théorie des
déformations, équation
(\ref{defos-equation-morphism-thickenings-ringed-topoi}), où
$\Ker(i^\sharp)$ est identifié à $\mathcal{G}$, tels que $(f')^\sharp$ induise
l'application donnée $c$. Nous dirons que
$(\Sh(\mathcal{C}'), \mathcal{O}')$ est une {\it solution} de
(\ref{cotangent-equation-to-solve-ringed-topoi}).

\begin{lemma}
\label{cotangent-lemma-find-obstruction-ringed-topoi}
Dans la situation précédente :
\begin{enumerate}
\item il existe un élément canonique
$\xi \in \Ext^2_\mathcal{O}(L_f, \mathcal{G})$
dont l'annulation est une condition nécessaire et suffisante pour l'existence
d'une solution de (\ref{cotangent-equation-to-solve-ringed-topoi}) ;
\item s'il existe une solution, l'ensemble des classes d'isomorphisme de
solutions est un espace principal homogène sous
$\Ext^1_\mathcal{O}(L_f, \mathcal{G})$ ;
\item étant donnée une solution $X'$, l'ensemble des automorphismes de $X'$
qui s'insèrent dans (\ref{cotangent-equation-to-solve-ringed-topoi}) est canoniquement
isomorphe à $\Ext^0_\mathcal{O}(L_f, \mathcal{G})$.
\end{enumerate}
\end{lemma}

\begin{proof}
Au moyen des identifications $\NL_f = \tau_{\geq -1}L_f$ (lemme
\ref{cotangent-lemma-compare-cotangent-complex-with-naive-ringed-topoi}) et
$H^0(L_f) = \Omega_f$
(lemme \ref{cotangent-lemma-H0-L-morphism-ringed-topoi}), nous avons déjà établi (2) et
(3) dans Théorie des déformations, lemmes
\ref{defos-lemma-huge-diagram-ringed-topoi} et
\ref{defos-lemma-choices-ringed-topoi}.

\medskip\noindent
Démonstration de (1). Pour accorder les notations à celles de Théorie des
déformations, section \ref{defos-section-deformations-ringed-topoi}, nous
écrirons
$\NL_f = \NL_{\mathcal{O}/\mathcal{O}_\mathcal{B}}$ et
$L_f = L_{\mathcal{O}/\mathcal{O}_\mathcal{B}}$, et de même pour les morphismes
$t$ et $t \circ f$. D'après Théorie des déformations, lemme
\ref{defos-lemma-parametrize-solutions-ringed-topoi}, il existe un élément
$$
\xi' \in
\Ext^1_\mathcal{O}(
Lf^*\NL_{\mathcal{O}_\mathcal{B}/\mathcal{O}_{\mathcal{B}'}}, \mathcal{G}) =
\Ext^1_\mathcal{O}(
Lf^*L_{\mathcal{O}_\mathcal{B}/\mathcal{O}_{\mathcal{B}'}}, \mathcal{G})
$$
tel qu'une solution existe si et seulement si cet élément appartient à l'image
de l'application
$$
\Ext^1_\mathcal{O}(
\NL_{\mathcal{O}/\mathcal{O}_{\mathcal{B}'}}, \mathcal{G}) =
\Ext^1_\mathcal{O}(
L_{\mathcal{O}/\mathcal{O}_{\mathcal{B}'}}, \mathcal{G})
\longrightarrow
\Ext^1_\mathcal{O}(
Lf^*L_{\mathcal{O}_\mathcal{B}/\mathcal{O}_{\mathcal{B}'}}, \mathcal{G})
$$
Le triangle distingué du lemme \ref{cotangent-lemma-triangle-ringed-topoi}, appliqué à
$f$ et $t$, donne naissance à une suite exacte longue
$$
\ldots \to
\Ext^1_\mathcal{O}(
L_{\mathcal{O}/\mathcal{O}_{\mathcal{B}'}}, \mathcal{G}) \to
\Ext^1_\mathcal{O}(
Lf^*L_{\mathcal{O}_\mathcal{B}/\mathcal{O}_{\mathcal{B}'}}, \mathcal{G})
\to
\Ext^1_\mathcal{O}(
L_{\mathcal{O}/\mathcal{O}_\mathcal{B}}, \mathcal{G})
$$
Il suffit donc de prendre pour $\xi$ l'image de $\xi'$.
\end{proof}












\section{Le complexe cotangent d'un morphisme de schémas}
\label{cotangent-section-cotangent-morphism-schemes}

\noindent
Comme annoncé plus haut, nous définissons comme suit le complexe cotangent d'un
morphisme de schémas.

\begin{definition}
\label{cotangent-definition-cotangent-morphism-schemes}
Soit $f : X \to Y$ un morphisme de schémas. Le {\it complexe cotangent
$L_{X/Y}$ de $X$ sur $Y$} est le complexe cotangent de $f$ considéré comme un
morphisme d'espaces annelés (définition
\ref{cotangent-definition-cotangent-complex-morphism-ringed-spaces}).
\end{definition}

\noindent
En particulier, les résultats de la section
\ref{cotangent-section-cotangent-morphism-ringed-spaces} s'appliquent aux complexes
cotangents des morphismes de schémas. Le lemme suivant montre que cette
définition est compatible avec celle donnée pour les homomorphismes d'anneaux ;
il implique aussi que $L_{X/Y}$ est un objet de $D_\QCoh(\mathcal{O}_X)$.

\begin{lemma}
\label{cotangent-lemma-morphism-affine-schemes}
Soit $f : X \to Y$ un morphisme de schémas. Soient
$U = \Spec(B) \subset X$ et $V = \Spec(A) \subset Y$ des ouverts affines tels
que $f(U) \subset V$. Il existe une application canonique
$$
\widetilde{L_{B/A}} \longrightarrow L_{X/Y}|_U
$$
de complexes qui est un isomorphisme dans $D(\mathcal{O}_U)$. Cette application
est compatible à la restriction à des ouverts affines plus petits de $X$ et
de $Y$.
\end{lemma}

\begin{proof}
D'après la remarque \ref{cotangent-remark-map-sections-over-U}, il existe une
application canonique de complexes
$L_{\mathcal{O}_X(U)/f^{-1}\mathcal{O}_Y(U)} \to L_{X/Y}(U)$
de $B = \mathcal{O}_X(U)$-modules, compatible aux restrictions ultérieures. En
utilisant l'application canonique $A \to f^{-1}\mathcal{O}_Y(U)$, nous
obtenons une application canonique
$L_{B/A} \to L_{\mathcal{O}_X(U)/f^{-1}\mathcal{O}_Y(U)}$
de complexes de $B$-modules. La propriété universelle du foncteur
$\widetilde{\ }$ (voir Schémas, lemme
\ref{schemes-lemma-compare-constructions}) fournit l'application de l'énoncé.
Pour vérifier qu'elle est un isomorphisme sur les faisceaux de cohomologie, il
suffit de vérifier qu'elle induit des isomorphismes sur les fibres. Cela
résulte immédiatement des lemmes \ref{cotangent-lemma-stalk-cotangent-complex} et
\ref{cotangent-lemma-localize} (et de la description des fibres de
$\mathcal{O}_X$ et $f^{-1}\mathcal{O}_Y$
en un point $\mathfrak p \in \Spec(B)$ comme $B_\mathfrak p$ et
$A_\mathfrak q$, où $\mathfrak q = A \cap \mathfrak p$ ; les références
utilisées sont Schémas, lemme \ref{schemes-lemma-spec-sheaves}, et Faisceaux,
lemme \ref{sheaves-lemma-stalk-pullback}).
\end{proof}

\begin{lemma}
\label{cotangent-lemma-scheme-over-ring}
Soit $\Lambda$ un anneau. Soit $X$ un schéma sur $\Lambda$. Alors
$$
L_{X/\Spec(\Lambda)} = L_{\mathcal{O}_X/\underline{\Lambda}}
$$
où $\underline{\Lambda}$ est le faisceau constant de valeur $\Lambda$ sur $X$.
\end{lemma}

\begin{proof}
Soit $p : X \to \Spec(\Lambda)$ le morphisme structural. Soit
$q : \Spec(\Lambda) \to (*, \Lambda)$ le morphisme évident. D'après le triangle
distingué du lemme \ref{cotangent-lemma-triangle-ringed-spaces}, il suffit de montrer que
$L_q = 0$. Pour cela, il suffit de montrer, pour
$\mathfrak p \in \Spec(\Lambda)$, que
$$
(L_q)_\mathfrak p =
L_{\mathcal{O}_{\Spec(\Lambda), \mathfrak p}/\Lambda} =
L_{\Lambda_\mathfrak p/\Lambda}
$$
(lemme \ref{cotangent-lemma-stalk-cotangent-complex}) est nul, ce qui résulte du lemme
\ref{cotangent-lemma-when-zero}.
\end{proof}






\section{Le complexe cotangent d'un schéma sur un anneau}
\label{cotangent-section-cotangent-schemes-variant}

\noindent
Soit $\Lambda$ un anneau et soit $X$ un schéma sur $\Lambda$. Écrivons
$L_{X/\Spec(\Lambda)} = L_{X/\Lambda}$, ce qui est justifié par le lemme
\ref{cotangent-lemma-scheme-over-ring}. Dans cette section, nous donnons de
$L_{X/\Lambda}$ une description analogue à celle du lemme
\ref{cotangent-lemma-compute-cotangent-complex}. Plus précisément, nous construisons une
catégorie $\mathcal{C}_{X/\Lambda}$ fibrée au-dessus de $X_{Zar}$ et la
munissons d'un faisceau $\mathcal{O}$ de $\Lambda$-algèbres (polynomiales) tel
que
$$
L_{X/\Lambda} =
L\pi_!(\Omega_{\mathcal{O}/\underline{\Lambda}} \otimes_\mathcal{O}
\underline{\mathcal{O}}_X).
$$
Nous utiliserons plus loin la catégorie $\mathcal{C}_{X/\Lambda}$ pour
construire une théorie naïve des obstructions pour le champ des faisceaux
cohérents.

\medskip\noindent
Soit $\Lambda$ un anneau. Soit $X$ un schéma sur $\Lambda$. Soit
$\mathcal{C}_{X/\Lambda}$ la catégorie dont les objets sont les diagrammes
commutatifs
\begin{equation}
\label{cotangent-equation-object}
\vcenter{
\xymatrix{
X \ar[d] & U \ar[l] \ar[d] \\
\Spec(\Lambda) & \mathbf{A} \ar[l]
}
}
\end{equation}
de schémas tels que
\begin{enumerate}
\item $U$ soit un sous-schéma ouvert de $X$,
\item il existe un isomorphisme $\mathbf{A} = \Spec(P)$, où $P$ est une
algèbre polynomiale sur $\Lambda$ (en un certain ensemble de variables).
\end{enumerate}
Autrement dit, $\mathbf{A}$ est un espace affine (de dimension infinie) sur
$\Spec(\Lambda)$. Les morphismes sont donnés par des diagrammes commutatifs.
Rappelons que $X_{Zar}$ désigne le petit site de Zariski de $X$. Il existe un
foncteur d'oubli
$$
u : \mathcal{C}_{X/\Lambda} \to X_{Zar},\ (U \to \mathbf{A}) \mapsto U
$$
Remarquons que la catégorie fibre au-dessus de $U$ est canoniquement
équivalente à la catégorie $\mathcal{C}_{\mathcal{O}_X(U)/\Lambda}$ introduite
dans la section \ref{cotangent-section-compute-L-pi-shriek}.

\begin{lemma}
\label{cotangent-lemma-category-fibred}
Dans la situation précédente, la catégorie $\mathcal{C}_{X/\Lambda}$ est
fibrée au-dessus de $X_{Zar}$.
\end{lemma}

\begin{proof}
Étant donnés un objet $U \to \mathbf{A}$ de $\mathcal{C}_{X/\Lambda}$ et un
morphisme $U' \to U$ de $X_{Zar}$, considérons l'objet
$U' \to \mathbf{A}$ de $\mathcal{C}_{X/\Lambda}$, où
$U' \to \mathbf{A}$ est la composée de $U' \to U$ et de
$U \to \mathbf{A}$. Le morphisme
$(U' \to \mathbf{A}) \to (U \to \mathbf{A})$ of
de $\mathcal{C}_{X/\Lambda}$ est fortement cartésien au-dessus de $X_{Zar}$.
\end{proof}

\noindent
Munissons $\mathcal{C}_{X/\Lambda}$ de la topologie héritée de $X_{Zar}$ (voir
Champs, section \ref{stacks-section-topology}). Le foncteur $u$ définit un
morphisme de topos
$\pi : \Sh(\mathcal{C}_{X/\Lambda}) \to \Sh(X_{Zar})$. Le site
$\mathcal{C}_{X/\Lambda}$ est muni de plusieurs faisceaux d'anneaux :
\begin{enumerate}
\item le faisceau $\mathcal{O}$ donné par la règle
$(U \to \mathbf{A}) \mapsto \Gamma(\mathbf{A}, \mathcal{O}_\mathbf{A})$.
\item le faisceau $\underline{\mathcal{O}}_X = \pi^{-1}\mathcal{O}_X$ donné
par la règle $(U \to \mathbf{A}) \mapsto \mathcal{O}_X(U)$ ;
\item le faisceau constant $\underline{\Lambda}$.
\end{enumerate}
Nous obtenons des morphismes de topos annelés
\begin{equation}
\label{cotangent-equation-pi-schemes}
\vcenter{
\xymatrix{
(\Sh(\mathcal{C}_{X/\Lambda}), \underline{\mathcal{O}}_X) \ar[r]_i \ar[d]_\pi &
(\Sh(\mathcal{C}_{X/\Lambda}), \mathcal{O}) \\
(\Sh(X_{Zar}), \mathcal{O}_X)
}
}
\end{equation}
Le morphisme $i$ est l'identité sur les topos sous-jacents, et
$i^\sharp : \mathcal{O} \to \underline{\mathcal{O}}_X$ est l'homomorphisme
évident. Le morphisme $\pi$ est un cas particulier de Cohomologie sur les
sites, situation \ref{sites-cohomology-situation-fibred-category}. Les
foncteurs dérivés suivants joueront un rôle important :
$
Li^* : D(\mathcal{O}) \longrightarrow D(\underline{\mathcal{O}}_X)
$
adjoint à gauche de
$Ri_* = i_* : D(\underline{\mathcal{O}}_X) \to D(\mathcal{O})$, et
$
L\pi_! : D(\underline{\mathcal{O}}_X) \longrightarrow D(\mathcal{O}_X)
$
adjoint à gauche de
$\pi^* = \pi^{-1} : D(\mathcal{O}_X) \to D(\underline{\mathcal{O}}_X)$. Nos
travaux précédents permettent de calculer $L\pi_!$.

\begin{remark}
\label{cotangent-remark-compute-L-pi-shriek}
Dans la situation précédente, pour tout ouvert $U \subset X$, soit
$P_{\bullet, U}$ la résolution standard de $\mathcal{O}_X(U)$ sur $\Lambda$.
Posons $\mathbf{A}_{n, U} = \Spec(P_{n, U})$. Alors
$\mathbf{A}_{\bullet, U}$
est un objet cosimplicial de la catégorie fibre
$\mathcal{C}_{\mathcal{O}_X(U)/\Lambda}$ de $\mathcal{C}_{X/\Lambda}$ au-dessus
de $U$. De plus, comme expliqué dans la remarque \ref{cotangent-remark-resolution},
$\mathbf{A}_{\bullet, U}$ est un objet cosimplicial de
$\mathcal{C}_{\mathcal{O}_X(U)/\Lambda}$ comme dans Cohomologie sur les sites,
lemme \ref{sites-cohomology-lemma-compute-by-cosimplicial-resolution}. La
construction $U \mapsto \mathbf{A}_{\bullet, U}$ étant fonctorielle en $U$,
tout faisceau (abélien) $\mathcal{F}$ sur $\mathcal{C}_{X/\Lambda}$ donne un
complexe de préfaisceaux
$$
U \longmapsto \mathcal{F}(\mathbf{A}_{\bullet, U})
$$
dont les groupes de cohomologie calculent l'homologie de $\mathcal{F}$ sur la
catégorie fibre. D'après Cohomologie sur les sites, lemme
\ref{sites-cohomology-lemma-compute-left-derived-pi-shriek}, nous en concluons
que la faisceautisation calcule $L_n\pi_!(\mathcal{F})$. Autrement dit, le
complexe de faisceaux dont le terme de degré $-n$ est la faisceautisation de
$U \mapsto \mathcal{F}(\mathbf{A}_{n, U})$ calcule $L\pi_!(\mathcal{F})$.
\end{remark}

\noindent
Nous pouvons maintenant énoncer le résultat principal de cette section.

\begin{lemma}
\label{cotangent-lemma-cotangent-morphism-schemes}
Dans la situation précédente, il existe un isomorphisme canonique
$$
L_{X/\Lambda} = 
L\pi_!(Li^*\Omega_{\mathcal{O}/\underline{\Lambda}}) =
L\pi_!(i^*\Omega_{\mathcal{O}/\underline{\Lambda}}) =
L\pi_!(\Omega_{\mathcal{O}/\underline{\Lambda}}
\otimes_\mathcal{O} \underline{\mathcal{O}}_X)
$$
dans $D(\mathcal{O}_X)$.
\end{lemma}

\begin{proof}
Observons d'abord que, pour tout objet $(U \to \mathbf{A})$ de
$\mathcal{C}_{X/\Lambda}$, la valeur du faisceau $\mathcal{O}$ est une algèbre
polynomiale sur $\Lambda$. Ainsi,
$\Omega_{\mathcal{O}/\underline{\Lambda}}$ est un $\mathcal{O}$-module plat,
et les deuxième et troisième égalités de l'énoncé sont vérifiées.

\medskip\noindent
D'après la remarque \ref{cotangent-remark-compute-L-pi-shriek}, l'objet
$L\pi_!(\Omega_{\mathcal{O}/\underline{\Lambda}}
\otimes_\mathcal{O} \underline{\mathcal{O}}_X)$
se calcule comme la faisceautisation du complexe de préfaisceaux
$$
U \mapsto
\left(\Omega_{\mathcal{O}/\underline{\Lambda}}
\otimes_\mathcal{O} \underline{\mathcal{O}}_X\right)(\mathbf{A}_{\bullet, U})
=
\Omega_{P_{\bullet, U}/\Lambda} \otimes_{P_{\bullet, U}} \mathcal{O}_X(U) =
L_{\mathcal{O}_X(U)/\Lambda}
$$
avec les notations de la remarque \ref{cotangent-remark-compute-L-pi-shriek}. La remarque
\ref{cotangent-remark-map-sections-over-U} montre alors que
$L\pi_!(\Omega_{\mathcal{O}/\underline{\Lambda}}
\otimes_\mathcal{O} \underline{\mathcal{O}}_X)$
calcule le complexe cotangent de l'homomorphisme d'anneaux
$\underline{\Lambda} \to \mathcal{O}_X$ sur $X$. C'est le résultat voulu
d'après le lemme \ref{cotangent-lemma-scheme-over-ring}.
\end{proof}








\section{Le complexe cotangent d'un morphisme d'espaces algébriques}
\label{cotangent-section-cotangent-morphism-spaces}

\noindent
Nous définissons le complexe cotangent d'un morphisme d'espaces algébriques au
moyen du morphisme associé entre les petits sites étales.

\begin{definition}
\label{cotangent-definition-cotangent-morphism-spaces}
Soit $S$ un schéma. Soit $f : X \to Y$ un morphisme d'espaces algébriques sur
$S$. Le {\it complexe cotangent $L_{X/Y}$ de $X$ sur $Y$} est le complexe
cotangent du morphisme de topos annelés $f_{small}$ entre les petits sites
étales de $X$ et de $Y$ (voir Propriétés des espaces, lemme
\ref{spaces-properties-lemma-morphism-ringed-topoi}, et définition
\ref{cotangent-definition-cotangent-complex-morphism-ringed-topoi}).
\end{definition}

\noindent
En particulier, les résultats de la section
\ref{cotangent-section-cotangent-morphism-ringed-topoi} s'appliquent aux complexes
cotangents des morphismes d'espaces algébriques. Les lemmes suivants montrent
que cette définition est compatible avec celles données pour les
homomorphismes d'anneaux et pour les schémas, et que $L_{X/Y}$ est un objet de
$D_\QCoh(\mathcal{O}_X)$.

\begin{lemma}
\label{cotangent-lemma-etale-localization}
Soit $S$ un schéma. Considérons un diagramme commutatif
$$
\xymatrix{
U \ar[d]_p \ar[r]_g & V \ar[d]^q \\
X \ar[r]^f & Y
}
$$
d'espaces algébriques sur $S$, où $p$ et $q$ sont étales. Il existe alors une
identification canonique $L_{X/Y}|_{U_\etale} = L_{U/V}$ dans
$D(\mathcal{O}_U)$.
\end{lemma}

\begin{proof}
La formation du complexe cotangent commute à l'image inverse (lemme
\ref{cotangent-lemma-pullback-cotangent-morphism-topoi}), et nous avons
$p_{small}^{-1}\mathcal{O}_X = \mathcal{O}_U$ et
$g_{small}^{-1}\mathcal{O}_{V_\etale} =
p_{small}^{-1}f_{small}^{-1}\mathcal{O}_{Y_\etale}$
parce que $q_{small}^{-1}\mathcal{O}_{Y_\etale} =
\mathcal{O}_{V_\etale}$
(Propriétés des espaces, lemme
\ref{spaces-properties-lemma-etale-exact-pullback}).
En suivant les définitions, nous en concluons que
$L_{X/Y}|_{U_\etale} = L_{U/V}$.
\end{proof}

\begin{lemma}
\label{cotangent-lemma-compare-spaces-schemes}
Soit $S$ un schéma. Soit $f : X \to Y$ un morphisme d'espaces algébriques sur
$S$. Supposons que $X$ et $Y$ soient représentables par des schémas $X_0$ et
$Y_0$. Il existe alors une identification canonique
$L_{X/Y} = \epsilon^*L_{X_0/Y_0}$ in $D(\mathcal{O}_X)$
où $\epsilon$ est celui de Catégories dérivées des espaces, section
\ref{spaces-perfect-section-derived-quasi-coherent-etale}, et où
$L_{X_0/Y_0}$ est celui de la définition
\ref{cotangent-definition-cotangent-morphism-schemes}.
\end{lemma}

\begin{proof}
Soit $f_0 : X_0 \to Y_0$ le morphisme de schémas correspondant à $f$. Il
existe une application canonique
$\epsilon^{-1}f_0^{-1}\mathcal{O}_{Y_0} \to f_{small}^{-1}\mathcal{O}_Y$
compatible à
$\epsilon^\sharp : \epsilon^{-1}\mathcal{O}_{X_0} \to \mathcal{O}_X$
parce qu'il existe un diagramme commutatif
$$
\xymatrix{
X_{0, Zar} \ar[d]_{f_0} & X_\etale \ar[l]^\epsilon \ar[d]^f \\
Y_{0, Zar} & Y_\etale \ar[l]_\epsilon
}
$$
voir Catégories dérivées des espaces, remarque
\ref{spaces-perfect-remark-match-total-direct-images}.
Nous obtenons ainsi une application canonique
$$
\epsilon^{-1}L_{X_0/Y_0} =
\epsilon^{-1}L_{\mathcal{O}_{X_0}/f_0^{-1}\mathcal{O}_{Y_0}} =
L_{\epsilon^{-1}\mathcal{O}_{X_0}/\epsilon^{-1}f_0^{-1}\mathcal{O}_{Y_0}}
\longrightarrow
L_{\mathcal{O}_X/f^{-1}_{small}\mathcal{O}_Y} = L_{X/Y}
$$
par la fonctorialité étudiée dans la section
\ref{cotangent-section-cotangent-complex} et le lemme
\ref{cotangent-lemma-pullback-cotangent-morphism-topoi}. Pour voir que l'application
induite $\epsilon^*L_{X_0/Y_0} \to L_{X/Y}$ est un isomorphisme, nous pouvons
le vérifier sur les fibres aux points géométriques (Propriétés des espaces,
théorème \ref{spaces-properties-theorem-exactness-stalks}). Nous utiliserons
le lemme \ref{cotangent-lemma-stalk-cotangent-complex} pour calculer les fibres. Soit
$\overline{x} : \Spec(k) \to X_0$ un point géométrique au-dessus de
$x \in X_0$, et soit $\overline{y} = f \circ \overline{x}$ au-dessus de
$y \in Y_0$. Alors
$$
L_{X/Y, \overline{x}} =
L_{\mathcal{O}_{X, \overline{x}}/\mathcal{O}_{Y, \overline{y}}}
$$
et
$$
(\epsilon^*L_{X_0/Y_0})_{\overline{x}} =
L_{X_0/Y_0, x} \otimes_{\mathcal{O}_{X_0, x}}
\mathcal{O}_{X, \overline{x}} =
L_{\mathcal{O}_{X_0, x}/\mathcal{O}_{Y_0, y}}
\otimes_{\mathcal{O}_{X_0, x}} \mathcal{O}_{X, \overline{x}}
$$
Nous omettons quelques détails (indication : utiliser le fait que la fibre
d'une image inverse est la fibre au point image ; voir Sites, lemme
\ref{sites-lemma-point-morphism-sites}, ainsi que le résultat correspondant
pour les modules, Modules sur les sites, lemme
\ref{sites-modules-lemma-pullback-stalk}). Remarquons que
$\mathcal{O}_{X, \overline{x}}$ est le hensélisé strict de
$\mathcal{O}_{X_0, x}$, et de même pour $\mathcal{O}_{Y, \overline{y}}$
(Propriétés des espaces, lemme
\ref{spaces-properties-lemma-describe-etale-local-ring}). Le résultat découle
donc du lemme \ref{cotangent-lemma-cotangent-complex-henselization}.
\end{proof}

\begin{lemma}
\label{cotangent-lemma-space-over-ring}
Soit $\Lambda$ un anneau. Soit $X$ un espace algébrique sur $\Lambda$. Alors
$$
L_{X/\Spec(\Lambda)} = L_{\mathcal{O}_X/\underline{\Lambda}}
$$
où $\underline{\Lambda}$ est le faisceau constant de valeur $\Lambda$ sur
$X_\etale$.
\end{lemma}

\begin{proof}
Soit $p : X \to \Spec(\Lambda)$ le morphisme structural. Soit
$q : \Spec(\Lambda)_\etale \to (*, \Lambda)$ le morphisme évident. D'après le
triangle distingué du lemme \ref{cotangent-lemma-triangle-ringed-topoi}, il suffit de
montrer que $L_q = 0$. Pour cela, il suffit de montrer (Propriétés des espaces,
théorème \ref{spaces-properties-theorem-exactness-stalks}) que, pour tout point
géométrique $\overline{t} : \Spec(k) \to \Spec(\Lambda)$,
$$
(L_q)_{\overline{t}} =
L_{\mathcal{O}_{\Spec(\Lambda)_\etale, \overline{t}}/\Lambda}
$$
(lemme \ref{cotangent-lemma-stalk-cotangent-complex}) est nul. Comme
$\mathcal{O}_{\Spec(\Lambda)_\etale, \overline{t}}$ est le hensélisé strict
d'un anneau local de $\Lambda$ (Propriétés des espaces, lemme
\ref{spaces-properties-lemma-describe-etale-local-ring}), cela résulte du lemme
\ref{cotangent-lemma-when-zero}.
\end{proof}










\section{Le complexe cotangent d'un espace algébrique sur un anneau}
\sectionmark{Complexe cotangent sur un anneau}
\label{cotangent-section-cotangent-spaces-variant}

\noindent
Soit $\Lambda$ un anneau et soit $X$ un espace algébrique sur $\Lambda$.
Écrivons $L_{X/\Spec(\Lambda)} = L_{X/\Lambda}$, ce qui est justifié par le
lemme \ref{cotangent-lemma-space-over-ring}. Dans cette section, nous donnons de
$L_{X/\Lambda}$ une description analogue à celle du lemme
\ref{cotangent-lemma-compute-cotangent-complex}. Plus précisément, nous construisons une
catégorie $\mathcal{C}_{X/\Lambda}$ fibrée au-dessus de $X_\etale$ et la
munissons d'un faisceau $\mathcal{O}$ de $\Lambda$-algèbres (polynomiales) tel
que
$$
L_{X/\Lambda} =
L\pi_!(\Omega_{\mathcal{O}/\underline{\Lambda}} \otimes_\mathcal{O}
\underline{\mathcal{O}}_X).
$$
Nous utiliserons plus loin la catégorie $\mathcal{C}_{X/\Lambda}$ pour
construire une théorie naïve des obstructions pour le champ des faisceaux
cohérents.

\medskip\noindent
Soit $\Lambda$ un anneau. Soit $X$ un espace algébrique sur $\Lambda$. Soit
$\mathcal{C}_{X/\Lambda}$ la catégorie dont les objets sont les diagrammes
commutatifs
\begin{equation}
\label{cotangent-equation-object-space}
\vcenter{
\xymatrix{
X \ar[d] & U \ar[l] \ar[d] \\
\Spec(\Lambda) & \mathbf{A} \ar[l]
}
}
\end{equation}
de schémas tels que
\begin{enumerate}
\item $U$ soit un schéma,
\item $U \to X$ soit étale,
\item il existe un isomorphisme $\mathbf{A} = \Spec(P)$, où $P$ est une
algèbre polynomiale sur $\Lambda$ (en un certain ensemble de variables).
\end{enumerate}
Autrement dit, $\mathbf{A}$ est un espace affine (de dimension infinie) sur
$\Spec(\Lambda)$. Les morphismes sont donnés par des diagrammes commutatifs.
Rappelons que $X_\etale$ désigne le petit site étale de $X$, dont les objets
sont les schémas étales sur $X$. Il existe un foncteur d'oubli
$$
u : \mathcal{C}_{X/\Lambda} \to X_\etale,
\quad
(U \to \mathbf{A}) \mapsto U
$$
Remarquons que la catégorie fibre au-dessus de $U$ est canoniquement
équivalente à la catégorie $\mathcal{C}_{\mathcal{O}_X(U)/\Lambda}$ introduite
dans la section \ref{cotangent-section-compute-L-pi-shriek}.

\begin{lemma}
\label{cotangent-lemma-category-fibred-space}
Dans la situation précédente, la catégorie $\mathcal{C}_{X/\Lambda}$ est
fibrée au-dessus de $X_\etale$.
\end{lemma}

\begin{proof}
Étant donnés un objet $U \to \mathbf{A}$ de $\mathcal{C}_{X/\Lambda}$ et un
morphisme $U' \to U$ de $X_\etale$, considérons l'objet
$U' \to \mathbf{A}$ de $\mathcal{C}_{X/\Lambda}$, où
$U' \to \mathbf{A}$ est la composée de $U' \to U$ et de
$U \to \mathbf{A}$. Le morphisme
$(U' \to \mathbf{A}) \to (U \to \mathbf{A})$ de
$\mathcal{C}_{X/\Lambda}$ est fortement cartésien au-dessus de $X_\etale$.
\end{proof}

\noindent
Munissons $\mathcal{C}_{X/\Lambda}$ de la topologie héritée de $X_\etale$
(voir Champs, section \ref{stacks-section-topology}). Le foncteur $u$ définit
un morphisme de topos
$\pi : \Sh(\mathcal{C}_{X/\Lambda}) \to \Sh(X_\etale)$. Le site
$\mathcal{C}_{X/\Lambda}$ est muni de plusieurs faisceaux d'anneaux :
\begin{enumerate}
\item le faisceau $\mathcal{O}$ donné par la règle
$(U \to \mathbf{A}) \mapsto \Gamma(\mathbf{A}, \mathcal{O}_\mathbf{A})$.
\item le faisceau $\underline{\mathcal{O}}_X = \pi^{-1}\mathcal{O}_X$ donné
par la règle $(U \to \mathbf{A}) \mapsto \mathcal{O}_X(U)$ ;
\item le faisceau constant $\underline{\Lambda}$.
\end{enumerate}
Nous obtenons des morphismes de topos annelés
\begin{equation}
\label{cotangent-equation-pi-spaces}
\vcenter{
\xymatrix{
(\Sh(\mathcal{C}_{X/\Lambda}), \underline{\mathcal{O}}_X) \ar[r]_i \ar[d]_\pi &
(\Sh(\mathcal{C}_{X/\Lambda}), \mathcal{O}) \\
(\Sh(X_\etale), \mathcal{O}_X)
}
}
\end{equation}
Le morphisme $i$ est l'identité sur les topos sous-jacents, et
$i^\sharp : \mathcal{O} \to \underline{\mathcal{O}}_X$ est l'homomorphisme
évident. Le morphisme $\pi$ est un cas particulier de Cohomologie sur les
sites, situation \ref{sites-cohomology-situation-fibred-category}. Les
foncteurs dérivés suivants joueront un rôle important :
$
Li^* : D(\mathcal{O}) \longrightarrow D(\underline{\mathcal{O}}_X)
$
adjoint à gauche de
$Ri_* = i_* : D(\underline{\mathcal{O}}_X) \to D(\mathcal{O})$, et
$
L\pi_! : D(\underline{\mathcal{O}}_X) \longrightarrow D(\mathcal{O}_X)
$
adjoint à gauche de
$\pi^* = \pi^{-1} : D(\mathcal{O}_X) \to D(\underline{\mathcal{O}}_X)$. Nos
travaux précédents permettent de calculer $L\pi_!$.

\begin{remark}
\label{cotangent-remark-compute-L-pi-shriek-spaces}
Dans la situation précédente, pour tout objet $U \to X$ de $X_\etale$, soit
$P_{\bullet, U}$ la résolution standard de $\mathcal{O}_X(U)$ sur $\Lambda$.
Posons $\mathbf{A}_{n, U} = \Spec(P_{n, U})$. Alors
$\mathbf{A}_{\bullet, U}$ est un objet cosimplicial de la catégorie fibre
$\mathcal{C}_{\mathcal{O}_X(U)/\Lambda}$ de $\mathcal{C}_{X/\Lambda}$ au-dessus
de $U$. De plus, comme expliqué dans la remarque \ref{cotangent-remark-resolution},
$\mathbf{A}_{\bullet, U}$ est un objet cosimplicial de
$\mathcal{C}_{\mathcal{O}_X(U)/\Lambda}$ comme dans Cohomologie sur les sites,
lemme \ref{sites-cohomology-lemma-compute-by-cosimplicial-resolution}. La
construction $U \mapsto \mathbf{A}_{\bullet, U}$ étant fonctorielle en $U$,
tout faisceau (abélien) $\mathcal{F}$ sur $\mathcal{C}_{X/\Lambda}$ donne un
complexe de préfaisceaux
$$
U \longmapsto \mathcal{F}(\mathbf{A}_{\bullet, U})
$$
dont les groupes de cohomologie calculent l'homologie de $\mathcal{F}$ sur la
catégorie fibre. D'après Cohomologie sur les sites, lemme
\ref{sites-cohomology-lemma-compute-left-derived-pi-shriek}, nous en concluons
que la faisceautisation calcule $L_n\pi_!(\mathcal{F})$. Autrement dit, le
complexe de faisceaux dont le terme de degré $-n$ est la faisceautisation de
$U \mapsto \mathcal{F}(\mathbf{A}_{n, U})$ calcule $L\pi_!(\mathcal{F})$.
\end{remark}

\noindent
Nous pouvons maintenant énoncer le résultat principal de cette section.

\begin{lemma}
\label{cotangent-lemma-cotangent-morphism-spaces}
Dans la situation précédente, il existe un isomorphisme canonique
$$
L_{X/\Lambda} = 
L\pi_!(Li^*\Omega_{\mathcal{O}/\underline{\Lambda}}) =
L\pi_!(i^*\Omega_{\mathcal{O}/\underline{\Lambda}}) =
L\pi_!(\Omega_{\mathcal{O}/\underline{\Lambda}}
\otimes_\mathcal{O} \underline{\mathcal{O}}_X)
$$
dans $D(\mathcal{O}_X)$.
\end{lemma}

\begin{proof}
Observons d'abord que, pour tout objet $(U \to \mathbf{A})$ de
$\mathcal{C}_{X/\Lambda}$, la valeur du faisceau $\mathcal{O}$ est une algèbre
polynomiale sur $\Lambda$. Ainsi,
$\Omega_{\mathcal{O}/\underline{\Lambda}}$ est un $\mathcal{O}$-module plat,
et les deuxième et troisième égalités de l'énoncé sont vérifiées.

\medskip\noindent
D'après la remarque \ref{cotangent-remark-compute-L-pi-shriek-spaces}, l'objet
$L\pi_!(\Omega_{\mathcal{O}/\underline{\Lambda}}
\otimes_\mathcal{O} \underline{\mathcal{O}}_X)$
se calcule comme la faisceautisation du complexe de préfaisceaux
$$
U \mapsto
\left(\Omega_{\mathcal{O}/\underline{\Lambda}}
\otimes_\mathcal{O} \underline{\mathcal{O}}_X\right)(\mathbf{A}_{\bullet, U})
=
\Omega_{P_{\bullet, U}/\Lambda} \otimes_{P_{\bullet, U}} \mathcal{O}_X(U) =
L_{\mathcal{O}_X(U)/\Lambda}
$$
avec les notations de la remarque \ref{cotangent-remark-compute-L-pi-shriek-spaces}. La
remarque \ref{cotangent-remark-map-sections-over-U} montre alors que
$L\pi_!(\Omega_{\mathcal{O}/\underline{\Lambda}}
\otimes_\mathcal{O} \underline{\mathcal{O}}_X)$
calcule le complexe cotangent de l'homomorphisme d'anneaux
$\underline{\Lambda} \to \mathcal{O}_X$ sur $X_\etale$. C'est le résultat
voulu d'après le lemme \ref{cotangent-lemma-space-over-ring}.
\end{proof}






\section{Produits fibrés d'espaces algébriques et complexe cotangent}
\sectionmark{Produits fibrés et complexe cotangent}
\label{cotangent-section-fibre-product}

\noindent
Soit $S$ un schéma. Soient $X \to B$ et $Y \to B$ des morphismes d'espaces
algébriques sur $S$. Considérons le produit fibré $X \times_B Y$, muni des
morphismes de projection $p : X \times_B Y \to X$ et
$q : X \times_B Y \to Y$. Dans cette section, nous étudions
$L_{X \times_B Y/B}$. L'essentiel des renseignements recherchés est contenu
dans le diagramme suivant
\begin{equation}
\label{cotangent-equation-fibre-product}
\vcenter{
\xymatrix{
Lp^*L_{X/B} \ar[r] &
L_{X \times_B Y/Y} \ar[r] &
E \\
Lp^*L_{X/B} \ar[r] \ar@{=}[u] &
L_{X \times_B Y/B} \ar[r] \ar[u] &
L_{X \times_B Y/X} \ar[u] \\
 &
Lq^*L_{Y/B} \ar[u] \ar@{=}[r] &
Lq^*L_{Y/B} \ar[u]
}
}
\end{equation}
Explication : la ligne médiane est le triangle fondamental du lemme
\ref{cotangent-lemma-triangle-ringed-topoi} associé aux morphismes
$X \times_B Y \to X \to B$. La colonne médiane est le triangle fondamental
associé aux morphismes $X \times_B Y \to Y \to B$. Ensuite, $E$ est un objet
de $D(\mathcal{O}_{X \times_B Y})$ qui « complète » le coin supérieur droit,
c'est-à-dire qui fait de la ligne supérieure et de la colonne droite des
triangles distingués. Un tel $E$ existe d'après Catégories dérivées,
proposition \ref{derived-proposition-9}, appliquée au carré inférieur gauche
(en plaçant $0$ à la position manquante). Plus explicitement, nous pourrions
par exemple définir $E$ comme le cône (Catégories dérivées, définition
\ref{derived-definition-cone}) de l'application de complexes
$$
Lp^*L_{X/B} \oplus Lq^*L_{Y/B} \longrightarrow L_{X \times_B Y/B}
$$
et obtenir les deux applications de but $E$ en appliquant TR3. Dans le cas
Tor-indépendant, l'objet $E$ est nul.

\begin{lemma}
\label{cotangent-lemma-fibre-product-tor-independent}
Dans la situation précédente, si $X$ et $Y$ sont Tor-indépendants sur $B$,
alors l'objet $E$ de (\ref{cotangent-equation-fibre-product}) est nul. Dans ce cas, nous
avons
$$
L_{X \times_B Y/B} = Lp^*L_{X/B} \oplus Lq^*L_{Y/B}
$$
\end{lemma}

\begin{proof}
Choisissons un schéma $W$ et un morphisme étale surjectif $W \to B$.
Choisissons un schéma $U$ et un morphisme étale surjectif
$U \to X \times_B W$. Choisissons un schéma $V$ et un morphisme étale
surjectif $V \to Y \times_B W$. Alors
$U \times_W V \to X \times_B Y$ est également étale surjectif. Il suffit donc
de montrer que la restriction de $E$ à $U \times_W V$ est nulle. D'après le
lemme \ref{cotangent-lemma-compare-spaces-schemes} et Catégories dérivées des espaces,
lemme \ref{spaces-perfect-lemma-tor-independent}, nous sommes ramenés au cas
des schémas. En choisissant des ouverts affines convenables, nous nous ramenons
au cas des schémas affines. Le lemme \ref{cotangent-lemma-morphism-affine-schemes} nous
ramène alors au cas d'un produit tensoriel d'anneaux, c'est-à-dire au lemme
\ref{cotangent-lemma-tensor-product-tor-independent}.
\end{proof}

\noindent
Dans le cas général, nous pouvons dire ce qui suit de l'objet $E$.

\begin{lemma}
\label{cotangent-lemma-fibre-product}
Soit $S$ un schéma. Soient $X \to B$ et $Y \to B$ des morphismes d'espaces
algébriques sur $S$. L'objet $E$ de (\ref{cotangent-equation-fibre-product}) vérifie
$H^i(E) = 0$ pour $i = 0, -1$, et, pour tout point géométrique
$(\overline{x}, \overline{y}) : \Spec(k) \to X \times_B Y$, nous avons
$$
H^{-2}(E)_{(\overline{x}, \overline{y})} =
\text{Tor}_1^R(A, B) \otimes_{A \otimes_R B} C
$$
où $R = \mathcal{O}_{B, \overline{b}}$, $A = \mathcal{O}_{X, \overline{x}}$,
$B = \mathcal{O}_{Y, \overline{y}}$ et
$C = \mathcal{O}_{X \times_B Y, (\overline{x}, \overline{y})}$.
\end{lemma}

\begin{proof}
La formation du complexe cotangent commute au passage aux fibres et aux images
inverses ; voir les lemmes \ref{cotangent-lemma-stalk-cotangent-complex} et
\ref{cotangent-lemma-pullback-cotangent-morphism-topoi}. Remarquons que $C$ est un
hensélisé de $A \otimes_R B$.
$L_{C/R} = L_{A \otimes_R B/R} \otimes_{A \otimes_R B} C$
d'après les résultats de la section \ref{cotangent-section-localization}. Ainsi, la fibre
de $E$ en notre point géométrique est le cône de l'application
$L_{A/R} \otimes C \to L_{A \otimes_R B/R} \otimes C$. Les assertions du
lemme résultent donc du cas des anneaux, c'est-à-dire du lemme
\ref{cotangent-lemma-tensor-product}.
\end{proof}












% OK, commencer ici.
%

\chapter{Problèmes de déformation}


\phantomsection
\label{examples-defos-section-phantom}


\section{Introduction}
\label{examples-defos-section-introduction}

\noindent
Le but de ce chapitre est de traiter des exemples de la théorie générale
développée dans les chapitres Théorie formelle des déformations,
Théorie des déformations et Le complexe cotangent.

\medskip\noindent
La section 3 de l'article \cite{Sch} de Schlessinger étudie elle aussi
quelques exemples.






\section{Exemples de problèmes de déformation}
\label{examples-defos-section-examples}

\noindent
Liste des sujets qui devraient figurer ici :
\begin{enumerate}
\item Déformations de schémas :
\begin{enumerate}
\item La condition de Rim-Schlessinger.
\item Calcul de l'espace tangent.
\item Calcul des déformations infinitésimales.
\item La catégorie de déformations d'une hypersurface affine.
\end{enumerate}
\item Déformations de faisceaux (par exemple, fixer $X/S$, un point $s$
de type fini de $S$ et un faisceau quasi-cohérent $\mathcal{F}_s$ sur $X_s$).
\item Déformations d'espaces algébriques (très semblables aux déformations
de schémas ; peut-être même plus faciles ?).
\item Déformations d'applications (par exemple de morphismes de schémas ; on peut fixer
à la fois la source et le but, ou seulement l'un des deux).
\item Ajouter d'autres sujets ici.
\end{enumerate}





\section{Plan général}
\label{examples-defos-section-general}

\noindent
Cette section expose la démarche suivie pour étudier les exemples qui viennent.

\medskip\noindent
Étape I. Dans chaque section, on fixe un anneau noethérien $\Lambda$ et
un morphisme fini d'anneaux $\Lambda \to k$, où $k$ est un corps.
Comme d'habitude, on prend $\mathcal{C}_\Lambda = \mathcal{C}_{\Lambda, k}$
pour catégorie de base ; voir Théorie formelle des déformations,
définition \ref{formal-defos-definition-CLambda}.

\medskip\noindent
Étape II. Dans chaque section, on définit une catégorie $\mathcal{F}$
cofibrée en groupoïdes sur $\mathcal{C}_\Lambda$. Il nous arrivera
de considérer à la place un foncteur
$F : \mathcal{C}_\Lambda \to \textit{Ensembles}$.

\medskip\noindent
Étape III. On explique dans quelle mesure $\mathcal{F}$ satisfait
la condition de Rim-Schlessinger (RS) étudiée dans Théorie formelle
des déformations, section \ref{formal-defos-section-RS-condition}.
De même, on pourra étudier dans quelle mesure $\mathcal{F}$
satisfait (S1) et (S2), ou dans quelle mesure $F$ satisfait
les conditions correspondantes (H1) et (H2) de Schlessinger.
Voir Théorie formelle des déformations, section
\ref{formal-defos-section-schlessinger-conditions}.

\medskip\noindent
Étape IV. Soit $x_0$ un objet de $\mathcal{F}(k)$, autrement dit un objet
de $\mathcal{F}$ sur $k$. Dans ce chapitre, on emploiera la notation
$$
\Deformationcategory_{x_0} = \mathcal{F}_{x_0}
$$
pour désigner la catégorie de prédéformations construite dans Théorie
formelle des déformations, remarque
\ref{formal-defos-remark-localize-cofibered-groupoid}.
Si $\mathcal{F}$ satisfait (RS), alors
$\Deformationcategory_{x_0}$ est une catégorie de déformations
(Théorie formelle des déformations, lemme
\ref{formal-defos-lemma-localize-RS})
et satisfait (S1) et (S2)
(Théorie formelle des déformations, lemme
\ref{formal-defos-lemma-RS-implies-S1-S2}).
Si (S1) et (S2) sont satisfaites, une question importante
est de savoir si l'espace tangent
$$
T\Deformationcategory_{x_0} = T_{x_0}\mathcal{F} = T\mathcal{F}_{x_0}
$$
(voir Théorie formelle des déformations, remarque
\ref{formal-defos-remark-tangent-space-cofibered-groupoid} et
définition \ref{formal-defos-definition-tangent-space})
est de dimension finie. En effet, cela garantit que
$\Deformationcategory_{x_0}$ possède un objet formel versel
(Théorie formelle des déformations, lemme
\ref{formal-defos-lemma-versal-object-existence}).

\medskip\noindent
Étape V. Si $\mathcal{F}$ franchit l'étape IV, la question suivante est de savoir
si le $k$-espace vectoriel
$$
\text{Inf}(\Deformationcategory_{x_0}) = \text{Inf}_{x_0}(\mathcal{F})
$$
des automorphismes infinitésimaux de $x_0$ est de dimension finie.
Dans l'affirmative, cela implique que
$\Deformationcategory_{x_0}$ admet une présentation par un
groupoïde lisse proreprésentable en foncteurs sur $\mathcal{C}_\Lambda$ ; voir
Théorie formelle des déformations, théorème
\ref{formal-defos-theorem-presentation-deformation-groupoid}.





\section{Modules projectifs de type fini}
\label{examples-defos-section-finite-projective-modules}

\noindent
Cette section n'est qu'un échauffement. Bien entendu, les modules projectifs
de type fini ne devraient avoir aucun « module » de paramètres.

\begin{example}[Modules projectifs de type fini]
\label{examples-defos-example-finite-projective-modules}
Soit $\mathcal{F}$ la catégorie définie comme suit :
\begin{enumerate}
\item un objet est un couple $(A, M)$ formé d'un
objet $A$ de $\mathcal{C}_\Lambda$ et d'un $A$-module
projectif de type fini $M$ ;
\item un morphisme $(f, g) : (B, N) \to (A, M)$ est formé d'un
morphisme $f : B \to A$ dans $\mathcal{C}_\Lambda$ et
d'une application $g : N \to M$ qui est $f$-linéaire et induit
un isomorphisme $N \otimes_{B, f} A \cong M$.
\end{enumerate}
Le foncteur $p : \mathcal{F} \to \mathcal{C}_\Lambda$ envoie $(A, M)$ sur $A$
et $(f, g)$ sur $f$. Il est clair que $p$ est cofibré en groupoïdes.
Étant donné un $k$-espace vectoriel de dimension finie $V$,
soit $x_0 = (k, V)$ l'objet correspondant de $\mathcal{F}(k)$.
On pose
$$
\Deformationcategory_V = \mathcal{F}_{x_0}
$$
\end{example}

\noindent
Puisque tout module projectif de type fini sur un anneau local est libre
de type fini (Algèbre, lemme \ref{algebra-lemma-finite-projective}), on a
$$
\begin{matrix}
\text{classes d'isomorphisme} \\
\text{des objets de }\mathcal{F}(A)
\end{matrix}
= \coprod\nolimits_{n \geq 0} \{*\}
$$
Bien que cela signifie que la théorie des déformations de $\mathcal{F}$
est essentiellement triviale, on parcourt néanmoins les étapes décrites
à la section \ref{examples-defos-section-general} afin de donner un exemple simple.

\begin{lemma}
\label{examples-defos-lemma-finite-projective-modules-RS}
L'exemple \ref{examples-defos-example-finite-projective-modules}
satisfait la condition de Rim-Schlessinger (RS).
En particulier, $\Deformationcategory_V$ est une catégorie de déformations
pour tout espace vectoriel $V$ de dimension finie sur $k$.
\end{lemma}

\begin{proof}
Soient $A_1 \to A$ et $A_2 \to A$ des morphismes de $\mathcal{C}_\Lambda$.
Supposons $A_2 \to A$ surjectif. D'après Théorie formelle
des déformations, lemme
\ref{formal-defos-lemma-RS-2-categorical}
il suffit de montrer que le foncteur
$\mathcal{F}(A_1 \times_A A_2) \to
\mathcal{F}(A_1) \times_{\mathcal{F}(A)} \mathcal{F}(A_2)$
est une équivalence de catégories.

\medskip\noindent
Il faut donc montrer que la catégorie des modules projectifs de type fini
sur $A_1 \times_A A_2$ est équivalente au produit fibré
des catégories des modules projectifs de type fini sur $A_1$ et $A_2$
au-dessus de la catégorie des modules projectifs de type fini sur $A$.
C'est un cas particulier de Compléments d'algèbre, lemme
\ref{more-algebra-lemma-finitely-presented-module-over-fibre-product}.
Rappelons que le foncteur inverse envoie le triplet
$(M_1, M_2, \varphi)$, où
$M_1$ est un $A_1$-module projectif de type fini,
$M_2$ est un $A_2$-module projectif de type fini et
$\varphi : M_1 \otimes_{A_1} A \to M_2 \otimes_{A_2} A$
est un isomorphisme de $A$-modules, sur le $A_1 \times_A A_2$-module
projectif de type fini $M_1 \times_\varphi M_2$.
\end{proof}

\begin{lemma}
\label{examples-defos-lemma-finite-projective-modules-TI}
Dans l'exemple \ref{examples-defos-example-finite-projective-modules},
soit $V$ un $k$-espace vectoriel de dimension finie. Alors
$$
T\Deformationcategory_V = (0)
\quad\text{et}\quad
\text{Inf}(\Deformationcategory_V) = \text{End}_k(V)
$$
sont de dimension finie.
\end{lemma}

\begin{proof}
Avec $\mathcal{F}$ comme dans l'exemple \ref{examples-defos-example-finite-projective-modules},
posons $x_0 = (k, V) \in \Ob(\mathcal{F}(k))$.
Rappelons que $T\Deformationcategory_V = T_{x_0}\mathcal{F}$
est l'ensemble des classes d'isomorphisme
de couples $(x, \alpha)$ formés d'un objet $x$ de $\mathcal{F}$
sur l'anneau des nombres duaux $k[\epsilon]$ et d'un morphisme
$\alpha : x \to x_0$ de $\mathcal{F}$ au-dessus de $k[\epsilon] \to k$.

\medskip\noindent
À isomorphisme près, il existe un unique couple $(M, \alpha)$ formé d'un
module projectif de type fini $M$ sur $k[\epsilon]$
et d'une application $k[\epsilon]$-linéaire $\alpha : M \to V$
qui induit un isomorphisme $M \otimes_{k[\epsilon]} k \to V$.
Par exemple, si $V = k^{\oplus n}$, on prend
$M = k[\epsilon]^{\oplus n}$ muni de l'application évidente $\alpha$.

\medskip\noindent
De même, $\text{Inf}(\Deformationcategory_V) = \text{Inf}_{x_0}(\mathcal{F})$
est l'ensemble des automorphismes
de la déformation triviale $x'_0$ de $x_0$ sur $k[\epsilon]$.
Voir Théorie formelle des déformations, définition
\ref{formal-defos-definition-infinitesimal-auts} pour plus de détails.

\medskip\noindent
Pour $(M, \alpha)$ comme au deuxième paragraphe, on voit qu'un élément de
$\text{Inf}_{x_0}(\mathcal{F})$ est un automorphisme $\gamma : M \to M$ tel que
$\gamma \bmod \epsilon = \text{id}$. On peut donc écrire
$\gamma = \text{id}_M + \epsilon \psi$, où
$\psi : M/\epsilon M \to M/\epsilon M$ est $k$-linéaire.
Au moyen de $\alpha$, on peut voir $\psi$ comme un élément de
$\text{End}_k(V)$, ce qui achève la démonstration.
\end{proof}



\section{Représentations d'un groupe}
\label{examples-defos-section-representations}

\noindent
La théorie des déformations des représentations peut être très intéressante.

\begin{example}[Représentations d'un groupe]
\label{examples-defos-example-representations}
Soit $\Gamma$ un groupe.
Soit $\mathcal{F}$ la catégorie définie comme suit :
\begin{enumerate}
\item un objet est un triplet $(A, M, \rho)$ formé d'un
objet $A$ de $\mathcal{C}_\Lambda$, d'un $A$-module projectif de type fini $M$
et d'un homomorphisme $\rho : \Gamma \to \text{GL}_A(M)$ ;
\item un morphisme $(f, g) : (B, N, \tau) \to (A, M, \rho)$ est formé d'un
morphisme $f : B \to A$ dans $\mathcal{C}_\Lambda$ et
d'une application $g : N \to M$ qui est $f$-linéaire, $\Gamma$-équivariante
et induit un isomorphisme $N \otimes_{B, f} A \cong M$.
\end{enumerate}
Le foncteur $p : \mathcal{F} \to \mathcal{C}_\Lambda$ envoie $(A, M, \rho)$
sur $A$ et $(f, g)$ sur $f$. Il est clair que $p$ est cofibré en groupoïdes.
Étant donné un $k$-espace vectoriel de dimension finie $V$ et une représentation
$\rho_0 : \Gamma \to \text{GL}_k(V)$,
soit $x_0 = (k, V, \rho_0)$ l'objet correspondant de $\mathcal{F}(k)$.
On pose
$$
\Deformationcategory_{V, \rho_0} = \mathcal{F}_{x_0}
$$
\end{example}

\noindent
Puisque tout module projectif de type fini sur un anneau local est libre
de type fini (Algèbre, lemme \ref{algebra-lemma-finite-projective}),
on a
$$
\begin{matrix}
\text{classes d'isomorphisme} \\
\text{des objets de }\mathcal{F}(A)
\end{matrix}
=
\coprod\nolimits_{n \geq 0}\quad
\begin{matrix}
\text{GL}_n(A)\text{-orbites de conjugaison des}\\
\text{homomorphismes }\rho : \Gamma \to \text{GL}_n(A)
\end{matrix}
$$
C'est déjà plus intéressant que l'étude menée à la
section \ref{examples-defos-section-finite-projective-modules}.

\begin{lemma}
\label{examples-defos-lemma-representations-RS}
L'exemple \ref{examples-defos-example-representations}
satisfait la condition de Rim-Schlessinger (RS).
En particulier, $\Deformationcategory_{V, \rho_0}$ est une catégorie de déformations
pour toute représentation de dimension finie
$\rho_0 : \Gamma \to \text{GL}_k(V)$.
\end{lemma}

\begin{proof}
Soient $A_1 \to A$ et $A_2 \to A$ des morphismes de $\mathcal{C}_\Lambda$.
Supposons $A_2 \to A$ surjectif. D'après Théorie formelle
des déformations, lemme
\ref{formal-defos-lemma-RS-2-categorical}
il suffit de montrer que le foncteur
$\mathcal{F}(A_1 \times_A A_2) \to
\mathcal{F}(A_1) \times_{\mathcal{F}(A)} \mathcal{F}(A_2)$
est une équivalence de catégories.

\medskip\noindent
Considérons un objet
$$
((A_1, M_1, \rho_1), (A_2, M_2, \rho_2), (\text{id}_A, \varphi))
$$
de la catégorie $\mathcal{F}(A_1) \times_{\mathcal{F}(A)} \mathcal{F}(A_2)$.
Alors, comme on l'a vu dans la démonstration du lemme \ref{examples-defos-lemma-finite-projective-modules-RS},
on peut considérer le $A_1 \times_A A_2$-module projectif
de type fini $M_1 \times_\varphi M_2$.
Puisque $\varphi$ est compatible aux actions données, on obtient
$$
\rho_1 \times \rho_2 : \Gamma \longrightarrow
\text{GL}_{A_1 \times_A A_2}(M_1 \times_\varphi M_2)
$$
Alors $(M_1 \times_\varphi M_2, \rho_1 \times \rho_2)$
est un objet de $\mathcal{F}(A_1 \times_A A_2)$.
Cette construction détermine un quasi-inverse de notre foncteur.
\end{proof}

\begin{lemma}
\label{examples-defos-lemma-representations-TI}
Dans l'exemple \ref{examples-defos-example-representations}, soit
$\rho_0 : \Gamma \to \text{GL}_k(V)$
une représentation de dimension finie. Alors
$$
T\Deformationcategory_{V, \rho_0} = \Ext^1_{k[\Gamma]}(V, V) =
H^1(\Gamma, \text{End}_k(V))
\quad\text{et}\quad
\text{Inf}(\Deformationcategory_{V, \rho_0}) = H^0(\Gamma, \text{End}_k(V))
$$
Ainsi, $\text{Inf}(\Deformationcategory_{V, \rho_0})$
est toujours de dimension finie,
et $T\Deformationcategory_{V, \rho_0}$ est de dimension finie
si $\Gamma$ est de type fini.
\end{lemma}

\begin{proof}
Traitons d'abord les automorphismes infinitésimaux.
Soit $M = V \otimes_k k[\epsilon]$, muni de l'action induite
$\rho_0' : \Gamma \to \text{GL}_n(M)$.
Un automorphisme infinitésimal, c'est-à-dire un élément de
$\text{Inf}(\Deformationcategory_{V, \rho_0})$,
est alors donné par un automorphisme
$\gamma = \text{id} + \epsilon \psi : M \to M$
comme dans la démonstration du lemme \ref{examples-defos-lemma-finite-projective-modules-TI},
où, de plus, $\psi$ doit commuter
à l'action de $\Gamma$ (donnée par $\rho_0$).
On voit donc que
$$
\text{Inf}(\Deformationcategory_{V, \rho_0}) = H^0(\Gamma, \text{End}_k(V))
$$
comme l'affirme le lemme.

\medskip\noindent
Ensuite, soit $(k[\epsilon], M, \rho)$ un objet de $\mathcal{F}$
sur $k[\epsilon]$, et soit $\alpha : M \to V$ une application $\Gamma$-équivariante
induisant un isomorphisme $M/\epsilon M \to V$.
Puisque $M$ est libre comme $k[\epsilon]$-module, on obtient
une extension de $\Gamma$-modules
$$
0 \to V \to M \xrightarrow{\alpha} V \to 0
$$
On omet la construction détaillée de l'application de gauche.
Réciproquement, si l'on dispose d'une extension de $\Gamma$-modules comme
ci-dessus, on peut l'utiliser pour définir une structure de $k[\epsilon]$-module
sur $M$ et obtenir un objet de $\mathcal{F}(k[\epsilon])$
ainsi qu'une application $\alpha$ comme ci-dessus.
Il s'ensuit que
$$
T\Deformationcategory_{V, \rho_0} = \Ext^1_{k[\Gamma]}(V, V)
$$
comme l'affirme le lemme. Cet espace est égal à
$H^1(\Gamma, \text{End}_k(V))$ d'après Cohomologie étale,
lemme \ref{etale-cohomology-lemma-ext-modules-hom}.

\medskip\noindent
L'assertion sur les dimensions résulte de Cohomologie étale,
lemme
\ref{etale-cohomology-lemma-finite-dim-group-cohomology}.
\end{proof}

\noindent
Dans l'exemple \ref{examples-defos-example-representations}, si $\Gamma$ est de type fini
et si $(V, \rho_0)$ est une représentation de dimension finie de $\Gamma$
sur $k$, alors $\Deformationcategory_{V, \rho_0}$
admet une présentation par un groupoïde lisse proreprésentable en foncteurs
sur $\mathcal{C}_\Lambda$
et possède a fortiori un objet formel versel (minimal). Cela résulte
des lemmes \ref{examples-defos-lemma-representations-RS} et \ref{examples-defos-lemma-representations-TI}
ainsi que de la discussion générale de la section \ref{examples-defos-section-general}.

\begin{lemma}
\label{examples-defos-lemma-representations-hull}
Dans l'exemple \ref{examples-defos-example-representations}, supposons $\Gamma$ de type fini.
Soit $\rho_0 : \Gamma \to \text{GL}_k(V)$ une représentation de dimension finie.
Supposons que $\Lambda$ soit un anneau local complet de corps résiduel $k$
(cas classique). Alors le foncteur
$$
F : \mathcal{C}_\Lambda \longrightarrow \textit{Ensembles},\quad
A \longmapsto \Ob(\Deformationcategory_{V, \rho_0}(A))/\cong
$$
des classes d'isomorphisme d'objets possède une enveloppe. Si
$H^0(\Gamma, \text{End}_k(V)) = k$, alors $F$ est
proreprésentable.
\end{lemma}

\begin{proof}
L'existence d'une enveloppe résulte des lemmes \ref{examples-defos-lemma-representations-RS} et
\ref{examples-defos-lemma-representations-TI} et
Théorie formelle des déformations, lemme \ref{formal-defos-lemma-RS-implies-S1-S2}
et remarque \ref{formal-defos-remark-compose-minimal-into-iso-classes}.

\medskip\noindent
Supposons $H^0(\Gamma, \text{End}_k(V)) = k$. Pour voir que $F$
est proreprésentable, il suffit de montrer que $F$ est un
foncteur de déformations ; voir Théorie formelle des déformations, théorème
\ref{formal-defos-theorem-Schlessinger-prorepresentability}.
Autrement dit, il faut montrer que $F$ satisfait (RS).
On peut pour cela utiliser le critère de Théorie formelle des déformations, lemme
\ref{formal-defos-lemma-RS-associated-functor}.
La surjectivité requise des groupes d'automorphismes résultera de ce que
l'on montre que
$$
A \cdot \text{id}_M =
\text{End}_{A[\Gamma]}(M)
$$
pour tout objet $(A, M, \rho)$ de $\mathcal{F}$ tel que
$M \otimes_A k$ soit isomorphe à $V$ comme représentation de $\Gamma$.
Puisque le membre de gauche est contenu dans celui de droite,
il suffit de montrer
$\text{length}_A \text{End}_{A[\Gamma]}(M) \leq \text{length}_A A$.
Choisissons des idéaux deux à deux distincts
$(0) = I_n \subset \ldots \subset I_1 \subset A$
avec $n = \text{length}(A)$. En filtrant $M$
de façon correspondante, on voit qu'il suffit de prouver que $\Hom_{A[\Gamma]}(M, I_tM/I_{t + 1}M)$
est de longueur $1$. Puisque $I_tM/I_{t + 1}M \cong M \otimes_A k$
et que toute application de $A[\Gamma]$-modules $M \to M \otimes_A k$ se factorise
de manière unique par l'application quotient $M \to M \otimes_A k$
et donne ainsi un élément de
$$
\text{End}_{A[\Gamma]}(M \otimes_A k) = \text{End}_{k[\Gamma]}(V) = k
$$
on conclut.
\end{proof}



\section{Représentations continues}
\label{examples-defos-section-continuous-representations}

\noindent
Une démarche particulièrement intéressante consiste à prendre un groupe de Galois
infini et à étudier la théorie des déformations de ses représentations ; voir
\cite{Mazur-deforming}.

\begin{example}[Représentations d'un groupe topologique]
\label{examples-defos-example-continuous-representations}
Soit $\Gamma$ un groupe topologique.
Soit $\mathcal{F}$ la catégorie définie comme suit :
\begin{enumerate}
\item un objet est un triplet $(A, M, \rho)$ formé d'un
objet $A$ de $\mathcal{C}_\Lambda$, d'un $A$-module projectif de type fini $M$
et d'un homomorphisme continu $\rho : \Gamma \to \text{GL}_A(M)$,
où $\text{GL}_A(M)$ est muni de la topologie discrète\footnote{Une autre possibilité
serait d'exiger que le $A$-module $M$, muni de l'action de $G$ donnée par $\rho$,
soit un $A\text{-}G$-module au sens de Cohomologie étale, définition
\ref{etale-cohomology-definition-G-module-continuous}. Cependant,
puisque $M$ est un $A$-module de type fini, ces conditions sont équivalentes.} ;
\item un morphisme $(f, g) : (B, N, \tau) \to (A, M, \rho)$ est formé d'un
morphisme $f : B \to A$ dans $\mathcal{C}_\Lambda$ et
d'une application $g : N \to M$ qui est $f$-linéaire, $\Gamma$-équivariante
et induit un isomorphisme $N \otimes_{B, f} A \cong M$.
\end{enumerate}
Le foncteur $p : \mathcal{F} \to \mathcal{C}_\Lambda$ envoie $(A, M, \rho)$
sur $A$ et $(f, g)$ sur $f$. Il est clair que $p$ est cofibré en groupoïdes.
Étant donné un $k$-espace vectoriel de dimension finie $V$ et une
représentation continue $\rho_0 : \Gamma \to \text{GL}_k(V)$,
soit $x_0 = (k, V, \rho_0)$ l'objet correspondant de $\mathcal{F}(k)$.
On pose
$$
\Deformationcategory_{V, \rho_0} = \mathcal{F}_{x_0}
$$
\end{example}

\noindent
Puisque tout module projectif de type fini sur un anneau local est libre
de type fini (Algèbre, lemme \ref{algebra-lemma-finite-projective}),
on a
$$
\begin{matrix}
\text{classes d'isomorphisme} \\
\text{des objets de }\mathcal{F}(A)
\end{matrix}
=
\coprod\nolimits_{n \geq 0}\quad
\begin{matrix}
\text{GL}_n(A)\text{-orbites de conjugaison des}\\
\text{homomorphismes continus }\rho : \Gamma \to \text{GL}_n(A)
\end{matrix}
$$

\begin{lemma}
\label{examples-defos-lemma-continuous-representations-RS}
L'exemple \ref{examples-defos-example-continuous-representations}
satisfait la condition de Rim-Schlessinger (RS).
En particulier, $\Deformationcategory_{V, \rho_0}$ est une catégorie de déformations
pour toute représentation continue de dimension finie
$\rho_0 : \Gamma \to \text{GL}_k(V)$.
\end{lemma}

\begin{proof}
La démonstration est exactement la même que celle du
lemme \ref{examples-defos-lemma-representations-RS}.
\end{proof}

\begin{lemma}
\label{examples-defos-lemma-continuous-representations-TI}
Dans l'exemple \ref{examples-defos-example-continuous-representations}, soit
$\rho_0 : \Gamma \to \text{GL}_k(V)$ une représentation continue
de dimension finie. Alors
$$
T\Deformationcategory_{V, \rho_0} = H^1(\Gamma, \text{End}_k(V))
\quad\text{et}\quad
\text{Inf}(\Deformationcategory_{V, \rho_0}) = H^0(\Gamma, \text{End}_k(V))
$$
Ainsi, $\text{Inf}(\Deformationcategory_{V, \rho_0})$
est toujours de dimension finie,
et $T\Deformationcategory_{V, \rho_0}$ est de dimension finie
si $\Gamma$ est topologiquement de type fini.
\end{lemma}

\begin{proof}
La démonstration est exactement la même que celle du
lemme \ref{examples-defos-lemma-representations-TI}.
\end{proof}

\noindent
Dans l'exemple \ref{examples-defos-example-continuous-representations}, si $\Gamma$
est topologiquement de type fini
et si $(V, \rho_0)$ est une représentation continue de dimension finie de $\Gamma$
sur $k$, alors $\Deformationcategory_{V, \rho_0}$
admet une présentation par un groupoïde lisse proreprésentable en foncteurs
sur $\mathcal{C}_\Lambda$
et possède a fortiori un objet formel versel (minimal). Cela résulte
des lemmes \ref{examples-defos-lemma-continuous-representations-RS} et
\ref{examples-defos-lemma-continuous-representations-TI}
ainsi que de la discussion générale de la section \ref{examples-defos-section-general}.

\begin{lemma}
\label{examples-defos-lemma-continuous-representations-hull}
Dans l'exemple \ref{examples-defos-example-continuous-representations}, supposons que $\Gamma$
soit topologiquement de type fini.
Soit $\rho_0 : \Gamma \to \text{GL}_k(V)$ une représentation de dimension finie.
Supposons que $\Lambda$ soit un anneau local complet de corps résiduel $k$
(cas classique). Alors le foncteur
$$
F : \mathcal{C}_\Lambda \longrightarrow \textit{Ensembles},\quad
A \longmapsto \Ob(\Deformationcategory_{V, \rho_0}(A))/\cong
$$
des classes d'isomorphisme d'objets possède une enveloppe. Si
$H^0(\Gamma, \text{End}_k(V)) = k$, alors $F$ est
proreprésentable.
\end{lemma}

\begin{proof}
La démonstration est exactement la même que celle du
lemme \ref{examples-defos-lemma-representations-hull}.
\end{proof}



\section{Algèbres graduées}
\label{examples-defos-section-graded-algebras}

\noindent
Nous utiliserons l'exemple de cette section dans la démonstration que le champ des
schémas propres polarisés est un champ algébrique. Pour cette raison, nous
considérerons des algèbres graduées commutatives dont les composantes homogènes sont
des modules projectifs de type fini (parfois dits « localement finis »).

\begin{example}[Algèbres graduées]
\label{examples-defos-example-graded-algebras}
Soit $\mathcal{F}$ la catégorie définie comme suit :
\begin{enumerate}
\item un objet est un couple $(A, P)$ formé d'un
objet $A$ de $\mathcal{C}_\Lambda$ et d'une $A$-algèbre graduée $P$
tel que $P_d$ soit un $A$-module projectif de type fini pour tout $d \geq 0$ ;
\item un morphisme $(f, g) : (B, Q) \to (A, P)$ est formé d'un
morphisme $f : B \to A$ dans $\mathcal{C}_\Lambda$ et
d'une application $g : Q \to P$ qui est $f$-linéaire et induit un
isomorphisme $Q \otimes_{B, f} A \cong P$.
\end{enumerate}
Le foncteur $p : \mathcal{F} \to \mathcal{C}_\Lambda$ envoie $(A, P)$
sur $A$ et $(f, g)$ sur $f$. Il est clair que $p$ est cofibré en groupoïdes.
Étant donnée une $k$-algèbre graduée $P$ telle que $\dim_k(P_d) < \infty$ pour tout
$d \geq 0$, soit $x_0 = (k, P)$ l'objet correspondant de $\mathcal{F}(k)$.
On pose
$$
\Deformationcategory_P = \mathcal{F}_{x_0}
$$
\end{example}

\begin{lemma}
\label{examples-defos-lemma-graded-algebras-RS}
L'exemple \ref{examples-defos-example-graded-algebras}
satisfait la condition de Rim-Schlessinger (RS).
En particulier, $\Deformationcategory_P$ est une catégorie de déformations
pour toute $k$-algèbre graduée $P$.
\end{lemma}

\begin{proof}
Soient $A_1 \to A$ et $A_2 \to A$ des morphismes de $\mathcal{C}_\Lambda$.
Supposons $A_2 \to A$ surjectif. D'après Théorie formelle
des déformations, lemme
\ref{formal-defos-lemma-RS-2-categorical}
il suffit de montrer que le foncteur
$\mathcal{F}(A_1 \times_A A_2) \to
\mathcal{F}(A_1) \times_{\mathcal{F}(A)} \mathcal{F}(A_2)$
est une équivalence de catégories.

\medskip\noindent
Considérons un objet
$$
((A_1, P_1), (A_2, P_2), (\text{id}_A, \varphi))
$$
de la catégorie $\mathcal{F}(A_1) \times_{\mathcal{F}(A)} \mathcal{F}(A_2)$.
Considérons alors $P_1 \times_\varphi P_2$. Comme
$\varphi : P_1 \otimes_{A_1} A \to P_2 \otimes_{A_2} A$
est un isomorphisme d'algèbres graduées, on voit que les composantes homogènes
de $P_1 \times_\varphi P_2$ sont des $A_1 \times_A A_2$-modules projectifs de type fini ;
voir la démonstration du lemme \ref{examples-defos-lemma-finite-projective-modules-RS}.
Ainsi, $P_1 \times_\varphi P_2$ est un objet de $\mathcal{F}(A_1 \times_A A_2)$.
Cette construction détermine un quasi-inverse de notre foncteur,
ce qui achève la démonstration.
\end{proof}

\begin{lemma}
\label{examples-defos-lemma-graded-algebras-TI}
Dans l'exemple \ref{examples-defos-example-graded-algebras}, soit $P$ une $k$-algèbre graduée.
Alors
$$
T\Deformationcategory_P
\quad\text{et}\quad
\text{Inf}(\Deformationcategory_P) = \text{Der}_k(P, P)
$$
sont de dimension finie si $P$ est de type fini sur $k$.
\end{lemma}

\begin{proof}
Traitons d'abord les automorphismes infinitésimaux.
Soit $Q = P \otimes_k k[\epsilon]$.
Un élément de $\text{Inf}(\Deformationcategory_P)$
est alors donné par un automorphisme
$\gamma = \text{id} + \epsilon \delta : Q \to Q$
comme ci-dessus, où maintenant $\delta : P \to P$.
Le fait que $\gamma$ soit gradué implique que
$\delta$ est homogène de degré $0$.
Le fait que $\gamma$ soit $k$-linéaire implique que
$\delta$ est $k$-linéaire.
Le fait que $\gamma$ soit multiplicatif implique que
$\delta$ est une $k$-dérivation.
Réciproquement, toute $k$-dérivation $\delta : P \to P$
homogène de degré $0$ fournit un automorphisme
$\gamma = \text{id} + \epsilon \delta$ comme ci-dessus.
On voit donc que
$$
\text{Inf}(\Deformationcategory_P) = \text{Der}_k(P, P)
$$
comme l'affirme le lemme.
Manifestement, si $P$ est engendrée par les composantes $P_i$ de degrés
$0 \leq i \leq N$, alors $\delta$ est déterminée par
les applications linéaires $\delta_i : P_i \to P_i$ pour
$0 \leq i \leq N$, et l'on voit que
$$
\dim_k \text{Der}_k(P, P) < \infty
$$
comme voulu.

\medskip\noindent
Pour achever la démonstration du lemme, montrons que l'espace des déformations
est de dimension finie. À cette fin, choisissons
une présentation
$$
k[X_1, \ldots, X_n]/(F_1, \ldots, F_m) \longrightarrow P
$$
de $k$-algèbres graduées, où $\deg(X_i) = d_i$ et où
$F_j$ est homogène de degré $e_j$.
Soit $Q$ une $k[\epsilon]$-algèbre graduée quelconque,
libre de type fini en chaque degré, munie d'un isomorphisme
$\alpha : Q/\epsilon Q \to P$ tel que $(Q, \alpha)$ définisse
un élément de $T\Deformationcategory_P$.
Choisissons un élément homogène $q_i \in Q$ de degré $d_i$
dont l'image est celle de $X_i$ dans $P$.
On obtient alors
$$
k[\epsilon][X_1, \ldots, X_n] \longrightarrow Q,\quad
X_i \longmapsto q_i
$$
et, puisque $P = Q/\epsilon Q$, cette application est surjective d'après le lemme de Nakayama.
Une chasse au diagramme élémentaire montre que l'on peut choisir des éléments homogènes
$F_{\epsilon, j} \in k[\epsilon][X_1, \ldots, X_n]$ de degré $e_j$
dont l'image dans $Q$ est nulle et dont l'image est $F_j$ dans $k[X_1, \ldots, X_n]$.
Alors
$$
k[\epsilon][X_1, \ldots, X_n]/(F_{\epsilon, 1}, \ldots, F_{\epsilon, m})
\longrightarrow Q
$$
est une présentation de $Q$, par platitude de $Q$ sur $k[\epsilon]$.
Écrivons
$$
F_{\epsilon, j} =  F_j + \epsilon G_j
$$
Le vecteur $(G_1, \ldots, G_m)$ comporte une certaine ambiguïté.
D'abord, en choisissant autrement les $F_{\epsilon, j}$,
on peut modifier $G_j$ par un élément arbitraire de degré $e_j$
du noyau de $k[X_1, \ldots, X_n] \to P$.
Ainsi, au lieu de $(G_1, \ldots, G_m)$, nous retenons
l'élément
$$
(g_1, \ldots, g_m) \in P_{e_1} \oplus \ldots \oplus P_{e_m}
$$
où $g_j$ est l'image de $G_j$ dans $P_{e_j}$.
En outre, si l'on remplace le choix de $q_i$ par $q_i + \epsilon p_i$,
où $p_i$ est de degré $d_i$, un calcul (omis) montre
que $g_j$ devient
$$
g_j^{new} = g_j - \sum\nolimits_{i = 1}^n p_i \partial F_j / \partial X_i
$$
On conclut que la classe d'isomorphisme de $Q$ est déterminée par
l'image du vecteur $(G_1, \ldots, G_m)$ dans le $k$-espace vectoriel
$$
W  = \Coker(P_{d_1} \oplus \ldots \oplus P_{d_n}
\xrightarrow{(\frac{\partial F_j}{\partial X_i})}
P_{e_1} \oplus \ldots \oplus P_{e_m})
$$
On obtient ainsi une injection
$$
T\Deformationcategory_P \longrightarrow W
$$
Puisque $W$ est manifestement de dimension finie, le lemme en résulte.
\end{proof}

\noindent
Dans l'exemple \ref{examples-defos-example-graded-algebras}, si $P$ est une $k$-algèbre graduée
de type fini, alors $\Deformationcategory_P$
admet une présentation par un groupoïde lisse proreprésentable en foncteurs
sur $\mathcal{C}_\Lambda$
et possède a fortiori un objet formel versel (minimal). Cela résulte
des lemmes \ref{examples-defos-lemma-graded-algebras-RS} et
\ref{examples-defos-lemma-graded-algebras-TI}
ainsi que de la discussion générale de la section \ref{examples-defos-section-general}.

\begin{lemma}
\label{examples-defos-lemma-graded-algebras-hull}
Dans l'exemple \ref{examples-defos-example-graded-algebras}, supposons que $P$ soit une
$k$-algèbre graduée de type fini. Supposons que $\Lambda$ soit un anneau local complet
de corps résiduel $k$
(cas classique). Alors le foncteur
$$
F : \mathcal{C}_\Lambda \longrightarrow \textit{Ensembles},\quad
A \longmapsto \Ob(\Deformationcategory_P(A))/\cong
$$
des classes d'isomorphisme d'objets possède une enveloppe.
\end{lemma}

\begin{proof}
Cela résulte immédiatement des lemmes \ref{examples-defos-lemma-graded-algebras-RS} et
\ref{examples-defos-lemma-graded-algebras-TI} ainsi que de
Théorie formelle des déformations, lemme \ref{formal-defos-lemma-RS-implies-S1-S2}
et remarque \ref{formal-defos-remark-compose-minimal-into-iso-classes}.
\end{proof}







\section{Anneaux}
\label{examples-defos-section-rings}

\noindent
La théorie des déformations des anneaux est la même que celle des
schémas affines. Pour les anneaux et les schémas, lorsque nous parlons de déformations,
nous entendons des déformations {\it plates}.

\begin{example}[Anneaux]
\label{examples-defos-example-rings}
Soit $\mathcal{F}$ la catégorie définie comme suit :
\begin{enumerate}
\item un objet est un couple $(A, P)$ formé d'un
objet $A$ de $\mathcal{C}_\Lambda$ et d'une $A$-algèbre plate $P$ ;
\item un morphisme $(f, g) : (B, Q) \to (A, P)$ est formé d'un
morphisme $f : B \to A$ dans $\mathcal{C}_\Lambda$ et
d'une application $g : Q \to P$ qui est $f$-linéaire et induit un
isomorphisme $Q \otimes_{B, f} A \cong P$.
\end{enumerate}
Le foncteur $p : \mathcal{F} \to \mathcal{C}_\Lambda$ envoie $(A, P)$
sur $A$ et $(f, g)$ sur $f$. Il est clair que $p$ est cofibré en groupoïdes.
Étant donnée une $k$-algèbre $P$, soit $x_0 = (k, P)$ l'objet correspondant
de $\mathcal{F}(k)$. On pose
$$
\Deformationcategory_P = \mathcal{F}_{x_0}
$$
\end{example}

\begin{lemma}
\label{examples-defos-lemma-rings-RS}
L'exemple \ref{examples-defos-example-rings}
satisfait la condition de Rim-Schlessinger (RS).
En particulier, $\Deformationcategory_P$ est une catégorie de déformations
pour toute $k$-algèbre $P$.
\end{lemma}

\begin{proof}
Soient $A_1 \to A$ et $A_2 \to A$ des morphismes de $\mathcal{C}_\Lambda$.
Supposons $A_2 \to A$ surjectif. D'après Théorie formelle
des déformations, lemme
\ref{formal-defos-lemma-RS-2-categorical}
il suffit de montrer que le foncteur
$\mathcal{F}(A_1 \times_A A_2) \to
\mathcal{F}(A_1) \times_{\mathcal{F}(A)} \mathcal{F}(A_2)$
est une équivalence de catégories.
C'est un cas particulier de Compléments sur l'algèbre, lemme
\ref{more-algebra-lemma-properties-algebras-over-fibre-product}.
\end{proof}

\begin{lemma}
\label{examples-defos-lemma-rings-TI}
Dans l'exemple \ref{examples-defos-example-rings}, soit $P$ une $k$-algèbre. Alors
$$
T\Deformationcategory_P = \text{Ext}^1_P(\NL_{P/k}, P)
\quad\text{et}\quad
\text{Inf}(\Deformationcategory_P) = \text{Der}_k(P, P)
$$
\end{lemma}

\begin{proof}
Rappelons que $\text{Inf}(\Deformationcategory_P)$ est l'ensemble des
automorphismes de la déformation triviale
$P[\epsilon] = P \otimes_k k[\epsilon]$ de $P$ sur $k[\epsilon]$
qui sont égaux à l'identité modulo $\epsilon$.
D'après Théorie des déformations, lemme \ref{defos-lemma-huge-diagram},
cet ensemble est égal à $\Hom_P(\Omega_{P/k}, P)$, qui est à son tour
égal à $\text{Der}_k(P, P)$ d'après
Algèbre, lemme \ref{algebra-lemma-universal-omega}.

\medskip\noindent
Rappelons que $T\Deformationcategory_P$ est l'ensemble des classes d'isomorphisme
des déformations plates $Q$ de $P$ sur $k[\epsilon]$, plus précisément
l'ensemble des classes d'isomorphisme de $\Deformationcategory_P(k[\epsilon])$.
Rappelons qu'une $k[\epsilon]$-algèbre $Q$ telle que $Q/\epsilon Q = P$
est plate sur $k[\epsilon]$ si et seulement si
$$
0 \to P \xrightarrow{\epsilon} Q \to P \to 0
$$
est exacte. Cela est démontré dans Compléments sur les morphismes, lemme
\ref{more-morphisms-lemma-deform}, et, plus généralement, dans
Théorie des déformations, lemme \ref{defos-lemma-deform-module}.
Nous pouvons donc appliquer
Théorie des déformations, lemme \ref{defos-lemma-choices},
pour voir que l'ensemble des classes d'isomorphisme de telles
déformations est égal à $\text{Ext}^1_P(\NL_{P/k}, P)$.
\end{proof}

\begin{lemma}
\label{examples-defos-lemma-smooth}
Dans l'exemple \ref{examples-defos-example-rings}, soit $P$ une $k$-algèbre lisse. Alors
$T\Deformationcategory_P = (0)$.
\end{lemma}

\begin{proof}
D'après le lemme \ref{examples-defos-lemma-rings-TI}, il faut montrer que
$\text{Ext}^1_P(\NL_{P/k}, P) = (0)$.
Puisque $k \to P$ est lisse, $\NL_{P/k}$ est quasi-isomorphe au
complexe constitué d'un
$P$-module projectif de type fini placé en degré $0$.
\end{proof}

\begin{lemma}
\label{examples-defos-lemma-finite-type-rings-TI}
Dans le lemme \ref{examples-defos-lemma-rings-TI}, si $P$ est une $k$-algèbre de type fini, alors
\begin{enumerate}
\item $\text{Inf}(\Deformationcategory_P)$ est de dimension finie si et seulement si
$\dim(P) = 0$ ;
\item $T\Deformationcategory_P$ est de dimension finie si
$\Spec(P) \to \Spec(k)$ est lisse sauf en un nombre fini de points.
\end{enumerate}
\end{lemma}

\begin{proof}
Preuve de (1). Considérons $\text{Der}_k(P, P)$ comme un $P$-module.
S'il est de dimension finie sur $k$, il est alors de longueur finie
comme $P$-module ; son support est donc contenu dans un nombre fini de
points fermés de $\Spec(P)$
(Algèbre, lemme \ref{algebra-lemma-simple-pieces}).
Puisque $\text{Der}_k(P, P) = \Hom_P(\Omega_{P/k}, P)$,
nous voyons que
$\text{Der}_k(P, P)_\mathfrak p = \text{Der}_k(P_\mathfrak p, P_\mathfrak p)$
pour tout idéal premier $\mathfrak p \subset P$
(cela utilise Algèbre, lemmes
\ref{algebra-lemma-differentials-localize},
\ref{algebra-lemma-differentials-finitely-presented}, et
\ref{algebra-lemma-hom-from-finitely-presented}).
Soit $\mathfrak p$ un idéal premier minimal de $P$
correspondant à une composante irréductible de dimension $d > 0$.
Alors $P_\mathfrak p$ est un anneau local artinien
essentiellement de type fini sur $k$, avec corps résiduel,
et $\Omega_{P_\mathfrak p/k}$ est non nul, par exemple d'après
Algèbre, lemme \ref{algebra-lemma-characterize-smooth-over-field}.
Tout module fini non nul sur un anneau local artinien
possède à la fois un sous-module et un module quotient isomorphes au corps résiduel.
Nous trouvons ainsi que
$\text{Der}_k(P_\mathfrak p, P_\mathfrak p) =
\Hom_{P_\mathfrak p}(\Omega_{P_\mathfrak p/k}, P_\mathfrak p)$
est lui aussi non nul. En combinant tout ce qui précède, nous obtenons (1).

\medskip\noindent
Preuve de (2). Pour un idéal premier $\mathfrak p$ de $P$, nous utiliserons l'égalité
$\NL_{P_\mathfrak p/k} = (\NL_{P/k})_\mathfrak p$
(Algèbre, lemme \ref{algebra-lemma-localize-NL})
ainsi que
l'égalité
$\text{Ext}_P^1(\NL_{P/k}, P)_\mathfrak p =
\text{Ext}_{P_\mathfrak p}^1(\NL_{P_\mathfrak p/k}, P_\mathfrak p)$
(Compléments sur l'algèbre, lemme
\ref{more-algebra-lemma-pseudo-coherence-and-base-change-ext}).
Étant donné un idéal premier $\mathfrak p \subset P$,
le morphisme $k \to P$ est lisse en $\mathfrak p$ si et seulement si
$(\NL_{P/k})_\mathfrak p$ est quasi-isomorphe
à un module projectif de type fini placé en degré $0$ (cela résulte
immédiatement de la définition d'un homomorphisme d'anneaux lisse, mais aussi
du résultat plus fort d'Algèbre, lemme \ref{algebra-lemma-smooth-at-point}).

\medskip\noindent
Supposons que $P$ soit lisse sur $k$ en dehors d'un nombre fini d'idéaux premiers.
Ces idéaux premiers « mauvais » sont alors les idéaux maximaux
$\mathfrak m_1, \ldots, \mathfrak m_n \subset P$, d'après
Algèbre, lemme \ref{algebra-lemma-finite-type-algebra-finite-nr-primes},
et parce que les idéaux premiers « mauvais » forment un fermé de $\Spec(P)$.
Pour $\mathfrak p \not \in \{\mathfrak m_1, \ldots, \mathfrak m_n\}$,
nous avons $\text{Ext}^1_P(\NL_{P/k}, P)_\mathfrak p = 0$ d'après ce qui précède.
Ainsi, $\text{Ext}^1_P(\NL_{P/k}, P)$ est un $P$-module fini
dont le support est contenu dans $\{\mathfrak m_1, \ldots, \mathfrak m_r\}$.
D'après Algèbre, proposition
\ref{algebra-proposition-minimal-primes-associated-primes}
par exemple, nous trouvons que la dimension sur $k$ de
$\text{Ext}^1_P(\NL_{P/k}, P)$ est une combinaison linéaire entière finie
des $\dim_k \kappa(\mathfrak m_i)$ ; elle est donc finie d'après
le théorème des zéros de Hilbert
(Algèbre, théorème \ref{algebra-theorem-nullstellensatz}).
\end{proof}

\noindent
Dans l'exemple \ref{examples-defos-example-rings}, soit $P$ une
$k$-algèbre de type fini. Alors $\Deformationcategory_P$
admet une présentation par un groupoïde lisse proreprésentable en foncteurs
sur $\mathcal{C}_\Lambda$ si et seulement si $\dim(P) = 0$.
De plus, $\Deformationcategory_P$ possède un objet
formel versel si $\Spec(P) \to \Spec(k)$ n'a qu'un nombre fini de
points singuliers. Cela résulte des lemmes \ref{examples-defos-lemma-rings-RS} et
\ref{examples-defos-lemma-finite-type-rings-TI},
ainsi que de la discussion générale de la section \ref{examples-defos-section-general}.

\begin{lemma}
\label{examples-defos-lemma-rings-hull}
Dans l'exemple \ref{examples-defos-example-rings}, supposons que $P$ soit une
$k$-algèbre de type fini telle que $\Spec(P) \to \Spec(k)$ soit lisse sauf
en un nombre fini de points.
Supposons que $\Lambda$ soit un anneau local complet de corps résiduel $k$
(cas classique). Alors le foncteur
$$
F : \mathcal{C}_\Lambda \longrightarrow \textit{Ensembles},\quad
A \longmapsto \Ob(\Deformationcategory_P(A))/\cong
$$
des classes d'isomorphisme d'objets possède une enveloppe.
\end{lemma}

\begin{proof}
Cela résulte immédiatement des lemmes \ref{examples-defos-lemma-rings-RS} et
\ref{examples-defos-lemma-finite-type-rings-TI}, ainsi que de
Théorie formelle des déformations, lemme \ref{formal-defos-lemma-RS-implies-S1-S2}
et remarque \ref{formal-defos-remark-compose-minimal-into-iso-classes}.
\end{proof}

\begin{lemma}
\label{examples-defos-lemma-localization}
Dans l'exemple \ref{examples-defos-example-rings}, soit $P$ une $k$-algèbre.
Soit $S \subset P$ une partie multiplicative. Il existe un foncteur naturel
$$
\Deformationcategory_P \longrightarrow \Deformationcategory_{S^{-1}P}
$$
de catégories de déformations.
\end{lemma}

\begin{proof}
Étant donnée une déformation de $P$, nous pouvons la localiser
pour obtenir une déformation de la localisation ; cela est
clair, et nous invitons le lecteur à omettre la démonstration. Plus précisément,
soit $(A, Q) \to (k, P)$ un morphisme de $\mathcal{F}$, c'est-à-dire
un objet de $\Deformationcategory_P$. Soit $S_Q \subset Q$ l'image
réciproque de $S$. Alors
Ainsi, $(A, S_Q^{-1}Q) \to (k, S^{-1}P)$
est l'objet recherché de $\Deformationcategory_{S^{-1}P}$.
\end{proof}

\begin{lemma}
\label{examples-defos-lemma-henselization}
Dans l'exemple \ref{examples-defos-example-rings}, soit $P$ une $k$-algèbre.
Soit $J \subset P$ un idéal.
Notons $(P^h, J^h)$ l'hensélisation du couple $(P, J)$.
Il existe un foncteur naturel
$$
\Deformationcategory_P \longrightarrow \Deformationcategory_{P^h}
$$
de catégories de déformations.
\end{lemma}

\begin{proof}
Étant donnée une déformation de $P$, nous pouvons en prendre l'hensélisation
pour obtenir une déformation de l'hensélisé ; cela est
clair, et nous invitons le lecteur à omettre la démonstration. Plus précisément,
soit $(A, Q) \to (k, P)$ un morphisme de $\mathcal{F}$, c'est-à-dire
un objet de $\Deformationcategory_P$. Notons $J_Q \subset Q$ l'image
réciproque de $J$ dans $Q$. Soit $(Q^h, J_Q^h)$ l'hensélisation
du couple $(Q, J_Q)$. Rappelons que $Q \to Q^h$ est plat
(Compléments sur l'algèbre, lemme \ref{more-algebra-lemma-henselization-flat})
et que $Q^h$ est donc plat sur $A$.
D'après Compléments sur l'algèbre, lemme \ref{more-algebra-lemma-henselization-integral},
nous voyons que l'application $Q^h \to P^h$ induit un isomorphisme
$Q^h \otimes_A k = Q^h \otimes_Q P = P^h$.
Ainsi, $(A, Q^h) \to (k, P^h)$ est l'objet recherché de
$\Deformationcategory_{P^h}$.
\end{proof}

\begin{lemma}
\label{examples-defos-lemma-strict-henselization}
Dans l'exemple \ref{examples-defos-example-rings}, soit $P$ une $k$-algèbre.
Supposons que $P$ soit un anneau local et soit $P^{sh}$ un hensélisé strict de $P$.
Il existe un foncteur naturel
$$
\Deformationcategory_P \longrightarrow \Deformationcategory_{P^{sh}}
$$
de catégories de déformations.
\end{lemma}

\begin{proof}
Étant donnée une déformation de $P$, nous pouvons en prendre l'hensélisé strict
pour obtenir une déformation de l'hensélisé strict ; cela est
clair, et nous invitons le lecteur à omettre la démonstration. Plus précisément,
soit $(A, Q) \to (k, P)$ un morphisme de $\mathcal{F}$, c'est-à-dire
un objet de $\Deformationcategory_P$. Puisque le noyau de la surjection
$Q \to P$ est nilpotent, nous trouvons que $Q$ est un anneau local de
même corps résiduel que $P$. Soit $Q^{sh}$ l'hensélisé strict
de $Q$. Rappelons que $Q \to Q^{sh}$ est plat
(Compléments sur l'algèbre, lemme \ref{more-algebra-lemma-dumb-properties-henselization})
et que $Q^{sh}$ est donc plat sur $A$.
D'après Algèbre, lemme \ref{algebra-lemma-quotient-strict-henselization},
nous voyons que l'application $Q^{sh} \to P^{sh}$ induit un isomorphisme
$Q^{sh} \otimes_A k = Q^{sh} \otimes_Q P = P^{sh}$.
Ainsi, $(A, Q^{sh}) \to (k, P^{sh})$ est l'objet recherché de
$\Deformationcategory_{P^{sh}}$.
\end{proof}

\begin{lemma}
\label{examples-defos-lemma-completion}
Dans l'exemple \ref{examples-defos-example-rings}, soit $P$ une $k$-algèbre.
Supposons $P$ noethérien et soit $J \subset P$ un idéal.
Notons $P^\wedge$ le complété $J$-adique.
Il existe un foncteur naturel
$$
\Deformationcategory_P \longrightarrow \Deformationcategory_{P^\wedge}
$$
de catégories de déformations.
\end{lemma}

\begin{proof}
Étant donnée une déformation de $P$, nous pouvons en prendre le complété
pour obtenir une déformation du complété ; cela est
clair, et nous invitons le lecteur à omettre la démonstration. Plus précisément,
soit $(A, Q) \to (k, P)$ un morphisme de $\mathcal{F}$, c'est-à-dire
un objet de $\Deformationcategory_P$. Observons que $Q$ est un anneau
noethérien : le noyau de l'homomorphisme surjectif d'anneaux $Q \to P$ est
nilpotent et de type fini, et $P$ est noethérien ; appliquons
Algèbre, lemme \ref{algebra-lemma-completion-Noetherian}.
Notons $J_Q \subset Q$ l'image
réciproque de $J$ dans $Q$. Soit $Q^\wedge$ le complété $J_Q$-adique de $Q$.
Rappelons que $Q \to Q^\wedge$ est plat
(Algèbre, lemme \ref{algebra-lemma-completion-flat})
et que $Q^\wedge$ est donc plat sur $A$.
L'application induite $Q^\wedge \to P^\wedge$ induit un isomorphisme
$Q^\wedge \otimes_A k = Q^\wedge \otimes_Q P = P^\wedge$ d'après
Algèbre, lemme \ref{algebra-lemma-completion-tensor}, par exemple.
Ainsi, $(A, Q^\wedge) \to (k, P^\wedge)$
est l'objet recherché de $\Deformationcategory_{P^\wedge}$.
\end{proof}

\begin{lemma}
\label{examples-defos-lemma-power-series-rings-TI}
Dans le lemme \ref{examples-defos-lemma-rings-TI}, si $P = k[[x_1, \ldots, x_n]]/(f)$
pour un certain $f \in (x_1, \ldots, x_n)^2$ non nul, alors
\begin{enumerate}
\item $\text{Inf}(\Deformationcategory_P)$ est de dimension finie
si et seulement si $n = 1$ ;
\item $T\Deformationcategory_P$ est de dimension finie si
$$
\sqrt{(f, \partial f/\partial x_1, \ldots,  \partial f/\partial x_n)} =
(x_1, \ldots, x_n)
$$
\end{enumerate}
\end{lemma}

\begin{proof}
Preuve de (1). Considérons les dérivations $\partial/\partial x_i$ de
$k[[x_1, \ldots, x_n]]$ sur $k$. Posons $f_i = \partial f/\partial x_i$.
La dérivation
$$
\theta = \sum h_i \partial/\partial x_i
$$
de $k[[x_1, \ldots, x_n]]$
induit une dérivation de $P = k[[x_1, \ldots, x_n]]/(f)$
si et seulement si
$\sum h_i f_i \in (f)$. De plus, la dérivation induite de $P$
est nulle si et seulement si $h_i \in (f)$ pour $i = 1, \ldots, n$.
Nous obtenons ainsi
$$
\Ker((f_1, \ldots, f_n) : P^{\oplus n} \longrightarrow P) \subset
\text{Der}_k(P, P)
$$
Le membre de gauche n'est un $k$-espace vectoriel de dimension finie que si
$n = 1$ ; nous omettons la démonstration. Nous laissons aussi au lecteur le soin de voir
que le membre de droite est de dimension finie si $n = 1$.
Cela démontre (1).

\medskip\noindent
Preuve de (2). Soit $Q$ une déformation plate de $P$ sur $k[\epsilon]$,
comme dans la démonstration du lemme \ref{examples-defos-lemma-rings-TI}. Choisissons des relèvements $q_i \in Q$
de l'image de $x_i$ dans $P$. Alors $Q$ est un anneau local complet
dont l'idéal maximal est engendré par $q_1, \ldots, q_n$ et $\epsilon$
(nous omettons un court argument). Nous obtenons ainsi une surjection
$$
k[\epsilon][[x_1, \ldots, x_n]] \longrightarrow Q,\quad
x_i \longmapsto q_i
$$
Choisissons un élément de la forme
$f + \epsilon g \in k[\epsilon][[x_1, \ldots, x_n]]$
qui s'envoie sur zéro dans $Q$. Observons que $g$ est bien défini modulo $(f)$.
Puisque $Q$ est plat sur $k[\epsilon]$, nous obtenons
$$
Q = k[\epsilon][[x_1, \ldots, x_n]]/(f + \epsilon g)
$$
Enfin, modifier le choix de $q_i$ revient à
remplacer les coordonnées $x_i$ par $x_i + \epsilon h_i$
pour certains $h_i \in k[[x_1, \ldots, x_n]]$. Alors
$f + \epsilon g$ devient $f + \epsilon (g + \sum h_i f_i)$,
où $f_i = \partial f/\partial x_i$. Nous voyons donc que la
classe d'isomorphisme de la déformation $Q$ est déterminée
par un élément de
$$
k[[x_1, \ldots, x_n]]/
(f, \partial f/\partial x_1, \ldots,  \partial f/\partial x_n)
$$
Ce quotient est de dimension finie sur $k$ si et seulement si
son support est le point fermé de $k[[x_1, \ldots, x_n]]$,
si et seulement si
$\sqrt{(f, \partial f/\partial x_1, \ldots,  \partial f/\partial x_n)} =
(x_1, \ldots, x_n)$.
\end{proof}






\section{Schémas}
\label{examples-defos-section-schemes}

\noindent
La théorie des déformations des schémas.

\begin{example}[Schémas]
\label{examples-defos-example-schemes}
Soit $\mathcal{F}$ la catégorie définie comme suit :
\begin{enumerate}
\item un objet est un couple $(A, X)$ formé d'un
objet $A$ de $\mathcal{C}_\Lambda$ et d'un schéma $X$ plat sur $A$ ;
\item un morphisme $(f, g) : (B, Y) \to (A, X)$ est formé d'un
morphisme $f : B \to A$ dans $\mathcal{C}_\Lambda$ et
d'un morphisme $g : X \to Y$ tel que
$$
\xymatrix{
X \ar[r]_g \ar[d] & Y \ar[d] \\
\Spec(A) \ar[r]^f & \Spec(B)
}
$$
soit un diagramme commutatif cartésien de schémas.
\end{enumerate}
Le foncteur $p : \mathcal{F} \to \mathcal{C}_\Lambda$ envoie $(A, X)$
sur $A$ et $(f, g)$ sur $f$. Il est clair que $p$ est cofibré en groupoïdes.
Étant donné un schéma $X$ sur $k$, soit $x_0 = (k, X)$ l'objet correspondant
de $\mathcal{F}(k)$. On pose
$$
\Deformationcategory_X = \mathcal{F}_{x_0}
$$
\end{example}

\begin{lemma}
\label{examples-defos-lemma-schemes-RS}
L'exemple \ref{examples-defos-example-schemes}
satisfait la condition de Rim-Schlessinger (RS).
En particulier, $\Deformationcategory_X$ est une catégorie de déformations
pour tout schéma $X$ sur $k$.
\end{lemma}

\begin{proof}
Soient $A_1 \to A$ et $A_2 \to A$ des morphismes de $\mathcal{C}_\Lambda$.
Supposons $A_2 \to A$ surjectif. D'après Théorie formelle
des déformations, lemme
\ref{formal-defos-lemma-RS-2-categorical}
il suffit de montrer que le foncteur
$\mathcal{F}(A_1 \times_A A_2) \to
\mathcal{F}(A_1) \times_{\mathcal{F}(A)} \mathcal{F}(A_2)$
est une équivalence de catégories.
Observons que
$$
\xymatrix{
\Spec(A) \ar[r] \ar[d] & \Spec(A_2) \ar[d] \\
\Spec(A_1) \ar[r] &
\Spec(A_1 \times_A A_2)
}
$$
est un diagramme cocartésien comme dans Compléments sur les morphismes, lemme
\ref{more-morphisms-lemma-pushout-along-thickening}.
Le lemme est donc un cas particulier de Compléments sur les morphismes, lemme
\ref{more-morphisms-lemma-equivalence-categories-schemes-over-pushout-flat}.
\end{proof}

\begin{lemma}
\label{examples-defos-lemma-schemes-TI}
Dans l'exemple \ref{examples-defos-example-schemes}, soit $X$ un schéma sur $k$. Alors
$$
\text{Inf}(\Deformationcategory_X) =
\text{Ext}^0_{\mathcal{O}_X}(\NL_{X/k}, \mathcal{O}_X) =
\Hom_{\mathcal{O}_X}(\Omega_{X/k}, \mathcal{O}_X) =
\text{Der}_k(\mathcal{O}_X, \mathcal{O}_X)
$$
et
$$
T\Deformationcategory_X =
\text{Ext}^1_{\mathcal{O}_X}(\NL_{X/k}, \mathcal{O}_X)
$$
\end{lemma}

\begin{proof}
Rappelons que $\text{Inf}(\Deformationcategory_X)$ est l'ensemble des
automorphismes de la déformation triviale
$X' = X \times_{\Spec(k)} \Spec(k[\epsilon])$ de $X$ sur $k[\epsilon]$
qui sont égaux à l'identité modulo $\epsilon$.
D'après Théorie des déformations, lemme \ref{defos-lemma-deform},
cet ensemble est égal à $\text{Ext}^0_{\mathcal{O}_X}(\NL_{X/k}, \mathcal{O}_X)$.
L'égalité $\text{Ext}^0_{\mathcal{O}_X}(\NL_{X/k}, \mathcal{O}_X) =
\Hom_{\mathcal{O}_X}(\Omega_{X/k}, \mathcal{O}_X)$ résulte de
Compléments sur les morphismes, lemme
\ref{more-morphisms-lemma-netherlander-quasi-coherent}.
L'égalité
$\Hom_{\mathcal{O}_X}(\Omega_{X/k}, \mathcal{O}_X) =
\text{Der}_k(\mathcal{O}_X, \mathcal{O}_X)$
résulte de Morphismes, lemme
\ref{morphisms-lemma-universal-derivation-universal}.

\medskip\noindent
Rappelons que $T_{x_0}\Deformationcategory_X$ est l'ensemble des classes d'isomorphisme
des déformations plates $X'$ de $X$ sur $k[\epsilon]$, plus précisément
l'ensemble des classes d'isomorphisme de $\Deformationcategory_X(k[\epsilon])$.
La seconde assertion du lemme résulte donc de
Théorie des déformations, lemme \ref{defos-lemma-deform}.
\end{proof}

\begin{lemma}
\label{examples-defos-lemma-proper-schemes-TI}
Dans le lemme \ref{examples-defos-lemma-schemes-TI}, si $X$ est propre sur $k$, alors
$\text{Inf}(\Deformationcategory_X)$ et $T\Deformationcategory_X$ sont
de dimension finie.
\end{lemma}

\begin{proof}
D'après le lemme, il faut montrer que
$\Ext^1_{\mathcal{O}_X}(\NL_{X/k}, \mathcal{O}_X)$ et
$\Ext^0_{\mathcal{O}_X}(\NL_{X/k}, \mathcal{O}_X)$ sont de
dimension finie. D'après Compléments sur les morphismes, lemme
\ref{more-morphisms-lemma-netherlander-fp}
et le fait que $X$ est noethérien, nous voyons que
$\NL_{X/k}$ a des faisceaux de cohomologie cohérents, nuls sauf
en degrés $0$ et $-1$.
D'après Catégories dérivées de schémas, lemme \ref{perfect-lemma-ext-finite},
les groupes $\Ext$ affichés sont des $k$-espaces vectoriels de dimension finie,
ce qui achève la démonstration.
\end{proof}

\noindent
Dans l'exemple \ref{examples-defos-example-schemes}, si $X$ est un schéma propre sur $k$,
alors $\Deformationcategory_X$
admet une présentation par un groupoïde lisse proreprésentable en foncteurs
sur $\mathcal{C}_\Lambda$
et possède a fortiori un objet formel versel (minimal). Cela résulte
des lemmes \ref{examples-defos-lemma-schemes-RS} et
\ref{examples-defos-lemma-proper-schemes-TI},
ainsi que de la discussion générale de la section \ref{examples-defos-section-general}.

\begin{lemma}
\label{examples-defos-lemma-schemes-hull}
Dans l'exemple \ref{examples-defos-example-schemes}, supposons que $X$ soit un $k$-schéma propre.
Supposons que $\Lambda$ soit un anneau local complet de corps résiduel $k$
(cas classique). Alors le foncteur
$$
F : \mathcal{C}_\Lambda \longrightarrow \textit{Ensembles},\quad
A \longmapsto \Ob(\Deformationcategory_X(A))/\cong
$$
des classes d'isomorphisme d'objets possède une enveloppe. Si
$\text{Der}_k(\mathcal{O}_X, \mathcal{O}_X) = 0$, alors
$F$ est proreprésentable.
\end{lemma}

\begin{proof}
L'existence d'une enveloppe résulte immédiatement des
lemmes \ref{examples-defos-lemma-schemes-RS} et \ref{examples-defos-lemma-proper-schemes-TI}, ainsi que de
Théorie formelle des déformations, lemme \ref{formal-defos-lemma-RS-implies-S1-S2}
et remarque \ref{formal-defos-remark-compose-minimal-into-iso-classes}.

\medskip\noindent
Supposons $\text{Der}_k(\mathcal{O}_X, \mathcal{O}_X) = 0$. Alors
$\Deformationcategory_X$ et $F$ sont équivalents d'après
Théorie formelle des déformations, lemme \ref{formal-defos-lemma-infdef-trivial}.
Ainsi, $F$ est un foncteur de déformations (car $\Deformationcategory_X$ est une
catégorie de déformations) d'espace tangent de dimension finie, et nous pouvons appliquer
Théorie formelle des déformations, théorème
\ref{formal-defos-theorem-Schlessinger-prorepresentability}.
\end{proof}

\begin{lemma}
\label{examples-defos-lemma-open}
Dans l'exemple \ref{examples-defos-example-schemes}, soit $X$ un schéma sur $k$.
Soit $U \subset X$ un sous-schéma ouvert.
Il existe un foncteur naturel
$$
\Deformationcategory_X \longrightarrow \Deformationcategory_U
$$
de catégories de déformations.
\end{lemma}

\begin{proof}
Étant donnée une déformation de $X$, nous pouvons en prendre l'ouvert correspondant
pour obtenir une déformation de $U$. Nous omettons les détails.
\end{proof}

\begin{lemma}
\label{examples-defos-lemma-affine}
Dans l'exemple \ref{examples-defos-example-schemes}, soit $X = \Spec(P)$ un
schéma affine sur $k$. Avec $\Deformationcategory_P$ comme dans
l'exemple \ref{examples-defos-example-rings}, il existe une équivalence naturelle
$$
\Deformationcategory_X \longrightarrow \Deformationcategory_P
$$
de catégories de déformations.
\end{lemma}

\begin{proof}
Le foncteur envoie $(A, Y)$ sur $\Gamma(Y, \mathcal{O}_Y)$.
Cette construction convient parce que
toute déformation de $X$ est affine d'après
Compléments sur les morphismes, lemme \ref{more-morphisms-lemma-thickening-affine-scheme}.
\end{proof}

\begin{lemma}
\label{examples-defos-lemma-local-ring}
Dans l'exemple \ref{examples-defos-example-schemes}, soit $X$ un schéma sur $k$.
Soit $p \in X$ un point. Avec $\Deformationcategory_{\mathcal{O}_{X, p}}$
comme dans l'exemple \ref{examples-defos-example-rings}, il existe un foncteur naturel
$$
\Deformationcategory_X
\longrightarrow
\Deformationcategory_{\mathcal{O}_{X, p}}
$$
de catégories de déformations.
\end{lemma}

\begin{proof}
Choisissons un ouvert affine $U = \Spec(P) \subset X$ contenant $p$.
Alors $\mathcal{O}_{X, p}$ est une localisation de $P$. Composons
les foncteurs des
lemmes \ref{examples-defos-lemma-open}, \ref{examples-defos-lemma-affine} et \ref{examples-defos-lemma-localization}.
\end{proof}

\begin{situation}
\label{examples-defos-situation-glueing}
Soit $\Lambda \to k$ comme dans la section \ref{examples-defos-section-general}.
Soit $X$ un schéma sur $k$ qui possède un recouvrement ouvert affine
$X = U_1 \cup U_2$, où $U_{12} = U_1 \cap U_2$ est lui aussi affine.
Écrivons $U_1 = \Spec(P_1)$, $U_2 = \Spec(P_2)$ et $U_{12} = \Spec(P_{12})$.
Soient $\Deformationcategory_X$, $\Deformationcategory_{U_1}$,
$\Deformationcategory_{U_2}$ et $\Deformationcategory_{U_{12}}$
comme dans l'exemple \ref{examples-defos-example-schemes}, et soient
$\Deformationcategory_{P_1}$, $\Deformationcategory_{P_2}$ et
$\Deformationcategory_{P_{12}}$ comme dans l'exemple \ref{examples-defos-example-rings}.
\end{situation}

\begin{lemma}
\label{examples-defos-lemma-glueing}
Dans la situation \ref{examples-defos-situation-glueing},
il existe une équivalence
$$
\Deformationcategory_X =
\Deformationcategory_{P_1}
\times_{\Deformationcategory_{P_{12}}}
\Deformationcategory_{P_2}
$$
de catégories de déformations ; voir les exemples \ref{examples-defos-example-schemes} et
\ref{examples-defos-example-rings}.
\end{lemma}

\begin{proof}
Il suffit de montrer que les foncteurs du lemme \ref{examples-defos-lemma-open}
définissent une équivalence
$$
\Deformationcategory_X \longrightarrow
\Deformationcategory_{U_1}
\times_{\Deformationcategory_{U_{12}}}
\Deformationcategory_{U_2}
$$
car nous pouvons alors appliquer le lemme \ref{examples-defos-lemma-affine} pour passer aux anneaux.
Pour cela, construisons un quasi-inverse. Notons
$F_i : \Deformationcategory_{U_i} \to \Deformationcategory_{U_{12}}$
le foncteur du lemme \ref{examples-defos-lemma-open}.
Un objet du membre de droite est donné par un $A$ dans $\mathcal{C}_\Lambda$,
des objets $(A, V_1) \to (k, U_1)$ et $(A, V_2) \to (k, U_2)$, et
un morphisme
$$
g : F_1(A, V_1) \to F_2(A, V_2)
$$
Ici, $F_i(A, V_i) = (A, V_{i, 3 - i})$, où $V_{i, 3 - i} \subset V_i$
est le sous-schéma ouvert dont le changement de base à $k$ est $U_{12} \subset U_i$.
Le morphisme $g$ définit un isomorphisme
$V_{1, 2} \to V_{2, 1}$ de schémas sur $A$, compatible
avec $\text{id} : U_{12} \to U_{12}$ sur $k$.
Ainsi, $(\{1, 2\}, V_i, V_{i, 3 - i}, g, g^{-1})$ est une donnée
de recollement comme dans Schémas, section \ref{schemes-section-glueing-schemes}.
Soit $Y$ le schéma recollé ; voir Schémas, lemme \ref{schemes-lemma-glue}.
Alors $Y$ est un schéma sur $A$, et les
compatibilités mentionnées ci-dessus montrent
qu'il existe un isomorphisme canonique
$Y \times_{\Spec(A)} \Spec(k) = X$.
Ainsi, $(A, Y) \to (k, X)$ est un objet de $\Deformationcategory_X$.
Nous omettons de vérifier que cette construction définit un foncteur
quasi-inverse de celui qui était donné.
\end{proof}






\section{Morphismes de schémas}
\label{examples-defos-section-schemes-morphisms}

\noindent
La théorie des déformations des morphismes de schémas.
Il ne s'agit bien sûr que d'un exemple de
déformations de diagrammes de schémas.

\begin{example}[Morphismes de schémas]
\label{examples-defos-example-schemes-morphisms}
Soit $\mathcal{F}$ la catégorie définie comme suit :
\begin{enumerate}
\item un objet est un couple $(A, X \to Y)$ formé d'un
objet $A$ de $\mathcal{C}_\Lambda$ et d'un morphisme
$X \to Y$ de schémas sur $A$, où $X$ et $Y$ sont tous deux plats sur $A$ ;
\item un morphisme $(f, g, h) : (A', X' \to Y') \to (A, X \to Y)$ est formé d'un
morphisme $f : A' \to A$ dans $\mathcal{C}_\Lambda$ et
de morphismes de schémas $g : X \to X'$ et $h : Y \to Y'$ tels que
$$
\xymatrix{
X \ar[r]_g \ar[d] & X' \ar[d] \\
Y \ar[r]_h \ar[d] & Y' \ar[d] \\
\Spec(A) \ar[r]^f & \Spec(A')
}
$$
soit un diagramme commutatif de schémas dont les deux carrés sont cartésiens.
\end{enumerate}
Le foncteur $p : \mathcal{F} \to \mathcal{C}_\Lambda$ envoie $(A, X \to Y)$
sur $A$ et $(f, g, h)$ sur $f$. Il est clair que $p$ est cofibré en groupoïdes.
Étant donné un morphisme de schémas $X \to Y$ sur $k$, soit $x_0 = (k, X \to Y)$
l'objet correspondant de $\mathcal{F}(k)$. On pose
$$
\Deformationcategory_{X \to Y} = \mathcal{F}_{x_0}
$$
\end{example}

\begin{lemma}
\label{examples-defos-lemma-schemes-morphisms-RS}
L'exemple \ref{examples-defos-example-schemes-morphisms}
satisfait la condition de Rim-Schlessinger (RS).
En particulier, $\Deformationcategory_{X \to Y}$ est une catégorie de déformations
pour tout morphisme de schémas $X \to Y$ sur $k$.
\end{lemma}

\begin{proof}
Soient $A_1 \to A$ et $A_2 \to A$ des morphismes de $\mathcal{C}_\Lambda$.
Supposons $A_2 \to A$ surjectif. D'après Théorie formelle
des déformations, lemme
\ref{formal-defos-lemma-RS-2-categorical}
il suffit de montrer que le foncteur
$\mathcal{F}(A_1 \times_A A_2) \to
\mathcal{F}(A_1) \times_{\mathcal{F}(A)} \mathcal{F}(A_2)$
est une équivalence de catégories.
Observons que
$$
\xymatrix{
\Spec(A) \ar[r] \ar[d] & \Spec(A_2) \ar[d] \\
\Spec(A_1) \ar[r] &
\Spec(A_1 \times_A A_2)
}
$$
est un diagramme cocartésien comme dans Compléments sur les morphismes, lemme
\ref{more-morphisms-lemma-pushout-along-thickening}.
Le lemme résulte donc immédiatement de
Compléments sur les morphismes, lemme
\ref{more-morphisms-lemma-equivalence-categories-schemes-over-pushout-flat}
car ce résultat décrit la catégorie des schémas plats sur $A_1 \times_A A_2$
comme le produit fibré de la catégorie des schémas plats sur $A_1$
et de la catégorie des schémas plats sur $A_2$ au-dessus de la catégorie des
schémas plats sur $A$.
\end{proof}

\begin{lemma}
\label{examples-defos-lemma-schemes-morphisms-TI}
Dans l'exemple \ref{examples-defos-example-schemes}, soit $f : X \to Y$ un morphisme de schémas
sur $k$. Il existe une suite exacte canonique de $k$-espaces vectoriels
$$
\xymatrix{
0 \ar[r] &
\text{Inf}(\Deformationcategory_{X \to Y}) \ar[r] &
\text{Inf}(\Deformationcategory_X \times \Deformationcategory_Y) \ar[r] &
\text{Der}_k(\mathcal{O}_Y, f_*\mathcal{O}_X) \ar[lld] \\
& T\Deformationcategory_{X \to Y} \ar[r] &
T(\Deformationcategory_X \times \Deformationcategory_Y) \ar[r] &
\text{Ext}^1_{\mathcal{O}_X}(Lf^*\NL_{Y/k}, \mathcal{O}_X)
}
$$
\end{lemma}

\begin{proof}
Le morphisme évident de catégories de déformations
$\Deformationcategory_{X \to Y} \to
\Deformationcategory_X \times \Deformationcategory_Y$
donne deux des flèches de la suite exacte du lemme.
Rappelons que $\text{Inf}(\Deformationcategory_{X \to Y})$
est l'ensemble des automorphismes de la déformation triviale
$$
f' :  X' = X \times_{\Spec(k)} \Spec(k[\epsilon])
\xrightarrow{f \times \text{id}}
Y' = Y \times_{\Spec(k)} \Spec(k[\epsilon])
$$
de $X \to Y$ sur $k[\epsilon]$ qui sont égaux à l'identité modulo $\epsilon$.
Il s'agit clairement des couples
$(\alpha, \beta) \in
\text{Inf}(\Deformationcategory_X \times \Deformationcategory_Y)$
d'automorphismes infinitésimaux de $X$ et $Y$ compatibles avec $f'$, c'est-à-dire
tels que $f' \circ \alpha = \beta \circ f'$.
D'après Théorie des déformations, lemme \ref{defos-lemma-huge-diagram-ringed-spaces},
pour un couple arbitraire $(\alpha, \beta)$, la différence entre
le morphisme $f' : X' \to Y'$ et le morphisme
$\beta^{-1} \circ f' \circ \alpha : X' \to Y'$ définit un élément
de
$$
\text{Der}_k(\mathcal{O}_Y, f_*\mathcal{O}_X) =
\Hom_{\mathcal{O}_Y}(\Omega_{Y/k}, f_*\mathcal{O}_X)
$$
L'égalité résulte de Compléments sur les morphismes, lemme
\ref{more-morphisms-lemma-netherlander-quasi-coherent}.
Cela définit la dernière flèche horizontale supérieure et montre l'exactitude
aux deux premières places. Pour l'application
$$
\text{Der}_k(\mathcal{O}_Y, f_*\mathcal{O}_X)
\to
T\Deformationcategory_{X \to Y}
$$
interprétons les éléments de la source comme des morphismes
$f_\epsilon : X' \to Y'$ sur $\Spec(k[\epsilon])$
égaux à $f$ modulo $\epsilon$,
à l'aide de Théorie des déformations, lemme \ref{defos-lemma-huge-diagram-ringed-spaces}.
Nous envoyons $f_\epsilon$ sur la classe d'isomorphisme de
$(f_\epsilon : X' \to Y')$ dans $T\Deformationcategory_{X \to Y}$.
Notons que $(f_\epsilon : X' \to Y')$ est isomorphe à la
déformation triviale $(f' : X' \to Y')$ exactement lorsque
$f_\epsilon  = \beta^{-1} \circ f \circ \alpha$ pour un certain
couple $(\alpha, \beta)$, ce qui implique l'exactitude à la troisième place.
Clairement, si une déformation du premier ordre
$(f_\epsilon : X_\epsilon \to Y_\epsilon)$
s'envoie sur zéro dans $T(\Deformationcategory_X \times \Deformationcategory_Y)$,
nous pouvons alors choisir des isomorphismes $X' \to X_\epsilon$ et $Y' \to Y_\epsilon$
et nous concluons que nous sommes dans l'image de la flèche sud-ouest.
Nous avons donc l'exactitude à la quatrième place.
Enfin, étant données deux déformations du premier ordre $X_\epsilon$, $Y_\epsilon$
de $X$, $Y$, il existe une obstruction dans
$$
ob(X_\epsilon, Y_\epsilon) \in
\text{Ext}^1_{\mathcal{O}_X}(Lf^*\NL_{Y/k}, \mathcal{O}_X)
$$
qui s'annule si et seulement si $f : X \to Y$ se relève en
$X_\epsilon \to Y_\epsilon$ ; voir
Théorie des déformations, lemme \ref{defos-lemma-huge-diagram-ringed-spaces}.
Cela achève la démonstration.
\end{proof}

\begin{lemma}
\label{examples-defos-lemma-proper-schemes-morphisms-TI}
Dans le lemme \ref{examples-defos-lemma-schemes-morphisms-TI}, si $X$ et $Y$ sont tous deux
propres sur $k$, alors
$\text{Inf}(\Deformationcategory_{X \to Y})$ et
$T\Deformationcategory_{X \to Y}$ sont de dimension finie.
\end{lemma}

\begin{proof}
Démonstration omise. Indication : raisonner comme dans le lemme \ref{examples-defos-lemma-proper-schemes-TI}
et utiliser la suite exacte du lemme.
\end{proof}

\noindent
Dans l'exemple \ref{examples-defos-example-schemes-morphisms},
si $X \to Y$ est un morphisme de schémas propres sur $k$,
alors $\Deformationcategory_{X \to Y}$
admet une présentation par un groupoïde lisse proreprésentable en foncteurs
sur $\mathcal{C}_\Lambda$
et possède a fortiori un objet formel versel (minimal). Cela résulte
des lemmes \ref{examples-defos-lemma-schemes-morphisms-RS} et
\ref{examples-defos-lemma-proper-schemes-morphisms-TI},
ainsi que de la discussion générale de la section \ref{examples-defos-section-general}.

\begin{lemma}
\label{examples-defos-lemma-schemes-morphisms-hull}
Dans l'exemple \ref{examples-defos-example-schemes-morphisms}, supposons que $X \to Y$
soit un morphisme de $k$-schémas propres.
Supposons que $\Lambda$ soit un anneau local complet de corps résiduel $k$
(cas classique). Alors le foncteur
$$
F : \mathcal{C}_\Lambda \longrightarrow \textit{Ensembles},\quad
A \longmapsto \Ob(\Deformationcategory_{X \to Y}(A))/\cong
$$
des classes d'isomorphisme d'objets possède une enveloppe. Si
$\text{Der}_k(\mathcal{O}_X, \mathcal{O}_X) =
\text{Der}_k(\mathcal{O}_Y, \mathcal{O}_Y) = 0$, alors
$F$ est proreprésentable.
\end{lemma}

\begin{proof}
L'existence d'une enveloppe résulte immédiatement des
lemmes \ref{examples-defos-lemma-schemes-morphisms-RS} et
\ref{examples-defos-lemma-proper-schemes-morphisms-TI}, ainsi que de
Théorie formelle des déformations, lemme \ref{formal-defos-lemma-RS-implies-S1-S2}
et remarque \ref{formal-defos-remark-compose-minimal-into-iso-classes}.

\medskip\noindent
Supposons $\text{Der}_k(\mathcal{O}_X, \mathcal{O}_X) =
\text{Der}_k(\mathcal{O}_Y, \mathcal{O}_Y) = 0$. Alors
la suite exacte du lemme \ref{examples-defos-lemma-schemes-morphisms-TI},
combinée au lemme \ref{examples-defos-lemma-schemes-TI},
montre que $\text{Inf}(\Deformationcategory_{X \to Y}) = 0$.
Alors $\Deformationcategory_{X \to Y}$ et $F$ sont équivalents d'après
Théorie formelle des déformations, lemme \ref{formal-defos-lemma-infdef-trivial}.
Ainsi, $F$ est un foncteur de déformations (car
$\Deformationcategory_{X \to Y}$ est une
catégorie de déformations) d'espace tangent de dimension finie, et nous pouvons appliquer
Théorie formelle des déformations, théorème
\ref{formal-defos-theorem-Schlessinger-prorepresentability}.
\end{proof}

\begin{lemma}
\label{examples-defos-lemma-schemes-morphisms-smooth-to-base}
\begin{reference}
Cette question est traitée dans \cite[Section 5.3]{Ravi-Murphys-Law} et
\cite[Theorem 3.3]{Ran-deformations}.
\end{reference}
Dans l'exemple \ref{examples-defos-example-schemes}, soit $f : X \to Y$ un morphisme de schémas
sur $k$. Si $f_*\mathcal{O}_X = \mathcal{O}_Y$ et $R^1f_*\mathcal{O}_X = 0$,
alors le morphisme de catégories de déformations
$$
\Deformationcategory_{X \to Y} \to \Deformationcategory_X
$$
est une équivalence.
\end{lemma}

\begin{proof}
Construisons un quasi-inverse du foncteur d'oubli du lemme.
Plus précisément, supposons que $(A, U)$ soit un objet de $\Deformationcategory_X$.
L'application donnée $X \to U$ est un épaississement d'ordre fini, que nous pouvons utiliser
pour identifier les espaces topologiques sous-jacents à $U$ et $X$ ; voir
Compléments sur les morphismes, section \ref{more-morphisms-section-thickenings}.
Nous pouvons donc considérer, et considérerons, $\mathcal{O}_U$ comme un faisceau de
$A$-algèbres sur $X$ ; de plus, la platitude de $U \to \Spec(A)$
signifie que $\mathcal{O}_U$ est plat comme faisceau de $A$-modules.
En particulier, nous avons une filtration
$$
0 = \mathfrak m_A^n\mathcal{O}_U \subset
\mathfrak m_A^{n - 1}\mathcal{O}_U \subset \ldots \subset
\mathfrak m_A^2\mathcal{O}_U \subset
\mathfrak m_A\mathcal{O}_U \subset \mathcal{O}_U
$$
dont les sous-quotients sont égaux à
$\mathcal{O}_X \otimes_k \mathfrak m_A^i/\mathfrak m_A^{i + 1}$
par platitude ; voir Compléments sur les morphismes, lemme \ref{more-morphisms-lemma-deform},
ou, plus généralement, Théorie des déformations, lemme \ref{defos-lemma-deform-module}.
Posons
$$
\mathcal{O}_V = f_*\mathcal{O}_U
$$
considéré comme un faisceau de $A$-algèbres sur $Y$. Puisque
$R^1f_*\mathcal{O}_X = 0$, la description ci-dessus montre que
$R^1f_*(\mathfrak m_A^i\mathcal{O}_U/\mathfrak m_A^{i + 1}\mathcal{O}_U) = 0$
pour tout $i$. Il s'ensuit que les suites
$$
0 \to
(f_*\mathcal{O}_X) \otimes_k \mathfrak m_A^i/\mathfrak m_A^{i + 1} \to
f_*(\mathcal{O}_U/\mathfrak m_A^{i + 1}\mathcal{O}_U) \to
f_*(\mathcal{O}_U/\mathfrak m_A^i\mathcal{O}_U) \to 0
$$
sont exactes pour tout $i$. En parcourant en sens inverse les références données ci-dessus
(et en raisonnant par récurrence), nous trouvons que $\mathcal{O}_V$ est un faisceau plat
de $A$-algèbres tel que
$\mathcal{O}_V/\mathfrak m_A\mathcal{O}_V = \mathcal{O}_Y$.
À l'aide de Compléments sur les morphismes, lemme
\ref{more-morphisms-lemma-first-order-thickening}
nous trouvons que $(Y, \mathcal{O}_V)$ est un schéma, que nous notons $V$.
L'égalité $\mathcal{O}_V = f_*\mathcal{O}_U$ définit un
morphisme d'espaces annelés $U \to V$ dont on voit aisément qu'il est
un morphisme de schémas. La platitude déjà établie
achève la démonstration.
\end{proof}







\section{Espaces algébriques}
\label{examples-defos-section-algebraic-spaces}

\noindent
La théorie des déformations des espaces algébriques.

\begin{example}[Espaces algébriques]
\label{examples-defos-example-spaces}
Soit $\mathcal{F}$ la catégorie définie comme suit :
\begin{enumerate}
\item un objet est un couple $(A, X)$ formé d'un
objet $A$ de $\mathcal{C}_\Lambda$ et d'un espace algébrique
$X$ plat sur $A$ ;
\item un morphisme $(f, g) : (B, Y) \to (A, X)$ est formé d'un
morphisme $f : B \to A$ dans $\mathcal{C}_\Lambda$ et
d'un morphisme $g : X \to Y$ d'espaces algébriques sur $\Lambda$
tel que
$$
\xymatrix{
X \ar[r]_g \ar[d] & Y \ar[d] \\
\Spec(A) \ar[r]^f & \Spec(B)
}
$$
soit un diagramme commutatif cartésien d'espaces algébriques.
\end{enumerate}
Le foncteur $p : \mathcal{F} \to \mathcal{C}_\Lambda$ envoie $(A, X)$
sur $A$ et $(f, g)$ sur $f$. Il est clair que $p$ est cofibré en groupoïdes.
Étant donné un espace algébrique $X$ sur $k$, soit
$x_0 = (k, X)$ l'objet correspondant de $\mathcal{F}(k)$. On pose
$$
\Deformationcategory_X = \mathcal{F}_{x_0}
$$
\end{example}

\begin{lemma}
\label{examples-defos-lemma-spaces-RS}
L'exemple \ref{examples-defos-example-spaces}
satisfait la condition de Rim-Schlessinger (RS).
En particulier, $\Deformationcategory_X$ est une catégorie de déformations
pour tout espace algébrique $X$ sur $k$.
\end{lemma}

\begin{proof}
Soient $A_1 \to A$ et $A_2 \to A$ des morphismes de $\mathcal{C}_\Lambda$.
Supposons $A_2 \to A$ surjectif. D'après Théorie formelle
des déformations, lemme
\ref{formal-defos-lemma-RS-2-categorical}
il suffit de montrer que le foncteur
$\mathcal{F}(A_1 \times_A A_2) \to
\mathcal{F}(A_1) \times_{\mathcal{F}(A)} \mathcal{F}(A_2)$
est une équivalence de catégories.
Observons que
$$
\xymatrix{
\Spec(A) \ar[r] \ar[d] & \Spec(A_2) \ar[d] \\
\Spec(A_1) \ar[r] &
\Spec(A_1 \times_A A_2)
}
$$
est un diagramme cocartésien comme dans Sommes amalgamées d'espaces, lemme
\ref{spaces-pushouts-lemma-pushout-along-thickening}.
Le lemme est donc un cas particulier de Sommes amalgamées d'espaces, lemme
\ref{spaces-pushouts-lemma-equivalence-categories-spaces-pushout-flat}.
\end{proof}

\begin{lemma}
\label{examples-defos-lemma-spaces-TI}
Dans l'exemple \ref{examples-defos-example-spaces}, soit $X$ un espace algébrique sur $k$. Alors
$$
\text{Inf}(\Deformationcategory_X) =
\text{Ext}^0_{\mathcal{O}_X}(\NL_{X/k}, \mathcal{O}_X) =
\Hom_{\mathcal{O}_X}(\Omega_{X/k}, \mathcal{O}_X) =
\text{Der}_k(\mathcal{O}_X, \mathcal{O}_X)
$$
et
$$
T\Deformationcategory_X =
\text{Ext}^1_{\mathcal{O}_X}(\NL_{X/k}, \mathcal{O}_X)
$$
\end{lemma}

\begin{proof}
Rappelons que $\text{Inf}(\Deformationcategory_X)$ est l'ensemble des
automorphismes de la déformation triviale
$X' = X \times_{\Spec(k)} \Spec(k[\epsilon])$ de $X$ sur $k[\epsilon]$
qui sont égaux à l'identité modulo $\epsilon$.
D'après Théorie des déformations, lemme \ref{defos-lemma-deform-spaces},
cet ensemble est égal à $\text{Ext}^0_{\mathcal{O}_X}(\NL_{X/k}, \mathcal{O}_X)$.
L'égalité $\text{Ext}^0_{\mathcal{O}_X}(\NL_{X/k}, \mathcal{O}_X) =
\Hom_{\mathcal{O}_X}(\Omega_{X/k}, \mathcal{O}_X)$ résulte de
Compléments sur les morphismes d'espaces, lemme
\ref{spaces-more-morphisms-lemma-netherlander-quasi-coherent}.
L'égalité
$\Hom_{\mathcal{O}_X}(\Omega_{X/k}, \mathcal{O}_X) =
\text{Der}_k(\mathcal{O}_X, \mathcal{O}_X)$
résulte de Compléments sur les morphismes d'espaces, définition
\ref{spaces-more-morphisms-definition-sheaf-differentials}, et de
Modules sur les sites, définition
\ref{sites-modules-definition-module-differentials}.

\medskip\noindent
Rappelons que $T_{x_0}\Deformationcategory_X$ est l'ensemble des classes d'isomorphisme
des déformations plates $X'$ de $X$ sur $k[\epsilon]$, plus précisément
l'ensemble des classes d'isomorphisme de $\Deformationcategory_X(k[\epsilon])$.
La seconde assertion du lemme résulte donc de
Théorie des déformations, lemme \ref{defos-lemma-deform-spaces}.
\end{proof}

\begin{lemma}
\label{examples-defos-lemma-proper-spaces-TI}
Dans le lemme \ref{examples-defos-lemma-spaces-TI}, si $X$ est propre sur $k$, alors
$\text{Inf}(\Deformationcategory_X)$ et $T\Deformationcategory_X$ sont
de dimension finie.
\end{lemma}

\begin{proof}
D'après le lemme, il faut montrer que
$\Ext^1_{\mathcal{O}_X}(\NL_{X/k}, \mathcal{O}_X)$ et
$\Ext^0_{\mathcal{O}_X}(\NL_{X/k}, \mathcal{O}_X)$ sont de
dimension finie. D'après Compléments sur les morphismes d'espaces, lemme
\ref{spaces-more-morphisms-lemma-netherlander-fp}
et le fait que $X$ est noethérien, nous voyons que
$\NL_{X/k}$ a des faisceaux de cohomologie cohérents, nuls sauf
en degrés $0$ et $-1$.
D'après Catégories dérivées d'espaces, lemme \ref{spaces-perfect-lemma-ext-finite},
les groupes $\Ext$ affichés sont des $k$-espaces vectoriels de dimension finie,
ce qui achève la démonstration.
\end{proof}

\noindent
Dans l'exemple \ref{examples-defos-example-spaces}, si $X$ est un espace algébrique propre sur $k$,
alors $\Deformationcategory_X$
admet une présentation par un groupoïde lisse proreprésentable en foncteurs
sur $\mathcal{C}_\Lambda$
et possède a fortiori un objet formel versel (minimal). Cela résulte
des lemmes \ref{examples-defos-lemma-spaces-RS} et
\ref{examples-defos-lemma-proper-spaces-TI},
ainsi que de la discussion générale de la section \ref{examples-defos-section-general}.

\begin{lemma}
\label{examples-defos-lemma-spaces-hull}
Dans l'exemple \ref{examples-defos-example-spaces}, supposons que $X$ soit un espace algébrique propre sur $k$.
Supposons que $\Lambda$ soit un anneau local complet de corps résiduel $k$
(cas classique). Alors le foncteur
$$
F : \mathcal{C}_\Lambda \longrightarrow \textit{Ensembles},\quad
A \longmapsto \Ob(\Deformationcategory_X(A))/\cong
$$
des classes d'isomorphisme d'objets possède une enveloppe. Si
$\text{Der}_k(\mathcal{O}_X, \mathcal{O}_X) = 0$, alors
$F$ est proreprésentable.
\end{lemma}

\begin{proof}
L'existence d'une enveloppe résulte immédiatement des
lemmes \ref{examples-defos-lemma-spaces-RS} et \ref{examples-defos-lemma-proper-spaces-TI}, ainsi que de
Théorie formelle des déformations, lemme \ref{formal-defos-lemma-RS-implies-S1-S2}
et remarque \ref{formal-defos-remark-compose-minimal-into-iso-classes}.

\medskip\noindent
Supposons $\text{Der}_k(\mathcal{O}_X, \mathcal{O}_X) = 0$. Alors
$\Deformationcategory_X$ et $F$ sont équivalents d'après
Théorie formelle des déformations, lemme \ref{formal-defos-lemma-infdef-trivial}.
Ainsi, $F$ est un foncteur de déformations (car $\Deformationcategory_X$ est une
catégorie de déformations) d'espace tangent de dimension finie, et nous pouvons appliquer
Théorie formelle des déformations, théorème
\ref{formal-defos-theorem-Schlessinger-prorepresentability}.
\end{proof}





\section{Déformations des complétés}
\label{examples-defos-section-compare}

\noindent
Dans cette section, nous comparons les problèmes de déformation posés
par une algèbre et par son complété.
Nous discutons d'abord de la « possibilité de relèvement ».

\begin{lemma}
\label{examples-defos-lemma-lift-equivalence-module-derived}
Soit $A' \to A$ une surjection d'anneaux de noyau nilpotent.
Soit $A' \to P'$ un homomorphisme d'anneaux plat.
Posons $P = P' \otimes_{A'} A$.
Soit $M$ un module plat sur $A$, muni d'une structure de $P$-module.
Alors les assertions suivantes sont équivalentes :
\begin{enumerate}
\item il existe un module plat sur $A'$, muni d'une structure de $P'$-module, noté $M'$, tel que
$M' \otimes_{P'} P = M$ ;
\item il existe un objet $K' \in D^-(P')$ tel que
$K' \otimes_{P'}^\mathbf{L} P = M$.
\end{enumerate}
\end{lemma}

\begin{proof}
Supposons que $M'$ soit comme dans (1). Alors
$$
M = M' \otimes_P P' = M' \otimes_{A'} A =
M' \otimes_A^\mathbf{L} A' = M' \otimes_{P'}^\mathbf{L} P
$$
Les deux premières égalités sont claires, la troisième vaut parce que
$M'$ est plat sur $A'$, et la quatrième résulte de
Compléments sur l'algèbre, lemme \ref{more-algebra-lemma-base-change-comparison}.
Ainsi, (2) est vérifiée. Réciproquement, supposons $K'$ comme dans (2).
Nous pouvons supposer, et supposons, $M$ non nul.
Soit $t$ le plus grand entier tel que $H^t(K')$ soit non nul
(il existe puisque $M$ est non nul).
Alors $H^t(K') \otimes_{P'} P = H^t(K' \otimes_{P'}^\mathbf{L} P)$
est nul si $t > 0$. Comme le noyau de $P' \to P$ est nilpotent,
cela implique $H^t(K') = 0$ par le lemme de Nakayama, contradiction.
Ainsi, $t = 0$ (le cas $t < 0$ est lui aussi absurde).
Alors $M' = H^0(K')$ est un $P'$-module tel que $M = M' \otimes_{P'} P$,
et la suite spectrale de Tor donne une application injective
$$
\text{Tor}_1^{P'}(M', P) \to H^{-1}(M' \otimes_{P'}^\mathbf{L} P) = 0
$$
Par la référence ci-dessus sur le changement de base dérivé,
$0 = \text{Tor}_1^{P'}(M', P) = \text{Tor}_1^{A'}(M', A)$.
Nous concluons que $M'$ est plat sur $A'$ d'après
Algèbre, lemme \ref{algebra-lemma-what-does-it-mean}.
\end{proof}

\begin{lemma}
\label{examples-defos-lemma-lift-equivalence-module}
Considérons un diagramme commutatif d'anneaux noethériens
$$
\xymatrix{
A' \ar[d] \ar[r] &
P' \ar[d] \ar[r] &
Q' \ar[d] \\
A \ar[r] &
P \ar[r] &
Q
}
$$
dont les carrés sont cartésiens, les flèches horizontales plates et les
flèches verticales surjectives, de noyaux nilpotents.
Soit $J' \subset P'$ un idéal tel que $P'/J' = Q'/J'Q'$.
Soit $M$ un module plat sur $A$, muni d'une structure de $P$-module.
Supposons que, pour tout $g \in J'$, il existe un module plat sur $A'$ muni d'une structure de $(P')_g$-module
relevant $M_g$. Alors les assertions suivantes sont équivalentes :
\begin{enumerate}
\item $M$ possède un relèvement plat sur $A'$ muni d'une structure de $P'$-module ;
\item $M \otimes_P Q$ possède un relèvement plat sur $A'$ muni d'une structure de $Q'$-module.
\end{enumerate}
\end{lemma}

\begin{proof}
Soit $I = \Ker(A' \to A)$. Par récurrence sur l'entier $n > 1$
tel que $I^n = 0$, nous nous ramenons au cas où $I$ est un idéal
de carré nul ; nous omettons les détails.
Nous traduisons la condition de possibilité de relèvement de
$M$ en le problème de trouver un objet de $D^-(P')$ comme dans le
lemme \ref{examples-defos-lemma-lift-equivalence-module-derived}.
L'obstruction à cette construction est l'élément
$$
\omega(M) \in \text{Ext}^2_P(M, M \otimes_P^\mathbf{L} IP) =
\text{Ext}^2_P(M, M \otimes_P IP)
$$
construit dans
Théorie des déformations, lemme \ref{defos-lemma-canonical-class-algebra}.
L'égalité dans la formule affichée vaut car
$M \otimes_P^\mathbf{L} IP = M \otimes_P IP$
puisque $M$ et $P$ sont plats sur $A$\footnote{Choisissons une résolution
$F_\bullet \to I$ par des $A$-modules libres. Puisque $A \to P$ est plat,
$P \otimes_A F_\bullet$ est une résolution libre de $IP$.
Ainsi, $M \otimes_P^\mathbf{L} IP$ est représenté par
$M \otimes_P P \otimes_A F_\bullet = M \otimes_A F_\bullet$.
Ce complexe n'a de cohomologie qu'en degré $0$, car $M$ est plat sur $A$.}.
De même, l'obstruction au relèvement de $M \otimes_P Q$ est
l'élément
$$
\omega(M \otimes_P Q) \in
\text{Ext}^2_Q(M \otimes_P Q, (M \otimes_P Q) \otimes_Q IQ)
$$
qui est l'image de $\omega(M)$ par la fonctorialité
de la construction $\omega(-)$ de
Théorie des déformations, lemme \ref{defos-lemma-canonical-class-algebra}.
D'après Compléments sur l'algèbre, lemme \ref{more-algebra-lemma-base-change-RHom},
nous avons
$$
\text{Ext}^2_Q(M \otimes_P Q, (M \otimes_P Q) \otimes_Q IQ) =
\text{Ext}^2_P(M, M \otimes_P IP) \otimes_P Q
$$
où nous utilisons le fait que $P$ est noethérien et $M$ fini.
Notre hypothèse sur $P' \to Q'$ garantit que, pour un $P$-module $E$,
l'application $E \to E \otimes_P Q$ est bijective sur la partie annulée par une puissance de $J'$ ; voir
Compléments sur l'algèbre, lemme
\ref{more-algebra-lemma-neighbourhood-equivalence}.
Nous concluons donc qu'il suffit de montrer que $\omega(M)$
est annulé par une puissance de $J'$. Autrement dit, il suffit de montrer que
$\omega(M)$ s'annule dans
$$
\text{Ext}^2_P(M, M \otimes_P IP)_g =
\text{Ext}^2_{P_g}(M_g, M_g \otimes_{P_g} IP_g)
$$
pour tout $g \in J'$. Cependant, par la compatibilité de la formation de $\omega(M)$
au changement de base, nous concluons à nouveau que cela est vrai puisque $M_g$
est supposé posséder un relèvement (il faut bien sûr utiliser de nouveau toute la
chaîne d'équivalences).
\end{proof}

\begin{lemma}
\label{examples-defos-lemma-lift-equivalence}
Soit $A' \to A$ un homomorphisme surjectif d'anneaux noethériens de noyau nilpotent.
Soit $A \to B$ un homomorphisme d'anneaux plat de type fini.
Soit $\mathfrak b \subset B$ un idéal tel que
$\Spec(B) \to \Spec(A)$ soit syntomique sur le complémentaire de $V(\mathfrak b)$.
Alors $B$ possède un relèvement plat sur $A'$ si et seulement si le complété
$\mathfrak b$-adique $B^\wedge$ possède un relèvement plat sur $A'$.
\end{lemma}

\begin{proof}
Choisissons une surjection de $A$-algèbres $P = A[x_1, \ldots, x_n] \to B$.
Soit $\mathfrak p \subset P$ l'image réciproque de $\mathfrak b$.
Posons $P' = A'[x_1, \ldots, x_n]$ et notons $\mathfrak p' \subset P'$
l'image réciproque de $\mathfrak p$. (Bien sûr, $\mathfrak p$
et $\mathfrak p'$ ne désignent pas ici des idéaux premiers.)
Nous noterons $P^\wedge$ et $(P')^\wedge$ les complétés respectifs.

\medskip\noindent
Supposons que $A' \to B'$ soit un relèvement plat de $A \to B$, autrement dit
que $A' \to B'$ soit plat et qu'il existe un isomorphisme de $A$-algèbres
$B = B' \otimes_{A'} A$. Nous pouvons alors choisir un homomorphisme de $A'$-algèbres
$P' \to B'$ relevant la surjection donnée $P \to B$.
Par le lemme de Nakayama (Algèbre, lemme \ref{algebra-lemma-NAK}),
nous trouvons que $B'$ est un quotient de $P'$. En particulier, nous trouvons
que nous pouvons rendre $B'$ plat sur $A'$ comme $P'$-module,
relevant $B$, qui est plat sur $A$, comme $P$-module.
Réciproquement, si nous pouvons relever $B$ en un $P'$-module $M'$ plat sur $A'$,
alors $M'$ est un module cyclique $M' \cong P'/J'$ (en utilisant de nouveau Nakayama),
et, en posant $B' = P'/J'$, nous obtenons un relèvement plat de $B$ comme algèbre.

\medskip\noindent
Posons $C = B^\wedge$ et $\mathfrak c = \mathfrak bC$.
Supposons que $A' \to C'$ soit un relèvement plat de $A \to C$.
Alors $C'$ est complet pour l'image réciproque $\mathfrak c'$
de $\mathfrak c$
(Algèbre, lemme \ref{algebra-lemma-complete-modulo-nilpotent}).
Choisissons un homomorphisme de $A'$-algèbres $P' \to C'$ relevant
l'homomorphisme de $A$-algèbres $P \to C$. Ces homomorphismes se prolongent aux
complétés et donnent des surjections $P^\wedge \to C$ et $(P')^\wedge \to C'$
(pour la seconde, en utilisant de nouveau le lemme de Nakayama).
En particulier, nous trouvons que $C'$ est plat sur $A'$ et peut être muni d'une structure de
$(P')^\wedge$-module relevant $C$, plat sur $A$, comme $P^\wedge$-module.
Réciproquement, si nous pouvons relever $C$ en un $(P')^\wedge$-module $N'$ plat sur $A'$,
alors $N'$ est un module cyclique $N' \cong (P')^\wedge/\tilde J$
(en utilisant de nouveau Nakayama) et, en posant $C' = (P')^\wedge/\tilde J$,
nous obtenons un relèvement plat de $C$ comme algèbre.

\medskip\noindent
Observons que $P' \to (P')^\wedge$ est un homomorphisme d'anneaux plat qui
induit un isomorphisme $P'/\mathfrak p' = (P')^\wedge/\mathfrak p'(P')^\wedge$.
Nous concluons que notre lemme résulte du
lemme \ref{examples-defos-lemma-lift-equivalence-module}, pourvu que nous puissions
montrer que $B_g$ se relève en un module plat sur $A'$ muni d'une structure de $P'_g$-module pour
$g \in \mathfrak p'$. Or l'homomorphisme d'anneaux $A \to B_g$ est syntomique
et se relève donc en une $A'$-algèbre plate $B'$ d'après
Lissage des homomorphismes d'anneaux, proposition \ref{smoothing-proposition-lift-smooth}.
Puisque $A' \to P'_g$ est lisse, nous pouvons relever $P_g \to B_g$
en une application surjective $P'_g \to B'$ comme ci-dessus, ce qui donne le résultat voulu.
\end{proof}

\noindent
Notation. Soit $A \to B$ un homomorphisme d'anneaux. Soit $N$ un $B$-module.
Nous notons $\text{Exal}_A(B, N)$ l'ensemble des classes d'isomorphisme
des extensions
$$
0 \to N \to C \to B \to 0
$$
de $A$-algèbres telles que $N$ soit un idéal de carré nul dans $C$.
Étant donnée une seconde extension $0 \to N \to C' \to B \to 0$ de ce type, un isomorphisme
est un isomorphisme de $A$-algèbres $C \to C'$ tel que le diagramme
$$
\xymatrix{
0 \ar[r] &
N \ar[r] \ar[d]_{\text{id}} &
C \ar[r] \ar[d] &
B \ar[r] \ar[d]_{\text{id}} & 0 \\
0 \ar[r] &
N \ar[r] &
C' \ar[r] &
B \ar[r] & 0
}
$$
soit commutatif. La correspondance $N \mapsto \text{Exal}_A(B, N)$
est un foncteur qui transforme les produits en produits.
C'est donc un foncteur additif, et $\text{Exal}_A(B, N)$
possède une structure naturelle de $B$-module. En fait, d'après
Théorie des déformations, lemme \ref{defos-lemma-choices},
nous avons $\text{Exal}_A(B, N) = \text{Ext}^1_B(\NL_{B/A}, N)$.

\begin{lemma}
\label{examples-defos-lemma-first-order-completion}
Soit $k$ un corps. Soit $B$ une $k$-algèbre de type fini.
Soit $J \subset B$ un idéal tel que
$\Spec(B) \to \Spec(k)$ soit lisse sur le complémentaire de $V(J)$.
Soit $N$ un $B$-module fini.
Alors il existe une bijection canonique
$$
\text{Exal}_k(B, N) \to \text{Exal}_k(B^\wedge, N^\wedge)
$$
Ici, $B^\wedge$ et $N^\wedge$ sont les complétés $J$-adiques.
\end{lemma}

\begin{proof}
L'application est donnée par complétion : étant donné $0 \to N \to C \to B \to 0$
dans $\text{Exal}_k(B, N)$, nous l'envoyons sur le complété $C^\wedge$
de $C$ pour l'image réciproque de $J$. Comparer avec
la démonstration du lemme \ref{examples-defos-lemma-completion}.

\medskip\noindent
Puisque $k \to B$ est de présentation finie, le complexe
$\NL_{B/k}$ peut être représenté par un complexe
$N^{-1} \to N^0$, où $N^i$ est un $B$-module fini ; voir
Algèbre, section \ref{algebra-section-netherlander}, et
en particulier
Algèbre, lemme \ref{algebra-lemma-NL-homotopy}.
Comme $B$ est noethérien, cela signifie que $\NL_{B/k}$
est pseudo-cohérent. Pour $g \in J$, la $k$-algèbre $B_g$
est lisse et donc $(\NL_{B/k})_g = \NL_{B_g/k}$
est quasi-isomorphe à un $B$-module projectif de type fini placé en degré $0$.
Ainsi, $\text{Ext}^i_B(\NL_{B/k}, N)_g = 0$ pour $i \geq 1$
et tout $B$-module $N$. D'après
Compléments sur l'algèbre, lemme \ref{more-algebra-lemma-ext-annihilated-into},
nous concluons que
$$
\text{Ext}^1_B(\NL_{B/k}, N) \longrightarrow
\lim_n \text{Ext}^1_B(\NL_{B/k}, N/J^n N)
$$
est un isomorphisme pour tout $B$-module fini $N$.

\medskip\noindent
Injectivité de l'application.
Supposons que $0 \to N \to C \to B \to 0$ appartienne à $\text{Exal}_k(B, N)$
et s'envoie sur zéro dans $\text{Exal}_k(B^\wedge, N^\wedge)$.
Choisissons un scindage $C^\wedge = B^\wedge \oplus N^\wedge$.
Alors l'application induite $C \to C^\wedge \to N^\wedge$
donne des applications $C \to N/J^nN$ pour tout $n$.
Nous voyons donc que notre élément appartient au noyau des applications
$$
\text{Ext}^1_B(\NL_{B/k}, N) \to
\text{Ext}^1_B(\NL_{B/k}, N/J^n N)
$$
pour tout $n$. Le paragraphe précédent montre que
notre élément est nul.

\medskip\noindent
Surjectivité de l'application. Soit $0 \to N^\wedge \to C' \to B^\wedge \to 0$
un élément de $\text{Exal}_k(B^\wedge, N^\wedge)$.
Par image réciproque selon $B \to B^\wedge$, nous obtenons un élément
$0 \to N^\wedge \to C'' \to B \to 0$ dans
$\text{Exal}_k(B, N^\wedge)$.
nous avons
$$
\text{Ext}^1_B(\NL_{B/k}, N^\wedge) =
\text{Ext}^1_B(\NL_{B/k}, N) \otimes_B B^\wedge =
\text{Ext}^1_B(\NL_{B/k}, N)
$$
La première égalité vaut puisque $N^\wedge = N \otimes_B B^\wedge$
(Algèbre, lemme \ref{algebra-lemma-completion-tensor})
et d'après
Compléments sur l'algèbre, lemme \ref{more-algebra-lemma-pseudo-coherence-and-ext}.
La seconde égalité vaut parce que $\text{Ext}^1_B(\NL_{B/k}, N)$
est annulé par une puissance de $J$ (voir ci-dessus), que $B \to B^\wedge$ est plat et induit
un isomorphisme $B/J \to B^\wedge/JB^\wedge$, et d'après
Compléments sur l'algèbre, lemme \ref{more-algebra-lemma-neighbourhood-equivalence}.
Nous pouvons donc trouver un $C \in \text{Exal}_k(B, N)$ s'envoyant sur $C''$ dans
$\text{Exal}_k(B, N^\wedge)$.
Ainsi,
$$
0 \to N^\wedge \to C' \to B^\wedge \to 0
\quad\text{et}\quad
0 \to N^\wedge \to C^\wedge \to B^\wedge \to 0
$$
ces suites sont deux éléments de $\text{Exal}_k(B^\wedge, N^\wedge)$
qui s'envoient sur le même élément de $\text{Exal}_k(B, N^\wedge)$.
En prenant leur différence, nous obtenons un élément
$0 \to N^\wedge \to C' \to B^\wedge \to 0$ de
$\text{Exal}_k(B^\wedge, N^\wedge)$
dont l'image dans $\text{Exal}_k(B, N^\wedge)$ est nulle.
Cela signifie qu'il existe
$$
\xymatrix{
0 \ar[r] &
N^\wedge \ar[r] &
C' \ar[r] &
B^\wedge \ar[r] & 0 \\
& & B \ar[u]^\sigma \ar[ru]
}
$$
Soit $J' \subset C'$ l'image réciproque de $JB^\wedge \subset B^\wedge$.
Pour achever la démonstration, il suffit de noter que
$\sigma$ est continue pour la topologie $J$-adique sur $B$
et la topologie $J'$-adique sur $C'$, et que $C'$ est complet pour la topologie $J'$-adique d'après
Algèbre, lemme \ref{algebra-lemma-complete-modulo-nilpotent}
(nous utilisons aussi ici que $C'$ est noethérien ; un petit détail est omis).
En effet, cela signifie que $\sigma$ se factorise par le
complété $B^\wedge$ et que $C' = 0$ dans $\text{Exal}_k(B^\wedge, N^\wedge)$.
\end{proof}

\begin{lemma}
\label{examples-defos-lemma-smooth-completion}
Dans l'exemple \ref{examples-defos-example-rings}, soit $P$ une $k$-algèbre.
Soit $J \subset P$ un idéal.
Notons $P^\wedge$ le complété $J$-adique. Si
\begin{enumerate}
\item $k \to P$ est de type fini ;
\item $\Spec(P) \to \Spec(k)$ est lisse sur le complémentaire de $V(J)$.
\end{enumerate}
alors le foncteur entre catégories de déformations du
lemme \ref{examples-defos-lemma-completion}
$$
\Deformationcategory_P \longrightarrow \Deformationcategory_{P^\wedge}
$$
est lisse et induit un isomorphisme sur les espaces tangents.
\end{lemma}

\begin{proof}
Nous savons que $\Deformationcategory_P$ et $\Deformationcategory_{P^\wedge}$
sont des catégories de déformations d'après le lemme \ref{examples-defos-lemma-rings-RS}.
Il suffit donc de vérifier que
notre foncteur identifie les espaces tangents et les possibilités
de relèvement ; voir Théorie formelle des déformations, lemme
\ref{formal-defos-lemma-easy-check-smooth}.
La propriété concernant les possibilités de relèvement est démontrée dans le
lemme \ref{examples-defos-lemma-lift-equivalence},
et l'isomorphisme sur les espaces tangents est le cas particulier du
lemme \ref{examples-defos-lemma-first-order-completion} où $N = B$.
\end{proof}



\section{Déformations des localisations}
\label{examples-defos-section-compare-localization}

\noindent
Dans cette section, nous comparons les problèmes de déformation posés
par une algèbre et par sa localisation en une partie multiplicative.
Nous discutons d'abord de la « possibilité de relèvement ».

\begin{lemma}
\label{examples-defos-lemma-lift-equivalence-localization}
Soit $A' \to A$ un homomorphisme surjectif d'anneaux noethériens de noyau nilpotent.
Soit $A \to B$ un homomorphisme d'anneaux plat de type fini.
Soit $S \subset B$ une partie multiplicative telle que,
si $\Spec(B) \to \Spec(A)$ n'est pas syntomique en $\mathfrak q$,
alors $S \cap \mathfrak q = \emptyset$.
Alors $B$ possède un relèvement plat sur $A'$ si et seulement si
$S^{-1}B$ possède un relèvement plat sur $A'$.
\end{lemma}

\begin{proof}
Cette démonstration est la même que celle du
lemme \ref{examples-defos-lemma-lift-equivalence}, mais plus facile. Nous suggérons au
lecteur de l'omettre.
Choisissons une surjection de $A$-algèbres $P = A[x_1, \ldots, x_n] \to B$.
Soit $S_P \subset P$ l'image réciproque de $S$.
Posons $P' = A'[x_1, \ldots, x_n]$ et notons $S_{P'} \subset P'$
l'image réciproque de $S_P$.

\medskip\noindent
Supposons que $A' \to B'$ soit un relèvement plat de $A \to B$, autrement dit
que $A' \to B'$ soit plat et qu'il existe un isomorphisme de $A$-algèbres
$B = B' \otimes_{A'} A$. Nous pouvons alors choisir un homomorphisme de $A'$-algèbres
$P' \to B'$ relevant la surjection donnée $P \to B$.
Par le lemme de Nakayama (Algèbre, lemme \ref{algebra-lemma-NAK}),
nous trouvons que $B'$ est un quotient de $P'$. En particulier, nous trouvons
que nous pouvons rendre $B'$ plat sur $A'$ comme $P'$-module,
relevant $B$, qui est plat sur $A$, comme $P$-module.
Réciproquement, si nous pouvons relever $B$ en un $P'$-module $M'$ plat sur $A'$,
alors $M'$ est un module cyclique $M' \cong P'/J'$ (en utilisant de nouveau Nakayama),
et, en posant $B' = P'/J'$, nous obtenons un relèvement plat de $B$ comme algèbre.

\medskip\noindent
Posons $C = S^{-1}B$. Supposons que $A' \to C'$ soit un relèvement plat de $A \to C$.
Les éléments de $C'$ dont l'image dans $C$ est inversible sont inversibles.
Choisissons un homomorphisme de $A'$-algèbres $P' \to C'$ relevant
l'homomorphisme de $A$-algèbres $P \to C$. D'après la remarque précédente,
ces homomorphismes se prolongent aux localisations et donnent des surjections
$S_P^{-1}P \to C$ et $S_{P'}^{-1}P' \to C'$
(pour la seconde, utiliser le lemme de Nakayama).
En particulier, nous trouvons que $C'$ est plat sur $A'$ et peut être muni d'une structure de
$S_{P'}^{-1}P'$-module relevant $C$, plat sur $A$, comme
$S_P^{-1}P$-module. Réciproquement, si nous pouvons relever $C$ en un
$S_{P'}^{-1}P'$-module $N'$ plat sur $A'$, alors $N'$
est un module cyclique $N' \cong S_{P'}^{-1}P'/\tilde J$
(en utilisant de nouveau Nakayama) et, en posant $C' = S_{P'}^{-1}P'/\tilde J$,
nous obtenons un relèvement plat de $C$ comme algèbre.

\medskip\noindent
Le lieu syntomique d'un morphisme de schémas est ouvert par définition.
Soit $J_B \subset B$ un idéal définissant l'ensemble des points
de $\Spec(B)$ où $\Spec(B) \to \Spec(A)$ n'est pas syntomique.
Notons $J_P \subset P$ et $J_{P'} \subset P'$ les idéaux
correspondants. Observons que $P' \to S_{P'}^{-1}P'$ est un homomorphisme d'anneaux plat qui
induit un isomorphisme $P'/J_{P'} = S_{P'}^{-1}P'/J_{P'}S_{P'}^{-1}P'$
grâce à notre hypothèse sur $S$ dans le lemme ; plus précisément, l'hypothèse
du lemme signifie exactement que $B/J_B = S^{-1}(B/J_B)$.
Nous concluons que notre lemme résulte du
lemme \ref{examples-defos-lemma-lift-equivalence-module}, pourvu que nous puissions
montrer que $B_g$ se relève en un module plat sur $A'$ muni d'une structure de $P'_g$-module pour
$g \in J_B$. Or l'homomorphisme d'anneaux $A \to B_g$ est syntomique
et se relève donc en une $A'$-algèbre plate $B'$ d'après
Lissage des homomorphismes d'anneaux, proposition \ref{smoothing-proposition-lift-smooth}.
Puisque $A' \to P'_g$ est lisse, nous pouvons relever $P_g \to B_g$
en une application surjective $P'_g \to B'$ comme ci-dessus, ce qui donne le résultat voulu.
\end{proof}

\begin{lemma}
\label{examples-defos-lemma-first-order-localization}
Soit $k$ un corps. Soit $B$ une $k$-algèbre de type fini.
Soit $S \subset B$ une partie multiplicative telle que,
si $\Spec(B) \to \Spec(k)$ n'est pas lisse en $\mathfrak q$,
alors $S \cap \mathfrak q = \emptyset$.
Soit $N$ un $B$-module fini.
Alors il existe une bijection canonique
$$
\text{Exal}_k(B, N) \to \text{Exal}_k(S^{-1}B, S^{-1}N)
$$
\end{lemma}

\begin{proof}
Cette démonstration est la même que celle du
lemme \ref{examples-defos-lemma-first-order-completion}, mais plus facile. Nous suggérons au
lecteur de l'omettre.
L'application est donnée par localisation : étant donné $0 \to N \to C \to B \to 0$
dans $\text{Exal}_k(B, N)$, nous l'envoyons sur la localisation $S_C^{-1}C$
de $C$ en l'image réciproque $S_C \subset C$ de $S$.
Comparer avec la démonstration du lemme \ref{examples-defos-lemma-localization}.

\medskip\noindent
Le lieu de lissité d'un morphisme de schémas est ouvert par définition.
Soit $J \subset B$ un idéal définissant l'ensemble des points
de $\Spec(B)$ où $\Spec(B) \to \Spec(A)$ n'est pas lisse.
Puisque $k \to B$ est de présentation finie, le complexe
$\NL_{B/k}$ peut être représenté par un complexe
$N^{-1} \to N^0$, où $N^i$ est un $B$-module fini ; voir
Algèbre, section \ref{algebra-section-netherlander}, et
en particulier
Algèbre, lemme \ref{algebra-lemma-NL-homotopy}.
Comme $B$ est noethérien, cela signifie que $\NL_{B/k}$
est pseudo-cohérent. Pour $g \in J$, la $k$-algèbre $B_g$
est lisse et donc $(\NL_{B/k})_g = \NL_{B_g/k}$
est quasi-isomorphe à un $B$-module projectif de type fini placé en degré $0$.
Ainsi, $\text{Ext}^i_B(\NL_{B/k}, N)_g = 0$ pour $i \geq 1$
et tout $B$-module $N$. Enfin, nous avons
$$
\text{Ext}^1_{S^{-1}B}(\NL_{S^{-1}B/k}, S^{-1}N) =
\text{Ext}^1_B(\NL_{B/k}, N) \otimes_B S^{-1}B =
\text{Ext}^1_B(\NL_{B/k}, N)
$$
La première égalité résulte de
Compléments sur l'algèbre, lemme \ref{more-algebra-lemma-base-change-RHom},
et d'Algèbre, lemme \ref{algebra-lemma-localize-NL}.
La seconde vaut parce que $\text{Ext}^1_B(\NL_{B/k}, N)$ est annulé par une puissance de $J$
et que les éléments de $S$ agissent inversiblement sur les modules annulés par une puissance de $J$.
Cela achève la démonstration grâce à la description de $\text{Exal}_A(B, N)$
comme $\text{Ext}^1_B(\NL_{B/A}, N)$ donnée juste avant le
lemme \ref{examples-defos-lemma-first-order-completion}.
\end{proof}

\begin{lemma}
\label{examples-defos-lemma-smooth-localization}
Dans l'exemple \ref{examples-defos-example-rings}, soit $P$ une $k$-algèbre.
Soit $S \subset P$ une partie multiplicative. Si
\begin{enumerate}
\item $k \to P$ est de type fini ;
\item $\Spec(P) \to \Spec(k)$ est lisse en tout point de
$V(g)$ pour tout $g \in S$.
\end{enumerate}
alors le foncteur entre catégories de déformations du
lemme \ref{examples-defos-lemma-localization}
$$
\Deformationcategory_P \longrightarrow \Deformationcategory_{S^{-1}P}
$$
est lisse et induit un isomorphisme sur les espaces tangents.
\end{lemma}

\begin{proof}
Nous savons que $\Deformationcategory_P$ et $\Deformationcategory_{S^{-1}P}$
sont des catégories de déformations d'après le lemme \ref{examples-defos-lemma-rings-RS}.
Il suffit donc de vérifier que
notre foncteur identifie les espaces tangents et les possibilités
de relèvement ; voir Théorie formelle des déformations, lemme
\ref{formal-defos-lemma-easy-check-smooth}.
La propriété concernant les possibilités de relèvement est démontrée dans le
lemme \ref{examples-defos-lemma-lift-equivalence-localization},
et l'isomorphisme sur les espaces tangents est le cas particulier du
lemme \ref{examples-defos-lemma-first-order-localization} où $N = B$.
\end{proof}



\section{Déformations des hensélisés}
\label{examples-defos-section-compare-henselization}

\noindent
Dans cette section, nous comparons les problèmes de déformation posés
par une algèbre et par son complété.
Nous discutons d'abord de la « possibilité de relèvement ».

\begin{lemma}
\label{examples-defos-lemma-lift-equivalence-henselization}
Soit $A' \to A$ un homomorphisme surjectif d'anneaux noethériens de noyau nilpotent.
Soit $A \to B$ un homomorphisme d'anneaux plat de type fini.
Soit $\mathfrak b \subset B$ un idéal tel que
$\Spec(B) \to \Spec(A)$ soit syntomique sur le complémentaire de $V(\mathfrak b)$.
Soit $(B^h, \mathfrak b^h)$ l'hensélisation du couple $(B, \mathfrak b)$.
Alors $B$ possède un relèvement plat sur $A'$ si et seulement si $B^h$ possède un relèvement plat sur $A'$.
\end{lemma}

\begin{proof}[Première démonstration]
Cette démonstration triche quelque peu. En effet, si $B$ possède un relèvement plat $B'$, alors,
en prenant l'hensélisé $(B')^h$, nous obtenons un relèvement plat de $B^h$
(comparer avec la démonstration du lemme \ref{examples-defos-lemma-henselization}).
Réciproquement, supposons que $C'$ soit un relèvement plat sur $A'$ de $(B')^h$.
Soit alors $\mathfrak c' \subset C'$ l'image réciproque de
l'idéal $\mathfrak b^h$. Le complété $(C')^\wedge$ de
$C'$ pour $\mathfrak c'$ est alors un relèvement de $B^\wedge$ (détails omis).
Nous voyons donc que $B$ possède un relèvement plat d'après le
lemme \ref{examples-defos-lemma-lift-equivalence}.
\end{proof}

\begin{proof}[Seconde démonstration]
Choisissons une surjection de $A$-algèbres $P = A[x_1, \ldots, x_n] \to B$.
Soit $\mathfrak p \subset P$ l'image réciproque de $\mathfrak b$.
Posons $P' = A'[x_1, \ldots, x_n]$ et notons $\mathfrak p' \subset P'$
l'image réciproque de $\mathfrak p$. (Bien sûr, $\mathfrak p$
et $\mathfrak p'$ ne désignent pas ici des idéaux premiers.)
Nous noterons $P^h$ et $(P')^h$ les hensélisés respectifs.
Nous utiliserons la fonctorialité de l'hensélisation ainsi que le fait que
l'hensélisé d'un quotient est le quotient correspondant
de l'hensélisé ; voir
Compléments sur l'algèbre, lemmes
\ref{more-algebra-lemma-irreducible-henselian-pair-connected} et
\ref{more-algebra-lemma-henselization-integral}.

\medskip\noindent
Supposons que $A' \to B'$ soit un relèvement plat de $A \to B$, autrement dit
que $A' \to B'$ soit plat et qu'il existe un isomorphisme de $A$-algèbres
$B = B' \otimes_{A'} A$. Nous pouvons alors choisir un homomorphisme de $A'$-algèbres
$P' \to B'$ relevant la surjection donnée $P \to B$.
Par le lemme de Nakayama (Algèbre, lemme \ref{algebra-lemma-NAK}),
nous trouvons que $B'$ est un quotient de $P'$. En particulier, nous trouvons
que nous pouvons rendre $B'$ plat sur $A'$ comme $P'$-module,
relevant $B$, qui est plat sur $A$, comme $P$-module.
Réciproquement, si nous pouvons relever $B$ en un $P'$-module $M'$ plat sur $A'$,
alors $M'$ est un module cyclique $M' \cong P'/J'$ (en utilisant de nouveau Nakayama),
et, en posant $B' = P'/J'$, nous obtenons un relèvement plat de $B$ comme algèbre.

\medskip\noindent
Posons $C = B^h$ et $\mathfrak c = \mathfrak bC$.
Supposons que $A' \to C'$ soit un relèvement plat de $A \to C$.
Alors $C'$ est hensélien pour l'image réciproque
$\mathfrak c'$ de $\mathfrak c$
(d'après Compléments sur l'algèbre, lemme \ref{more-algebra-lemma-henselian-henselian-pair},
et parce que le noyau de $C' \to C$ est nilpotent).
Choisissons un homomorphisme de $A'$-algèbres $P' \to C'$ relevant
l'homomorphisme de $A$-algèbres $P \to C$. Ces homomorphismes se prolongent aux
hensélisés et donnent des surjections $P^h \to C$ et $(P')^h \to C'$
(pour la seconde, en utilisant de nouveau le lemme de Nakayama).
En particulier, nous trouvons que $C'$ est plat sur $A'$ et peut être muni d'une structure de
$(P')^h$-module relevant $C$, plat sur $A$, comme $P^h$-module.
Réciproquement, si nous pouvons relever $C$ en un $(P')^h$-module $N'$ plat sur $A'$,
alors $N'$ est un module cyclique $N' \cong (P')^h/\tilde J$
(en utilisant de nouveau Nakayama) et, en posant $C' = (P')^h/\tilde J$,
nous obtenons un relèvement plat de $C$ comme algèbre.

\medskip\noindent
Observons que $P' \to (P')^h$ est un homomorphisme d'anneaux plat qui
induit un isomorphisme $P'/\mathfrak p' = (P')^h/\mathfrak p'(P')^h$
(Compléments sur l'algèbre, lemme \ref{more-algebra-lemma-henselization-flat}).
Nous concluons que notre lemme résulte du
lemme \ref{examples-defos-lemma-lift-equivalence-module}, pourvu que nous puissions
montrer que $B_g$ se relève en un module plat sur $A'$ muni d'une structure de $P'_g$-module pour
$g \in \mathfrak p'$. Or l'homomorphisme d'anneaux $A \to B_g$ est syntomique
et se relève donc en une $A'$-algèbre plate $B'$ d'après
Lissage des homomorphismes d'anneaux, proposition \ref{smoothing-proposition-lift-smooth}.
Puisque $A' \to P'_g$ est lisse, nous pouvons relever $P_g \to B_g$
en une application surjective $P'_g \to B'$ comme ci-dessus, ce qui donne le résultat voulu.
\end{proof}

\begin{lemma}
\label{examples-defos-lemma-first-order-henselization}
Soit $k$ un corps. Soit $B$ une $k$-algèbre de type fini.
Soit $J \subset B$ un idéal tel que
$\Spec(B) \to \Spec(k)$ soit lisse sur le complémentaire de $V(J)$.
Soit $N$ un $B$-module fini.
Alors il existe une bijection canonique
$$
\text{Exal}_k(B, N) \to \text{Exal}_k(B^h, N^h)
$$
Ici, $(B^h, J^h)$ est l'hensélisation de $(B, J)$
et $N^h = N \otimes_B B^h$.
\end{lemma}

\begin{proof}
Cette démonstration est la même que celle du
lemme \ref{examples-defos-lemma-first-order-completion}, mais plus facile. Nous suggérons au
lecteur de l'omettre.
L'application est donnée par hensélisation : étant donné $0 \to N \to C \to B \to 0$
dans $\text{Exal}_k(B, N)$, nous l'envoyons sur
l'hensélisé $C^h$
de $C$ pour l'image réciproque $J_C \subset C$ de $J$.
Comparer avec la démonstration du lemme \ref{examples-defos-lemma-henselization}.

\medskip\noindent
Puisque $k \to B$ est de présentation finie, le complexe
$\NL_{B/k}$ peut être représenté par un complexe
$N^{-1} \to N^0$, où $N^i$ est un $B$-module fini ; voir
Algèbre, section \ref{algebra-section-netherlander}, et
en particulier
Algèbre, lemme \ref{algebra-lemma-NL-homotopy}.
Comme $B$ est noethérien, cela signifie que $\NL_{B/k}$
est pseudo-cohérent. Pour $g \in J$, la $k$-algèbre $B_g$
est lisse et donc $(\NL_{B/k})_g = \NL_{B_g/k}$
est quasi-isomorphe à un $B$-module projectif de type fini placé en degré $0$.
Ainsi, $\text{Ext}^i_B(\NL_{B/k}, N)_g = 0$ pour $i \geq 1$
et tout $B$-module $N$. Enfin, nous avons
\begin{align*}
\text{Ext}^1_{B^h}(\NL_{B^h/k}, N^h)
& =
\text{Ext}^1_{B^h}(\NL_{B/k} \otimes_B B^h, N \otimes_B B^h) \\
& =
\text{Ext}^1_B(\NL_{B/k}, N) \otimes_B B^h \\
& =
\text{Ext}^1_B(\NL_{B/k}, N)
\end{align*}
La première égalité résulte de
Compléments sur l'algèbre, lemme \ref{more-algebra-lemma-henselization-NL}
(ou plutôt de son analogue pour les hensélisations de couples).
La deuxième résulte de
Compléments sur l'algèbre, lemme \ref{more-algebra-lemma-base-change-RHom}.
La troisième vaut parce que $\text{Ext}^1_B(\NL_{B/k}, N)$ est annulé par une puissance de $J$,
que l'application $B \to B^h$ est plate et induit un isomorphisme
$B/J \to B^h/JB^h$ (Compléments sur l'algèbre, lemme
\ref{more-algebra-lemma-henselization-flat}), et d'après
Compléments sur l'algèbre, lemme \ref{more-algebra-lemma-neighbourhood-equivalence}.
Cela achève la démonstration grâce à la description de $\text{Exal}_A(B, N)$
comme $\text{Ext}^1_B(\NL_{B/A}, N)$ donnée juste avant le
lemme \ref{examples-defos-lemma-first-order-completion}.
\end{proof}

\begin{lemma}
\label{examples-defos-lemma-smooth-henselization}
Dans l'exemple \ref{examples-defos-example-rings}, soit $P$ une $k$-algèbre.
Soit $J \subset P$ un idéal.
Notons $(P^h, J^h)$ l'hensélisation du couple $(P, J)$. Si
\begin{enumerate}
\item $k \to P$ est de type fini ;
\item $\Spec(P) \to \Spec(k)$ est lisse sur le complémentaire de $V(J)$,
\end{enumerate}
alors le foncteur entre catégories de déformations du
lemme \ref{examples-defos-lemma-henselization}
$$
\Deformationcategory_P \longrightarrow \Deformationcategory_{P^h}
$$
est lisse et induit un isomorphisme sur les espaces tangents.
\end{lemma}

\begin{proof}
Nous savons que $\Deformationcategory_P$ et $\Deformationcategory_{P^h}$
sont des catégories de déformations d'après le lemme \ref{examples-defos-lemma-rings-RS}.
Il suffit donc de vérifier que
notre foncteur identifie les espaces tangents et les possibilités
de relèvement ; voir Théorie formelle des déformations, lemme
\ref{formal-defos-lemma-easy-check-smooth}.
La propriété concernant les possibilités de relèvement est démontrée dans le
lemme \ref{examples-defos-lemma-lift-equivalence-henselization},
et l'isomorphisme sur les espaces tangents est le cas particulier du
lemme \ref{examples-defos-lemma-first-order-henselization} où $N = B$.
\end{proof}



\section{Application aux singularités isolées}
\label{examples-defos-section-isolated}

\noindent
Nous appliquons la discussion précédente à l'étude de la théorie des déformations
d'une algèbre de type fini qui ne possède qu'un nombre fini de points singuliers.

\begin{lemma}
\label{examples-defos-lemma-isolated}
Dans l'exemple \ref{examples-defos-example-rings}, soit $P$ une $k$-algèbre.
Supposons que $k \to P$ soit de type fini et que $\Spec(P) \to \Spec(k)$
soit lisse sauf aux idéaux maximaux
$\mathfrak m_1, \ldots, \mathfrak m_n$ de $P$.
Soient $P_{\mathfrak m_i}$, $P_{\mathfrak m_i}^h$, $P_{\mathfrak m_i}^\wedge$
l'anneau local, l'hensélisé et le complété.
Alors les applications de catégories de déformations
$$
\Deformationcategory_P \to
\prod \Deformationcategory_{P_{\mathfrak m_i}} \to
\prod \Deformationcategory_{P_{\mathfrak m_i}^h} \to
\prod \Deformationcategory_{P_{\mathfrak m_i}^\wedge}
$$
sont lisses et induisent des isomorphismes sur leurs espaces tangents
de dimension finie.
\end{lemma}

\begin{proof}
L'espace tangent est de dimension finie d'après le
lemme \ref{examples-defos-lemma-finite-type-rings-TI}.
Les foncteurs entre les catégories sont construits
dans les lemmes \ref{examples-defos-lemma-localization}, \ref{examples-defos-lemma-henselization} et
\ref{examples-defos-lemma-completion} (nous omettons certaines vérifications du type :
le complété de l'hensélisé est le complété).

\medskip\noindent
Posons $J = \mathfrak m_1 \cap \ldots \cap \mathfrak m_n$ et appliquons le
lemme \ref{examples-defos-lemma-smooth-completion} pour obtenir que
$\Deformationcategory_P \to \Deformationcategory_{P^\wedge}$
est lisse et induit un isomorphisme sur les espaces tangents,
où $P^\wedge$ est le complété $J$-adique de $P$.
Cependant, puisque $P^\wedge = \prod P_{\mathfrak m_i}^\wedge$,
nous voyons que l'application $\Deformationcategory_P \to
\prod \Deformationcategory_{P_{\mathfrak m_i}^\wedge}$
est lisse et induit un isomorphisme sur les espaces tangents.

\medskip\noindent
Soit $(P^h, J^h)$ l'hensélisation du couple $(P, J)$.
Alors $P^h = \prod P_{\mathfrak m_i}^h$ (considérer les idempotents
et utiliser Compléments sur l'algèbre, lemme
\ref{more-algebra-lemma-characterize-henselian-pair}).
Nous pouvons donc appliquer le lemme \ref{examples-defos-lemma-smooth-henselization}
et conclure comme dans le cas du complété.

\medskip\noindent
Pour obtenir le dernier cas, il suffit de montrer que
$\Deformationcategory_{P_{\mathfrak m_i}} \to
\Deformationcategory_{P_{\mathfrak m_i}^\wedge}$
est lisse et induit des isomorphismes sur les espaces tangents pour chaque $i$ séparément.
Pour cela, nous pouvons remplacer $P$ par une localisation principale
dont l'unique point singulier est un idéal maximal $\mathfrak m$
(correspondant à $\mathfrak m_i$ dans le $P$ initial).
Nous pouvons alors appliquer le
lemme \ref{examples-defos-lemma-smooth-localization}
à la partie multiplicative $S = P \setminus \mathfrak m$ pour conclure.
Les détails mineurs sont omis.
\end{proof}






\section{Problèmes de déformation non obstrués}
\label{examples-defos-section-unobstructed}

\noindent
Soit $p : \mathcal{F} \to \mathcal{C}_\Lambda$ une
catégorie cofibrée en groupoïdes. Rappelons que nous disons que $\mathcal{F}$
est {\it lisse} ou {\it non obstruée} si $p$ est lisse.
Cela signifie qu'étant données une surjection $\varphi : A' \to A$ dans
$\mathcal{C}_\Lambda$ et $x \in \Ob(\mathcal{F}(A))$,
il existe un morphisme $f : x' \to x$ dans $\mathcal{F}$
tel que $p(f) = \varphi$.
Voir Théorie formelle des déformations, section \ref{formal-defos-section-smooth}.
Dans cette section, nous donnons quelques exemples géométriquement significatifs.

\begin{lemma}
\label{examples-defos-lemma-lci-unobstructed}
Dans l'exemple \ref{examples-defos-example-rings}, soit $P$ une intersection complète
locale sur $k$ (Algèbre, définition \ref{algebra-definition-lci-field}).
Alors $\Deformationcategory_P$ est non obstruée.
\end{lemma}

\begin{proof}
Soit $(A, Q) \to (k, P)$ un objet de $\Deformationcategory_P$.
Alors nous voyons que $A \to Q$ est un homomorphisme d'anneaux syntomique d'après
Algèbre, définition \ref{algebra-definition-lci}.
Ainsi, pour toute surjection $A' \to A$ dans $\mathcal{C}_\Lambda$,
nous voyons qu'il existe un morphisme $(A', Q') \to (A, Q)$
relevant $A' \to A$ d'après
Lissage des homomorphismes d'anneaux, proposition \ref{smoothing-proposition-lift-smooth}.
Cela démontre le lemme.
\end{proof}

\begin{lemma}
\label{examples-defos-lemma-glueing-smooth}
Dans la situation \ref{examples-defos-situation-glueing}, si $U_{12} \to \Spec(k)$ est lisse,
alors le morphisme
$$
\Deformationcategory_X
\longrightarrow
\Deformationcategory_{U_1} \times \Deformationcategory_{U_2} =
\Deformationcategory_{P_1} \times \Deformationcategory_{P_2}
$$
est lisse. Si, de plus,
$U_1$ est une intersection complète locale sur $k$, alors
$$
\Deformationcategory_X
\longrightarrow
\Deformationcategory_{U_2} = \Deformationcategory_{P_2}
$$
est lisse.
\end{lemma}

\begin{proof}
Les signes d'égalité sont justifiés par le lemme \ref{examples-defos-lemma-affine}.
Considérons $\mathcal{C}_\Lambda$ comme une catégorie de
déformations sur $\mathcal{C}_\Lambda$, comme dans
Théorie formelle des déformations, section \ref{formal-defos-section-smooth}.
Alors
$$
\Deformationcategory_{P_1} \times \Deformationcategory_{P_2} =
\Deformationcategory_{P_1}
\times_{\mathcal{C}_\Lambda}
\Deformationcategory_{P_2},
$$
voir Théorie formelle des déformations, remarques
\ref{formal-defos-remarks-cofibered-groupoids}
(\ref{formal-defos-item-product}).
En utilisant le
lemme \ref{examples-defos-lemma-glueing},
la première assertion revient à dire que le foncteur
$$
\Deformationcategory_{P_1}
\times_{\Deformationcategory_{P_{12}}}
\Deformationcategory_{P_2}
\longrightarrow
\Deformationcategory_{P_1}
\times_{\mathcal{C}_\Lambda}
\Deformationcategory_{P_2}
$$
est lisse. Cela résulte de Théorie formelle des déformations, lemme
\ref{formal-defos-lemma-map-fibre-products-smooth}, pourvu que
nous puissions montrer que $T\Deformationcategory_{P_{12}} = (0)$.
Cette annulation résulte du lemme \ref{examples-defos-lemma-smooth},
puisque $P_{12}$ est lisse sur $k$.
Pour la seconde assertion, il suffit de montrer que
$\Deformationcategory_{P_1} \to \mathcal{C}_\Lambda$
est lisse ; voir Théorie formelle des déformations, lemme
\ref{formal-defos-lemma-smooth-properties}.
Autrement dit, il faut montrer que $\Deformationcategory_{P_1}$
est non obstruée, ce qui est le lemme \ref{examples-defos-lemma-lci-unobstructed}.
\end{proof}

\begin{lemma}
\label{examples-defos-lemma-curve-isolated}
Dans l'exemple \ref{examples-defos-example-schemes}, soit $X$ un schéma sur $k$. Supposons que
\begin{enumerate}
\item $X$ soit séparé, de type fini sur $k$ et $\dim(X) \leq 1$ ;
\item $X \to \Spec(k)$ soit lisse sauf aux
points fermés $p_1, \ldots, p_n \in X$.
\end{enumerate}
Soient $\mathcal{O}_{X, p_1}$, $\mathcal{O}_{X, p_1}^h$,
$\mathcal{O}_{X, p_1}^\wedge$ l'anneau local, l'hensélisé et le complété.
Considérons les applications de catégories de déformations
$$
\Deformationcategory_X
\longrightarrow
\prod \Deformationcategory_{\mathcal{O}_{X, p_i}}
\longrightarrow
\prod \Deformationcategory_{\mathcal{O}_{X, p_i}^h}
\longrightarrow
\prod \Deformationcategory_{\mathcal{O}_{X, p_i}^\wedge}
$$
La première flèche est lisse, et les deuxième et troisième flèches
sont lisses et induisent des isomorphismes sur les espaces tangents.
\end{lemma}

\begin{proof}
Choisissons un ouvert affine $U_2 \subset X$ contenant
$p_1, \ldots, p_n$ et le point générique de toute composante
irréductible de $X$. Cela est possible d'après
Variétés, lemme \ref{varieties-lemma-dim-1-quasi-projective},
et Propriétés, lemme \ref{properties-lemma-ample-finite-set-in-affine}.
Alors $X \setminus U_2$ est fini, et nous pouvons choisir un ouvert affine
$U_1 \subset X \setminus \{p_1, \ldots, p_n\}$ tel que
$X = U_1 \cup U_2$. Posons $U_{12} = U_1 \cap U_2$.
Alors $U_1$ et $U_{12}$ sont des schémas affines lisses sur $k$.
Nous concluons que
$$
\Deformationcategory_X \longrightarrow \Deformationcategory_{U_2}
$$
est lisse d'après le lemme \ref{examples-defos-lemma-glueing-smooth}.
Le résultat découle des lemmes \ref{examples-defos-lemma-affine} et \ref{examples-defos-lemma-isolated}.
\end{proof}

\begin{lemma}
\label{examples-defos-lemma-curve-isolated-lci}
Dans l'exemple \ref{examples-defos-example-schemes}, soit $X$ un schéma sur $k$. Supposons que
\begin{enumerate}
\item $X$ soit séparé, de type fini sur $k$ et $\dim(X) \leq 1$ ;
\item $X$ soit une intersection complète locale sur $k$ ;
\item $X \to \Spec(k)$ soit lisse sauf en un nombre fini de points.
\end{enumerate}
Alors $\Deformationcategory_X$ est non obstruée.
\end{lemma}

\begin{proof}
Soient $p_1, \ldots, p_n \in X$ les points où $X \to \Spec(k)$
n'est pas lisse. Choisissons un ouvert affine $U_2 \subset X$ contenant
$p_1, \ldots, p_n$ et le point générique de toute composante
irréductible de $X$. Cela est possible d'après
Variétés, lemme \ref{varieties-lemma-dim-1-quasi-projective},
et Propriétés, lemme \ref{properties-lemma-ample-finite-set-in-affine}.
Alors $X \setminus U_2$ est fini, et nous pouvons choisir un ouvert affine
$U_1 \subset X \setminus \{p_1, \ldots, p_n\}$ tel que
$X = U_1 \cup U_2$. Posons $U_{12} = U_1 \cap U_2$.
Alors $U_1$ et $U_{12}$ sont des schémas affines lisses sur $k$.
Nous concluons que
$$
\Deformationcategory_X \longrightarrow \Deformationcategory_{U_2}
$$
est lisse d'après le lemme \ref{examples-defos-lemma-glueing-smooth}.
Le résultat découle des lemmes \ref{examples-defos-lemma-affine} et \ref{examples-defos-lemma-lci-unobstructed}.
\end{proof}




\section{Lissages}
\label{examples-defos-section-smoothing}

\noindent
Soit donné un schéma ou espace algébrique de type fini $X$ sur un corps $k$.
Il est souvent utile de trouver un morphisme plat de type fini $Y \to \Spec(k[[t]])$
dont la fibre générique est lisse et la fibre spéciale isomorphe à $X$.
Un tel objet est appelé un lissage de $X$. Dans cette section, nous trouverons
un lissage en dimension $1$ pour un espace séparé $X$ qui possède des singularités
isolées d'intersection complète locale.

\begin{lemma}
\label{examples-defos-lemma-criterion-smoothing}
Soit $k$ un corps. Posons $S = \Spec(k[[t]])$ et
$S_n = \Spec(k[t]/(t^n))$. Soit $Y \to S$ un morphisme propre et plat
de schémas dont la fibre spéciale $X$ est de Cohen-Macaulay et
équidimensionnelle de dimension $d$. Notons $X_n = Y \times_S S_n$.
Si, pour un certain $n \geq 1$, le $d$-ième idéal de Fitting de $\Omega_{X_n/S_n}$
contient $t^{n - 1}$, alors la fibre générique de $Y \to S$ est lisse.
\end{lemma}

\begin{proof}
D'après Compléments sur les morphismes, lemme
\ref{more-morphisms-lemma-flat-finite-presentation-CM-open}
nous voyons que $Y \to S$ est un morphisme de Cohen-Macaulay.
D'après Morphismes, lemme
\ref{morphisms-lemma-flat-finite-presentation-CM-fibres-relative-dimension}
nous voyons que $Y \to S$ est de dimension relative $d$.
D'après Diviseurs, lemme \ref{divisors-lemma-d-fitting-ideal-omega-smooth},
le $d$-ième idéal de Fitting $\mathcal{I} \subset \mathcal{O}_Y$
de $\Omega_{Y/S}$ définit le lieu singulier du morphisme $Y \to S$.
Autrement dit, $V(\mathcal{I}) \subset Y$ est le sous-ensemble fermé
des points où $Y \to S$ n'est pas lisse.
D'après Diviseurs, lemme \ref{divisors-lemma-base-change-and-fitting-ideal-omega},
la formation de cet idéal de Fitting commute au changement de base.
Par hypothèse, nous voyons que $t^{n - 1}$ est une section de
$\mathcal{I} + t^n\mathcal{O}_Y$. Ainsi, pour tout
$x \in X = V(t) \subset Y$, nous concluons que
$t^{n - 1} \in \mathcal{I}_x$, où $\mathcal{I}_x$ est le germe en $x$.
Cela implique que $V(\mathcal{I}) \subset V(t)$ dans un
voisinage ouvert de $X$ dans $Y$. Puisque $Y \to S$
est propre, il s'ensuit que $V(\mathcal{I}) \subset V(t)$,
comme souhaité.
\end{proof}

\begin{lemma}
\label{examples-defos-lemma-jouanolou-type-thing}
Soit $k$ un corps. Soient $1 \leq c \leq n$ des entiers.
Soient $f_1, \ldots, f_c \in k[x_1, \ldots x_n]$ des éléments.
Soient $a_{ij}$, $0 \leq i \leq n$, $1 \leq j \leq c$ des
indéterminées. Considérons
$$
g_j = f_j + a_{0j} + a_{1j}x_1 + \ldots + a_{nj}x_n \in
k[a_{ij}][x_1, \ldots, x_n]
$$
Notons $Y \subset \mathbf{A}^{n + c(n + 1)}_k$
le sous-schéma fermé défini par $g_1, \ldots, g_c$.
Notons $\pi : Y \to \mathbf{A}^{c(n + 1)}_k$ la projection
sur l'espace affine d'indéterminées $a_{ij}$.
Alors il existe un ouvert de Zariski non vide
de $\mathbf{A}^{c(n + 1)}_k$ au-dessus duquel $\pi$ est lisse.
\end{lemma}

\begin{proof}
Rappelons que l'ensemble des points où $\pi$ est lisse est ouvert.
Son complémentaire, c'est-à-dire le lieu singulier, est donc fermé.
D'après le théorème de Chevalley (sous la forme de
Morphismes, lemme \ref{morphisms-lemma-chevalley}),
l'image du lieu singulier est constructible.
Par conséquent, si le point générique de $\mathbf{A}^{c(n + 1)}_k$
n'appartient pas à l'image du lieu singulier, alors
le lemme en découle (par exemple d'après Topologie, lemme
\ref{topology-lemma-generic-point-in-constructible}).
Il faut donc montrer qu'il n'existe aucun point
$y \in Y$ où $\pi$ ne soit pas lisse et qui s'envoie sur
le point générique de $\mathbf{A}^{c(n + 1)}_k$.
Considérons la matrice des dérivées partielles
$$
(\frac{\partial g_j}{\partial x_i}) =
(\frac{\partial f_j}{\partial x_i} + a_{ij})
$$
L'image de cette matrice dans $\kappa(y)$ doit être de rang $< c$,
car sinon $\pi$ serait lisse en $y$ ; voir la discussion dans
Lissage des homomorphismes d'anneaux, section \ref{smoothing-section-singular-ideal}.
Nous pouvons donc trouver $\lambda_1, \ldots, \lambda_c \in \kappa(y)$,
non tous nuls, tels que le vecteur $(\lambda_1, \ldots, \lambda_c)$
appartienne au noyau de cette matrice.
Après renumérotation, nous pouvons supposer $\lambda_1 \not = 0$.
En divisant par $\lambda_1$, nous pouvons supposer que notre vecteur est
de la forme $(1, \lambda_2, \ldots, \lambda_c)$.
Nous obtenons alors
$$
a_{i1} = -
\frac{\partial f_j}{\partial x_1} -
\sum\nolimits_{j = 2, \ldots, c} \lambda_j(\frac{\partial f_j}{\partial x_i} + a_{ij})
$$
dans $\kappa(y)$ pour $i = 1, \ldots, n$. De plus, puisque $y \in Y$, nous avons aussi
l'égalité
$$
a_{0j} = -f_j - a_{1j}x_1 - \ldots - a_{nj}x_n
$$
dans $\kappa(y)$. Cela signifie que le sous-corps de $\kappa(y)$
engendré par $a_{ij}$ est contenu dans le sous-corps de $\kappa(y)$
engendré par les images de $x_1, \ldots, x_n, \lambda_2, \ldots, \lambda_c$,
et par les $a_{ij}$ autres que $a_{i1}$ et $a_{0j}$.
En comptant, nous voyons que son degré de transcendance est
au plus $c(n + 1) - 1$. Ainsi, $y$ ne peut pas s'envoyer sur le point générique,
comme souhaité.
\end{proof}

\begin{lemma}
\label{examples-defos-lemma-smoothing-affine-lci}
Soit $k$ un corps. Soit $A$ une intersection complète globale
sur $k$. Il existe un homomorphisme d'anneaux plat de type fini
$k[[t]] \to B$ tel que $B/tB \cong A$ et que
$B[1/t]$ soit lisse sur $k((t))$.
\end{lemma}

\begin{proof}
Écrivons $A = k[x_1, \ldots, x_n]/(f_1, \ldots, f_c)$ comme dans
Algèbre, définition \ref{algebra-definition-lci-field}.
Nous allons choisir
$a_{ij} \in (t) \subset k[[t]]$ et poser
$$
g_j = f_j + a_{0j} + a_{1j}x_1 + \ldots + a_{nj}x_n \in
k[[t]][x_1, \ldots, x_n]
$$
Après ce choix, nous prenons
$B = k[[t]][x_1, \ldots, x_n]/(g_1, \ldots, g_c)$.
Nous affirmons que $k[[t]] \to B$ est plat en tout idéal premier
au-dessus de $(t)$. En effet, les éléments $f_1, \ldots, f_c$
forment une suite régulière dans l'anneau local en tout idéal premier
$\mathfrak p$ de $k[x_1, \ldots, x_n]$ contenant $f_1, \ldots, f_c$
(Algèbre, lemme \ref{algebra-lemma-lci}). Ainsi, $g_1, \ldots, g_c$
est localement un relèvement d'une suite régulière, et nous pouvons appliquer
Algèbre, lemme \ref{algebra-lemma-grothendieck-regular-sequence}.
La platitude aux idéaux premiers au-dessus de $(0) \subset k[[t]]$ est automatique,
car $k((t)) = k[[t]]_{(0)}$ est un corps. Ainsi, $B$ est plat
sur $k[[t]]$.

\medskip\noindent
Il reste seulement à montrer que, pour des choix convenables
des $a_{ij}$, la fibre générique $B_{(0)}$ est lisse sur
$k((t))$. Pour cela, nous devons montrer que nous pouvons choisir
nos $a_{ij}$ de sorte que le morphisme induit
$$
(a_{ij}) : \Spec(k[[t]]) \longrightarrow \mathbf{A}^{c(n + 1)}_k
$$
s'envoie dans l'ouvert de Zariski non vide du
lemme \ref{examples-defos-lemma-jouanolou-type-thing}.
Cela est clair, car aucun polynôme non nul en les
$a_{ij}$ ne s'annule sur $(t)^{\oplus c(n + 1)}$.
(Nous laissons cela en exercice au lecteur.)
\end{proof}

\begin{lemma}
\label{examples-defos-lemma-smoothing-artinian-lci}
Soit $k$ un corps. Soit $A$ une $k$-algèbre de dimension finie
qui est une intersection complète locale sur $k$. Alors il existe
une $k[[t]]$-algèbre finie et plate $B$ telle que $B/tB \cong A$
et $B[1/t]$ soit \'etale sur $k((t))$.
\end{lemma}

\begin{proof}
Puisque $A$ est artinien
(Algèbre, lemme \ref{algebra-lemma-finite-dimensional-algebra}),
nous pouvons écrire $A$ comme un produit d'anneaux locaux artiniens
(Algèbre, lemme \ref{algebra-lemma-artinian-finite-length}).
Il suffit donc de démontrer le lemme lorsque $A$ est local
(nous utilisons ici que la propriété d'être une intersection complète locale
est préservée par localisation principale ; voir
Algèbre, lemme \ref{algebra-lemma-localize-lci}).
Dans ce cas, $A$ est une intersection complète globale.
Considérons l'algèbre $B$ construite dans le
lemme \ref{examples-defos-lemma-smoothing-affine-lci}.
Alors $k[[t]] \to B$ est quasi-fini en l'unique idéal premier de $B$
au-dessus de $(t)$ (Algèbre, définition \ref{algebra-definition-quasi-finite}).
Observons que $k[[t]]$ est un anneau local hensélien
(Algèbre, lemme \ref{algebra-lemma-complete-henselian}).
Ainsi, $B = B' \times C$, où $B'$ est fini sur $k[[t]]$
et $C$ n'a aucun idéal premier au-dessus de $(t)$ ; voir
Algèbre, lemme \ref{algebra-lemma-characterize-henselian}.
Alors $B'$ est l'anneau recherché
(rappelons que \'etale signifie la même chose que
lisse de dimension relative $0$).
\end{proof}

\begin{lemma}
\label{examples-defos-lemma-smoothing-at-lci-point}
Soit $k$ un corps. Soit $A$ une $k$-algèbre. Supposons que
\begin{enumerate}
\item $A$ soit un anneau local essentiellement de type fini sur $k$ ;
\item $A$ soit une intersection complète sur $k$
(Algèbre, définition \ref{algebra-definition-lci-local-ring}).
\end{enumerate}
Posons $d = \dim(A) + \text{trdeg}_k(\kappa)$, où $\kappa$
est le corps résiduel de $A$. Alors il existe un entier $n$
et un homomorphisme d'anneaux plat, essentiellement de type fini,
$k[[t]] \to B$ tel que $B/tB \cong A$ et que $t^n$ appartienne au
$d$-ième idéal de Fitting de $\Omega_{B/k[[t]]}$.
\end{lemma}

\begin{proof}
D'après Algèbre, lemme \ref{algebra-lemma-lci-local}, nous pouvons écrire $A$ comme la
localisation en un idéal premier $\mathfrak p$ d'une intersection complète globale $P$
sur $k$. Observons que $\dim(P) = d$ d'après
Algèbre, lemme \ref{algebra-lemma-dimension-at-a-point-finite-type-field}.
D'après le lemme \ref{examples-defos-lemma-smoothing-affine-lci}, nous pouvons trouver un
homomorphisme d'anneaux plat de type fini $k[[t]] \to Q$ tel que $P \cong Q/tQ$ et
tel que $k((t)) \to Q[1/t]$ soit lisse. Il résulte de la construction
de $Q$ dans le lemme que $k[[t]] \to Q$ est un morphisme
d'intersection complète globale, de dimension relative $d$ ; autrement dit,
Algèbre, lemme \ref{algebra-lemma-syntomic}, nous dit que $Q$ ou une
localisation principale convenable de $Q$ est une telle intersection complète globale.
Ainsi, d'après Diviseurs, lemme \ref{divisors-lemma-d-fitting-ideal-omega-smooth},
le $d$-ième idéal de Fitting $I \subset Q$ de $\Omega_{Q/k[[t]]}$
définit le lieu singulier de $\Spec(Q) \to \Spec(k[[t]])$.
Ainsi, $t^n \in I$ pour un certain $n$.
Soit $\mathfrak q \subset Q$
l'image réciproque de $\mathfrak p$. Posons $B = Q_\mathfrak q$.
Le lemme est démontré.
\end{proof}

\begin{lemma}
\label{examples-defos-lemma-smoothing-proper-curve-isolated-lci}
Soit $X$ un schéma sur un corps $k$. Supposons que
\begin{enumerate}
\item $X$ soit propre sur $k$ ;
\item $X$ soit une intersection complète locale sur $k$ ;
\item $X$ soit de dimension $\leq 1$ ;
\item $X \to \Spec(k)$ soit lisse sauf en un nombre fini de points.
\end{enumerate}
Alors il existe un morphisme projectif plat $Y \to \Spec(k[[t]])$
dont la fibre générique est lisse et la fibre spéciale
est isomorphe à $X$.
\end{lemma}

\begin{proof}
Observons que $X$ est de Cohen-Macaulay ; voir
Algèbre, lemme \ref{algebra-lemma-lci-CM}.
Ainsi, $X = X' \amalg X''$ avec $\dim(X') = 0$,
et $X''$ est équidimensionnel de dimension $1$ ; voir Morphismes, lemme
\ref{morphisms-lemma-flat-finite-presentation-CM-fibres-relative-dimension}.
Puisque $X'$ est fini sur $k$ (Variétés, lemme
\ref{varieties-lemma-algebraic-scheme-dim-0})
nous pouvons trouver $Y' \to \Spec(k[[t]])$ dont la fibre
spéciale est $X'$ et la fibre générique est lisse d'après le
lemme \ref{examples-defos-lemma-smoothing-artinian-lci}.
Il suffit donc de démontrer le lemme pour $X''$.
Après avoir remplacé $X$ par $X''$, le schéma $X$ est
de Cohen-Macaulay et équidimensionnel de dimension $1$.

\medskip\noindent
Nous allons utiliser la théorie des déformations dans la situation $\Lambda = k \to k$.
Soient $p_1, \ldots, p_r \in X$ les points singuliers fermés de $X$, c'est-à-dire
les points où $X \to \Spec(k)$ n'est pas lisse. Pour chaque $i$, choisissons
un entier $n_i$ et un homomorphisme d'anneaux plat, essentiellement de type fini,
$$
k[[t]] \longrightarrow B_i
$$
tel que $B_i/tB_i \cong \mathcal{O}_{X, p_i}$ et que
$t^{n_i}$ appartienne au $1$er idéal de Fitting de $\Omega_{B_i/k[[t]]}$.
Cela est possible d'après le lemme \ref{examples-defos-lemma-smoothing-at-lci-point}.
Observons que le système $(B_i/t^nB_i)$ définit un objet formel de
$\Deformationcategory_{\mathcal{O}_{X, p_i}}$ sur $k[[t]]$.
D'après le lemme \ref{examples-defos-lemma-curve-isolated}, l'application
$$
\Deformationcategory_X
\longrightarrow
\prod\nolimits_{i = 1, \ldots, r} \Deformationcategory_{\mathcal{O}_{X, p_i}}
$$
est une application lisse entre catégories de déformations. Ainsi, d'après
Théorie formelle des déformations, lemme
\ref{formal-defos-lemma-smooth-morphism-essentially-surjective}
il existe un objet formel $(X_n)$ dans $\Deformationcategory_X$
qui s'envoie sur l'objet formel $\prod_i (B_i/t^n)$ par la flèche ci-dessus.
D'après Compléments sur les morphismes d'espaces, lemme
\ref{spaces-more-morphisms-lemma-formal-algebraic-space-proper-reldim-1}
il existe un schéma projectif $Y$ sur $k[[t]]$ et des isomorphismes compatibles
$Y \times_{\Spec(k[[t]])} \Spec(k[t]/(t^n)) \cong X_n$.
D'après Compléments sur les morphismes, lemme
\ref{more-morphisms-lemma-check-flatness-on-infinitesimal-nbhds}
nous voyons que $Y \to \Spec(k[[t]])$ est plat.
Puisque $X$ est de Cohen-Macaulay et équidimensionnel de dimension $1$,
nous pouvons appliquer le lemme \ref{examples-defos-lemma-criterion-smoothing}
pour vérifier que $Y$ a une fibre générique lisse\footnote{Attention : en général, il n'est
{\bf pas} vrai que l'anneau local de $Y$ au point
$p_i$ soit isomorphe à $B_i$. Nous savons seulement que cela devient vrai après
quotient par $t^n$ des deux côtés !}.
Choisissons $n$ strictement supérieur au maximum des entiers $n_i$ trouvés ci-dessus.
Si nous pouvons montrer que $t^{n - 1}$ appartient au premier idéal de Fitting de
$\Omega_{X_n/S_n}$, où $S_n = \Spec(k[t]/(t^n))$, la démonstration est achevée.
Pour cela, il suffit de le démontrer dans chacun
des anneaux locaux de $X_n$ aux points fermés $p$.
Cependant, si $p$ correspond à un point lisse de $X \to \Spec(k)$,
alors $\Omega_{X_n/S_n, p}$ est libre de rang $1$ et le premier idéal de Fitting
est égal à l'anneau local. Si $p = p_i$ pour un certain $i$, alors
$$
\Omega_{X_n/S_n, p_i} =
\Omega_{(B_i/t^nB_i)/(k[t]/(t^n))} =
\Omega_{B_i/k[[t]]}/t^n\Omega_{B_i/k[[t]]}
$$
Puisque la formation des idéaux de Fitting commute au changement de base
(nous avons déjà utilisé ce fait, mais, dans ce cadre algébrique,
il résulte de Compléments sur l'algèbre, lemme
\ref{more-algebra-lemma-fitting-ideal-basics}),
et puisque $n - 1 \geq n_i$, nous voyons que $t^{n - 1}$ appartient
à l'idéal de Fitting de ce module sur $B_i/t^nB_i$, comme souhaité.
\end{proof}

\begin{lemma}
\label{examples-defos-lemma-smoothing-curve-isolated-lci}
Soit $k$ un corps et soit $X$ un schéma sur $k$. Supposons que
\begin{enumerate}
\item $X$ soit séparé, de type fini sur $k$ et $\dim(X) \leq 1$ ;
\item $X$ soit une intersection complète locale sur $k$ ;
\item $X \to \Spec(k)$ soit lisse sauf en un nombre fini de points.
\end{enumerate}
Alors il existe un morphisme plat, séparé et de type fini $Y \to \Spec(k[[t]])$
dont la fibre générique est lisse et la fibre spéciale
est isomorphe à $X$.
\end{lemma}

\begin{proof}
Si $X$ est réduit, nous pouvons choisir une immersion
$X \subset \overline{X}$ comme dans
Variétés, lemme \ref{varieties-lemma-reduced-dim-1-projective-completion}.
En écrivant $X = \overline{X} \setminus \{x_1, \ldots, x_n\}$,
nous voyons que $\mathcal{O}_{\overline{X}, x_i}$ est un anneau de
valuation discrète, donc en particulier une intersection complète locale
(Algèbre, définition \ref{algebra-definition-lci-local-ring}).
Ainsi, $\overline{X}$ est une intersection complète locale
sur $k$, car cette propriété vaut sur l'ouvert $X$ et
aux points $x_i$ d'après Algèbre, lemme \ref{algebra-lemma-lci-local}.
Nous pouvons donc appliquer le lemme \ref{examples-defos-lemma-smoothing-proper-curve-isolated-lci}
pour trouver un morphisme projectif plat $\overline{Y} \to \Spec(k[[t]])$
dont la fibre générique est lisse et la fibre spéciale
est $\overline{X}$. Nous retirons alors $x_1, \ldots, x_n$
de $\overline{Y}$ pour obtenir $Y$.

\medskip\noindent
Dans le cas général, écrivons $X = X' \amalg X''$, où
$\dim(X') = 0$ et $X''$ est équidimensionnel de dimension $1$.
Alors $X''$ est réduit et le premier paragraphe s'y applique.
D'autre part, $X'$ peut être traité
comme dans la démonstration du lemme \ref{examples-defos-lemma-smoothing-proper-curve-isolated-lci}.
Certains détails sont omis.
\end{proof}









% OK, commencer ici.
%

\chapter{Champs algébriques}



\phantomsection
\label{algebraic-section-phantom}


\section{Introduction}
\label{algebraic-section-introduction}

\noindent
C'est ici que nous définissons les champs algébriques et formulons quelques
observations très élémentaires. La philosophie générale consistera à n'imposer
aucune condition de séparation et à ajouter les conditions nécessaires pour que
les lemmes, propositions et théorèmes soient vrais/démontrables. Ainsi, les notions
étudiées ici diffèrent légèrement de celles que l'on trouve ailleurs dans la littérature, par exemple
\cite{LM-B}.

\medskip\noindent
Ce chapitre ne constitue pas une introduction aux champs algébriques.
Pour une présentation informelle des champs algébriques, voir
Introduction aux champs algébriques, section
\ref{stacks-introduction-section-introduction}.


\section{Conventions}
\label{algebraic-section-conventions}

\noindent
Les conventions utilisées dans ce chapitre sont les mêmes que dans le
chapitre consacré aux espaces algébriques. Nous les rappelons ici par commodité.

\medskip\noindent
Nous travaillons dans un grand site fppf convenable $\Sch_{fppf}$
comme dans Topologies, Définition \ref{topologies-definition-big-fppf-site}.
Ainsi, sauf mention explicite du contraire, tous les schémas seront des objets
de $\Sch_{fppf}$. Nous examinons les changements qu'entraîne le remplacement du grand
site fppf dans la
section \ref{algebraic-section-change-big-site}.

\medskip\noindent
Nous travaillerons toujours relativement à une base $S$ contenue dans $\Sch_{fppf}$.
Nous travaillerons alors avec le grand site fppf $(\Sch/S)_{fppf}$, voir
Topologies, Définition \ref{topologies-definition-big-small-fppf}.
On retrouve le cas absolu en prenant
$S = \Spec(\mathbf{Z})$.

\medskip\noindent
Si $U, T$ sont des schémas sur $S$, nous notons
$U(T)$ l'ensemble de ses points à valeurs dans $T$ {\it au-dessus de} $S$.
En formule : $U(T) = \Mor_S(T, U)$.

\medskip\noindent
Notons que tout recouvrement fpqc est un épimorphisme effectif
universel, voir
Descente, Lemme \ref{descent-lemma-fpqc-universal-effective-epimorphisms}.
Par conséquent, la topologie sur $\Sch_{fppf}$
est moins fine que la topologie canonique et tous les préfaisceaux représentables
sont des faisceaux.








\section{Notations}
\label{algebraic-section-notation}

\noindent
Nous employons les lettres $S, T, U, V, X, Y$ pour désigner des schémas.
Nous employons les lettres $\mathcal{X}, \mathcal{Y}, \mathcal{Z}$ pour désigner des
catégories (fibrées, fibrées en groupoïdes, champs, ...)
au-dessus de $(\Sch/S)_{fppf}$. Nous employons les lettres minuscules
$f$, $g$ pour les foncteurs tels que $f : \mathcal{X} \to \mathcal{Y}$
au-dessus de $(\Sch/S)_{fppf}$.
Nous employons les majuscules $F$, $G$, $H$ pour les espaces algébriques sur $S$, et plus
généralement pour les préfaisceaux d'ensembles sur $(\Sch/S)_{fppf}$.
(Dans les chapitres ultérieurs, nous emploierons de nouveau aussi $X$, $Y$, etc.
pour les espaces algébriques.)

\medskip\noindent
Ces choix visent à distinguer clairement les différents
types d'objets de ce chapitre, afin d'en établir les fondements.









\section{Catégories fibrées en groupoïdes représentables}
\label{algebraic-section-representable}

\noindent
Soit $S$ un schéma contenu dans $\Sch_{fppf}$.
L'objet d'étude fondamental de ce chapitre sera une
catégorie fibrée en groupoïdes
$p : \mathcal{X} \to (\Sch/S)_{fppf}$, voir
Catégories, Définition \ref{categories-definition-fibred-groupoids}.
Nous dirons souvent simplement « soit $\mathcal{X}$ une catégorie fibrée
en groupoïdes au-dessus de $(\Sch/S)_{fppf}$ » pour désigner
cette situation. Un $1$-morphisme $\mathcal{X} \to \mathcal{Y}$ de catégories
fibrées en groupoïdes au-dessus de $(\Sch/S)_{fppf}$ sera un $1$-morphisme
dans la $2$-catégorie des catégories fibrées en groupoïdes au-dessus de
$(\Sch/S)_{fppf}$, voir
Catégories,
Définition \ref{categories-definition-categories-fibred-in-groupoids-over-C}.
Il s'agit simplement d'un foncteur $\mathcal{X} \to \mathcal{Y}$ au-dessus de
$(\Sch/S)_{fppf}$.
Rappelons qu'il s'agit en fait d'une $(2, 1)$-catégorie et que tous les $2$-produits fibrés
existent.

\medskip\noindent
Soit $\mathcal{X}$ une catégorie fibrée en groupoïdes au-dessus de
$(\Sch/S)_{fppf}$. Rappelons que $\mathcal{X}$
est dite {\it représentable} s'il existe un
schéma $U \in \Ob((\Sch/S)_{fppf})$ et une
équivalence
$$
j : \mathcal{X} \longrightarrow (\Sch/U)_{fppf}
$$
de catégories au-dessus de $(\Sch/S)_{fppf}$, voir
Catégories,
Définition \ref{categories-definition-representable-fibred-category}.
Nous dirons parfois que $\mathcal{X}$ est
{\it représentable par un schéma} pour distinguer ce cas de celui où
$\mathcal{X}$ est représentable par un espace algébrique (voir
ci-dessous).

\medskip\noindent
Si $\mathcal{X}, \mathcal{Y}$ sont fibrées en groupoïdes et
représentables par $U, V$, alors on a
\begin{equation}
\label{algebraic-equation-morphisms-schemes}
\Mor_{\textit{Cat}/(\Sch/S)_{fppf}}(\mathcal{X}, \mathcal{Y})
\Big/
2\text{-isomorphisme}
=
\Mor_{\Sch/S}(U, V)
\end{equation}
voir
Catégories,
Lemme \ref{categories-lemma-morphisms-representable-fibred-categories}.
Plus précisément, tout $1$-morphisme $\mathcal{X} \to \mathcal{Y}$
donne lieu à un morphisme $U \to V$. Réciproquement, étant donné un morphisme
de schémas $U \to V$ au-dessus de $S$, il existe un $1$-morphisme
$\phi : \mathcal{X} \to \mathcal{Y}$ qui donne lieu à $U \to V$
et qui est unique à un unique $2$-isomorphisme près.






\section{Le lemme de 2-Yoneda}
\label{algebraic-section-2-yoneda}

\noindent
Soient $U \in \Ob((\Sch/S)_{fppf})$ et $\mathcal{X}$ une
catégorie fibrée en groupoïdes au-dessus de $(\Sch/S)_{fppf}$.
Nous utiliserons fréquemment le lemme de $2$-Yoneda, voir
Catégories, Lemme \ref{categories-lemma-yoneda-2category}.
Techniquement, il affirme qu'il existe une équivalence de catégories
$$
\Mor_{\textit{Cat}/(\Sch/S)_{fppf}}(
(\Sch/U)_{fppf}, \mathcal{X})
\longrightarrow
\mathcal{X}_U, \quad
f \longmapsto f(U/U).
$$
Il affirme que les $1$-morphismes $(\Sch/U)_{fppf} \to \mathcal{X}$
correspondent aux objets $x$ de la catégorie fibre $\mathcal{X}_U$.
En effet, un $1$-morphisme $f : (\Sch/U)_{fppf} \to \mathcal{X}$
donne l'objet $x = f(U/U) \in \Ob(\mathcal{X}_U)$.
Réciproquement, étant donnés un choix d'images inverses pour $\mathcal{X}$ comme dans
Catégories,
Définition \ref{categories-definition-pullback-functor-fibred-category},
et un objet $x$ de $\mathcal{X}_U$, on obtient un foncteur
$(\Sch/U)_{fppf} \to \mathcal{X}$ défini sur les objets par la règle
$$
(\varphi : V \to U) \longmapsto \varphi^*x
$$
Par abus de notation, nous employons
$x : (\Sch/U)_{fppf} \to \mathcal{X}$
pour désigner ce foncteur. Il vérifie bien $x(U/U) = x$
et, de plus, pour tout autre foncteur $f$ tel que $f(U/U) = x$, il existe
un unique $2$-isomorphisme $x \to f$. Autrement dit, le foncteur $x$
est déterminé par l'objet $x$ à un unique $2$-isomorphisme près.

\medskip\noindent
Nous l'utiliserons désormais sans autre mention.





\section{Morphismes représentables de catégories fibrées en groupoïdes}
\sectionmark{Morphismes représentables de catégories fibrées}
\label{algebraic-section-representable-morphism}


\noindent
Soient $\mathcal{X}$, $\mathcal{Y}$ des catégories fibrées en groupoïdes
au-dessus de $(\Sch/S)_{fppf}$. Soit $f : \mathcal{X} \to \mathcal{Y}$
un {\it $1$-morphisme représentable}, voir
Catégories, Définition
\ref{categories-definition-representable-map-categories-fibred-in-groupoids}.
Cela signifie que, pour tout $U \in \Ob((\Sch/S)_{fppf})$ et
tout $y \in \Ob(\mathcal{Y}_U)$, le $2$-produit fibré
$(\Sch/U)_{fppf} \times_{y, \mathcal{Y}} \mathcal{X}$
est représentable comme catégorie fibrée en groupoïdes au-dessus de
$(\Sch/U)_{fppf}$. Choisissons un objet représentant $f_y : V_y \to U$
et une équivalence
$$
(\Sch/V_y)_{fppf}
\longrightarrow
(\Sch/U)_{fppf} \times_{y, \mathcal{Y}} \mathcal{X}.
$$
Le morphisme $f_y$ correspond à la projection
$(\Sch/V_y)_{fppf} \to
(\Sch/U)_{fppf} \times_\mathcal{Y} \mathcal{Y}
\to (\Sch/U)_{fppf}$
Voir section \ref{algebraic-section-representable},
Équation (\ref{algebraic-equation-morphisms-schemes}). Nous représentons ceci par le diagramme
\begin{equation}
\label{algebraic-equation-representable}
\vcenter{
\xymatrix{
V_y \ar@{~>}[r] \ar[d]_{f_y} &
(\Sch/V_y)_{fppf} \ar[d] \ar[r] &
\mathcal{X} \ar[d]^f \\
U \ar@{~>}[r] &
(\Sch/U)_{fppf} \ar[r]^-y &
\mathcal{Y}
}
}
\end{equation}
où les flèches ondulées représentent le plongement de $2$-Yoneda.
Voici quelques lemmes sur cette notion, valables dans une grande généralité
(plus précisément, ils valent pour les catégories fibrées en groupoïdes au-dessus de toute
catégorie de base qui possède des produits fibrés).

\begin{lemma}
\label{algebraic-lemma-morphism-schemes-gives-representable-transformation}
Soit $f : X \to Y$ un morphisme de $(\Sch/S)_{fppf}$.
Alors le $1$-morphisme induit par $f$
$$
(\Sch/X)_{fppf} \longrightarrow (\Sch/Y)_{fppf}
$$
est un $1$-morphisme représentable.
\end{lemma}

\begin{proof}
C'est formel et ne repose que sur le fait que
la catégorie $(\Sch/S)_{fppf}$ possède des produits fibrés.
\end{proof}

\begin{lemma}
\label{algebraic-lemma-representable-morphism-equivalent}
Soit $S$ un objet de $\Sch_{fppf}$.
Considérons un diagramme $2$-commutatif
$$
\xymatrix{
\mathcal{X}' \ar[r] \ar[d]_{f'} & \mathcal{X} \ar[d]^f \\
\mathcal{Y}' \ar[r] & \mathcal{Y}
}
$$
de $1$-morphismes de catégories fibrées en groupoïdes au-dessus de
$(\Sch/S)_{fppf}$.
Supposons que les flèches horizontales soient des équivalences.
Alors $f$ est représentable si et seulement si $f'$ est représentable.
\end{lemma}

\begin{proof}
Omis.
\end{proof}

\begin{lemma}
\label{algebraic-lemma-composition-representable-transformations}
Soit $S$ un schéma contenu dans $\Sch_{fppf}$.
Soient $\mathcal{X}, \mathcal{Y}, \mathcal{Z}$
des catégories fibrées en groupoïdes au-dessus de $(\Sch/S)_{fppf}$
Soient $f : \mathcal{X} \to \mathcal{Y}$, $g : \mathcal{Y} \to \mathcal{Z}$
des $1$-morphismes représentables. Alors
$$
g \circ f : \mathcal{X} \longrightarrow \mathcal{Z}
$$
est un $1$-morphisme représentable.
\end{lemma}

\begin{proof}
C'est entièrement formel et vaut dans toute catégorie.
\end{proof}

\begin{lemma}
\label{algebraic-lemma-base-change-representable-transformations}
Soit $S$ un schéma contenu dans $\Sch_{fppf}$.
Soient $\mathcal{X}, \mathcal{Y}, \mathcal{Z}$
des catégories fibrées en groupoïdes au-dessus de $(\Sch/S)_{fppf}$
Soit $f : \mathcal{X} \to \mathcal{Y}$ un $1$-morphisme représentable.
Soit $g : \mathcal{Z} \to \mathcal{Y}$ un $1$-morphisme quelconque.
Considérons le diagramme de produit fibré
$$
\xymatrix{
\mathcal{Z} \times_{g, \mathcal{Y}, f} \mathcal{X} \ar[r]_-{g'} \ar[d]_{f'} &
\mathcal{X} \ar[d]^f \\
\mathcal{Z} \ar[r]^g & \mathcal{Y}
}
$$
Alors le changement de base $f'$ est un $1$-morphisme représentable.
\end{lemma}

\begin{proof}
C'est entièrement formel et vaut dans toute catégorie.
\end{proof}

\begin{lemma}
\label{algebraic-lemma-product-representable-transformations}
Soit $S$ un schéma contenu dans $\Sch_{fppf}$.
Soient $\mathcal{X}_i, \mathcal{Y}_i$ des catégories fibrées en groupoïdes au-dessus de
$(\Sch/S)_{fppf}$, $i = 1, 2$.
Soient $f_i : \mathcal{X}_i \to \mathcal{Y}_i$, $i = 1, 2$
des $1$-morphismes représentables.
Alors
$$
f_1 \times f_2 :
\mathcal{X}_1 \times \mathcal{X}_2
\longrightarrow
\mathcal{Y}_1 \times \mathcal{Y}_2
$$
est un $1$-morphisme représentable.
\end{lemma}

\begin{proof}
Écrivons $f_1 \times f_2$ comme le composé
$\mathcal{X}_1 \times \mathcal{X}_2 \to
\mathcal{Y}_1 \times \mathcal{X}_2 \to
\mathcal{Y}_1 \times \mathcal{Y}_2$.
La première flèche est le changement de base de $f_1$ par l'application
$\mathcal{Y}_1 \times \mathcal{X}_2 \to \mathcal{Y}_1$, et la seconde flèche
est le changement de base de $f_2$ par l'application
$\mathcal{Y}_1 \times \mathcal{Y}_2 \to \mathcal{Y}_2$.
Ce lemme est donc une conséquence formelle des Lemmes
\ref{algebraic-lemma-composition-representable-transformations}
et \ref{algebraic-lemma-base-change-representable-transformations}.
\end{proof}



\section{Catégories fibrées en groupoïdes scindées}
\label{algebraic-section-split}

\noindent
Soit $S$ un schéma contenu dans $\Sch_{fppf}$.
Rappelons qu'étant donné un « préfaisceau de groupoïdes »
$$
F : (\Sch/S)_{fppf}^{opp} \longrightarrow \textit{Groupoïdes}
$$
on obtient une catégorie fibrée en groupoïdes $\mathcal{S}_F$ au-dessus de
$(\Sch/S)_{fppf}$, voir
Catégories, Exemple \ref{categories-example-functor-groupoids}.
Toute catégorie fibrée en groupoïdes isomorphe (!) à l'une d'elles
est appelée une {\it catégorie fibrée en groupoïdes scindée}.
Toute catégorie fibrée en groupoïdes est équivalente à une catégorie scindée.

\medskip\noindent
Si $F$ est un préfaisceau d'ensembles, alors $\mathcal{S}_F$ est
fibrée en ensembles, voir
Catégories,
Définition \ref{categories-definition-category-fibred-sets},
et
Catégories, Exemple \ref{categories-example-presheaf}.
La règle $F \mapsto \mathcal{S}_F$ est en un certain sens pleinement fidèle
sur les préfaisceaux, voir
Catégories, Lemme \ref{categories-lemma-2-category-fibred-sets}.
Si $F, G$ sont des préfaisceaux, alors
$$
\mathcal{S}_{F \times G}
=
\mathcal{S}_F \times_{(\Sch/S)_{fppf}} \mathcal{S}_G
$$
et si $F \to H$ et $G \to H$ sont des applications de préfaisceaux d'ensembles, alors
$$
\mathcal{S}_{F \times_H G} =
\mathcal{S}_F \times_{\mathcal{S}_H} \mathcal{S}_G
$$
où les membres de droite sont des $2$-produits fibrés. Cela résulte immédiatement
des définitions, puisque les catégories fibres de
$\mathcal{S}_F, \mathcal{S}_G, \mathcal{S}_H$ ne possèdent que des morphismes identité.

\medskip\noindent
Un cas encore plus particulier est celui où $F = h_X$ est un préfaisceau
représentable. Dans ce cas, on a
$\mathcal{S}_{h_X} = (\Sch/X)_{fppf}$, voir
Catégories,
Exemple \ref{categories-example-fibred-category-from-functor-of-points}.

\medskip\noindent
Nous emploierons désormais la notation $\mathcal{S}_F$ sans autre mention.





\section{Catégories fibrées en groupoïdes représentables par des espaces algébriques}
\label{algebraic-section-representable-by-algebraic-spaces}

\noindent
La représentabilité par des espaces algébriques, que nous étudions dans cette section,
est une notion légèrement plus faible que la représentabilité.
On aurait pu éviter cette discussion si nous avions travaillé avec une catégorie
$\textit{Espaces}_{fppf}$ d'espaces algébriques plutôt qu'avec la catégorie
$\Sch_{fppf}$. Il nous semble toutefois naturel de considérer la
catégorie des schémas comme la collection naturelle d'« objets tests » au-dessus
desquels sont définies les catégories fibres d'un champ algébrique.

\medskip\noindent
Par analogie avec Catégories, Définitions
\ref{categories-definition-representable-fibred-category},
nous posons la définition suivante.

\begin{definition}
\label{algebraic-definition-representable-by-algebraic-space}
Soit $S$ un schéma contenu dans $\Sch_{fppf}$.
Une catégorie fibrée en groupoïdes $p : \mathcal{X} \to (\Sch/S)_{fppf}$
est dite {\it représentable par un espace algébrique sur $S$}
s'il existe un espace algébrique $F$ sur $S$ et une équivalence
$j : \mathcal{X} \to \mathcal{S}_F$
de catégories au-dessus de $(\Sch/S)_{fppf}$.
\end{definition}

\noindent
Nous poursuivons notre abus de notation en omettant l'équivalence $j$
chaque fois que nous rencontrons une telle situation.
Il résulte formellement de ce qui précède que, si $\mathcal{X}$ est
représentable (par un schéma), elle est représentable par un
espace algébrique. Voici l'analogue de
Catégories,
Lemme \ref{categories-lemma-characterize-representable-fibred-category}.

\begin{lemma}
\label{algebraic-lemma-characterize-representable-by-space}
Soit $S$ un schéma contenu dans $\Sch_{fppf}$.
Soit $p : \mathcal{X} \to (\Sch/S)_{fppf}$
une catégorie fibrée en groupoïdes.
Alors $\mathcal{X}$ est représentable par un espace algébrique sur $S$
si et seulement si les conditions suivantes sont satisfaites :
\begin{enumerate}
\item $\mathcal{X}$ est fibrée en setoïdes\footnote{Cela signifie qu'elle
est fibrée en groupoïdes et que les objets des catégories fibres
n'ont pas d'automorphismes non triviaux, voir Catégories,
Définition \ref{categories-definition-category-fibred-setoids}.}, et
\item le préfaisceau $U \mapsto \Ob(\mathcal{X}_U)/\!\!\cong$ est
un espace algébrique.
\end{enumerate}
\end{lemma}

\begin{proof}
Omis, mais voir Catégories,
Lemme \ref{categories-lemma-characterize-representable-fibred-category}.
\end{proof}

\noindent
Si $\mathcal{X}, \mathcal{Y}$ sont fibrées en groupoïdes et
représentables par des espaces algébriques $F, G$ sur $S$, alors on a
\begin{equation}
\label{algebraic-equation-morphisms-spaces}
\Mor_{\textit{Cat}/(\Sch/S)_{fppf}}(\mathcal{X}, \mathcal{Y})
\Big/
2\text{-isomorphisme}
=
\Mor_{\Sch/S}(F, G)
\end{equation}
voir
Catégories, Lemme \ref{categories-lemma-2-category-fibred-setoids}.
Plus précisément, tout $1$-morphisme $\mathcal{X} \to \mathcal{Y}$
donne lieu à un morphisme $F \to G$. Réciproquement, étant donné un morphisme
de faisceaux $F \to G$ au-dessus de $S$, il existe un $1$-morphisme
$\phi : \mathcal{X} \to \mathcal{Y}$ qui donne lieu à $F \to G$
et qui est unique à un unique $2$-isomorphisme près.



\section{Morphismes représentables par des espaces algébriques}
\label{algebraic-section-morphisms-representable-by-algebraic-spaces}

\noindent
Par analogie avec Catégories, Définition
\ref{categories-definition-representable-map-categories-fibred-in-groupoids},
nous posons la définition suivante.

\begin{definition}
\label{algebraic-definition-representable-by-algebraic-spaces}
Soit $S$ un schéma contenu dans $\Sch_{fppf}$.
Un $1$-morphisme $f : \mathcal{X} \to \mathcal{Y}$ de
catégories fibrées en groupoïdes au-dessus de $(\Sch/S)_{fppf}$
est dit {\it représentable par des espaces algébriques} si,
pour tout $U \in \Ob((\Sch/S)_{fppf})$
et tout $y : (\Sch/U)_{fppf} \to \mathcal{Y}$,
la catégorie fibrée en groupoïdes
$$
(\Sch/U)_{fppf} \times_{y, \mathcal{Y}} \mathcal{X}
$$
au-dessus de $(\Sch/U)_{fppf}$
est représentable par un espace algébrique sur $U$.
\end{definition}

\noindent
Choisissons un espace algébrique $F_y$ sur $U$ qui représente
$(\Sch/U)_{fppf} \times_{y, \mathcal{Y}} \mathcal{X}$.
On peut considérer $F_y$ comme un espace algébrique sur $S$
muni d'un morphisme canonique $f_y : F_y \to U$
au-dessus de $S$, voir
Espaces, section \ref{spaces-section-change-base-scheme}.
Voici le diagramme
\begin{equation}
\label{algebraic-equation-representable-by-algebraic-spaces}
\vcenter{
\xymatrix{
F_y \ar[d]_{f_y} &
(\Sch/U)_{fppf} \times_{y, \mathcal{Y}} \mathcal{X}
\ar@{~>}[l] \ar[d]_{\text{pr}_0} \ar[r]_-{\text{pr}_1} &
\mathcal{X} \ar[d]^f \\
U &
(\Sch/U)_{fppf} \ar@{~>}[l] \ar[r]^-y &
\mathcal{Y}
}
}
\end{equation}
où les flèches ondulées représentent la construction qui associe
à un champ fibré en setoïdes le faisceau associé des classes d'isomorphie
de ses objets. Le carré de droite est
$2$-commutatif et est un carré $2$-cartésien.

\medskip\noindent
Voici l'analogue de Catégories,
Lemme \ref{categories-lemma-criterion-representable-map-stack-in-groupoids}.

\begin{lemma}
\label{algebraic-lemma-criterion-map-representable-spaces-fibred-in-groupoids}
Soit $S$ un schéma contenu dans $\Sch_{fppf}$.
Soit $f : \mathcal{X} \to \mathcal{Y}$ un $1$-morphisme
de catégories fibrées en groupoïdes au-dessus de $(\Sch/S)_{fppf}$.
Pour que $f$ soit représentable par des espaces algébriques, il faut et il
suffit que les conditions suivantes soient satisfaites :
\begin{enumerate}
\item pour tout schéma $U/S$, le
foncteur $f_U : \mathcal{X}_U \longrightarrow \mathcal{Y}_U$
entre les catégories fibres est fidèle, et
\item pour tout $U$ et tout $y \in \Ob(\mathcal{Y}_U)$, le préfaisceau
$$
(h : V \to U)
\longmapsto
\{(x, \phi) \mid x \in \Ob(\mathcal{X}_V), \phi : h^*y \to f(x)\}/\cong
$$
est un espace algébrique sur $U$.
\end{enumerate}
Nous avons choisi ici des images inverses pour $\mathcal{Y}$.
\end{lemma}

\begin{proof}
Cela résulte de la description des catégories fibres des $2$-produits fibrés
$(\Sch/U)_{fppf} \times_{y, \mathcal{Y}} \mathcal{X}$ donnée dans
Catégories, Lemme \ref{categories-lemma-identify-fibre-product},
combinée au
Lemme \ref{algebraic-lemma-characterize-representable-by-space}.
\end{proof}

\noindent
Voici quelques lemmes sur cette notion, valables dans une grande généralité.

\begin{lemma}
\label{algebraic-lemma-representable-by-spaces-morphism-equivalent}
Soit $S$ un objet de $\Sch_{fppf}$.
Considérons un diagramme $2$-commutatif
$$
\xymatrix{
\mathcal{X}' \ar[r] \ar[d]_{f'} & \mathcal{X} \ar[d]^f \\
\mathcal{Y}' \ar[r] & \mathcal{Y}
}
$$
de $1$-morphismes de catégories fibrées en groupoïdes au-dessus de
$(\Sch/S)_{fppf}$.
Supposons que les flèches horizontales soient des équivalences.
Alors $f$ est représentable par des espaces algébriques
si et seulement si $f'$ est représentable par des espaces algébriques.
\end{lemma}

\begin{proof}
Omis.
\end{proof}

\begin{lemma}
\label{algebraic-lemma-morphism-spaces-gives-representable-by-spaces}
Soit $S$ un objet de $\Sch_{fppf}$.
Soit $f : \mathcal{X} \to \mathcal{Y}$
un $1$-morphisme de catégories fibrées en groupoïdes au-dessus de $S$.
Si $\mathcal{X}$ et $\mathcal{Y}$ sont représentables par des
espaces algébriques sur $S$, alors le $1$-morphisme $f$
est représentable par des espaces algébriques.
\end{lemma}

\begin{proof}
Omis. Cela ne repose que sur le fait que
la catégorie des espaces algébriques sur $S$ possède des produits fibrés,
voir Espaces, Lemme \ref{spaces-lemma-fibre-product-spaces}.
\end{proof}

\begin{lemma}
\label{algebraic-lemma-map-presheaves-representable-by-algebraic-spaces}
Soit $S$ un objet de $\Sch_{fppf}$.
Soit $a : F \to G$ un morphisme de préfaisceaux d'ensembles sur $(\Sch/S)_{fppf}$.
Notons $a' : \mathcal{S}_F  \to \mathcal{S}_G$ le morphisme associé
de catégories fibrées en ensembles.
Alors $a$ est représentable par des espaces algébriques (voir
Amorçage,
Définition \ref{bootstrap-definition-morphism-representable-by-spaces})
si et seulement si $a'$ est représentable par des espaces algébriques.
\end{lemma}

\begin{proof}
Omis.
\end{proof}

\begin{lemma}
\label{algebraic-lemma-map-fibred-setoids-representable-algebraic-spaces}
Soit $S$ un objet de $\Sch_{fppf}$.
Soit $f : \mathcal{X} \to \mathcal{Y}$ un $1$-morphisme de
catégories fibrées en setoïdes au-dessus de $(\Sch/S)_{fppf}$.
Soient $F$ et $G$ les préfaisceaux qui associent respectivement à $T$
l'ensemble des classes d'isomorphie des objets de
$\mathcal{X}_T$ et celui des objets de $\mathcal{Y}_T$.
Soit $a : F \to G$ le morphisme de préfaisceaux correspondant à $f$.
Alors $a$ est représentable par des espaces algébriques (voir
Amorçage,
Définition \ref{bootstrap-definition-morphism-representable-by-spaces})
si et seulement si $f$ est représentable par des espaces algébriques.
\end{lemma}

\begin{proof}
Omis. Indication : combiner les
Lemmes \ref{algebraic-lemma-representable-by-spaces-morphism-equivalent}
et \ref{algebraic-lemma-map-presheaves-representable-by-algebraic-spaces}.
\end{proof}

\begin{lemma}
\label{algebraic-lemma-base-change-representable-by-spaces}
Soit $S$ un schéma contenu dans $\Sch_{fppf}$.
Soient $\mathcal{X}, \mathcal{Y}, \mathcal{Z}$
des catégories fibrées en groupoïdes au-dessus de $(\Sch/S)_{fppf}$.
Soit $f : \mathcal{X} \to \mathcal{Y}$ un $1$-morphisme
représentable par des espaces algébriques.
Soit $g : \mathcal{Z} \to \mathcal{Y}$ un $1$-morphisme quelconque.
Considérons le diagramme de produit fibré
$$
\xymatrix{
\mathcal{Z} \times_{g, \mathcal{Y}, f} \mathcal{X} \ar[r]_-{g'} \ar[d]_{f'} &
\mathcal{X} \ar[d]^f \\
\mathcal{Z} \ar[r]^g & \mathcal{Y}
}
$$
Alors le changement de base $f'$ est un $1$-morphisme représentable par des
espaces algébriques.
\end{lemma}

\begin{proof}
C'est formel.
\end{proof}

\begin{lemma}
\label{algebraic-lemma-base-change-by-space-representable-by-space}
Soit $S$ un schéma contenu dans $\Sch_{fppf}$.
Soient $\mathcal{X}, \mathcal{Y}, \mathcal{Z}$
des catégories fibrées en groupoïdes au-dessus de $(\Sch/S)_{fppf}$
Soient $f : \mathcal{X} \to \mathcal{Y}$,
$g : \mathcal{Z} \to \mathcal{Y}$ des $1$-morphismes.
Supposons que
\begin{enumerate}
\item $f$ soit représentable par des espaces algébriques, et
\item $\mathcal{Z}$ soit représentable par un espace algébrique sur $S$.
\end{enumerate}
Alors le $2$-produit fibré
$\mathcal{Z} \times_{g, \mathcal{Y}, f} \mathcal{X}$
est représentable par un espace algébrique.
\end{lemma}

\begin{proof}
C'est une reformulation de
Amorçage, Lemme \ref{bootstrap-lemma-representable-by-spaces-over-space}.
Notons d'abord que
$\mathcal{Z} \times_{g, \mathcal{Y}, f} \mathcal{X}$
est fibrée en setoïdes au-dessus de $(\Sch/S)_{fppf}$.
Elle est donc équivalente à $\mathcal{S}_F$ pour un certain préfaisceau
$F$ sur $(\Sch/S)_{fppf}$, voir
Catégories, Lemme \ref{categories-lemma-setoid-fibres}.
De plus, soit $G$ un espace algébrique qui représente
$\mathcal{Z}$. Le $1$-morphisme
$\mathcal{Z} \times_{g, \mathcal{Y}, f} \mathcal{X} \to \mathcal{Z}$
est représentable par des espaces algébriques d'après le
Lemme \ref{algebraic-lemma-base-change-representable-by-spaces}.
Et $\mathcal{Z} \times_{g, \mathcal{Y}, f} \mathcal{X} \to \mathcal{Z}$
correspond à un morphisme $F \to G$ d'après
Catégories, Lemme \ref{categories-lemma-2-category-fibred-setoids}.
Alors $F \to G$ est représentable par des espaces algébriques d'après le
Lemme \ref{algebraic-lemma-map-fibred-setoids-representable-algebraic-spaces}.
Par conséquent,
Amorçage, Lemme \ref{bootstrap-lemma-representable-by-spaces-over-space}
implique que $F$ est un espace algébrique, comme souhaité.
\end{proof}

\noindent
Soient $S$, $\mathcal{X}$, $\mathcal{Y}$, $\mathcal{Z}$, $f$, $g$ comme dans le
Lemme \ref{algebraic-lemma-base-change-by-space-representable-by-space}.
Soient $F$ et $G$ des espaces algébriques sur $S$ tels que
$F$ représente $\mathcal{Z} \times_{g, \mathcal{Y}, f} \mathcal{X}$
et $G$ représente $\mathcal{Z}$. Le $1$-morphisme
$f' : \mathcal{Z} \times_{g, \mathcal{Y}, f} \mathcal{X} \to \mathcal{Z}$
correspond à un morphisme $f' : F \to G$ d'espaces algébriques
d'après (\ref{algebraic-equation-morphisms-spaces}).
On a donc le diagramme suivant
\begin{equation}
\label{algebraic-equation-representable-by-algebraic-spaces-on-space}
\vcenter{
\xymatrix{
F \ar[d]_{f'} &
\mathcal{Z} \times_{g, \mathcal{Y}, f} \mathcal{X}
\ar@{~>}[l] \ar[d] \ar[r] &
\mathcal{X} \ar[d]^f \\
G &
\mathcal{Z} \ar@{~>}[l] \ar[r]^-g &
\mathcal{Y}
}
}
\end{equation}
où les flèches ondulées représentent la construction qui associe
à un champ fibré en setoïdes le faisceau associé des classes d'isomorphie
de ses objets.

\begin{lemma}
\label{algebraic-lemma-composition-representable-by-spaces}
Soit $S$ un schéma contenu dans $\Sch_{fppf}$.
Soient $\mathcal{X}, \mathcal{Y}, \mathcal{Z}$
des catégories fibrées en groupoïdes au-dessus de $(\Sch/S)_{fppf}$.
Si $f : \mathcal{X} \to \mathcal{Y}$, $g : \mathcal{Y} \to \mathcal{Z}$
sont des $1$-morphismes représentables par des espaces algébriques, alors
$$
g \circ f : \mathcal{X} \longrightarrow \mathcal{Z}
$$
est un $1$-morphisme représentable par des espaces algébriques.
\end{lemma}

\begin{proof}
Cela résulte du
Lemme \ref{algebraic-lemma-base-change-by-space-representable-by-space}.
Les détails sont omis.
\end{proof}

\begin{lemma}
\label{algebraic-lemma-product-representable-by-spaces}
Soit $S$ un schéma contenu dans $\Sch_{fppf}$.
Soient $\mathcal{X}_i, \mathcal{Y}_i$ des catégories fibrées en groupoïdes au-dessus de
$(\Sch/S)_{fppf}$, $i = 1, 2$.
Soient $f_i : \mathcal{X}_i \to \mathcal{Y}_i$, $i = 1, 2$
des $1$-morphismes représentables par des espaces algébriques.
Alors
$$
f_1 \times f_2 :
\mathcal{X}_1 \times \mathcal{X}_2
\longrightarrow
\mathcal{Y}_1 \times \mathcal{Y}_2
$$
est un $1$-morphisme représentable par des espaces algébriques.
\end{lemma}

\begin{proof}
Écrivons $f_1 \times f_2$ comme le composé
$\mathcal{X}_1 \times \mathcal{X}_2 \to
\mathcal{Y}_1 \times \mathcal{X}_2 \to
\mathcal{Y}_1 \times \mathcal{Y}_2$.
La première flèche est le changement de base de $f_1$ par l'application
$\mathcal{Y}_1 \times \mathcal{X}_2 \to \mathcal{Y}_1$, et la seconde flèche
est le changement de base de $f_2$ par l'application
$\mathcal{Y}_1 \times \mathcal{Y}_2 \to \mathcal{Y}_2$.
Ce lemme est donc une conséquence formelle des Lemmes
\ref{algebraic-lemma-composition-representable-by-spaces}
et \ref{algebraic-lemma-base-change-representable-by-spaces}.
\end{proof}

\begin{lemma}
\label{algebraic-lemma-get-a-stack}
\begin{reference}
Lemme figurant dans un courriel de Matthew Emerton daté du 15 juin 2016
\end{reference}
Soit $S$ un schéma contenu dans $\Sch_{fppf}$.
Soient $\mathcal{X} \to \mathcal{Z}$ et $\mathcal{Y} \to \mathcal{Z}$
des $1$-morphismes de catégories fibrées en groupoïdes au-dessus de $(\Sch/S)_{fppf}$.
Si $\mathcal{X} \to \mathcal{Z}$ est représentable par des espaces algébriques
et si $\mathcal{Y}$ est un champ en groupoïdes, alors
$\mathcal{X} \times_\mathcal{Z} \mathcal{Y}$ est un champ en groupoïdes.
\end{lemma}

\begin{proof}
La propriété, pour un morphisme, d'être représentable par des espaces algébriques
est préservée par changement de base
(Lemme \ref{algebraic-lemma-base-change-by-space-representable-by-space}) ;
ainsi, en passant au changement de base
$\mathcal{X} \times_\mathcal{Z} \mathcal{Y}$ au-dessus de $\mathcal{Y}$,
on peut se ramener au cas d'un morphisme de catégories
fibrées en groupoïdes $\mathcal{X} \to \mathcal{Y}$
qui est représentable par des espaces algébriques et
dont le but est un champ en groupoïdes ; il s'agit alors de démontrer
que $\mathcal{X}$ est également un champ en groupoïdes.
Cela résulte de Champs, Lemme
\ref{stacks-lemma-relative-sheaf-over-stack-is-stack},
dont les hypothèses sont satisfaites en vertu du
Lemme \ref{algebraic-lemma-criterion-map-representable-spaces-fibred-in-groupoids}.
\end{proof}







\section{Propriétés des morphismes représentables par des espaces algébriques}
\label{algebraic-section-representable-properties}

\noindent
Voici la définition qui permet de procéder.

\begin{definition}
\label{algebraic-definition-relative-representable-property}
Soit $S$ un schéma appartenant à $\Sch_{fppf}$.
Soit $f : \mathcal{X} \to \mathcal{Y}$ un $1$-morphisme
de catégories fibrées en groupoïdes au-dessus de $(\Sch/S)_{fppf}$.
Supposons que $f$ soit représentable par des espaces algébriques.
Soit $\mathcal{P}$ une propriété des morphismes d'espaces algébriques qui
\begin{enumerate}
\item est préservée par tout changement de base, et
\item est locale pour la topologie fppf sur la base, voir
Descente et espaces algébriques,
Définition \ref{spaces-descent-definition-property-morphisms-local}.
\end{enumerate}
Dans ce cas, on dit que $f$ a la {\it propriété $\mathcal{P}$} si, pour tout
$U \in \Ob((\Sch/S)_{fppf})$ et
tout $y \in \mathcal{Y}_U$, le morphisme d'espaces algébriques qui en résulte
$f_y : F_y \to U$, voir le
diagramme (\ref{algebraic-equation-representable-by-algebraic-spaces}),
a la propriété $\mathcal{P}$.
\end{definition}

\noindent
Il importe de noter que nous n'utiliserons cette définition que pour des
propriétés de morphismes stables par changement de base et
locales pour la topologie fppf sur le but. Ce n'est
pas que la définition n'ait autrement aucun sens ; c'est plutôt
que nous pourrions vouloir donner une autre définition qui soit
mieux adaptée à la propriété envisagée.

\begin{lemma}
\label{algebraic-lemma-property-morphism-equivalent}
Soit $S$ un objet de $\Sch_{fppf}$.
Soit $\mathcal{P}$ comme dans la
Définition \ref{algebraic-definition-relative-representable-property}.
Considérons un diagramme $2$-commutatif
$$
\xymatrix{
\mathcal{X}' \ar[r] \ar[d]_{f'} & \mathcal{X} \ar[d]^f \\
\mathcal{Y}' \ar[r] & \mathcal{Y}
}
$$
de $1$-morphismes de catégories fibrées en groupoïdes au-dessus de
$(\Sch/S)_{fppf}$.
Supposons que les flèches horizontales soient des équivalences et que $f$
(ou, de manière équivalente, $f'$) soit représentable par des espaces algébriques.
Alors $f$ a $\mathcal{P}$ si et seulement si $f'$ a $\mathcal{P}$.
\end{lemma}

\begin{proof}
Cet énoncé a bien un sens en vertu du
Lemme \ref{algebraic-lemma-representable-by-spaces-morphism-equivalent}.
Démonstration omise.
\end{proof}

\noindent
Voici un contrôle de cohérence.

\begin{lemma}
\label{algebraic-lemma-map-presheaves-representable-by-spaces-transformation-property}
Soit $S$ un schéma appartenant à $\Sch_{fppf}$.
Soit $a : F \to G$ une application de préfaisceaux sur $(\Sch/S)_{fppf}$.
Soit $\mathcal{P}$ comme dans la
Définition \ref{algebraic-definition-relative-representable-property}.
Supposons que $a$ soit représentable par des espaces algébriques.
Alors $a : F \to G$ a la propriété $\mathcal{P}$ (voir
Amorçage, Définition \ref{bootstrap-definition-property-transformation})
si et seulement si le morphisme correspondant
$\mathcal{S}_F \to \mathcal{S}_G$ de catégories fibrées en groupoïdes
a la propriété $\mathcal{P}$.
\end{lemma}

\begin{proof}
L'énoncé du lemme a bien un sens en vertu du
Lemme \ref{algebraic-lemma-map-presheaves-representable-by-algebraic-spaces}.
Démonstration omise.
\end{proof}

\begin{lemma}
\label{algebraic-lemma-map-fibred-setoids-property}
Soit $S$ un objet de $\Sch_{fppf}$. Soit $\mathcal{P}$ comme dans la
Définition \ref{algebraic-definition-relative-representable-property}.
Soit $f : \mathcal{X} \to \mathcal{Y}$ un $1$-morphisme de
catégories fibrées en setoïdes au-dessus de $(\Sch/S)_{fppf}$.
Soient $F$ et $G$ les préfaisceaux qui associent respectivement à $T$
l'ensemble des classes d'isomorphie des objets de
$\mathcal{X}_T$ et celui des objets de $\mathcal{Y}_T$.
Soit $a : F \to G$ l'application de préfaisceaux correspondant à $f$.
Alors $a$ a $\mathcal{P}$ si et seulement si $f$ a $\mathcal{P}$.
\end{lemma}

\begin{proof}
L'énoncé du lemme a bien un sens en vertu du
Lemme \ref{algebraic-lemma-map-fibred-setoids-representable-algebraic-spaces}.
Le lemme résulte de la combinaison des
Lemmes \ref{algebraic-lemma-property-morphism-equivalent}
et \ref{algebraic-lemma-map-presheaves-representable-by-spaces-transformation-property}.
\end{proof}

\begin{lemma}
\label{algebraic-lemma-composition-representable-transformations-property}
Soit $S$ un schéma appartenant à $\Sch_{fppf}$.
Soient $\mathcal{X}$, $\mathcal{Y}$, $\mathcal{Z}$ des catégories fibrées
en groupoïdes au-dessus de $(\Sch/S)_{fppf}$.
Soit $\mathcal{P}$ une propriété comme dans la
Définition \ref{algebraic-definition-relative-representable-property}
qui est stable par composition.
Soient $f : \mathcal{X} \to \mathcal{Y}$ et
$g : \mathcal{Y} \to \mathcal{Z}$ des $1$-morphismes
représentables par des espaces algébriques.
Si $f$ et $g$ ont la propriété $\mathcal{P}$, il en est de même de
$g \circ f : \mathcal{X} \to \mathcal{Z}$.
\end{lemma}

\begin{proof}
L'énoncé du lemme a bien un sens en vertu du
Lemme \ref{algebraic-lemma-composition-representable-by-spaces}.
Démonstration omise.
\end{proof}

\begin{lemma}
\label{algebraic-lemma-base-change-representable-transformations-property}
Soit $S$ un schéma appartenant à $\Sch_{fppf}$.
Soient $\mathcal{X}, \mathcal{Y}, \mathcal{Z}$
des catégories fibrées en groupoïdes au-dessus de $(\Sch/S)_{fppf}$.
Soit $\mathcal{P}$ une propriété comme dans la
Définition \ref{algebraic-definition-relative-representable-property}.
Soit $f : \mathcal{X} \to \mathcal{Y}$ un $1$-morphisme
représentable par des espaces algébriques.
Soit $g : \mathcal{Z} \to \mathcal{Y}$ un $1$-morphisme quelconque.
Considérons le diagramme de $2$-produit fibré
$$
\xymatrix{
\mathcal{Z} \times_{g, \mathcal{Y}, f} \mathcal{X} \ar[r]_-{g'} \ar[d]_{f'} &
\mathcal{X} \ar[d]^f \\
\mathcal{Z} \ar[r]^g & \mathcal{Y}
}
$$
Si $f$ a $\mathcal{P}$, alors le changement de base $f'$
a $\mathcal{P}$.
\end{lemma}

\begin{proof}
L'énoncé du lemme a bien un sens en vertu du
Lemme \ref{algebraic-lemma-base-change-representable-by-spaces}.
Démonstration omise.
\end{proof}

\begin{lemma}
\label{algebraic-lemma-descent-representable-transformations-property}
Soit $S$ un schéma appartenant à $\Sch_{fppf}$.
Soient $\mathcal{X}, \mathcal{Y}, \mathcal{Z}$
des catégories fibrées en groupoïdes au-dessus de $(\Sch/S)_{fppf}$.
Soit $\mathcal{P}$ une propriété comme dans la
Définition \ref{algebraic-definition-relative-representable-property}.
Soit $f : \mathcal{X} \to \mathcal{Y}$ un $1$-morphisme
représentable par des espaces algébriques.
Soit $g : \mathcal{Z} \to \mathcal{Y}$ un $1$-morphisme quelconque.
Considérons le diagramme de produit fibré
$$
\xymatrix{
\mathcal{Z} \times_{g, \mathcal{Y}, f} \mathcal{X} \ar[r]_-{g'} \ar[d]_{f'} &
\mathcal{X} \ar[d]^f \\
\mathcal{Z} \ar[r]^g & \mathcal{Y}
}
$$
Supposons que, pour tout schéma $U$ et tout objet $x$ de $\mathcal{Y}_U$,
il existe un recouvrement fppf $\{U_i \to U\}$ tel que $x|_{U_i}$
appartienne à l'image essentielle du foncteur
$g : \mathcal{Z}_{U_i} \to \mathcal{Y}_{U_i}$.
Dans ce cas, si $f'$ a $\mathcal{P}$, alors $f$ a $\mathcal{P}$.
\end{lemma}

\begin{proof}
Démonstration omise. Indication : comparer avec la démonstration de
Espaces,
Lemme \ref{spaces-lemma-descent-representable-transformations-property}.
\end{proof}

\begin{lemma}
\label{algebraic-lemma-product-representable-transformations-property}
Soit $S$ un schéma appartenant à $\Sch_{fppf}$.
Soit $\mathcal{P}$ une propriété comme dans la
Définition \ref{algebraic-definition-relative-representable-property}
qui est stable par composition.
Soient $\mathcal{X}_i, \mathcal{Y}_i$ des catégories fibrées en groupoïdes au-dessus de
$(\Sch/S)_{fppf}$, $i = 1, 2$.
Soient $f_i : \mathcal{X}_i \to \mathcal{Y}_i$, $i = 1, 2$,
des $1$-morphismes représentables par des espaces algébriques.
Si $f_1$ et $f_2$ ont la propriété $\mathcal{P}$, il en est de même de
$
f_1 \times f_2 :
\mathcal{X}_1 \times \mathcal{X}_2
\to
\mathcal{Y}_1 \times \mathcal{Y}_2
$.
\end{lemma}

\begin{proof}
L'énoncé du lemme a bien un sens en vertu du
Lemme \ref{algebraic-lemma-product-representable-by-spaces}.
Démonstration omise.
\end{proof}

\begin{lemma}
\label{algebraic-lemma-representable-transformations-property-implication}
Soit $S$ un schéma appartenant à $\Sch_{fppf}$.
Soient $\mathcal{X}$, $\mathcal{Y}$ des catégories fibrées en groupoïdes
au-dessus de $(\Sch/S)_{fppf}$.
Soit $f : \mathcal{X} \to \mathcal{Y}$ un $1$-morphisme représentable
par des espaces algébriques.
Soient $\mathcal{P}$, $\mathcal{P}'$ des propriétés comme dans la
Définition \ref{algebraic-definition-relative-representable-property}.
Supposons que, pour tout morphisme d'espaces algébriques $a : F \to G$,
on ait $\mathcal{P}(a) \Rightarrow \mathcal{P}'(a)$.
Si $f$ a la propriété $\mathcal{P}$, alors
$f$ a la propriété $\mathcal{P}'$.
\end{lemma}

\begin{proof}
C'est formel.
\end{proof}

\begin{lemma}
\label{algebraic-lemma-open-fibred-category-is-full}
Soit $S$ un schéma appartenant à $\Sch_{fppf}$.
Soit $j : \mathcal X \to \mathcal Y$ un $1$-morphisme de
catégories fibrées en groupoïdes au-dessus de $(\Sch/S)_{fppf}$.
Supposons que $j$ soit représentable par des espaces algébriques et soit un monomorphisme
(voir la
Définition \ref{algebraic-definition-relative-representable-property}
et
Descente et espaces algébriques, Lemme
\ref{spaces-descent-lemma-descending-property-monomorphism}).
Alors $j$ est pleinement fidèle sur les catégories fibres.
\end{lemma}

\begin{proof}
Nous avons vu dans le
Lemme \ref{algebraic-lemma-criterion-map-representable-spaces-fibred-in-groupoids}
que $j$ est fidèle sur les catégories fibres. Considérons un schéma $U$,
deux objets $u, v$ de $\mathcal{X}_U$ et un isomorphisme
$t : j(u) \to j(v)$ dans $\mathcal{Y}_U$. Il faut construire un
isomorphisme dans $\mathcal{X}_U$ entre $u$ et $v$.
Par le lemme de $2$-Yoneda (voir section \ref{algebraic-section-2-yoneda}),
nous considérons $u$, $v$ comme des $1$-morphismes
$u, v : (\Sch/U)_{fppf} \to \mathcal{X}$
et considérons le $2$-produit fibré
$$
(\Sch/U)_{fppf} \times_{j \circ v, \mathcal{Y}} \mathcal{X}.
$$
Par hypothèse, celui-ci est représentable par un espace algébrique
$F_{j \circ v}$ au-dessus de $U$, et le morphisme
$F_{j \circ v} \to U$ est un monomorphisme.
Mais puisque $(1_U, v, 1_{j(v)})$ donne un $1$-morphisme de
$(\Sch/U)_{fppf}$ dans le $2$-produit fibré affiché,
on voit que $F_{j \circ v} = U$ (on utilise ici
que si $V \to U$ est un monomorphisme d'espaces algébriques qui possède une
section, alors $V = U$). Par conséquent, le $1$-morphisme de projection sur
la première coordonnée
$$
(\Sch/U)_{fppf} \times_{j \circ v, \mathcal{Y}} \mathcal{X}
\to (\Sch/U)_{fppf}
$$
est une équivalence de catégories fibres.
Puisque $(1_U, u, t)$ et $(1_U, v, 1_{j(v)})$ donnent deux
objets de $((\Sch/U)_{fppf} \times_{j \circ v, \mathcal{Y}}
\mathcal{X})_U$ qui ont la même première coordonnée, il doit
exister entre eux un $2$-morphisme dans le $2$-produit fibré.
C'est, par définition, un morphisme $\tilde t : u \to v$ tel que
$j(\tilde t) = t$.
\end{proof}

\noindent
Voici une caractérisation des catégories fibrées en groupoïdes
dont la diagonale est représentable par des espaces algébriques.

\begin{lemma}
\label{algebraic-lemma-representable-diagonal}
Soit $S$ un schéma appartenant à $\Sch_{fppf}$.
Soit $\mathcal{X}$ une catégorie fibrée en groupoïdes au-dessus de
$(\Sch/S)_{fppf}$. Les assertions suivantes sont équivalentes :
\begin{enumerate}
\item la diagonale $\mathcal{X} \to \mathcal{X} \times \mathcal{X}$
est représentable par des espaces algébriques,
\item pour tout schéma $U$ au-dessus de $S$ et tous
$x, y \in \Ob(\mathcal{X}_U)$, le faisceau
$\mathit{Isom}(x, y)$ est un espace algébrique au-dessus de $U$,
\item pour tout schéma $U$ au-dessus de $S$ et tout $x \in \Ob(\mathcal{X}_U)$,
le $1$-morphisme associé $x : (\Sch/U)_{fppf} \to \mathcal{X}$
est représentable par des espaces algébriques,
\item pour toute paire de schémas $T_1, T_2$ au-dessus de $S$ et tous
$x_i \in \Ob(\mathcal{X}_{T_i})$, $i = 1, 2$, le $2$-produit fibré
$(\Sch/T_1)_{fppf} \times_{x_1, \mathcal{X}, x_2}
(\Sch/T_2)_{fppf}$
est représentable par un espace algébrique,
\item pour toute catégorie fibrée en groupoïdes représentable $\mathcal{U}$
au-dessus de $(\Sch/S)_{fppf}$, tout $1$-morphisme
$\mathcal{U} \to \mathcal{X}$ est représentable par des espaces algébriques,
\item pour toute paire $\mathcal{T}_1, \mathcal{T}_2$ de catégories fibrées
en groupoïdes représentables au-dessus de $(\Sch/S)_{fppf}$ et tous
$1$-morphismes $x_i : \mathcal{T}_i \to \mathcal{X}$, $i = 1, 2$, le
$2$-produit fibré $\mathcal{T}_1 \times_{x_1, \mathcal{X}, x_2} \mathcal{T}_2$
est représentable par un espace algébrique,
\item pour toute catégorie fibrée en groupoïdes $\mathcal{U}$
au-dessus de $(\Sch/S)_{fppf}$ qui est
représentable par un espace algébrique, tout $1$-morphisme
$\mathcal{U} \to \mathcal{X}$ est représentable par des espaces algébriques,
\item pour toute paire $\mathcal{T}_1, \mathcal{T}_2$ de catégories fibrées
en groupoïdes au-dessus de $(\Sch/S)_{fppf}$ qui sont représentables
par des espaces algébriques, et tous $1$-morphismes
$x_i : \mathcal{T}_i \to \mathcal{X}$, le
$2$-produit fibré $\mathcal{T}_1 \times_{x_1, \mathcal{X}, x_2} \mathcal{T}_2$
est représentable par un espace algébrique.
\end{enumerate}
\end{lemma}

\begin{proof}
L'équivalence de (1) et (2) résulte du
Lemme \ref{stacks-lemma-isom-as-2-fibre-product} de Champs
et des définitions.
Démontrons l'équivalence de (1) et (3).
Posons $\mathcal{C} = (\Sch/S)_{fppf}$ pour la catégorie de base.
Nous utiliserons certaines des observations de la démonstration de l'énoncé analogue
Catégories, Lemme \ref{categories-lemma-representable-diagonal-groupoids}.
Nous emploierons le symbole $\cong$ au sens d'« équivalence de catégories fibrées
en groupoïdes au-dessus de $\mathcal{C} = (\Sch/S)_{fppf}$ ».
Supposons (1). Soient $U$ et $x$ comme dans (3). Pour tout schéma $V$
et tout $y \in \Ob(\mathcal{X}_V)$, on voit (comparer la référence ci-dessus) que
$$
\mathcal{C}/U
\times_{x, \mathcal{X}, y}
\mathcal{C}/V
\cong
(\mathcal{C}/U \times_S V)
\times_{(x, y), \mathcal{X} \times \mathcal{X}, \Delta}
\mathcal{X}
$$
qui est représentable par un espace algébrique par hypothèse. Réciproquement,
supposons (3). Considérons un schéma $U$ quelconque au-dessus de $S$ et une paire $(x, x')$
d'objets de $\mathcal{X}$ au-dessus de $U$. Il faut montrer que
$\mathcal{X} \times_{\Delta, \mathcal{X} \times \mathcal{X}, (x, x')} U$
est représentable par un espace algébrique. C'est clair, car
(comparer la référence ci-dessus)
$$
\mathcal{X}
\times_{\Delta, \mathcal{X} \times \mathcal{X}, (x, x')}
\mathcal{C}/U
\cong
(\mathcal{C}/U \times_{x, \mathcal{X}, x'} \mathcal{C}/U)
\times_{\mathcal{C}/U \times_S U, \Delta}
\mathcal{C}/U
$$
et le membre de droite est représentable par un espace algébrique par hypothèse
et parce que la catégorie des espaces algébriques au-dessus de $S$ admet des produits fibrés
et contient $U$ et $S$.

\medskip\noindent
Les équivalences
(3) $\Leftrightarrow$ (4),
(5) $\Leftrightarrow$ (6)
et
(7) $\Leftrightarrow$ (8)
sont formelles. Les équivalences
(3) $\Leftrightarrow$ (5) et
(4) $\Leftrightarrow$ (6)
résultent du
Lemme \ref{algebraic-lemma-representable-by-spaces-morphism-equivalent}.
Supposons (3), et soit $\mathcal{U} \to \mathcal{X}$ comme dans (7).
Pour démontrer (7), il faut montrer que, pour tout schéma $V$ et tout $1$-morphisme
$y : (\Sch/V)_{fppf} \to \mathcal{X}$, le $2$-produit fibré
$(\Sch/V)_{fppf} \times_{y, \mathcal{X}} \mathcal{U}$
est représentable par un espace algébrique. La propriété (3) nous dit
que $y$ est représentable par des espaces algébriques ; par conséquent, le
Lemme \ref{algebraic-lemma-base-change-by-space-representable-by-space}
donne l'assertion voulue. Enfin, (7) implique directement (3).
\end{proof}

\noindent
Dans la situation du lemme, pour tout $1$-morphisme
$x : (\Sch/U)_{fppf} \to \mathcal{X}$ comme dans le lemme, il est légitime
de dire que $x$ a la propriété $\mathcal{P}$, pour toute propriété comme dans la
Définition \ref{algebraic-definition-relative-representable-property}.
Cela vaut en particulier pour
$\mathcal{P} = $ « surjectif »,
$\mathcal{P} = $ « lisse » et
$\mathcal{P} = $ « \'etale »,
voir
Descente et espaces algébriques,
Lemmes \ref{spaces-descent-lemma-descending-property-surjective},
\ref{spaces-descent-lemma-descending-property-smooth} et
\ref{spaces-descent-lemma-descending-property-etale}.
Nous utiliserons ces trois cas dans les définitions
des champs algébriques ci-dessous.








\section{Champs en groupoïdes}
\label{algebraic-section-stacks}

\noindent
Soit $S$ un schéma appartenant à $\Sch_{fppf}$.
Rappelons qu'une catégorie $p : \mathcal{X} \to (\Sch/S)_{fppf}$
au-dessus de $(\Sch/S)_{fppf}$ est appelée un
{\it champ en groupoïdes} (voir
Champs, Définition \ref{stacks-definition-stack-in-groupoids})
si et seulement si
\begin{enumerate}
\item $p : \mathcal{X} \to (\Sch/S)_{fppf}$ est fibrée
en groupoïdes au-dessus de $(\Sch/S)_{fppf}$,
\item pour tout $U \in \Ob((\Sch/S)_{fppf})$ et
tous $x, y\in \Ob(\mathcal{X}_U)$, le préfaisceau
$\mathit{Isom}(x, y)$ est un faisceau sur le site $(\Sch/U)_{fppf}$, et
\item pour tout recouvrement $\mathcal{U} = \{U_i \to U\}$ de
$(\Sch/S)_{fppf}$, toutes les données de descente $(x_i, \phi_{ij})$
relatives à $\mathcal{U}$ sont effectives.
\end{enumerate}
Pour des exemples, voir
Exemples de champs,
section \ref{examples-stacks-section-examples-stacks-in-groupoids}
et suivantes.










\section{Champs algébriques}
\label{algebraic-section-algebraic-stacks}

\noindent
Voici la définition d'un champ algébrique. Remarquons que la condition
(2) permet de donner un sens à la condition de la partie (3) selon laquelle
$(\Sch/U)_{fppf} \to \mathcal{X}$
est lisse et surjective ; voir la discussion qui suit le
Lemme \ref{algebraic-lemma-representable-diagonal}.

\begin{definition}
\label{algebraic-definition-algebraic-stack}
Soit $S$ un schéma de base appartenant à $\Sch_{fppf}$.
Un {\it champ algébrique au-dessus de $S$} est une catégorie
$$
p : \mathcal{X} \to (\Sch/S)_{fppf}
$$
au-dessus de $(\Sch/S)_{fppf}$ qui possède les propriétés suivantes :
\begin{enumerate}
\item La catégorie $\mathcal{X}$ est un champ en groupoïdes au-dessus de
$(\Sch/S)_{fppf}$.
\item La diagonale
$\Delta : \mathcal{X} \to \mathcal{X} \times \mathcal{X}$
est représentable par des espaces algébriques.
\item Il existe un schéma $U \in \Ob((\Sch/S)_{fppf})$
et un $1$-morphisme $(\Sch/U)_{fppf} \to \mathcal{X}$
qui est surjectif et lisse\footnote{Dans les chapitres ultérieurs, nous noterons
simplement ce morphisme $U \to \mathcal{X}$, comme il est d'usage dans la littérature. Une autre
bonne possibilité consisterait à formuler cette condition par l'existence d'une
catégorie fibrée en groupoïdes représentable $\mathcal{U}$ et d'un $1$-morphisme
lisse et surjectif $\mathcal{U} \to \mathcal{X}$.}.
\end{enumerate}
\end{definition}

\noindent
Il existe quelques différences avec d'autres définitions de la littérature.

\medskip\noindent
La première est que nous exigeons que $\mathcal{X}$ soit un champ en groupoïdes
pour la topologie fppf, tandis que de nombreuses références utilisent la topologie \'etale.
La topologie fppf nous semble en quelque sorte être la topologie naturelle
avec laquelle travailler. En définitive, la $2$-catégorie de champs algébriques
ainsi obtenue est la même. Ceci est expliqué dans
Critères de représentabilité, section \ref{criteria-section-stacks-etale}.

\medskip\noindent
La deuxième est que nous exigeons seulement que l'application diagonale de $\mathcal{X}$ soit
représentable par des espaces algébriques, alors que la plupart des références imposent d'autres
conditions. Notre point de vue consiste à essayer de démontrer un certain
nombre des résultats qui suivent en supposant seulement que la diagonale
de $\mathcal{X}$ soit représentable par des espaces algébriques, puis à ajouter simplement
une hypothèse supplémentaire chaque fois que cela est nécessaire. Cela présente en outre
l'avantage que tout espace algébrique (tel que défini dans
Espaces, Définition \ref{spaces-definition-algebraic-space})
donne lieu à un champ algébrique.

\medskip\noindent
La troisième est que certains articles exigent l'existence d'un
schéma $U$ et d'un morphisme surjectif et \'etale $U \to \mathcal{X}$.
Dans l'article fondateur \cite{DM}, où les champs algébriques furent introduits,
Deligne et Mumford ont utilisé cette définition et montré que
le champ de modules des courbes stables de genre $g > 1$ est un champ algébrique
qui admet un recouvrement \'etale par un schéma. Michael Artin, voir
\cite{ArtinVersal}, a compris que de nombreux
résultats naturels sur les champs algébriques s'étendent au cas où l'on
suppose seulement l'existence d'un recouvrement lisse par un schéma. D'où notre choix ci-dessus.
Pour distinguer les deux cas, la littérature emploie les termes « champ de Deligne-Mumford »
et « champ d'Artin ». Nous réserverons le terme
« champ d'Artin » à un usage ultérieur (insérer ici une référence future), et continuerons
d'employer « champ algébrique » ; en revanche, nous emploierons « champ de Deligne-Mumford »
pour désigner les champs algébriques qui admettent un recouvrement \'etale par un
schéma.

\begin{definition}
\label{algebraic-definition-deligne-mumford}
Soit $S$ un schéma appartenant à $\Sch_{fppf}$.
Soit $\mathcal{X}$ un champ algébrique au-dessus de $S$.
On dit que $\mathcal{X}$ est un {\it champ de Deligne-Mumford} s'il existe
un schéma $U$ et un morphisme \'etale surjectif
$(\Sch/U)_{fppf} \to \mathcal{X}$.
\end{definition}

\noindent
Nous comparerons plus loin notre notion de champ de Deligne-Mumford avec
la notion définie dans l'article de Deligne et Mumford
(voir insérer ici une référence future).

\medskip\noindent
La catégorie des champs algébriques au-dessus de $S$ forme une $2$-catégorie.
En voici la définition précise.

\begin{definition}
\label{algebraic-definition-morphism-algebraic-stacks}
Soit $S$ un schéma appartenant à $\Sch_{fppf}$.
La {\it $2$-catégorie des champs algébriques au-dessus de $S$} est la
sous-$2$-catégorie de la $2$-catégorie des catégories fibrées en
groupoïdes au-dessus de $(\Sch/S)_{fppf}$ (voir
Catégories,
Définition \ref{categories-definition-categories-fibred-in-groupoids-over-C})
définie comme suit :
\begin{enumerate}
\item Ses objets sont les catégories fibrées en groupoïdes
au-dessus de $(\Sch/S)_{fppf}$ qui sont des champs algébriques au-dessus de $S$.
\item Ses $1$-morphismes $f : \mathcal{X} \to \mathcal{Y}$ sont
les foncteurs quelconques de catégories au-dessus de $(\Sch/S)_{fppf}$, comme dans
Catégories, Définition \ref{categories-definition-categories-over-C}.
\item Ses $2$-morphismes sont les transformations entre foncteurs
au-dessus de $(\Sch/S)_{fppf}$, comme dans
Catégories, Définition \ref{categories-definition-categories-over-C}.
\end{enumerate}
\end{definition}

\noindent
Autrement dit, cette $2$-catégorie est la sous-$2$-catégorie pleine de
$\textit{Cat}/(\Sch/S)_{fppf}$ dont les objets sont les champs algébriques.
Remarquons que tout $2$-morphisme est automatiquement un isomorphisme.
Il s'agit donc en fait d'une $(2, 1)$-catégorie et non pas seulement d'une $2$-catégorie.

\medskip\noindent
Nous verrons plus loin (insérer ici une référence future) que cette $2$-catégorie
admet des $2$-produits fibrés.

\medskip\noindent
Comme dans la remarque ci-dessus, la $2$-catégorie des champs algébriques au-dessus de $S$ est une
sous-$2$-catégorie pleine de la $2$-catégorie des catégories fibrées en groupoïdes
au-dessus de $(\Sch/S)_{fppf}$. Il se trouve qu'elle est stable par
équivalences. En voici l'énoncé précis.

\begin{lemma}
\label{algebraic-lemma-equivalent}
Soit $S$ un schéma appartenant à $\Sch_{fppf}$.
Soient $\mathcal{X}$, $\mathcal{Y}$ des catégories au-dessus de $(\Sch/S)_{fppf}$.
Supposons que $\mathcal{X}$, $\mathcal{Y}$ soient équivalentes comme catégories au-dessus de
$(\Sch/S)_{fppf}$. Alors $\mathcal{X}$ est un champ algébrique si et
seulement si $\mathcal{Y}$ est un champ algébrique. De même, $\mathcal{X}$
est un champ de Deligne-Mumford si et seulement si $\mathcal{Y}$ est un champ de Deligne-Mumford.
\end{lemma}

\begin{proof}
Supposons que $\mathcal{X}$ soit un champ algébrique (resp.\ un champ de Deligne-Mumford). Le
Lemme \ref{stacks-lemma-stack-in-groupoids-equivalent} de Champs
implique alors que $\mathcal{Y}$ est un champ en groupoïdes au-dessus de
$\Sch_{fppf}$. Choisissons une équivalence $f : \mathcal{X} \to \mathcal{Y}$
au-dessus de $\Sch_{fppf}$. Elle donne un diagramme $2$-commutatif
$$
\xymatrix{
\mathcal{X} \ar[r]_f \ar[d]_{\Delta_\mathcal{X}} &
\mathcal{Y} \ar[d]^{\Delta_\mathcal{Y}} \\
\mathcal{X} \times \mathcal{X} \ar[r]^{f \times f} &
\mathcal{Y} \times \mathcal{Y}
}
$$
dont les flèches horizontales sont des équivalences. Cela implique que
$\Delta_\mathcal{Y}$ est représentable par des espaces algébriques d'après le
Lemme \ref{algebraic-lemma-representable-by-spaces-morphism-equivalent}.
Enfin, soit $U$ un schéma au-dessus de $S$, et soit
$x : (\Sch/U)_{fppf} \to \mathcal{X}$ un $1$-morphisme qui
est surjectif et lisse (resp.\ \'etale). En considérant le diagramme
$$
\xymatrix{
(\Sch/U)_{fppf} \ar[r]_{\text{id}} \ar[d]_x &
(\Sch/U)_{fppf} \ar[d]^{f \circ x} \\
\mathcal{X} \ar[r]^f &
\mathcal{Y}
}
$$
et en appliquant le
Lemme \ref{algebraic-lemma-property-morphism-equivalent},
on conclut que $f \circ x$ est surjectif et lisse (resp.\ \'etale),
comme voulu.
\end{proof}





\section{Champs algébriques et espaces algébriques}
\label{algebraic-section-stacks-spaces}

\noindent
Dans cette section, nous examinons quelques critères simples qui impliquent qu'un
champ algébrique est un espace algébrique. Le résultat principal affirme que c'est
exactement le cas lorsque les objets des catégories fibres n'ont aucun
automorphisme non trivial. Ce n'est pas une trivialité ! Avant d'y venir,
effectuons d'abord un contrôle de cohérence.

\begin{lemma}
\label{algebraic-lemma-representable-algebraic}
Soit $S$ un schéma appartenant à $\Sch_{fppf}$.
\begin{enumerate}
\item Une catégorie fibrée en groupoïdes
$p : \mathcal{X} \to (\Sch/S)_{fppf}$
qui est représentable par un espace algébrique est un champ de Deligne-Mumford.
\item Si $F$ est un espace algébrique au-dessus de $S$, alors la
catégorie fibrée en groupoïdes associée
$p : \mathcal{S}_F \to (\Sch/S)_{fppf}$
est un champ de Deligne-Mumford.
\item Si $X \in \Ob((\Sch/S)_{fppf})$, alors
$(\Sch/X)_{fppf} \to (\Sch/S)_{fppf}$ est
un champ de Deligne-Mumford.
\end{enumerate}
\end{lemma}

\begin{proof}
Il est clair que (2) implique (3).
Les parties (1) et (2) sont équivalentes par le Lemme \ref{algebraic-lemma-equivalent}.
Il suffit donc de démontrer (2).
Tout d'abord, remarquons que $\mathcal{S}_F$ est un champ en ensembles puisque
$F$ est un faisceau (Champs, Lemme
\ref{stacks-lemma-stack-in-setoids-characterize}).
A fortiori, c'est un champ en groupoïdes. Ensuite, le morphisme diagonal
$\mathcal{S}_F \to \mathcal{S}_F  \times \mathcal{S}_F$
est le même que le morphisme $\mathcal{S}_F \to \mathcal{S}_{F \times F}$
qui provient de la diagonale de $F$. Il est donc représentable
par des espaces algébriques d'après le
Lemme \ref{algebraic-lemma-morphism-spaces-gives-representable-by-spaces}.
En fait, il est même représentable (par des schémas), puisque la diagonale d'un
espace algébrique est représentable, mais nous n'en avons pas besoin.
Soit $U$ un schéma et soit $h_U \to F$ un morphisme \'etale surjectif.
Nous pouvons le considérer comme un morphisme \'etale surjectif d'espaces algébriques.
Ainsi, par le
Lemme
\ref{algebraic-lemma-map-presheaves-representable-by-spaces-transformation-property},
le $1$-morphisme correspondant $(\Sch/U)_{fppf} \to \mathcal{S}_F$
est surjectif et \'etale.
\end{proof}

\noindent
Le résultat suivant affirme qu'un champ de Deligne-Mumford dont l'inertie
est triviale « est » un espace algébrique. Ce lemme sera rendu caduc par la
Proposition plus forte \ref{algebraic-proposition-algebraic-stack-no-automorphisms}
ci-dessous, qui affirme que cela vaut plus généralement pour les champs algébriques…

\begin{lemma}
\label{algebraic-lemma-algebraic-stack-no-automorphisms}
Soit $S$ un schéma appartenant à $\Sch_{fppf}$.
Soit $\mathcal{X}$ un champ algébrique au-dessus de $S$.
Les assertions suivantes sont équivalentes :
\begin{enumerate}
\item $\mathcal{X}$ est un champ de Deligne-Mumford et est un champ en setoïdes,
\item $\mathcal{X}$ est un champ de Deligne-Mumford tel que le
$1$-morphisme canonique $\mathcal{I}_\mathcal{X} \to \mathcal{X}$
soit une équivalence, et
\item $\mathcal{X}$ est représentable par un espace algébrique.
\end{enumerate}
\end{lemma}

\begin{proof}
L'équivalence de (1) et (2) résulte du
Lemme \ref{stacks-lemma-characterize-stack-in-setoids} de Champs.
L'implication (3) $\Rightarrow$ (1) résulte du
Lemme \ref{algebraic-lemma-representable-algebraic}.
Enfin, supposons (1). Par le
Lemme \ref{stacks-lemma-stack-in-setoids-characterize} de Champs,
il existe un faisceau $F$ sur $(\Sch/S)_{fppf}$
et une équivalence $j : \mathcal{X} \to \mathcal{S}_F$. D'après le
Lemme \ref{algebraic-lemma-map-presheaves-representable-by-algebraic-spaces},
le fait que $\Delta_\mathcal{X}$ soit représentable par des espaces
algébriques signifie que $\Delta_F : F \to F \times F$
est représentable par des espaces algébriques.
Soit $U$ un schéma et soit $x : (\Sch/U)_{fppf} \to \mathcal{X}$
un morphisme \'etale surjectif. La composée
$j \circ x : (\Sch/U)_{fppf} \to \mathcal{S}_F$
correspond à un morphisme $h_U \to F$ de faisceaux. Par le
Lemme \ref{bootstrap-lemma-representable-diagonal} de Amorçage,
ce morphisme est représentable par des espaces algébriques.
Ainsi, par le
Lemme \ref{algebraic-lemma-map-fibred-setoids-property},
on conclut que $h_U \to F$ est surjectif et \'etale.
Enfin, on applique le
Théorème \ref{bootstrap-theorem-bootstrap} de Amorçage
pour voir que $F$ est un espace algébrique.
\end{proof}

\begin{proposition}
\label{algebraic-proposition-algebraic-stack-no-automorphisms}
Soit $S$ un schéma appartenant à $\Sch_{fppf}$.
Soit $\mathcal{X}$ un champ algébrique au-dessus de $S$.
Les assertions suivantes sont équivalentes :
\begin{enumerate}
\item $\mathcal{X}$ est un champ en setoïdes,
\item le $1$-morphisme canonique $\mathcal{I}_\mathcal{X} \to \mathcal{X}$
est une équivalence, et
\item $\mathcal{X}$ est représentable par un espace algébrique.
\end{enumerate}
\end{proposition}

\begin{proof}
L'équivalence de (1) et (2) résulte du
Lemme \ref{stacks-lemma-characterize-stack-in-setoids} de Champs.
L'implication (3) $\Rightarrow$ (1) résulte du
Lemme \ref{algebraic-lemma-algebraic-stack-no-automorphisms}.
Enfin, supposons (1). Par le
Lemme \ref{stacks-lemma-stack-in-setoids-characterize} de Champs,
il existe une équivalence $j : \mathcal{X} \to \mathcal{S}_F$
où $F$ est un faisceau sur $(\Sch/S)_{fppf}$.  D'après le
Lemme \ref{algebraic-lemma-map-presheaves-representable-by-algebraic-spaces},
le fait que $\Delta_\mathcal{X}$ soit représentable par des espaces
algébriques signifie que $\Delta_F : F \to F \times F$
est représentable par des espaces algébriques.
Soit $U$ un schéma et soit $x : (\Sch/U)_{fppf} \to \mathcal{X}$
un morphisme lisse surjectif. La composée
$j \circ x : (\Sch/U)_{fppf} \to \mathcal{S}_F$
correspond à un morphisme $h_U \to F$ de faisceaux. Par le
Lemme \ref{bootstrap-lemma-representable-diagonal} de Amorçage,
ce morphisme est représentable par des espaces algébriques.
Ainsi, par le
Lemme \ref{algebraic-lemma-map-fibred-setoids-property},
on conclut que $h_U \to F$ est surjectif et lisse.
En particulier, il est surjectif, plat et localement de présentation finie
(par le
Lemme \ref{algebraic-lemma-representable-transformations-property-implication}
et le fait qu'un morphisme lisse d'espaces algébriques est plat et
localement de présentation finie, voir
Morphismes d'espaces algébriques,
Lemmes \ref{spaces-morphisms-lemma-smooth-locally-finite-presentation} et
\ref{spaces-morphisms-lemma-smooth-flat}).
Enfin, on applique le
Théorème \ref{bootstrap-theorem-final-bootstrap} de Amorçage
pour voir que $F$ est un espace algébrique.
\end{proof}







\section{2-produits fibrés de champs algébriques}
\label{algebraic-section-2-fibre-products}

\noindent
La $2$-catégorie des champs algébriques admet des produits et des $2$-produits fibrés.
Le premier lemme est en réalité un cas particulier du
Lemme \ref{algebraic-lemma-2-fibre-product},
mais sa démonstration est légèrement plus simple.

\begin{lemma}
\label{algebraic-lemma-product-spaces}
Soit $S$ un schéma appartenant à $\Sch_{fppf}$.
Soient $\mathcal{X}$, $\mathcal{Y}$ des champs algébriques au-dessus de $S$.
Alors $\mathcal{X} \times_{(\Sch/S)_{fppf}} \mathcal{Y}$
est un champ algébrique, et est un produit dans la $2$-catégorie des
champs algébriques au-dessus de $S$.
\end{lemma}

\begin{proof}
Un objet de $\mathcal{X} \times_{(\Sch/S)_{fppf}} \mathcal{Y}$
au-dessus de $T$ est simplement une paire $(x, y)$, où $x$ est un objet de $\mathcal{X}_T$
et $y$ un objet de $\mathcal{Y}_T$. Il résulte donc immédiatement
des définitions que
$\mathcal{X} \times_{(\Sch/S)_{fppf}} \mathcal{Y}$ est un
champ en groupoïdes. Si $(x, y)$ et $(x', y')$ sont
deux objets de $\mathcal{X} \times_{(\Sch/S)_{fppf}} \mathcal{Y}$
au-dessus de $T$, alors
$$
\mathit{Isom}((x, y), (x', y')) =
\mathit{Isom}(x, x') \times \mathit{Isom}(y, y').
$$
Il résulte donc des équivalences du
Lemme \ref{algebraic-lemma-representable-diagonal}
et du fait que la catégorie des espaces algébriques admet des produits
que la diagonale de $\mathcal{X} \times_{(\Sch/S)_{fppf}} \mathcal{Y}$
est représentable par des espaces algébriques.
Enfin, supposons que $U, V \in \Ob((\Sch/S)_{fppf})$,
et soient $x, y$ des morphismes lisses surjectifs
$x : (\Sch/U)_{fppf} \to \mathcal{X}$,
$y : (\Sch/V)_{fppf} \to \mathcal{Y}$.
Remarquons que
$$
(\Sch/U \times_S V)_{fppf} =
(\Sch/U)_{fppf}
\times_{(\Sch/S)_{fppf}} (\Sch/V)_{fppf}.
$$
L'objet $(\text{pr}_U^*x, \text{pr}_V^*y)$ de
$\mathcal{X} \times_{(\Sch/S)_{fppf}} \mathcal{Y}$ au-dessus de
$(\Sch/U \times_S V)_{fppf}$ définit ainsi un $1$-morphisme
$$
(\Sch/U \times_S V)_{fppf}
\longrightarrow
\mathcal{X} \times_{(\Sch/S)_{fppf}} \mathcal{Y}
$$
qui est la composée de changements de base de $x$ et de $y$, donc
est surjectif et lisse, voir les
Lemmes \ref{algebraic-lemma-base-change-representable-transformations-property} et
\ref{algebraic-lemma-composition-representable-transformations-property}.
On conclut que $\mathcal{X} \times_{(\Sch/S)_{fppf}} \mathcal{Y}$
est bien un champ algébrique. Nous omettons de vérifier qu'il
est effectivement un produit.
\end{proof}

\begin{lemma}
\label{algebraic-lemma-2-fibre-product-general}
Soit $S$ un schéma appartenant à $\Sch_{fppf}$.
Soit $\mathcal{Z}$ un champ en groupoïdes au-dessus de $(\Sch/S)_{fppf}$
dont la diagonale est représentable par des espaces algébriques.
Soient $\mathcal{X}$, $\mathcal{Y}$ des champs algébriques au-dessus de $S$.
Soient $f : \mathcal{X} \to \mathcal{Z}$, $g : \mathcal{Y} \to \mathcal{Z}$
des $1$-morphismes de champs en groupoïdes. Alors le $2$-produit fibré
$\mathcal{X} \times_{f, \mathcal{Z}, g} \mathcal{Y}$ est un champ algébrique.
\end{lemma}

\begin{proof}
Il faut vérifier les conditions (1), (2) et (3) de la
Définition \ref{algebraic-definition-algebraic-stack}.
La première condition résulte du
Lemme \ref{stacks-lemma-2-product-stacks-in-groupoids} de Champs.

\medskip\noindent
La deuxième condition à vérifier est que les faisceaux $\mathit{Isom}$
soient représentables par des espaces algébriques. Pour cela, supposons que
$T$ soit un schéma au-dessus de $S$, et que $u, v$ soient des objets de
$(\mathcal{X} \times_{f, \mathcal{Z}, g} \mathcal{Y})_T$.
D'après notre construction des $2$-produits fibrés (qui remonte jusqu'à
Catégories, Lemme \ref{categories-lemma-2-product-categories-over-C}),
nous pouvons écrire $u = (x, y, \alpha)$ et $v = (x', y', \alpha')$.
Ici, $\alpha : f(x) \to g(y)$, et de même pour $\alpha'$.
Il est alors clair que
$$
\xymatrix{
\mathit{Isom}(u, v) \ar[d] \ar[rr] & &
\mathit{Isom}(y, y') \ar[d]^{\phi \mapsto g(\phi) \circ \alpha} \\
\mathit{Isom}(x, x') \ar[rr]^-{\psi \mapsto \alpha' \circ f(\psi)} & &
\mathit{Isom}(f(x), g(y'))
}
$$
est un diagramme cartésien de faisceaux sur $(\Sch/T)_{fppf}$.
Puisque, par hypothèse, les faisceaux
$\mathit{Isom}(y, y')$, $\mathit{Isom}(x, x')$, $\mathit{Isom}(f(x), g(y'))$
sont des espaces algébriques (voir le
Lemme \ref{algebraic-lemma-representable-diagonal}),
on voit que $\mathit{Isom}(u, v)$
est un espace algébrique.

\medskip\noindent
Soient $U, V \in \Ob((\Sch/S)_{fppf})$,
et soient $x, y$ des morphismes lisses surjectifs
$x : (\Sch/U)_{fppf} \to \mathcal{X}$,
$y : (\Sch/V)_{fppf} \to \mathcal{Y}$.
Considérons le morphisme
$$
(\Sch/U)_{fppf}
\times_{f \circ x, \mathcal{Z}, g \circ y}
(\Sch/V)_{fppf}
\longrightarrow
\mathcal{X} \times_{f, \mathcal{Z}, g} \mathcal{Y}.
$$
Comme la diagonale de $\mathcal{Z}$ est représentable par des espaces algébriques,
la source de cette flèche est représentable par un espace algébrique $F$, voir le
Lemme \ref{algebraic-lemma-representable-diagonal}.
De plus, le morphisme est la composée
de changements de base de $x$ et de $y$, donc est surjectif et lisse, voir les
Lemmes \ref{algebraic-lemma-base-change-representable-transformations-property} et
\ref{algebraic-lemma-composition-representable-transformations-property}.
En choisissant un schéma $W$ et un morphisme \'etale surjectif $W \to F$,
on voit que la composée du $1$-morphisme affiché
avec le $1$-morphisme correspondant
$$
(\Sch/W)_{fppf}
\longrightarrow
(\Sch/U)_{fppf}
\times_{f \circ x, \mathcal{Z}, g \circ y}
(\Sch/V)_{fppf}
$$
est surjective et lisse, ce qui démontre la dernière condition.
\end{proof}

\begin{lemma}
\label{algebraic-lemma-2-fibre-product}
Soit $S$ un schéma appartenant à $\Sch_{fppf}$.
Soient $\mathcal{X}, \mathcal{Y}, \mathcal{Z}$ des champs algébriques au-dessus de $S$.
Soient $f : \mathcal{X} \to \mathcal{Z}$, $g : \mathcal{Y} \to \mathcal{Z}$
des $1$-morphismes de champs algébriques. Alors le $2$-produit fibré
$\mathcal{X} \times_{f, \mathcal{Z}, g} \mathcal{Y}$ est un champ algébrique.
C'est aussi le $2$-produit fibré dans la $2$-catégorie des champs algébriques
au-dessus de $(\Sch/S)_{fppf}$.
\end{lemma}

\begin{proof}
Le fait que $\mathcal{X} \times_{f, \mathcal{Z}, g} \mathcal{Y}$ soit un
champ algébrique résulte du
Lemme plus fort \ref{algebraic-lemma-2-fibre-product-general}.
Le fait que $\mathcal{X} \times_{f, \mathcal{Z}, g} \mathcal{Y}$
soit un $2$-produit fibré dans la $2$-catégorie des champs algébriques au-dessus de $S$
résulte formellement du fait que la $2$-catégorie des champs algébriques
au-dessus de $S$ est une sous-$2$-catégorie pleine de la $2$-catégorie des champs en
groupoïdes au-dessus de $(\Sch/S)_{fppf}$.
\end{proof}








\section{Champs algébriques, nouvelle présentation}
\label{algebraic-section-overhaul}

\noindent
Quelques résultats fondamentaux sur les champs algébriques.

\begin{lemma}
\label{algebraic-lemma-lift-morphism-presentations}
Soit $S$ un schéma appartenant à $\Sch_{fppf}$.
Soit $f : \mathcal{X} \to \mathcal{Y}$ un $1$-morphisme de champs
algébriques au-dessus de $S$.
Soit $V \in \Ob((\Sch/S)_{fppf})$.
Soit $y : (\Sch/V)_{fppf} \to \mathcal{Y}$ surjectif et lisse.
Alors il existe un objet $U \in \Ob((\Sch/S)_{fppf})$
et un diagramme $2$-commutatif
$$
\xymatrix{
(\Sch/U)_{fppf} \ar[r]_a \ar[d]_x &
(\Sch/V)_{fppf} \ar[d]^y \\
\mathcal{X} \ar[r]^f & \mathcal{Y}
}
$$
où $x$ est surjectif et lisse.
\end{lemma}

\begin{proof}
Choisissons d'abord $W \in \Ob((\Sch/S)_{fppf})$ et un $1$-morphisme
lisse surjectif $z : (\Sch/W)_{fppf} \to \mathcal{X}$.
Comme $\mathcal{Y}$ est un champ algébrique, on peut choisir une équivalence
$$
j :
\mathcal{S}_F
\longrightarrow
(\Sch/W)_{fppf}
\times_{f \circ z, \mathcal{Y}, y}
(\Sch/V)_{fppf}
$$
où $F$ est un espace algébrique. Par le
Lemme \ref{algebraic-lemma-base-change-representable-transformations-property},
le morphisme
$\mathcal{S}_F \to (\Sch/W)_{fppf}$ est surjectif et lisse
comme changement de base de $y$. Ainsi, par le
Lemme \ref{algebraic-lemma-composition-representable-transformations-property},
on voit que $\mathcal{S}_F \to \mathcal{X}$ est surjectif et lisse.
Choisissons un objet $U \in \Ob((\Sch/S)_{fppf})$
et un morphisme \'etale surjectif $U \to F$. En appliquant alors
une fois encore le Lemme \ref{algebraic-lemma-composition-representable-transformations-property},
on obtient les propriétés voulues.
\end{proof}

\noindent
Ce lemme généralise la
Proposition \ref{algebraic-proposition-algebraic-stack-no-automorphisms}.

\begin{lemma}
\label{algebraic-lemma-characterize-representable-by-algebraic-spaces}
Soit $S$ un schéma appartenant à $\Sch_{fppf}$.
Soit $f : \mathcal{X} \to \mathcal{Y}$ un $1$-morphisme de champs
algébriques au-dessus de $S$. Les assertions suivantes sont équivalentes :
\begin{enumerate}
\item pour tout $U \in \Ob((\Sch/S)_{fppf})$,
le foncteur $f : \mathcal{X}_U \to \mathcal{Y}_U$ est fidèle,
\item le foncteur $f$ est fidèle, et
\item $f$ est représentable par des espaces algébriques.
\end{enumerate}
\end{lemma}

\begin{proof}
Les parties (1) et (2) sont équivalentes par les propriétés générales des $1$-morphismes
de catégories fibrées en groupoïdes, voir
Catégories, Lemme \ref{categories-lemma-equivalence-fibred-categories}.
On voit que (3) implique (2) par le
Lemme \ref{algebraic-lemma-criterion-map-representable-spaces-fibred-in-groupoids}.
Enfin, supposons (2).
Soit $U$ un schéma. Soit $y \in \Ob(\mathcal{Y}_U)$.
Il faut démontrer que
$$
\mathcal{W} = (\Sch/U)_{fppf} \times_{y, \mathcal{Y}} \mathcal{X}
$$
est représentable par un espace algébrique au-dessus de $U$. Puisque
$(\Sch/U)_{fppf}$ est un champ algébrique, il résulte du
Lemme \ref{algebraic-lemma-2-fibre-product}
que $\mathcal{W}$ est un champ algébrique.
D'autre part, la description explicite des objets de $\mathcal{W}$
comme triplets $(V, x, \alpha : y(V) \to f(x))$ et le fait que $f$ soit
fidèle montrent que les catégories fibres de $\mathcal{W}$ sont des setoïdes. Par conséquent, la
Proposition \ref{algebraic-proposition-algebraic-stack-no-automorphisms}
garantit que $\mathcal{W}$ est représentable par un espace algébrique.
\end{proof}

\begin{lemma}
\label{algebraic-lemma-smooth-surjective-morphism-implies-algebraic}
Soit $S$ un schéma appartenant à $\Sch_{fppf}$.
Soit $u : \mathcal{U} \to \mathcal{X}$ un $1$-morphisme de
champs en groupoïdes au-dessus de $(\Sch/S)_{fppf}$. Si
\begin{enumerate}
\item $\mathcal{U}$ est représentable par un espace algébrique, et
\item $u$ est représentable par des espaces algébriques, surjectif et lisse,
\end{enumerate}
alors $\mathcal X$ est un champ algébrique au-dessus de $S$.
\end{lemma}

\begin{proof}
Il faut montrer que $\Delta : \mathcal{X} \to \mathcal{X} \times \mathcal{X}$
est représentable par des espaces algébriques, voir la
Définition \ref{algebraic-definition-algebraic-stack}.
Étant donnés deux schémas $T_1$, $T_2$ au-dessus de $S$, notons
$\mathcal{T}_i = (\Sch/T_i)_{fppf}$ les catégories fibrées représentables
associées. Supposons donnés des $1$-morphismes
$f_i : \mathcal{T}_i \to \mathcal{X}$.
D'après le
Lemme \ref{algebraic-lemma-representable-diagonal},
il suffit de démontrer que le $2$-produit
fibré $\mathcal{T}_1 \times_\mathcal{X} \mathcal{T}_2$
est représentable par un espace algébrique. Par le
Lemme de Champs
\ref{stacks-lemma-2-fibre-product-stacks-in-setoids-over-stack-in-groupoids},
il s'agit en tout cas d'un champ en setoïdes. Ainsi,
$\mathcal{T}_1 \times_\mathcal{X} \mathcal{T}_2$ correspond
à un faisceau $F$ sur $(\Sch/S)_{fppf}$, voir
Champs, Lemme \ref{stacks-lemma-stack-in-setoids-characterize}.
Soit $U$ l'espace algébrique qui représente $\mathcal{U}$.
Par hypothèse,
$$
\mathcal{T}_i' = \mathcal{U} \times_{u, \mathcal{X}, f_i} \mathcal{T}_i
$$
est représentable par un espace algébrique $T'_i$ au-dessus de $S$. Par conséquent,
$\mathcal{T}_1' \times_\mathcal{U} \mathcal{T}_2'$ est représentable
par l'espace algébrique $T'_1 \times_U T'_2$.
Considérons le diagramme commutatif
$$
\xymatrix{
&
\mathcal{T}_1 \times_{\mathcal X} \mathcal{T}_2 \ar[rr]\ar'[d][dd] & &
\mathcal{T}_1 \ar[dd] \\
\mathcal{T}_1' \times_\mathcal{U} \mathcal{T}_2' \ar[ur]\ar[rr]\ar[dd] & &
\mathcal{T}_1' \ar[ur]\ar[dd] \\
&
\mathcal{T}_2 \ar'[r][rr] & &
\mathcal X \\
\mathcal{T}_2' \ar[rr]\ar[ur] & &
\mathcal{U} \ar[ur] }
$$
Dans ce diagramme, le carré inférieur, le carré droit, le carré arrière et
le carré avant sont des $2$-produits fibrés. Un argument formel montre alors
que $\mathcal{T}_1' \times_\mathcal{U} \mathcal{T}_2' \to
\mathcal{T}_1 \times_{\mathcal X} \mathcal{T}_2$
est le « changement de base » de $\mathcal{U} \to \mathcal{X}$ ; plus précisément,
le diagramme
$$
\xymatrix{
\mathcal{T}_1' \times_\mathcal{U} \mathcal{T}_2' \ar[d] \ar[r] &
\mathcal{U} \ar[d] \\
\mathcal{T}_1 \times_{\mathcal X} \mathcal{T}_2 \ar[r] &
\mathcal{X}
}
$$
est un carré $2$-cartésien.
Ainsi, $T'_1 \times_U T'_2 \to F$ est représentable par des espaces algébriques,
lisse et surjectif, voir les
Lemmes \ref{algebraic-lemma-map-fibred-setoids-representable-algebraic-spaces},
\ref{algebraic-lemma-base-change-representable-by-spaces},
\ref{algebraic-lemma-map-fibred-setoids-property} et
\ref{algebraic-lemma-base-change-representable-transformations-property}.
Par conséquent, $F$ est un espace algébrique par le
Théorème \ref{bootstrap-theorem-final-bootstrap} de Amorçage,
ce qui conclut.
\end{proof}

\noindent
Une application du
Lemme \ref{algebraic-lemma-smooth-surjective-morphism-implies-algebraic}
est que ce qui est un espace algébrique au-dessus d'un champ algébrique
est un champ algébrique. C'est l'analogue du
Lemme \ref{bootstrap-lemma-representable-by-spaces-over-space} de Amorçage.
En fait, il suffit de supposer que le morphisme
$\mathcal{X} \to \mathcal{Y}$ est « algébrique », comme nous le verrons dans
Critères de représentabilité,
Lemme \ref{criteria-lemma-algebraic-morphism-to-algebraic}.

\begin{lemma}
\label{algebraic-lemma-representable-morphism-to-algebraic}
Soit $S$ un schéma appartenant à $\Sch_{fppf}$.
Soit $\mathcal{X} \to \mathcal{Y}$ un morphisme de champs en groupoïdes
au-dessus de $(\Sch/S)_{fppf}$. Supposons que
\begin{enumerate}
\item $\mathcal{X} \to \mathcal{Y}$ soit représentable par des espaces algébriques, et
\item $\mathcal{Y}$ soit un champ algébrique au-dessus de $S$.
\end{enumerate}
Alors $\mathcal{X}$ est un champ algébrique au-dessus de $S$.
\end{lemma}

\begin{proof}
Soit $\mathcal{V} \to \mathcal{Y}$ un $1$-morphisme lisse surjectif
d'un champ en groupoïdes représentable vers $\mathcal{Y}$. Il existe par la
Définition \ref{algebraic-definition-algebraic-stack}.
Alors le $2$-produit fibré
$\mathcal{U} = \mathcal{V} \times_{\mathcal Y} \mathcal X$
est représentable par un espace algébrique d'après le
Lemme \ref{algebraic-lemma-base-change-by-space-representable-by-space}.
Le $1$-morphisme $\mathcal{U} \to \mathcal X$ est représentable par des espaces
algébriques, lisse et surjectif, voir les
Lemmes \ref{algebraic-lemma-base-change-representable-by-spaces} et
\ref{algebraic-lemma-base-change-representable-transformations-property}.
Par le
Lemme \ref{algebraic-lemma-smooth-surjective-morphism-implies-algebraic},
on conclut que $\mathcal{X}$ est un champ algébrique.
\end{proof}

\begin{lemma}
\label{algebraic-lemma-open-fibred-category-is-algebraic}
\begin{reference}
La suppression de l'hypothèse selon laquelle $j$ est un monomorphisme a été signalée
dans un courriel de Matthew Emerton daté du 15 juin 2016.
\end{reference}
Soit $S$ un schéma appartenant à $\Sch_{fppf}$.
Soit $j : \mathcal X \to \mathcal Y$ un $1$-morphisme de
catégories fibrées en groupoïdes au-dessus de $(\Sch/S)_{fppf}$.
Supposons que $j$ soit représentable par des espaces algébriques.
Alors, si $\mathcal{Y}$ est un champ en groupoïdes
(resp.\ un champ algébrique), il en est de même de $\mathcal{X}$.
\end{lemma}

\begin{proof}
L'assertion concernant les champs algébriques résulte de celle concernant les
champs en groupoïdes par le Lemme \ref{algebraic-lemma-representable-morphism-to-algebraic}.
Si $j$ est représentable par des espaces algébriques, alors $j$ est
fidèle sur les catégories fibres et, pour chaque $U$ et chaque
$y \in \Ob(\mathcal{Y}_U)$, le préfaisceau
$$
(h : V \to U)
\longmapsto
\{(x, \phi) \mid x \in \Ob(\mathcal{X}_V), \phi : h^*y \to f(x)\}/\cong
$$
est un espace algébrique au-dessus de $U$. Voir le
Lemme \ref{algebraic-lemma-criterion-map-representable-spaces-fibred-in-groupoids}.
En particulier, ce préfaisceau est un faisceau, et la conclusion résulte du
Lemme \ref{stacks-lemma-relative-sheaf-over-stack-is-stack} de Champs.
\end{proof}



\section{D'un champ algébrique à une présentation}
\label{algebraic-section-stack-to-presentation}

\noindent
Étant donné un champ algébrique au-dessus de $S$, on obtient un groupoïde en espaces algébriques
au-dessus de $S$ dont le champ quotient associé est le champ algébrique.

\medskip\noindent
Rappelons que si $(U, R, s, t, c)$ est un groupoïde en espaces algébriques au-dessus de $S$,
alors $[U/R]$ désigne le champ quotient associé à cette donnée, voir
Groupoïdes in Espaces,
Définition \ref{spaces-groupoids-definition-quotient-stack}.
En général, $[U/R]$ n'est {\bf pas} un champ algébrique. En particulier, le
champ $[U/R]$ qui apparaît dans le lemme suivant n'est en général pas
algébrique.

\begin{lemma}
\label{algebraic-lemma-map-space-into-stack}
Soit $S$ un schéma appartenant à $\Sch_{fppf}$.
Soit $\mathcal{X}$ un champ algébrique au-dessus de $S$.
Soit $\mathcal{U}$ un champ algébrique au-dessus de $S$ qui
est représentable par un espace algébrique.
Soit $f : \mathcal{U} \to \mathcal{X}$ un 1-morphisme. Alors
\begin{enumerate}
\item le $2$-produit fibré
$\mathcal{R} = \mathcal{U} \times_{f, \mathcal{X}, f} \mathcal{U}$
est représentable par un espace algébrique,
\item il existe une équivalence canonique
$$
\mathcal{U} \times_{f, \mathcal{X}, f} \mathcal{U}
\times_{f, \mathcal{X}, f} \mathcal{U} =
\mathcal{R} \times_{\text{pr}_1, \mathcal{U}, \text{pr}_0} \mathcal{R},
$$
\item la projection $\text{pr}_{02}$ induit, via (2), un $1$-morphisme
$$
\text{pr}_{02} :
\mathcal{R} \times_{\text{pr}_1, \mathcal{U}, \text{pr}_0} \mathcal{R}
\longrightarrow
\mathcal{R}
$$
\item soient $U$, $R$ les espaces algébriques représentant
$\mathcal{U}, \mathcal{R}$, et soient $t, s : R \to U$ et
$c : R \times_{s, U, t} R \to R$ les morphismes correspondant
aux $1$-morphismes
$\text{pr}_0, \text{pr}_1 : \mathcal{R} \to \mathcal{U}$
et au
$\text{pr}_{02} :
\mathcal{R} \times_{\text{pr}_1, \mathcal{U}, \text{pr}_0} \mathcal{R} \to
\mathcal{R}$ ci-dessus ; alors le quintuplet $(U, R, s, t, c)$ est un groupoïde en
espaces algébriques au-dessus de $S$,
\item le morphisme $f$ induit un $1$-morphisme canonique
$f_{can} : [U/R] \to \mathcal{X}$
de champs en groupoïdes au-dessus de $(\Sch/S)_{fppf}$, et
\item le $1$-morphisme $f_{can} : [U/R] \to \mathcal{X}$ est pleinement fidèle.
\end{enumerate}
\end{lemma}

\begin{proof}
Démonstration de (1). Par définition, $\Delta_\mathcal{X}$ est représentable
par des espaces algébriques ; le
Lemme \ref{algebraic-lemma-representable-diagonal}
s'applique donc et montre que $\mathcal{U} \to \mathcal{X}$ est représentable
par des espaces algébriques. Le résultat découle alors du
Lemme \ref{algebraic-lemma-base-change-by-space-representable-by-space}.

\medskip\noindent
Soit $T$ un schéma au-dessus de $S$. Par construction du $2$-produit fibré (voir
Catégories, Lemme \ref{categories-lemma-2-product-categories-over-C}),
on voit que les objets de la catégorie fibre $\mathcal{R}_T$
sont les triplets $(a, b, \alpha)$ où $a, b \in \Ob(\mathcal{U}_T)$
et où $\alpha : f(a) \to f(b)$
est un morphisme de la catégorie fibre $\mathcal{X}_T$.

\medskip\noindent
Démonstration de (2). L'équivalence s'obtient en appliquant successivement
Catégories, Lemmes \ref{categories-lemma-associativity-2-fibre-product} et
\ref{categories-lemma-2-fibre-product-erase-factor}.
Identifions
$\mathcal{U} \times_\mathcal{X} \mathcal{U} \times_\mathcal{X} \mathcal{U}$
à
$(\mathcal{U} \times_\mathcal{X} \mathcal{U})
\times_\mathcal{X} \mathcal{U}$.
Si $T$ est un schéma au-dessus de $S$, alors, sur les catégories fibres au-dessus de $T$,
cette équivalence envoie l'objet
$((a, b, \alpha), c, \beta)$ du membre de gauche
sur l'objet $((a, b, \alpha), (b, c, \beta))$ du membre de droite.

\medskip\noindent
Démonstration de (3). Le $1$-morphisme $\text{pr}_{02}$ est construit dans la démonstration de
Catégories, Lemme \ref{categories-lemma-triple-2-fibre-product-pr02}.
À l'aide de la description ci-dessus des objets de la catégorie fibre,
on voit que $((a, b, \alpha), (b, c, \beta))$
est envoyé sur $(a, c, \beta \circ \alpha)$.

\medskip\noindent
Malheureusement, cela n'est {\it pas compatible} avec nos conventions sur les
groupoïdes, selon lesquelles on a toujours $j = (t, s) : R \to U$, et l'on « considère »
un point à valeurs dans $T$, noté $r$, de $R$ comme un morphisme $r : s(r) \to t(r)$.
Cela n'affecte toutefois pas la démonstration de (4), puisque l'opposé d'un
groupoïde est un groupoïde. Mais, dans la démonstration de (5), cela explique
les inverses dans la formule affichée ci-dessous.

\medskip\noindent
Démonstration de (4). Rappelons que le faisceau $U$ est isomorphe au faisceau
$T \mapsto \Ob(\mathcal{U}_T)/\!\cong$, et
de même pour $R$, voir le
Lemme \ref{algebraic-lemma-characterize-representable-by-space}.
Il résulte du
Lemme de Catégories
\ref{categories-lemma-category-fibred-setoids-presheaves-products}
que cette description est compatible aux $2$-produits fibrés ;
on obtient donc une identification analogue de
$\mathcal{R} \times_{\text{pr}_1, \mathcal{U}, \text{pr}_0} \mathcal{R}$
et de $R \times_{s, U, t} R$.
Les morphismes $t, s : R \to U$ et $c : R \times_{s, U, t} R \to R$
s'obtiennent à partir de l'égalité générale (\ref{algebraic-equation-morphisms-spaces}).
Explicitement, ces applications sont les transformations de foncteurs qui proviennent
de l'action de $\text{pr}_0$, $\text{pr}_0$, $\text{pr}_{02}$
sur les classes d'isomorphie d'objets des catégories fibres.
Ainsi, pour montrer que l'on obtient un groupoïde en espaces
algébriques, il suffit de montrer que, pour tout schéma $T$ au-dessus de $S$,
la structure
$$
(\Ob(\mathcal{U}_T)/\!\cong,
\Ob(\mathcal{R}_T)/\!\cong,
\text{pr}_1, \text{pr}_0, \text{pr}_{02})
$$
est un groupoïde, ce qui est clair d'après notre description ci-dessus des objets de
$\mathcal{R}_T$.

\medskip\noindent
Démonstration de (5). Nous appliquerons finalement
Groupoïdes in Espaces,
Lemme \ref{spaces-groupoids-lemma-quotient-stack-2-coequalizer}
pour obtenir le foncteur $[U/R] \to \mathcal{X}$.
Considérons le $1$-morphisme $f : \mathcal{U} \to \mathcal{X}$.
Nous avons une $2$-flèche $\tau : f \circ \text{pr}_1 \to f \circ \text{pr}_0$
par définition de $\mathcal{R}$ comme $2$-produit fibré.
En effet, sur un objet $(a, b, \alpha)$ de $\mathcal{R}$ au-dessus de $T$, c'est
l'application $\alpha^{-1} : b \to a$. Nous affirmons que
$$
\tau \circ \text{id}_{\text{pr}_{02}} =
(\tau \star \text{id}_{\text{pr}_0})
\circ
(\tau \star \text{id}_{\text{pr}_1}).
$$
Cette identité dit qu'étant donné un objet
$((a, b, \alpha), (b, c, \beta))$ de
$\mathcal{R} \times_{\text{pr}_1, \mathcal{U}, \text{pr}_0} \mathcal{R}$
au-dessus de $T$, la composée de
$$
\xymatrix{
c \ar[r]^{\beta^{-1}} & b \ar[r]^{\alpha^{-1}} & a
}
$$
est la même que la flèche $(\beta \circ \alpha)^{-1} : a \to c$. C'est
manifestement vrai, d'où l'affirmation. On voit ainsi que toutes les
hypothèses du
Lemme de Groupoïdes in Espaces
\ref{spaces-groupoids-lemma-quotient-stack-2-coequalizer}
sont satisfaites pour la structure
$(\mathcal{U}, \mathcal{R}, \text{pr}_0, \text{pr}_1, \text{pr}_{02})$
ainsi que le $1$-morphisme $f$ et le $2$-morphisme $\tau$.
Cependant, pour appliquer le lemme, il faut démontrer cela
pour la structure $(\mathcal{S}_U, \mathcal{S}_R, s, t, c)$
munie de morphismes appropriés.

\medskip\noindent
Il devrait exister un argument général de théorie abstraite des catégories
qui transfère ces données de l'une à l'autre, mais il semble
assez long. Nous employons plutôt l'astuce suivante.
Choisissons un quasi-inverse $j^{-1} : \mathcal{S}_U \to \mathcal{U}$
de l'équivalence canonique $j : \mathcal{U} \to \mathcal{S}_U$ qui provient
de $U(T) = \Ob(\mathcal{U}_T)/\!\!\cong$.
Cela signifie simplement que, pour tout schéma $T/S$ et tout
objet $a \in \mathcal{U}_T$, nous avons choisi un
élément particulier de sa classe d'isomorphie, à savoir $j^{-1}(j(a))$.
Grâce à $j^{-1}$, on peut donc considérer $\mathcal{S}_U$
comme une sous-catégorie de $\mathcal{U}$. Cette sous-catégorie une fois choisie,
on peut considérer les objets $(a, b, \alpha)$ de $\mathcal{R}_T$
tels que $a, b$ soient des objets de $(\mathcal{S}_U)_T$, c'est-à-dire tels
que $j^{-1}(j(a)) = a$ et $j^{-1}(j(b)) = b$. Il est alors clair qu'ils
forment une sous-catégorie de $\mathcal{R}$ qui s'envoie isomorphiquement
sur $\mathcal{S}_R$ par l'équivalence canonique
$\mathcal{R} \to \mathcal{S}_R$. De plus, ceci est manifestement compatible
avec la formation du $2$-produit fibré
$\mathcal{R} \times_{\text{pr}_1, \mathcal{U}, \text{pr}_0} \mathcal{R}$.
On voit donc que l'on peut simplement restreindre
$f$ à $\mathcal{S}_U$ et restreindre $\tau$ en une transformation
entre foncteurs $\mathcal{S}_R \to \mathcal{X}$. Il est dès lors clair que
l'égalité affichée du
Lemme de Groupoïdes in Espaces
\ref{spaces-groupoids-lemma-quotient-stack-2-coequalizer}
est satisfaite, puisqu'elle l'est même comme égalité de transformations de foncteurs
$\mathcal{R} \times_{\text{pr}_1, \mathcal{U}, \text{pr}_0} \mathcal{R}
\to \mathcal{X}$ avant restriction à la sous-catégorie
$\mathcal{S}_{R \times_{s, U, t} R}$.

\medskip\noindent
Ceci démontre que le
Lemme de Groupoïdes in Espaces
\ref{spaces-groupoids-lemma-quotient-stack-2-coequalizer}
s'applique et fournit le morphisme de champs recherché
$f_{can} : [U/R] \to \mathcal{X}$. Précisons brièvement comment
$f_{can}$ est défini dans ce cas particulier.
Sur un objet $a$ de $\mathcal{S}_U$ au-dessus de $T$,
on a $f_{can}(a) = f(a)$, où l'on considère
$\mathcal{S}_U \subset \mathcal{U}$ par le plongement choisi ci-dessus.
Si $a, b$ sont des objets de $\mathcal{S}_U$ au-dessus de $T$, alors un morphisme
$\varphi : a \to b$ dans $[U/R]$ est, par définition, un objet de la
forme $\varphi = (b, a, \alpha)$ de $\mathcal{R}$ au-dessus de $T$. (Noter
l'interversion.) Et la règle de la démonstration du
Lemme de Groupoïdes in Espaces
\ref{spaces-groupoids-lemma-quotient-stack-2-coequalizer}
est
\begin{equation}
\label{algebraic-equation-on-morphisms}
f_{can}(\varphi) = \Big(f(a) \xrightarrow{\alpha^{-1}} f(b)\Big).
\end{equation}
Démonstration de (6). $[U/R]$ et $\mathcal{X}$ sont tous deux des champs.
Ainsi, étant donnés un schéma $T/S$ et des objets $a, b$ de $[U/R]$
au-dessus de $T$, on obtient une transformation de faisceaux fppf
$$
\mathit{Isom}(a, b) \longrightarrow \mathit{Isom}(f_{can}(a), f_{can}(b))
$$
sur $(\Sch/T)_{fppf}$. Il faut montrer qu'il s'agit d'un
isomorphisme. On peut travailler localement pour la topologie fppf sur $T$, et donc supposer que
$a, b$ proviennent de morphismes $a, b : T \to U$. Par le plongement
$\mathcal{S}_U \subset \mathcal{U}$ ci-dessus, on peut aussi considérer $a, b$ comme
des objets de $\mathcal{U}$ au-dessus de $T$. Dans
Groupoïdes in Espaces,
Lemme \ref{spaces-groupoids-lemma-quotient-stack-morphisms},
nous avons vu que le faisceau de gauche est représenté par l'espace algébrique
$$
R \times_{(t, s), U \times_S U, (b, a)} T
$$
au-dessus de $T$. D'autre part, le membre de droite est, par le
Lemme \ref{stacks-lemma-isom-as-2-fibre-product} de Champs,
égal au faisceau associé au champ en setoïdes suivant :
$$
\mathcal{X}
\times_{\mathcal{X} \times \mathcal{X}, (f \circ b, f \circ a)} T =
\mathcal{X}
\times_{\mathcal{X} \times \mathcal{X}, (f, f)}
(\mathcal{U} \times \mathcal{U})
\times_{\mathcal{U} \times \mathcal{U}, (b, a)} T =
\mathcal{R}
\times_{(\text{pr}_0, \text{pr}_1), \mathcal{U} \times \mathcal{U}, (b, a)} T
$$
qui est représentable par le produit fibré affiché ci-dessus.
À ce stade, nous avons montré que les deux faisceaux $\mathit{Isom}$
sont isomorphes. Notre $1$-morphisme $f_{can} : [U/R] \to \mathcal{X}$ induit
cet isomorphisme sur les faisceaux $\mathit{Isom}$ d'après
l'Équation (\ref{algebraic-equation-on-morphisms}).
\end{proof}

\noindent
On peut utiliser le lemme très abstrait précédent pour produire des
présentations.

\begin{lemma}
\label{algebraic-lemma-stack-presentation}
Soit $S$ un schéma appartenant à $\Sch_{fppf}$.
Soit $\mathcal{X}$ un champ algébrique au-dessus de $S$.
Soit $U$ un espace algébrique au-dessus de $S$.
Soit $f : \mathcal{S}_U \to \mathcal{X}$ un morphisme lisse surjectif.
Soient $(U, R, s, t, c)$ le groupoïde en espaces algébriques
et $f_{can} : [U/R] \to \mathcal{X}$ le résultat de l'application du
Lemme \ref{algebraic-lemma-map-space-into-stack}
à $U$ et à $f$. Alors
\begin{enumerate}
\item les morphismes $s$, $t$ sont lisses, et
\item le $1$-morphisme $f_{can} : [U/R] \to \mathcal{X}$
est une équivalence.
\end{enumerate}
\end{lemma}

\begin{proof}
Les morphismes $s, t$ sont lisses par les
Lemmes \ref{algebraic-lemma-property-morphism-equivalent} et
\ref{algebraic-lemma-map-presheaves-representable-by-spaces-transformation-property}.
Comme le $1$-morphisme $f$ est lisse et
surjectif, il est clair que, pour tout schéma $T$ et tout objet
$a \in \Ob(\mathcal{X}_T)$, il existe un morphisme lisse et surjectif
$T' \to T$ tel que $a|_T'$ provienne d'un objet de
$[U/R]_{T'}$. Puisque $f_{can} : [U/R] \to \mathcal{X}$
est pleinement fidèle, on en déduit que
$[U/R] \to \mathcal{X}$ est essentiellement surjectif, car
les données de descente sur les objets sont effectives des deux côtés, voir
Champs, Lemme \ref{stacks-lemma-characterize-essentially-surjective-when-ff}.
\end{proof}

\begin{remark}
\label{algebraic-remark-flat-fp-presentation}
Si l'on suppose seulement que le morphisme $f : \mathcal{S}_U \to \mathcal{X}$ du
Lemme \ref{algebraic-lemma-stack-presentation}
est surjectif, plat et localement de présentation finie, alors
$f_{can} : [U/R] \to \mathcal{X}$ sera encore une
équivalence. Dans ce cas, les morphismes $s$, $t$ seront plats et
localement de présentation finie, mais ne seront bien sûr pas lisses en général.
\end{remark}

\noindent
Le Lemme \ref{algebraic-lemma-stack-presentation}
suggère les définitions suivantes.

\begin{definition}
\label{algebraic-definition-smooth-groupoid}
Soit $S$ un schéma. Soit $B$ un espace algébrique au-dessus de $S$.
Soit $(U, R, s, t, c)$ un groupoïde en espaces algébriques au-dessus de $B$.
On dit que $(U, R, s, t, c)$ est un {\it groupoïde lisse}\footnote{Cette terminologie
peut prêter quelque peu à confusion : elle n'implique pas que $[U/R]$ soit lisse
au-dessus de quoi que ce soit.}
si $s, t : R \to U$ sont des morphismes lisses d'espaces algébriques.
\end{definition}

\begin{definition}
\label{algebraic-definition-presentation}
Soit $\mathcal{X}$ un champ algébrique au-dessus de $S$.
Une {\it présentation} de $\mathcal{X}$ est la donnée d'un groupoïde lisse
$(U, R, s, t, c)$ en espaces algébriques au-dessus de $S$, et d'une
équivalence $f : [U/R] \to \mathcal{X}$.
\end{definition}

\noindent
Nous avons vu ci-dessus que tout champ algébrique admet une présentation.
Notre tâche suivante consiste à montrer que tout groupoïde lisse en espaces
algébriques au-dessus de $S$ donne lieu à un champ algébrique.


\section{Le champ algébrique associé à un groupoïde lisse}
\label{algebraic-section-smooth-groupoid-gives-algebraic-stack}

\noindent
Dans cette section, nous partons d'un groupoïde lisse en espaces algébriques
et montrons que le champ quotient associé est un champ algébrique.

\begin{lemma}
\label{algebraic-lemma-diagonal-quotient-stack}
Soit $S$ un schéma appartenant à $\Sch_{fppf}$.
Soit $(U, R, s, t, c)$ un groupoïde en espaces algébriques au-dessus de $S$.
Alors la diagonale de $[U/R]$ est représentable par des espaces algébriques.
\end{lemma}

\begin{proof}
Il suffit de montrer que les faisceaux $\mathit{Isom}$ sont des espaces
algébriques, voir le
Lemme \ref{algebraic-lemma-representable-diagonal}.
Cela résulte du
Lemme \ref{bootstrap-lemma-quotient-stack-isom} de Amorçage.
\end{proof}

\begin{lemma}
\label{algebraic-lemma-smooth-quotient-smooth-presentation}
Soit $S$ un schéma appartenant à $\Sch_{fppf}$.
Soit $(U, R, s, t, c)$ un groupoïde lisse en espaces algébriques au-dessus de $S$.
Alors le morphisme $\mathcal{S}_U \to [U/R]$ est lisse et surjectif.
\end{lemma}

\begin{proof}
Soit $T$ un schéma et soit $x : (\Sch/T)_{fppf} \to [U/R]$
un $1$-morphisme. Il faut montrer que la projection
$$
\mathcal{S}_U \times_{[U/R]} (\Sch/T)_{fppf}
\longrightarrow
(\Sch/T)_{fppf}
$$
est surjective et lisse. Nous savons déjà que le membre de gauche
est représentable par un espace algébrique $F$, voir les
Lemmes \ref{algebraic-lemma-diagonal-quotient-stack} et
\ref{algebraic-lemma-representable-diagonal}.
Il faut donc montrer que le morphisme correspondant $F \to T$
d'espaces algébriques est surjectif et lisse.
Puisqu'il s'agit de propriétés de morphismes d'espaces algébriques
qui sont locales pour la topologie fppf sur le but, nous
pouvons les vérifier localement pour la topologie fppf sur $T$. Par construction, il existe
un recouvrement fppf $\{T_i \to T\}$ de $T$ tel que
$x|_{(\Sch/T_i)_{fppf}}$ provienne d'un morphisme
$x_i : T_i \to U$. (Remarquons que $F \times_T T_i$ représente le
$2$-produit fibré $\mathcal{S}_U \times_{[U/R]} (\Sch/T_i)_{fppf}$,
de sorte que tout est compatible au changement de base par $T_i \to T$.)
On peut donc supposer que $x$ provient de $x : T \to U$.
Dans ce cas, on voit que
$$
\mathcal{S}_U \times_{[U/R]} (\Sch/T)_{fppf}
=
(\mathcal{S}_U \times_{[U/R]} \mathcal{S}_U)
\times_{\mathcal{S}_U} (\Sch/T)_{fppf}
=
\mathcal{S}_R \times_{\mathcal{S}_U} (\Sch/T)_{fppf}
$$
La première égalité résulte du
Lemme \ref{categories-lemma-2-fibre-product-erase-factor} de Catégories,
et la seconde du
Lemme de Groupoïdes in Espaces
\ref{spaces-groupoids-lemma-quotient-stack-2-cartesian}.
Manifestement, le dernier $2$-produit fibré est représenté par l'espace
algébrique $F = R \times_{s, U, x} T$, et la projection
$R \times_{s, U, x} T \to T$ est lisse comme changement de base du
morphisme lisse d'espaces algébriques $s : R \to U$.
Elle est aussi surjective, car $s$ possède une section (à savoir l'identité
$e : U \to R$ du groupoïde).
Cela démontre le lemme.
\end{proof}

\noindent
Voici le résultat principal de cette section.

\begin{theorem}
\label{algebraic-theorem-smooth-groupoid-gives-algebraic-stack}
Soit $S$ un schéma appartenant à $\Sch_{fppf}$.
Soit $(U, R, s, t, c)$ un groupoïde lisse en espaces algébriques au-dessus de $S$.
Alors le champ quotient $[U/R]$ est un champ algébrique au-dessus de $S$.
\end{theorem}

\begin{proof}
Vérifions les trois conditions de la
Définition \ref{algebraic-definition-algebraic-stack}.
Par construction, $[U/R]$ est un champ en groupoïdes,
ce qui est la première condition.

\medskip\noindent
La deuxième condition résulte du résultat plus fort qu'est le
Lemme \ref{algebraic-lemma-diagonal-quotient-stack}.

\medskip\noindent
Enfin, il faut montrer qu'il existe un schéma $W$ au-dessus de $S$
et un $1$-morphisme lisse surjectif
$(\Sch/W)_{fppf} \longrightarrow \mathcal{X}$.
Choisissons d'abord $W \in \Ob((\Sch/S)_{fppf})$ et un
morphisme \'etale surjectif $W \to U$. Remarquons que cela
donne un morphisme \'etale surjectif $\mathcal{S}_W \to \mathcal{S}_U$
de catégories fibrées en ensembles, voir le
Lemme
\ref{algebraic-lemma-map-presheaves-representable-by-spaces-transformation-property}.
Bien entendu, $\mathcal{S}_W \to \mathcal{S}_U$ est alors aussi surjectif et
lisse, voir le
Lemme \ref{algebraic-lemma-representable-transformations-property-implication}.
Ainsi, $\mathcal{S}_W \to \mathcal{S}_U \to [U/R]$ est surjectif
et lisse par combinaison des
Lemmes \ref{algebraic-lemma-smooth-quotient-smooth-presentation} et
\ref{algebraic-lemma-composition-representable-transformations-property}.
\end{proof}





\section{Changement de grand site}
\label{algebraic-section-change-big-site}

\noindent
Dans cette section, nous examinons brièvement ce qui se passe lorsque l'on change de grand site.
En substance, on peut toujours agrandir à volonté le grand site ; on peut donc
supposer que tout ensemble de schémas que l'on souhaite considérer appartient au grand
site fppf au-dessus duquel on considère notre espace algébrique.
Nous invitons le lecteur à sauter cette section.

\medskip\noindent
Les images réciproques de champs sont définies dans
Champs, section \ref{stacks-section-inverse-image}.

\begin{lemma}
\label{algebraic-lemma-change-big-site}
Supposons donnés de grands sites $\Sch_{fppf}$ et $\Sch'_{fppf}$.
Supposons que $\Sch_{fppf}$ soit contenu dans $\Sch'_{fppf}$,
voir Topologies, section \ref{topologies-section-change-alpha}.
Soit $S$ un objet de $\Sch_{fppf}$.
Soit $f : (\Sch'/S)_{fppf} \to (\Sch/S)_{fppf}$ le morphisme
de sites correspondant au foncteur d'inclusion
$u : (\Sch/S)_{fppf} \to (\Sch'/S)_{fppf}$.
Soit $\mathcal{X}$ un champ en groupoïdes au-dessus de $(\Sch/S)_{fppf}$.
\begin{enumerate}
\item si $\mathcal{X}$ est représentable par un
$X \in \Ob((\Sch/S)_{fppf})$, alors
$f^{-1}\mathcal{X}$ est aussi représentable ; il est même représenté par le
même schéma $X$, considéré maintenant comme objet de $(\Sch'/S)_{fppf}$,
\item si $\mathcal{X}$ est représentable par
$F \in \Sh((\Sch/S)_{fppf})$, qui est
un espace algébrique, alors $f^{-1}\mathcal{X}$ est représentable
par l'espace algébrique $f^{-1}F$,
\item si $\mathcal{X}$ est un champ algébrique, alors $f^{-1}\mathcal{X}$
est un champ algébrique, et
\item si $\mathcal{X}$ est un champ de Deligne-Mumford, alors $f^{-1}\mathcal{X}$
est aussi un champ de Deligne-Mumford.
\end{enumerate}
\end{lemma}

\begin{proof}
Démontrons (3). Par le
Lemme \ref{algebraic-lemma-stack-presentation},
on peut écrire $\mathcal{X} = [U/R]$ pour un groupoïde lisse
en espaces algébriques $(U, R, s, t, c)$. Par le
Lemme de Groupoïdes in Espaces
\ref{spaces-groupoids-lemma-quotient-stack-change-big-site},
on voit que $f^{-1}[U/R] = [f^{-1}U/f^{-1}R]$.
Bien entendu, $(f^{-1}U, f^{-1}R, f^{-1}s, f^{-1}t, f^{-1}c)$
est aussi un groupoïde lisse en espaces algébriques. Cela démontre (3).

\medskip\noindent
Les autres cas (1), (2), (4) signifient chacun que $\mathcal{X}$ possède
une présentation $[U/R]$ d'un type particulier, et se traduisent donc par le
même type de présentation de $f^{-1}\mathcal{X} = [f^{-1}U/f^{-1}R]$.
Le lemme est ainsi démontré.
\end{proof}

\noindent
Il n'est pas vrai, en général, que la restriction d'un espace algébrique
sur le grand site le plus gros soit un espace algébrique sur le plus petit (pour de simples
raisons de cardinalité). On ne peut donc utiliser un lemme simple de ce
genre que pour agrandir la catégorie de base, jamais pour la rétrécir.

\begin{lemma}
\label{algebraic-lemma-fully-faithful}
Supposons que $\Sch_{fppf}$ soit contenu dans $\Sch'_{fppf}$.
Soit $S$ un objet de $\Sch_{fppf}$. Notons
$\textit{Algebraic-Champs}/S$ la $2$-catégorie des champs algébriques au-dessus de $S$
définie à l'aide de $\Sch_{fppf}$. De même, notons
$\textit{Algebraic-Champs}'/S$ la $2$-catégorie des champs algébriques au-dessus de $S$
définie à l'aide de $\Sch'_{fppf}$. La règle
$\mathcal{X} \mapsto f^{-1}\mathcal{X}$ du
Lemme \ref{algebraic-lemma-change-big-site}
définit un foncteur de $2$-catégories
$$
\textit{Algebraic-Champs}/S \longrightarrow \textit{Algebraic-Champs}'/S
$$
qui définit des équivalences de catégories de morphismes
$$
\Mor_{\textit{Algebraic-Champs}/S}(\mathcal{X}, \mathcal{Y})
\longrightarrow
\Mor_{\textit{Algebraic-Champs}'/S}(f^{-1}\mathcal{X}, f^{-1}\mathcal{Y})
$$
pour tous objets $\mathcal{X}, \mathcal{Y}$ de
$\textit{Algebraic-Champs}/S$. Un objet
$\mathcal{X}'$ de $\textit{Algebraic-Champs}'/S$
est équivalent à $f^{-1}\mathcal{X}$ pour un certain
$\mathcal{X}$ de $\textit{Algebraic-Champs}/S$
si et seulement s'il possède une présentation $\mathcal{X} = [U'/R']$
où $U', R'$ sont isomorphes à $f^{-1}U$, $f^{-1}R$ pour certains
$U, R \in \textit{Espaces}/S$.
\end{lemma}

\begin{proof}
L'assertion concernant les catégories de morphismes résulte de l'énoncé plus général
Champs, Lemme \ref{stacks-lemma-bigger-site}.
La caractérisation de l'« image essentielle » découle de la description
de $f^{-1}$ dans la démonstration du
Lemme \ref{algebraic-lemma-change-big-site}.
\end{proof}


\section{Changement de schéma de base}
\label{algebraic-section-change-base-scheme}

\noindent
Dans cette section, nous examinons brièvement ce qui se passe lorsque l'on change de schéma de base.
En substance, étant donné un morphisme $S \to S'$ de schémas de base, tout champ
algébrique au-dessus de $S$ peut être considéré comme un champ algébrique au-dessus de $S'$.

\begin{lemma}
\label{algebraic-lemma-category-of-spaces-over-smaller-base-scheme}
Soit $\Sch_{fppf}$ un grand site fppf.
Soit $S \to S'$ un morphisme de ce site.
Les constructions A et B de
Champs, section \ref{stacks-section-localize}
ci-dessus donnent des isomorphismes de $2$-catégories
$$
\left\{
\begin{matrix}
2\text{-catégorie des champs}\\
\text{algébriques }\mathcal{X}\text{ au-dessus de }S
\end{matrix}
\right\}
\leftrightarrow
\left\{
\begin{matrix}
2\text{-catégorie des paires }(\mathcal{X}', f)\text{ formées d'un}\\
\text{champ algébrique }\mathcal{X}'\text{ au-dessus de }S'\text{ et d'un morphisme}\\
f : \mathcal{X}' \to (\Sch/S)_{fppf}\text{ de champs algébriques au-dessus de }S'
\end{matrix}
\right\}
$$
\end{lemma}

\begin{proof}
L'énoncé a bien un sens, car le foncteur
$j : (\Sch/S)_{fppf} \to (\Sch/S')_{fppf}$
est le foncteur de localisation associé à l'objet $S/S'$
de $(\Sch/S')_{fppf}$. Par le
Lemme \ref{stacks-lemma-localize-stacks} de Champs,
il suffit de montrer que les constructions A et B
préservent les sous-catégories de champs algébriques.
Par exemple, si $\mathcal{X} = [U/R]$, alors la construction A
appliquée à $\mathcal{X}$ produit simplement
$\mathcal{X}' = \mathcal{X}$. Réciproquement, si $\mathcal{X}' = [U'/R']$,
le morphisme $p$ induit des morphismes d'espaces algébriques
$U' \to S$ et $R' \to S$, et alors $\mathcal{X} = [U'/R']$,
mais considéré maintenant comme un champ au-dessus de $S$. Le lemme est donc clair.
\end{proof}

\begin{definition}
\label{algebraic-definition-viewed-as}
Soit $\Sch_{fppf}$ un grand site fppf.
Soit $S \to S'$ un morphisme de ce site.
Si $p : \mathcal{X} \to (\Sch/S)_{fppf}$
est un champ algébrique au-dessus de $S$, alors
$\mathcal{X}$ {\it considéré comme champ algébrique au-dessus de $S'$}
est le champ algébrique
$$
\mathcal{X} \longrightarrow (\Sch/S')_{fppf}
$$
obtenu en appliquant la construction A du
Lemme \ref{algebraic-lemma-category-of-spaces-over-smaller-base-scheme}
à $\mathcal{X}$.
\end{definition}

\noindent
Réciproquement, que se passe-t-il si l'on part d'un champ algébrique $\mathcal{X}'$
au-dessus de $S'$ et que l'on souhaite obtenir un champ algébrique au-dessus de $S$ ?
On considère alors le $2$-produit fibré
$$
\mathcal{X}'_S
=
(\Sch/S)_{fppf} \times_{(\Sch/S')_{fppf}} \mathcal{X}'
$$
qui est un champ algébrique au-dessus de $S'$ d'après le
Lemme \ref{algebraic-lemma-2-fibre-product}.
De plus, il est muni d'un $1$-morphisme naturel
$p : \mathcal{X}'_S \to (\Sch/S)_{fppf}$ et, par le
Lemme \ref{algebraic-lemma-category-of-spaces-over-smaller-base-scheme},
il correspond donc canoniquement à un champ algébrique au-dessus de $S$.

\begin{definition}
\label{algebraic-definition-change-of-base}
Soit $\Sch_{fppf}$ un grand site fppf.
Soit $S \to S'$ un morphisme de ce site.
Soit $\mathcal{X}'$ un champ algébrique au-dessus de $S'$.
Le {\it changement de base de $\mathcal{X}'$} est le
champ algébrique $\mathcal{X}'_S$ au-dessus de $S$ décrit ci-dessus.
\end{definition}








% OK, start here.
%

\chapter{Exemples de champs}



\phantomsection
\label{examples-stacks-section-phantom}


\section{Introduction}
\label{examples-stacks-section-introduction}

\noindent
Nous étudions ici des exemples de champs en géométrie algébrique.
Certains sont des champs algébriques, d'autres ne le sont pas.
Nous déterminerons lesquels sont des champs algébriques dans un chapitre ultérieur.
Dans le présent chapitre, nous nous préoccupons donc principalement des conditions
de descente. Voir par exemple \cite{Vis2}.

\medskip\noindent
Certaines notations, conventions et certains termes de ce chapitre sont peu naturels
et pourront sembler inversés au lecteur plus expérimenté. Cela est intentionnel.
On trouvera une explication dans Quot, section \ref{quot-section-conventions}.




\section{Notations}
\label{examples-stacks-section-notation}

\noindent
Dans ce chapitre, nous fixons un grand site fppf convenable $\Sch_{fppf}$
comme dans Topologies, définition \ref{topologies-definition-big-fppf-site}.
Ainsi, sauf indication explicite contraire, tous les schémas seront des objets
de $\Sch_{fppf}$.
Nous travaillerons toujours relativement à une base $S$ contenue dans $\Sch_{fppf}$.
Nous travaillerons alors avec le grand site fppf $(\Sch/S)_{fppf}$,
voir Topologies, définition \ref{topologies-definition-big-small-fppf}.
On retrouve le cas absolu en prenant
$S = \Spec(\mathbf{Z})$.




\section{Exemples de champs}
\label{examples-stacks-section-examples-stacks}

\noindent
Nous donnons d'abord quelques exemples importants de champs sur
$(\Sch/S)_{fppf}$.




\section{Faisceaux quasi-cohérents}
\label{examples-stacks-section-stack-of-quasi-coherent-sheaves}

\noindent
Nous définissons une catégorie $\QCohstack$ comme suit :
\begin{enumerate}
\item Un objet de $\QCohstack$ est un couple $(X, \mathcal{F})$,
où $X/S$ est un objet de $(\Sch/S)_{fppf}$ et $\mathcal{F}$
est un $\mathcal{O}_X$-module quasi-cohérent ;
\item un morphisme $(f, \varphi) : (Y, \mathcal{G}) \to (X, \mathcal{F})$
est un couple formé d'un morphisme $f : Y \to X$ de schémas sur $S$
et d'un $f$-morphisme (voir
Faisceaux, section \ref{sheaves-section-ringed-spaces-functoriality-modules})
$\varphi : \mathcal{F} \to \mathcal{G}$ ;
\item le composé des morphismes
$$
(Z, \mathcal{H}) \xrightarrow{(g, \psi)}
(Y, \mathcal{G}) \xrightarrow{(f, \phi)} (X, \mathcal{F})
$$
est $(f \circ g, \psi \circ \phi)$, où $\psi \circ \phi$ est
le composé des $f$-morphismes.
\end{enumerate}
Ainsi, $\QCohstack$ est une catégorie et
$$
p : \QCohstack \to (\Sch/S)_{fppf},
\quad
(X, \mathcal{F}) \mapsto X
$$
est un foncteur. Remarquons que la catégorie fibre de $\QCohstack$ au-dessus
d'un schéma $X$ est l'opposée de la catégorie $\QCoh(\mathcal{O}_X)$
des $\mathcal{O}_X$-modules quasi-cohérents.
Pour un usage ultérieur, remarquons qu'étant donnés
$(X, \mathcal{F}), (Y, \mathcal{G}) \in \Ob(\QCohstack)$,
nous avons
\begin{equation}
\label{examples-stacks-equation-morphisms-qcoh}
\Mor_{\QCohstack}((Y, \mathcal{G}), (X, \mathcal{F}))
=
\coprod\nolimits_{f \in \Mor_S(Y, X)}
\Mor_{\QCoh(\mathcal{O}_Y)}(f^*\mathcal{F}, \mathcal{G}).
\end{equation}
Voir la discussion des $f$-morphismes de modules dans
Faisceaux, section \ref{sheaves-section-ringed-spaces-functoriality-modules}.

\medskip\noindent
La catégorie $\QCohstack$ n'est pas un champ sur $(\Sch/S)_{fppf}$,
car sa collection d'objets est une classe propre. Nous verrons néanmoins
qu'elle satisfait tous les axiomes d'un champ. Nous contournerons
la difficulté ensembliste dans la
section \ref{examples-stacks-section-stack-of-finitely-generated-quasi-coherent-sheaves}.

\begin{lemma}
\label{examples-stacks-lemma-quasi-coherent-strongly-cartesian}
Un morphisme $(f, \varphi) : (Y, \mathcal{G}) \to (X, \mathcal{F})$
de $\QCohstack$ est fortement cartésien si et seulement si le
morphisme $\varphi$ induit un isomorphisme $f^*\mathcal{F} \to \mathcal{G}$.
\end{lemma}

\begin{proof}
Soit $(X, \mathcal{F}) \in \Ob(\QCohstack)$.
Soit $f : Y \to X$ un morphisme de $(\Sch/S)_{fppf}$.
Il existe un $f$-morphisme canonique $c : \mathcal{F} \to f^*\mathcal{F}$
et nous obtenons donc un morphisme
$(f, c) : (Y, f^*\mathcal{F}) \to (X, \mathcal{F})$.
Nous affirmons que $(f, c)$ est fortement cartésien.
En effet, pour tout objet $(Z, \mathcal{H})$ de $\QCohstack$, nous avons
\begin{align*}
\Mor_{\QCohstack}((Z, \mathcal{H}), (Y, f^*\mathcal{F}))
& =
\coprod\nolimits_{g \in \Mor_S(Z, Y)}
\Mor_{\QCoh(\mathcal{O}_Z)}(g^*f^*\mathcal{F}, \mathcal{H}) \\
& =
\coprod\nolimits_{g \in \Mor_S(Z, Y)}
\Mor_{\QCoh(\mathcal{O}_Z)}((f \circ g)^*\mathcal{F}, \mathcal{H}) \\
& =
\Mor_{\QCohstack}((Z, \mathcal{H}), (X, \mathcal{F}))
\times_{\Mor_S(Z, X)} \Mor_S(Z, Y),
\end{align*}
où nous avons utilisé deux fois l'équation (\ref{examples-stacks-equation-morphisms-qcoh}).
Cela prouve que la condition de
Catégories, définition \ref{categories-definition-cartesian-over-C},
est satisfaite pour $(f, c)$, d'où notre assertion. Par
Catégories, lemme \ref{categories-lemma-composition-cartesian},
les isomorphismes sont fortement cartésiens et
les composés de morphismes fortement cartésiens le sont aussi,
ce qui prouve l'implication « si » du lemme. Réciproquement,
étant donnés $(X, \mathcal{F})$ et $f : Y \to X$, s'il existe un
morphisme fortement cartésien relevant $f$ et de but $(X, \mathcal{F})$,
alors il est nécessairement isomorphe à $(f, c)$ (voir la discussion qui suit
Catégories, définition \ref{categories-definition-cartesian-over-C}).
L'implication « seulement si » du lemme en résulte.
\end{proof}

\begin{lemma}
\label{examples-stacks-lemma-stack-of-quasi-coherent-sheaves}
Le foncteur $p : \QCohstack \to (\Sch/S)_{fppf}$
satisfait les conditions (1), (2) et (3) de
Champs, définition \ref{stacks-definition-stack}.
\end{lemma}

\begin{proof}
Il résulte clairement du
lemme \ref{examples-stacks-lemma-quasi-coherent-strongly-cartesian}
que $\QCohstack$ est une catégorie fibrée sur $(\Sch/S)_{fppf}$.
Pour un recouvrement $\mathcal{U} = \{X_i \to X\}_{i \in I}$ de
$(\Sch/S)_{fppf}$, le foncteur
$$
\QCoh(\mathcal{O}_X) \longrightarrow DD(\mathcal{U})
$$
est pleinement fidèle et essentiellement surjectif, voir
Descente, proposition \ref{descent-proposition-fpqc-descent-quasi-coherent}.
Par conséquent,
Champs, lemme \ref{stacks-lemma-stack-equivalences},
montre que $\QCohstack$ satisfait tous les
axiomes d'un champ.
\end{proof}





\section{Le champ des faisceaux quasi-cohérents de type fini}
\label{examples-stacks-section-stack-of-finitely-generated-quasi-coherent-sheaves}

\noindent
On peut en fait obtenir un champ de faisceaux quasi-cohérents
si l'on ne considère que les modules quasi-cohérents de type fini.
Notons
$$
p_{fg} : \QCohstack_{fg} \to (\Sch/S)_{fppf}
$$
la sous-catégorie pleine de $\QCohstack$ sur $(\Sch/S)_{fppf}$
formée des couples $(T, \mathcal{F})$ tels que $\mathcal{F}$
soit un $\mathcal{O}_T$-module quasi-cohérent de type fini.

\begin{lemma}
\label{examples-stacks-lemma-stack-of-finite-type-quasi-coherent-sheaves}
Le foncteur $p_{fg} : \QCohstack_{fg} \to (\Sch/S)_{fppf}$
satisfait les conditions (1), (2) et (3) de
Champs, définition \ref{stacks-definition-stack}.
\end{lemma}

\begin{proof}
Nous allons pour cela vérifier les hypothèses (1), (2) et (3) de
Champs, lemme \ref{stacks-lemma-substack}.
D'après le
lemme \ref{examples-stacks-lemma-quasi-coherent-strongly-cartesian},
un morphisme $(Y, \mathcal{G}) \to (X, \mathcal{F})$ est
fortement cartésien si et seulement s'il induit un isomorphisme
$f^*\mathcal{F} \to \mathcal{G}$. D'après
Modules, lemme \ref{modules-lemma-pullback-finite-type},
l'image inverse d'un $\mathcal{O}_X$-module de type fini est de type
fini. L'hypothèse (1) de
Champs, lemme \ref{stacks-lemma-substack},
est donc satisfaite. L'hypothèse (2) est immédiate.
Enfin, pour prouver l'hypothèse (3), il faut montrer ceci :
si $\mathcal{F}$ est un $\mathcal{O}_X$-module quasi-cohérent
et si $\{f_i : X_i \to X\}$ est un recouvrement fppf tel que chaque
$f_i^*\mathcal{F}$ soit de type fini, alors $\mathcal{F}$ est de
type fini. En restreignant $\mathcal{F}$ à
un ouvert affine de $X$, on se ramène à l'énoncé algébrique suivant :
soient $R \to S$ un homomorphisme d'anneaux de présentation finie et fidèlement plat,
et $M$ un $R$-module. Si $M \otimes_R S$ est un
$S$-module de type fini, alors $M$ est un $R$-module de type fini.
Une forme plus forte de ce fait algébrique se trouve dans
Algèbre, lemme \ref{algebra-lemma-descend-properties-modules}.
\end{proof}

\begin{lemma}
\label{examples-stacks-lemma-finite-type}
Soit $(X, \mathcal{O}_X)$ un espace annelé.
\begin{enumerate}
\item La catégorie des $\mathcal{O}_X$-modules de type fini possède un
ensemble de classes d'isomorphisme.
\item La catégorie des $\mathcal{O}_X$-modules quasi-cohérents
de type fini possède un ensemble de classes d'isomorphisme.
\end{enumerate}
\end{lemma}

\begin{proof}
La partie (2) résulte de la partie (1), car la catégorie de (2) est une sous-catégorie
pleine de celle de (1). Considérons un recouvrement ouvert quelconque
$\mathcal{U} : X = \bigcup_{i \in I} U_i$. Notons $j_i : U_i \to X$
les morphismes d'inclusion. Considérons une application $r : I \to \mathbf{N}$.
Si $\mathcal{F}$ est un $\mathcal{O}_X$-module dont la restriction à
$U_i$ est engendrée par au plus $r(i)$ sections de $\mathcal{F}(U_i)$,
alors $\mathcal{F}$ est un quotient du faisceau
$$
\mathcal{H}_{\mathcal{U}, r} =
\bigoplus\nolimits_{i \in I} j_{i, !}\mathcal{O}_{U_i}^{\oplus r(i)}.
$$
Par définition, si $\mathcal{F}$ est de type fini, il existe
un recouvrement ouvert $\mathcal{U}$ dont l'ensemble d'indices est $I = X$
et pour lequel cette condition est vérifiée. Il suffit donc de montrer qu'il existe
un ensemble de choix possibles pour $\mathcal{U}$ (c'est évident),
un ensemble de choix possibles pour $r : I \to \mathbf{N}$ (c'est évident) et,
pour chaque $\mathcal{U}$ et chaque $r$, un ensemble de modules quotients possibles
de $\mathcal{H}_{\mathcal{U}, r}$. Autrement dit, il suffit de montrer
que, pour un $\mathcal{O}_X$-module $\mathcal{H}$ donné, les quotients ont
au plus un ensemble de classes d'isomorphisme.
Cette dernière assertion devient évidente
si l'on considère les noyaux d'un morphisme quotient
$\mathcal{H} \to \mathcal{F}$
comme paramétrés par une partie de l'ensemble des parties de
$\prod_{U \subset X\text{ ouvert}} \mathcal{H}(U)$.
\end{proof}

\begin{lemma}
\label{examples-stacks-lemma-stack-fg-quasi-coherent}
Il existe une sous-catégorie
$\QCohstack_{fg, small} \subset \QCohstack_{fg}$
ayant les propriétés suivantes :
\begin{enumerate}
\item le foncteur d'inclusion
$\QCohstack_{fg, small} \to \QCohstack_{fg}$ est
pleinement fidèle et essentiellement surjectif ;
\item le foncteur
$p_{fg, small} : \QCohstack_{fg, small} \to (\Sch/S)_{fppf}$
fait de $\QCohstack_{fg, small}$ un champ sur $(\Sch/S)_{fppf}$.
\end{enumerate}
\end{lemma}

\begin{proof}
Nous avons vu dans les
lemmes \ref{examples-stacks-lemma-stack-of-finite-type-quasi-coherent-sheaves} et
\ref{examples-stacks-lemma-finite-type}
que $p_{fg} : \QCohstack_{fg} \to (\Sch/S)_{fppf}$
satisfait les conditions (1), (2) et (3) de
Champs, définition \ref{stacks-definition-stack},
ainsi que la condition supplémentaire (4) de
Champs, remarque \ref{stacks-remark-stack-make-small}.
La discussion de cette remarque fournit donc $\QCohstack_{fg, small}$.
\end{proof}

\noindent
Nous effectuerons souvent le remplacement
$$
\QCohstack_{fg} \leadsto \QCohstack_{fg, small}
$$
sans le signaler davantage et, par abus de notation, nous désignerons
simplement ce remplacement par $\QCohstack_{fg}$.

\begin{remark}
\label{examples-stacks-remark-higher-rank}
Toute la discussion de cette section reste valable
si l'on veut considérer les
faisceaux quasi-cohérents qui sont localement engendrés par au plus $\kappa$
sections, pour un cardinal infini $\kappa$, par exemple $\kappa = \aleph_0$.
\end{remark}





\section{Revêtements étales finis}
\label{examples-stacks-section-finite-etale}

\noindent
Nous définissons une catégorie $\textit{F\'Et}$ comme suit :
\begin{enumerate}
\item Un objet de $\textit{F\'Et}$ est un morphisme étale fini $Y \to X$
de schémas (d'après nos conventions, il s'agit d'un morphisme étale fini
dans $(\Sch/S)_{fppf}$) ;
\item un morphisme $(b, a) : (Y \to X) \to (Y' \to X')$ de $\textit{F\'Et}$
est un diagramme commutatif
$$
\xymatrix{
Y \ar[d] \ar[r]_b & Y' \ar[d] \\
X \ar[r]_a & X'
}
$$
dans la catégorie des schémas.
\end{enumerate}
Ainsi, $\textit{F\'Et}$ est une catégorie et
$$
p : \textit{F\'Et} \to (\Sch/S)_{fppf},
\quad
(Y \to X) \mapsto X
$$
est un foncteur. La catégorie fibre de $\textit{F\'Et}$ au-dessus
d'un schéma $X$ est précisément la catégorie $\textit{F\'Et}_X$ étudiée dans
Groupes fondamentaux, section \ref{pione-section-finite-etale}.

\begin{lemma}
\label{examples-stacks-lemma-finite-etale-stack}
Le foncteur
$$
p : \textit{F\'Et} \longrightarrow (\Sch/S)_{fppf}
$$
définit un champ sur $(\Sch/S)_{fppf}$.
\end{lemma}

\begin{proof}
La descente fppf des morphismes étales finis résulte de
Descente, lemmes \ref{descent-lemma-affine},
\ref{descent-lemma-descending-property-finite} et
\ref{descent-lemma-descending-property-etale}.
Nous omettons les détails.
\end{proof}






\section{Espaces algébriques}
\label{examples-stacks-section-stack-of-spaces}

\noindent
Nous définissons une catégorie $\Spacesstack$ comme suit :
\begin{enumerate}
\item Un objet de $\Spacesstack$ est un morphisme $X \to U$
d'espaces algébriques sur $S$, où $U$ est représentable par un objet de
$(\Sch/S)_{fppf}$ ;
\item un morphisme $(f, g) : (X \to U) \to (Y \to V)$
est un diagramme commutatif
$$
\xymatrix{
X \ar[d] \ar[r]_f & Y \ar[d] \\
U \ar[r]^g & V
}
$$
de morphismes d'espaces algébriques sur $S$.
\end{enumerate}
Ainsi, $\Spacesstack$ est une catégorie et
$$
p : \Spacesstack \to (\Sch/S)_{fppf},
\quad
(X \to U) \mapsto U
$$
est un foncteur. La catégorie fibre de $\Spacesstack$ au-dessus
d'un schéma $U$ est précisément la catégorie $\textit{Espaces}/U$ des
espaces algébriques sur $U$ (voir
Topologies sur les espaces, section \ref{spaces-topologies-section-procedure}).
Nous considérerons donc parfois un objet de $\Spacesstack$ comme un
couple $X/U$ formé d'un schéma $U$ et d'un espace algébrique $X$ sur $U$.
Pour un usage ultérieur, remarquons qu'étant donnés
$(X/U), (Y/V) \in \Ob(\Spacesstack)$,
nous avons
\begin{equation}
\label{examples-stacks-equation-morphisms-spaces}
\Mor_{\Spacesstack}(X/U, Y/V)
=
\coprod\nolimits_{g \in \Mor_S(U, V)}
\Mor_{\textit{Espaces}/U}(X, U \times_{g, V} Y).
\end{equation}
La catégorie $\Spacesstack$ est presque, mais pas tout à fait, un champ
sur $(\Sch/S)_{fppf}$. La difficulté est de nature ensembliste,
comme nous l'expliquerons ci-dessous.

\begin{lemma}
\label{examples-stacks-lemma-spaces-strongly-cartesian}
Un morphisme $(f, g) : X/U \to Y/V$
de $\Spacesstack$ est fortement cartésien si et seulement si le
morphisme $f$ induit un isomorphisme $X \to U \times_{g, V} Y$.
\end{lemma}

\begin{proof}
Soit $Y/V \in \Ob(\Spacesstack)$.
Soit $g : U \to V$ un morphisme de $(\Sch/S)_{fppf}$.
La projection $p : U \times_{g, V} Y \to Y$
donne un morphisme
$(p, g) : U \times_{g, V} Y/U \to Y/V$ de $\Spacesstack$.
Nous affirmons que $(p, g)$ est fortement cartésien.
En effet, pour tout objet $Z/W$ de $\Spacesstack$, nous avons
\begin{align*}
\Mor_{\Spacesstack}(Z/W, U \times_{g, V} Y/U)
& =
\coprod\nolimits_{h \in \Mor_S(W, U)}
\Mor_{\textit{Espaces}/W}(Z, W \times_{h, U} U \times_{g, V} Y) \\
& =
\coprod\nolimits_{h \in \Mor_S(W, U)}
\Mor_{\textit{Espaces}/W}(Z, W \times_{g \circ h, V} Y) \\
& =
\Mor_{\Spacesstack}(Z/W, Y/V)
\times_{\Mor_S(W, V)} \Mor_S(W, U),
\end{align*}
où nous avons utilisé deux fois l'équation (\ref{examples-stacks-equation-morphisms-spaces}).
Cela prouve que la condition de
Catégories, définition \ref{categories-definition-cartesian-over-C},
est satisfaite pour $(p, g)$, d'où notre assertion. Par
Catégories, lemme \ref{categories-lemma-composition-cartesian},
les isomorphismes sont fortement cartésiens et
les composés de morphismes fortement cartésiens le sont aussi,
ce qui prouve l'implication « si » du lemme. Réciproquement,
étant donnés $Y/V$ et $g : U \to V$, s'il existe un
morphisme fortement cartésien relevant $g$ et de but $Y/V$,
alors il est nécessairement isomorphe à $(p, g)$ (voir la discussion qui suit
Catégories, définition \ref{categories-definition-cartesian-over-C}).
L'implication « seulement si » du lemme en résulte.
\end{proof}

\begin{lemma}
\label{examples-stacks-lemma-pre-stack-of-spaces}
Le foncteur $p : \Spacesstack \to (\Sch/S)_{fppf}$
satisfait les conditions (1) et (2) de
Champs, définition \ref{stacks-definition-stack}.
\end{lemma}

\begin{proof}
Il résulte du
lemme \ref{examples-stacks-lemma-spaces-strongly-cartesian}
que $\Spacesstack$ est une catégorie fibrée sur $(\Sch/S)_{fppf}$,
ce qui prouve (1).
Supposons que $\{U_i \to U\}_{i \in I}$ soit un recouvrement de
$(\Sch/S)_{fppf}$. Soient $X, Y$ des espaces algébriques sur
$U$. Enfin, supposons donnés des morphismes $\varphi_i : X_{U_i} \to Y_{U_i}$
de $\textit{Espaces}/U_i$ tels que les restrictions de $\varphi_i$ et $\varphi_j$
soient les mêmes morphismes $X_{U_i \times_U U_j} \to Y_{U_i \times_U U_j}$
d'espaces algébriques sur $U_i \times_U U_j$.
Pour prouver (2), il faut montrer qu'il existe un unique morphisme
$\varphi  : X \to Y$ sur $U$ dont le changement de base à $U_i$ est
égal à $\varphi_i$. Comme un morphisme de $X$ vers $Y$ équivaut à
un morphisme de faisceaux, cela résulte directement de
Sites, lemme \ref{sites-lemma-glue-maps}.
\end{proof}

\begin{remark}
\label{examples-stacks-remark-stack-spaces}
Si l'on néglige les difficultés ensemblistes\footnote{La difficulté ne tient pas
à ce que $\Spacesstack$ soit une classe propre, puisque notre définition d'un
espace algébrique sur $S$ n'autorise qu'un ensemble de
classes d'isomorphisme d'espaces algébriques sur $S$. Elle tient plutôt à ce que
des réunions disjointes arbitraires d'espaces algébriques puissent devenir trop grandes
et sortir ainsi de l'« univers partiel » d'ensembles que nous avons choisi.},
$\Spacesstack$ satisfait aussi
la descente des objets et est donc un champ. Il faut en effet montrer qu'étant donnés
\begin{enumerate}
\item un recouvrement fppf $\{U_i \to U\}_{i \in I}$ ;
\item pour chaque $i \in I$, un espace algébrique $X_i/U_i$ ;
\item pour chaque $i, j \in I$, un isomorphisme
$\varphi_{ij} : X_i \times_U U_j \to U_i \times_U X_j$ d'espaces algébriques
sur $U_i \times_U U_j$, satisfaisant la condition de cocycle sur
$U_i \times_U U_j \times_U U_k$,
\end{enumerate}
il existe un espace algébrique $X/U$ et des isomorphismes
$X_{U_i} \cong X_i$ sur $U_i$ qui redonnent les isomorphismes $\varphi_{ij}$.
D'abord, d'après
Sites, lemme \ref{sites-lemma-glue-sheaves},
il existe sur $(\Sch/U)_{fppf}$ un faisceau $X$ qui redonne
les $X_i$ et les $\varphi_{ij}$. Puis, d'après
Amorçage, lemme \ref{bootstrap-lemma-locally-algebraic-space},
$X$ est un espace algébrique (si l'on néglige la condition ensembliste
de ce lemme).
Nous emploierons cet argument dans la section suivante pour montrer que,
si l'on ne considère que les espaces algébriques de type fini, on obtient
un champ.
\end{remark}





\section{Le champ des espaces algébriques de type fini}
\label{examples-stacks-section-stack-of-finite-type-spaces}


\noindent
On peut en fait obtenir un champ d'espaces
si l'on ne considère que les espaces de type fini.
Notons
$$
p_{ft} : \Spacesstack_{ft} \to (\Sch/S)_{fppf}
$$
la sous-catégorie pleine de $\Spacesstack$ sur $(\Sch/S)_{fppf}$
formée des couples $X/U$ tels que $X \to U$
soit un morphisme de type fini.

\begin{lemma}
\label{examples-stacks-lemma-stack-of-finite-type-spaces}
Le foncteur
$p_{ft} : \Spacesstack_{ft} \to (\Sch/S)_{fppf}$
satisfait les conditions (1), (2) et (3) de
Champs, définition \ref{stacks-definition-stack}.
\end{lemma}

\begin{proof}
Nous allons donner tous les détails, jusqu'à l'excès (ce qui risque
de masquer l'idée de la démonstration).

\medskip\noindent
Nous avons vu dans le
lemme \ref{examples-stacks-lemma-spaces-strongly-cartesian}
qu'un morphisme $(f, g) : X/U \to Y/V$ de $\Spacesstack$ est
fortement cartésien si le morphisme induit $f : X \to U \times_V Y$
est un isomorphisme. Si $Y \to V$ est de type fini,
alors $U \times_V Y \to U$ l'est aussi, voir
Morphismes d'espaces,
lemme \ref{spaces-morphisms-lemma-base-change-finite-type}.
Ainsi, si $(f, g) : X/U \to Y/V$ dans $\Spacesstack$ est
fortement cartésien dans $\Spacesstack$ et si $Y/V$ est un objet
de $\Spacesstack_{ft}$, alors $X/U$ est automatiquement aussi un
objet de $\Spacesstack_{ft}$ et, bien entendu, $(f, g)$ est
également fortement cartésien dans $\Spacesstack_{ft}$. On en
conclut que $\Spacesstack_{ft}$ est une catégorie fibrée sur
$(\Sch/S)_{fppf}$. Cela prouve (1).

\medskip\noindent
L'argument précédent montre aussi que le foncteur d'inclusion
$\Spacesstack_{ft} \to \Spacesstack$ transforme les
morphismes fortement cartésiens en morphismes fortement cartésiens.
Autrement dit, $\Spacesstack_{ft} \to \Spacesstack$ est
un $1$-morphisme de catégories fibrées sur $(\Sch/S)_{fppf}$.

\medskip\noindent
Soit $U \in \Ob((\Sch/S)_{fppf})$.
Soient $X, Y$ des espaces algébriques de type fini sur $U$. D'après
Champs, lemme \ref{stacks-lemma-presheaf-mor-map-fibred-categories},
nous obtenons un morphisme de préfaisceaux
$$
\mathit{Mor}_{\Spacesstack_{ft}}(X, Y)
\longrightarrow
\mathit{Mor}_{\Spacesstack}(X, Y)
$$
qui est un isomorphisme, puisque $\Spacesstack_{ft}$ est une sous-catégorie pleine de
$\Spacesstack$. Le membre de gauche est donc un faisceau, car nous avons montré,
au lemme \ref{examples-stacks-lemma-pre-stack-of-spaces},
que celui de droite en est un. Cela prouve (2).

\medskip\noindent
Pour prouver la condition (3) de
Champs, définition \ref{stacks-definition-stack},
il faut montrer ceci : étant donnés
\begin{enumerate}
\item un recouvrement $\{U_i \to U\}_{i \in I}$ de $(\Sch/S)_{fppf}$ ;
\item pour chaque $i \in I$, un espace algébrique $X_i$ de type fini sur $U_i$ ;
\item pour chaque $i, j \in I$, un isomorphisme
$\varphi_{ij} : X_i \times_U U_j \to U_i \times_U X_j$ d'espaces algébriques
sur $U_i \times_U U_j$, satisfaisant la condition de cocycle sur
$U_i \times_U U_j \times_U U_k$,
\end{enumerate}
il existe un espace algébrique $X$ de type fini sur $U$ et des isomorphismes
$X_{U_i} \cong X_i$ sur $U_i$ qui redonnent les isomorphismes $\varphi_{ij}$.
Cela résulte de
Amorçage, lemme \ref{bootstrap-lemma-descend-algebraic-space}, partie (2). D'après
Descente sur les espaces, lemme
\ref{spaces-descent-lemma-descending-property-locally-finite-presentation},
le morphisme $X \to U$ est de type fini, ce qui achève la démonstration.
\end{proof}

\begin{lemma}
\label{examples-stacks-lemma-stack-ft-spaces}
Il existe une sous-catégorie
$\Spacesstack_{ft, small} \subset \Spacesstack_{ft}$
ayant les propriétés suivantes :
\begin{enumerate}
\item le foncteur d'inclusion
$\Spacesstack_{ft, small} \to \Spacesstack_{ft}$ est
pleinement fidèle et essentiellement surjectif ;
\item le foncteur
$p_{ft, small} : \Spacesstack_{ft, small} \to (\Sch/S)_{fppf}$
fait de $\Spacesstack_{ft, small}$ un champ sur
$(\Sch/S)_{fppf}$.
\end{enumerate}
\end{lemma}

\begin{proof}
Nous avons vu dans le
lemme \ref{examples-stacks-lemma-stack-of-finite-type-spaces}
que $p_{ft} : \Spacesstack_{ft} \to (\Sch/S)_{fppf}$
satisfait les conditions (1), (2) et (3) de
Champs, définition \ref{stacks-definition-stack}.
La condition supplémentaire (4) de
Champs, remarque \ref{stacks-remark-stack-make-small},
est satisfaite parce que tout espace algébrique $X$ sur $S$ est de la
forme $U/R$, avec $U, R \in \Ob((\Sch/S)_{fppf})$, voir
Espaces, lemme \ref{spaces-lemma-space-presentation}.
Il n'existe donc qu'un ensemble de classes d'isomorphisme d'objets.
La discussion de cette remarque fournit alors $\Spacesstack_{ft, small}$.
\end{proof}

\noindent
Nous effectuerons souvent le remplacement
$$
\Spacesstack_{ft} \leadsto \Spacesstack_{ft, small}
$$
sans le signaler davantage et, par abus de notation, nous désignerons
simplement ce remplacement par $\Spacesstack_{ft}$.

\begin{remark}
\label{examples-stacks-remark-higher-cardinality-spaces}
Toute la discussion de cette section reste valable
si l'on veut considérer les espaces algébriques $X/U$ qui sont
localement de type fini et tels que l'image inverse dans $X$ d'un ouvert affine
de $U$ puisse être recouverte par un nombre dénombrable d'ouverts affines.
Au besoin, on peut aussi introduire la notion de morphisme de
$\kappa$-type (c'est-à-dire imposer une borne au nombre de générateurs des
extensions d'anneaux et une borne au cardinal de la famille d'ouverts affines au-dessus
d'un ouvert affine donné de la base), où $\kappa$ est un cardinal, puis
construire un champ
$$
\Spacesstack_\kappa \longrightarrow (\Sch/S)_{fppf}
$$
exactement de la même manière que ci-dessus (à condition de veiller à ce que
$\Sch$ soit assez grande en fonction de $\kappa$).
\end{remark}





\section{Exemples de champs en groupoïdes}
\label{examples-stacks-section-examples-stacks-in-groupoids}

\noindent
Les exemples précédents sont des champs qui ne sont pas des champs en
groupoïdes. Dans le reste de ce chapitre, nous donnons
des exemples issus de la géométrie algébrique de champs en groupoïdes.



\section{Le champ associé à un faisceau}
\label{examples-stacks-section-stack-associated-to-sheaf}

\noindent
Soit $F : (\Sch/S)_{fppf}^{opp} \to \textit{Ensembles}$ un préfaisceau.
Nous obtenons une catégorie fibrée en ensembles
$$
p_F : \mathcal{S}_F \to (\Sch/S)_{fppf},
$$
voir
Catégories, exemple \ref{categories-example-presheaf}.
C'est un champ en ensembles si et seulement si $F$ est un faisceau, voir
Champs, lemme \ref{stacks-lemma-stack-in-setoids-characterize}.



\section{Le champ en groupoïdes des faisceaux quasi-cohérents de type fini}
\label{examples-stacks-section-stack-in-groupoids-of-quasi-coherent-sheaves}

\noindent
Soit $p : \QCohstack_{fg} \to (\Sch/S)_{fppf}$ le champ
introduit à la
section \ref{examples-stacks-section-stack-of-finitely-generated-quasi-coherent-sheaves}
(avec l'abus de notation qui y a été introduit).
Nous pouvons le transformer en un champ en groupoïdes
$p' : \QCohstack_{fg}' \to (\Sch/S)_{fppf}$ au moyen
du procédé de
Catégories, lemme \ref{categories-lemma-fibred-gives-fibred-groupoids},
voir
Champs, lemme \ref{stacks-lemma-stack-gives-stack-groupoids}.
Dans ce cas particulier, cela signifie simplement que $\QCohstack_{fg}'$ a
les mêmes objets que $\QCohstack_{fg}$, mais que ses morphismes sont les
couples $(f, g) : (U, \mathcal{F}) \to (U', \mathcal{F}')$
où $g$ est un isomorphisme $g : f^*\mathcal{F}' \to \mathcal{F}$.


\section{Le champ en groupoïdes des espaces algébriques de type fini}
\label{examples-stacks-section-stack-in-groupoids-of-finite-type-spaces}

\noindent
Soit $p : \Spacesstack_{ft} \to (\Sch/S)_{fppf}$ le champ
introduit à la
section \ref{examples-stacks-section-stack-of-finite-type-spaces}
(avec l'abus de notation qui y a été introduit).
Nous pouvons le transformer en un champ en groupoïdes
$p' : \Spacesstack_{ft}' \to (\Sch/S)_{fppf}$ au moyen
du procédé de
Catégories, lemme \ref{categories-lemma-fibred-gives-fibred-groupoids},
voir
Champs, lemme \ref{stacks-lemma-stack-gives-stack-groupoids}.
Dans ce cas particulier, cela signifie simplement que $\Spacesstack_{ft}'$
a les mêmes objets que $\Spacesstack_{ft}$, c'est-à-dire les morphismes de type fini
$X \to U$, où $X$ est un espace algébrique sur $S$ et $U$ un schéma
sur $S$. Mais les morphismes $(f, g) : X/U \to Y/V$ sont maintenant les
diagrammes commutatifs
$$
\xymatrix{
X \ar[d] \ar[r]_f & Y \ar[d] \\
U \ar[r]^g & V
}
$$
qui sont cartésiens.



\section{Champs quotients}
\label{examples-stacks-section-quotient-stacks}

\noindent
Soit $(U, R, s, t, c)$ un groupoïde en espaces algébriques sur $S$.
Dans ce cas, le champ quotient
$$
[U/R] \longrightarrow (\Sch/S)_{fppf}
$$
est par construction un champ en groupoïdes, voir
Groupoïdes dans les espaces,
définition \ref{spaces-groupoids-definition-quotient-stack}.
Les faisceaux $\mathit{Isom}$ sont même
représentables par des espaces algébriques, voir
Amorçage, lemme \ref{bootstrap-lemma-quotient-stack-isom}.
Ces champs quotients sont d'une importance fondamentale dans la théorie des
champs algébriques.

\medskip\noindent
Un cas particulier de la construction précédente est le champ quotient
$$
[X/G] \longrightarrow (\Sch/S)_{fppf}
$$
associé à une donnée $(B, G/B, m, X/B, a)$. Ici,
\begin{enumerate}
\item $B$ est un espace algébrique sur $S$ ;
\item $(G, m)$ est un espace algébrique en groupes sur $B$ ;
\item $X$ est un espace algébrique sur $B$ ;
\item $a : G \times_B X \to X$ est une action de $G$ sur $X$ au-dessus de $B$.
\end{enumerate}
En effet, d'après
Groupoïdes dans les espaces,
définition \ref{spaces-groupoids-definition-quotient-stack},
le champ en groupoïdes $[X/G]$ est le
champ quotient $[X/G \times_B X]$ ci-dessus. Il convient
d'expliciter la catégorie $[X/G]$. Nous le ferons à la
section \ref{examples-stacks-section-group-quotient-stacks}.



\section{Classification des torseurs}
\label{examples-stacks-section-torsors}

\noindent
Nous allons expliquer avec soin plusieurs variantes de ce que peut
signifier l'étude du champ des torseurs sous un espace algébrique en groupes $G$
ou sous un faisceau de groupes $\mathcal{G}$.



\subsection{Torseurs sous un faisceau de groupes}
\label{examples-stacks-subsection-torsors-sheaf}

\noindent
Soit $\mathcal{G}$ un faisceau de groupes sur $(\Sch/S)_{fppf}$.
Pour $U \in \Ob((\Sch/S)_{fppf})$, nous notons
$\mathcal{G}|_U$ la restriction de $\mathcal{G}$ à $(\Sch/U)_{fppf}$.
Nous définissons une catégorie $\mathcal{G}\textit{-Torsors}$ comme suit :
\begin{enumerate}
\item Un objet de $\mathcal{G}\textit{-Torsors}$ est un couple
$(U, \mathcal{F})$, où $U$ est un objet de $(\Sch/S)_{fppf}$
et $\mathcal{F}$ un torseur sous $\mathcal{G}|_U$, voir
Cohomologie sur les sites, définition \ref{sites-cohomology-definition-torsor}.
\item Un morphisme $(U, \mathcal{F}) \to (V, \mathcal{H})$ est donné
par un couple $(f, \alpha)$, où $f : U \to V$ est un morphisme de schémas
sur $S$, et $\alpha : f^{-1}\mathcal{H} \to \mathcal{F}$ un
isomorphisme de torseurs sous $\mathcal{G}|_U$.
\end{enumerate}
Ainsi, $\mathcal{G}\textit{-Torsors}$ est une catégorie et
$$
p : \mathcal{G}\textit{-Torsors} \longrightarrow (\Sch/S)_{fppf},
\quad
(U, \mathcal{F}) \longmapsto U
$$
est un foncteur. La catégorie fibre de $\mathcal{G}\textit{-Torsors}$
au-dessus de $U$ est la catégorie des torseurs sous $\mathcal{G}|_U$, qui est un groupoïde.

\begin{lemma}
\label{examples-stacks-lemma-torsors-sheaf-stack-in-groupoids}
Moyennant un remplacement comme dans
Champs, remarque \ref{stacks-remark-stack-make-small},
le foncteur
$$
p : \mathcal{G}\textit{-Torsors} \longrightarrow (\Sch/S)_{fppf}
$$
définit un champ en groupoïdes sur $(\Sch/S)_{fppf}$.
\end{lemma}

\begin{proof}
La partie la plus difficile de la démonstration consiste à montrer
l'effectivité de la descente des objets.
Soit $\{U_i \to U\}_{i \in I}$ un recouvrement de $(\Sch/S)_{fppf}$.
Supposons donné, pour chaque $i$, un torseur $\mathcal{F}_i$
sous $\mathcal{G}|_{U_i}$ et, pour chaque $i, j \in I$, un isomorphisme
$\varphi_{ij} :
\mathcal{F}_i|_{U_i \times_U U_j} \to \mathcal{F}_j|_{U_i \times_U U_j}$
de torseurs sous $\mathcal{G}|_{U_i \times_U U_j}$,
satisfaisant une condition de cocycle convenable sur $U_i \times_U U_j \times_U U_k$.
Alors, d'après
Sites, section \ref{sites-section-glueing-sheaves},
nous obtenons sur $(\Sch/U)_{fppf}$ un faisceau $\mathcal{F}$
dont la restriction à chaque $U_i$ redonne $\mathcal{F}_i$ ainsi que
la donnée de descente. Par l'équivalence de catégories de
Sites, lemme \ref{sites-lemma-mapping-property-glue},
les morphismes d'action $\mathcal{G}|_{U_i} \times \mathcal{F}_i \to \mathcal{F}_i$
se recollent en un morphisme $a : \mathcal{G}|_U \times \mathcal{F} \to \mathcal{F}$.
Il reste à montrer que $a$ est une action et que $\mathcal{F}$ devient
un torseur sous $\mathcal{G}|_U$. Ces deux propriétés se vérifient localement
et résultent donc des propriétés correspondantes des actions
$\mathcal{G}|_{U_i} \times \mathcal{F}_i \to \mathcal{F}_i$.
Cela prouve l'effectivité de la descente des objets dans
$\mathcal{G}\textit{-Torsors}$.
Nous omettons certains détails.
\end{proof}



\subsection{Variante pour les torseurs sous un faisceau}
\label{examples-stacks-subsection-variant-torsor-sheaf}

\noindent
On peut généraliser légèrement la construction de la
sous-section \ref{examples-stacks-subsection-torsors-sheaf}.
Soient en effet $\mathcal{G} \to \mathcal{B}$ un morphisme de faisceaux
sur $(\Sch/S)_{fppf}$ et
$$
m :
\mathcal{G} \times_\mathcal{B} \mathcal{G}
\longrightarrow
\mathcal{G}
$$
une loi de groupe sur $\mathcal{G}/\mathcal{B}$. Autrement dit, le couple
$(\mathcal{G}, m)$ est un objet en groupes du topos
$\Sh((\Sch/S)_{fppf})/\mathcal{B}$. Voir
Sites, section \ref{sites-section-localize-topoi},
pour les localisations de topos.
Dans cette situation, nous pouvons définir une catégorie
$\mathcal{G}/\mathcal{B}\textit{-Torsors}$ comme suit
(nous utilisons le plongement de Yoneda pour considérer les schémas comme des faisceaux) :
\begin{enumerate}
\item Un objet de $\mathcal{G}/\mathcal{B}\textit{-Torsors}$ est un triplet
$(U, b, \mathcal{F})$, où
\begin{enumerate}
\item $U$ est un objet de $(\Sch/S)_{fppf}$ ;
\item $b : U \to \mathcal{B}$ est une section de $\mathcal{B}$ sur $U$ ;
\item $\mathcal{F}$ est un torseur sous $U \times_{b, \mathcal{B}}\mathcal{G}$
sur $U$.
\end{enumerate}
\item Un morphisme $(U, b, \mathcal{F}) \to (U', b', \mathcal{F}')$ est donné
par un couple $(f, g)$, où $f : U \to U'$ est un morphisme de schémas
sur $S$ tel que $b = b' \circ f$, et où
$g : f^{-1}\mathcal{F}' \to \mathcal{F}$ est un
isomorphisme de torseurs sous $U \times_{b, \mathcal{B}} \mathcal{G}$.
\end{enumerate}
Ainsi, $\mathcal{G}/\mathcal{B}\textit{-Torsors}$ est une catégorie et
$$
p :
\mathcal{G}/\mathcal{B}\textit{-Torsors}
\longrightarrow
(\Sch/S)_{fppf},
\quad
(U, b, \mathcal{F}) \longmapsto U
$$
est un foncteur. La catégorie fibre de
$\mathcal{G}/\mathcal{B}\textit{-Torsors}$
au-dessus de $U$ est la réunion disjointe, lorsque $b : U \to \mathcal{B}$ varie,
des catégories de torseurs sous $U \times_{b, \mathcal{B}} \mathcal{G}$ ;
c'est donc un groupoïde.

\medskip\noindent
Dans le cas particulier $\mathcal{B} = S$, on retrouve la catégorie
$\mathcal{G}\textit{-Torsors}$ introduite dans la
sous-section \ref{examples-stacks-subsection-torsors-sheaf}.

\begin{lemma}
\label{examples-stacks-lemma-variant-torsors-sheaf-stack-in-groupoids}
Moyennant un remplacement comme dans
Champs, remarque \ref{stacks-remark-stack-make-small},
le foncteur
$$
p :
\mathcal{G}/\mathcal{B}\textit{-Torsors}
\longrightarrow
(\Sch/S)_{fppf}
$$
définit un champ en groupoïdes sur $(\Sch/S)_{fppf}$.
\end{lemma}

\begin{proof}
Cette démonstration reprend celle du
lemme \ref{examples-stacks-lemma-torsors-sheaf-stack-in-groupoids}.
Le lecteur est invité à lire d'abord cette dernière, dont
les notations sont moins lourdes.
La partie la plus difficile de la démonstration consiste à montrer
l'effectivité de la descente des objets. Soit $\{U_i \to U\}_{i \in I}$
un recouvrement de $(\Sch/S)_{fppf}$.
Supposons donné, pour chaque $i$, un couple $(b_i, \mathcal{F}_i)$
formé d'un morphisme $b_i : U_i \to \mathcal{B}$ et d'un
torseur $\mathcal{F}_i$ sous $U_i \times_{b_i, \mathcal{B}} \mathcal{G}$,
et supposons que, pour chaque $i, j \in I$,
on ait $b_i|_{U_i \times_U U_j} = b_j|_{U_i \times_U U_j}$ et
que soit donné un isomorphisme
$\varphi_{ij} :
\mathcal{F}_i|_{U_i \times_U U_j} \to \mathcal{F}_j|_{U_i \times_U U_j}$
de torseurs sous $(U_i \times_U U_j) \times_\mathcal{B} \mathcal{G}$,
satisfaisant une condition de cocycle convenable sur $U_i \times_U U_j \times_U U_k$.
Alors, d'après
Sites, section \ref{sites-section-glueing-sheaves},
nous obtenons sur $(\Sch/U)_{fppf}$ un faisceau $\mathcal{F}$
dont la restriction à chaque $U_i$ redonne $\mathcal{F}_i$ ainsi que
la donnée de descente. L'axiome de faisceau pour $\mathcal{B}$ montre
que les morphismes $b_i$ proviennent d'un unique morphisme $b : U \to \mathcal{B}$.
Par l'équivalence de catégories de
Sites, lemme \ref{sites-lemma-mapping-property-glue},
les morphismes d'action
$(U_i \times_{b_i, \mathcal{B}} \mathcal{G}) \times_{U_i} \mathcal{F}_i
\to \mathcal{F}_i$
se recollent en un morphisme
$(U \times_{b, \mathcal{B}} \mathcal{G}) \times \mathcal{F} \to \mathcal{F}$.
Il reste à montrer que c'est une action et que $\mathcal{F}$ devient
un torseur sous $U \times_{b, \mathcal{B}} \mathcal{G}$.
Ces deux propriétés se vérifient localement
et résultent donc des propriétés correspondantes des actions
sur les $\mathcal{F}_i$.
Cela prouve l'effectivité de la descente des objets dans
$\mathcal{G}/\mathcal{B}\textit{-Torsors}$.
Nous omettons certains détails.
\end{proof}



\subsection{Espaces principaux homogènes}
\label{examples-stacks-subsection-principal-homogeneous-spaces}

\noindent
Soit $B$ un espace algébrique sur $S$.
Soit $G$ un espace algébrique en groupes sur $B$.
Nous définissons une catégorie $G\textit{-Principal}$ comme suit :
\begin{enumerate}
\item Un objet de $G\textit{-Principal}$ est un triplet $(U, b, X)$, où
\begin{enumerate}
\item $U$ est un objet de $(\Sch/S)_{fppf}$ ;
\item $b : U \to B$ est un morphisme sur $S$ ;
\item $X$ est un espace principal homogène sous $G_U$ sur $U$, où
$G_U = U \times_{b, B} G$.
\end{enumerate}
Voir
Groupoïdes dans les espaces,
définition \ref{spaces-groupoids-definition-principal-homogeneous-space}.
\item Un morphisme $(U, b, X) \to (U', b', X')$ est donné
par un couple $(f, g)$, où $f : U \to U'$ est un morphisme de schémas
sur $B$, et $g : X \to U \times_{f, U'} X'$ un
isomorphisme d'espaces principaux homogènes sous $G_U$.
\end{enumerate}
Ainsi, $G\textit{-Principal}$ est une catégorie et
$$
p : G\textit{-Principal} \longrightarrow (\Sch/S)_{fppf},
\quad
(U, b, X) \longmapsto U
$$
est un foncteur. La catégorie fibre de $G\textit{-Principal}$
au-dessus de $U$ est la réunion disjointe, lorsque $b : U \to B$ varie,
des catégories d'espaces principaux homogènes sous $U \times_{b, B} G$ ;
c'est donc un groupoïde.

\medskip\noindent
Dans le cas particulier $S = B$, les objets sont simplement les couples
$(U, X)$, où $U$ est un schéma sur $S$ et $X$ un espace principal homogène
sous $G_U$ sur $U$. De plus, les morphismes sont simplement les diagrammes
cartésiens
$$
\xymatrix{
X \ar[d] \ar[r]_g & X' \ar[d] \\
U \ar[r]^f & U'
}
$$
où $g$ est $G$-équivariant.

\begin{remark}
\label{examples-stacks-remark-principal-stack-in-groupoids}
Nous conjecturons que, moyennant un remplacement comme dans
Champs, remarque \ref{stacks-remark-stack-make-small},
le foncteur
$$
p : G\textit{-Principal} \longrightarrow (\Sch/S)_{fppf}
$$
définit un champ en groupoïdes sur $(\Sch/S)_{fppf}$. Cela résulterait
de l'énoncé suivant : étant donnés
\begin{enumerate}
\item un recouvrement $\{U_i \to U\}_{i \in I}$ de $(\Sch/S)_{fppf}$ ;
\item un espace algébrique en groupes $H$ sur $U$ ;
\item pour chaque $i$, un espace principal homogène $X_i$ sous $H_{U_i}$
sur $U_i$ ;
\item des isomorphismes $H$-équivariants
$\varphi_{ij} : X_{i, U_i \times_U U_j} \to X_{j, U_i \times_U U_j}$
satisfaisant la condition de cocycle,
\end{enumerate}
il existe un espace principal homogène $X$ sous $H$ sur $U$
qui redonne $(X_i, \varphi_{ij})$. La méthode employée dans la démonstration de
Amorçage, lemme \ref{bootstrap-lemma-descent-torsor},
ramène ce problème à une question ensembliste ; le lecteur qui néglige
les questions ensemblistes « saura » donc que le résultat est vrai. On trouve dans
\url{https://math.columbia.edu/~dejong/wordpress/?p=591}
une suggestion pour aborder ce problème.
\end{remark}



\subsection{Variante pour les espaces principaux homogènes}
\label{examples-stacks-subsection-variant-principal-homogeneous-spaces}

\noindent
Soit $S$ un schéma. Soit $B = S$.
Soit $G$ un schéma en groupes sur $B = S$.
Dans cette situation, nous pouvons définir une sous-catégorie pleine
$G\textit{-Principal-Schémas} \subset G\textit{-Principal}$
dont les objets sont les couples $(U, X)$, où $U$ est un objet de
$(\Sch/S)_{fppf}$ et $X \to U$ un espace principal homogène
sous $G$ sur $U$ qui est représentable, c'est-à-dire un schéma.

\medskip\noindent
En général, $G\textit{-Principal-Schémas}$ n'est pas
un champ en groupoïdes sur $(\Sch/S)_{fppf}$. Cela tient
à ce qu'il existe effectivement, en général, des espaces principaux homogènes
qui ne sont pas des schémas ; la descente des objets n'est donc pas effective
en général.



\subsection{Torseurs pour la topologie fppf}
\label{examples-stacks-subsection-fppf-torsors}

\noindent
Soit $B$ un espace algébrique sur $S$.
Soit $G$ un espace algébrique en groupes sur $B$.
Nous définissons une catégorie $G\textit{-Torsors}$ comme suit :
\begin{enumerate}
\item Un objet de $G\textit{-Torsors}$ est un triplet $(U, b, X)$, où
\begin{enumerate}
\item $U$ est un objet de $(\Sch/S)_{fppf}$ ;
\item $b : U \to B$ est un morphisme ;
\item $X$ est un $G_U$-torseur fppf sur $U$, où $G_U = U \times_{b, B} G$.
\end{enumerate}
Voir
Groupoïdes dans les espaces,
définition \ref{spaces-groupoids-definition-principal-homogeneous-space}.
\item Un morphisme $(U, b, X) \to (U', b', X')$ est donné
par un couple $(f, g)$, où $f : U \to U'$ est un morphisme de schémas
sur $B$, et $g : X \to U \times_{f, U'} X'$ un
isomorphisme de $G_U$-torseurs.
\end{enumerate}
Ainsi, $G\textit{-Torsors}$ est une catégorie et
$$
p : G\textit{-Torsors} \longrightarrow (\Sch/S)_{fppf},
\quad
(U, a, X) \longmapsto U
$$
est un foncteur. La catégorie fibre de $G\textit{-Torsors}$
au-dessus de $U$ est la réunion disjointe, lorsque $b : U \to B$ varie,
des catégories de $U \times_{b, B} G$-torseurs fppf ;
c'est donc un groupoïde.

\medskip\noindent
Dans le cas particulier $S = B$, les objets sont simplement les couples
$(U, X)$, où $U$ est un schéma sur $S$ et $X$ un
$G_U$-torseur fppf sur $U$. De plus, les morphismes sont simplement les diagrammes
cartésiens
$$
\xymatrix{
X \ar[d] \ar[r]_g & X' \ar[d] \\
U \ar[r]^f & U'
}
$$
où $g$ est $G$-équivariant.

\begin{lemma}
\label{examples-stacks-lemma-torsors-stack-in-groupoids}
Moyennant un remplacement comme dans
Champs, remarque \ref{stacks-remark-stack-make-small},
le foncteur
$$
p : G\textit{-Torsors} \longrightarrow (\Sch/S)_{fppf}
$$
définit un champ en groupoïdes sur $(\Sch/S)_{fppf}$.
\end{lemma}

\begin{proof}
La partie la plus difficile de la démonstration consiste à montrer l'effectivité
de la descente des objets ; c'est le contenu de
Amorçage, lemme \ref{bootstrap-lemma-descent-torsor}.
Nous omettons la démonstration des axiomes (1) et (2) de
Champs, définition \ref{stacks-definition-stack-in-groupoids}.
\end{proof}

\begin{lemma}
\label{examples-stacks-lemma-compare-torsors}
Soit $B$ un espace algébrique sur $S$. Soit $G$ un espace algébrique en
groupes sur $B$. Notons $\mathcal{G}$, resp.\ $\mathcal{B}$, l'espace algébrique
$G$, resp.\ $B$, considéré comme un faisceau sur $(\Sch/S)_{fppf}$.
Le foncteur
$$
G\textit{-Torsors} \longrightarrow \mathcal{G}/\mathcal{B}\textit{-Torsors}
$$
qui associe au triplet $(U, b, X)$ le triplet
$(U, b, \mathcal{X})$, où $\mathcal{X}$ est $X$ considéré comme un faisceau,
est une équivalence de champs en groupoïdes sur $(\Sch/S)_{fppf}$.
\end{lemma}

\begin{proof}
Nous allons utiliser
Champs, lemme \ref{stacks-lemma-characterize-essentially-surjective-when-ff},
pour le prouver. Le foncteur est pleinement fidèle, puisque la catégorie des
espaces algébriques sur $S$ est une sous-catégorie pleine de la catégorie des
faisceaux sur $(\Sch/S)_{fppf}$.
En outre, tous les objets (des deux côtés) sont des torseurs localement triviaux,
de sorte que la condition (2) du lemme cité ci-dessus est satisfaite.
Le foncteur est donc une équivalence.
\end{proof}



\subsection{Variante pour les torseurs en topologie fppf}
\label{examples-stacks-subsection-variant-fppf-torsors}

\noindent
Soit $S$ un schéma. Soit $B = S$.
Soit $G$ un schéma en groupes sur $B = S$.
Dans cette situation, nous pouvons définir une sous-catégorie pleine
$G\textit{-Torsors-Schémas} \subset G\textit{-Torsors}$
dont les objets sont les couples $(U, X)$, où $U$ est un objet de
$(\Sch/S)_{fppf}$ et $X \to U$ un
$G$-torseur fppf sur $U$ qui est représentable, c'est-à-dire un schéma.

\medskip\noindent
En général, $G\textit{-Torsors-Schémas}$ n'est pas
un champ en groupoïdes sur $(\Sch/S)_{fppf}$. Cela tient
à ce qu'il existe effectivement, en général, des $G$-torseurs fppf
qui ne sont pas des schémas ; la descente des objets n'est donc pas effective
en général.




\section{Quotients par des actions de groupes}
\label{examples-stacks-section-group-quotient-stacks}

\noindent
Nous avons maintenant introduit assez de notations pour décrire
plus précisément les champs $[X/G]$ de la
section \ref{examples-stacks-section-quotient-stacks}.

\begin{situation}
\label{examples-stacks-situation-quotient-stack}
Dans cette situation,
\begin{enumerate}
\item $S$ est un schéma contenu dans $\Sch_{fppf}$ ;
\item $B$ est un espace algébrique sur $S$ ;
\item $(G, m)$ est un espace algébrique en groupes sur $B$ ;
\item $\pi : X \to B$ est un espace algébrique sur $B$ ;
\item $a : G \times_B X \to X$ est une action de $G$ sur $X$ au-dessus de $B$.
\end{enumerate}
\end{situation}

\noindent
Dans cette situation, nous construisons comme suit une catégorie $[[X/G]]$\footnote{La notation
$[[X/G]]$ avec doubles crochets sert à distinguer cette catégorie du
champ $[X/G]$ introduit précédemment. Dans la
proposition \ref{examples-stacks-proposition-equal-quotient-stacks},
nous montrerons que les deux sont canoniquement équivalents.
Nous utiliserons ensuite la notation $[X/G]$ pour désigner indifféremment l'un ou l'autre.} :
\begin{enumerate}
\item Un objet de $[[X/G]]$ consiste en un quadruplet
$(U, b, P, \varphi : P \to X)$, où
\begin{enumerate}
\item $U$ est un objet de $(\Sch/S)_{fppf}$ ;
\item $b : U \to B$ est un morphisme sur $S$ ;
\item $P$ est un $G_U$-torseur fppf sur $U$, où $G_U = U \times_{b, B} G$ ;
\item $\varphi : P \to X$ est un morphisme $G$-équivariant qui s'insère
dans le diagramme commutatif
$$
\xymatrix{
P \ar[d] \ar[r]_{\varphi} & X \ar[d] \\
U \ar[r]^b & B
}
$$
\end{enumerate}
\item Un morphisme de $[[X/G]]$ est un couple
$(f, g) : (U, b, P, \varphi) \to (U', b', P', \varphi')$,
où $f : U \to U'$ est un morphisme de schémas sur $B$
et $g : P \to P'$ un morphisme $G$-équivariant au-dessus de $f$
qui induit un isomorphisme $P \cong U \times_{f, U'} P'$ et vérifie
$\varphi = \varphi' \circ g$.
Autrement dit, $(f, g)$ s'insère dans le diagramme commutatif
suivant :
$$
\xymatrix{
P \ar[d] \ar[rrrd]_\varphi \ar[r]^g & P' \ar[d] \ar[rrd]^{\varphi'} \\
U \ar[rrrd]_b \ar[r]^f & U' \ar[rrd]^{b'} & & X \ar[d] \\
& & & B
}
$$
\end{enumerate}
Ainsi, $[[X/G]]$ est une catégorie et
$$
p : [[X/G]] \longrightarrow (\Sch/S)_{fppf},
\quad
(U, b, P, \varphi) \longmapsto U
$$
est un foncteur. La catégorie fibre de $[[X/G]]$
au-dessus de $U$ est la réunion disjointe, lorsque $b \in \Mor_S(U, B)$ varie,
des $U \times_{b, B} G$-torseurs fppf $P$ munis d'un morphisme
$G$-équivariant vers $X$. Les catégories fibres de $[[X/G]]$ sont donc des groupoïdes.

\medskip\noindent
Remarquons que le foncteur
$$
[[X/G]] \longrightarrow G\textit{-Torsors},
\quad
(U, b, P, \varphi) \longmapsto (U, b, P)
$$
est un $1$-morphisme de catégories sur $(\Sch/S)_{fppf}$.

\begin{lemma}
\label{examples-stacks-lemma-group-quotient-stack-in-groupoids}
Moyennant un remplacement comme dans
Champs, remarque \ref{stacks-remark-stack-make-small},
le foncteur
$$
p : [[X/G]] \longrightarrow (\Sch/S)_{fppf}
$$
définit un champ en groupoïdes sur $(\Sch/S)_{fppf}$.
\end{lemma}

\begin{proof}
La partie la plus difficile de la démonstration consiste à montrer l'effectivité
de la descente des objets. Supposons que $\{U_i \to U\}_{i \in I}$ soit un recouvrement de
$(\Sch/S)_{fppf}$. Soient
$\xi_i = (U_i, b_i, P_i, \varphi_i)$ des objets de $[[X/G]]$ au-dessus de $U_i$,
et soit $\varphi_{ij} : \text{pr}_0^*\xi_i \to \text{pr}_1^*\xi_j$
une donnée de descente. En appliquant le foncteur
$[[X/G]] \to G\textit{-Torsors}$ ci-dessus, nous obtenons en particulier une donnée de descente
sur les triplets $(U_i, b_i, P_i)$ pour le champ en groupoïdes
$G\textit{-Torsors}$. Nous avons vu que
$G\textit{-Torsors}$ est un champ en groupoïdes
(lemme \ref{examples-stacks-lemma-torsors-stack-in-groupoids}).
Nous pouvons donc supposer que $b_i = b|_{U_i}$ pour un morphisme $b : U \to B$, et
que $P_i = U_i \times_U P$ pour un $G_U = U \times_{b, B} G$-torseur fppf
$P$ sur $U$. Les morphismes $\varphi_i$ sont compatibles
avec la donnée de descente canonique sur les restrictions $U_i \times_U P$
et définissent donc un morphisme $\varphi : P \to X$. (On peut par exemple
utiliser
Sites, lemme \ref{sites-lemma-mapping-property-glue},
ou
Descente sur les espaces,
lemme \ref{spaces-descent-lemma-fpqc-universal-effective-epimorphisms},
pour obtenir $\varphi$.)
Cela prouve l'effectivité de la descente des objets.
Nous omettons la démonstration des axiomes (1) et (2) de
Champs, définition \ref{stacks-definition-stack-in-groupoids}.
\end{proof}

\begin{proposition}
\label{examples-stacks-proposition-equal-quotient-stacks}
Dans la
situation \ref{examples-stacks-situation-quotient-stack},
il existe une équivalence canonique
$$
[X/G] \longrightarrow [[X/G]]
$$
de champs en groupoïdes sur $(\Sch/S)_{fppf}$.
\end{proposition}

\begin{proof}
Nous donnons tous les détails afin de nous assurer que les définitions
s'accordent exactement comme il le faut.
Rappelons que $[X/G]$ est le champ quotient
associé au groupoïde en espaces algébriques
$(X, G \times_B X, s, t, c)$, voir
Groupoïdes dans les espaces,
définition \ref{spaces-groupoids-definition-quotient-stack}.
Cela signifie que $[X/G]$ est la champification de la
catégorie fibrée en groupoïdes $[X/_{\!p}G]$ associée au foncteur
$$
(\Sch/S)_{fppf} \longrightarrow \textit{Groupoïdes},
\quad
U \longmapsto (X(U), G(U) \times_{B(U)} X(U), s, t, c),
$$
où $s(g, x) = x$, $t(g, x) = a(g, x)$ et
$c((g, x), (g', x')) = (m(g, g'), x')$. D'après la construction de
Catégories, exemple \ref{categories-example-functor-groupoids},
un objet de $[X/_{\!p}G]$ est un couple $(U, x)$ avec $x \in X(U)$,
et un morphisme $(f, g) : (U, x) \to (U', x')$ de $[X/_{\!p}G]$
est donné par un morphisme de schémas $f : U \to U'$ et un élément
$g \in G(U)$ tel que $a(g, x) = x' \circ f$.
Nous pouvons donc définir un $1$-morphisme de champs en groupoïdes
$$
F_p : [X/_{\!p}G] \longrightarrow [[X/G]]
$$
par les règles suivantes : sur les objets, nous posons
$$
F_p(U, x) =
(U, \pi \circ x, G \times_{B, \pi \circ x} U, a \circ (\text{id}_G \times x)).
$$
Cette définition a un sens, car le diagramme
$$
\xymatrix{
G \times_{B, \pi \circ x} U \ar[d] \ar[r]_{\text{id}_G \times x} &
G \times_{B, \pi} X \ar[r]_-a &
X \ar[d]^\pi \\
U \ar[rr]^{\pi \circ x} & & B
}
$$
est commutatif, et les deux flèches horizontales sont $G$-équivariantes si l'on considère
les produits fibrés comme des $G$-torseurs triviaux sur $U$, resp.\ sur $X$.
Sur un morphisme $(f, g) : (U, x) \to (U', x')$, nous posons
$F_p(f, g) = (f, R_{g^{-1}})$, où $R_{g^{-1}}$ désigne la translation à droite
par l'inverse de $g$. Plus précisément, le morphisme
$F_p(f, g) : F_p(U, x) \to F_p(U', x')$ est donné par le diagramme cartésien
$$
\xymatrix{
G \times_{B, \pi \circ x} U \ar[d] \ar[r]_{R_{g^{-1}}} &
G \times_{B, \pi \circ x'} U' \ar[d] \\
U \ar[r]^f & U'
}
$$
où, sur les points à valeurs dans $T$, $R_{g^{-1}}$ est défini par
$$
R_{g^{-1}}(g', u) = (m(g', i(g(u))), f(u)).
$$
Pour vérifier que cela convient, il faut établir l'égalité
$$
a \circ (\text{id}_G \times x)
=
a \circ (\text{id}_G \times x') \circ R_{g^{-1}},
$$
qui est vraie puisque le membre de droite, appliqué au point à valeurs dans $T$
$(g', u)$, donne bien l'égalité souhaitée
\begin{align*}
a((\text{id}_G \times x')(m(g', i(g(u))), f(u)))
& =
a(m(g', i(g(u))), x'(f(u))) \\
& =
a(g', a(i(g(u)), x'(f(u)))) \\
& =
a(g', x(u)),
\end{align*}
car $a(g, x) = x' \circ f$ et donc $a(i(g), x' \circ f) = x$.

\medskip\noindent
La propriété universelle de la champification donnée dans
Champs, lemme \ref{stacks-lemma-stackify-groupoids-universal-property},
fournit une extension canonique $F : [X/G] \to [[X/G]]$ du $1$-morphisme
$F_p$ ci-dessus. Montrons d'abord que $F$ est pleinement fidèle.
Comme la source et le but sont des champs en groupoïdes,
il suffit pour cela de montrer que $F$ identifie les faisceaux $\mathit{Isom}$.
Choisissons un schéma $U$ et des objets $\xi, \xi'$ de
$[X/G]$ au-dessus de $U$. Nous voulons montrer que
$$
F :
\mathit{Isom}_{[X/G]}(\xi, \xi')
\longrightarrow
\mathit{Isom}_{[[X/G]]}(F(\xi), F(\xi'))
$$
est un isomorphisme de faisceaux. Il suffit de travailler localement
sur $U$ ; nous pouvons donc supposer que $\xi, \xi'$ proviennent d'objets
$(U, x)$, $(U, x')$ de $[X/_{\!p}G]$ au-dessus de $U$. Cela résulte directement
de la construction de la champification et est aussi détaillé dans
Groupoïdes dans les espaces,
section \ref{spaces-groupoids-section-explicit-quotient-stacks}.
Soit en utilisant directement la description ci-dessus des morphismes de
$[X/_{\!p}G]$, soit en appliquant
Groupoïdes dans les espaces,
lemme \ref{spaces-groupoids-lemma-quotient-stack-morphisms},
nous voyons que, dans ce cas,
$$
\mathit{Isom}_{[X/G]}(\xi, \xi') =
U \times_{(x, x'), X \times_S X, (s, t)} (G \times_B X).
$$
Un point à valeurs dans $T$ de ce produit fibré correspond à un couple
$(u, g)$, où $u \in U(T)$ et $g \in G(T)$, tel que
$a(g, x \circ u) = x' \circ u$. (Cela implique en particulier
$\pi \circ x \circ u = \pi \circ x' \circ u$.)
D'autre part, un point à valeurs dans $T$
de $\mathit{Isom}_{[[X/G]]}(F(\xi), F(\xi'))$ correspond par définition
à un morphisme $u : T \to U$ tel que
$\pi \circ x \circ u = \pi \circ x' \circ u : T \to B$, ainsi qu'à un isomorphisme
$$
R :
G \times_{B, \pi \circ x \circ u} T
\longrightarrow
G \times_{B, \pi \circ x' \circ u} T
$$
de $G_T$-torseurs triviaux compatible avec les morphismes donnés vers $X$.
Puisque les torseurs sont triviaux, on a $R = R_{g^{-1}}$
(multiplication à droite) pour un certain $g \in G(T)$. La compatibilité avec les morphismes
$a \circ (1_G, x \circ u), a \circ (1_G, x' \circ u) : G \times_B T \to X$
équivaut à la condition $a(g, x \circ u) = x' \circ u$.
Nous obtenons donc l'égalité voulue des faisceaux $\mathit{Isom}$.

\medskip\noindent
Sachant maintenant que $F$ est pleinement fidèle, nous pouvons appliquer
Champs, lemme \ref{stacks-lemma-characterize-essentially-surjective-when-ff}.
Pour montrer que $F$ est une équivalence, il suffit donc
de montrer que les objets de $[[X/G]]$ appartiennent localement pour la topologie fppf à l'image essentielle
de $F$. Cela est clair, puisque les torseurs fppf sont localement triviaux pour la topologie fppf,
et la conclusion en résulte.
\end{proof}

\begin{lemma}
\label{examples-stacks-lemma-classifying-stacks}
\begin{slogan}
Le champ classifiant d'un schéma en groupes ou d'un espace algébrique en groupes.
\end{slogan}
Soit $S$ un schéma. Soit $B$ un espace algébrique sur $S$.
Soit $G$ un espace algébrique en groupes sur $B$. Alors les champs
en groupoïdes
$$
[B/G],\quad
[[B/G]],\quad
G\textit{-Torsors},\quad
\mathcal{G}/\mathcal{B}\textit{-Torsors}
$$
sont tous canoniquement équivalents.
Si $G \to B$ est plat et localement
de présentation finie, ils sont aussi équivalents à
$G\textit{-Principal}$.
\end{lemma}

\begin{proof}
L'équivalence
$G\textit{-Torsors} \to \mathcal{G}/\mathcal{B}\textit{-Torsors}$
est donnée au lemme \ref{examples-stacks-lemma-compare-torsors}.
L'équivalence $[B/G] \to [[B/G]]$ est donnée par la
proposition \ref{examples-stacks-proposition-equal-quotient-stacks}.
En développant la définition de $[[B/G]]$ donnée à la
section \ref{examples-stacks-section-group-quotient-stacks},
nous voyons que $[[B//G]] = G\textit{-Torsors}$.

\medskip\noindent
Supposons enfin $G \to B$ plat et localement de présentation finie.
Pour montrer que le foncteur naturel
$G\textit{-Torsors} \to G\textit{-Principal}$ est une équivalence,
il suffit de montrer que, pour un schéma $U$ sur $B$,
un espace principal homogène $X \to U$ sous $G_U$
est localement trivial pour la topologie fppf. Par notre définition des espaces principaux homogènes
(Groupoïdes dans les espaces,
définition \ref{spaces-groupoids-definition-principal-homogeneous-space}),
il existe un recouvrement fpqc $\{U_i \to U\}$ tel que
$U_i \times_U X \cong G \times_B U_i$ comme espaces algébriques sur $U_i$.
Il en résulte que $X \to U$ est surjectif, plat et localement de présentation
finie, voir
Descente sur les espaces, lemmes
\ref{spaces-descent-lemma-descending-property-surjective},
\ref{spaces-descent-lemma-descending-property-flat} et
\ref{spaces-descent-lemma-descending-property-locally-finite-presentation}.
Choisissons un schéma $W$ et un morphisme étale surjectif $W \to X$.
Il résulte alors de ce qui précède que $\{W \to U\}$ est un recouvrement fppf
tel que $X_W \to W$ admette une section. Ainsi, $X$ est un $G_U$-torseur fppf.
\end{proof}

\begin{remark}
\label{examples-stacks-remark-X-mod-G-group}
Soit $S$ un schéma.
Soit $G$ un groupe abstrait.
Soit $X$ un espace algébrique sur $S$.
Soit $G \to \text{Aut}_S(X)$ un homomorphisme de groupes.
Dans cette situation, nous pouvons définir $[[X/G]]$ de façon analogue
à ce qui précède, comme suit :
\begin{enumerate}
\item Un objet de $[[X/G]]$ consiste en un triplet
$(U, P, \varphi : P \to X)$, où
\begin{enumerate}
\item $U$ est un objet de $(\Sch/S)_{fppf}$ ;
\item $P$ est un faisceau sur $(\Sch/U)_{fppf}$ muni
d'une action de $G$ qui en fait un torseur sous le faisceau constant
de valeur $G$ ;
\item $\varphi : P \to X$ est un morphisme $G$-équivariant de faisceaux.
\end{enumerate}
\item Un morphisme
$(f, g) : (U, P, \varphi) \to (U', P', \varphi')$
est donné par un morphisme de schémas $f : T \to T'$
et un isomorphisme $G$-équivariant
$g : P \to f^{-1}P'$ tel que $\varphi = \varphi' \circ g$.
\end{enumerate}
Exactement comme ci-dessus, nous obtenons un foncteur
$$
[[X/G]] \longrightarrow (\Sch/S)_{fppf}
$$
qui fait de $[[X/G]]$ un champ en groupoïdes sur $(\Sch/S)_{fppf}$.
Le faisceau constant $\underline{G}$ est représentable (pourvu que le cardinal de $G$ ne soit
pas trop grand) par $G_S$ sur $(\Sch/S)_{fppf}$,
et cette version de $[[X/G]]$ est équivalente au champ
$[[X/G_S]]$ introduit ci-dessus.
\end{remark}






\section{Le champ de Picard}
\label{examples-stacks-section-picard-stack}

\noindent
Dans cette section, nous introduisons le champ de Picard en toute généralité.
Dans le chapitre sur Quot et Hilb, nous montrerons qu'il s'agit, sous des hypothèses
convenables, d'un champ algébrique, voir
Quot, section \ref{quot-section-picard-stack}.

\medskip\noindent
Soit $S$ un schéma.
Soit $\pi : X \to B$ un morphisme d'espaces algébriques sur $S$.
Nous définissons une catégorie $\Picardstack_{X/B}$ comme suit :
\begin{enumerate}
\item Un objet est un triplet $(U, b, \mathcal{L})$, où
\begin{enumerate}
\item $U$ est un objet de $(\Sch/S)_{fppf}$ ;
\item $b : U \to B$ est un morphisme sur $S$ ;
\item $\mathcal{L}$ est un faisceau inversible sur le changement de base
$X_U = U \times_{b, B} X$.
\end{enumerate}
\item Un morphisme $(f, g) : (U, b, \mathcal{L}) \to (U', b', \mathcal{L}')$
est donné par un morphisme de schémas $f : U \to U'$ sur $B$ et un
isomorphisme $g : f^*\mathcal{L}' \to \mathcal{L}$.
\end{enumerate}
Le composé de
$(f, g) : (U, b, \mathcal{L}) \to (U', b', \mathcal{L}')$
avec
$(f', g') : (U', b', \mathcal{L}') \to (U'', b'', \mathcal{L}'')$
est $(f' \circ f, g \circ f^*(g'))$.
Nous obtenons ainsi une catégorie $\Picardstack_{X/B}$ et
$$
p : \Picardstack_{X/B} \longrightarrow (\Sch/S)_{fppf},
\quad
(U, b, \mathcal{L}) \longmapsto U
$$
est un foncteur. La catégorie fibre de $\Picardstack_{X/B}$ au-dessus de $U$
est la réunion disjointe, lorsque $b \in \Mor_S(U, B)$ varie, des catégories
de faisceaux inversibles sur $X_U = U \times_{b, B} X$. Les catégories
fibres sont donc des groupoïdes.

\begin{lemma}
\label{examples-stacks-lemma-picard-stack}
Moyennant un remplacement comme dans
Champs, remarque \ref{stacks-remark-stack-make-small},
le foncteur
$$
\Picardstack_{X/B} \longrightarrow (\Sch/S)_{fppf}
$$
définit un champ en groupoïdes sur $(\Sch/S)_{fppf}$.
\end{lemma}

\begin{proof}
Comme d'habitude, la partie la plus difficile consiste à montrer l'effectivité de la descente des objets.
Soit pour cela $\{U_i \to U\}$ un recouvrement de $(\Sch/S)_{fppf}$.
Soit $\xi_i = (U_i, b_i, \mathcal{L}_i)$ un objet de
$\Picardstack_{X/B}$ au-dessus de $U_i$, et soit
$\varphi_{ij} : \text{pr}_0^*\xi_i \to \text{pr}_1^*\xi_j$
une donnée de descente. Elle implique en particulier que les morphismes
$b_i$ sont les restrictions d'un morphisme $b : U \to B$.
Écrivons $X_U = U \times_{b, B} X$ et
$X_i = U_i \times_{b_i, B} X =
U_i \times_U U \times_{b, B} X = U_i \times_U X_U$.
Le faisceau $\mathcal{L}_i$ est un $\mathcal{O}_{X_i}$-module inversible.
La famille $\{X_i \to X_U\}$ est elle aussi un recouvrement fppf.
De plus, la donnée de descente $\varphi_{ij}$ fournit une
donnée de descente sur les faisceaux inversibles $\mathcal{L}_i$ relativement
au recouvrement fppf $\{X_i \to X_U\}$.
Par conséquent, d'après
Descente sur les espaces,
proposition \ref{spaces-descent-proposition-fpqc-descent-quasi-coherent},
nous obtenons un unique faisceau inversible $\mathcal{L}$ sur $X_U$
qui redonne $\mathcal{L}_i$ et les données de descente sur $X_i$.
Le triplet $(U, b, \mathcal{L})$ est donc l'objet de
$\Picardstack_{X/B}$ au-dessus de $U$ que nous cherchions.
Nous omettons les détails.
\end{proof}





\section{Exemples de champs d'inertie}
\label{examples-stacks-section-examples-inertia}

\noindent
Voici quelques exemples de champs d'inertie.

\begin{example}
\label{examples-stacks-example-inertia-stack-of-X-mod-G}
Soit $S$ un schéma. Soit $G$ un groupe commutatif.
Soit $X \to S$ un schéma sur $S$.
Soit $a : G \times X \to X$ une action de $G$ sur $X$.
Pour $g \in G$, nous notons $g : X \to X$ l'automorphisme correspondant.
Dans ce cas, le champ d'inertie de $[X/G]$ (voir
remarque \ref{examples-stacks-remark-X-mod-G-group})
est donné par
$$
I_{[X/G]} = \coprod\nolimits_{g\in G} [X^g/G],
$$
où, pour un élément $g$ de $G$, le symbole $X^g$ désigne le
schéma $X^g = \{x \in X \mid g(x) = x\}$. Plus précisément,
$X^g$ est le produit
fibré
$$
X^g =  X \times_{(1, 1), X \times_S X, (g, 1)} X.
$$
En effet, pour tout $S$-schéma $T$, un
point à valeurs dans $T$ du champ d'inertie de $[X/G]$ consiste en un
triplet $(P/T, \phi, \alpha)$ formé d'un $G$-torseur fppf
$P \to T$, d'un morphisme $G$-équivariant
$\phi : P \to X$ et
d'un automorphisme $\alpha$ de $P$ sur $T$ tel que
$\phi \circ \alpha = \phi$.
Puisque $G$ est un faisceau de groupes \emph{commutatifs},
$\alpha$ est donné, localement pour la topologie fppf sur $T$,
par la multiplication par un élément $g$ de $G$.
La condition $\phi \circ \alpha = \phi$ signifie que $\phi$
se factorise par l'inclusion de $X^g$
dans $X$, c'est-à-dire que $\phi$ est le composé de cette inclusion avec un
morphisme $P \to X^g$.
La discussion précédente permet de définir un morphisme de catégories fibrées
$I_{[X/G]} \to \coprod_{g\in G} [X^g/G]$ sur les points à valeurs dans $T$ comme ci-dessus.
Nous omettons la vérification qu'il s'agit d'une équivalence.
\end{example}

\begin{example}
\label{examples-stacks-example-inertia-stack-of-picard}
Soit $f : X \to S$ un morphisme de schémas.
Supposons que, pour tout $T \to S$, le changement de base $f_T : X_T \to T$
ait la propriété que le morphisme $\mathcal{O}_T \to f_{T, *}\mathcal{O}_{X_T}$
est un isomorphisme. (Cela implique que $f$ est
{\it cohomologiquement plat en dimension $0$} (insérer ici une référence ultérieure),
mais c'est une condition plus forte.) Considérons le champ de Picard $\Picardstack_{X/S}$, voir
section \ref{examples-stacks-section-picard-stack}.
Les points de son champ d'inertie au-dessus d'un
$S$-schéma $T$ sont les couples $(\mathcal{L}, \alpha)$,
où $\mathcal{L}$ est un faisceau inversible
sur $X_T$ et $\alpha$ un automorphisme de ce faisceau inversible.
Autrement dit, nous pouvons considérer $\alpha$ comme un élément de
$H^0(X_T, \mathcal{O}_{X_T})^\times = H^0(T, \mathcal{O}_T^*)$
d'après notre hypothèse. Or $H^0(T, \mathcal{O}_T^*) = \mathbf{G}_{m, S}(T)$,
voir Groupoïdes, exemple \ref{groupoids-example-multiplicative-group}.
Le champ d'inertie de $\Picardstack_{X/S}$ est donc
$$
I_{\Picardstack_{X/S}} = \mathbf{G}_{m, S} \times_S \Picardstack_{X/S}
$$
en tant que champ sur $(\Sch/S)_{fppf}$.
\end{example}






\section{Champs de Hilbert finis}
\label{examples-stacks-section-hilbert-d-stack}

\noindent
Nous formulons ce qui suit dans une généralité peut-être un peu plus grande
que strictement nécessaire. Fixons un $1$-morphisme
$$
F : \mathcal{X} \longrightarrow \mathcal{Y}
$$
de champs en groupoïdes sur $(\Sch/S)_{fppf}$. Pour chaque entier
$d \geq 1$, considérons une catégorie $\mathcal{H}_d(\mathcal{X}/\mathcal{Y})$
définie comme suit :
\begin{enumerate}
\item Un objet est un quintuplet $(U, Z, y, x, \alpha)$, où $U, Z$ sont des objets de
$(\Sch/S)_{fppf}$ et où $Z$ est fini et localement libre de degré
$d$ sur $U$, avec
$y \in \Ob(\mathcal{Y}_U)$, $x \in \Ob(\mathcal{X}_Z)$
et un isomorphisme $\alpha : y|_Z \to F(x)$\footnote{
Cela signifie que ces données fournissent, par le $2$-lemme de Yoneda
(Catégories, lemme \ref{categories-lemma-yoneda-2category}), un
diagramme $2$-commutatif
$$
\xymatrix{
(\Sch/Z)_{fppf} \ar[r]_-x \ar[d] & \mathcal{X} \ar[d]^F \\
(\Sch/U)_{fppf} \ar[r]^-y & \mathcal{Y}
}
$$
de champs en groupoïdes sur $(\Sch/S)_{fppf}$.
On peut aussi représenter $\alpha$ comme un $2$-morphisme
$$
\xymatrix{
(\Sch/Z)_{fppf}
\rrtwocell^{y \circ (Z \to U)}_{F \circ x}{\alpha} & &
\mathcal{Y}.
}
$$
}.
\item Un morphisme $(U, Z, y, x, \alpha) \to (U', Z', y', x', \alpha')$ est
donné par un morphisme de schémas $f : U \to U'$, un morphisme de schémas
$g : Z \to Z'$ qui induit un isomorphisme $Z \to Z' \times_U U'$,
et des isomorphismes $b : y \to f^*y'$, $a : x \to g^*x'$ induisant un diagramme
commutatif
$$
\xymatrix{
y|_Z \ar[rr]_\alpha \ar[d]_{b|_Z} & &
F(x) \ar[d]^{F(a)} \\
f^*y'|_Z \ar[rr]^{\alpha'} & &
F(g^*x') \\
}
$$
\end{enumerate}
Les définitions montrent immédiatement qu'il existe un foncteur d'oubli
canonique
$$
p :
\mathcal{H}_d(\mathcal{X}/\mathcal{Y})
\longrightarrow
(\Sch/S)_{fppf}
$$
qui associe au quintuplet $(U, Z, y, x, \alpha)$ le schéma $U$
et au morphisme
$(f, g, b, a) : (U, Z, y, x, \alpha) \to (U', Z', y', x', \alpha')$
le morphisme $f : U \to U'$.

\begin{lemma}
\label{examples-stacks-lemma-hilbert-d-stack}
La catégorie $\mathcal{H}_d(\mathcal{X}/\mathcal{Y})$, munie
du foncteur $p$ ci-dessus, définit un champ en groupoïdes sur
$(\Sch/S)_{fppf}$.
\end{lemma}

\begin{proof}
Comme d'habitude, la partie la plus difficile consiste à montrer l'effectivité de la descente des objets.
Soit pour cela $\{U_i \to U\}$ un recouvrement de $(\Sch/S)_{fppf}$.
Soit $\xi_i = (U_i, Z_i, y_i, x_i, \alpha_i)$ un objet de
$\mathcal{H}_d(\mathcal{X}/\mathcal{Y})$ au-dessus de $U_i$, et soit
$\varphi_{ij} : \text{pr}_0^*\xi_i \to \text{pr}_1^*\xi_j$
une donnée de descente. Remarquons d'abord que $\varphi_{ij}$
induit une donnée de descente $(Z_i/U_i, \varphi_{ij})$, qui est effective d'après
Descente, lemme \ref{descent-lemma-affine}.
On obtient ainsi un schéma $Z/U$ fini et localement libre de degré $d$ d'après
Descente, lemme \ref{descent-lemma-descending-property-finite-locally-free}.
Désormais, nous identifions $Z_i$ à $Z \times_U U_i$.
Ensuite, les objets $y_i$ des catégories fibres $\mathcal{Y}_{U_i}$
descendent en un objet $y$ de $\mathcal{Y}_U$, puisque $\mathcal{Y}$ est un
champ en groupoïdes. De même, les objets $x_i$ des catégories fibres
$\mathcal{X}_{Z_i}$ descendent en un objet $x$ de $\mathcal{X}_Z$, puisque
$\mathcal{X}$ est un champ en groupoïdes. Enfin, les isomorphismes donnés
$$
\alpha_i :
(y|_Z)_{Z_i} = y_i|_{Z_i}
\longrightarrow
F(x_i) = F(x|_{Z_i})
$$
se recollent en un morphisme $\alpha : y|_Z \to F(x)$, car $\mathcal{Y}$
est un champ et $\mathit{Isom}_\mathcal{Y}(y|_Z, F(x))$ est donc
un faisceau. Nous omettons les détails.
\end{proof}

\begin{definition}
\label{examples-stacks-definition-hilbert-d-stack}
Nous appellerons $\mathcal{H}_d(\mathcal{X}/\mathcal{Y})$
le {\it champ de Hilbert fini de degré $d$ de $\mathcal{X}$ sur $\mathcal{Y}$}
construit ci-dessus. Si $\mathcal{Y} = S$, nous écrivons
$\mathcal{H}_d(\mathcal{X}) = \mathcal{H}_d(\mathcal{X}/\mathcal{Y})$.
Si $\mathcal{X} = \mathcal{Y} = S$, nous le notons $\mathcal{H}_d$.
\end{definition}

\noindent
Étant donné $F : \mathcal{X} \to \mathcal{Y}$ comme ci-dessus, nous avons les
$1$-morphismes naturels suivants de champs en groupoïdes sur
$(\Sch/S)_{fppf}$ :
\begin{equation}
\label{examples-stacks-equation-diagram-hilbert-d-stack}
\vcenter{
\xymatrix{
\mathcal{H}_d(\mathcal{X}) \ar[rd] &
\mathcal{H}_d(\mathcal{X}/\mathcal{Y}) \ar[d] \ar[l] \ar[r] &
\mathcal{Y} \\
& \mathcal{H}_d
}
}
\end{equation}
Chacune des flèches est donnée par un « foncteur d'oubli ».

\begin{lemma}
\label{examples-stacks-lemma-faithful-hilbert}
Le $1$-morphisme
$\mathcal{H}_d(\mathcal{X}/\mathcal{Y}) \to \mathcal{H}_d(\mathcal{X})$
est fidèle.
\end{lemma}

\begin{proof}
Pour vérifier que
$\mathcal{H}_d(\mathcal{X}/\mathcal{Y}) \to \mathcal{H}_d(\mathcal{X})$
est fidèle, il suffit de le vérifier sur les catégories fibres.
Soient $\xi = (U, Z, y, x, \alpha)$ et $\xi' = (U, Z', y', x', \alpha')$
deux objets de $\mathcal{H}_d(\mathcal{X}/\mathcal{Y})$ au-dessus du
schéma $U$. Soient $(g, b, a), (g', b', a') : \xi \to \xi'$ deux morphismes
de la catégorie fibre de $\mathcal{H}_d(\mathcal{X}/\mathcal{Y})$ au-dessus de $U$.
Les images de ces morphismes dans $\mathcal{H}_d(\mathcal{X})$ coïncident
si et seulement si $g = g'$ et $a = a'$. Le diagramme commutatif
$$
\xymatrix{
y|_Z \ar[rr]_\alpha \ar[d]_{b|_Z, \ b'|_Z} & &
F(x) \ar[d]^{F(a) = F(a')} \\
y'|_Z \ar[rr]^-{\alpha'} & &
F(g^*x') = F((g')^*x') \\
}
$$
implique alors que $b|_Z = b'|_Z$. Puisque $Z \to U$ est fini et localement libre de degré
$d$, la famille $\{Z \to U\}$ est un recouvrement fppf ; ainsi $b = b'$.
\end{proof}
















% OK, start here.
%

\chapter{Faisceaux sur les champs algébriques}



\phantomsection
\label{stacks-sheaves-section-phantom}


\section{Introduction}
\label{stacks-sheaves-section-introduction}

\noindent
Il existe une multitude de façons d'envisager les faisceaux sur les champs algébriques.
Dans ce chapitre, nous présentons une approche particulièrement bien
adaptée à nos fondements pour les champs algébriques. Chaque fois que nous introduirons
un type de faisceaux, nous préciserons sa relation exacte avec les
notions analogues de la littérature.
Le but de ce chapitre est d'énoncer les résultats
qui sont soit évidents, soit faciles à démontrer,
et de réserver les constructions plus délicates pour plus tard.

\medskip\noindent
En fait, il s'avère que, pour élaborer une théorie pleinement développée des
faisceaux \'etales constructibles, ou encore une étude adéquate des
catégories dérivées de complexes de $\mathcal{O}$-modules dont les
faisceaux de cohomologie sont quasi-cohérents, il faut un travail considérable ; voir
\cite{olsson_sheaves}. Nous y reviendrons dans
Cohomologie des champs, section \ref{stacks-cohomology-section-introduction}.

\medskip\noindent
Dans la littérature et les articles de recherche consacrés aux faisceaux sur les champs
algébriques, le site lisse-\'etale d'un champ algébrique joue souvent un rôle majeur.
Il s'agit toutefois d'un objet problématique : un morphisme de
champs algébriques n'induit pas, en général, de morphisme de topos lisses-\'etales. Nous avons
donc choisi d'éviter toute mention du site lisse-\'etale
aussi longtemps que possible. Les arguments qui utilisent traditionnellement le site lisse-\'etale
seront remplacés par des arguments faisant appel à un recouvrement de {\v C}ech
dans le site $\mathcal{X}_{smooth}$ défini ci-dessous.

\medskip\noindent
Certaines notations, conventions et certains termes de ce chapitre sont peu commodes
et pourront sembler contraires aux usages au lecteur expérimenté. C'est intentionnel.
Voir Quotients, section \ref{quot-section-conventions}, pour une
explication.




\section{Conventions}
\label{stacks-sheaves-section-conventions}

\noindent
Les conventions employées dans ce chapitre sont les mêmes que dans le
chapitre sur les champs algébriques ; voir
Champs algébriques, section \ref{algebraic-section-conventions}.
Nous les répétons ici par commodité.

\medskip\noindent
Nous travaillons dans un grand site fppf convenable $\Sch_{fppf}$, comme dans
Topologies, Définition \ref{topologies-definition-big-fppf-site}.
Ainsi, sauf indication contraire explicite, tous les schémas seront des objets
de $\Sch_{fppf}$. Nous indiquerons ailleurs ce qui change lorsqu'on remplace le
grand site fppf (insérer ici une référence future).

\medskip\noindent
Nous travaillerons toujours relativement à une base $S$ appartenant à $\Sch_{fppf}$.
Nous utiliserons alors le grand site fppf $(\Sch/S)_{fppf}$ ; voir
Topologies, Définition \ref{topologies-definition-big-small-fppf}.
On retrouve le cas absolu en prenant
$S = \Spec(\mathbf{Z})$.





\section{Préfaisceaux}
\label{stacks-sheaves-section-presheaves}

\noindent
Dans cette section, nous définissons les préfaisceaux sur les catégories fibrées en
groupoïdes au-dessus de $(\Sch/S)_{fppf}$, mais l'essentiel de la discussion vaut
pour les catégories au-dessus d'une catégorie de base quelconque. Cette section sert aussi à
introduire les notations que nous emploierons par la suite.

\begin{definition}
\label{stacks-sheaves-definition-presheaves}
Soit $p : \mathcal{X} \to (\Sch/S)_{fppf}$ une catégorie fibrée en
groupoïdes.
\begin{enumerate}
\item Un {\it préfaisceau sur $\mathcal{X}$} est un préfaisceau sur la
catégorie sous-jacente à $\mathcal{X}$.
\item Un {\it morphisme de préfaisceaux sur $\mathcal{X}$} est un morphisme de
préfaisceaux sur la catégorie sous-jacente à $\mathcal{X}$.
\end{enumerate}
Nous notons $\textit{PSh}(\mathcal{X})$ la catégorie des préfaisceaux sur
$\mathcal{X}$.
\end{definition}

\noindent
Cela définit les préfaisceaux d'ensembles. On peut bien entendu parler aussi de
préfaisceaux d'ensembles pointés, de groupes abéliens, de groupes, de monoïdes, d'anneaux,
de modules sur un anneau fixé, d'algèbres de Lie sur un corps fixé, etc.
La catégorie des {\it préfaisceaux abéliens}, c'est-à-dire des préfaisceaux de groupes
abéliens, est notée $\textit{PAb}(\mathcal{X})$.

\medskip\noindent
Soit $f : \mathcal{X} \to \mathcal{Y}$ un $1$-morphisme de catégories
fibrées en groupoïdes au-dessus de $(\Sch/S)_{fppf}$. Rappelons que cela
signifie simplement que $f$ est un foncteur au-dessus de $(\Sch/S)_{fppf}$.
Les résultats de
Sites, section \ref{sites-section-more-functoriality-PSh},
nous fournissent une paire de foncteurs adjoints\footnote{Ces foncteurs
seront notés $f^{-1}$ et $f_*$ une fois le
Lemme \ref{stacks-sheaves-lemma-functoriality-sheaves}
démontré.}
\begin{equation}
\label{stacks-sheaves-equation-pushforward-pullback}
f^p : \textit{PSh}(\mathcal{Y}) \longrightarrow \textit{PSh}(\mathcal{X})
\quad\text{et}\quad
{}_pf : \textit{PSh}(\mathcal{X}) \longrightarrow \textit{PSh}(\mathcal{Y}).
\end{equation}
L'adjonction s'écrit
$$
\Mor_{\textit{PSh}(\mathcal{X})}(f^p\mathcal{G}, \mathcal{F})
=
\Mor_{\textit{PSh}(\mathcal{Y})}(\mathcal{G}, {}_pf\mathcal{F})
$$
où $\mathcal{F} \in \Ob(\textit{PSh}(\mathcal{X}))$ et
$\mathcal{G} \in \Ob(\textit{PSh}(\mathcal{Y}))$. Nous appelons
$f^p\mathcal{G}$ l'{\it image réciproque} de $\mathcal{G}$. Il résulte
des définitions que
$$
f^p\mathcal{G}(x) = \mathcal{G}(f(x))
$$
pour tout $x \in \Ob(\mathcal{X})$. Le préfaisceau ${}_pf\mathcal{F}$
est appelé l'{\it image directe} de $\mathcal{F}$. Il est décrit
par la formule
$$
({}_pf\mathcal{F})(y) = \lim_{f(x) \to y} \mathcal{F}(x).
$$
Le reste de cette section devrait probablement être déplacé dans le chapitre
sur les sites et peut, en tout cas, être omis en première lecture.

\begin{lemma}
\label{stacks-sheaves-lemma-1-morphisms-presheaves}
Soient $f : \mathcal{X} \to \mathcal{Y}$ et $g : \mathcal{Y} \to \mathcal{Z}$
des $1$-morphismes de catégories fibrées en groupoïdes au-dessus de
$(\Sch/S)_{fppf}$. Alors $(g \circ f)^p = f^p \circ g^p$ et
il existe un isomorphisme canonique
${}_p(g \circ f) \to {}_pg \circ {}_pf$
compatible aux adjonctions de $(f^p, {}_pf)$, $(g^p, {}_pg)$ et
$((g \circ f)^p, {}_p(g \circ f))$.
\end{lemma}

\begin{proof}
Soit $\mathcal{H}$ un préfaisceau sur $\mathcal{Z}$. Alors
$(g \circ f)^p\mathcal{H} = f^p (g^p\mathcal{H})$ est donné
par les égalités
$$
(g \circ f)^p\mathcal{H}(x) = \mathcal{H}((g \circ f)(x))
= \mathcal{H}(g(f(x))) = f^p (g^p\mathcal{H})(x).
$$
Nous omettons de vérifier la compatibilité avec les applications de restriction.

\medskip\noindent
Définissons ensuite la transformation ${}_p(g \circ f) \to {}_pg \circ {}_pf$.
Soit $\mathcal{F}$ un préfaisceau sur $\mathcal{X}$.
Si $z$ est un objet de $\mathcal{Z}$, on obtient une
catégorie $\mathcal{J}$ de quadruplets
$(x, f(x) \to y, y, g(y) \to z)$ et une catégorie $\mathcal{I}$
de couples $(x, g(f(x)) \to z)$. Il existe un foncteur canonique
$\mathcal{J} \to \mathcal{I}$ qui envoie l'objet
$(x, \alpha : f(x) \to y, y, \beta : g(y) \to z)$ sur
$(x, \beta \circ f(\alpha) : g(f(x)) \to z)$. Il fournit la flèche dans
\begin{align*}
({}_p(g \circ f)\mathcal{F})(z) & =
\lim_{g(f(x)) \to z} \mathcal{F}(x) \\
& = \lim_\mathcal{I} \mathcal{F} \\
& \to \lim_\mathcal{J} \mathcal{F} \\
& = \lim_{g(y) \to z}
\Big(\lim_{f(x) \to y} \mathcal{F}(x)\Big) \\
& =
({}_pg \circ {}_pf\mathcal{F})(x)
\end{align*}
par
Catégories, Lemme \ref{categories-lemma-functorial-limit}.
Nous omettons de vérifier la compatibilité avec les applications de restriction.
Au lieu de cette construction directe, on peut définir
${}_p(g \circ f) \cong {}_pg \circ {}_pf$
comme l'unique application compatible aux propriétés d'adjonction. Cette méthode
présente aussi l'avantage de dispenser de démontrer la compatibilité.

\medskip\noindent
La compatibilité aux adjonctions de $(f^p, {}_pf)$, $(g^p, {}_pg)$ et
$((g \circ f)^p, {}_p(g \circ f))$ signifie que, pour les préfaisceaux
$\mathcal{H}$ et $\mathcal{F}$ ci-dessus, le diagramme suivant est commutatif :
$$
\xymatrix{
\Mor_{\textit{PSh}(\mathcal{X})}(f^pg^p\mathcal{H}, \mathcal{F})
\ar@{=}[r] \ar@{=}[d] &
\Mor_{\textit{PSh}(\mathcal{Y})}(g^p\mathcal{H}, {}_pf\mathcal{F})
\ar@{=}[r] &
\Mor_{\textit{PSh}(\mathcal{Y})}(\mathcal{H}, {}_pg{}_pf\mathcal{F})
\\
\Mor_{\textit{PSh}(\mathcal{X})}((g \circ f)^p\mathcal{G}, \mathcal{F})
\ar@{=}[rr] & &
\Mor_{\textit{PSh}(\mathcal{Y})}(\mathcal{G}, {}_p(g \circ f)\mathcal{F})
\ar[u]
}
$$
Démonstration omise.
\end{proof}

\begin{lemma}
\label{stacks-sheaves-lemma-2-morphisms-presheaves}
Soient $f, g : \mathcal{X} \to \mathcal{Y}$ des $1$-morphismes de catégories
fibrées en groupoïdes au-dessus de $(\Sch/S)_{fppf}$. Soit $t : f \to g$
un $2$-morphisme de catégories fibrées en groupoïdes au-dessus de
$(\Sch/S)_{fppf}$. À $t$ sont associés des
isomorphismes canoniques de foncteurs
$$
t^p : g^p \longrightarrow f^p
\quad\text{et}\quad
{}_pt : {}_pf \longrightarrow {}_pg
$$
compatibles aux adjonctions de $(f^p, {}_pf)$ et
$(g^p, {}_pg)$, ainsi qu'aux
compositions verticale et horizontale des $2$-morphismes.
\end{lemma}

\begin{proof}
Soit $\mathcal{G}$ un préfaisceau sur $\mathcal{Y}$. Alors
$t^p : g^p\mathcal{G} \to f^p\mathcal{G}$ est donné par la famille
d'applications
$$
g^p\mathcal{G}(x) = \mathcal{G}(g(x))
\xrightarrow{\mathcal{G}(t_x)}
\mathcal{G}(f(x)) = f^p\mathcal{G}(x)
$$
paramétrée par $x \in \Ob(\mathcal{X})$. Cela a un sens puisque
$t_x : f(x) \to g(x)$ et que $\mathcal{G}$ est un foncteur contravariant.
Nous omettons de vérifier la compatibilité avec les applications
de restriction.

\medskip\noindent
Pour définir la transformation ${}_pt$, posons, pour $y \in \Ob(\mathcal{Y})$,
que ${}_y^f\mathcal{I}$, resp.\ ${}_y^g\mathcal{I}$, est la catégorie
des couples $(x, \psi : f(x) \to y)$, resp.\ $(x, \psi : g(x) \to y)$ ; voir
Sites, section \ref{sites-section-more-functoriality-PSh}.
Remarquons que $t$ définit un foncteur
${}_yt : {}_y^g\mathcal{I} \to {}_y^f\mathcal{I}$
donné par la règle
$$
(x, g(x) \to y) \longmapsto (x, f(x) \xrightarrow{t_x} g(x) \to y).
$$
Pour un préfaisceau $\mathcal{F}$ sur $\mathcal{X}$, remarquons que le composé
de ${}_yt$ avec $\mathcal{F} : {}_y^f\mathcal{I}^{opp} \to \textit{Ensembles}$,
$(x, f(x) \to y) \mapsto \mathcal{F}(x)$, est égal à
$\mathcal{F} : {}_y^g\mathcal{I}^{opp} \to \textit{Ensembles}$. Ainsi, d'après
Catégories, Lemme \ref{categories-lemma-functorial-limit},
nous obtenons pour tout $y \in \Ob(\mathcal{Y})$ une application canonique
$$
({}_pf\mathcal{F})(y) = \lim_{{}_y^f\mathcal{I}} \mathcal{F}
\longrightarrow
\lim_{{}_y^g\mathcal{I}} \mathcal{F} = ({}_pg\mathcal{F})(y)
$$
Nous omettons de vérifier la compatibilité avec les applications
de restriction. Au lieu de cette construction directe, on peut définir
${}_pt$ comme l'unique application compatible aux propriétés d'adjonction
des couples $(f^p, {}_pf)$ et $(g^p, {}_pg)$ (voir ci-dessous). Cette méthode
présente aussi l'avantage de dispenser de démontrer la compatibilité.

\medskip\noindent
La compatibilité aux adjonctions de $(f^p, {}_pf)$ et $(g^p, {}_pg)$ signifie
que, pour les préfaisceaux $\mathcal{G}$ et $\mathcal{F}$ ci-dessus, nous avons
un diagramme commutatif
$$
\xymatrix{
\Mor_{\textit{PSh}(\mathcal{X})}(f^p\mathcal{G}, \mathcal{F})
\ar@{=}[r] \ar[d]_{- \circ t^p} &
\Mor_{\textit{PSh}(\mathcal{Y})}(\mathcal{G}, {}_pf\mathcal{F})
\ar[d]^{{}_pt \circ -} \\
\Mor_{\textit{PSh}(\mathcal{X})}(g^p\mathcal{G}, \mathcal{F})
\ar@{=}[r] &
\Mor_{\textit{PSh}(\mathcal{Y})}(\mathcal{G}, {}_pg\mathcal{F})
}
$$
Démonstration omise. Indication : reprendre la démonstration de
Sites, Lemme \ref{sites-lemma-adjoints-pu},
et déduire la compatibilité de la description explicite des
applications horizontales et verticales du diagramme.

\medskip\noindent
Nous omettons de vérifier la compatibilité avec les compositions verticale et
horizontale. Indication : pour $t^p$, la démonstration est immédiate, et
on peut en déduire que cela vaut pour les applications ${}_pt$ par compatibilité
à l'adjonction.
\end{proof}







\section{Faisceaux}
\label{stacks-sheaves-section-sheaves}

\noindent
Commençons par une observation importante et triviale
(surtout pour les lecteurs que les questions de théorie des ensembles
ne préoccupent pas).

\medskip\noindent
Considérons un grand site fppf $\Sch_{fppf}$ comme dans
Topologies, Définition \ref{topologies-definition-big-fppf-site},
et notons $\Sch_\alpha$ sa catégorie sous-jacente.
Outre qu'elle est la catégorie sous-jacente à un site fppf,
la catégorie $\Sch_\alpha$ peut aussi servir de catégorie
sous-jacente à un grand site de Zariski, un grand site \'etale, un grand site lisse
et un grand site syntomique ; voir
Topologies, Remarque \ref{topologies-remark-choice-sites}.
Nous notons ces sites $\Sch_{Zar}$, $\Sch_\etale$,
$\Sch_{smooth}$ et $\Sch_{syntomic}$.
Dans cette situation, puisque nous avons défini
le grand site de Zariski $(\Sch/S)_{Zar}$ de $S$,
le grand site \'etale $(\Sch/S)_\etale$ de $S$,
le grand site lisse $(\Sch/S)_{smooth}$ de $S$,
le grand site syntomique $(\Sch/S)_{syntomic}$ de $S$ et
le grand site fppf $(\Sch/S)_{fppf}$ de $S$
comme les localisations (voir
Sites, section \ref{sites-section-localize})
$\Sch_{Zar}/S$, $\Sch_\etale/S$,
$\Sch_{smooth}/S$, $\Sch_{syntomic}/S$ et
$\Sch_{fppf}/S$
de ces grands sites (absolus), tous ont la
même catégorie sous-jacente, à savoir $\Sch_\alpha/S$.

\medskip\noindent
Il s'ensuit que, si l'on se donne une catégorie
$p : \mathcal{X} \to (\Sch/S)_{fppf}$ fibrée en groupoïdes, alors
$\mathcal{X}$ hérite de topologies de Zariski, \'etale, lisse, syntomique et
fppf ; voir
Champs, Définition \ref{stacks-definition-topology-inherited}.

\begin{definition}
\label{stacks-sheaves-definition-inherited-topologies}
Soit $\mathcal{X}$ une catégorie fibrée en groupoïdes au-dessus de
$(\Sch/S)_{fppf}$.
\begin{enumerate}
\item Le {\it site de Zariski associé}, noté $\mathcal{X}_{Zar}$,
est la structure de site sur $\mathcal{X}$ héritée de
$(\Sch/S)_{Zar}$.
\item Le {\it site \'etale associé}, noté $\mathcal{X}_\etale$,
est la structure de site sur $\mathcal{X}$ héritée de
$(\Sch/S)_\etale$.
\item Le {\it site lisse associé}, noté $\mathcal{X}_{smooth}$,
est la structure de site sur $\mathcal{X}$ héritée de
$(\Sch/S)_{smooth}$.
\item Le {\it site syntomique associé}, noté $\mathcal{X}_{syntomic}$,
est la structure de site sur $\mathcal{X}$ héritée de
$(\Sch/S)_{syntomic}$.
\item Le {\it site fppf associé}, noté $\mathcal{X}_{fppf}$,
est la structure de site sur $\mathcal{X}$ héritée de
$(\Sch/S)_{fppf}$.
\end{enumerate}
\end{definition}

\noindent
La discussion précédente donne un sens à cette définition. Si $\mathcal{X}$
est un champ algébrique, la littérature appelle $\mathcal{X}_{fppf}$ (ou un
site qui lui est équivalent) le {\it grand site fppf} de $\mathcal{X}$, et de même
pour les autres. Nous emploierons parfois cette terminologie afin de
distinguer cette construction d'autres constructions.

\begin{remark}
\label{stacks-sheaves-remark-ambiguity}
Nous n'employons cette notation que lorsque le symbole $\mathcal{X}$ désigne une
catégorie fibrée en groupoïdes, et non un schéma, un espace algébrique, etc.
Nous éviterons ainsi toute confusion avec le petit site \'etale d'un
schéma ou d'un espace algébrique, que l'on note $X_\etale$ (auquel cas
la lettre capitale est romaine et non calligraphique).
\end{remark}

\noindent
Ces topologies étant définies, nous pouvons préciser ce qu'est
un faisceau sur $\mathcal{X}$, c'est-à-dire définir les topos correspondants.

\begin{definition}
\label{stacks-sheaves-definition-sheaves}
Soit $\mathcal{X}$ une catégorie fibrée en groupoïdes au-dessus de
$(\Sch/S)_{fppf}$. Soit $\mathcal{F}$ un préfaisceau sur $\mathcal{X}$.
\begin{enumerate}
\item Nous disons que $\mathcal{F}$ est un {\it faisceau de Zariski}, ou un
{\it faisceau pour la topologie de Zariski}, si $\mathcal{F}$
est un faisceau sur le site de Zariski associé $\mathcal{X}_{Zar}$.
\item Nous disons que $\mathcal{F}$ est un {\it faisceau \'etale}, ou un
{\it faisceau pour la topologie \'etale}, si $\mathcal{F}$
est un faisceau sur le site \'etale associé $\mathcal{X}_\etale$.
\item Nous disons que $\mathcal{F}$ est un {\it faisceau lisse}, ou un
{\it faisceau pour la topologie lisse}, si $\mathcal{F}$
est un faisceau sur le site lisse associé $\mathcal{X}_{smooth}$.
\item Nous disons que $\mathcal{F}$ est un {\it faisceau syntomique}, ou un
{\it faisceau pour la topologie syntomique}, si $\mathcal{F}$
est un faisceau sur le site syntomique associé $\mathcal{X}_{syntomic}$.
\item Nous disons que $\mathcal{F}$ est un {\it faisceau fppf}, ou un {\it faisceau},
ou encore un {\it faisceau pour la topologie fppf}, si $\mathcal{F}$
est un faisceau sur le site fppf associé $\mathcal{X}_{fppf}$.
\end{enumerate}
Un morphisme de faisceaux est simplement un morphisme de préfaisceaux. Nous notons
ces catégories de faisceaux
$\Sh(\mathcal{X}_{Zar})$,
$\Sh(\mathcal{X}_\etale)$,
$\Sh(\mathcal{X}_{smooth})$,
$\Sh(\mathcal{X}_{syntomic})$ et
$\Sh(\mathcal{X}_{fppf})$.
\end{definition}

\noindent
On peut bien entendu parler aussi de faisceaux d'ensembles pointés, de groupes abéliens,
de groupes, de monoïdes, d'anneaux, de modules sur un anneau fixé et d'algèbres de Lie sur
un corps fixé, etc. La catégorie des {\it faisceaux abéliens}, c'est-à-dire des faisceaux
de groupes abéliens, est notée $\textit{Ab}(\mathcal{X}_{fppf})$,
et de même pour les autres topologies. Si $\mathcal{X}$ est un champ
algébrique, alors $\Sh(\mathcal{X}_{fppf})$ est équivalente (aux
questions ensemblistes près) à ce que la littérature appelle
la {\it catégorie des faisceaux sur le grand site fppf de $\mathcal{X}$}. Il en va de même
pour les autres topologies. Nous emploierons parfois cette terminologie afin de
distinguer cette construction d'autres constructions.

\medskip\noindent
Puisque les topologies sont énumérées par finesse croissante, nous avons
les inclusions strictement pleines suivantes :
$$
\Sh(\mathcal{X}_{fppf}) \subset
\Sh(\mathcal{X}_{syntomic}) \subset
\Sh(\mathcal{X}_{smooth}) \subset
\Sh(\mathcal{X}_\etale) \subset
\Sh(\mathcal{X}_{Zar}) \subset \textit{PSh}(\mathcal{X})
$$
Nous écrivons parfois
$\Sh(\mathcal{X}_{fppf}) = \Sh(\mathcal{X})$
et
$\textit{Ab}(\mathcal{X}_{fppf}) = \textit{Ab}(\mathcal{X})$
conformément à notre convention selon laquelle un faisceau sur $\mathcal{X}$
est un faisceau fppf sur $\mathcal{X}$.

\medskip\noindent
Dans ce cadre, la fonctorialité de ces topos est immédiate et
elle est, de plus, compatible aux foncteurs d'inclusion ci-dessus.

\begin{lemma}
\label{stacks-sheaves-lemma-functoriality-sheaves}
Soit $f : \mathcal{X} \to \mathcal{Y}$ un $1$-morphisme de catégories
fibrées en groupoïdes au-dessus de $(\Sch/S)_{fppf}$. Soit
$\tau \in \{Zar, \etale, smooth, syntomic, fppf\}$.
Les foncteurs ${}_pf$ et $f^p$ de (\ref{stacks-sheaves-equation-pushforward-pullback})
transforment les faisceaux pour $\tau$ en faisceaux pour $\tau$ et définissent un morphisme
de topos
$f : \Sh(\mathcal{X}_\tau) \to \Sh(\mathcal{Y}_\tau)$.
\end{lemma}

\begin{proof}
Cela résulte immédiatement de
Champs, Lemme \ref{stacks-lemma-topology-inherited-functorial}.
\end{proof}

\noindent
Autrement dit, l'image directe et l'image réciproque des préfaisceaux définies dans la
section \ref{stacks-sheaves-section-presheaves}
donnent aussi l'{\it image directe} et l'{\it image réciproque} des faisceaux pour $\tau$.
Compte tenu de ce qui précède, nous pouvons écrire $f^p = f^{-1}$
et ${}_pf = f_*$ sans risque de confusion.

\begin{definition}
\label{stacks-sheaves-definition-morphism}
Soit $f : \mathcal{X} \to \mathcal{Y}$ un morphisme de catégories
fibrées en groupoïdes au-dessus de $(\Sch/S)_{fppf}$. Nous notons
$$
f = (f^{-1}, f_*) :
\Sh(\mathcal{X}_{fppf})
\longrightarrow
\Sh(\mathcal{Y}_{fppf})
$$
le {\it morphisme de topos fppf associé} construit ci-dessus.
De même pour les topos de Zariski, \'etale, lisse et syntomique associés.
\end{definition}

\noindent
Comme expliqué dans
Sites, section \ref{sites-section-sheaves-algebraic-structures},
la même formule (appliquée au faisceau d'ensembles sous-jacent) définit
l'image directe et l'image réciproque des faisceaux (pour l'une de nos topologies)
d'ensembles pointés, de groupes abéliens, de groupes, de monoïdes, d'anneaux, de modules
sur un anneau fixé, d'algèbres de Lie sur un corps fixé, etc.








\section{Calcul de l'image directe}
\label{stacks-sheaves-section-pushforward}

\noindent
Soit $f : \mathcal{X} \to \mathcal{Y}$ un $1$-morphisme de catégories
fibrées en groupoïdes au-dessus de $(\Sch/S)_{fppf}$. Soit $\mathcal{F}$
un préfaisceau sur $\mathcal{X}$. Soit $y \in \Ob(\mathcal{Y})$.
On peut calculer $f_*\mathcal{F}(y)$ comme suit. Supposons que
$y$ soit au-dessus du schéma $V$ et, au moyen du lemme de $2$-Yoneda, considérons
$y$ comme un $1$-morphisme. Considérons la projection
$$
\text{pr} :
(\Sch/V)_{fppf} \times_{y, \mathcal{Y}} \mathcal{X}
\longrightarrow
\mathcal{X}
$$
On a alors une identification canonique
\begin{equation}
\label{stacks-sheaves-equation-pushforward}
f_*\mathcal{F}(y) = \Gamma\Big(
(\Sch/V)_{fppf} \times_{y, \mathcal{Y}} \mathcal{X},
\ \text{pr}^{-1}\mathcal{F}\Big)
\end{equation}
En effet, les objets du $2$-produit fibré sont les triplets
$(h : U \to V, x, f(x) \to h^*y)$. En omettant $h$ dans la
notation, cela équivaut à la donnée d'un objet
$x$ de $\mathcal{X}$ et d'un morphisme $\alpha : f(x) \to y$ de $\mathcal{Y}$.
Comme $f_*\mathcal{F}(y) = \lim_{f(x) \to y} \mathcal{F}(x)$ par définition,
l'égalité en résulte.

\medskip\noindent
Nous en déduisons le résultat de « changement de base » suivant pour les
images directes. Il est trivial et repose sur le fait que
nous utilisons des « grands » sites.

\begin{lemma}
\label{stacks-sheaves-lemma-base-change}
Soit $S$ un schéma. Soit
$$
\xymatrix{
\mathcal{Y}' \times_\mathcal{Y} \mathcal{X} \ar[r]_{g'} \ar[d]_{f'} &
\mathcal{X} \ar[d]^f \\
\mathcal{Y}' \ar[r]^g & \mathcal{Y}
}
$$
un diagramme $2$-cartésien de catégories fibrées en groupoïdes au-dessus de $S$.
On a alors un isomorphisme canonique
$$
g^{-1}f_*\mathcal{F} \longrightarrow f'_*(g')^{-1}\mathcal{F}
$$
fonctoriel en le préfaisceau $\mathcal{F}$ sur $\mathcal{X}$.
\end{lemma}

\begin{proof}
Pour tout objet $y'$ de $\mathcal{Y}'$ au-dessus de $V$,
il existe une équivalence
$$
(\Sch/V)_{fppf} \times_{g(y'), \mathcal{Y}} \mathcal{X}
=
(\Sch/V)_{fppf} \times_{y', \mathcal{Y}'}
(\mathcal{Y}' \times_\mathcal{Y} \mathcal{X})
$$
D'où, par (\ref{stacks-sheaves-equation-pushforward}), une bijection
$g^{-1}f_*\mathcal{F}(y') \to f'_*(g')^{-1}\mathcal{F}(y')$.
Nous omettons de vérifier la compatibilité avec les applications
de restriction.
\end{proof}

\noindent
Dans le cas d'un morphisme représentable de catégories fibrées en groupoïdes,
la formule (\ref{stacks-sheaves-equation-pushforward}) se simplifie. Nous conseillons au
lecteur d'omettre le reste de cette section.

\begin{lemma}
\label{stacks-sheaves-lemma-representable}
Soit $f : \mathcal{X} \to \mathcal{Y}$ un $1$-morphisme de catégories
fibrées en groupoïdes au-dessus de $(\Sch/S)_{fppf}$. Les conditions suivantes sont
équivalentes :
\begin{enumerate}
\item $f$ est représentable, et
\item pour tout $y \in \Ob(\mathcal{Y})$, le foncteur
$\mathcal{X}^{opp} \to \textit{Ensembles}$,
$x \mapsto \Mor_\mathcal{Y}(f(x), y)$
est représentable.
\end{enumerate}
\end{lemma}

\begin{proof}
Selon la discussion de
Champs algébriques, section \ref{algebraic-section-representable-morphism},
$f$ est représentable si et seulement si, pour tout
$y \in \Ob(\mathcal{Y})$
au-dessus de $U$, le $2$-produit fibré
$(\Sch/U)_{fppf} \times_{y, \mathcal{Y}} \mathcal{X}$
est représentable, c'est-à-dire de la forme $(\Sch/V_y)_{fppf}$ pour un
schéma $V_y$ au-dessus de $U$. Les objets de ce $2$-produit fibré sont les triplets
$(h : V \to U, x, \alpha : f(x) \to h^*y)$ où $\alpha$ est
au-dessus de $\text{id}_V$. En omettant $h$ dans la notation, cela
équivaut à la donnée d'un objet $x$ de $\mathcal{X}$ et d'un
morphisme $f(x) \to y$. Ainsi, le $2$-produit fibré est
représenté par $V_y$ et $f(x_y) \to y$, où $x_y$ est un objet
de $\mathcal{X}$ au-dessus de $V_y$, si et seulement si le foncteur de (2) est représenté
par $x_y$, l'objet universel étant une application $f(x_y) \to y$.
\end{proof}

\noindent
Soit
$$
\xymatrix{
\mathcal{X} \ar[rr]_f \ar[rd]_p & &  \mathcal{Y} \ar[ld]^q \\
& (\Sch/S)_{fppf}
}
$$
un $1$-morphisme de catégories fibrées en groupoïdes. Supposons $f$
représentable. Pour tout $y \in \Ob(\mathcal{Y})$, choisissons
un objet $u(y) \in \Ob(\mathcal{X})$ représentant le foncteur
$x \mapsto \Mor_\mathcal{Y}(f(x), y)$ du
Lemme \ref{stacks-sheaves-lemma-representable}
(ce qui est possible par l'axiome du choix).
Ces objets sont munis de morphismes canoniques $f(u(y)) \to y$ par
construction.
Pour tout morphisme $\beta : y' \to y$ de $\mathcal{Y}$, on obtient un unique
morphisme $u(\beta) : u(y') \to u(y)$ de $\mathcal{X}$ tel que le diagramme
$$
\xymatrix{
f(u(y')) \ar[d] \ar[rr]_{f(u(\beta))} & & f(u(y)) \ar[d] \\
y' \ar[rr] & & y
}
$$
soit commutatif. Autrement dit, $u : \mathcal{Y} \to \mathcal{X}$ est un foncteur.
On peut en fait en dire un peu plus. Supposons en effet que
$V' = q(y')$, $V = q(y)$, $U' = p(u(y'))$ et $U = p(u(y))$. Alors
$$
\xymatrix{
U' \ar[rr]_{p(u(\beta))} \ar[d] & & U \ar[d] \\
V' \ar[rr]^{q(\beta)} & & V
}
$$
est un carré cartésien. Cela tient à ce que $U' \to U$ représente
le changement de base
$(\Sch/V')_{fppf} \times_{y', \mathcal{Y}} \mathcal{X} \to
(\Sch/V)_{fppf} \times_{y, \mathcal{Y}} \mathcal{X}$
de $V' \to V$.

\begin{lemma}
\label{stacks-sheaves-lemma-representable-pushforward}
Soit $f : \mathcal{X} \to \mathcal{Y}$ un $1$-morphisme représentable de
catégories fibrées en groupoïdes au-dessus de $(\Sch/S)_{fppf}$. Soit
$\tau \in \{Zar, \etale, smooth, syntomic, fppf\}$.
Alors le foncteur $u : \mathcal{Y}_\tau \to \mathcal{X}_\tau$ est continu
et définit un morphisme de sites $\mathcal{X}_\tau \to \mathcal{Y}_\tau$
qui induit le même morphisme de topos
$\Sh(\mathcal{X}_\tau) \to \Sh(\mathcal{Y}_\tau)$
que le morphisme $f$ construit dans le
Lemme \ref{stacks-sheaves-lemma-functoriality-sheaves}.
De plus, $f_*\mathcal{F}(y) = \mathcal{F}(u(y))$ pour tout préfaisceau
$\mathcal{F}$ sur $\mathcal{X}$.
\end{lemma}

\begin{proof}
Soit $\{y_i \to y\}$ un recouvrement pour $\tau$ dans $\mathcal{Y}$. Par définition,
cela signifie simplement que $\{q(y_i) \to q(y)\}$ est un recouvrement pour $\tau$ de
schémas. La dernière remarque précédant le lemme montre que
$\{p(u(y_i)) \to p(u(y))\}$ est le changement de base du recouvrement pour $\tau$
$\{q(y_i) \to q(y)\}$ par $p(u(y)) \to q(y)$ ; c'est donc lui-même un
recouvrement pour $\tau$ par les axiomes d'un site. Ainsi $\{u(y_i) \to u(y)\}$
est un recouvrement pour $\tau$ de $\mathcal{X}$. Cela démontre que $u$ est
continu.

\medskip\noindent
Employons les notations $u_p, u_s, u^p, u^s$ de
Sites, sections \ref{sites-section-functoriality-PSh} et
\ref{sites-section-continuous-functors}.
Si nous démontrons la dernière assertion du lemme, alors
$f_* = u^p = u^s$ (par la continuité de $u$ établie ci-dessus) ; l'adjonction donne donc
$f^{-1} = u_s$, ce qui prouve que $u_s$ est exact, puis que $u$ détermine
un morphisme de sites ; l'égalité sera elle aussi claire.
Pour voir que $f_*\mathcal{F}(y) = \mathcal{F}(u(y))$, remarquons que, par
définition,
$$
f_*\mathcal{F}(y) = ({}_pf\mathcal{F})(y) =
\lim_{f(x) \to y} \mathcal{F}(x).
$$
Puisque $u(y)$ est un objet final de la catégorie sur laquelle la limite est
prise, on conclut.
\end{proof}





\section{Le faisceau structural}
\label{stacks-sheaves-section-structure-sheaf}

\noindent
Soit $\tau \in \{Zar, \etale, smooth, syntomic, fppf\}$.
Soit $p : \mathcal{X} \to (\Sch/S)_{fppf}$ une catégorie
fibrée en groupoïdes. La 2-catégorie des catégories fibrées en groupoïdes au-dessus de
$(\Sch/S)_{fppf}$ possède un objet final, à savoir
$\text{id} : (\Sch/S)_{fppf} \to (\Sch/S)_{fppf}$
et $p$ est un $1$-morphisme de $\mathcal{X}$ vers cet objet final.
Ainsi, tout préfaisceau $\mathcal{G}$ sur $(\Sch/S)_{fppf}$ fournit un
préfaisceau $p^{-1}\mathcal{G}$ sur $\mathcal{X}$ défini par la règle
$p^{-1}\mathcal{G}(x) = \mathcal{G}(p(x))$. De plus, la discussion de la
section \ref{stacks-sheaves-section-sheaves}
montre que $p^{-1}\mathcal{G}$ est un faisceau pour $\tau$ dès que
$\mathcal{G}$ est un faisceau pour $\tau$.

\medskip\noindent
Rappelons que le site $(\Sch/S)_{fppf}$ est un site annelé,
de faisceau structural $\mathcal{O}$ défini par la règle
$$
(\Sch/S)^{opp} \longrightarrow \textit{Rings},
\quad
U/S \longmapsto \Gamma(U, \mathcal{O}_U)
$$
voir
Descente, Définition \ref{descent-definition-structure-sheaf}.

\begin{definition}
\label{stacks-sheaves-definition-structure-sheaf}
Soit $p : \mathcal{X} \to (\Sch/S)_{fppf}$ une catégorie
fibrée en groupoïdes. Le
{\it faisceau structural de $\mathcal{X}$} est le faisceau d'anneaux
$\mathcal{O}_\mathcal{X} = p^{-1}\mathcal{O}$.
\end{definition}

\noindent
Pour un objet $x$ de $\mathcal{X}$ au-dessus de $U$, nous avons
$\mathcal{O}_\mathcal{X}(x) = \mathcal{O}(U) = \Gamma(U, \mathcal{O}_U)$.
Il va sans dire que $\mathcal{O}_\mathcal{X}$ est aussi un faisceau de Zariski, \'etale,
lisse et syntomique ; chacun des sites
$\mathcal{X}_{Zar}$, $\mathcal{X}_\etale$, $\mathcal{X}_{smooth}$,
$\mathcal{X}_{syntomic}$ et $\mathcal{X}_{fppf}$ est donc un site annelé.
Cette construction est elle aussi fonctorielle.

\begin{lemma}
\label{stacks-sheaves-lemma-functoriality-structure-sheaf}
Soit $f : \mathcal{X} \to \mathcal{Y}$ un $1$-morphisme de catégories
fibrées en groupoïdes au-dessus de $(\Sch/S)_{fppf}$. Soit
$\tau \in \{Zar, \etale, smooth, syntomic, fppf\}$.
Il existe une identification canonique
$f^{-1}\mathcal{O}_\mathcal{Y} = \mathcal{O}_\mathcal{X}$
qui fait de
$f : \Sh(\mathcal{X}_\tau) \to \Sh(\mathcal{Y}_\tau)$
un morphisme de topos annelés.
\end{lemma}

\begin{proof}
Notons $p : \mathcal{X} \to (\Sch/S)_{fppf}$ et
$q : \mathcal{Y} \to (\Sch/S)_{fppf}$ les foncteurs structuraux.
Alors $p = q \circ f$, donc $p^{-1} = f^{-1} \circ q^{-1}$ d'après le
Lemme \ref{stacks-sheaves-lemma-1-morphisms-presheaves}. Puisque
$\mathcal{O}_\mathcal{X} = p^{-1}\mathcal{O}$ et
$\mathcal{O}_\mathcal{Y} = q^{-1}\mathcal{O}$
le résultat en découle.
\end{proof}

\begin{remark}
\label{stacks-sheaves-remark-flat}
Dans la situation du
Lemme \ref{stacks-sheaves-lemma-functoriality-structure-sheaf},
le morphisme de topos annelés
$f : \Sh(\mathcal{X}_\tau) \to \Sh(\mathcal{Y}_\tau)$
est plat, comme le montre l'égalité
$f^{-1}\mathcal{O}_\mathcal{X} = \mathcal{O}_\mathcal{Y}$.
Cela est quelque peu contre-intuitif, car, par exemple, une immersion
fermée de champs algébriques n'est généralement pas plate (comme morphisme de
champs algébriques).
Toutefois, il se produit exactement la même chose pour une immersion
fermée $i : X \to Y$ de schémas : dans ce cas, le morphisme associé
de grands sites pour $\tau$
$i : (\Sch/X)_\tau \to (\Sch/Y)_\tau$
est lui aussi plat.
\end{remark}




\section{Faisceaux de modules}
\label{stacks-sheaves-section-modules}

\noindent
Puisque nous disposons d'un faisceau structural, nous disposons de modules.

\begin{definition}
\label{stacks-sheaves-definition-modules}
Soit $\mathcal{X}$ une catégorie fibrée en groupoïdes au-dessus de
$(\Sch/S)_{fppf}$.
\begin{enumerate}
\item Un {\it préfaisceau de modules sur $\mathcal{X}$} est un
préfaisceau de $\mathcal{O}_\mathcal{X}$-modules. La catégorie des
préfaisceaux de modules est notée $\textit{PMod}(\mathcal{O}_\mathcal{X})$.
\item Nous disons qu'un préfaisceau de modules $\mathcal{F}$ est un
{\it $\mathcal{O}_\mathcal{X}$-module}, ou plus précisément un
{\it faisceau de $\mathcal{O}_\mathcal{X}$-modules}, si $\mathcal{F}$
est un faisceau fppf. La catégorie des $\mathcal{O}_\mathcal{X}$-modules
est notée $\textit{Mod}(\mathcal{O}_\mathcal{X})$.
\end{enumerate}
\end{definition}

\noindent
Dans la littérature, ces (pré)faisceaux de modules apparaissent comme des {\it (pré)faisceaux
de $\mathcal{O}_\mathcal{X}$-modules sur le grand site fppf de $\mathcal{X}$}.
Nous emploierons parfois cette terminologie pour distinguer ces
catégories d'autres catégories. Nous rencontrerons aussi des préfaisceaux de modules qui
sont des faisceaux pour les topologies de Zariski, \'etale, lisse ou syntomique
(sans être nécessairement des faisceaux fppf). Au besoin, nous les noterons
$\textit{Mod}(\mathcal{X}_\etale, \mathcal{O}_\mathcal{X})$
et de même pour les autres topologies.

\medskip\noindent
Étudions maintenant la fonctorialité, d'abord pour les préfaisceaux de modules. Soit
$$
\xymatrix{
\mathcal{X} \ar[rr]_f \ar[rd]_p & &  \mathcal{Y} \ar[ld]^q \\
& (\Sch/S)_{fppf}
}
$$
un $1$-morphisme de catégories fibrées en groupoïdes.
Les foncteurs $f^{-1}$ et $f_*$ sur les préfaisceaux abéliens s'étendent en des foncteurs
\begin{equation}
\label{stacks-sheaves-equation-functoriality-presheaves-modules}
f^{-1} :
\textit{PMod}(\mathcal{O}_\mathcal{Y})
\longrightarrow
\textit{PMod}(\mathcal{O}_\mathcal{X})
\quad\text{et}\quad
f_* :
\textit{PMod}(\mathcal{O}_\mathcal{X})
\longrightarrow
\textit{PMod}(\mathcal{O}_\mathcal{Y})
\end{equation}
Cela est immédiat pour $f^{-1}$, car
$f^{-1}\mathcal{G}(x) = \mathcal{G}(f(x))$, qui est un module sur
$\mathcal{O}_\mathcal{Y}(f(x)) = \mathcal{O}(q(f(x))) = \mathcal{O}(p(x)) =
\mathcal{O}_\mathcal{X}(x)$. On peut aussi le déduire de ce que
$f^{-1}\mathcal{O}_\mathcal{Y} = \mathcal{O}_\mathcal{X}$
et de ce que $f^{-1}$ commute aux limites (sur les préfaisceaux).
Comme $f_*$ est un adjoint à droite, il commute à toutes les limites
(sur les préfaisceaux), en particulier aux produits. Nous pouvons donc étendre
$f_*$ en un foncteur sur les préfaisceaux de modules comme dans la démonstration de
Modules sur les sites, Lemme \ref{sites-modules-lemma-pushforward-module}.
Nous affirmons que les foncteurs (\ref{stacks-sheaves-equation-functoriality-presheaves-modules})
forment une paire de foncteurs adjoints :
$$
\Mor_{\textit{PMod}(\mathcal{O}_\mathcal{X})}(
f^{-1}\mathcal{G}, \mathcal{F})
=
\Mor_{\textit{PMod}(\mathcal{O}_\mathcal{Y})}(
\mathcal{G}, f_*\mathcal{F}).
$$
Comme $f^{-1}\mathcal{O}_\mathcal{Y} = \mathcal{O}_\mathcal{X}$,
cela résulte de
Modules sur les sites, Lemme \ref{sites-modules-lemma-adjoint-push-pull-modules},
en munissant $\mathcal{X}$ et $\mathcal{Y}$ de la topologie
chaotique.

\medskip\noindent
Étudions ensuite la fonctorialité pour les modules, c'est-à-dire pour les faisceaux de modules
de la topologie fppf. Notons encore $f$ le morphisme induit de topos
annelés ; voir le
Lemme \ref{stacks-sheaves-lemma-functoriality-structure-sheaf}
(pour l'instant appliqué aux topologies fppf). Remarquons que les foncteurs
$f^{-1}$ et $f_*$ de (\ref{stacks-sheaves-equation-functoriality-presheaves-modules})
préservent les sous-catégories de faisceaux de modules ; voir le
Lemme \ref{stacks-sheaves-lemma-functoriality-sheaves}.
Il s'ensuit immédiatement que
\begin{equation}
\label{stacks-sheaves-equation-functoriality-sheaves-modules}
f^{-1} :
\textit{Mod}(\mathcal{O}_\mathcal{Y})
\longrightarrow
\textit{Mod}(\mathcal{O}_\mathcal{X})
\quad\text{et}\quad
f_* :
\textit{Mod}(\mathcal{O}_\mathcal{X})
\longrightarrow
\textit{Mod}(\mathcal{O}_\mathcal{Y})
\end{equation}
forment une paire de foncteurs adjoints :
$$
\Mor_{\textit{Mod}(\mathcal{O}_\mathcal{X})}(
f^{-1}\mathcal{G}, \mathcal{F})
=
\Mor_{\textit{Mod}(\mathcal{O}_\mathcal{Y})}(
\mathcal{G}, f_*\mathcal{F}).
$$
D'après l'unicité des adjoints, nous concluons que
$f^* = f^{-1}$, où $f^*$ est défini dans
Modules sur les sites, section \ref{sites-modules-section-functoriality-modules},
pour le morphisme de topos annelés $f$ ci-dessus. Nous aurions bien entendu pu
le voir directement, car
$f^*(-) = f^{-1}(-) \otimes_{f^{-1}\mathcal{O}_\mathcal{Y}}
\mathcal{O}_\mathcal{X}$, et puisque
$f^{-1}\mathcal{O}_\mathcal{Y} = \mathcal{O}_\mathcal{X}$.

\medskip\noindent
Il en va de même pour les faisceaux de modules des topologies de Zariski, \'etale, lisse et
syntomique.



\section{Catégories représentables}
\label{stacks-sheaves-section-representable}

\noindent
Dans cette courte section, nous comparons nos définitions à la situation
où les champs algébriques considérés sont représentables.

\begin{lemma}
\label{stacks-sheaves-lemma-compare-with-scheme}
Soit $S$ un schéma. Soit $\mathcal{X}$ une catégorie fibrée
en groupoïdes au-dessus de $(\Sch/S)$. Supposons que $\mathcal{X}$ soit représentable
par un schéma $X$. Pour $\tau \in \{Zar,\linebreak[0] \etale,\linebreak[0]
smooth,\linebreak[0] syntomic,\linebreak[0] fppf\}$
il existe une équivalence canonique
$$
(\mathcal{X}_\tau, \mathcal{O}_\mathcal{X}) =
((\Sch/X)_\tau, \mathcal{O}_X)
$$
de sites annelés.
\end{lemma}

\begin{proof}
Il suffit de choisir une équivalence
$(\Sch/X)_\tau \to \mathcal{X}$ de catégories fibrées en groupoïdes
au-dessus de $(\Sch/S)_{fppf}$ et d'utiliser la fonctorialité de
la construction $\mathcal{X} \leadsto \mathcal{X}_\tau$.
\end{proof}

\begin{lemma}
\label{stacks-sheaves-lemma-compare-with-morphism-of-schemes}
Soit $S$ un schéma. Soit $f : \mathcal{X} \to \mathcal{Y}$ un morphisme
de catégories fibrées en groupoïdes au-dessus de $S$.
Supposons que $\mathcal{X}$ et $\mathcal{Y}$ soient représentables par des schémas
$X$ et $Y$. Soit $f : X \to Y$ le morphisme de schémas correspondant
à $f$. Pour $\tau \in \{Zar,\linebreak[0] \etale,\linebreak[0]
smooth,\linebreak[0] syntomic,\linebreak[0] fppf\}$
le morphisme de topos annelés
$f : (\Sh(\mathcal{X}_\tau), \mathcal{O}_\mathcal{X}) \to
(\Sh(\mathcal{Y}_\tau), \mathcal{O}_\mathcal{Y})$
coïncide avec le morphisme de topos annelés
$f : (\Sh((\Sch/X)_\tau), \mathcal{O}_X) \to 
(\Sh((\Sch/Y)_\tau), \mathcal{O}_Y)$ au moyen des identifications du
Lemme \ref{stacks-sheaves-lemma-compare-with-scheme}.
\end{lemma}

\begin{proof}
Il suffit de dérouler les définitions.
\end{proof}




\section{Restriction}
\label{stacks-sheaves-section-restriction}


\noindent
Une observation triviale mais utile est que la localisation
d'une catégorie fibrée en groupoïdes en un objet
est équivalente au grand site du schéma au-dessus duquel il se trouve.

\begin{lemma}
\label{stacks-sheaves-lemma-localizing}
Soit $p : \mathcal{X} \to (\Sch/S)_{fppf}$ une catégorie fibrée
en groupoïdes. Soit $\tau \in \{Zar, \etale, smooth, syntomic, fppf\}$.
Soit $x \in \Ob(\mathcal{X})$ au-dessus de $U = p(x)$.
Le foncteur $p$ induit une équivalence de sites
$\mathcal{X}_\tau/x \to (\Sch/U)_\tau$.
\end{lemma}

\begin{proof}
Cas particulier de Champs, Lemme \ref{stacks-lemma-localizing}.
\end{proof}

\noindent
Nous utilisons le lemme précédent pour définir l'image réciproque et la restriction
d'un (pré)faisceau à un schéma.

\begin{definition}
\label{stacks-sheaves-definition-pullback}
Soit $p : \mathcal{X} \to (\Sch/S)_{fppf}$ une catégorie fibrée
en groupoïdes. Soit $x \in \Ob(\mathcal{X})$ au-dessus de $U = p(x)$.
Soit $\mathcal{F}$ un préfaisceau sur $\mathcal{X}$.
\begin{enumerate}
\item L'{\it image réciproque $x^{-1}\mathcal{F}$ de $\mathcal{F}$} est la
restriction $\mathcal{F}|_{(\mathcal{X}/x)}$, considérée comme un préfaisceau sur
$(\Sch/U)_{fppf}$ au moyen de l'équivalence
$\mathcal{X}/x \to (\Sch/U)_{fppf}$ du
Lemme \ref{stacks-sheaves-lemma-localizing}.
\item La {\it restriction de $\mathcal{F}$ à $U_\etale$}
est $x^{-1}\mathcal{F}|_{U_\etale}$, notée abusivement
$\mathcal{F}|_{U_\etale}$.
\end{enumerate}
\end{definition}

\noindent
Cette notation a un sens, car, à l'objet $x$, le lemme de $2$-Yoneda, voir
Champs algébriques, section \ref{algebraic-section-2-yoneda},
associe un $1$-morphisme $x : (\Sch/U)_{fppf} \to \mathcal{X}/x$
quasi-inverse de $p : \mathcal{X}/x \to (\Sch/U)_{fppf}$.
Ainsi, $x^{-1}\mathcal{F}$ est bien l'image réciproque de $\mathcal{F}$ par ce
$1$-morphisme. En particulier, d'après ce qui précède, si $\mathcal{F}$
est un faisceau (ou un faisceau de Zariski, \'etale, lisse ou syntomique), alors
$x^{-1}\mathcal{F}$ est un faisceau sur $(\Sch/U)_{fppf}$ (ou sur
$(\Sch/U)_{Zar}$, $(\Sch/U)_\etale$,
$(\Sch/U)_{smooth}$, $(\Sch/U)_{syntomic}$).

\medskip\noindent
Soit $p : \mathcal{X} \to (\Sch/S)_{fppf}$ une catégorie fibrée
en groupoïdes. Soit $\varphi : x \to y$ un morphisme de $\mathcal{X}$
au-dessus du morphisme de schémas $a : U \to V$.
Rappelons que $a$ induit un morphisme de petits sites \'etales
$a_{small} : U_\etale \to V_\etale$ ; voir
Cohomologie \'etale, section \ref{etale-cohomology-section-functoriality}.
Soit $\mathcal{F}$ un préfaisceau sur $\mathcal{X}$.
Soient $\mathcal{F}|_{U_\etale}$ et
$\mathcal{F}|_{V_\etale}$ les restrictions de $\mathcal{F}$
par $x$ et $y$. Il existe une application naturelle de {\it comparaison}
\begin{equation}
\label{stacks-sheaves-equation-comparison-push}
c_\varphi :
\mathcal{F}|_{V_\etale}
\longrightarrow
a_{small, *}(\mathcal{F}|_{U_\etale})
\end{equation}
de préfaisceaux sur $U_\etale$. En effet, si $V' \to V$ est \'etale,
posons $U' = V' \times_V U$ et définissons $c_\varphi$ sur les sections au-dessus de $V'$
par
$$
\xymatrix{
a_{small, *}(\mathcal{F}|_{U_\etale})(V') &
\mathcal{F}|_{U_\etale}(U') \ar@{=}[l] &
\mathcal{F}(x') \ar@{=}[l] \\
\mathcal{F}|_{V_\etale}(V') \ar@{=}[rr] \ar[u]^{c_\varphi} &
&
\mathcal{F}(y') \ar[u]_{\mathcal{F}(\varphi')}
}
$$
Ici, $\varphi' : x' \to y'$ est un morphisme de $\mathcal{X}$
s'inscrivant dans un diagramme commutatif
$$
\vcenter{
\xymatrix{
x' \ar[r] \ar[d]_{\varphi'} & x \ar[d]^\varphi \\
y' \ar[r] & y
}
}
\quad\text{au-dessus de}\quad
\vcenter{
\xymatrix{
U' \ar[r] \ar[d] & U \ar[d]^a \\
V' \ar[r] & V
}
}
$$
L'existence et l'unicité de $\varphi'$ résultent des axiomes
d'une catégorie fibrée en groupoïdes.
Nous omettons de vérifier que $c_\varphi$ ainsi défini est bien un morphisme
de préfaisceaux (c'est-à-dire compatible avec les applications de restriction) et qu'il
est fonctoriel en $\mathcal{F}$. Si $\mathcal{F}$ est un faisceau pour la
topologie \'etale, nous obtenons une application de {\it comparaison}
\begin{equation}
\label{stacks-sheaves-equation-comparison}
c_\varphi : a_{small}^{-1}(\mathcal{F}|_{V_\etale})
\longrightarrow
\mathcal{F}|_{U_\etale}
\end{equation}
que l'on note également $c_\varphi$, comme indiqué (c'est l'abus de notation usuel
qui consiste à ne pas distinguer des applications adjointes).

\begin{lemma}
\label{stacks-sheaves-lemma-comparison}
Soit $\mathcal{F}$ un faisceau \'etale sur $\mathcal{X} \to (\Sch/S)_{fppf}$.
\begin{enumerate}
\item Si $\varphi : x \to y$ et $\psi : y \to z$
sont des morphismes de $\mathcal{X}$ au-dessus de $a : U \to V$ et
$b : V \to W$, alors le composé
$$
a_{small}^{-1}(b_{small}^{-1} (\mathcal{F}|_{W_\etale}))
\xrightarrow{a_{small}^{-1}c_\psi}
a_{small}^{-1}(\mathcal{F}|_{V_\etale})
\xrightarrow{c_\varphi}
\mathcal{F}|_{U_\etale}
$$
est égal à $c_{\psi \circ \varphi}$ au moyen de l'identification
$$
(b \circ a)_{small}^{-1}(\mathcal{F}|_{W_\etale}) =
a_{small}^{-1}(b_{small}^{-1} (\mathcal{F}|_{W_\etale})).
$$
\item Si $\varphi : x \to y$ est au-dessus d'un morphisme \'etale de schémas
$a : U \to V$, alors (\ref{stacks-sheaves-equation-comparison}) est un isomorphisme.
\item Supposons que $f : \mathcal{Y} \to \mathcal{X}$ soit un $1$-morphisme de
catégories fibrées en groupoïdes au-dessus de $(\Sch/S)_{fppf}$, et que $y$ soit
un objet de $\mathcal{Y}$ au-dessus du schéma $U$, d'image
$x = f(y)$. Il existe alors une identification canonique
$f^{-1}\mathcal{F}|_{U_\etale} = \mathcal{F}|_{U_\etale}$.
\item De plus, étant donné $\psi : y' \to y$ dans $\mathcal{Y}$ au-dessus de
$a : U' \to U$, l'application de comparaison
$c_\psi : a_{small}^{-1}(f^{-1}\mathcal{F}|_{U_\etale}) \to
f^{-1}\mathcal{F}|_{U'_\etale}$ est égale à
l'application de comparaison $c_{f(\psi)} : a_{small}^{-1}\mathcal{F}|_{U_\etale}
\to \mathcal{F}|_{U'_\etale}$ au moyen des identifications de (3).
\end{enumerate}
\end{lemma}

\begin{proof}
La vérification de ces propriétés est omise.
\end{proof}

\noindent
Passons à la restriction des (pré)faisceaux de modules.

\begin{lemma}
\label{stacks-sheaves-lemma-localizing-structure-sheaf}
Soit $p : \mathcal{X} \to (\Sch/S)_{fppf}$ une catégorie fibrée
en groupoïdes. Soit $\tau \in \{Zar, \etale, smooth, syntomic, fppf\}$.
Soit $x \in \Ob(\mathcal{X})$ au-dessus de $U = p(x)$.
L'équivalence du
Lemme \ref{stacks-sheaves-lemma-localizing}
s'étend en une équivalence de sites annelés
$(\mathcal{X}_\tau/x, \mathcal{O}_\mathcal{X}|_x) \to
((\Sch/U)_\tau, \mathcal{O})$.
\end{lemma}

\begin{proof}
Cela résulte immédiatement de la construction des faisceaux structuraux.
\end{proof}

\noindent
Soit $\mathcal{X}$ une catégorie fibrée en groupoïdes au-dessus de $(\Sch/S)_{fppf}$.
Soit $\mathcal{F}$ un (pré)faisceau de modules sur $\mathcal{X}$ comme dans la
Définition \ref{stacks-sheaves-definition-modules}.
Soit $x$ un objet de $\mathcal{X}$ au-dessus de $U$. Alors le
Lemme \ref{stacks-sheaves-lemma-localizing-structure-sheaf}
garantit que la restriction
$x^{-1}\mathcal{F}$ est un (pré)faisceau de modules sur $(\Sch/U)_{fppf}$.
Dans ce cas, nous écrirons parfois $x^*\mathcal{F} = x^{-1}\mathcal{F}$.
De même, si $\mathcal{F}$ est un faisceau pour la topologie de Zariski, \'etale, lisse
ou syntomique, alors $x^{-1}\mathcal{F}$ en est un aussi. De plus, la
restriction
$\mathcal{F}|_{U_\etale} = x^{-1}\mathcal{F}|_{U_\etale}$
à $U$ est un préfaisceau de $\mathcal{O}_{U_\etale}$-modules.
Si $\mathcal{F}$ est un faisceau pour la topologie \'etale, alors
$\mathcal{F}|_{U_\etale}$ est un faisceau de modules. De plus,
si $\varphi : x \to y$ est un morphisme de $\mathcal{X}$ au-dessus de
$a : U \to V$, alors l'application de comparaison (\ref{stacks-sheaves-equation-comparison})
est compatible à $a_{small}^\sharp$ (voir
Descente, Remarque \ref{descent-remark-change-topologies-ringed})
et induit une application de {\it comparaison}
\begin{equation}
\label{stacks-sheaves-equation-comparison-modules}
c_\varphi : a_{small}^*(\mathcal{F}|_{V_\etale})
\longrightarrow
\mathcal{F}|_{U_\etale}
\end{equation}
de $\mathcal{O}_{U_\etale}$-modules.
Remarquons que les propriétés (1), (2), (3) et (4) du
Lemme \ref{stacks-sheaves-lemma-comparison}
valent également pour les faisceaux \'etales de modules.
Nous l'utiliserons dans la suite sans autre mention.

\begin{lemma}
\label{stacks-sheaves-lemma-enough-points}
Soit $p : \mathcal{X} \to (\Sch/S)_{fppf}$ une catégorie fibrée
en groupoïdes. Soit $\tau \in \{Zar, \etale, smooth, syntomic, fppf\}$.
Le site $\mathcal{X}_\tau$ a suffisamment de points.
\end{lemma}

\begin{proof}
D'après
Sites, Lemme \ref{sites-lemma-enough-points-local},
il faut montrer qu'il existe une famille d'objets $x$ de $\mathcal{X}$
telle que $\mathcal{X}_\tau/x$ ait suffisamment de points et que les faisceaux
$h_x^\#$ recouvrent l'objet final de la catégorie des faisceaux.
D'après le
Lemme \ref{stacks-sheaves-lemma-localizing}
et
Cohomologie \'etale, Lemme \ref{etale-cohomology-lemma-points-fppf},
$\mathcal{X}_\tau/x$ a suffisamment de points pour tout objet
$x$, ce qui conclut.
\end{proof}







\section{Restriction aux espaces algébriques}
\label{stacks-sheaves-section-restriction-algebraic-spaces}

\noindent
Dans cette section, nous considérons les faisceaux sur des catégories représentables par des
espaces algébriques. Le lemme suivant est l'analogue de
Topologies, Lemme \ref{topologies-lemma-at-the-bottom-etale}
pour les espaces algébriques.

\begin{lemma}
\label{stacks-sheaves-lemma-compare}
Soit $S$ un schéma. Soit $\mathcal{X} \to (\Sch/S)_{fppf}$ une catégorie
fibrée en groupoïdes. Supposons que $\mathcal{X}$ soit représentable par un espace
algébrique $F$. Il existe alors un foncteur continu et cocontinu
$
F_\etale \to \mathcal{X}_\etale
$
qui induit un morphisme de sites annelés
$$
\pi_F :
(\mathcal{X}_\etale, \mathcal{O}_\mathcal{X})
\longrightarrow
(F_\etale, \mathcal{O}_F)
$$
et un morphisme de topos annelés
$$
i_F :
(\Sh(F_\etale), \mathcal{O}_F)
\longrightarrow
(\Sh(\mathcal{X}_\etale), \mathcal{O}_\mathcal{X})
$$
tel que $\pi_F \circ i_F = \text{id}$. De plus, $\pi_{F, *} = i_F^{-1}$.
\end{lemma}

\begin{proof}
Choisissons une équivalence $j : \mathcal{S}_F \to \mathcal{X}$ ; voir
Champs algébriques, sections \ref{algebraic-section-split} et
\ref{algebraic-section-representable-by-algebraic-spaces}.
Un objet de $F_\etale$ est un schéma $U$ muni d'un
morphisme \'etale $\varphi : U \to F$. Alors $\varphi$ est un objet
de $\mathcal{S}_F$ au-dessus de $U$. Ainsi $j(\varphi)$ est un objet de
$\mathcal{X}$ au-dessus de $U$. De cette façon, $j$ induit un foncteur
$u : F_\etale \to \mathcal{X}$. Il est clair que
$u$ est continu et cocontinu pour la topologie \'etale sur
$\mathcal{X}$. Puisque $j$ est une équivalence, le foncteur $u$ est pleinement
fidèle. En outre, les produits fibrés et les égalisateurs existent dans $F_\etale$
et $u$ y commute, car ils se calculent au niveau
des schémas sous-jacents dans $F_\etale$. Ainsi,
Sites, Lemmes \ref{sites-lemma-when-shriek},
\ref{sites-lemma-preserve-equalizers} et
\ref{sites-lemma-back-and-forth}
s'appliquent. En particulier, $u$ définit un morphisme de topos
$i_F : \Sh(F_\etale) \to \Sh(\mathcal{X}_\etale)$
et il existe un adjoint à gauche $i_{F, !}$ de $i_F^{-1}$ qui commute
aux produits fibrés et aux égalisateurs.

\medskip\noindent
Nous affirmons que $i_{F, !}$ est exact. Si tel est le cas, nous pouvons définir
$\pi_F$ par les règles $\pi_F^{-1} = i_{F, !}$ et $\pi_{F, *} = i_F^{-1}$,
et tout est clair. Pour démontrer l'assertion, remarquons que nous savons déjà
que $i_{F, !}$
est exact à droite et préserve les produits fibrés. Il suffit donc de montrer
que $i_{F, !}* = *$, où $*$ désigne l'objet final de la catégorie
des faisceaux d'ensembles. Soit $U$ un schéma et soit
$\varphi : U \to F$ surjectif et \'etale. Posons $R = U \times_F U$.
Alors
$$
\xymatrix{
h_R \ar@<1ex>[r] \ar@<-1ex>[r] & h_U \ar[r] & {*}
}
$$
est un diagramme coégalisateur dans $\Sh(F_\etale)$. En utilisant
l'exactitude à droite de $i_{F, !}$, l'égalité $i_{F, !} = (u_p\ )^\#$ et le
Sites, Lemme \ref{sites-lemma-pullback-representable-presheaf}
nous voyons que
$$
\xymatrix{
h_{u(R)} \ar@<1ex>[r] \ar@<-1ex>[r] & h_{u(U)} \ar[r] & i_{F, !}{*}
}
$$
est un diagramme coégalisateur dans $\Sh(\mathcal{X}_\etale)$. Puisque
$j$ est une équivalence et que $F = U/R$, il s'ensuit que
le coégalisateur dans $\Sh(\mathcal{X}_\etale)$ des
deux applications $h_{u(R)} \to h_{u(U)}$ est $*$. Nous omettons de démontrer que
ces morphismes sont compatibles aux faisceaux structuraux.
\end{proof}

\begin{remark}
\label{stacks-sheaves-remark-compare}
Les constructions du Lemme \ref{stacks-sheaves-lemma-compare} sont compatibles avec la localisation \'etale ;
en voici une formulation précise. Soit $S$ un schéma.
Soit $f : \mathcal{X} \to \mathcal{Y}$ un morphisme
de catégories fibrées en groupoïdes au-dessus de $(\Sch/S)_{fppf}$. Supposons que
$\mathcal{X}$ et $\mathcal{Y}$ soient représentables par des espaces algébriques $F$ et $G$,
et que le morphisme induit $f : F \to G$ d'espaces algébriques soit \'etale.
Notons $f_{small} : F_\etale \to G_\etale$
le morphisme correspondant de topos annelés. Alors
$$
\xymatrix{
(\Sh(F_\etale), \mathcal{O}_F) \ar[rr]_{f_{small}} \ar[d]_{i_F} & &
(\Sh(G_\etale), \mathcal{O}_G) \ar[d]^{i_G} \\
(\Sh(\mathcal{X}_\etale), \mathcal{O}_\mathcal{X})
\ar[d]_{\pi_F} \ar[rr]_f & &
(\Sh(\mathcal{Y}_\etale), \mathcal{O}_\mathcal{Y}) \ar[d]^{\pi_G} \\
(\Sh(F_\etale), \mathcal{O}_F) \ar[rr]^{f_{small}} & &
(\Sh(G_\etale), \mathcal{O}_G)
}
$$
est un diagramme commutatif de topos annelés. Nous omettons les détails.
\end{remark}

\noindent
Supposons que $\mathcal{X}$ soit un champ algébrique représenté par
l'espace algébrique $F$.
Soit $j : \mathcal{S}_F \to \mathcal{X}$ une équivalence, et notons
$u : F_\etale \to \mathcal{X}_\etale$ le
foncteur de la démonstration du Lemme \ref{stacks-sheaves-lemma-compare} ci-dessus.
Pour tout faisceau $\mathcal{F}$ sur $\mathcal{X}_\etale$, nous avons
$$
\pi_{F, *}\mathcal{F}(U) = i_F^{-1}\mathcal{F}(U) = \mathcal{F}(u(U)).
$$
C'est pourquoi nous considérons souvent $i_F^{-1}$ comme un {\it foncteur de restriction},
par analogie avec la
Définition \ref{stacks-sheaves-definition-pullback}
et avec la restriction d'un faisceau du grand site \'etale d'un
schéma au petit site \'etale de ce schéma. Nous employons souvent la notation
\begin{equation}
\label{stacks-sheaves-equation-restrict}
\mathcal{F}|_{F_\etale} = i_F^{-1}\mathcal{F} = \pi_{F, *}\mathcal{F}
\end{equation}
dans cette situation.

\begin{lemma}
\label{stacks-sheaves-lemma-compare-morphism}
Soit $S$ un schéma. Soit $f : \mathcal{X} \to \mathcal{Y}$ un morphisme
de catégories fibrées en groupoïdes au-dessus de $(\Sch/S)_{fppf}$. Supposons que
$\mathcal{X}$ et $\mathcal{Y}$ soient représentables par des espaces algébriques $F$ et $G$.
Notons $f : F \to G$ le morphisme induit d'espaces algébriques, et
$f_{small} : F_\etale \to G_\etale$
le morphisme correspondant de topos annelés. Alors
$$
\xymatrix{
(\Sh(\mathcal{X}_\etale), \mathcal{O}_\mathcal{X})
\ar[d]_{\pi_F} \ar[rr]_f & &
(\Sh(\mathcal{Y}_\etale), \mathcal{O}_\mathcal{Y}) \ar[d]^{\pi_G} \\
(\Sh(F_\etale), \mathcal{O}_F) \ar[rr]^{f_{small}} & &
(\Sh(G_\etale), \mathcal{O}_G)
}
$$
est un diagramme commutatif de topos annelés.
\end{lemma}

\begin{proof}
Cela ressemble à
Topologies, Lemme \ref{topologies-lemma-morphism-big-small-etale} (3),
mais une petite difficulté vient de ce que $F \to G$ n'est pas nécessairement
représentable par des schémas. En particulier, nous n'obtenons pas un diagramme commutatif
de sites annelés, mais seulement un diagramme commutatif de topos annelés.

\medskip\noindent
Avant d'aborder la démonstration proprement dite, choisissons des équivalences
$j : \mathcal{S}_F \to \mathcal{X}$ et
$j' : \mathcal{S}_G \to \mathcal{Y}$ qui induisent des foncteurs
$u : F_\etale \to \mathcal{X}$ et
$u' : G_\etale \to \mathcal{Y}$ comme dans la démonstration du
Lemme \ref{stacks-sheaves-lemma-compare}. Grâce à la 2-fonctorialité des
faisceaux sur les catégories fibrées en groupoïdes au-dessus de $\Sch_{fppf}$
(voir la discussion de la section \ref{stacks-sheaves-section-presheaves}),
nous pouvons supposer que $\mathcal{X} = \mathcal{S}_F$ et
$\mathcal{Y} = \mathcal{S}_G$, et que $f : \mathcal{S}_F \to \mathcal{S}_G$
est le foncteur associé au morphisme $f : F \to G$. Par conséquent,
nous omettrons $u$ et $u'$ de la notation : étant donné un objet
$U \to F$ de $F_\etale$, nous notons $U/F$
l'objet correspondant de $\mathcal{X}$. De même pour $G$.

\medskip\noindent
Soit $\mathcal{G}$ un faisceau sur $\mathcal{X}_\etale$.
Pour démontrer (2), calculons $\pi_{G, *}f_*\mathcal{G}$ et
$f_{small, *}\pi_{F, *}\mathcal{G}$. Pour cela, soit $V \to G$ un objet
de $G_\etale$. Alors
$$
\pi_{G, *}f_*\mathcal{G}(V) = f_*\mathcal{G}(V/G) =
\Gamma\Big(
(\Sch/V)_{fppf} \times_{\mathcal{Y}} \mathcal{X},
\ \text{pr}^{-1}\mathcal{G}\Big)
$$
voir (\ref{stacks-sheaves-equation-pushforward}). Le produit fibré figurant dans la formule est
$$
(\Sch/V)_{fppf} \times_{\mathcal{Y}} \mathcal{X} =
(\Sch/V)_{fppf} \times_{\mathcal{S}_G} \mathcal{S}_F =
\mathcal{S}_{V \times_G F}
$$
c'est-à-dire la catégorie scindée fibrée en groupoïdes associée à
l'espace algébrique $V \times_G F$. Et $\text{pr}^{-1}\mathcal{G}$ est un
faisceau sur $\mathcal{S}_{V \times_G F}$ pour la topologie \'etale.

\medskip\noindent
En particulier, si $V \times_G F$ est représentable, c'est-à-dire s'il est un schéma,
alors $\pi_{G, *}f_*\mathcal{G}(V) = \mathcal{G}(V \times_G F/F)$, et
également
$$
f_{small, *}\pi_{F, *}\mathcal{G}(V) =
\pi_{F, *}\mathcal{G}(V \times_G F) =
\mathcal{G}(V \times_G F/F)
$$
ce qui démontre l'égalité voulue dans ce cas particulier.

\medskip\noindent
En général, choisissons un schéma $U$ et un morphisme \'etale surjectif
$U \to V \times_G F$. Posons $R = U \times_{V \times_G F} U$. Alors
$U/V \times_G F$ et $R/V \times_G F$ sont des objets de la catégorie
produit fibré ci-dessus. Puisque $\text{pr}^{-1}\mathcal{G}$ est un
faisceau pour la topologie \'etale sur $\mathcal{S}_{V \times_G F}$,
le diagramme
$$
\xymatrix{
\Gamma\Big(
(\Sch/V)_{fppf} \times_{\mathcal{Y}} \mathcal{X},
\ \text{pr}^{-1}\mathcal{G}\Big)
\ar[r] &
\text{pr}^{-1}\mathcal{G}(U/V \times_G F) \ar@<1ex>[r] \ar@<-1ex>[r] &
\text{pr}^{-1}\mathcal{G}(R/V \times_G F)
}
$$
est un diagramme égalisateur. Remarquons que
$\text{pr}^{-1}\mathcal{G}(U/V \times_G F) = \mathcal{G}(U/F)$
et $\text{pr}^{-1}\mathcal{G}(R/V \times_G F) = \mathcal{G}(R/F)$
par définition des images réciproques. De plus, d'après les résultats de
Propriétés des espaces, section \ref{spaces-properties-section-etale-site}
(en particulier,
Propriétés des espaces,
Remarque \ref{spaces-properties-remark-explain-equivalence} et
Lemme \ref{spaces-properties-lemma-functoriality-etale-site}),
nous voyons qu'il existe un diagramme égalisateur
$$
\xymatrix{
f_{small, *}\pi_{F, *}\mathcal{G}(V)
\ar[r] &
\pi_{F, *}\mathcal{G}(U/F) \ar@<1ex>[r] \ar@<-1ex>[r] &
\pi_{F, *}\mathcal{G}(R/F)
}
$$
Comme nous avons aussi $\pi_{F, *}\mathcal{G}(U/F) = \mathcal{G}(U/F)$
et $\pi_{F, *}\mathcal{G}(U/F) = \mathcal{G}(U/F)$,
nous obtenons une identification canonique
$f_{small, *}\pi_{F, *}\mathcal{G}(V) = \pi_{G, *}f_*\mathcal{G}(V)$.
Nous omettons de démontrer la compatibilité avec les applications de restriction
et la fonctorialité en $\mathcal{G}$.
\end{proof}

\noindent
Soient $f : \mathcal{X} \to \mathcal{Y}$ et $f : F \to G$ comme dans la
seconde partie du lemme précédent. Le lemme et
(\ref{stacks-sheaves-equation-restrict}) donnent
\begin{equation}
\label{stacks-sheaves-equation-compare-big-small}
(f_*\mathcal{F})|_{G_\etale} =
f_{small, *}(\mathcal{F}|_{F_\etale})
\end{equation}
pour tout faisceau $\mathcal{F}$ sur $\mathcal{X}_\etale$.
De plus, si $\mathcal{F}$ est un faisceau de $\mathcal{O}$-modules, alors
(\ref{stacks-sheaves-equation-compare-big-small}) est un isomorphisme de
$\mathcal{O}_G$-modules sur $G_\etale$.

\medskip\noindent
Enfin, supposons donné un diagramme $2$-commutatif
$$
\xymatrix{
\mathcal{U} \ar[r]^a \ar[dr]_f \drtwocell<\omit>{<-2>\varphi} &
\mathcal{V} \ar[d]^g \\
& \mathcal{X}
}
$$
de $1$-morphismes de catégories fibrées en groupoïdes au-dessus de $(\Sch/S)_{fppf}$,
que $\mathcal{F}$ soit un faisceau sur $\mathcal{X}_\etale$,
et que $\mathcal{U}, \mathcal{V}$ soient représentables par des espaces algébriques
$U, V$. Nous obtenons alors une application de comparaison
\begin{equation}
\label{stacks-sheaves-equation-comparison-algebraic-spaces}
c_\varphi : a_{small}^{-1}(g^{-1}\mathcal{F}|_{V_\etale})
\longrightarrow
f^{-1}\mathcal{F}|_{U_\etale}
\end{equation}
où $a : U \to V$ désigne le morphisme d'espaces algébriques correspondant
à $a$. C'est l'analogue de (\ref{stacks-sheaves-equation-comparison}). Nous définissons
$c_\varphi$ comme l'adjointe de l'application
$$
g^{-1}\mathcal{F}|_{V_\etale}
\longrightarrow
a_{small, *}(f^{-1}\mathcal{F}|_{U_\etale}) =
(a_*f^{-1}\mathcal{F})|_{V_\etale}
$$
(égalité donnée par (\ref{stacks-sheaves-equation-compare-big-small}))
qui est la restriction à $V$, au sens de (\ref{stacks-sheaves-equation-restrict}), de l'application
$$
g^{-1}\mathcal{F} \to a_*a^{-1}g^{-1}\mathcal{F} = a_*f^{-1}\mathcal{F}
$$
où la dernière égalité utilise la $2$-commutativité du diagramme ci-dessus.
Si $\mathcal{F}$ est un faisceau de $\mathcal{O}_\mathcal{X}$-modules,
$c_\varphi$ induit une application de {\it comparaison}
\begin{equation}
\label{stacks-sheaves-equation-comparison-algebraic-spaces-modules}
c_\varphi : a_{small}^*(g^*\mathcal{F}|_{V_\etale})
\longrightarrow
f^*\mathcal{F}|_{U_\etale}
\end{equation}
de $\mathcal{O}_{U_\etale}$-modules. C'est l'analogue de
(\ref{stacks-sheaves-equation-comparison-modules}).
Remarquons que les propriétés (1), (2), (3) et (4) du
Lemme \ref{stacks-sheaves-lemma-comparison}
valent aussi dans ce cadre.












\section{Modules quasi-cohérents}
\label{stacks-sheaves-section-quasi-coherent}

\noindent
Nous pouvons maintenant appliquer la définition générale d'un module
quasi-cohérent à la situation étudiée dans ce chapitre.

\begin{definition}
\label{stacks-sheaves-definition-quasi-coherent}
Soit $p : \mathcal{X} \to (\Sch/S)_{fppf}$ une catégorie fibrée
en groupoïdes. Un {\it module quasi-cohérent sur $\mathcal{X}$}, ou un
{\it $\mathcal{O}_\mathcal{X}$-module quasi-cohérent}, est un
module quasi-cohérent sur le site annelé
$(\mathcal{X}_{fppf}, \mathcal{O}_\mathcal{X})$ au sens de
Modules sur les sites, Définition \ref{sites-modules-definition-site-local}.
La catégorie des faisceaux quasi-cohérents sur $\mathcal{X}$
est notée $\QCoh(\mathcal{O}_\mathcal{X})$.
\end{definition}

\noindent
Si $\mathcal{X}$ est un champ algébrique, cette définition coïncide avec toutes les
définitions de la littérature, au sens où $\QCoh(\mathcal{O}_\mathcal{X})$
est équivalente (aux questions ensemblistes près) à toute variante de cette catégorie
définie dans la littérature. Par exemple, nous comparerons notre définition à
celle de \cite[définition 6.1]{olsson_sheaves} dans
Cohomologie des champs, Lemme \ref{stacks-sheaves-lemma-quasi-coherent}.
Nous verrons également plus loin d'autres constructions de cette catégorie.

\medskip\noindent
En général (comme pour les morphismes de schémas), l'image directe
d'un faisceau quasi-cohérent par un $1$-morphisme n'est pas quasi-cohérente.
L'image réciproque, en revanche, préserve la quasi-cohérence.

\begin{lemma}
\label{stacks-sheaves-lemma-pullback-quasi-coherent}
Soit $f : \mathcal{X} \to \mathcal{Y}$ un $1$-morphisme de catégories
fibrées en groupoïdes au-dessus de $(\Sch/S)_{fppf}$.
Le foncteur d'image réciproque
$f^* = f^{-1} : \textit{Mod}(\mathcal{O}_\mathcal{Y}) \to
\textit{Mod}(\mathcal{O}_\mathcal{X})$
préserve les faisceaux quasi-cohérents.
\end{lemma}

\begin{proof}
C'est un fait général ; voir
Modules sur les sites, Lemme \ref{sites-modules-lemma-local-pullback}.
\end{proof}

\noindent
Les faisceaux quasi-cohérents admettent une caractérisation très simple
en termes de leurs images réciproques. Voir aussi le
Lemme \ref{stacks-sheaves-lemma-quasi-coherent}
pour une caractérisation en termes de restrictions.

\begin{lemma}
\label{stacks-sheaves-lemma-characterize-quasi-coherent}
Soit $p : \mathcal{X} \to (\Sch/S)_{fppf}$ une catégorie
fibrée en groupoïdes. Soit $\mathcal{F}$
un faisceau de $\mathcal{O}_\mathcal{X}$-modules. Alors $\mathcal{F}$
est quasi-cohérent si et seulement si $x^*\mathcal{F}$ est un faisceau
quasi-cohérent sur $(\Sch/U)_{fppf}$ pour tout objet $x$ de
$\mathcal{X}$, où $U = p(x)$.
\end{lemma}

\begin{proof}
D'après le
Lemme \ref{stacks-sheaves-lemma-pullback-quasi-coherent},
la condition est nécessaire. Réciproquement, puisque $x^*\mathcal{F}$
n'est autre que la restriction à $\mathcal{X}_{fppf}/x$, la condition est
suffisante par la définition même d'un faisceau quasi-cohérent
(et par le fait que la quasi-cohérence est une propriété intrinsèque
des faisceaux de modules ; voir
Modules sur les sites, section \ref{sites-modules-section-intrinsic}).
\end{proof}

\begin{lemma}
\label{stacks-sheaves-lemma-characterize-quasi-coherent-bis}
Soit $p : \mathcal{X} \to (\Sch/S)_{fppf}$ une catégorie
fibrée en groupoïdes. Soit $\mathcal{F}$ un préfaisceau de
modules sur $\mathcal{X}$. Les conditions suivantes sont équivalentes :
\begin{enumerate}
\item $\mathcal{F}$ est un objet de
$\textit{Mod}(\mathcal{X}_{Zar}, \mathcal{O}_\mathcal{X})$
et $\mathcal{F}$ est un module quasi-cohérent sur
$(\mathcal{X}_{Zar}, \mathcal{O}_\mathcal{X})$ au sens de
Modules sur les sites, Définition \ref{sites-modules-definition-site-local},
\item $\mathcal{F}$ est un objet de
$\textit{Mod}(\mathcal{X}_\etale, \mathcal{O}_\mathcal{X})$
et $\mathcal{F}$ est un module quasi-cohérent sur
$(\mathcal{X}_\etale, \mathcal{O}_\mathcal{X})$ au sens de
Modules sur les sites, Définition \ref{sites-modules-definition-site-local}, et
\item $\mathcal{F}$ est un module quasi-cohérent sur $\mathcal{X}$
au sens de la Définition \ref{stacks-sheaves-definition-quasi-coherent}.
\end{enumerate}
\end{lemma}

\begin{proof}
Supposons que l'une des conditions (1), (2) ou (3) soit satisfaite.
Soit $x$ un objet de $\mathcal{X}$ au-dessus du schéma $U$.
Rappelons que $x^*\mathcal{F} = x^{-1}\mathcal{F}$ n'est autre que la
restriction à $\mathcal{X}/x = (\Sch/U)_\tau$, où
$\tau = fppf$, $\tau = \etale$ ou $\tau = Zar$ ; voir la
section \ref{stacks-sheaves-section-restriction}.
D'après la définition des modules quasi-cohérents sur un site annelé,
cette restriction est quasi-cohérente dès que $\mathcal{F}$ l'est.
D'après Descente, Proposition \ref{descent-proposition-equivalence-quasi-coherent},
$x^*\mathcal{F}$ est le faisceau associé à
un $\mathcal{O}_U$-module quasi-cohérent ; c'est donc
un module quasi-cohérent pour les topologies fppf, \'etale et de
Zariski ; nous utilisons ici également
Descente, Lemme \ref{descent-lemma-sheaf-condition-holds} et
Définition \ref{descent-definition-structure-sheaf}.
Comme cela vaut pour tout objet $x$ de $\mathcal{X}$,
$\mathcal{F}$ est un faisceau pour chacune des trois topologies.
De plus, $\mathcal{F}$ est quasi-cohérent pour chacune
des trois topologies ; cela résulte directement de la définition de la
quasi-cohérence et du fait que $x$ est un objet arbitraire de $\mathcal{X}$.
\end{proof}







\section{Modules localement quasi-cohérents}
\label{stacks-sheaves-section-locally-quasi-coherent}

\noindent
Bien qu'il existe une variante pour la topologie de Zariski, la
topologie \'etale semble être la topologie naturelle à employer dans la
définition suivante.

\begin{definition}
\label{stacks-sheaves-definition-locally-quasi-coherent}
Soit $p : \mathcal{X} \to (\Sch/S)_{fppf}$ une catégorie
fibrée en groupoïdes. Soit $\mathcal{F}$
un préfaisceau de $\mathcal{O}_\mathcal{X}$-modules.
On dit que $\mathcal{F}$ est {\it localement quasi-cohérent}\footnote{Cette
notation n'est pas standard.} si
$\mathcal{F}$ est un faisceau pour la topologie \'etale et si,
pour tout objet $x$ de $\mathcal{X}$, la restriction
$x^*\mathcal{F}|_{U_\etale}$ est un faisceau
quasi-cohérent. Ici $U = p(x)$.
\end{definition}

\noindent
Nous notons $\textit{LQCoh}(\mathcal{O}_\mathcal{X})$ la catégorie des
modules localement quasi-cohérents. Nous avons alors le diagramme suivant
de catégories de modules
$$
\xymatrix{
\QCoh(\mathcal{O}_\mathcal{X}) \ar[r] \ar[d] &
\textit{Mod}(\mathcal{O}_\mathcal{X}) \ar[d] \\
\textit{LQCoh}(\mathcal{O}_\mathcal{X}) \ar[r] &
\textit{Mod}(\mathcal{X}_\etale, \mathcal{O}_\mathcal{X})
}
$$
où les flèches sont des plongements strictement pleins.
Il se trouve que de nombreux résultats sur les faisceaux quasi-cohérents ont un
analogue pour les modules localement quasi-cohérents. De plus, à bien des
égards (comme nous le verrons plus loin), c'est une catégorie naturelle à considérer.
Par exemple, les faisceaux quasi-cohérents sont exactement les
modules localement quasi-cohérents qui sont « cartésiens », c'est-à-dire qui satisfont
la seconde condition du lemme ci-dessous.

\begin{lemma}
\label{stacks-sheaves-lemma-quasi-coherent}
Soit $p : \mathcal{X} \to (\Sch/S)_{fppf}$ une catégorie
fibrée en groupoïdes. Soit $\mathcal{F}$
un préfaisceau de $\mathcal{O}_\mathcal{X}$-modules. Alors $\mathcal{F}$
est quasi-cohérent si et seulement si les deux conditions suivantes sont satisfaites :
\begin{enumerate}
\item $\mathcal{F}$ est localement quasi-cohérent, et
\item pour tout morphisme $\varphi : x \to y$ de $\mathcal{X}$ au-dessus de
$f : U \to V$, le morphisme de comparaison
$c_\varphi : f_{small}^*\mathcal{F}|_{V_\etale} \to
\mathcal{F}|_{U_\etale}$ de
(\ref{stacks-sheaves-equation-comparison-modules}) est un isomorphisme.
\end{enumerate}
\end{lemma}

\begin{proof}
Supposons $\mathcal{F}$ quasi-cohérent. Alors $\mathcal{F}$ est un faisceau
pour la topologie fppf, donc un faisceau pour la topologie \'etale. De plus,
toute image réciproque de $\mathcal{F}$ sur un topos annelé est quasi-cohérente ; ainsi
les restrictions $x^*\mathcal{F}|_{U_\etale}$ sont quasi-cohérentes.
Cela prouve que $\mathcal{F}$ est localement quasi-cohérent.
Soit $y$ un objet de $\mathcal{X}$ tel que $V = p(y)$.
Nous avons vu que $\mathcal{X}/y = (\Sch/V)_{fppf}$. D'après
Descente, Proposition \ref{descent-proposition-equivalence-quasi-coherent},
il s'ensuit que $y^*\mathcal{F}$ est le module quasi-cohérent
associé à un module quasi-cohérent (usuel) $\mathcal{F}_V$ sur
le schéma $V$. Les morphismes de comparaison
(\ref{stacks-sheaves-equation-comparison-modules}) sont donc bien des isomorphismes.

\medskip\noindent
Réciproquement, supposons que $\mathcal{F}$ satisfasse (1) et (2).
Soit $y$ un objet de $\mathcal{X}$ tel que $V = p(y)$. Notons
$\mathcal{F}_V$ le module quasi-cohérent sur
le schéma $V$ correspondant à la restriction
$y^*\mathcal{F}|_{V_\etale}$, laquelle est quasi-cohérente par
l'hypothèse (1), voir
Descente, Proposition \ref{descent-proposition-equivalence-quasi-coherent}.
La condition (2) signifie alors que les restrictions
$x^*\mathcal{F}|_{U_\etale}$, pour $x$ au-dessus de $y$, sont chacune
isomorphes au faisceau \'etale associé à l'image réciproque de $\mathcal{F}_V$
par le morphisme de schémas correspondant $U \to V$.
Ainsi $y^*\mathcal{F}$ est le faisceau sur $(\Sch/V)_{fppf}$
associé à $\mathcal{F}_V$. Il est donc quasi-cohérent (encore d'après
Descente, Proposition \ref{descent-proposition-equivalence-quasi-coherent}) ;
ainsi, nous voyons que $\mathcal{F}$ est quasi-cohérent sur $\mathcal{X}$ d'après le
Lemme \ref{stacks-sheaves-lemma-characterize-quasi-coherent}.
\end{proof}

\begin{lemma}
\label{stacks-sheaves-lemma-pullback-lqc}
Soit $f : \mathcal{X} \to \mathcal{Y}$ un $1$-morphisme de catégories
fibrées en groupoïdes sur $(\Sch/S)_{fppf}$. Le foncteur image réciproque
$f^* = f^{-1} :
\textit{Mod}(\mathcal{Y}_\etale, \mathcal{O}_\mathcal{Y})
\to
\textit{Mod}(\mathcal{X}_\etale, \mathcal{O}_\mathcal{X})$
préserve les faisceaux localement quasi-cohérents.
\end{lemma}

\begin{proof}
Soit $\mathcal{G}$ localement quasi-cohérent sur $\mathcal{Y}$.
Choisissons un objet $x$ de $\mathcal{X}$ au-dessus du schéma $U$.
La restriction $x^*f^*\mathcal{G}|_{U_\etale}$ est égale à
$(f \circ x)^*\mathcal{G}|_{U_\etale}$
et est donc un faisceau quasi-cohérent par hypothèse sur $\mathcal{G}$.
\end{proof}

\begin{lemma}
\label{stacks-sheaves-lemma-lqc-colimits}
Soit $p : \mathcal{X} \to (\Sch/S)_{fppf}$ une catégorie fibrée en
groupoïdes.
\begin{enumerate}
\item La catégorie $\textit{LQCoh}(\mathcal{O}_\mathcal{X})$
admet des limites inductives, qui coïncident avec les limites inductives dans la catégorie
$\textit{Mod}(\mathcal{X}_\etale, \mathcal{O}_\mathcal{X})$.
\item La catégorie $\textit{LQCoh}(\mathcal{O}_\mathcal{X})$
est abélienne, les noyaux et conoyaux étant calculés dans
$\textit{Mod}(\mathcal{X}_\etale, \mathcal{O}_\mathcal{X})$,
autrement dit, le foncteur d'inclusion est exact.
\item Étant donnée une suite exacte courte
$0 \to \mathcal{F}_1 \to \mathcal{F}_2 \to \mathcal{F}_3 \to 0$ de
$\textit{Mod}(\mathcal{X}_\etale, \mathcal{O}_\mathcal{X})$
si deux de ces trois modules sont localement quasi-cohérents, le troisième l'est aussi.
\item Étant donnés $\mathcal{F}, \mathcal{G}$ dans
$\textit{LQCoh}(\mathcal{O}_\mathcal{X})$
le produit tensoriel $\mathcal{F} \otimes_{\mathcal{O}_\mathcal{X}} \mathcal{G}$
dans $\textit{Mod}(\mathcal{X}_\etale, \mathcal{O}_\mathcal{X})$
est un objet de $\textit{LQCoh}(\mathcal{O}_\mathcal{X})$.
\item Étant donnés $\mathcal{F}, \mathcal{G}$ dans
$\textit{LQCoh}(\mathcal{O}_\mathcal{X})$
avec $\mathcal{F}$ de présentation finie sur
$\mathcal{X}_\etale$, le faisceau
$\SheafHom_{\mathcal{O}_\mathcal{X}}(\mathcal{F}, \mathcal{G})$
dans $\textit{Mod}(\mathcal{X}_\etale, \mathcal{O}_\mathcal{X})$
est un objet de $\textit{LQCoh}(\mathcal{O}_\mathcal{X})$.
\end{enumerate}
\end{lemma}

\begin{proof}
Dans les arguments ci-dessous, $x$ désigne un objet arbitraire de $\mathcal{X}$
au-dessus du schéma $U$. Pour montrer qu'un objet $\mathcal{H}$
de $\textit{Mod}(\mathcal{X}_\etale, \mathcal{O}_\mathcal{X})$ appartient
à $\textit{LQCoh}(\mathcal{O}_\mathcal{X})$, nous montrerons que la
restriction $x^*\mathcal{H}|_{U_\etale} = \mathcal{H}|_{U_\etale}$
est un objet quasi-cohérent de $\textit{Mod}(U_\etale, \mathcal{O}_U)$.

\medskip\noindent
Preuve de (1). Soit $\mathcal{I} \to \textit{LQCoh}(\mathcal{O}_\mathcal{X})$,
$i \mapsto \mathcal{F}_i$, un diagramme.
Considérons l'objet $\mathcal{F} = \colim_i \mathcal{F}_i$ de
$\textit{Mod}(\mathcal{X}_\etale, \mathcal{O}_\mathcal{X})$.
Le foncteur image réciproque $x^*$ commute à toutes les limites inductives,
puisqu'il est adjoint à gauche. Ainsi
$x^*\mathcal{F} = \colim_i x^*\mathcal{F}_i$. De même,
$x^*\mathcal{F}|_{U_\etale} =
\colim_i x^*\mathcal{F}_i|_{U_\etale}$.
Or, par hypothèse, chaque $x^*\mathcal{F}_i|_{U_\etale}$
est quasi-cohérent. Par conséquent, $\colim_i x^*\mathcal{F}_i|_{U_\etale}$
est quasi-cohérent d'après
Descente, Lemme \ref{descent-lemma-properties-quasi-coherent}.
Ainsi $x^*\mathcal{F}|_{U_\etale}$ est quasi-cohérent, comme voulu.

\medskip\noindent
Preuve de (2). Il résulte de (1) que les conoyaux existent dans
$\textit{LQCoh}(\mathcal{O}_\mathcal{X})$ et coïncident avec ceux calculés
dans $\textit{Mod}(\mathcal{X}_\etale, \mathcal{O}_\mathcal{X})$.
Soit $\varphi : \mathcal{F} \to \mathcal{G}$ un morphisme de
$\textit{LQCoh}(\mathcal{O}_\mathcal{X})$ et soit
$\mathcal{K} = \Ker(\varphi)$ calculé dans
$\textit{Mod}(\mathcal{X}_\etale, \mathcal{O}_\mathcal{X})$.
Si nous montrons que $\mathcal{K}$ est un module localement quasi-cohérent,
la preuve de (2) sera achevée. Remarquons pour cela que les noyaux
se calculent dans la catégorie des préfaisceaux (sans prendre le faisceau associé).
Ainsi $\mathcal{K}|_{U_\etale}$ est le noyau du morphisme
$\mathcal{F}|_{U_\etale} \to \mathcal{G}|_{U_\etale}$,
c'est-à-dire le noyau d'un morphisme de faisceaux quasi-cohérents sur $U_\etale$,
donc est quasi-cohérent d'après
Descente, Lemme \ref{descent-lemma-properties-quasi-coherent}.
Cela prouve (2).

\medskip\noindent
Preuve de (3). Soit
$0 \to \mathcal{F}_1 \to \mathcal{F}_2 \to \mathcal{F}_3 \to 0$
une suite exacte courte de
$\textit{Mod}(\mathcal{X}_\etale, \mathcal{O}_\mathcal{X})$.
Puisque nous employons la topologie \'etale, la restriction
$0 \to \mathcal{F}_1|_{U_\etale} \to
\mathcal{F}_2|_{U_\etale} \to \mathcal{F}_3|_{U_\etale} \to 0$
est elle aussi une suite exacte courte. L'assertion (3) résulte donc de l'énoncé
correspondant de
Descente, Lemme \ref{descent-lemma-properties-quasi-coherent}.

\medskip\noindent
Preuve de (4). Soient $\mathcal{F}$ et $\mathcal{G}$ dans
$\textit{LQCoh}(\mathcal{O}_\mathcal{X})$. Comme la restriction à $U_\etale$
est donnée par l'image réciproque suivant le morphisme de topos annelés
$U_\etale \to (\Sch/U)_\etale \to \mathcal{X}_\etale$
nous voyons que la restriction du produit tensoriel
$\mathcal{F} \otimes_{\mathcal{O}_\mathcal{X}} \mathcal{G}$
à $U_\etale$ est égale à
$\mathcal{F}|_{U_\etale} \otimes_{\mathcal{O}_U} \mathcal{G}|_{U_\etale}$,
voir Modules sur les sites, Lemme \ref{sites-modules-lemma-tensor-product-pullback}.
Comme $\mathcal{F}|_{U_\etale}$ et $\mathcal{G}|_{U_\etale}$
sont quasi-cohérents, leur produit tensoriel l'est aussi, voir
Descente, Lemme \ref{descent-lemma-properties-quasi-coherent}.

\medskip\noindent
Preuve de (5). Soient $\mathcal{F}$ et $\mathcal{G}$ dans
$\textit{LQCoh}(\mathcal{O}_\mathcal{X})$, avec $\mathcal{F}$
de présentation finie. Puisque
$(\Sch/U)_\etale = \mathcal{X}_\etale/x$
est une localisation de $\mathcal{X}_\etale$ en un objet,
nous voyons que la restriction de
$\SheafHom_{\mathcal{O}_\mathcal{X}}(\mathcal{F}, \mathcal{G})$
à $(\Sch/U)_\etale$ est égale à
$$
\mathcal{H} =
\SheafHom_{\mathcal{O}|_{(\Sch/U)_\etale}}(
\mathcal{F}|_{(\Sch/U)_\etale}, \mathcal{G}|_{(\Sch/U)_\etale})
$$
par Modules sur les sites, Lemme \ref{sites-modules-lemma-internal-hom-restriction}.
Le morphisme de topos annelés
$(U_\etale, \mathcal{O}_U) \to ((\Sch/U)_\etale, \mathcal{O})$
est plat, car l'image réciproque de $\mathcal{O}$ est $\mathcal{O}_U$.
L'image réciproque de $\mathcal{H}$ par ce morphisme est donc égale à
$\SheafHom_{\mathcal{O}_U}(\mathcal{F}|_{U_\etale}, \mathcal{G}|_{U_\etale})$
par Modules sur les sites, Lemme \ref{sites-modules-lemma-pullback-internal-hom}.
Autrement dit, la restriction de
$\SheafHom_{\mathcal{O}_\mathcal{X}}(\mathcal{F}, \mathcal{G})$
à $U_\etale$ est
$\SheafHom_{\mathcal{O}_U}(\mathcal{F}|_{U_\etale}, \mathcal{G}|_{U_\etale})$.
Comme $\mathcal{F}|_{U_\etale}$ et $\mathcal{G}|_{U_\etale}$
sont quasi-cohérents, il en va de même de
$\SheafHom_{\mathcal{O}_U}(\mathcal{F}|_{U_\etale}, \mathcal{G}|_{U_\etale})$,
voir Descente, Lemme \ref{descent-lemma-properties-quasi-coherent}.
On conclut comme précédemment.
\end{proof}

\noindent
Dans le degré de généralité considéré ici, la catégorie des faisceaux quasi-cohérents
n'est pas abélienne. Voir Exemples, section \ref{examples-section-nonabelian-QCoh}.
Voici ce que nous pouvons démontrer sans travail supplémentaire.

\begin{lemma}
\label{stacks-sheaves-lemma-qc-colimits}
Soit $p : \mathcal{X} \to (\Sch/S)_{fppf}$ une catégorie
fibrée en groupoïdes.
\begin{enumerate}
\item La catégorie $\QCoh(\mathcal{O}_\mathcal{X})$
admet des limites inductives, qui coïncident avec les limites inductives dans les catégories
$\textit{Mod}(\mathcal{X}_{Zar}, \mathcal{O}_\mathcal{X})$,
$\textit{Mod}(\mathcal{X}_\etale, \mathcal{O}_\mathcal{X})$,
$\textit{Mod}(\mathcal{O}_\mathcal{X})$, et
$\textit{LQCoh}(\mathcal{O}_\mathcal{X})$.
\item Étant donnés $\mathcal{F}, \mathcal{G}$ dans $\QCoh(\mathcal{O}_\mathcal{X})$,
les produits tensoriels $\mathcal{F} \otimes_{\mathcal{O}_\mathcal{X}} \mathcal{G}$
calculés dans $\textit{Mod}(\mathcal{X}_{Zar}, \mathcal{O}_\mathcal{X})$,
$\textit{Mod}(\mathcal{X}_\etale, \mathcal{O}_\mathcal{X})$, ou
$\textit{Mod}(\mathcal{O}_\mathcal{X})$ coïncident, et leur valeur commune
est un objet de $\QCoh(\mathcal{O}_\mathcal{X})$.
\item Étant donnés $\mathcal{F}, \mathcal{G}$ dans $\QCoh(\mathcal{O}_\mathcal{X})$,
avec $\mathcal{F}$ localement libre de type fini (pour la topologie fppf, ou, de manière équivalente, \'etale, ou
encore, de manière équivalente, de Zariski), les Hom internes
$\SheafHom_{\mathcal{O}_\mathcal{X}}(\mathcal{F}, \mathcal{G})$
calculés dans $\textit{Mod}(\mathcal{X}_{Zar}, \mathcal{O}_\mathcal{X})$,
$\textit{Mod}(\mathcal{X}_\etale, \mathcal{O}_\mathcal{X})$, ou
$\textit{Mod}(\mathcal{O}_\mathcal{X})$ coïncident, et leur valeur commune
est un objet de $\QCoh(\mathcal{O}_\mathcal{X})$.
\end{enumerate}
\end{lemma}

\begin{proof}
Soit $x$ un objet arbitraire de $\mathcal{X}$ au-dessus du schéma $U$.
Soit $\tau \in \{Zariski, \etale, fppf\}$. Pour montrer qu'un objet $\mathcal{H}$
de $\textit{Mod}(\mathcal{X}_\tau, \mathcal{O}_\mathcal{X})$ appartient
à $\QCoh(\mathcal{O}_\mathcal{X})$, il suffit de montrer que la
restriction $x^*\mathcal{H}$ (section \ref{stacks-sheaves-section-restriction})
est un objet quasi-cohérent de $\textit{Mod}((\Sch/U)_\tau, \mathcal{O})$.
Voir les Lemmes \ref{stacks-sheaves-lemma-characterize-quasi-coherent} et
\ref{stacks-sheaves-lemma-characterize-quasi-coherent-bis}. Il en va de même pour la propriété d'être
localement libre de type fini. Rappelons que $(\Sch/U)_\tau = \mathcal{X}_\tau/x$
est une localisation de $\mathcal{X}_\tau$ en un objet. La restriction
commute donc aux limites inductives, aux produits tensoriels et à la formation du Hom interne
(voir Modules sur les sites, Lemmes
\ref{sites-modules-lemma-exactness-pushforward-pullback},
\ref{sites-modules-lemma-tensor-product-pullback}, et
\ref{sites-modules-lemma-internal-hom-restriction}).
Le lemme se ramène ainsi à Descente, Lemme \ref{descent-lemma-qc-colimits}.
\end{proof}






\section{Champification et faisceaux}
\label{stacks-sheaves-section-stackification}

\noindent
Il se trouve que la catégorie des faisceaux sur une catégorie fibrée en
groupoïdes ne « connaît » que la champification.

\begin{lemma}
\label{stacks-sheaves-lemma-stackification}
Soit $f : \mathcal{X} \to \mathcal{Y}$ un $1$-morphisme de catégories
fibrées en groupoïdes sur $(\Sch/S)_{fppf}$. Si
$f$ induit une équivalence des champifications, alors le morphisme
de topos
$f : \Sh(\mathcal{X}_{fppf}) \to \Sh(\mathcal{Y}_{fppf})$
est une équivalence.
\end{lemma}

\begin{proof}
Nous pouvons supposer que $\mathcal{Y}$ est la champification de $\mathcal{X}$.
Nous affirmons que $f : \mathcal{X} \to \mathcal{Y}$ est un foncteur
cocontinu spécial, voir
Sites, Définition \ref{sites-definition-special-cocontinuous-functor},
ce qui prouvera le lemme. D'après
Champs, Lemme \ref{stacks-lemma-topology-inherited-functorial},
le foncteur $f$ est continu et cocontinu. D'après
Champs, Lemme \ref{stacks-lemma-stackify},
nous voyons que les conditions (3), (4) et (5) de
Sites, Lemme \ref{sites-lemma-equivalence},
sont satisfaites.
\end{proof}

\begin{lemma}
\label{stacks-sheaves-lemma-stackification-quasi-coherent}
Soit $f : \mathcal{X} \to \mathcal{Y}$ un $1$-morphisme de catégories
fibrées en groupoïdes sur $(\Sch/S)_{fppf}$. Si
$f$ induit une équivalence des champifications, alors $f^*$
induit des équivalences
$\textit{Mod}(\mathcal{O}_\mathcal{X}) \to
\textit{Mod}(\mathcal{O}_\mathcal{Y})$
et
$\QCoh(\mathcal{O}_\mathcal{X}) \to
\QCoh(\mathcal{O}_\mathcal{Y})$.
\end{lemma}

\begin{proof}
Nous pouvons supposer que $\mathcal{Y}$ est la champification de $\mathcal{X}$.
La première assertion résulte clairement du
Lemme \ref{stacks-sheaves-lemma-stackification}
et de
$\mathcal{O}_\mathcal{X} = f^{-1}\mathcal{O}_\mathcal{Y}$.
Les images réciproques de faisceaux quasi-cohérents sont quasi-cohérentes, voir
Lemme \ref{stacks-sheaves-lemma-pullback-quasi-coherent}.
Il suffit donc de montrer que, si $f^*\mathcal{G}$ est
quasi-cohérent, alors $\mathcal{G}$ l'est.
Pour cela, soit $y$ un objet de $\mathcal{Y}$.
En explicitant la condition que $\mathcal{Y}$ est la champification
de $\mathcal{X}$, nous voyons qu'il existe un recouvrement fppf $\{y_i \to y\}$
dans $\mathcal{Y}$ tel que $y_i \cong f(x_i)$ pour un
objet $x_i$ de $\mathcal{X}$. Disons que $x_i$ et $y_i$ sont au-dessus du schéma $U_i$.
La quasi-cohérence de $f^*\mathcal{G}$ signifie alors que $x_i^*f^*\mathcal{G}$
est quasi-cohérent. Comme $x_i^*f^*\mathcal{G}$ est isomorphe à
$y_i^*\mathcal{G}$ (en tant que faisceaux sur $(\Sch/U_i)_{fppf}$), nous
voyons que $y_i^*\mathcal{G}$ est quasi-cohérent.
Il résulte de
Modules sur les sites, Lemme \ref{sites-modules-lemma-local-final-object}
que la restriction de $\mathcal{G}$ à $\mathcal{Y}/y$ est
quasi-cohérente. Ainsi $\mathcal{G}$ est quasi-cohérent d'après le
Lemme \ref{stacks-sheaves-lemma-characterize-quasi-coherent}.
\end{proof}





\section{Faisceaux quasi-cohérents et présentations}
\label{stacks-sheaves-section-quasi-coherent-presentation}

\noindent
Commençons par comparer les faisceaux quasi-cohérents aux notions définies
précédemment pour les schémas et les espaces algébriques.

\begin{lemma}
\label{stacks-sheaves-lemma-compare-locally-quasi-coherent}
Soit $S$ un schéma. Soit $\mathcal{X} \to (\Sch/S)_{fppf}$ une catégorie
fibrée en groupoïdes représentable par un espace algébrique $F$.
Si $\mathcal{F}$ appartient à $\textit{LQCoh}(\mathcal{O}_\mathcal{X})$,
alors la restriction $\mathcal{F}|_{F_\etale}$ (\ref{stacks-sheaves-equation-restrict})
est quasi-cohérente.
\end{lemma}

\begin{proof}
Soit $U$ un schéma \'etale sur $F$. Alors
$\mathcal{F}|_{U_\etale} = (\mathcal{F}|_{F_\etale})|_{U_\etale}$.
C'est clair, mais voir aussi la Remarque \ref{stacks-sheaves-remark-compare}.
L'assertion résulte donc des définitions.
\end{proof}

\begin{lemma}
\label{stacks-sheaves-lemma-compare-quasi-coherent}
Soit $S$ un schéma. Soit $\mathcal{X} \to (\Sch/S)_{fppf}$ une catégorie
fibrée en groupoïdes représentable par un espace algébrique $F$.
Le foncteur (\ref{stacks-sheaves-equation-restrict}) définit une équivalence
$$
\QCoh(\mathcal{O}_\mathcal{X}) \to \QCoh(\mathcal{O}_F),\quad
\mathcal{F} \longmapsto \mathcal{F}|_{F_\etale}
$$
dont un quasi-inverse est donné par $\mathcal{G} \mapsto \pi_F^*\mathcal{G}$.
Cette équivalence est compatible aux images réciproques pour les morphismes entre
catégories fibrées en groupoïdes représentables par des espaces algébriques.
\end{lemma}

\begin{proof}
D'après le Lemme \ref{stacks-sheaves-lemma-characterize-quasi-coherent-bis}, nous pouvons travailler avec
la topologie \'etale. Nous emploierons sans autre mention les notations et résultats du
Lemme \ref{stacks-sheaves-lemma-compare}.
Rappelons que le foncteur de restriction
$\textit{Mod}(\mathcal{X}_\etale, \mathcal{O}_\mathcal{X}) \to
\textit{Mod}(F_\etale, \mathcal{O}_F)$,
$\mathcal{F} \mapsto \mathcal{F}|_{F_\etale}$ est donné par $i_F^*$. D'après le
Lemme \ref{stacks-sheaves-lemma-compare-locally-quasi-coherent}, ou d'après
Modules sur les sites, Lemme \ref{sites-modules-lemma-local-pullback}
nous voyons que $\mathcal{F}|_{F_\etale}$ est quasi-cohérent si
$\mathcal{F}$ est quasi-cohérent.
Nous obtenons donc le foncteur indiqué dans l'énoncé du
lemme et, dans l'autre sens, un foncteur $\pi_F^*$.
Comme $\pi_F \circ i_F = \text{id}$, nous avons
$i_F^*\pi_F^*\mathcal{G} = \mathcal{G}$.

\medskip\noindent
Pour $\mathcal{F}$ dans
$\textit{Mod}(\mathcal{X}_\etale, \mathcal{O}_\mathcal{X})$
il existe un morphisme canonique $\pi_F^*(\mathcal{F}|_{F_\etale}) \to \mathcal{F}$,
à savoir le morphisme adjoint à l'identification
$\mathcal{F}|_{F_\etale} = \pi_{F, *}\mathcal{F}$.
Nous montrerons que ce morphisme est un isomorphisme si $\mathcal{F}$ est
un module quasi-cohérent sur $\mathcal{X}$.
Choisissons un schéma $U$ et un morphisme \'etale surjectif $U \to F$.
Notons $x : U \to \mathcal{X}$ l'objet correspondant de $\mathcal{X}$
au-dessus de $U$. Il suffit de montrer que
$\pi_F^*(\mathcal{F}|_{F_\etale}) \to \mathcal{F}$ est un isomorphisme
après restriction à $\mathcal{X}_\etale/x = (\Sch/U)_\etale$.
Comme $U \to F$ est \'etale, la Remarque \ref{stacks-sheaves-remark-compare} donne
$$
\pi_F^*(\mathcal{F}|_{F_\etale})|_{\mathcal{X}_\etale/x} =
\pi_U^*(\mathcal{F}|_{U_\etale})
$$
et montre que la restriction du morphisme
$\pi_F^*(\mathcal{F}|_{F_\etale}) \to \mathcal{F}$
à $\mathcal{X}_\etale/x = (\Sch/U)_\etale$
est égale au morphisme correspondant
$\pi_U^*(\mathcal{F}|_{U_\etale}) \to \mathcal{F}|_{(\Sch/U)_\etale}$.
Puisque nous avons démontré le résultat pour les schémas dans
Descente, section \ref{descent-section-quasi-coherent-sheaves}\footnote{En effet,
si $U$ est un schéma et $\mathcal{F}$ est quasi-cohérent sur
$(\Sch/U)_\etale$, alors $\mathcal{F} = \mathcal{H}^a$ pour un
module quasi-cohérent $\mathcal{H}$ sur le schéma $U$, d'après
Descente, Proposition \ref{descent-proposition-equivalence-quasi-coherent}.
Autrement dit, $\mathcal{F} = (\text{id}_{\etale,Zar})^*\mathcal{H}$
d'après Descente, Remarque \ref{descent-remark-change-topologies-ringed-sites},
avec les notations de Descente, Lemme \ref{descent-lemma-compare-sites}.
Or nous avons
$\text{id}_{\etale,Zar} = \pi_U \circ \text{id}_{small,\etale,Zar}$
et nous voyons donc que $\mathcal{F} = \pi_U^*\mathcal{G}$, où
$\mathcal{G} = (\text{id}_{small,\etale,Zar})^*\mathcal{H}$ est
quasi-cohérent. Alors
$\pi_U^*i_U^*\mathcal{F} = \pi_U^*i_U^*\pi_U^*\mathcal{G} =
\pi_U^*\mathcal{G} = \mathcal{F}$, comme voulu.}
nous concluons.

\medskip\noindent
La compatibilité aux images réciproques résulte du fait que
le quasi-inverse est donné par $\pi_F^*$ et du
diagramme commutatif de topos annelés du
Lemme \ref{stacks-sheaves-lemma-compare-morphism}.
\end{proof}

\noindent
Dans
Groupoïdes dans les espaces, Définition
\ref{spaces-groupoids-definition-groupoid-module}
nous avons défini la notion de module quasi-cohérent
sur un groupoïde arbitraire. La proposition (formelle) suivante affirme
que nous pouvons étudier les faisceaux quasi-cohérents sur les champs quotients
en termes de modules quasi-cohérents sur les présentations.

\begin{proposition}
\label{stacks-sheaves-proposition-quasi-coherent}
Soit $(U, R, s, t, c)$ un groupoïde en espaces algébriques sur $S$.
Soit $\mathcal{X} = [U/R]$ le champ quotient.
La catégorie des modules quasi-cohérents sur $\mathcal{X}$
est équivalente à la catégorie des modules quasi-cohérents
sur $(U, R, s, t, c)$.
\end{proposition}

\begin{proof}
Nous allons construire des foncteurs quasi-inverses
$$
\QCoh(\mathcal{O}_\mathcal{X})
\longleftrightarrow
\QCoh(U, R, s, t, c).
$$
où $\QCoh(U, R, s, t, c)$ désigne la catégorie des modules quasi-cohérents
sur le groupoïde $(U, R, s, t, c)$.

\medskip\noindent
Soit $\mathcal{F}$ un objet de $\QCoh(\mathcal{O}_\mathcal{X})$.
Notons $\mathcal{U}$ et $\mathcal{R}$ les catégories fibrées en groupoïdes
correspondant à $U$ et $R$. Notons $x$ l'objet (définissant) de $\mathcal{X}$
au-dessus de $U$. Rappelons que nous avons un diagramme $2$-commutatif
$$
\xymatrix{
\mathcal{R} \ar[r]_s \ar[d]_t & \mathcal{U} \ar[d]^x \\
\mathcal{U} \ar[r]^x & \mathcal{X}
}
$$
Voir Groupoïdes dans les espaces, Lemme
\ref{spaces-groupoids-lemma-quotient-stack-2-arrow}.
D'après le Lemme \ref{stacks-sheaves-lemma-2-morphisms-presheaves},
la $2$-flèche inhérente au diagramme induit un isomorphisme
$\alpha : t^*x^*\mathcal{F} \to s^*x^*\mathcal{F}$
qui satisfait la condition de cocycle sur
$\mathcal{R} \times_{s, \mathcal{U}, t} \mathcal{R}$ ;
c'est une conséquence de Groupoïdes dans les espaces, Lemme
\ref{spaces-groupoids-lemma-quotient-stack-cocycle}.
Ainsi, si nous posons $\mathcal{G} = x^*\mathcal{F}|_{U_\etale}$,
l'équivalence de catégories du
Lemme \ref{stacks-sheaves-lemma-compare-quasi-coherent}
(employée plusieurs fois de façon compatible aux images réciproques)
donne un isomorphisme
$\alpha : t_{small}^*\mathcal{G} \to s_{small}^*\mathcal{G}$
satisfaisant la condition de cocycle sur $R \times_{s, U, t} R$, c'est-à-dire que
$(\mathcal{G}, \alpha)$ est un objet de $\QCoh(U, R, s, t, c)$.
La règle $\mathcal{F} \mapsto (\mathcal{G}, \alpha)$ est notre
foncteur de gauche à droite.

\medskip\noindent
Construction du foncteur dans l'autre sens.
Soit $(\mathcal{G}, \alpha)$ un objet de $\QCoh(U, R, s, t, c)$.
D'après le Lemme \ref{stacks-sheaves-lemma-stackification-quasi-coherent},
le morphisme de champification $[U/_{\!p}R] \to [U/R]$ (voir
Groupoïdes dans les espaces, Définition
\ref{spaces-groupoids-definition-quotient-stack})
induit une équivalence des catégories de faisceaux quasi-cohérents.
Il suffit donc de construire un module quasi-cohérent $\mathcal{F}$
sur $[U/_{\!p}R]$.

\medskip\noindent
Rappelons qu'un objet $x = (T, u)$ de $[U/_{\!p}R]$
est donné par un schéma $T$ et un morphisme $u : T \to U$. Un morphisme
$(T, u) \to (T', u')$ est donné par une paire $(f, r)$, où $f : T \to T'$
et $r : T \to R$, avec $s \circ r = u$ et $t \circ r = u' \circ f$.
Appelons {\it morphisme spécial} tout morphisme de la forme
$(f, e \circ u' \circ f) : (T, u' \circ f) \to (T', u')$.
La catégorie des $(T, u)$, munie des morphismes spéciaux, n'est autre que la
catégorie des schémas sur $U$.

\medskip\noindent
Avec ces notations, pour un objet $(T, u)$ de $[U/_{\!p}R]$, posons
$$
\mathcal{F}(T, u) : = \Gamma(T, u_{small}^*\mathcal{G}).
$$
À un morphisme $(f, r) : (T, u) \to (T', u')$ est associé un morphisme
\begin{align*}
\mathcal{F}(T', u') & = \Gamma(T', (u')_{small}^*\mathcal{G}) \\
& \to \Gamma(T, f_{small}^*(u')_{small}^*\mathcal{G}) =
\Gamma(T, (u' \circ f)_{small}^*\mathcal{G}) \\
& = \Gamma(T, (t \circ r)_{small}^*\mathcal{G}) =
\Gamma(T, r_{small}^*t_{small}^*\mathcal{G}) \\
& \to \Gamma(T, r_{small}^*s_{small}^*\mathcal{G}) =
\Gamma(T, (s \circ r)_{small}^*\mathcal{G}) \\
& = \Gamma(T, u_{small}^*\mathcal{G}) \\
& = \mathcal{F}(T, u)
\end{align*}
où la première flèche est l'image réciproque suivant $f$ et la seconde est
$\alpha$. Remarquons que si $(f, r)$ est un morphisme spécial, ce
morphisme est simplement l'image réciproque suivant $f$, car $e_{small}^*\alpha = \text{id}$ par
les axiomes d'un faisceau de modules quasi-cohérents sur un groupoïde.
La condition de cocycle implique que $\mathcal{F}$ est un préfaisceau
de modules (détails omis). Nous voyons que la restriction de $\mathcal{F}$
à $(\Sch/T)_{fppf}$ est quasi-cohérente par la description simple des
morphismes de restriction de $\mathcal{F}$ dans le cas d'un morphisme spécial.
Ainsi $\mathcal{F}$ est un faisceau quasi-cohérent sur $[U/_{\!p}R]$
(Lemme \ref{stacks-sheaves-lemma-characterize-quasi-coherent}).

\medskip\noindent
Nous omettons de vérifier que les foncteurs construits ci-dessus sont
quasi-inverses l'un de l'autre.
\end{proof}

\noindent
Nous terminons cette section par un lemme technique sur les morphismes ayant pour source un faisceau
quasi-cohérent. C'est un analogue de
Schémas, Lemme \ref{schemes-lemma-compare-constructions}.
Nous verrons plus loin (Critères de représentabilité, Théorème
\ref{criteria-theorem-flat-groupoid-gives-algebraic-stack})
que les hypothèses faites sur le groupoïde impliquent que $\mathcal{X}$ est
un champ algébrique.

\begin{lemma}
\label{stacks-sheaves-lemma-map-from-quasi-coherent}
Soit $(U, R, s, t, c)$ un groupoïde en espaces algébriques sur $S$.
Supposons $s, t$ plats et localement de présentation finie.
Soit $\mathcal{X} = [U/R]$ le champ quotient. Notons
$x$ l'objet de $\mathcal{X}$ au-dessus de $U$.
Soit $\mathcal{F}$ un
$\mathcal{O}_\mathcal{X}$-module quasi-cohérent, et soit $\mathcal{H}$ un objet quelconque
de $\textit{Mod}(\mathcal{O}_\mathcal{X})$.
Le morphisme
$$
\Hom_{\mathcal{O}_\mathcal{X}}(\mathcal{F}, \mathcal{H})
\longrightarrow
\Hom_{\mathcal{O}_U}(x^*\mathcal{F}|_{U_\etale},
x^*\mathcal{H}|_{U_\etale}),
\quad
\phi \longmapsto x^*\phi|_{U_\etale}
$$
est injectif, et son image est exactement formée des
$\varphi : x^*\mathcal{F}|_{U_\etale} \to
x^*\mathcal{H}|_{U_\etale}$ qui donnent lieu à un
diagramme commutatif
$$
\xymatrix{
s_{small}^*(x^*\mathcal{F}|_{U_\etale})
\ar[r] \ar[d]^{s_{small}^*\varphi} &
(x \circ s)^*\mathcal{F}|_{R_\etale} =
(x \circ t)^*\mathcal{F}|_{R_\etale} &
t_{small}^*(x^*\mathcal{F}|_{U_\etale})
\ar[l] \ar[d]_{t_{small}^*\varphi} \\
s_{small}^*(x^*\mathcal{H}|_{U_\etale})
\ar[r] &
(x \circ s)^*\mathcal{H}|_{R_\etale} =
(x \circ t)^*\mathcal{H}|_{R_\etale} &
t_{small}^*(x^*\mathcal{H}|_{U_\etale})
\ar[l]
}
$$
de modules sur $R_\etale$,
où les flèches horizontales sont les morphismes de comparaison
(\ref{stacks-sheaves-equation-comparison-algebraic-spaces-modules}).
\end{lemma}

\begin{proof}
D'après le
Lemme \ref{stacks-sheaves-lemma-stackification-quasi-coherent}
le morphisme de champification $[U/_{\!p}R] \to [U/R]$ (voir
Groupoïdes dans les espaces, Définition
\ref{spaces-groupoids-definition-quotient-stack})
induit une équivalence des catégories de faisceaux quasi-cohérents
et de $\mathcal{O}$-modules fppf.
Il suffit donc de démontrer le lemme pour $\mathcal{X} = [U/_{\!p}R]$.
D'après la Proposition \ref{stacks-sheaves-proposition-quasi-coherent}
et sa preuve, il existe un module quasi-cohérent
$(\mathcal{G}, \alpha)$ sur $(U, R, s, t, c)$ tel que
$\mathcal{F}$ soit donné par la règle
$\mathcal{F}(T, u) = \Gamma(T, u^*\mathcal{G})$.
En particulier, $x^*\mathcal{F}|_{U_\etale} = \mathcal{G}$
et il est clair que le morphisme de l'énoncé du
lemme est injectif. De plus, étant donnés un morphisme
$\varphi : \mathcal{G} \to x^*\mathcal{H}|_{U_\etale}$
et un objet quelconque
$y = (T, u)$ de $[U/_{\!p}R]$, nous pouvons considérer le morphisme
$$
\mathcal{F}(y) = \Gamma(T, u^*\mathcal{G})
\xrightarrow{u_{small}^*\varphi}
\Gamma(T, u_{small}^*x^*\mathcal{H}|_{U_\etale})
\rightarrow
\Gamma(T, y^*\mathcal{H}|_{T_\etale}) = \mathcal{H}(y)
$$
où la seconde flèche est le morphisme de comparaison
(\ref{stacks-sheaves-equation-comparison-modules}) pour le faisceau $\mathcal{H}$.
Cette construction est compatible aux morphismes de restriction des
faisceaux $\mathcal{F}$ et $\mathcal{G}$ pour les morphismes de
$[U/_{\!p}R]$ si la condition de cocycle du
lemme est satisfaite. Preuve omise. Indication : les morphismes de restriction
de $\mathcal{F}$ sont explicités en termes de $(\mathcal{G}, \alpha)$
dans la preuve de la
Proposition \ref{stacks-sheaves-proposition-quasi-coherent}.
\end{proof}







\section{Faisceaux quasi-cohérents sur les champs algébriques}
\label{stacks-sheaves-section-quasi-coherent-algebraic-stacks}

\noindent
Soit $\mathcal{X}$ un champ algébrique sur $S$. D'après
Champs algébriques, Lemme \ref{algebraic-lemma-stack-presentation}
il existe une équivalence $[U/R] \to \mathcal{X}$,
où $(U, R, s, t, c)$ est un groupoïde lisse en espaces algébriques.
Alors
$$
\QCoh(\mathcal{O}_\mathcal{X})
\cong
\QCoh(\mathcal{O}_{[U/R]})
\cong
\QCoh(U, R, s, t, c)
$$
où la seconde équivalence est la
Proposition \ref{stacks-sheaves-proposition-quasi-coherent}.
Ainsi, la catégorie des faisceaux quasi-cohérents sur un champ algébrique
est équivalente à la catégorie des modules quasi-cohérents sur un groupoïde
lisse en espaces algébriques. En particulier, d'après
Groupoïdes dans les espaces, Lemme \ref{spaces-groupoids-lemma-abelian}
nous voyons que $\QCoh(\mathcal{O}_\mathcal{X})$ est abélienne !

\medskip\noindent
Notre cadre actuel présente un aspect légèrement déconcertant :
le plongement pleinement fidèle
$$
\QCoh(\mathcal{O}_\mathcal{X})
\longrightarrow
\textit{Mod}(\mathcal{O}_\mathcal{X})
$$
n'est en général {\bf pas} exact. Toutefois, il en va exactement de même
pour les schémas : pour la plupart des schémas $X$, le plongement
$$
\QCoh(\mathcal{O}_X) \cong
\QCoh((\Sch/X)_{fppf}, \mathcal{O}_X) \longrightarrow
\textit{Mod}((\Sch/X)_{fppf}, \mathcal{O}_X)
$$
n'est pas exact, voir
Descente, Lemme \ref{descent-lemma-equivalence-quasi-coherent-limits}.
Incidemment, l'exemple donné dans la preuve de
Descente, Lemme \ref{descent-lemma-equivalence-quasi-coherent-limits}
montre qu'en général le plongement strictement plein
$\QCoh(\mathcal{O}_\mathcal{X}) \to
\textit{LQCoh}(\mathcal{O}_\mathcal{X})$ n'est pas exact non plus.

\medskip\noindent
Nous réunissons tous les résultats obtenus jusqu'ici en un seul énoncé.

\begin{lemma}
\label{stacks-sheaves-lemma-quasi-coherent-algebraic-stack}
Soit $\mathcal{X}$ un champ algébrique sur $S$.
\begin{enumerate}
\item Si $[U/R] \to \mathcal{X}$ est une présentation de $\mathcal{X}$,
il existe une équivalence canonique
$\QCoh(\mathcal{O}_\mathcal{X}) \cong
\QCoh(U, R, s, t, c)$.
\item La catégorie $\QCoh(\mathcal{O}_\mathcal{X})$ est abélienne.
\item Le foncteur d'inclusion $\QCoh(\mathcal{O}_\mathcal{X}) \to
\textit{Mod}(\mathcal{O}_\mathcal{X})$ est exact à droite, mais
n'est en général {\bf pas} exact.
\item La catégorie $\QCoh(\mathcal{O}_\mathcal{X})$
admet des limites inductives, qui coïncident avec les limites inductives dans la catégorie
$\textit{Mod}(\mathcal{O}_\mathcal{X})$.
\item Étant donnés $\mathcal{F}, \mathcal{G}$ dans
$\QCoh(\mathcal{O}_\mathcal{X})$
le produit tensoriel $\mathcal{F} \otimes_{\mathcal{O}_\mathcal{X}} \mathcal{G}$
dans $\textit{Mod}(\mathcal{O}_\mathcal{X})$
est un objet de $\QCoh(\mathcal{O}_\mathcal{X})$.
\item Étant donnés $\mathcal{F}, \mathcal{G}$ dans
$\QCoh(\mathcal{O}_\mathcal{X})$
avec $\mathcal{F}$ localement libre de type fini, le faisceau
$\SheafHom_{\mathcal{O}_\mathcal{X}}(\mathcal{F}, \mathcal{G})$
dans $\textit{Mod}(\mathcal{O}_\mathcal{X})$
est un objet de $\QCoh(\mathcal{O}_\mathcal{X})$.
\item Étant donnée une suite exacte courte
$0 \to \mathcal{F}_1 \to \mathcal{F}_2 \to \mathcal{F}_3 \to 0$
dans $\textit{Mod}(\mathcal{O}_\mathcal{X})$,
si $\mathcal{F}_1$ et $\mathcal{F}_3$ sont quasi-cohérents, alors
$\mathcal{F}_2$ est quasi-cohérent.
\end{enumerate}
\end{lemma}

\begin{proof}
Les propriétés (4), (5) et (6) ont été démontrées dans le
Lemme \ref{stacks-sheaves-lemma-qc-colimits}.
L'assertion (1) est la
Proposition \ref{stacks-sheaves-proposition-quasi-coherent}.
L'assertion (2) résulte de (1) et de
Groupoïdes dans les espaces, Lemme \ref{spaces-groupoids-lemma-abelian}
comme expliqué ci-dessus. L'exactitude à droite du foncteur d'inclusion
dans (3) résulte de (4) ; comparer avec
Homologie, Lemme \ref{homology-lemma-exact-functor}.
Pour la non-exactitude du foncteur d'inclusion dans (3), voir
Descente, Lemme \ref{descent-lemma-equivalence-quasi-coherent-limits}.
Pour (7), remarquons qu'il suffit de vérifier que la restriction
de $\mathcal{F}_2$ au grand site d'un schéma est quasi-cohérente
(Lemme \ref{stacks-sheaves-lemma-characterize-quasi-coherent}) ;
cela résulte donc de l'assertion correspondante de
Descente, Lemme \ref{descent-lemma-equivalence-quasi-coherent-limits}.
\end{proof}

\noindent
Construisons maintenant le cohérateur des modules sur un champ algébrique.

\begin{proposition}
\label{stacks-sheaves-proposition-coherator}
Soit $\mathcal{X}$ un champ algébrique sur $S$.
\begin{enumerate}
\item La catégorie $\QCoh(\mathcal{O}_\mathcal{X})$ est une catégorie
abélienne de Grothendieck. Par conséquent, $\QCoh(\mathcal{O}_\mathcal{X})$
a suffisamment d'injectifs et admet toutes les limites.
\item Le foncteur d'inclusion
$\QCoh(\mathcal{O}_\mathcal{X}) \to
\textit{Mod}(\mathcal{O}_\mathcal{X})$ admet un adjoint à droite\footnote{Ce
foncteur est parfois appelé le {\it cohérateur}.}
$$
Q :
\textit{Mod}(\mathcal{O}_\mathcal{X})
\longrightarrow
\QCoh(\mathcal{O}_\mathcal{X})
$$
tel que, pour tout faisceau quasi-cohérent $\mathcal{F}$, le morphisme d'adjonction
$Q(\mathcal{F}) \to \mathcal{F}$ soit un isomorphisme.
\end{enumerate}
\end{proposition}

\begin{proof}
Cette preuve reprend celle du cas des schémas, voir
Propriétés, Proposition \ref{properties-proposition-coherator}
et celle du cas des espaces algébriques, voir
Propriétés des espaces, Proposition
\ref{spaces-properties-proposition-coherator}.
Nous conseillons au lecteur de lire d'abord l'une de ces preuves.

\medskip\noindent
L'assertion (1) signifie que $\QCoh(\mathcal{O}_\mathcal{X})$ (a) admet toutes les limites inductives,
(b) que les limites inductives filtrantes sont exactes, et (c) qu'elle admet un générateur, voir
Injectifs, section \ref{injectives-section-grothendieck-conditions}.
D'après le Lemme \ref{stacks-sheaves-lemma-quasi-coherent-algebraic-stack},
les limites inductives dans $\QCoh(\mathcal{O}_X)$ existent et coïncident
avec celles de $\textit{Mod}(\mathcal{O}_X)$. D'après
Modules sur les sites, Lemme \ref{sites-modules-lemma-limits-colimits}
les limites inductives filtrantes sont exactes. Ainsi (a) et (b) sont satisfaites.

\medskip\noindent
Choisissons une présentation $\mathcal{X} = [U/R]$ telle que $(U, R, s, t, c)$
soit un groupoïde lisse en espaces algébriques ; en particulier, $s$ et $t$
sont des morphismes plats d'espaces algébriques. D'après le
Lemme \ref{stacks-sheaves-lemma-quasi-coherent-algebraic-stack}
ci-dessus, nous avons
$\QCoh(\mathcal{O}_\mathcal{X}) = \QCoh(U, R, s, t, c)$.
D'après Groupoïdes dans les espaces, Lemme \ref{spaces-groupoids-lemma-set-generators},
il existe un ensemble $T$ et une famille $(\mathcal{F}_t)_{t \in T}$ de
faisceaux quasi-cohérents sur $\mathcal{X}$ tels que tout faisceau quasi-cohérent
sur $\mathcal{X}$ soit la limite inductive filtrante de ses sous-faisceaux
isomorphes à l'un des $\mathcal{F}_t$.
Ainsi $\bigoplus_t \mathcal{F}_t$ est
un générateur de $\QCoh(\mathcal{O}_X)$, et (c) est satisfait.
Les assertions sur les limites et les injectifs valent dans toute
catégorie abélienne de Grothendieck, voir
Injectifs, Théorème
\ref{injectives-theorem-injective-embedding-grothendieck} et
Lemme \ref{injectives-lemma-grothendieck-products}.

\medskip\noindent
Preuve de (2). Pour construire $Q$, nous employons le procédé général suivant.
Étant donné un objet $\mathcal{F}$ de $\textit{Mod}(\mathcal{O}_\mathcal{X})$,
considérons le foncteur
$$
\QCoh(\mathcal{O}_\mathcal{X})^{opp}
\longrightarrow
\textit{Ensembles},
\quad
\mathcal{G}
\longmapsto
\Hom_\mathcal{X}(\mathcal{G}, \mathcal{F})
$$
Ce foncteur transforme les limites inductives en limites projectives ;
il est donc représentable, voir
Injectifs, Lemme \ref{injectives-lemma-grothendieck-brown}.
Il existe donc un faisceau quasi-cohérent $Q(\mathcal{F})$
et un isomorphisme fonctoriel
$\Hom_\mathcal{X}(\mathcal{G}, \mathcal{F}) =
\Hom_\mathcal{X}(\mathcal{G}, Q(\mathcal{F}))$
pour $\mathcal{G}$ dans $\QCoh(\mathcal{O}_\mathcal{X})$.
D'après le lemme de Yoneda
(Catégories, Lemme \ref{categories-lemma-yoneda})
la construction $\mathcal{F} \leadsto Q(\mathcal{F})$ est
fonctorielle en $\mathcal{F}$. Par construction, $Q$ est un adjoint
à droite du foncteur d'inclusion.
Le fait que $Q(\mathcal{F}) \to \mathcal{F}$ soit un isomorphisme
lorsque $\mathcal{F}$ est quasi-cohérent est une conséquence formelle du fait
que le foncteur d'inclusion
$\QCoh(\mathcal{O}_\mathcal{X}) \to
\textit{Mod}(\mathcal{O}_\mathcal{X})$
est pleinement fidèle.
\end{proof}








\section{Cohomologie}
\label{stacks-sheaves-section-cohomology-general}

\noindent
Soit $S$ un schéma et soit $\mathcal{X}$ une catégorie fibrée en groupoïdes
sur $(\Sch/S)_{fppf}$. Pour tout $\tau \in \{Zariski, \etale, smooth,
syntomic, fppf\}$, les catégories $\textit{Ab}(\mathcal{X}_\tau)$ et
$\textit{Mod}(\mathcal{X}_\tau, \mathcal{O}_\mathcal{X})$ ont
suffisamment d'injectifs, voir
Injectifs, Théorèmes \ref{injectives-theorem-sheaves-injectives} et
\ref{injectives-theorem-sheaves-modules-injectives}.
Nous pouvons donc employer le formalisme de
Cohomologie sur les sites, section \ref{sites-cohomology-section-cohomology-sheaves}
pour définir les groupes de cohomologie
$$
H^p(\mathcal{X}_\tau, \mathcal{F}) = H^p_\tau(\mathcal{X}, \mathcal{F})
\quad\text{et}\quad
H^p(x, \mathcal{F}) = H^p_\tau(x, \mathcal{F})
$$
pour tout $x \in \Ob(\mathcal{X})$ et tout objet $\mathcal{F}$ de
$\textit{Ab}(\mathcal{X}_\tau)$ ou
$\textit{Mod}(\mathcal{X}_\tau, \mathcal{O}_\mathcal{X})$. De plus, si
$f : \mathcal{X} \to \mathcal{Y}$ est un $1$-morphisme de catégories fibrées
en groupoïdes sur $(\Sch/S)_{fppf}$, on obtient les images directes
supérieures $R^if_*\mathcal{F}$ dans $\textit{Ab}(\mathcal{Y}_\tau)$ ou
$\textit{Mod}(\mathcal{Y}_\tau, \mathcal{O}_\mathcal{Y})$.
Bien entendu, comme expliqué dans
Cohomologie sur les sites, section \ref{sites-cohomology-section-derived-functors}
il existe aussi des versions dérivées de $H^p(-)$ et de $R^if_*$.

\begin{lemma}
\label{stacks-sheaves-lemma-cohomology-restriction}
Soit $S$ un schéma. Soit $\mathcal{X}$ une catégorie fibrée en groupoïdes
sur $(\Sch/S)_{fppf}$. Soit $\tau \in \{Zariski, \etale, smooth,
syntomic, fppf\}$. Soit $x \in \Ob(\mathcal{X})$ un objet au-dessus
du schéma $U$. Soit $\mathcal{F}$
un objet de $\textit{Ab}(\mathcal{X}_\tau)$ ou de
$\textit{Mod}(\mathcal{X}_\tau, \mathcal{O}_\mathcal{X})$. Alors
$$
H^p_\tau(x, \mathcal{F}) = H^p((\Sch/U)_\tau, x^{-1}\mathcal{F})
$$
et, si $\tau = \etale$, nous avons aussi
$$
H^p_\etale(x, \mathcal{F}) =
H^p(U_\etale, \mathcal{F}|_{U_\etale}).
$$
\end{lemma}

\begin{proof}
La première assertion résulte de
Cohomologie sur les sites, Lemme \ref{sites-cohomology-lemma-cohomology-of-open}
et de l'équivalence du
Lemme \ref{stacks-sheaves-lemma-localizing-structure-sheaf}.
La seconde assertion résulte de la première, combinée avec
Cohomologie \'etale, Lemme
\ref{etale-cohomology-lemma-compare-cohomology-big-small}.
\end{proof}














\section{Faisceaux injectifs}
\label{stacks-sheaves-section-lower-shriek}

\noindent
L'image directe d'un faisceau abélien injectif ou d'un module injectif est injective.

\begin{lemma}
\label{stacks-sheaves-lemma-pushforward-injective}
Soit $f : \mathcal{X} \to \mathcal{Y}$ un $1$-morphisme de catégories
fibrées en groupoïdes sur $(\Sch/S)_{fppf}$. Soit
$\tau \in \{Zar, \etale, smooth, syntomic, fppf\}$.
\begin{enumerate}
\item $f_*\mathcal{I}$ est injectif dans $\textit{Ab}(\mathcal{Y}_\tau)$
pour $\mathcal{I}$ injectif dans $\textit{Ab}(\mathcal{X}_\tau)$, et
\item $f_*\mathcal{I}$ est injectif dans
$\textit{Mod}(\mathcal{Y}_\tau, \mathcal{O}_\mathcal{Y})$
pour $\mathcal{I}$ injectif dans
$\textit{Mod}(\mathcal{X}_\tau, \mathcal{O}_\mathcal{X})$.
\end{enumerate}
\end{lemma}

\begin{proof}
Cela résulte formellement du fait que $f^{-1}$ est un adjoint à gauche
exact de $f_*$, voir
Homologie, Lemme \ref{homology-lemma-adjoint-preserve-injectives}.
\end{proof}

\noindent
Dans la suite de cette section, nous prouvons que l'image réciproque $f^{-1}$ admet un
adjoint à gauche $f_!$ sur les faisceaux abéliens et les modules. Si $f$ est représentable (par
des schémas ou des espaces algébriques), il se trouvera que $f_!$ est exact
et que $f^{-1}$ préserve les injectifs. Nous démontrons d'abord quelques
lemmes préliminaires sur les produits fibrés et les égalisateurs dans les catégories
fibrées en groupoïdes, ainsi que sur leur comportement à l'égard des morphismes.

\begin{lemma}
\label{stacks-sheaves-lemma-fibre-products}
Soit $p : \mathcal{X} \to (\Sch/S)_{fppf}$
une catégorie fibrée en groupoïdes.
\begin{enumerate}
\item La catégorie $\mathcal{X}$ admet des produits fibrés.
\item Si les préfaisceaux $\mathit{Isom}$ de $\mathcal{X}$
sont représentables par des espaces algébriques, alors $\mathcal{X}$ admet des égalisateurs.
\item Si $\mathcal{X}$ est un champ algébrique (ou, plus généralement,
un champ quotient), alors $\mathcal{X}$ admet des égalisateurs.
\end{enumerate}
\end{lemma}

\begin{proof}
L'assertion (1) résulte de
Catégories, Lemme \ref{categories-lemma-fibred-groupoids-fibre-product-goes-up}
puisque $(\Sch/S)_{fppf}$ admet des produits fibrés.

\medskip\noindent
Soient $a, b : x \to y$ des morphismes de $\mathcal{X}$.
Posons $U = p(x)$ et $V = p(y)$. La catégorie des schémas admet des égalisateurs ;
on peut donc prendre pour $W \to U$ l'égalisateur de $p(a)$ et $p(b)$.
Notons $c : z \to x$ un morphisme de $\mathcal{X}$ au-dessus de $W \to U$.
L'égalisateur de $a$ et $b$, s'il existe, est celui de $a \circ c$
et $b \circ c$. Nous pouvons donc supposer que $p(a) = p(b) = f : U \to V$.
Comme $\mathcal{X}$ est fibrée en groupoïdes, il existe un unique automorphisme
$i : x \to x$ dans la catégorie fibre de $\mathcal{X}$ au-dessus de $U$ tel que
$a \circ i = b$. De nouveau, l'égalisateur de $a$ et $b$ est celui de
$\text{id}_x$ et $i$. Rappelons que $\mathit{Isom}_\mathcal{X}(x)$
est le préfaisceau sur $(\Sch/U)_{fppf}$ qui associe à
$T/U$ l'ensemble des automorphismes de $x|_T$ dans la catégorie fibre
de $\mathcal{X}$ au-dessus de $T$, voir
Champs, Définition \ref{stacks-definition-mor-presheaf}.
Si $\mathit{Isom}_\mathcal{X}(x)$ est représentable par un espace algébrique
$G \to U$, alors $\text{id}_x$ et $i$ définissent des morphismes
$e, i : U \to G$ au-dessus de $U$. Posons $M = U \times_{e, G, i} U$, qui, d'après
Morphismes d'espaces, Lemme \ref{spaces-morphisms-lemma-section-immersion}
est un schéma. Il est alors clair que $x|_M \to x$ est l'égalisateur des
morphismes $\text{id}_x$ et $i$ dans $\mathcal{X}$.
Cela prouve (2).

\medskip\noindent
Si $\mathcal{X} = [U/R]$ pour un groupoïde en espaces algébriques
$(U, R, s, t, c)$ sur $S$, l'hypothèse de (2) est satisfaite d'après
Amorçage, Lemme \ref{bootstrap-lemma-quotient-stack-isom}.
Si $\mathcal{X}$ est un champ algébrique, nous pouvons choisir une
présentation $[U/R] \cong \mathcal{X}$ d'après
Champs algébriques, Lemme \ref{algebraic-lemma-stack-presentation}.
\end{proof}

\begin{lemma}
\label{stacks-sheaves-lemma-fibre-products-morphism}
Soit $f : \mathcal{X} \to \mathcal{Y}$ un $1$-morphisme de catégories
fibrées en groupoïdes sur $(\Sch/S)_{fppf}$.
\begin{enumerate}
\item Le foncteur $f$ transforme les produits fibrés en produits fibrés.
\item Si $f$ est fidèle, alors $f$ transforme les égalisateurs en égalisateurs.
\end{enumerate}
\end{lemma}

\begin{proof}
D'après
Catégories, Lemme \ref{categories-lemma-fibred-groupoids-fibre-product-goes-up}
un produit fibré dans $\mathcal{X}$ est tout carré commutatif au-dessus
d'un diagramme de produit fibré dans $(\Sch/S)_{fppf}$. Il en va de même pour
$\mathcal{Y}$. L'assertion (1) est donc claire.

\medskip\noindent
Soit $x \to x'$ l'égalisateur de deux morphismes $a, b : x' \to x''$
dans $\mathcal{X}$. Montrons que $f(x) \to f(x')$ est l'égalisateur
de $f(a)$ et $f(b)$. Soit $y \to f(x')$ un morphisme de $\mathcal{Y}$
égalisant $f(a)$ et $f(b)$. Disons que $x, x', x''$ sont au-dessus des schémas
$U, U', U''$ et que $y$ est au-dessus de $V$. Notons $h : V \to U'$ l'image
de $y \to f(x')$ dans la catégorie des schémas. D'après les axiomes des catégories
fibrées, le morphisme $y \to f(x')$ est isomorphe à $f(h^*x') \to f(x')$.
Comme $f$ est fidèle, nous voyons donc que
$h^*x' \to x'$ égalise $a$ et $b$. Nous obtenons ainsi un unique morphisme
$h^*x' \to x$ dont l'image $y = f(h^*x') \to f(x)$ est le morphisme cherché
dans $\mathcal{Y}$.
\end{proof}

\begin{lemma}
\label{stacks-sheaves-lemma-fibre-products-preserve-properties}
Soient $f : \mathcal{X} \to \mathcal{Y}$ et $g : \mathcal{Z} \to \mathcal{Y}$
des $1$-morphismes fidèles de catégories
fibrées en groupoïdes sur $(\Sch/S)_{fppf}$.
\begin{enumerate}
\item le foncteur $\mathcal{X} \times_\mathcal{Y} \mathcal{Z} \to \mathcal{Y}$
est fidèle, et
\item si $\mathcal{X}, \mathcal{Z}$ admettent des égalisateurs, il en va de même de
$\mathcal{X} \times_\mathcal{Y} \mathcal{Z}$.
\end{enumerate}
\end{lemma}

\begin{proof}
Nous considérons les objets de $\mathcal{X} \times_\mathcal{Y} \mathcal{Z}$ comme des
quadruplets $(U, x, z, \alpha)$, où $\alpha : f(x) \to g(z)$ est un
isomorphisme au-dessus de $U$, voir
Catégories, Lemme \ref{categories-lemma-2-product-categories-over-C}.
Un morphisme $(U, x, z, \alpha) \to (U', x', z', \alpha')$ est une
paire de morphismes $a : x \to x'$ et $b : z \to z'$ compatibles
à $\alpha$ et $\alpha'$. Il est donc clair que, si $f$ et
$g$ sont fidèles, le foncteur
$\mathcal{X} \times_\mathcal{Y} \mathcal{Z} \to \mathcal{Y}$
est fidèle. Supposons maintenant que
$(a, b), (a', b') : (U, x, z, \alpha) \to (U', x', z', \alpha')$
soient deux morphismes du $2$-produit fibré. Considérons l'égalisateur
$x'' \to x$ de $a$ et $a'$, et l'égalisateur $z'' \to z$ de $b$ et $b'$.
Comme $f$ commute aux égalisateurs (d'après le
Lemme \ref{stacks-sheaves-lemma-fibre-products-morphism})
nous voyons que $f(x'') \to f(x)$ est l'égalisateur de $f(a)$ et $f(a')$.
De même, $g(z'') \to g(z)$ est l'égalisateur de $g(b)$ et $g(b')$.
Considérons le diagramme
$$
\xymatrix{
f(x'') \ar[r] \ar@{..>}[d]_{\alpha''}&
f(x) \ar[d]_\alpha
\ar@<0.5ex>[r]^{f(a)}
\ar@<-0.5ex>[r]_{f(a')}
 &
f(x') \ar[d]^{\alpha'} \\
g(z'') \ar[r] &
g(z)
\ar@<0.5ex>[r]^{g(b)}
\ar@<-0.5ex>[r]_{g(b')}
 &
g(z')
}
$$
Il est clair que la flèche pointillée existe et est un isomorphisme.
Toutefois, il n'est pas assuré a priori que l'image de $\alpha''$
dans la catégorie des schémas soit l'identité de sa source. D'autre
part, l'existence de $\alpha''$ signifie que nous pouvons supposer que $x''$
et $z''$ sont définis sur le même schéma et que les morphismes
$x'' \to x$ et $z'' \to z$ ont la même image dans la catégorie des schémas.
En refaisant le diagramme ci-dessus, nous voyons que la flèche pointillée se
projette alors bien sur un morphisme identité, ce qui conclut. Quelques détails sont omis.
\end{proof}

\noindent
Puisque nous travaillons avec de grands sites, nous obtenons le résultat quelque peu
contre-intuitif suivant (qui vaut aussi pour les morphismes entre grands sites
de schémas). Attention : ce résultat est faux si l'on omet l'hypothèse
que $f$ est fidèle.

\begin{lemma}
\label{stacks-sheaves-lemma-pullback-injective}
Soit $f : \mathcal{X} \to \mathcal{Y}$ un $1$-morphisme de catégories
fibrées en groupoïdes sur $(\Sch/S)_{fppf}$. Soit
$\tau \in \{Zar, \etale, smooth, syntomic, fppf\}$.
Le foncteur
$f^{-1} : \textit{Ab}(\mathcal{Y}_\tau) \to \textit{Ab}(\mathcal{X}_\tau)$
admet un adjoint à gauche
$f_! : \textit{Ab}(\mathcal{X}_\tau) \to \textit{Ab}(\mathcal{Y}_\tau)$.
Si $f$ est fidèle et si $\mathcal{X}$ admet des égalisateurs, alors
\begin{enumerate}
\item $f_!$ est exact, et
\item $f^{-1}\mathcal{I}$ est injectif dans $\textit{Ab}(\mathcal{X}_\tau)$
pour $\mathcal{I}$ injectif dans $\textit{Ab}(\mathcal{Y}_\tau)$.
\end{enumerate}
\end{lemma}

\begin{proof}
D'après
Champs, Lemme \ref{stacks-lemma-topology-inherited-functorial}
le foncteur $f$ est continu et cocontinu. Par conséquent, d'après
Modules sur les sites, Lemme \ref{sites-modules-lemma-g-shriek-adjoint}
le foncteur
$f^{-1} : \textit{Ab}(\mathcal{Y}_\tau) \to \textit{Ab}(\mathcal{X}_\tau)$
admet un adjoint à gauche
$f_! : \textit{Ab}(\mathcal{X}_\tau) \to \textit{Ab}(\mathcal{Y}_\tau)$.
Pour (1), nous appliquons
Modules sur les sites, Lemme \ref{sites-modules-lemma-exactness-lower-shriek}
et, pour vérifier les hypothèses de ce lemme, nous employons les
Lemmes \ref{stacks-sheaves-lemma-fibre-products} et
\ref{stacks-sheaves-lemma-fibre-products-morphism}
ci-dessus. L'assertion (2) en résulte formellement, voir
Homologie, Lemme \ref{homology-lemma-adjoint-preserve-injectives}.
\end{proof}

\begin{lemma}
\label{stacks-sheaves-lemma-pullback-injective-modules}
Soit $f : \mathcal{X} \to \mathcal{Y}$ un $1$-morphisme de catégories
fibrées en groupoïdes sur $(\Sch/S)_{fppf}$. Soit
$\tau \in \{Zar, \etale, smooth, syntomic, fppf\}$.
Le foncteur
$f^* : \textit{Mod}(\mathcal{Y}_\tau, \mathcal{O}_\mathcal{Y}) \to
\textit{Mod}(\mathcal{X}_\tau, \mathcal{O}_\mathcal{X})$
admet un adjoint à gauche
$f_! : \textit{Mod}(\mathcal{X}_\tau, \mathcal{O}_\mathcal{X}) \to
\textit{Mod}(\mathcal{Y}_\tau, \mathcal{O}_\mathcal{Y})$, qui
coïncide avec le foncteur $f_!$ du Lemme \ref{stacks-sheaves-lemma-pullback-injective}
sur les faisceaux abéliens sous-jacents.
Si $f$ est fidèle et si $\mathcal{X}$ admet des égalisateurs, alors
\begin{enumerate}
\item $f_!$ est exact, et
\item $f^{-1}\mathcal{I}$ est injectif dans
$\textit{Mod}(\mathcal{X}_\tau, \mathcal{O}_\mathcal{X})$
pour $\mathcal{I}$ injectif dans
$\textit{Mod}(\mathcal{Y}_\tau, \mathcal{O}_\mathcal{X})$.
\end{enumerate}
\end{lemma}

\begin{proof}
Rappelons que $f$ est un foncteur continu et cocontinu de sites
et que $f^{-1}\mathcal{O}_\mathcal{Y} = \mathcal{O}_\mathcal{X}$. Ainsi
Modules sur les sites, Lemme \ref{sites-modules-lemma-lower-shriek-modules}
implique que $f^*$ admet un adjoint à gauche $f_!^{Mod}$.
Soit $x$ un objet de $\mathcal{X}$ au-dessus du schéma $U$.
Alors $f$ induit une équivalence de sites annelés
$$
\mathcal{X}/x \longrightarrow \mathcal{Y}/f(x)
$$
puisque les deux membres sont équivalents à $(\Sch/U)_\tau$, voir
Lemme \ref{stacks-sheaves-lemma-localizing-structure-sheaf}.
Modules sur les sites, Remarque \ref{sites-modules-remark-when-shriek-equal}
montre que $f_!$ coïncide avec le foncteur sur les faisceaux abéliens.

\medskip\noindent
Supposons maintenant que $\mathcal{X}$ admette des égalisateurs et que $f$ soit fidèle.
Le Lemme \ref{stacks-sheaves-lemma-pullback-injective}
nous dit que $f_!$ est exact. Enfin,
Homologie, Lemme \ref{homology-lemma-adjoint-preserve-injectives}
implique l'assertion sur les images réciproques de modules injectifs.
\end{proof}




\section{Le complexe de {\v C}ech}
\label{stacks-sheaves-section-cech}

\noindent
Pour calculer la cohomologie d'un faisceau sur un champ algébrique, nous la comparons
à la cohomologie du faisceau restreint à des recouvrements du
champ algébrique donné.

\medskip\noindent
Dans toute cette section, la situation sera la suivante. On se donne
un $1$-morphisme de catégories fibrées en groupoïdes
\begin{equation}
\label{stacks-sheaves-equation-covering}
\vcenter{
\xymatrix{
\mathcal{U} \ar[rr]_f \ar[rd]_q & &  \mathcal{X} \ar[ld]^p \\
& (\Sch/S)_{fppf}
}
}
\end{equation}
Nous considérerons $\mathcal{U}$ comme un « recouvrement » de $\mathcal{X}$.
Nous voulons donc considérer l'objet simplicial
$$
\xymatrix{
\mathcal{U} \times_\mathcal{X} \mathcal{U} \times_\mathcal{X} \mathcal{U}
\ar@<1ex>[r]
\ar@<0ex>[r]
\ar@<-1ex>[r]
&
\mathcal{U} \times_\mathcal{X} \mathcal{U}
\ar@<0.5ex>[r]
\ar@<-0.5ex>[r]
&
\mathcal{U}
}
$$
dans la catégorie des catégories fibrées en groupoïdes sur
$(\Sch/S)_{fppf}$. Cependant, comme il s'agit d'une $(2, 1)$-catégorie et
non d'une catégorie, précisons explicitement ce que nous entendons. Soit
$\mathcal{U}_n$ la catégorie dont les objets sont les
$(u_0, \ldots, u_n, x, \alpha_0, \ldots, \alpha_n)$
où $\alpha_i : f(u_i) \to x$ est un isomorphisme dans $\mathcal{X}$.
Nous notons $f_n : \mathcal{U}_n \to \mathcal{X}$ le $1$-morphisme
qui associe à $(u_0, \ldots, u_n, x, \alpha_0, \ldots, \alpha_n)$
l'objet $x$. Remarquons que $\mathcal{U}_0 = \mathcal{U}$ et $f_0 = f$.
À un morphisme $\varphi : [m] \to [n]$, nous associons le $1$-morphisme
$\mathcal{U}_\varphi : \mathcal{U}_n \longrightarrow \mathcal{U}_n$
donné sur les objets par
$$
(u_0, \ldots, u_n, x, \alpha_0, \ldots, \alpha_n)
\longmapsto
(u_{\varphi(0)}, \ldots, u_{\varphi(m)}, x,
\alpha_{\varphi(0)}, \ldots, \alpha_{\varphi(m)})
$$
Tous ces $1$-morphismes se composent strictement comme il convient
(sans qu'il faille de $2$-morphismes), et chacun de ces $1$-morphismes est un $1$-morphisme
sur $\mathcal{X}$. Nous notons $\mathcal{U}_\bullet$ cet objet simplicial.
Si $\mathcal{F}$ est un préfaisceau d'ensembles sur $\mathcal{X}$, nous obtenons un
ensemble cosimplicial
$$
\xymatrix{
\Gamma(\mathcal{U}_0, f_0^{-1}\mathcal{F})
\ar@<0.5ex>[r]
\ar@<-0.5ex>[r]
&
\Gamma(\mathcal{U}_1, f_1^{-1}\mathcal{F})
\ar@<1ex>[r]
\ar@<0ex>[r]
\ar@<-1ex>[r]
&
\Gamma(\mathcal{U}_2, f_2^{-1}\mathcal{F})
}
$$
Ici, les flèches sont les morphismes d'image réciproque suivant les morphismes donnés de
l'objet simplicial.
Si $\mathcal{F}$ est un préfaisceau de groupes abéliens, on obtient un groupe
abélien cosimplicial.

\medskip\noindent
Soit $\mathcal{U} \to \mathcal{X}$ comme ci-dessus et soit $\mathcal{F}$
un préfaisceau abélien sur $\mathcal{X}$.
Le {\it complexe de {\v C}ech} associé à cette situation est noté
$\check{\mathcal{C}}^\bullet(\mathcal{U} \to \mathcal{X}, \mathcal{F})$.
C'est le complexe de cochaînes associé au groupe abélien cosimplicial
ci-dessus, voir
Simplicial, section \ref{simplicial-section-dold-kan-cosimplicial}.
Ses termes sont
$$
\check{\mathcal{C}}^n(\mathcal{U} \to \mathcal{X}, \mathcal{F})
= \Gamma(\mathcal{U}_n, f_n^{-1}\mathcal{F}).
$$
Les différentielles sont les morphismes
$$
d^n = \sum\nolimits_{i = 0}^{n + 1} (-1)^i \delta^{n + 1}_i :
\Gamma(\mathcal{U}_n, f_n^{-1}\mathcal{F})
\longrightarrow
\Gamma(\mathcal{U}_{n + 1}, f_{n + 1}^{-1}\mathcal{F})
$$
où $\delta^{n + 1}_i$ correspond au morphisme
$[n] \to [n + 1]$ qui omet l'indice $i$. Remarquons que le morphisme
$\Gamma(\mathcal{X}, \mathcal{F}) \to
\Gamma(\mathcal{U}_0, f_0^{-1}\mathcal{F}_0)$
est dans le noyau de la différentielle $d^0$. Nous définissons donc
le {\it complexe de {\v C}ech augmenté} comme le complexe
$$
\ldots \to 0 \to
\Gamma(\mathcal{X}, \mathcal{F}) \to
\Gamma(\mathcal{U}_0, f_0^{-1}\mathcal{F}_0) \to
\Gamma(\mathcal{U}_1, f_1^{-1}\mathcal{F}_1) \to \ldots
$$
où $\Gamma(\mathcal{X}, \mathcal{F})$ est placé en degré $-1$.
Le complexe de {\v C}ech augmenté est acyclique si et seulement si le morphisme canonique
$$
\Gamma(\mathcal{X}, \mathcal{F})[0]
\longrightarrow
\check{\mathcal{C}}^\bullet(\mathcal{U} \to \mathcal{X}, \mathcal{F})
$$
est un quasi-isomorphisme de complexes.

\begin{lemma}
\label{stacks-sheaves-lemma-generalities}
Généralités sur les complexes de {\v C}ech.
\begin{enumerate}
\item Si
$$
\xymatrix{
\mathcal{V} \ar[d]_g \ar[r]_h & \mathcal{U} \ar[d]^f \\
\mathcal{Y} \ar[r]^e & \mathcal{X}
}
$$
est un diagramme $2$-commutatif de catégories fibrées en groupoïdes sur
$(\Sch/S)_{fppf}$, alors il existe un morphisme de complexes de {\v C}ech
$$
\check{\mathcal{C}}^\bullet(\mathcal{U} \to \mathcal{X}, \mathcal{F})
\longrightarrow
\check{\mathcal{C}}^\bullet(\mathcal{V} \to \mathcal{Y}, e^{-1}\mathcal{F})
$$
\item si $h$ et $e$ sont des équivalences, le morphisme de (1) est un isomorphisme,
\item si $f, f' : \mathcal{U} \to \mathcal{X}$ sont $2$-isomorphes, alors
les complexes de {\v C}ech associés sont isomorphes.
\end{enumerate}
\end{lemma}

\begin{proof}
Dans la situation de (1), soit $t : f \circ h \to e \circ g$ un $2$-morphisme.
Le morphisme de complexes est donné en degré $n$ par
l'image réciproque suivant les $1$-morphismes
$\mathcal{V}_n \to \mathcal{U}_n$ définis par la règle
$$
(v_0, \ldots, v_n, y, \beta_0, \ldots, \beta_n)
\longmapsto
(h(v_0), \ldots, h(v_n), e(y),
e(\beta_0) \circ t_{v_0}, \ldots, e(\beta_n) \circ t_{v_n}).
$$
Pour (2), remarquons que l'image réciproque sur les sections globales est un isomorphisme
pour tout préfaisceau d'ensembles lorsqu'elle est prise suivant une équivalence
de catégories. L'assertion (3) résulte de la combinaison de (1) et (2).
\end{proof}

\begin{lemma}
\label{stacks-sheaves-lemma-homotopy}
S'il existe un $1$-morphisme $s : \mathcal{X} \to \mathcal{U}$
tel que $f \circ s$ soit $2$-isomorphe à $\text{id}_\mathcal{X}$,
alors le complexe de {\v C}ech augmenté est homotope à zéro.
\end{lemma}

\begin{proof}
Posons $\mathcal{U}' = \mathcal{U} \times_\mathcal{X} \mathcal{X}$,
égal au produit fibré décrit dans
Catégories, Lemme \ref{categories-lemma-2-product-categories-over-C}.
Posons $f' : \mathcal{U}' \to \mathcal{X}$ égal à la seconde projection.
Alors $\mathcal{U} \to \mathcal{U}'$, $u \mapsto (u, f(x), 1)$
est une équivalence sur $\mathcal{X}$ ; nous pouvons donc remplacer
$(\mathcal{U}, f)$ par $(\mathcal{U}', f')$ d'après le
Lemme \ref{stacks-sheaves-lemma-generalities}.
L'avantage est que $f'$ admet maintenant une section $s'$ telle
que $f' \circ s' = \text{id}_\mathcal{X}$ strictement. En effet, si
$t : s \circ f \to \text{id}_\mathcal{X}$ est un $2$-isomorphisme,
nous pouvons poser $s'(x) = (s(x), x, t_x)$. Nous pouvons donc supposer que
$f \circ s = \text{id}_\mathcal{X}$.

\medskip\noindent
Lorsque $f \circ s = \text{id}_\mathcal{X}$, le résultat découle
de principes généraux. Explicitons l'homotopie. Pour
$n \geq 0$, définissons $s_n : \mathcal{U}_n \to \mathcal{U}_{n + 1}$
comme le $1$-morphisme donné sur les objets par la règle
$$
(u_0, \ldots, u_n, x, \alpha_0, \ldots, \alpha_n)
\longmapsto
(u_0, \ldots, u_n, s(x), x,
\alpha_0, \ldots, \alpha_n, \text{id}_x).
$$
Définissons
$$
h^{n + 1} :
\Gamma(\mathcal{U}_{n + 1}, f_{n + 1}^{-1}\mathcal{F})
\longrightarrow
\Gamma(\mathcal{U}_n, f_n^{-1}\mathcal{F})
$$
comme l'image réciproque suivant $s_n$. Posons aussi $s_{-1} = s$ et
$h^0 : \Gamma(\mathcal{U}_0, f_0^{-1}\mathcal{F}) \to
\Gamma(\mathcal{X}, \mathcal{F})$ égal à l'image réciproque suivant $s_{-1}$.
Alors la famille de morphismes $\{h^n\}_{n \geq 0}$ est une homotopie entre
$1$ et $0$ sur le complexe de {\v C}ech augmenté.
\end{proof}









\section{Le complexe de {\v C}ech relatif}
\label{stacks-sheaves-section-sheaf-cech-complex}

\noindent
Soit $f : \mathcal{U} \to \mathcal{X}$ un $1$-morphisme de catégories
fibrées en groupoïdes sur $(\Sch/S)_{fppf}$ comme dans
(\ref{stacks-sheaves-equation-covering}). Considérons l'objet simplicial
associé $\mathcal{U}_\bullet$ et les morphismes
$f_n : \mathcal{U}_n \to \mathcal{X}$. Soit
$\tau \in \{Zar, \etale, smooth, syntomic, fppf\}$.
Enfin, supposons que $\mathcal{F}$ soit un faisceau (d'ensembles)
sur $\mathcal{X}_\tau$. Alors
$$
\xymatrix{
f_{0, *}f_0^{-1}\mathcal{F}
\ar@<0.5ex>[r]
\ar@<-0.5ex>[r]
&
f_{1, *}f_1^{-1}\mathcal{F}
\ar@<1ex>[r]
\ar@<0ex>[r]
\ar@<-1ex>[r]
&
f_{2, *}f_2^{-1}\mathcal{F}
}
$$
est un faisceau cosimplicial sur $\mathcal{X}_\tau$, où nous employons les morphismes d'image réciproque
introduits dans
Sites, section \ref{sites-section-pullback}.
Si $\mathcal{F}$ est un faisceau abélien, les $f_{n, *}f_n^{-1}\mathcal{F}$
forment un faisceau abélien cosimplicial sur $\mathcal{X}_\tau$.
Le complexe associé (voir
Simplicial, section \ref{simplicial-section-dold-kan-cosimplicial})
$$
\ldots \to 0 \to
f_{0, *}f_0^{-1}\mathcal{F} \to
f_{1, *}f_1^{-1}\mathcal{F} \to
f_{2, *}f_2^{-1}\mathcal{F} \to \ldots
$$
est appelé le {\it complexe de {\v C}ech relatif} associé à la situation.
Nous noterons ce complexe $\mathcal{K}^\bullet(f, \mathcal{F})$.
Le {\it complexe de {\v C}ech relatif augmenté} est le complexe
$$
\ldots \to 0 \to
\mathcal{F} \to
f_{0, *}f_0^{-1}\mathcal{F} \to
f_{1, *}f_1^{-1}\mathcal{F} \to
f_{2, *}f_2^{-1}\mathcal{F} \to \ldots
$$
où $\mathcal{F}$ est en degré $-1$. Le complexe de {\v C}ech relatif augmenté
est acyclique si et seulement si le morphisme
$\mathcal{F}[0] \to \mathcal{K}^\bullet(f, \mathcal{F})$
est un quasi-isomorphisme de complexes de faisceaux.

\begin{remark}
\label{stacks-sheaves-remark-cech-complex-presheaves}
Nous pouvons aussi définir le complexe $\mathcal{K}^\bullet(f, \mathcal{F})$
si $\mathcal{F}$ est un préfaisceau, mais nous ne pouvons alors employer
Sites, section \ref{sites-section-pullback}
pour définir les morphismes d'image réciproque. Pour expliquer ceux-ci, supposons donné
un diagramme commutatif
$$
\xymatrix{
\mathcal{V} \ar[rd]_g \ar[rr]_h & &  \mathcal{U} \ar[ld]^f \\
& \mathcal{X}
}
$$
de catégories fibrées en groupoïdes sur $(\Sch/S)_{fppf}$
et un préfaisceau $\mathcal{G}$ sur $\mathcal{U}$.
Nous pouvons définir le morphisme d'image réciproque $f_*\mathcal{G} \to g_*h^{-1}\mathcal{G}$
comme la composée
$$
f_*\mathcal{G} \longrightarrow
f_*h_*h^{-1}\mathcal{G} = g_*h^{-1}\mathcal{G}
$$
où le morphisme provient du morphisme d'adjonction
$\mathcal{G} \to h_*h^{-1}\mathcal{G}$. Cela fonctionne parce que, dans notre situation,
les foncteurs $h_*$ et $h^{-1}$ sont adjoints sur les préfaisceaux (et coïncident avec
leurs homologues sur les faisceaux). Voir les
sections \ref{stacks-sheaves-section-presheaves} et \ref{stacks-sheaves-section-sheaves}.
\end{remark}

\begin{lemma}
\label{stacks-sheaves-lemma-generalities-sheafified}
Généralités sur les complexes de {\v C}ech relatifs.
\begin{enumerate}
\item Si
$$
\xymatrix{
\mathcal{V} \ar[d]_g \ar[r]_h & \mathcal{U} \ar[d]^f \\
\mathcal{Y} \ar[r]^e & \mathcal{X}
}
$$
est un diagramme $2$-commutatif de catégories fibrées en groupoïdes sur
$(\Sch/S)_{fppf}$, alors il existe un morphisme
$e^{-1}\mathcal{K}^\bullet(f, \mathcal{F}) \to
\mathcal{K}^\bullet(g, e^{-1}\mathcal{F})$.
\item si $h$ et $e$ sont des équivalences, le morphisme de (1) est un isomorphisme,
\item si $f, f' : \mathcal{U} \to \mathcal{X}$ sont $2$-isomorphes, alors
les complexes de {\v C}ech relatifs associés sont isomorphes.
\end{enumerate}
\end{lemma}

\begin{proof}
Littéralement la même que la preuve du
Lemme \ref{stacks-sheaves-lemma-generalities}
en employant les morphismes d'image réciproque de la
Remarque \ref{stacks-sheaves-remark-cech-complex-presheaves}.
\end{proof}

\begin{lemma}
\label{stacks-sheaves-lemma-homotopy-sheafified}
S'il existe un $1$-morphisme $s : \mathcal{X} \to \mathcal{U}$
tel que $f \circ s$ soit $2$-isomorphe à $\text{id}_\mathcal{X}$,
alors le complexe de {\v C}ech relatif augmenté est homotope à zéro.
\end{lemma}

\begin{proof}
Littéralement la même que la preuve du
Lemme \ref{stacks-sheaves-lemma-homotopy}.
\end{proof}

\begin{remark}
\label{stacks-sheaves-remark-cech-complex-sections}
« Calculons » la valeur du complexe de {\v C}ech relatif en un
objet $x$ de $\mathcal{X}$. Écrivons $p(x) = U$.
Considérons le diagramme de $2$-produit fibré (qui sert à introduire
la notation $g : \mathcal{V} \to \mathcal{Y}$)
$$
\xymatrix{
\mathcal{V} \ar@{=}[r] \ar[d]_g &
(\Sch/U)_{fppf} \times_{x, \mathcal{X}} \mathcal{U} \ar[r] \ar[d] &
\mathcal{U} \ar[d]^f \\
\mathcal{Y} \ar@{=}[r] &
(\Sch/U)_{fppf} \ar[r]^-x & \mathcal{X}
}
$$
Remarquons que le morphisme $\mathcal{V}_n \to \mathcal{U}_n$ de la preuve du
Lemme \ref{stacks-sheaves-lemma-generalities}
induit une équivalence
$\mathcal{V}_n =
(\Sch/U)_{fppf} \times_{x, \mathcal{X}} \mathcal{U}_n$.
Il résulte donc de
(\ref{stacks-sheaves-equation-pushforward})
que
$$
\Gamma(x, \mathcal{K}^\bullet(f, \mathcal{F})) =
\check{\mathcal{C}}^\bullet(\mathcal{V} \to \mathcal{Y}, x^{-1}\mathcal{F})
$$
En d'autres termes : la valeur du complexe de {\v C}ech relatif en un objet $x$ de
$\mathcal{X}$ est le complexe de {\v C}ech du changement de base de $f$ à
$\mathcal{X}/x \cong (\Sch/U)_{fppf}$. Cela implique par exemple que le
Lemme \ref{stacks-sheaves-lemma-homotopy}
implique le
Lemme \ref{stacks-sheaves-lemma-homotopy-sheafified}
et, plus généralement, que les résultats sur le complexe de {\v C}ech (usuel) impliquent
les résultats correspondants pour le complexe de {\v C}ech relatif.
\end{remark}

\begin{lemma}
\label{stacks-sheaves-lemma-base-change-cech-complex}
Soit
$$
\xymatrix{
\mathcal{V} \ar[d]_g \ar[r]_h & \mathcal{U} \ar[d]^f \\
\mathcal{Y} \ar[r]^e & \mathcal{X}
}
$$
un $2$-produit fibré de catégories fibrées en groupoïdes sur
$(\Sch/S)_{fppf}$ et soit $\mathcal{F}$ un préfaisceau abélien
sur $\mathcal{X}$. Alors le morphisme
$e^{-1}\mathcal{K}^\bullet(f, \mathcal{F}) \to
\mathcal{K}^\bullet(g, e^{-1}\mathcal{F})$
du
Lemme \ref{stacks-sheaves-lemma-generalities-sheafified}
est un isomorphisme de complexes de préfaisceaux abéliens.
\end{lemma}

\begin{proof}
Soit $y$ un objet de $\mathcal{Y}$ au-dessus du schéma $T$.
Posons $x = e(y)$. Nous allons montrer que le morphisme induit un isomorphisme
sur les sections au-dessus de $y$. Remarquons que
$$
\Gamma(y, e^{-1}\mathcal{K}^\bullet(f, \mathcal{F})) =
\Gamma(x, \mathcal{K}^\bullet(f, \mathcal{F})) =
\check{\mathcal{C}}^\bullet(
(\Sch/T)_{fppf} \times_{x, \mathcal{X}} \mathcal{U} \to
(\Sch/T)_{fppf}, x^{-1}\mathcal{F})
$$
d'après la
Remarque \ref{stacks-sheaves-remark-cech-complex-sections}. D'autre part,
$$
\Gamma(y, \mathcal{K}^\bullet(g, e^{-1}\mathcal{F})) =
\check{\mathcal{C}}^\bullet(
(\Sch/T)_{fppf} \times_{y, \mathcal{Y}} \mathcal{V} \to
(\Sch/T)_{fppf}, y^{-1}e^{-1}\mathcal{F})
$$
également d'après la
Remarque \ref{stacks-sheaves-remark-cech-complex-sections}.
Remarquons que $y^{-1}e^{-1}\mathcal{F} = x^{-1}\mathcal{F}$
et que, puisque le diagramme est $2$-cartésien, le $1$-morphisme
$$
(\Sch/T)_{fppf} \times_{y, \mathcal{Y}} \mathcal{V} \to
(\Sch/T)_{fppf} \times_{x, \mathcal{X}} \mathcal{U}
$$
est une équivalence. Le morphisme sur les sections au-dessus de $y$ est donc un
isomorphisme d'après le
Lemme \ref{stacks-sheaves-lemma-generalities}.
\end{proof}

\noindent
L'exactitude peut se vérifier sur un « recouvrement ».

\begin{lemma}
\label{stacks-sheaves-lemma-check-exactness-covering}
Soit $f : \mathcal{U} \to \mathcal{X}$ un $1$-morphisme de catégories fibrées
en groupoïdes sur $(\Sch/S)_{fppf}$. Soit
$\tau \in \{Zar, \etale, smooth, syntomic, fppf\}$.
Soit
$$
\mathcal{F} \to \mathcal{G} \to \mathcal{H}
$$
un complexe dans $\textit{Ab}(\mathcal{X}_\tau)$. Supposons que
\begin{enumerate}
\item pour tout objet $x$ de $\mathcal{X}$, il existe un recouvrement
$\{x_i \to x\}$ dans $\mathcal{X}_\tau$ tel que chaque $x_i$ soit isomorphe
à $f(u_i)$ pour un objet $u_i$ de $\mathcal{U}$, et
\item $f^{-1}\mathcal{F} \to f^{-1}\mathcal{G} \to f^{-1}\mathcal{H}$ soit exact.
\end{enumerate}
Alors la suite $\mathcal{F} \to \mathcal{G} \to \mathcal{H}$
est exacte.
\end{lemma}

\begin{proof}
Soit $x$ un objet de $\mathcal{X}$ au-dessus du schéma $T$.
Considérons la suite
$x^{-1}\mathcal{F} \to x^{-1}\mathcal{G} \to x^{-1}\mathcal{H}$
de faisceaux abéliens sur $(\Sch/T)_\tau$. Il suffit de montrer
que cette suite est exacte. Par hypothèse, il existe un $\tau$-recouvrement
$\{T_i \to T\}$ tel que $x|_{T_i}$ soit isomorphe à $f(u_i)$ pour
un objet $u_i$ de $\mathcal{U}$ au-dessus de $T_i$ et que, de plus, la suite
$u_i^{-1}f^{-1}\mathcal{F} \to u_i^{-1}f^{-1}\mathcal{G} \to
u_i^{-1}f^{-1}\mathcal{H}$ de faisceaux abéliens sur $(\Sch/T_i)_\tau$
soit exacte. Comme
$u_i^{-1}f^{-1}\mathcal{F} = x^{-1}\mathcal{F}|_{(\Sch/T_i)_\tau}$
nous concluons que la suite
$x^{-1}\mathcal{F} \to x^{-1}\mathcal{G} \to x^{-1}\mathcal{H}$
devient exacte après localisation en chacun des membres d'un recouvrement ;
la suite est donc exacte.
\end{proof}

\begin{proposition}
\label{stacks-sheaves-proposition-exactness-cech-complex}
Soit $f : \mathcal{U} \to \mathcal{X}$ un $1$-morphisme de catégories fibrées
en groupoïdes sur $(\Sch/S)_{fppf}$. Soit
$\tau \in \{Zar, \etale, smooth, syntomic, fppf\}$.
Si
\begin{enumerate}
\item $\mathcal{F}$ est un faisceau abélien sur $\mathcal{X}_\tau$, et
\item pour tout objet $x$ de $\mathcal{X}$, il existe un recouvrement
$\{x_i \to x\}$ dans $\mathcal{X}_\tau$ tel que chaque $x_i$ soit isomorphe
à $f(u_i)$ pour un objet $u_i$ de $\mathcal{U}$,
\end{enumerate}
alors le complexe de {\v C}ech relatif augmenté
$$
\ldots \to 0 \to
\mathcal{F} \to
f_{0, *}f_0^{-1}\mathcal{F} \to
f_{1, *}f_1^{-1}\mathcal{F} \to
f_{2, *}f_2^{-1}\mathcal{F} \to \ldots
$$
est exact dans $\textit{Ab}(\mathcal{X}_\tau)$.
\end{proposition}

\begin{proof}
D'après le
Lemme \ref{stacks-sheaves-lemma-check-exactness-covering}
il suffit de vérifier l'exactitude après image réciproque sur $\mathcal{U}$.
D'après le
Lemme \ref{stacks-sheaves-lemma-base-change-cech-complex}
l'image réciproque du complexe de {\v C}ech relatif augmenté est isomorphe
au complexe de {\v C}ech relatif augmenté du morphisme
$\mathcal{U} \times_\mathcal{X} \mathcal{U} \to \mathcal{U}$
et d'un faisceau abélien sur $\mathcal{U}_\tau$. Comme il existe une section
$\Delta_{\mathcal{U}/\mathcal{X}} : \mathcal{U} \to
\mathcal{U} \times_\mathcal{X} \mathcal{U}$, l'exactitude résulte du
Lemme \ref{stacks-sheaves-lemma-homotopy-sheafified}.
\end{proof}

\noindent
Nous pouvons ainsi construire la suite spectrale de {\v C}ech vers la cohomologie
comme suit. Nous donnons d'abord une version technique précise. Dans la section suivante,
nous donnerons une version qui ne concerne que les champs algébriques.

\begin{lemma}
\label{stacks-sheaves-lemma-cech-to-cohomology}
Soit $f : \mathcal{U} \to \mathcal{X}$ un $1$-morphisme de catégories fibrées
en groupoïdes sur $(\Sch/S)_{fppf}$. Soit
$\tau \in \{Zar, \etale, smooth, syntomic, fppf\}$.
Supposons que
\begin{enumerate}
\item $\mathcal{F}$ soit un faisceau abélien sur $\mathcal{X}_\tau$,
\item pour tout objet $x$ de $\mathcal{X}$, il existe un recouvrement
$\{x_i \to x\}$ dans $\mathcal{X}_\tau$ tel que chaque $x_i$ soit isomorphe
à $f(u_i)$ pour un objet $u_i$ de $\mathcal{U}$,
\item la catégorie $\mathcal{U}$ admette des égalisateurs, et
\item le foncteur $f$ soit fidèle.
\end{enumerate}
Alors il existe une suite spectrale de groupes abéliens du premier quadrant
$$
E_1^{p, q} = H^q((\mathcal{U}_p)_\tau, f_p^{-1}\mathcal{F})
\Rightarrow
H^{p + q}(\mathcal{X}_\tau, \mathcal{F})
$$
qui converge vers la cohomologie de $\mathcal{F}$ pour la topologie $\tau$.
\end{lemma}

\begin{proof}
Avant de commencer la preuve, faisons quelques remarques. D'après le
Lemme \ref{stacks-sheaves-lemma-fibre-products-preserve-properties}
(et par récurrence), toutes les catégories fibrées en groupoïdes $\mathcal{U}_p$
admettent des égalisateurs et tous les morphismes $f_p : \mathcal{U}_p \to \mathcal{X}$
sont fidèles. Soit $\mathcal{I}$ un objet injectif
de $\textit{Ab}(\mathcal{X}_\tau)$. D'après le
Lemme \ref{stacks-sheaves-lemma-pullback-injective}
nous voyons que $f_p^{-1}\mathcal{I}$ est un objet injectif de
$\textit{Ab}((\mathcal{U}_p)_\tau)$.
Ainsi $f_{p, *}f_p^{-1}\mathcal{I}$ est un objet injectif de
$\textit{Ab}(\mathcal{X}_\tau)$ d'après le
Lemme \ref{stacks-sheaves-lemma-pushforward-injective}.
Par conséquent, la
Proposition \ref{stacks-sheaves-proposition-exactness-cech-complex}
montre que le complexe de {\v C}ech relatif augmenté
$$
\ldots \to 0 \to
\mathcal{I} \to
f_{0, *}f_0^{-1}\mathcal{I} \to
f_{1, *}f_1^{-1}\mathcal{I} \to
f_{2, *}f_2^{-1}\mathcal{I} \to \ldots
$$
est un complexe exact dans $\textit{Ab}(\mathcal{X}_\tau)$ dont tous les
termes sont injectifs. La prise des sections globales de ce complexe est exacte,
et nous voyons que le complexe de {\v C}ech
$\check{\mathcal{C}}^\bullet(\mathcal{U} \to \mathcal{X}, \mathcal{I})$
est quasi-isomorphe à $\Gamma(\mathcal{X}_\tau, \mathcal{I})[0]$.

\medskip\noindent
Après ces préliminaires, considérons les deux suites spectrales
associées au complexe double (voir
Homologie, section \ref{homology-section-double-complex})
$$
\check{\mathcal{C}}^\bullet(\mathcal{U} \to \mathcal{X}, \mathcal{I}^\bullet)
$$
où $\mathcal{F} \to \mathcal{I}^\bullet$ est une résolution injective
dans $\textit{Ab}(\mathcal{X}_\tau)$.
La discussion ci-dessus montre que
Homologie, Lemme \ref{homology-lemma-double-complex-gives-resolution}
s'applique et montre que
$\Gamma(\mathcal{X}_\tau, \mathcal{I}^\bullet)$
est quasi-isomorphe au complexe total associé au complexe double.
D'après les remarques ci-dessus, le complexe $f_p^{-1}\mathcal{I}^\bullet$ est une
résolution injective de $f_p^{-1}\mathcal{F}$. L'autre suite
spectrale est donc celle indiquée dans le lemme.
\end{proof}

\noindent
Il existe bien sûr aussi une version pour les modules.

\begin{lemma}
\label{stacks-sheaves-lemma-cech-to-cohomology-modules}
Soit $f : \mathcal{U} \to \mathcal{X}$ un $1$-morphisme de catégories fibrées
en groupoïdes sur $(\Sch/S)_{fppf}$. Soit
$\tau \in \{Zar, \etale, smooth, syntomic, fppf\}$.
Supposons que
\begin{enumerate}
\item $\mathcal{F}$ soit un objet de
$\textit{Mod}(\mathcal{X}_\tau, \mathcal{O}_\mathcal{X})$,
\item pour tout objet $x$ de $\mathcal{X}$, il existe un recouvrement
$\{x_i \to x\}$ dans $\mathcal{X}_\tau$ tel que chaque $x_i$ soit isomorphe
à $f(u_i)$ pour un objet $u_i$ de $\mathcal{U}$,
\item la catégorie $\mathcal{U}$ admette des égalisateurs, et
\item le foncteur $f$ soit fidèle.
\end{enumerate}
Alors il existe une suite spectrale du premier quadrant de
$\Gamma(\mathcal{O}_\mathcal{X})$-modules
$$
E_1^{p, q} = H^q((\mathcal{U}_p)_\tau, f_p^*\mathcal{F})
\Rightarrow
H^{p + q}(\mathcal{X}_\tau, \mathcal{F})
$$
qui converge vers la cohomologie de $\mathcal{F}$ pour la topologie $\tau$.
\end{lemma}

\begin{proof}
La preuve de ce lemme est identique à celle du
Lemme \ref{stacks-sheaves-lemma-cech-to-cohomology}
si ce n'est qu'elle emploie une résolution injective dans
$\textit{Mod}(\mathcal{X}_\tau, \mathcal{O}_\mathcal{X})$
et qu'elle emploie le
Lemme \ref{stacks-sheaves-lemma-pullback-injective-modules}
au lieu du
Lemme \ref{stacks-sheaves-lemma-pullback-injective}.
\end{proof}

\noindent
Voici un lemme qui ramène un type plus usuel de recouvrement aux
types de recouvrements rencontrés ci-dessus.

\begin{lemma}
\label{stacks-sheaves-lemma-surjective-flat-locally-finite-presentation}
Soit $f : \mathcal{X} \to \mathcal{Y}$ un $1$-morphisme de
catégories fibrées en groupoïdes sur $(\Sch/S)_{fppf}$.
\begin{enumerate}
\item Supposons $f$ représentable par des espaces algébriques, surjectif,
plat et localement de présentation finie. Alors, pour tout objet $y$ de
$\mathcal{Y}$, il existe un recouvrement fppf $\{y_i \to y\}$ et des objets
$x_i$ de $\mathcal{X}$ tels que $f(x_i) \cong y_i$ dans $\mathcal{Y}$.
\item Supposons $f$ représentable par des espaces algébriques, surjectif
et lisse. Alors, pour tout objet $y$ de
$\mathcal{Y}$, il existe un recouvrement \'etale $\{y_i \to y\}$ et des objets
$x_i$ de $\mathcal{X}$ tels que $f(x_i) \cong y_i$ dans $\mathcal{Y}$.
\end{enumerate}
\end{lemma}

\begin{proof}
Preuve de (1). Supposons que $y$ soit au-dessus du schéma $V$.
Nous pouvons voir $y$ comme un morphisme $(\Sch/V)_{fppf} \to \mathcal{Y}$.
Par définition, le $2$-produit fibré
$\mathcal{X} \times_\mathcal{Y} (\Sch/V)_{fppf}$
est représenté par un espace algébrique $W$, et le morphisme
$W \to V$ est surjectif, plat et localement de présentation finie.
Choisissons un schéma $U$ et un morphisme \'etale surjectif $U \to W$.
Alors $U \to V$ est aussi surjectif, plat et localement de présentation finie
(voir Morphismes d'espaces, Lemmes
\ref{spaces-morphisms-lemma-etale-flat},
\ref{spaces-morphisms-lemma-etale-locally-finite-presentation},
\ref{spaces-morphisms-lemma-composition-surjective},
\ref{spaces-morphisms-lemma-composition-finite-presentation}, et
\ref{spaces-morphisms-lemma-composition-flat}).
Ainsi $\{U \to V\}$ est un recouvrement fppf. Notons $x$ l'objet de
$\mathcal{X}$ au-dessus de $U$ correspondant au $1$-morphisme
$(\Sch/U)_{fppf} \to \mathcal{X}$. Alors $\{f(x) \to y\}$ est
le recouvrement fppf cherché de $\mathcal{Y}$.

\medskip\noindent
Preuve de (2). Supposons que $y$ soit au-dessus du schéma $V$.
Nous pouvons voir $y$ comme un morphisme $(\Sch/V)_{fppf} \to \mathcal{Y}$.
Par définition, le $2$-produit fibré
$\mathcal{X} \times_\mathcal{Y} (\Sch/V)_{fppf}$
est représenté par un espace algébrique $W$, et le morphisme
$W \to V$ est surjectif et lisse.
Choisissons un schéma $U$ et un morphisme \'etale surjectif $U \to W$.
Alors $U \to V$ est aussi surjectif et lisse
(voir Morphismes d'espaces, Lemmes
\ref{spaces-morphisms-lemma-etale-smooth},
\ref{spaces-morphisms-lemma-composition-surjective}, et
\ref{spaces-morphisms-lemma-composition-smooth}).
Ainsi $\{U \to V\}$ est un recouvrement lisse. D'après
Compléments sur les morphismes, Lemme \ref{more-morphisms-lemma-etale-dominates-smooth}
il existe un recouvrement \'etale $\{V_i \to V\}$ tel que
chaque $V_i \to V$ se factorise par $U$. Notons $x_i$ l'objet de
$\mathcal{X}$ au-dessus de $V_i$ correspondant au $1$-morphisme
$$
(\Sch/V_i)_{fppf} \to (\Sch/U)_{fppf} \to \mathcal{X}.
$$
Alors $\{f(x_i) \to y\}$ est
le recouvrement \'etale cherché de $\mathcal{Y}$.
\end{proof}

\begin{lemma}
\label{stacks-sheaves-lemma-cech-to-cohomology-relative}
Soient $f : \mathcal{U} \to \mathcal{X}$ et
$g : \mathcal{X} \to \mathcal{Y}$
des $1$-morphismes composables de catégories fibrées
en groupoïdes sur $(\Sch/S)_{fppf}$. Soit
$\tau \in \{Zar, \etale, smooth, syntomic, \linebreak[0] fppf\}$.
Supposons que
\begin{enumerate}
\item $\mathcal{F}$ soit un faisceau abélien sur $\mathcal{X}_\tau$,
\item pour tout objet $x$ de $\mathcal{X}$, il existe un recouvrement
$\{x_i \to x\}$ dans $\mathcal{X}_\tau$ tel que chaque $x_i$ soit isomorphe
à $f(u_i)$ pour un objet $u_i$ de $\mathcal{U}$,
\item la catégorie $\mathcal{U}$ admette des égalisateurs, et
\item le foncteur $f$ soit fidèle.
\end{enumerate}
Alors il existe une suite spectrale du premier quadrant de faisceaux abéliens
sur $\mathcal{Y}_\tau$
$$
E_1^{p, q} = R^q(g \circ f_p)_*f_p^{-1}\mathcal{F}
\Rightarrow
R^{p + q}g_*\mathcal{F}
$$
où toutes les images directes supérieures sont calculées pour la topologie $\tau$.
\end{lemma}

\begin{proof}
Remarquons que les hypothèses sur $f : \mathcal{U} \to \mathcal{X}$
et $\mathcal{F}$ sont identiques à celles du
Lemme \ref{stacks-sheaves-lemma-cech-to-cohomology}.
Les remarques préliminaires faites dans la preuve de ce lemme
s'appliquent donc ici aussi. Elles impliquent en particulier que
$$
0 \to g_*\mathcal{I} \to
(g \circ f_0)_*f_0^{-1}\mathcal{I} \to
(g \circ f_1)_*f_1^{-1}\mathcal{I} \to \ldots
$$
est exacte si $\mathcal{I}$ est un objet injectif de
$\textit{Ab}(\mathcal{X}_\tau)$.
Cela étant, considérons les deux suites spectrales de
Homologie, section \ref{homology-section-double-complex}
associées au complexe double $\mathcal{C}^{\bullet, \bullet}$ de termes
$$
\mathcal{C}^{p, q} = (g \circ f_p)_*\mathcal{I}^q
$$
où $\mathcal{F} \to \mathcal{I}^\bullet$ est une résolution injective
dans $\textit{Ab}(\mathcal{X}_\tau)$. La première suite spectrale implique, par
Homologie, Lemme \ref{homology-lemma-double-complex-gives-resolution},
que $g_*\mathcal{I}^\bullet$ est quasi-isomorphe au complexe total
associé à $\mathcal{C}^{\bullet, \bullet}$.
Comme $f_p^{-1}\mathcal{I}^\bullet$ est une résolution injective de
$f_p^{-1}\mathcal{F}$ (voir
Lemme \ref{stacks-sheaves-lemma-pullback-injective})
la seconde suite spectrale a pour termes
$E_1^{p, q} = R^q(g \circ f_p)_*f_p^{-1}\mathcal{F}$, comme dans l'énoncé
du lemme.
\end{proof}

\begin{lemma}
\label{stacks-sheaves-lemma-cech-to-cohomology-relative-modules}
Soient $f : \mathcal{U} \to \mathcal{X}$ et
$g : \mathcal{X} \to \mathcal{Y}$
des $1$-morphismes composables de catégories fibrées
en groupoïdes sur $(\Sch/S)_{fppf}$. Soit
$\tau \in \{Zar, \etale, smooth, syntomic, \linebreak[0] fppf\}$.
Supposons que
\begin{enumerate}
\item $\mathcal{F}$ soit un objet de
$\textit{Mod}(\mathcal{X}_\tau, \mathcal{O}_\mathcal{X})$,
\item pour tout objet $x$ de $\mathcal{X}$, il existe un recouvrement
$\{x_i \to x\}$ dans $\mathcal{X}_\tau$ tel que chaque $x_i$ soit isomorphe
à $f(u_i)$ pour un objet $u_i$ de $\mathcal{U}$,
\item la catégorie $\mathcal{U}$ admette des égalisateurs, et
\item le foncteur $f$ soit fidèle.
\end{enumerate}
Alors il existe une suite spectrale du premier quadrant dans
$\textit{Mod}(\mathcal{Y}_\tau, \mathcal{O}_\mathcal{Y})$
$$
E_1^{p, q} = R^q(g \circ f_p)_*f_p^{-1}\mathcal{F}
\Rightarrow
R^{p + q}g_*\mathcal{F}
$$
où toutes les images directes supérieures sont calculées pour la topologie $\tau$.
\end{lemma}

\begin{proof}
La preuve est identique à celle du
Lemme \ref{stacks-sheaves-lemma-cech-to-cohomology-relative}
si ce n'est qu'elle emploie une résolution injective dans
$\textit{Mod}(\mathcal{X}_\tau, \mathcal{O}_\mathcal{X})$
et qu'elle emploie le
Lemme \ref{stacks-sheaves-lemma-pullback-injective-modules}
au lieu du
Lemme \ref{stacks-sheaves-lemma-pullback-injective}.
\end{proof}






\section{Cohomologie des champs algébriques}
\label{stacks-sheaves-section-cohomology}

\noindent
Soit $\mathcal{X}$ un champ algébrique sur $S$. Dans les sections précédentes,
nous avons vu comment définir les faisceaux pour les topologies \'etale, ..., fppf
sur $\mathcal{X}$. En fait, nous avons construit un site
$\mathcal{X}_\tau$ pour chaque $\tau \in \{Zar, \etale, smooth, syntomic,
fppf\}$. Sur ces sites, on dispose de la notion de faisceau abélien $\mathcal{F}$.
Dans le chapitre sur la cohomologie des sites, nous avons expliqué comment définir la
cohomologie. En réunissant ces constructions, définissons les
{\it sections globales dérivées}, ou la {\it cohomologie totale},
$$
R\Gamma_{Zar}(\mathcal{X}, \mathcal{F}),
R\Gamma_\etale(\mathcal{X}, \mathcal{F}), \ldots,
R\Gamma_{fppf}(\mathcal{X}, \mathcal{F})
$$
comme $\Gamma(\mathcal{X}_\tau, \mathcal{I}^\bullet)$, où
$\mathcal{F} \to \mathcal{I}^\bullet$ est une résolution injective
dans $\textit{Ab}(\mathcal{X}_\tau)$. Le $i$-ème groupe de cohomologie de $\mathcal{F}$
est la $i$-ème cohomologie de la cohomologie totale. Nous le noterons
ainsi :
$$
H^i_{Zar}(\mathcal{X}, \mathcal{F}),
H^i_\etale(\mathcal{X}, \mathcal{F}), \ldots,
H^i_{fppf}(\mathcal{X}, \mathcal{F}).
$$
Il se trouvera que $H^i_\etale = H^i_{smooth}$ en vertu de
Compléments sur les morphismes, Lemme \ref{more-morphisms-lemma-etale-dominates-smooth}.

\medskip\noindent
Si $\mathcal{F}$ est un préfaisceau de $\mathcal{O}_\mathcal{X}$-modules
qui est un faisceau pour la topologie $\tau$, nous employons des résolutions
injectives dans $\textit{Mod}(\mathcal{X}_\tau, \mathcal{O}_\mathcal{X})$
pour calculer sa cohomologie totale, resp.\ ses groupes de cohomologie ;
le résultat est quasi-isomorphe, resp.\ isomorphe à la cohomologie de
$\mathcal{F}$ considéré comme faisceau de groupes abéliens, d'après le très général
Cohomologie sur les sites, Lemme
\ref{sites-cohomology-lemma-cohomology-modules-abelian-agree}.

\medskip\noindent
Jusqu'ici, notre seul outil de calcul des groupes de cohomologie est le résultat sur les
complexes de {\v C}ech démontré ci-dessus. Reformulons-le ici dans le langage
des champs algébriques pour les topologies \'etale et fppf. Soit
$f : \mathcal{U} \to \mathcal{X}$ un $1$-morphisme de champs algébriques.
Rappelons que
$$
f_p : \mathcal{U}_p =
\mathcal{U} \times_\mathcal{X} \ldots \times_\mathcal{X} \mathcal{U}
\longrightarrow
\mathcal{X}
$$
est le morphisme structural, avec $(p + 1)$ facteurs. Rappelons aussi
qu'un faisceau sur $\mathcal{X}$ est un faisceau pour la topologie fppf. Remarquons
que, si $\mathcal{U}$ est un espace algébrique, alors
$f : \mathcal{U} \to \mathcal{X}$ est représentable par des espaces algébriques,
voir
Champs algébriques, Lemme \ref{algebraic-lemma-representable-diagonal}.
La proposition s'applique donc en particulier à un recouvrement lisse du
champ algébrique $\mathcal{X}$ par un schéma.

\begin{proposition}
\label{stacks-sheaves-proposition-smooth-covering-compute-cohomology}
\begin{slogan}
La cohomologie d'un champ se calcule sur le nerf de Čech d'une présentation
\end{slogan}
Soit $f : \mathcal{U} \to \mathcal{X}$ un $1$-morphisme de champs algébriques.
\begin{enumerate}
\item Soit $\mathcal{F}$ un faisceau abélien \'etale sur $\mathcal{X}$.
Supposons $f$ représentable par des espaces algébriques, surjectif et lisse.
Alors il existe une suite spectrale
$$
E_1^{p, q} = H^q_\etale(\mathcal{U}_p, f_p^{-1}\mathcal{F})
\Rightarrow
H^{p + q}_\etale(\mathcal{X}, \mathcal{F})
$$
\item Soit $\mathcal{F}$ un faisceau abélien sur $\mathcal{X}$.
Supposons $f$ représentable par des espaces algébriques, surjectif, plat
et localement de présentation finie. Alors il existe
une suite spectrale
$$
E_1^{p, q} = H^q_{fppf}(\mathcal{U}_p, f_p^{-1}\mathcal{F})
\Rightarrow
H^{p + q}_{fppf}(\mathcal{X}, \mathcal{F})
$$
\end{enumerate}
\end{proposition}

\begin{proof}
Pour le voir, vérifions les hypothèses (1) -- (4) du
Lemme \ref{stacks-sheaves-lemma-cech-to-cohomology}.
Le $1$-morphisme $f$ est fidèle d'après
Champs algébriques, Lemme
\ref{algebraic-lemma-characterize-representable-by-algebraic-spaces}.
Cela prouve (4).
L'hypothèse (3) résulte du fait que $\mathcal{U}$ est un champ
algébrique, voir
Lemme \ref{stacks-sheaves-lemma-fibre-products}.
Pour (2), appliquons le
Lemme \ref{stacks-sheaves-lemma-surjective-flat-locally-finite-presentation}.
La condition (1) est satisfaite par hypothèse.
\end{proof}










\section{Images directes supérieures et champs algébriques}
\label{stacks-sheaves-section-higher-direct-images}

\noindent
Soit $g : \mathcal{X} \to \mathcal{Y}$ un $1$-morphisme de champs algébriques
sur $S$. Dans les sections précédentes, nous avons construit un morphisme de topos
annelés $g : \Sh(\mathcal{X}_\tau) \to \Sh(\mathcal{Y}_\tau)$
pour chaque $\tau \in \{Zar, \etale, smooth, syntomic, fppf\}$.
Dans le chapitre sur la cohomologie des sites, nous avons expliqué comment
définir les images directes supérieures. Ainsi l'{\it image directe totale}
$Rg_*\mathcal{F}$ est définie comme $g_*\mathcal{I}^\bullet$, où
$\mathcal{F} \to \mathcal{I}^\bullet$ est une résolution injective dans
$\textit{Ab}(\mathcal{X}_\tau)$. La $i$-ème image directe supérieure
$R^ig_*\mathcal{F}$ est la $i$-ème cohomologie de l'image directe totale.
Important : le choix de la topologie $\tau$ importe ici !

\medskip\noindent
Si $\mathcal{F}$ est un préfaisceau de $\mathcal{O}_\mathcal{X}$-modules
qui est un faisceau pour la topologie $\tau$, nous employons des résolutions
injectives dans $\textit{Mod}(\mathcal{X}_\tau, \mathcal{O}_\mathcal{X})$
pour calculer l'image directe totale et les images directes supérieures.

\medskip\noindent
Jusqu'ici, notre seul outil de calcul des images directes supérieures de $g_*$
est le résultat sur les complexes de {\v C}ech démontré ci-dessus. Il requiert le choix
d'un « recouvrement » $f : \mathcal{U} \to \mathcal{X}$. Si $\mathcal{U}$
est un espace algébrique, alors $f : \mathcal{U} \to \mathcal{X}$
est représentable par des espaces algébriques, voir
Champs algébriques, Lemme \ref{algebraic-lemma-representable-diagonal}.
La proposition s'applique donc en particulier à un recouvrement lisse du
champ algébrique $\mathcal{X}$ par un schéma.

\begin{proposition}
\label{stacks-sheaves-proposition-smooth-covering-compute-direct-image}
Soient $f : \mathcal{U} \to \mathcal{X}$ et $g : \mathcal{X} \to \mathcal{Y}$
des $1$-morphismes composables de champs algébriques.
\begin{enumerate}
\item Supposons $f$ représentable par des espaces algébriques, surjectif et
lisse.
\begin{enumerate}
\item Si $\mathcal{F}$ appartient à $\textit{Ab}(\mathcal{X}_\etale)$,
il existe une suite spectrale
$$
E_1^{p, q} = R^q(g \circ f_p)_*f_p^{-1}\mathcal{F}
\Rightarrow
R^{p + q}g_*\mathcal{F}
$$
dans $\textit{Ab}(\mathcal{Y}_\etale)$, les images directes supérieures étant
calculées pour la topologie \'etale.
\item Si $\mathcal{F}$ appartient à
$\textit{Mod}(\mathcal{X}_\etale, \mathcal{O}_\mathcal{X})$, alors
il existe une suite spectrale
$$
E_1^{p, q} = R^q(g \circ f_p)_*f_p^{-1}\mathcal{F}
\Rightarrow
R^{p + q}g_*\mathcal{F}
$$
dans $\textit{Mod}(\mathcal{Y}_\etale, \mathcal{O}_\mathcal{Y})$.
\end{enumerate}
\item Supposons $f$ représentable par des espaces algébriques, surjectif,
plat et localement de présentation finie.
\begin{enumerate}
\item Si $\mathcal{F}$ appartient à $\textit{Ab}(\mathcal{X})$, il existe
une suite spectrale
$$
E_1^{p, q} = R^q(g \circ f_p)_*f_p^{-1}\mathcal{F}
\Rightarrow
R^{p + q}g_*\mathcal{F}
$$
dans $\textit{Ab}(\mathcal{Y})$, les images directes supérieures étant
calculées pour la topologie fppf.
\item Si $\mathcal{F}$ appartient à $\textit{Mod}(\mathcal{O}_\mathcal{X})$, alors
il existe une suite spectrale
$$
E_1^{p, q} = R^q(g \circ f_p)_*f_p^{-1}\mathcal{F}
\Rightarrow
R^{p + q}g_*\mathcal{F}
$$
dans $\textit{Mod}(\mathcal{O}_\mathcal{Y})$.
\end{enumerate}
\end{enumerate}
\end{proposition}

\begin{proof}
Pour le voir, vérifions les hypothèses (1) -- (4) du
Lemme \ref{stacks-sheaves-lemma-cech-to-cohomology-relative} et du
Lemme \ref{stacks-sheaves-lemma-cech-to-cohomology-relative-modules}.
Le $1$-morphisme $f$ est fidèle d'après
Champs algébriques, Lemme
\ref{algebraic-lemma-characterize-representable-by-algebraic-spaces}.
Cela prouve (4).
L'hypothèse (3) résulte du fait que $\mathcal{U}$ est un champ
algébrique, voir
Lemme \ref{stacks-sheaves-lemma-fibre-products}.
Pour (2), appliquons le
Lemme \ref{stacks-sheaves-lemma-surjective-flat-locally-finite-presentation}.
La condition (1) est satisfaite par hypothèse dans les quatre cas.
\end{proof}

\noindent
Voici une description des images directes supérieures pour un
morphisme de champs algébriques.

\begin{lemma}
\label{stacks-sheaves-lemma-pushforward-restriction}
Soit $S$ un schéma. Soit $f : \mathcal{X} \to \mathcal{Y}$ un
$1$-morphisme de champs algébriques\footnote{Ce résultat devrait valoir
pour tout $1$-morphisme de catégories fibrées en groupoïdes sur
$(\Sch/S)_{fppf}$.} sur $S$.
Soit $\tau \in \{Zariski,\linebreak[0] \etale,\linebreak[0]
smooth,\linebreak[0] syntomic,\linebreak[0] fppf\}$.
Soit $\mathcal{F}$
un objet de $\textit{Ab}(\mathcal{X}_\tau)$ ou de
$\textit{Mod}(\mathcal{X}_\tau, \mathcal{O}_\mathcal{X})$.
Alors le faisceau $R^if_*\mathcal{F}$ est le faisceau associé au
préfaisceau
$$
y \longmapsto
H^i_\tau\Big((\Sch/V)_{fppf} \times_{y, \mathcal{Y}} \mathcal{X},
\ \text{pr}^{-1}\mathcal{F}\Big)
$$
Ici, $y$ est un objet de $\mathcal{Y}$ au-dessus du schéma $V$.
\end{lemma}

\begin{proof}
Choisissons une résolution injective $\mathcal{F}[0] \to \mathcal{I}^\bullet$.
D'après la formule d'image directe (\ref{stacks-sheaves-equation-pushforward}),
$R^if_*\mathcal{F}$ est le faisceau associé au préfaisceau qui associe
à $y$ la cohomologie du complexe
$$
\begin{matrix}
\Gamma\Big((\Sch/V)_{fppf} \times_{y, \mathcal{Y}} \mathcal{X},
\ \text{pr}^{-1}\mathcal{I}^{i - 1}\Big) \\
\downarrow \\
\Gamma\Big((\Sch/V)_{fppf} \times_{y, \mathcal{Y}} \mathcal{X},
\ \text{pr}^{-1}\mathcal{I}^i\Big) \\
\downarrow \\
\Gamma\Big((\Sch/V)_{fppf} \times_{y, \mathcal{Y}} \mathcal{X},
\ \text{pr}^{-1}\mathcal{I}^{i + 1}\Big)
\end{matrix}
$$
Comme $\text{pr}^{-1}$ est exact, il suffit de montrer que
$\text{pr}^{-1}$ préserve les injectifs. Cela résulte des
Lemmes \ref{stacks-sheaves-lemma-pullback-injective} et
\ref{stacks-sheaves-lemma-pullback-injective-modules}
ainsi que du fait que $\text{pr}$ est un morphisme représentable de
champs algébriques (de sorte que $\text{pr}$ est fidèle d'après
Champs algébriques, Lemme
\ref{algebraic-lemma-characterize-representable-by-algebraic-spaces}
et que
$(\Sch/V)_{fppf} \times_{y, \mathcal{Y}} \mathcal{X}$
admet des égalisateurs d'après le
Lemme \ref{stacks-sheaves-lemma-fibre-products}).
\end{proof}

\noindent
Voici un résultat immédiat de changement de base.

\begin{lemma}
\label{stacks-sheaves-lemma-base-change-higher-direct-images}
Soit $S$ un schéma. Soit
$\tau \in \{Zariski,\linebreak[0] \etale,\linebreak[0]
smooth,\linebreak[0] syntomic,\linebreak[0] fppf\}$. Soit
$$
\xymatrix{
\mathcal{Y}' \times_\mathcal{Y} \mathcal{X} \ar[r]_{g'} \ar[d]_{f'} &
\mathcal{X} \ar[d]^f \\
\mathcal{Y}' \ar[r]^g & \mathcal{Y}
}
$$
un diagramme $2$-cartésien de champs algébriques sur $S$. Alors le morphisme de
changement de base est un isomorphisme
$$
g^{-1}Rf_*\mathcal{F} \longrightarrow Rf'_*(g')^{-1}\mathcal{F}
$$
fonctoriel en $\mathcal{F}$ dans $\textit{Ab}(\mathcal{X}_\tau)$
ou en $\mathcal{F}$ dans $\textit{Mod}(\mathcal{X}_\tau, \mathcal{O}_\mathcal{X})$.
\end{lemma}

\begin{proof}
L'isomorphisme $g^{-1}f_*\mathcal{F} = f'_*(g')^{-1}\mathcal{F}$ est le
Lemme \ref{stacks-sheaves-lemma-base-change} (et vaut pour des préfaisceaux arbitraires).
Pour les images directes totales, il existe un morphisme de changement de base parce que les
morphismes $g$ et $g'$ sont plats, voir
Cohomologie sur les sites, section \ref{sites-cohomology-section-base-change-map}.
Pour voir que ce morphisme est un quasi-isomorphisme, nous pouvons employer le fait que, pour
un objet $y'$ de $\mathcal{Y}'$ au-dessus d'un schéma $V$, il existe une équivalence
$$
(\Sch/V)_{fppf} \times_{g(y'), \mathcal{Y}} \mathcal{X}
=
(\Sch/V)_{fppf} \times_{y', \mathcal{Y}'}
(\mathcal{Y}' \times_\mathcal{Y} \mathcal{X})
$$
Nous concluons que le morphisme induit
$g^{-1}R^if_*\mathcal{F} \to R^if'_*(g')^{-1}\mathcal{F}$
est un isomorphisme d'après le
Lemme \ref{stacks-sheaves-lemma-pushforward-restriction}.
\end{proof}












\section{Comparaison}
\label{stacks-sheaves-section-compare}

\noindent
Dans cette section, nous réunissons quelques résultats comparant la cohomologie définie
au moyen des champs et celle définie au moyen des espaces algébriques.

\begin{lemma}
\label{stacks-sheaves-lemma-compare-injectives}
Soit $S$ un schéma. Soit $\mathcal{X}$ un champ algébrique sur $S$
représentable par l'espace algébrique $F$.
\begin{enumerate}
\item Si $\mathcal{I}$ est injectif dans $\textit{Ab}(\mathcal{X}_\etale)$, alors
$\mathcal{I}|_{F_\etale}$ est injectif dans $\textit{Ab}(F_\etale)$,
\item Si $\mathcal{I}^\bullet$ est un complexe K-injectif dans
$\textit{Ab}(\mathcal{X}_\etale)$, alors $\mathcal{I}^\bullet|_{F_\etale}$
est un complexe K-injectif dans $\textit{Ab}(F_\etale)$.
\end{enumerate}
Le même énoncé ne vaut pas pour les modules.
\end{lemma}

\begin{proof}
Cela résulte formellement du fait que le foncteur de restriction
$\pi_{F, *} = i_F^{-1}$ (voir Lemme \ref{stacks-sheaves-lemma-compare})
est adjoint à droite du foncteur exact $\pi_F^{-1}$, voir
Homologie, Lemme \ref{homology-lemma-adjoint-preserve-injectives} et
Catégories dérivées, Lemme \ref{derived-lemma-adjoint-preserve-K-injectives}.
Pour voir que le lemme ne vaut pas pour les modules, nous renvoyons le
lecteur à Cohomologie \'etale, Lemme
\ref{etale-cohomology-lemma-compare-injectives}.
\end{proof}

\begin{lemma}
\label{stacks-sheaves-lemma-compare-morphism-cohomology}
Soit $S$ un schéma. Soit $f : \mathcal{X} \to \mathcal{Y}$ un morphisme
de champs algébriques sur $S$. Supposons $\mathcal{X}$ et $\mathcal{Y}$
représentables par les espaces algébriques $F$ et $G$. Notons $f : F \to G$ le
morphisme induit d'espaces algébriques.
\begin{enumerate}
\item Pour tout $\mathcal{F} \in \textit{Ab}(\mathcal{X}_\etale)$,
nous avons
$$
(Rf_*\mathcal{F})|_{G_\etale} =
Rf_{small, *}(\mathcal{F}|_{F_\etale})
$$
dans $D(G_\etale)$.
\item Pour tout objet $\mathcal{F}$ de
$\textit{Mod}(\mathcal{X}_\etale, \mathcal{O}_\mathcal{X})$
nous avons
$$
(Rf_*\mathcal{F})|_{G_\etale} =
Rf_{small, *}(\mathcal{F}|_{F_\etale})
$$
dans $D(\mathcal{O}_G)$.
\end{enumerate}
\end{lemma}

\begin{proof}
L'assertion (1) résulte immédiatement du
Lemme \ref{stacks-sheaves-lemma-compare-injectives}
et de (\ref{stacks-sheaves-equation-compare-big-small}),
après choix d'une résolution injective de $\mathcal{F}$.

\medskip\noindent
L'assertion (2) se démontre comme suit. Dans le Lemme \ref{stacks-sheaves-lemma-compare-morphism},
nous avons vu que $\pi_G \circ f = f_{small} \circ \pi_F$ en tant que morphismes
de sites annelés. Nous obtenons donc
$R\pi_{G, *} \circ Rf_* = Rf_{small, *} \circ R\pi_{F, *}$
d'après Cohomologie sur les sites, Lemme
\ref{sites-cohomology-lemma-derived-pushforward-composition}.
Comme les foncteurs de restriction $\pi_{F, *}$ et $\pi_{G, *}$
sont exacts, nous concluons.
\end{proof}

\begin{lemma}
\label{stacks-sheaves-lemma-compare-representable-morphism-cohomology}
Soit $S$ un schéma. Considérons un carré qui est un $2$-produit fibré
$$
\xymatrix{
\mathcal{X}' \ar[r]_{g'} \ar[d]_{f'} & \mathcal{X} \ar[d]^f \\
\mathcal{Y}' \ar[r]^g & \mathcal{Y}
}
$$
de champs algébriques sur $S$. Supposons $f$ représentable par des espaces
algébriques et $\mathcal{Y}'$ représentable par un espace algébrique $G'$.
Alors $\mathcal{X}'$ est représentable par un espace algébrique $F'$ et,
en notant $f' : F' \to G'$ le morphisme induit d'espaces algébriques,
nous avons
$$
g^{-1}(Rf_*\mathcal{F})|_{G'_\etale} =
Rf'_{small, *}((g')^{-1}\mathcal{F}|_{F'_\etale})
$$
pour tout $\mathcal{F}$ dans $\textit{Ab}(\mathcal{X}_\etale)$
ou dans
$\textit{Mod}(\mathcal{X}_\etale, \mathcal{O}_\mathcal{X})$.
\end{lemma}

\begin{proof}
Cela résulte formellement de la combinaison des
Lemmes \ref{stacks-sheaves-lemma-base-change-higher-direct-images} et
\ref{stacks-sheaves-lemma-compare-morphism-cohomology}.
\end{proof}











\section{Changement de topologie}
\label{stacks-sheaves-section-change-topology}

\noindent
Voici un lemme technique affirmant que la cohomologie fppf
d'un faisceau localement quasi-cohérent est égale à sa cohomologie
\'etale, pourvu que les morphismes de comparaison soient des isomorphismes
pour les morphismes de $\mathcal{X}$ au-dessus de morphismes plats.

\begin{lemma}
\label{stacks-sheaves-lemma-lqc-flat-base-change-fppf-sheaf}
Soit $S$ un schéma. Soit $\mathcal{X}$ un champ algébrique sur $S$.
Soit $\mathcal{F}$ un préfaisceau de $\mathcal{O}_\mathcal{X}$-modules.
Supposons que
\begin{enumerate}
\item[(a)] $\mathcal{F}$ soit localement quasi-cohérent, et
\item[(b)] pour tout morphisme $\varphi : x \to y$ de $\mathcal{X}$ au-dessus
d'un morphisme de schémas $f : U \to V$ plat et
localement de présentation finie, le morphisme de comparaison
$c_\varphi : f_{small}^*\mathcal{F}|_{V_\etale} \to
\mathcal{F}|_{U_\etale}$ de
(\ref{stacks-sheaves-equation-comparison-modules}) soit un isomorphisme.
\end{enumerate}
Alors $\mathcal{F}$ est un faisceau pour la topologie fppf.
\end{lemma}

\begin{proof}
Soit $\{x_i \to x\}$ un recouvrement fppf de $\mathcal{X}$ au-dessus du
recouvrement fppf $\{f_i : U_i \to U\}$ de schémas sur $S$.
Par hypothèse, la restriction $\mathcal{G} = \mathcal{F}|_{U_\etale}$
est quasi-cohérente et les morphismes de comparaison
$f_{i, small}^*\mathcal{G} \to \mathcal{F}|_{U_{i, \etale}}$
sont des isomorphismes. La condition de faisceau pour $\mathcal{F}$
et le recouvrement $\{x_i \to x\}$ équivaut donc à celle de
$\mathcal{G}^a$ sur $(\Sch/U)_{fppf}$ pour le recouvrement $\{U_i \to U\}$,
laquelle est satisfaite d'après
Descente, Lemme \ref{descent-lemma-sheaf-condition-holds}.
\end{proof}

\begin{lemma}
\label{stacks-sheaves-lemma-compare-fppf-etale}
Soit $S$ un schéma. Soit $\mathcal{X}$ un champ algébrique sur $S$.
Soit $\mathcal{F}$ un $\mathcal{O}_\mathcal{X}$-module préfaisceau tel que
\begin{enumerate}
\item[(a)] $\mathcal{F}$ soit localement quasi-cohérent, et
\item[(b)] pour tout morphisme $\varphi : x \to y$ de $\mathcal{X}$ au-dessus
d'un morphisme de schémas $f : U \to V$ plat et
localement de présentation finie, le morphisme de comparaison
$c_\varphi : f_{small}^*\mathcal{F}|_{V_\etale} \to
\mathcal{F}|_{U_\etale}$ de
(\ref{stacks-sheaves-equation-comparison-modules}) soit un isomorphisme.
\end{enumerate}
Alors $\mathcal{F}$ est un $\mathcal{O}_\mathcal{X}$-module et
les assertions suivantes valent :
\begin{enumerate}
\item Si $\epsilon : \mathcal{X}_{fppf} \to \mathcal{X}_\etale$
est le morphisme de comparaison, alors
$R\epsilon_*\mathcal{F} = \epsilon_*\mathcal{F}$.
\item Les groupes de cohomologie $H^p_{fppf}(\mathcal{X}, \mathcal{F})$ sont égaux
aux groupes de cohomologie calculés pour la topologie \'etale sur $\mathcal{X}$.
Il en va de même des groupes $H^p_{fppf}(x, \mathcal{F})$ et des
versions dérivées $R\Gamma(\mathcal{X}, \mathcal{F})$ et
$R\Gamma(x, \mathcal{F})$.
\item Si $f : \mathcal{X} \to \mathcal{Y}$ est un $1$-morphisme de
catégories fibrées en groupoïdes sur $(\Sch/S)_{fppf}$, alors
$R^if_*\mathcal{F}$ est égal au faisceau fppf associé à l'image
directe supérieure calculée en cohomologie \'etale.
Il en va de même pour l'image réciproque dérivée.
\end{enumerate}
\end{lemma}

\begin{proof}
L'assertion que $\mathcal{F}$ est un $\mathcal{O}_\mathcal{X}$-module
résulte du
Lemme \ref{stacks-sheaves-lemma-lqc-flat-base-change-fppf-sheaf}.
Remarquons que $\epsilon$ est un morphisme de sites donné par le foncteur
identité de $\mathcal{X}$. Le faisceau $R^p\epsilon_*\mathcal{F}$ est donc
le faisceau associé au préfaisceau
$x \mapsto H^p_{fppf}(x, \mathcal{F})$, voir
Cohomologie sur les sites, Lemme \ref{sites-cohomology-lemma-higher-direct-images}.
Pour prouver (1), il suffit de montrer que
$H^p_{fppf}(x, \mathcal{F}) = 0$ pour $p > 0$
chaque fois que $x$ est au-dessus d'un schéma affine $U$. D'après le
Lemme \ref{stacks-sheaves-lemma-cohomology-restriction}
nous avons
$H^p_{fppf}(x, \mathcal{F}) = H^p((\Sch/U)_{fppf}, x^{-1}\mathcal{F})$.
En combinant
Descente, Lemme \ref{descent-lemma-quasi-coherent-and-flat-base-change}
avec Cohomologie des schémas, Lemme
\ref{coherent-lemma-quasi-coherent-affine-cohomology-zero}
nous voyons que ces groupes de cohomologie sont nuls.

\medskip\noindent
Nous avons vu ci-dessus que $\epsilon_*\mathcal{F}$ et $\mathcal{F}$ sont les
faisceaux sur $\mathcal{X}_\etale$ et $\mathcal{X}_{fppf}$
correspondant au même préfaisceau sur $\mathcal{X}$ (et cela vaut plus
généralement pour tout faisceau fppf sur $\mathcal{X}$).
Nous identifions souvent abusivement $\mathcal{F}$ et $\epsilon_*\mathcal{F}$ ;
c'est en ce sens qu'il faut comprendre les assertions (2) et (3) du
lemme. Ainsi (2) résulte formellement de (1) et de la suite spectrale de
Leray, voir
Cohomologie sur les sites, Lemme \ref{sites-cohomology-lemma-apply-Leray}.

\medskip\noindent
Enfin, prouvons (3). Le faisceau $R^if_*\mathcal{F}$
(resp.\ $Rf_{\etale, *}\mathcal{F}$)
est le faisceau associé au préfaisceau
$$
y \longmapsto
H^i_\tau\Big((\Sch/V)_{fppf} \times_{y, \mathcal{Y}} \mathcal{X},
\ \text{pr}^{-1}\mathcal{F}\Big)
$$
où $\tau$ vaut $fppf$ (resp.\ $\etale$), voir le
Lemme \ref{stacks-sheaves-lemma-pushforward-restriction}.
Remarquons que $\text{pr}^{-1}\mathcal{F}$ satisfait aussi les propriétés (a) et (b)
(d'après les Lemmes \ref{stacks-sheaves-lemma-pullback-lqc} et \ref{stacks-sheaves-lemma-comparison}) ;
ces deux préfaisceaux sont donc égaux par (2).
Cela implique immédiatement (3).
\end{proof}

\noindent
Nous emploierons le lemme suivant pour comparer la cohomologie \'etale des faisceaux
sur les champs algébriques à la cohomologie du topos lisse-\'etale.

\begin{lemma}
\label{stacks-sheaves-lemma-cohomology-on-subcategory}
Soit $S$ un schéma. Soit $\mathcal{X}$ un champ algébrique sur $S$.
Soit $\tau = \etale$ (resp.\ $\tau = fppf$). Soit
$\mathcal{X}' \subset \mathcal{X}$ une sous-catégorie pleine possédant les
propriétés suivantes :
\begin{enumerate}
\item si $x \to x'$ est un morphisme de $\mathcal{X}$ au-dessus d'un morphisme
lisse (resp.\ plat et localement de présentation finie) de
schémas et $x' \in \Ob(\mathcal{X}')$, alors $x \in \Ob(\mathcal{X}')$, et
\item il existe un objet $x \in \Ob(\mathcal{X}')$ au-dessus
d'un schéma $U$ tel que le $1$-morphisme associé
$x : (\Sch/U)_{fppf} \to \mathcal{X}$ soit lisse et surjectif.
\end{enumerate}
On obtient un site $\mathcal{X}'_\tau$ en déclarant qu'un recouvrement de $\mathcal{X}'$
est toute famille de morphismes $\{x_i \to x\}$ dans $\mathcal{X}'$ qui est un
recouvrement dans $\mathcal{X}_\tau$. Alors le foncteur d'inclusion
$\mathcal{X}' \to \mathcal{X}_\tau$ est pleinement fidèle, cocontinu et
continu, et définit donc un morphisme de topos
$$
g : \Sh(\mathcal{X}'_\tau) \longrightarrow \Sh(\mathcal{X}_\tau)
$$
et $H^p(\mathcal{X}'_\tau, g^{-1}\mathcal{F}) =
H^p(\mathcal{X}_\tau, \mathcal{F})$ pour tout $p \geq 0$ et tout
$\mathcal{F} \in \textit{Ab}(\mathcal{X}_\tau)$.
\end{lemma}

\begin{proof}
L'hypothèse (1) implique que, si $\{x_i \to x\}$ est un recouvrement
de $\mathcal{X}_\tau$ et $x \in \Ob(\mathcal{X}')$, alors
$x_i \in \Ob(\mathcal{X}')$. Ainsi $\mathcal{X}' \to \mathcal{X}$
est continu et cocontinu, puisque les recouvrements des objets de
$\mathcal{X}'_\tau$ coïncident avec leurs recouvrements comme objets de
$\mathcal{X}_\tau$. Nous obtenons le morphisme $g$, et le foncteur
$g^{-1}$ s'identifie au foncteur de restriction, voir
Sites, Lemme \ref{sites-lemma-when-shriek}.

\medskip\noindent
En particulier, si $\{x_i \to x\}$ est un recouvrement dans $\mathcal{X}'_\tau$,
alors, pour tout faisceau abélien $\mathcal{F}$ sur $\mathcal{X}$,
$$
\check H^p(\{x_i \to x\}, g^{-1}\mathcal{F}) =
\check H^p(\{x_i \to x\}, \mathcal{F})
$$
Ainsi, si $\mathcal{I}$ est un faisceau abélien injectif sur $\mathcal{X}_\tau$,
les groupes de cohomologie de {\v C}ech supérieurs sont nuls
(Cohomologie sur les sites,
Lemme \ref{sites-cohomology-lemma-injective-trivial-cech}).
Ainsi $H^p(x, g^{-1}\mathcal{I}) = 0$ pour tout objet $x$
de $\mathcal{X}'$
(Cohomologie sur les sites,
Lemme \ref{sites-cohomology-lemma-cech-vanish-collection}).
Autrement dit, les faisceaux abéliens injectifs sur $\mathcal{X}_\tau$
sont acycliques à droite pour le foncteur $H^0(x, g^{-1}-)$.
Il s'ensuit que $H^p(x, g^{-1}\mathcal{F}) = H^p(x, \mathcal{F})$
pour tout $\mathcal{F} \in \textit{Ab}(\mathcal{X})$ et tout
$x \in \Ob(\mathcal{X}')$.

\medskip\noindent
Choisissons un objet $x \in \mathcal{X}'$ au-dessus d'un schéma $U$
comme dans l'hypothèse (2). En particulier, $\mathcal{X}/x \to \mathcal{X}$
est un morphisme de champs algébriques représentable par des espaces algébriques,
surjectif et lisse. (Remarquons que $\mathcal{X}/x$ est équivalent à
$(\Sch/U)_{fppf}$, voir le Lemme \ref{stacks-sheaves-lemma-localizing}.)
Le morphisme de faisceaux
$$
h_x \longrightarrow *
$$
dans $\Sh(\mathcal{X}_\tau)$ est surjectif. En effet, pour tout objet $x'$
de $\mathcal{X}$, il existe un $\tau$-recouvrement $\{x'_i \to x'\}$
tel qu'il existe des morphismes $x'_i \to x$, voir
Lemme \ref{stacks-sheaves-lemma-surjective-flat-locally-finite-presentation}.
Comme $g$ est exact, le morphisme de faisceaux
$$
g^{-1}h_x \longrightarrow * = g^{-1}*
$$
dans $\Sh(\mathcal{X}'_\tau)$ est aussi surjectif. Soit $h_{x, n}$ le
produit de $(n + 1)$ facteurs $h_x \times \ldots \times h_x$.
Nous avons alors les suites spectrales
\begin{equation}
\label{stacks-sheaves-equation-spectral-sequence-one}
E_1^{p, q} = H^q(h_{x, p}, \mathcal{F}) \Rightarrow
H^{p + q}(\mathcal{X}_\tau, \mathcal{F})
\end{equation}
et
\begin{equation}
\label{stacks-sheaves-equation-spectral-sequence-two}
E_1^{p, q} = H^q(g^{-1}h_{x, p}, g^{-1}\mathcal{F}) \Rightarrow
H^{p + q}(\mathcal{X}'_\tau, g^{-1}\mathcal{F})
\end{equation}
voir Cohomologie sur les sites,
Lemme \ref{sites-cohomology-lemma-cech-to-cohomology-sheaf-sets}.

\medskip\noindent
Cas I : $\mathcal{X}$ admet un objet final $x$ qui appartient aussi à
$\mathcal{X}'$. Ce cas résulte immédiatement de la discussion
du deuxième paragraphe ci-dessus.

\medskip\noindent
Cas II : $\mathcal{X}$ est représentable par un espace algébrique $F$.
Dans ce cas, les faisceaux $h_{x, n}$ sont représentables par un
objet $x_n$ de $\mathcal{X}$. (En effet, si $\mathcal{S}_F = \mathcal{X}$
et si $x : U \to F$ est l'objet donné, alors $h_{x, n}$ est représenté
par l'objet $U \times_F \ldots \times_F U \to F$ de $\mathcal{S}_F$.)
Il s'ensuit que $H^q(h_{x, p}, \mathcal{F}) = H^q(x_p, \mathcal{F})$.
Les morphismes $x_n \to x$ sont au-dessus de morphismes lisses de schémas, donc
$x_n \in \mathcal{X}'$ pour tout $n$. Ainsi
$H^q(g^{-1}h_{x, p}, g^{-1}\mathcal{F}) = H^q(x_p, g^{-1}\mathcal{F})$.
Ainsi, dans les deux suites spectrales
(\ref{stacks-sheaves-equation-spectral-sequence-one}) et
(\ref{stacks-sheaves-equation-spectral-sequence-two}) ci-dessus, les termes $E_1^{p, q}$ coïncident
d'après la discussion du deuxième paragraphe. Le lemme en résulte aussi
dans le Cas II.

\medskip\noindent
Cas III : $\mathcal{X}$ est un champ algébrique. Nous affirmons que, dans ce cas,
les groupes de cohomologie $H^q(h_{x, p}, \mathcal{F})$ et
$H^q(g^{-1}h_{x, n}, g^{-1}\mathcal{F})$ coïncident d'après le Cas II.
Une fois ceci prouvé, le résultat suivra comme précédemment.

\medskip\noindent
Considérons en effet la catégorie $\mathcal{X}/h_{x, n}$, voir
Sites, Lemme \ref{sites-lemma-localize-topos-site}.
Comme $h_{x, n}$ est le produit de $(n + 1)$ facteurs $h_x$, un
objet de cette catégorie est un $(n + 2)$-uplet
$(y, s_0, \ldots, s_n)$, où $y$ est un objet de $\mathcal{X}$ et chaque
$s_i : y \to x$ est un morphisme de $\mathcal{X}$.
C'est une catégorie sur $(\Sch/S)_{fppf}$. Il existe une équivalence
$$
\mathcal{X}/h_{x, n}
\longrightarrow
(\Sch/U)_{fppf} \times_\mathcal{X} \ldots \times_\mathcal{X} (\Sch/U)_{fppf}
=: \mathcal{U}_n
$$
sur $(\Sch/S)_{fppf}$. En effet, si $x : (\Sch/U)_{fppf} \to \mathcal{X}$ désigne aussi
le $1$-morphisme associé à $x$ et
$p : \mathcal{X} \to (\Sch/S)_{fppf}$ le foncteur structural,
alors nous pouvons voir $(y, s_0, \ldots, s_n)$ comme
$(y, f_0, \ldots, f_n, \alpha_0, \ldots, \alpha_n)$
où $y$ est un objet de $\mathcal{X}$, $f_i : p(y) \to p(x)$ un
morphisme de schémas et $\alpha_i : y \to x(f_i)$ un isomorphisme.
La catégorie des $2n+3$-uplets
$(y, f_0, \ldots, f_n, \alpha_0, \ldots, \alpha_n)$
est une réalisation du produit fibré à $(n + 1)$ facteurs $\mathcal{U}_n$
de champs algébriques affiché ci-dessus, comme discuté dans la
section \ref{stacks-sheaves-section-cech}.
D'après Cohomologie sur les sites, Lemme
\ref{sites-cohomology-lemma-cohomology-on-sheaf-sets}
nous avons
$$
H^p(\mathcal{U}_n, \mathcal{F}|_{\mathcal{U}_n}) =
H^p(\mathcal{X}/h_{x, n}, \mathcal{F}|_{\mathcal{X}/h_{x, n}}) =
H^p(h_{x, n}, \mathcal{F}).
$$
Enfin, examinons l'analogue « primé ». La catégorie
$\mathcal{X}'/h_{x, n}$ correspond, par l'équivalence ci-dessus,
à la sous-catégorie pleine $\mathcal{U}'_n \subset \mathcal{U}_n$
formée des uplets
$(y, f_0, \ldots, f_n, \alpha_0, \ldots, \alpha_n)$
avec $y \in \mathcal{X}'$. La propriété (1) de l'énoncé du
lemme vaut donc bien
pour l'inclusion $\mathcal{U}'_n \subset \mathcal{U}_n$.
Pour la propriété (2), choisissons un objet $\xi = (y, s_0, \ldots, s_n)$
au-dessus d'un schéma $W$ tel que $(\Sch/W)_{fppf} \to \mathcal{U}_n$
soit lisse et surjectif (c'est possible puisque $\mathcal{U}_n$ est
un champ algébrique). Alors
$(\Sch/W)_{fppf} \to \mathcal{U}_n \to (\Sch/U)_{fppf}$
est lisse comme composée de changements de base du morphisme
$x : (\Sch/U)_{fppf} \to \mathcal{X}$, voir
Champs algébriques, Lemmes
\ref{algebraic-lemma-base-change-representable-transformations-property} et
\ref{algebraic-lemma-composition-representable-transformations-property}.
Ainsi l'axiome (1) pour $\mathcal{X}$ implique que $y$ est un objet de
$\mathcal{X}'$, donc $\xi$ est un objet de $\mathcal{U}'_n$.
En employant encore
$$
H^p(\mathcal{U}'_n, \mathcal{F}|_{\mathcal{U}'_n}) =
H^p(\mathcal{X}'/h_{x, n}, \mathcal{F}|_{\mathcal{X}'/h_{x, n}}) =
H^p(g^{-1}h_{x, n}, g^{-1}\mathcal{F}).
$$
nous pouvons maintenant appliquer le Cas II à
$\mathcal{U}'_n \subset \mathcal{U}_n$
et conclure.
\end{proof}








\section{Restriction aux affines}
\label{stacks-sheaves-section-alternative}

\noindent
Dans cette section, pour une catégorie $\mathcal{X}$ fibrée en groupoïdes sur
$(\Sch/S)_{fppf}$, nous considérons la sous-catégorie pleine $\mathcal{X}_{affine}$
de $\mathcal{X}$ formée des objets $x$ au-dessus de schémas affines $U$.
Nous verrons que, pour toute topologie $\tau$ plus fine que la topologie de Zariski,
les catégories de faisceaux sur $\mathcal{X}$ et sur $\mathcal{X}_{affine, \tau}$
coïncident.

\begin{definition}
\label{stacks-sheaves-definition-alternative}
Soit $p : \mathcal{X} \to (\Sch/S)_{fppf}$ une catégorie fibrée en groupoïdes.
Le {\it site affine associé} est la sous-catégorie pleine
$\mathcal{X}_{affine}$ de $\mathcal{X}$ dont les objets sont les
$x \in \Ob(\mathcal{X})$ au-dessus d'un schéma $U$ tel que $U$
soit affine. La topologie sur $\mathcal{X}_{affine}$ sera la topologie chaotique,
c'est-à-dire que les faisceaux sur $\mathcal{X}_{affine}$ sont les préfaisceaux.
\end{definition}

\noindent
Ainsi le foncteur $p : \mathcal{X} \to (\Sch/S)_{fppf}$ se restreint en un foncteur
$$
p : \mathcal{X}_{affine} \longrightarrow (\textit{Aff}/S)_{fppf}
$$
où la notation du membre de droite est celle introduite dans
Topologies, Définition \ref{topologies-definition-big-small-fppf}.
Il est clair que $\mathcal{X}_{affine}$ est fibrée en groupoïdes sur
$(\textit{Aff}/S)_{fppf}$. Il s'ensuit que $\mathcal{X}_{affine}$
hérite de topologies de Zariski, \'etale, lisse, syntomique et fppf
à partir de $(\textit{Aff}/S)_{Zar}$, $(\textit{Aff}/S)_\etale$,
$(\textit{Aff}/S)_{smooth}$, $(\textit{Aff}/S)_{syntomic}$ et
$(\textit{Aff}/S)_{fppf}$ ; voir
Champs, Définition \ref{stacks-definition-topology-inherited}.

\begin{definition}
\label{stacks-sheaves-definition-alternative-inherited-topologies}
Soit $p : \mathcal{X} \to (\Sch/S)_{fppf}$ une catégorie fibrée en groupoïdes.
\begin{enumerate}
\item Le {\it site affine de Zariski associé} $\mathcal{X}_{affine, Zar}$
est la structure de site sur $\mathcal{X}_{affine}$ héritée de
$(\textit{Aff}/S)_{Zar}$.
\item Le {\it site affine \'etale associé} $\mathcal{X}_{affine, \etale}$
est la structure de site sur $\mathcal{X}_{affine}$ héritée de
$(\textit{Aff}/S)_\etale$.
\item Le {\it site affine lisse associé}
$\mathcal{X}_{affine, smooth}$
est la structure de site sur $\mathcal{X}_{affine}$ héritée de
$(\textit{Aff}/S)_{smooth}$.
\item Le {\it site affine syntomique associé} $\mathcal{X}_{affine, syntomic}$
est la structure de site sur $\mathcal{X}_{affine}$ héritée de
$(\textit{Aff}/S)_{syntomic}$.
\item Le {\it site affine fppf associé} $\mathcal{X}_{affine, fppf}$
est la structure de site sur $\mathcal{X}_{affine}$ héritée de
$(\textit{Aff}/S)_{fppf}$.
\end{enumerate}
\end{definition}

\noindent
Cette définition a un sens d'après la discussion ci-dessus. Pour chaque
$\tau \in \{Zariski,\allowbreak \etale,\allowbreak smooth,\allowbreak
syntomic,\allowbreak fppf\}$, une
famille de morphismes $\{x_i \to x\}_{i \in I}$ de but fixé
dans $\mathcal{X}_{affine}$ est un recouvrement de $\mathcal{X}_{affine, \tau}$
si et seulement si la famille $\{p(x_i) \to p(x)\}_{i \in I}$
de morphismes de schémas affines est un $\tau$-recouvrement standard au sens de
Topologies, Définitions
\ref{topologies-definition-standard-Zariski},
\ref{topologies-definition-standard-etale},
\ref{topologies-definition-standard-smooth},
\ref{topologies-definition-standard-syntomic} et
\ref{topologies-definition-standard-fppf}.

\begin{lemma}
\label{stacks-sheaves-lemma-alternative}
Soit $p : \mathcal{X} \to (\Sch/S)_{fppf}$ une catégorie fibrée en groupoïdes.
Soit $\tau \in \{Zariski, \etale, smooth, syntomic, fppf\}$. Le foncteur
$\mathcal{X}_{affine, \tau} \to \mathcal{X}_\tau$ est un foncteur
cocontinu spécial. Il induit donc une équivalence de topos
de $\Sh(\mathcal{X}_{affine, \tau})$ vers $\Sh(\mathcal{X}_\tau)$.
\end{lemma}

\begin{proof}
Omis. Indication : la preuve est exactement la même que celle de
Topologies, Lemmes
\ref{topologies-lemma-affine-big-site-Zariski},
\ref{topologies-lemma-affine-big-site-etale},
\ref{topologies-lemma-affine-big-site-smooth},
\ref{topologies-lemma-affine-big-site-syntomic} et
\ref{topologies-lemma-affine-big-site-fppf}.
\end{proof}

\noindent
Soit $p : \mathcal{X} \to (\Sch/S)_{fppf}$ une catégorie fibrée en groupoïdes.
Notons $\mathcal{O}$ la restriction de $\mathcal{O}_\mathcal{X}$
à $\mathcal{X}_{affine}$. Alors $\mathcal{O}$ est un faisceau pour les topologies
de Zariski, \'etale, lisse, syntomique et fppf sur
$\mathcal{X}_{affine}$. De plus, l'équivalence de topos du
Lemme \ref{stacks-sheaves-lemma-alternative} s'étend en une équivalence
\begin{equation}
\label{stacks-sheaves-equation-alternative-ringed}
(\Sh(\mathcal{X}_{affine, \tau}), \mathcal{O})
\longrightarrow
(\Sh(\mathcal{X}_\tau), \mathcal{O}_\mathcal{X})
\end{equation}
de topos annelés pour $\tau \in \{Zariski, \etale, smooth, syntomic, fppf\}$.









\section{Modules quasi-cohérents et affines}
\label{stacks-sheaves-section-alternative-quasi-coherent}

\noindent
Soit $p : \mathcal{X} \to (\Sch/S)_{fppf}$ une catégorie fibrée en groupoïdes.
Dans la section \ref{stacks-sheaves-section-alternative}, nous lui avons associé un site annelé
$(\mathcal{X}_{affine}, \mathcal{O})$.

\begin{lemma}
\label{stacks-sheaves-lemma-quasi-coherent-alternative}
Soit $p : \mathcal{X} \to (\Sch/S)_{fppf}$ une catégorie fibrée en groupoïdes.
Soit $\mathcal{F}$ un $\mathcal{O}$-module sur $\mathcal{X}_{affine}$.
Les conditions suivantes sont équivalentes :
\begin{enumerate}
\item pour tout morphisme $x \to x'$ de $\mathcal{X}_{affine}$, le morphisme
$\mathcal{F}(x') \otimes_{\mathcal{O}(x')} \mathcal{O}(x) \to \mathcal{F}(x)$
est un isomorphisme,
\item $\mathcal{F}$ est un module quasi-cohérent sur
$(\mathcal{X}_{affine}, \mathcal{O})$ au sens de
Modules sur les sites, Définition \ref{sites-modules-definition-site-local},
\item $\mathcal{F}$ est un faisceau pour la topologie de Zariski sur
$\mathcal{X}_{affine}$ et un module quasi-cohérent sur
$(\mathcal{X}_{affine, Zar}, \mathcal{O})$ au sens de
Modules sur les sites, Définition \ref{sites-modules-definition-site-local},
\item comme en (3) pour la topologie \'etale,
\item comme en (3) pour la topologie lisse,
\item comme en (3) pour la topologie syntomique,
\item comme en (3) pour la topologie fppf, et
\item $\mathcal{F}$ correspond à un module quasi-cohérent
sur $\mathcal{X}$ par l'équivalence (\ref{stacks-sheaves-equation-alternative-ringed}).
\end{enumerate}
\end{lemma}

\begin{proof}
Pour interpréter (2), rappelons que $\mathcal{X}_{affine}$ est le site
obtenu en munissant la catégorie $\mathcal{X}_{affine}$ de la topologie
chaotique (Définition \ref{stacks-sheaves-definition-alternative}) ; un faisceau
de $\mathcal{O}$-modules $\mathcal{F}$ est donc la même chose qu'un
préfaisceau de $\mathcal{O}$-modules. Les conditions (1) et (2)
sont équivalentes d'après Modules sur les sites, Lemme
\ref{sites-modules-lemma-chaotic-quasi-coherent}.
Remarquons que, pour $\tau \in \{Zariski, \etale, smooth, syntomic, fppf\}$,
le préfaisceau $\mathcal{F}$ est un $\tau$-faisceau si et seulement si,
pour tout $x \in \Ob(\mathcal{X}_{affine})$, sa restriction
à $\mathcal{X}_{affine}/x$ est un $\tau$-faisceau. Posons $U = p(x)$.
Comme dans la discussion de la section \ref{stacks-sheaves-section-restriction},
l'objet $x$ de $\mathcal{X}_{affine}$ induit une équivalence
$\mathcal{X}_{affine, \etale}/x \to (\textit{Aff}/U)_\etale$ de sites.
Ainsi, l'équivalence de (1) avec (3) -- (7)
résulte de Descente, Lemme \ref{descent-lemma-quasi-coherent-alternative},
appliqué à chacun de ces sites. L'équivalence de (8) et (7)
résulte immédiatement du fait qu'« être quasi-cohérent » est une
propriété intrinsèque des faisceaux de modules, voir
Modules sur les sites, section \ref{sites-modules-section-intrinsic}.
\end{proof}

\begin{lemma}
\label{stacks-sheaves-lemma-locally-quasi-coherent-alternative}
Soit $p : \mathcal{X} \to (\Sch/S)_{fppf}$ une catégorie fibrée en groupoïdes.
Soit $\mathcal{F}$ un $\mathcal{O}$-module sur $\mathcal{X}_{affine}$.
Les conditions suivantes sont équivalentes :
\begin{enumerate}
\item pour tout morphisme $x \to x'$ de $\mathcal{X}_{affine}$ tel que
$p(x) \to p(x')$ soit un morphisme \'etale (de schémas affines), le morphisme
$\mathcal{F}(x') \otimes_{\mathcal{O}(x')} \mathcal{O}(x) \to \mathcal{F}(x)$
est un isomorphisme,
\item $\mathcal{F}$ est un faisceau pour la topologie \'etale sur
$\mathcal{X}_{affine}$ et, pour tout objet $x$ de $\mathcal{X}_{affine}$,
la restriction $x^*\mathcal{F}|_{U_{affine, \etale}}$ est quasi-cohérente,
où $U = p(x)$,
\item $\mathcal{F}$ correspond à un module localement quasi-cohérent
sur $\mathcal{X}$ par l'équivalence (\ref{stacks-sheaves-equation-alternative-ringed})
pour la topologie \'etale.
\end{enumerate}
\end{lemma}

\begin{proof}
Pour interpréter la condition (2), rappelons que $U_{affine, \etale}$
est la sous-catégorie pleine de $U_\etale$ formée des objets affines, voir
Topologies, Définition \ref{topologies-definition-big-small-etale}.
Comme dans la discussion de la section \ref{stacks-sheaves-section-restriction},
l'objet $x$ de $\mathcal{X}_{affine}$ induit une équivalence
$\mathcal{X}_{affine, \etale}/x \to (\textit{Aff}/U)_\etale$ de sites.
Alors $x^*\mathcal{F}$ est le faisceau de modules sur $(\textit{Aff}/U)_\etale$
correspondant à la restriction
$\mathcal{F}|_{\mathcal{X}_{affine, \etale}/x}$.
Enfin, au moyen du foncteur d'inclusion continu et cocontinu
$U_{affine, \etale} \to (\textit{Aff}/U)_\etale$, nous pouvons encore
restreindre et obtenir $x^*\mathcal{F}|_{U_{affine, \etale}}$.

\medskip\noindent
L'équivalence de (1) et (2) résulte des remarques ci-dessus et de
Descente, Lemme \ref{descent-lemma-quasi-coherent-alternative-small}
appliqué à la restriction de $\mathcal{F}$ à $U_{affine, \etale}$
pour tout objet $x$ de $\mathcal{X}$ au-dessus d'un schéma affine $U$.
L'équivalence de (2) et (3) résulte immédiatement des définitions
et du fait que les modules quasi-cohérents sur $U_{affine, \etale}$
et $U_\etale$ se correspondent (encore d'après
Descente, Lemme \ref{descent-lemma-quasi-coherent-alternative-small}
par exemple).
\end{proof}













\section{Objets quasi-cohérents de la catégorie dérivée}
\label{stacks-sheaves-section-QC}

\noindent
Depuis des décennies, les géomètres algébristes considèrent des invariants de
foncteurs non représentables $X$ (à valeurs dans les ensembles ou les groupoïdes) sur $\Sch/S$.
Par exemple,
avant l'invention de la notion de champ, Mumford a défini \cite{mumford_picard}
le groupoïde de Picard $\Pic(X)$ du foncteur de modules $X$ des courbes elliptiques
comme la $2$-limite $Pic(U)$ sur la catégorie de tous les schémas $U$
munis d'un morphisme vers $X$ (c'est-à-dire d'une famille de courbes elliptiques).
De même, Beilinson-Drinfeld ont défini \cite{BVGD} la catégorie
$\QCoh(X)$ d'un ind-schéma $X = \colim X_i$ comme la $2$-limite
$\lim \QCoh(X_i)$.
Cette stratégie suffit à définir des invariants $1$-catégoriques
comme $\QCoh(-)$, mais non des invariants catégoriques dérivés
(tels que la catégorie dérivée quasi-cohérente), car les $2$-limites
de catégories triangulées se comportent mal.
Avec l'avènement des techniques de catégories supérieures
et de la géométrie algébrique dérivée, ce problème se
résout élégamment : on peut définir une
$\infty$-catégorie dérivée quasi-cohérente $\mathcal{D}_{qc}(X)$ du foncteur $X$
comme la limite $\lim \mathcal{D}_{qc}(U)$, où $U$
parcourt les affines dérivés sur X (voir \cite{lurie-thesis}).

\medskip\noindent
Le but de cette section est d'associer une catégorie triangulée
$\mathit{QC}(X)$ à un foncteur $X$ (à valeurs dans les ensembles ou les groupoïdes)
comme ci-dessus. En fait, la construction vaut pour toute catégorie
$p : \mathcal{X} \to (\Sch/S)_{fppf}$ fibrée en groupoïdes
(et pas seulement pour les catégories scindées). Dans les bons cas, la catégorie $\mathit{QC}(\mathcal{X})$
coïncide avec la catégorie homotopique de
$\mathcal{D}_{qc}(\mathcal{X})$,
mais expliquer cette comparaison dépasse le cadre de ce document.
Les principales propriétés de la construction sont les suivantes :
\begin{enumerate}
\item[(a)] $\mathit{QC}(\mathcal{X})$ est par construction une sous-catégorie pleine de
$D(\mathcal{X}_{affine}, \mathcal{O})$,
\item[(b)] $\mathit{QC}(\mathcal{X})$ coïncide avec $D_\QCoh(\mathcal{O}_X)$
lorsque $\mathcal{X}$ est représentable par l'espace algébrique $X$,
\item[(c)] $\mathit{QC}(\mathcal{X})$ coïncide avec
$D_\QCoh(\mathcal{O}_\mathcal{X})$ lorsque $\mathcal{X}$ est un champ algébrique,
\item[(d)] lorsque $X = \text{Spf}(A)$ est un espace algébrique formel affine
associé à un anneau noethérien $A$ muni de la topologie $I$-adique
pour un idéal $I$, la catégorie triangulée
$\mathit{QC}(X)$ coïncide avec la sous-catégorie pleine
$D_{comp}(A, I) \subset D(A)$ des objets dérivés complets.
\end{enumerate}
Ces résultats sont démontrés dans la
Proposition \ref{stacks-sheaves-proposition-QC-compare},
Catégories dérivées des champs, Proposition
\ref{stacks-perfect-proposition-QC-compare} et
Proposition \ref{stacks-sheaves-proposition-QC-derived-complete}.

\medskip\noindent
Pour motiver la définition précise de $\mathit{QC}(\mathcal{X})$,
nous renvoyons le lecteur à la caractérisation, dans le
Lemme \ref{stacks-sheaves-lemma-quasi-coherent-alternative},
des modules quasi-cohérents sur $\mathcal{X}$ comme préfaisceaux de
$\mathcal{O}$-modules sur $\mathcal{X}_{affine}$ qui
satisfont une propriété de changement de base.

\begin{definition}
\label{stacks-sheaves-definition-QC}
Soit $p : \mathcal{X} \to (\Sch/S)_{fppf}$ une catégorie fibrée en groupoïdes.
Soit $\mathcal{O}$ le faisceau d'anneaux sur $\mathcal{X}_{affine}$ introduit
dans la section \ref{stacks-sheaves-section-alternative}. Nous définissons la
{\it catégorie triangulée des objets quasi-cohérents de la catégorie dérivée}
par la formule
$$
\mathit{QC}(\mathcal{X}) = \mathit{QC}(\mathcal{X}_{affine}, \mathcal{O})
$$
où le membre de droite est défini dans
Cohomologie sur les sites, Définition \ref{sites-cohomology-definition-cartesian}.
\end{definition}

\noindent
Cela a un sens, car $\mathcal{X}_{affine}$ est une catégorie considérée
comme un site en la munissant de la topologie chaotique, et
$\mathcal{O}$ est un faisceau d'anneaux sur cette catégorie, exactement comme requis dans
Cohomologie sur les sites, Définition \ref{sites-cohomology-definition-cartesian}.

\medskip\noindent
La relation entre cette définition et la catégorie des modules quasi-cohérents
sur $\mathcal{X}$ n'est pas si claire en général ! Supposons par exemple
que $M$ soit un objet de $\mathit{QC}(\mathcal{X})$. Les faisceaux de cohomologie
$H^i(M)$ de $M$ sont alors des (pré)faisceaux de $\mathcal{O}$-modules sur
$\mathcal{X}_{affine}$, mais ils ne sont pas quasi-cohérents en général. Toutefois, le dernier
faisceau de cohomologie non nul est quasi-cohérent.

\begin{lemma}
\label{stacks-sheaves-lemma-QC-bounded-above}
Dans la situation de la Définition \ref{stacks-sheaves-definition-QC}, supposons que
$M$ soit un objet de $\mathit{QC}(\mathcal{X})$ et que $b \in \mathbf{Z}$
vérifie $H^i(M) = 0$ pour tout $i > b$.
Alors $H^b(M)$ est un module quasi-cohérent sur
$(\mathcal{X}_{affine}, \mathcal{O})$ ; voir le
Lemme \ref{stacks-sheaves-lemma-quasi-coherent-alternative}.
\end{lemma}

\begin{proof}
Cas particulier de
Cohomologie sur les sites, Lemme \ref{sites-cohomology-lemma-QC-bounded-above}.
\end{proof}

\begin{lemma}
\label{stacks-sheaves-lemma-QC-compare}
Soit $S$ un schéma. Soit $\mathcal{X} \to (\Sch/S)_{fppf}$ une catégorie
fibrée en groupoïdes. Le morphisme de comparaison
$\epsilon : \mathcal{X}_{affine, \etale} \to \mathcal{X}_{affine}$
satisfait les hypothèses et les conclusions de Cohomologie sur les sites, Lemme
\ref{sites-cohomology-lemma-cartesion-plus-topology}.
\end{lemma}

\begin{proof}
L'hypothèse (1) vaut par définition de $\mathcal{X}_{affine}$.
Pour la condition (2), nous employons le fait que, pour $x \in \Ob(\mathcal{X})$
au-dessus du schéma affine $U = p(x)$, il existe une équivalence
$\mathcal{X}_{affine, \etale}/x = (\textit{Aff}/U)_\etale$ compatible
aux faisceaux structuraux ; voir la
discussion de la section \ref{stacks-sheaves-section-restriction}.
Il suffit donc de montrer ceci : pour un schéma affine $U = \Spec(R)$
et un complexe de $R$-modules $M^\bullet$, la cohomologie totale
du complexe de modules sur $(\textit{Aff}/U)_\etale$
associé à $M^\bullet$ est quasi-isomorphe à $M^\bullet$.
Cela résulte de la combinaison de :
Catégories dérivées des schémas, Lemme
\ref{perfect-lemma-affine-compare-bounded} (cohomologie totale des complexes
de modules sur les affines pour la topologie de Zariski),
Catégories dérivées des espaces, Remarque
\ref{spaces-perfect-remark-match-total-direct-images}
(coïncidence de la cohomologie totale pour les petites
topologies de Zariski et \'etale des
complexes quasi-cohérents de modules), et
Cohomologie \'etale, Lemme \ref{etale-cohomology-lemma-compare-cohomology}
(pour voir que la cohomologie \'etale d'un complexe de modules
sur le grand site \'etale d'un schéma
se calcule après restriction au petit site \'etale).
\end{proof}

\noindent
Si nous appliquons la définition lorsque notre catégorie fibrée en groupoïdes
$\mathcal{X}$ est représentable par un espace algébrique $X$, nous
retrouvons $D_\QCoh(\mathcal{O}_X)$. Nous énoncerons et démontrerons plus loin le
résultat analogue pour les champs algébriques (insérer ici une référence future).

\begin{proposition}
\label{stacks-sheaves-proposition-QC-compare}
Soit $S$ un schéma. Soit $\mathcal{X} \to (\Sch/S)_{fppf}$ une catégorie
fibrée en groupoïdes. Supposons $\mathcal{X}$ représentable par un espace
algébrique $X$. Alors $\mathit{QC}(\mathcal{X})$ est canoniquement équivalente à
$D_\QCoh(\mathcal{O}_X)$.
\end{proposition}

\begin{proof}
Notons $X_{affine}$ la catégorie des schémas affines \'etales sur $X$,
munie de la topologie chaotique et de son faisceau structural $\mathcal{O}_X$, voir
Catégories dérivées des espaces, section \ref{spaces-perfect-section-QC}.
Le foncteur $u : X_\etale \to \mathcal{X}_\etale$ du Lemme \ref{stacks-sheaves-lemma-compare}
induit un foncteur $X_{affine} \to \mathcal{X}_{affine}$.
Celui-ci est compatible aux faisceaux structuraux et fournit un foncteur
$$
G :
\mathit{QC}(\mathcal{X}) = \mathit{QC}(\mathcal{X}_{affine}, \mathcal{O})
\longrightarrow
\mathit{QC}(X_{affine}, \mathcal{O}_X)
$$
Voir
Cohomologie sur les sites, Lemme \ref{sites-cohomology-lemma-cartesian-functorial}.
D'après Catégories dérivées des espaces, Lemme
\ref{spaces-perfect-lemma-DQCoh-alternative-small}
la catégorie triangulée $\mathit{QC}(X_{affine}, \mathcal{O}_X)$
est équivalente à $D_\QCoh(\mathcal{O}_X)$. Il suffit donc de
prouver que $G$ est une équivalence.

\medskip\noindent
Considérons les morphismes de comparaison plats
$\epsilon_\mathcal{X} : \mathcal{X}_{affine, \etale} \to \mathcal{X}_{affine}$
et $\epsilon_X : X_{affine, \etale} \to X_{affine}$
de sites annelés. Le Lemme \ref{stacks-sheaves-lemma-QC-compare} et
(la preuve de) Catégories dérivées des espaces, Lemme
\ref{spaces-perfect-lemma-DQCoh-alternative-small}
montrent que les foncteurs
$\epsilon_\mathcal{X}^*$ et $\epsilon_X^*$
identifient $\mathit{QC}(\mathcal{X}_{affine}, \mathcal{O})$
et $\mathit{QC}(X_{affine}, \mathcal{O}_X)$
à des sous-catégories
$Q_\mathcal{X} \subset D(\mathcal{X}_{affine, \etale}, \mathcal{O})$
et
$Q_X \subset D(X_{affine, \etale}, \mathcal{O}_X)$.
Sous ces identifications,
le foncteur $G$ du premier paragraphe est induit par le foncteur
$$
Li_X^* = R\pi_{X, *}:
D(\mathcal{X}_{affine, \etale}, \mathcal{O})
\longrightarrow
D(X_{affine, \etale}, \mathcal{O}_X)
$$
où $i_X$ et $\pi_X$ sont les morphismes du Lemme \ref{stacks-sheaves-lemma-compare},
les sites \'etales étant remplacés par les sites affines correspondants.
Le lecteur peut montrer que ce remplacement est licite soit
en redémontrant directement le lemme pour les sites affines, soit
en employant les équivalences de topos
$\Sh(\mathcal{X}_{affine, \etale}) = \Sh(\mathcal{X}_\etale)$ et
$\Sh(X_{affine, \etale}) = \Sh(X_\etale)$.
Le lemme nous dit aussi que $Li_X^*$ admet un adjoint à gauche
$$
L\pi_X^*:
D(X_{affine, \etale}, \mathcal{O}_X)
\longrightarrow
D(\mathcal{X}_{affine, \etale}, \mathcal{O})
$$
et, de plus, $Li_X^* \circ L\pi_X^* = \text{id}$
puisque $\pi_X \circ i_X$ est l'identité. Il suffit donc
de montrer que (a) $L\pi_X^*$ envoie $Q_X$ dans $Q_\mathcal{X}$
et (b) le noyau de $Li_X^*$ est $0$.
Voir Catégories dérivées, Lemme
\ref{derived-lemma-fully-faithful-adjoint-kernel-zero}.

\medskip\noindent
Preuve de (a). D'après Catégories dérivées des espaces, Lemme
\ref{spaces-perfect-lemma-DQCoh-alternative-small}
nous avons $Q_X = D_\QCoh(X_{affine, \etale}, \mathcal{O}_X)$.
Soit $K$ un objet de $Q_X$. Soit $x$ un objet de
$\mathcal{X}_{affine, \etale}$ au-dessus du schéma
affine $U = p(x)$. Notons $f : U \to X$ le morphisme correspondant
à $x$. Alors
$$
R\Gamma(x, L\pi_X^*K) = R\Gamma(U, Lf^*K)
$$
Cela résulte de la transitivité des images réciproques ; voir la discussion de la
section \ref{stacks-sheaves-section-restriction-algebraic-spaces}.
Supposons ensuite que $x \to x'$ soit un morphisme de $\mathcal{X}_{affine, \etale}$
au-dessus du morphisme $h : U \to U'$ de schémas affines.
Comme précédemment, notons $f : U \to X$ et $f' : U' \to X$ les morphismes
correspondant à $x$ et $x'$, de sorte que $f = f' \circ h$.
Alors
\begin{align*}
R\Gamma(x, L\pi_X^*K)
& =
R\Gamma(U, Lf^*K) \\
& =
R\Gamma(U, Lh^*L(f')^*K) \\
& =
R\Gamma(U', L(f')^*K) \otimes_{\mathcal{O}(U')}^\mathbf{L} \mathcal{O}(U) \\
& =
R\Gamma(x', L\pi_X^*K) \otimes_{\mathcal{O}(x')}^\mathbf{L} \mathcal{O}(x)
\end{align*}
et nous obtenons donc (a) par la note de l'énoncé de
Cohomologie sur les sites, Lemme
\ref{sites-cohomology-lemma-cartesion-plus-topology}. La troisième égalité est
Catégories dérivées des schémas, Lemme
\ref{perfect-lemma-quasi-coherence-pullback}.

\medskip\noindent
Preuve de (b). Soit $M$ un objet de $Q_\mathcal{X}$
tel que $Li_X^*M = 0$. Soit $x'$ un objet de $\mathcal{X}_{affine, \etale}$
au-dessus du schéma affine $U' = p(x')$, et supposons que le morphisme
correspondant $f' : U' \to X$ soit \'etale. Alors $f' : U' \to X$ est un objet de
$X_{affine, \etale}$, et la condition $Li_X^*M = 0$ implique
$M|_{U'_\etale} = 0$. En particulier, $R\Gamma(x', M) = 0$.
Cependant, pour tout objet $x$ du site
$\mathcal{X}_{affine, \etale}$, il existe un recouvrement $\{x_i \to x\}$
tel que, pour chaque $i$, il existe un morphisme $x_i \to x'_i$, où
$x'_i$ correspond à un objet de $X_{affine, \etale}$.
Puisque $M$ appartient à $Q_\mathcal{X}$, nous avons
$$
R\Gamma(x_i, M) =
R\Gamma(x_i', M) \otimes_{\mathcal{O}(x_i')}^\mathbf{L} \mathcal{O}(x_i) =
0
$$
et nous concluons que $M$ est nul, comme voulu.
\end{proof}

\noindent
Pour montrer que la construction produit une catégorie intéressante dans un autre
cas, énonçons et démontrons une caractérisation de
$\mathit{QC}(\text{Spf}(A))$ pour le spectre formel d'un
anneau adique noethérien $A$.

\begin{proposition}
\label{stacks-sheaves-proposition-QC-derived-complete}
Soit $S$ un schéma. Soit $X = \text{Spf}(A)$, où $A$ est une
$S$-algèbre topologique adique noethérienne d'idéal de définition $I$, voir
Compléments d'algèbre, Définition \ref{more-algebra-definition-topological-ring}
et Espaces formels, Définition
\ref{formal-spaces-definition-affine-formal-spectrum}.
Soit $p : \mathcal{X} \to (\Sch/S)_{fppf}$
la catégorie fibrée en ensembles associée au foncteur $X$, voir
Catégories, Exemple \ref{categories-example-presheaf}.
Alors $\mathit{QC}(\mathcal{X})$ est canoniquement équivalente à la
catégorie $D_{comp}(A, I)$ des objets de $D(A)$ qui sont
dérivés complets pour $I$.
\end{proposition}

\begin{proof}
Rappelons que $X = \colim \Spec(A/I^n)$ comme faisceau fppf.
Un objet de $\mathcal{X}_{affine}$ est la même chose
qu'un schéma affine $U = \Spec(R)$ muni d'un morphisme $f : U \to X$.
D'après Espaces formels, Lemme \ref{formal-spaces-lemma-factor-through-thickening},
il existe $n \geq 1$ tel que
$f$ se factorise par le monomorphisme $\Spec(A/I^n) \to X$.
Considérons la sous-catégorie pleine
$\mathcal{C} \subset \mathcal{X}_{affine}$ formée des
objets $\Spec(A/I^n) \to X$.
D'après les remarques précédentes et Faisceaux différentiels gradués, Lemme
\ref{sdga-lemma-cartesian-cofinal-subcategory}
la restriction à $\mathcal{C}$ est une équivalence exacte
$\mathit{QC}(\mathcal{X}) \to
\mathit{QC}(\mathcal{C}, \mathcal{O}|_\mathcal{C})$.
Pour simplifier, supposons que
$I^n \not = I^{n + 1}$ pour tout $n \geq 1$.
Alors $(\mathcal{C}, \mathcal{O}|_\mathcal{C})$ est isomorphe,
comme site annelé, au site annelé $(\mathbf{N}, (A/I^n))$, voir
Faisceaux différentiels gradués, section
\ref{sdga-section-qc-inverse-system}.
Nous concluons donc par Faisceaux différentiels gradués, Proposition
\ref{sdga-proposition-derived-complete}.
\end{proof}

\noindent
Le lemme suivant servira à comparer $\mathit{QC}(\mathcal{X})$
à $D_\QCoh(\mathcal{O}_\mathcal{X})$ lorsque $\mathcal{X}$ est un champ algébrique.

\begin{lemma}
\label{stacks-sheaves-lemma-QC-compare-fppf}
Soit $S$ un schéma. Soit $\mathcal{X} \to (\Sch/S)_{fppf}$ une catégorie
fibrée en groupoïdes. Le morphisme de comparaison
$\epsilon : \mathcal{X}_{affine, fppf} \to \mathcal{X}_{affine}$
satisfait les hypothèses et les conclusions de Cohomologie sur les sites, Lemme
\ref{sites-cohomology-lemma-cartesion-plus-topology}.
\end{lemma}

\begin{proof}
La preuve est exactement la même que celle du Lemme \ref{stacks-sheaves-lemma-QC-compare}.
L'hypothèse (1) vaut par définition de $\mathcal{X}_{affine}$.
Pour la condition (2), nous employons le fait que, pour $x \in \Ob(\mathcal{X})$
au-dessus du schéma affine $U = p(x)$, il existe une équivalence
$\mathcal{X}_{affine, \etale}/x = (\textit{Aff}/U)_\etale$ compatible
aux faisceaux structuraux ; voir la
discussion de la section \ref{stacks-sheaves-section-restriction}.
Il suffit donc de montrer ceci : pour un schéma affine $U = \Spec(R)$
et un complexe de $R$-modules $M^\bullet$, la cohomologie totale
du complexe de modules sur $(\textit{Aff}/U)_{fppf}$
associé à $M^\bullet$ est quasi-isomorphe à $M^\bullet$.
C'est Cohomologie \'etale, Lemme
\ref{etale-cohomology-lemma-cohomological-descent-complex-modules}.
\end{proof}














% OK, start here.
%

\chapter{Critères de représentabilité}


\phantomsection
\label{criteria-section-phantom}





\section{Introduction}
\label{criteria-section-introduction}

\noindent
Le but de ce chapitre est de trouver des critères garantissant qu'un
champ en groupoïdes sur la catégorie des schémas munie de la topologie fppf
est un champ algébrique. Historiquement, cela consistait souvent à démontrer que
certains foncteurs étaient représentables ; voir les exposés de Grothendieck
\cite{Gr-I},
\cite{Gr-II},
\cite{Gr-III},
\cite{Gr-IV},
\cite{Gr-V}, et
\cite{Gr-VI}.
Cela explique le titre de ce chapitre. Les travaux d'Artin constituent une autre
source importante de ces résultats ; voir
\cite{ArtinI},
\cite{ArtinII},
\cite{Artin-Theorem-Representability},
\cite{Artin-Construction-Techniques},
\cite{Artin-Algebraic-Spaces},
\cite{Artin-Algebraic-Approximation},
\cite{Artin-Implicit-Function}, et
\cite{ArtinVersal}.

\medskip\noindent
Certaines notations, conventions et certains termes de ce chapitre sont peu
commodes et peuvent sembler inversés au lecteur expérimenté. C'est intentionnel.
On trouvera une explication dans Quot, section
\ref{quot-section-conventions}.



\section{Conventions}
\label{criteria-section-conventions}

\noindent
Les conventions employées dans ce chapitre sont les mêmes que dans le
chapitre consacré aux champs algébriques ; voir
Champs algébriques, section \ref{algebraic-section-conventions}.




\section{Ce que nous savons déjà}
\label{criteria-section-done-so-far}

\noindent
L'analogue de ce chapitre pour les espaces algébriques est le chapitre intitulé
« Amorçage » ; voir
Amorçage, section \ref{bootstrap-section-introduction}.
Ce chapitre contient déjà certains résultats de représentabilité.
En outre, nous avons déjà établi dans le chapitre sur les champs algébriques
une partie des résultats préliminaires qui y sont traités.
En voici la liste :
\begin{enumerate}
\item Les morphismes de préfaisceaux représentables par des espaces algébriques sont étudiés dans
Amorçage, section
\ref{bootstrap-section-morphism-representable-by-spaces}.
Dans
Champs algébriques, section
\ref{algebraic-section-morphisms-representable-by-algebraic-spaces}
nous étudions la notion de $1$-morphisme de catégories fibrées en groupoïdes
représentable par des espaces algébriques.
\item Les propriétés des morphismes de préfaisceaux représentables par des
espaces algébriques sont étudiées dans
Amorçage, section
\ref{bootstrap-section-representable-by-spaces-properties}.
Dans
Champs algébriques, section
\ref{algebraic-section-representable-properties}
nous étudions les propriétés des $1$-morphismes de catégories fibrées en groupoïdes
représentables par des espaces algébriques.
\item Nous avons démontré que, si $F$ est un faisceau dont la diagonale est représentable
par des espaces algébriques et qui possède un recouvrement \'etale par un espace
algébrique, alors $F$ est un espace algébrique ; voir
Amorçage, théorème \ref{bootstrap-theorem-bootstrap}.
(Il s'agit d'une version faible du résultat de l'alinéa suivant.)
\item
\label{criteria-item-bootstrap-final}
Nous avons démontré que, si $F$ est un faisceau et s'il existe un espace
algébrique $U$ et un morphisme $U \to F$ représentable par des espaces
algébriques, surjectif, plat et localement de présentation finie, alors
$F$ est un espace algébrique ; voir
Amorçage, théorème \ref{bootstrap-theorem-final-bootstrap}.
\item Nous avons aussi démontré l'analogue « lisse » de
(\ref{criteria-item-bootstrap-final}) pour les champs
algébriques : si $\mathcal{X}$ est un champ en groupoïdes sur
$(\Sch/S)_{fppf}$ et s'il existe un champ en groupoïdes
$\mathcal{U}$ sur $(\Sch/S)_{fppf}$ représentable
par un espace algébrique ainsi qu'un $1$-morphisme $u : \mathcal{U} \to \mathcal{X}$
représentable par des espaces algébriques, surjectif et lisse,
alors $\mathcal{X}$ est un champ algébrique ; voir
Champs algébriques, lemme
\ref{algebraic-lemma-smooth-surjective-morphism-implies-algebraic}.
\end{enumerate}
Notre première tâche consiste maintenant à démontrer l'analogue de
(\ref{criteria-item-bootstrap-final}) pour les champs
algébriques en général ; c'est le
théorème \ref{criteria-theorem-bootstrap}.



\section{Morphismes de champs en groupoïdes}
\label{criteria-section-1-morphisms}

\noindent
Cette section est préliminaire et peut être omise en première lecture.

\begin{lemma}
\label{criteria-lemma-etale-permanence}
Soient $\mathcal{X} \to \mathcal{Y} \to \mathcal{Z}$ des
$1$-morphismes de catégories fibrées en groupoïdes sur
$(\Sch/S)_{fppf}$.
Si $\mathcal{X} \to \mathcal{Z}$ et $\mathcal{Y} \to \mathcal{Z}$ sont
représentables par des espaces algébriques et étales, il en est de même de
$\mathcal{X} \to \mathcal{Y}$.
\end{lemma}

\begin{proof}
Soit $\mathcal{U}$ une catégorie fibrée en groupoïdes représentable sur $S$.
Soit $f : \mathcal{U} \to \mathcal{Y}$ un $1$-morphisme. Il faut montrer que
$\mathcal{X} \times_\mathcal{Y} \mathcal{U}$ est représentable par un
espace algébrique et étale sur $\mathcal{U}$.
Considérons la composée $h : \mathcal{U} \to \mathcal{Z}$. Alors
$$
\mathcal{X} \times_\mathcal{Z} \mathcal{U}
\longrightarrow
\mathcal{Y} \times_\mathcal{Z} \mathcal{U}
$$
est un $1$-morphisme entre deux catégories fibrées en groupoïdes, toutes deux
représentables par des espaces algébriques et étales sur $\mathcal{U}$.
Ainsi, d'après
Propriétés des espaces algébriques, lemme \ref{spaces-properties-lemma-etale-permanence},
il est représenté par un morphisme étale d'espaces algébriques.
Enfin, on obtient le résultat voulu puisque le morphisme $f$ induit
un morphisme $\mathcal{U} \to \mathcal{Y} \times_\mathcal{Z} \mathcal{U}$
et que l'on a
$$
\mathcal{X} \times_\mathcal{Y} \mathcal{U} =
(\mathcal{X} \times_\mathcal{Z} \mathcal{U})
\times_{(\mathcal{Y} \times_\mathcal{Z} \mathcal{U})}
\mathcal{U}.
$$
\end{proof}

\begin{lemma}
\label{criteria-lemma-stack-in-setoids-descent}
Soient $\mathcal{X}, \mathcal{Y}, \mathcal{Z}$ des champs en groupoïdes
sur $(\Sch/S)_{fppf}$. Supposons que $\mathcal{X} \to \mathcal{Y}$
et $\mathcal{Z} \to \mathcal{Y}$ soient des $1$-morphismes.
Si
\begin{enumerate}
\item $\mathcal{Y}$ et $\mathcal{Z}$ sont représentables par des espaces algébriques
$Y$ et $Z$ sur $S$,
\item le morphisme associé d'espaces algébriques $Y \to Z$ est surjectif,
plat et localement de présentation finie, et
\item $\mathcal{Y} \times_\mathcal{Z} \mathcal{X}$ est un champ en
groupoïdes discrets,
\end{enumerate}
alors $\mathcal{X}$ est un champ en groupoïdes discrets.
\end{lemma}

\begin{proof}
C'est un cas particulier de
Champs, lemme \ref{stacks-lemma-stack-in-setoids-descent}.
\end{proof}

\noindent
Le lemme suivant est l'analogue de
Champs algébriques, lemme
\ref{algebraic-lemma-smooth-surjective-morphism-implies-algebraic}
et sera supplanté par le
théorème \ref{criteria-theorem-bootstrap}, plus fort.

\begin{lemma}
\label{criteria-lemma-flat-finite-presentation-surjective-diagonal}
Soit $S$ un schéma.
Soit $u : \mathcal{U} \to \mathcal{X}$ un $1$-morphisme de
champs en groupoïdes sur $(\Sch/S)_{fppf}$. Si
\begin{enumerate}
\item $\mathcal{U}$ est représentable par un espace algébrique, et
\item $u$ est représentable par des espaces algébriques, surjectif, plat et
localement de présentation finie,
\end{enumerate}
alors
$\Delta : \mathcal{X} \to \mathcal{X} \times \mathcal{X}$
est représentable par des espaces algébriques.
\end{lemma}

\begin{proof}
Étant donnés deux schémas $T_1$, $T_2$ sur $S$, notons
$\mathcal{T}_i = (\Sch/T_i)_{fppf}$ les catégories fibrées représentables
associées. Supposons donnés des $1$-morphismes
$f_i : \mathcal{T}_i \to \mathcal{X}$.
D'après
Champs algébriques, lemme \ref{algebraic-lemma-representable-diagonal},
il suffit de démontrer que le $2$-produit
fibré $\mathcal{T}_1 \times_\mathcal{X} \mathcal{T}_2$
est représentable par un espace algébrique. D'après
Champs, lemme
\ref{stacks-lemma-2-fibre-product-stacks-in-setoids-over-stack-in-groupoids}
c'est en tout cas un champ en groupoïdes discrets. Ainsi,
$\mathcal{T}_1 \times_\mathcal{X} \mathcal{T}_2$ correspond
à un faisceau $F$ sur $(\Sch/S)_{fppf}$ ; voir
Champs, lemme \ref{stacks-lemma-stack-in-setoids-characterize}.
Soit $U$ l'espace algébrique qui représente $\mathcal{U}$.
Par hypothèse,
$$
\mathcal{T}_i' = \mathcal{U} \times_{u, \mathcal{X}, f_i} \mathcal{T}_i
$$
est représentable par un espace algébrique $T'_i$ sur $S$. Par conséquent,
$\mathcal{T}_1' \times_\mathcal{U} \mathcal{T}_2'$ est représentable
par l'espace algébrique $T'_1 \times_U T'_2$.
Considérons le diagramme commutatif
$$
\xymatrix{
&
\mathcal{T}_1 \times_{\mathcal X} \mathcal{T}_2 \ar[rr]\ar'[d][dd] & &
\mathcal{T}_1 \ar[dd] \\
\mathcal{T}_1' \times_\mathcal{U} \mathcal{T}_2' \ar[ur]\ar[rr]\ar[dd] & &
\mathcal{T}_1' \ar[ur]\ar[dd] \\
&
\mathcal{T}_2 \ar'[r][rr] & &
\mathcal X \\
\mathcal{T}_2' \ar[rr]\ar[ur] & &
\mathcal{U} \ar[ur] }
$$
Dans ce diagramme, les carrés inférieur, droit, arrière et
avant sont des $2$-produits fibrés. Un argument formel montre alors
que $\mathcal{T}_1' \times_\mathcal{U} \mathcal{T}_2' \to
\mathcal{T}_1 \times_{\mathcal X} \mathcal{T}_2$
est le « changement de base » de $\mathcal{U} \to \mathcal{X}$ ; plus précisément,
le diagramme
$$
\xymatrix{
\mathcal{T}_1' \times_\mathcal{U} \mathcal{T}_2' \ar[d] \ar[r] &
\mathcal{U} \ar[d] \\
\mathcal{T}_1 \times_{\mathcal X} \mathcal{T}_2 \ar[r] &
\mathcal{X}
}
$$
est un carré $2$-cartésien.
Ainsi, $T'_1 \times_U T'_2 \to F$ est représentable par des espaces algébriques,
plat, localement de présentation finie et surjectif ; voir
Champs algébriques, lemmes
\ref{algebraic-lemma-map-fibred-setoids-representable-algebraic-spaces},
\ref{algebraic-lemma-base-change-representable-by-spaces},
\ref{algebraic-lemma-map-fibred-setoids-property}, et
\ref{algebraic-lemma-base-change-representable-transformations-property}.
Par conséquent, $F$ est un espace algébrique d'après
Amorçage, théorème \ref{bootstrap-theorem-final-bootstrap},
ce qui conclut.
\end{proof}

\begin{lemma}
\label{criteria-lemma-second-diagonal}
Soit $\mathcal{X}$ une catégorie fibrée en groupoïdes sur $(\Sch/S)_{fppf}$.
Les assertions suivantes sont équivalentes :
\begin{enumerate}
\item $\Delta_\Delta : \mathcal{X} \to
\mathcal{X} \times_{\mathcal{X} \times \mathcal{X}} \mathcal{X}$
est représentable par des espaces algébriques,
\item pour tout $1$-morphisme $\mathcal{V} \to \mathcal{X} \times \mathcal{X}$
où $\mathcal{V}$ est représentable (par un schéma), le produit fibré
$\mathcal{Y} =
\mathcal{X} \times_{\Delta, \mathcal{X} \times \mathcal{X}} \mathcal{V}$
a une diagonale représentable par des espaces algébriques.
\end{enumerate}
\end{lemma}

\begin{proof}
Bien que cela demande un peu de réflexion, l'argument est entièrement formel.
Rappelons en effet que
$\mathcal{X} \times_{\mathcal{X} \times \mathcal{X}} \mathcal{X} =
\mathcal{I}_\mathcal{X}$ est l'inertie de $\mathcal{X}$ et que
$\Delta_\Delta$ est la section unité de $\mathcal{I}_\mathcal{X}$ ; voir
Catégories, section \ref{categories-section-inertia}.
La condition (1) signifie donc ceci : étant donnés un schéma $V$, un objet $x$ de
$\mathcal{X}$ au-dessus de $V$ et un morphisme $\alpha : x \to x$ de $\mathcal{X}_V$,
la condition « $\alpha = \text{id}_x$ » définit un espace algébrique sur $V$.
(Autrement dit, il existe un monomorphisme d'espaces algébriques $W \to V$
tel qu'un morphisme de schémas $f : T \to V$ se factorise par $W$
si et seulement si $f^*\alpha = \text{id}_{f^*x}$.)

\medskip\noindent
Soient d'autre part $V$ un schéma et $x, y$ des objets de
$\mathcal{X}$ au-dessus de $V$. Alors $(x, y)$ définit un morphisme
$\mathcal{V} = (\Sch/V)_{fppf} \to \mathcal{X} \times \mathcal{X}$.
Soient ensuite $h : V' \to V$ un morphisme de schémas et
$\alpha : h^*x \to h^*y$ et $\beta : h^*x \to h^*y$ des morphismes
de $\mathcal{X}_{V'}$. Alors $(\alpha, \beta)$ définit un morphisme
$\mathcal{V}' = (\Sch/V)_{fppf} \to \mathcal{Y} \times \mathcal{Y}$.
La condition (2) affirme alors que, quels que soient les choix ci-dessus, la
condition « $\alpha = \beta$ » définit un espace algébrique sur $V$.

\medskip\noindent
Pour voir l'équivalence, étant donné $(\alpha, \beta)$ comme en (2), on constate que
(1) implique que « $\alpha^{-1} \circ \beta = \text{id}_{h^*x}$ »
définit un espace algébrique. L'implication (2) $\Rightarrow$ (1)
s'obtient en prenant $h = \text{id}_V$ et $\beta = \text{id}_x$.
\end{proof}











\section{Commutation aux limites sur les objets}
\label{criteria-section-limit-preserving}

\noindent
Soit $S$ un schéma. Soit $p : \mathcal{X} \to \mathcal{Y}$ un $1$-morphisme
de catégories fibrées en groupoïdes sur $(\Sch/S)_{fppf}$. Nous dirons que
$p$ {\it commute aux limites sur les objets} si la condition suivante est satisfaite :
pour toute donnée formée de
\begin{enumerate}
\item un schéma affine $U = \lim_{i \in I} U_i$ écrit comme
limite projective filtrante de schémas affines $U_i$ sur $S$,
\item un objet $y_i$ de $\mathcal{Y}$ au-dessus de $U_i$ pour un certain $i$,
\item un objet $x$ de $\mathcal{X}$ au-dessus de $U$, et
\item un isomorphisme $\gamma : p(x) \to y_i|_U$,
\end{enumerate}
il existe un $i' \geq i$, un objet $x_{i'}$ de
$\mathcal{X}$ au-dessus de $U_{i'}$, un isomorphisme
$\beta : x_{i'}|_U \to x$ et un isomorphisme
$\gamma_{i'} : p(x_{i'}) \to y_i|_{U_{i'}}$
tels que
\begin{equation}
\label{criteria-equation-limit-preserving}
\vcenter{
\xymatrix{
p(x_{i'}|_U) \ar[d]_{p(\beta)} \ar[rr]_{\gamma_{i'}|_U} & &
(y_i|_{U_{i'}})|_U \ar@{=}[d] \\
p(x) \ar[rr]^\gamma & & y_i|_U
}
}
\end{equation}
soit commutatif. Dans cette situation, on dit que « $(i', x_{i'}, \beta, \gamma_{i'})$
est une {\it solution} du problème posé par les données (1), (2), (3), (4) ».
La motivation de cette définition vient de
Limits of Espaces,
lemme \ref{spaces-limits-lemma-characterize-relative-limit-preserving}.

\begin{lemma}
\label{criteria-lemma-base-change-limit-preserving}
Soient $p : \mathcal{X} \to \mathcal{Y}$ et $q : \mathcal{Z} \to \mathcal{Y}$
des $1$-morphismes de catégories fibrées en groupoïdes sur $(\Sch/S)_{fppf}$.
Si $p : \mathcal{X} \to \mathcal{Y}$ commute aux limites sur les objets, il en
est de même du changement de base
$p' : \mathcal{X} \times_\mathcal{Y} \mathcal{Z} \to \mathcal{Z}$
de $p$ par $q$.
\end{lemma}

\begin{proof}
L'assertion est formelle. Soit $U = \lim_{i \in I} U_i$ la limite projective filtrante
de schémas affines $U_i$ sur $S$, soit $z_i$ un objet de $\mathcal{Z}$
au-dessus de $U_i$ pour un certain $i$, soit $w$ un objet de
$\mathcal{X} \times_\mathcal{Y} \mathcal{Z}$ au-dessus de $U$, et soit
$\delta : p'(w) \to z_i|_U$ un isomorphisme.
On peut écrire
$w = (U, x, z, \alpha)$, où $x$ est un objet de $\mathcal{X}$ au-dessus de $U$,
$z$ un objet de $\mathcal{Z}$ au-dessus de $U$ et
$\alpha : p(x) \to q(z)$ un isomorphisme. Puisque $p'(w) = z$, on a
$\delta : z \to z_i|_U$. Posons $y_i = q(z_i)$ et
$\gamma = q(\delta) \circ \alpha : p(x) \to y_i|_U$.
Comme $p$ commute aux limites sur les objets, il existe $i' \geq i$
et un objet $x_{i'}$ de $\mathcal{X}$ au-dessus de $U_{i'}$, ainsi que des
isomorphismes $\beta : x_{i'}|_U \to x$ et
$\gamma_{i'} : p(x_{i'}) \to y_i|_{U_{i'}}$ tels que
(\ref{criteria-equation-limit-preserving}) soit commutatif. Considérons alors l'objet
$w_{i'} = (U_{i'}, x_{i'}, z_i|_{U_{i'}}, \gamma_{i'})$ de
$\mathcal{X} \times_\mathcal{Y} \mathcal{Z}$ au-dessus de $U_{i'}$
et définissons les isomorphismes
$$
w_{i'}|_U = (U, x_{i'}|_U, z_i|_U, \gamma_{i'}|_U)
\xrightarrow{(\beta, \delta^{-1})}
(U, x, z, \alpha) = w
$$
et
$$
p'(w_{i'}) = z_i|_{U_{i'}} \xrightarrow{\text{id}} z_i|_{U_{i'}}.
$$
Ceux-ci se combinent en une solution du problème.
\end{proof}

\begin{lemma}
\label{criteria-lemma-composition-limit-preserving}
Soient $p : \mathcal{X} \to \mathcal{Y}$ et $q : \mathcal{Y} \to \mathcal{Z}$
des $1$-morphismes de catégories fibrées en groupoïdes sur $(\Sch/S)_{fppf}$.
Si $p$ et $q$ commutent aux limites sur les objets, il en est de même de la composée
$q \circ p$.
\end{lemma}

\begin{proof}
L'assertion est formelle. Soit $U = \lim_{i \in I} U_i$ la limite projective filtrante
de schémas affines $U_i$ sur $S$, soit $z_i$ un objet de $\mathcal{Z}$
au-dessus de $U_i$ pour un certain $i$, soit $x$ un objet de $\mathcal{X}$ au-dessus de $U$,
et soit $\gamma : q(p(x)) \to z_i|_U$ un isomorphisme. Comme $q$
commute aux limites sur les objets, il existe $i' \geq i$, un objet
$y_{i'}$ de $\mathcal{Y}$ au-dessus de $U_{i'}$, un isomorphisme
$\beta : y_{i'}|_U \to p(x)$ et un isomorphisme
$\gamma_{i'} : q(y_{i'}) \to z_i|_{U_{i'}}$
tels que (\ref{criteria-equation-limit-preserving}) soit commutatif. Comme $p$
commute aux limites sur les objets, il existe $i'' \geq i'$, un objet
$x_{i''}$ de $\mathcal{X}$ au-dessus de $U_{i''}$, un isomorphisme
$\beta' : x_{i''}|_U \to x$ et un isomorphisme
$\gamma'_{i''} : p(x_{i''}) \to y_{i'}|_{U_{i''}}$
tels que (\ref{criteria-equation-limit-preserving}) soit commutatif.
La solution consiste à prendre $x_{i''}$ au-dessus de $U_{i''}$ avec l'isomorphisme
$$
q(p(x_{i''})) \xrightarrow{q(\gamma'_{i''})}
q(y_{i'})|_{U_{i''}} \xrightarrow{\gamma_{i'}|_{U_{i''}}}
z_i|_{U_{i''}}
$$
et l'isomorphisme $\beta' : x_{i''}|_U \to x$. Nous omettons de vérifier
que (\ref{criteria-equation-limit-preserving}) est commutatif.
\end{proof}

\begin{lemma}
\label{criteria-lemma-representable-by-spaces-limit-preserving}
Soit $p : \mathcal{X} \to \mathcal{Y}$ un $1$-morphisme de catégories
fibrées en groupoïdes sur $(\Sch/S)_{fppf}$. Si $p$ est
représentable par des espaces algébriques, les assertions suivantes sont équivalentes :
\begin{enumerate}
\item $p$ commute aux limites sur les objets, et
\item $p$ est localement de présentation finie (voir
Champs algébriques,
définition \ref{algebraic-definition-relative-representable-property}).
\end{enumerate}
\end{lemma}

\begin{proof}
Supposons (2). Soit $U = \lim_{i \in I} U_i$ la limite projective filtrante
de schémas affines $U_i$ sur $S$, soit $y_i$ un objet de $\mathcal{Y}$
au-dessus de $U_i$ pour un certain $i$, soit $x$ un objet de $\mathcal{X}$ au-dessus de $U$,
et soit $\gamma : p(x) \to y_i|_U$ un isomorphisme. Notons
$X_{y_i}$ un espace algébrique sur $U_i$ représentant le $2$-produit
fibré
$$
(\Sch/U_i)_{fppf} \times_{y_i, \mathcal{Y}, p} \mathcal{X}.
$$
Notons que $\xi = (U, U \to U_i, x, \gamma^{-1})$ définit un objet de
ce $2$-produit fibré au-dessus de $U$. Par le lemme de $2$-Yoneda, $\xi$ correspond
à un morphisme $f_\xi : U \to X_{y_i}$ sur $U_i$. D'après
Limits of Espaces, proposition
\ref{spaces-limits-proposition-characterize-locally-finite-presentation}
il existe $i' \geq i$ et un morphisme $f_{i'} : U_{i'} \to X_{y_i}$
tels que $f_\xi$ soit la composée de $f_{i'}$ et du morphisme de
projection $U \to U_{i'}$. En outre, le lemme de $2$-Yoneda montre que
$f_{i'}$ correspond à un objet
$\xi_{i'} = (U_{i'}, U_{i'} \to U_i, x_{i'}, \alpha)$
du $2$-produit fibré affiché au-dessus de $U_{i'}$, dont la restriction à
$U$ redonne $\xi$. En particulier, on obtient un isomorphisme
$\gamma : x_{i'}|U \to x$. Notons que $\alpha : y_i|_{U_{i'}} \to p(x_{i'})$.
Ainsi, prendre $x_{i'}$, l'isomorphisme
$\gamma : x_{i'}|U \to x$ et l'isomorphisme
$\beta = \alpha^{-1} : p(x_{i'}) \to y_i|_{U_{i'}}$
fournit une solution du problème.

\medskip\noindent
Supposons (1). Choisissons un schéma $T$ et un $1$-morphisme
$y : (\Sch/T)_{fppf} \to \mathcal{Y}$. Soit
$X_y$ un espace algébrique sur $T$ représentant le $2$-produit fibré
$(\Sch/T)_{fppf} \times_{y, \mathcal{Y}, p} \mathcal{X}$.
Il faut montrer que $X_y \to T$ est localement de présentation finie.
Pour cela, nous utiliserons le critère de
Limits of Espaces, remarque \ref{spaces-limits-remark-limit-preserving}.
Considérons un schéma affine $U = \lim_{i \in I} U_i$ écrit comme
limite projective filtrante de schémas affines sur $T$.
Choisissons $i \in I$ et posons $y_i = y|_{U_i}$. Notons aussi $i'$ un élément
de $I$ supérieur ou égal à $i$. D'après le lemme de $2$-Yoneda, les
morphismes $U \to X_y$ sur $T$ correspondent bijectivement
aux classes d'isomorphisme de couples $(x, \alpha)$, où $x$ est un objet
de $\mathcal{X}$ au-dessus de $U$ et $\alpha : y|_U \to p(x)$ un isomorphisme.
Bien entendu, donner $\alpha$ revient, à inversion près, à donner
un isomorphisme $\gamma : p(x) \to y_i|_U$.
Il en va de même pour les morphismes $U_{i'} \to X_y$ sur $T$. Ainsi, (1) garantit
que l'application canonique
$$
\colim_{i' \geq i} X_y(U_{i'}) \longrightarrow X_y(U)
$$
est surjective dans cette situation. Il résulte de
Limits of Espaces, lemme \ref{spaces-limits-lemma-surjection-is-enough},
que $X_y \to T$ est localement de présentation finie.
\end{proof}

\begin{lemma}
\label{criteria-lemma-open-immersion-limit-preserving}
Soit $p : \mathcal{X} \to \mathcal{Y}$ un $1$-morphisme de catégories
fibrées en groupoïdes sur $(\Sch/S)_{fppf}$. Supposons $p$ représentable
par des espaces algébriques et une immersion ouverte. Alors $p$ commute aux
limites sur les objets.
\end{lemma}

\begin{proof}
Cela résulte du
lemme \ref{criteria-lemma-representable-by-spaces-limit-preserving}
et, par le principe général de
Champs algébriques, lemme
\ref{algebraic-lemma-representable-transformations-property-implication})
du fait qu'une immersion ouverte d'espaces algébriques est
localement de présentation finie ; voir
Morphismes d'espaces algébriques, lemme
\ref{spaces-morphisms-lemma-open-immersion-locally-finite-presentation}.
\end{proof}

\noindent
Soit $S$ un schéma. Dans le lemme suivant, nous aurons besoin de la notion de
{\it taille} d'un espace algébrique $X$ sur $S$. Étant donné un cardinal
$\kappa$, nous dirons que $X$ vérifie $\text{size}(X) \leq \kappa$ si et seulement
s'il existe un schéma $U$ tel que $\text{size}(U) \leq \kappa$ (voir
Ensembles, section \ref{sets-section-categories-schemes}) et un morphisme
étale surjectif $U \to X$.

\begin{lemma}
\label{criteria-lemma-check-representable-limit-preserving}
Soit $S$ un schéma.
Soit $\kappa = \text{size}(T)$ pour un certain $T \in \Ob((\Sch/S)_{fppf})$.
Soit $f : \mathcal{X} \to \mathcal{Y}$ un $1$-morphisme
de catégories fibrées en groupoïdes sur $(\Sch/S)_{fppf}$
tel que
\begin{enumerate}
\item $\mathcal{Y} \to (\Sch/S)_{fppf}$ commute aux limites sur les objets,
\item pour tout schéma affine $V$ localement de présentation finie sur $S$ et
tout $y \in \Ob(\mathcal{Y}_V)$, le produit fibré
$(\Sch/V)_{fppf} \times_{y, \mathcal{Y}} \mathcal{X}$ est représentable
par un espace algébrique de taille $\leq \kappa$\footnote{Ceux qui négligent les
questions ensemblistes peuvent supprimer la condition sur la taille.},
\item $\mathcal{X}$ et $\mathcal{Y}$ sont des champs pour la topologie de Zariski.
\end{enumerate}
Alors $f$ est représentable par des espaces algébriques.
\end{lemma}

\begin{proof}
Soient $V$ un schéma sur $S$ et $y \in \mathcal{Y}_V$. Il faut démontrer que
$(\Sch/V)_{fppf} \times_{y, \mathcal{Y}} \mathcal{X}$ est représentable
par un espace algébrique.

\medskip\noindent
Cas I : $V$ est affine et s'envoie dans un ouvert affine $\Spec(\Lambda) \subset S$.
On peut alors écrire $V = \lim V_i$, chaque $V_i$ étant affine et de
présentation finie sur $\Spec(\Lambda)$ ; voir
Algèbre, lemme \ref{algebra-lemma-ring-colimit-fp}.
L'hypothèse (1) montre alors que $y$ provient d'un objet $y_i$ au-dessus de $V_i$ pour un certain $i$.
Par l'hypothèse (3), le produit fibré
$(\Sch/V_i)_{fppf} \times_{y_i, \mathcal{Y}} \mathcal{X}$ est représentable
par un espace algébrique $Z_i$. Alors
$(\Sch/V)_{fppf} \times_{y, \mathcal{Y}} \mathcal{X}$ est représentable
par $Z \times_{V_i} V$.

\medskip\noindent
Cas II : $V$ est quelconque. Choisissons un recouvrement ouvert affine
$V = \bigcup_{i \in I} V_i$ tel que chaque $V_i$ s'envoie dans un ouvert affine
de $S$. Nous affirmons d'abord
que $\mathcal{Z} = (\Sch/V)_{fppf} \times_{y, \mathcal{Y}} \mathcal{X}$
est un champ en groupoïdes discrets pour la topologie de Zariski. En effet, c'est un champ en
groupoïdes pour la topologie de Zariski d'après
Champs, lemme \ref{stacks-lemma-2-product-stacks-in-groupoids}.
Supposons ensuite que $z$ soit un objet de $\mathcal{Z}$ au-dessus d'un schéma $T$.
Notons $g : T \to V$ le morphisme correspondant à la
projection de $z$ dans $(\Sch/V)_{fppf}$. Considérons le faisceau de Zariski
$\mathit{I} = \mathit{Isom}_{\mathcal{Z}}(z, z)$. Le cas I montre que
$\mathit{I}|_{g^{-1}(V_i)} = *$ (le faisceau réduit à un élément). Ainsi,
$\mathcal{I} = *$. Par conséquent, $\mathcal{Z}$ est fibrée en groupoïdes discrets. Pour achever
la démonstration, il faut montrer que le faisceau de Zariski
$Z : T \mapsto \Ob(\mathcal{Z}_T)/\cong$ est un espace algébrique ; voir
Champs algébriques, lemme
\ref{algebraic-lemma-characterize-representable-by-space}.
Il existe une application $p : Z \to V$ (une transformation de foncteurs) et le cas I
montre que $Z_i = p^{-1}(V_i)$ est un espace algébrique. Les morphismes
$Z_i \to Z$ sont représentables par des immersions ouvertes et
$\coprod Z_i \to Z$ est surjectif (pour la topologie de Zariski).
Ainsi, $Z$ est un faisceau pour la topologie fppf d'après
Amorçage, lemme \ref{bootstrap-lemma-glueing-sheaves}.
Le lemme de Espaces, \ref{spaces-lemma-glueing-algebraic-spaces},
s'applique donc et montre que $Z$ est un espace algébrique\footnote{
Pour vérifier la condition ensembliste de ce lemme,
raisonnons comme suit. Choisissons d'abord le recouvrement ouvert de sorte que
$|I| \leq \text{size}(V)$. Choisissons ensuite des schémas $U_i$ de taille
$\leq \max(\kappa, \text{size}(V))$ et des morphismes étales surjectifs
$U_i \to Z_i$ ; cela est possible par l'hypothèse (2) et
Ensembles, lemme \ref{sets-lemma-bound-size-fibre-product}
(nous omettons les détails). Alors
Ensembles, lemme \ref{sets-lemma-what-is-in-it},
implique que $\coprod U_i$ est un objet de $(\Sch/S)_{fppf}$.
Par conséquent, $\coprod Z_i$ est un espace algébrique d'après
Espaces, lemme \ref{spaces-lemma-coproduct-algebraic-spaces}.
}.
\end{proof}

\begin{lemma}
\label{criteria-lemma-check-property-limit-preserving}
Soit $S$ un schéma. Soit $f : \mathcal{X} \to \mathcal{Y}$ un $1$-morphisme
de catégories fibrées en groupoïdes sur $(\Sch/S)_{fppf}$. Soit $\mathcal{P}$
une propriété des morphismes d'espaces algébriques comme dans
Champs algébriques, définition
\ref{algebraic-definition-relative-representable-property}. Si
\begin{enumerate}
\item $f$ est représentable par des espaces algébriques,
\item $\mathcal{Y} \to (\Sch/S)_{fppf}$ commute aux limites sur les objets,
\item pour tout schéma affine $V$ localement de présentation finie sur $S$ et
tout $y \in \mathcal{Y}_V$, le morphisme d'espaces algébriques qui en résulte
$f_y : F_y \to V$, voir Champs algébriques, équation
(\ref{algebraic-equation-representable-by-algebraic-spaces}),
a la propriété $\mathcal{P}$.
\end{enumerate}
Alors $f$ a la propriété $\mathcal{P}$.
\end{lemma}

\begin{proof}
Soient $V$ un schéma sur $S$ et $y \in \mathcal{Y}_V$. Il faut montrer
que $F_y \to V$ a la propriété $\mathcal{P}$. Comme $\mathcal{P}$ est
locale pour la topologie fppf sur la base, on peut supposer que $V$ est un schéma affine qui
s'envoie dans un ouvert affine $\Spec(\Lambda) \subset S$. On peut donc écrire
$V = \lim V_i$, chaque $V_i$ étant affine et de présentation finie sur
$\Spec(\Lambda)$ ; voir Algèbre, lemme \ref{algebra-lemma-ring-colimit-fp}.
L'hypothèse (2) montre alors que $y$ provient d'un objet $y_i$ au-dessus de $V_i$ pour un certain $i$.
Par l'hypothèse (3), le morphisme $F_{y_i} \to V_i$ a la propriété $\mathcal{P}$.
Comme $\mathcal{P}$ est stable par tout changement de base et que
$F_y = F_{y_i} \times_{V_i} V$, on conclut que $F_y \to V$ a la propriété
$\mathcal{P}$, comme voulu.
\end{proof}



\section{Formellement lisse sur les objets}
\label{criteria-section-formally-smooth}

\noindent
Soit $S$ un schéma. Soit $p : \mathcal{X} \to \mathcal{Y}$ un $1$-morphisme
de catégories fibrées en groupoïdes sur $(\Sch/S)_{fppf}$. Nous dirons que
$p$ est {\it formellement lisse sur les objets} si la condition suivante est satisfaite :
pour toute donnée formée de
\begin{enumerate}
\item un épaississement du premier ordre $U \subset U'$ de schémas affines sur $S$,
\item un objet $y'$ de $\mathcal{Y}$ au-dessus de $U'$,
\item un objet $x$ de $\mathcal{X}$ au-dessus de $U$, et
\item un isomorphisme $\gamma : p(x) \to y'|_U$,
\end{enumerate}
il existe un objet $x'$ de
$\mathcal{X}$ au-dessus de $U'$, muni d'un isomorphisme
$\beta : x'|_U \to x$ et d'un isomorphisme $\gamma' : p(x') \to y'$
tels que
\begin{equation}
\label{criteria-equation-formally-smooth}
\vcenter{
\xymatrix{
p(x'|_U) \ar[d]_{p(\beta)} \ar[rr]_{\gamma'|_U} & &
y'|_U \ar@{=}[d] \\
p(x) \ar[rr]^\gamma & & y'|_U
}
}
\end{equation}
soit commutatif. Dans cette situation, on dit que « $(x', \beta, \gamma')$
est une {\it solution} du problème posé par les données (1), (2), (3), (4) ».

\begin{lemma}
\label{criteria-lemma-base-change-formally-smooth}
Soient $p : \mathcal{X} \to \mathcal{Y}$ et $q : \mathcal{Z} \to \mathcal{Y}$
des $1$-morphismes de catégories fibrées en groupoïdes sur $(\Sch/S)_{fppf}$.
Si $p : \mathcal{X} \to \mathcal{Y}$ est formellement lisse sur les objets, il en
est de même du changement de base
$p' : \mathcal{X} \times_\mathcal{Y} \mathcal{Z} \to \mathcal{Z}$
de $p$ par $q$.
\end{lemma}

\begin{proof}
L'assertion est formelle. Soit $U \subset U'$ un épaississement du premier ordre
de schémas affines sur $S$, soit $z'$ un objet de $\mathcal{Z}$
au-dessus de $U'$, soit $w$ un objet de
$\mathcal{X} \times_\mathcal{Y} \mathcal{Z}$ au-dessus de $U$, et soit
$\delta : p'(w) \to z'|_U$ un isomorphisme.
On peut écrire
$w = (U, x, z, \alpha)$, où $x$ est un objet de $\mathcal{X}$ au-dessus de $U$,
$z$ un objet de $\mathcal{Z}$ au-dessus de $U$ et
$\alpha : p(x) \to q(z)$ un isomorphisme. Puisque $p'(w) = z$, on a
$\delta : z \to z|_U$. Posons $y' = q(z')$ et
$\gamma = q(\delta) \circ \alpha : p(x) \to y'|_U$.
Comme $p$ est formellement lisse sur les objets, il existe un
objet $x'$ de $\mathcal{X}$ au-dessus de $U'$, ainsi que des
isomorphismes $\beta : x'|_U \to x$ et $\gamma' : p(x') \to y'$ tels que
(\ref{criteria-equation-formally-smooth}) soit commutatif. Considérons alors l'objet
$w = (U', x', z', \gamma')$ de $\mathcal{X} \times_\mathcal{Y} \mathcal{Z}$
au-dessus de $U'$ et définissons les isomorphismes
$$
w'|_U = (U, x'|_U, z'|_U, \gamma'|_U)
\xrightarrow{(\beta, \delta^{-1})}
(U, x, z, \alpha) = w
$$
et
$$
p'(w') = z' \xrightarrow{\text{id}} z'.
$$
Ceux-ci se combinent en une solution du problème.
\end{proof}

\begin{lemma}
\label{criteria-lemma-composition-formally-smooth}
Soient $p : \mathcal{X} \to \mathcal{Y}$ et $q : \mathcal{Y} \to \mathcal{Z}$
des $1$-morphismes de catégories fibrées en groupoïdes sur $(\Sch/S)_{fppf}$.
Si $p$ et $q$ sont formellement lisses sur les objets, il en est de même de la composée
$q \circ p$.
\end{lemma}

\begin{proof}
L'assertion est formelle. Soit $U \subset U'$ un épaississement du premier ordre
de schémas affines sur $S$, soit $z'$ un objet de $\mathcal{Z}$
au-dessus de $U'$, soit $x$ un objet de $\mathcal{X}$ au-dessus de $U$,
et soit $\gamma : q(p(x)) \to z'|_U$ un isomorphisme. Comme $q$ est
formellement lisse sur les objets, il existe un objet
$y'$ de $\mathcal{Y}$ au-dessus de $U'$, un isomorphisme
$\beta : y'|_U \to p(x)$ et un isomorphisme $\gamma' : q(y') \to z'$
tels que (\ref{criteria-equation-formally-smooth}) soit commutatif. Comme $p$ est
formellement lisse sur les objets, il existe un objet
$x'$ de $\mathcal{X}$ au-dessus de $U'$, un isomorphisme
$\beta' : x'|_U \to x$ et un isomorphisme $\gamma'' : p(x') \to y'$
tels que (\ref{criteria-equation-formally-smooth}) soit commutatif.
La solution consiste à prendre $x'$ au-dessus de $U'$ avec l'isomorphisme
$$
q(p(x')) \xrightarrow{q(\gamma'')} q(y') \xrightarrow{\gamma'} z'
$$
et l'isomorphisme $\beta' : x'|_U \to x$. Nous omettons de vérifier
que (\ref{criteria-equation-formally-smooth}) est commutatif.
\end{proof}

\noindent
Notons que la classe des morphismes formellement lisses d'espaces algébriques est
stable par tout changement de base et locale sur le but pour la
topologie fpqc ; voir
Compléments sur les morphismes d'espaces algébriques,
lemmes \ref{spaces-more-morphisms-lemma-base-change-formally-smooth} et
\ref{spaces-more-morphisms-lemma-descending-property-formally-smooth}.
La condition (2) du lemme ci-dessous a donc un sens.

\begin{lemma}
\label{criteria-lemma-representable-by-spaces-formally-smooth}
Soit $p : \mathcal{X} \to \mathcal{Y}$ un $1$-morphisme de catégories
fibrées en groupoïdes sur $(\Sch/S)_{fppf}$. Si $p$ est
représentable par des espaces algébriques, les assertions suivantes sont équivalentes :
\begin{enumerate}
\item $p$ est formellement lisse sur les objets, et
\item $p$ est formellement lisse (voir
Champs algébriques,
définition \ref{algebraic-definition-relative-representable-property}).
\end{enumerate}
\end{lemma}

\begin{proof}
Supposons (2). Soit $U \subset U'$ un épaississement du premier ordre
de schémas affines sur $S$, soit $y'$ un objet de $\mathcal{Y}$
au-dessus de $U'$, soit $x$ un objet de $\mathcal{X}$ au-dessus de $U$,
et soit $\gamma : p(x) \to y'|_U$ un isomorphisme. Notons
$X_{y'}$ un espace algébrique sur $U'$ représentant le $2$-produit
fibré
$$
(\Sch/U')_{fppf} \times_{y', \mathcal{Y}, p} \mathcal{X}.
$$
Notons que $\xi = (U, U \to U', x, \gamma^{-1})$ définit un objet de
ce $2$-produit fibré au-dessus de $U$. Par le lemme de $2$-Yoneda, $\xi$ correspond
à un morphisme $f_\xi : U \to X_{y'}$ sur $U'$. Comme $X_{y'} \to U'$ est
formellement lisse par hypothèse, il existe un morphisme
$f' : U' \to X_{y'}$ tel que $f_\xi$ soit la composée de $f'$
et du morphisme $U \to U'$. En outre, le lemme de $2$-Yoneda montre que
$f'$ correspond à un objet $\xi' = (U', U' \to U', x', \alpha)$ du
$2$-produit fibré affiché au-dessus de $U'$, dont la restriction à
$U$ redonne $\xi$. En particulier, on obtient un isomorphisme
$\gamma : x'|U \to x$. Notons que $\alpha : y' \to p(x')$.
Ainsi, prendre $x'$, l'isomorphisme
$\gamma : x'|U \to x$ et l'isomorphisme
$\beta = \alpha^{-1} : p(x') \to y'$
fournit une solution du problème.

\medskip\noindent
Supposons (1). Choisissons un schéma $T$ et un $1$-morphisme
$y : (\Sch/T)_{fppf} \to \mathcal{Y}$. Soit
$X_y$ un espace algébrique sur $T$ représentant le $2$-produit fibré
$(\Sch/T)_{fppf} \times_{y, \mathcal{Y}, p} \mathcal{X}$.
Il faut montrer que $X_y \to T$ est formellement lisse.
Il suffit donc de montrer que, pour tout épaississement du premier ordre
$U \subset U'$ de schémas affines sur $T$,
$X_y(U') \to X_y(U')$ est surjective (il s'agit de morphismes dans la
catégorie des espaces algébriques sur $T$). Posons $y' = y|_{U'}$.
D'après le lemme de $2$-Yoneda, les morphismes $U \to X_y$ sur $T$ correspondent bijectivement
aux classes d'isomorphisme de couples $(x, \alpha)$, où $x$ est un objet
de $\mathcal{X}$ au-dessus de $U$ et $\alpha : y|_U \to p(x)$ un isomorphisme.
Bien entendu, donner $\alpha$ revient, à inversion près, à donner
un isomorphisme $\gamma : p(x) \to y'|_U$.
Il en va de même pour les morphismes $U' \to X_y$ sur $T$. Ainsi, (1) garantit
la surjectivité de $X_y(U') \to X_y(U')$
dans cette situation, ce qui conclut.
\end{proof}







\section{Surjectif sur les objets}
\label{criteria-section-formally-surjective}

\noindent
Soit $S$ un schéma. Soit $p : \mathcal{X} \to \mathcal{Y}$ un $1$-morphisme
de catégories fibrées en groupoïdes sur $(\Sch/S)_{fppf}$. Nous dirons que
$p$ est {\it surjectif sur les objets} si la condition suivante est satisfaite :
pour toute donnée formée de
\begin{enumerate}
\item un corps $k$ sur $S$, et
\item un objet $y$ de $\mathcal{Y}$ au-dessus de $\Spec(k)$,
\end{enumerate}
il existe une extension de corps $K/k$ sur $S$ et un
objet $x$ de $\mathcal{X}$ au-dessus de $\Spec(K)$
tel que $p(x) \cong y|_{\Spec(K)}$.

\begin{lemma}
\label{criteria-lemma-base-change-surjective}
Soient $p : \mathcal{X} \to \mathcal{Y}$ et $q : \mathcal{Z} \to \mathcal{Y}$
des $1$-morphismes de catégories fibrées en groupoïdes sur $(\Sch/S)_{fppf}$.
Si $p : \mathcal{X} \to \mathcal{Y}$ est surjectif sur les objets, il en
est de même du changement de base
$p' : \mathcal{X} \times_\mathcal{Y} \mathcal{Z} \to \mathcal{Z}$
de $p$ par $q$.
\end{lemma}

\begin{proof}
L'assertion est formelle. Soit $z$ un objet de $\mathcal{Z}$ au-dessus d'un corps $k$.
Comme $p$ est surjectif sur les objets, il existe une extension $K/k$,
un objet $x$ de $\mathcal{X}$ au-dessus de $K$ et un isomorphisme
$\alpha : p(x) \to q(z)|_{\Spec(K)}$. Alors
$w = (\Spec(K), x, z|_{\Spec(K)}, \alpha)$ est un objet de
$\mathcal{X} \times_\mathcal{Y} \mathcal{Z}$ au-dessus de $K$ tel que
$p'(w) = z|_{\Spec(K)}$.
\end{proof}

\begin{lemma}
\label{criteria-lemma-composition-surjective}
Soient $p : \mathcal{X} \to \mathcal{Y}$ et $q : \mathcal{Y} \to \mathcal{Z}$
des $1$-morphismes de catégories fibrées en groupoïdes sur $(\Sch/S)_{fppf}$.
Si $p$ et $q$ sont surjectifs sur les objets, il en est de même de la composée
$q \circ p$.
\end{lemma}

\begin{proof}
L'assertion est formelle. Soit $z$ un objet de $\mathcal{Z}$ au-dessus d'un corps $k$.
Comme $q$ est surjectif sur les objets, il existe une extension de corps $K/k$
et un objet $y$ de $\mathcal{Y}$ au-dessus de $K$ tels que
$q(y) \cong x|_{\Spec(K)}$. Comme $p$ est surjectif sur les objets, il
existe une extension de corps $L/K$ et un objet $x$ de $\mathcal{X}$
au-dessus de $L$ tel que $p(x) \cong y|_{\Spec(L)}$. Alors l'extension de corps
$L/k$ et l'objet $x$ de $\mathcal{X}$ au-dessus de $L$ vérifient
$q(p(x)) \cong z|_{\Spec(L)}$, comme voulu.
\end{proof}

\begin{lemma}
\label{criteria-lemma-representable-by-spaces-surjective}
Soit $p : \mathcal{X} \to \mathcal{Y}$ un $1$-morphisme de catégories
fibrées en groupoïdes sur $(\Sch/S)_{fppf}$. Si $p$ est
représentable par des espaces algébriques, les assertions suivantes sont équivalentes :
\begin{enumerate}
\item $p$ est surjectif sur les objets, et
\item $p$ est surjectif (voir
Champs algébriques,
définition \ref{algebraic-definition-relative-representable-property}).
\end{enumerate}
\end{lemma}

\begin{proof}
Supposons (2). Soit $k$ un corps et soit $y$ un objet de
$\mathcal{Y}$ au-dessus de $k$. Notons $X_y$ un espace algébrique sur $k$
représentant le $2$-produit fibré
$$
(\Sch/\Spec(k))_{fppf} \times_{y, \mathcal{Y}, p} \mathcal{X}.
$$
Puisque $p$ est supposé surjectif, $X_y$ n'est pas vide.
Il existe donc une extension de corps $K/k$ et, à valeurs dans $K$, un point
$x$ de $X_y$. Par le lemme de $2$-Yoneda, cela correspond à un objet
$x$ de $\mathcal{X}$ au-dessus de $K$, muni d'un isomorphisme
$p(x) \cong y|_{\Spec(K)}$ ; ainsi, (1) est satisfaite.

\medskip\noindent
Supposons (1). Choisissons un schéma $T$ et un $1$-morphisme
$y : (\Sch/T)_{fppf} \to \mathcal{Y}$. Soit
$X_y$ un espace algébrique sur $T$ représentant le $2$-produit fibré
$(\Sch/T)_{fppf} \times_{y, \mathcal{Y}, p} \mathcal{X}$.
Il faut montrer que $X_y \to T$ est surjectif. D'après
Morphismes d'espaces algébriques, définition \ref{spaces-morphisms-definition-surjective},
il faut montrer que $|X_y| \to |T|$ est surjectif.
Cela signifie exactement que, pour tout corps $k$ sur $T$ et tout
morphisme $t : \Spec(k) \to T$, il existe une extension de corps
$K/k$ et un morphisme $x : \Spec(K) \to X_y$ tels que
$$
\xymatrix{
\Spec(K) \ar[d] \ar[r]_x & X_y \ar[d] \\
\Spec(k) \ar[r]^t & T
}
$$
soit commutatif. D'après le lemme de $2$-Yoneda, cela signifie exactement qu'il faut trouver
$k \subset K$ et un objet $x$ de $\mathcal{X}$ au-dessus de $K$ tels que
$p(x) \cong t^*y|_{\Spec(K)}$. Ainsi, (1) garantit bien cette
existence, ce qui conclut.
\end{proof}












\section{Morphismes algébriques}
\label{criteria-section-algebraic}

\noindent
La notion suivante est parfois utile.

\begin{definition}
\label{criteria-definition-algebraic}
Soit $S$ un schéma. Soit $F : \mathcal{X} \to \mathcal{Y}$ un
$1$-morphisme de champs en groupoïdes sur $(\Sch/S)_{fppf}$.
On dit que $F$ est {\it algébrique} si, pour tout schéma $T$ et tout
objet $\xi$ de $\mathcal{Y}$ au-dessus de $T$, le $2$-produit fibré
$$
(\Sch/T)_{fppf} \times_{\xi, \mathcal{Y}} \mathcal{X}
$$
est un champ algébrique sur $S$.
\end{definition}

\noindent
Avec cette terminologie, on obtient le résultat suivant, qui généralise
Champs algébriques, lemme
\ref{algebraic-lemma-representable-morphism-to-algebraic}.

\begin{lemma}
\label{criteria-lemma-algebraic-morphism-to-algebraic}
Soit $S$ un schéma.
Soit $F : \mathcal{X} \to \mathcal{Y}$ un $1$-morphisme de
champs en groupoïdes sur $(\Sch/S)_{fppf}$. Si
\begin{enumerate}
\item $\mathcal{Y}$ est un champ algébrique, et
\item $F$ est algébrique (au sens ci-dessus),
\end{enumerate}
alors $\mathcal{X}$ est un champ algébrique.
\end{lemma}

\begin{proof}
D'après l'hypothèse (1), il existe un schéma $T$ et un objet
$\xi$ de $\mathcal{Y}$ au-dessus de $T$ tels que le
$1$-morphisme correspondant $\xi : (\Sch/T)_{fppf} \to \mathcal{Y}$
soit lisse et surjectif. Alors
$\mathcal{U} = (\Sch/T)_{fppf} \times_{\xi, \mathcal{Y}} \mathcal{X}$
est un champ algébrique par l'hypothèse (2).
Choisissons un schéma $U$ et un $1$-morphisme lisse surjectif
$(\Sch/U)_{fppf} \to \mathcal{U}$.
La projection $\mathcal{U} \longrightarrow \mathcal{X}$,
qui est le changement de base du morphisme
$\xi : (\Sch/T)_{fppf} \to \mathcal{Y}$,
est surjective et lisse ; voir
Champs algébriques, lemme
\ref{algebraic-lemma-base-change-representable-transformations-property}.
La composée
$(\Sch/U)_{fppf} \to \mathcal{U} \to \mathcal{X}$
est alors surjective et lisse comme composée de
morphismes surjectifs et lisses ; voir
Champs algébriques, lemme
\ref{algebraic-lemma-composition-representable-transformations-property}.
Ainsi, $\mathcal{X}$ est un champ algébrique d'après
Champs algébriques, lemme
\ref{algebraic-lemma-smooth-surjective-morphism-implies-algebraic}.
\end{proof}

\begin{lemma}
\label{criteria-lemma-map-from-algebraic}
Soit $S$ un schéma. Soit $F : \mathcal{X} \to \mathcal{Y}$ un $1$-morphisme
de champs en groupoïdes sur $(\Sch/S)_{fppf}$. Si $\mathcal{X}$ est un
champ algébrique et si $\Delta : \mathcal{Y} \to \mathcal{Y} \times \mathcal{Y}$
est représentable par des espaces algébriques, alors $F$ est algébrique.
\end{lemma}

\begin{proof}
Choisissons un champ en groupoïdes représentable $\mathcal{U}$ et un
$1$-morphisme lisse surjectif $\mathcal{U} \to \mathcal{X}$. Soit $T$ un schéma et
soit $\xi$ un objet de $\mathcal{Y}$ au-dessus de $T$. Le morphisme de
$2$-produits fibrés
$$
(\Sch/T)_{fppf} \times_{\xi, \mathcal{Y}} \mathcal{U}
\longrightarrow
(\Sch/T)_{fppf} \times_{\xi, \mathcal{Y}} \mathcal{X}
$$
est représentable par des espaces algébriques, surjectif et lisse comme
changement de base de $\mathcal{U} \to \mathcal{X}$ ; voir
Champs algébriques,
lemmes \ref{algebraic-lemma-base-change-representable-by-spaces} et
\ref{algebraic-lemma-base-change-representable-transformations-property}.
La condition imposée à la diagonale de $\mathcal{Y}$ montre que
la source de ce morphisme est représentable par un espace algébrique ; voir
Champs algébriques, lemme \ref{algebraic-lemma-representable-diagonal}.
La cible est donc un champ algébrique d'après
Champs algébriques,
lemme \ref{algebraic-lemma-smooth-surjective-morphism-implies-algebraic}.
\end{proof}

\begin{lemma}
\label{criteria-lemma-diagonals-and-algebraic-morphisms}
Soit $S$ un schéma. Soit $F : \mathcal{X} \to \mathcal{Y}$ un $1$-morphisme
de champs en groupoïdes sur $(\Sch/S)_{fppf}$.
Si $F$ est algébrique et si
$\Delta : \mathcal{Y} \to \mathcal{Y} \times \mathcal{Y}$
est représentable par des espaces algébriques, alors
$\Delta : \mathcal{X} \to \mathcal{X} \times \mathcal{X}$
est représentable par des espaces algébriques.
\end{lemma}

\begin{proof}
Supposons $F$ algébrique et
$\Delta : \mathcal{Y} \to \mathcal{Y} \times \mathcal{Y}$
représentable par des espaces algébriques.
Prenons un schéma $U$ sur $S$ et deux objets $x_1, x_2$ de
$\mathcal{X}$ au-dessus de $U$.
Il faut montrer que $\mathit{Isom}(x_1, x_2)$ est un espace algébrique
sur $U$ ; voir
Champs algébriques, lemme \ref{algebraic-lemma-representable-diagonal}.
Posons $y_i = F(x_i)$. On dispose d'un morphisme de faisceaux d'ensembles
$$
f : \mathit{Isom}(x_1, x_2) \to \mathit{Isom}(y_1, y_2)
$$
dont la cible est un espace algébrique par hypothèse.
Il suffit donc de montrer que $f$ est représentable par des
espaces algébriques ; voir Amorçage, lemme
\ref{bootstrap-lemma-representable-by-spaces-over-space}.
On peut donc choisir un schéma $V$ sur $U$ et un
isomorphisme $\beta : y_{1, V} \to y_{2, V}$ ;
il faut montrer que le foncteur
$$
(\Sch/V)_{fppf} \to \textit{Ensembles},\quad
T/V \mapsto \{\alpha : x_{1, T} \to x_{2, T}
\text{ dans }\mathcal{X}_T \mid F(\alpha) = \beta|_T\}
$$
est un espace algébrique. Considérons les objets
$z_1 = (V, x_{1, V}, \text{id})$ et
$z_2 = (V, x_{2, V}, \beta)$ de
$$
\mathcal{Z} = (\Sch/V)_{fppf} \times_{y_{1, V}, \mathcal{Y}} \mathcal{X}
$$
On vérifie immédiatement que
le foncteur ci-dessus est égal à $\mathit{Isom}(z_1, z_2)$
sur $(\Sch/V)_{fppf}$. C'est donc un espace algébrique
par l'hypothèse que $F$ est algébrique (et par la définition
des champs algébriques).
\end{proof}
















\section{Espaces de sections}
\label{criteria-section-spaces-sections}

\noindent
Étant donnés des morphismes $W \to Z \to U$, on peut considérer le foncteur qui associe
à un schéma $U'$ sur $U$ l'ensemble des sections $\sigma : Z_{U'} \to W_{U'}$
du changement de base $W_{U'} \to Z_{U'}$ du morphisme $W \to Z$.
Dans cette section, nous démontrons quelques lemmes préliminaires sur ce foncteur.

\begin{lemma}
\label{criteria-lemma-surjection-space-of-sections}
Soit $Z \to U$ un morphisme fini de schémas.
Soit $W$ un espace algébrique et soit $W \to Z$ un
morphisme étale surjectif. Il existe alors un morphisme
étale surjectif $U' \to U$ et une section
$$
\sigma : Z_{U'} \to W_{U'}
$$
du morphisme $W_{U'} \to Z_{U'}$.
\end{lemma}

\begin{proof}
On peut choisir un schéma séparé $W'$ et un morphisme étale surjectif
$W' \to W$. Après avoir remplacé $W$ par $W'$, on peut donc supposer que $W$
est un schéma séparé. Écrivons $f : W \to Z$ et $\pi : Z \to U$.
Notons que $f \circ \pi : W \to U$ est séparé puisque
$W$ est séparé (voir
Schémas, lemme \ref{schemes-lemma-compose-after-separated}).
Soit $u \in U$ un point. Il suffit évidemment
de trouver un voisinage étale $(U', u')$ de $(U, u)$ tel qu'une
section $\sigma$ existe sur $U'$. Soient $z_1, \ldots, z_r$
les points de $Z$ situés au-dessus de $u$. Pour chaque $i$, choisissons un point
$w_i \in W$ d'image $z_i$. On peut choisir un voisinage étale
$(U', u') \to (U, u)$ tel que les conclusions de
Compléments sur les morphismes, lemme
\ref{more-morphisms-lemma-etale-splits-off-quasi-finite-part-technical-variant}
soient satisfaites à la fois pour $Z \to U$ et les points $z_1, \ldots, z_r$,
et pour $W \to U$ et les points $w_1, \ldots, w_r$. Après avoir
remplacé $(U, u)$ par $(U', u')$ et renommé les objets, on peut donc supposer que
toutes les extensions de corps $\kappa(z_i)/\kappa(u)$ et
$\kappa(w_i)/\kappa(u)$ sont radicielles et, de plus,
qu'il existe des décompositions en sommes disjointes
$$
Z = V_1 \amalg \ldots \amalg V_r \amalg A, \quad
W = W_1 \amalg \ldots \amalg W_r \amalg B
$$
en sous-schémas ouverts et fermés,
avec $z_i \in V_i$, $w_i \in W_i$, et $V_i \to U$, $W_i \to U$ finis.
Après avoir remplacé $U$ par $U \setminus \pi(A)$, on peut supposer que
$A = \emptyset$, c'est-à-dire $Z = V_1 \amalg \ldots \amalg V_r$.
Après avoir remplacé $W_i$ par $W_i \cap f^{-1}(V_i)$ et
$B$ par $B \cup \bigcup W_i \cap f^{-1}(Z \setminus V_i)$,
on peut supposer que $f$ envoie $W_i$ dans $V_i$.
Alors $f_i = f|_{W_i} : W_i \to V_i$ est un morphisme entre schémas finis sur $U$,
donc il est fini (voir
Morphismes, lemme \ref{morphisms-lemma-finite-permanence}).
Il est aussi étale (par hypothèse),
$f_i^{-1}(\{z_i\}) = w_i$, et induit un isomorphisme des corps
résiduels $\kappa(z_i) = \kappa(w_i)$ (car les deux sont des extensions radicielles
de $\kappa(u)$, tandis que $\kappa(w_i)/\kappa(z_i)$
est séparable puisque $f$ est étale). D'après
Morphismes étales, lemme \ref{etale-lemma-finite-etale-one-point}
le morphisme $f_i$ est donc un isomorphisme dans un voisinage $V_i'$ de
$z_i$. Comme $\pi : Z \to U$ est fermé, quitte à rétrécir $U$, on peut supposer
que $W_i \to V_i$ est un isomorphisme. Cela démontre le lemme.
\end{proof}

\begin{lemma}
\label{criteria-lemma-space-of-sections}
Soit $Z \to U$ un morphisme fini localement libre de schémas.
Soit $W$ un espace algébrique et soit $W \to Z$ un morphisme étale.
Alors le foncteur
$$
F : (\Sch/U)_{fppf}^{opp} \longrightarrow \textit{Ensembles},
$$
défini par la règle
$$
U' \longmapsto
F(U') =
\{\sigma : Z_{U'} \to W_{U'}\text{ section de }W_{U'} \to Z_{U'}\}
$$
est un espace algébrique et le morphisme $F \to U$ est étale.
\end{lemma}

\begin{proof}
Supposons d'abord que $W \to Z$ soit aussi séparé.
Soit $U'$ un schéma sur $U$ et soit $\sigma \in F(U')$. D'après
Morphismes d'espaces algébriques, lemme \ref{spaces-morphisms-lemma-section-immersion},
le morphisme $\sigma$ est une immersion fermée.
De plus, $\sigma$ est étale d'après
Propriétés des espaces algébriques, lemme \ref{spaces-properties-lemma-etale-permanence}.
Ainsi, $\sigma$ est aussi une immersion ouverte ; voir
Morphismes d'espaces algébriques,
lemme \ref{spaces-morphisms-lemma-etale-universally-injective-open}.
Autrement dit, $Z_\sigma = \sigma(Z_{U'}) \subset W_{U'}$ est
un sous-espace ouvert tel que le morphisme $Z_\sigma \to Z_{U'}$
soit un isomorphisme. En particulier, le morphisme $Z_\sigma \to U'$
est fini. On obtient donc une transformation de foncteurs
$$
F \longrightarrow (W/U)_{fin}, \quad
\sigma \longmapsto (U' \to U, Z_\sigma)
$$
où $(W/U)_{fin}$ est la partie finie du morphisme $W \to U$
introduite dans
Compléments sur les groupoïdes en espaces algébriques, section
\ref{spaces-more-groupoids-section-finite}.
Il est clair que cette transformation de foncteurs est injective (puisqu'on peut
retrouver $\sigma$ à partir de $Z_\sigma$ comme inverse de l'isomorphisme
$Z_\sigma \to Z_{U'}$). D'après
Compléments sur les groupoïdes en espaces algébriques, proposition
\ref{spaces-more-groupoids-proposition-finite-algebraic-space}
nous savons que $(W/U)_{fin}$ est un espace algébrique étale sur $U$.
Pour achever la preuve dans ce cas, il suffit donc de montrer que
$F \to (W/U)_{fin}$ est représentable et est une immersion ouverte.
Pour cela, supposons donnés un morphisme de schémas $U' \to U$
et un sous-espace ouvert $Z' \subset W_{U'}$ tels que $Z' \to U'$
soit fini. Il suffit alors de montrer qu'il existe un
sous-schéma ouvert $U'' \subset U'$ tel qu'un morphisme
$T \to U'$ se factorise par $U''$ si et seulement si $Z' \times_{U'} T$
s'envoie isomorphiquement sur $Z \times_{U'} T$. Cela résulte de
Compléments sur les morphismes d'espaces algébriques, lemme
\ref{spaces-more-morphisms-lemma-where-isomorphism}
(on utilise ici que $Z \to B$ est plat et localement de présentation finie,
ainsi que fini).
Le lemme est donc démontré lorsque $W \to Z$ est séparé
et étale.

\medskip\noindent
Dans le cas général, choisissons un schéma séparé $W'$ et un morphisme
étale surjectif $W' \to W$. Notons que les morphismes $W' \to W$ et
$W \to Z$ sont séparés puisque leur source est séparée. Notons $F'$ le
foncteur associé à $W' \to Z \to U$ comme dans le lemme. Dans le premier
paragraphe de la preuve, nous avons montré que $F'$ est représentable par un
espace algébrique étale sur $U$. D'après le
lemme \ref{criteria-lemma-surjection-space-of-sections},
l'application de foncteurs $F' \to F$ est surjective pour la topologie étale
sur $\Sch/U$. De plus, si $U'$ et $\sigma : Z_{U'} \to W_{U'}$
définissent un élément $\xi \in F(U')$, alors le produit fibré
$$
F'' = F' \times_{F, \xi} U'
$$
est le foncteur sur $\Sch/U'$ associé aux morphismes
$$
W'_{U'} \times_{W_{U'}, \sigma} Z_{U'} \to Z_{U'} \to U'.
$$
Comme le premier morphisme est séparé, étant le changement de base d'un morphisme
séparé, le résultat du premier paragraphe montre que $F''$ est un espace algébrique étale sur $U'$.
Il s'ensuit que $F' \to F$ est une transformation de foncteurs
étale surjective, représentable par des espaces algébriques. Ainsi,
$F$ est un espace algébrique d'après
Amorçage, théorème \ref{bootstrap-theorem-final-bootstrap}.
Puisque $F' \to F$ est un morphisme étale surjectif d'espaces algébriques,
$F \to U$ est étale puisque $F' \to U$ l'est.
\end{proof}









\section{Morphismes relatifs}
\label{criteria-section-relative-morphisms}

\noindent
Nous poursuivons la discussion commencée dans
Compléments sur les morphismes, section \ref{more-morphisms-section-relative-morphisms}.

\medskip\noindent
Soit $S$ un schéma. Soient $Z \to B$ et $X \to B$ des morphismes
d'espaces algébriques sur $S$. Pour tout schéma $T$, on peut considérer les couples
$(a, b)$, où $a : T \to B$
est un morphisme et $b : T \times_{a, B} Z \to T \times_{a, B} X$
un morphisme sur $T$. On a le diagramme
\begin{equation}
\label{criteria-equation-hom}
\vcenter{
\xymatrix{
T \times_{a, B} Z \ar[rd] \ar[rr]_b & &
T \times_{a, B} X \ar[ld] & Z \ar[rd] & & X \ar[ld] \\
& T \ar[rrr]^a & & & B
}
}
\end{equation}
Bien entendu, on peut aussi voir $b$ comme un morphisme
$b : T \times_{a, B} Z \to X$ tel que
$$
\xymatrix{
T \times_{a, B} Z \ar[r] \ar[d] \ar@/^1pc/[rrr]_-b &
Z \ar[rd] & & X \ar[ld] \\
T \ar[rr]^a & & B
}
$$
soit commutatif. Dans cette situation, on peut définir un foncteur
\begin{equation}
\label{criteria-equation-hom-functor}
\mathit{Mor}_B(Z, X) : (\Sch/S)^{opp} \longrightarrow \textit{Ensembles},
\quad
T \longmapsto \{(a, b)\text{ comme ci-dessus}\}
\end{equation}
Parfois, on le considère comme un foncteur défini sur la catégorie
des schémas sur $B$ ; dans ce cas, on omet $a$ de la notation.

\begin{lemma}
\label{criteria-lemma-hom-functor-sheaf}
Soit $S$ un schéma. Soient $Z \to B$ et $X \to B$ des morphismes
d'espaces algébriques sur $S$. Alors
\begin{enumerate}
\item $\mathit{Mor}_B(Z, X)$ est un faisceau sur
$(\Sch/S)_{fppf}$.
\item Si $T$ est un espace algébrique sur $S$, il existe une
bijection canonique
$$
\Mor_{\Sh((\Sch/S)_{fppf})}(T, \mathit{Mor}_B(Z, X))
=
\{(a, b)\text{ comme dans }(\ref{criteria-equation-hom})\}
$$
\end{enumerate}
\end{lemma}

\begin{proof}
Soit $T$ un espace algébrique sur $S$. Soit $\{T_i \to T\}$ un
recouvrement fppf de $T$ (comme dans
Topologies on Espaces, section \ref{spaces-topologies-section-fppf}).
Supposons que $(a_i, b_i) \in \mathit{Mor}_B(Z, X)(T_i)$ et
que $(a_i, b_i)|_{T_i \times_T T_j} = (a_j, b_j)|_{T_i \times_T T_j}$
pour tous $i, j$. D'après
Descente et espaces algébriques,
lemme \ref{spaces-descent-lemma-fpqc-universal-effective-epimorphisms},
il existe un unique morphisme $a : T \to B$ tel que $a_i$ soit la
composée de $T_i \to T$ et de $a$. Alors
$\{T_i \times_{a_i, B} Z \to T \times_{a, B} Z\}$ est aussi un recouvrement
fppf, et le même lemme implique qu'il existe un unique morphisme
$b : T \times_{a, B} Z \to T \times_{a, B} X$ tel que $b_i$ soit la
composée de $T_i \times_{a_i, B} Z \to T \times_{a, B} Z$ et de $b$. Ainsi,
$(a, b) \in \mathit{Mor}_B(Z, X)(T)$ se restreint à $(a_i, b_i)$
sur $T_i$ pour tout $i$.

\medskip\noindent
Notons que le résultat du paragraphe précédent implique en particulier (1).

\medskip\noindent
Soit $T$ un espace algébrique sur $S$. Pour démontrer (2), nous allons
construire des applications réciproques entre les ensembles affichés. Dans la
suite, le mot « couple » désigne un couple $(a, b)$ s'insérant
dans (\ref{criteria-equation-hom}).

\medskip\noindent
Soit $v : T \to \mathit{Mor}_B(Z, X)$ une transformation naturelle.
Choisissons un schéma $U$ et un morphisme étale surjectif $p : U \to T$.
Alors $v(p) \in \mathit{Mor}_B(Z, X)(U)$ correspond à un couple $(a_U, b_U)$
sur $U$. Soit $R = U \times_T U$, de projections $t, s : R \to U$.
Comme $v$ est une transformation de foncteurs, les images inverses de
$(a_U, b_U)$ par $s$ et $t$ coïncident. Puisque $\{U \to T\}$ est un
recouvrement fppf, le résultat du premier paragraphe permet donc de
déduire qu'il existe un unique couple $(a, b)$ sur $T$.

\medskip\noindent
Réciproquement, soit $(a, b)$ un couple sur $T$.
Soient $U \to T$, $R = U \times_T U$ et $t, s : R \to U$ comme
ci-dessus. La restriction $(a, b)|_U$ donne alors naissance à une
transformation de foncteurs $v : h_U \to \mathit{Mor}_B(Z, X)$ par le
lemme de Yoneda
(Catégories, lemme \ref{categories-lemma-yoneda}).
Comme les deux images inverses $s^*(a, b)|_U$ et $t^*(a, b)|_U$
sont égales, $v$ coégalise les deux applications
$h_t, h_s : h_R \to h_U$. Puisque $T = U/R$ est le faisceau quotient fppf d'après
Espaces, lemme \ref{spaces-lemma-space-presentation},
et puisque $\mathit{Mor}_B(Z, X)$ est un faisceau fppf d'après (1), on conclut
que $v$ se factorise par une application $T \to \mathit{Mor}_B(Z, X)$.

\medskip\noindent
Nous omettons de vérifier que les deux constructions ci-dessus sont
réciproques.
\end{proof}

\begin{lemma}
\label{criteria-lemma-base-change-hom-functor}
Soit $S$ un schéma. Soient $Z \to B$, $X \to B$ et $B' \to B$
des morphismes d'espaces algébriques sur $S$. Posons $Z' = B' \times_B Z$
et $X' = B' \times_B X$. Alors
$$
\mathit{Mor}_{B'}(Z', X')
=
B' \times_B \mathit{Mor}_B(Z, X)
$$
dans $\Sh((\Sch/S)_{fppf})$.
\end{lemma}

\begin{proof}
L'égalité des foncteurs résulte immédiatement des définitions.
L'égalité des faisceaux s'en déduit, car les deux membres sont des
faisceaux d'après le
lemme \ref{criteria-lemma-hom-functor-sheaf},
et parce qu'un produit fibré de faisceaux est égal au
produit fibré correspondant des préfaisceaux (c'est-à-dire des foncteurs).
\end{proof}

\begin{lemma}
\label{criteria-lemma-etale-covering-hom-functor}
Soit $S$ un schéma. Soient $Z \to B$ et $X' \to X \to B$ des morphismes
d'espaces algébriques sur $S$. Supposons que
\begin{enumerate}
\item $X' \to X$ soit étale, et
\item $Z \to B$ soit fini localement libre.
\end{enumerate}
Alors $\mathit{Mor}_B(Z, X') \to \mathit{Mor}_B(Z, X)$ est représentable
par des espaces algébriques et étale. Si $X' \to X$ est aussi surjectif,
alors $\mathit{Mor}_B(Z, X') \to \mathit{Mor}_B(Z, X)$ est surjectif.
\end{lemma}

\begin{proof}
Soit $U$ un schéma et soit $\xi = (a, b)$ un élément de
$\mathit{Mor}_B(Z, X)(U)$. Il faut montrer que le foncteur
$$
h_U \times_{\xi, \mathit{Mor}_B(Z, X)} \mathit{Mor}_B(Z, X')
$$
est représentable par un espace algébrique étale sur $U$. Posons
$Z_U = U \times_{a, B} Z$ et $W = Z_U \times_{b, X} X'$.
Alors $W \to Z_U \to U$ est comme dans le
lemme \ref{criteria-lemma-space-of-sections},
et le faisceau $F$ qui y est défini s'identifie au produit fibré
affiché ci-dessus. D'où la première assertion du lemme.
La seconde assertion résulte de celle-ci et du
lemme \ref{criteria-lemma-surjection-space-of-sections},
qui garantit que $F \to U$ est surjectif dans la situation ci-dessus.
\end{proof}

\begin{proposition}
\label{criteria-proposition-hom-functor-algebraic-space}
Soit $S$ un schéma. Soient $Z \to B$ et $X \to B$ des morphismes
d'espaces algébriques sur $S$. Si $Z \to B$ est fini localement libre,
alors $\mathit{Mor}_B(Z, X)$ est un espace algébrique.
\end{proposition}

\begin{proof}
Choisissons un schéma $B' = \coprod B'_i$, somme disjointe de
schémas affines $B'_i$, et un morphisme étale surjectif $B' \to B$.
On peut aussi supposer que $B'_i \times_B Z$ est le spectre d'un anneau
qui est un $\Gamma(B'_i, \mathcal{O}_{B'_i})$-module libre de type fini.
D'après le
lemme \ref{criteria-lemma-base-change-hom-functor}
et
Espaces, lemme
\ref{spaces-lemma-base-change-representable-transformations-property}
le morphisme $\mathit{Mor}_{B'}(Z', X') \to \mathit{Mor}_B(Z, X)$
est étale surjectif. D'après
Amorçage, théorème \ref{bootstrap-theorem-final-bootstrap},
il suffit donc de démontrer la proposition lorsque $B = B'$ est une somme disjointe de
schémas affines $B'_i$ telle que chaque $B'_i \times_B Z$ soit fini et libre
sur $B'_i$. Il suffit alors de démontrer le résultat après restriction
à chaque $B'_i$. On peut donc supposer que $B$ est affine et que
$\Gamma(Z, \mathcal{O}_Z)$ est un $\Gamma(B, \mathcal{O}_B)$-module libre de type fini.

\medskip\noindent
Choisissons un schéma $X'$ qui soit une somme disjointe de schémas affines et
un morphisme étale surjectif $X' \to X$. D'après le
lemme \ref{criteria-lemma-etale-covering-hom-functor},
le morphisme $\mathit{Mor}_B(Z, X') \to \mathit{Mor}_B(Z, X)$
est représentable par des espaces algébriques, étale et surjectif.
D'après
Amorçage, théorème \ref{bootstrap-theorem-final-bootstrap},
il suffit donc de démontrer la proposition lorsque $X$ est une somme disjointe
de schémas affines. On se ramène ainsi au cas traité dans le paragraphe
suivant.

\medskip\noindent
Supposons que $X = \coprod_{i \in I} X_i$ soit une somme disjointe de schémas
affines, que $B$ soit affine et que $\Gamma(Z, \mathcal{O}_Z)$ soit un
$\Gamma(B, \mathcal{O}_B)$-module libre de type fini. Pour toute partie finie
$E \subset I$, posons
$$
F_E = \mathit{Mor}_B(Z, \coprod\nolimits_{i \in E} X_i).
$$
D'après Compléments sur les morphismes,
lemme \ref{more-morphisms-lemma-hom-from-finite-free-into-affine},
$F_E$ est un espace algébrique. Considérons le morphisme
$$
\coprod\nolimits_{E \subset I\text{ fini}} F_E
\longrightarrow
\mathit{Mor}_B(Z, X)
$$
Chacun des morphismes
$F_E \to \mathit{Mor}_B(Z, X)$ est une immersion ouverte, car son image est
simplement le lieu paramétrant les couples $(a, b)$ pour lesquels $b$ se factorise par
le sous-schéma ouvert $\coprod\nolimits_{i \in E} X_i$ de $X$. De plus,
si $T$ est quasi-compact, alors, pour tout couple $(a, b)$, l'image
de $b$ est contenue dans $\coprod\nolimits_{i \in E} X_i$ pour une certaine partie
finie $E \subset I$. La flèche affichée est donc en fait un
recouvrement ouvert, ce qui conclut\footnote{Sous réserve de
quelques arguments ensemblistes. Il faut en effet montrer que
$\coprod F_E$ est un espace algébrique. Cela résulte des inégalités
$|I| \leq \text{size}(X)$ et $\text{size}(F_E) \leq \text{size}(X)$,
la seconde découlant de la description explicite de $F_E$ dans la preuve de
Compléments sur les morphismes,
lemme \ref{more-morphisms-lemma-hom-from-finite-free-into-affine}.
Nous omettons certains détails.}, d'après
Espaces, lemme \ref{spaces-lemma-glueing-algebraic-spaces}.
\end{proof}










\section{Restriction des scalaires}
\label{criteria-section-restriction-of-scalars}

\noindent
Supposons que $X \to Z \to B$ soient des morphismes d'espaces algébriques sur $S$.
Pour tout schéma $T$, on peut considérer les couples $(a, b)$, où $a : T \to B$
est un morphisme et $b : T \times_{a, B} Z \to X$ un morphisme sur $Z$.
On a le diagramme
\begin{equation}
\label{criteria-equation-pairs}
\vcenter{
\xymatrix{
& X \ar[d] \\
T \times_{a, B} Z \ar[d] \ar[ru]^b \ar[r] & Z \ar[d] \\
T \ar[r]^a & B
}
}
\end{equation}
Dans cette situation, on peut définir un
foncteur
\begin{equation}
\label{criteria-equation-restriction-of-scalars}
\text{Res}_{Z/B}(X) : (\Sch/S)^{opp} \longrightarrow \textit{Ensembles},
\quad
T \longmapsto \{(a, b)\text{ comme ci-dessus}\}
\end{equation}
Parfois, on le considère comme un foncteur défini sur la catégorie
des schémas sur $B$ ; dans ce cas, on omet $a$ de la notation.

\begin{lemma}
\label{criteria-lemma-restriction-of-scalars-sheaf}
Soit $S$ un schéma. Soient $X \to Z \to B$ des morphismes
d'espaces algébriques sur $S$. Alors
\begin{enumerate}
\item $\text{Res}_{Z/B}(X)$ est un faisceau sur
$(\Sch/S)_{fppf}$.
\item Si $T$ est un espace algébrique sur $S$, il existe une
bijection canonique
$$
\Mor_{\Sh((\Sch/S)_{fppf})}(T, \text{Res}_{Z/B}(X))
=
\{(a, b)\text{ comme dans }(\ref{criteria-equation-pairs})\}
$$
\end{enumerate}
\end{lemma}

\begin{proof}
Soit $T$ un espace algébrique sur $S$. Soit $\{T_i \to T\}$ un
recouvrement fppf de $T$ (comme dans
Topologies on Espaces, section \ref{spaces-topologies-section-fppf}).
Supposons que $(a_i, b_i) \in \text{Res}_{Z/B}(X)(T_i)$ et
que $(a_i, b_i)|_{T_i \times_T T_j} = (a_j, b_j)|_{T_i \times_T T_j}$
pour tous $i, j$. D'après
Descente et espaces algébriques,
lemme \ref{spaces-descent-lemma-fpqc-universal-effective-epimorphisms},
il existe un unique morphisme $a : T \to B$ tel que $a_i$ soit la
composée de $T_i \to T$ et de $a$. Alors
$\{T_i \times_{a_i, B} Z \to T \times_{a, B} Z\}$ est aussi un recouvrement
fppf, et le même lemme implique qu'il existe un unique morphisme
$b : T \times_{a, B} Z \to X$ tel que $b_i$ soit la composée
de $T_i \times_{a_i, B} Z \to T \times_{a, B} Z$ et de $b$. Ainsi,
$(a, b) \in \text{Res}_{Z/B}(X)(T)$ se restreint à $(a_i, b_i)$
sur $T_i$ pour tout $i$.

\medskip\noindent
Notons que le résultat du paragraphe précédent implique en particulier (1).

\medskip\noindent
Soit $T$ un espace algébrique sur $S$. Pour démontrer (2), nous allons
construire des applications réciproques entre les ensembles affichés. Dans la
suite, le mot « couple » désigne un couple $(a, b)$ s'insérant
dans (\ref{criteria-equation-pairs}).

\medskip\noindent
Soit $v : T \to \text{Res}_{Z/B}(X)$ une transformation naturelle.
Choisissons un schéma $U$ et un morphisme étale surjectif $p : U \to T$.
Alors $v(p) \in \text{Res}_{Z/B}(X)(U)$ correspond à un couple $(a_U, b_U)$
sur $U$. Soit $R = U \times_T U$, de projections $t, s : R \to U$.
Comme $v$ est une transformation de foncteurs, les images inverses de
$(a_U, b_U)$ par $s$ et $t$ coïncident. Puisque $\{U \to T\}$ est un
recouvrement fppf, le résultat du premier paragraphe permet donc de
déduire qu'il existe un unique couple $(a, b)$ sur $T$.

\medskip\noindent
Réciproquement, soit $(a, b)$ un couple sur $T$.
Soient $U \to T$, $R = U \times_T U$ et $t, s : R \to U$ comme
ci-dessus. La restriction $(a, b)|_U$ donne alors naissance à une
transformation de foncteurs $v : h_U \to \text{Res}_{Z/B}(X)$ par le
lemme de Yoneda
(Catégories, lemme \ref{categories-lemma-yoneda}).
Comme les deux images inverses $s^*(a, b)|_U$ et $t^*(a, b)|_U$
sont égales, $v$ coégalise les deux applications
$h_t, h_s : h_R \to h_U$. Puisque $T = U/R$ est le faisceau quotient fppf d'après
Espaces, lemme \ref{spaces-lemma-space-presentation},
et puisque $\text{Res}_{Z/B}(X)$ est un faisceau fppf d'après (1), on conclut
que $v$ se factorise par une application $T \to \text{Res}_{Z/B}(X)$.

\medskip\noindent
Nous omettons de vérifier que les deux constructions ci-dessus sont
réciproques.
\end{proof}

\noindent
Bien entendu, le faisceau $\text{Res}_{Z/B}(X)$ est muni d'une transformation
naturelle de foncteurs $\text{Res}_{Z/B}(X) \to B$. Nous l'utiliserons sans autre
mention dans la suite.

\begin{lemma}
\label{criteria-lemma-etale-base-change-restriction-of-scalars}
Soit $S$ un schéma. Soient $X \to Z \to B$ et $B' \to B$
des morphismes d'espaces algébriques sur $S$.
Posons $Z' = B' \times_B Z$ et $X' = B' \times_B X$. Alors
$$
\text{Res}_{Z'/B'}(X')
=
B' \times_B \text{Res}_{Z/B}(X)
$$
dans $\Sh((\Sch/S)_{fppf})$.
\end{lemma}

\begin{proof}
L'égalité des foncteurs résulte immédiatement des définitions.
L'égalité des faisceaux s'en déduit, car les deux membres sont des
faisceaux d'après le
lemme \ref{criteria-lemma-restriction-of-scalars-sheaf},
et parce qu'un produit fibré de faisceaux est égal au
produit fibré correspondant des préfaisceaux (c'est-à-dire des foncteurs).
\end{proof}

\begin{lemma}
\label{criteria-lemma-etale-covering-restriction-of-scalars}
Soit $S$ un schéma. Soient $X' \to X \to Z \to B$ des morphismes
d'espaces algébriques sur $S$. Supposons que
\begin{enumerate}
\item $X' \to X$ soit étale, et
\item $Z \to B$ soit fini localement libre.
\end{enumerate}
Alors $\text{Res}_{Z/B}(X') \to \text{Res}_{Z/B}(X)$ est représentable
par des espaces algébriques et étale. Si $X' \to X$ est aussi surjectif,
alors $\text{Res}_{Z/B}(X') \to \text{Res}_{Z/B}(X)$ est surjectif.
\end{lemma}

\begin{proof}
Soit $U$ un schéma et soit $\xi = (a, b)$ un élément de
$\text{Res}_{Z/B}(X)(U)$. Il faut montrer que le foncteur
$$
h_U \times_{\xi, \text{Res}_{Z/B}(X)} \text{Res}_{Z/B}(X')
$$
est représentable par un espace algébrique étale sur $U$. Posons
$Z_U = U \times_{a, B} Z$ et $W = Z_U \times_{b, X} X'$.
Alors $W \to Z_U \to U$ est comme dans le
lemme \ref{criteria-lemma-space-of-sections},
et le faisceau $F$ qui y est défini s'identifie au produit fibré
affiché ci-dessus. D'où la première assertion du lemme.
La seconde assertion résulte de celle-ci et du
lemme \ref{criteria-lemma-surjection-space-of-sections},
qui garantit que $F \to U$ est surjectif dans la situation ci-dessus.
\end{proof}

\noindent
On peut maintenant utiliser les lemmes ci-dessus pour montrer que $\text{Res}_{Z/B}(X)$
est un espace algébrique dès que $Z \to B$ est fini localement libre, d'une manière
presque identique à la preuve que $\mathit{Mor}_B(Z, X)$ est un
espace algébrique ; voir la
proposition \ref{criteria-proposition-hom-functor-algebraic-space}.
Nous déduirons plutôt directement ce résultat du lemme suivant
et du fait que $\mathit{Mor}_B(Z, X)$ est un espace algébrique.

\begin{lemma}
\label{criteria-lemma-fibre-diagram}
Soit $S$ un schéma. Soient $X \to Z \to B$ des morphismes
d'espaces algébriques sur $S$. Le diagramme suivant
$$
\xymatrix{
\mathit{Mor}_B(Z, X) \ar[r] & \mathit{Mor}_B(Z, Z) \\
\text{Res}_{Z/B}(X) \ar[r] \ar[u] & B \ar[u]_{\text{id}_Z}
}
$$
est un diagramme cartésien de faisceaux sur $(\Sch/S)_{fppf}$.
\end{lemma}

\begin{proof}
La preuve est omise. Indication : exercice sur le point de vue fonctoriel en
géométrie algébrique.
\end{proof}

\begin{proposition}
\label{criteria-proposition-restriction-of-scalars-algebraic-space}
Soit $S$ un schéma. Soient $X \to Z \to B$ des morphismes
d'espaces algébriques sur $S$. Si $Z \to B$ est fini localement libre,
alors $\text{Res}_{Z/B}(X)$ est un espace algébrique.
\end{proposition}

\begin{proof}
D'après la
proposition \ref{criteria-proposition-hom-functor-algebraic-space},
les foncteurs $\mathit{Mor}_B(Z, X)$ et $\mathit{Mor}_B(Z, Z)$
sont des espaces algébriques. L'assertion résulte donc du diagramme cartésien
du
lemme \ref{criteria-lemma-fibre-diagram}
et du fait que les produits fibrés d'espaces algébriques existent et
sont donnés par le produit fibré dans la catégorie sous-jacente des
faisceaux d'ensembles (voir
Espaces, lemme
\ref{spaces-lemma-fibre-product-spaces-over-sheaf-with-representable-diagonal}).
\end{proof}






\section{Champs de Hilbert finis}
\label{criteria-section-finite-hilbert-stacks}

\noindent
Dans cette section, nous démontrons quelques résultats sur les champs de
Hilbert finis $\mathcal{H}_d(\mathcal{X}/\mathcal{Y})$
introduits dans
Exemples de champs, section \ref{examples-stacks-section-hilbert-d-stack}.

\begin{lemma}
\label{criteria-lemma-map-hilbert}
Considérons un diagramme $2$-commutatif
$$
\xymatrix{
\mathcal{X}' \ar[r]_G \ar[d]_{F'} & \mathcal{X} \ar[d]^F \\
\mathcal{Y}' \ar[r]^H & \mathcal{Y}
}
$$
de champs en groupoïdes sur $(\Sch/S)_{fppf}$, muni d'un
$2$-isomorphisme $\gamma : H \circ F' \to F \circ G$. Dans cette situation, on
obtient un $1$-morphisme canonique
$\mathcal{H}_d(\mathcal{X}'/\mathcal{Y}') \to
\mathcal{H}_d(\mathcal{X}/\mathcal{Y})$.
Ce morphisme est compatible avec les $1$-morphismes d'oubli de
Exemples de champs,
équation (\ref{examples-stacks-equation-diagram-hilbert-d-stack}).
\end{lemma}

\begin{proof}
On envoie l'objet $(U, Z, y', x', \alpha')$ sur l'objet
$(U, Z, H(y'), G(x'), \gamma \star \text{id}_H \star \alpha')$,
où $\star$ désigne la composition horizontale des $2$-morphismes ; voir
Catégories, définition \ref{categories-definition-horizontal-composition}.
Au morphisme
$(f, g, b, a) :
(U_1, Z_1, y_1', x_1', \alpha_1') \to (U_2, Z_2, y_2', x_2', \alpha_2')$
on associe
$(f, g, H(b), G(a))$.
Nous omettons de vérifier que cela définit un foncteur entre catégories sur
$(\Sch/S)_{fppf}$.
\end{proof}

\begin{lemma}
\label{criteria-lemma-cartesian-map-hilbert}
Dans la situation du
lemme \ref{criteria-lemma-map-hilbert},
supposons que le carré donné soit $2$-cartésien. Alors le diagramme
$$
\xymatrix{
\mathcal{H}_d(\mathcal{X}'/\mathcal{Y}') \ar[r] \ar[d] &
\mathcal{H}_d(\mathcal{X}/\mathcal{Y}) \ar[d] \\
\mathcal{Y}' \ar[r] &
\mathcal{Y}
}
$$
est $2$-cartésien.
\end{lemma}

\begin{proof}
Le lemme \ref{criteria-lemma-map-hilbert} fournit un diagramme $2$-commutatif
et, par conséquent,
un $1$-morphisme (c'est-à-dire un foncteur)
$$
\mathcal{H}_d(\mathcal{X}'/\mathcal{Y}')
\longrightarrow
\mathcal{Y}' \times_\mathcal{Y} \mathcal{H}_d(\mathcal{X}/\mathcal{Y})
$$
Indiquons pourquoi ce foncteur est essentiellement surjectif. Un objet
de la catégorie du membre de droite consiste en un schéma $U$ sur $S$,
un objet $y'$ de $\mathcal{Y}'_U$, un objet $(U, Z, y, x, \alpha)$
de $\mathcal{H}_d(\mathcal{X}/\mathcal{Y})$ au-dessus de $U$ et un isomorphisme
$H(y') \to y$ dans $\mathcal{Y}_U$. L'hypothèse signifie exactement qu'il
existe un objet $x'$ de $\mathcal{X}'_Z$ muni d'isomorphismes
$G(x') \cong x$ et $\alpha' : y'|_Z \to F'(x')$ compatibles
avec $\alpha$. Alors $(U, Z, y', x', \alpha')$ est un
objet de $\mathcal{H}_d(\mathcal{X}'/\mathcal{Y}')$ au-dessus de $U$.
Nous omettons les détails.
\end{proof}

\begin{lemma}
\label{criteria-lemma-etale-covering-hilbert}
Dans la situation du
lemme \ref{criteria-lemma-map-hilbert},
supposons que
\begin{enumerate}
\item $\mathcal{Y}' = \mathcal{Y}$ et $H = \text{id}_\mathcal{Y}$,
\item $G$ soit représentable par des espaces algébriques et étale.
\end{enumerate}
Alors $\mathcal{H}_d(\mathcal{X}'/\mathcal{Y}) \to
\mathcal{H}_d(\mathcal{X}/\mathcal{Y})$ est représentable par
des espaces algébriques et étale.
Si $G$ est en outre surjectif, alors
$\mathcal{H}_d(\mathcal{X}'/\mathcal{Y}) \to
\mathcal{H}_d(\mathcal{X}/\mathcal{Y})$ est surjectif.
\end{lemma}

\begin{proof}
Soit $U$ un schéma et soit $\xi = (U, Z, y, x, \alpha)$ un objet de
$\mathcal{H}_d(\mathcal{X}/\mathcal{Y})$ au-dessus de $U$.
Il faut montrer que le $2$-produit fibré
\begin{equation}
\label{criteria-equation-to-show}
(\Sch/U)_{fppf}
\times_{\xi, \mathcal{H}_d(\mathcal{X}/\mathcal{Y})}
\mathcal{H}_d(\mathcal{X}'/\mathcal{Y})
\end{equation}
est représentable par un espace algébrique étale sur $U$.
Un objet de celui-ci au-dessus de $U'$ correspond à un objet
$x'$ de la catégorie fibre de $\mathcal{X}'$ au-dessus de $Z_{U'}$
tel que $G(x') \cong x|_{Z_{U'}}$.
Par hypothèse, le $2$-produit fibré
$$
(\Sch/Z)_{fppf} \times_{x, \mathcal{X}} \mathcal{X}'
$$
est représentable par un espace algébrique $W$ tel que la projection
$W \to Z$ soit étale. Alors (\ref{criteria-equation-to-show})
est représenté par l'espace algébrique $F$ qui paramètre les sections de
$W \to Z$ sur $U$, introduit dans le
lemme \ref{criteria-lemma-space-of-sections}.
Comme $F \to U$ est étale, on conclut que
$\mathcal{H}_d(\mathcal{X}'/\mathcal{Y}) \to
\mathcal{H}_d(\mathcal{X}/\mathcal{Y})$ est représentable par
des espaces algébriques et étale.
Enfin, si $\mathcal{X}' \to \mathcal{X}$ est aussi surjectif,
alors $W \to Z$ est surjectif, et donc $F \to U$ l'est d'après le
lemme \ref{criteria-lemma-surjection-space-of-sections}.
Ainsi, dans ce cas,
$\mathcal{H}_d(\mathcal{X}'/\mathcal{Y}) \to
\mathcal{H}_d(\mathcal{X}/\mathcal{Y})$ est aussi surjectif.
\end{proof}

\begin{lemma}
\label{criteria-lemma-etale-map-hilbert}
Dans la situation du
lemme \ref{criteria-lemma-map-hilbert},
supposons que $G$ et $H$ soient représentables par des espaces algébriques et étales.
Alors $\mathcal{H}_d(\mathcal{X}'/\mathcal{Y}') \to
\mathcal{H}_d(\mathcal{X}/\mathcal{Y})$ est représentable par
des espaces algébriques et étale.
Si, de plus, $H$ est surjectif et le foncteur induit
$\mathcal{X}' \to \mathcal{Y}' \times_\mathcal{Y} \mathcal{X}$
est surjectif, alors
$\mathcal{H}_d(\mathcal{X}'/\mathcal{Y}') \to
\mathcal{H}_d(\mathcal{X}/\mathcal{Y})$ est surjectif.
\end{lemma}

\begin{proof}
Posons $\mathcal{X}'' = \mathcal{Y}' \times_\mathcal{Y} \mathcal{X}$. D'après le
lemme \ref{criteria-lemma-etale-permanence},
le $1$-morphisme $\mathcal{X}' \to \mathcal{X}''$ est représentable par
des espaces algébriques et étale (en particulier, la condition de la seconde
assertion du lemme selon laquelle $\mathcal{X}' \to \mathcal{X}''$ est surjectif
a un sens). On obtient un diagramme $2$-commutatif
$$
\xymatrix{
\mathcal{X}' \ar[r] \ar[d] &
\mathcal{X}'' \ar[r] \ar[d] &
\mathcal{X} \ar[d] \\
\mathcal{Y}' \ar[r] &
\mathcal{Y}' \ar[r] &
\mathcal{Y}
}
$$
Il résulte du
lemme \ref{criteria-lemma-cartesian-map-hilbert}
que $\mathcal{H}_d(\mathcal{X}''/\mathcal{Y}')$ est le changement de base
de $\mathcal{H}_d(\mathcal{X}/\mathcal{Y})$ par $\mathcal{Y}' \to \mathcal{Y}$.
En particulier,
$\mathcal{H}_d(\mathcal{X}''/\mathcal{Y}') \to
\mathcal{H}_d(\mathcal{X}/\mathcal{Y})$ est
représentable par des espaces algébriques et étale ; voir
Champs algébriques, lemme
\ref{algebraic-lemma-base-change-representable-transformations-property}.
De plus, il est aussi surjectif si $H$ l'est.
Par conséquent, il suffit de montrer que
le résultat vaut pour le carré de gauche du diagramme.
On se ramène ainsi au cas où $\mathcal{Y}' = \mathcal{Y}$,
qui fait l'objet du
lemme \ref{criteria-lemma-etale-covering-hilbert}.
\end{proof}

\begin{lemma}
\label{criteria-lemma-relative-hilbert}
Soit $F : \mathcal{X} \to \mathcal{Y}$ un $1$-morphisme de champs en groupoïdes
sur $(\Sch/S)_{fppf}$. Supposons que
$\Delta : \mathcal{Y} \to \mathcal{Y} \times \mathcal{Y}$
soit représentable par des espaces algébriques. Alors
$$
\mathcal{H}_d(\mathcal{X}/\mathcal{Y})
\longrightarrow
\mathcal{H}_d(\mathcal{X}) \times \mathcal{Y}
$$
qui figure dans
Exemples de champs, équation
(\ref{examples-stacks-equation-diagram-hilbert-d-stack}),
est représentable par des espaces algébriques.
\end{lemma}

\begin{proof}
Soit $U$ un schéma et soit $\xi = (U, Z, p, x, 1)$ un objet de
$\mathcal{H}_d(\mathcal{X}) = \mathcal{H}_d(\mathcal{X}/S)$ au-dessus de $U$.
Ici, $p$ est simplement le morphisme structural de $U$.
La cinquième composante $1$ existe et est unique
puisque tout est au-dessus de $S$.
Soit aussi $y$ un objet de $\mathcal{Y}$ au-dessus de $U$.
Il faut montrer que le $2$-produit fibré
\begin{equation}
\label{criteria-equation-res-isom}
(\Sch/U)_{fppf}
\times_{\xi \times y, \mathcal{H}_d(\mathcal{X}) \times \mathcal{Y}}
\mathcal{H}_d(\mathcal{X}/\mathcal{Y})
\end{equation}
est représentable par un espace algébrique. Pour l'expliquer,
introduisons
$$
I = \mathit{Isom}_\mathcal{Y}(y|_Z, F(x))
$$
qui est un espace algébrique sur $Z$ par hypothèse. Soit $a : U' \to U$
un schéma sur $U$. Que signifie se donner un objet de la catégorie
fibre de (\ref{criteria-equation-res-isom}) au-dessus de $U'$ ? Cela signifie que l'on
dispose d'un objet $\xi' = (U', Z', y', x', \alpha')$ de
$\mathcal{H}_d(\mathcal{X}/\mathcal{Y})$ au-dessus de $U'$ et d'isomorphismes
$(U', Z', p', x', 1) \cong (U, Z, p, x, 1)|_{U'}$ et
$y' \cong y|_{U'}$. Ainsi, $\xi'$ est isomorphe à
$(U', U' \times_{a, U} Z, a^*y, x|_{U' \times_{a, U} Z}, \alpha)$
pour un certain morphisme
$$
\alpha :
a^*y|_{U' \times_{a, U} Z}
\longrightarrow
F(x|_{U' \times_{a, U} Z})
$$
dans la catégorie fibre de $\mathcal{Y}$ au-dessus de $U' \times_{a, U} Z$. Ainsi,
on peut voir $\alpha$ comme un morphisme $b : U' \times_{a, U} Z \to I$.
On voit de cette manière que (\ref{criteria-equation-res-isom})
est représenté par $\text{Res}_{Z/U}(I)$, qui est un espace algébrique d'après la
proposition \ref{criteria-proposition-restriction-of-scalars-algebraic-space}.
\end{proof}

\noindent
Le lemme suivant est une généralisation (partielle) du
lemme \ref{criteria-lemma-etale-covering-hilbert}.

\begin{lemma}
\label{criteria-lemma-representable-on-top}
Soient $F : \mathcal{X} \to \mathcal{Y}$ et $G : \mathcal{X}' \to \mathcal{X}$
des $1$-morphismes de champs en groupoïdes sur $(\Sch/S)_{fppf}$.
Si $G$ est représentable par des espaces algébriques, alors le $1$-morphisme
$$
\mathcal{H}_d(\mathcal{X}'/\mathcal{Y})
\longrightarrow
\mathcal{H}_d(\mathcal{X}/\mathcal{Y})
$$
est représentable par des espaces algébriques.
\end{lemma}

\begin{proof}
Soit $U$ un schéma et soit $\xi = (U, Z, y, x, \alpha)$ un objet de
$\mathcal{H}_d(\mathcal{X}/\mathcal{Y})$ au-dessus de $U$.
Il faut montrer que le $2$-produit fibré
\begin{equation}
\label{criteria-equation-to-show-again}
(\Sch/U)_{fppf}
\times_{\xi, \mathcal{H}_d(\mathcal{X}/\mathcal{Y})}
\mathcal{H}_d(\mathcal{X}'/\mathcal{Y})
\end{equation}
est représentable par un espace algébrique sur $U$.
Un objet de celui-ci au-dessus de $a : U' \to U$ correspond à un objet
$x'$ de $\mathcal{X}'$ au-dessus de $U' \times_{a, U} Z$ tel que
$G(x') \cong x|_{U' \times_{a, U} Z}$. Par hypothèse, le $2$-produit fibré
$$
(\Sch/Z)_{fppf} \times_{x, \mathcal{X}} \mathcal{X}'
$$
est représentable par un espace algébrique $W$ sur $Z$. Il s'ensuit que
(\ref{criteria-equation-to-show-again}) est représenté par $\text{Res}_{Z/U}(W)$,
qui est un espace algébrique d'après la
proposition \ref{criteria-proposition-restriction-of-scalars-algebraic-space}.
\end{proof}

\begin{lemma}
\label{criteria-lemma-limit-preserving}
Soit $F : \mathcal{X} \to \mathcal{Y}$ un $1$-morphisme de champs en groupoïdes
sur $(\Sch/S)_{fppf}$. Supposons que $F$ soit représentable par des espaces
algébriques et localement de présentation finie. Alors
$$
p : \mathcal{H}_d(\mathcal{X}/\mathcal{Y}) \to \mathcal{Y}
$$
commute aux limites sur les objets.
\end{lemma}

\begin{proof}
Cela signifie qu'il faut montrer ceci : étant donnés
\begin{enumerate}
\item un schéma affine $U = \lim_i U_i$ écrit comme la
limite projective filtrante de schémas affines $U_i$ sur $S$,
\item un objet $y_i$ de $\mathcal{Y}$ au-dessus de $U_i$ pour un certain $i$, et
\item un objet $\Xi = (U, Z, y, x, \alpha)$ de
$\mathcal{H}_d(\mathcal{X}/\mathcal{Y})$
au-dessus de $U$ tel que $y = y_i|_U$,
\end{enumerate}
il existe $i' \geq i$ et un objet
$\Xi_{i'} = (U_{i'}, Z_{i'}, y_{i'}, x_{i'}, \alpha_{i'})$ de
$\mathcal{H}_d(\mathcal{X}/\mathcal{Y})$ au-dessus de $U_{i'}$ tel que
$\Xi_{i'}|_U = \Xi$ et $y_{i'} = y_i|_{U_{i'}}$.
En effet, ces deux dernières égalités assurent la commutativité de
(\ref{criteria-equation-limit-preserving}).

\medskip\noindent
Soit $X_{y_i} \to U_i$ un espace algébrique représentant le $2$-produit
fibré
$$
(\Sch/U_i)_{fppf} \times_{y_i, \mathcal{Y}, F} \mathcal{X}.
$$
Notons que $X_{y_i} \to U_i$ est localement de présentation finie par notre
hypothèse sur $F$. Écrivons $\Xi $. Il est clair que
$\xi = (Z, Z \to U_i, x, \alpha)$ est un objet du $2$-produit fibré
affiché ci-dessus ; ainsi, $\xi$ donne naissance à un morphisme
$f_\xi : Z \to X_{y_i}$ d'espaces algébriques sur $U_i$
(puisque $X_{y_i}$ est le foncteur des classes d'isomorphisme d'objets de
$(\Sch/U_i)_{fppf} \times_{y, \mathcal{Y}, F} \mathcal{X}$ ; voir
Champs algébriques,
lemme \ref{algebraic-lemma-characterize-representable-by-space}).
D'après
Limits, lemmes \ref{limits-lemma-descend-finite-presentation} et
\ref{limits-lemma-descend-finite-locally-free}
il existe $i' \geq i$ et un morphisme fini localement libre
$Z_{i'} \to U_{i'}$ de degré $d$ dont le changement de base à $U$ est $Z$. D'après
Limits of Espaces, proposition
\ref{spaces-limits-proposition-characterize-locally-finite-presentation}
on peut, après avoir remplacé $i'$ par un indice plus grand, supposer qu'il existe un
morphisme $f_{i'} : Z_{i'} \to X_{y_i}$ tel que
$$
\xymatrix{
Z \ar[d] \ar[r] \ar@/^3ex/[rr]^{f_\xi} &
Z_{i'} \ar[d] \ar[r]_{f_{i'}} & X_{y_i} \ar[d] \\
U \ar[r] & U_{i'} \ar[r] & U_i
}
$$
soit commutatif. Posons
$\Xi_{i'} = (U_{i'}, Z_{i'}, y_{i'}, x_{i'}, \alpha_{i'})$
où
\begin{enumerate}
\item $y_{i'}$ est l'objet de $\mathcal{Y}$ au-dessus de $U_{i'}$
qui est l'image inverse de $y_i$ sur $U_{i'}$,
\item $x_{i'}$ est l'objet de $\mathcal{X}$ au-dessus de $Z_{i'}$ qui correspond,
par le lemme de $2$-Yoneda, au $1$-morphisme
$$
(\Sch/Z_{i'})_{fppf} \to
\mathcal{S}_{X_{y_i}} \to
(\Sch/U_i)_{fppf} \times_{y_i, \mathcal{Y}, F} \mathcal{X} \to
\mathcal{X}
$$
où la flèche du milieu est l'équivalence qui définit $X_{y_i}$
(notations comme dans
Champs algébriques, sections
\ref{algebraic-section-representable-by-algebraic-spaces} et
\ref{algebraic-section-split}).
\item $\alpha_{i'} : y_{i'}|_{Z_{i'}} \to F(x_{i'})$ est l'isomorphisme
provenant de la $2$-commutativité du diagramme
$$
\xymatrix{
(\Sch/Z_{i'})_{fppf} \ar[r] \ar[rd] &
(\Sch/U_i)_{fppf} \times_{y_i, \mathcal{Y}, F} \mathcal{X}
\ar[r] \ar[d] &
\mathcal{X} \ar[d]^F \\
& (\Sch/U_{i'})_{fppf} \ar[r] & \mathcal{Y}
}
$$
\end{enumerate}
Rappelons que $f_\xi : Z \to X_{y_i}$ était le morphisme correspondant à
l'objet $\xi = (Z, Z \to U_i, x, \alpha)$ de
$(\Sch/U_i)_{fppf} \times_{y_i, \mathcal{Y}, F} \mathcal{X}$
au-dessus de $Z$. Par construction, $f_{i'}$ est le morphisme correspondant à
l'objet $\xi_{i'} = (Z_{i'}, Z_{i'} \to U_i, x_{i'}, \alpha_{i'})$.
Comme $f_\xi = f_{i'} \circ (Z \to Z_{i'})$, on voit que
l'objet $\xi_{i'} = (Z_{i'}, Z_{i'} \to U_i, x_{i'}, \alpha_{i'})$ a pour image inverse
$\xi$ sur $Z$. Ainsi, $x_{i'}$ a pour image inverse $x$ et $\alpha_{i'}$
a pour image inverse $\alpha$. Cela signifie que $\Xi_{i'}$ a pour image inverse
$\Xi$ sur $U$, ce qui conclut.
\end{proof}










\section{Le champ de Hilbert fini d'un point}
\label{criteria-section-hilbert-point}

\noindent
Soit $d \geq 1$ un entier. Dans
Exemples de champs, définition \ref{examples-stacks-definition-hilbert-d-stack},
nous avons défini un champ en groupoïdes $\mathcal{H}_d$.
Dans cette section, nous montrons que $\mathcal{H}_d$ est un
champ algébrique. Nous supposerons partout que
$S = \Spec(\mathbf{Z})$.
Le cas général s'en déduira par changement de base.
Rappelons que la catégorie fibre de $\mathcal{H}_d$ au-dessus d'un schéma $T$
est la catégorie des morphismes finis localement libres $\pi : Z \to T$ de
degré $d$. Au lieu de les classifier directement, nous étudions d'abord
les faisceaux quasi-cohérents d'algèbres $\pi_*\mathcal{O}_Z$.

\medskip\noindent
Soit $R$ un anneau. Adoptons temporairement la définition suivante :
une {\it algèbre libre de rang $d$ sur $R$}
est la donnée d'une structure de $R$-algèbre commutative $m$ sur $R^{\oplus d}$
dont $e_1 = (1, 0, \ldots, 0)$ est l'unité\footnote{Il est peut-être préférable
de la voir comme un couple formé d'une application de multiplication
$m : R^{\oplus d} \otimes_R R^{\oplus d} \to R^{\oplus d}$ et
d'un homomorphisme d'anneaux $\psi : R \to R^{\oplus d}$ satisfaisant à un certain nombre d'axiomes.}.
On considère $m$ comme une application $R$-linéaire
$$
m : R^{\oplus d} \otimes_R R^{\oplus d} \longrightarrow R^{\oplus d}
$$
telle que $m(e_1, x) = m(x, e_1) = x$ et que $m$ définisse une
structure d'anneau commutative et associative. Si l'on écrit
$m(e_i, e_j) = \sum a_{ij}^ke_k$, cela revient
aux conditions
$$
\left\{
\begin{matrix}
\sum_l a_{ij}^la_{lk}^m = \sum_l a_{il}^ma_{jk}^l & \forall i, j, k, m \\
a_{ij}^k = a_{ji}^k & \forall i, j, k \\
a_{i1}^j = \delta_{ij} & \forall i, j
\end{matrix}
\right.
$$
où $\delta_{ij}$ est le symbole $\delta$ de Kronecker. Définissons donc
$$
R_{univ} = \mathbf{Z}[a_{ij}^k]/J
$$
où $J$ est l'idéal engendré par les relations affichées ci-dessus.
Notons
$$
m_{univ} :
R_{univ}^{\oplus d} \otimes_{R_{univ}} R_{univ}^{\oplus d}
\longrightarrow
R_{univ}^{\oplus d}
$$
l'algèbre libre de rang $d$, notée $m$, sur $R_{univ}$ dont les constantes
de structure sont les classes des $a_{ij}^k$ modulo $J$.
Il est alors clair que, pour toute algèbre libre de rang $d$, notée $m$, sur un anneau
$R$, il existe un unique homomorphisme de $\mathbf{Z}$-algèbres
$\psi : R_{univ} \to R$ tel que $\psi_*m_{univ} = m$ (cela signifie que
$m$ s'obtient en appliquant le foncteur de changement de base
$- \otimes_{R_{univ}} R$ à $m_{univ}$). Autrement dit, en posant
$X = \Spec(R_{univ})$, on obtient une identification canonique
$$
X(T) = \{\text{algèbres libres de rang }d\text{, de multiplication }m\text{ sur }R\}
$$
pour $T = \Spec(R)$ variable. Par localisation de Zariski, on obtient
l'identification suivante, apparemment plus générale,
\begin{equation}
\label{criteria-equation-objects}
X(T) = \{\text{algèbres libres de rang }d\text{, de multiplication }
m\text{ sur }\Gamma(T, \mathcal{O}_T)\}
\end{equation}
pour tout schéma $T$.

\medskip\noindent
Parlons maintenant brièvement des {\it isomorphismes d'algèbres libres de rang $d$
sur $R$}. Supposons que $m$ et $m'$ soient deux algèbres libres de rang $d$
sur un anneau $R$. Un {\it isomorphisme de $m$ vers $m'$} est la donnée d'une
application $R$-linéaire inversible
$$
\varphi : R^{\oplus d} \longrightarrow R^{\oplus d}
$$
telle que $\varphi(e_1) = e_1$ et que
$$
m \circ \varphi \otimes \varphi = \varphi \circ m'.
$$
Notons que l'on peut composer ces isomorphismes, de sorte que la collection des
algèbres libres de rang $d$ sur $R$ devienne une catégorie.
On obtient ainsi un foncteur
\begin{equation}
\label{criteria-equation-FAd}
FA_d : \Sch_{fppf}^{opp} \longrightarrow \textit{Groupoïdes}
\end{equation}
de la catégorie des schémas vers les groupoïdes : à un schéma $T$, on associe
l'ensemble des algèbres libres de rang $d$ sur $\Gamma(T, \mathcal{O}_T)$,
muni d'une structure
de catégorie au moyen de la notion d'isomorphisme qui vient d'être définie.

\medskip\noindent
Ce qui précède suggère de considérer le foncteur en groupes $G$
qui associe à tout schéma $T$ le groupe
$$
G(T) = \{g \in \text{GL}_d(\Gamma(T, \mathcal{O}_T)) \mid g(e_1) = e_1\}
$$
Il est clair que $G \subset \text{GL}_d$ (voir
Groupoïdes, exemple \ref{groupoids-example-general-linear-group})
est le sous-schéma en groupes fermé défini par les équations
$x_{11} = 1$ et $x_{i1} = 0$ pour $i > 1$. Ainsi, $G$ est un
schéma en groupes affine et lisse sur $\Spec(\mathbf{Z})$. Considérons
l'action
$$
a : G \times_{\Spec(\mathbf{Z})} X \longrightarrow X
$$
qui associe à un point à valeurs dans $T$, noté $(g, m)$, où $T = \Spec(R)$,
du membre de gauche l'algèbre libre de rang $d$ sur $R$
donnée par
$$
a(g, m) = g^{-1} \circ m \circ g \otimes g.
$$
Notons que cela signifie que $g$ définit un isomorphisme $m \to a(g, m)$
d'algèbres libres de rang $d$ sur $R$. Nous omettons de vérifier que
$a$ définit bien une action du schéma en groupes $G$ sur le schéma $X$.

\begin{lemma}
\label{criteria-lemma-represent-FAd}
Le foncteur en groupoïdes $FA_d$ défini dans (\ref{criteria-equation-FAd})
est isomorphe (!) au foncteur en groupoïdes qui associe
à un schéma $T$ la catégorie dont
\begin{enumerate}
\item l'ensemble des objets est $X(T)$,
\item l'ensemble des morphismes est $G(T) \times X(T)$,
\item $s : G(T) \times X(T) \to X(T)$ est la projection,
\item $t : G(T) \times X(T) \to X(T)$ est $a(T)$, et
\item la composition $G(T) \times X(T) \times_{s, X(T), t} G(T) \times X(T)
\to G(T) \times X(T)$ est donnée par $((g, m), (g', m')) \mapsto (gg', m')$.
\end{enumerate}
\end{lemma}

\begin{proof}
Nous avons vu la règle sur les objets dans (\ref{criteria-equation-objects}).
Nous avons aussi vu ci-dessus que $g \in G(T)$ peut être considéré comme
un morphisme de $m$ vers $a(g, m)$ pour toute algèbre libre de rang $d$, notée $m$.
Réciproquement, tout morphisme $m \to m'$ est donné par une application linéaire
inversible $\varphi$ qui correspond à un élément $g \in G(T)$ tel
que $m' = a(g, m)$.
\end{proof}

\noindent
En fait, le groupoïde $(X, G \times X, s, t, c)$ décrit dans le
lemme ci-dessus est le groupoïde associé à l'action $a : G \times X \to X$
tel qu'il est défini dans
Groupoïdes, lemme \ref{groupoids-lemma-groupoid-from-action}.
Comme $G$ est lisse sur $\Spec(\mathbf{Z})$,
les deux morphismes $s, t : G \times X \to X$ sont
lisses : par symétrie, il suffit de le montrer pour l'un d'eux, et
$s$ est le changement de base de $G \to \Spec(\mathbf{Z})$.
Ainsi, $(G \times X, X, s, t, c)$ est un groupoïde lisse en schémas,
et le champ quotient $[X/G]$ est un champ algébrique d'après
Champs algébriques,
théorème \ref{algebraic-theorem-smooth-groupoid-gives-algebraic-stack}.

\begin{proposition}
\label{criteria-proposition-finite-hilbert-point}
Le champ $\mathcal{H}_d$ est équivalent au champ quotient
$[X/G]$ décrit ci-dessus. En particulier, $\mathcal{H}_d$ est un
champ algébrique.
\end{proposition}

\begin{proof}
Notons que, d'après
Groupoïdes in Espaces, définition
\ref{spaces-groupoids-definition-quotient-stack}
le champ quotient $[X/G]$ est le champ associé à la
catégorie fibrée en groupoïdes associée au « préfaisceau en groupoïdes »
qui associe à un schéma $T$ le groupoïde
$$
(X(T), G(T) \times X(T), s, t, c).
$$
Comme ce « préfaisceau en groupoïdes » est isomorphe à $FA_d$ d'après le
lemme \ref{criteria-lemma-represent-FAd},
il suffit de montrer que $\mathcal{H}_d$ est le champ associé
à (la catégorie fibrée en groupoïdes associée au
« préfaisceau en groupoïdes ») $FA_d$. Pour cela, définissons d'abord un
foncteur
$$
\Spec : FA_d \longrightarrow \mathcal{H}_d
$$
Rappelons que la catégorie fibre de $\mathcal{H}_d$ au-dessus d'un schéma $T$
est la catégorie des morphismes finis localement libres $Z \to T$ de degré $d$.
Étant donnés un schéma $T$ et une algèbre libre de rang $d$
sur $\Gamma(T, \mathcal{O}_T)$, notée $m$, on peut lui associer l'objet
$$
Z = \underline{\Spec}_T(\mathcal{A})
$$
de $\mathcal{H}_{d, T}$,
où $\mathcal{A} = \mathcal{O}_T^{\oplus d}$ est munie d'une
structure de $\mathcal{O}_T$-algèbre au moyen de $m$. De plus, si $m'$ est
une seconde algèbre libre de rang $d$ sur $\Gamma(T, \mathcal{O}_T)$
et si $\varphi : m \to m'$ est un isomorphisme entre elles, alors
l'application $\mathcal{O}_T$-linéaire induite
$\varphi : \mathcal{O}_T^{\oplus d} \to \mathcal{O}_T^{\oplus d}$
induit un isomorphisme
$$
\varphi : \mathcal{A}' \longrightarrow \mathcal{A}
$$
de $\mathcal{O}_T$-algèbres quasi-cohérentes. Par conséquent,
$$
\underline{\Spec}_T(\varphi) :
\underline{\Spec}_T(\mathcal{A})
\longrightarrow
\underline{\Spec}_T(\mathcal{A}')
$$
est un morphisme dans la catégorie fibre $\mathcal{H}_{d, T}$. Nous omettons de
vérifier que cette construction est compatible au changement de base ; on
obtient donc bien un foncteur $\Spec : FA_d \to \mathcal{H}_d$,
comme annoncé ci-dessus.

\medskip\noindent
Pour montrer que $\Spec : FA_d \to \mathcal{H}_d$ induit une équivalence
entre le champ associé à $FA_d$ et $\mathcal{H}_d$, il suffit de
vérifier que
\begin{enumerate}
\item $\mathit{Isom}(m, m') = \mathit{Isom}(\Spec(m), \Spec(m'))$
pour tous $m, m' \in FA_d(T)$.
\item pour tout schéma $T$ et tout objet $Z \to T$ de $\mathcal{H}_{d, T}$,
il existe un recouvrement $\{T_i \to T\}$ tel que $Z|_{T_i}$ soit
isomorphe à $\Spec(m)$ pour un certain $m \in FA_d(T_i)$.
\end{enumerate}
Voir
Champs, lemme \ref{stacks-lemma-stackify-groupoids}.
La première assertion résulte de l'observation que tout isomorphisme
$$
\underline{\Spec}_T(\mathcal{A})
\longrightarrow
\underline{\Spec}_T(\mathcal{A}')
$$
est nécessairement donné par une matrice globale inversible $g$ lorsque
$\mathcal{A} = \mathcal{A}' = \mathcal{O}_T^{\oplus d}$ comme modules.
Pour démontrer la seconde assertion, soit $\pi : Z \to T$ un morphisme fini
localement libre de degré $d$. Alors $\mathcal{A}$ est un faisceau
de $\mathcal{O}_T$-modules localement libre de rang $d$.
Considérons l'élément $1 \in \Gamma(T, \mathcal{A})$. Cet élément est
non nul dans $\mathcal{A} \otimes_{\mathcal{O}_{T, t}} \kappa(t)$
pour tout $t \in T$, puisque le schéma
$Z_t = \Spec(\mathcal{A} \otimes_{\mathcal{O}_{T, t}} \kappa(t))$
est non vide, étant de degré $d > 0$ sur $\kappa(t)$. Ainsi,
$1 : \mathcal{O}_T \to \mathcal{A}$ peut localement servir de premier
élément d'une base (on peut, par exemple, utiliser les parties (1) et (2) du
lemme \ref{algebra-lemma-cokernel-flat} de Algèbre
pour le voir). Après localisation sur
$T$, on peut donc supposer qu'il existe un isomorphisme
$\varphi : \mathcal{A} \to \mathcal{O}_T^{\oplus d}$
tel que $1 \in \Gamma(\mathcal{A})$ corresponde au premier élément de la base.
Dans cette situation, l'application de multiplication
$\mathcal{A} \otimes_{\mathcal{O}_T} \mathcal{A} \to \mathcal{A}$
se traduit, par $\varphi$, en une algèbre libre de rang $d$, notée $m$, sur
$\Gamma(T, \mathcal{O}_T)$. Cela achève la preuve.
\end{proof}




\section{Champs de Hilbert finis d'espaces}
\label{criteria-section-spaces-hilbert}

\noindent
Le champ de Hilbert fini d'un espace algébrique est un champ algébrique.

\begin{lemma}
\label{criteria-lemma-hilbert-stack-of-space}
Soit $S$ un schéma.
Soit $X$ un espace algébrique sur $S$.
Alors $\mathcal{H}_d(X)$ est un champ algébrique.
\end{lemma}

\begin{proof}
Le $1$-morphisme
$$
\mathcal{H}_d(X) \longrightarrow \mathcal{H}_d
$$
est représentable par des espaces algébriques d'après le
lemme \ref{criteria-lemma-representable-on-top}.
Le champ $\mathcal{H}_d$ est algébrique d'après la
proposition \ref{criteria-proposition-finite-hilbert-point}.
Par conséquent, $\mathcal{H}_d(X)$ est un champ algébrique d'après
Champs algébriques,
lemme \ref{algebraic-lemma-representable-morphism-to-algebraic}.
\end{proof}

\noindent
Ce lemme permet d'amorcer le raisonnement.

\begin{lemma}
\label{criteria-lemma-hilbert-stack-relative-space}
Soit $S$ un schéma. Soit $F : \mathcal{X} \to \mathcal{Y}$ un $1$-morphisme
de champs en groupoïdes sur $(\Sch/S)_{fppf}$ tel que
\begin{enumerate}
\item $\mathcal{X}$ soit représentable par un espace algébrique, et
\item $F$ soit représentable par des espaces algébriques, surjectif, plat et
localement de présentation finie.
\end{enumerate}
Alors $\mathcal{H}_d(\mathcal{X}/\mathcal{Y})$ est un champ algébrique.
\end{lemma}

\begin{proof}
Choisissons un champ en groupoïdes représentable $\mathcal{U}$ sur $S$ et un
$1$-morphisme $f : \mathcal{U} \to \mathcal{H}_d(\mathcal{X})$
qui soit représentable par des espaces algébriques, lisse et surjectif.
C'est possible car $\mathcal{H}_d(\mathcal{X})$ est un champ algébrique d'après le
lemme \ref{criteria-lemma-hilbert-stack-of-space}.
Considérons le $2$-produit fibré
$$
\mathcal{W} =
\mathcal{H}_d(\mathcal{X}/\mathcal{Y})
\times_{\mathcal{H}_d(\mathcal{X}), f}
\mathcal{U}.
$$
Puisque $\mathcal{U}$ est représentable (c'est en particulier un champ en
groupoïdes discrets), il résulte de
Exemples de champs, lemme \ref{examples-stacks-lemma-faithful-hilbert}
et de
Champs, lemme \ref{stacks-lemma-2-fibre-product-gives-stack-in-setoids}
que $\mathcal{W}$ est un champ en groupoïdes discrets. Le $1$-morphisme
$\mathcal{W} \to \mathcal{H}_d(\mathcal{X}/\mathcal{Y})$ est
représentable par des espaces algébriques, lisse et surjectif comme changement
de base du morphisme $f$ (voir
Champs algébriques,
lemmes \ref{algebraic-lemma-base-change-representable-by-spaces} et
\ref{algebraic-lemma-base-change-representable-transformations-property}).
Ainsi, si l'on montre que $\mathcal{W}$ est représentable par un espace algébrique,
le lemme résultera de
Champs algébriques,
lemme \ref{algebraic-lemma-smooth-surjective-morphism-implies-algebraic}.

\medskip\noindent
La diagonale de $\mathcal{Y}$ est représentable par des espaces algébriques d'après le
lemme \ref{criteria-lemma-flat-finite-presentation-surjective-diagonal}.
On peut appliquer le
lemme \ref{criteria-lemma-relative-hilbert}
pour voir que le $1$-morphisme
$$
\mathcal{H}_d(\mathcal{X}/\mathcal{Y})
\longrightarrow
\mathcal{H}_d(\mathcal{X}) \times \mathcal{Y}
$$
est représentable par des espaces algébriques. Considérons le $2$-produit fibré
$$
\mathcal{V} =
\mathcal{H}_d(\mathcal{X}/\mathcal{Y})
\times_{(\mathcal{H}_d(\mathcal{X}) \times \mathcal{Y}), f \times F}
(\mathcal{U} \times \mathcal{X}).
$$
Le morphisme de projection $\mathcal{V} \to \mathcal{U} \times \mathcal{X}$
est représentable par des espaces algébriques comme changement de base du dernier
morphisme affiché. Par conséquent, $\mathcal{V}$ est un espace algébrique (voir
Amorçage, lemme \ref{bootstrap-lemma-representable-by-spaces-over-space}
ou
Champs algébriques,
lemme \ref{algebraic-lemma-base-change-by-space-representable-by-space}).
Le $1$-morphisme $\mathcal{V} \to \mathcal{U}$ s'insère dans le diagramme
$2$-cartésien suivant
$$
\xymatrix{
\mathcal{V} \ar[d] \ar[r] & \mathcal{X} \ar[d]^F \\
\mathcal{W} \ar[r] & \mathcal{Y}
}
$$
car
$$
\mathcal{H}_d(\mathcal{X}/\mathcal{Y})
\times_{(\mathcal{H}_d(\mathcal{X}) \times \mathcal{Y}), f \times F}
(\mathcal{U} \times \mathcal{X})
=
(\mathcal{H}_d(\mathcal{X}/\mathcal{Y})
\times_{\mathcal{H}_d(\mathcal{X}), f}
\mathcal{U}) \times_{\mathcal{Y}, F} \mathcal{X}.
$$
Ainsi, $\mathcal{V} \to \mathcal{W}$ est représentable par des espaces algébriques,
surjectif, plat et localement de présentation finie comme changement de base
de $F$. Il en est donc de même pour le morphisme correspondant entre les faisceaux
d'ensembles associés à $\mathcal{V}$ et à $\mathcal{W}$, voir
Champs algébriques, lemme \ref{algebraic-lemma-map-fibred-setoids-property}.
On en conclut que le faisceau associé à $\mathcal{W}$ est un espace
algébrique d'après
Amorçage, théorème \ref{bootstrap-theorem-final-bootstrap}.
\end{proof}




\section{Lieu LCI dans le champ de Hilbert}
\label{criteria-section-lci}

\noindent
On se reportera à
Exemples de champs, section \ref{examples-stacks-section-hilbert-d-stack}
pour les notations. Fixons un $1$-morphisme $F : \mathcal{X} \longrightarrow \mathcal{Y}$
de champs en groupoïdes sur $(\Sch/S)_{fppf}$. Supposons que
$F$ soit représentable par des espaces algébriques. Fixons $d \geq 1$. Considérons un
objet $(U, Z, y, x, \alpha)$ de $\mathcal{H}_d(\mathcal{X}/\mathcal{Y})$.
On en déduit un $1$-morphisme
$$
(\Sch/Z)_{fppf}
\longrightarrow
(\Sch/U)_{fppf} \times_{y, \mathcal{Y}, F} \mathcal{X}
$$
(d'après la propriété universelle des $2$-produits fibrés) qui est représenté par
un morphisme d'espaces algébriques sur $U$.
En effet, puisque $F$ est représentable par des espaces algébriques, on peut choisir
un espace algébrique $X_y$ sur $U$ qui représente le $2$-produit fibré
$(\Sch/U)_{fppf} \times_{y, \mathcal{Y}, F} \mathcal{X}$.
Comme $\alpha : y|_Z \to F(x)$ est un isomorphisme, on voit que
$\xi = (Z, Z \to U, x, \alpha)$ est un objet du $2$-produit fibré
$(\Sch/U)_{fppf} \times_{y, \mathcal{Y}, F} \mathcal{X}$ au-dessus de $Z$.
Ainsi, $\xi$ donne naissance à un morphisme $x_\alpha : Z \to X_y$ d'espaces algébriques
sur $U$, puisque $X_y$ est le foncteur des classes d'isomorphisme d'objets de
$(\Sch/U)_{fppf} \times_{y, \mathcal{Y}, F} \mathcal{X}$, voir
Champs algébriques,
lemme \ref{algebraic-lemma-characterize-representable-by-space}.
Voici un diagramme
\begin{equation}
\label{criteria-equation-relative-map}
\vcenter{
\xymatrix{
Z \ar[r]_{x_\alpha} \ar[rd] & X_y \ar[d] \\
& U
}
}
\quad\quad
\vcenter{
\xymatrix{
(\Sch/Z)_{fppf} \ar[rd] \ar[r]_-{x, \alpha} &
(\Sch/U)_{fppf} \times_{y, \mathcal{Y}, F} \mathcal{X} \ar[r] \ar[d] &
\mathcal{X} \ar[d]^F \\
& (\Sch/U)_{fppf} \ar[r]^y & \mathcal{Y}
}
}
\end{equation}
Remarquons que si
$(f, g, b, a) : (U, Z, y, x, \alpha) \to (U', Z', y', x', \alpha')$
est un morphisme entre objets de $\mathcal{H}_d$, alors le morphisme
$x'_{\alpha'} : Z' \to X'_{y'}$ est le changement de base du morphisme
$x_\alpha$ par le morphisme $g : U' \to U$ (détails omis).

\medskip\noindent
Supposons maintenant, de plus, que $F$ soit plat et localement de présentation finie.
Dans cette situation, définissons une sous-catégorie pleine
$$
\mathcal{H}_{d, lci}(\mathcal{X}/\mathcal{Y}) \subset
\mathcal{H}_d(\mathcal{X}/\mathcal{Y})
$$
formée des objets $(U, Z, y, x, \alpha)$ de
$\mathcal{H}_d(\mathcal{X}/\mathcal{Y})$ tels
que le morphisme correspondant $x_\alpha : Z \to X_y$ soit non ramifié
et localement d'intersection complète (voir
Morphismes d'espaces algébriques, définition \ref{spaces-morphisms-definition-unramified}
et
Compléments sur les morphismes d'espaces algébriques,
définition \ref{spaces-more-morphisms-definition-lci}
pour les définitions correspondantes).

\begin{lemma}
\label{criteria-lemma-lci-locus-stack-in-groupoids}
Soit $S$ un schéma. Fixons un $1$-morphisme
$F : \mathcal{X} \longrightarrow \mathcal{Y}$
de champs en groupoïdes sur $(\Sch/S)_{fppf}$.
Supposons que $F$ soit représentable par des espaces algébriques, plat et localement
de présentation finie. Alors $\mathcal{H}_{d, lci}(\mathcal{X}/\mathcal{Y})$
est un champ en groupoïdes et le foncteur d'inclusion
$$
\mathcal{H}_{d, lci}(\mathcal{X}/\mathcal{Y})
\longrightarrow
\mathcal{H}_d(\mathcal{X}/\mathcal{Y})
$$
est représentable et est une immersion ouverte.
\end{lemma}

\begin{proof}
Soit $\Xi = (U, Z, y, x, \alpha)$ un objet de $\mathcal{H}_d$. Il résulte
de la remarque qui suit
(\ref{criteria-equation-relative-map})
que l'image inverse de $\Xi$ par $U' \to U$ appartient à
$\mathcal{H}_{d, lci}(\mathcal{X}/\mathcal{Y})$ si et seulement si le changement
de base de $x_\alpha$ est non ramifié et localement d'intersection complète.
Notons que $Z \to U$ est fini localement libre (donc plat, localement de
présentation finie et universellement fermé) et que $X_y \to U$ est
plat et localement de présentation finie par notre hypothèse sur $F$. Alors
Compléments sur les morphismes d'espaces algébriques, lemmes
\ref{spaces-more-morphisms-lemma-where-unramified} et
\ref{spaces-more-morphisms-lemma-where-lci}
entraînent qu'il existe un sous-schéma ouvert $W \subset U$ tel qu'un morphisme
$U' \to U$ se factorise par $W$ si et seulement si le changement de base de
$x_\alpha$ par $U' \to U$ est non ramifié et localement d'intersection
complète. Cela implique que
$$
(\Sch/U)_{fppf}
\times_{\Xi, \mathcal{H}_d(\mathcal{X}/\mathcal{Y})}
\mathcal{H}_{d, lci}(\mathcal{X}/\mathcal{Y})
$$
est représentable par $W$. La dernière assertion du lemme
en résulte. La première assertion (à savoir que
$\mathcal{H}_{d, lci}(\mathcal{X}/\mathcal{Y})$ est un champ en groupoïdes)
résulte alors de ce qui précède et de
Champs algébriques,
lemme \ref{algebraic-lemma-open-fibred-category-is-algebraic}.
\end{proof}

\noindent
Les morphismes localement d'intersection complète sont « localement sans obstruction ».
Cette assertion vaut dans une bien plus grande généralité que le cas particulier
dont nous avons besoin ici.

\begin{lemma}
\label{criteria-lemma-lci-unobstructed}
Soit $U \subset U'$ un épaississement du premier ordre de schémas affines.
Soit $X'$ un espace algébrique plat sur $U'$. Posons $X = U \times_{U'} X'$.
Soit $Z \to U$ fini localement libre de degré $d$. Enfin, soit
$f : Z \to X$ non ramifié et localement d'intersection complète.
Alors il existe un diagramme commutatif
$$
\xymatrix{
(Z \subset Z') \ar[rd] \ar[rr]_{(f, f')} & & (X \subset X') \ar[ld] \\
& (U \subset U')
}
$$
d'espaces algébriques sur $U'$ tel que $Z' \to U'$ soit fini localement libre
de degré $d$ et que $Z = U \times_{U'} Z'$.
\end{lemma}

\begin{proof}
D'après
Compléments sur les morphismes d'espaces algébriques,
lemme \ref{spaces-more-morphisms-lemma-unramified-lci},
le faisceau conormal $\mathcal{C}_{Z/X}$ du morphisme non ramifié $Z \to X$
est un $\mathcal{O}_Z$-module localement libre de type fini et, d'après
Compléments sur les morphismes d'espaces algébriques,
lemme \ref{spaces-more-morphisms-lemma-transitivity-conormal-lci},
on a une suite exacte
$$
0 \to i^*\mathcal{C}_{X/X'} \to
\mathcal{C}_{Z/X'} \to
\mathcal{C}_{Z/X} \to 0
$$
de faisceaux conormaux. Comme $Z$ est affine, cette suite est scindée. Choisissons
un scindage
$$
\mathcal{C}_{Z/X'} = i^*\mathcal{C}_{X/X'} \oplus \mathcal{C}_{Z/X}
$$
Soit $Z \subset Z''$ l'épaississement universel du premier ordre de $Z$
sur $X'$ (voir
Compléments sur les morphismes d'espaces algébriques,
section \ref{spaces-more-morphisms-section-universal-thickening}).
Notons $\mathcal{I} \subset \mathcal{O}_{Z''}$ le faisceau quasi-cohérent
d'idéaux correspondant à $Z \subset Z''$. Par définition,
$\mathcal{C}_{Z/X'}$ est $\mathcal{I}$ considéré comme faisceau sur $Z$.
Le scindage ci-dessus détermine donc un scindage
$$
\mathcal{I} = i^*\mathcal{C}_{X/X'} \oplus \mathcal{C}_{Z/X}
$$
Soit $Z' \subset Z''$ le sous-schéma fermé défini par
$\mathcal{C}_{Z/X} \subset \mathcal{I}$, considéré comme faisceau quasi-cohérent
d'idéaux sur $Z''$. Il est clair que $Z'$ est un épaississement du premier ordre
de $Z$ et que l'on obtient un diagramme commutatif d'épaississements du premier ordre
comme dans l'énoncé du lemme.

\medskip\noindent
Puisque $X' \to U'$ est plat et que $X = U \times_{U'} X'$, on voit que
$\mathcal{C}_{X/X'}$ est l'image inverse de $\mathcal{C}_{U/U'}$ sur $X$, voir
Compléments sur les morphismes d'espaces algébriques, lemme \ref{spaces-more-morphisms-lemma-deform}.
Notons que, par construction, $\mathcal{C}_{Z/Z'} = i^*\mathcal{C}_{X/X'}$ ;
on en conclut que $\mathcal{C}_{Z/Z'}$ est isomorphe à l'image inverse
de $\mathcal{C}_{U/U'}$ sur $Z$. En appliquant
Compléments sur les morphismes d'espaces algébriques, lemme \ref{spaces-more-morphisms-lemma-deform}
une nouvelle fois (ou son analogue pour les schémas, voir
Compléments sur les morphismes, lemme \ref{more-morphisms-lemma-deform}),
on en conclut que $Z' \to U'$ est plat et que $Z = U \times_{U'} Z'$.
Enfin,
Compléments sur les morphismes, lemme \ref{more-morphisms-lemma-deform-property}
montre que $Z' \to U'$ est fini localement libre de degré $d$.
\end{proof}

\begin{lemma}
\label{criteria-lemma-lci-formally-smooth}
Soit $F : \mathcal{X} \to \mathcal{Y}$ un $1$-morphisme de champs en groupoïdes
sur $(\Sch/S)_{fppf}$. Supposons que $F$ soit représentable par des espaces
algébriques, plat et localement de présentation finie. Alors
$$
p : \mathcal{H}_{d, lci}(\mathcal{X}/\mathcal{Y}) \to \mathcal{Y}
$$
est formellement lisse sur les objets.
\end{lemma}

\begin{proof}
Il s'agit de montrer ceci : étant donnés
\begin{enumerate}
\item un objet $(U, Z, y, x, \alpha)$ de
$\mathcal{H}_{d, lci}(\mathcal{X}/\mathcal{Y})$ au-dessus d'un schéma affine $U$,
\item un épaississement du premier ordre $U \subset U'$, et
\item un objet $y'$ de $\mathcal{Y}$ au-dessus de $U'$ tel que $y'|_U = y$,
\end{enumerate}
il existe alors un objet $(U', Z', y', x', \alpha')$ de
$\mathcal{H}_{d, lci}(\mathcal{X}/\mathcal{Y})$ au-dessus de $U'$ tel que
$Z = U \times_{U'} Z'$, que $x = x'|_Z$ et que
$\alpha = \alpha'|_U$. En effet, les deux dernières égalités assureront
la commutativité de (\ref{criteria-equation-formally-smooth}).

\medskip\noindent
Considérons le morphisme $x_\alpha : Z \to X_y$ construit dans
l'équation (\ref{criteria-equation-relative-map}). Notons de même $X'_{y'}$
l'espace algébrique sur $U'$ qui représente le $2$-produit fibré
$(\Sch/U')_{fppf} \times_{y', \mathcal{Y}, F} \mathcal{X}$.
Par hypothèse, le morphisme $X'_{y'} \to U'$ est plat (et localement de
présentation finie). Comme $y'|_U = y$, on voit que $X_y = U \times_{U'} X'_{y'}$.
On peut donc appliquer le
lemme \ref{criteria-lemma-lci-unobstructed}
pour trouver $Z' \to U'$ fini localement libre de degré $d$ tel que
$Z = U \times_{U'} Z'$ et tel que $Z' \to X'_{y'}$ prolonge $x_\alpha$.
Par construction, le morphisme $Z' \to X'_{y'}$ correspond à un couple
$(x', \alpha')$. Il est clair que $(U', Z', y', x', \alpha')$
est un objet de $\mathcal{H}_d(\mathcal{X}/\mathcal{Y})$ au-dessus de $U'$
tel que $Z = U \times_{U'} Z'$, que $x = x'|_Z$ et que
$\alpha = \alpha'|_U$. Comme le
lemme \ref{criteria-lemma-lci-locus-stack-in-groupoids}
montre que $\mathcal{H}_{d, lci}(\mathcal{X}/\mathcal{Y}) \subset
\mathcal{H}_d(\mathcal{X}/\mathcal{Y})$ est un « sous-champ ouvert »,
il s'ensuit que $(U', Z', y', x', \alpha')$ est un objet de
$\mathcal{H}_{d, lci}(\mathcal{X}/\mathcal{Y})$, comme voulu.
\end{proof}

\begin{lemma}
\label{criteria-lemma-lci-surjective}
Soit $F : \mathcal{X} \to \mathcal{Y}$ un $1$-morphisme de champs en groupoïdes
sur $(\Sch/S)_{fppf}$. Supposons que $F$ soit représentable par des espaces
algébriques, plat, surjectif et localement de présentation finie. Alors
$$
\coprod\nolimits_{d \geq 1} \mathcal{H}_{d, lci}(\mathcal{X}/\mathcal{Y})
\longrightarrow
\mathcal{Y}
$$
est surjectif sur les objets.
\end{lemma}

\begin{proof}
Il suffit de montrer ceci : pour tout corps $k$
et tout objet $y$ de $\mathcal{Y}$ au-dessus de $\Spec(k)$, il existe
un entier $d \geq 1$ et un objet $(U, Z, y, x, \alpha)$ de
$\mathcal{H}_{d, lci}(\mathcal{X}/\mathcal{Y})$ avec $U = \Spec(k)$.
En effet, dans ce cas, on voit que $p$ est surjectif sur les objets au
sens fort où aucune extension du corps n'est nécessaire.

\medskip\noindent
Notons $X_y$ l'espace algébrique sur $U = \Spec(k)$
qui représente le $2$-produit fibré
$(\Sch/U')_{fppf} \times_{y', \mathcal{Y}, F} \mathcal{X}$.
Par hypothèse, le morphisme $X_y \to \Spec(k)$ est surjectif et
localement de présentation finie (et plat). En particulier, $X_y$ est
non vide. Choisissons un schéma affine non vide $V$ et un morphisme étale
$V \to X_y$. Notons que $V \to \Spec(k)$ est plat, surjectif
et localement de présentation finie (d'après
Morphismes d'espaces algébriques,
définition \ref{spaces-morphisms-definition-locally-finite-presentation}).
Choisissons un point fermé $v \in V$ où $V \to \Spec(k)$ est Cohen-Macaulay
(c'est-à-dire que $V$ est Cohen-Macaulay en $v$), voir
Compléments sur les morphismes,
lemme \ref{more-morphisms-lemma-flat-finite-presentation-CM-open}.
En appliquant
Compléments sur les morphismes,
lemme \ref{more-morphisms-lemma-slice-CM},
on trouve une immersion régulière $Z \to V$ telle que $Z = \{v\}$.
Cela implique que $Z \to V$ est une immersion fermée. De plus, il s'ensuit que
$Z \to \Spec(k)$ est fini (par exemple d'après
Algèbre, lemme \ref{algebra-lemma-isolated-point}).
Ainsi, $Z \to \Spec(k)$ est fini localement libre d'un certain degré $d$.
Le morphisme $Z \to X_y$ est alors non ramifié comme composé
d'une immersion fermée et d'un morphisme étale
(voir
Morphismes d'espaces algébriques, lemmes \ref{spaces-morphisms-lemma-composition-unramified},
\ref{spaces-morphisms-lemma-etale-unramified} et
\ref{spaces-morphisms-lemma-immersion-unramified}).
Enfin, $Z \to X_y$ est localement d'intersection complète
comme composé d'une immersion régulière de schémas et d'un morphisme
étale d'espaces algébriques (voir
Compléments sur les morphismes, lemme \ref{more-morphisms-lemma-regular-immersion-lci}
et
Morphismes d'espaces algébriques, lemmes \ref{spaces-morphisms-lemma-etale-smooth} et
\ref{spaces-morphisms-lemma-smooth-syntomic}, ainsi que
Compléments sur les morphismes d'espaces algébriques,
lemmes \ref{spaces-more-morphisms-lemma-flat-lci} et
\ref{spaces-more-morphisms-lemma-composition-lci}).
Le morphisme $Z \to X_y$ correspond à un objet $x$ de $\mathcal{X}$
au-dessus de $Z$, muni d'un isomorphisme $\alpha : y|_Z \to F(x)$.
On obtient un objet $(U, Z, y, x, \alpha)$ de
$\mathcal{H}_d(\mathcal{X}/\mathcal{Y})$. D'après ce qui précède au sujet
du morphisme $Z \to X_y$, on voit qu'il s'agit en fait d'un objet de la
sous-catégorie $\mathcal{H}_{d, lci}(\mathcal{X}/\mathcal{Y})$, ce qui conclut.
\end{proof}




















\section{Amorçage des champs algébriques}
\label{criteria-section-bootstrap}

\noindent
Le théorème suivant est l'un des principaux résultats de ce chapitre.

\begin{theorem}
\label{criteria-theorem-bootstrap}
\begin{slogan}
Théorème d'Artin sur la représentabilité des groupoïdes plats
\end{slogan}
Soit $S$ un schéma. Soit $F : \mathcal{X} \to \mathcal{Y}$
un $1$-morphisme de champs en groupoïdes sur $(\Sch/S)_{fppf}$. Si
\begin{enumerate}
\item $\mathcal{X}$ est représentable par un espace algébrique, et
\item $F$ est représentable par des espaces algébriques, surjectif, plat et
localement de présentation finie,
\end{enumerate}
alors $\mathcal{Y}$ est un champ algébrique.
\end{theorem}

\begin{proof}
D'après le
lemme \ref{criteria-lemma-flat-finite-presentation-surjective-diagonal},
on voit que la diagonale de $\mathcal{Y}$ est représentable par des espaces
algébriques. Il suffit donc de vérifier l'existence d'un $1$-morphisme
$f : \mathcal{V} \to \mathcal{Y}$ de champs en groupoïdes sur
$(\Sch/S)_{fppf}$, avec $\mathcal{V}$ représentable et
$f$ surjectif et lisse. D'après le
lemme \ref{criteria-lemma-hilbert-stack-relative-space},
on sait que
$$
\coprod\nolimits_{d \geq 1} \mathcal{H}_d(\mathcal{X}/\mathcal{Y})
$$
est un champ algébrique. Il résulte du
lemme \ref{criteria-lemma-lci-locus-stack-in-groupoids}
et de
Champs algébriques,
lemme \ref{algebraic-lemma-open-fibred-category-is-algebraic}
que
$$
\coprod\nolimits_{d \geq 1} \mathcal{H}_{d, lci}(\mathcal{X}/\mathcal{Y})
$$
est lui aussi un champ algébrique. Choisissons un champ en groupoïdes représentable
$\mathcal{V}$ sur $(\Sch/S)_{fppf}$ ainsi qu'un
$1$-morphisme surjectif et lisse
$$
\mathcal{V}
\longrightarrow
\coprod\nolimits_{d \geq 1} \mathcal{H}_{d, lci}(\mathcal{X}/\mathcal{Y}).
$$
Nous affirmons que le composé
$$
\mathcal{V}
\longrightarrow
\coprod\nolimits_{d \geq 1} \mathcal{H}_{d, lci}(\mathcal{X}/\mathcal{Y})
\longrightarrow
\mathcal{Y}
$$
est lisse et surjectif, ce qui achève la preuve du théorème. En effet,
la lissité résultera des
lemmes \ref{criteria-lemma-limit-preserving} et \ref{criteria-lemma-lci-formally-smooth},
et la surjectivité du
lemme \ref{criteria-lemma-lci-surjective}.
Détaillons cela dans le paragraphe suivant.

\medskip\noindent
Par construction, $\mathcal{V} \to
\coprod\nolimits_{d \geq 1} \mathcal{H}_{d, lci}(\mathcal{X}/\mathcal{Y})$
est représentable par des espaces algébriques, surjectif et lisse (donc
aussi localement de présentation finie et formellement lisse d'après le principe
général
Champs algébriques, lemme
\ref{algebraic-lemma-representable-transformations-property-implication}
et
Compléments sur les morphismes d'espaces algébriques,
lemme \ref{spaces-more-morphisms-lemma-smooth-formally-smooth}).
En appliquant les
lemmes \ref{criteria-lemma-representable-by-spaces-limit-preserving},
\ref{criteria-lemma-representable-by-spaces-formally-smooth} et
\ref{criteria-lemma-representable-by-spaces-surjective}
on voit que $\mathcal{V} \to
\coprod\nolimits_{d \geq 1} \mathcal{H}_{d, lci}(\mathcal{X}/\mathcal{Y})$
commute aux limites sur les objets, est formellement lisse sur les objets et
est surjectif sur les objets. Le $1$-morphisme
$\coprod\nolimits_{d \geq 1} \mathcal{H}_{d, lci}(\mathcal{X}/\mathcal{Y})
\to \mathcal{Y}$ possède les propriétés suivantes :
\begin{enumerate}
\item commute aux limites sur les objets : cela résulte du
lemme \ref{criteria-lemma-limit-preserving}
pour $\mathcal{H}_d(\mathcal{X}/\mathcal{Y}) \to \mathcal{Y}$,
que l'on combine aux lemmes
\ref{criteria-lemma-lci-locus-stack-in-groupoids},
\ref{criteria-lemma-open-immersion-limit-preserving} et
\ref{criteria-lemma-composition-limit-preserving}
pour l'obtenir pour $\mathcal{H}_{d, lci}(\mathcal{X}/\mathcal{Y}) \to \mathcal{Y}$ ;
\item est formellement lisse sur les objets d'après le
lemme \ref{criteria-lemma-lci-formally-smooth} ;

\item est surjectif sur les objets d'après le
lemme \ref{criteria-lemma-lci-surjective}.
\end{enumerate}
En utilisant les
lemmes \ref{criteria-lemma-composition-limit-preserving},
\ref{criteria-lemma-composition-formally-smooth} et
\ref{criteria-lemma-composition-surjective}
on en conclut que le composé $\mathcal{V} \to \mathcal{Y}$
commute aux limites sur les objets, est formellement lisse sur les objets et
est surjectif sur les objets.
En utilisant les
lemmes \ref{criteria-lemma-representable-by-spaces-limit-preserving},
\ref{criteria-lemma-representable-by-spaces-formally-smooth} et
\ref{criteria-lemma-representable-by-spaces-surjective}
on voit que $\mathcal{V} \to \mathcal{Y}$ est
localement de présentation finie, formellement lisse et surjectif.
Enfin, en utilisant (par l'intermédiaire du principe général
Champs algébriques,
lemme \ref{algebraic-lemma-representable-transformations-property-implication})
le critère infinitésimal de relèvement
(Compléments sur les morphismes d'espaces algébriques, lemme
\ref{spaces-more-morphisms-lemma-smooth-formally-smooth})
on voit que $\mathcal{V} \to \mathcal{Y}$ est lisse, ce qui conclut.
\end{proof}








\section{Applications}
\label{criteria-section-applications}

\noindent
Notre première tâche consiste à montrer que le champ quotient $[U/R]$ associé à
un « groupoïde plat et localement de présentation finie » est un champ algébrique.
Voir
Groupoïdes in Espaces,
définition \ref{spaces-groupoids-definition-quotient-stack}
pour la définition du champ quotient.
Le lemme préliminaire suivant est l'analogue de
Champs algébriques,
lemme \ref{algebraic-lemma-smooth-quotient-smooth-presentation}.

\begin{lemma}
\label{criteria-lemma-flat-quotient-flat-presentation}
Soit $S$ un schéma appartenant à $\Sch_{fppf}$.
Soit $(U, R, s, t, c)$ un groupoïde en espaces algébriques sur $S$.
Supposons que $s, t$ soient plats et localement de présentation finie.
Alors le morphisme $\mathcal{S}_U \to [U/R]$ est plat, localement de
présentation finie et surjectif.
\end{lemma}

\begin{proof}
Soit $T$ un schéma et soit $x : (\Sch/T)_{fppf} \to [U/R]$
un $1$-morphisme. Il faut montrer que la projection
$$
\mathcal{S}_U \times_{[U/R]} (\Sch/T)_{fppf}
\longrightarrow
(\Sch/T)_{fppf}
$$
est surjective, plate et localement de présentation finie.
On sait déjà que le membre de gauche
est représentable par un espace algébrique $F$, voir
Champs algébriques, lemmes \ref{algebraic-lemma-diagonal-quotient-stack} et
\ref{algebraic-lemma-representable-diagonal}.
Il faut donc montrer que le morphisme correspondant $F \to T$ d'espaces
algébriques est surjectif, localement de présentation finie et plat.
Puisqu'il s'agit de propriétés de morphismes d'espaces
algébriques qui sont locales sur le but pour la topologie fppf, on
peut les vérifier localement pour la topologie fppf sur $T$. Par construction, il existe
un recouvrement fppf $\{T_i \to T\}$ de $T$ tel que
$x|_{(\Sch/T_i)_{fppf}}$ provienne d'un morphisme
$x_i : T_i \to U$. (Notons que $F \times_T T_i$ représente le
$2$-produit fibré $\mathcal{S}_U \times_{[U/R]} (\Sch/T_i)_{fppf}$,
de sorte que tout est compatible au changement de base par $T_i \to T$.)
On peut donc supposer que $x$ provient de $x : T \to U$.
Dans ce cas, on voit que
$$
\mathcal{S}_U \times_{[U/R]} (\Sch/T)_{fppf}
=
(\mathcal{S}_U \times_{[U/R]} \mathcal{S}_U)
\times_{\mathcal{S}_U} (\Sch/T)_{fppf}
=
\mathcal{S}_R \times_{\mathcal{S}_U} (\Sch/T)_{fppf}
$$
La première égalité résulte de
Catégories, lemme \ref{categories-lemma-2-fibre-product-erase-factor},
et la seconde de
Groupoïdes in Espaces,
lemme \ref{spaces-groupoids-lemma-quotient-stack-2-cartesian}.
Manifestement, le dernier $2$-produit fibré est représenté par l'espace
algébrique $F = R \times_{s, U, x} T$, et la projection
$R \times_{s, U, x} T \to T$ est plate et localement de présentation finie
comme changement de base du morphisme d'espaces algébriques
plat et localement de présentation finie $s : R \to U$.
Elle est aussi surjective puisque $s$ admet une section (à savoir la section unité
$e : U \to R$ du groupoïde).
Cela démontre le lemme.
\end{proof}

\noindent
Voici le premier résultat principal de cette section.

\begin{theorem}
\label{criteria-theorem-flat-groupoid-gives-algebraic-stack}
Soit $S$ un schéma appartenant à $\Sch_{fppf}$.
Soit $(U, R, s, t, c)$ un groupoïde en espaces algébriques sur $S$.
Supposons que $s, t$ soient plats et localement de présentation finie.
Alors le champ quotient $[U/R]$ est un champ algébrique sur $S$.
\end{theorem}

\begin{proof}
Vérifions les deux conditions du
théorème \ref{criteria-theorem-bootstrap}
pour le morphisme
$$
(\Sch/U)_{fppf} \longrightarrow [U/R].
$$
La première est immédiate (car $U$ est un espace algébrique).
La seconde est le
lemme \ref{criteria-lemma-flat-quotient-flat-presentation}.
\end{proof}









\section{Quand un champ quotient est-il algébrique ?}
\label{criteria-section-quotient-algebraic}

\noindent
Dans
Groupoïdes in Espaces, section \ref{spaces-groupoids-section-stacks},
nous avons défini le champ quotient $[U/R]$ associé à un groupoïde
$(U, R, s, t, c)$ en espaces algébriques. Notons que $[U/R]$ est un champ
en groupoïdes dont la diagonale est représentable par des espaces algébriques (voir
Amorçage, lemme \ref{bootstrap-lemma-quotient-stack-isom}
et
Champs algébriques, lemme \ref{algebraic-lemma-representable-diagonal})
et qu'il existe un espace algébrique $U$ et un $1$-morphisme
$(\Sch/U)_{fppf} \to [U/R]$ qui est une « surjection fppf »
au sens où il induit un morphisme entre les préfaisceaux des classes d'isomorphisme
d'objets qui devient surjectif après passage aux faisceaux associés.
Cependant, $[U/R]$ n'est pas en général un champ
algébrique. Cela ne contredit pas le
théorème \ref{criteria-theorem-bootstrap},
car le $1$-morphisme $(\Sch/U)_{fppf} \to [U/R]$ peut ne pas
être plat et localement de présentation finie.

\medskip\noindent
La manière la plus simple de construire des exemples de champs quotients non algébriques
consiste à considérer des quotients de la forme $[S/G]$, où $S$ est un schéma et $G$
un schéma en groupes sur $S$ qui agit trivialement sur $S$. En effet, nous verrons
ci-dessous
(lemme \ref{criteria-lemma-BG-algebraic})
que, si $[S/G]$ est algébrique, alors $G \to S$ doit être plat et localement
de présentation finie. On trouvera un exemple explicite dans
Exemples, section \ref{examples-section-not-algebraic-stack}.

\begin{lemma}
\label{criteria-lemma-quotient-algebraic}
Soit $S$ un schéma et soit $B$ un espace algébrique sur $S$.
Soit $(U, R, s, t, c)$ un groupoïde en espaces algébriques sur $B$.
Le champ quotient $[U/R]$ est un champ algébrique si et seulement s'il
existe un morphisme d'espaces algébriques $g : U' \to U$ tel que
\begin{enumerate}
\item le composé
$U' \times_{g, U, t} R \to R \xrightarrow{s} U$ est un morphisme surjectif de
faisceaux, et
\item les morphismes $s', t' : R' \to U'$ sont plats et localement de présentation
finie, où $(U', R', s', t', c')$ est la restriction de
$(U, R, s, t, c)$ par $g$.
\end{enumerate}
\end{lemma}

\begin{proof}
Supposons d'abord que $g : U' \to U$ satisfasse à (1) et (2). La propriété (1)
implique que $[U'/R'] \to [U/R]$ est une équivalence, voir
Groupoïdes in Espaces,
lemme \ref{spaces-groupoids-lemma-quotient-stack-restrict-equivalence}.
D'après le
théorème \ref{criteria-theorem-flat-groupoid-gives-algebraic-stack},
le champ quotient $[U'/R']$ est un champ algébrique. Par conséquent,
$[U/R]$ est lui aussi un champ algébrique, voir
Champs algébriques, lemme \ref{algebraic-lemma-equivalent}.

\medskip\noindent
Réciproquement, supposons que $[U/R]$ soit un champ algébrique. On peut choisir un
schéma $W$ et un $1$-morphisme lisse et surjectif
$$
f : (\Sch/W)_{fppf} \longrightarrow [U/R].
$$
D'après le lemme de $2$-Yoneda
(Champs algébriques, section \ref{algebraic-section-2-yoneda}),
cela correspond à un objet $\xi$ de $[U/R]$ au-dessus de $W$.
D'après la description de $[U/R]$ dans
Groupoïdes in Espaces, lemme \ref{spaces-groupoids-lemma-quotient-stack-objects},
on peut trouver un morphisme de schémas surjectif, plat et localement de présentation finie
$b : U' \to W$ tel que $\xi' = b^*\xi$ corresponde à un morphisme
$g : U' \to U$. Notons que le $1$-morphisme
$$
f' : (\Sch/U')_{fppf} \longrightarrow [U/R].
$$
correspondant à $\xi'$ est surjectif, plat et localement de
présentation finie, voir
Champs algébriques, lemme
\ref{algebraic-lemma-composition-representable-transformations-property}.
Par conséquent,
$(\Sch/U')_{fppf} \times_{[U/R]} (\Sch/U')_{fppf}$
qui est représenté par l'espace algébrique
$$
\mathit{Isom}_{[U/R]}(\text{pr}_0^*\xi', \text{pr}_1^*\xi') =
(U' \times_S U')
\times_{(g \circ \text{pr}_0, g \circ \text{pr}_1), U \times_S U} R = R'
$$
(voir
Groupoïdes in Espaces, lemme
\ref{spaces-groupoids-lemma-quotient-stack-morphisms}
pour la première égalité ; la seconde est la définition de la restriction)
est plat et localement de présentation finie sur $U'$ à la fois par $s'$ et par $t'$
(par changement de base, voir
Champs algébriques, lemme
\ref{algebraic-lemma-base-change-representable-transformations-property}).
D'après cette description de $R'$ et
Champs algébriques, lemme \ref{algebraic-lemma-map-space-into-stack},
on obtient un $1$-morphisme canonique pleinement fidèle $[U'/R'] \to [U/R]$.
Ce $1$-morphisme est essentiellement surjectif parce que $f'$ est plat,
localement de présentation finie et surjectif (voir
Champs, lemme \ref{stacks-lemma-characterize-essentially-surjective-when-ff}) ;
une autre manière de le démontrer consiste à utiliser
Champs algébriques, remarque
\ref{algebraic-remark-flat-fp-presentation}.
Enfin, on peut utiliser
Groupoïdes in Espaces, lemme
\ref{spaces-groupoids-lemma-quotient-stack-restrict-equivalence}
pour conclure que le composé
$U' \times_{g, U, t} R \to R \xrightarrow{s} U$ est un morphisme surjectif de faisceaux.
\end{proof}

\begin{lemma}
\label{criteria-lemma-group-quotient-algebraic}
Soit $S$ un schéma et soit $B$ un espace algébrique sur $S$.
Soit $G$ un espace algébrique en groupes sur $B$.
Soit $X$ un espace algébrique sur $B$ et soit $a : G \times_B X \to X$
une action de $G$ sur $X$ au-dessus de $B$.
Le champ quotient $[X/G]$ est un champ algébrique si et seulement s'il
existe un morphisme d'espaces algébriques $\varphi : X' \to X$ tel que
\begin{enumerate}
\item $G \times_B X' \to X$, $(g, x') \mapsto a(g, \varphi(x'))$, soit un
morphisme surjectif de faisceaux, et
\item les deux projections $X'' \to X'$ de l'espace algébrique $X''$
défini par la règle
$$
T \longmapsto \{(x'_1, g, x'_2) \in (X' \times_B G \times_B X')(T)
\mid \varphi(x'_1) = a(g, \varphi(x'_2))\}
$$
soient plates et localement de présentation finie.
\end{enumerate}
\end{lemma}

\begin{proof}
Ce lemme est un cas particulier du
lemme \ref{criteria-lemma-quotient-algebraic}.
En effet, le champ quotient $[X/G]$ est, d'après
Groupoïdes in Espaces, définition \ref{spaces-groupoids-definition-quotient-stack},
égal au champ quotient $[X/G \times_B X]$ du groupoïde
en espaces algébriques $(X, G \times_B X, s, t, c)$ associé à
l'action du groupe dans
Groupoïdes in Espaces, lemme \ref{spaces-groupoids-lemma-groupoid-from-action}.
Une petite observation est nécessaire pour obtenir la condition (1).
En effet, le morphisme $s : G \times_B X \to X$ est la seconde projection
et le morphisme $t :  G \times_B X \to X$ est le morphisme d'action $a$.
Ainsi, le morphisme $h : U' \times_{g, U, t} R \to R \xrightarrow{s} U$ du
lemme \ref{criteria-lemma-quotient-algebraic}
correspond au morphisme
$$
X' \times_{\varphi, X, a} (G \times_B X) \xrightarrow{\text{pr}_1} X
$$
dans la situation présente. Toutefois, grâce à la symétrie fournie par
le passage à l'inverse dans $G$, ce morphisme est isomorphe au morphisme
$$
(G \times_B X) \times_{\text{pr}_1, X, \varphi} X' \xrightarrow{a} X
$$
de l'énoncé du lemme. Les détails sont omis.
\end{proof}

\begin{lemma}
\label{criteria-lemma-BG-algebraic}
\begin{slogan}
Les gerbes sont algébriques si et seulement si les groupes associés sont plats
et localement de présentation finie
\end{slogan}
Soit $S$ un schéma et soit $B$ un espace algébrique sur $S$.
Soit $G$ un espace algébrique en groupes sur $B$.
Munissons $B$ de l'action triviale de $G$.
Alors le champ quotient $[B/G]$ est un champ algébrique
si et seulement si $G$ est plat et localement de présentation finie sur $B$.
\end{lemma}

\begin{proof}
Si $G$ est plat et localement de présentation finie sur $B$, alors
$[B/G]$ est un champ algébrique d'après le
théorème \ref{criteria-theorem-flat-groupoid-gives-algebraic-stack}.

\medskip\noindent
Réciproquement, supposons que $[B/G]$ soit un champ algébrique. D'après le
lemme \ref{criteria-lemma-group-quotient-algebraic}
et puisque l'action est triviale, on voit
qu'il existe un espace algébrique $B'$ et un morphisme
$B' \to B$ tels que (1) $B' \to B$ soit un morphisme surjectif
de faisceaux et (2) que les projections
$$
B' \times_B G \times_B B' \to B'
$$
soient plates et localement de présentation finie. Notons que le changement de base
$B' \times_B G \times_B B' \to G \times_B B'$ de $B' \to B$
est lui aussi un morphisme surjectif de faisceaux. Il résulte donc de
Descente et espaces algébriques, lemme \ref{spaces-descent-lemma-curiosity},
que la projection $G \times_B B' \to B'$ est plate et localement
de présentation finie. D'après (1), on peut trouver un recouvrement fppf
$\{B_i \to B\}$ tel que $B_i \to B$ se factorise par $B' \to B$.
Ainsi, $G \times_B B_i \to B_i$ est plat et localement de présentation finie
par changement de base. D'après
Descente et espaces algébriques, lemmes
\ref{spaces-descent-lemma-descending-property-flat} et
\ref{spaces-descent-lemma-descending-property-locally-finite-presentation}
on en conclut que $G \to B$ est plat et localement de présentation finie.
\end{proof}

\noindent
Nous verrons plus loin que le champ quotient d'un $S$-espace lisse
par un espace algébrique en groupes $G$ est lisse, même lorsque $G$ ne l'est pas
(Morphismes of Champs, lemme
\ref{stacks-morphisms-lemma-smooth-quotient-stack}).




\section{Champs algébriques pour la topologie étale}
\label{criteria-section-stacks-etale}

\noindent
Soit $S$ un schéma. Au lieu de travailler avec des champs en groupoïdes sur
le grand site fppf $(\Sch/S)_{fppf}$, on pourrait travailler avec des champs en groupoïdes
sur le grand site étale $(\Sch/S)_\etale$. Tout le contenu de
Champs algébriques, sections
\ref{algebraic-section-representable},
\ref{algebraic-section-2-yoneda},
\ref{algebraic-section-representable-morphism},
\ref{algebraic-section-split},
\ref{algebraic-section-representable-by-algebraic-spaces},
\ref{algebraic-section-morphisms-representable-by-algebraic-spaces},
\ref{algebraic-section-representable-properties} et
\ref{algebraic-section-stacks}
s'applique aux catégories fibrées en groupoïdes sur $(\Sch/S)_\etale$.
On obtient ainsi une seconde notion de champ algébrique en travaillant pour la
topologie étale. Cette notion est a priori plus faible que celle introduite
dans Champs algébriques, définition \ref{algebraic-definition-algebraic-stack},
puisqu'un champ pour la topologie fppf est certainement un champ pour la topologie
étale. Toutefois, les deux notions sont équivalentes, comme le montre le
lemme suivant.

\begin{lemma}
\label{criteria-lemma-stacks-etale}
La catégorie sous-jacente commune à $\Sch_{fppf}$
et à $\Sch_\etale$ sera notée $\Sch_\alpha$ (voir
Faisceaux on Champs, section \ref{stacks-sheaves-section-sheaves} et
Topologies, remarque \ref{topologies-remark-choice-sites}). Soit $S$ un objet
de $\Sch_\alpha$. Soit
$$
p : \mathcal{X} \to \Sch_\alpha/S
$$
une catégorie fibrée en groupoïdes ayant les propriétés suivantes :
\begin{enumerate}
\item $\mathcal{X}$ est un champ en groupoïdes sur $(\Sch/S)_\etale$,
\item la diagonale $\Delta : \mathcal{X} \to \mathcal{X} \times \mathcal{X}$
est représentable par des espaces algébriques\footnote{Ici, on peut entendre
soit des faisceaux pour la topologie étale dont la diagonale est représentable et qui
admettent un recouvrement étale surjectif par un schéma, soit des espaces algébriques
au sens de
Espaces algébriques, définition \ref{spaces-definition-algebraic-space}.
En effet, d'après Amorçage, lemme \ref{bootstrap-lemma-spaces-etale},
il n'y a aucune différence.}, et
\item il existe $U \in \Ob(\Sch_\alpha/S)$
et un $1$-morphisme $(\Sch/U)_\etale \to \mathcal{X}$
qui est surjectif et lisse.
\end{enumerate}
Alors $\mathcal{X}$ est un champ algébrique au sens de
Champs algébriques, définition \ref{algebraic-definition-algebraic-stack}.
\end{lemma}

\begin{proof}
Notons que les propriétés (2) et (3) du lemme et les propriétés
(2) et (3) correspondantes de
Champs algébriques, définition \ref{algebraic-definition-algebraic-stack},
sont indépendantes de la topologie. En effet, ces propriétés
ne font intervenir que la notion de $2$-produit fibré de catégories fibrées en
groupoïdes, les $1$- et $2$-morphismes de catégories fibrées en groupoïdes, la
notion de $1$-morphisme de catégories fibrées en groupoïdes représentable
par des espaces algébriques et ce que signifie, pour un tel $1$-morphisme, être
surjectif et lisse.
Il suffit donc de montrer que tout champ en groupoïdes pour la topologie étale,
noté $\mathcal{X}$, qui possède les propriétés (2) et (3), est aussi un champ en groupoïdes pour la topologie fppf.

\medskip\noindent
D'après (2), soit $R$ un espace algébrique représentant
$$
(\Sch_\alpha/U) \times_\mathcal{X} (\Sch_\alpha/U)
$$
D'après (3), les projections $s, t : R \to U$ sont lisses. Exactement comme dans la preuve de
Champs algébriques, lemme \ref{algebraic-lemma-map-space-into-stack},
il existe un groupoïde en espaces $(U, R, s, t, c)$ et, canoniquement,
un $1$-morphisme pleinement fidèle $[U/R]_\etale \to \mathcal{X}$,
où $[U/R]_\etale$ est le champ associé pour la topologie étale au préfaisceau
en groupoïdes
$$
T \longmapsto (U(T), R(T), s(T), t(T), c(T))
$$
Affirmation : si $V \to T$ est un morphisme lisse et surjectif d'un espace algébrique
$V$ vers un schéma $T$, alors il existe un recouvrement étale $\{T_i \to T\}$
plus fin que le recouvrement $\{V \to T\}$. Cela résulte de
Compléments sur les morphismes, lemme \ref{more-morphisms-lemma-etale-dominates-smooth}
ou du résultat plus général
Faisceaux on Champs, lemme
\ref{stacks-sheaves-lemma-surjective-flat-locally-finite-presentation}.
En utilisant l'affirmation et en raisonnant exactement comme dans
Champs algébriques, lemme \ref{algebraic-lemma-stack-presentation},
on voit que $[U/R]_\etale \to \mathcal{X}$ est une
équivalence.

\medskip\noindent
Notons ensuite $[U/R]$ le champ quotient pour la topologie fppf,
qui est un champ algébrique d'après
Champs algébriques, théorème
\ref{algebraic-theorem-smooth-groupoid-gives-algebraic-stack}.
On a donc des $1$-morphismes
$$
U \to [U/R]_\etale \to [U/R].
$$
Les morphismes $U \to [U/R]_\etale \cong \mathcal{X}$ et
$U \to [U/R]$ sont tous deux surjectifs et lisses (le premier par hypothèse
et le second d'après le théorème) et, dans les deux cas, les
produits fibrés $U \times_\mathcal{X} U$ et $U \times_{[U/R]} U$
sont représentés par $R$. Par conséquent, le $1$-morphisme
$[U/R]_\etale \to [U/R]$ est pleinement fidèle (puisque les morphismes
dans les champs quotients sont donnés par des morphismes vers $R$, voir
Groupoïdes in Espaces, section
\ref{spaces-groupoids-section-explicit-quotient-stacks}).

\medskip\noindent
Enfin, pour tout schéma $T$ et tout morphisme $t : T \to [U/R]$, le produit fibré
$V = T \times_{U/R} U$ est un espace algébrique surjectif et lisse sur $T$.
D'après l'affirmation ci-dessus, il existe un recouvrement étale $\{T_i \to T\}_{i \in I}$
et des morphismes $T_i \to V$ au-dessus de $T$. Cela prouve que l'objet
$t$ de $[U/R]$ au-dessus de $T$ provient de $U$ localement pour la topologie étale. On en conclut que
$[U/R]_\etale \to [U/R]$ est une équivalence de champs en
groupoïdes sur $(\Sch/S)_\etale$ d'après
Champs, lemme \ref{stacks-lemma-characterize-essentially-surjective-when-ff}.
Cela achève la preuve.
\end{proof}














% OK, commencer ici.
%

\chapter{Axiomes d'Artin}


\phantomsection
\label{artin-section-phantom}





\section{Introduction}
\label{artin-section-introduction}

\noindent
Dans ce chapitre, nous étudions les axiomes d'Artin relatifs à la
représentabilité des foncteurs par des espaces algébriques. Nous renvoyons aux articles
\cite{ArtinI}, \cite{ArtinII}, \cite{ArtinVersal}.

\medskip\noindent
Certaines notations, conventions et certains termes employés dans ce chapitre sont peu commodes
et peuvent sembler aller à rebours des usages au lecteur expérimenté. Ce choix est délibéré.
Voir Quot, section \ref{quot-section-conventions}, pour une
explication.

\medskip\noindent
Soit $S$ un schéma de base localement noethérien. Soit
$$
p : \mathcal{X} \longrightarrow (\Sch/S)_{fppf}
$$
une catégorie fibrée en groupoïdes. Soit $x_0$ un objet de $\mathcal{X}$
au-dessus d'un corps $k$ de type fini sur $S$. Dans tout ce chapitre, un rôle
important est joué par la catégorie de prédéformations
(voir Théorie formelle des déformations,
définition \ref{formal-defos-definition-predeformation-category})
$$
\mathcal{F}_{\mathcal{X}, k, x_0}
\longrightarrow
\{\text{algèbres locales artiniennes sur }S\text{ de corps résiduel }k\}
$$
associée à $x_0$ au-dessus de $k$. Nous introduisons la condition de Rim--Schlessinger (RS)
pour $\mathcal{X}$ et montrons qu'elle entraîne que
$\mathcal{F}_{\mathcal{X}, k, x_0}$ est une catégorie de déformations, c'est-à-dire que
$\mathcal{F}_{\mathcal{X}, k, x_0}$ satisfait elle-même (RS).
Nous étudions comment $\mathcal{F}_{\mathcal{X}, k, x_0}$
se transforme lorsqu'on remplace $k$ par une extension finie,
ainsi que ses espaces tangents.

\medskip\noindent
Nous étudions ensuite les objets formels $\xi = (\xi_n)$ de $\mathcal{X}$, c'est-à-dire les
systèmes projectifs d'objets situés au-dessus des quotients $R/\mathfrak m^n$,
où $R$ est une $S$-algèbre locale noethérienne complète dont le corps résiduel
est de type fini sur $S$. Cela revient à se donner un objet formel
de $\mathcal{F}_{\mathcal{X}, k, x_0}$ pour certains $x_0$ et $k$.
Un objet formel est dit effectif s'il existe un objet de
$\mathcal{X}$ au-dessus de $R$ qui engendre ce système projectif.
Un objet formel de $\mathcal{X}$ est dit versel s'il engendre un
objet formel versel de $\mathcal{F}_{\mathcal{X}, k, x_0}$.
Enfin, étant donnés un $S$-schéma $U$ de type fini, un objet $x$
de $\mathcal{X}$ au-dessus de $U$ et un point fermé $u_0 \in U$, on dit que
$x$ est versel en $u_0$ si l'objet formel induit sur l'anneau local complété
$\mathcal{O}_{U, u_0}^\wedge$ est versel.

\medskip\noindent
Après ces préliminaires, nous pouvons énoncer le célèbre théorème
d'Artin : notre $\mathcal{X}$ est un champ algébrique si les conditions suivantes sont satisfaites :
\begin{enumerate}
\item $\mathcal{O}_{S, s}$ est un G-anneau pour tout $s \in S$\footnote{Cela
signifie que $S$ est un G-schéma ; voir
Propriétés, section \ref{properties-section-G-scheme}.},
\item $\Delta : \mathcal{X} \to \mathcal{X} \times \mathcal{X}$
est représentable par des espaces algébriques,
\item $\mathcal{X}$ est un champ pour la topologie \'etale,
\item $\mathcal{X}$ commute aux limites,
\item $\mathcal{X}$ satisfait (RS),
\item les espaces tangents et les espaces d'automorphismes infinitésimaux
des catégories de déformations $\mathcal{F}_{\mathcal{X}, k, x_0}$
sont de dimension finie,
\item les objets formels sont effectifs,
\item $\mathcal{X}$ satisfait l'ouverture de la versalité.
\end{enumerate}
Il s'agit du lemme \ref{artin-lemma-diagonal-representable} ; voir aussi la
proposition \ref{artin-proposition-second-diagonal-representable}
pour une légère amélioration. Une proposition analogue
caractérise les foncteurs $F : (\Sch/S)_{fppf}^{opp} \to \textit{Ensembles}$
qui sont des espaces algébriques ; voir la section \ref{artin-section-algebraic-spaces}.

\medskip\noindent
Voici les grandes lignes de la démonstration du théorème d'Artin.
Nous montrons d'abord qu'il existe suffisamment d'objets formels versels
en utilisant (RS) et la finitude des dimensions des espaces tangents et d'automorphismes ; voir
par exemple Théorie formelle des déformations, lemme
\ref{formal-defos-lemma-minimal-groupoid-in-functors-construction}.
Ces objets formels sont effectifs par hypothèse.
Les objets formels effectifs peuvent être ``approchés'' par des objets $x$ au-dessus de
$S$-schémas $U$ de type fini ; voir le lemme \ref{artin-lemma-approximate}.
Cette approximation utilise le fait que les anneaux locaux de $S$ sont des G-anneaux et
que $\mathcal{X}$ commute aux limites ; c'est peut-être la partie la plus difficile
de la démonstration, car elle s'appuie sur la désingularisation générale de N\'eron pour
approcher des solutions formelles d'équations algébriques sur un G-anneau local
noethérien par des solutions dans son hensélisé.
L'ouverture de la versalité permet ensuite, quitte à restreindre $U$,
de supposer que $x$ est versel en tout point fermé de $U$.
Après ces étapes, nous montrons que $U \to \mathcal{X}$
est un morphisme lisse. En prenant suffisamment de morphismes $U \to \mathcal{X}$,
nous obtenons un ``atlas lisse'' de $\mathcal{X}$, ce qui montre
que $\mathcal{X}$ est un champ algébrique.

\medskip\noindent
Lorsqu'on vérifie les axiomes d'Artin pour une catégorie $\mathcal{X}$ fibrée
en groupoïdes, l'étape la plus difficile consiste souvent à établir l'ouverture
de la versalité. Dans la suite, supposons que $\mathcal{X}/S$
satisfasse déjà les autres conditions énumérées ci-dessus.
Nous présentons dans ce chapitre deux méthodes qui permettent
de démontrer que $\mathcal{X}$ satisfait l'ouverture de la versalité :
\begin{enumerate}
\item La première consiste à supposer une condition de Rim--Schlessinger
plus forte, appelée (RS*), ainsi qu'une version renforcée de
l'effectivité formelle, qui exige essentiellement que les objets au-dessus de
systèmes projectifs d'épaississements soient effectifs. Sous
ces hypothèses, l'ouverture de la versalité en découle automatiquement ;
voir le lemme \ref{artin-lemma-SGE-implies-openness-versality}.
Remarquons que nous utilisons ici de manière essentielle le fait
que $\mathcal{X}$ est définie sur la catégorie de tous les schémas
au-dessus de $S$, et non seulement sur celle des schémas noethériens !
\item La seconde méthode, suivant Artin, consiste à munir $\mathcal{X}$
d'une théorie des obstructions. Si cette théorie des obstructions
``commute aux produits'' en un sens convenable, alors
$\mathcal{X}$ satisfait l'ouverture de la versalité ; voir le
lemme \ref{artin-lemma-get-openness-obstruction-theory}.
\end{enumerate}
Les théories des obstructions admettent de nombreuses axiomatisations
et la littérature en propose effectivement maintes variantes, souvent adaptées à des champs de modules particuliers.
dans la littérature. Nous présentons une variante fondée sur la catégorie dérivée
(qui apparaît souvent naturellement dans les calculs de théorie des déformations
effectués dans la littérature) dans le lemme \ref{artin-lemma-dual-openness}.

\medskip\noindent
Dans la section \ref{artin-section-algebraic-spaces-noetherian},
nous indiquons les modifications nécessaires pour traiter
les foncteurs définis sur la catégorie
$(\textit{Noetherian}/S)_\etale$ des schémas localement noethériens
au-dessus de $S$.

\medskip\noindent
Dans la dernière section de ce chapitre, nous appliquons les axiomes d'Artin
pour démontrer le théorème d'Artin sur l'existence des contractions ; voir la
section \ref{artin-section-contractions}. Ce théorème affirme, en gros, qu'étant donnés un
espace algébrique $X'$ séparé et de type fini sur $S$,
une partie fermée $T' \subset |X'|$ et une modification formelle
$$
\mathfrak{f} : X'_{/T'} \longrightarrow \mathfrak{X}
$$
où $\mathfrak{X}$ est un espace algébrique formel noethérien au-dessus de $S$,
il existe un morphisme propre $f : X' \to X$ qui
``réalise la contraction''. Cela signifie qu'il existe une
identification $\mathfrak{X} = X_{/T}$ telle que
$\mathfrak{f} = f_{/T'} : X'_{/T'} \to X_{/T}$, où $T = f(T')$,
et que $f$ soit en outre un isomorphisme au-dessus de $X \setminus T$. La démonstration consiste
à définir un foncteur $F$ sur la catégorie des schémas localement noethériens
au-dessus de $S$ et à vérifier les axiomes d'Artin pour $F$. Fait remarquable, dans cette
application des axiomes d'Artin, l'ouverture de la versalité n'est pas le point le plus
difficile à démontrer ; c'est plutôt la preuve que $F$ commute aux limites qui exige
de nombreux travaux et résultats préliminaires.




\section{Conventions}
\label{artin-section-conventions}

\noindent
Les conventions employées dans ce chapitre sont les mêmes que dans le
chapitre consacré aux champs algébriques ; voir
Champs algébriques, section \ref{algebraic-section-conventions}.
Dans ce chapitre, le schéma de base $S$ sera souvent localement noethérien
(bien que nous répétions toujours cette hypothèse dans les énoncés
des résultats).





\section{Catégories de prédéformations}
\label{artin-section-predeformation-categories}

\noindent
Soit $S$ un schéma de base localement noethérien. Soit
$$
p : \mathcal{X} \longrightarrow (\Sch/S)_{fppf}
$$
une catégorie fibrée en groupoïdes. Soit $k$ un corps
et soit $\Spec(k) \to S$ un morphisme de type fini (voir
Morphismes, lemme \ref{morphisms-lemma-point-finite-type}). Nous dirons parfois
simplement que {\it $k$ est un corps de type fini sur $S$}. Soit
$x_0$ un objet de $\mathcal{X}$ au-dessus de $\Spec(k)$.
À partir de $S$, $\mathcal{X}$, $k$ et $x_0$, nous allons construire une
catégorie de prédéformations, au sens de
Théorie formelle des déformations,
définition \ref{formal-defos-definition-predeformation-category}.
Cette construction sera analogue à celle de
Théorie formelle des déformations,
remarque \ref{formal-defos-remark-localize-cofibered-groupoid}.

\medskip\noindent
D'abord, d'après Morphismes, lemme \ref{morphisms-lemma-point-finite-type},
nous pouvons choisir un ouvert affine $\Spec(\Lambda) \subset S$ tel que
$\Spec(k) \to S$ se factorise par $\Spec(\Lambda)$ et que l'homomorphisme
d'anneaux associé $\Lambda \to k$ soit fini. Nous obtenons ainsi la catégorie
$\mathcal{C}_\Lambda$ ; voir
Théorie formelle des déformations, définition \ref{formal-defos-definition-CLambda}.
À équivalence canonique près, la catégorie $\mathcal{C}_\Lambda$
ne dépend pas du choix de l'ouvert affine $\Spec(\Lambda)$ de $S$.
En effet, $\mathcal{C}_\Lambda$ est équivalente à l'opposée
de la catégorie des factorisations
\begin{equation}
\label{artin-equation-factor}
\Spec(k) \to \Spec(A) \to S
\end{equation}
du morphisme structural telles que $A$ soit un anneau local artinien et
que $\Spec(k) \to \Spec(A)$ corresponde à un homomorphisme d'anneaux $A \to k$
qui identifie $k$ au corps résiduel de $A$.

\medskip\noindent
Notons $\mathcal{F} = \mathcal{F}_{\mathcal{X}, k, x_0}$ la
catégorie dont
\begin{enumerate}
\item les objets sont les morphismes $x_0 \to x$ de $\mathcal{X}$ tels que
$p(x) = \Spec(A)$, où $A$ est un anneau local artinien, et que
$p(x_0) \to p(x) \to S$ soit une factorisation comme en (\ref{artin-equation-factor}) ;
\item les morphismes $(x_0 \to x) \to (x_0 \to x')$ sont les diagrammes
commutatifs
$$
\xymatrix{
x & & x' \ar[ll] \\
& x_0 \ar[lu] \ar[ru]
}
$$
dans $\mathcal{X}$. (Noter le renversement des flèches.)
\end{enumerate}
Si $x_0 \to x$ est un objet de $\mathcal{F}$, en écrivant $p(x) = \Spec(A)$,
nous obtenons un objet $A$ de $\mathcal{C}_\Lambda$. Nous dirons souvent que
$x_0 \to x$, ou $x$, est au-dessus de $A$. Un morphisme de $\mathcal{F}$ entre les objets
$x_0 \to x$ au-dessus de $A$ et $x_0 \to x'$ au-dessus de $A'$
correspond à un morphisme $x' \to x$ de $\mathcal{X}$, donc à un morphisme
$p(x' \to x) : \Spec(A') \to \Spec(A)$, lequel correspond à son tour à un
homomorphisme d'anneaux $A \to A'$. Comme $\mathcal{X}$ est une catégorie
au-dessus de la catégorie des schémas sur $S$, nous voyons que $A \to A'$ est un
homomorphisme de $\Lambda$-algèbres. Nous obtenons donc un foncteur
\begin{equation}
\label{artin-equation-predeformation-category}
p : \mathcal{F} = \mathcal{F}_{\mathcal{X}, k, x_0}
\longrightarrow
\mathcal{C}_\Lambda.
\end{equation}
Nous noterons $\mathcal{F}(A)$ la catégorie fibre
au-dessus d'un objet $A$ de $\mathcal{C}_\Lambda$. Un objet de $\mathcal{F}(A)$
est simplement un morphisme $x_0 \to x$ de $\mathcal{X}$ tel que
$x$ soit au-dessus de $\Spec(A)$ et $x_0 \to x$ au-dessus de $\Spec(k) \to \Spec(A)$.

\begin{lemma}
\label{artin-lemma-predeformation-category}
Le foncteur $p : \mathcal{F} \to \mathcal{C}_\Lambda$ défini ci-dessus
est une catégorie de prédéformations.
\end{lemma}

\begin{proof}
Il faut montrer (a) que $\mathcal{F}$ est cofibrée en groupoïdes au-dessus de
$\mathcal{C}_\Lambda$ et (b) que $\mathcal{F}(k)$ est équivalente à une
catégorie possédant un unique objet et un unique morphisme.

\medskip\noindent
Démonstration de (a). Les catégories fibres de $\mathcal{F}$
au-dessus de $\mathcal{C}_\Lambda$ sont des groupoïdes puisque les catégories fibres
de $\mathcal{X}$ le sont. Soient $A \to A'$ un morphisme de
$\mathcal{C}_\Lambda$ et $x_0 \to x$ un objet de $\mathcal{F}(A)$.
Comme $\mathcal{X}$ est fibrée en groupoïdes, il existe un morphisme
$x' \to x$ au-dessus de $\Spec(A') \to \Spec(A)$. Puisque le composé
$A \to A' \to k$ est égal à l'homomorphisme donné $A \to k$, l'unicité
des images inverses à isomorphisme près montre que l'image inverse par $\Spec(k) \to \Spec(A')$
de $x'$ est $x_0$ ; autrement dit, il existe un morphisme $x_0 \to x'$
au-dessus de $\Spec(k) \to \Spec(A')$, compatible avec
$x_0 \to x$ et $x' \to x$. Ainsi, $\mathcal{F}$ admet des
images directes. On conclut par le dual de
Catégories, lemme \ref{categories-lemma-fibred-groupoids}.

\medskip\noindent
Démonstration de (b). Si $A = k$, alors $\Spec(k) = \Spec(A)$ et, puisque $\mathcal{X}$
est fibrée en groupoïdes au-dessus de $(\Sch/S)_{fppf}$, pour tout objet
$x_0 \to x$ de $\mathcal{F}(k)$, le morphisme $x_0 \to x$ est un isomorphisme.
Tout objet de $\mathcal{F}(k)$ est donc isomorphe à $x_0 \to x_0$.
Enfin, l'unique endomorphisme de $x_0 \to x_0$ dans $\mathcal{F}$ est
l'identité.
\end{proof}

\noindent
Soit $S$ un schéma de base localement noethérien. Soit
$F : \mathcal{X} \to \mathcal{Y}$ un $1$-morphisme entre catégories
fibrées en groupoïdes au-dessus de $(\Sch/S)_{fppf}$. Soit $k$ un corps
de type fini sur $S$. Soit $x_0$ un objet de $\mathcal{X}$ au-dessus
de $\Spec(k)$. Posons $y_0 = F(x_0)$ ; c'est un objet de $\mathcal{Y}$
au-dessus de $\Spec(k)$. Alors $F$ induit un foncteur
\begin{equation}
\label{artin-equation-functoriality}
F :
\mathcal{F}_{\mathcal{X}, k, x_0}
\longrightarrow
\mathcal{F}_{\mathcal{Y}, k, y_0}
\end{equation}
entre catégories cofibrées au-dessus de $\mathcal{C}_\Lambda$. Plus précisément, à l'objet
$x_0 \to x$ de $\mathcal{F}_{\mathcal{X}, k, x_0}(A)$, nous associons
l'objet $F(x_0) \to F(x)$ de $\mathcal{F}_{\mathcal{Y}, k, y_0}(A)$.

\begin{lemma}
\label{artin-lemma-formally-smooth-on-deformation-categories}
Soit $S$ un schéma localement noethérien. Soit $F : \mathcal{X} \to \mathcal{Y}$
un $1$-morphisme de catégories fibrées en groupoïdes au-dessus de $(\Sch/S)_{fppf}$.
Supposons que l'une des conditions suivantes soit satisfaite :
\begin{enumerate}
\item $F$ est formellement lisse sur les objets (Critères de représentabilité,
section \ref{criteria-section-formally-smooth}) ;
\item $F$ est représentable par des espaces algébriques et formellement lisse ;
\item $F$ est représentable par des espaces algébriques et lisse.
\end{enumerate}
Alors, pour tout corps $k$ de type fini sur $S$ et tout objet
$x_0$ de $\mathcal{X}$ au-dessus de $k$, le foncteur (\ref{artin-equation-functoriality})
est lisse au sens de
Théorie formelle des déformations, définition
\ref{formal-defos-definition-smooth-morphism}.
\end{lemma}

\begin{proof}
Le cas (1) résulte directement des définitions.
L'hypothèse (2) entraîne (1) d'après
Critères de représentabilité, lemme
\ref{criteria-lemma-representable-by-spaces-formally-smooth}.
L'hypothèse (3) entraîne (2) d'après
Compléments sur les morphismes d'espaces, lemme
\ref{spaces-more-morphisms-lemma-smooth-formally-smooth}
et le principe de
Champs algébriques, lemme
\ref{algebraic-lemma-representable-transformations-property-implication}.
\end{proof}

\begin{lemma}
\label{artin-lemma-fibre-product-deformation-categories}
Soit $S$ un schéma localement noethérien. Soit
$$
\xymatrix{
\mathcal{W} \ar[d] \ar[r] & \mathcal{Z} \ar[d] \\
\mathcal{X} \ar[r] & \mathcal{Y}
}
$$
un $2$-produit fibré de catégories fibrées en groupoïdes au-dessus de
$(\Sch/S)_{fppf}$. Soient $k$ un corps de type fini sur $S$ et
$w_0$ un objet de $\mathcal{W}$ au-dessus de $k$. Notons $x_0, z_0, y_0$
les images de $w_0$ par les morphismes du diagramme. Alors
$$
\xymatrix{
\mathcal{F}_{\mathcal{W}, k, w_0} \ar[d] \ar[r] &
\mathcal{F}_{\mathcal{Z}, k, z_0} \ar[d] \\
\mathcal{F}_{\mathcal{X}, k, x_0} \ar[r] & \mathcal{F}_{\mathcal{Y}, k, y_0}
}
$$
est un produit fibré de catégories de prédéformations.
\end{lemma}

\begin{proof}
Cela résulte directement des définitions. Nous omettons les détails.
\end{proof}






\section{Sommes amalgamées et champs}
\label{artin-section-pushouts}

\noindent
Dans cette section, nous montrons que les champs algébriques se comportent bien
à l'égard de certaines sommes amalgamées. Les résultats de cette section sont valables sur
un schéma de base quelconque.

\medskip\noindent
Le lemme suivant reste vrai lorsque $Y$, $X'$, $X$, $Y'$ sont
des espaces algébriques ; voir (insérer ici une référence future).

\begin{lemma}
\label{artin-lemma-pushout}
\begin{slogan}
Les champs algébriques satisfont la condition (forte) de Rim--Schlessinger
\end{slogan}
Soit $S$ un schéma. Soit
$$
\xymatrix{
X \ar[r] \ar[d] & X' \ar[d] \\
Y \ar[r] & Y'
}
$$
une somme amalgamée dans la catégorie des schémas sur $S$, où $X \to X'$
est un épaississement et $X \to Y$ est affine ; voir
Compléments sur les morphismes, lemme \ref{more-morphisms-lemma-pushout-along-thickening}.
Soit $\mathcal{Z}$ un champ algébrique sur $S$.
Alors le foncteur entre catégories fibres
$$
\mathcal{Z}_{Y'}
\longrightarrow
\mathcal{Z}_Y \times_{\mathcal{Z}_X} \mathcal{Z}_{X'}
$$
est une équivalence de catégories.
\end{lemma}

\begin{proof}
Soit $y'$ un objet du membre de gauche. Le faisceau
$\mathit{Isom}(y', y')$ sur la catégorie des schémas au-dessus de $Y'$
est représentable par un espace algébrique $I$ sur $Y'$ ; voir
Champs algébriques, lemme \ref{algebraic-lemma-representable-diagonal}.
Le foncteur du lemme est donc pleinement fidèle, puisque
$Y'$ est la somme amalgamée dans la catégorie des espaces algébriques aussi bien
que dans celle des schémas ; voir
Sommes amalgamées d'espaces, lemme
\ref{spaces-pushouts-lemma-pushout-along-thickening-schemes}.

\medskip\noindent
Soit $(y, x', f)$ un objet du membre de droite. Ici, $f : y|_X \to x'|_X$
est un isomorphisme. Pour achever la démonstration, nous devons construire un objet $y'$ de
$\mathcal{Z}_{Y'}$ dont les restrictions à $Y$ et à $X'$ s'identifient à $y$ et à $x'$
de manière compatible avec $f$. Il suffit en fait de construire $y'$
localement pour la topologie fppf sur $Y'$ ; voir
Champs, lemme \ref{stacks-lemma-characterize-essentially-surjective-when-ff}.
Choisissons un champ algébrique représentable
$\mathcal{W}$ et un morphisme lisse surjectif $\mathcal{W} \to \mathcal{Z}$.
Alors
$$
(\Sch/Y)_{fppf} \times_{y, \mathcal{Z}} \mathcal{W}
\quad\text{et}\quad
(\Sch/X')_{fppf} \times_{x', \mathcal{Z}} \mathcal{W}
$$
sont des champs algébriques représentables par des espaces algébriques $V$ et $U'$
lisses sur $Y$ et $X'$. L'isomorphisme $f$ induit un isomorphisme
$\varphi : V \times_Y X \to U' \times_{X'} X$ sur $X$. D'après
Sommes amalgamées d'espaces, lemmes
\ref{spaces-pushouts-lemma-pushout-along-thickening} et
\ref{spaces-pushouts-lemma-equivalence-categories-spaces-pushout-flat}
nous voyons que la somme amalgamée $V' = V \amalg_{V \times_Y X} U'$ est
un espace algébrique lisse sur $Y'$ dont les changements de base à
$Y$ et $X'$ redonnent $V$ et $U'$ de manière compatible avec $\varphi$.

\medskip\noindent
Soit $W$ l'espace algébrique qui représente $\mathcal{W}$.
Les projections $V \to W$ et $U' \to W$ coïncident comme morphismes
sur $V \times_Y X \cong U' \times_{X'} X$ ; la propriété universelle
de la somme amalgamée détermine donc un morphisme d'espaces algébriques
$V' \to W$. Choisissons un schéma $Y_1'$ et un morphisme \'etale surjectif
$Y_1' \to V'$. Posons $Y_1 = Y \times_{Y'} Y_1'$,
$X_1' = X' \times_{Y'} Y_1'$ et $X_1 = X \times_{Y'} Y_1'$.
Le composé
$$
(\Sch/Y_1') \to (\Sch/V') \to (\Sch/W) = \mathcal{W} \to \mathcal{Z}
$$
correspond, par le lemme de Yoneda $2$-catégorique, à un objet $y_1'$ de $\mathcal{Z}$
au-dessus de $Y_1'$ dont les restrictions à $Y_1$ et à $X_1'$ s'identifient à $y|_{Y_1}$
et à $x'|_{X_1'}$ de manière compatible avec $f|_{X_1}$. Nous avons ainsi
construit localement pour la topologie lisse sur $Y'$ l'objet voulu, ce qui achève la démonstration.
\end{proof}








\section{La condition de Rim--Schlessinger}
\label{artin-section-RS}

\noindent
La définition suivante est motivée par le
lemme \ref{artin-lemma-pushout}
ainsi que par
Théorie formelle des déformations, définition \ref{formal-defos-definition-RS} et
lemme \ref{formal-defos-lemma-RS-2-categorical}.

\begin{definition}
\label{artin-definition-RS}
Soit $S$ un schéma localement noethérien. Soit $\mathcal{Z}$ une catégorie
fibrée en groupoïdes au-dessus de $(\Sch/S)_{fppf}$. Nous disons que $\mathcal{Z}$
satisfait la {\it condition (RS)} si, pour toute somme amalgamée
$$
\xymatrix{
X \ar[r] \ar[d] & X' \ar[d] \\
Y \ar[r] & Y' = Y \amalg_X X'
}
$$
dans la catégorie des schémas sur $S$, telle que
\begin{enumerate}
\item $X$, $X'$, $Y$, $Y'$ soient les spectres d'anneaux locaux artiniens ;
\item $X$, $X'$, $Y$, $Y'$ soient de type fini sur $S$ ;
\item $X \to X'$ (et donc $Y \to Y'$) soit une immersion fermée,
\end{enumerate}
le foncteur entre catégories fibres
$$
\mathcal{Z}_{Y'}
\longrightarrow
\mathcal{Z}_Y \times_{\mathcal{Z}_X} \mathcal{Z}_{X'}
$$
est une équivalence de catégories.
\end{definition}

\noindent
Si $A$ est un anneau local artinien de corps résiduel $k$, alors
tout morphisme $\Spec(A) \to S$ est affine et de type fini si et
seulement si le morphisme induit $\Spec(k) \to S$ est de type fini ; voir
Morphismes, lemmes \ref{morphisms-lemma-Artinian-affine} et
\ref{morphisms-lemma-artinian-finite-type}.

\begin{lemma}
\label{artin-lemma-algebraic-stack-RS}
Soit $\mathcal{X}$ un champ algébrique sur une base localement noethérienne
$S$. Alors $\mathcal{X}$ satisfait (RS).
\end{lemma}

\begin{proof}
Cela résulte immédiatement des définitions et du lemme \ref{artin-lemma-pushout}.
\end{proof}

\begin{lemma}
\label{artin-lemma-fibre-product-RS}
Soit $S$ un schéma. Soient $p : \mathcal{X} \to \mathcal{Y}$ et
$q : \mathcal{Z} \to \mathcal{Y}$ des $1$-morphismes de catégories
fibrées en groupoïdes au-dessus de $(\Sch/S)_{fppf}$. Si $\mathcal{X}$, $\mathcal{Y}$
et $\mathcal{Z}$ satisfont (RS), alors
$\mathcal{X} \times_\mathcal{Y} \mathcal{Z}$ la satisfait également.
\end{lemma}

\begin{proof}
C'est formel. Soit
$$
\xymatrix{
X \ar[r] \ar[d] & X' \ar[d] \\
Y \ar[r] & Y' = Y \amalg_X X'
}
$$
un diagramme comme dans la définition \ref{artin-definition-RS}. Il faut montrer que
$$
(\mathcal{X} \times_{\mathcal{Y}} \mathcal{Z})_{Y'}
\longrightarrow
(\mathcal{X} \times_{\mathcal{Y}} \mathcal{Z})_Y
\times_{(\mathcal{X} \times_{\mathcal{Y}} \mathcal{Z})_X}
(\mathcal{X} \times_{\mathcal{Y}} \mathcal{Z})_{X'}
$$
est une équivalence. En utilisant la définition du $2$-produit fibré,
cela devient
\begin{equation}
\label{artin-equation-RS-fibre-product}
\mathcal{X}_{Y'} \times_{\mathcal{Y}_{Y'}} \mathcal{Z}_{Y'}
\longrightarrow
(\mathcal{X}_Y \times_{\mathcal{Y}_Y} \mathcal{Z}_Y)
\times_{(\mathcal{X}_X \times_{\mathcal{Y}_X} \mathcal{Z}_X)}
(\mathcal{X}_{X'} \times_{\mathcal{Y}_{X'}} \mathcal{Z}_{X'}).
\end{equation}
Par hypothèse, chacun des foncteurs
$$
\mathcal{X}_{Y'} \to \mathcal{X}_Y \times_{\mathcal{Y}_Y} \mathcal{Z}_Y,
\quad
\mathcal{Y}_{Y'} \to \mathcal{X}_X \times_{\mathcal{Y}_X} \mathcal{Z}_X,
\quad
\mathcal{Z}_{Y'} \to
\mathcal{X}_{X'} \times_{\mathcal{Y}_{X'}} \mathcal{Z}_{X'}
$$
est une équivalence. Un objet du membre de droite de
(\ref{artin-equation-RS-fibre-product}) est un système
$$
((x_Y, z_Y, \phi_Y), (x_{X'}, z_{X'}, \phi_{X'}), (\alpha, \beta)).
$$
Alors $(x_Y, x_{Y'}, \alpha)$ est isomorphe à l'image d'un objet
$x_{Y'}$ de $\mathcal{X}_{Y'}$, et $(z_Y, z_{Y'}, \beta)$ est isomorphe
à l'image d'un objet $z_{Y'}$ de $\mathcal{Z}_{Y'}$. Le couple de
morphismes $(\phi_Y, \phi_{X'})$ correspond à un morphisme $\psi$
entre les images de $x_{Y'}$ et $z_{Y'}$ dans $\mathcal{Y}_{Y'}$.
Alors $(x_{Y'}, z_{Y'}, \psi)$ est un objet du membre de gauche de
(\ref{artin-equation-RS-fibre-product}) qui s'envoie sur l'objet donné du
membre de droite. Cela montre que (\ref{artin-equation-RS-fibre-product}) est
essentiellement surjectif. Nous omettons la preuve de sa pleine fidélité.
\end{proof}





\section{Catégories de déformations}
\label{artin-section-deformation-categories}

\noindent
Rapprochons les notations introduites ci-dessus de celles du
chapitre ``Théorie formelle des déformations''.

\begin{lemma}
\label{artin-lemma-deformation-category}
Soit $S$ un schéma localement noethérien. Soit $\mathcal{X}$ une catégorie
fibrée en groupoïdes au-dessus de $(\Sch/S)_{fppf}$ satisfaisant (RS). Pour tout corps
$k$ de type fini sur $S$ et tout objet $x_0$ de $\mathcal{X}$ au-dessus
de $k$, la catégorie de prédéformations
$p : \mathcal{F}_{\mathcal{X}, k, x_0} \to \mathcal{C}_\Lambda$
(\ref{artin-equation-predeformation-category}) est une catégorie de déformations ; voir
Théorie formelle des déformations, définition
\ref{formal-defos-definition-deformation-category}.
\end{lemma}

\begin{proof}
Posons $\mathcal{F} = \mathcal{F}_{\mathcal{X}, k, x_0}$.
Soient $f_1 : A_1 \to A$ et $f_2 : A_2 \to A$ des homomorphismes d'anneaux dans
$\mathcal{C}_\Lambda$, avec $f_2$ surjectif. Il faut montrer que
le foncteur
$$
\mathcal{F}(A_1 \times_A A_2)
\longrightarrow
\mathcal{F}(A_1) \times_{\mathcal{F}(A)} \mathcal{F}(A_2)
$$
est une équivalence ; voir
Théorie formelle des déformations, lemme \ref{formal-defos-lemma-RS-2-categorical}.
Posons $X = \Spec(A)$, $X' = \Spec(A_2)$, $Y = \Spec(A_1)$ et
$Y' = \Spec(A_1 \times_A A_2)$. Notons que $Y' = Y \amalg_X X'$ dans la
catégorie des schémas ; voir
Compléments sur les morphismes, lemme \ref{more-morphisms-lemma-pushout-along-thickening}.
Dans le diagramme de foncteurs entre catégories fibres
$$
\xymatrix{
\mathcal{X}_{Y'} \ar[r] \ar[d] &
\mathcal{X}_Y \times_{\mathcal{X}_X} \mathcal{X}_{X'} \ar[d] \\
\mathcal{X}_{\Spec(k)} \ar@{=}[r] & \mathcal{X}_{\Spec(k)}
}
$$
la flèche horizontale supérieure est une équivalence d'après la
définition \ref{artin-definition-RS}.
Comme $\mathcal{F}(B)$ est la catégorie des objets de $\mathcal{X}_{\Spec(B)}$
munis d'une identification à $x_0$ au-dessus de $k$, le résultat en découle.
\end{proof}

\begin{remark}
\label{artin-remark-deformation-category-implies}
Soit $S$ un schéma localement noethérien. Soit $\mathcal{X}$ une catégorie fibrée
en groupoïdes au-dessus de $(\Sch/S)_{fppf}$. Soient $k$ un corps de type fini sur
$S$ et $x_0$ un objet
de $\mathcal{X}$ au-dessus de $k$. Soit $p : \mathcal{F} \to \mathcal{C}_\Lambda$
comme en (\ref{artin-equation-predeformation-category}). Si $\mathcal{F}$
est une catégorie de déformations, c'est-à-dire si $\mathcal{F}$ satisfait la
condition (RS) de Rim--Schlessinger, alors $\mathcal{F}$ satisfait
les conditions (S1) et (S2) de Schlessinger d'après
Théorie formelle des déformations, lemme \ref{formal-defos-lemma-RS-implies-S1-S2}.
Soit $\overline{\mathcal{F}}$ le foncteur des classes d'isomorphisme ; voir
Théorie formelle des déformations, remarques
\ref{formal-defos-remarks-cofibered-groupoids}
(\ref{formal-defos-item-associated-functor-isomorphism-classes}).
Alors $\overline{\mathcal{F}}$ satisfait également (S1) et (S2) ; voir
Théorie formelle des déformations, lemme
\ref{formal-defos-lemma-S1-S2-associated-functor}.
C'est en particulier le cas dans la situation du
lemme \ref{artin-lemma-deformation-category}.
\end{remark}




\section{Changement de corps}
\label{artin-section-change-of-field}

\noindent
Cette section est l'analogue de
Théorie formelle des déformations, section \ref{formal-defos-section-change-of-field}.
Comme cela y est signalé, pour étudier le comportement par changement de corps,
nous devons écrire $\mathcal{C}_{\Lambda, k}$ plutôt que $\mathcal{C}_\Lambda$.
Dans le lemme suivant, nous utilisons la notation $\mathcal{F}_{l/k}$
introduite dans Théorie formelle des déformations, situation
\ref{formal-defos-situation-change-of-fields}.

\begin{lemma}
\label{artin-lemma-change-of-field}
Soit $S$ un schéma localement noethérien. Soit $\mathcal{X}$ une catégorie
fibrée en groupoïdes sur $(\Sch/S)_{fppf}$. Soit $k$ un
corps de type fini sur $S$ et soit $l/k$ une extension finie.
Soit $x_0$ un objet de $\mathcal{F}$ au-dessus de $\Spec(k)$.
Notons $x_{l, 0}$ la restriction de $x_0$ à $\Spec(l)$.
Alors il existe un foncteur canonique
$$
(\mathcal{F}_{\mathcal{X}, k , x_0})_{l/k}
\longrightarrow
\mathcal{F}_{\mathcal{X}, l, x_{l, 0}}
$$
de catégories cofibrées en groupoïdes sur $\mathcal{C}_{\Lambda, l}$.
Si $\mathcal{X}$ satisfait à (RS), alors ce foncteur est une équivalence.
\end{lemma}

\begin{proof}
Considérons une factorisation
$$
\Spec(l) \to \Spec(B) \to S
$$
comme dans (\ref{artin-equation-factor}). Par définition, nous avons
$$
(\mathcal{F}_{\mathcal{X}, k , x_0})_{l/k}(B) =
\mathcal{F}_{\mathcal{X}, k, x_0}(B \times_l k)
$$
voir Théorie formelle des déformations, situation
\ref{formal-defos-situation-change-of-fields}. Ainsi, un objet de cette catégorie
est un morphisme $x_0 \to x$ de $\mathcal{X}$ au-dessus du morphisme
$\Spec(k) \to \Spec(B \times_l k)$. En choisissant un foncteur image réciproque pour $\mathcal{X}$,
on peut associer à $x_0 \to x$ le morphisme $x_{l, 0} \to x_B$,
où $x_B$ est la restriction de $x$ à $\Spec(B)$ (par le morphisme
$\Spec(B) \to \Spec(B \times_l k)$ provenant de $B \times_l k \subset B$).
Cette construction est fonctorielle en $B$ et compatible avec les morphismes.

\medskip\noindent
Supposons maintenant que $\mathcal{X}$ satisfait à (RS). Considérons les diagrammes
$$
\vcenter{
\xymatrix{
l & B \ar[l] \\
k \ar[u] & B \times_l k \ar[l] \ar[u]
}
}
\quad\text{et}\quad
\vcenter{
\xymatrix{
\Spec(l) \ar[d] \ar[r] & \Spec(B) \ar[d] \\
\Spec(k) \ar[r] & \Spec(B \times_l k)
}
}
$$
Le diagramme de gauche est un produit fibré d'anneaux. Celui de
droite est une somme amalgamée dans la catégorie des schémas, voir
Compléments sur les morphismes, lemme \ref{more-morphisms-lemma-pushout-along-thickening}.
Tous ces schémas sont de type fini sur $S$ (voir les remarques suivant la
définition \ref{artin-definition-RS}). Ainsi, (RS) s'applique et donne une équivalence
de catégories fibres
$$
\mathcal{X}_{\Spec(B \times_l k)}
\longrightarrow
\mathcal{X}_{\Spec(k)}
\times_{\mathcal{X}_{\Spec(l)}}
\mathcal{X}_{\Spec(B)}
$$
Cela entraîne que le foncteur défini ci-dessus donne une équivalence de
catégories fibres. Par conséquent, le foncteur est une équivalence de catégories
cofibrées en groupoïdes par (le dual de)
Catégories, lemme \ref{categories-lemma-equivalence-fibred-categories}.
\end{proof}








\section{Espaces tangents}
\label{artin-section-tangent-spaces}

\noindent
Soit $S$ un schéma localement noethérien. Soit $\mathcal{X}$ une catégorie
fibrée en groupoïdes sur $(\Sch/S)_{fppf}$. Soit $k$ un corps de type
fini sur $S$ et soit $x_0$ un objet de $\mathcal{X}$ au-dessus de $k$.
Dans Théorie formelle des déformations, section \ref{formal-defos-section-tangent-spaces},
nous avons défini l'{\it espace tangent}
\begin{equation}
\label{artin-equation-tangent-space}
T\mathcal{F}_{\mathcal{X}, k, x_0} =
\left\{
\begin{matrix}
\text{classes d'isomorphisme de morphismes}\\
x_0 \to x\text{ au-dessus de }\Spec(k) \to \Spec(k[\epsilon])
\end{matrix}
\right\}
\end{equation}
de la catégorie de prédéformations $\mathcal{F}_{\mathcal{X}, k, x_0}$.
Dans Théorie formelle des déformations, section
\ref{formal-defos-section-infinitesimal-automorphisms},
nous avons défini
\begin{equation}
\label{artin-equation-infinitesimal-automorphisms}
\text{Inf}(\mathcal{F}_{\mathcal{X}, k, x_0}) =
\Ker\left(
\text{Aut}_{\Spec(k[\epsilon])}(x'_0) \to \text{Aut}_{\Spec(k)}(x_0)
\right)
\end{equation}
où $x_0'$ est l'image réciproque de $x_0$ sur $\Spec(k[\epsilon])$.
Si $\mathcal{X}$ satisfait à la condition de Rim-Schlessinger (RS), alors
$T\mathcal{F}_{\mathcal{X}, k, x_0}$ est naturellement muni d'une
structure d'espace vectoriel sur $k$ par Théorie formelle des déformations, lemme
\ref{formal-defos-lemma-tangent-space-vector-space}
(les hypothèses sont vérifiées d'après le lemme \ref{artin-lemma-deformation-category} et la
remarque \ref{artin-remark-deformation-category-implies}). De plus,
Théorie formelle des déformations, lemme \ref{formal-defos-lemma-infaut-vector-space},
montre que $\text{Inf}(\mathcal{F}_{\mathcal{X}, k, x_0})$ possède une
structure naturelle d'espace vectoriel sur $k$ telle que l'addition correspond à la
composition des automorphismes. Il est naturel
de demander que ces espaces vectoriels soient de dimension finie.

\medskip\noindent
Le lemme suivant affirme que cela est vrai si
$\mathcal{X}$ est localement de type fini sur $S$ (voir
Morphismes de champs, section \ref{stacks-morphisms-section-finite-type}).

\begin{lemma}
\label{artin-lemma-finite-dimension}
Soit $S$ un schéma localement noethérien. Supposons que
\begin{enumerate}
\item $\mathcal{X}$ est un champ algébrique,
\item $U$ est un schéma localement de type fini sur $S$, et
\item $(\Sch/U)_{fppf} \to \mathcal{X}$ est un morphisme lisse et
surjectif.
\end{enumerate}
Alors, pour tout $\mathcal{F} = \mathcal{F}_{\mathcal{X}, k, x_0}$ comme dans la
section \ref{artin-section-predeformation-categories},
l'espace tangent $T\mathcal{F}$ et l'espace des automorphismes infinitésimaux
$\text{Inf}(\mathcal{F})$ sont de dimension finie sur $k$.
\end{lemma}

\begin{proof}
Posons $\mathcal{U} = (\Sch/U)_{fppf}$. Selon notre définition
des champs algébriques, le $1$-morphisme $\mathcal{U} \to \mathcal{X}$
est représentable par des espaces algébriques. En particulier, le
2-produit fibré
$$
\mathcal{U}_{x_0} = (\Sch/\Spec(k))_{fppf} \times_\mathcal{X} \mathcal{U}
$$
est représenté par un espace algébrique $U_{x_0}$ sur $\Spec(k)$. Alors
$U_{x_0} \to \Spec(k)$ est lisse et surjectif (en particulier, $U_{x_0}$
est non vide). D'après Espaces sur les corps, lemme
\ref{spaces-over-fields-lemma-smooth-separable-closed-points-dense},
il existe une extension finie $l/k$ et un point
$\Spec(l) \to U_{x_0}$ au-dessus de $k$. Nous avons
$$
(\mathcal{F}_{\mathcal{X}, k , x_0})_{l/k} =
\mathcal{F}_{\mathcal{X}, l, x_{l, 0}}
$$
d'après le lemme \ref{artin-lemma-change-of-field} et le fait que $\mathcal{X}$
satisfait à (RS). Nous voyons donc que
$$
T\mathcal{F} \otimes_k l \cong T\mathcal{F}_{\mathcal{X}, l, x_{l, 0}}
\quad\text{et}\quad
\text{Inf}(\mathcal{F}) \otimes_k l \cong
\text{Inf}(\mathcal{F}_{\mathcal{X}, l, x_{l, 0}})
$$
d'après
Théorie formelle des déformations, lemmes
\ref{formal-defos-lemma-tangent-space-change-of-field} et
\ref{formal-defos-lemma-inf-aut-change-of-field}
(ces lemmes sont applicables d'après les
lemmes \ref{artin-lemma-algebraic-stack-RS} et
\ref{artin-lemma-deformation-category} et la
remarque \ref{artin-remark-deformation-category-implies}).
Il suffit donc de démontrer que $T\mathcal{F}_{\mathcal{X}, l, x_{l, 0}}$
et $\text{Inf}(\mathcal{F}_{\mathcal{X}, l, x_{l, 0}})$
sont de dimension finie sur $l$. Notons que $x_{l, 0}$ provient d'un point
$u_0$ de $\mathcal{U}$ au-dessus de $l$.

\medskip\noindent
Interrompons le cours de l'argument pour montrer que le lemme concernant les
automorphismes infinitésimaux se déduit du lemme concernant les espaces tangents.
En effet, posons
$\mathcal{R} = \mathcal{U} \times_\mathcal{X} \mathcal{U}$.
Soit $r_0$ le point à valeurs dans $l$, $(u_0, u_0, \text{id}_{x_0})$, de
$\mathcal{R}$. En combinant le
lemme \ref{artin-lemma-fibre-product-deformation-categories} et
Théorie formelle des déformations, lemme
\ref{formal-defos-lemma-deformation-functor-diagonal},
nous obtenons
$$
\text{Inf}(\mathcal{F}_{\mathcal{X}, l, x_{l, 0}})
\subset
T\mathcal{F}_{\mathcal{R}, l, r_0}
$$
Notons que $\mathcal{R}$ est un champ algébrique, voir
Champs algébriques, lemme \ref{algebraic-lemma-2-fibre-product-general}.
De plus, $\mathcal{R}$ est représenté par un espace algébrique $R$
lisse sur $U$ (par l'une ou l'autre projection, voir
Champs algébriques, lemme \ref{algebraic-lemma-stack-presentation}).
Par conséquent, choisissons un schéma $U'$ et un morphisme \'etale surjectif
$U' \to R$ ; nous voyons que $U'$ est lisse sur $U$, donc localement de
type fini sur $S$. Comme $(\Sch/U')_{fppf} \to \mathcal{R}$ est
surjectif et lisse, nous avons ramené la question au cas
des espaces tangents.

\medskip\noindent
Le foncteur (\ref{artin-equation-functoriality})
$$
\mathcal{F}_{\mathcal{U}, l, u_0}
\longrightarrow
\mathcal{F}_{\mathcal{X}, l, x_{l, 0}}
$$
est lisse d'après le lemme \ref{artin-lemma-formally-smooth-on-deformation-categories}.
L'application induite sur les espaces tangents
$$
T\mathcal{F}_{\mathcal{U}, l, u_0}
\longrightarrow
T\mathcal{F}_{\mathcal{X}, l, x_{l, 0}}
$$
est $l$-linéaire (d'après
Théorie formelle des déformations, lemme
\ref{formal-defos-lemma-k-linear-differential})
et surjective (puisque les morphismes lisses de catégories de prédéformations induisent
des applications surjectives sur les espaces tangents d'après
Théorie formelle des déformations, lemme
\ref{formal-defos-lemma-smooth-morphism-essentially-surjective}).
Il suffit donc de prouver que l'espace tangent de l'espace de déformations
associé au champ algébrique représentable $\mathcal{U}$
au point $u_0$ est de dimension finie. Soit $\Spec(R) \subset U$ un ouvert affine
tel que $u_0 : \Spec(l) \to U$ se factorise par $\Spec(R)$
et tel que $\Spec(R) \to S$ se factorise par $\Spec(\Lambda) \subset S$.
Soit $\mathfrak m_R \subset R$ le noyau du morphisme de $\Lambda$-algèbres
$\varphi_0 : R \to l$ correspondant à $u_0$. Comme $R$ est de type
fini sur l'anneau noethérien $\Lambda$, c'est un anneau noethérien. Ainsi,
$\mathfrak m_R = (f_1, \ldots, f_n)$ est un idéal de type fini.
Nous avons
$$
T\mathcal{F}_{\mathcal{U}, l, u_0}
=
\{\varphi : R \to l[\epsilon] \mid
\varphi \text{ est un morphisme de } \Lambda\text{-algèbres et }
\varphi \bmod \epsilon = \varphi_0\}
$$
Un élément du membre de droite est déterminé par ses valeurs sur
$f_1, \ldots, f_n$ ; le membre de droite est donc de dimension au plus $n$, ce qui conclut.
Nous omettons quelques détails.
\end{proof}

\begin{lemma}
\label{artin-lemma-fibre-product-tangent-spaces}
Soit $S$ un schéma localement noethérien. Soient $p : \mathcal{X} \to \mathcal{Y}$
et $q : \mathcal{Z} \to \mathcal{Y}$ des $1$-morphismes de catégories
fibrées en groupoïdes sur $(\Sch/S)_{fppf}$. Supposons que $\mathcal{X}$,
$\mathcal{Y}$ et $\mathcal{Z}$ satisfassent à (RS).
Soit $k$ un corps de type fini sur $S$ et soit $w_0$ un objet de
$\mathcal{W} = \mathcal{X} \times_\mathcal{Y} \mathcal{Z}$ au-dessus de $k$.
Notons $x_0, y_0, z_0$ les objets de $\mathcal{X}, \mathcal{Y}, \mathcal{Z}$
déduits de $w_0$. Alors il existe une suite exacte à $6$ termes
$$
\xymatrix{
0 \ar[r] &
\text{Inf}(\mathcal{F}_{\mathcal{W}, k, w_0}) \ar[r] &
\text{Inf}(\mathcal{F}_{\mathcal{X}, k, x_0}) \oplus
\text{Inf}(\mathcal{F}_{\mathcal{Z}, k, z_0}) \ar[r] &
\text{Inf}(\mathcal{F}_{\mathcal{Y}, k, y_0}) \ar[lld] \\
 &
T\mathcal{F}_{\mathcal{W}, k, w_0} \ar[r] &
T\mathcal{F}_{\mathcal{X}, k, x_0} \oplus
T\mathcal{F}_{\mathcal{Z}, k, z_0} \ar[r] &
T\mathcal{F}_{\mathcal{Y}, k, y_0}
}
$$
d'espaces vectoriels sur $k$.
\end{lemma}

\begin{proof}
D'après le lemme \ref{artin-lemma-fibre-product-RS}, $\mathcal{W}$
satisfait à (RS), de sorte que l'énoncé du lemme a un sens. Pour démontrer
le lemme, appliquons les lemmes \ref{artin-lemma-fibre-product-deformation-categories} et
\ref{artin-lemma-deformation-category},
ainsi que Théorie formelle des déformations, lemme
\ref{formal-defos-lemma-deformation-categories-fiber-product-morphisms}.
\end{proof}






\section{Objets formels}
\label{artin-section-formal-objects}

\noindent
Dans cette section, nous transposons au présent cadre certaines des notions déjà définies
dans le chapitre « Théorie formelle des déformations ». Dans la suite, nous dirons que
« $R$ est une $S$-algèbre » pour signifier
que $R$ est un anneau muni d'un morphisme de schémas $\Spec(R) \to S$.

\begin{definition}
\label{artin-definition-formal-objects}
Soit $S$ un schéma localement noethérien. Soit
$p : \mathcal{X} \to (\Sch/S)_{fppf}$ une catégorie fibrée en groupoïdes.
\begin{enumerate}
\item Un {\it objet formel} $\xi = (R, \xi_n, f_n)$ de $\mathcal{X}$ est constitué
d'une $S$-algèbre locale noethérienne complète $R$, d'objets $\xi_n$ de
$\mathcal{X}$ au-dessus de $\Spec(R/\mathfrak m_R^n)$, et de morphismes
$f_n : \xi_n \to \xi_{n + 1}$ de $\mathcal{X}$ au-dessus de
$\Spec(R/\mathfrak m^n) \to \Spec(R/\mathfrak m^{n + 1})$
tels que $R/\mathfrak m$ soit un corps de type fini sur $S$.
\item Un {\it morphisme d'objets formels}
$a : \xi = (R, \xi_n, f_n) \to \eta = (T, \eta_n, g_n)$
est la donnée de morphismes $a_n : \xi_n \to \eta_n$ tels que, pour tout $n$,
le diagramme
$$
\xymatrix{
\xi_n \ar[r]_{f_n} \ar[d]_{a_n} & \xi_{n + 1} \ar[d]^{a_{n + 1}} \\
\eta_n \ar[r]^{g_n} & \eta_{n + 1}
}
$$
soit commutatif. En appliquant le foncteur $p$, on obtient une famille compatible
de morphismes $\Spec(R/\mathfrak m_R^n) \to \Spec(T/\mathfrak m_T^n)$, donc
un morphisme $a_0 : \Spec(R) \to \Spec(T)$ au-dessus de $S$. On dit que
$a$ {\it est au-dessus de} $a_0$.
\end{enumerate}
\end{definition}

\noindent
On obtient ainsi une catégorie d'objets formels de $\mathcal{X}$.

\begin{remark}
\label{artin-remark-formal-objects-match}
Soit $S$ un schéma localement noethérien. Soit
$p : \mathcal{X} \to (\Sch/S)_{fppf}$ une catégorie fibrée en groupoïdes.
Soit $\xi = (R, \xi_n, f_n)$ un objet formel. Posons $k = R/\mathfrak m$ et
$x_0 = \xi_1$. L'objet formel $\xi$ définit un objet formel
$\xi$ de la catégorie de prédéformations $\mathcal{F}_{\mathcal{X}, k, x_0}$.
Cela résulte immédiatement de la
Définition \ref{artin-definition-formal-objects} ci-dessus, de
Théorie formelle des déformations, Définition
\ref{formal-defos-definition-formal-objects},
et de notre construction de la catégorie de prédéformations
$\mathcal{F}_{\mathcal{X}, k, x_0}$ dans la
section \ref{artin-section-predeformation-categories}.
\end{remark}

\noindent
Si $F : \mathcal{X} \to \mathcal{Y}$ est un $1$-morphisme de catégories fibrées
en groupoïdes sur $(\Sch/S)_{fppf}$, alors $F$ induit aussi un foncteur entre
les catégories d'objets formels.

\begin{lemma}
\label{artin-lemma-smooth-lift-formal}
Soit $S$ un schéma localement noethérien. Soit $F : \mathcal{X} \to \mathcal{Y}$
un $1$-morphisme de catégories fibrées en groupoïdes sur $(\Sch/S)_{fppf}$.
Soit $\eta = (R, \eta_n, g_n)$ un objet formel de $\mathcal{Y}$
et soit $\xi_1$ un objet de $\mathcal{X}$ tel que $F(\xi_1) \cong \eta_1$.
Si $F$ est formellement lisse sur les objets (voir
Critères de représentabilité, section \ref{criteria-section-formally-smooth}),
alors il existe un objet formel $\xi = (R, \xi_n, f_n)$ de $\mathcal{X}$
tel que $F(\xi) \cong \eta$.
\end{lemma}

\begin{proof}
Notons que chacun des morphismes
$\Spec(R/\mathfrak m^n) \to \Spec(R/\mathfrak m^{n + 1})$ est un épaississement
du premier ordre de schémas affines sur $S$. L'hypothèse sur $F$ signifie donc
que l'on peut relever successivement $\xi_1$ en des objets $\xi_2, \xi_3, \ldots$
de $\mathcal{X}$ munis d'isomorphismes compatibles
$\eta_n|_{\Spec(R/\mathfrak m^{n - 1})} \cong \eta_{n - 1}$
et $F(\eta_n) \cong \xi_n$.
\end{proof}

\noindent
Soit $S$ un schéma localement noethérien. Soit
$p : \mathcal{X} \to (\Sch/S)_{fppf}$ une catégorie fibrée en groupoïdes.
Supposons que $x$ soit un objet de $\mathcal{X}$ au-dessus de $R$, où $R$ est une
$S$-algèbre locale noethérienne complète dont le corps résiduel est de type fini
sur $S$. On peut alors considérer le système des restrictions
$\xi_n = x|_{\Spec(R/\mathfrak m^n)}$ muni des morphismes naturels
$\xi_1 \to \xi_2 \to \ldots$ provenant de la transitivité de la restriction.
Ainsi, $\xi = (R, \xi_n, \xi_n \to \xi_{n + 1})$ est un objet formel de
$\mathcal{X}$. Cette construction est fonctorielle en l'objet $x$.
On obtient donc un foncteur
\begin{equation}
\label{artin-equation-approximation}
\left\{
\begin{matrix}
\text{objets }x\text{ de }\mathcal{X} \text{ tels que }p(x) = \Spec(R) \\
\text{où }R\text{ est un anneau local noethérien complet}\\
\text{et }R/\mathfrak m\text{ est de type fini sur }S
\end{matrix}
\right\}
\longrightarrow
\left\{
\begin{matrix}
\text{objets formels de }\mathcal{X}
\end{matrix}
\right\}
\end{equation}
Plus précisément, le membre de gauche est la sous-catégorie pleine de $\mathcal{X}$
constituée des objets indiqués, et le membre de droite est la catégorie
des objets formels de $\mathcal{X}$ au sens de la
Définition \ref{artin-definition-formal-objects}.

\begin{definition}
\label{artin-definition-effective}
Soit $S$ un schéma localement noethérien. Soit $\mathcal{X}$ une catégorie
fibrée en groupoïdes sur $(\Sch/S)_{fppf}$. Un objet formel
$\xi = (R, \xi_n, f_n)$ de $\mathcal{X}$ est dit {\it effectif}
s'il appartient à l'image essentielle du foncteur
(\ref{artin-equation-approximation}).
\end{definition}

\noindent
Si la catégorie fibrée en groupoïdes est un champ algébrique, alors tout
objet formel est effectif, comme il résulte du lemme suivant.

\begin{lemma}
\label{artin-lemma-effective}
Soit $S$ un schéma localement noethérien. Soit $\mathcal{X}$ un champ algébrique
sur $S$. Le foncteur (\ref{artin-equation-approximation}) est une équivalence.
\end{lemma}

\begin{proof}
Cas I : $\mathcal{X}$ est représentable (par un schéma). Écrivons
$\mathcal{X} = (\Sch/X)_{fppf}$ pour un schéma $X$ sur $S$.
En explicitant les définitions, il faut démontrer ceci : étant donnée
une $S$-algèbre locale noethérienne complète $R$ telle que $R/\mathfrak m$ soit
de type fini sur $S$, l'application
$$
\Mor_S(\Spec(R), X) \longrightarrow \lim \Mor_S(\Spec(R/\mathfrak m^n), X)
$$
est bijective. Cela résulte de Espaces formels, Lemme
\ref{formal-spaces-lemma-map-into-scheme}.

\medskip\noindent
Cas II. $\mathcal{X}$ est représentable par un espace algébrique. Écrivons
que $\mathcal{X}$ est représentable par $X$. Il faut à nouveau montrer que
$$
\Mor_S(\Spec(R), X) \longrightarrow \lim \Mor_S(\Spec(R/\mathfrak m^n), X)
$$
est bijective pour $R$ comme ci-dessus. C'est Espaces formels, Lemme
\ref{formal-spaces-lemma-map-into-algebraic-space}.

\medskip\noindent
Cas III : cas général d'un champ algébrique. Remarquons de manière générale que
les membres de gauche et de droite de (\ref{artin-equation-approximation}) sont des
catégories fibrées en groupoïdes sur la catégorie des schémas affines
sur $S$ qui sont les spectres d'anneaux locaux noethériens complets
dont le corps résiduel est de type fini sur $S$. Nous verrons aussi dans la
démonstration ci-dessous qu'elles sont des champs pour une certaine topologie sur cette
catégorie.

\medskip\noindent
Démontrons d'abord la pleine fidélité. Soit $R$ une algèbre locale noethérienne complète
sur $S$ telle que $k = R/\mathfrak m$ soit de type fini sur $S$.
Soient $x, x'$ des objets de $\mathcal{X}$ au-dessus de $R$. Comme $\mathcal{X}$ est
un champ algébrique, $\mathit{Isom}(x, x')$ est représentable par un
espace algébrique $I$ sur $\Spec(R)$, voir
Champs algébriques, Lemme \ref{algebraic-lemma-representable-diagonal}.
L'application du Cas II à $I$ sur $\Spec(R)$ implique immédiatement que
(\ref{artin-equation-approximation}) est pleinement fidèle sur les catégories fibres au-dessus de
$\Spec(R)$. Le foncteur est donc pleinement fidèle d'après
Catégories, Lemme \ref{categories-lemma-equivalence-fibred-categories}.

\medskip\noindent
Surjectivité essentielle. Soit $\xi = (R, \xi_n, f_n)$ un objet formel de
$\mathcal{X}$. Choisissons un schéma $U$ sur $S$ et un morphisme lisse surjectif
$f : (\Sch/U)_{fppf} \to \mathcal{X}$. Pour tout $n$, considérons le produit fibré
$$
(\Sch/\Spec(R/\mathfrak m^n))_{fppf}
\times_{\xi_n, \mathcal{X}, f}
(\Sch/U)_{fppf}
$$
Par hypothèse, il est représentable par un espace algébrique $V_n$, lisse et
surjectif sur $\Spec(R/\mathfrak m^n)$. Les morphismes
$f_n : \xi_n \to \xi_{n + 1}$ induisent des carrés cartésiens
$$
\xymatrix{
V_{n + 1} \ar[d] & V_n \ar[d] \ar[l] \\
\Spec(R/\mathfrak m^{n + 1}) & \Spec(R/\mathfrak m^n) \ar[l]
}
$$
d'espaces algébriques. D'après Espaces sur les corps, Lemme
\ref{spaces-over-fields-lemma-smooth-separable-closed-points-dense},
on peut trouver une extension finie séparable $k'/k$ et un point
$v'_1 : \Spec(k') \to V_1$ au-dessus de $k$. Soit $R \subset R'$ l'extension finie \'etale
dont l'extension de corps résiduels est $k'/k$ (elle existe et
est unique d'après
Algèbre, Lemmes \ref{algebra-lemma-henselian-cat-finite-etale} et
\ref{algebra-lemma-complete-henselian}).
Par le critère infinitésimal de lissité (voir
Compléments sur les morphismes d'espaces, Lemme
\ref{spaces-more-morphisms-lemma-smooth-formally-smooth}),
appliqué à $V_n \to \Spec(R/\mathfrak m^n)$ pour $n = 2, 3, 4, \ldots$,
on peut trouver successivement des morphismes
$v'_n : \Spec(R'/(\mathfrak m')^n) \to V_n$ au-dessus de $\Spec(R/\mathfrak m^n)$
s'inscrivant dans des diagrammes commutatifs
$$
\xymatrix{
\Spec(R'/(\mathfrak m')^{n + 1}) \ar[d]_{v'_{n + 1}} &
\Spec(R'/(\mathfrak m')^n) \ar[d]^{v'_n} \ar[l] \\
V_{n + 1} & V_n \ar[l]
}
$$
En composant avec les morphismes de projection $V_n \to U$, on obtient un système
compatible de morphismes $u'_n : \Spec(R'/(\mathfrak m')^n) \to U$.
D'après le Cas I, la famille $(u'_n)$ provient d'un unique
morphisme $u' : \Spec(R') \to U$. Notons $x'$ l'objet de $\mathcal{X}$
au-dessus de $\Spec(R')$ que l'on obtient en appliquant le $1$-morphisme $f$ à $u'$.
Par construction, il existe un morphisme d'objets formels
$$
(\ref{artin-equation-approximation})(x') =
(R', x'|_{\Spec(R'/(\mathfrak m')^n)}, \ldots)
\longrightarrow
(R, \xi_n, f_n)
$$
au-dessus de $\Spec(R') \to \Spec(R)$. Notons que $R' \otimes_R R'$ est un produit fini
de spectres d'anneaux locaux noethériens complets auxquels s'applique notre présente
discussion. Notons $p_0, p_1 : \Spec(R' \otimes_R R') \to \Spec(R')$
les deux projections. Par la pleine fidélité démontrée ci-dessus, il existe
un isomorphisme canonique $\varphi : p_0^*x' \to p_1^*x'$, puisque l'on dispose de
tels isomorphismes au-dessus de
$\Spec((R' \otimes_R R')/\mathfrak m^n(R' \otimes_R R'))$.
Nous omettons de démontrer que l'isomorphisme $\varphi$ satisfait la condition de cocycle
(voir Champs, Définition \ref{stacks-definition-descent-data}).
Puisque $\{\Spec(R') \to \Spec(R)\}$ est un recouvrement fppf, on conclut
que $x'$ provient par descente d'un objet $x$ de $\mathcal{X}$ au-dessus de $\Spec(R)$.
Nous omettons de démontrer que $x_n$ est la restriction de $x$ à
$\Spec(R/\mathfrak m^n)$.
\end{proof}

\begin{lemma}
\label{artin-lemma-fibre-product-effective}
Soit $S$ un schéma. Soient $p : \mathcal{X} \to \mathcal{Y}$ et
$q : \mathcal{Z} \to \mathcal{Y}$ des $1$-morphismes de catégories
fibrées en groupoïdes sur $(\Sch/S)_{fppf}$. Si le foncteur
(\ref{artin-equation-approximation}) est une équivalence pour 
$\mathcal{X}$, $\mathcal{Y}$ et $\mathcal{Z}$, alors il est 
une équivalence pour $\mathcal{X} \times_\mathcal{Y} \mathcal{Z}$.
\end{lemma}

\begin{proof}
Les membres de gauche et de droite de (\ref{artin-equation-approximation})
pour $\mathcal{X} \times_\mathcal{Y} \mathcal{Z}$ sont simplement les $2$-produits
fibrés des membres de gauche et de droite de (\ref{artin-equation-approximation})
pour $\mathcal{X}$ et $\mathcal{Z}$ au-dessus de $\mathcal{Y}$.
Le résultat en découle, car la formation des $2$-produits fibrés est compatible
avec les équivalences de catégories, voir
Catégories, Lemme \ref{categories-lemma-equivalence-2-fibre-product}.
\end{proof}






\section{Approximation}
\label{artin-section-approximation}

\noindent
Une idée fondamentale due à Michael Artin est que l'on peut approcher les
objets d'un champ commutant aux limites. Plus précisément, étant donné un objet $x$
du champ sur un anneau local noethérien complet, on peut trouver
un objet $x_A$ sur un anneau algébrique qui soit « proche de » $x$.
Ici, un anneau algébrique désigne une $S$-algèbre de type fini, et proche
signifie proche pour la topologie adique. Nous présentons dans cette section une forme simple,
mais générale, de ce résultat.

\medskip\noindent
Pour formuler le résultat, il faut réunir quelques définitions provenant de
différents endroits du projet Champs. Tout d'abord, dans
Critères de représentabilité, section \ref{criteria-section-limit-preserving},
nous avons introduit la notion de {\it commutation aux limites sur les objets} pour les $1$-morphismes
de catégories fibrées en groupoïdes sur la catégorie des schémas.
Dans Compléments d'algèbre, Définition \ref{more-algebra-definition-G-ring},
nous avons défini la notion d'{\it anneau G}. Soit $S$ un schéma localement noethérien.
Soit $A$ une $S$-algèbre. On dit que $A$ est {\it de type fini sur $S$}
ou que c'est une {\it $S$-algèbre de type fini} si $\Spec(A) \to S$ est de type fini.
Dans ce cas, $A$ est un anneau noethérien. Enfin, étant donnés un anneau $A$ et un idéal
$I$, on note $\text{Gr}_I(A) = \bigoplus I^n/I^{n + 1}$.

\begin{lemma}
\label{artin-lemma-approximate}
Soit $S$ un schéma localement noethérien. Soit
$p : \mathcal{X} \to (\Sch/S)_{fppf}$ une catégorie
fibrée en groupoïdes. Soit $x$ un objet de
$\mathcal{X}$ au-dessus de $\Spec(R)$, où $R$ est un anneau local noethérien complet
de corps résiduel $k$ de type fini sur $S$. Soit $s \in S$
l'image de $\Spec(k) \to S$. Supposons que (a) $\mathcal{O}_{S, s}$ soit
un anneau G et (b) $p$ commute aux limites sur les objets. Alors, pour tout
entier $N \geq 1$, il existe
\begin{enumerate}
\item une $S$-algèbre de type fini $A$,
\item un idéal maximal $\mathfrak m_A \subset A$,
\item un objet $x_A$ de $\mathcal{X}$ au-dessus de $\Spec(A)$,
\item un $S$-isomorphisme $R/\mathfrak m_R^N \cong A/\mathfrak m_A^N$,
\item un isomorphisme
$x|_{\Spec(R/\mathfrak m_R^N)} \cong x_A|_{\Spec(A/\mathfrak m_A^N)}$
compatible avec (4), et
\item un isomorphisme
$\text{Gr}_{\mathfrak m_R}(R) \cong \text{Gr}_{\mathfrak m_A}(A)$
de $k$-algèbres graduées.
\end{enumerate}
\end{lemma}

\begin{proof}
Choisissons un ouvert affine $\Spec(\Lambda) \subset S$ tel que $k$ soit une
$\Lambda$-algèbre finie, voir
Morphismes, Lemme \ref{morphisms-lemma-point-finite-type}.
On peut, et on va, remplacer $S$ par $\Spec(\Lambda)$.

\medskip\noindent
On peut écrire $R$ comme une limite inductive filtrante $R = \colim C_j$, où chaque
$C_j$ est une $\Lambda$-algèbre de type fini (voir
Algèbre, Lemme \ref{algebra-lemma-ring-colimit-fp}).
D'après l'hypothèse (b), l'objet $x$ est isomorphe à la restriction d'un
objet au-dessus de l'un des $C_j$. On peut donc choisir une
$\Lambda$-algèbre de type fini $C$, un homomorphisme de $\Lambda$-algèbres $C \to R$ et un objet
$x_C$ de $\mathcal{X}$ au-dessus de $\Spec(C)$ tels que $x = x_C|_{\Spec(R)}$.
Le choix de $C$ est une commodité de notation dont on pourrait se passer.
Pour la suite, écrivons $C = \Lambda[y_1, \ldots, y_u]/(f_1, \ldots, f_v)$
et notons $\overline{a}_i \in R$ l'image de $y_i$ par
l'homomorphisme $C \to R$. Posons $\mathfrak m_C = C \cap \mathfrak m_R$.

\medskip\noindent
Choisissons une surjection de $\Lambda$-algèbres $\Lambda[x_1, \ldots, x_s] \to k$
et notons $\mathfrak m'$ son noyau.
Par la propriété universelle des anneaux de polynômes, on peut la relever
en un homomorphisme de $\Lambda$-algèbres $\Lambda[x_1, \ldots, x_s] \to R$.
Ajoutons quelques variables (c'est-à-dire augmentons un peu $s$) envoyées sur des générateurs
de $\mathfrak m_R$. Cela fait, on voit que
$\Lambda[x_1, \ldots, x_s] \to R/\mathfrak m_R^2$ est surjectif.
On en déduit que
\begin{equation}
\label{artin-equation-surjection}
P = \Lambda[x_1, \ldots, x_s]_{\mathfrak m'}^\wedge \longrightarrow R
\end{equation}
est un homomorphisme surjectif d'anneaux locaux noethériens complets, voir par exemple
Théorie formelle des déformations, Lemme
\ref{formal-defos-lemma-surjective-cotangent-space}.

\medskip\noindent
Choisissons des relèvements $a_i \in P$ des $\overline{a}_i$ trouvés ci-dessus.
Choisissons des générateurs $b_1, \ldots, b_r \in P$ du noyau de
(\ref{artin-equation-surjection}).
Choisissons des $c_{ji} \in P$ tels que
$$
f_j(a_1, \ldots, a_u) = \sum c_{ji} b_i
$$
dans $P$, ce qui est possible en vertu des choix déjà effectués. Choisissons des générateurs
$$
k_1, \ldots, k_t \in
\Ker(P^{\oplus r} \xrightarrow{(b_1, \ldots, b_r)} P)
$$
et écrivons $k_i = (k_{i1}, \ldots, k_{ir})$ et $K = (k_{ij})$
de sorte que
$$
P^{\oplus t} \xrightarrow{K}
P^{\oplus r} \xrightarrow{(b_1, \ldots, b_r)}
P \to R \to 0
$$
soit une suite exacte de $P$-modules. En particulier, on a
$\sum k_{ij} b_j = 0$. Quitte à augmenter $N$, on peut
supposer que $N - 1$ convient dans le lemme d'Artin-Rees pour les deux premiers morphismes de cette
suite exacte (voir Compléments d'algèbre, section
\ref{more-algebra-section-artin-rees} pour la terminologie).

\medskip\noindent
Par hypothèse, $\mathcal{O}_{S, s} = \Lambda_{\Lambda \cap \mathfrak m'}$ est
un anneau G. Ainsi, d'après Compléments d'algèbre, Proposition
\ref{more-algebra-proposition-finite-type-over-G-ring},
l'anneau $\Lambda[x_1, \ldots, x_s]_{\mathfrak m'}$ est un anneau $G$.
Par conséquent, d'après Lissage des homomorphismes d'anneaux, Théorème
\ref{smoothing-theorem-approximation-property-variant},
il existe un homomorphisme d'anneaux \'etale
$$
\Lambda[x_1, \ldots, x_s]_{\mathfrak m'} \to B,
$$
un idéal maximal $\mathfrak m_B$ de $B$ au-dessus de $\mathfrak m'$, et
des éléments $a'_i, b'_i, c'_{ij}, k'_{ij} \in B'$ tels que
\begin{enumerate}
\item $\kappa(\mathfrak m') = \kappa(\mathfrak m_B)$, ce qui implique
que $\Lambda[x_1, \ldots, x_s]_{\mathfrak m'} \subset B_{\mathfrak m_B}
\subset P$ et que $P$ s'identifie au complété de $B$ en
$\mathfrak m_B$, voir la remarque précédant Lissage des homomorphismes d'anneaux, Théorème
\ref{smoothing-theorem-approximation-property-variant},
\item $a_i - a'_i, b_i - b'_i, c_{ij} - c'_{ij}, k_{ij} - k'_{ij} \in
(\mathfrak m')^N P$, et
\item $f_j(a'_1, \ldots, a'_u) = \sum c'_{ji} b'_i$ et $\sum k'_{ij}b'_j = 0$.
\end{enumerate}
Posons $A = B/(b'_1, \ldots, b'_r)$ et notons $\mathfrak m_A$ l'image de
$\mathfrak m_B$ dans $A$. (Notons que $A$ est essentiellement de type
fini sur $\Lambda$ ; à la fin de la démonstration, nous expliquerons comment obtenir
un $A$ qui soit de type fini sur $\Lambda$.) Il existe un homomorphisme d'anneaux
$C \to A$ envoyant $y_i \mapsto a'_i$, car les $a'_i$ satisfont
les équations voulues modulo $(b'_1, \ldots, b'_r)$.
Notons que $A/\mathfrak m_A^N = R/\mathfrak m_R^N$ en tant que quotients de
$P = B^\wedge$ d'après la propriété (2) ci-dessus. Posons $x_A = x_C|_{\Spec(A)}$.
Comme les homomorphismes
$$
C \to A \to A/\mathfrak m_A^N \cong R/\mathfrak m_R^N
\quad\text{et}\quad
C \to R \to R/\mathfrak m_R^N
$$
sont égaux, on voit que $x_A$ et $x$ coïncident modulo $\mathfrak m_R^N$
via l'isomorphisme $A/\mathfrak m_A^N = R/\mathfrak m_R^N$. À ce
stade, nous avons établi les propriétés (1) -- (5) de l'énoncé du lemme.
Pour voir (6), notons que
$$
P^{\oplus t} \xrightarrow{K}
P^{\oplus r} \xrightarrow{(b_1, \ldots, b_r)}
P
\quad\text{et}\quad
P^{\oplus t} \xrightarrow{K'}
P^{\oplus r} \xrightarrow{(b'_1, \ldots, b'_r)}
P
$$
sont deux complexes de $P$-modules congrus modulo
$(\mathfrak m')^N$, le premier étant exact. D'après notre choix de $N$
ci-dessus, on voit d'après
Compléments d'algèbre, Lemme \ref{more-algebra-lemma-approximate-complex-graded},
que $R = P/(b_1, \ldots, b_r)$ et
$P/(b'_1, \ldots, b'_r) = B^\wedge/(b'_1, \ldots, b'_r) = A^\wedge$
ont des algèbres graduées associées isomorphes, ce qu'il fallait démontrer.

\medskip\noindent
Ce dernier paragraphe de la démonstration sert à régler le fait que $A$ est
essentiellement de type fini sur $S$, et non encore de type fini.
La construction ci-dessus donne $A = B/(b'_1, \ldots, b'_r)$ et
$\mathfrak m_A \subset A$, avec $B$ \'etale sur
$\Lambda[x_1, \ldots, x_s]_{\mathfrak m'}$. Ainsi, $A$ est de type
fini sur l'anneau noethérien $\Lambda[x_1, \ldots, x_s]_{\mathfrak m'}$.
On peut donc écrire $A = (A_0)_{\mathfrak m'}$ pour une
$\Lambda[x_1, \ldots, x_n]$-algèbre de type fini $A_0$. Alors
$A = \colim (A_0)_f$, où
$f \in \Lambda[x_1, \ldots, x_n] \setminus \mathfrak m'$, voir
Algèbre, Lemme \ref{algebra-lemma-localization-colimit}.
Puisque $p : \mathcal{X} \to (\Sch/S)_{fppf}$ commute aux limites sur les
objets, on voit que
$x_A$ provient d'un objet $x_{(A_0)_f}$ au-dessus de $\Spec((A_0)_f)$ pour
un $f$ comme ci-dessus. Après avoir remplacé $A$ par $(A_0)_f$, $x_A$ par
$x_{(A_0)_f}$ et $\mathfrak m_A$ par $(A_0)_f \cap \mathfrak m_A$,
la démonstration est achevée.
\end{proof}





\section{Commutation aux limites}
\label{artin-section-limits}

\noindent
Le morphisme $p : \mathcal{X} \to (\Sch/S)_{fppf}$ commute aux limites
sur les objets, au sens de Critères de représentabilité, section
\ref{criteria-section-limit-preserving}, si le foncteur de la définition
ci-dessous est essentiellement surjectif. Toutefois, l'exemple
de Exemples, section \ref{examples-section-limit-preserving},
montre que cela n'équivaut pas à la commutation aux limites.

\begin{definition}
\label{artin-definition-limit-preserving}
Soit $S$ un schéma. Soit $\mathcal{X}$ une catégorie fibrée en groupoïdes
sur $(\Sch/S)_{fppf}$. On dit que $\mathcal{X}$ {\it commute aux limites}
si, pour tout schéma affine $T$ sur $S$ qui est une limite $T = \lim T_i$
d'un système projectif filtrant de schémas affines $T_i$ sur $S$, on a
une équivalence
$$
\colim \mathcal{X}_{T_i} \longrightarrow \mathcal{X}_T
$$
de catégories fibres.
\end{definition}

\noindent
Explicitons ce que cela signifie. Premièrement, pour des objets $x, y$ de $\mathcal{X}$
au-dessus de $T_i$, on doit avoir
$$
\Mor_{\mathcal{X}_T}(x|_T, y|_T) =
\colim_{i' \geq i} \Mor_{\mathcal{X}_{T_{i'}}}(x|_{T_{i'}}, y|_{T_{i'}})
$$
et, deuxièmement, tout objet de $\mathcal{X}_T$ est isomorphe à la restriction
d'un objet au-dessus de $T_i$ pour un certain $i$. Notons que la première condition signifie
que les préfaisceaux $\mathit{Isom}_\mathcal{X}(x, y)$ (voir
Champs, Définition \ref{stacks-definition-mor-presheaf})
commutent aux limites.

\begin{lemma}
\label{artin-lemma-fibre-product-limit-preserving}
Soit $S$ un schéma. Soient $p : \mathcal{X} \to \mathcal{Y}$ et
$q : \mathcal{Z} \to \mathcal{Y}$ des $1$-morphismes de catégories
fibrées en groupoïdes sur $(\Sch/S)_{fppf}$.
\begin{enumerate}
\item Si $\mathcal{X} \to (\Sch/S)_{fppf}$ et
$\mathcal{Z} \to (\Sch/S)_{fppf}$ commutent aux limites sur les objets et
$\mathcal{Y}$ commute aux limites, alors
$\mathcal{X} \times_\mathcal{Y} \mathcal{Z} \to (\Sch/S)_{fppf}$ commute
aux limites sur les objets.
\item Si $\mathcal{X}$, $\mathcal{Y}$
et $\mathcal{Z}$ commutent aux limites, alors il en va de même de
$\mathcal{X} \times_\mathcal{Y} \mathcal{Z}$.
\end{enumerate}
\end{lemma}

\begin{proof}
C'est formel. Preuve de (1). Soit $T = \lim_{i \in I} T_i$ la limite
du système projectif filtrant de schémas affines $T_i$ sur $S$. Montrons que le foncteur
$\colim \mathcal{X}_{T_i} \to \mathcal{X}_T$ est essentiellement surjectif.
Rappelons qu'un objet du produit fibré au-dessus de $T$ est un quadruplet
$(T, x, z, \alpha)$, où $x$ est un objet de $\mathcal{X}$ au-dessus de $T$,
$z$ est un objet de $\mathcal{Z}$ au-dessus de $T$, et
$\alpha : p(x) \to q(z)$ est un morphisme dans la catégorie fibre de
$\mathcal{Y}$ au-dessus de $T$. D'après l'hypothèse sur $\mathcal{X}$ et $\mathcal{Z}$,
on peut trouver un $i$ et des objets $x_i$ et $z_i$ au-dessus de $T_i$ tels que
$x_i|_T \cong T$ et $z_i|_T \cong z$. Alors $\alpha$ correspond à
un isomorphisme $p(x_i)|_T \to q(z_i)|_T$ qui provient d'un isomorphisme
$\alpha_{i'} : p(x_i)|_{T_{i'}} \to q(z_i)|_{T_{i'}}$ d'après notre hypothèse sur
$\mathcal{Y}$. Après avoir remplacé $i$ par $i'$, $x_i$ par $x_i|_{T_{i'}}$ et
$z_i$ par $z_i|_{T_{i'}}$, on voit que $(T_i, x_i, z_i, \alpha_i)$
est un objet du produit fibré au-dessus de $T_i$ dont la restriction est
un objet isomorphe à $(T, x, z, \alpha)$ au-dessus de $T$, comme voulu.

\medskip\noindent
Nous omettons les arguments montrant que $\colim \mathcal{X}_{T_i} \to \mathcal{X}_T$
est pleinement fidèle dans (2).
\end{proof}

\begin{lemma}
\label{artin-lemma-limit-preserving-algebraic-space}
Soit $S$ un schéma. Soit $\mathcal{X}$ un champ algébrique sur $S$.
Les conditions suivantes sont alors équivalentes :
\begin{enumerate}
\item $\mathcal{X}$ est un champ en setoïdes et
$\mathcal{X} \to (\Sch/S)_{fppf}$ commute aux limites sur les objets,
\item $\mathcal{X}$ est un champ en setoïdes et commute aux limites,
\item $\mathcal{X}$ est représentable par un espace algébrique
localement de présentation finie.
\end{enumerate}
\end{lemma}

\begin{proof}
Sous chacune des trois hypothèses, $\mathcal{X}$ est représentable
par un espace algébrique $X$ sur $S$, voir Champs algébriques, Proposition
\ref{algebraic-proposition-algebraic-stack-no-automorphisms}.
Il est clair que (1) et (2) sont équivalentes, puisqu'un foncteur entre
setoïdes est une équivalence si et seulement s'il est surjectif sur les classes
d'isomorphisme. Enfin, (1) et (3) sont équivalentes d'après
Limites d'espaces, Proposition
\ref{spaces-limits-proposition-characterize-locally-finite-presentation}.
\end{proof}

\begin{lemma}
\label{artin-lemma-diagonal}
Soit $S$ un schéma. Soit $\mathcal{X}$ une catégorie fibrée
en groupoïdes sur $(\Sch/S)_{fppf}$. Supposons que
$\Delta : \mathcal{X} \to \mathcal{X} \times \mathcal{X}$ soit
représentable par des espaces algébriques et que $\mathcal{X}$ commute aux limites.
Alors $\Delta$ est localement de type fini.
\end{lemma}

\begin{proof}
Appliquons Critères de représentabilité, Lemme
\ref{criteria-lemma-check-property-limit-preserving}.
Soit $V$ un schéma affine $V$ localement de présentation finie sur $S$,
et soit $\theta$ un objet de $\mathcal{X} \times \mathcal{X}$
au-dessus de $V$. Soit $F_\theta$ un espace algébrique représentant
$\mathcal{X} \times_{\Delta, \mathcal{X} \times \mathcal{X}, \theta}
(\Sch/V)_{fppf}$, et soit $f_\theta : F_\theta \to V$ le morphisme canonique
(voir Champs algébriques, section
\ref{algebraic-section-morphisms-representable-by-algebraic-spaces}).
Il suffit de montrer que
$F_\theta \to V$ possède les propriétés correspondantes. D'après les
Lemmes \ref{artin-lemma-fibre-product-limit-preserving} et
\ref{artin-lemma-limit-preserving-algebraic-space},
on voit que $F_\theta \to S$ est localement de présentation finie.
Il s'ensuit que $F_\theta \to V$ est localement de type fini
d'après Morphismes d'espaces, Lemme
\ref{spaces-morphisms-lemma-permanence-finite-type}.
\end{proof}
 
 
 
 
 
 
\section{Versalité}
\label{artin-section-versality}

\noindent
Dans la section précédente, nous avons expliqué comment approcher les objets sur
des anneaux locaux complets par des objets algébriques. Mais, pour montrer
qu'un champ $\mathcal{X}$ est un champ algébrique, nous devons trouver des
$1$-morphismes lisses de schémas vers $\mathcal{X}$. Comme nous n'allons pas
supposer a priori que $\mathcal{X}$ possède une diagonale représentable, nous
ne pouvons même pas parler de morphismes lisses vers $\mathcal{X}$. Nous allons
plutôt introduire, en empruntant la terminologie de la théorie des déformations,
les objets versels.

\begin{definition}
\label{artin-definition-versal-formal-object}
Soit $S$ un schéma localement noethérien. Soit
$p : \mathcal{X} \to (\Sch/S)_{fppf}$ une catégorie fibrée en groupoïdes.
Soit $\xi = (R, \xi_n, f_n)$ un objet formel. Posons $k = R/\mathfrak m$ et
$x_0 = \xi_1$. Nous dirons que $\xi$ est {\it versel} si $\xi$,
considéré comme objet formel de $\mathcal{F}_{\mathcal{X}, k, x_0}$
(Remarque \ref{artin-remark-formal-objects-match}), est versel au sens de
Théorie des déformations formelles, Définition \ref{formal-defos-definition-versal}.
\end{definition}

\noindent
Explicitons brièvement ce que cela signifie. Avec les notations de la définition,
supposons donnés des morphismes $\xi_1 = x_0 \to y \to z$ de $\mathcal{X}$ au-dessus des
immersions fermées
$\Spec(k) \to \Spec(A) \to \Spec(B)$,
où $A, B$ sont des anneaux locaux artiniens de corps résiduel $k$.
Supposons donnés un entier $n \geq 1$ et un diagramme commutatif
$$
\vcenter{
\xymatrix{
& y \ar[ld] \\
\xi_n & \xi_1 \ar[u] \ar[l]
}
}
\quad\text{au-dessus de}\quad
\vcenter{
\xymatrix{
& \Spec(A) \ar[ld] \\
\Spec(R/\mathfrak m^n) & \Spec(k) \ar[u] \ar[l]
}
}
$$
La versalité signifie que, pour toutes données comme ci-dessus,
il existe un entier $m \geq n$ et un diagramme commutatif
$$
\vcenter{
\xymatrix{
& & z \ar[lldd] \\
& & y \ar[ld] \ar[u] \\
\xi_m & \xi_n \ar[l] & \xi_1 \ar[u] \ar[l]
}
}
\quad\text{au-dessus de}
\vcenter{
\xymatrix{
& & \Spec(B) \ar[lldd] \\
& & \Spec(A) \ar[ld] \ar[u] \\
\Spec(R/\mathfrak m^m) &
\Spec(R/\mathfrak m^n) \ar[l] &
\Spec(k) \ar[u] \ar[l]
}
}
$$
Comparer avec Théorie des déformations formelles, Remarque
\ref{formal-defos-remark-versal-object}.

\medskip\noindent
Soit $S$ un schéma localement noethérien. Soit $U$ un schéma sur $S$
dont le morphisme structural $U \to S$ est localement de type fini. Soit
$u_0 \in U$ un point de type fini de $U$ ; voir
Morphismes, Définition \ref{morphisms-definition-finite-type-point}.
Posons $k = \kappa(u_0)$.
Notons que le composé $\Spec(k) \to S$ est aussi de type fini ;
voir Morphismes, Lemme \ref{morphisms-lemma-composition-finite-type}.
Soit $p : \mathcal{X} \to (\Sch/S)_{fppf}$ une catégorie fibrée en groupoïdes.
Soit $x$ un objet de $\mathcal{X}$ au-dessus de $U$. Notons $x_0$
l'image réciproque de $x$ par $u_0$. Par le lemme de $2$-Yoneda, $x$ correspond
à un $1$-morphisme
$$
x : (\Sch/U)_{fppf} \longrightarrow \mathcal{X},
$$
voir Champs algébriques, section \ref{algebraic-section-2-yoneda}. Nous obtenons un
morphisme de catégories de prédéformation
\begin{equation}
\label{artin-equation-hat-x}
\hat x :
\mathcal{F}_{(\Sch/U)_{fppf}, k, u_0}
\longrightarrow
\mathcal{F}_{\mathcal{X}, k, x_0},
\end{equation}
au-dessus de $\mathcal{C}_\Lambda$ ; voir (\ref{artin-equation-functoriality}).

\begin{definition}
\label{artin-definition-versal}
Soit $S$ un schéma localement noethérien.
Soit $\mathcal{X}$ une catégorie fibrée en groupoïdes sur $(\Sch/S)_{fppf}$.
Soit $U$ un schéma localement de type fini sur $S$.
Soit $x$ un objet de $\mathcal{X}$ au-dessus de $U$.
Soit $u_0$ un point de type fini de $U$.
On dit que $x$ est {\it versel} en $u_0$ si le morphisme $\hat x$
(\ref{artin-equation-hat-x}) est lisse ; voir Théorie des déformations formelles, Définition
\ref{formal-defos-definition-smooth-morphism}.
\end{definition}

\noindent
Cette définition coïncide avec notre notion de versalité pour les objets formels de
$\mathcal{X}$.

\begin{lemma}
\label{artin-lemma-versality-matches}
Avec les notations de la Définition \ref{artin-definition-versal},
soit $R = \mathcal{O}_{U, u_0}^\wedge$.
Soit $\xi$ l'objet formel de $\mathcal{X}$
sur $R$ associé à $x|_{\Spec(R)}$ ; voir (\ref{artin-equation-approximation}).
Alors
$$
x\text{ est versel en }u_0
\Leftrightarrow
\xi\text{ est versel}
$$
\end{lemma}

\begin{proof}
Observons que $\mathcal{O}_{U, u_0}$ est une $S$-algèbre locale noethérienne
de corps résiduel $k$. Ainsi, $R = \mathcal{O}_{U, u_0}^\wedge$ est un objet de
$\mathcal{C}_\Lambda^\wedge$ ; voir Théorie des déformations formelles, Définition
\ref{formal-defos-definition-completion-CLambda}.
Rappelons que $\xi$ est versel si
$\underline{\xi} : \underline{R}|_{\mathcal{C}_\Lambda} \to
\mathcal{F}_{\mathcal{X}, k, x_0}$
est lisse, et que $x$ est versel en $u_0$ si
$\hat x : \mathcal{F}_{(\Sch/U)_{fppf}, k, u_0}
\to \mathcal{F}_{\mathcal{X}, k, x_0}$ est lisse.
Il existe une identification de catégories de prédéformation
$$
\underline{R}|_{\mathcal{C}_\Lambda}
=
\mathcal{F}_{(\Sch/U)_{fppf}, k, u_0},
$$
voir Théorie des déformations formelles, Remarque
\ref{formal-defos-remark-formal-objects-yoneda} pour les notations.
En effet, pour toute $S$-algèbre locale artinienne $A$ dont le corps résiduel
est identifié à $k$, nous avons
$$
\Mor_{\mathcal{C}_\Lambda^\wedge}(R, A) =
\{\varphi \in \Mor_S(\Spec(A), U) \mid \varphi|_{\Spec(k)} = u_0\}
$$
En explicitant les définitions, on vérifie que le morphisme obtenu
$$
\underline{R}|_{\mathcal{C}_\Lambda} =
\mathcal{F}_{(\Sch/U)_{fppf}, k, u_0}
\xrightarrow{\hat x}
\mathcal{F}_{\mathcal{X}, k, x_0},
$$
est égal à $\underline{\xi}$, ce qui démontre le lemme.
\end{proof}

\noindent
Voici une vérification de cohérence.

\begin{lemma}
\label{artin-lemma-versal-implies-smooth}
Soit $S$ un schéma localement noethérien. Soit $f : U \to V$
un morphisme de schémas localement de type fini sur $S$.
Soit $u_0 \in U$ un point de type fini. Les conditions suivantes sont équivalentes :
\begin{enumerate}
\item $f$ est lisse en $u_0$,
\item $f$, considéré comme objet de $(\Sch/V)_{fppf}$ au-dessus de $U$, est
versel en $u_0$.
\end{enumerate}
\end{lemma}

\begin{proof}
C'est une reformulation de Compléments sur les morphismes, Lemme
\ref{more-morphisms-lemma-lifting-along-artinian-at-point}.
\end{proof}

\noindent
Il s'avère que cette notion se comporte bien par
extension de corps.

\begin{lemma}
\label{artin-lemma-versal-change-of-field}
Soient $S$, $\mathcal{X}$, $U$, $x$, $u_0$ comme dans la
Définition \ref{artin-definition-versal}. Soit $l$ un corps et soit
$u_{l, 0} : \Spec(l) \to U$ un morphisme d'image $u_0$ tel que
l'extension $l/k = \kappa(u_0)$ soit finie. Posons $x_{l, 0} = x_0|_{\Spec(l)}$.
Si $\mathcal{X}$ vérifie (RS) et si $x$ est versel en $u_0$, alors
$$
\mathcal{F}_{(\Sch/U)_{fppf}, l, u_{l, 0}}
\longrightarrow
\mathcal{F}_{\mathcal{X}, l, x_{l, 0}}
$$
est lisse.
\end{lemma}

\begin{proof}
Notons que $(\Sch/U)_{fppf}$ vérifie (RS) d'après le
Lemme \ref{artin-lemma-algebraic-stack-RS}.
Le foncteur du lemme est donc le foncteur
$$
(\mathcal{F}_{(\Sch/U)_{fppf}, k , u_0})_{l/k}
\longrightarrow
(\mathcal{F}_{\mathcal{X}, k , x_0})_{l/k}
$$
associé à $\hat x$ ; voir le Lemme \ref{artin-lemma-change-of-field}.
Le lemme résulte donc de
Théorie des déformations formelles, Lemme
\ref{formal-defos-lemma-change-of-fields-smooth}.
\end{proof}

\noindent
Le lemme suivant est une autre vérification de cohérence. Il
exprime essentiellement que, si $x$ est versel en $u_0$ comme dans la
Définition \ref{artin-definition-versal},
alors $x$, considéré comme morphisme de $U$ vers $\mathcal{X}$, est
lisse après tout changement de base par un schéma.

\begin{lemma}
\label{artin-lemma-base-change-versal}
Soient $S$, $\mathcal{X}$, $U$, $x$, $u_0$ comme dans la
Définition \ref{artin-definition-versal}. Supposons que
\begin{enumerate}
\item $\Delta : \mathcal{X} \to \mathcal{X} \times \mathcal{X}$
soit représentable par des espaces algébriques,
\item $\Delta$ soit localement de type fini
(par exemple si $\mathcal{X}$ commute aux limites), et
\item $\mathcal{X}$ vérifie (RS).
\end{enumerate}
Soit $V$ un schéma localement de type fini sur $S$
et soit $y$ un objet de $\mathcal{X}$ au-dessus de $V$.
Formons le $2$-produit fibré
$$
\xymatrix{
\mathcal{Z} \ar[r] \ar[d] & (\Sch/U)_{fppf} \ar[d]^x \\
(\Sch/V)_{fppf} \ar[r]^y & \mathcal{X}
}
$$
Soit $Z$ l'espace algébrique qui représente $\mathcal{Z}$,
et soit $z_0 \in |Z|$ un point de type fini au-dessus de $u_0$.
Si $x$ est versel en $u_0$, alors
le morphisme $Z \to V$ est lisse en $z_0$.
\end{lemma}

\begin{proof}
(La remarque entre parenthèses dans l'énoncé résulte du
Lemme \ref{artin-lemma-diagonal}.)
L'espace $Z$ existe par l'hypothèse (1) et Champs algébriques, Lemme
\ref{algebraic-lemma-representable-diagonal}. Par l'hypothèse (2), nous voyons que
$Z \to V \times_S U$ est localement de type fini.
Choisissons un schéma $W$, un point fermé $w_0 \in W$ et
un morphisme \'etale $W \to Z$ envoyant $w_0$ sur $z_0$ ; voir
Morphismes d'espaces, Définition
\ref{spaces-morphisms-definition-finite-type-point}.
Alors $W$ est localement de type fini sur $S$ et
$w_0$ est un point de type fini de $W$.
Posons $l = \kappa(z_0)$. Notons $z_{l, 0}$, $v_{l, 0}$,
$u_{l, 0}$ et $x_{l, 0}$ les objets de
$\mathcal{Z}$, $(\Sch/V)_{fppf}$, $(\Sch/U)_{fppf}$
et $\mathcal{X}$ au-dessus de $\Spec(l)$ obtenus par image réciproque sur $\Spec(l) = w_0$.
Considérons
$$
\xymatrix{
\mathcal{F}_{(\Sch/W)_{fppf}, l, w_0} \ar[r] &
\mathcal{F}_{\mathcal{Z}, l, z_{l, 0}} \ar[d] \ar[r] &
\mathcal{F}_{(\Sch/U)_{fppf}, l, u_{l, 0}} \ar[d] \\
& \mathcal{F}_{(\Sch/V)_{fppf}, l, v_{l, 0}} \ar[r] &
\mathcal{F}_{\mathcal{X}, l, x_{l, 0}}
}
$$
D'après le Lemme \ref{artin-lemma-fibre-product-deformation-categories},
le carré est un produit fibré de catégories de prédéformation.
D'après le Lemme \ref{artin-lemma-versal-change-of-field},
la flèche verticale de droite est lisse.
D'après Théorie des déformations formelles, Lemme
\ref{formal-defos-lemma-smooth-properties},
la flèche verticale de gauche est lisse.
D'après le Lemme \ref{artin-lemma-formally-smooth-on-deformation-categories},
la flèche horizontale de gauche est lisse.
Nous en concluons que l'application
$$
\mathcal{F}_{(\Sch/W)_{fppf}, l, w_0} \to
\mathcal{F}_{(\Sch/V)_{fppf}, l, v_{l, 0}}
$$
est lisse d'après Théorie des déformations formelles, Lemme
\ref{formal-defos-lemma-smooth-properties}.
Nous en déduisons que $W \to V$ est lisse en $w_0$ d'après
Compléments sur les morphismes, Lemme
\ref{more-morphisms-lemma-lifting-along-artinian-at-point}.
Cela signifie précisément que $Z \to V$ est lisse en $z_0$,
ce qui achève la démonstration.
\end{proof}

\noindent
Reformulons le résultat d'approximation au moyen des
objets versels.

\begin{lemma}
\label{artin-lemma-approximate-versal}
Soit $S$ un schéma localement noethérien. Soit
$p : \mathcal{X} \to (\Sch/S)_{fppf}$ une catégorie fibrée en groupoïdes.
Soit $\xi = (R, \xi_n, f_n)$ un objet formel de $\mathcal{X}$ dont
$\xi_1$ est au-dessus de $\Spec(k) \to S$, d'image $s \in S$. Supposons que
\begin{enumerate}
\item $\xi$ soit versel,
\item $\xi$ soit effectif,
\item $\mathcal{O}_{S, s}$ soit un anneau G, et
\item $p : \mathcal{X} \to (\Sch/S)_{fppf}$ commute aux limites sur les objets.
\end{enumerate}
Alors il existe un morphisme de type fini $U \to S$, un point de type fini
$u_0 \in U$ de corps résiduel $k$, et un objet $x$ de $\mathcal{X}$
au-dessus de $U$, tel que $x$ soit versel en $u_0$ et que
$x|_{\Spec(\mathcal{O}_{U, u_0}/\mathfrak m_{u_0}^n)} \cong \xi_n$.
\end{lemma}

\begin{proof}
Choisissons un objet $x_R$ de $\mathcal{X}$ au-dessus de $\Spec(R)$ dont l'objet
formel associé est $\xi$. Posons $N = 2$ et appliquons le Lemme \ref{artin-lemma-approximate}.
Nous obtenons $A, \mathfrak m_A, x_A, \ldots$.
Soit $\eta = (A^\wedge, \eta_n, g_n)$ l'objet formel associé à
$x_A|_{\Spec(A^\wedge)}$. Nous avons un diagramme
$$
\vcenter{
\xymatrix{
& \eta \ar[d] \\
\xi \ar[r] \ar@{..>}[ru] & \xi_2 = \eta_2
}
}
\quad\text{au-dessus de}\quad
\vcenter{
\xymatrix{
& A^\wedge \ar[d] \\
R \ar[r] \ar@{..>}[ru] & R/\mathfrak m_R^2 = A/\mathfrak m_A^2
}
}
$$
La versalité de $\xi$ signifie précisément que nous pouvons construire les
flèches pointillées dans les diagrammes, car nous pouvons construire successivement
des morphismes $\xi \to \eta_3$, $\xi \to \eta_4$, et ainsi de suite, d'après
Théorie des déformations formelles, Remarque \ref{formal-defos-remark-versal-object}.
L'homomorphisme d'anneaux correspondant $R \to A^\wedge$ est surjectif d'après
Théorie des déformations formelles, Lemme
\ref{formal-defos-lemma-surjective-cotangent-space}.
D'autre part, nous avons
$\dim_k \mathfrak m_R^n/\mathfrak m_R^{n + 1} =
\dim_k \mathfrak m_A^n/\mathfrak m_A^{n + 1}$ pour tout $n$, par construction.
Ainsi, $R/\mathfrak m_R^n$ et $A/\mathfrak m_A^n$ ont la même longueur (finie)
comme $\Lambda$-modules, par additivité de la longueur et d'après
Théorie des déformations formelles, Lemme \ref{formal-defos-lemma-length}.
Il s'ensuit que $R/\mathfrak m_R^n \to A/\mathfrak m_A^n$ est un isomorphisme
pour tout $n$ ; par conséquent, $R \to A^\wedge$ est un isomorphisme. Ainsi, $\eta$ est
isomorphe à un objet versel, donc est lui-même versel. D'après le
Lemme \ref{artin-lemma-versality-matches},
nous concluons que $x_A$ est versel au point $u_0$ de
$U = \Spec(A)$ correspondant à $\mathfrak m_A$.
\end{proof}

\begin{example}
\label{artin-example-approximate-versal-implies}
Dans cet exemple, nous montrons que l'anneau local $\mathcal{O}_{S, s}$ doit être
un anneau G pour que le résultat du Lemme \ref{artin-lemma-approximate-versal}
soit vrai. En effet, soit $\Lambda$ un anneau noethérien et soit $\mathfrak m$
un idéal maximal de $\Lambda$. Posons $R = \Lambda_\mathfrak m^\wedge$. Soit
$\Lambda \to C \to R$ une factorisation où $C$ est de type fini sur
$\Lambda$. Posons $S = \Spec(\Lambda)$, $U = S \setminus \{\mathfrak m\}$ et
$S' = U \amalg \Spec(C)$. Considérons le foncteur
$F : (\Sch/S)_{fppf}^{opp} \to \textit{Ensembles}$ défini par la règle
$$
F(T) = 
\left\{
\begin{matrix}
* & \text{si }T \to S\text{ se factorise par }S' \\
\emptyset & \text{sinon}
\end{matrix}
\right.
$$
Soit $\mathcal{X} = \mathcal{S}_F$ la catégorie fibrée en ensembles associée
à $F$ ; voir Champs algébriques, section \ref{algebraic-section-split}.
Alors $\mathcal{X} \to (\Sch/S)_{fppf}$ commute aux limites sur les objets, et
il existe un objet formel effectif et versel $\xi$ sur $R$.
Ainsi, si la conclusion du Lemme \ref{artin-lemma-approximate-versal} vaut
pour $\mathcal{X}$, alors il existe un homomorphisme d'anneaux de type fini $\Lambda \to A$
et un idéal maximal $\mathfrak m_A$ au-dessus de $\mathfrak m$ tels que
\begin{enumerate}
\item $\kappa(\mathfrak m) = \kappa(\mathfrak m_A)$,
\item $\Lambda \to A$ et $\mathfrak m_A$ vérifient la condition (4) du
Lemme \ref{algebra-lemma-smooth-test-artinian} d'Algèbre, et
\item il existe un homomorphisme de $\Lambda$-algèbres $C \to A$.
\end{enumerate}
Ainsi, $\Lambda \to A$ est lisse en $\mathfrak m_A$ d'après le lemme cité.
Après une localisation puis un passage au quotient convenables de $A$, nous pouvons supposer que $\Lambda \to A$ est \'etale en
$\mathfrak m_A$ ; voir par exemple
Compléments sur les morphismes, Lemme \ref{more-morphisms-lemma-slice-smooth},
ou le vérifier directement. Écrivons $C = \Lambda[y_1, \ldots, y_n]/(f_1, \ldots, f_m)$.
Alors $C \to R$ correspond à une solution dans $R$ du système d'équations
$f_1 = \ldots = f_m = 0$ ; voir Lissage des homomorphismes d'anneaux, section
\ref{smoothing-section-approximation-G-rings}.
Par conséquent, si la conclusion du
Lemme \ref{artin-lemma-approximate-versal} vaut pour tout $\mathcal{X}$ comme
ci-dessus, alors tout système d'équations qui possède une solution dans $R$ possède une
solution dans l'hensélisé de $\Lambda_{\mathfrak m}$.
Autrement dit, la propriété d'approximation vaut pour
$\Lambda_{\mathfrak m}^h$. Cela implique que $\Lambda_{\mathfrak m}^h$
est un anneau G (insérer ici une référence ultérieure ; voir aussi la discussion dans
Lissage des homomorphismes d'anneaux, section \ref{smoothing-section-introduction}),
ce qui implique à son tour que $\Lambda_{\mathfrak m}$ est un anneau G.
\end{example}






\section{Ouverture de la versalité}
\label{artin-section-openness-versality}

\noindent
Nous abordons maintenant l'ouverture de la versalité.

\begin{definition}
\label{artin-definition-openness-versality}
Soit $S$ un schéma localement noethérien.
\begin{enumerate}
\item Soit $\mathcal{X}$ une catégorie
fibrée en groupoïdes sur $(\Sch/S)_{fppf}$. Nous disons que $\mathcal{X}$ vérifie
{\it l'ouverture de la versalité} si, étant donnés un schéma $U$ localement de type fini
sur $S$, un objet $x$ de $\mathcal{X}$ au-dessus de $U$ et un point de type fini
$u_0 \in U$ tel que $x$ soit versel en $u_0$, il existe un voisinage ouvert
$u_0 \in U' \subset U$ tel que $x$ soit versel en tout point de type fini
de $U'$.
\item Soit $f : \mathcal{X} \to \mathcal{Y}$ un $1$-morphisme de catégories
fibrées en groupoïdes sur $(\Sch/S)_{fppf}$. Nous disons que $f$ vérifie
{\it l'ouverture de la versalité} si, étant donnés un schéma $U$ localement de type fini
sur $S$ et un objet $y$ de $\mathcal{Y}$ au-dessus de $U$, l'ouverture
de la versalité vaut pour
$(\Sch/U)_{fppf} \times_\mathcal{Y} \mathcal{X}$.
\end{enumerate}
\end{definition}

\noindent
L'ouverture de la versalité est souvent la condition la plus difficile à vérifier. L'exemple suivant
montre néanmoins qu'il est nécessaire de l'imposer.

\begin{example}
\label{artin-example-versality}
Soit $k$ un corps et posons $\Lambda = k[s, t]$. Considérons le foncteur
$F : \Lambda\text{-algèbres} \longrightarrow \textit{Ensembles}$
défini par la règle
$$
F(A) =
\left\{
\begin{matrix}
* & \text{s'il existe }f_1, \ldots, f_n \in A\text{ tels que } \\
  & A = (s, t, f_1, \ldots, f_n)\text{ et } f_i s = 0\ \forall i \\
\emptyset & \text{sinon}
\end{matrix}
\right.
$$
Géométriquement, $F(A) = *$ signifie qu'il existe un voisinage ouvert quasi-compact
$W$ de $V(s, t) \subset \Spec(A)$ tel que $s|_W = 0$.
Soit $\mathcal{X} \subset (\Sch/\Spec(\Lambda))_{fppf}$ la sous-catégorie pleine
formée des schémas $T$ qui admettent un recouvrement ouvert affine
$T = \bigcup \Spec(A_j)$ tel que $F(A_j) = *$ pour tout $j$. Alors $\mathcal{X}$
vérifie [0], [1], [2], [3] et [4], mais pas [5]. En effet, au-dessus de
$U = \Spec(k[s, t]/(s))$,
il existe un objet $x$ qui est versel en $u_0 = (s, t)$, mais qui ne l'est
en aucun autre point. Nous omettons les détails.
\end{example}

\noindent
Soit $S$ un schéma localement noethérien.
Soit $f : \mathcal{X} \to \mathcal{Y}$ un $1$-morphisme de catégories
fibrées en groupoïdes sur $(\Sch/S)_{fppf}$. Considérons la propriété suivante :
\begin{equation}
\label{artin-equation-smooth}
\begin{matrix}
\text{pour tout corps }k\text{ de type fini sur }S
\text{ et tout }x_0 \in \Ob(\mathcal{X}_{\Spec(k)})\text{,}\\
\text{ l'application }
\mathcal{F}_{\mathcal{X}, k, x_0} \to \mathcal{F}_{\mathcal{Y}, k, f(x_0)}
\text{ de catégories de prédéformation est lisse}
\end{matrix}
\end{equation}
Formulons quelques lemmes concernant cette notion. Nous la relions d'abord à
(l'ouverture de) la versalité.

\begin{lemma}
\label{artin-lemma-versal-smooth}
Soit $S$ un schéma localement noethérien. Soit $\mathcal{X}$ une catégorie
fibrée en groupoïdes sur $(\Sch/S)_{fppf}$. Soit $U$ un schéma localement
de type fini sur $S$. Soit $x$ un objet de $\mathcal{X}$ au-dessus de $U$.
Supposons que $x$ soit versel en tout point de type fini de $U$ et que
$\mathcal{X}$ vérifie (RS). Alors $x : (\Sch/U)_{fppf} \to \mathcal{X}$
vérifie (\ref{artin-equation-smooth}).
\end{lemma}

\begin{proof}
Soit $\Spec(l) \to U$ un morphisme où $l$ est de type fini sur $S$.
Alors l'image $u_0 \in U$ est un point de type fini de $U$, et
$l/\kappa(u_0)$ est une extension finie ; voir la discussion dans
Morphismes, section \ref{morphisms-section-points-finite-type}.
Par conséquent,
$\mathcal{F}_{(\Sch/U)_{fppf}, l, u_{l, 0}} \to
\mathcal{F}_{\mathcal{X}, l, x_{l, 0}}$
est lisse d'après le Lemme \ref{artin-lemma-versal-change-of-field}.
\end{proof}

\begin{lemma}
\label{artin-lemma-composition-smooth}
Soit $S$ un schéma localement noethérien. Soient $f : \mathcal{X} \to \mathcal{Y}$
et $g : \mathcal{Y} \to \mathcal{Z}$ deux $1$-morphismes composables de
catégories fibrées en groupoïdes sur $(\Sch/S)_{fppf}$. Si $f$ et $g$
vérifient (\ref{artin-equation-smooth}), alors $g \circ f$ la vérifie aussi.
\end{lemma}

\begin{proof}
Cela résulte formellement de Théorie des déformations formelles, Lemme
\ref{formal-defos-lemma-smooth-properties}.
\end{proof}

\begin{lemma}
\label{artin-lemma-base-change-smooth}
Soit $S$ un schéma localement noethérien. Soient $f : \mathcal{X} \to \mathcal{Y}$
et $\mathcal{Z} \to \mathcal{Y}$ des $1$-morphismes de
catégories fibrées en groupoïdes sur $(\Sch/S)_{fppf}$. Si $f$
vérifie (\ref{artin-equation-smooth}), alors la projection
$\mathcal{X} \times_\mathcal{Y} \mathcal{Z} \to \mathcal{Z}$ la vérifie aussi.
\end{lemma}

\begin{proof}
Cela résulte immédiatement du
Lemme \ref{artin-lemma-fibre-product-deformation-categories}
et de
Théorie des déformations formelles, Lemme
\ref{formal-defos-lemma-smooth-properties}.
\end{proof}

\begin{lemma}
\label{artin-lemma-smooth-smooth}
Soit $S$ un schéma localement noethérien. Soit $f : \mathcal{X} \to \mathcal{Y}$
un $1$-morphisme de catégories fibrées en groupoïdes sur $(\Sch/S)_{fppf}$.
Si $f$ est formellement lisse sur les objets, alors $f$ vérifie
(\ref{artin-equation-smooth}). Si $f$ est représentable par des espaces algébriques
et lisse, alors $f$ vérifie (\ref{artin-equation-smooth}).
\end{lemma}

\begin{proof}
C'est une reformulation du Lemme \ref{artin-lemma-formally-smooth-on-deformation-categories}.
\end{proof}

\begin{lemma}
\label{artin-lemma-implies-smooth}
Soit $S$ un schéma localement noethérien. Soit $f : \mathcal{X} \to \mathcal{Y}$
un $1$-morphisme de catégories fibrées en groupoïdes sur $(\Sch/S)_{fppf}$.
Supposons que
\begin{enumerate}
\item $f$ soit représentable par des espaces algébriques,
\item $f$ vérifie (\ref{artin-equation-smooth}),
\item $\mathcal{X} \to (\Sch/S)_{fppf}$ commute aux limites sur les objets, et
\item $\mathcal{Y}$ commute aux limites.
\end{enumerate}
Alors $f$ est lisse.
\end{lemma}

\begin{proof}
L'ingrédient clé de la démonstration est Compléments sur les morphismes, Lemme
\ref{more-morphisms-lemma-lifting-along-artinian-at-point},
qui affirme presque qu'un morphisme de schémas de type fini sur $S$
vérifiant (\ref{artin-equation-smooth}) est un morphisme lisse. Les autres
arguments de la démonstration se réduisent essentiellement à des vérifications de compatibilité.

\medskip\noindent
Soit $V$ un schéma sur $S$ et soit $y$ un objet de $\mathcal{Y}$ au-dessus de
$V$. Soit $Z$ un espace algébrique représentant le $2$-produit fibré
$\mathcal{Z} = \mathcal{X} \times_{f, \mathcal{X}, y} (\Sch/V)_{fppf}$.
Nous devons montrer que le morphisme de projection $Z \to V$ est lisse ; voir
Champs algébriques, Définition
\ref{algebraic-definition-relative-representable-property}.
Il suffit même de le faire lorsque $V$ est un schéma affine
localement de présentation finie sur $S$ ; voir
Critères de représentabilité, Lemme
\ref{criteria-lemma-check-property-limit-preserving}.
Alors $(\Sch/V)_{fppf}$ commute aux limites d'après le
Lemme \ref{artin-lemma-limit-preserving-algebraic-space}.
Par conséquent, $Z \to S$ est localement de présentation finie d'après les
Lemmes \ref{artin-lemma-fibre-product-limit-preserving} et
\ref{artin-lemma-limit-preserving-algebraic-space}.
Choisissons un schéma $W$ et un morphisme \'etale surjectif $W \to Z$.
Alors $W$ est localement de présentation finie sur $S$.

\medskip\noindent
Puisque $f$ vérifie (\ref{artin-equation-smooth}), il en est de même de
$\mathcal{Z} \to (\Sch/V)_{fppf}$ ; voir le Lemme \ref{artin-lemma-base-change-smooth}.
Ensuite, $(\Sch/W)_{fppf} \to \mathcal{Z}$ vérifie
(\ref{artin-equation-smooth}) d'après le Lemme \ref{artin-lemma-smooth-smooth}.
Ainsi, le composé
$$
(\Sch/W)_{fppf} \to \mathcal{Z} \to (\Sch/V)_{fppf}
$$
vérifie (\ref{artin-equation-smooth}) d'après le Lemme \ref{artin-lemma-composition-smooth}.
Compléments sur les morphismes, Lemme
\ref{more-morphisms-lemma-lifting-along-artinian-at-point},
montre que le composé $W \to Z \to V$ est lisse en tout point de type fini
$w_0$ de $W$. Le lieu de lissité étant ouvert, nous en concluons
que $W \to V$ est un morphisme lisse de schémas d'après
Morphismes, Lemme \ref{morphisms-lemma-enough-finite-type-points}.
Nous en concluons que $Z \to V$ est un morphisme lisse
d'espaces algébriques, par définition.
\end{proof}

\noindent
Le lemme ci-dessous indique comment nous utiliserons l'ouverture de la versalité.

\begin{lemma}
\label{artin-lemma-get-smooth}
Soit $S$ un schéma localement noethérien. Soit
$p : \mathcal{X} \to (\Sch/S)_{fppf}$ une catégorie fibrée en groupoïdes.
Soit $k$ un corps de type fini sur $S$ et soit $x_0$ un objet de
$\mathcal{X}$ au-dessus de $\Spec(k)$, d'image $s \in S$. Supposons que
\begin{enumerate}
\item $\Delta : \mathcal{X} \to \mathcal{X} \times \mathcal{X}$ soit
représentable par des espaces algébriques,
\item $\mathcal{X}$ vérifie les axiomes [1], [2], [3] (voir la
section \ref{artin-section-axioms}),
\item tout objet formel de $\mathcal{X}$ soit effectif,
\item l'ouverture de la versalité vaille pour $\mathcal{X}$, et
\item $\mathcal{O}_{S, s}$ soit un anneau G.
\end{enumerate}
Alors il existe un morphisme de type fini $U \to S$ et un objet
$x$ de $\mathcal{X}$ au-dessus de $U$ tels que
$$
x : (\Sch/U)_{fppf} \longrightarrow \mathcal{X}
$$
soit lisse et qu'il existe un point de type fini $u_0 \in U$
de corps résiduel $k$ tel que $x|_{u_0} \cong x_0$.
\end{lemma}

\begin{proof}
Par l'axiome [2], le Lemme \ref{artin-lemma-deformation-category} et la
Remarque \ref{artin-remark-deformation-category-implies},
nous voyons que $\mathcal{F}_{\mathcal{X}, k, x_0}$ vérifie (S1) et (S2).
Comme l'espace tangent est en outre de dimension finie par l'axiome [3],
nous déduisons de Théorie des déformations formelles, Lemme
\ref{formal-defos-lemma-versal-object-existence},
que $\mathcal{F}_{\mathcal{X}, k, x_0}$ possède un objet formel versel $\xi$.
L'hypothèse (3) affirme que $\xi$ est effectif. Par l'axiome [1] et le
Lemme \ref{artin-lemma-approximate-versal},
il existe un morphisme de type fini $U \to S$, un objet $x$ de
$\mathcal{X}$ au-dessus de $U$, et un point de type fini $u_0$ de $U$, de corps
résiduel $k$, tels que $x$ soit versel en $u_0$ et que
$x|_{\Spec(k)} \cong x_0$. Par l'ouverture de la versalité, quitte à rétrécir
$U$, nous pouvons supposer que $x$ est versel en tout point de type fini de $U$.
Nous affirmons que
$$
x : (\Sch/U)_{fppf} \longrightarrow \mathcal{X}
$$
est lisse, ce qui démontre le lemme. En effet, d'après le Lemme \ref{artin-lemma-versal-smooth},
$x$ vérifie (\ref{artin-equation-smooth}),
et le Lemme \ref{artin-lemma-implies-smooth}
achève alors la démonstration.
\end{proof}






\section{Axiomes}
\label{artin-section-axioms}

\noindent
Soit $S$ un schéma localement noethérien. Soit
$p : \mathcal{X} \to (\Sch/S)_{fppf}$ une catégorie fibrée en groupoïdes.
Voici les axiomes que nous considérerons pour $\mathcal{X}$.
\begin{enumerate}
\item[{[-1]}] une condition ensembliste\footnote{La condition est la
suivante : la borne supérieure de tous les cardinaux
$|\Ob(\mathcal{X}_{\Spec(k)})/\cong|$ et
$|\text{Flèches}(\mathcal{X}_{\Spec(k)})|$, où $k$ parcourt les corps de
type fini sur $S$, est inférieure ou égale ($\leq$) au cardinal d'un
objet de $(\Sch/S)_{fppf}$.} que peuvent ignorer les
lecteurs que les questions ensemblistes n'intéressent pas,
\item[{[0]}] $\mathcal{X}$ est un champ en groupoïdes pour la topologie \'etale,
\item[{[1]}] $\mathcal{X}$ commute aux limites,
\item[{[2]}] $\mathcal{X}$ vérifie la condition de Rim--Schlessinger (RS),
\item[{[3]}] les espaces $T\mathcal{F}_{\mathcal{X}, k, x_0}$ et
$\text{Inf}(\mathcal{F}_{\mathcal{X}, k, x_0})$
sont de dimension finie
pour tout $k$ et tout $x_0$ ; voir
(\ref{artin-equation-tangent-space}) et
(\ref{artin-equation-infinitesimal-automorphisms}),
\item[{[4]}] le foncteur (\ref{artin-equation-approximation}) est une équivalence,
\item[{[5]}] $\mathcal{X}$ et
$\Delta : \mathcal{X} \to \mathcal{X} \times \mathcal{X}$ vérifient
l'ouverture de la versalité.
\end{enumerate}
 
 
 
 
 
 




\section{Axiomes pour les foncteurs}
\label{artin-section-axioms-functors}

\noindent
Soit $S$ un schéma. Soit $F : (\Sch/S)_{fppf}^{opp} \to \textit{Ensembles}$ un
foncteur. Notons $\mathcal{X} = \mathcal{S}_F$ la catégorie fibrée en ensembles
associée à $F$, voir Champs algébriques, section \ref{algebraic-section-split}.
Dans cette section, nous traduisons les résultats précédents, tels
qu'ils s'appliquent à $\mathcal{X}$, en énoncés portant sur $F$.

\medskip\noindent
Soit $S$ un schéma localement noethérien. Soit
$F : (\Sch/S)_{fppf}^{opp} \to \textit{Ensembles}$ un foncteur. Soit $k$ un
corps de type fini sur $S$. Soit $x_0 \in F(\Spec(k))$.
La catégorie de prédéformations associée (\ref{artin-equation-predeformation-category})
correspond au foncteur
$$
F_{k, x_0} : \mathcal{C}_\Lambda \longrightarrow \textit{Ensembles},
\quad
A \longmapsto \{x \in F(\Spec(A)) \mid x|_{\Spec(k)} = x_0 \}.
$$
Rappelons que nous ne distinguons pas
les catégories cofibrées en ensembles sur $\mathcal{C}_\Lambda$
des foncteurs $\mathcal{C}_\Lambda \to \textit{Ensembles}$,
voir Théorie formelle des déformations, Remarques
\ref{formal-defos-remarks-cofibered-groupoids}
(\ref{formal-defos-item-convention-cofibered-sets}).
Étant donnée une transformation de foncteurs $a : F \to G$, en posant
$y_0 = a(x_0)$, nous obtenons un morphisme
$$
F_{k, x_0} \longrightarrow G_{k, y_0}
$$
voir (\ref{artin-equation-functoriality}).
Le Lemme \ref{artin-lemma-formally-smooth-on-deformation-categories} nous dit que si
$a : F \to G$ est formellement lisse (au sens de
Compléments sur les morphismes d'espaces, Définition
\ref{spaces-more-morphisms-definition-formally-smooth-etale-unramified}), alors
$F_{k, x_0} \longrightarrow G_{k, y_0}$ est lisse au sens de
Théorie formelle des déformations, Remarque
\ref{formal-defos-remark-compare-smooth-schlessinger}.

\medskip\noindent
Le Lemme \ref{artin-lemma-pushout} dit que si $Y' = Y \amalg_X X'$ dans la
catégorie des schémas sur $S$, où $X \to X'$ est un épaississement et
$X \to Y$ est affine, alors l'application
$$
F(Y \amalg_X X') \to F(Y) \times_{F(X)} F(X')
$$
est une bijection, pourvu que $F$ soit un espace algébrique.
Nous disons qu'un foncteur $F$ vérifie la {\it condition de Rim-Schlessinger}
ou que $F$ {\it vérifie (RS)} si, pour toute
somme amalgamée $Y' = Y \amalg_X X'$, où $Y, X, X'$ sont les spectres d'anneaux
artiniens locaux de type fini sur $S$, l'application
$$
F(Y \amalg_X X') \to F(Y) \times_{F(X)} F(X')
$$
est une bijection. Ainsi, tout espace algébrique vérifie (RS).

\medskip\noindent
Le Lemme \ref{artin-lemma-deformation-category} dit que, si
un foncteur $F$ vérifie (RS), alors tous les $F_{k, x_0}$
sont des foncteurs de déformations au sens de
Théorie formelle des déformations, Définition
\ref{formal-defos-definition-deformation-category}, c'est-à-dire qu'ils vérifient
la condition (RS) au sens de
Théorie formelle des déformations, Remarque
\ref{formal-defos-remark-compare-schlessinger-H4}.
En particulier, l'espace tangent
$$
TF_{k, x_0} = \{x \in F(\Spec(k[\epsilon])) \mid x|_{\Spec(k)} = x_0\}
$$
est muni d'une structure de $k$-espace vectoriel par Théorie formelle des déformations,
Lemme \ref{formal-defos-lemma-tangent-space-vector-space}.

\medskip\noindent
Le Lemme \ref{artin-lemma-finite-dimension} dit qu'un espace algébrique $F$
localement de type fini sur $S$ donne naissance à des foncteurs de déformations
$F_{k, x_0}$ dont les espaces tangents $TF_{k, x_0}$ sont de dimension finie.

\medskip\noindent
Un {\it objet formel}\footnote{C'est ce qu'Artin appelle une déformation formelle.}
$\xi = (R, \xi_n)$ de $F$ consiste en une algèbre noethérienne
locale complète sur $S$, notée $R$, dont le corps résiduel est de type fini
sur $S$, munie d'éléments $\xi_n \in F(\Spec(R/\mathfrak m^n))$
tels que $\xi_{n + 1}|_{\Spec(R/\mathfrak m^n)} = \xi_n$. Un objet
formel $\xi$ définit un objet formel $\xi$ de $F_{R/\mathfrak m, \xi_1}$.
Nous disons que $\xi$ est {\it versel} si et seulement s'il est versel au sens de
Théorie formelle des déformations, Définition \ref{formal-defos-definition-versal}.
Un objet formel $\xi = (R, \xi_n)$ est dit {\it effectif}
s'il existe $x \in F(\Spec(R))$ tel que
$\xi_n = x|_{\Spec(R/\mathfrak m^n)}$ pour tout $n \geq 1$.
Le Lemme \ref{artin-lemma-effective} dit que, si $F$ est un espace algébrique,
alors tout objet formel est effectif.

\medskip\noindent
Soit $U$ un schéma localement de type fini sur $S$ et soit $x \in F(U)$.
Soit $u_0 \in U$ un point de type fini. Nous disons que $x$ est versel en $u_0$
si et seulement si
$\xi = (\mathcal{O}_{U, u_0}^\wedge,
x|_{\Spec(\mathcal{O}_{U, u_0}/\mathfrak m_{u_0}^n)})$
est un objet formel versel au sens décrit ci-dessus.

\medskip\noindent
Soit $S$ un schéma localement noethérien. Soit
$F : (\Sch/S)_{fppf}^{opp} \to \textit{Ensembles}$ un foncteur.
Voici les axiomes que nous considérerons pour $F$.
\begin{enumerate}
\item[{[-1]}] une condition ensembliste\footnote{La condition est la
suivante : la borne supérieure de tous les cardinaux
$|F(\Spec(k))|$, où $k$ parcourt les corps de type
fini sur $S$, est inférieur ou égal ($\leq$) au cardinal d'un certain
objet de $(\Sch/S)_{fppf}$.} que peuvent ignorer les
lecteurs que les questions ensemblistes n'intéressent pas,
\item[{[0]}] $F$ est un faisceau pour la topologie \'etale,
\item[{[1]}] $F$ commute aux limites,
\item[{[2]}] $F$ vérifie la condition de Rim-Schlessinger (RS),
\item[{[3]}] tout espace tangent $TF_{k, x_0}$ est de dimension finie,
\item[{[4]}] tout objet formel est effectif,
\item[{[5]}] $F$ vérifie l'ouverture de la versalité.
\end{enumerate}
Ici, l'expression {\it commute aux limites} désigne la notion définie dans
Limites d'espaces, Définition
\ref{spaces-limits-definition-locally-finite-presentation}, et
{\it ouverture de la versalité} signifie ce qui suit : étant donnés un schéma $U$
localement de type fini sur $S$, un élément $x \in F(U)$ et
un point de type fini $u_0 \in U$ tel que $x$ soit versel en $u_0$,
il existe alors un voisinage ouvert $u_0 \in U' \subset U$
tel que $x$ soit versel en tout point de type fini de $U'$.






\section{Espaces algébriques}
\label{artin-section-algebraic-spaces}

\noindent
Voici notre premier résultat principal sur les espaces algébriques.

\begin{proposition}
\label{artin-proposition-spaces-diagonal-representable}
Soit $S$ un schéma localement noethérien. Soit
$F : (\Sch/S)_{fppf}^{opp} \to \textit{Ensembles}$ un foncteur. Supposons que
\begin{enumerate}
\item $\Delta : F \to F \times F$ soit représentable par des espaces algébriques,
\item $F$ vérifie les axiomes [-1], [0], [1], [2], [3], [4], [5]
(voir section \ref{artin-section-axioms-functors}), et
\item $\mathcal{O}_{S, s}$ soit un anneau G pour tout point $s$ de type fini de $S$.
\end{enumerate}
Alors $F$ est un espace algébrique.
\end{proposition}

\begin{proof}
Le Lemme \ref{artin-lemma-get-smooth} s'applique à $F$. Nous pouvons donc
choisir, pour tout corps $k$ de type fini sur $S$ et tout $x_0 \in F(\Spec(k))$,
un schéma affine $U_{k, x_0}$ de type fini sur $S$ et un morphisme lisse
$U_{k, x_0} \to F$ tel qu'il existe un point de type fini
$u_{k, x_0} \in U_{k, x_0}$ de corps résiduel $k$ tel que $x_0$
soit l'image de $u_{k, x_0}$. Alors
$$
U = \coprod\nolimits_{k, x_0} U_{k, x_0} \longrightarrow F
$$
est lisse\footnote{Remarque ensembliste : ce coproduit est (isomorphe à)
un objet de $(\Sch/S)_{fppf}$, puisque l'axiome [-1] nous donne une borne
sur l'ensemble d'indices ; voir Ensembles, Lemme \ref{sets-lemma-what-is-in-it}.}.
Pour achever la démonstration, il suffit de montrer que cette application est surjective ;
voir Fondements, Lemme \ref{bootstrap-lemma-spaces-etale-smooth-cover}
(c'est ici que nous utilisons l'axiome [0]). Par Critères de représentabilité, Lemme
\ref{criteria-lemma-check-property-limit-preserving},
il suffit de montrer que $U \times_F V \to V$ est surjectif pour les morphismes
$V \to F$ où $V$ est un schéma affine localement de présentation finie
sur $S$. Comme $U \times_F V \to V$ est lisse, son image est ouverte. Il
suffit donc de montrer que l'image de $U \times_F V \to V$ contient tous les
points de type fini de $V$ ; voir
Morphismes, Lemme \ref{morphisms-lemma-enough-finite-type-points}.
Soit $v_0 \in V$ un point de type fini. Alors $k = \kappa(v_0)$ est
un corps de type fini sur $S$. Notons $x_0$ le composé
$\Spec(k) \xrightarrow{v_0} V \to F$. Alors
$(u_{k, x_0}, v_0) : \Spec(k) \to U \times_F V$ est un point s'envoyant sur
$v_0$, ce qui donne le résultat.
\end{proof}

\begin{lemma}
\label{artin-lemma-monomorphism}
Soit $S$ un schéma localement noethérien. Soit $a : F \to G$ une transformation
de foncteurs $(\Sch/S)_{fppf}^{opp} \to \textit{Ensembles}$.
Supposons que
\begin{enumerate}
\item $a$ soit injective,
\item $F$ vérifie les axiomes [0], [1], [2], [4] et [5],
\item $\mathcal{O}_{S, s}$ soit un anneau G pour tout point $s$ de type fini de $S$,
\item $G$ soit un espace algébrique localement de type fini sur $S$,
\end{enumerate}
Alors $F$ est un espace algébrique.
\end{lemma}

\begin{proof}
Par le Lemme \ref{artin-lemma-finite-dimension}, le foncteur $G$ vérifie [3].
Comme $F \to G$ est injectif, nous en concluons que $F$ vérifie aussi [3].
De plus, comme $F \to G$ est injectif, étant donnés des schémas
$U$, $V$ et des morphismes $U \to F$ et $V \to F$, nous avons
$U \times_F V = U \times_G V$. Ainsi, $\Delta : F \to F \times F$ est
représentable (par des schémas), puisque c'est le cas de $G$ par hypothèse.
La Proposition \ref{artin-proposition-spaces-diagonal-representable}
s'applique donc\footnote{La condition
ensembliste [-1] est satisfaite par $F$, puisqu'elle l'est par $G$. Détails
omis.}.
\end{proof}












\section{Champs algébriques}
\label{artin-section-algebraic-stacks}

\noindent
La Proposition \ref{artin-proposition-second-diagonal-representable} est notre premier
résultat principal sur les champs algébriques.

\begin{lemma}
\label{artin-lemma-diagonal-representable}
Soit $S$ un schéma localement noethérien. Soit
$p : \mathcal{X} \to (\Sch/S)_{fppf}$ une catégorie fibrée en groupoïdes.
Supposons que
\begin{enumerate}
\item $\Delta : \mathcal{X} \to \mathcal{X} \times \mathcal{X}$
soit représentable par des espaces algébriques,
\item $\mathcal{X}$ vérifie les axiomes [-1], [0], [1], [2], [3] (voir
section \ref{artin-section-axioms}),
\item tout objet formel de $\mathcal{X}$ soit effectif,
\item $\mathcal{X}$ vérifie l'ouverture de la versalité, et
\item $\mathcal{O}_{S, s}$ soit un anneau G pour tout point $s$ de type fini de $S$.
\end{enumerate}
Alors $\mathcal{X}$ est un champ algébrique.
\end{lemma}

\begin{proof}
Le Lemme \ref{artin-lemma-get-smooth} s'applique à $\mathcal{X}$. Nous pouvons donc
choisir, pour tout corps $k$ de type fini sur $S$ et toute
classe d'isomorphisme d'un objet $x_0 \in \Ob(\mathcal{X}_{\Spec(k)})$,
un schéma affine $U_{k, x_0}$ de type fini sur $S$ et un morphisme lisse
$(\Sch/U_{k, x_0})_{fppf} \to \mathcal{X}$ tel qu'il existe un point de type
fini $u_{k, x_0} \in U_{k, x_0}$ de corps résiduel $k$ tel que $x_0$
soit l'image de $u_{k, x_0}$. Alors
$$
(\Sch/U)_{fppf} \to \mathcal{X},
\quad\text{avec}\quad
U = \coprod\nolimits_{k, x_0} U_{k, x_0}
$$
est lisse\footnote{Remarque ensembliste : ce coproduit est (isomorphe à)
un objet de $(\Sch/S)_{fppf}$, puisque l'axiome [-1] nous donne une borne
sur l'ensemble d'indices ; voir Ensembles, Lemme \ref{sets-lemma-what-is-in-it}.}.
Pour achever la démonstration, il suffit de montrer que cette application est surjective ;
voir Critères de représentabilité, Lemme \ref{criteria-lemma-stacks-etale}
(c'est ici que nous utilisons l'axiome [0]). Par Critères de représentabilité, Lemme
\ref{criteria-lemma-check-property-limit-preserving},
il suffit de montrer que
$(\Sch/U)_{fppf} \times_\mathcal{X} (\Sch/V)_{fppf} \to (\Sch/V)_{fppf}$
est surjectif pour les morphismes $y : (\Sch/V)_{fppf} \to \mathcal{X}$ où
$V$ est un schéma affine localement de présentation finie
sur $S$. Par l'hypothèse (1), le produit fibré
$(\Sch/U)_{fppf} \times_\mathcal{X} (\Sch/V)_{fppf}$ est représentable
par un espace algébrique $W$. Alors $W \to V$ est lisse, donc son image est
ouverte. Il suffit donc de montrer que l'image de $W \to V$ contient tous les
points de type fini de $V$ ; voir
Morphismes, Lemme \ref{morphisms-lemma-enough-finite-type-points}.
Soit $v_0 \in V$ un point de type fini. Alors $k = \kappa(v_0)$ est
un corps de type fini sur $S$. Notons $x_0 = y|_{\Spec(k)}$
l'image réciproque de $y$ par $v_0$. Alors $(u_{k, x_0}, v_0)$ fournit
un morphisme $\Spec(k) \to W$ dont le composé avec $W \to V$
est $v_0$, ce qui donne le résultat.
\end{proof}

\begin{proposition}
\label{artin-proposition-second-diagonal-representable}
Soit $S$ un schéma localement noethérien. Soit
$p : \mathcal{X} \to (\Sch/S)_{fppf}$ une catégorie fibrée en groupoïdes.
Supposons que
\begin{enumerate}
\item $\Delta_\Delta : \mathcal{X} \to
\mathcal{X} \times_{\mathcal{X} \times \mathcal{X}} \mathcal{X}$
soit représentable par des espaces algébriques,
\item $\mathcal{X}$ vérifie les axiomes [-1], [0], [1], [2], [3], [4] et [5]
(voir section \ref{artin-section-axioms}),
\item $\mathcal{O}_{S, s}$ soit un anneau G pour tout point $s$ de type fini de $S$.
\end{enumerate}
Alors $\mathcal{X}$ est un champ algébrique.
\end{proposition}

\begin{proof}
Nous montrons d'abord que $\Delta : \mathcal{X} \to \mathcal{X} \times \mathcal{X}$
est représentable par des espaces algébriques. Pour cela, il suffit de montrer
que
$$
\mathcal{Y} =
\mathcal{X} \times_{\Delta, \mathcal{X} \times \mathcal{X}, y} (\Sch/V)_{fppf}
$$
est représentable par un espace algébrique pour tout schéma affine $V$ localement
de présentation finie sur $S$ et tout objet $y$ de
$\mathcal{X} \times \mathcal{X}$ au-dessus de $V$ ; voir
Critères de représentabilité, Lemme
\ref{criteria-lemma-check-representable-limit-preserving}\footnote{La
condition ensembliste de Critères de représentabilité, Lemme
\ref{criteria-lemma-check-representable-limit-preserving},
sera satisfaite : la taille de l'espace algébrique $Y$ représentant $\mathcal{Y}$ est
convenablement bornée. En effet, $Y \to S$ sera localement de type fini et $Y$
vérifiera l'axiome [-1]. Détails omis.}.
Observons que $\mathcal{Y}$ est fibrée en catégories discrètes
(Champs, Lemme \ref{stacks-lemma-isom-as-2-fibre-product})
et soit $Y : (\Sch/S)_{fppf}^{opp} \to \textit{Ensembles}$ le foncteur
$T \mapsto \Ob(\mathcal{Y}_T)/\cong$ des classes d'isomorphisme.
Nous appliquerons la
Proposition \ref{artin-proposition-spaces-diagonal-representable}
pour voir que $Y$ est un espace algébrique.

\medskip\noindent
Notons que
$\Delta_\mathcal{Y} : \mathcal{Y} \to \mathcal{Y} \times \mathcal{Y}$
(et donc aussi $Y \to Y \times Y$)
est représentable par des espaces algébriques par la condition (1) et
Critères de représentabilité, Lemme \ref{criteria-lemma-second-diagonal}.
Observons que $Y$ est un faisceau pour la topologie \'etale par
Champs, Lemmes \ref{stacks-lemma-stack-in-setoids-characterize} et
\ref{stacks-lemma-2-fibre-product-gives-stack-in-setoids}, c'est-à-dire que
l'axiome [0] est satisfait. De plus, $Y$ commute aux limites par
le Lemme \ref{artin-lemma-fibre-product-limit-preserving}, c'est-à-dire que [1] est satisfait.
Notons que $Y$ vérifie (RS), c'est-à-dire l'axiome [2], par les
Lemmes \ref{artin-lemma-algebraic-stack-RS} et
\ref{artin-lemma-fibre-product-RS}. L'axiome [3] pour $Y$ résulte
des Lemmes \ref{artin-lemma-finite-dimension} et
\ref{artin-lemma-fibre-product-tangent-spaces}.
L'axiome [4] résulte des Lemmes \ref{artin-lemma-effective} et
\ref{artin-lemma-fibre-product-effective}.
L'axiome [5] pour $Y$ résulte directement de l'ouverture de la versalité
pour $\Delta_\mathcal{X}$, qui fait partie de l'axiome [5] pour $\mathcal{X}$.
Ainsi, toutes les hypothèses de la
Proposition \ref{artin-proposition-spaces-diagonal-representable}
sont satisfaites et $Y$ est un espace algébrique.

\medskip\noindent
À ce stade, il résulte du Lemme \ref{artin-lemma-diagonal-representable}
que $\mathcal{X}$ est un champ algébrique.
\end{proof}





\section{Condition de Rim-Schlessinger forte}
\label{artin-section-RS-star}

\noindent
Dans la suite de ce chapitre, la version strictement plus forte suivante
des conditions de Rim-Schlessinger jouera un rôle important.

\begin{definition}
\label{artin-definition-RS-star}
Soit $S$ un schéma. Soit $\mathcal{X}$ une catégorie
fibrée en groupoïdes sur $(\Sch/S)_{fppf}$. Nous disons que $\mathcal{X}$
vérifie la {\it condition (RS*)} si, pour tout diagramme de produit fibré
$$
\xymatrix{
B' \ar[r] & B \\
A' = A \times_B B' \ar[u] \ar[r] & A \ar[u]
}
$$
de $S$-algèbres, où $B' \to B$ est surjectif et de noyau de carré nul,
le foncteur entre les catégories fibres
$$
\mathcal{X}_{\Spec(A')}
\longrightarrow
\mathcal{X}_{\Spec(A)} \times_{\mathcal{X}_{\Spec(B)}} \mathcal{X}_{\Spec(B')}
$$
est une équivalence de catégories.
\end{definition}

\noindent
Faisons quelques observations :
pour $A \to B \leftarrow B'$ comme dans la Définition \ref{artin-definition-RS-star},
\begin{enumerate}
\item nous avons $\Spec(A') = \Spec(A) \amalg_{\Spec(B)} \Spec(B')$
dans la catégorie des schémas ; voir
Compléments sur les morphismes, Lemme \ref{more-morphisms-lemma-pushout-along-thickening},
et
\item si $\mathcal{X}$ est un champ algébrique, alors $\mathcal{X}$ vérifie
(RS*) par le Lemme \ref{artin-lemma-algebraic-stack-RS-star}.
\end{enumerate}
Si $S$ est localement noethérien, alors
\begin{enumerate}
\item[(3)] si $A$, $B$, $B'$ sont de type fini sur $S$ et si
$B$ est fini sur $A$, alors $A'$ est de type fini sur
$S$\footnote{Si $\Spec(A)$ s'envoie dans un ouvert affine de $S$,
cela résulte de
Compléments d'algèbre, Lemme \ref{more-algebra-lemma-fibre-product-finite-type}.
Le cas général résulte du
Lemme \ref{more-algebra-lemma-diagram-localize} de Compléments d'algèbre.}, et
\item[(4)] si $\mathcal{X}$ vérifie (RS*), alors $\mathcal{X}$ vérifie (RS),
car (RS) recouvre exactement les cas de (RS*) où
$A$, $B$, $B'$ sont artiniens locaux.
\end{enumerate}

\begin{lemma}
\label{artin-lemma-algebraic-stack-RS-star}
Soit $\mathcal{X}$ un champ algébrique sur une base $S$.
Alors $\mathcal{X}$ vérifie (RS*).
\end{lemma}

\begin{proof}
Cela résulte du Lemme \ref{artin-lemma-pushout} ; voir les
remarques qui suivent la Définition \ref{artin-definition-RS-star}.
\end{proof}

\begin{lemma}
\label{artin-lemma-fibre-product-RS-star}
Soit $S$ un schéma. Soient $p : \mathcal{X} \to \mathcal{Y}$ et
$q : \mathcal{Z} \to \mathcal{Y}$ des $1$-morphismes de catégories
fibrées en groupoïdes sur $(\Sch/S)_{fppf}$. Si $\mathcal{X}$, $\mathcal{Y}$
et $\mathcal{Z}$ vérifient (RS*), alors il en va de même de
$\mathcal{X} \times_\mathcal{Y} \mathcal{Z}$.
\end{lemma}

\begin{proof}
La démonstration est exactement la même que celle du
Lemme \ref{artin-lemma-fibre-product-RS}.
\end{proof}








\section{Versalité et générisations}
\label{artin-section-generalize-versality}

\noindent
Nous montrons que la versalité se conserve par générisation
pour les champs qui vérifient (RS*) et commutent aux limites.
Nous suggérons d'omettre cette section en première lecture.

\begin{lemma}
\label{artin-lemma-single-point}
Soit $S$ un schéma localement noethérien. Soit $\mathcal{X}$ une catégorie
fibrée en groupoïdes sur $(\Sch/S)_{fppf}$ vérifiant (RS*).
Soit $x$ un objet de $\mathcal{X}$ au-dessus d'un schéma affine $U$
de type fini sur $S$. Soit $u \in U$ un point de type fini tel que
$x$ ne soit pas versel en $u$. Il existe alors un morphisme $x \to y$
de $\mathcal{X}$ au-dessus de $U \to T$ satisfaisant aux conditions suivantes :
\begin{enumerate}
\item le morphisme $U \to T$ est un épaississement du premier ordre,
\item nous avons une suite exacte courte
$$
0 \to \kappa(u) \to \mathcal{O}_T \to \mathcal{O}_U \to 0
$$
\item il n'existe {\bf pas} de couple $(W, \alpha)$
formé d'un voisinage ouvert $W \subset T$ de $u$
et d'un morphisme $\beta : y|_W \to x$ tel que le composé
$$
x|_{U \cap W} \xrightarrow{\text{restriction de }x \to y}
y|_W \xrightarrow{\beta} x
$$
soit le morphisme canonique $x|_{U \cap W} \to x$.
\end{enumerate}
\end{lemma}

\begin{proof}
Posons $R = \mathcal{O}_{U, u}^\wedge$. Soit $k = \kappa(u)$
le corps résiduel de $R$. Soit $\xi$ l'objet formel
de $\mathcal{X}$ au-dessus de $R$ associé à $x$. Comme $x$ n'est pas
versel en $u$, nous voyons que $\xi$ n'est pas versel ; voir le
Lemme \ref{artin-lemma-versality-matches}. D'après la discussion qui suit la
Définition \ref{artin-definition-versal-formal-object},
cela signifie que nous pouvons trouver des
morphismes $\xi_1 \to x_A \to x_B$ de $\mathcal{X}$ au-dessus des
immersions fermées $\Spec(k) \to \Spec(A) \to \Spec(B)$,
où $A, B$ sont des anneaux artiniens locaux de corps résiduel $k$,
un entier $n \geq 1$ et un diagramme commutatif
$$
\vcenter{
\xymatrix{
& x_A \ar[ld] \\
\xi_n & \xi_1 \ar[u] \ar[l]
}
}
\quad\text{au-dessus de}\quad
\vcenter{
\xymatrix{
& \Spec(A) \ar[ld] \\
\Spec(R/\mathfrak m^n) & \Spec(k) \ar[u] \ar[l]
}
}
$$
tels qu'il n'existe {\bf pas} d'entier $m \geq n$ et de diagramme commutatif
$$
\vcenter{
\xymatrix{
& & x_B \ar[lldd] \\
& & x_A \ar[ld] \ar[u] \\
\xi_m & \xi_n \ar[l] & \xi_1 \ar[u] \ar[l]
}
}
\quad\text{au-dessus de}
\vcenter{
\xymatrix{
& & \Spec(B) \ar[lldd] \\
& & \Spec(A) \ar[ld] \ar[u] \\
\Spec(R/\mathfrak m^m) &
\Spec(R/\mathfrak m^n) \ar[l] &
\Spec(k) \ar[u] \ar[l]
}
}
$$
Nous pouvons en outre supposer que $B \to A$ est une petite
extension, c'est-à-dire que le noyau $I$ de la surjection $B \to A$
est isomorphe à $k$ comme $A$-module.
Cela résulte de Théorie formelle des déformations, Remarque
\ref{formal-defos-remark-versal-object}.
Nous définissons alors simplement
$$
T = U \amalg_{\Spec(A)} \Spec(B)
$$
Grâce à la propriété (RS*), nous trouvons $y$ au-dessus de $T$ dont la restriction à
$\Spec(B)$ est $x_B$ et dont la restriction à $U$ est $x$
(cela donne la flèche $x \to y$ au-dessus de $U \to T$).
Pour achever la démonstration, vérifions les conditions (1), (2) et (3).

\medskip\noindent
Par construction de la somme amalgamée, nous avons un diagramme commutatif
$$
\xymatrix{
0 \ar[r] &
I \ar[r] &
B \ar[r] &
A \ar[r] &
0 \\
0 \ar[r] &
I \ar[r] \ar[u] &
\Gamma(T, \mathcal{O}_T) \ar[r] \ar[u] &
\Gamma(U, \mathcal{O}_U) \ar[r] \ar[u] &
0
}
$$
à lignes exactes. Cela démontre immédiatement (1) et (2).
Pour achever la démonstration, nous allons raisonner par l'absurde.
Supposons donné un couple $(W, \beta)$ comme en (3).
Comme $\Spec(B) \to T$ se factorise par $W$, nous obtenons le morphisme
$$
x_B \to y|_W \xrightarrow{\beta} x
$$
Comme $B$ est artinien local de corps résiduel $k = \kappa(u)$,
nous voyons que $x_B \to x$ est au-dessus d'un morphisme $\Spec(B) \to U$
qui se factorise par $\Spec(\mathcal{O}_{U, u}/\mathfrak m_u^m)$
pour un certain $m \geq n$. Autrement dit, $x_B \to x$ se factorise
par $\xi_m$, ce qui donne un morphisme $x_B \to \xi_m$.
La condition de compatibilité imposée au morphisme $\alpha$
dans la condition (3) se traduit par la commutativité du diagramme
$$
\xymatrix{
x_B \ar[d] & x_A \ar[d] \ar[l] \\
\xi_m & \xi_n \ar[l]
}
$$
Cela donne la contradiction recherchée.
\end{proof}

\begin{lemma}
\label{artin-lemma-generalization-versality}
Soit $S$ un schéma localement noethérien. Soit $\mathcal{X}$ une catégorie fibrée
en groupoïdes sur $(\Sch/S)_{fppf}$. Supposons que
\begin{enumerate}
\item $\Delta : \mathcal{X} \to \mathcal{X} \times \mathcal{X}$ soit
représentable par des espaces algébriques,
\item $\mathcal{X}$ vérifie (RS*),
\item $\mathcal{X}$ commute aux limites.
\end{enumerate}
Soit $x$ un objet de $\mathcal{X}$ au-dessus d'un schéma $U$ de type fini sur
$S$. Soit $u \leadsto u_0$ une spécialisation de points de type fini de $U$
telle que $x$ soit versel en $u_0$. Alors $x$ est versel en $u$.
\end{lemma}

\begin{proof}
Après avoir rétréci $U$, nous pouvons supposer que $U$ est affine et que $U$ s'envoie dans un
ouvert affine $\Spec(\Lambda)$ de $S$. Si $x$ n'est pas versel en $u$, nous
pouvons choisir $x \to y$ au-dessus de $U \to T$ comme dans le
Lemme \ref{artin-lemma-single-point}. Écrivons $U = \Spec(R_0)$ et $T = \Spec(R)$.
Le morphisme $U \to T$ correspond à un homomorphisme surjectif d'anneaux
$R \to R_0$ dont le noyau est un idéal de carré nul.
Par l'hypothèse (3), $y$ provient d'un objet $x'$ au-dessus de
$U' = \Spec(R')$, pour une certaine $\Lambda$-sous-algèbre de type fini
$R' \subset R$. En agrandissant $R'$, nous pouvons supposer, et supposons désormais, que
$R' \to R_0$ est surjectif, de sorte que $U \subset U'$ est un épaississement du premier ordre.
Nous avons donc maintenant
$$
x \to y \to x'
\text{ au-dessus de }
U \to T \to U'
$$
Par l'hypothèse (1), il existe un espace algébrique $Z$ sur $S$ représentant
$$
(\Sch/U)_{fppf} \times_{x, \mathcal{X}, x'} (\Sch/U')_{fppf}
$$
voir Champs algébriques, Lemme \ref{algebraic-lemma-representable-diagonal}.
Par construction des $2$-produits fibrés, un point de $Z$ à valeurs dans $V$
correspond à un triplet $(a, a', \alpha)$ formé de morphismes
$a : V \to U$, $a' : V \to U'$ et d'un morphisme $\alpha : a^*x \to (a')^*x'$.
Nous obtenons un diagramme commutatif
$$
\xymatrix{
U \ar[rd] \ar[rdd] \ar[rrd] \\
& Z \ar[r]_{p'} \ar[d]^p & U' \ar[d] \\
& U \ar[r] & S
}
$$
Le morphisme $i : U \to Z$ provient de l'isomorphisme $x \to x'|_U$.
Soit $z_0 = i(u_0) \in Z$. Par le Lemme \ref{artin-lemma-base-change-versal},
nous voyons que $Z \to U'$ est lisse en $z_0$. En remplaçant $U$ par un
voisinage ouvert affine de $u_0$, $U'$ par l'ouvert correspondant
et $Z$ par l'intersection des images réciproques
de ces ouverts par $p$ et $p'$, nous nous ramenons à la situation où
$Z \to U'$ est lisse le long de $i(U)$. Comme $u \leadsto u_0$, le point
$u$ appartient à cet ouvert. La condition (3) du Lemme \ref{artin-lemma-single-point}
est manifestement préservée lorsque l'on rétrécit $U$ (tous les schémas $U$, $T$, $U'$
ont le même espace topologique sous-jacent).
Comme $U \to U'$ est un épaississement du premier ordre de schémas affines,
nous pouvons choisir un morphisme $i' : U' \to Z$
tel que $p' \circ i' = \text{id}_{U'}$ et
dont la restriction à $U$ soit $i$
(Compléments sur les morphismes d'espaces, Lemme
\ref{spaces-more-morphisms-lemma-smooth-formally-smooth}).
En prenant par $i'$ l'image réciproque du morphisme universel $p^*x \to (p')^*x'$,
nous obtenons un morphisme
$$
x' \to x
$$
au-dessus de $p \circ i' : U' \to U$ tel que le composé
$$
x \to x' \to x
$$
soit l'identité. Rappelons que nous avons $y \to x'$ au-dessus du
morphisme $T \to U'$. En composant, nous obtenons un morphisme
$y \to x$ dont l'existence contredit la condition
(3) du Lemme \ref{artin-lemma-single-point}.
Cette contradiction achève la démonstration.
\end{proof}






\section{Effectivité formelle forte}
\label{artin-section-strong-formal-effectiveness}

\noindent
Dans cette section, nous montrons comment une version forte de l'effectivité
des objets formels implique l'ouverture de la versalité. La démonstration de
\cite[théorème 1.1]{Bhatt-Algebraize} montre que les espaces algébriques quasi-compacts et
quasi-séparés vérifient l'effectivité formelle forte
décrite dans la Remarque \ref{artin-remark-strong-effectiveness}. En outre, la
théorie que nous développons n'est pas vide : nous l'utiliserons plus loin pour démontrer l'ouverture de la versalité
pour le champ des faisceaux cohérents et pour le champ de modules de complexes ; voir
Quot, Théorèmes
\ref{quot-theorem-coherent-algebraic-general} et
\ref{quot-theorem-complexes-algebraic}.

\begin{lemma}
\label{artin-lemma-infinite-sequence-pre}
Soit $S$ un schéma localement noethérien. Soit $\mathcal{X}$ une catégorie
fibrée en groupoïdes sur $(\Sch/S)_{fppf}$ vérifiant (RS*).
Soit $x$ un objet de
$\mathcal{X}$ au-dessus d'un schéma affine $U$ de type fini sur $S$.
Soient $u_n \in U$, $n \geq 1$, des points de type fini tels que
(a) il n'existe aucune spécialisation $u_n \leadsto u_m$ pour $n \not = m$, et
(b) $x$ ne soit versel en $u_n$ pour tout $n$. Il existe alors des morphismes
$$
x \to x_1 \to x_2 \to \ldots
\quad\text{dans }\mathcal{X}\text{ au-dessus de }\quad
U \to U_1 \to U_2 \to \ldots
$$
sur $S$ tels que
\begin{enumerate}
\item pour tout $n$, le morphisme $U \to U_n$ soit un épaississement du premier
ordre,
\item pour tout $n$, on ait une suite exacte courte
$$
0 \to \kappa(u_n) \to \mathcal{O}_{U_n} \to \mathcal{O}_{U_{n - 1}} \to 0
$$
avec $U_0 = U$ pour $n = 1$,
\item pour tout $n$, il n'existe {\bf pas} de couple $(W, \alpha)$
formé d'un voisinage ouvert $W \subset U_n$ de $u_n$
et d'un morphisme $\alpha : x_n|_W \to x$
tel que le composé
$$
x|_{U \cap W} \xrightarrow{\text{restriction de }x \to x_n}
x_n|_W \xrightarrow{\alpha} x
$$
soit le morphisme canonique $x|_{U \cap W} \to x$.
\end{enumerate}
\end{lemma}

\begin{proof}
Comme il n'y a pas de spécialisations entre les points $u_n$ (et que les
$u_n$ sont donc en particulier deux à deux distincts), pour tout $n$
nous pouvons trouver un ouvert $U' \subset U$
tel que $u_n \in U'$ et $u_i \not \in U'$ pour $i = 1, \ldots, n - 1$.
Par le Lemme \ref{artin-lemma-single-point}, pour tout $n \geq 1$, nous pouvons trouver
$$
x \to y_n
\quad\text{dans }\mathcal{X}\text{ au-dessus de}\quad
U \to T_n
$$
tels que
\begin{enumerate}
\item le morphisme $U \to T_n$ soit un épaississement du premier ordre,
\item on ait une suite exacte courte
$$
0 \to \kappa(u_n) \to \mathcal{O}_{T_n} \to \mathcal{O}_U \to 0
$$
\item il n'existe {\bf pas} de couple $(W, \alpha)$
formé d'un voisinage ouvert $W \subset T_n$ de $u_n$
et d'un morphisme $\beta : y_n|_W \to x$ tel que le composé
$$
x|_{U \cap W} \xrightarrow{\text{restriction de }x \to y_n}
y_n|_W \xrightarrow{\beta} x
$$
soit le morphisme canonique $x|_{U \cap W} \to x$.
\end{enumerate}
Nous pouvons donc définir par récurrence
$$
U_1 = T_1, \quad
U_{n + 1} = U_n \amalg_U T_{n + 1}
$$
En posant $x_1 = y_1$ et en utilisant (RS*), nous trouvons par récurrence
$x_{n + 1}$ au-dessus de $U_{n + 1}$ dont la restriction à
$U_n$ est $x_n$ et dont la restriction à $T_{n + 1}$ est $y_{n + 1}$.
La propriété (1) de $U \to U_n$ résulte de la construction
de la somme amalgamée dans Compléments sur les morphismes, Lemme
\ref{more-morphisms-lemma-pushout-along-thickening}.
La propriété (2) pour $U_n$ résulte de même de
la propriété (2) pour $T_n$ et de la construction de la somme amalgamée.
Après restriction à un voisinage ouvert $U'$ de $u_n$
comme ci-dessus, la propriété (3) pour $(U_n, x_n)$ résulte de la propriété (3)
pour $(T_n, y_n)$, simplement parce que les sous-schémas ouverts correspondants
de $T_n$ et $U_n$ sont isomorphes. Nous omettons certains détails.
\end{proof}

\begin{remark}[Effectivité forte]
\label{artin-remark-strong-effectiveness}
Soit $S$ un schéma localement noethérien.
Soit $\mathcal{X}$ une catégorie fibrée en groupoïdes sur $(\Sch/S)_{fppf}$.
Supposons donnés
\begin{enumerate}
\item un ouvert affine $\Spec(\Lambda) \subset S$,
\item un système projectif $(R_n)$ de $\Lambda$-algèbres
à morphismes de transition surjectifs et dont les noyaux sont localement nilpotents,
\item un système $(\xi_n)$ d'objets de $\mathcal{X}$ au-dessus
du système $(\Spec(R_n))$.
\end{enumerate}
Dans cette situation, posons $R = \lim R_n$. Nous disons que
$(\xi_n)$ est {\it effectif} s'il existe un objet
$\xi$ de $\mathcal{X}$ au-dessus de $\Spec(R)$ dont la restriction
à $\Spec(R_n)$ redonne le système $(\xi_n)$.
\end{remark}

\noindent
Il n'est pas vrai que tout champ algébrique $\mathcal{X}$
sur $S$ vérifie un axiome d'effectivité forte de la forme suivante :
tout système $(\xi_n)$ comme dans la Remarque \ref{artin-remark-strong-effectiveness}
est effectif. Un exemple est donné dans
Exemples, section \ref{examples-section-non-formal-effectiveness}.

\begin{lemma}
\label{artin-lemma-SGE-implies-openness-versality}
Soit $S$ un schéma localement noethérien. Soit $\mathcal{X}$ une catégorie fibrée
en groupoïdes sur $(\Sch/S)_{fppf}$. Supposons que
\begin{enumerate}
\item $\Delta : \mathcal{X} \to \mathcal{X} \times \mathcal{X}$ soit
représentable par des espaces algébriques,
\item $\mathcal{X}$ vérifie (RS*),
\item $\mathcal{X}$ commute aux limites,
\item les systèmes $(\xi_n)$ comme dans la Remarque \ref{artin-remark-strong-effectiveness},
où $\Ker(R_m \to R_n)$ est un idéal de carré nul pour tous $m \geq n$,
soient effectifs.
\end{enumerate}
Alors $\mathcal{X}$ vérifie l'ouverture de la versalité.
\end{lemma}

\begin{proof}
Choisissons un schéma $U$ localement de type fini sur $S$,
un point $u_0$ de type fini de $U$ et un objet $x$ de $\mathcal{X}$
au-dessus de $U$, tel que $x$ soit versel en $u_0$. Après avoir rétréci
$U$, nous pouvons supposer que $U$ est affine et que $U$ s'envoie dans un ouvert affine
$\Spec(\Lambda)$ de $S$. Soit $E \subset U$ l'ensemble des points de type fini
$u$ tels que $x$ ne soit pas versel en $u$. D'après le
Lemme \ref{artin-lemma-generalization-versality},
si $u \in E$, alors $u_0$ n'est pas une spécialisation de $u$.
Si l'ouverture de la versalité n'est pas satisfaite, alors $u_0$ appartient à l'adhérence
$\overline{E}$ de $E$. Par le
Lemme \ref{properties-lemma-countable-dense-subset} de Propriétés,
nous pouvons choisir un sous-ensemble dénombrable $E' \subset E$ de même adhérence
que $E$. Par Propriétés, Lemme \ref{properties-lemma-maximal-points},
nous pouvons supposer qu'il n'existe aucune spécialisation entre les points de $E'$.
Observons que $E'$ doit être (dénombrable) infini, puisque $u_0$
n'est la spécialisation d'aucun point de $E'$, comme nous venons de le voir.
Nous pouvons donc écrire $E' = \{u_1, u_2, u_3, \ldots\}$ ; il
n'existe aucune spécialisation entre les $u_i$, et $u_0$ appartient à l'adhérence
de $E'$.

\medskip\noindent
Choisissons $x \to x_1 \to x_2 \to \ldots$ au-dessus de
$U \to U_1 \to U_2 \to \ldots$ comme dans le Lemme \ref{artin-lemma-infinite-sequence-pre}.
Écrivons $U_n = \Spec(R_n)$ et $U = \Spec(R_0)$.
Posons $R = \lim R_n$. Observons que $R \to R_0$ est surjectif
et que son noyau est un idéal de carré nul. Par l'hypothèse (4),
nous obtenons $\xi$ au-dessus de $\Spec(R)$ dont le changement de base à $R_n$ est $x_n$.
Par l'hypothèse (3), $\xi$ provient d'un objet $\xi'$ au-dessus de
$U' = \Spec(R')$, pour une certaine $\Lambda$-sous-algèbre de type fini
$R' \subset R$. En agrandissant $R'$, nous pouvons supposer, et supposons désormais, que
$R' \to R_0$ est surjectif, de sorte que $U \subset U'$ est un épaississement du premier ordre.
Nous avons donc maintenant
$$
x \to x_1 \to x_2 \to \ldots \to \xi'
\text{ au-dessus de }
U \to U_1 \to U_2 \to \ldots \to U'
$$
Par l'hypothèse (1), il existe un espace algébrique $Z$ sur $S$ représentant
$$
(\Sch/U)_{fppf} \times_{x, \mathcal{X}, \xi'} (\Sch/U')_{fppf}
$$
voir Champs algébriques, Lemme \ref{algebraic-lemma-representable-diagonal}.
Par construction des $2$-produits fibrés, un point de $Z$ à valeurs dans $T$
correspond à un triplet $(a, a', \alpha)$ formé de morphismes
$a : T \to U$, $a' : T \to U'$ et d'un morphisme $\alpha : a^*x \to (a')^*\xi'$.
Nous obtenons un diagramme commutatif
$$
\xymatrix{
U \ar[rd] \ar[rdd] \ar[rrd] \\
& Z \ar[r]_{p'} \ar[d]^p & U' \ar[d] \\
& U \ar[r] & S
}
$$
Le morphisme $i : U \to Z$ provient de l'isomorphisme $x \to \xi'|_U$.
Soit $z_0 = i(u_0) \in Z$. Par le Lemme \ref{artin-lemma-base-change-versal},
nous voyons que $Z \to U'$ est lisse en $z_0$. En remplaçant $U$ par un
voisinage ouvert affine de $u_0$, $U'$ par l'ouvert correspondant
et $Z$ par l'intersection des images réciproques
de ces ouverts par $p$ et $p'$, nous nous ramenons à la situation où
$Z \to U'$ est lisse le long de $i(U)$. Notons que cela
implique aussi de remplacer la suite $u_n$ par une sous-suite, à savoir
celle correspondant aux indices tels que $u_n$ appartienne à l'ouvert. De plus, la condition
(3) du Lemme \ref{artin-lemma-infinite-sequence-pre}
est manifestement préservée lorsque l'on rétrécit $U$
(tous les schémas $U$, $U_n$, $U'$ ont le même espace
topologique sous-jacent).
Comme $U \to U'$ est un épaississement du premier ordre de schémas affines,
nous pouvons choisir un morphisme $i' : U' \to Z$
tel que $p' \circ i' = \text{id}_{U'}$ et
dont la restriction à $U$ soit $i$
(Compléments sur les morphismes d'espaces, Lemme
\ref{spaces-more-morphisms-lemma-smooth-formally-smooth}).
En prenant l'image réciproque du morphisme universel
$p^*x \to (p')^*\xi'$ par $i'$, nous obtenons un morphisme
$$
\xi' \to x
$$
au-dessus de $p \circ i' : U' \to U$ tel que le composé
$$
x \to \xi' \to x
$$
soit l'identité. Rappelons que nous avons $x_1 \to \xi'$ au-dessus du
morphisme $U_1 \to U'$. En composant, nous obtenons un morphisme
$x_1 \to x$ dont l'existence contredit la condition
(3) du Lemme \ref{artin-lemma-infinite-sequence-pre}.
Cette contradiction achève la démonstration.
\end{proof}

\begin{remark}
\label{artin-remark-trade-openness-versality-diagonal-with-strong-effectiveness}
Il existe une manière de déduire l'ouverture de la versalité de la diagonale
d'une catégorie fibrée en groupoïdes à partir d'un axiome d'effectivité formelle
forte.
Soit $S$ un schéma localement noethérien. Soit $\mathcal{X}$ une catégorie fibrée
en groupoïdes sur $(\Sch/S)_{fppf}$. Supposons que
\begin{enumerate}
\item $\Delta_\Delta : \mathcal{X} \to
\mathcal{X} \times_{\mathcal{X} \times \mathcal{X}} \mathcal{X}$
soit représentable par des espaces algébriques,
\item $\mathcal{X}$ vérifie (RS*),
\item $\mathcal{X}$ commute aux limites,
\item pour tout système projectif $(R_n)$ de $S$-algèbres
comme dans la Remarque \ref{artin-remark-strong-effectiveness},
où $\Ker(R_m \to R_n)$ est un idéal de carré nul pour tous $m \geq n$,
le foncteur
$$
\mathcal{X}_{\Spec(\lim R_n)} \longrightarrow
\lim_n \mathcal{X}_{\Spec(R_n)}
$$
soit pleinement fidèle.
\end{enumerate}
Alors $\Delta : \mathcal{X} \to \mathcal{X} \times \mathcal{X}$
vérifie l'ouverture de la versalité. Cela résulte de l'application du
Lemme \ref{artin-lemma-SGE-implies-openness-versality}
aux produits fibrés de la forme
$\mathcal{X} \times_{\Delta, \mathcal{X} \times \mathcal{X}, y}
(\Sch/V)_{fppf}$, pour tout schéma affine $V$ localement
de présentation finie sur $S$ et tout objet $y$ de
$\mathcal{X} \times \mathcal{X}$ au-dessus de $V$.
Si jamais nous en avons besoin, nous transformerons cette remarque en
un lemme et en donner une démonstration détaillée.
\end{remark}











\section{Déformations infinitésimales}
\label{artin-section-inf}

\noindent
Dans cette section, nous étudions une généralisation de la notion
d'espace tangent introduite à la section \ref{artin-section-tangent-spaces}.
Pour le faire de façon naturelle, nous empruntons certaines notations à
Théorie formelle des déformations, sections
\ref{formal-defos-section-tangent-spaces-functors},
\ref{formal-defos-section-lifts} et
\ref{formal-defos-section-infinitesimal-automorphisms}.

\medskip\noindent
Soit $S$ un schéma. Soit $\mathcal{X}$ une catégorie fibrée en groupoïdes
au-dessus de $(\Sch/S)_{fppf}$. Étant donnés un homomorphisme $A' \to A$ de $S$-algèbres
et un objet $x$ de $\mathcal{X}$ au-dessus de $\Spec(A)$, nous notons
$\textit{Lift}(x, A')$ la catégorie des relèvements de $x$ à $\Spec(A')$.
Un objet de $\textit{Lift}(x, A')$ est un morphisme $x \to x'$ de $\mathcal{X}$
au-dessus de $\Spec(A) \to \Spec(A')$, et les morphismes de $\textit{Lift}(x, A')$
sont définis par des diagrammes commutatifs. L'ensemble des classes d'isomorphisme de
$\textit{Lift}(x, A')$ est noté $\text{Lift}(x, A')$. Voir
Théorie formelle des déformations, définition \ref{formal-defos-definition-lifts} et
remarque \ref{formal-defos-remark-omit-arrow}.
Si $A' \to A$ est surjectif de noyau localement nilpotent, nous appelons tout élément
$x'$ de $\text{Lift}(x, A')$ une {\it déformation (infinitésimale)} de $x$.
Dans ce cas, le {\it groupe des automorphismes infinitésimaux de $x'$ au-dessus de $x$}
est le noyau
$$
\text{Inf}(x'/x) =
\Ker\left(
\text{Aut}_{\mathcal{X}_{\Spec(A')}}(x') \to
\text{Aut}_{\mathcal{X}_{\Spec(A)}}(x)\right)
$$
Remarquons qu'un élément de $\text{Inf}(x'/x)$ revient à un relèvement
de $\text{id}_x$ au-dessus de $\Spec(A')$ pour (la catégorie fibrée en ensembles associée
à) $\mathit{Aut}_\mathcal{X}(x')$. Comparer à
Théorie formelle des déformations, définition
\ref{formal-defos-definition-relative-infinitesimal-auts} et
Théorie formelle des déformations, remarque
\ref{formal-defos-remark-infaut-lifting-equalities}.

\medskip\noindent
Si $M$ est un $A$-module, nous notons $A[M]$ la $A$-algèbre dont le
$A$-module sous-jacent est $A \oplus M$ et dont la multiplication est donnée par
$(a, m) \cdot (a', m') = (aa', am' + a'm)$. Lorsque $M = A$, c'est l'anneau
des nombres duaux sur $A$, que nous notons, comme il est d'usage, $A[\epsilon]$.
Il existe un homomorphisme de $A$-algèbres $A[M] \to A$. L'image inverse de $x$ sur $\Spec(A[M])$
est appelée la {\it déformation triviale} de $x$ sur $\Spec(A[M])$.

\begin{lemma}
\label{artin-lemma-functoriality}
Soit $S$ un schéma. Soit $f : \mathcal{X} \to \mathcal{Y}$ un $1$-morphisme
de catégories fibrées en groupoïdes au-dessus de $(\Sch/S)_{fppf}$. Soit
$$
\xymatrix{
B' \ar[r] & B \\
A' \ar[u] \ar[r] & A \ar[u]
}
$$
un diagramme commutatif de $S$-algèbres. Soient $x$ un objet de $\mathcal{X}$
au-dessus de $\Spec(A)$, $y$ un objet de $\mathcal{Y}$ au-dessus de $\Spec(B)$,
et $\phi : f(x)|_{\Spec(B)} \to y$ un morphisme de $\mathcal{Y}$
au-dessus de $\Spec(B)$. Il existe alors un foncteur canonique
$$
\textit{Lift}(x, A') \longrightarrow \textit{Lift}(y, B')
$$
entre catégories de relèvements, induit par $f$ et $\phi$. La construction est
compatible, de manière évidente, aux compositions de $1$-morphismes de catégories fibrées en
groupoïdes.
\end{lemma}

\begin{proof}
Le lemme résulte immédiatement des définitions.
\end{proof}

\noindent
Soit $S$ un schéma de base. Soit $\mathcal{X}$ une catégorie fibrée
en groupoïdes au-dessus de $(\Sch/S)_{fppf}$. Nous définissons une catégorie dont les objets sont
les couples $(x, A' \to A)$ tels que
\begin{enumerate}
\item $A' \to A$ soit une surjection de $S$-algèbres dont le noyau
est un idéal de carré nul ;
\item $x$ soit un objet de $\mathcal{X}$ au-dessus de $\Spec(A)$.
\end{enumerate}
Un morphisme $(y, B' \to B) \to (x, A' \to A)$ est donné par un diagramme
commutatif
$$
\xymatrix{
B' \ar[r] & B \\
A' \ar[u] \ar[r] & A \ar[u]
}
$$
de $S$-algèbres, ainsi que par un morphisme $x|_{\Spec(B)} \to y$ au-dessus de
$\Spec(B)$. Appelons-la catégorie des {\it situations de déformation}.

\begin{lemma}
\label{artin-lemma-properties-lift-RS-star}
Soit $S$ un schéma. Soit $\mathcal{X}$ une catégorie
fibrée en groupoïdes au-dessus de $(\Sch/S)_{fppf}$. Supposons que $\mathcal{X}$ satisfasse
la condition (RS*). Soient $A$ une $S$-algèbre et $x$ un objet de
$\mathcal{X}$ au-dessus de $\Spec(A)$.
\begin{enumerate}
\item Il existe un foncteur $A$-linéaire
$\text{Inf}_x : \text{Mod}_A \to \text{Mod}_A$
tel que, pour toute situation de déformation $(x, A' \to A)$ et tout relèvement $x'$,
il existe un isomorphisme $\text{Inf}_x(I) \to \text{Inf}(x'/x)$, où
$I = \Ker(A' \to A)$.
\item Il existe un foncteur $A$-linéaire
$T_x : \text{Mod}_A \to \text{Mod}_A$
tel que
\begin{enumerate}
\item pour tout $M$ dans $\text{Mod}_A$, il existe une bijection
$T_x(M) \to \text{Lift}(x, A[M])$ ;
\item pour toute situation de déformation $(x, A' \to A)$, il existe une action
$$
T_x(I) \times \text{Lift}(x, A') \to \text{Lift}(x, A')
$$
où $I = \Ker(A' \to A)$. Elle est simplement transitive si
$\text{Lift}(x, A') \not = \emptyset$.
\end{enumerate}
\end{enumerate}
\end{lemma}

\begin{proof}
Nous définissons $\text{Inf}_x$ comme le foncteur
$$
\text{Mod}_A \longrightarrow \textit{Ensembles},\quad
M \longrightarrow
\text{Inf}(x'_M/x) = \text{Lift}(\text{id}_x, A[M])
$$
qui associe à $M$ le groupe des automorphismes infinitésimaux
de la déformation triviale $x'_M$ de $x$ sur $\Spec(A[M])$,
ou, ce qui revient au même, le groupe des relèvements de $\text{id}_x$ dans
$\mathit{Aut}_\mathcal{X}(x'_M)$.
Nous définissons $T_x$ comme le foncteur
$$
\text{Mod}_A \longrightarrow \textit{Ensembles},\quad
M \longrightarrow \text{Lift}(x, A[M])
$$
des classes d'isomorphisme de déformations infinitésimales de $x$ sur
$\Spec(A[M])$. Nous appliquons Théorie formelle des déformations, lemme
\ref{formal-defos-lemma-linear-functor}
à $\text{Inf}_x$ et $T_x$. Ce lemme s'applique, puisque
(RS*) nous dit que
$$
\textit{Lift}(x, A[M \times N]) =
\textit{Lift}(x, A[M]) \times \textit{Lift}(x, A[N])
$$
comme catégories (et que les déformations triviales se correspondent également).

\medskip\noindent
Soit $(x, A' \to A)$ une situation de déformation. Considérons l'homomorphisme d'anneaux
$g : A' \times_A A' \to A[I]$ défini par la
règle $g(a_1, a_2) = \overline{a_1} \oplus a_2 - a_1$.
Il existe un isomorphisme
$$
A' \times_A A' \longrightarrow A' \times_A A[I]
$$
donné par $(a_1, a_2) \mapsto (a_1, g(a_1, a_2))$.
Cet isomorphisme commute aux projections sur $A'$ par le premier
facteur, et donc aux projections sur $A$. En appliquant deux fois (RS*),
nous obtenons ainsi des équivalences de catégories
\begin{align*}
\textit{Lift}(x, A') \times \textit{Lift}(x, A')
& =
\textit{Lift}(x, A' \times_A A') \\
& =
\textit{Lift}(x, A' \times_A A[I]) \\
& =
\textit{Lift}(x, A') \times \textit{Lift}(x, A[I])
\end{align*}
En utilisant ces applications et la projection sur le dernier facteur du dernier produit,
nous obtenons des ``applications différence''
$$
\text{Inf}(x'/x) \times  \text{Inf}(x'/x)
\longrightarrow
\text{Inf}_x(I)
\quad\text{et}\quad
\text{Lift}(x, A') \times \text{Lift}(x, A')
\longrightarrow
T_x(I)
$$
Ces applications différence satisfont la règle de transitivité
``$(x'_1 - x'_2) + (x'_2 - x'_3) = x'_1 - x'_3$'', car le diagramme
$$
\xymatrix{
A' \times_A A' \times_A A'
\ar[rrrrr]_-{(a_1, a_2, a_3) \mapsto (g(a_1, a_2), g(a_2, a_3))}
\ar[rrrrrd]_{(a_1, a_2, a_3) \mapsto g(a_1, a_3)} & & & & &
A[I] \times_A A[I] = A[I \times I] \ar[d]^{+} \\
& & & & & A[I]
}
$$
est commutatif. En inversant la chaîne d'équivalences ci-dessus, nous obtenons
une action libre et transitive dès que $\text{Inf}(x'/x)$,
resp.\ $\text{Lift}(x, A')$, est non vide. Remarquons que $\text{Inf}(x'/x)$
est toujours non vide, puisque c'est un groupe.
\end{proof}

\begin{remark}[Fonctorialité]
\label{artin-remark-functoriality}
Hypothèses et notations comme dans le lemme \ref{artin-lemma-properties-lift-RS-star}.
Supposons donné un homomorphisme d'anneaux $A \to B$ et posons $y = x|_{\Spec(B)}$.
Soient $M \in \text{Mod}_A$, $N \in \text{Mod}_B$
et $M \to N$ une application $A$-linéaire. Il existe alors des applications canoniques
$\text{Inf}_x(M) \to \text{Inf}_y(N)$ et
$T_x(M) \to T_y(N)$, simplement parce qu'il existe un foncteur d'image inverse
$$
\textit{Lift}(x, A[M]) \to \textit{Lift}(y, B[N])
$$
provenant de l'homomorphisme d'anneaux $A[M] \to B[N]$. De même, tout morphisme de
situations de déformation $(y, B' \to B) \to (x, A' \to A)$ induit un foncteur
d'image inverse $\textit{Lift}(x, A') \to \textit{Lift}(y, B')$. Comme la
construction de l'action, de l'addition et de la multiplication scalaire
sur $\text{Inf}_x$ et $T_x$ n'emploie que des morphismes dans les catégories de relèvements
(voir la démonstration de
Théorie formelle des déformations, lemme
\ref{formal-defos-lemma-linear-functor})
nous voyons que les constructions précédentes sont fonctorielles. Autrement dit, nous
obtenons des applications $A$-linéaires
$$
\text{Inf}_x(M) \to \text{Inf}_y(N)
\quad\text{et}\quad
T_x(M) \to T_y(N)
$$
telles que les diagrammes
$$
\vcenter{
\xymatrix{
\text{Inf}_y(J) \ar[r] & \text{Inf}(y'/y) \\
\text{Inf}_x(I) \ar[r] \ar[u] & \text{Inf}(x'/x) \ar[u]
}
}
\quad\text{et}\quad
\vcenter{
\xymatrix{
T_y(J) \times \text{Lift}(y, B') \ar[r] & \text{Lift}(y, B') \\
T_x(I) \times \text{Lift}(x, A') \ar[r] \ar[u] & \text{Lift}(x, A') \ar[u]
}
}
$$
sont commutatifs. Ici, $I = \Ker(A' \to A)$, $J = \Ker(B' \to B)$,
$x'$ est un relèvement de $x$ à $A'$ (qui n'existe pas nécessairement) et
$y' = x'|_{\Spec(B')}$.
\end{remark}

\begin{remark}[Automorphismes]
\label{artin-remark-automorphisms}
Hypothèses et notations comme dans le lemme \ref{artin-lemma-properties-lift-RS-star}.
Soient $x', x''$ des relèvements de $x$ à $A'$. Nous disposons alors d'une application
de composition
$$
\text{Inf}(x'/x) \times
\Mor_{\textit{Lift}(x, A')}(x', x'') \times \text{Inf}(x''/x)
\longrightarrow
\Mor_{\textit{Lift}(x, A')}(x', x'').
$$
Comme $\textit{Lift}(x, A')$ est un groupoïde, si
$\Mor_{\textit{Lift}(x, A')}(x', x'')$ est non vide, cela définit
une action à gauche simplement transitive de $\text{Inf}(x'/x)$ sur
$\Mor_{\textit{Lift}(x, A')}(x', x'')$ et une action à droite
simplement transitive de $\text{Inf}(x''/x)$. Or le lemme donne
$\text{Inf}(x'/x) = \text{Inf}_x(I) = \text{Inf}(x''/x)$.
Nous affirmons que, par ces identifications, les deux actions décrites ci-dessus coïncident.
En effet, ou bien $x' \not \cong x''$, auquel cas l'affirmation est claire, ou bien
$x' \cong x''$ ; dans ce cas, nous pouvons supposer $x'' = x'$, et le
résultat découle alors de la commutativité de $\text{Inf}(x'/x)$.
En particulier, nous obtenons une action bien définie
$$
\text{Inf}_x(I) \times \Mor_{\textit{Lift}(x, A')}(x', x'')
\longrightarrow
\Mor_{\textit{Lift}(x, A')}(x', x'')
$$
qui est simplement transitive dès que $\Mor_{\textit{Lift}(x, A')}(x', x'')$
est non vide.
\end{remark}

\begin{remark}
\label{artin-remark-short-exact-sequence-thickenings}
Soit $S$ un schéma. Soit $\mathcal{X}$ une catégorie
fibrée en groupoïdes au-dessus de $(\Sch/S)_{fppf}$. Soit $A$ une $S$-algèbre. Il
existe une notion de {\it suite exacte courte}
$$
(x, A_1' \to A) \to (x, A_2' \to A) \to (x, A_3' \to A)
$$
de situations de déformation : nous demandons que les applications correspondantes entre
les noyaux $I_i = \Ker(A_i' \to A)$ donnent une suite exacte courte
$$
0 \to I_3 \to I_2 \to I_1 \to 0
$$
de $A$-modules. Remarquons que, dans ce cas, l'application $A_3' \to A_1'$
se factorise par $A$ ; il existe donc un isomorphisme canonique
$A_1' = A[I_1]$.
\end{remark}

\begin{lemma}
\label{artin-lemma-ses-inf-and-T}
Soit $S$ un schéma. Soient $p : \mathcal{X} \to \mathcal{Y}$
et $q : \mathcal{Z} \to \mathcal{Y}$ des $1$-morphismes de catégories
fibrées en groupoïdes au-dessus de $(\Sch/S)_{fppf}$. Supposons que $\mathcal{X}$,
$\mathcal{Y}$ et $\mathcal{Z}$ satisfassent (RS*).
Soient $A$ une $S$-algèbre et $w$ un objet de
$\mathcal{W} = \mathcal{X} \times_\mathcal{Y} \mathcal{Z}$ au-dessus de $A$.
Notons $x, y, z$ les objets de $\mathcal{X}, \mathcal{Y}, \mathcal{Z}$
issus de $w$. Pour tout $A$-module $M$, il existe une suite exacte à $6$ termes
$$
\xymatrix{
0 \ar[r] &
\text{Inf}_w(M) \ar[r] &
\text{Inf}_x(M) \oplus \text{Inf}_z(M) \ar[r] &
\text{Inf}_y(M) \ar[lld] \\
 &
T_w(M) \ar[r] &
T_x(M) \oplus T_z(M) \ar[r] &
T_y(M)
}
$$
de $A$-modules.
\end{lemma}

\begin{proof}
D'après le lemme \ref{artin-lemma-fibre-product-RS-star}, $\mathcal{W}$
satisfait (RS*) ; par conséquent, $T_w(M)$ et $\text{Inf}_w(M)$ sont définis.
Les flèches horizontales sont définies par la fonctorialité du
lemme \ref{artin-lemma-functoriality}.

\medskip\noindent
Définition du ``morphisme de connexion'' $\delta : \text{Inf}_y(M) \to T_w(M)$.
Choisissons des isomorphismes $p(x) \to y$ et $y \to q(z)$ tels que
$w = (x, z, p(x) \to y \to q(z))$ dans la description du
$2$-produit fibré de
Catégories, lemme \ref{categories-lemma-2-product-fibred-categories},
et plus précisément
Catégories, lemme \ref{categories-lemma-2-product-categories-over-C}.
Notons $x', y', z', w'$ les déformations triviales de
$x, y, z, w$ sur $A[M]$. Par image inverse, nous obtenons des isomorphismes
$y' \to p(x')$ et $q(z') \to y'$. Un élément $\alpha \in \text{Inf}_y(M)$
revient à un automorphisme $\alpha : y' \to y'$
au-dessus de $A[M]$ dont la restriction à $y$ au-dessus de $A$ est l'identité.
En posant
$$
\delta(\alpha) =
(x', z', p(x') \to y' \xrightarrow{\alpha} y' \to q(z'))
$$
nous obtenons un objet de $T_w(M)$. C'est une application de $A$-modules
d'après Théorie formelle des déformations, lemme
\ref{formal-defos-lemma-morphism-linear-functors}.

\medskip\noindent
La suite de la démonstration est exactement la même que celle de
Théorie formelle des déformations, lemme
\ref{formal-defos-lemma-deformation-categories-fiber-product-morphisms}.
\end{proof}

\begin{remark}[Compatibilité avec les espaces tangents précédents]
\label{artin-remark-compare-deformation-spaces}
Soit $S$ un schéma localement noethérien. Soit $\mathcal{X}$ une catégorie fibrée
en groupoïdes au-dessus de $(\Sch/S)_{fppf}$. Supposons que $\mathcal{X}$ satisfasse (RS*).
Soient $k$ un corps de type fini sur $S$ et $x_0$ un objet de
$\mathcal{X}$ au-dessus de $\Spec(k)$. Nous avons alors les égalités d'espaces
vectoriels sur $k$
$$
T\mathcal{F}_{\mathcal{X}, k, x_0} = T_{x_0}(k)
\quad\text{et}\quad
\text{Inf}(\mathcal{F}_{\mathcal{X}, k, x_0}) =
\text{Inf}_{x_0}(k)
$$
où les espaces figurant aux membres de gauche des égalités sont
donnés en (\ref{artin-equation-tangent-space}) et
(\ref{artin-equation-infinitesimal-automorphisms})
où ceux des membres de droite sont donnés par le
lemme \ref{artin-lemma-properties-lift-RS-star}.
\end{remark}

\begin{remark}[Élément canonique]
\label{artin-remark-canonical-element}
Hypothèses et notations comme dans le lemme \ref{artin-lemma-properties-lift-RS-star}.
Choisissons un ouvert affine $\Spec(\Lambda) \subset S$ tel que $\Spec(A) \to S$
corresponde à un homomorphisme d'anneaux $\Lambda \to A$. Considérons l'homomorphisme d'anneaux
$$
A \longrightarrow A[\Omega_{A/\Lambda}],
\quad
a \longmapsto (a, \text{d}_{A/\Lambda}(a))
$$
En prenant l'image inverse de $x$ par le morphisme correspondant
$\Spec(A[\Omega_{A/\Lambda}]) \to \Spec(A)$, nous obtenons une
déformation $x_{can}$ de $x$ sur $A[\Omega_{A/\Lambda}]$. Nous l'appelons
l'{\it élément canonique}
$$
x_{can} \in T_x(\Omega_{A/\Lambda}) = \text{Lift}(x, A[\Omega_{A/\Lambda}]).
$$
Supposons ensuite que $\Lambda$ soit noethérien et que $\Lambda \to A$
soit de type fini. Soit
$k = \kappa(\mathfrak p)$ le corps résiduel en un point de type fini $u_0$
de $U = \Spec(A)$. Posons $x_0 = x|_{u_0}$. D'après (RS*) et l'égalité
$A[k] = A \times_k k[k]$, l'espace $T_x(k)$ est l'espace tangent au
foncteur de déformations $\mathcal{F}_{\mathcal{X}, k, x_0}$. Par l'identification
$$
T\mathcal{F}_{U, k, u_0} =
\text{Der}_\Lambda(A, k) = \Hom_A(\Omega_{A/\Lambda}, k)
$$
(voir Théorie formelle des déformations, exemple
\ref{formal-defos-example-tangent-space-prorepresentable-functor})
et par la fonctorialité de $T_x$, l'élément canonique fournit l'application
sur les espaces tangents induite par l'objet $x$ au-dessus de $U$. Plus précisément,
$\theta \in T\mathcal{F}_{U, k, u_0}$ s'envoie sur $T_x(\theta)(x_{can})$
dans $T_x(k) = T\mathcal{F}_{\mathcal{X}, k, x_0}$.
\end{remark}

\begin{remark}[Automorphisme canonique]
\label{artin-remark-canonical-isomorphism}
Soit $S$ un schéma localement noethérien. Soit $\mathcal{X}$ une catégorie
fibrée en groupoïdes au-dessus de $(\Sch/S)_{fppf}$. Supposons que $\mathcal{X}$ satisfasse
la condition (RS*). Soit $A$ une $S$-algèbre telle que
$\Spec(A) \to S$ se factorise par un ouvert affine, et soient $x, y$ des objets de
$\mathcal{X}$ au-dessus de $\Spec(A)$. Soient en outre $A \to B$ un homomorphisme d'anneaux et
$\alpha : x|_{\Spec(B)} \to y|_{\Spec(B)}$ un morphisme de
$\mathcal{X}$ au-dessus de $\Spec(B)$. Considérons l'homomorphisme d'anneaux
$$
B \longrightarrow B[\Omega_{B/A}],
\quad
b \longmapsto (b, \text{d}_{B/A}(b))
$$
En prenant l'image inverse de $\alpha$ par le morphisme correspondant
$\Spec(B[\Omega_{B/A}]) \to \Spec(B)$, nous obtenons un
morphisme $\alpha_{can}$ entre les images inverses de $x$ et $y$ au-dessus de
$B[\Omega_{B/A}]$. D'autre part, nous pouvons prendre l'image inverse de $\alpha$
par le morphisme $\Spec(B[\Omega_{B/A}]) \to \Spec(B)$ correspondant
à l'injection de $B$ dans le premier facteur de $B[\Omega_{B/A}]$.
D'après la remarque \ref{artin-remark-automorphisms},
nous pouvons prendre la différence
$$
\varphi(x, y, \alpha) = \alpha_{can} - \alpha|_{\Spec(B[\Omega_{B/A}])} \in
\text{Inf}_{x|_{\Spec(B)}}(\Omega_{B/A}).
$$
Nous l'appellerons l'{\it automorphisme canonique}. Il dépend
de toutes les données $A$, $x$, $y$, $A \to B$ et $\alpha$.
\end{remark}





\section{Théories de l'obstruction}
\label{artin-section-obstruction-theory}

\noindent
Dans cette section, nous décrivons ce qu'est une théorie de l'obstruction.
Contrairement aux espaces de déformations infinitésimales et d'automorphismes
infinitésimaux, une théorie de l'obstruction est une donnée supplémentaire.
La formulation est motivée par les résultats du
Lemme \ref{artin-lemma-properties-lift-RS-star}
et de la Remarque \ref{artin-remark-functoriality}.

\begin{definition}
\label{artin-definition-obstruction-theory}
Soit $S$ une base localement noethérienne. Soit $\mathcal{X}$ une catégorie fibrée
en groupoïdes sur $(\Sch/S)_{fppf}$. Une {\it théorie de l'obstruction} est
constituée des données suivantes :
\begin{enumerate}
\item pour toute $S$-algèbre $A$ telle que $\Spec(A) \to S$
se factorise par un ouvert affine, et pour tout objet $x$ de $\mathcal{X}$ sur
$\Spec(A)$, un foncteur $A$-linéaire
$$
\mathcal{O}_x : \text{Mod}_A \to \text{Mod}_A
$$
de {\it modules d'obstruction},
\item pour $(x, A)$ comme en (1), un homomorphisme d'anneaux $A \to B$,
$M \in \text{Mod}_A$, $N \in \text{Mod}_B$, et un morphisme $A$-linéaire
$M \to N$, un morphisme $A$-linéaire induit $\mathcal{O}_x(M) \to \mathcal{O}_y(N)$,
où $y = x|_{\Spec(B)}$, et
\item pour toute situation de déformation $(x, A' \to A)$, un
élément d'{\it obstruction}
$o_x(A') \in \mathcal{O}_x(I)$, où $I = \Ker(A' \to A)$.
\end{enumerate}
Ces données sont soumises aux conditions suivantes :
\begin{enumerate}
\item[(i)] les morphismes de fonctorialité font des modules d'obstruction un foncteur
de la catégorie des triplets $(x, A, M)$ à valeurs dans les ensembles,
\item[(ii)] pour tout morphisme de situations de déformation
$(y, B' \to B) \to (x, A' \to A)$, l'élément $o_x(A')$ s'envoie
sur $o_y(B')$, et
\item[(iii)] on a
$$
\text{Lift}(x, A') \not = \emptyset
\Leftrightarrow
o_x(A') = 0
$$
pour toute situation de déformation $(x, A' \to A)$.
\end{enumerate}
\end{definition}

\noindent
Cette dernière condition explique la terminologie. Le module $\mathcal{O}_x(A')$
est appelé le {\it module d'obstruction}. L'élément $o_x(A')$ est
l'{\it obstruction}.
La plupart des théories de l'obstruction possèdent des propriétés supplémentaires et,
pour qu'elles soient utiles, il faut imposer des conditions supplémentaires.
De plus, il ne s'agit ici que d'une définition possible ; on pourrait par exemple
n'y considérer que les situations de déformation de type fini sur $S$.

\medskip\noindent
L'une des principales raisons d'introduire les théories de l'obstruction est de vérifier
l'ouverture de la versalité. Un exemple de résultat de cette nature est le
Lemme \ref{artin-lemma-get-openness-obstruction-theory} ci-dessous.
L'idée initiale de cette méthode est due à Artin ; voir
les articles d'Artin mentionnés dans l'introduction. Elle a été reprise,
par exemple, dans les travaux de Flenner \cite{Flenner},
Hall \cite{Hall-coherent},
Hall et Rydh \cite{rydh_axioms},
Olsson \cite{olsson_deformation},
Olsson et Starr \cite{olsson-starr}, et
Lieblich \cite{lieblich-complexes} (références citées sans ordre particulier).
En outre, pour des catégories fibrées en groupoïdes particulières, les auteurs
développent souvent un peu de théorie adaptée au problème considéré.
Nous développerons cette théorie plus loin (insérer ici une référence future).

\begin{lemma}
\label{artin-lemma-get-openness-obstruction-theory}
\begin{reference}
Il s'agit de \cite[théorème 4.4]{Hall-coherent}.
\end{reference}
Soit $S$ un schéma localement noethérien. Soit $\mathcal{X}$ une catégorie fibrée
en groupoïdes sur $(\Sch/S)_{fppf}$. Supposons que
\begin{enumerate}
\item $\Delta : \mathcal{X} \to \mathcal{X} \times \mathcal{X}$ soit
représentable par des espaces algébriques,
\item $\mathcal{X}$ vérifie (RS*),
\item $\mathcal{X}$ commute aux limites,
\item il existe une théorie de l'obstruction\footnote{L'analyse de la démonstration
montre qu'il suffit en fait de vérifier
la fonctorialité (ii) des classes d'obstruction dans la
Définition \ref{artin-definition-obstruction-theory}
pour les morphismes $(y, B' \to B) \to (x, A' \to A)$
avec $B = A$ et $y = x$.},
\item pour un objet $x$ de $\mathcal{X}$ sur $\Spec(A)$
et des $A$-modules $M_n$, $n \geq 1$, on ait
\begin{enumerate}
\item $T_x(\prod M_n) = \prod T_x(M_n)$,
\item $\mathcal{O}_x(\prod M_n) \to \prod \mathcal{O}_x(M_n)$
est injectif.
\end{enumerate}
\end{enumerate}
Alors $\mathcal{X}$ vérifie l'ouverture de la versalité.
\end{lemma}

\begin{proof}
Nous le démontrons en vérifiant la condition (4) du
Lemme \ref{artin-lemma-SGE-implies-openness-versality}.
Soient $(\xi_n)$ et $(R_n)$ comme dans la Remarque \ref{artin-remark-strong-effectiveness},
de sorte que $\Ker(R_m \to R_n)$ soit un idéal de carré nul
pour tous $m \geq n$. Posons $A = R_1$ et $x = \xi_1$.
Notons $M_n = \Ker(R_n \to R_1)$.
Alors $M_n$ est un $A$-module. Posons $R = \lim R_n$.
Soit
$$
\tilde R = \{(r_1, r_2, r_3 \ldots) \in \prod R_n
\text{ tels que tous aient la même image dans }A\}
$$
Alors $\tilde R \to A$ est surjectif de noyau $M = \prod M_n$.
Il existe un morphisme $R \to \tilde R$ et un morphisme
$\tilde R \to A[M]$, $(r_1, r_2, r_3, \ldots) \mapsto
(r_1, r_2 - r_1, r_3 - r_2, \ldots)$.
Ils donnent ensemble une suite exacte courte
$$
(x, R \to A) \to (x, \tilde R \to A) \to (x, A[M])
$$
de situations de déformation ; voir la
Remarque \ref{artin-remark-short-exact-sequence-thickenings}.
La suite associée des noyaux
$0 \to \lim M_n \to M \to M \to 0$
est la suite canonique qui calcule la limite
du système de modules $(M_n)$.

\medskip\noindent
Soit $o_x(\tilde R) \in \mathcal{O}_x(M)$ l'élément d'obstruction.
Puisque nous disposons des relèvements $\xi_n$, nous voyons que $o_x(\tilde R)$
s'envoie sur zéro dans $\mathcal{O}_x(M_n)$. D'après l'hypothèse (5)(b),
on a $o_x(\tilde R) = 0$. Choisissons un relèvement $\tilde \xi$
de $x$ à $\Spec(\tilde R)$. Soit $\tilde \xi_n$ la
restriction de $\tilde \xi$ à $\Spec(R_n)$. Il existe
des éléments $t_n \in T_x(M_n)$ tels que
$t_n \cdot \tilde \xi_n = \xi_n$, d'après la partie (2)(b) du
Lemme \ref{artin-lemma-properties-lift-RS-star}.
D'après l'hypothèse (5)(a), nous pouvons trouver $t \in T_x(M)$
qui s'envoie sur $t_n$ dans $T_x(M_n)$. Après avoir remplacé
$\tilde \xi$ par $t \cdot \tilde \xi$, nous obtenons que
$\tilde \xi$ se restreint à $\xi_n$ sur $\Spec(R_n)$ pour tout $n$.
En particulier, puisque $\xi_{n + 1}$ se restreint à $\xi_n$
sur $\Spec(R_n)$, la restriction $\overline{\xi}$ de $\tilde \xi$
à $\Spec(A[M])$ est telle que sa restriction à
$\Spec(A[M_n])$ est la déformation triviale pour tout $n$.
Par conséquent, l'hypothèse (5)(a) montre que $\overline{\xi}$
est la déformation triviale de $x$. L'axiome (RS*)
appliqué à $R = \tilde R \times_{A[M]} A$
implique que $\tilde \xi$ est l'image réciproque
d'une déformation $\xi$ de $x$ sur $R$. Ceci achève la démonstration.
\end{proof}

\begin{example}
\label{artin-example-global-sections}
Soit $S = \Spec(\Lambda)$ pour un anneau noethérien $\Lambda$.
Soit $W \to S$ un morphisme de schémas. Soit $\mathcal{F}$
un $\mathcal{O}_W$-module quasi-cohérent et plat sur $S$.
Considérons le foncteur
$$
F : (\Sch/S)_{fppf}^{opp} \longrightarrow \textit{Ensembles},
\quad
T/S \longrightarrow H^0(W_T, \mathcal{F}_T)
$$
où $W_T = T \times_S W$ est le changement de base et $\mathcal{F}_T$
l'image réciproque de $\mathcal{F}$ sur $W_T$. Si $T = \Spec(A)$,
nous écrirons $W_T = W_A$, etc. Soit $\mathcal{X} \to (\Sch/S)_{fppf}$
la catégorie fibrée en groupoïdes associée à $F$. Alors
$\mathcal{X}$ possède une théorie de l'obstruction. Plus précisément,
\begin{enumerate}
\item étant donnés $A$, une algèbre sur $\Lambda$, et
$x \in H^0(W_A, \mathcal{F}_A)$, nous posons
$\mathcal{O}_x(M) = H^1(W_A, \mathcal{F}_A \otimes_A M)$,
\item étant donnée une situation de déformation $(x, A' \to A)$, nous définissons
$o_x(A') \in \mathcal{O}_x(A)$ comme l'image de $x$ par l'homomorphisme cobord
$$
H^0(W_A, \mathcal{F}_A) \longrightarrow H^1(W_A, \mathcal{F}_A \otimes_A I)
$$
provenant de la suite exacte courte de modules
$$
0 \to \mathcal{F}_A \otimes_A I \to
\mathcal{F}_{A'} \to \mathcal{F}_A \to 0.
$$
\end{enumerate}
Nous avons omis certains détails, notamment la construction de la suite
exacte courte ci-dessus (elle utilise le fait que $W_A$ et $W_{A'}$ ont le même
espace topologique sous-jacent), ainsi que l'explication du fait que la platitude
de $\mathcal{F}$ sur $S$ implique l'exactitude de la suite ci-dessus.
\end{example}

\begin{example}[Exemple clé]
\label{artin-example-key}
Soit $S = \Spec(\Lambda)$ pour un anneau noethérien $\Lambda$.
Supposons que $\mathcal{X} = (\Sch/X)_{fppf}$, avec $X = \Spec(R)$ et
$R = \Lambda[x_1, \ldots, x_n]/J$. Le complexe cotangent
naïf $\NL_{R/\Lambda}$ est (canoniquement) homotopiquement équivalent à
$$
J/J^2
\longrightarrow
\bigoplus\nolimits_{i = 1, \ldots, n} R\text{d}x_i,
$$
voir Algèbre, Lemme \ref{algebra-lemma-NL-homotopy}.
Considérons une situation de déformation $(x, A' \to A)$. Notons $I$ le noyau de
$A' \to A$. L'objet $x$ correspond à $(a_1, \ldots, a_n)$,
où $a_i \in A$ et $f(a_1, \ldots, a_n) = 0$ dans $A$ pour tout $f \in J$.
Posons
\begin{align*}
\mathcal{O}_x(A')
& =
\Hom_R(J/J^2, I)/\Hom_R(R^{\oplus n}, I) \\
& =
\Ext^1_R(\NL_{R/\Lambda}, I) \\
& =
\Ext^1_A(\NL_{R/\Lambda} \otimes_R A, I).
\end{align*}
Choisissons des relèvements $a_i' \in A'$ des $a_i$ dans $A$. Alors $o_x(A')$
est la classe du morphisme $J/J^2 \to I$ défini en envoyant $f \in J$ sur
$f(a_1', \ldots, a'_n) \in I$. Nous omettons de vérifier que
$o_x(A')$ est indépendant des choix. Il est clair que, si $o_x(A') = 0$,
alors le morphisme se relève. Enfin, la fonctorialité est immédiate.
Nous obtenons ainsi une théorie de l'obstruction. Observons que $o_x(A')$
peut se décrire un peu plus canoniquement comme le composé
$$
\NL_{R/\Lambda} \to \NL_{A/\Lambda} \to \NL_{A/A'} = I[1]
$$
dans $D(A)$ ; voir Algèbre, Lemme \ref{algebra-lemma-NL-surjection}
pour la dernière identification.
\end{example}









\section{Théories naïves de l'obstruction}
\label{artin-section-naive-obstruction-theory}

\noindent
Le titre de cette section renvoie au fait que nous y utiliserons le
complexe cotangent naïf. Soit $(x, A' \to A)$
une situation de déformation pour une catégorie fibrée en groupoïdes donnée sur un
schéma localement noethérien $S$. L'Exemple clé \ref{artin-example-key}
laisse penser que toute théorie de l'obstruction devrait être étroitement liée aux
morphismes dans $D(A)$ aboutissant au complexe cotangent naïf de $A$.
En précisant cette idée, nous obtenons dans le
Lemme \ref{artin-lemma-characterize-versal} un critère de versalité qui conduit à un critère
d'ouverture de la versalité dans le Lemme \ref{artin-lemma-openness}. Nous introduisons la notion de
théorie naïve de l'obstruction dans la
Définition \ref{artin-definition-naive-obstruction-theory} afin de formaliser
un peu plus cette idée.

\medskip\noindent
Dans ce qui suit, nous utiliserons le complexe cotangent naïf tel qu'il est
défini dans Algèbre, section \ref{algebra-section-netherlander}.
En particulier, si $A' \to A$ est une surjection de $\Lambda$-algèbres
de noyau $I$ de carré nul, il existe des morphismes
$$
\NL_{A'/\Lambda} \to \NL_{A/\Lambda} \to \NL_{A/A'}
$$
dont le composé est homotopiquement équivalent à zéro (voir
Algèbre, Remarque \ref{algebra-remark-composition-homotopy-equivalent-to-zero}).
Ils ne forment pas en général un triangle distingué, car nous utilisons
le complexe cotangent naïf et non le complexe cotangent proprement dit.
Il existe une équivalence d'homotopie $\NL_{A/A'} \to I[1]$ (le complexe
constitué de $I$ placé en degré $-1$ ; voir
Algèbre, Lemme \ref{algebra-lemma-NL-surjection}).
Enfin, remarquons qu'il existe un morphisme canonique
$\NL_{A/\Lambda} \to \Omega_{A/\Lambda}$.

\begin{lemma}
\label{artin-lemma-compute-ext-into-field}
Soit $A \to k$ un homomorphisme d'anneaux, où $k$ est un corps. Soit $E \in D^-(A)$.
Alors $\Ext^i_A(E, k) = \Hom_k(H^{-i}(E \otimes^\mathbf{L} k), k)$.
\end{lemma}

\begin{proof}
Omis. Indication : remplacer $E$ par un complexe de $A$-modules libres borné
supérieurement et calculer les deux membres.
\end{proof}

\begin{lemma}
\label{artin-lemma-construct-essential-surjection}
Soient $\Lambda \to A \to k$ des homomorphismes de type fini d'anneaux noethériens,
avec $k = \kappa(\mathfrak p)$ pour un idéal premier $\mathfrak p$ de $A$. Soit
$\xi : E \to \NL_{A/\Lambda}$ un morphisme de $D^{-}(A)$ tel que
$H^{-1}(\xi \otimes^{\mathbf{L}} k)$ ne soit pas surjectif.
Alors il existe une surjection $A' \to A$ de $\Lambda$-algèbres
telle que
\begin{enumerate}
\item[(a)] $I = \Ker(A' \to A)$ soit de carré nul et isomorphe à $k$
comme $A$-module,
\item[(b)] $\Omega_{A'/\Lambda} \otimes k = \Omega_{A/\Lambda} \otimes k$, et
\item[(c)] $E \to \NL_{A/A'}$ soit nul.
\end{enumerate}
\end{lemma}

\begin{proof}
Soit $f \in A$, $f \not \in \mathfrak p$. Supposons que $A'' \to A_f$
vérifie (a), (b), (c) pour le morphisme induit
$E \otimes_A A_f \to \NL_{A_f/\Lambda}$ ; voir
Algèbre, Lemme \ref{algebra-lemma-localize-NL}.
On peut alors poser $A' = A'' \times_{A_f} A$ et obtenir une solution.
Il est en effet clair que $A' \to A$ vérifie (a), puisque
$\Ker(A' \to A) = \Ker(A'' \to A) = I$. Choisissons
$f'' \in A''$ relevant $f$. Alors le localisé de $A'$ en
$(f'', f)$ est isomorphe à $A''$
(par exemple, d'après
Compléments sur l'algèbre, Lemme \ref{more-algebra-lemma-diagram-localize}).
Ainsi, (b) et (c) sont également clairs pour $A'$.
Nous voyons de cette manière que nous pouvons remplacer $A$ par le localisé
$A_f$ (un nombre fini de fois).
En particulier, après un tel remplacement, nous pouvons supposer que $\mathfrak p$
est un idéal maximal de $A$ ; voir
Morphismes, Lemme \ref{morphisms-lemma-point-finite-type}.

\medskip\noindent
Choisissons une présentation $A = \Lambda[x_1, \ldots, x_n]/J$. Alors
$\NL_{A/\Lambda}$ est (canoniquement) homotopiquement équivalent à
$$
J/J^2
\longrightarrow
\bigoplus\nolimits_{i = 1, \ldots, n} A\text{d}x_i,
$$
voir Algèbre, Lemme \ref{algebra-lemma-NL-homotopy}. Quitte à localiser
si nécessaire (en utilisant le lemme de Nakayama), nous pouvons choisir des générateurs
$f_1, \ldots, f_m$ de $J$ tels que les $f_j \otimes 1$ forment une base de
$J/J^2 \otimes_A k$. De plus, après renumérotation, nous pouvons supposer que les
images de $\text{d}f_1, \ldots, \text{d}f_r$ forment une
base de l'image de $J/J^2 \otimes k \to \bigoplus k\text{d}x_i$
et que $\text{d}f_{r + 1}, \ldots, \text{d}f_m$ s'envoient sur zéro dans
$\bigoplus k\text{d}x_i$. Avec ces choix, l'espace
$$
H^{-1}(\NL_{A/\Lambda} \otimes^{\mathbf{L}}_A k) =
H^{-1}(\NL_{A/\Lambda} \otimes_A k)
$$
a pour base $f_{r + 1} \otimes 1, \ldots, f_m \otimes 1$. En changeant
encore une fois de base, nous pouvons supposer que l'image de $H^{-1}(\xi \otimes^{\mathbf{L}} k)$
est contenue dans le sous-espace vectoriel sur $k$ engendré par
$f_{r + 1} \otimes 1, \ldots, f_{m - 1} \otimes 1$.
Posons
$$
A' = \Lambda[x_1, \ldots, x_n]/(f_1, \ldots, f_{m - 1}, \mathfrak pf_m)
$$
Par construction, $A' \to A$ vérifie (a). Puisque $\text{d}f_m$ s'envoie
sur zéro dans $\bigoplus k\text{d}x_i$, nous voyons que (b) est satisfaite. Enfin, par
construction, le morphisme induit $E \to \NL_{A/A'} = I[1]$ induit le morphisme nul
$H^{-1}(E \otimes_A^\mathbf{L} k) \to I \otimes_A k$. D'après le
Lemme \ref{artin-lemma-compute-ext-into-field},
nous voyons que le composé est nul.
\end{proof}

\noindent
Le lemme suivant est notre principal résultat technique.

\begin{lemma}
\label{artin-lemma-characterize-versal}
Soit $S$ un schéma localement noethérien. Soit $\mathcal{X}$ une catégorie
fibrée en groupoïdes sur $(\Sch/S)_{fppf}$ vérifiant (RS*).
Soit $U = \Spec(A)$ un
schéma affine de type fini sur $S$ dont le morphisme se factorise par un ouvert affine
$\Spec(\Lambda)$. Soit $x$ un objet de $\mathcal{X}$ sur $U$.
Soit $\xi : E \to \NL_{A/\Lambda}$ un morphisme de $D^{-}(A)$. Supposons que
\begin{enumerate}
\item[(i)] pour toute situation de déformation $(x, A' \to A)$, on ait :
$x$ se relève à $\Spec(A')$ si et seulement si
$E \to \NL_{A/\Lambda} \to \NL_{A/A'}$ est nul, et
\item[(ii)] il existe un isomorphisme de foncteurs
$T_x(-) \to \Ext^0_A(E, -)$
tel que $E \to \NL_{A/\Lambda} \to \Omega^1_{A/\Lambda}$
corresponde à l'élément canonique (voir la
Remarque \ref{artin-remark-canonical-element}).
\end{enumerate}
Soit $u_0 \in U$ un point de type fini de corps résiduel
$k = \kappa(u_0)$. Considérons les assertions suivantes :
\begin{enumerate}
\item $x$ est versel en $u_0$, et
\item $\xi : E \to \NL_{A/\Lambda}$ induit une surjection
$H^{-1}(E \otimes_A^{\mathbf{L}} k) \to
H^{-1}(\NL_{A/\Lambda} \otimes_A^{\mathbf{L}} k)$
et une injection
$H^0(E \otimes_A^{\mathbf{L}} k) \to
H^0(\NL_{A/\Lambda} \otimes_A^{\mathbf{L}} k)$.
\end{enumerate}
Alors on a toujours (2) $\Rightarrow$ (1), et on a (1) $\Rightarrow$ (2)
si $u_0$ est un point fermé.
\end{lemma}

\begin{proof}
Soit $\mathfrak p = \Ker(A \to k)$ l'idéal premier correspondant à $u_0$.

\medskip\noindent
Supposons que $x$ soit versel en $u_0$ et que $u_0$ soit un point fermé de $U$.
Si $H^{-1}(\xi \otimes_A^{\mathbf{L}} k)$ n'est pas surjectif, alors
soit $A' \to A$ une extension de noyau $I$ comme dans le
Lemme \ref{artin-lemma-construct-essential-surjection}.
Puisque $u_0$ est un point fermé, nous voyons que $I$ est un $A$-module de type fini,
et donc que $A'$ est une $\Lambda$-algèbre de type fini (ceci est faux si
$u_0$ n'est pas fermé). En particulier, $A'$ est noethérien.
D'après la propriété (c) de $A'$ et la propriété (i) de $\xi$, nous voyons que $x$ se relève en
un objet $x'$ sur $A'$.
Soit $\mathfrak p' \subset A'$ le noyau de l'homomorphisme surjectif vers $k$.
D'après Artin--Rees (Algèbre, Lemme \ref{algebra-lemma-Artin-Rees}),
il existe un $n > 1$ tel que $(\mathfrak p')^n \cap I = 0$.
Nous voyons alors que
$$
B' = A'/(\mathfrak p')^n \longrightarrow A/\mathfrak p^n = B
$$
est une petite extension essentielle d'anneaux locaux artiniens ; voir
Théorie formelle des déformations, Lemme
\ref{formal-defos-lemma-essential-surjection}.
D'autre part, puisque $x$ est versel en $u_0$ et que $x'|_{\Spec(B')}$
est un relèvement de $x|_{\Spec(B)}$, il existe un entier
$m \geq n$ et un morphisme $q : A/\mathfrak p^m \to B'$
tel que le composé
$A/\mathfrak p^m \to B' \to B$ soit le morphisme quotient.
Puisque $B'$ est tel que la puissance $n$-ième de son idéal maximal est nulle, ce
$q$ se factorise par $B$, ce qui contredit le fait que $B' \to B$ est une
surjection essentielle. Cette contradiction montre que
$H^{-1}(\xi \otimes_A^{\mathbf{L}} k)$
est surjectif.

\medskip\noindent
Supposons que $x$ soit versel en $u_0$. D'après le Lemme \ref{artin-lemma-compute-ext-into-field},
le morphisme $H^0(\xi \otimes_A^{\mathbf{L}} k)$ est dual du morphisme
$\Ext^0_A(\NL_{A/\Lambda}, k) \to \text{Ext}^0_A(E, k)$. Remarquons que
$$
\Ext^0_A(\NL_{A/\Lambda}, k) = \text{Der}_\Lambda(A, k)
\quad\text{et}\quad
T_x(k) = \Ext^0_A(E, k)
$$
La condition (ii) assure que le morphisme
$\Ext^0_A(\NL_{A/\Lambda}, k) \to \text{Ext}^0_A(E, k)$
envoie un vecteur tangent $\theta$ à $U$ en $u_0$ sur la déformation
infinitésimale correspondante de $x_0$ ; voir la Remarque \ref{artin-remark-canonical-element}.
Par conséquent, si $x$ est versel, ce morphisme est surjectif ; voir
Théorie formelle des déformations, Lemme \ref{formal-defos-lemma-versal-criterion}.
Ainsi, $H^0(\xi \otimes_A^{\mathbf{L}} k)$ est injectif.
Ceci achève la démonstration de (1) $\Rightarrow$ (2) lorsque $u_0$ est un
point fermé.

\medskip\noindent
Pour la suite de la démonstration, supposons que $H^{-1}(E \otimes_A^\mathbf{L} k) \to
H^{-1}(\NL_{A/\Lambda} \otimes_A^\mathbf{L} k)$
soit surjectif et que
$H^0(E \otimes_A^\mathbf{L} k) \to
H^0(\NL_{A/\Lambda} \otimes_A^\mathbf{L} k)$
soit injectif. Posons $R = A_\mathfrak p^\wedge$ et soit $\eta$ l'objet
formel sur $R$ associé à $x|_{\Spec(R)}$.
Le morphisme $d\underline{\eta}$ sur les espaces tangents est surjectif
car il s'identifie au dual du morphisme injectif
$H^0(E \otimes_A^{\mathbf{L}} k) \to
H^0(\NL_{A/\Lambda} \otimes_A^{\mathbf{L}} k)$
(voir le paragraphe précédent). D'après le
Lemme \ref{formal-defos-lemma-versal-criterion} de la Théorie formelle des déformations,
il suffit de démontrer ce qui suit :
Soit $C' \to C$ une petite extension d'algèbres locales artiniennes
de type fini sur $\Lambda$, de corps résiduel $k$. Soit $R \to C$ un
homomorphisme de $\Lambda$-algèbres compatible avec les identifications des corps résiduels.
Soit $y = x|_{\Spec(C)}$ et soit $y'$ un relèvement de $y$ à $C'$.
Il faut montrer que l'on peut relever l'homomorphisme de $\Lambda$-algèbres $R \to C$ en un homomorphisme $R \to C'$.

\medskip\noindent
Observons qu'il suffit de relever l'homomorphisme de $\Lambda$-algèbres $A \to C$.
Soit $I = \Ker(C' \to C)$. Remarquons que $I$ est de dimension $1$ sur $k$ comme espace vectoriel.
L'obstruction $ob$ au relèvement de $A \to C$ est un élément de
$\Ext^1_A(\NL_{A/\Lambda}, I)$ ; voir l'Exemple \ref{artin-example-key}.
D'après le Lemme \ref{artin-lemma-compute-ext-into-field} et notre hypothèse, le morphisme
$\xi$ induit une injection
$$
\Ext^1_A(\NL_{A/\Lambda}, I)
\longrightarrow
\Ext^1_A(E, I)
$$
D'après la construction de $ob$ et (i), l'image de $ob$ dans $\Ext^1_A(E, I)$
est l'obstruction au relèvement de $x$ à $A \times_C C'$. D'après (RS*), le fait que
$y/C$ ait pour relèvement $y'/C'$ implique que $x$ se relève à $A \times_C C'$. Ainsi,
$ob = 0$, et la démonstration est achevée.
\end{proof}

\noindent
Le lemme clé ci-dessus nous permet de conclure à l'ouverture de la
versalité dans certains cas.

\begin{lemma}
\label{artin-lemma-openness}
Soit $S$ un schéma localement noethérien. Soit $\mathcal{X}$ une catégorie
fibrée en groupoïdes sur $(\Sch/S)_{fppf}$ vérifiant (RS*).
Soit $U = \Spec(A)$ un schéma affine de type fini sur $S$ dont le morphisme se factorise
par un ouvert affine $\Spec(\Lambda)$. Soit $x$ un objet de $\mathcal{X}$
sur $U$. Soit $\xi : E \to \NL_{A/\Lambda}$ un morphisme de $D^{-}(A)$.
Supposons que
\begin{enumerate}
\item[(i)] pour toute situation de déformation $(x, A' \to A)$, on ait :
$x$ se relève à $\Spec(A')$ si et seulement si
$E \to \NL_{A/\Lambda} \to \NL_{A/A'}$ est nul,
\item[(ii)] il existe un isomorphisme de foncteurs
$T_x(-) \to \Ext^0_A(E, -)$
tel que $E \to \NL_{A/\Lambda} \to \Omega^1_{A/\Lambda}$
corresponde à l'élément canonique (voir la
Remarque \ref{artin-remark-canonical-element}),
\item[(iii)] les groupes de cohomologie de $E$ soient des $A$-modules de type fini.
\end{enumerate}
Si $x$ est versel en un point fermé $u_0 \in U$,
alors il existe un voisinage ouvert $u_0 \in U' \subset U$
tel que $x$ soit versel en tout point de type fini de $U'$.
\end{lemma}

\begin{proof}
Soit $C$ le cône de $\xi$, de sorte que nous ayons un triangle distingué
$$
E \to \NL_{A/\Lambda} \to C \to E[1]
$$
dans $D^{-}(A)$. D'après le Lemme \ref{artin-lemma-characterize-versal},
l'hypothèse selon laquelle $x$ est versel en $u_0$ implique que
$H^{-1}(C \otimes^\mathbf{L} k) = 0$. D'après
Compléments sur l'algèbre, Lemme \ref{more-algebra-lemma-cut-complex-in-two},
il existe un $f \in A$ qui n'appartient pas à l'idéal premier correspondant à
$u_0$ tel que $H^{-1}(C \otimes^\mathbf{L}_A M) = 0$ pour
tout $A_f$-module $M$. En utilisant de nouveau le
Lemme \ref{artin-lemma-characterize-versal},
nous obtenons la versalité en tout point de type fini de
l'ouvert $D(f) \subset U$.
\end{proof}

\noindent
Les lemmes techniques ci-dessus suggèrent la définition suivante.

\begin{definition}
\label{artin-definition-naive-obstruction-theory}
Soit $S$ une base localement noethérienne. Soit $\mathcal{X}$ une catégorie fibrée
en groupoïdes sur $(\Sch/S)_{fppf}$. Supposons que $\mathcal{X}$
vérifie (RS*). Une {\it théorie naïve de l'obstruction} est
constituée des données suivantes :
\begin{enumerate}
\item
\label{artin-item-map}
pour toute $S$-algèbre $A$ telle que $\Spec(A) \to S$
se factorise par un ouvert affine $\Spec(\Lambda) \subset S$, et pour tout objet $x$
de $\mathcal{X}$ sur $\Spec(A)$, on se donne un objet $E_x \in D^-(A)$
et un morphisme $\xi_x : E \to \NL_{A/\Lambda}$,
\item
\label{artin-item-inf}
étant donné $(x, A)$ comme dans (\ref{artin-item-map}), il existe des transformations de
foncteurs
$$
\text{Inf}_x( - ) \to \Ext^{-1}_A(E_x, -)
\quad\text{et}\quad
T_x(-) \to \Ext^0_A(E_x, -)
$$
\item
\label{artin-item-functoriality}
pour $(x, A)$ comme dans (\ref{artin-item-map}) et un homomorphisme d'anneaux $A \to B$,
en posant $y = x|_{\Spec(B)}$, il existe un morphisme de fonctorialité
$E_x \to E_y$ dans $D(A)$.
\end{enumerate}
Ces données sont soumises aux conditions suivantes :
\begin{enumerate}
\item[(i)]
dans la situation de (\ref{artin-item-functoriality}), le diagramme
$$
\xymatrix{
E_y \ar[r]_{\xi_y} & \NL_{B/\Lambda} \\
E_x \ar[u] \ar[r]^{\xi_x} & \NL_{A/\Lambda} \ar[u]
}
$$
est commutatif dans $D(A)$,
\item[(ii)]
étant donnés $(x, A)$ comme dans (\ref{artin-item-map}) et $A \to B \to C$,
en posant $y = x|_{\Spec(B)}$ et $z = x|_{\Spec(C)}$, le
composé des morphismes de fonctorialité $E_x \to E_y$ et $E_y \to E_z$ est
le morphisme de fonctorialité $E_x \to E_z$,
\item[(iii)]
les morphismes de (\ref{artin-item-inf}) sont des isomorphismes
compatibles avec les morphismes de fonctorialité
et avec les morphismes de la Remarque \ref{artin-remark-functoriality},
\item[(iv)]
le composé $E_x \to \NL_{A/\Lambda} \to \Omega_{A/\Lambda}$
correspond à l'élément canonique de
$T_x(\Omega_{A/\Lambda}) = \Ext^0(E_x, \Omega_{A/\Lambda})$ ; voir la
Remarque \ref{artin-remark-canonical-element},
\item[(v)]
étant donnée une situation de déformation $(x, A' \to A)$ avec $I = \Ker(A' \to A)$,
le composé $E_x \to \NL_{A/\Lambda} \to \NL_{A/A'}$ est nul dans
$$
\Hom_A(E_x, \NL_{A/\Lambda}) = \Ext^0_A(E_x, \NL_{A/A'}) =
\Ext^1_A(E_x, I)
$$
si et seulement si $x$ se relève à $A'$.
\end{enumerate}
\end{definition}

\noindent
Nous voyons ainsi, en particulier, que nous obtenons une théorie de l'obstruction
comme dans la section \ref{artin-section-obstruction-theory} en posant
$\mathcal{O}_x( - ) = \Ext^1_A(E_x, -)$.

\begin{lemma}
\label{artin-lemma-naive-obstruction-theory-qis}
Soient $S$ et $\mathcal{X}$ comme dans la
Définition \ref{artin-definition-naive-obstruction-theory},
et munissons $\mathcal{X}$ d'une théorie naïve de l'obstruction.
Soient $A \to B$ et $y \to x$ comme dans (\ref{artin-item-functoriality}).
Soit $k$ une $B$-algèbre qui est un corps. Alors le morphisme de fonctorialité
$E_x \to E_y$ induit des bijections
$$
H^i(E_x \otimes_A^{\mathbf{L}} k) \to H^i(E_y \otimes_B^{\mathbf{L}} k)
$$
pour $i = 0, 1$.
\end{lemma}

\begin{proof}
Soit $z = x|_{\Spec(k)}$. Alors (RS*) implique que
$$
\textit{Lift}(x, A[k]) = \textit{Lift}(z, k[k])
\quad\text{et}\quad
\textit{Lift}(y, B[k]) = \textit{Lift}(z, k[k])
$$
car $A[k] = A \times_k k[k]$ et $B[k] = B \times_k k[k]$.
Les propriétés d'une théorie naïve de l'obstruction impliquent donc que le
morphisme de fonctorialité $E_x \to E_y$ induit des bijections
$\Ext^i_A(E_x, k) \to \text{Ext}^i_B(E_y, k)$
pour $i = -1, 0$. D'après le Lemme \ref{artin-lemma-compute-ext-into-field}, nos morphismes
$H^i(E_x \otimes_A^{\mathbf{L}} k) \to H^i(E_y \otimes_B^{\mathbf{L}} k)$,
$i = 0, 1$, induisent des isomorphismes sur les espaces vectoriels duaux ; ce sont donc des isomorphismes.
\end{proof}

\begin{lemma}
\label{artin-lemma-naive-obstruction-theory-gives-openness}
Soit $S$ un schéma localement noethérien. Soit
$p : \mathcal{X} \to (\Sch/S)_{fppf}^{opp}$ une catégorie fibrée en groupoïdes.
Supposons que $\mathcal{X}$ vérifie (RS*)
et que $\mathcal{X}$ possède une théorie naïve de l'obstruction.
Alors l'ouverture de la versalité est satisfaite pour $\mathcal{X}$, pourvu que les
complexes $E_x$ de la Définition \ref{artin-definition-naive-obstruction-theory}
aient des groupes de cohomologie de type fini pour les couples $(A, x)$ où
$A$ est de type fini sur $S$.
\end{lemma}

\begin{proof}
Soit $U$ un schéma localement de type fini sur $S$, soit $x$ un objet de
$\mathcal{X}$ sur $U$, et soit $u_0$ un point de type fini de $U$ tel que
$x$ soit versel en $u_0$. Nous pouvons d'abord remplacer $U$ par un ouvert affine
tel que $u_0$ y soit fermé et que le morphisme $U \to S$ se factorise par un ouvert
affine $\Spec(\Lambda)$. Écrivons $U = \Spec(A)$. Soit
$\xi_x : E_x \to \NL_{A/\Lambda}$ le morphisme d'obstruction.
Nous pouvons alors appliquer le Lemme \ref{artin-lemma-openness} pour conclure.
\end{proof}
 
 
 
 
 
 
 
 
 
\section{Une notion duale}
\label{artin-section-dual}

\noindent
Soit $(x, A' \to A)$ une situation de déformation pour une catégorie donnée
$\mathcal{X}$ fibrée en groupoïdes au-dessus d'un schéma localement noethérien $S$.
Supposons que $\mathcal{X}$ soit munie d'une théorie des obstructions ; voir la
définition \ref{artin-definition-obstruction-theory}. En pratique,
on dispose souvent d'un complexe $K^\bullet$ de $A$-modules et d'isomorphismes de
foncteurs
$$
\text{Inf}_x(-) \to H^0(K^\bullet \otimes_A^\mathbf{L} -),\quad
T_x(-) \to H^1(K^\bullet \otimes_A^\mathbf{L} -),\quad
\mathcal{O}_x(-) \to H^2(K^\bullet \otimes_A^\mathbf{L} -)
$$
Dans cette section, nous formalisons quelque peu cette situation et montrons comment elle conduit,
dans certains cas, à vérifier l'ouverture de la versalité.

\begin{example}
\label{artin-example-global-sections-dual}
Soient $\Lambda, S, W, \mathcal{F}$ comme dans
l'exemple \ref{artin-example-global-sections}.
Supposons que $W \to S$ soit propre et que $\mathcal{F}$ soit cohérent. D'après
Cohomologie des schémas, remarque
\ref{coherent-remark-explain-perfect-direct-image}
il existe un complexe fini de $\Lambda$-modules projectifs de type fini
$N^\bullet$ qui calcule universellement la cohomologie de $\mathcal{F}$.
En particulier, les espaces d'obstructions de l'exemple \ref{artin-example-global-sections}
sont $\mathcal{O}_x(M) = H^1(N^\bullet \otimes_\Lambda M)$.
Ainsi, avec $K^\bullet = N^\bullet \otimes_\Lambda A[-1]$, nous obtenons
$\mathcal{O}_x(M) = H^2(K^\bullet \otimes_A^\mathbf{L} M)$.
\end{example}

\begin{situation}
\label{artin-situation-dual}
Soit $S$ un schéma localement noethérien. Soit $\mathcal{X}$ une catégorie
fibrée en groupoïdes au-dessus de $(\Sch/S)_{fppf}$. Supposons que
$\mathcal{X}$ satisfasse (RS*), de sorte que le foncteur $T_x(-)$ soit défini ; voir le
lemme \ref{artin-lemma-properties-lift-RS-star}.
Soit $U = \Spec(A)$ un schéma affine de type fini sur $S$ dont l'image
est contenue dans un ouvert affine $\Spec(\Lambda)$. Soit $x$ un objet de $\mathcal{X}$
au-dessus de $U$. Supposons données
\begin{enumerate}
\item un complexe de $A$-modules $K^\bullet$ ;
\item une transformation de foncteurs
$T_x(-) \to H^1(K^\bullet \otimes_A^\mathbf{L} -)$ ;
\item pour toute situation de déformation $(x, A' \to A)$ de noyau
$I = \Ker(A' \to A)$, un élément
$o_x(A') \in H^2(K^\bullet \otimes_A^\mathbf{L} I)$
\end{enumerate}
qui satisfont les conditions (minimales) suivantes :
\begin{enumerate}
\item[(i)] la transformation
$T_x(-) \to H^1(K^\bullet \otimes_A^\mathbf{L} -)$
est un isomorphisme ;
\item[(ii)] pour tout morphisme $(x, A'' \to A) \to (x, A' \to A)$ de situations
de déformation, l'élément $o_x(A')$ s'envoie sur l'élément $o_x(A'')$
par l'application
$H^2(K^\bullet \otimes_A^\mathbf{L} I) \to
H^2(K^\bullet \otimes_A^\mathbf{L} I')$
où $I' = \Ker(A'' \to A)$ ;
\item[(iii)] $x$ se relève en un objet au-dessus de $\Spec(A')$ si et seulement si
$o_x(A') = 0$.
\end{enumerate}
On pourrait aussi incorporer les automorphismes infinitésimaux, mais
nous ne le faisons pas afin d'obtenir le résultat le plus précis possible.
\end{situation}

\noindent
Dans la situation \ref{artin-situation-dual}, le complexe
$K^\bullet \otimes_A^\mathbf{L} \NL_{A/\Lambda}$ jouera un rôle important. Supposons donné un
élément $\xi \in H^1(K^\bullet \otimes_A^\mathbf{L} \NL_{A/\Lambda})$.
Alors (1), pour toute surjection $A' \to A$ de $\Lambda$-algèbres de noyau
$I$ de carré nul, l'application canonique $\NL_{A/\Lambda} \to \NL_{A/A'} = I[1]$
envoie $\xi$ sur un élément $\xi_{A'} \in H^2(K^\bullet \otimes_A^\mathbf{L} I)$,
et (2) l'application $\NL_{A/\Lambda} \to \Omega_{A/\Lambda}$ envoie
$\xi$ sur un élément $\xi_{can}$ de
$H^1(K^\bullet \otimes_A^\mathbf{L} \Omega_{A/\Lambda})$.

\begin{lemma}
\label{artin-lemma-dual-obstruction}
Dans la situation \ref{artin-situation-dual}, supposons en outre que
\begin{enumerate}
\item[(iv)] pour toute suite exacte courte de situations de déformation
comme dans la remarque \ref{artin-remark-short-exact-sequence-thickenings} et tout
relèvement $x'_2 \in \text{Lift}(x, A_2')$,
$o_x(A_3') \in H^2(K^\bullet \otimes_A^\mathbf{L} I_3)$
soit égal à $\partial\theta$, où
$\theta \in H^1(K^\bullet \otimes_A^\mathbf{L} I_1)$
est l'élément correspondant à $x'_2|_{\Spec(A_1')}$ par
$A_1' = A[I_1]$ et l'application donnée
$T_x(-) \to H^1(K^\bullet \otimes_A^\mathbf{L} -)$.
\end{enumerate}
Alors il existe un élément
$\xi \in H^1(K^\bullet \otimes_A^\mathbf{L} \NL_{A/\Lambda})$
tel que
\begin{enumerate}
\item pour toute situation de déformation $(x, A' \to A)$, on ait
$\xi_{A'} = o_x(A')$ ;
\item $\xi_{can}$ corresponde à l'élément canonique de la
remarque \ref{artin-remark-canonical-element} par la transformation donnée
$T_x(-) \to H^1(K^\bullet \otimes_A^\mathbf{L} -)$.
\end{enumerate}
\end{lemma}

\begin{proof}
Choisissons un homomorphisme $\alpha : \Lambda[x_1, \ldots, x_n] \to A$ de noyau $J$.
Posons $P = \Lambda[x_1, \ldots, x_n]$. Dans la suite de cette démonstration, nous travaillons avec
$$
\NL(\alpha) = (J/J^2 \longrightarrow \bigoplus A \text{d}x_i)
$$
ce qui est permis par
Algèbre, lemme \ref{algebra-lemma-NL-homotopy},
et
Compléments d'algèbre, lemme \ref{more-algebra-lemma-derived-tor-homotopy}.
Considérons l'élément
$o_x(P/J^2) \in H^2(K^\bullet \otimes_A^\mathbf{L} J/J^2)$ ainsi que
le quotient
$$
C = (P/J^2 \times \bigoplus A \text{d}x_i)/(J/J^2)
$$
où $J/J^2$ est plongé diagonalement. Remarquons que $C \to A$ est une surjection
de noyau $\bigoplus A\text{d}x_i$. Il existe en outre une section
$A \to C$ de $C \to A$, qui envoie la classe de $f \in P$ sur la classe
de $(f, \text{d}f)$ dans la somme amalgamée. Pour la suite, notons $x_C$
l'image inverse de $x$ par le morphisme correspondant $\Spec(C) \to \Spec(A)$.
Nous avons donc $o_x(C) = 0$.
Il s'ensuit que $o_x(P/J^2)$ s'envoie sur zéro dans
$H^2(K^\bullet \otimes_A^\mathbf{L} \bigoplus A\text{d}x_i)$.
Il existe par conséquent un élément
$\xi \in H^1(K^\bullet \otimes_A^\mathbf{L} \NL(\alpha))$
qui s'envoie sur $o_x(P/J^2)$.

\medskip\noindent
Remarquons que, pour toute situation de déformation $(x, A' \to A)$, il existe
un homomorphisme de $\Lambda$-algèbres $P/J^2 \to A'$ compatible aux augmentations
vers $A$. Par conséquent, l'élément
$\xi$ satisfait par construction la première propriété du lemme
et la propriété (ii) de la situation \ref{artin-situation-dual}.

\medskip\noindent
Remarquons que notre choix de $\xi$ est bien défini à l'ajout près d'un
élément de $H^1(K^\bullet \otimes_A^\mathbf{L} \bigoplus A\text{d}x_i)$.
Nous allons montrer qu'en modifiant $\xi$ par un élément de ce
groupe, on peut faire en sorte que la seconde assertion du lemme soit satisfaite.
Soit $C' \subset C$ l'image de $P/J^2$ par l'homomorphisme de
$\Lambda$-algèbres $P/J^2 \to C$ (inclusion du premier facteur).
Observons que
$\Ker(C' \to A) = \Im(J/J^2 \to \bigoplus A\text{d}x_i)$.
Posons $\overline{C} = A[\Omega_{A/\Lambda}]$. L'application
$P/J^2 \times \bigoplus A \text{d}x_i \to \overline{C}$,
$(f, \sum f_i \text{d}x_i) \mapsto (f \bmod J, \sum f_i \text{d}x_i)$
se factorise par une application surjective $C \to \overline{C}$. Alors
$$
(x, \overline{C} \to A) \to (x, C \to A) \to (x, C' \to A)
$$
est une suite exacte courte de situations de déformation. Le
scindage associé $\overline{C} = A[\Omega_{A/\Lambda}]$ (issu de la
remarque \ref{artin-remark-short-exact-sequence-thickenings}) est égal au
scindage donné ci-dessus. En outre, la section $A \to C$ composée avec l'application
$C \to \overline{C}$
est l'application $(1, \text{d}) : A \to A[\Omega_{A/\Lambda}]$ de la
remarque \ref{artin-remark-canonical-element}.
Ainsi, la restriction de $x_C$ est l'élément canonique $x_{can}$ de
$T_x(\Omega_{A/\Lambda}) = \text{Lift}(x, A[\Omega_{A/\Lambda}])$.
La condition (iv) montre que $o_x(P/J^2)$ s'envoie sur $\partial x_{can}$
dans
$$
H^1(K^\bullet \otimes_A^\mathbf{L} \Im(J/J^2 \to \bigoplus A\text{d}x_i))
$$
Par construction, $\xi$ s'envoie sur $o_x(P/J^2)$. Il en résulte que
$x_{can}$ et $\xi_{can}$ s'envoient sur le même élément dans le
groupe affiché, ce qui signifie (d'après la suite exacte longue de cohomologie)
qu'ils diffèrent d'un élément de
$H^1(K^\bullet \otimes_A^\mathbf{L} \bigoplus A\text{d}x_i)$
comme souhaité.
\end{proof}

\begin{lemma}
\label{artin-lemma-dual-openness}
Dans la situation \ref{artin-situation-dual}, supposons que la condition (iv) du
lemme \ref{artin-lemma-dual-obstruction} soit satisfaite et que $K^\bullet$ soit un
objet parfait de $D(A)$. Si $x$ est versel en un point fermé
$u_0 \in U$, il existe alors un voisinage ouvert
$u_0 \in U' \subset U$ tel que $x$ soit versel en tout point de type fini
de $U'$.
\end{lemma}

\begin{proof}
Nous pouvons supposer que $K^\bullet$ est un complexe fini de
$A$-modules projectifs de type fini. Le produit tensoriel dérivé par $K^\bullet$
revient donc au simple produit tensoriel par $K^\bullet$. Soit
$E^\bullet$ le complexe parfait dual de $K^\bullet$ ; voir
Compléments d'algèbre, lemme \ref{more-algebra-lemma-dual-perfect-complex}.
(Ainsi $E^n = \Hom_A(K^{-n}, A)$, les différentielles étant les transposées des
différentielles de $K^\bullet$.) Notons $E \in D^{-}(A)$ l'objet
représenté par le complexe $E^\bullet[-1]$.
Soit $\xi \in H^1(\text{Tot}(K^\bullet \otimes_A \NL_{A/\Lambda}))$
l'élément construit dans le lemme \ref{artin-lemma-dual-obstruction},
et notons $\xi : E = E^\bullet[-1] \to \NL_{A/\Lambda}$ l'application
correspondante (loc. cit.). Nous affirmons que le couple $(E, \xi)$ satisfait toutes les
hypothèses du lemme \ref{artin-lemma-openness}, ce qui achève la démonstration.

\medskip\noindent
En effet, l'hypothèse (i) du lemme \ref{artin-lemma-openness} découle de la conclusion
(1) du lemme \ref{artin-lemma-dual-obstruction}
et de l'égalité $H^2(K^\bullet \otimes_A^\mathbf{L} -) =
\Ext^1(E, -)$, d'après la référence citée. L'hypothèse (ii) du
lemme \ref{artin-lemma-openness} découle de la conclusion (2) du
lemme \ref{artin-lemma-dual-obstruction}
et de l'égalité $H^1(K^\bullet \otimes_A^\mathbf{L} -) =
\Ext^0(E, -)$, d'après la même référence. L'hypothèse (iii) du lemme \ref{artin-lemma-openness}
est claire.
\end{proof}







\section{Foncteurs commutant aux limites sur les schémas noethériens}
\sectionmark{Foncteurs commutant aux limites}
\label{artin-section-noetherian}

\noindent
Il est parfois commode de considérer des foncteurs ou des champs définis seulement
sur la sous-catégorie pleine des schémas (localement) noethériens. Dans cette section,
nous examinons cette question dans le cas des espaces algébriques.

\medskip\noindent
Soit $S$ un schéma localement noethérien. Soyons un peu pédants afin
d'ajuster correctement nos catégories ; ceux qui font abstraction des questions
ensemblistes peuvent simplement remplacer les ensembles de schémas choisis par la collection
de tous les schémas dans ce qui suit. Comme dans
Topologies, remarque \ref{topologies-remark-choice-sites},
nous choisissons une catégorie $\Sch_\alpha$ de schémas contenant
$S$ de sorte que nous obtenions les gros sites $(\Sch/S)_{Zar}$,
$(\Sch/S)_\etale$, $(\Sch/S)_{smooth}$, $(\Sch/S)_{syntomic}$
et $(\Sch/S)_{fppf}$, ayant tous la même
catégorie sous-jacente $\Sch_\alpha/S$. Notons
$$
\textit{Noetherian}_\alpha \subset \Sch_\alpha
$$
la sous-catégorie pleine formée des schémas localement noethériens.
Cela détermine une sous-catégorie pleine
$$
\textit{Noetherian}_\alpha/S \subset \Sch_\alpha/S
$$
Pour $\tau \in \{Zariski, \etale, smooth, syntomic, fppf\}$, on a
\begin{enumerate}
\item si $f : X \to Y$ est un morphisme de $\Sch_\alpha/S$
où $Y$ appartient à $\textit{Noetherian}_\alpha/S$ et si
$f$ est localement de type fini, alors $X$ appartient à $\textit{Noetherian}_\alpha/S$,
\item pour des morphismes $f : X \to Y$ et $g : Z \to Y$ de
$\textit{Noetherian}_\alpha/S$, avec $f$ localement de type fini,
le produit fibré $X \times_Y Z$ dans $\textit{Noetherian}_\alpha/S$
existe et coïncide avec le produit fibré dans $\Sch_\alpha/S$,
\item si $\{X_i \to X\}_{i \in I}$ est un recouvrement de
$(\Sch/S)_\tau$ et si $X$ appartient à $\textit{Noetherian}_\alpha/S$,
alors les objets $X_i$ appartiennent à $\textit{Noetherian}_\alpha/S$,
\item la catégorie $\textit{Noetherian}_\alpha/S$ munie
de l'ensemble des recouvrements de $(\Sch/S)_\tau$ dont les objets
appartiennent à $\textit{Noetherian}_\alpha/S$ est un site
que nous noterons $(\textit{Noetherian}/S)_\tau$,
\item le foncteur d'inclusion
$(\textit{Noetherian}/S)_\tau \to (\Sch/S)_\tau$
est pleinement fidèle, continu et cocontinu.
\end{enumerate}
D'après Sites, lemmes \ref{sites-lemma-cocontinuous-morphism-topoi} et
\ref{sites-lemma-when-shriek}, nous obtenons un morphisme de topos
$$
g_\tau : \Sh((\textit{Noetherian}/S)_\tau) \longrightarrow \Sh((\Sch/S)_\tau)
$$
dont le foncteur image inverse est la restriction des faisceaux suivant
le foncteur d'inclusion $(\textit{Noetherian}/S)_\tau \to (\Sch/S)_\tau$.

\begin{remark}[Avertissement]
\label{artin-remark-no-fibre-products}
Le site $(\textit{Noetherian}/S)_\tau$ ne possède pas de produits fibrés.
Il faut donc être prudent lorsque l'on travaille avec des faisceaux. Par exemple,
le foncteur d'inclusion continu
$(\textit{Noetherian}/S)_\tau \to (\Sch/S)_\tau$
ne définit pas un morphisme de sites. Voir
Exemples, section \ref{examples-section-sheaves-locally-Noetherian},
pour un exemple lorsque $\tau = fppf$.
\end{remark}

\noindent
Soit $F : (\textit{Noetherian}/S)_\tau^{opp} \to \textit{Ensembles}$
un foncteur. On dit que $F$ {\it commute aux limites} si, pour toute
limite projective filtrante de schémas affines $X = \lim X_i$ de
$(\textit{Noetherian}/S)_\tau$, on a $F(X) = \colim F(X_i)$.

\begin{lemma}
\label{artin-lemma-canonical-extension}
Soit $\tau \in \{Zariski, \etale, smooth, syntomic, fppf\}$.
La restriction suivant le foncteur d'inclusion
$(\textit{Noetherian}/S)_\tau \to (\Sch/S)_\tau$
définit une équivalence de catégories entre
\begin{enumerate}
\item la catégorie des faisceaux commutant aux limites sur
$(\Sch/S)_\tau$ et
\item la catégorie des faisceaux commutant aux limites sur
$(\textit{Noetherian}/S)_\tau$.
\end{enumerate}
\end{lemma}

\begin{proof}
Soit $F : (\textit{Noetherian}/S)_\tau^{opp} \to \textit{Ensembles}$
un foncteur qui commute aux limites et qui est un faisceau.
D'après Topologies, lemmes
\ref{topologies-lemma-extend} et \ref{topologies-lemma-extend-sheaf-general},
il existe un unique foncteur
$F' : (\Sch/S)_\tau^{opp} \to \textit{Ensembles}$
qui commute aux limites, est un faisceau et dont la restriction est $F$.
En fait, la construction de $F'$ dans
Topologies, lemme \ref{topologies-lemma-extend},
est fonctorielle en $F$, et cette construction est un foncteur quasi-inverse
de la restriction. Nous omettons quelques détails.
\end{proof}

\begin{lemma}
\label{artin-lemma-representable-limit-preserving}
Soit $X$ un objet de $(\textit{Noetherian}/S)_\tau$. Si le foncteur
des points $h_X : (\textit{Noetherian}/S)_\tau^{opp} \to \textit{Ensembles}$
commute aux limites, alors $X$ est localement de présentation finie sur $S$.
\end{lemma}

\begin{proof}
Soit $V \subset X$ un ouvert affine dont l'image est contenue dans un
ouvert affine $U \subset S$. On peut écrire $V = \lim V_i$ comme limite projective filtrante de
schémas affines $V_i$ de présentation finie sur $U$, voir
Algèbre, lemme \ref{algebra-lemma-ring-colimit-fp}.
Par hypothèse, la flèche $V \to X$ se factorise sous la forme $V \to V_i \to X$
pour un certain $i$. Quitte à augmenter $i$, on peut supposer que $V_i \to X$
se factorise par $V$ : en effet, l'image réciproque de $V \subset X$ dans $V_i$
est égale à $V_i$ pour un indice assez grand, d'après
Limites, lemme \ref{limits-lemma-descend-opens}.
Alors le morphisme identité $V \to V$ se factorise par $V_i$ pour un certain $i$
dans la catégorie des schémas sur $U$. Ainsi $V \to U$ est de présentation finie ;
le fait algébrique correspondant est que, si $B$ est une $A$-algèbre
telle que $\text{id} : B \to B$ se factorise par une
$A$-algèbre de présentation finie, alors $B$ est de présentation finie sur $A$ (joli exercice).
Par conséquent, $X$ est localement de présentation finie sur $S$.
\end{proof}

\noindent
Le lemme suivant possède une variante pour les morphismes de foncteurs
représentables par des espaces algébriques.

\begin{lemma}
\label{artin-lemma-representable}
Soit $\tau \in \{Zariski, \etale, smooth, syntomic, fppf\}$.
Soient $F', G' : (\Sch/S)_\tau^{opp} \to \textit{Ensembles}$ des foncteurs commutant aux limites
et qui sont des faisceaux. Soit $a' : F' \to G'$ un morphisme de foncteurs.
Notons $a : F \to G$ la restriction de $a' : F' \to G'$ à
$(\textit{Noetherian}/S)_\tau$. Les assertions suivantes sont équivalentes :
\begin{enumerate}
\item $a'$ est représentable (comme morphisme de foncteurs, voir
Catégories, définition \ref{categories-definition-representable-morphism}),
\item pour tout objet $V$ de $(\textit{Noetherian}/S)_\tau$
et tout morphisme $V \to G$, le produit fibré
$F \times_G V : (\textit{Noetherian}/S)_\tau^{opp} \to \textit{Ensembles}$
est un foncteur représentable,
\item même assertion qu'en (2), mais seulement pour $V$ affine de type fini sur $S$
dont l'image est contenue dans un ouvert affine de $S$.
\end{enumerate}
\end{lemma}

\begin{proof}
Supposons (1). D'après Limites d'espaces, lemme
\ref{spaces-limits-lemma-locally-finite-presentation-permanence},
le morphisme de foncteurs $a'$ commute aux limites\footnote{Cela
a un sens même si $\tau \not = fppf$, car la catégorie sous-jacente
à $(\Sch/S)_\tau$ est égale à la catégorie sous-jacente
à $(\Sch/S)_{fppf}$ et l'énoncé ne fait pas intervenir la topologie.}.
Soit $\xi : V \to G$ comme en (2). Notons $V' = V$, mais considéré comme un
objet de $(\Sch/S)_\tau$. Puisque $G$ est la restriction de
$G'$ à $(\textit{Noetherian}/S)_\tau$, on voit que
$\xi \in G(V)$ correspond à $\xi' \in G'(V')$.
Par hypothèse, $V' \times_{\xi', G'} F'$ est représentable
par un schéma $U'$. Le morphisme de schémas $U' \to V'$ correspondant à
la projection $V' \times_{\xi', G'} F' \to V'$ est localement de
présentation finie d'après
Limites d'espaces, lemme
\ref{spaces-limits-lemma-base-change-locally-finite-presentation} et
Limites, proposition
\ref{limits-proposition-characterize-locally-finite-presentation}.
Ainsi $U'$ est un schéma localement noethérien et, par conséquent, $U'$ est
isomorphe à un objet $U$ de $(\textit{Noetherian}/S)_\tau$.
Alors $U$ représente $F \times_G V$, comme voulu.

\medskip\noindent
L'implication (2) $\Rightarrow$ (3) est immédiate. Supposons (3).
Nous allons démontrer (1). Soit $T$ un objet de $(\Sch/S)_\tau$
et soit $T \to G'$ un morphisme. Il faut montrer que
le foncteur $F' \times_{G'} T$ est représentable par un schéma $X$
sur $T$. Soit $\mathcal{B}$ l'ensemble des ouverts affines de
$T$ dont l'image est contenue dans un ouvert affine de $S$. Ils forment une base
de la topologie de $T$. Nous montrerons ci-dessous que, pour $W \in \mathcal{B}$,
le produit fibré $F' \times_{G'} W$ est représentable
par un schéma $X_W$ sur $W$. Si $W_1 \subset W_2$ dans $\mathcal{B}$, alors
nous obtenons un isomorphisme $X_{W_1} \to X_{W_2} \times_{W_2} W_1$, car
$X_{W_1}$ et $X_{W_2} \times_{W_2} W_1$ représentent tous deux le foncteur
$F' \times_{G'} W_1$.
Ces isomorphismes sont canoniques et satisfont la condition de cocycle
mentionnée dans Constructions, lemme
\ref{constructions-lemma-relative-glueing}.
Nous pouvons donc recoller les schémas $X_W$ en un schéma $X$ sur $T$.
La compatibilité des isomorphismes de recollement avec les morphismes
$X_W \to F'$ nous fournit un morphisme $X \to F'$.
Le morphisme ainsi obtenu $X \to F' \times_{G'} T$ est un
isomorphisme, comme on peut le vérifier localement sur $T$ (la source
et le but de cette flèche étant des faisceaux pour la topologie de Zariski).

\medskip\noindent
Soit $W$ un schéma affine dont l'image est contenue dans un ouvert affine $U \subset S$.
Soit $W \to G'$ un morphisme.
Toujours sous l'hypothèse (3), il faut montrer que $F' \times_{G'} W$
est représentable par un schéma.
On peut écrire $W = \lim V'_i$ comme limite projective filtrante de schémas
affines $V'_i$ de présentation finie sur $U$, voir
Algèbre, lemme \ref{algebra-lemma-ring-colimit-fp}.
Puisque $V'_i$ est de type fini sur un schéma noethérien,
on voit que $V'_i$ est un schéma noethérien.
Notons $V_i = V'_i$, mais considéré comme un objet de
$(\textit{Noetherian}/S)_\tau$. Puisque $G'$
commute aux limites, on peut choisir un $i$ et un morphisme
$V'_i \to G'$ tels que $W \to G'$ soit le composé
$W \to V'_i \to G'$. Puisque $G$ est la restriction de $G'$
à $(\textit{Noetherian}/S)_\tau$, le morphisme $V'_i \to G'$
est la même chose qu'un morphisme $V_i \to G$ (voir ci-dessus).
Par hypothèse (3), le foncteur $F \times_G V_i$ est représentable par un objet
$X_i$ de $(\textit{Noetherian}/S)_\tau$.
Le foncteur $F \times_G V_i$ commute aux limites,
car il est la restriction de $F' \times_{G'} V'_i$,
et ce foncteur commute aux limites d'après
Limites d'espaces, lemme
\ref{spaces-limits-lemma-fibre-product-locally-finite-presentation},
l'hypothèse que $F'$ et $G'$ commutent aux limites, ainsi que
Limites, remarque \ref{limits-remark-limit-preserving}, qui
nous dit que le foncteur des points de $V'_i$ commute aux limites.
D'après le lemme \ref{artin-lemma-representable-limit-preserving},
on conclut que $X_i$ est localement de présentation finie sur $S$.
Notons $X'_i = X_i$, mais considéré comme un objet de
$(\Sch/S)_\tau$. Alors on voit que $F' \times_{G'} V'_i$
et le foncteur des points $h_{X'_i}$ sont tous deux des prolongements
de $h_{X_i} : (\textit{Noetherian}/S)_\tau^{opp} \to \textit{Ensembles}$
en des faisceaux sur $(\Sch/S)_\tau$ qui commutent aux limites.
Par l'équivalence de catégories du lemme \ref{artin-lemma-canonical-extension},
on en déduit que $X'_i$ représente $F' \times_{G'} V'_i$.
Enfin,
$$
F' \times_{G'} W = F' \times_{G'} V'_i \times_{V'_i} W =
X'_i \times_{V'_i} W
$$
est représentable, comme voulu.
\end{proof}






\section{Espaces algébriques dans le cadre noethérien}
\label{artin-section-algebraic-spaces-noetherian}

\noindent
Soit $S$ un schéma localement noethérien. Soit
$(\textit{Noetherian}/S)_\etale \subset (\Sch/S)_\etale$
le site étudié à la section \ref{artin-section-noetherian}.
Soit $F : (\textit{Noetherian}/S)_\etale^{opp} \to \textit{Ensembles}$
un foncteur, c'est-à-dire que $F$ est un préfaisceau sur $(\textit{Noetherian}/S)_\etale$.
Dans ce cadre, tous les axiomes [-1],  [0], [1], [2], [3], [4], [5] de la
section \ref{artin-section-axioms-functors} ont un sens. Nous allons les passer
en revue un à un et préciser exactement ce que nous entendons.

\medskip\noindent
Axiome [-1]. Il s'agit d'une condition ensembliste que peuvent ignorer
les lecteurs que les questions ensemblistes n'intéressent pas.
Elle a un sens pour $F$, puisqu'elle concerne l'évaluation de
$F$ sur les spectres de corps de type fini sur $S$, qui sont
des objets de $(\textit{Noetherian}/S)_\etale$.

\medskip\noindent
Axiome [0]. C'est l'axiome selon lequel $F$ est un faisceau
sur $(\textit{Noetherian}/S)_\etale^{opp}$, c'est-à-dire qu'il vérifie
la condition de faisceau pour les recouvrements \'etales.

\medskip\noindent
Axiome [1]. C'est l'axiome selon lequel $F$ commute aux limites au sens défini
à la section \ref{artin-section-noetherian} : pour toute
limite projective filtrante de schémas affines $X = \lim X_i$ de
$(\textit{Noetherian}/S)_\etale$, on a $F(X) = \colim F(X_i)$.

\medskip\noindent
Axiome [2]. C'est l'axiome selon lequel $F$ vérifie la condition
de Rim-Schlessinger (RS). D'après sa définition, donnée à la
Définition \ref{artin-definition-RS}, et la discussion de la
section \ref{artin-section-axioms-functors},
cela signifie que, pour toute somme amalgamée $Y' = Y \amalg_X X'$
de schémas de type fini sur $S$, où $Y, X, X'$
sont les spectres d'anneaux locaux artiniens, l'application
$$
F(Y \amalg_X X') \to F(Y) \times_{F(X)} F(X')
$$
est bijective. Cette condition a un sens, car les schémas
$X$, $X'$, $Y$ et $Y'$ appartiennent à $(\text{Noetherian}/S)_\etale$,
puisqu'ils sont de type fini sur $S$.

\medskip\noindent
Axiome [3]. C'est l'axiome selon lequel tout espace tangent $TF_{k, x_0}$
est de dimension finie. Cela a un sens, car les espaces tangents $TF_{k, x_0}$
sont construits à partir des évaluations de $F$ en $\Spec(k)$ et
en $\Spec(k[\epsilon])$, où $k$ est un corps de type fini sur $S$,
et sont donc obtenus par évaluation sur des objets de la catégorie
$(\textit{Noetherian}/S)_\etale$.

\medskip\noindent
Axiome [4]. C'est l'axiome selon lequel tout objet formel est effectif.
D'après la discussion des
sections \ref{artin-section-formal-objects} et \ref{artin-section-axioms-functors},
cela ne fait intervenir que l'évaluation de notre foncteur sur des schémas noethériens,
et cette condition a donc un sens pour $F$.

\medskip\noindent
Axiome [5]. C'est l'axiome selon lequel $F$ vérifie l'ouverture de la versalité.
Rappelons que cela signifie ce qui suit : étant donnés un schéma $U$
localement de type fini sur $S$, un élément $x \in F(U)$ et
un point de type fini $u_0 \in U$ tel que $x$ soit versel en $u_0$,
il existe alors un voisinage ouvert $u_0 \in U' \subset U$
tel que $x$ soit versel en tout point de type fini de $U'$.
Comme précédemment, sa vérification ne fait intervenir que l'évaluation
de notre foncteur sur des schémas noethériens.

\begin{proposition}
\label{artin-proposition-spaces-diagonal-representable-noetherian}
Soit $S$ un schéma localement noethérien. Soit
$F : (\textit{Noetherian}/S)_\etale^{opp} \to \textit{Ensembles}$
un foncteur. Supposons que
\begin{enumerate}
\item $\Delta : F \to F \times F$ soit représentable
(comme transformation de foncteurs, voir
Catégories, Définition \ref{categories-definition-representable-morphism}),
\item $F$ vérifie les axiomes [-1], [0], [1], [2], [3], [4], [5]
(voir ci-dessus), et
\item $\mathcal{O}_{S, s}$ soit un anneau G pour tout point $s$ de type fini de $S$.
\end{enumerate}
Alors il existe un unique espace algébrique
$F' : (\Sch/S)_{fppf}^{opp} \to \textit{Ensembles}$
dont la restriction à $(\textit{Noetherian}/S)_\etale$ est $F$
(voir la démonstration pour plus de précisions).
\end{proposition}

\begin{proof}
Rappelons que les sites $(\Sch/S)_{fppf}$ et $(\Sch/S)_\etale$ ont la même
catégorie sous-jacente ; voir la discussion de la section \ref{artin-section-noetherian}.
De même, les sites $(\textit{Noetherian}/S)_\etale$ et
$(\textit{Noetherian}/S)_{fppf}$ ont la même catégorie sous-jacente.
Par les axiomes [0] et [1], le foncteur $F$ est un faisceau et
commute aux limites.
Soit $F' : (\Sch/S)_\etale^{opp} \to \textit{Ensembles}$
l'unique prolongement de $F$ qui soit un faisceau (pour la topologie \'etale)
et qui commute aux limites ; voir le
Lemme \ref{artin-lemma-canonical-extension}.
Alors $F'$ vérifie les axiomes [0] et [1] tels qu'ils sont énoncés à la
section \ref{artin-section-axioms-functors}.
D'après le Lemme \ref{artin-lemma-representable},
$\Delta' : F' \to F' \times F'$ est représentable (par des schémas).
D'autre part, il est immédiat que
$F'$ vérifie les axiomes [-1], [2], [3], [4], [5] de la
section \ref{artin-section-axioms-functors},
puisque chacun de ces axiomes ne fait intervenir que l'évaluation de $F'$ sur des objets
de $(\textit{Noetherian}/S)_\etale$ et que nous avons supposé les
conditions correspondantes pour $F$.
Ainsi, $F'$ est un espace algébrique d'après la
Proposition \ref{artin-proposition-spaces-diagonal-representable}.
\end{proof}












\section{Théorème d'Artin sur les contractions}
\label{artin-section-contractions}

\noindent
Dans cette section, nous utiliserons librement le langage des espaces algébriques
formels, voir Espaces formels, section \ref{formal-spaces-section-introduction}.
Le théorème d'Artin sur les contractions est l'un des deux principaux théorèmes de
l'article d'Artin \cite{ArtinII} ; le premier est son théorème sur les dilatations,
que nous avons énoncé et démontré dans Algébrisation des espaces formels,
section \ref{restricted-section-dilatations}.

\begin{situation}
\label{artin-situation-contractions}
Soit $S$ un schéma localement noethérien. Soit $X'$ un espace algébrique
localement de type fini sur $S$. Soit $T' \subset |X'|$ une partie
fermée. Soit $U' \subset X'$ le sous-espace ouvert tel que
$|U'| = |X'| \setminus T'$. Soit $W$ un espace algébrique formel
localement noethérien sur $S$, tel que $W_{red}$ soit localement de type fini
sur $S$. Enfin, soit
$$
g : X'_{/T'} \longrightarrow W
$$
une modification formelle, voir
Algébrisation des espaces formels, définition
\ref{restricted-definition-formal-modification}.
Rappelons que $X'_{/T'}$ désigne le complété formel de $X'$ le long de
$T'$, voir Espaces formels, section \ref{formal-spaces-section-completion}.
\end{situation}

\noindent
Dans la situation ci-dessus, notre but est de démontrer qu'il existe un
morphisme propre $f : X' \to X$ d'espaces algébriques sur $S$,
une partie fermée $T \subset |X|$ et un isomorphisme
$a : X_{/T} \to W$ d'espaces algébriques formels tels que
\begin{enumerate}
\item $T'$ soit l'image inverse de $T$ par $|f| : |X'| \to |X|$,
\item $f : X' \to X$ envoie isomorphiquement $U'$ sur
un sous-espace ouvert $U$ de $X$, et
\item $g = a \circ f_{/T}$, où $f_{/T} : X'_{/T'} \to X_{/T}$
est le morphisme induit.
\end{enumerate}
Nous dirons que $(f : X' \to X, T, a)$ est une {\it solution}.

\medskip\noindent
Nous suivrons la stratégie d'Artin en construisant un foncteur $F$ sur
la catégorie des schémas localement noethériens sur $S$, en montrant que $F$ est un
espace algébrique à l'aide de la
proposition \ref{artin-proposition-spaces-diagonal-representable-noetherian},
et en démontrant que le choix $X = F$ convient.

\begin{remark}
\label{artin-remark-G-rings}
En particulier, nous ne pouvons pas démontrer que le résultat souhaité est vrai pour
toute Situation \ref{artin-situation-contractions}, car nous devrons
supposer que les anneaux locaux de $S$ sont des G-anneaux. Si vous pouvez démontrer le
résultat en général ou si vous disposez d'un contre-exemple, veuillez
nous le signaler à l'adresse
\href{mailto:stacks.project@gmail.com}{stacks.project@gmail.com}.
\end{remark}

\noindent
Dans la Situation \ref{artin-situation-contractions},
soit $V$ un schéma localement noethérien sur $S$.
La valeur de notre foncteur $F$ en $V$ sera l'ensemble de tous les triplets
$$
(Z, u' : V \setminus Z \to U', \hat x : V_{/Z} \to W)
$$
satisfaisant aux conditions suivantes :
\begin{enumerate}
\item $Z \subset V$ est une partie fermée,
\item $u' : V \setminus Z \to U'$ est un morphisme sur $S$,
\item $\hat x : V_{/Z} \to W$ est un morphisme adique d'espaces algébriques
formels sur $S$,
\item $u'$ et $\hat x$ sont compatibles (voir ci-dessous).
\end{enumerate}
La condition de compatibilité est la suivante : par changement de base de la
modification formelle $g$, nous obtenons une modification formelle
$$
X'_{/T'} \times_{g, W, \hat x} V_{/Z} \longrightarrow V_{/Z}
$$
Voir Algébrisation des espaces formels, lemme
\ref{restricted-lemma-base-change-formal-modification}.
D'après le théorème principal sur les dilatations
(Algébrisation des espaces formels, théorème
\ref{restricted-theorem-dilatations}), il existe un unique
morphisme propre $V' \to V$ d'espaces algébriques, qui est un isomorphisme au-dessus de
$V \setminus Z$, tel que $V'_{/Z} \to V_{/Z}$ soit isomorphe à la
flèche affichée. Autrement dit, pour un certain morphisme
$\hat x' : V'_{/Z} \to X'_{/T'}$, nous avons un diagramme cartésien
$$
\xymatrix{
V'_{/Z} \ar[r] \ar[d]_{\hat x'} & V_{/Z} \ar[d]^{\hat x} \\
X'_{/T'} \ar[r]^g & W
}
$$
d'espaces algébriques formels. Nous considérerons
$V \setminus Z$ comme un sous-espace ouvert de $V'$ sans plus le mentionner.
La condition de compatibilité exige qu'il existe un
morphisme $x' : V' \to X'$ dont les restrictions soient $u'$ et $\hat x$
sur $V \setminus Z \subset V'$ et $V'_{/Z}$.
Autrement dit, il faut que le diagramme
$$
\xymatrix{
V \setminus Z \ar[r] \ar[d]_{u'} &
V' \ar[d]^{x'} &
V'_{/Z} \ar[l] \ar[d]^{\hat x'} \ar[r] &
V_{/Z} \ar[d]^{\hat x} \\
U' \ar[r] &
X' &
X'_{/T'} \ar[r]^g \ar[l] &
W
}
$$
soit commutatif. Observons que, d'après Algébrisation des espaces formels,
lemme \ref{restricted-lemma-faithful}, le morphisme $x'$ est unique
s'il existe. Nous indiquerons cette situation en disant
``{\it $V' \to V$, $\hat x'$ et $x'$ témoignent de la compatibilité
entre $u'$ et $\hat x$}''.

\begin{remark}
\label{artin-remark-how-to-think-compatibility}
Dans la Situation \ref{artin-situation-contractions}, soit $V$ un schéma localement noethérien
sur $S$. Soit $(Z, u', \hat x)$ un triplet satisfaisant aux conditions (1), (2) et
(3) ci-dessus. Nous voulons expliquer une manière d'envisager la condition de compatibilité
(4). Elle ne sera pas mathématiquement précise, car nous allons utiliser
une catégorie fictive $\textit{An}_S$ d'espaces analytiques sur $S$
et un foncteur d'analytification fictif
$$
\left\{
\begin{matrix}
\text{espaces algébriques formels} \\
\text{localement noethériens sur }S
\end{matrix}
\right\}
\longrightarrow
\textit{An}_S,
\quad\quad
Y \longmapsto Y^{an}
$$
Par exemple, si $Y = \text{Spf}(k[[t]])$ sur $S = \Spec(k)$, alors $Y^{an}$
doit être considéré comme un disque unité ouvert. Si $Y = \Spec(k)$, alors $Y^{an}$
est réduit à un seul point. La catégorie $\textit{An}_S$ devrait posséder des immersions
ouvertes et fermées, et nous devrions pouvoir prendre l'ouvert complémentaire
d'une partie fermée. Pour $Y$ donné, le morphisme $Y_{red} \to Y$ devrait induire une
immersion fermée $Y_{red}^{an} \to Y^{an}$. Nous posons
$Y^{rig} = Y^{an} \setminus Y_{red}^{an}$ égal à son ouvert complémentaire.
Si $Y$ est un espace algébrique et si $Z \subset Y$ est fermé, alors
le morphisme $Y_{/Z} \to Y$ devrait induire une immersion ouverte
$Y_{/Z}^{an} \to Y^{an}$, laquelle devrait à son tour induire une immersion ouverte
$$
can : (Y_{/Z})^{rig} \longrightarrow (Y \setminus Z)^{an}
$$
En outre, étant donnée une modification formelle $g : Y' \to Y$ d'espaces algébriques
formels localement noethériens, le morphisme induit $g^{rig} : (Y')^{rig} \to Y^{rig}$
devrait être un isomorphisme. Étant donnés $\text{An}_S$ et le foncteur
d'analytification, nous pouvons considérer l'exigence que le diagramme
$$
\xymatrix{
(V_{/Z})^{rig} \ar[rr]_{can} \ar[d]_{(g^{rig})^{-1} \circ \hat x^{an}} & &
(V \setminus Z)^{an} \ar[d]^{(u')^{an}} \\
(X'_{/T'})^{rig} \ar[rr]^{can} & & (X' \setminus T')^{an}
}
$$
soit commutatif. Cela a un sens puisque $g^{rig} : (X'_{T'})^{rig} \to W^{rig}$
est un isomorphisme et que $U' = X' \setminus T'$. Enfin, sous certaines hypothèses
de fidélité du foncteur d'analytification, cette exigence sera
équivalente à la condition de compatibilité formulée ci-dessus.
Nous espérons que cela incitera le lecteur à considérer la compatibilité
de $u'$ et de $\hat x$ comme l'exigence que certaines applications soient égales,
plutôt que comme la demande d'existence d'un certain diagramme commutatif.
\end{remark}

\begin{lemma}
\label{artin-lemma-functor}
Dans la Situation \ref{artin-situation-contractions}, la règle $F$ qui associe
à un schéma localement noethérien $V$ sur $S$ l'ensemble des triplets
$(Z, u', \hat x)$ satisfaisant à la condition de compatibilité et qui associe à
un morphisme $\varphi : V_2 \to V_1$ de schémas localement noethériens sur $S$
l'application
$$
F(\varphi) : F(V_1) \longrightarrow F(V_2)
$$
envoyant un élément $(Z_1, u'_1, \hat x_1)$ de $F(V_1)$ sur
$(Z_2, u'_2, \hat x_2)$ dans $F(V_2)$ défini par
\begin{enumerate}
\item $Z_2 \subset V_2$ est l'image inverse de $Z_1$ par $\varphi$,
\item $u'_2$ est le composé de $u'_1$ et de
$\varphi|_{V_2 \setminus Z_2} : V_2 \setminus Z_2 \to V_1 \setminus Z_1$,
\item $\hat x_2$ est le composé de $\hat x_1$ et de
$\varphi_{/Z_2} : V_{2, /Z_2} \to V_{1, /Z_1}$
\end{enumerate}
est un foncteur contravariant.
\end{lemma}

\begin{proof}
Pour vérifier la condition de compatibilité entre $u'_2$ et $\hat x_2$, soient
$V'_1 \to V_1$, $\hat x'_1$ et $x'_1$ témoignant de la compatibilité entre
$u'_1$ et $\hat x_1$. Posons $V'_2 = V_2 \times_{V_1} V'_1$, posons
$\hat x'_2$ égal au composé de $\hat x'_1$ et de
$V'_{2, /Z_2} \to V'_{1, /Z_1}$, et posons $x'_2$
égal au composé de $x'_1$ et de $V'_2 \to V'_1$.
Alors $V'_2 \to V_2$, $\hat x'_2$ et $x'_2$ témoignent de la compatibilité entre
$u'_2$ et $\hat x_2$. Nous omettons la vérification détaillée.
\end{proof}

\begin{lemma}
\label{artin-lemma-solution}
Dans la Situation \ref{artin-situation-contractions}, s'il existe une solution
$(f : X' \to X, T, a)$, alors il existe une bijection fonctorielle
$F(V) = \Mor_S(V, X)$ sur la catégorie des
schémas localement noethériens $V$ sur $S$.
\end{lemma}

\begin{proof}
Soit $V$ un schéma localement noethérien sur $S$.
Soit $x : V \to X$ un morphisme sur $S$.
Nous obtenons alors un élément $(Z, u', \hat x)$ de $F(V)$ comme suit :
\begin{enumerate}
\item $Z \subset V$ est l'image inverse de $T$ par $x$,
\item $u' : V \setminus Z \to U' = U$ est la restriction de
$x$ à $V \setminus Z$,
\item $\hat x : V_{/Z} \to W$ est le composé de
$x_{/Z} : V_{/Z} \to X_{/T}$ avec l'isomorphisme $a : X_{/T} \to W$.
\end{enumerate}
Ce triplet satisfait à la condition de compatibilité, car nous
pouvons prendre $V' = V \times_{x, X} X'$ et prendre pour $\hat x'$
le complété de la projection $x' : V' \to X'$.

\medskip\noindent
Réciproquement, supposons donné un élément $(Z, u', \hat x)$ de $F(V)$. Nous affirmons
qu'il existe un unique morphisme $x : V \to X$ compatible avec $u'$ et $\hat x$.
En effet, soient $V' \to V$, $\hat x'$ et $x'$ témoignant de la
compatibilité entre $u'$ et $\hat x$. Alors
Algébrisation des espaces formels, proposition
\ref{restricted-proposition-glue-modification}
est exactement le résultat dont nous avons besoin pour trouver
un unique morphisme $x : V \to X$ coïncidant avec
$\hat x$ sur $V_{/Z}$ et avec $x'$ sur $V'$ (et a fortiori
coïncidant avec $u'$ sur $V \setminus Z$).

\medskip\noindent
Nous omettons la vérification que les deux constructions ci-dessus définissent des
bijections réciproques entre leurs domaines respectifs.
\end{proof}

\begin{lemma}
\label{artin-lemma-functor-is-solution}
Dans la Situation \ref{artin-situation-contractions}, s'il existe un
espace algébrique $X$ localement de type fini sur $S$ et une
bijection fonctorielle $F(V) = \Mor_S(V, X)$ sur la catégorie des
schémas localement noethériens $V$ sur $S$, alors $X$ est une solution.
\end{lemma}

\begin{proof}
Nous devons construire un morphisme propre $f : X' \to X$, une partie fermée
$T \subset |X|$ et un isomorphisme $a : X_{/T} \to W$ ayant les propriétés
(1), (2), (3) énumérées juste après la Situation \ref{artin-situation-contractions}.

\medskip\noindent
La discussion de cette démonstration est quelque peu pédante, car nous voulons
faire coïncider soigneusement les catégories sous-jacentes. Dans ce paragraphe,
nous expliquons comment le lecteur aventureux peut procéder avec moins de timidité.
En effet, le lecteur peut étendre notre définition du foncteur $F$
à tous les espaces algébriques localement noethériens sur $S$.
Ce faisant, le lecteur peut alors conclure que $F$ et $X$ coïncident en tant que
foncteurs sur la catégorie de ces espaces algébriques, c'est-à-dire que
$X$ représente $F$. On considère alors l'objet universel
$(T, u', \hat x)$ dans $F(X)$. Le lecteur constatera alors que,
pour le triplet $X'' \to X$, $\hat x'$, $x'$ témoignant de la compatibilité
entre $u'$ et $\hat x$, le morphisme $x' : X'' \to X'$ est un isomorphisme,
ce qui produira $f : X' \to X$ en inversant $x'$. Enfin, nous disposons déjà de
$T \subset |X|$, et le lecteur peut montrer que $\hat x$ est un isomorphisme
qui peut servir de dernier ingrédient, à savoir $a$.

\medskip\noindent
Notons $h_X(-) = \Mor_S(-, X)$ le foncteur des
points de $X$ restreint à la catégorie
$(\textit{Noetherian}/S)_\etale$ de la section \ref{artin-section-noetherian}.
D'après Limites d'espaces, remarque \ref{spaces-limits-remark-limit-preserving},
les espaces algébriques $X$ et $X'$ commutent aux limites. Il en est donc de même
des restrictions $h_X$ et $h_{X'}$.
Pour construire $f$, il suffit donc de construire une
transformation $h_{X'} \to h_X = F$, voir le lemme \ref{artin-lemma-canonical-extension}.
Soit donc $V \to S$ un objet de $(\textit{Noetherian}/S)_\etale$
et soit $\tilde x : V \to X'$ un élément de $h_{X'}(V)$.
Nous obtenons alors un élément $(Z, u', \hat x)$ de $F(V)$ comme suit :
\begin{enumerate}
\item $Z \subset V$ est l'image inverse de $T'$ par $\tilde x$,
\item $u' : V \setminus Z \to U'$ est la restriction de
$\tilde x$ à $V \setminus Z$,
\item $\hat x : V_{/Z} \to W$ est le composé de
$x_{/Z} : V_{/Z} \to X'_{/T'}$ avec $g : X'_{/T'} \to W$.
\end{enumerate}
Ce triplet satisfait à la condition de compatibilité : tout d'abord, nous
obtenons gratuitement $V' \to V$ et $\hat x' : V'_{/Z'} \to X'_{/T'}$
(voir la discussion précédant le lemme \ref{artin-lemma-functor}).
Il suffit ensuite de définir $x' : V' \to X'$ comme le composé
de $V' \to V$ et du morphisme $\tilde x : V \to X'$.
Nous omettons de vérifier que cette construction convient.

\medskip\noindent
Si $\xi : V \to X$ est un morphisme \'etale, où $V$ est un schéma, alors
nous obtenons $\xi = (Z, u', \hat x) \in F(V) = h_X(V) = X(V)$.
Bien entendu, si $\varphi : V' \to V$ est encore un morphisme \'etale de schémas,
alors l'image inverse de $(Z, u', \hat x)$ dans $F(V')$ correspond à
$\xi \circ \varphi$.
La partie fermée $T \subset |X|$ est simplement définie par la condition que
le morphisme $\xi : V \to X$ associé à $\xi = (Z, u', \hat x)$
ait pour image inverse de $T$ la partie $Z$.

\medskip\noindent
Considérons un schéma noethérien $V$ sur $S$ et un morphisme
$\xi : V \to X$ correspondant à $(Z, u', \hat x)$ comme ci-dessus.
Alors $\xi(V)$ est ensemblistement contenue dans $T$
si et seulement si $V = Z$ (comme espaces topologiques). Il s'ensuit que
$X_{/T}$ coïncide avec $W$ en tant que foncteur. On obtient ainsi l'isomorphisme
$a : X_{/T} \to W$. (Nous avons omis ici un petit détail :
pour les espaces algébriques formels localement noethériens $X_{/T}$ et
$W$, il suffit de vérifier que l'on obtient un isomorphisme après évaluation
sur les schémas localement noethériens sur $S$.)

\medskip\noindent
Nous omettons la démonstration des conditions (1), (2) et (3).
\end{proof}

\begin{remark}
\label{artin-remark-diagonal}
Plaçons-nous dans la Situation \ref{artin-situation-contractions}.
Soit $V$ un schéma localement noethérien sur $S$.
Soit $(Z_i, u'_i, \hat x_i) \in F(V)$ pour $i = 1, 2$. Supposons que $V'_i \to V$,
$\hat x'_i$ et $x'_i$ témoignent de la compatibilité entre $u'_i$ et
$\hat x_i$ pour $i = 1, 2$.

\medskip\noindent
Posons $V' = V'_1 \times_V V'_2$. Notons $E' \to V'$ le noyau du couple
de morphismes
$$
V' \to V'_1 \xrightarrow{x'_1} X'
\quad\text{et}\quad
V' \to V'_2 \xrightarrow{x'_2} X'
$$
Posons $Z = Z_1 \cap Z_2$. Notons $E_W \to V_{/Z}$
le noyau du couple de morphismes
$$
V_{/Z} \to V_{/Z_1} \xrightarrow{\hat x_1} W
\quad\text{et}\quad
V_{/Z} \to V_{/Z_2} \xrightarrow{\hat x_2} W
$$
Remarquons que $E' \to V$ est séparé et localement de type fini
et que $E_W$ est un espace algébrique formel localement noethérien
séparé sur $V$.
Les compatibilités entre les divers morphismes en jeu montrent que
\begin{enumerate}
\item $\Im(E' \to V) \cap (Z_1 \cup Z_2)$
est contenue dans $Z = Z_1 \cap Z_2$,
\item le morphisme $E' \times_V (V \setminus Z) \to V \setminus Z$
est un monomorphisme et est le noyau du couple formé par les restrictions
de $u'_1$ et $u'_2$ à $V \setminus (Z_1 \cup Z_2)$,
\item le morphisme $E'_{/Z} \to V_{/Z}$ se factorise par $E_W$
et le diagramme
$$
\xymatrix{
E'_{/Z} \ar[r] \ar[d] & X'_{/T'} \ar[d]^g \\
E_W \ar[r] & W
}
$$
est cartésien. En particulier, le morphisme $E'_{/Z} \to E_W$
est une modification formelle, car il s'obtient par changement de base de $g$,
\item les données $E'$, $(E' \to V)^{-1}Z$ et $E'_{/Z} \to E_W$
forment un triplet comme dans la Situation \ref{artin-situation-contractions}
dont le schéma de base est le schéma localement noethérien $V$,
\item pour tout morphisme $\varphi : A \to V$
de schémas localement noethériens, les conditions suivantes sont équivalentes :
\begin{enumerate}
\item $(Z_1, u'_1, \hat x_1)$ et $(Z_2, u'_2, \hat x_2)$
induisent le même élément de $F(A)$,
\item $A \setminus \varphi^{-1}(Z) \to V \setminus Z$
se factorise par $E' \times_V (V \setminus Z)$
et $A_{/\varphi^{-1}(Z)} \to V_{/Z}$
se factorise par $E_W$.
\end{enumerate}
\end{enumerate}
Nous concluons, à l'aide des
lemmes \ref{artin-lemma-solution} et \ref{artin-lemma-functor-is-solution},
que, s'il existe une solution $E \to V$
pour le triplet de (4), alors $E$ représente
$F \times_{\Delta, F \times F} V$ sur la catégorie des
schémas localement noethériens sur $V$.
\end{remark}

\begin{lemma}
\label{artin-lemma-closed-immersion}
Dans la Situation \ref{artin-situation-contractions}, supposons donnée une partie
fermée $Z \subset S$ telle que
\begin{enumerate}
\item l'image inverse de $Z$ dans $X'$ soit $T'$,
\item $U' \to S \setminus Z$ soit une immersion fermée,
\item $W \to S_{/Z}$ soit une immersion fermée.
\end{enumerate}
Alors il existe une solution $(f : X' \to X, T, a)$
et, de plus, $X \to S$ est une immersion fermée.
\end{lemma}

\begin{proof}
Supposons qu'il existe un sous-schéma fermé $X \subset S$ tel que
$X \cap (S \setminus Z) = U'$ et $X_{/Z} = W$. Alors
$X$ représente le foncteur $F$ (nous omettons certains détails) et est donc
une solution. Trouver $X$ est manifestement une question locale sur $S$.
On se ramène ainsi au cas traité au paragraphe suivant.

\medskip\noindent
Supposons que $S = \Spec(A)$ soit affine. Soit $I \subset A$ l'idéal
radical définissant $Z$.
Écrivons $I = (f_1, \ldots, f_r)$. Par hypothèse, les assertions suivantes sont vraies :
\begin{enumerate}
\item l'immersion fermée $U' \to S \setminus Z$ détermine des
idéaux $J_i \subset A[1/f_i]$ tels que $J_i$ et $J_j$
engendrent le même idéal de $A[1/f_if_j]$,
\item l'immersion fermée $W \to S_{/Z}$ est l'application
$\text{Spf}(A^\wedge/J') \to \text{Spf}(A^\wedge)$ pour un certain
idéal $J' \subset A^\wedge$ du complété $I$-adique $A^\wedge$ de $A$.
\end{enumerate}
Pour achever la démonstration, il nous faut trouver un idéal $J \subset A$
tel que $J_i = J[1/f_i]$ et $J' = JA^\wedge$. D'après la
proposition \ref{more-algebra-proposition-equivalence} de Compléments d'algèbre,
il suffit de montrer que $J_i$ et $J'$ engendrent le même idéal
de $A^\wedge[1/f_i]$ pour tout $i$.

\medskip\noindent
Rappelons que $A' = H^0(X', \mathcal{O})$ est une $A$-algèbre finie
dont la formation commute au changement de base plat
(Cohomologie des espaces, lemmes
\ref{spaces-cohomology-lemma-proper-over-affine-cohomology-finite} et
\ref{spaces-cohomology-lemma-flat-base-change-cohomology}). Posons
$J'' = \Ker(A \to A')$\footnote{Contrairement à ce à quoi le lecteur pourrait
s'attendre, les idéaux $J$ et $J''$ ne coïncident pas en général.}.
L'égalité $J_i = J''A[1/f_i]$ résulte
du changement de base au spectre de $A[1/f_i]$.
Considérons le diagramme commutatif
$$
\xymatrix{
X' \ar[d] &
X'_{/T'} \times_{S_{/Z}} \text{Spf}(A^\wedge) \ar[l] \ar[d] &
X'_{/T'} \times_W \text{Spf}(A^\wedge/J') \ar@{=}[l] \ar[d] \\
\Spec(A) &
\text{Spf}(A^\wedge) \ar[l] &
\text{Spf}(A^\wedge/J') \ar[l]
}
$$
La flèche verticale médiane est le complété de la flèche verticale de gauche
le long des parties fermées évidentes. D'après le théorème des
fonctions formelles, on a
$$
(A')^\wedge = \Gamma(X' \times_S \Spec(A^\wedge), \mathcal{O}) =
\lim H^0(X' \times_S \Spec(A/I^n), \mathcal{O})
$$
Voir Cohomologie des espaces, théorème
\ref{spaces-cohomology-theorem-formal-functions}.
Le diagramme montre que $J'$ s'envoie sur zéro dans $(A')^\wedge$.
Ainsi $J' \subset J'' A^\wedge$. Considérons les flèches
$$
X'_{/T'} \to
\text{Spf}(A^\wedge/J''A^\wedge) \to
\text{Spf}(A^\wedge/J') = W
$$
Nous savons que le composé $g$ est une modification formelle
(donc, en particulier, rig-\'etale et rig-surjectif) et que la seconde
flèche est une immersion fermée (donc, en particulier, un monomorphisme adique).
Par conséquent, $X'_{/T'} \to \text{Spf}(A^\wedge/J''A^\wedge)$ est
rig-surjectif et rig-\'etale ; voir
Algébrisation des espaces formels, lemmes
\ref{restricted-lemma-rig-surjective-alternative-permanence} et
\ref{restricted-lemma-rig-etale-alternative-permanence}.
En appliquant les lemmes d'Algébrisation des espaces formels
\ref{restricted-lemma-rig-etale-descent} et
\ref{restricted-lemma-permanence-rig-surjective},
nous concluons que $\text{Spf}(A^\wedge/J''A^\wedge) \to W$
est rig-\'etale et rig-surjectif.
D'après Algébrisation des espaces formels, lemme
\ref{restricted-lemma-closed-immersion-rig-smooth-rig-surjective},
nous concluons que $I^n J'' A^\wedge \subset J'$ pour un certain $n > 0$.
Il s'ensuit que $J'' A^\wedge[1/f_i] = J' A^\wedge[1/f_i]$ et
nous en déduisons $J_i A^\wedge[1/f_i] = J' A^\wedge[1/f_i]$ pour tout
$i$, comme souhaité.
\end{proof}

\begin{lemma}
\label{artin-lemma-diagonal-contractions}
Dans la Situation \ref{artin-situation-contractions}, supposons que $X' \to S$
et $W \to S$ soient séparés. Alors la diagonale $\Delta : F \to F \times F$
est représentable par des immersions fermées.
\end{lemma}

\begin{proof}
Il suffit de combiner le lemme \ref{artin-lemma-closed-immersion}
avec la discussion de la remarque \ref{artin-remark-diagonal}.
\end{proof}

\begin{lemma}
\label{artin-lemma-sheaf}
Dans la Situation \ref{artin-situation-contractions}, le foncteur
$F$ vérifie la propriété de faisceau pour tout recouvrement \'etale
de schémas localement noethériens sur $S$.
\end{lemma}

\begin{proof}
Démonstration omise. Indication : les morphismes peuvent se définir localement pour la topologie \'etale.
\end{proof}

\begin{lemma}
\label{artin-lemma-limit-preserving}
Dans la Situation \ref{artin-situation-contractions}, le foncteur $F$ commute aux limites :
pour toute limite projective $V = \lim V_\lambda$ d'un système filtrant de schémas affines noethériens
sur $S$, nous avons $F(V) = \colim F(V_\lambda)$.
\end{lemma}

\begin{proof}
Cette démonstration est absurdement longue. Elle repose en grande partie sur des arguments classiques
sur les limites et la localisation \'etale. Nous invitons le lecteur à passer directement
à la dernière partie de la démonstration, où survient quelque chose d'intéressant.

\medskip\noindent
Soit $V = \lim_{\lambda \in \Lambda} V_i$ la limite projective d'un système filtrant de schémas
sur $S$, où $V$ et $V_\lambda$ sont noethériens et où les morphismes de transition
sont affines. Voir Limites, section \ref{limits-section-limits}, pour les résultats sur les
limites de schémas. Nous allons démontrer que $\colim F(V_\lambda) \to F(V)$
est bijective.

\medskip\noindent
Démonstration de l'injectivité : notations.
Soit $\lambda \in \Lambda$ et soient
$\xi_{\lambda, 1}, \xi_{\lambda, 2} \in F(V_\lambda)$
des éléments qui induisent le même élément de $F(V)$.
Écrivons $\xi_{\lambda, 1} = (Z_{\lambda, 1}, u'_{\lambda, 1},
\hat x_{\lambda, 1})$ et
$\xi_{\lambda, 2} = (Z_{\lambda, 2}, u'_{\lambda, 2}, \hat x_{\lambda, 2})$.

\medskip\noindent
Démonstration de l'injectivité : coïncidence des $Z_{\lambda, i}$.
Puisque $Z_{\lambda, 1}$ et $Z_{\lambda, 2}$ ont la même image inverse fermée
dans $V$, quitte à augmenter $i$, nous pouvons supposer que
$Z_{\lambda, 1} = Z_{\lambda, 2}$ ; voir
Limites, lemme \ref{limits-lemma-inverse-limit-top}, et
Topologie, lemme \ref{topology-lemma-describe-limits}.
Notons leur valeur commune $Z_\lambda \subset V_\lambda$ et, pour
$\mu \geq \lambda$,
notons $Z_\mu \subset V_\mu$ son image inverse dans $V_\mu$, et
notons $Z$ son image inverse dans $V$. Nous utiliserons plus bas l'égalité
$Z = \lim_{\mu \geq \lambda} Z_\mu$ en tant que schémas si l'on considère $Z$
et les $Z_\mu$ comme des sous-schémas fermés réduits.

\medskip\noindent
Démonstration de l'injectivité : coïncidence des $u'_{\lambda, i}$.
Puisque $U'$ est localement de type fini sur $S$ et que
les restrictions de $u'_{\lambda, 1}$ et $u'_{\lambda, 2}$
à $V \setminus Z$ coïncident, quitte à augmenter $\lambda$, nous pouvons supposer que
$u'_{\lambda, 1} = u'_{\lambda, 2}$ ; voir Limites, proposition
\ref{limits-proposition-characterize-locally-finite-presentation}.
Notons leur valeur commune $u'_\lambda$ et notons
$u'$ sa restriction à $V \setminus Z$.

\medskip\noindent
Démonstration de l'injectivité : reformulation.
À ce stade, nous avons
$\xi_{\lambda, 1} = (Z_\lambda, u'_\lambda, \hat x_{\lambda, 1})$ et
$\xi_{\lambda, 2} = (Z_\lambda, u'_\lambda, \hat x_{\lambda, 2})$.
Le principal problème de cette partie de la démonstration
consiste à montrer que les morphismes
$\hat x_{\lambda, 1}$ et $\hat x_{\lambda, 2}$ coïncident après avoir
éventuellement augmenté $\lambda$.

\medskip\noindent
Démonstration de l'injectivité : coïncidence des $\hat x_{\lambda, i}|_{Z_\lambda}$.
Considérons les morphismes
$\hat x_{\lambda, 1}|_{Z_\lambda}, \hat x_{\lambda, 2}|_{Z_\lambda} :
Z_\lambda \to W_{red}$.
Ces morphismes induisent le même morphisme $Z \to W_{red}$.
Puisque $W_{red}$ est un schéma localement de type fini sur $S$,
la proposition de Limites
\ref{limits-proposition-characterize-locally-finite-presentation}
montre que, quitte à remplacer $\lambda$ par un indice plus grand,
nous pouvons supposer que
$\hat x_{\lambda, 1}|_{Z_\lambda} = \hat x_{\lambda, 2}|_{Z_\lambda} :
Z_\lambda \to W_{red}$.

\medskip\noindent
Démonstration de l'injectivité : fin.
Nous allons maintenant appliquer la discussion de la
remarque \ref{artin-remark-diagonal} à $V_\lambda$ et aux deux éléments
$\xi_{\lambda, 1}, \xi_{\lambda, 2} \in F(V_\lambda)$.
Nous obtenons ainsi
\begin{enumerate}
\item $e_\lambda : E_\lambda' \to V_\lambda$
séparé et localement de type fini,
\item $e_\lambda^{-1}(V_\lambda \setminus Z_\lambda) \to
V_\lambda \setminus Z_\lambda$ est un isomorphisme,
\item un monomorphisme $E_{W, \lambda} \to V_{\lambda, /Z_\lambda}$
qui est l'égalisateur de $\hat x_{\lambda, 1}$ et $\hat x_{\lambda, 2}$,
\item une modification formelle $E'_{\lambda, /Z_\lambda} \to E_{W, \lambda}$
\end{enumerate}
L'assertion (2) découle de l'assertion (2) de la Remarque \ref{artin-remark-diagonal}
et du travail préparatoire effectué ci-dessus, qui donne
$u'_{\lambda, 1} = u'_{\lambda, 2} = u'_\lambda$.
Puisque $Z_\lambda = (V_{\lambda, /Z_\lambda})_{red}$ se factorise par
$E_{W, \lambda}$ car
$\hat x_{\lambda, 1}|_{Z_\lambda} = \hat x_{\lambda, 2}|_{Z_\lambda}$,
il résulte de
Espaces formels, Lemme \ref{formal-spaces-lemma-monomorphism-iso-over-red},
que $E_{W, \lambda} \to V_{\lambda, /Z_\lambda}$ est une immersion fermée.
Alors il résulte de l'assertion (4) de la Remarque \ref{artin-remark-diagonal}
et du Lemme \ref{artin-lemma-closed-immersion} appliqué au triplet
$E_\lambda'$, $e_\lambda^{-1}(Z_\lambda)$,
$E'_{\lambda, /Z_\lambda} \to E_{W, \lambda}$ au-dessus de $V_\lambda$ que
l'on obtient une immersion fermée $E_\lambda \to V_\lambda$
qui est une solution pour ce triplet.
Ensuite, nous utilisons l'assertion (5) de la Remarque \ref{artin-remark-diagonal}
qui, combinée au Lemme \ref{artin-lemma-solution},
dit que $E_\lambda$ est l'``égalisateur'' de $\xi_{\lambda, 1}$
et $\xi_{\lambda, 2}$. En particulier, nous voyons que $V \to V_\lambda$
se factorise par $E_\lambda$. Puis, en utilisant encore Limites, Proposition
\ref{limits-proposition-characterize-locally-finite-presentation},
nous trouvons une fois de plus un $\mu \geq \lambda$ tel que $V_\mu \to V_\lambda$
se factorise par $E_\lambda$ et, par conséquent, les images réciproques de
$\xi_{\lambda, 1}$ et $\xi_{\lambda, 2}$ dans $V_\mu$ sont égales,
comme souhaité.

\medskip\noindent
Démonstration de la surjectivité : énoncé.
Soit $\xi = (Z, u', \hat x)$ un élément de $F(V)$.
Il faut trouver un $\lambda \in \Lambda$ et un élément
$\xi_\lambda \in F(V_\lambda)$ dont la restriction soit $\xi$.

\medskip\noindent
Démonstration de la surjectivité : la question est locale pour la topologie \'etale.
Grâce à l'unicité établie dans la partie précédente de la démonstration et à la
propriété de faisceau de $F$ du Lemme \ref{artin-lemma-sheaf}, le problème
est local sur $V$ pour la topologie \'etale. Plus précisément, soit $v \in V$.
Nous affirmons qu'il suffit de trouver un morphisme \'etale
$(\tilde V, \tilde v) \to (V, v)$ ainsi qu'un
$\lambda$, un morphisme \'etale $\tilde V_\lambda \to V_\lambda$,
et un élément $\tilde \xi_\lambda \in F(\tilde V_\lambda)$ tels que
$\tilde V = \tilde V_\lambda \times_{V_\lambda} V$
et $\xi|_U = \tilde \xi_\lambda|_U$. Nous omettons une démonstration détaillée de
cette affirmation\footnote{Pour la démontrer,
on réunit une famille de morphismes $\tilde V \to V$
en un recouvrement \'etale fini et on montre que les morphismes correspondants
$\tilde V_\lambda \to V_\lambda$ forment également un recouvrement \'etale (après
avoir augmenté $\lambda$). Ensuite, on utilise l'injectivité pour voir que
les éléments $\tilde \xi_\lambda$ se recollent (après avoir augmenté $\lambda$),
et la propriété de faisceau de $F$ pour descendre ces
éléments en un élément de $F(V_\lambda)$.}.

\medskip\noindent
Démonstration de la surjectivité : reformulation du problème.
Rappelons que tout morphisme \'etale $(\tilde V, \tilde v) \to (V, v)$
avec $\tilde V$ affine s'obtient par changement de base à partir d'un morphisme \'etale
$\tilde V_\lambda \to V_\lambda$ avec $\tilde V_\lambda$ affine
pour un certain $\lambda$, voir par exemple
Topologies, Lemme \ref{topologies-lemma-limit-fppf-topology}.
Une fois $\tilde V_\lambda$ donné, nous avons
$\tilde V = \lim_{\mu \geq \lambda} \tilde V_\lambda \times_{V_\lambda} V_\mu$.
Ainsi, étant donné $(\tilde V, \tilde v) \to (V, v)$ \'etale avec $\tilde V$ affine,
nous pouvons remplacer $(V, v)$ par $(\tilde V, \tilde v)$ et $\xi$ par
la restriction de $\xi$ à $\tilde V$.

\medskip\noindent
Démonstration de la surjectivité : réduction au cas où la base est affine. En particulier,
supposons que $\tilde S \subset S$ soit un sous-schéma ouvert affine tel
que $v \in V$ s'envoie sur un point de $\tilde S$. Alors, conformément
au paragraphe précédent, nous pouvons remplacer $V$ par $\tilde V = \tilde S \times_S V$.
Bien sûr, si nous procédons ainsi, il suffit de résoudre le problème
pour le foncteur $F$ restreint à la catégorie des schémas localement noethériens
sur $\tilde S$. Ce foncteur est évidemment celui
associé à toute la situation après changement de base à $\tilde S$.
Nous pouvons donc supposer, et nous supposerons, que $S = \Spec(R)$ est un schéma affine noethérien
dans la suite de la démonstration.

\medskip\noindent
Démonstration de la surjectivité : cas facile.
Si $v \in V \setminus Z$, nous pouvons prendre $\tilde V = V \setminus Z$.
Celui-ci descend en un sous-schéma ouvert $\tilde V_\lambda \subset V_\lambda$
pour un certain $\lambda$ par Limites, Lemme
\ref{limits-lemma-descend-opens}.
Ensuite, après avoir augmenté $\lambda$, nous pouvons supposer
qu'il existe un morphisme $u'_\lambda : \tilde V_\lambda \to U'$
dont la restriction est $u'$. En prenant
$\tilde \xi_\lambda = (\emptyset, u'_\lambda, \emptyset)$,
on obtient l'élément souhaité de $F(\tilde V_\lambda)$.

\medskip\noindent
Démonstration de la surjectivité : cas difficile et réduction au cas où $W$ est affine.
Le cas le plus difficile est celui où l'on considère $v \in Z \subset V$.
Nous affirmons que nous pouvons nous ramener au cas où $W$ est un
schéma formel affine ; nous invitons le lecteur à sauter cet argument\footnote{La démarche
d'Artin pour démontrer ce lemme consiste à contourner cette difficulté et
il peut par conséquent éviter de démontrer d'abord l'injectivité. En effet, Artin
travaille systématiquement avec des recouvrements \'etales affines finis de tous les espaces
considérés, en gardant la trace des applications qui les relient au cours de la démonstration.
Rétrospectivement, ce serait peut-être préférable à notre méthode.}.
En effet, nous pouvons choisir un morphisme \'etale
$\tilde W \to W$ où $\tilde W$ est un espace algébrique formel affine
tel que l'image de $v$ par $\hat x : V_{/Z} \to W$
appartienne à l'image de $\tilde W \to W$ (sur les réductions).
Alors les morphismes
$$
p : \tilde W \times_{W, g} X'_{/T'} \longrightarrow X'_{/T'}
$$
et
$$
q : \tilde W \times_{W, \hat x} V_{/Z} \to V_{/Z}
$$
sont des morphismes \'etales d'espaces algébriques formels localement noethériens.
D'après Algébrisation des espaces formels, Théorème
\ref{restricted-theorem-dilatations-general} (dans un cas facile),
il existe un morphisme $\tilde X' \to X'$ d'espaces algébriques
qui est localement de type fini, est un isomorphisme au-dessus de $U'$, et
tel que $\tilde X'_{/T'} \to X'_{/T'}$ soit isomorphe à $p$.
D'après Algébrisation des espaces formels, Lemme \ref{restricted-lemma-output-etale},
le morphisme $\tilde X' \to X'$ est \'etale. Notons
$\tilde T' \subset |\tilde X'|$ l'image réciproque de $T'$.
Notons $\tilde U' \subset \tilde X'$ le sous-espace ouvert complémentaire.
Notons $\tilde g' : \tilde X'_{/\tilde T'} \to \tilde W$
la modification formelle obtenue par changement de base de $g$ par
$\tilde W \to W$. Alors
$$
\tilde X',\ \tilde T',\ \tilde U',\ \tilde W,
\ \tilde g : \tilde X'_{/\tilde T'} \to \tilde W
$$
constitue un autre exemple de la Situation \ref{artin-situation-contractions}.
Notons $\tilde F$ le foncteur construit à partir de ce triplet.
Il existe une transformation de foncteurs
$$
\tilde F \longrightarrow F
$$
construite de manière évidente à l'aide des morphismes $\tilde X' \to X'$ et
$\tilde W \to W$ ; nous omettons les détails.

\medskip\noindent
Démonstration de la surjectivité : cas difficile et réduction au cas où $W$ est affine, partie 2.
D'après le même théorème que ci-dessus, il existe un morphisme
$\tilde V \to V$ d'espaces algébriques qui est localement de type fini,
est un isomorphisme au-dessus de $V \setminus Z$ et tel que
$\tilde V_{/Z} \to V_{/Z}$ soit isomorphe à $q$.
Notons $\tilde Z \subset \tilde V$ l'image réciproque de $Z$.
D'après Algébrisation des espaces formels, Lemmes
\ref{restricted-lemma-output-etale} et \ref{restricted-lemma-output-separated},
le morphisme $\tilde V \to V$ est \'etale et séparé.
En particulier, $\tilde V$ est un schéma (localement noethérien), voir par exemple
Morphismes d'espaces, Proposition
\ref{spaces-morphisms-proposition-locally-quasi-finite-separated-over-scheme}.
Nous avons le morphisme $u'$, que nous pouvons considérer comme un morphisme
$$
\tilde u' : \tilde V \setminus \tilde Z \longrightarrow \tilde U'
$$
où $\tilde U' \subset \tilde X'$ est l'ouvert qui s'envoie isomorphiquement
sur $U'$. Nous avons un morphisme
$$
\tilde {\hat x} :
\tilde V_{/\tilde Z} = \tilde W \times_{W, \hat x} V_{/Z}
\longrightarrow
\tilde W
$$
Ici, il s'agit simplement de la projection. Nous avons donc le triplet
$$
\tilde \xi = (\tilde Z, \tilde u', \tilde {\hat x}) \in \tilde F(\tilde V)
$$
Nous omettons de démontrer la condition de compatibilité ; indications : si $V' \to V$, $\hat x'$,
et $x'$ témoignent de la compatibilité entre $u'$ et $\hat x$, alors on
pose $\tilde V' = V' \times_V \tilde V$, qui est muni de morphismes
$\tilde{\hat x}'$ et $\tilde x'$, et on montre que cela convient.
L'image de $\tilde \xi$ par la transformation $\tilde F \to F$
est la restriction de $\xi$ à $\tilde V$.

\medskip\noindent
Démonstration de la surjectivité : cas difficile et réduction au cas où $W$ est affine, partie 3.
Par notre choix de $\tilde W \to W$, il existe un ouvert affine
$\tilde V_{open} \subset \tilde V$ (nous sommes à court de notation)
dont l'image dans $V$ contient le point choisi $v \in V$.
Or, grâce au cas étudié dans le paragraphe suivant et aux remarques faites
ci-dessus, nous pouvons descendre $\tilde \xi|_{\tilde V_{open}}$
en un élément $\tilde \xi_\lambda$ de $\tilde F$ sur
$\tilde V_{\lambda, open}$
pour un certain morphisme \'etale $\tilde V_{\lambda, open} \to V_\lambda$
dont le changement de base à $V$ est $\tilde V_{open}$.
En appliquant la transformation de foncteurs $\tilde F \to F$,
nous obtenons l'élément de $F(\tilde V_{\lambda, open})$
que nous cherchions. Cela nous ramène au cas traité
dans le paragraphe suivant.

\medskip\noindent
Démonstration de la surjectivité : le cas où $W$ est affine. Nous avons $v \in Z \subset V$
et $W$ est un espace algébrique formel affine. Rappelons que
$$
\xi = (Z, u', \hat x) \in F(V)
$$
Nous pouvons encore remplacer $V$ par un voisinage \'etale de $v$.
En particulier, nous pouvons supposer, et supposons, que $V$ et $V_\lambda$ sont affines.

\medskip\noindent
Démonstration de la surjectivité : descente de $Z$. Nous pouvons trouver un $\lambda$
et un sous-schéma fermé $Z_\lambda \subset V_\lambda$ tel que
$Z$ soit obtenu par changement de base de $Z_\lambda$ à $V$. Voir
Limites, Lemme \ref{limits-lemma-descend-finite-presentation}.
Attention : nous ne savons pas (et, en général, ce ne sera pas vrai)
que $Z_\lambda$ soit un sous-schéma fermé réduit de $V_\lambda$.
Pour $\mu \geq \lambda$, notons $Z_\mu \subset V_\mu$ l'image réciproque schématique
dans $V_\mu$. Nous utiliserons plus bas le fait que
$Z = \lim_{\mu \geq \lambda} Z_\mu$ comme schémas.

\medskip\noindent
Démonstration de la surjectivité : descente de $u'$.
Puisque $U'$ est localement de type fini sur $S$,
nous pouvons supposer, après avoir augmenté $\lambda$,
qu'il existe un morphisme
$u'_\lambda : V_\lambda \setminus Z_\lambda \to U'$
dont la restriction à $V \setminus Z$ est $u'$.
Voir Limites, Proposition
\ref{limits-proposition-characterize-locally-finite-presentation}.
Pour $\mu \geq \lambda$, nous noterons $u'_\mu$ la restriction
de $u'_\lambda$ à $V_\mu \setminus Z_\mu$.

\medskip\noindent
Démonstration de la surjectivité : descente d'un témoin.
Soient $V' \to V$, $\hat x'$ et $x'$ des données témoignant de la compatibilité entre
$u'$ et $\hat x$. En utilisant les mêmes références que ci-dessus, nous pouvons supposer
(après avoir augmenté $\lambda$) qu'il existe un morphisme
$V'_\lambda \to V_\lambda$ de type fini dont le changement de base à $V$
est $V' \to V$. Après avoir augmenté $\lambda$, nous pouvons supposer que
$V'_\lambda \to V_\lambda$ est propre
(Limites, Lemme \ref{limits-lemma-eventually-proper}).
Ensuite, nous pouvons supposer que $V'_\lambda \to V_\lambda$ est un isomorphisme
au-dessus de $V_\lambda \setminus Z_\lambda$
(Limites, Lemme \ref{limits-lemma-descend-isomorphism}).
Ensuite, nous pouvons supposer qu'il existe un morphisme $x'_\lambda : V'_\lambda \to X'$
dont la restriction à $V'$ est $x'$.
En augmentant encore $\lambda$, nous pouvons supposer que $x'_\lambda$
coïncide avec $u'_\lambda$ au-dessus de $V_\lambda \setminus Z_\lambda$.
Pour $\mu \geq \lambda$, nous noterons
$V'_\mu$ et $x'_\mu$ respectivement le changement de base de $V'_\lambda$ et la
restriction de $x'_\lambda$.

\medskip\noindent
Démonstration de la surjectivité : algèbre.
Écrivons $W = \text{Spf}(B)$, $V = \Spec(A)$ et,
pour $\mu \geq \lambda$, écrivons $V_\mu = \Spec(A_\mu)$.
Notons $I_\mu \subset A_\mu$ et $I \subset A$
les idéaux définissant $Z_\mu$ et $Z$.
Alors $I_\lambda A_\mu = I_\mu$ et $I_\lambda A = I$.
Le morphisme $\hat x$ détermine, et est déterminé par, un
homomorphisme continu d'anneaux
$$
(\hat x)^\sharp : B \longrightarrow A^\wedge
$$
où $A^\wedge$ est le complété $I$-adique de $A$.
Pour achever la démonstration, il faut montrer que cette application descend en une application à valeurs dans
$A_\mu^\wedge$ pour un $\mu$ suffisamment grand, où $A_\mu^\wedge$
est le complété $I_\mu$-adique de $A_\mu$.
C'est un fait non trivial ; Artin écrit dans son article
\cite{ArtinII} : ``Puisque les données (3.5) mettent en jeu les complétés $I$-adiques,
qui ne commutent pas aux limites inductives, la vérification est assez
délicate. C'est l'analogue algébrique d'une preuve de convergence en analyse.''

\medskip\noindent
Démonstration de la surjectivité : algèbre, encore des anneaux.
Posons
$$
C_\mu = \Gamma(V'_\mu, \mathcal{O})
\quad\text{et}\quad
C = \Gamma(V', \mathcal{O})
$$
Remarquons que les homomorphismes d'anneaux $A \to C$ et $A_\mu \to C_\mu$ sont finis puisque
$V' \to V$ et $V'_\mu \to V_\mu$ sont des morphismes propres, voir
Cohomologie des espaces, Lemme
\ref{spaces-cohomology-lemma-proper-over-affine-cohomology-finite}.
Puisque $V = \lim V_\mu$ et $V' = \lim V'_\mu$,
nous avons
$$
A = \colim A_\mu
\quad\text{et}\quad
C = \colim C_\mu
$$
par Limites, Lemme \ref{limits-lemma-descend-section}\footnote{Nous ne
savons pas que $C_\mu = C_\lambda \otimes_{A_\lambda} A_\mu$ puisque les
divers morphismes ne sont pas plats.}. Pour un élément
$a \in I$, resp.\ $a \in I_\mu$, les applications $A_a \to C_a$,
resp.\ $(A_\mu)_a \to (C_\mu)_a$ sont des isomorphismes par changement de base plat
(Cohomologie des espaces, Lemme
\ref{spaces-cohomology-lemma-flat-base-change-cohomology}).
Ainsi, le noyau et le conoyau de $A \to C$ sont à support
dans $V(I)$, et de même pour $A_\mu \to C_\mu$.
Nous en concluons que le noyau et le conoyau de $A \to C$
sont annulés par une puissance de $I$, et que le noyau et le conoyau de
$A_\mu \to C_\mu$ sont annulés par une puissance de $I_\mu$, voir
Algèbre, Lemme \ref{algebra-lemma-Noetherian-power-ideal-kills-module}.

\medskip\noindent
Preuve de la surjectivité : algèbre, d'autres homomorphismes d'anneaux.
Notons $Z_n \subset V$ le $n$-ième voisinage infinitésimal
de $Z$ et notons $Z_{\mu, n} \subset V_\mu$
le $n$-ième voisinage infinitésimal de $Z_\mu$.
D'après le théorème des fonctions formelles
(Cohomologie des espaces, Théorème
\ref{spaces-cohomology-theorem-formal-functions}),
nous avons
$$
C^\wedge = \lim_n H^0(V' \times_V Z_n, \mathcal{O})
\quad\text{et}\quad
C_\mu^\wedge =
\lim_n H^0(V'_\mu \times_{V_\mu} Z_{\mu, n}, \mathcal{O})
$$
où $C^\wedge$ et $C_\mu^\wedge$ sont les complétés pour
les topologies $I$-adique et $I_\mu$-adique.
En composant le complété du morphisme
$x'_\mu : V'_\mu \to X'$ avec le morphisme $g : X'_{/T'} \to W$, nous obtenons
$$
g \circ x'_{\mu, /Z_\mu} :
V'_{\mu, /Z_\mu} = \colim V_\mu' \times_{V_\mu} Z_{\mu, n}
\longrightarrow
W
$$
et donc, d'après la description ci-dessus du complété
$C_\mu^\wedge$, nous obtenons un homomorphisme continu d'anneaux
$$
(g \circ x'_{\mu, /Z_\mu})^\sharp : B \longrightarrow C_\mu^\wedge
$$
Le fait que $V' \to V$, $\hat x'$, $x'$ réalisent la compatibilité
entre $u'$ et $\hat x$ entraîne la
commutativité du diagramme suivant :
$$
\xymatrix{
C_\mu^\wedge \ar[r] &
C^\wedge \\
B \ar[u]^{(g \circ x'_{\mu, /Z_\mu})^\sharp} \ar[r]^{(\hat x)^\sharp} &
A^\wedge \ar[u]
}
$$

\medskip\noindent
Preuve de la surjectivité : encore des arguments algébriques. Rappelons que les
$A$-modules de type fini $\Ker(A \to C)$ et $\Coker(A \to C)$ sont annulés
par une puissance de $I$, et que, de même, les $A_\mu$-modules de type fini
$\Ker(A_\mu \to C_\mu)$ et $\Coker(A_\mu \to C_\mu)$ sont annulés
par une puissance de $I_\mu$. Il s'ensuit que ces modules
coïncident avec leurs complétés. Comme le passage au complété $I$-adique dans la catégorie des
$A$-modules de type fini est exact (voir
Algèbre, section \ref{algebra-section-completion-noetherian}),
il s'ensuit que
$$
\Coker(A^\wedge \to C^\wedge) = \Coker(A \to C)
$$
et de même pour les noyaux et pour les homomorphismes $A_\mu \to C_\mu$.
Bien entendu, nous avons aussi
$$
\Ker(A \to C) = \colim \Ker(A_\mu \to C_\mu)
\quad\text{et}\quad
\Coker(A \to C) = \colim \Coker(A_\mu \to C_\mu)
$$
Rappelons que $S = \Spec(R)$ est affine. Tous les homomorphismes d'anneaux
ci-dessus sont des homomorphismes de $R$-algèbres, puisque tous les morphismes
sont des morphismes au-dessus de $S$. D'après
Algébrisation des espaces formels, Lemme
\ref{restricted-lemma-Noetherian-finite-type-red},
nous voyons que $B$ est topologiquement de type fini sur $R$.
Supposons que l'algèbre $B$ soit topologiquement engendrée par $b_1, \ldots, b_n$.
Fixons un $\mu$ (par exemple $\lambda$) et considérons
les éléments
$$
\text{images de }
(g \circ x'_{\mu, /Z_\mu})^\sharp(b_1)
, \ldots,
(g \circ x'_{\mu, /Z_\mu})^\sharp(b_n)
\text{ dans }\Coker(A_\mu \to C_\mu)
$$
Les images de ces éléments dans $\Coker(\alpha)$ sont nulles
par commutativité du carré ci-dessus. Puisque
$\Coker(A \to C) = \colim \Coker(A_\mu \to C_\mu)$ et que ces
conoyaux coïncident avec leurs complétés,
nous voyons qu'après avoir augmenté $\mu$, nous pouvons supposer que toutes ces
images sont nulles. Cela signifie que l'image de l'homomorphisme continu
$(g \circ x'_{\mu, /Z_\mu})^\sharp$ est contenue
dans $\Im(A_\mu \to C_\mu)$.
Choisissons des éléments $a_{\mu, j} \in (A_\mu)^\wedge$ envoyés sur
$(g \circ x'_{\mu, /Z_\mu})^\sharp(b_1)$ dans $(C_\mu)^\wedge$.
Alors $a_{\mu, j} \in A_\mu^\wedge$ et $(\hat x)^\sharp(b_j) \in A^\wedge$
sont envoyés sur le même élément de $C^\wedge$ par commutativité du
carré ci-dessus. Puisque
$\Ker(A \to C) = \colim \Ker(A_\mu \to C_\mu)$ et que ces noyaux
coïncident avec leurs complétés, nous pouvons, après avoir augmenté
$\mu$, modifier nos choix des $a_{\mu, j}$ de sorte que
l'image de $a_{\mu, j}$ dans $A^\wedge$ soit égale à $(\hat x)^\sharp(b_j)$.

\medskip\noindent
Preuve de la surjectivité : derniers arguments algébriques.
Soit $\mathfrak b \subset B$ l'idéal des éléments topologiquement
nilpotents. Soit $J \subset R[x_1, \ldots, x_n]$ l'idéal
constitué des $h(x_1, \ldots, x_n)$ tels que
$h(b_1, \ldots, b_n) \in \mathfrak b$. Nous obtenons alors un homomorphisme continu
surjectif de $R$-algèbres topologiques
$$
\Phi : R[x_1, \ldots, x_n]^\wedge \longrightarrow B,\quad
x_j \longmapsto b_j
$$
où le complété du membre de gauche est pris pour la topologie $J$-adique.
Comme $R[x_1, \ldots, x_n]$ est noethérien, nous pouvons choisir
des générateurs $h_1, \ldots, h_m$ de $J$. Par commutativité
du carré ci-dessus, nous voyons que $h_j(a_{\mu, 1}, \ldots, a_{\mu, n})$ est
un élément de $A_\mu^\wedge$ dont l'image dans $A^\wedge$ appartient
à $IA^\wedge$. En effet, l'homomorphisme d'anneaux $(\hat x)^\sharp$ est continu
et $IA^\wedge$ est l'idéal des éléments topologiquement nilpotents
de $A^\wedge$, car $A^\wedge/IA^\wedge = A/I$ est réduit.
(Voir Algèbre, section \ref{algebra-section-completion-noetherian}
pour les résultats sur les complétés d'anneaux noethériens.)
Puisque $A/I = \colim A_\mu/I_\mu$, nous concluons qu'après avoir augmenté
$\mu$, nous pouvons supposer que $h_j(a_{\mu, 1}, \ldots, a_{\mu, n})$ appartient à
$I_\mu A_\mu^\wedge$. En particulier, les éléments
$h_j(a_{\mu, 1}, \ldots, a_{\mu, n})$ de $A_\mu^\wedge$
sont topologiquement nilpotents dans $A_\mu^\wedge$.
Nous obtenons donc un homomorphisme continu de $R$-algèbres
$$
\Psi : R[x_1, \ldots, x_n]^\wedge \longrightarrow A_\mu^\wedge,\quad
x_j \longmapsto a_{\mu, j}
$$
Pour conclure, il nous reste à voir si $\Ker(\Phi)$ est
annulé par $\Psi$. Ce n'est peut-être pas le cas, mais nous pouvons faire en sorte
qu'il le soit après avoir augmenté $\mu$. En effet, puisque
$R[x_1, \ldots, x_n]^\wedge$ est noethérien,
nous pouvons choisir des générateurs $g_1, \ldots, g_l$ de l'idéal
$\Ker(\Phi)$. Nous voyons alors que
$$
\Psi(g_1), \ldots, \Psi(g_l) \in
\Ker(A_\mu^\wedge \to C_\mu^\wedge) = \Ker(A_\mu \to C_\mu)
$$
s'envoient sur zéro dans $\Ker(A \to C) = \colim \Ker(A_\mu \to C_\mu)$.
Ainsi, en augmentant $\mu$ comme précédemment, nous obtenons le résultat voulu.

\medskip\noindent
Preuve de la surjectivité : conclusion. L'homomorphisme continu d'anneaux
$B \to (A_\mu)^\wedge$ construit ci-dessus détermine un morphisme
$\hat x_\mu : V_{\mu, /Z_\mu} \to W$.
La compatibilité de $\hat x_\mu$ et de $u'_\mu$ résulte
du fait que l'homomorphisme d'anneaux $B \to (A_\mu)^\wedge$
est, par construction, compatible avec l'homomorphisme d'anneaux $A_\mu \to C_\mu$.
En fait, la compatibilité sera réalisée par le morphisme propre
$V'_\mu \to V_\mu$ et les morphismes
$x'_\mu$ et $\hat x'_\mu = x'_{\mu, /Z_\mu}$ utilisés dans
la construction. Ceci achève la démonstration.
\end{proof}

\begin{lemma}
\label{artin-lemma-rs}
Dans la Situation \ref{artin-situation-contractions}, le foncteur $F$ vérifie
la condition de Rim-Schlessinger (RS).
\end{lemma}

\begin{proof}
Rappelons que la condition ne fait intervenir que l'évaluation $F(V)$ du foncteur
$F$ sur les schémas $V$ au-dessus de $S$ qui sont les spectres d'anneaux locaux artiniens,
et les applications de restriction $F(V_2) \to F(V_1)$ associées aux morphismes $V_1 \to V_2$
de schémas au-dessus de $S$ qui sont les spectres d'anneaux locaux artiniens.
Soit donc $V/S$ le spectre d'un anneau local artinien.
Si $\xi = (Z, u', \hat x) \in F(V)$, alors soit $Z = \emptyset$,
soit $Z = V$ (ensemblistement). Dans le premier cas, nous voyons que
$\hat x$ est un morphisme de l'espace algébrique formel vide
vers $W$. Dans le second cas, nous voyons que $u'$ est un morphisme du
schéma vide vers $X'$ et que $\hat x : V \to W$
est un morphisme vers $W$. Nous concluons que
$$
F(V) = U'(V) \amalg W(V)
$$
et, de plus, pour $V_1 \to V_2$ comme ci-dessus, l'application induite
$F(V_2) \to F(V_1)$ est compatible avec cette décomposition.
Il suffit donc de prouver que $U'$ et $W$ vérifient tous deux la
condition de Rim-Schlessinger. Pour $U'$, cela résulte du
Lemme \ref{artin-lemma-algebraic-stack-RS}.
Pour $W$, écrivons $W = \colim W_n$ comme dans
Espaces formels, Lemme
\ref{formal-spaces-lemma-structure-locally-noetherian}.
Écrivons $V = \Spec(A)$, où $(A, \mathfrak m)$ est un anneau local artinien.
Choisissons $n \geq 1$ tel que $\mathfrak m^n = 0$. Alors nous avons
$W(V) = W_n(V)$. Ainsi, la condition de Rim-Schlessinger
pour $W$ résulte de la condition de Rim-Schlessinger pour $W_n$, pour
tout $n$ (laquelle résulte à son tour du
Lemme \ref{artin-lemma-algebraic-stack-RS}).
\end{proof}

\begin{lemma}
\label{artin-lemma-finite-dim}
Dans la Situation \ref{artin-situation-contractions}, les espaces tangents du
foncteur $F$ sont de dimension finie.
\end{lemma}

\begin{proof}
Dans la preuve du Lemme \ref{artin-lemma-rs}, nous avons vu que
$F(V) = U'(V) \amalg W(V)$ si $V$ est le spectre d'un anneau
local artinien. Les espaces tangents se calculent entièrement à partir des évaluations
de $F$ sur de tels schémas au-dessus de $S$.
Il suffit donc de prouver que les espaces tangents
des foncteurs $U'$ et $W$ sont de dimension finie.
Pour $U'$, cela résulte du
Lemme \ref{artin-lemma-finite-dimension}.
Écrivons $W = \colim W_n$ comme dans la preuve du Lemme \ref{artin-lemma-rs}.
Nous voyons alors que les espaces tangents de $W$ sont égaux aux
espaces tangents de $W_2$, car, pour calculer l'espace tangent,
il nous suffit d'évaluer $W$ sur les spectres d'anneaux locaux
artiniens $(A, \mathfrak m)$ tels que $\mathfrak m^2 = 0$.
Nous voyons alors, de nouveau, que les espaces tangents de $W_2$ sont
de dimension finie d'après le
Lemme \ref{artin-lemma-finite-dimension}.
\end{proof}

\begin{lemma}
\label{artin-lemma-formal-object-effective}
Dans la Situation \ref{artin-situation-contractions}, supposons que $X' \to S$ soit séparé.
Alors tout objet formel de $F$ est effectif.
\end{lemma}

\begin{proof}
Un objet formel $\xi = (R, \xi_n)$ de $F$ consiste en une
$S$-algèbre locale noethérienne complète $R$ dont le corps résiduel est de type fini
sur $S$, munie d'éléments $\xi_n \in F(\Spec(R/\mathfrak m^n))$
pour tout $n$, tels que $\xi_{n + 1}|_{\Spec(R/\mathfrak m^n)} = \xi_n$.
D'après la discussion dans la preuve du Lemme \ref{artin-lemma-rs},
nous voyons que $\xi$ est soit un objet formel de $U'$, soit un objet
formel de $W$. Dans le premier cas, nous voyons que $\xi$ est effectif
d'après le Lemme \ref{artin-lemma-effective}. Le second cas est le cas intéressant.
Posons $V = \Spec(R)$. Nous allons construire un élément
$(Z, u', \hat x) \in F(V)$ dont l'image dans $F(\Spec(R/\mathfrak m^n))$
est $\xi_n$ pour tout $n \geq 1$.

\medskip\noindent
Nous pouvons voir la famille d'éléments $\xi_n$ comme un morphisme
$$
\xi : \text{Spf}(R) \longrightarrow W
$$
d'espaces algébriques formels localement noethériens au-dessus de $S$. Remarquons que $\xi$
n'est {\it pas}, en général, un morphisme adique. Pour y remédier, soit $I \subset R$
l'idéal correspondant au sous-espace formel fermé
$$
\text{Spf}(R) \times_{\xi, W} W_{red} \subset \text{Spf}(R)
$$
Notons que $I \subset \mathfrak m_R$. Posons $Z = V(I) \subset V = \Spec(R)$.
Puisque $R$ est complet pour la topologie $\mathfrak m_R$-adique, il est a fortiori
complet pour la topologie $I$-adique (Algèbre, Lemme \ref{algebra-lemma-complete-by-sub}).
De plus, nous affirmons que, pour tout $n \geq 1$, le morphisme
$$
\xi|_{\text{Spf}(R/I^n)} :
\text{Spf}(R/I^n)
\longrightarrow
W
$$
provient en fait d'un morphisme
$$
\xi'_n :  \Spec(R/I^n) \longrightarrow W
$$
En effet, cela résulte de l'écriture $W = \colim W_n$ comme dans la
preuve du Lemme \ref{artin-lemma-rs}, en remarquant que $\xi|_{\text{Spf}(R/I^n)}$
se factorise par $W_n$, et en appliquant Espaces formels, Lemme
\ref{formal-spaces-lemma-map-into-algebraic-space}
pour l'algébriser en un morphisme $\Spec(R/I^n) \to W_n$
comme voulu. Notons $\text{Spf}'(R) = V_{/Z}$ le spectre formel
de $R$ muni de la topologie $I$-adique -- autrement dit, le complété formel
de $V$ le long de $Z$. À l'aide des morphismes
$\xi'_n$, nous obtenons un morphisme adique
$$
\hat x = (\xi'_n) : \text{Spf}'(R) \longrightarrow W
$$
d'espaces algébriques formels localement noethériens au-dessus de $S$.
Considérons le changement de base
$$
\text{Spf}'(R) \times_{\hat x, W, g} X'_{/T'} \longrightarrow \text{Spf}'(R)
$$
C'est une modification formelle d'après
Algébrisation des espaces formels, Lemme
\ref{restricted-lemma-base-change-formal-modification}.
Ainsi, par le théorème principal sur les dilatations
(Algébrisation des espaces formels, Théorème \ref{restricted-theorem-dilatations}),
nous obtenons un morphisme propre
$$
V' \longrightarrow V = \Spec(R)
$$
qui est un isomorphisme au-dessus de $\Spec(R) \setminus V(I)$ et
dont le complété redonne la modification formelle ci-dessus, autrement dit
$$
V' \times_{\Spec(R)} \Spec(R/I^n) =
\Spec(R/I^n) \times_{\xi'_n, W, g} X'_{/T'}
$$
Cela nous donne en particulier un système compatible de morphismes
$$
V' \times_{\Spec(R)} \Spec(R/I^n) \longrightarrow X' \times_S \Spec(R/I^n)
$$
Par conséquent, d'après le théorème d'algébrisation de Grothendieck (sous la forme de
Compléments sur les morphismes d'espaces, Lemme
\ref{spaces-more-morphisms-lemma-algebraize-morphism})
nous obtenons un morphisme
$$
x' : V' \to X'
$$
sur $S$ qui redonne les morphismes affichés ci-dessus. Finalement, en
posant $u' : V \setminus Z \to X'$ égal à la restriction de $x'$ à
$V \setminus Z \subset V'$, on obtient la troisième composante de notre
élément cherché $(Z, u', \hat x) \in F(V)$.
\end{proof}

\begin{lemma}
\label{artin-lemma-openness-smoothness}
Soit $S$ un schéma localement noethérien. Soit $V$ un schéma localement
de type fini sur $S$. Soit $Z \subset V$ un fermé. Soit $W$ un
espace algébrique formel localement noethérien sur $S$ tel que
$W_{red}$ soit localement de type fini sur $S$. Soit $g : V_{/Z} \to W$
un morphisme adique d'espaces algébriques formels sur $S$. Soit $v \in V$
un point fermé tel que $g$ soit versel en $v$ (au sens de la
section \ref{artin-section-axioms-functors}).
Alors, après avoir remplacé $V$ par un voisinage ouvert de $v$, le
morphisme $g$ est lisse (voir la démonstration).
\end{lemma}

\begin{proof}
Puisque $g$ est adique, il est représentable par des espaces algébriques (Espaces formels,
section \ref{formal-spaces-section-adic}).
Ainsi, dire que $g$ est lisse signifie que $g$ doit être lisse
au sens de
Fondements, définition \ref{bootstrap-definition-property-transformation}.

\medskip\noindent
Écrivons $W = \colim W_n$ comme dans Espaces formels, lemme
\ref{formal-spaces-lemma-structure-locally-noetherian}.
Posons $V_n = V_{/Z} \times_{\hat x, W} W_n$.
Alors $V_n$ est un sous-schéma fermé d'ensemble sous-jacent $Z$.
Le fait que $V \to W$ soit lisse équivaut à la lissité
de tous les morphismes $V_n \to W_n$ (cela vaut parce que tout morphisme
$T \to W$, où $T$ est un schéma quasi-compact, se factorise par
$W_n$ pour un certain $n$). Nous savons que le morphisme $V_n \to W_n$
est lisse en $v$ d'après le
lemme \ref{artin-lemma-base-change-versal}\footnote{Le lemme s'applique puisque
la diagonale de $W$ est représentable par des espaces algébriques et
localement de type fini, voir Espaces formels, lemme
\ref{formal-spaces-lemma-diagonal-morphism-formal-algebraic-spaces},
et nous avons vu que $W$ vérifie (RS) dans la démonstration du lemme \ref{artin-lemma-rs}.}.
Bien sûr, cela signifie que, pour tout $n$, nous pouvons restreindre $V$
de sorte que $V_n \to W_n$ soit lisse. Le problème est de trouver
un ouvert qui convienne simultanément à tous les $n$.

\medskip\noindent
La question est locale sur $V$ ; nous pouvons donc supposer que $S = \Spec(R)$ et
$V = \Spec(A)$ sont affines.

\medskip\noindent
Dans ce paragraphe, nous nous ramenons au cas où $W$ est un espace algébrique
formel affine. Choisissons un schéma formel affine $W'$ et un morphisme \'etale
$W' \to W$ tels que l'image de $v$ dans $W_{red}$ appartienne à
l'image de $W'_{red} \to W_{red}$. Alors $V_{/Z} \times_{g, W} W' \to V_{/Z}$
est un morphisme adique \'etale d'espaces algébriques formels sur $S$,
et $V_{/Z} \times_{g, W} W'$ est un espace algébrique formel affine.
D'après Algébrisation des espaces formels,
lemme \ref{restricted-lemma-algebraize-rig-etale-affine},
il existe un morphisme \'etale $\varphi : V' \to V$ de schémas affines
tel que le complété de $V'$ le long de $Z' = \varphi^{-1}(Z)$
soit isomorphe à $V_{/Z} \times_{g, W} W'$ sur $V_{/Z}$.
Notons que $v$ est l'image d'un certain $v' \in V'$.
Comme la lissité est stable par changement de base, nous voyons que
$V'_n \to W'_n$ est lisse pour tout $n$. Dans le paragraphe suivant,
nous montrons qu'après avoir remplacé $V'$ par un voisinage ouvert de $v'$,
les morphismes $V'_n \to W'_n$ sont lisses pour tout $n$.
Alors, après avoir remplacé $V$ par l'image ouverte de $V' \to V$,
nous obtenons que $V_n \to W_n$ est lisse par descente \'etale de la lissité.
Nous omettons quelques détails.

\medskip\noindent
Supposons $S = \Spec(R)$, $V = \Spec(A)$, $Z = V(I)$ et $W = \text{Spf}(B)$.
Soit $v$ le point correspondant à l'idéal maximal $I \subset \mathfrak m \subset A$.
On se donne un homomorphisme continu adique de $R$-algèbres
$$
B \longrightarrow A^\wedge
$$
Soit $\mathfrak b \subset B$ l'idéal des éléments topologiquement nilpotents
(c'est le plus grand idéal de définition de l'anneau topologique adique
noethérien $B$). Notons que $\mathfrak b A^\wedge$ et
$IA^\wedge$ sont tous deux des idéaux de définition de l'anneau adique
noethérien $A^\wedge$. En outre, $\mathfrak m A^\wedge$ est un idéal maximal
de $A^\wedge$ contenant à la fois $\mathfrak b A^\wedge$ et $IA^\wedge$.
On sait que
$$
B_n = B/\mathfrak b^n \to A^\wedge/\mathfrak b^n A^\wedge = A_n
$$
est lisse en $\mathfrak m$ pour tout $n$. D'après la discussion précédente,
nous pouvons supposer et supposons que $B_1 \to A_1$ est un homomorphisme lisse.
Notons $\mathfrak m_1 \subset A_1$ l'idéal maximal correspondant
à $\mathfrak m$. Puisque la lissité implique la platitude, nous voyons que :
pour tout $n \geq 1$, l'application
$$
\mathfrak b^n/\mathfrak b^{n + 1} \otimes_{B_1} (A_1)_{\mathfrak m_1}
\longrightarrow
\left(\mathfrak b^nA^\wedge/\mathfrak b^{n + 1}A^\wedge\right)_{\mathfrak m_1}
$$
est un isomorphisme (voir
Algèbre, lemme \ref{algebra-lemma-what-does-it-mean-again}).
Considérons l'algèbre de Rees
$$
B' = \bigoplus\nolimits_{n \geq 0} \mathfrak b^n/\mathfrak b^{n + 1}
$$
qui est une algèbre graduée de type fini sur l'anneau noethérien $B_1$, et
l'algèbre de Rees
$$
A' = \bigoplus\nolimits_{n \geq 0}
\mathfrak b^nA^\wedge/\mathfrak b^{n + 1}A^\wedge
$$
qui est une algèbre graduée de type fini sur l'anneau noethérien $A_1$.
Considérons l'homomorphisme d'$A_1$-algèbres graduées
$$
\Psi : B' \otimes_{B_1} A_1 \longrightarrow A'
$$
D'après ce qui précède, cette application devient un isomorphisme après localisation en
l'idéal maximal $\mathfrak m_1$ de $A_1$.
Ainsi, $\Ker(\Psi)$, resp.\ $\Coker(\Psi)$, est un module de type fini
sur $B' \otimes_{B_1} A_1$, resp.\ $A'$, dont le localisé
en $\mathfrak m_1$ est nul. Il s'ensuit qu'après avoir remplacé
$A_1$ (et, de manière correspondante, $A$) par une localisation principale,
nous pouvons supposer que $\Psi$ est un isomorphisme. (C'est l'étape clé de la démonstration.)
En remontant alors le raisonnement, nous voyons que $B_n \to A_n$ est plat, voir
Algèbre, lemme \ref{algebra-lemma-what-does-it-mean-again}.
Par suite, $A_n \to B_n$ est lisse (comme homomorphisme plat à
fibres lisses, voir Algèbre, lemme \ref{algebra-lemma-flat-fibre-smooth}),
ce qui achève la démonstration.
\end{proof}

\begin{lemma}
\label{artin-lemma-openness-versality}
Dans la situation \ref{artin-situation-contractions}, le foncteur
$F$ satisfait à l'ouverture de la versalité.
\end{lemma}

\begin{proof}
Nous devons montrer ce qui suit. Étant donnés un schéma $V$ localement de type fini sur
$S$, un élément $\xi \in F(V)$ et un point de type fini $v_0 \in V$ tel que
$\xi$ soit versel en $v_0$, après avoir remplacé $V$ par un voisinage ouvert
de $v_0$, l'élément $\xi$ est versel en tout point de type fini de $V$.
Écrivons $\xi = (Z, u', \hat x)$.

\medskip\noindent
Premier cas : $v_0 \not \in Z$. Nous pouvons d'abord remplacer $V$ par
$V \setminus Z$. Nous voyons ainsi que $\xi = (\emptyset, u', \emptyset)$
et que le morphisme $u' : V \to X'$ est versel en $v_0$.
D'après Compléments sur les morphismes d'espaces, lemme
\ref{spaces-more-morphisms-lemma-lifting-along-artinian-at-point},
cela signifie que $u' : V \to X'$ est lisse en $v_0$.
Puisque l'ensemble des points où un morphisme est lisse est ouvert,
nous pouvons, après avoir restreint $V$, supposer que $u'$ est lisse.
Le même lemme nous dit alors que $\xi$ est versel en tout
point, comme souhaité.

\medskip\noindent
Second cas : $v_0 \in Z$. Écrivons $W = \colim W_n$ comme dans
Espaces formels, lemme \ref{formal-spaces-lemma-structure-locally-noetherian}.
D'après le lemme \ref{artin-lemma-openness-smoothness}, nous pouvons supposer que $\hat x : V_{/Z} \to W$
est un morphisme lisse d'espaces algébriques formels. Il en résulte immédiatement
que $\xi = (Z, u', \hat x)$ est versel en tout point de type fini de $Z$.
Soient $V' \to V$, $\hat x'$ et $x'$ réalisant la compatibilité entre $u'$
et $\hat x$. Nous voyons que $\hat x' : V'_{/Z} \to X'_{/T'}$ est lisse en tant que
changement de base de $\hat x$. Puisque $\hat x'$ est le complété de
$x' : V' \to X'$, cela entraîne que $x' : V' \to X'$ est lisse en tous les
points de $(V' \to V)^{-1}(Z) = |x'|^{-1}(T') \subset |V'|$
d'après Compléments sur les morphismes d'espaces, lemme déjà utilisé
\ref{spaces-more-morphisms-lemma-lifting-along-artinian-at-point}.
Puisque l'ensemble des points de lissité d'un morphisme est ouvert, nous voyons que
l'ensemble fermé des points $B \subset |V'|$ où $x'$ n'est pas lisse
ne rencontre pas $(V' \to V)^{-1}(Z)$. Comme $V' \to V$ est propre et
donc fermé, nous voyons que $(V' \to V)(B) \subset V$ est un sous-ensemble fermé
qui ne rencontre pas $Z$. Après avoir restreint $V$, nous pouvons donc supposer que
$B = \emptyset$, c'est-à-dire que $x'$ est lisse. D'après la discussion du paragraphe
précédent, cela signifie exactement que $\xi$ est versel en tout point de type fini
de $V$ qui n'est pas dans $Z$, ce qui achève la démonstration.
\end{proof}

\noindent
Voici le résultat final.

\begin{theorem}
\label{artin-theorem-contractions}
\begin{reference}
\cite[théorème 3.1]{ArtinII}
\end{reference}
Soit $S$ un schéma localement noethérien tel que $\mathcal{O}_{S, s}$
soit un G-anneau pour tout point de type fini $s \in S$. Soit $X'$ un espace algébrique
localement de type fini sur $S$. Soit $T' \subset |X'|$ un sous-ensemble fermé.
Soit $W$ un espace algébrique formel localement noethérien sur $S$,
où $W_{red}$ est localement de type fini sur $S$. Enfin, soit
$$
g : X'_{/T'} \longrightarrow W
$$
une modification formelle, voir Algébrisation des espaces formels, définition
\ref{restricted-definition-formal-modification}. Si $X'$ et $W$ sont
séparés\footnote{Voir la remarque \ref{artin-remark-separated-needed}.} sur $S$, alors
il existe un morphisme propre $f : X' \to X$ d'espaces algébriques sur $S$,
un sous-ensemble fermé $T \subset |X|$ et un isomorphisme $a : X_{/T} \to W$
d'espaces algébriques formels tels que
\begin{enumerate}
\item $T'$ soit l'image réciproque de $T$ par $|f| : |X'| \to |X|$,
\item $f : X' \to X$ induise un isomorphisme de $X' \setminus T'$ sur
$X \setminus T$, et
\item $g = a \circ f_{/T}$, où $f_{/T} : X'_{/T'} \to X_{/T}$
est le morphisme induit.
\end{enumerate}
Autrement dit, $(f : X' \to X, T, a)$ est une solution au sens défini plus haut dans
cette section.
\end{theorem}

\begin{proof}
Soit $F$ le foncteur construit dans cette section à partir de $X'$, $T'$, $W$ et $g$.
D'après le lemme \ref{artin-lemma-functor-is-solution}, il suffit de montrer que
$F$ est représenté par un espace algébrique $X$ localement de type fini sur $S$.
À cette fin, nous allons appliquer la
proposition \ref{artin-proposition-spaces-diagonal-representable-noetherian}.
En effet, d'après le lemme \ref{artin-lemma-diagonal-contractions},
la diagonale de $F$ est représentable par des immersions fermées
et, d'après les
lemmes \ref{artin-lemma-sheaf}, \ref{artin-lemma-limit-preserving},
\ref{artin-lemma-rs}, \ref{artin-lemma-finite-dim},
\ref{artin-lemma-formal-object-effective} et \ref{artin-lemma-openness-versality},
les axiomes [0], [1], [2], [3], [4] et [5] sont satisfaits.
\end{proof}

\begin{remark}
\label{artin-remark-separated-needed}
La démonstration du théorème \ref{artin-theorem-contractions} utilise en deux endroits le fait que $X'$ et $W$
sont séparés sur $S$. Premièrement, elle l'utilise pour montrer que
$\Delta : F \to F \times F$ est représentable par des espaces algébriques.
On peut éviter entièrement cet emploi de l'hypothèse en prouvant
que $\Delta$ est représentable par application du théorème dans le cas
séparé aux triplets
$E'$, $(E' \to V)^{-1}Z$ et $E'_{/Z} \to E_W$
trouvés dans la remarque \ref{artin-remark-diagonal} (c'est le procédé d'amorçage
usuel pour la diagonale). Ainsi, la démonstration du
lemme \ref{artin-lemma-formal-object-effective} est le seul
endroit de notre démonstration du théorème \ref{artin-theorem-contractions}
où nous ayons réellement besoin d'utiliser le fait que $X' \to S$ est séparé.
Le lecteur vérifiera que nous n'utilisons l'hypothèse que pour obtenir
le morphisme $x' : V' \to X'$. L'existence de $x'$ peut être établie,
à l'aide de résultats de la littérature, si $X' \to S$ est quasi-séparé, voir
Compléments sur les morphismes d'espaces, remarque
\ref{spaces-more-morphisms-remark-weaken-separation-axioms-question}.
Nous concluons que le théorème reste vrai tel quel lorsque l'on remplace
``séparé'' par ``quasi-séparé''. Si nous en avons un jour besoin,
nous l'énoncerons précisément et le démontrerons ici avec soin.
\end{remark}



















% OK, commençons ici.
%

\chapter{Espaces de Quot et de Hilbert}


\phantomsection
\label{quot-section-phantom}





\section{Introduction}
\label{quot-section-introduction}

\noindent
Tel qu'il avait été conçu initialement, ce chapitre devait traiter des
foncteurs de Quot et de Hilbert et démontrer que ce sont des espaces
algébriques sous réserve que certaines conditions techniques soient
satisfaites. Dans le cadre des schémas, ce sujet est traité dans les exposés
de Grothendieck au s\'eminaire Bourbaki ; voir
\cite{Gr-I},
\cite{Gr-II},
\cite{Gr-III},
\cite{Gr-IV},
\cite{Gr-V} et
\cite{Gr-VI}. Pour les schémas projectifs, les schémas de Quot et de Hilbert
se plongent dans des grassmanniennes d'espaces de sections de faisceaux
inversibles très amples convenables, ce qui fournit une méthode de
construction de ces schémas. Notre approche est différente : nous utilisons
les axiomes d'Artin pour démontrer que Quot et Hilb sont des espaces
algébriques.

\medskip\noindent
Après plus ample réflexion, il s'est révélé plus commode, pour le
développement de la théorie dans le projet Champs, de commencer par le champ
$\Cohstack_{X/B}$ des faisceaux cohérents (à support propre sur la base),
introduit dans \cite{lieblich_remarks}. Pour nous, $f : X \to B$ est un
morphisme d'espaces algébriques satisfaisant des conditions techniques
convenables, bien que ceci puisse être généralisé (voir ci-dessous). Étant
donnés des modules $\mathcal{F}$ et $\mathcal{G}$ sur $X$, sous des
hypothèses convenables, le foncteur
$T/B \mapsto \Hom_{X_T}(\mathcal{F}_T, \mathcal{G}_T)$ est un espace
algébrique $\mathit{Hom}(\mathcal{F}, \mathcal{G})$ sur $B$. Voir la
section \ref{quot-section-hom}. On montre à la section \ref{quot-section-isom} que le
sous-foncteur $\mathit{Isom}(\mathcal{F}, \mathcal{G})$ des isomorphismes est
un espace algébrique. Ceci est utilisé dans les sections suivantes pour
montrer que la diagonale du champ $\Cohstack_{X/B}$ est représentable. Nous
démontrons que $\Cohstack_{X/B}$ est un champ algébrique dans la section
\ref{quot-section-stack-coherent-sheaves} lorsque $X \to B$ est plat, et dans la
section \ref{quot-section-not-flat} en général. On trouvera dans l'introduction de
cette section des indications bibliographiques.

\medskip\noindent
Cela étant établi, il est assez immédiat de démontrer que
$\Quotfunctor_{\mathcal{F}/X/B}$, $\Hilbfunctor_{X/B}$ et
$\Picardfunctor_{X/B}$ sont des espaces algébriques et que
$\Picardstack_{X/B}$ est un champ algébrique. Voir les sections
\ref{quot-section-quot}, \ref{quot-section-hilb},
\ref{quot-section-picard-functor} et \ref{quot-section-picard-stack}.

\medskip\noindent
De la manière habituelle, nous en déduisons à la section
\ref{quot-section-relative-morphisms} que le foncteur
$\mathit{Mor}_B(Z, X)$ des morphismes relatifs est un espace algébrique (sous
des hypothèses convenables).

\medskip\noindent
Dans la section \ref{quot-section-stack-of-spaces}, nous démontrons que le champ en
groupoïdes
$$
\Spacesstack'_{fp, flat, proper}
$$
qui paramètre les familles plates d'espaces algébriques propres satisfait à
tous les axiomes d'Artin (y compris l'ouverture de la versalité), sauf à
l'effectivité formelle. Nous avons choisi à dessein cette notation fort
malcommode pour ce champ, car le lecteur doit user de ses propriétés avec
prudence.

\medskip\noindent
Dans la section \ref{quot-section-polarized}, nous démontrons que le champ
$\Polarizedstack$ qui paramètre les familles plates d'espaces algébriques
propres polarisés est un champ algébrique. Grâce à notre étude des familles
plates d'espaces algébriques propres, il suffit pour cela de démontrer
l'effectivité formelle pour les schémas polarisés, résultat souvent connu
sous le nom de théorème d'algébrisation de Grothendieck.

\medskip\noindent
Dans la section \ref{quot-section-curves}, nous démontrons que le champ
$\Curvesstack$ qui paramètre les familles de courbes est algébrique.

\medskip\noindent
Dans la section \ref{quot-section-moduli-complexes}, nous étudions les modules de
complexes sur un morphisme propre et obtenons un champ algébrique
$\Complexesstack_{X/B}$. L'idée de l'énoncé et celle de la démonstration sont
empruntées à \cite{lieblich-complexes}.

\medskip\noindent
Que ne trouve-t-on pas dans ce chapitre ? Il n'y est presque pas question des
propriétés que possèdent les espaces et champs de modules ainsi obtenus (au-delà
de leur algébricité) ; à ce sujet, nous renvoyons à Champs de modules,
section \ref{moduli-section-introduction}. Dans la plupart des résultats
examinés, on peut généraliser les constructions en considérant un morphisme
$\mathcal{X} \to \mathcal{B}$ de champs algébriques au lieu d'un morphisme
$X \to B$ d'espaces algébriques. Nous en traiterons (insérer ici une référence
future). Dans le cas des espaces de Hilbert, il existe une notion plus générale
de « champs de Hilbert », dont nous traiterons dans un chapitre séparé ; voir
(insérer ici une référence future).




\section{Conventions}
\label{quot-section-conventions}

\noindent
Nous avons placé à dessein ce chapitre, ainsi que les chapitres « Exemples de
champs », « Faisceaux sur les champs algébriques », « Critères de
représentabilité » et « Axiomes d'Artin », avant le développement général de
la théorie des champs algébriques. La raison en est qu'à partir du chapitre
suivant (voir Propriétés des champs, section
\ref{stacks-properties-section-conventions}), nous ne distinguerons plus un
schéma du champ algébrique auquel il donne naissance. Notre langage deviendra
donc plus souple et plus facile à lire pour un être humain, mais aussi moins
précis. Ces premiers chapitres, y compris le chapitre initial « Champs
algébriques », posent les fondements qui nous permettront plus tard de négliger
certaines distinctions très techniques entre diverses façons d'envisager les
champs algébriques. Cependant, notamment dans les chapitres « Axiomes d'Artin »
et « Critères de représentabilité », nous devons préciser exactement les objets
avec lesquels nous travaillons, puisque nous cherchons à démontrer que
certaines constructions produisent des champs algébriques ou des espaces
algébriques.

\medskip\noindent
Malheureusement, il en résulte que certaines notations, conventions et certains
termes sont malcommodes et peuvent paraître contraires aux usages au lecteur
plus expérimenté. Nous espérons qu'il nous le pardonnera !

\medskip\noindent
Nous supposerons constamment que tous les schémas sont contenus dans un grand
site fppf $\Sch_{fppf}$. En outre, tout anneau $A$ considéré possède la
propriété que $\Spec(A)$ est (isomorphe à) un objet de ce grand site.

\medskip\noindent
Soit $S$ un schéma et soit $X$ un espace algébrique sur $S$. Dans ce chapitre
et dans le suivant, nous écrirons $X \times_S X$ pour le produit de $X$ par
lui-même (dans la catégorie des espaces algébriques sur $S$), plutôt que
$X \times X$.














\section{Le foncteur Hom}
\label{quot-section-hom}

\noindent
Dans cette section, nous étudions le foncteur des homomorphismes défini
ci-dessous.

\begin{situation}
\label{quot-situation-hom}
Soit $S$ un schéma. Soit $f : X \to B$ un morphisme d'espaces algébriques sur
$S$. Soient $\mathcal{F}$ et $\mathcal{G}$ des $\mathcal{O}_X$-modules
quasi-cohérents. Pour tout schéma $T$ sur $B$, nous noterons $\mathcal{F}_T$
et $\mathcal{G}_T$ les changements de base de $\mathcal{F}$ et
$\mathcal{G}$ à $T$, autrement dit leurs images inverses par le morphisme de
projection $X_T = X \times_B T \to X$. Nous considérons le foncteur
\begin{equation}
\label{quot-equation-hom}
\mathit{Hom}(\mathcal{F}, \mathcal{G}) :
(\Sch/B)^{opp}
\longrightarrow
\textit{Ensembles},\quad
T
\longrightarrow
\Hom_{\mathcal{O}_{X_T}}(\mathcal{F}_T, \mathcal{G}_T)
\end{equation}
\end{situation}

\noindent
Dans la situation \ref{quot-situation-hom}, nous regarderons parfois
$\mathit{Hom}(\mathcal{F}, \mathcal{G})$ comme un foncteur
$(\Sch/S)^{opp} \to \textit{Ensembles}$ muni d'un morphisme
$\mathit{Hom}(\mathcal{F}, \mathcal{G}) \to B$. En effet, si $T$ est un
schéma sur $S$, un élément de
$\mathit{Hom}(\mathcal{F}, \mathcal{G})(T)$ consiste en un couple $(h, u)$,
où $h$ est un morphisme $h : T \to B$ et
$u : \mathcal{F}_T \to \mathcal{G}_T$ est un homomorphisme de
$\mathcal{O}_{X_T}$-modules, avec $X_T = T \times_{h, B} X$ et
$\mathcal{F}_T$, $\mathcal{G}_T$ les images inverses sur $X_T$. En
particulier, lorsque nous disons que
$\mathit{Hom}(\mathcal{F}, \mathcal{G})$ est un espace algébrique, nous
entendons que le foncteur correspondant
$(\Sch/S)^{opp} \to \textit{Ensembles}$ est un espace algébrique.

\begin{lemma}
\label{quot-lemma-hom-sheaf}
Dans la situation \ref{quot-situation-hom}, le foncteur
$\mathit{Hom}(\mathcal{F}, \mathcal{G})$
satisfait à la propriété de faisceau pour la topologie fpqc.
\end{lemma}

\begin{proof}
Soit $\{T_i \to T\}_{i \in I}$ un recouvrement fpqc de schémas sur $B$.
Posons $X_i = X_{T_i} = X \times_S T_i$ et $\mathcal{F}_i = u_{T_i}$,
ainsi que $\mathcal{G}_i = \mathcal{G}_{T_i}$. Remarquons que
$\{X_i \to X_T\}_{i \in I}$ est un recouvrement fpqc de $X_T$ ; voir
Topologies sur les espaces, lemme
\ref{spaces-topologies-lemma-fpqc}. Ainsi, une famille d'homomorphismes
$u_i : \mathcal{F}_i \to \mathcal{G}_i$ tels que $u_i$ et $u_j$
induisent le même homomorphisme sur $X_{T_i \times_T T_j}$ provient par
descente d'un unique homomorphisme
$u : \mathcal{F}_T \to \mathcal{G}_T$
(Descente sur les espaces, proposition
\ref{spaces-descent-proposition-fpqc-descent-quasi-coherent}).
\end{proof}

\begin{lemma}
\label{quot-lemma-extend-hom-to-spaces}
Dans la situation \ref{quot-situation-hom}. Soit $T$ un espace algébrique sur $S$.
Nous avons
$$
\Mor_{\Sh((\Sch/S)_{fppf})}(T, \mathit{Hom}(\mathcal{F}, \mathcal{G})) =
\{(h, u) \mid h : T \to B, u : \mathcal{F}_T \to \mathcal{G}_T\}
$$
où $\mathcal{F}_T, \mathcal{G}_T$ désignent les images inverses de
$\mathcal{F}$ et $\mathcal{G}$ sur l'espace algébrique
$X \times_{B, h} T$.
\end{lemma}

\begin{proof}
Choisissons un schéma $U$ et un morphisme \'etale surjectif $p : U \to T$.
Soit $R = U \times_T U$, de projections $t, s : R \to U$.

\medskip\noindent
Soit $v : T \to \mathit{Hom}(\mathcal{F}, \mathcal{G})$ une transformation
naturelle. Alors $v(p)$ correspond à un couple $(h_U, u_U)$ sur $U$. Comme
$v$ est une transformation de foncteurs, les images inverses de
$(h_U, u_U)$ par $s$ et $t$ coïncident. Puisque $T = U/R$ (Espaces, lemme
\ref{spaces-lemma-space-presentation}), nous obtenons un morphisme
$h : T \to B$ tel que $h_U = h \circ p$. Alors $\mathcal{F}_U$ est l'image
inverse de $\mathcal{F}_T$ sur $X_U$, et il en est de même pour
$\mathcal{G}_U$. Par conséquent, $u_U$ descend en un homomorphisme de
$\mathcal{O}_{X_T}$-modules $u : \mathcal{F}_T \to \mathcal{G}_T$, d'après
Descente sur les espaces, proposition
\ref{spaces-descent-proposition-fpqc-descent-quasi-coherent}.

\medskip\noindent
Réciproquement, soit $(h, u)$ un couple sur $T$. Nous obtenons alors une
transformation naturelle $v : T \to \mathit{Hom}(\mathcal{F}, \mathcal{G})$
en envoyant un morphisme $a : T' \to T$, où $T'$ est un schéma, sur
$(h \circ a, a^*u)$. Nous omettons de vérifier que cette construction et celle
du paragraphe précédent sont inverses l'une de l'autre.
\end{proof}

\begin{remark}
\label{quot-remark-hom-base-change}
Dans la situation \ref{quot-situation-hom}, soit $B' \to B$ un morphisme d'espaces
algébriques sur $S$. Posons $X' = X \times_B B'$ et désignons par
$\mathcal{F}'$, $\mathcal{G}'$ les images inverses de $\mathcal{F}$,
$\mathcal{G}$ sur $X'$. Nous obtenons alors un foncteur
$\mathit{Hom}(\mathcal{F}', \mathcal{G}') : (\Sch/B')^{opp} \to \textit{Ensembles}$
associé au changement de base $f' : X' \to B'$. Pour un schéma $T$ sur $B'$,
il est clair que
$$
\mathit{Hom}(\mathcal{F}', \mathcal{G}')(T) =
\mathit{Hom}(\mathcal{F}, \mathcal{G})(T)
$$
où, dans le membre de droite, on regarde $T$ comme un schéma sur $B$ au moyen
du composé $T \to B' \to B$. Cette remarque élémentaire sera parfois utile
pour changer l'espace algébrique de base.
\end{remark}

\begin{lemma}
\label{quot-lemma-hom-sheaf-in-X}
Dans la situation \ref{quot-situation-hom}, soit $\{X_i \to X\}_{i \in I}$ un
recouvrement fppf et, pour chaque $i, j \in I$, soit
$\{X_{ijk} \to X_i \times_X X_j\}$ un recouvrement fppf. Désignons par
$\mathcal{F}_i$, resp.\ $\mathcal{F}_{ijk}$, l'image inverse de
$\mathcal{F}$ sur $X_i$, resp.\ sur $X_{ijk}$. Définissons de même
$\mathcal{G}_i$ et $\mathcal{G}_{ijk}$. Pour tout schéma $T$ sur $B$, le
diagramme
$$
\xymatrix{
\mathit{Hom}(\mathcal{F}, \mathcal{G})(T) \ar[r] &
\prod\nolimits_i
\mathit{Hom}(\mathcal{F}_i, \mathcal{G}_i)(T)
\ar@<1ex>[r]^-{\text{pr}_0^*} \ar@<-1ex>[r]_-{\text{pr}_1^*}
&
\prod\nolimits_{i, j, k}
\mathit{Hom}(\mathcal{F}_{ijk}, \mathcal{G}_{ijk})(T)
}
$$
présente la première flèche comme l'égalisateur des deux autres.
\end{lemma}

\begin{proof}
Soit $u_i : \mathcal{F}_{i, T} \to \mathcal{G}_{i, T}$ un élément de
l'égalisateur de $\text{pr}_0^*$ et $\text{pr}_1^*$. Comme le changement de
base d'un recouvrement fppf est un recouvrement fppf (Topologies sur les
espaces, lemme \ref{spaces-topologies-lemma-fppf}), nous voyons que
$\{X_{i, T} \to X_T\}_{i \in I}$ et
$\{X_{ijk, T} \to X_{i, T} \times_{X_T} X_{j, T}\}$ sont des recouvrements
fppf. En appliquant Descente sur les espaces, proposition
\ref{spaces-descent-proposition-fpqc-descent-quasi-coherent}, nous concluons
d'abord que $u_i$ et $u_j$ induisent le même morphisme sur
$X_{i, T} \times_{X_T} X_{j, T}$ ; une seconde application montre alors qu'il
existe un unique morphisme $u : \mathcal{F}_T \to \mathcal{G}_T$ dont la
restriction est $u_i$ pour tout $i$. Ceci achève la démonstration.
\end{proof}

\begin{lemma}
\label{quot-lemma-hom-limits}
Dans la situation \ref{quot-situation-hom}. Si $\mathcal{F}$ est de présentation
finie et si $f$ est quasi-compact et quasi-séparé, alors
$\mathit{Hom}(\mathcal{F}, \mathcal{G})$ commute aux limites.
\end{lemma}

\begin{proof}
Soit $T = \lim_{i \in I} T_i$ une limite filtrante de $B$-schémas affines.
Nous devons montrer que
$$
\mathit{Hom}(\mathcal{F}, \mathcal{G})(T) =
\colim \mathit{Hom}(\mathcal{F}, \mathcal{G})(T_i)
$$
Choisissons $0 \in I$. Nous pouvons remplacer $B$ par $T_0$, $X$ par
$X_{T_0}$, $\mathcal{F}$ par $\mathcal{F}_{T_0}$, $\mathcal{G}$ par
$\mathcal{G}_{T_0}$ et $I$ par $\{i \in I \mid i \geq 0\}$. Voir la
remarque \ref{quot-remark-hom-base-change}. Nous pouvons donc supposer que
$B = \Spec(R)$ est affine.

\medskip\noindent
Lorsque $B$ est affine, $X$ est quasi-compact et quasi-séparé. Choisissons un
morphisme \'etale surjectif $U \to X$, où $U$ est un schéma affine (Propriétés
des espaces, lemme
\ref{spaces-properties-lemma-quasi-compact-affine-cover}). Comme $X$ est
quasi-séparé, le schéma $U \times_X U$ est quasi-compact, et nous pouvons
choisir un morphisme \'etale surjectif $V \to U \times_X U$, où $V$ est un
schéma affine. En appliquant le lemme \ref{quot-lemma-hom-sheaf-in-X}, nous voyons
que $\mathit{Hom}(\mathcal{F}, \mathcal{G})$ est l'égalisateur de deux
applications entre
$$
\mathit{Hom}(\mathcal{F}|_U, \mathcal{G}|_U)
\quad\text{et}\quad
\mathit{Hom}(\mathcal{F}|_V, \mathcal{G}|_V)
$$
Nous sommes ainsi ramenés au cas où $X$ est affine.

\medskip\noindent
Dans le cas affine, l'énoncé du lemme se ramène au problème suivant : étant
donnés un homomorphisme d'anneaux $R \to A$, deux $A$-modules $M$, $N$ et un
système inductif de $R$-algèbres $C = \colim C_i$. Quand l'application
$$
\colim \Hom_{A \otimes_R C_i}(M \otimes_R C_i, N \otimes_R C_i)
\longrightarrow
\Hom_{A \otimes_R C}(M \otimes_R C, N \otimes_R C)
$$
est-elle bijective ? D'après Algèbre, lemme
\ref{algebra-lemma-module-map-property-in-colimit}, elle l'est si
$M \otimes_R C$ est de présentation finie sur $A \otimes_R C$, c'est-à-dire
si $M$ est de présentation finie sur $A$.
\end{proof}

\begin{lemma}
\label{quot-lemma-hom-closed}
Soit $S$ un schéma. Soit $B$ un espace algébrique sur $S$. Soit
$i : X' \to X$ une immersion fermée d'espaces algébriques sur $B$. Soit
$\mathcal{F}$ un $\mathcal{O}_X$-module quasi-cohérent et soit
$\mathcal{G}'$ un $\mathcal{O}_{X'}$-module quasi-cohérent. Alors
$$
\mathit{Hom}(\mathcal{F}, i_*\mathcal{G}') =
\mathit{Hom}(i^*\mathcal{F}, \mathcal{G}')
$$
comme foncteurs sur $(\Sch/B)$.
\end{lemma}

\begin{proof}
Soit $g : T \to B$ un morphisme, où $T$ est un schéma. Désignons par
$i_T : X'_T \to X_T$ le changement de base de $i$. Désignons par
$h : X_T \to X$ et $h' : X'_T \to X'$ les projections. Observons que
$(h')^*i^*\mathcal{F} = i_T^*h^*\mathcal{F}$. Comme une immersion fermée est
affine (Morphismes d'espaces, lemme
\ref{spaces-morphisms-lemma-closed-immersion-affine}), nous avons
$h^*i_*\mathcal{G} = i_{T, *}(h')^*\mathcal{G}$ d'après Cohomologie des
espaces, lemme \ref{spaces-cohomology-lemma-affine-base-change}. Ainsi,
\begin{align*}
\mathit{Hom}(\mathcal{F}, i_*\mathcal{G}')(T)
& =
\Hom_{\mathcal{O}_{X_T}}(h^*\mathcal{F}, h^*i_*\mathcal{G}') \\
& =
\Hom_{\mathcal{O}_{X_T}}(h^*\mathcal{F}, i_{T, *}(h')^*\mathcal{G}) \\
& =
\Hom_{\mathcal{O}_{X'_T}}(i_T^*h^*\mathcal{F}, (h')^*\mathcal{G}) \\
& =
\Hom_{\mathcal{O}_{X'_T}}((h')^*i^*\mathcal{F}, (h')^*\mathcal{G}) \\
& =
\mathit{Hom}(i^*\mathcal{F}, \mathcal{G}')(T)
\end{align*}
comme voulu. L'égalité du milieu résulte de l'adjonction des foncteurs
$i_{T, *}$ et $i_T^*$.
\end{proof}

\begin{lemma}
\label{quot-lemma-cohomology-perfect-complex}
Soit $S$ un schéma. Soit $B$ un espace algébrique sur $S$. Soit $K$ un objet
pseudo-cohérent de $D(\mathcal{O}_B)$.
\begin{enumerate}
\item Si, pour tout $g : T \to B$ dans $(\Sch/B)$, le faisceau de cohomologie
$H^{-1}(Lg^*K)$ est nul, alors le foncteur
$$
(\Sch/B)^{opp} \longrightarrow \textit{Ensembles},\quad
(g : T \to B) \longmapsto H^0(T, H^0(Lg^*K))
$$
est un espace algébrique affine et de présentation finie sur $B$.
\item Si, pour tout $g : T \to B$ dans $(\Sch/B)$, les faisceaux de
cohomologie $H^i(Lg^*K)$ sont nuls pour $i < 0$, alors $K$ est parfait,
$K$ est localement d'amplitude de Tor contenue dans $[0, b]$, et le foncteur
$$
(\Sch/B)^{opp} \longrightarrow \textit{Ensembles},\quad
(g : T \to B) \longmapsto H^0(T, Lg^*K)
$$
est un espace algébrique affine et de présentation finie sur $B$.
\end{enumerate}
\end{lemma}

\begin{proof}
Sous les hypothèses de (2), nous avons
$H^0(T, Lg^*K) = H^0(T, H^0(Lg^*K))$. Montrons que la règle
$T \mapsto H^0(T, H^0(Lg^*K))$ satisfait à la propriété de faisceau pour la
topologie fppf. Supposons pour cela que nous disposions d'un recouvrement fppf
$\{h_i : T_i \to T\}$ d'un schéma $g : T \to B$ sur $B$. Posons
$g_i = g \circ h_i$. Remarquons que, puisque $h_i$ est plat, nous avons
$Lh_i^* = h_i^*$ et $h_i^*$ commute à la prise de la cohomologie. Par suite,
$$
H^0(T_i, H^0(Lg_i^*K)) =
H^0(T_i, H^0(h_i^*Lg^*K)) =
H^0(T, h_i^*H^0(Lg^*K))
$$
Il en est de même pour l'image inverse sur $T_i \times_T T_j$. Comme
$Lg^*K$ est un complexe pseudo-cohérent sur $T$ (Cohomologie sur les sites,
lemme \ref{sites-cohomology-lemma-pseudo-coherent-pullback}), le faisceau de
cohomologie $\mathcal{F} = H^0(Lg^*K)$ est quasi-cohérent (Catégories dérivées
des espaces, lemme \ref{spaces-perfect-lemma-pseudo-coherent}). Ainsi, d'après
Descente sur les espaces, proposition
\ref{spaces-descent-proposition-fpqc-descent-quasi-coherent}, nous avons
$$
H^0(T, \mathcal{F}) = \Ker(
\prod H^0(T_i, h_i^*\mathcal{F}) \to
\prod H^0(T_i \times_T T_j, (T_i \times_T T_j \to T)^*\mathcal{F}))
$$
Nous voyons ainsi que les règles de (1) et (2) satisfont à la propriété de
faisceau pour les recouvrements fppf. Nous pouvons donc appliquer Amorçage,
lemme \ref{bootstrap-lemma-locally-algebraic-space-finite-type}, pour voir
qu'il suffit de démontrer la représentabilité localement pour la topologie
\'etale sur $B$. De plus, on peut vérifier localement pour la topologie \'etale
sur $B$ si le résultat final est affine et de présentation finie ; voir
Morphismes d'espaces, lemmes \ref{spaces-morphisms-lemma-affine-local} et
\ref{spaces-morphisms-lemma-finite-presentation-local}. Nous pouvons donc
supposer que $B$ est un schéma affine.

\medskip\noindent
Supposons que $B = \Spec(A)$ soit un schéma affine. D'après les résultats de
Catégories dérivées des espaces, lemmes
\ref{spaces-perfect-lemma-pseudo-coherent},
\ref{spaces-perfect-lemma-derived-quasi-coherent-small-etale-site} et
\ref{spaces-perfect-lemma-descend-pseudo-coherent}, nous pouvons, dans le reste
de la démonstration, considérer $K$ comme un objet parfait de la catégorie
dérivée des complexes de modules sur $B$ pour la topologie de Zariski.
D'après Catégories dérivées des schémas, lemmes
\ref{perfect-lemma-pseudo-coherent},
\ref{perfect-lemma-affine-compare-bounded} et
\ref{perfect-lemma-pseudo-coherent-affine}, nous pouvons trouver un complexe
pseudo-cohérent $M^\bullet$ de $A$-modules tel que $K$ soit l'objet
correspondant de $D(\mathcal{O}_B)$. Notre hypothèse sur les images inverses
implique que $M^\bullet \otimes^\mathbf{L}_A \kappa(\mathfrak p)$ a son
$H^{-1}$ nul pour tout idéal premier $\mathfrak p \subset A$. D'après Plus
d'algèbre, lemme \ref{more-algebra-lemma-cut-complex-in-two}, nous pouvons
écrire
$$
M^\bullet =
\tau_{\geq 0}M^\bullet \oplus \tau_{\leq - 1}M^\bullet
$$
où $\tau_{\geq 0}M^\bullet$ est parfait d'amplitude de Tor contenue dans
$[0, b]$ pour un certain $b \geq 0$ (nous avons aussi utilisé ici Plus
d'algèbre, lemmes \ref{more-algebra-lemma-glue-perfect} et
\ref{more-algebra-lemma-glue-tor-amplitude}). Remarquons que, dans le cas (2),
nous voyons aussi que $\tau_{\leq - 1}M^\bullet = 0$ dans $D(A)$, de sorte que
$M^\bullet$ et $K$ sont parfaits d'amplitude de Tor contenue dans $[0, b]$.
Pour tout $B$-schéma $g : T \to B$, nous avons
$$
H^0(T, H^0(Lg^*K)) = H^0(T, H^0(Lg^*\tau_{\geq 0}K))
$$
(par le dual de Catégories dérivées, lemme
\ref{derived-lemma-negative-vanishing}) ; nous pouvons donc remplacer $K$ par
$\tau_{\geq 0}K$ et, corrélativement, $M^\bullet$ par
$\tau_{\geq 0}M^\bullet$. Autrement dit, nous pouvons supposer que
$M^\bullet$ est d'amplitude de Tor contenue dans $[0, b]$.

\medskip\noindent
Supposons $M^\bullet$ d'amplitude de Tor contenue dans $[0, b]$. Nous pouvons
supposer que $M^\bullet$ est un complexe borné supérieurement de $A$-modules
libres de type fini (d'après notre définition des complexes pseudo-cohérents ;
voir Plus d'algèbre, définition
\ref{more-algebra-definition-pseudo-coherent}, et la discussion qui suit la
définition). D'après Plus d'algèbre, lemme
\ref{more-algebra-lemma-last-one-flat}, nous voyons que
$M = \Coker(M^{- 1} \to M^0)$ est plat. D'après Algèbre, lemme
\ref{algebra-lemma-finite-projective}, nous voyons que $M$ est localement
libre de type fini. Par conséquent, $M^\bullet$ est quasi-isomorphe à
$$
M \to M^1 \to M^2 \to \ldots \to M^d \to 0 \ldots
$$
Remarquons qu'il s'agit d'un complexe K-plat (Cohomologie, lemme
\ref{cohomology-lemma-bounded-flat-K-flat}) ; l'image inverse dérivée de $K$
par un morphisme $T \to B$ est donc calculée par le complexe
$$
g^*\widetilde{M} \to g^*\widetilde{M^1} \to \ldots
$$
Il suffit ainsi de montrer que le foncteur
$$
(g : T \to B) \longmapsto
\Ker(
\Gamma(T,g^*\widetilde{M})
\to
\Gamma(T, g^*(\widetilde{M^1})
)
$$
est représentable par un schéma affine de présentation finie sur $B$.

\medskip\noindent
Nous pouvons encore remplacer $B$ par les membres d'un recouvrement ouvert
affine pour démontrer ce dernier énoncé. Nous pouvons donc supposer que $M$
est libre de type fini (rappelons que $M^1$ est libre de type fini dès le
départ). Écrivons $M = A^{\oplus n}$ et $M^1 = A^{\oplus m}$. Supposons que
l'application $M \to M^1$ soit donnée par la matrice $m \times n$
$(a_{ij})$ à coefficients dans $A$. Alors
$\widetilde{M} = \mathcal{O}_B^{\oplus n}$ et
$\widetilde{M^1} = \mathcal{O}_B^{\oplus m}$. Le foncteur ci-dessus est donc
égal au foncteur
$$
(g : T \to B) \longmapsto
\{(f_1, \ldots, f_n) \in \Gamma(T, \mathcal{O}_T) \mid
\sum g^\sharp(a_{ij})f_i = 0,\ j = 1, \ldots, m\}
$$
Il est clairement représentable par le schéma affine
$$
\Spec\left(A[x_1, \ldots, x_n]/(\sum a_{ij}x_i; j = 1, \ldots, m)\right)
$$
ce qui démontre le lemme.
\end{proof}

\noindent
Le foncteur $\mathit{Hom}(\mathcal{F}, \mathcal{G})$ est représentable dans
plusieurs situations. Tous nos résultats reposeront sur le cas fondamental
suivant. La démonstration ci-dessous de ce lemme est, en un certain sens, la
généralisation naturelle de celle de \cite[III, cor. 7.7.8]{EGA}.

\begin{lemma}
\label{quot-lemma-noetherian-hom}
Dans la situation \ref{quot-situation-hom}, supposons que
\begin{enumerate}
\item $B$ soit un espace algébrique noethérien,
\item $f$ soit localement de type fini et quasi-séparé,
\item $\mathcal{F}$ soit un $\mathcal{O}_X$-module de type fini, et
\item $\mathcal{G}$ soit un $\mathcal{O}_X$-module de type fini, plat sur
$B$, à support propre sur $B$.
\end{enumerate}
Alors le foncteur $\mathit{Hom}(\mathcal{F}, \mathcal{G})$ est un espace
algébrique affine et de présentation finie sur $B$.
\end{lemma}

\begin{proof}
Nous pouvons remplacer $X$ par un voisinage ouvert quasi-compact du support de
$\mathcal{G}$ ; nous pouvons donc supposer $X$ noethérien. Dans ce cas, $X$ et
$f$ sont quasi-compacts et quasi-séparés. Choisissons une approximation
$P \to \mathcal{F}$ par un complexe parfait $P$ du triplet
$(X, \mathcal{F}, -1)$ ; voir Catégories dérivées des espaces, définition
\ref{spaces-perfect-definition-approximation-holds} et théorème
\ref{spaces-perfect-theorem-approximation}). Alors l'application induite
$$
\Hom_{\mathcal{O}_X}(\mathcal{F}, \mathcal{G})
\longrightarrow
\Hom_{D(\mathcal{O}_X)}(P, \mathcal{G})
$$
est un isomorphisme, parce que $P \to \mathcal{F}$ induit un isomorphisme
$H^0(P) \to \mathcal{F}$ et que $H^i(P) = 0$ pour $i > 0$. De plus, pour tout
morphisme $g : T \to B$, désignons par $h : X_T = T \times_B X \to X$ la
projection et posons $P_T = Lh^*P$. Il est alors encore vrai que
$$
\Hom_{\mathcal{O}_{X_T}}(\mathcal{F}_T, \mathcal{G}_T)
\longrightarrow
\Hom_{D(\mathcal{O}_{X_T})}(P_T, \mathcal{G}_T)
$$
est un isomorphisme, puisque
$P_T = Lh^*P \to Lh^*\mathcal{F} \to \mathcal{F}_T$ induit un isomorphisme
$H^0(P_T) \to \mathcal{F}_T$ (car $h^*$ est exact à droite et
$H^i(P) = 0$ pour $i > 0$). Il suffit donc de démontrer le résultat pour le
foncteur
$$
T \longmapsto \Hom_{D(\mathcal{O}_{X_T})}(P_T, \mathcal{G}_T).
$$
D'après la suite spectrale de Leray (voir Cohomologie sur les sites, remarque
\ref{sites-cohomology-remark-before-Leray}), nous avons
$$
\Hom_{D(\mathcal{O}_{X_T})}(P_T, \mathcal{G}_T) =
H^0(X_T, R\SheafHom(P_T, \mathcal{G}_T)) =
H^0(T, Rf_{T, *}R\SheafHom(P_T, \mathcal{G}_T))
$$
où $f_T : X_T \to T$ est le changement de base de $f$. D'après Catégories
dérivées des espaces, lemme
\ref{spaces-perfect-lemma-base-change-RHom}, nous avons
$$
Rf_{T, *}R\SheafHom(P_T, \mathcal{G}_T) = Lg^*Rf_*R\SheafHom(P, \mathcal{G}).
$$
D'après Catégories dérivées des espaces, lemme
\ref{spaces-perfect-lemma-ext-perfect}, l'objet
$K = Rf_*R\SheafHom(P, \mathcal{G})$ de $D(\mathcal{O}_B)$ est parfait. Nous
pouvons donc appliquer le lemme \ref{quot-lemma-cohomology-perfect-complex}, pourvu
que nous démontrions que le faisceau de cohomologie $H^i(Lg^*K)$ est $0$ pour
tout $i < 0$ et tout $g : T \to B$ comme ci-dessus. Ceci résulte clairement de
la dernière formule affichée, puisque les faisceaux de cohomologie de
$Rf_{T, *}R\SheafHom(P_T, \mathcal{G}_T)$ sont nuls en degrés négatifs : en
effet, les faisceaux de cohomologie de
$R\SheafHom(P_T, \mathcal{G}_T)$ sont nuls en degrés négatifs, car $P_T$ est
parfait et ses faisceaux de cohomologie sont nuls en degrés positifs.
\end{proof}

\noindent
Voici une conséquence facile du lemme \ref{quot-lemma-noetherian-hom}.

\begin{proposition}
\label{quot-proposition-hom}
Dans la situation \ref{quot-situation-hom}, supposons que
\begin{enumerate}
\item $f$ soit de présentation finie, et
\item $\mathcal{G}$ soit un $\mathcal{O}_X$-module de présentation finie,
plat sur $B$, à support propre sur $B$.
\end{enumerate}
Alors le foncteur $\mathit{Hom}(\mathcal{F}, \mathcal{G})$ est un espace
algébrique affine sur $B$. Si $\mathcal{F}$ est de présentation finie, alors
$\mathit{Hom}(\mathcal{F}, \mathcal{G})$ est de présentation finie sur $B$.
\end{proposition}

\begin{proof}
D'après le lemme \ref{quot-lemma-hom-sheaf}, le foncteur
$\mathit{Hom}(\mathcal{F}, \mathcal{G})$ satisfait à la propriété de faisceau
pour les recouvrements fppf. Cela signifie que nous pouvons\footnote{Nous
omettons de vérifier la condition ensembliste (3) du lemme cité.} appliquer
Amorçage, lemme \ref{bootstrap-lemma-locally-algebraic-space}, pour vérifier
la représentabilité localement pour la topologie \'etale sur $B$. De plus, nous
pouvons vérifier localement pour la topologie \'etale sur $B$ si le résultat
final est affine ou de présentation finie ; voir Morphismes d'espaces, lemmes
\ref{spaces-morphisms-lemma-affine-local} et
\ref{spaces-morphisms-lemma-finite-presentation-local}. Nous pouvons donc
supposer que $B$ est un schéma affine.

\medskip\noindent
Supposons que $B$ soit un schéma affine. Comme $f$ est de présentation finie,
il en résulte que $X$ est quasi-compact et quasi-séparé. Nous pouvons donc
écrire $\mathcal{F} = \colim \mathcal{F}_i$ comme une colimite filtrante de
$\mathcal{O}_X$-modules de présentation finie (Limites d'espaces, lemme
\ref{spaces-limits-lemma-colimit-finitely-presented}). Il est clair que
$$
\mathit{Hom}(\mathcal{F}, \mathcal{G}) =
\lim \mathit{Hom}(\mathcal{F}_i, \mathcal{G})
$$
Si nous pouvons donc montrer que chaque
$\mathit{Hom}(\mathcal{F}_i, \mathcal{G})$ est représentable par un schéma
affine, nous verrons qu'il en va de même de
$\mathit{Hom}(\mathcal{F}, \mathcal{G})$. Utiliser les résultats de Limites,
section \ref{limits-section-limits}, et de Limites d'espaces, section
\ref{spaces-limits-section-limits}. Nous pouvons donc supposer que
$\mathcal{F}$ est de présentation finie.

\medskip\noindent
Écrivons $B = \Spec(R)$. Écrivons $R = \colim R_i$, chaque $R_i$ étant une
$\mathbf{Z}$-algèbre de type fini. Posons $B_i = \Spec(R_i)$. D'après les
résultats de Limites d'espaces, lemmes
\ref{spaces-limits-lemma-descend-finite-presentation} et
\ref{spaces-limits-lemma-descend-modules-finite-presentation}, nous pouvons
trouver un $i$, un morphisme d'espaces algébriques $X_i \to B_i$, et des
$\mathcal{O}_{X_i}$-modules de présentation finie $\mathcal{F}_i$ et
$\mathcal{G}_i$, tels que le changement de base de
$(X_i, \mathcal{F}_i, \mathcal{G}_i)$ à $B$ redonne
$(X, \mathcal{F}, \mathcal{G})$. D'après Limites d'espaces, lemme
\ref{spaces-limits-lemma-descend-flat}, nous pouvons, quitte à augmenter $i$,
supposer que $\mathcal{G}_i$ est plat sur $B_i$. D'après Limites d'espaces,
lemme \ref{spaces-limits-lemma-eventually-proper-support}, nous pouvons de
même supposer que le support schématique de $\mathcal{G}_i$ est propre sur
$B_i$. Nous pouvons alors appliquer le lemme \ref{quot-lemma-noetherian-hom} pour
voir que $H_i = \mathit{Hom}(\mathcal{F}_i, \mathcal{G}_i)$ est un espace
algébrique affine de présentation finie sur $B_i$. En effectuant l'image
inverse sur $B$ (à l'aide de la remarque \ref{quot-remark-hom-base-change}), nous
voyons que $H_i \times_{B_i} B = \mathit{Hom}(\mathcal{F}, \mathcal{G})$,
et nous avons gagné.
\end{proof}

\noindent
Vérification de cohérence : le faisceau $\mathit{Hom}$ joue le même rôle parmi
les espaces algébriques sur $S$.

\section{Le foncteur Isom}
\label{quot-section-isom}

\noindent
Dans la situation \ref{quot-situation-hom}, nous pouvons considérer le
sous-foncteur
$$
\mathit{Isom}(\mathcal{F}, \mathcal{G}) \subset
\mathit{Hom}(\mathcal{F}, \mathcal{G})
$$
dont la valeur sur un schéma $T$ sur $B$ est l'ensemble des homomorphismes
{\it inversibles} de $\mathcal{O}_{X_T}$-modules
$u : \mathcal{F}_T \to \mathcal{G}_T$.

\medskip\noindent
Nous regardons parfois $\mathit{Isom}(\mathcal{F}, \mathcal{G})$ comme un
foncteur $(\Sch/S)^{opp} \to \textit{Ensembles}$ muni d'un morphisme
$\mathit{Isom}(\mathcal{F}, \mathcal{G}) \to B$. En effet, si $T$ est un
schéma sur $S$, un élément de
$\mathit{Isom}(\mathcal{F}, \mathcal{G})(T)$ consiste en un couple $(h, u)$,
où $h$ est un morphisme $h : T \to B$ et
$u : \mathcal{F}_T \to \mathcal{G}_T$ est un isomorphisme de
$\mathcal{O}_{X_T}$-modules, avec $X_T = T \times_{h, B} X$ et
$\mathcal{F}_T$, $\mathcal{G}_T$ les images inverses sur $X_T$. En
particulier, lorsque nous disons que
$\mathit{Isom}(\mathcal{F}, \mathcal{G})$ est un espace algébrique, nous
entendons que le foncteur correspondant
$(\Sch/S)^{opp} \to \textit{Ensembles}$ est un espace algébrique.

\begin{lemma}
\label{quot-lemma-isom-sheaf}
Dans la situation \ref{quot-situation-hom}, le foncteur
$\mathit{Isom}(\mathcal{F}, \mathcal{G})$
satisfait à la propriété de faisceau pour la topologie fpqc.
\end{lemma}

\begin{proof}
Nous avons déjà vu que $\mathit{Hom}(\mathcal{F}, \mathcal{G})$ satisfait à
la propriété de faisceau. Il reste donc à montrer ceci : si
$\{T_i \to T\}_{i \in I}$ est un recouvrement fpqc de schémas sur $B$ et si
$u : \mathcal{F}_T \to \mathcal{G}_T$ est une application
$\mathcal{O}_{X_T}$-linéaire telle que $u_{T_i}$ soit un isomorphisme pour tout
$i$, alors $u$ est un isomorphisme. Puisque
$\{X_i \to X_T\}_{i \in I}$ est un recouvrement fpqc de $X_T$ ; voir
Topologies sur les espaces, lemme \ref{spaces-topologies-lemma-fpqc}, ceci
résulte de Descente sur les espaces, proposition
\ref{spaces-descent-proposition-fpqc-descent-quasi-coherent}.
\end{proof}

\noindent
Vérification de cohérence : le faisceau $\mathit{Isom}$ joue le même rôle parmi
les espaces algébriques sur $S$.

\begin{lemma}
\label{quot-lemma-extend-isom-to-spaces}
Dans la situation \ref{quot-situation-hom}. Soit $T$ un espace algébrique sur $S$.
Nous avons
$$
\Mor_{\Sh((\Sch/S)_{fppf})}(T, \mathit{Isom}(\mathcal{F}, \mathcal{G})) =
\{(h, u) \mid
h : T \to B, u : \mathcal{F}_T \to \mathcal{G}_T\text{ isomorphisme}\}
$$
où $\mathcal{F}_T, \mathcal{G}_T$ désignent les images inverses de
$\mathcal{F}$ et $\mathcal{G}$ sur l'espace algébrique
$X \times_{B, h} T$.
\end{lemma}

\begin{proof}
Observons que les membres de gauche et de droite de l'égalité sont des
sous-ensembles des membres de gauche et de droite de l'égalité du lemme
\ref{quot-lemma-extend-hom-to-spaces}. Nous omettons de vérifier que ces
sous-ensembles se correspondent par l'identification donnée dans la
démonstration de ce lemme.
\end{proof}

\begin{proposition}
\label{quot-proposition-isom}
Dans la situation \ref{quot-situation-hom}, supposons que
\begin{enumerate}
\item $f$ soit de présentation finie, et
\item $\mathcal{F}$ et $\mathcal{G}$ soient des $\mathcal{O}_X$-modules de
présentation finie, plats sur $B$, à support propre sur $B$.
\end{enumerate}
Alors le foncteur $\mathit{Isom}(\mathcal{F}, \mathcal{G})$ est un espace
algébrique affine de présentation finie sur $B$.
\end{proposition}

\begin{proof}
Nous emploierons les abréviations
$H = \mathit{Hom}(\mathcal{F}, \mathcal{G})$,
$I = \mathit{Hom}(\mathcal{F}, \mathcal{F})$,
$H' = \mathit{Hom}(\mathcal{G}, \mathcal{F})$ et
$I' = \mathit{Hom}(\mathcal{G}, \mathcal{G})$.
D'après la proposition \ref{quot-proposition-hom}, les foncteurs
$H$, $I$, $H'$, $I'$ sont des espaces algébriques et les morphismes
$H \to B$, $I \to B$, $H' \to B$ et $I' \to B$ sont affines et de
présentation finie. La composition des applications donne un morphisme
$$
c : H' \times_B H \longrightarrow I \times_B I',\quad
(u', u) \longmapsto (u \circ u', u' \circ u)
$$
d'espaces algébriques sur $B$. Comme $I \times_B I' \to B$ est séparé, la
section $\sigma : B \to I \times_B I'$ correspondant à
$(\text{id}_\mathcal{F}, \text{id}_\mathcal{G})$ est une immersion fermée
(Morphismes d'espaces, lemme
\ref{spaces-morphisms-lemma-section-immersion}). De plus, $\sigma$ est de
présentation finie (Morphismes d'espaces, lemme
\ref{spaces-morphisms-lemma-finite-presentation-permanence}). Par conséquent,
$$
\mathit{Isom}(\mathcal{F}, \mathcal{G}) =
(H' \times_B H) \times_{c, I \times_B I', \sigma} B
$$
est lui aussi un espace algébrique affine de présentation finie sur $B$.
Nous omettons quelques détails.
\end{proof}






\section{Le champ des faisceaux cohérents}
\label{quot-section-stack-coherent-sheaves}

\noindent
Dans cette section, nous démontrons que le champ des faisceaux cohérents sur
$X/B$ est algébrique sous des hypothèses convenables. C'est un cas particulier
de \cite[théorème 2.1.1]{lieblich_remarks}, qui traite du champ des faisceaux
cohérents sur un champ d'Artin au-dessus d'une base.

\begin{situation}
\label{quot-situation-coherent}
Soit $S$ un schéma. Soit $f : X \to B$ un morphisme d'espaces algébriques sur
$S$. Supposons que $f$ soit de présentation finie. Nous désignons par
$\Cohstack_{X/B}$ la catégorie dont les objets sont les triplets
$(T, g, \mathcal{F})$ où
\begin{enumerate}
\item $T$ est un schéma sur $S$,
\item $g : T \to B$ est un morphisme sur $S$, et, en posant
$X_T = T \times_{g, B} X$,
\item $\mathcal{F}$ est un $\mathcal{O}_{X_T}$-module quasi-cohérent de
présentation finie, plat sur $T$, à support propre sur $T$.
\end{enumerate}
Un morphisme $(T, g, \mathcal{F}) \to (T', g', \mathcal{F}')$ est donné par
un couple $(h, \varphi)$ où
\begin{enumerate}
\item $h : T \to T'$ est un morphisme de schémas sur $B$
(c'est-à-dire que $g' \circ h = g$), et
\item $\varphi : (h')^*\mathcal{F}' \to \mathcal{F}$ est un isomorphisme de
$\mathcal{O}_{X_T}$-modules, où $h' : X_T \to X_{T'}$ est le changement de
base de $h$.
\end{enumerate}
\end{situation}

\noindent
Ainsi, $\Cohstack_{X/B}$ est une catégorie et la règle
$$
p : \Cohstack_{X/B} \longrightarrow (\Sch/S)_{fppf},
\quad
(T, g, \mathcal{F}) \longmapsto T
$$
est un foncteur. Pour un schéma $T$ sur $S$, nous désignons par
$\Cohstack_{X/B, T}$ la catégorie fibre de $p$ en $T$. Ces catégories fibres
sont des groupoïdes.

\begin{lemma}
\label{quot-lemma-coherent-fibred-in-groupoids}
Dans la situation \ref{quot-situation-coherent}, le foncteur
$p : \Cohstack_{X/B} \longrightarrow (\Sch/S)_{fppf}$ est fibré en
groupoïdes.
\end{lemma}

\begin{proof}
Nous montrons que $p$ est fibré en groupoïdes en vérifiant les conditions (1)
et (2) de Catégories, définition
\ref{categories-definition-fibred-groupoids}. Étant donnés un objet
$(T', g', \mathcal{F}')$ de $\Cohstack_{X/B}$ et un morphisme
$h : T \to T'$ de schémas sur $S$, nous pouvons poser $g = h \circ g'$ et
$\mathcal{F} = (h')^*\mathcal{F}'$, où $h' : X_T \to X_{T'}$ est le
changement de base de $h$. Il est alors clair que nous obtenons un morphisme
$(T, g, \mathcal{F}) \to (T', g', \mathcal{F}')$ de $\Cohstack_{X/B}$
au-dessus de $h$. Cela démontre (1). Pour (2), supposons donnés des morphismes
$$
(h_1, \varphi_1) : (T_1, g_1, \mathcal{F}_1) \to (T, g, \mathcal{F})
\quad\text{et}\quad
(h_2, \varphi_2) : (T_2, g_2, \mathcal{F}_2) \to (T, g, \mathcal{F})
$$
de $\Cohstack_{X/B}$ et un morphisme $h : T_1 \to T_2$ tel que
$h_2 \circ h = h_1$. Nous pouvons alors prendre pour $\varphi$ le composé
$$
(h')^*\mathcal{F}_2
\xrightarrow{(h')^*\varphi_2^{-1}}
(h')^*(h_2)^*\mathcal{F} = (h_1)^*\mathcal{F}
\xrightarrow{\varphi_1}
\mathcal{F}_1
$$
et obtenir ainsi le morphisme
$(h, \varphi) : (T_1, g_1, \mathcal{F}_1) \to (T_2, g_2, \mathcal{F}_2)$
qui atteste la condition (2).
\end{proof}

\begin{lemma}
\label{quot-lemma-coherent-diagonal}
Dans la situation \ref{quot-situation-coherent}. Posons
$\mathcal{X} = \Cohstack_{X/B}$. Alors
$\Delta : \mathcal{X} \to \mathcal{X} \times \mathcal{X}$ est représentable
par des espaces algébriques.
\end{lemma}

\begin{proof}
Considérons deux objets $x = (T, g, \mathcal{F})$ et
$y = (T, h, \mathcal{G})$ de $\mathcal{X}$ sur un schéma $T$. Nous devons
montrer que $\mathit{Isom}_\mathcal{X}(x, y)$ est un espace algébrique sur
$T$ ; voir Champs algébriques, lemme
\ref{algebraic-lemma-representable-diagonal}. Si, pour $a : T' \to T$, les
restrictions $x|_{T'}$ et $y|_{T'}$ sont isomorphes dans la catégorie fibre
$\mathcal{X}_{T'}$, alors $g \circ a = h \circ a$. Il existe donc une
transformation de préfaisceaux
$$
\mathit{Isom}_\mathcal{X}(x, y) \longrightarrow \text{Égalisateur}(g, h)
$$
Comme la diagonale de $B$ est représentable (par des schémas), cet égalisateur
est un schéma. Nous pouvons donc remplacer $T$ par cet égalisateur et les
faisceaux $\mathcal{F}$ et $\mathcal{G}$ par leurs images inverses. Nous
pouvons ainsi supposer $g = h$. Dans ce cas,
$\mathit{Isom}_\mathcal{X}(x, y) = \mathit{Isom}(\mathcal{F}, \mathcal{G})$,
et le résultat découle de la proposition \ref{quot-proposition-isom}.
\end{proof}

\begin{lemma}
\label{quot-lemma-coherent-stack}
Dans la situation \ref{quot-situation-coherent}, le foncteur
$p : \Cohstack_{X/B} \longrightarrow (\Sch/S)_{fppf}$ est un champ en
groupoïdes.
\end{lemma}

\begin{proof}
Pour démontrer que $\Cohstack_{X/B}$ est un champ en groupoïdes, il faut
montrer que les préfaisceaux $\mathit{Isom}$ sont des faisceaux et que les
données de descente sont effectives. L'assertion relative à $\mathit{Isom}$
résulte du lemme \ref{quot-lemma-coherent-diagonal} ; voir Champs algébriques,
lemme \ref{algebraic-lemma-representable-diagonal}. Démontrons l'assertion
relative aux données de descente. Supposons que
$\{a_i : T_i \to T\}$ soit un recouvrement fppf de schémas sur $S$. Soit
$(\xi_i, \varphi_{ij})$ une donnée de descente pour $\{T_i \to T\}$ à valeurs
dans $\Cohstack_{X/B}$. Pour chaque $i$, écrivons
$\xi_i = (T_i, g_i, \mathcal{F}_i)$. Désignons par
$\text{pr}_0 : T_i \times_T T_j \to T_i$ et
$\text{pr}_1 : T_i \times_T T_j \to T_j$ les projections. La condition
$\xi_i|_{T_i \times_T T_j} = \xi_j|_{T_i \times_T T_j}$ implique notamment
que $g_i \circ \text{pr}_0 = g_j \circ \text{pr}_1$. Il existe donc un unique
morphisme $g : T \to B$ tel que $g_i = g \circ a_i$ ; voir Descente sur les
espaces, lemme
\ref{spaces-descent-lemma-fpqc-universal-effective-epimorphisms}.
Désignons par $X_T = T \times_{g, B} X$. Posons
$X_i = X_{T_i} = T_i \times_{g_i, B} X = T_i \times_{a_i, T} X_T$ et
$$
X_{ij} = X_{T_i} \times_{X_T} X_{T_j} = X_i \times_{X_T} X_j
$$
avec les projections $\text{pr}_i$ et $\text{pr}_j$ sur $X_i$ et $X_j$.
Observons que l'image inverse de $(T_i, g_i, \mathcal{F}_i)$ par
$\text{pr}_0 : T_i \times_T T_j \to T_i$ est donnée par
$(T_i \times_T T_j, g_i \circ \text{pr}_0, \text{pr}_i^*\mathcal{F}_i)$.
Ainsi, une donnée de descente pour $\{T_i \to T\}$ dans
$\Cohstack_{X/B}$ est donnée par les objets
$(T_i, g \circ a_i, \mathcal{F}_i)$ et, pour chaque couple $i, j$, par un
isomorphisme de $\mathcal{O}_{X_{ij}}$-modules
$$
\varphi_{ij} :
\text{pr}_i^*\mathcal{F}_i \longrightarrow \text{pr}_j^*\mathcal{F}_j
$$
satisfaisant à la condition de cocycle sur (l'image inverse de $X$ à)
$T_i \times_T T_j \times_T T_k$. Bien ; il suffit maintenant d'utiliser le
fait que $\{X_i \to X_T\}$ est un recouvrement fppf pour regarder
$(\mathcal{F}_i, \varphi_{ij})$ comme une donnée de descente pour ce
recouvrement. D'après Descente sur les espaces, proposition
\ref{spaces-descent-proposition-fpqc-descent-quasi-coherent}, cette donnée de
descente est effective et nous obtenons un faisceau quasi-cohérent
$\mathcal{F}$ sur $X_T$, dont la restriction à $X_i$ est $\mathcal{F}_i$.
D'après Morphismes d'espaces, lemme
\ref{spaces-morphisms-lemma-flat-permanence}, nous voyons que $\mathcal{F}$ est
plat sur $T$, et Descente sur les espaces, lemme
\ref{spaces-descent-lemma-finite-presentation-descends}, garantit que
$\mathcal{F}$ est de présentation finie comme $\mathcal{O}_{X_T}$-module.
Enfin, d'après Descente sur les espaces, lemme
\ref{spaces-descent-lemma-descending-property-proper}, nous voyons que le
support schématique de $\mathcal{F}$ est propre sur $T$, puisque nous avons
supposé le support schématique de $\mathcal{F}_i$ propre sur $T_i$
(remarquons que la formation du support schématique commute au changement de
base plat d'après Morphismes d'espaces, lemme
\ref{spaces-morphisms-lemma-flat-pullback-support}). Nous obtenons ainsi
l'objet recherché sur $T$.
\end{proof}

\begin{remark}
\label{quot-remark-coherent-base-change}
Dans la situation \ref{quot-situation-coherent}, la règle
$(T, g, \mathcal{F}) \mapsto (T, g)$ définit un $1$-morphisme
$$
\Cohstack_{X/B} \longrightarrow \mathcal{S}_B
$$
de champs en groupoïdes (voir le lemme \ref{quot-lemma-coherent-stack}, Champs
algébriques, section \ref{algebraic-section-split}, et Exemples de champs,
section \ref{examples-stacks-section-stack-associated-to-sheaf}). Soit
$B' \to B$ un morphisme d'espaces algébriques sur $S$. Soit
$\mathcal{S}_{B'} \to \mathcal{S}_B$ le $1$-morphisme associé de champs
fibrés en ensembles. Posons $X' = X \times_B B'$. Nous obtenons un champ en
groupoïdes $\Cohstack_{X'/B'} \to (\Sch/S)_{fppf}$ associé au changement de
base $f' : X' \to B'$. Dans cette situation, le diagramme
$$
\vcenter{
\xymatrix{
\Cohstack_{X'/B'} \ar[r] \ar[d] & \Cohstack_{X/B} \ar[d] \\
\mathcal{S}_{B'} \ar[r] & \mathcal{S}_B
}
}
\quad
\begin{matrix}
\text{ou, dans} \\
\text{une autre} \\
\text{notation}
\end{matrix}
\quad
\vcenter{
\xymatrix{
\Cohstack_{X'/B'} \ar[r] \ar[d] & \Cohstack_{X/B} \ar[d] \\
\Sch/B' \ar[r] & \Sch/B
}
}
$$
est un carré $2$-cartésien. Cette remarque élémentaire sera parfois utile pour
changer l'espace algébrique de base.
\end{remark}

\begin{lemma}
\label{quot-lemma-coherent-limits}
Dans la situation \ref{quot-situation-coherent}, supposons que $B \to S$ soit
localement de présentation finie. Alors
$p : \Cohstack_{X/B} \to (\Sch/S)_{fppf}$ commute aux limites (Axiomes
d'Artin, définition \ref{artin-definition-limit-preserving}).
\end{lemma}

\begin{proof}
Notons $B(T)$ la catégorie discrète dont les objets sont les $S$-morphismes
$T \to B$. Soit $T = \lim T_i$ une limite filtrante de schémas affines sur
$S$. Associer à un objet $(T, h, \mathcal{F})$ de
$\Cohstack_{X/B, T}$ l'objet $h$ de $B(T)$ donne un diagramme commutatif de
catégories fibres
$$
\xymatrix{
\colim \Cohstack_{X/B, T_i} \ar[r] \ar[d] &
\Cohstack_{X/B, T} \ar[d] \\
\colim B(T_i) \ar[r] & B(T)
}
$$
Nous devons montrer que la flèche horizontale supérieure est une équivalence.
Puisque nous avons supposé que $B$ est localement de présentation finie sur
$S$, il résulte de Limites d'espaces, remarque
\ref{spaces-limits-remark-limit-preserving}, que la flèche horizontale
inférieure est une équivalence. Cela signifie que nous pouvons supposer que
$T = \lim T_i$ soit une limite filtrante de schémas affines sur $B$.
Désignons par $g_i : T_i \to B$ et $g : T \to B$ les morphismes
correspondants. Posons $X_i = T_i \times_{g_i, B} X$ et
$X_T = T \times_{g, B} X$. Observons que $X_T = \colim X_i$ et que les
espaces algébriques $X_i$ et $X_T$ sont quasi-séparés et quasi-compacts
(puisqu'ils sont de présentation finie sur les schémas affines $T_i$ et $T$).
D'après Limites d'espaces, lemme
\ref{spaces-limits-lemma-descend-modules-finite-presentation}, nous avons
$$
\colim \textit{FP}(X_i) = \textit{FP}(X_T).
$$
où $\textit{FP}(W)$ abrège la catégorie des $\mathcal{O}_W$-modules de
présentation finie. Les résultats de Limites d'espaces, lemmes
\ref{spaces-limits-lemma-descend-flat} et
\ref{spaces-limits-lemma-eventually-proper-support}, nous disent que la même
assertion demeure vraie si nous remplaçons $\textit{FP}(X_i)$ et
$\textit{FP}(X_T)$ par la sous-catégorie pleine des objets plats sur $T_i$ et
$T$, dont le support schématique est propre sur $T_i$ et $T$. Cela démontre le
lemme.
\end{proof}

\begin{lemma}
\label{quot-lemma-coherent-RS-star}
Dans la situation \ref{quot-situation-coherent}. Soit
$$
\xymatrix{
Z \ar[r] \ar[d] & Z' \ar[d] \\
Y \ar[r] & Y'
}
$$
une somme amalgamée dans la catégorie des schémas sur $S$, où
$Z \to Z'$ est un épaississement et $Z \to Y$ est affine ; voir Plus sur les
morphismes, lemme
\ref{more-morphisms-lemma-pushout-along-thickening}. Alors le foncteur entre
catégories fibres
$$
\Cohstack_{X/B, Y'}
\longrightarrow
\Cohstack_{X/B, Y} \times_{\Cohstack_{X/B, Z}} \Cohstack_{X/B, Z'}
$$
est une équivalence.
\end{lemma}

\begin{proof}
Observons que l'application correspondante
$$
B(Y') \longrightarrow B(Y) \times_{B(Z)} B(Z')
$$
est une bijection ; voir Sommes amalgamées d'espaces, lemme
\ref{spaces-pushouts-lemma-pushout-along-thickening-schemes}. Le diagramme
commutatif
$$
\xymatrix{
\Cohstack_{X/B, Y'} \ar[r] \ar[d] &
\Cohstack_{X/B, Y} \times_{\Cohstack_{X/B, Z}} \Cohstack_{X/B, Z'}
\ar[d] \\
B(Y') \ar[r] & B(Y) \times_{B(Z)} B(Z')
}
$$
montre donc que nous pouvons supposer que $Y'$ est un schéma sur $B'$.
D'après la remarque \ref{quot-remark-coherent-base-change}, nous pouvons remplacer
$B$ par $Y'$ et $X$ par $X \times_B Y'$. Nous pouvons ainsi supposer
$B = Y'$. Dans ce cas, l'assertion résulte de Sommes amalgamées d'espaces,
lemme \ref{spaces-pushouts-lemma-space-over-pushout-flat-modules}.
\end{proof}

\begin{lemma}
\label{quot-lemma-coherent-over-first-order-thickening}
Soit
$$
\xymatrix{
X \ar[d] \ar[r]_i & X' \ar[d] \\
T \ar[r] & T'
}
$$
un carré cartésien d'espaces algébriques, où $T \to T'$ est un
épaississement du premier ordre. Soit $\mathcal{F}'$ un
$\mathcal{O}_{X'}$-module plat sur $T'$. Posons
$\mathcal{F} = i^*\mathcal{F}'$. Les conditions suivantes sont équivalentes :
\begin{enumerate}
\item $\mathcal{F}'$ est un $\mathcal{O}_{X'}$-module quasi-cohérent de
présentation finie,
\item $\mathcal{F}'$ est un $\mathcal{O}_{X'}$-module de présentation finie,
\item $\mathcal{F}$ est un $\mathcal{O}_X$-module quasi-cohérent de
présentation finie,
\item $\mathcal{F}$ est un $\mathcal{O}_X$-module de présentation finie.
\end{enumerate}
\end{lemma}

\begin{proof}
Rappelons qu'un module de présentation finie est quasi-cohérent ; cela donne
les équivalences entre (1) et (2), et entre (3) et (4). L'équivalence de (2)
et (4) est un cas particulier de Théorie des déformations, lemme
\ref{defos-lemma-deform-fp-module-ringed-topoi}.
\end{proof}

\begin{lemma}
\label{quot-lemma-coherent-tangent-space}
Dans la situation \ref{quot-situation-coherent}, supposons que $S$ soit un schéma
localement noethérien et que $B \to S$ soit localement de présentation finie.
Soit $k$ un corps de type fini sur $S$ et soit
$x_0 = (\Spec(k), g_0, \mathcal{G}_0)$ un objet de
$\mathcal{X} = \Cohstack_{X/B}$ sur $k$. Alors les espaces
$T\mathcal{F}_{\mathcal{X}, k, x_0}$ et
$\text{Inf}(\mathcal{F}_{\mathcal{X}, k, x_0})$ (Axiomes d'Artin, section
\ref{artin-section-tangent-spaces}) sont de dimension finie.
\end{lemma}

\begin{proof}
Observons que, d'après le lemme \ref{quot-lemma-coherent-RS-star}, notre champ en
groupoïdes $\mathcal{X}$ satisfait à la propriété (RS*) définie dans Axiomes
d'Artin, section \ref{artin-section-inf}. En particulier, $\mathcal{X}$
satisfait à (RS). Toutes les catégories de prédéformations associées sont donc
des catégories de déformations (Axiomes d'Artin, lemme
\ref{artin-lemma-deformation-category}), et l'énoncé a un sens.

\medskip\noindent
Nous montrons dans ce paragraphe que nous pouvons nous ramener au cas
$B = \Spec(k)$. Posons $X_0 = \Spec(k) \times_{g_0, B} X$ et désignons par
$\mathcal{X}_0 = \Cohstack_{X_0/k}$. Dans la remarque
\ref{quot-remark-coherent-base-change}, nous avons vu que $\mathcal{X}_0$ est le
$2$-produit fibré de $\mathcal{X}$ et $\Spec(k)$ sur $B$, comme catégories
fibrées en groupoïdes sur $(\Sch/S)_{fppf}$. Ainsi, d'après Axiomes d'Artin,
lemme \ref{artin-lemma-fibre-product-tangent-spaces}, nous sommes ramenés à
démontrer que $B$, $\Spec(k)$ et $\mathcal{X}_0$ possèdent des espaces
tangents et des espaces d'automorphismes infinitésimaux de dimension finie.
Les espaces tangents de $B$ et de $\Spec(k)$ sont de dimension finie d'après
Axiomes d'Artin, lemme \ref{artin-lemma-finite-dimension}, et leur
$\text{Inf}$ est évidemment nul. Il suffit donc de traiter $\mathcal{X}_0$.

\medskip\noindent
Soit $k[\epsilon]$ l'algèbre des nombres duaux sur $k$. Soit
$\Spec(k[\epsilon]) \to B$ le composé de $g_0 : \Spec(k) \to B$ et du
morphisme $\Spec(k[\epsilon]) \to \Spec(k)$ provenant de l'inclusion
$k \to k[\epsilon]$. Posons $X_0 = \Spec(k) \times_B X$ et
$X_\epsilon = \Spec(k[\epsilon]) \times_B X$. Observons que $X_\epsilon$ est
un épaississement du premier ordre de $X_0$, plat sur l'épaississement du
premier ordre $\Spec(k) \to \Spec(k[\epsilon])$. En développant les
définitions et en utilisant le lemme
\ref{quot-lemma-coherent-over-first-order-thickening}, nous voyons que
$T\mathcal{F}_{\mathcal{X}_0, k, x_0}$ est l'ensemble des relèvements de
$\mathcal{G}_0$ en un module plat sur $X_\epsilon$. D'après Théorie des
déformations, lemme \ref{defos-lemma-flat-ringed-topoi}, nous concluons que
$$
T\mathcal{F}_{\mathcal{X}_0, k, x_0} =
\Ext^1_{\mathcal{O}_{X_0}}(\mathcal{G}_0, \mathcal{G}_0)
$$
Nous avons utilisé ici l'identification $\epsilon k[\epsilon] \cong k$ de
$k[\epsilon]$-modules. En utilisant une fois encore Théorie des déformations,
lemme \ref{defos-lemma-flat-ringed-topoi}, nous voyons que
$$
\text{Inf}(\mathcal{F}_{\mathcal{X}, k, x_0}) =
\Ext^0_{\mathcal{O}_{X_0}}(\mathcal{G}_0, \mathcal{G}_0)
$$
Ces espaces sont de dimension finie sur $k$, car $\mathcal{G}_0$ est à support
propre sur $\Spec(k)$. En effet, $X_0$ est de présentation finie sur
$\Spec(k)$, donc noethérien. Comme $\mathcal{G}_0$ est de présentation finie,
c'est un $\mathcal{O}_{X_0}$-module cohérent. Nous pouvons donc appliquer
Catégories dérivées des espaces, lemme
\ref{spaces-perfect-lemma-ext-finite}, pour conclure à la finitude voulue.
\end{proof}

\begin{lemma}
\label{quot-lemma-coherent-existence}
Dans la situation \ref{quot-situation-coherent}, supposons que $S$ soit un schéma
localement noethérien et que $f : X \to B$ soit séparé. Soit
$\mathcal{X} = \Cohstack_{X/B}$. Alors le foncteur d'Axiomes d'Artin,
équation (\ref{artin-equation-approximation}), est une équivalence.
\end{lemma}

\begin{proof}
Soit $A$ une $S$-algèbre qui est un anneau local noethérien complet, d'idéal
maximal $\mathfrak m$, et dont le corps résiduel $k$ est de type fini sur $S$.
Nous devons montrer que la catégorie des objets sur $A$ est équivalente à la
catégorie des objets formels sur $A$. Comme nous savons que ceci vaut pour la
catégorie $\mathcal{S}_B$ fibrée en ensembles associée à $B$, d'après Axiomes
d'Artin, lemme \ref{artin-lemma-effective}, il suffit de le démontrer pour les
objets situés au-dessus d'un morphisme donné $\Spec(A) \to B$.

\medskip\noindent
Posons $X_A = \Spec(A) \times_B X$ et
$X_n = \Spec(A/\mathfrak m^n) \times_B X$. D'après le théorème d'existence de
Grothendieck (Plus sur les morphismes d'espaces, théorème
\ref{spaces-more-morphisms-theorem-grothendieck-existence}), la catégorie des
modules cohérents $\mathcal{F}$ sur $X_A$ à support propre sur $\Spec(A)$ est
équivalente à la catégorie des systèmes $(\mathcal{F}_n)$ de modules cohérents
$\mathcal{F}_n$ sur $X_n$ à support propre sur
$\Spec(A/\mathfrak m^n)$. L'équivalence envoie $\mathcal{F}$ sur le système
$(\mathcal{F} \otimes_A A/\mathfrak m^n)$. Voir la discussion dans Plus sur
les morphismes d'espaces, remarque
\ref{spaces-more-morphisms-remark-reformulate-existence-theorem}. Pour achever
la démonstration du lemme, il suffit de montrer que $\mathcal{F}$ est plat sur
$A$ si et seulement si tous les $\mathcal{F} \otimes_A A/\mathfrak m^n$ sont
plats sur $A/\mathfrak m^n$. Cela résulte de Plus sur les morphismes d'espaces,
lemme \ref{spaces-more-morphisms-lemma-flatness-over-Noetherian-ring}.
\end{proof}

\begin{lemma}
\label{quot-lemma-coherent-defo-thy}
Dans la situation \ref{quot-situation-coherent}, supposons que $S$ soit un schéma
localement noethérien, que $S = B$ et que $f : X \to B$ soit plat. Soit
$\mathcal{X} = \Cohstack_{X/B}$. Alors la versalité est ouverte pour
$\mathcal{X}$ (voir Axiomes d'Artin, définition
\ref{artin-definition-openness-versality}).
\end{lemma}

\begin{proof}[Première démonstration]
Cette démonstration repose sur le critère d'Axiomes d'Artin, lemme
\ref{artin-lemma-dual-openness}. Soit $U \to S$ un morphisme de schémas de
type fini, soit $x$ un objet de $\mathcal{X}$ sur $U$ et soit $u_0 \in U$ un
point de type fini tel que $x$ soit versel en $u_0$. En rétrécissant $U$, nous
pouvons supposer que $u_0$ est un point fermé (Morphismes, lemme
\ref{morphisms-lemma-point-finite-type}) et que $U = \Spec(A)$, le morphisme
$U \to S$ étant à valeurs dans un ouvert affine $\Spec(\Lambda)$ de $S$.
Soit $\mathcal{F}$ le module cohérent sur
$X_A = \Spec(A) \times_S X$, plat sur $A$, qui correspond à l'objet donné
$x$.

\medskip\noindent
D'après Théorie des déformations, lemme
\ref{defos-lemma-flat-ringed-topoi}, nous avons un isomorphisme de foncteurs
$$
T_x(M) = \Ext^1_{X_A}(\mathcal{F}, \mathcal{F} \otimes_A M)
$$
et, étant donnée une surjection $A' \to A$ de $\Lambda$-algèbres, de noyau de
carré nul $I$, nous avons une classe d'obstruction
$$
\xi_{A'} \in \Ext^2_{X_A}(\mathcal{F}, \mathcal{F} \otimes_A I)
$$
Ceci utilise le fait que, pour tout $A' \to A$ comme ci-dessus, le changement
de base $X_{A'} = \Spec(A') \times_B X$ est plat sur $A'$. De plus, la
construction de la classe d'obstruction est fonctorielle en la surjection
$A' \to A$ (pour $A$ fixé), d'après Théorie des déformations, lemme
\ref{defos-lemma-functorial-ringed-topoi}. Appliquons Catégories dérivées des
espaces, lemme \ref{spaces-perfect-lemma-compute-ext}, au calcul des groupes
Ext $\Ext^i_{X_A}(\mathcal{F}, \mathcal{F} \otimes_A M)$ pour $i \leq m$,
avec $m = 2$. Nous trouvons un objet parfait $K \in D(A)$ et des
isomorphismes fonctoriels
$$
H^i(K \otimes_A^\mathbf{L} M)
\longrightarrow
\Ext^i_{X_A}(\mathcal{F}, \mathcal{F} \otimes_A M)
$$
pour $i \leq m$, compatibles aux morphismes de bord. Cet objet $K$, avec les
identifications affichées ci-dessus, nous fournit une donnée comme dans
Axiomes d'Artin, situation \ref{artin-situation-dual}. Enfin, la condition (iv)
d'Axiomes d'Artin, lemme \ref{artin-lemma-dual-obstruction}, est satisfaite
d'après Théorie des déformations, lemme
\ref{defos-lemma-verify-iv-ringed-topoi}. Axiomes d'Artin, lemme
\ref{artin-lemma-dual-openness}, s'applique donc bien, ce qui démontre le
lemme.
\end{proof}

\begin{proof}[Deuxième démonstration]
Cette démonstration repose sur Axiomes d'Artin, lemme
\ref{artin-lemma-get-openness-obstruction-theory}. Les conditions (1), (2) et
(3) de ce lemme correspondent aux lemmes \ref{quot-lemma-coherent-diagonal},
\ref{quot-lemma-coherent-RS-star} et \ref{quot-lemma-coherent-limits}.

\medskip\noindent
Nous avons construit une théorie de l'obstruction dans le chapitre sur la
théorie des déformations. Plus précisément, étant données une $S$-algèbre $A$
et un objet $x$ de $\Cohstack_{X/B}$ sur $\Spec(A)$, donné par
$\mathcal{F}$ sur $X_A$, nous posons
$\mathcal{O}_x(M) = \Ext^2_{X_A}(\mathcal{F}, \mathcal{F} \otimes_A M)$ ; si
$A' \to A$ est une surjection de noyau $I$, nous prenons pour élément
d'obstruction l'élément
$$
o_x(A') = o(\mathcal{F}, \mathcal{F} \otimes_A I, 1) \in
\mathcal{O}_x(I) = \Ext^2_{X_A}(\mathcal{F}, \mathcal{F} \otimes_A I)
$$
de Théorie des déformations, lemme \ref{defos-lemma-flat-ringed-topoi}.
Toutes les propriétés d'une théorie de l'obstruction telle qu'elle est définie
dans Axiomes d'Artin, définition
\ref{artin-definition-obstruction-theory}, résultent de ce lemme, sauf la
fonctorialité des classes d'obstruction formulée dans la condition (ii) de la
définition. Mais, comme l'indique la note de bas de page jointe à l'hypothèse
(4) d'Axiomes d'Artin, lemme
\ref{artin-lemma-get-openness-obstruction-theory}, il suffit de vérifier la
fonctorialité des classes d'obstruction pour $A$ fixé, ce qui résulte de
Théorie des déformations, lemme
\ref{defos-lemma-functorial-ringed-topoi}. Théorie des déformations, lemme
\ref{defos-lemma-flat-ringed-topoi}, nous dit aussi que
$T_x(M) = \Ext^1_{X_A}(\mathcal{F}, \mathcal{F} \otimes_A M)$ pour tout
$A$-module $M$.

\medskip\noindent
Pour achever la démonstration, il suffit de montrer que
$T_x(\prod M_n) = \prod T_x(M_n)$ et
$\mathcal{O}_x(\prod M_n) = \prod \mathcal{O}_x(M)$. Appliquons Catégories
dérivées des espaces, lemme \ref{spaces-perfect-lemma-compute-ext}, au calcul
des groupes Ext $\Ext^i_{X_A}(\mathcal{F}, \mathcal{F} \otimes_A M)$ pour
$i \leq m$, avec $m = 2$. Nous trouvons un objet parfait $K \in D(A)$ et des
isomorphismes fonctoriels
$$
H^i(K \otimes_A^\mathbf{L} M)
\longrightarrow
\Ext^i_{X_A}(\mathcal{F}, \mathcal{F} \otimes_A M)
$$
pour $i = 1, 2$. Un argument immédiat montre que
$$
H^i(K \otimes_A^\mathbf{L} \prod M_n) =
\prod H^i(K \otimes_A^\mathbf{L} M_n)
$$
dès que $K$ est un objet pseudo-cohérent de $D(A)$. En fait, cette propriété
(pour tout $i$) caractérise les complexes pseudo-cohérents ; voir Plus
d'algèbre, lemme \ref{more-algebra-lemma-pseudo-coherent-tensor}.
\end{proof}

\begin{theorem}[Algébricité du champ des faisceaux cohérents ; cas plat]
\label{quot-theorem-coherent-algebraic}
Soit $S$ un schéma. Soit $f : X \to B$ un morphisme d'espaces algébriques sur
$S$. Supposons que $f$ soit de présentation finie, séparé et
plat\footnote{Cette hypothèse n'est pas nécessaire. Voir la section
\ref{quot-section-not-flat}.}. Alors $\Cohstack_{X/B}$ est un champ algébrique sur
$S$.
\end{theorem}

\begin{proof}
Posons $\mathcal{X} = \Cohstack_{X/B}$. Nous avons vu que $\mathcal{X}$ est un
champ en groupoïdes sur $(\Sch/S)_{fppf}$, à diagonale représentable par des
espaces algébriques (lemmes \ref{quot-lemma-coherent-stack} et
\ref{quot-lemma-coherent-diagonal}). Il suffit donc de trouver un schéma $W$ et un
morphisme lisse surjectif $W \to \mathcal{X}$.

\medskip\noindent
Soit $B'$ un schéma et soit $B' \to B$ un morphisme \'etale surjectif. Posons
$X' = B' \times_B X$ et désignons par $f' : X' \to B'$ la projection. Alors
$\mathcal{X}' = \Cohstack_{X'/B'}$ est égal au $2$-produit fibré de
$\mathcal{X}$ et de la catégorie fibrée en ensembles associée à $B'$, sur la
catégorie fibrée en ensembles associée à $B$ (remarque
\ref{quot-remark-coherent-base-change}). D'après les résultats de Champs
algébriques, section \ref{algebraic-section-representable-properties}, le
morphisme $\mathcal{X}' \to \mathcal{X}$ est \'etale et surjectif. Il suffit donc
de démontrer le résultat pour $\mathcal{X}'$. Autrement dit, nous pouvons
supposer que $B$ est un schéma.

\medskip\noindent
Supposons que $B$ soit un schéma. Dans ce cas, nous pouvons remplacer $S$ par
$B$ ; voir Champs algébriques, section
\ref{algebraic-section-change-base-scheme}. Nous pouvons donc supposer
$S = B$.

\medskip\noindent
Supposons $S = B$. Choisissons un recouvrement ouvert affine
$S = \bigcup U_i$. Désignons par $\mathcal{X}_i$ la restriction de
$\mathcal{X}$ à $(\Sch/U_i)_{fppf}$. Si nous pouvons trouver des schémas $W_i$
sur $U_i$ et des morphismes lisses surjectifs
$W_i \to \mathcal{X}_i$, alors, en posant $W = \coprod W_i$, nous obtenons un
morphisme lisse surjectif $W \to \mathcal{X}$. Nous pouvons donc supposer que
$S = B$ est affine.

\medskip\noindent
Supposons $S = B$ affine, disons $S = \Spec(\Lambda)$. Écrivons
$\Lambda = \colim \Lambda_i$ comme une colimite filtrante, chaque $\Lambda_i$
étant de type fini sur $\mathbf{Z}$. Pour un certain $i$, nous pouvons trouver
un morphisme d'espaces algébriques $X_i \to \Spec(\Lambda_i)$ qui soit de
présentation finie, séparé et plat, et dont le changement de base à $\Lambda$
soit $X$. Voir Limites d'espaces, lemmes
\ref{spaces-limits-lemma-descend-finite-presentation},
\ref{spaces-limits-lemma-descend-separated-morphism} et
\ref{spaces-limits-lemma-descend-flat}. Si nous montrons que
$\Cohstack_{X_i/\Spec(\Lambda_i)}$ est un champ algébrique, il résultera par
changement de base (remarque \ref{quot-remark-coherent-base-change} et Champs
algébriques, section \ref{algebraic-section-change-base-scheme}) que
$\mathcal{X}$ est un champ algébrique. Nous pouvons donc supposer que
$\Lambda$ est une $\mathbf{Z}$-algèbre de type fini.

\medskip\noindent
Supposons $S = B = \Spec(\Lambda)$ affine de type fini sur $\mathbf{Z}$. Dans
ce cas, nous vérifierons les conditions (1), (2), (3), (4) et (5) d'Axiomes
d'Artin, lemme \ref{artin-lemma-diagonal-representable}, pour conclure que
$\mathcal{X}$ est un champ algébrique. Remarquons que $\Lambda$ est un
anneau G ; voir Plus d'algèbre, proposition
\ref{more-algebra-proposition-ubiquity-G-ring}. Ainsi, tous les anneaux locaux
de $S$ sont des anneaux G. La condition (5) est donc satisfaite. D'après le
lemme \ref{quot-lemma-coherent-defo-thy}, la versalité est ouverte pour
$\mathcal{X}$, donc (4) est satisfaite. Pour vérifier (2), nous devons vérifier
les axiomes [-1], [0], [1], [2] et [3] d'Axiomes d'Artin, section
\ref{artin-section-axioms}. Nous omettons la vérification de [-1], et les
axiomes [0], [1], [2], [3] correspondent respectivement aux lemmes
\ref{quot-lemma-coherent-stack}, \ref{quot-lemma-coherent-limits},
\ref{quot-lemma-coherent-RS-star}, \ref{quot-lemma-coherent-tangent-space}. La condition
(3) résulte du lemme \ref{quot-lemma-coherent-existence}. Enfin, la condition (1)
est le lemme \ref{quot-lemma-coherent-diagonal}. Ceci achève la démonstration du
théorème.
\end{proof}









\section{Le champ des faisceaux cohérents dans le cas non plat}
\label{quot-section-not-flat}

\noindent
Dans le théorème \ref{quot-theorem-coherent-algebraic}, l'hypothèse que
$f : X \to B$ est plat n'est pas nécessaire. Dans cette section, nous donnons
une démonstration différente qui évite l'hypothèse de platitude et la
vérification de l'ouverture de la versalité, en utilisant les résultats de
Platitude sur les espaces, section \ref{spaces-flat-section-existence}, et
d'Axiomes d'Artin, section
\ref{artin-section-strong-formal-effectiveness}.

\medskip\noindent
Pour une autre approche de ce problème, le lecteur pourra consulter
\cite{ArtinI} et suivre la méthode exposée dans les articles
\cite{olsson-starr}, \cite{lieblich_remarks}, \cite{olsson_proper},
\cite{Hall-Rydh}, \cite{Hall-Rydh-Hilbert}, \cite{rydh_representability}.
Certains de ces articles traitent du cas plus général du champ des faisceaux
cohérents sur un champ algébrique au-dessus d'un champ algébrique, et d'autres
traitent de problèmes analogues pour les champs de Hilbert ou les foncteurs de
Quot. Notre stratégie consistera à démontrer l'algébricité de certains champs
de Hilbert et foncteurs de Quot comme conséquence de l'algébricité du champ des
faisceaux cohérents.

\begin{theorem}[Algébricité du champ des faisceaux cohérents ; cas général]
\label{quot-theorem-coherent-algebraic-general}
Soit $S$ un schéma. Soit $f : X \to B$ un morphisme d'espaces algébriques sur
$S$. Supposons que $f$ soit de présentation finie et séparé. Alors
$\Cohstack_{X/B}$ est un champ algébrique sur $S$.
\end{theorem}

\begin{proof}
Seule la dernière étape diffère de la démonstration du cas plat, mais nous
répétons ici tous les arguments afin de nous assurer que tout fonctionne.

\medskip\noindent
Posons $\mathcal{X} = \Cohstack_{X/B}$. Nous avons vu que $\mathcal{X}$ est un
champ en groupoïdes sur $(\Sch/S)_{fppf}$, à diagonale représentable par des
espaces algébriques (lemmes \ref{quot-lemma-coherent-stack} et
\ref{quot-lemma-coherent-diagonal}). Il suffit donc de trouver un schéma $W$ et un
morphisme lisse surjectif $W \to \mathcal{X}$.

\medskip\noindent
Soit $B'$ un schéma et soit $B' \to B$ un morphisme \'etale surjectif. Posons
$X' = B' \times_B X$ et désignons par $f' : X' \to B'$ la projection. Alors
$\mathcal{X}' = \Cohstack_{X'/B'}$ est égal au $2$-produit fibré de
$\mathcal{X}$ et de la catégorie fibrée en ensembles associée à $B'$, sur la
catégorie fibrée en ensembles associée à $B$ (remarque
\ref{quot-remark-coherent-base-change}). D'après les résultats de Champs
algébriques, section \ref{algebraic-section-representable-properties}, le
morphisme $\mathcal{X}' \to \mathcal{X}$ est \'etale et surjectif. Il suffit donc
de démontrer le résultat pour $\mathcal{X}'$. Autrement dit, nous pouvons
supposer que $B$ est un schéma.

\medskip\noindent
Supposons que $B$ soit un schéma. Dans ce cas, nous pouvons remplacer $S$ par
$B$ ; voir Champs algébriques, section
\ref{algebraic-section-change-base-scheme}. Nous pouvons donc supposer
$S = B$.

\medskip\noindent
Supposons $S = B$. Choisissons un recouvrement ouvert affine
$S = \bigcup U_i$. Désignons par $\mathcal{X}_i$ la restriction de
$\mathcal{X}$ à $(\Sch/U_i)_{fppf}$. Si nous pouvons trouver des schémas $W_i$
sur $U_i$ et des morphismes lisses surjectifs
$W_i \to \mathcal{X}_i$, alors, en posant $W = \coprod W_i$, nous obtenons un
morphisme lisse surjectif $W \to \mathcal{X}$. Nous pouvons donc supposer que
$S = B$ est affine.

\medskip\noindent
Supposons $S = B$ affine, disons $S = \Spec(\Lambda)$. Écrivons
$\Lambda = \colim \Lambda_i$ comme une colimite filtrante, chaque $\Lambda_i$
étant de type fini sur $\mathbf{Z}$. Pour un certain $i$, nous pouvons trouver
un morphisme d'espaces algébriques $X_i \to \Spec(\Lambda_i)$ qui soit séparé
et de présentation finie et dont le changement de base à $\Lambda$ soit $X$.
Voir Limites d'espaces, lemmes
\ref{spaces-limits-lemma-descend-finite-presentation} et
\ref{spaces-limits-lemma-descend-separated-morphism}. Si nous montrons que
$\Cohstack_{X_i/\Spec(\Lambda_i)}$ est un champ algébrique, il résultera par
changement de base (remarque \ref{quot-remark-coherent-base-change} et Champs
algébriques, section \ref{algebraic-section-change-base-scheme}) que
$\mathcal{X}$ est un champ algébrique. Nous pouvons donc supposer que
$\Lambda$ est une $\mathbf{Z}$-algèbre de type fini.

\medskip\noindent
Supposons $S = B = \Spec(\Lambda)$ affine de type fini sur $\mathbf{Z}$. Dans
ce cas, nous vérifierons les conditions (1), (2), (3), (4) et (5) d'Axiomes
d'Artin, lemme \ref{artin-lemma-diagonal-representable}, pour conclure que
$\mathcal{X}$ est un champ algébrique. Remarquons que $\Lambda$ est un anneau
G ; voir Plus d'algèbre, proposition
\ref{more-algebra-proposition-ubiquity-G-ring}. Ainsi, tous les anneaux locaux
de $S$ sont des anneaux G. La condition (5) est donc satisfaite. Pour vérifier
(2), nous devons vérifier les axiomes [-1], [0], [1], [2] et [3] d'Axiomes
d'Artin, section \ref{artin-section-axioms}. Nous omettons la vérification de
[-1], et les axiomes [0], [1], [2], [3] correspondent respectivement aux
lemmes \ref{quot-lemma-coherent-stack}, \ref{quot-lemma-coherent-limits},
\ref{quot-lemma-coherent-RS-star}, \ref{quot-lemma-coherent-tangent-space}. La condition
(3) est le lemme \ref{quot-lemma-coherent-existence}. La condition (1) est le lemme
\ref{quot-lemma-coherent-diagonal}.

\medskip\noindent
Il reste à démontrer la condition (4), qui est l'ouverture de la versalité.
Pour cela, nous utiliserons Axiomes d'Artin, lemme
\ref{artin-lemma-SGE-implies-openness-versality}. Nous avons déjà vu que
$\mathcal{X}$ a sa diagonale représentable par des espaces algébriques,
possède (RS*) et commute aux limites (voir les lemmes utilisés ci-dessus). Il
nous reste donc seulement à voir que $\mathcal{X}$ satisfait à l'effectivité
formelle forte formulée dans Axiomes d'Artin, lemme
\ref{artin-lemma-SGE-implies-openness-versality}. C'est le théorème de
Platitude sur les espaces \ref{spaces-flat-theorem-existence}, et la
démonstration est achevée.
\end{proof}

\section{Le foncteur des quotients}
\label{quot-section-functor-quotients}

\noindent
Dans cette section, nous donnons quelques généralités sur le foncteur
$Q_{\mathcal{F}/X/B}$ défini ci-dessous. La notation
$\Quotfunctor_{\mathcal{F}/X/B}$ est réservée à un sous-foncteur de
$\text{Q}_{\mathcal{F}/X/B}$. Nous recommandons au lecteur d'omettre cette
section lors d'une première lecture.

\begin{situation}
\label{quot-situation-q}
Soit $S$ un schéma. Soit $f : X \to B$ un morphisme d'espaces algébriques sur
$S$. Soit $\mathcal{F}$ un $\mathcal{O}_X$-module quasi-cohérent. Pour tout
schéma $T$ sur $B$, nous désignerons par $X_T$ le changement de base de $X$ à
$T$ et par $\mathcal{F}_T$ l'image inverse de $\mathcal{F}$ par le morphisme
de projection $X_T = X \times_B T \to X$. Pour un tel $T$, posons
$$
\text{Q}_{\mathcal{F}/X/B}(T) =
\left\{
\begin{matrix}
\text{quotients }\mathcal{F}_T \to \mathcal{Q}\text{ où }
\mathcal{Q}\text{ est un}\\
\text{module quasi-cohérent sur }
\mathcal{O}_{X_T}\text{, plat sur }T
\end{matrix}
\right\}
$$
Nous identifions deux quotients s'ils ont le même noyau. Supposons que
$T' \to T$ soit un morphisme de schémas sur $B$ et que
$\mathcal{F}_T \to \mathcal{Q}$ soit un élément de
$\text{Q}_{\mathcal{F}/X/B}(T)$. Alors l'image inverse
$\mathcal{Q}' = (X_{T'} \to X_T)^*\mathcal{Q}$ est un
$\mathcal{O}_{X_{T'}}$-module quasi-cohérent plat sur $T'$, d'après Morphismes
d'espaces, lemme \ref{spaces-morphisms-lemma-base-change-module-flat}. Nous
obtenons donc un foncteur
\begin{equation}
\label{quot-equation-q}
\text{Q}_{\mathcal{F}/X/B} : (\Sch/B)^{opp} \longrightarrow \textit{Ensembles}
\end{equation}
C'est le {\it foncteur des quotients de $\mathcal{F}/X/B$}. Nous définissons
un sous-foncteur
\begin{equation}
\label{quot-equation-q-fp}
\text{Q}^{fp}_{\mathcal{F}/X/B} : (\Sch/B)^{opp} \longrightarrow \textit{Ensembles}
\end{equation}
qui associe à $T$ le sous-ensemble de $\text{Q}_{\mathcal{F}/X/B}(T)$ formé
des quotients $\mathcal{F}_T \to \mathcal{Q}$ tels que $\mathcal{Q}$ soit de
présentation finie comme $\mathcal{O}_{X_T}$-module. C'est un sous-foncteur
d'après Propriétés des espaces, section
\ref{spaces-properties-section-properties-modules}.
\end{situation}

\noindent
Dans la situation \ref{quot-situation-q}, nous regardons parfois
$\text{Q}_{\mathcal{F}/X/B}$ comme un foncteur
$(\Sch/S)^{opp} \to \textit{Ensembles}$ muni d'un morphisme
$\text{Q}_{\mathcal{F}/X/S} \to B$. En effet, si $T$ est un schéma sur $S$,
un élément de $\text{Q}_{\mathcal{F}/X/B}(T)$ est un couple
$(h, \mathcal{Q})$, où $h$ un morphisme $h : T \to B$ et où $\mathcal{Q}$
est un quotient $\mathcal{F}_T \to \mathcal{Q}$ plat sur $T$, de présentation
finie sur $X_T = X \times_{B, h} T$. En particulier, lorsque nous disons que
$\text{Q}_{\mathcal{F}/X/S}$ est un espace algébrique, nous entendons que le
foncteur correspondant $(\Sch/S)^{opp} \to \textit{Ensembles}$ est un espace
algébrique. Des remarques analogues valent pour
$\text{Q}^{fp}_{\mathcal{F}/X/B}$.

\begin{remark}
\label{quot-remark-q-base-change}
Dans la situation \ref{quot-situation-q}, soit $B' \to B$ un morphisme d'espaces
algébriques sur $S$. Posons $X' = X \times_B B'$ et désignons par
$\mathcal{F}'$ l'image inverse de $\mathcal{F}$ sur $X'$. Nous avons donc le
foncteur $Q_{\mathcal{F}'/X'/B'}$ sur la catégorie des schémas sur $B'$. Pour
un schéma $T$ sur $B'$, il est clair que
$$
Q_{\mathcal{F}'/X'/B'}(T) = Q_{\mathcal{F}/X/B}(T)
$$
où, dans le membre de droite, nous regardons $T$ comme un schéma sur $B$ au
moyen du composé $T \to B' \to B$. Des remarques analogues valent pour
$\text{Q}^{fp}_{\mathcal{F}/X/B}$. Ces remarques élémentaires seront parfois
utiles pour changer l'espace algébrique de base.
\end{remark}

\begin{remark}
\label{quot-remark-q-sheaf}
Soient $S$ un schéma, $X$ un espace algébrique sur $S$ et $\mathcal{F}$ un
$\mathcal{O}_X$-module quasi-cohérent. Supposons que
$\{f_i : X_i \to X\}_{i \in I}$ soit un recouvrement fpqc et que, pour chaque
$i, j \in I$, nous soit donné un recouvrement fpqc
$\{X_{ijk} \to X_i \times_X X_j\}$. Dans cette situation, nous avons une
bijection
$$
\left\{
\begin{matrix}
\text{quotients }\mathcal{F} \to \mathcal{Q}\text{ où } \\
\mathcal{Q}\text{ est quasi-cohérent}\\
\end{matrix}
\right\}
\longrightarrow
\left\{
\begin{matrix}
\text{familles de quotients }f_i^*\mathcal{F} \to \mathcal{Q}_i
\text{ où } \\
\mathcal{Q}_i\text{ est quasi-cohérent et où }
\mathcal{Q}_i\text{ et }\mathcal{Q}_j\\
\text{induisent le même quotient sur }X_{ijk}
\end{matrix}
\right\}
$$
En effet, soit $(f_i^*\mathcal{F} \to \mathcal{Q}_i)_{i \in I}$ un élément du
membre de droite. Comme $\{X_{ijk} \to X_i \times_X X_j\}$ est un
recouvrement fpqc, nous voyons que les images inverses de $\mathcal{Q}_i$ et
$\mathcal{Q}_j$ induisent le même quotient de l'image inverse de
$\mathcal{F}$ sur $X_i \times_X X_j$ (par pleine fidélité dans Descente sur
les espaces, proposition
\ref{spaces-descent-proposition-fpqc-descent-quasi-coherent}). Nous obtenons
donc une donnée de descente pour les modules quasi-cohérents relativement à
$\{X_i \to X\}_{i \in I}$. D'après Descente sur les espaces, proposition
\ref{spaces-descent-proposition-fpqc-descent-quasi-coherent}, nous trouvons un
homomorphisme de $\mathcal{O}_X$-modules quasi-cohérents
$\mathcal{F} \to \mathcal{Q}$ dont la restriction à $X_i$ redonne les
homomorphismes donnés $f_i^*\mathcal{F} \to \mathcal{Q}_i$. Comme la famille
de morphismes $\{X_i \to X\}$ est conjointement surjective et plate, pour
tout point $x \in |X|$, il existe un $i$ et un point $x_i \in |X_i|$
s'envoyant sur $x$. Remarquons que l'homomorphisme induit sur les anneaux
locaux
$\mathcal{O}_{X, \overline{x}} \to \mathcal{O}_{X_i, \overline{x_i}}$
est fidèlement plat ; voir Morphismes d'espaces, section
\ref{spaces-morphisms-section-flat}. Nous voyons donc que
$\mathcal{F} \to \mathcal{Q}$ est surjectif.
\end{remark}

\begin{lemma}
\label{quot-lemma-q-sheaf}
Dans la situation \ref{quot-situation-q}. Les foncteurs
$\text{Q}_{\mathcal{F}/X/B}$ et $\text{Q}^{fp}_{\mathcal{F}/X/B}$ satisfont à
la propriété de faisceau pour la topologie fpqc.
\end{lemma}

\begin{proof}
Soit $\{T_i \to T\}_{i \in I}$ un recouvrement fpqc de schémas sur $S$.
Posons $X_i = X_{T_i} = X \times_S T_i$ et
$\mathcal{F}_i = \mathcal{F}_{T_i}$. Remarquons que
$\{X_i \to X_T\}_{i \in I}$ est un recouvrement fpqc de $X_T$ (Topologies
sur les espaces, lemme \ref{spaces-topologies-lemma-fpqc}) et que
$X_{T_i \times_T T_{i'}} = X_i \times_{X_T} X_{i'}$. Supposons que
$\mathcal{F}_i \to \mathcal{Q}_i$ soit une famille d'éléments de
$\text{Q}_{\mathcal{F}/X/B}(T_i)$ telle que $\mathcal{Q}_i$ et
$\mathcal{Q}_{i'}$ induisent le même élément de
$\text{Q}_{\mathcal{F}/X/B}(T_i \times_T T_{i'})$. D'après la remarque
\ref{quot-remark-q-sheaf}, nous obtenons une application surjective de
$\mathcal{O}_{X_T}$-modules quasi-cohérents
$\mathcal{F}_T \to \mathcal{Q}$ dont la restriction à $X_i$ redonne les
quotients donnés. D'après Morphismes d'espaces, lemme
\ref{spaces-morphisms-lemma-flat-permanence}, nous voyons que $\mathcal{Q}$ est
plat sur $T$. Enfin, dans le cas de $\text{Q}^{fp}_{\mathcal{F}/X/B}$,
c'est-à-dire si les $\mathcal{Q}_i$ sont de présentation finie, Descente sur
les espaces, lemme \ref{spaces-descent-lemma-finite-presentation-descends},
garantit que $\mathcal{Q}$ est de présentation finie comme
$\mathcal{O}_{X_T}$-module.
\end{proof}

\noindent
Vérification de cohérence : $\text{Q}_{\mathcal{F}/X/B}$ et
$\text{Q}^{fp}_{\mathcal{F}/X/B}$ jouent le même rôle parmi les espaces
algébriques sur $S$.

\begin{lemma}
\label{quot-lemma-extend-q-to-spaces}
Dans la situation \ref{quot-situation-q}. Soit $T$ un espace algébrique sur $S$.
Nous avons
$$
\Mor_{\Sh((\Sch/S)_{fppf})}(T,  \text{Q}_{\mathcal{F}/X/B}) =
\left\{
\begin{matrix}
(h, \mathcal{F}_T \to \mathcal{Q}) \text{ où }
h : T \to B \text{ et où}\\
\mathcal{Q}\text{ est quasi-cohérent et plat sur }T
\end{matrix}
\right\}
$$
où $\mathcal{F}_T$ désigne l'image inverse de $\mathcal{F}$ sur l'espace
algébrique $X \times_{B, h} T$. De même, nous avons
$$
\Mor_{\Sh((\Sch/S)_{fppf})}(T,  \text{Q}^{fp}_{\mathcal{F}/X/B}) =
\left\{
\begin{matrix}
(h, \mathcal{F}_T \to \mathcal{Q}) \text{ où }
h : T \to B \text{ et où}\\
\mathcal{Q}\text{ est de présentation finie et plat sur }T
\end{matrix}
\right\}
$$
\end{lemma}

\begin{proof}
Choisissons un schéma $U$ et un morphisme \'etale surjectif $p : U \to T$.
Soit $R = U \times_T U$, de projections $t, s : R \to U$.

\medskip\noindent
Soit $v : T \to \text{Q}_{\mathcal{F}/X/B}$ une transformation naturelle.
Alors $v(p)$ correspond à un couple
$(h_U, \mathcal{F}_U \to \mathcal{Q}_U)$ sur $U$. Comme $v$ est une
transformation de foncteurs, les images inverses de
$(h_U, \mathcal{F}_U \to \mathcal{Q}_U)$ par $s$ et $t$ coïncident. Puisque
$T = U/R$ (Espaces, lemme \ref{spaces-lemma-space-presentation}), nous
obtenons un morphisme $h : T \to B$ tel que $h_U = h \circ p$. D'après
Descente sur les espaces, proposition
\ref{spaces-descent-proposition-fpqc-descent-quasi-coherent}, le quotient
$\mathcal{Q}_U$ descend en un quotient $\mathcal{F}_T \to \mathcal{Q}$ sur
$X_T$. Comme $U \to T$ est surjectif et plat, il résulte de Morphismes
d'espaces, lemme \ref{spaces-morphisms-lemma-flat-permanence}, que
$\mathcal{Q}$ est plat sur $T$.

\medskip\noindent
Réciproquement, soit $(h, \mathcal{F}_T \to \mathcal{Q})$ un couple sur $T$.
Nous obtenons alors une transformation naturelle
$v : T \to \text{Q}_{\mathcal{F}/X/B}$ en envoyant un morphisme
$a : T' \to T$, où $T'$ est un schéma, sur
$(h \circ a, \mathcal{F}_{T'} \to a^*\mathcal{Q})$. Nous omettons de vérifier
que cette construction et celle du paragraphe précédent sont inverses l'une
de l'autre.

\medskip\noindent
Dans le cas de $\text{Q}^{fp}_{\mathcal{F}/X/B}$, ajoutons ceci : étant donné
un morphisme $h : T \to B$, un faisceau quasi-cohérent sur $X_T$ est de
présentation finie comme $\mathcal{O}_{X_T}$-module si et seulement si son
image inverse sur $X_U$ est de présentation finie comme
$\mathcal{O}_{X_U}$-module. Cela résulte du fait que $X_U \to X_T$ est
surjectif et \'etale, et de Descente sur les espaces, lemme
\ref{spaces-descent-lemma-finite-presentation-descends}.
\end{proof}

\begin{lemma}
\label{quot-lemma-q-sheaf-in-X}
Dans la situation \ref{quot-situation-q}, soit $\{X_i \to X\}_{i \in I}$ un
recouvrement fpqc et, pour chaque $i, j \in I$, soit
$\{X_{ijk} \to X_i \times_X X_j\}$ un recouvrement fpqc. Désignons par
$\mathcal{F}_i$, resp.\ $\mathcal{F}_{ijk}$, l'image inverse de
$\mathcal{F}$ sur $X_i$, resp.\ sur $X_{ijk}$. Pour tout schéma $T$ sur $B$,
le diagramme
$$
\xymatrix{
Q_{\mathcal{F}/X/B}(T) \ar[r] &
\prod\nolimits_i
Q_{\mathcal{F}_i/X_i/B}(T)
\ar@<1ex>[r]^-{\text{pr}_0^*} \ar@<-1ex>[r]_-{\text{pr}_1^*}
&
\prod\nolimits_{i, j, k}
Q_{\mathcal{F}_{ijk}/X_{ijk}/B}(T)
}
$$
présente la première flèche comme l'égalisateur des deux autres. La même
assertion vaut pour le foncteur $\text{Q}^{fp}_{\mathcal{F}/X/B}$.
\end{lemma}

\begin{proof}
Soit $\mathcal{F}_{i, T} \to \mathcal{Q}_i$ un élément de l'égalisateur de
$\text{pr}_0^*$ et $\text{pr}_1^*$. D'après la remarque
\ref{quot-remark-q-sheaf}, nous obtenons une surjection
$\mathcal{F}_T \to \mathcal{Q}$ de $\mathcal{O}_{X_T}$-modules
quasi-cohérents, dont la restriction à $X_{i, T}$ redonne
$\mathcal{F}_i \to \mathcal{Q}_i$. D'après Morphismes d'espaces, lemme
\ref{spaces-morphisms-lemma-flat-permanence}, nous voyons que $\mathcal{Q}$ est
plat sur $T$, comme voulu. Dans le cas du foncteur
$\text{Q}^{fp}_{\mathcal{F}/X/B}$, c'est-à-dire si $\mathcal{Q}_i$ est de
présentation finie, $\mathcal{Q}$ est lui aussi de présentation finie d'après
Descente sur les espaces, lemme
\ref{spaces-descent-lemma-finite-presentation-descends}.
\end{proof}

\begin{lemma}
\label{quot-lemma-q-limit-preserving}
Dans la situation \ref{quot-situation-q}, supposons en outre que (a) $f$ soit
quasi-compact et quasi-séparé, et que (b) $\mathcal{F}$ soit de présentation
finie. Alors le foncteur $\text{Q}^{fp}_{\mathcal{F}/X/B}$ commute aux limites
au sens suivant : si $T = \lim T_i$ est une limite filtrante de schémas
affines sur $B$, alors
$\text{Q}^{fp}_{\mathcal{F}/X/B}(T) =
\colim \text{Q}^{fp}_{\mathcal{F}/X/B}(T_i)$.
\end{lemma}

\begin{proof}
Soit $T = \lim T_i$ comme dans l'énoncé du lemme. Choisissons $i_0 \in I$ et
remplaçons $I$ par $\{i \in I \mid i \geq i_0\}$. Nous pouvons poser
$B = S = T_{i_0}$ et remplacer $X$ par $X_{T_0}$ et $\mathcal{F}$ par son
image inverse sur $X_{T_0}$. Alors $X_T = \lim X_{T_i}$ ; voir Limites
d'espaces, lemme
\ref{spaces-limits-lemma-directed-inverse-system-has-limit}. Soit
$\mathcal{F}_T \to \mathcal{Q}$ un élément de
$\text{Q}^{fp}_{\mathcal{F}/X/B}(T)$. D'après Limites d'espaces, lemme
\ref{spaces-limits-lemma-descend-modules-finite-presentation}, il existe un
$i$ et une application $\mathcal{F}_{T_i} \to \mathcal{Q}_i$ de
$\mathcal{O}_{X_{T_i}}$-modules de présentation finie dont l'image inverse sur
$X_T$ est l'application quotient donnée.

\medskip\noindent
Il nous reste à vérifier que, quitte à augmenter $i$, l'application
$\mathcal{F}_{T_i} \to \mathcal{Q}_i$ est surjective et $\mathcal{Q}_i$ est
plat sur $T_i$. Pour cela, choisissons un schéma affine $U$ et un morphisme
\'etale surjectif $U \to X$ (voir Propriétés des espaces, lemme
\ref{spaces-properties-lemma-quasi-compact-affine-cover}). Nous pouvons
vérifier la surjectivité et la platitude sur $T_i$ après image inverse par le
recouvrement \'etale $U_{T_i} \to X_{T_i}$ (par définition). Nous sommes ainsi
ramenés au cas où $X = \Spec(B_0)$ est un schéma affine de présentation finie
sur $B = S = T_0 = \Spec(A_0)$. En écrivant $T_i = \Spec(A_i)$, puis
$T = \Spec(A)$ avec $A = \colim A_i$, nous sommes parvenus au problème
d'algèbre suivant. Soit $M_i \to N_i$ une application de
$B_0 \otimes_{A_0} A_i$-modules de présentation finie telle que
$M_i \otimes_{A_i} A \to N_i \otimes_{A_i} A$ soit surjective et que
$N_i \otimes_{A_i} A$ soit plat sur $A$. Montrer qu'il existe
$i' \geq i$ tel que
$M_i \otimes_{A_i} A_{i'} \to N_i \otimes_{A_i} A_{i'}$ soit surjective et
que $N_i \otimes_{A_i} A_{i'}$ soit plat sur $A$. La première assertion
résulte d'Algèbre, lemme
\ref{algebra-lemma-module-map-property-in-colimit}, et la seconde d'Algèbre,
lemme \ref{algebra-lemma-flat-finite-presentation-limit-flat}.
\end{proof}

\begin{lemma}
\label{quot-lemma-q-RS-star}
Dans la situation \ref{quot-situation-q}. Soit
$$
\xymatrix{
Z \ar[r] \ar[d] & Z' \ar[d] \\
Y \ar[r] & Y'
}
$$
une somme amalgamée dans la catégorie des schémas sur $B$, où
$Z \to Z'$ est un épaississement et $Z \to Y$ est affine ; voir Plus sur les
morphismes, lemme
\ref{more-morphisms-lemma-pushout-along-thickening}. Alors l'application
naturelle
$$
Q_{\mathcal{F}/X/B}(Y') \longrightarrow
Q_{\mathcal{F}/X/B}(Y) \times_{Q_{\mathcal{F}/X/B}(Z)} Q_{\mathcal{F}/X/B}(Z')
$$
est bijective. Si $X \to B$ est localement de présentation finie, la même
assertion vaut pour $Q^{fp}_{\mathcal{F}/X/B}$.
\end{lemma}

\begin{proof}
Construisons une application inverse. Supposons donc donnés
$\mathcal{F}_Y \to \mathcal{A}$,
$\mathcal{F}_{Z'} \to \mathcal{B}'$ et un isomorphisme
$\mathcal{A}|_{X_Z} \to \mathcal{B}'|_{X_Z}$ compatible aux surjections
données. Nous appliquons alors Sommes amalgamées d'espaces, lemme
\ref{spaces-pushouts-lemma-space-over-pushout-flat-modules}, pour obtenir sur
$X_{Y'}$ un module quasi-cohérent $\mathcal{A}'$ plat sur $Y'$. Comme ce
faisceau est construit comme un produit fibré (voir la démonstration du lemme
cité), il existe une application canonique
$\mathcal{F}_{Y'} \to \mathcal{A}'$. On peut voir que cette application est
surjective parce qu'elle se factorise en
$$
\begin{matrix}
\mathcal{F}_{Y'} \\
\downarrow \\
(X_Y \to X_{Y'})_*\mathcal{F}_Y
\times_{(X_Z \to X_{Y'})_*\mathcal{F}_Z}
(X_{Z'} \to X_{Y'})_*\mathcal{F}_{Z'} \\
\downarrow \\
\mathcal{A}' =
(X_Y \to X_{Y'})_*\mathcal{A}
\times_{(X_Z \to X_{Y'})_*\mathcal{A}|_{X_Z}}
(X_{Z'} \to X_{Y'})_*\mathcal{B}'
\end{matrix}
$$
la première flèche étant surjective d'après Plus d'algèbre, lemme
\ref{more-algebra-lemma-module-over-fibre-product-bis}, et la seconde d'après
Plus d'algèbre, lemme
\ref{more-algebra-lemma-surjection-module-over-fibre-product}.

\medskip\noindent
Dans le cas de $Q^{fp}_{\mathcal{F}/X/B}$, il nous suffit de montrer que la
construction ci-dessus produit un module de présentation finie. Cela est
expliqué dans Plus d'algèbre, remarque
\ref{more-algebra-remark-relative-modules-over-fibre-product}, dans le cadre de
l'algèbre commutative. Le présent cas des modules sur des espaces algébriques
s'en déduit par localisation \'etale.
\end{proof}

\begin{remark}[Obstructions aux quotients]
\label{quot-remark-q-obs}
Dans la situation \ref{quot-situation-q}, {\bf supposons} que $\mathcal{F}$ soit
plat sur $B$. Soit $T \subset T'$ un épaississement du premier ordre de
schémas sur $B$, de faisceau d'idéaux $\mathcal{J}$. Alors
$X_T \subset X_{T'}$ est un épaississement du premier ordre d'espaces
algébriques, dont le faisceau d'idéaux $\mathcal{I}$ est un quotient de
$f_T^*\mathcal{J}$. Nous considérerons les faisceaux sur $X_{T'}$, resp.\ sur
$T'$, comme des faisceaux sur $X_T$, resp.\ sur $T$, au moyen de l'équivalence
fondamentale décrite dans Plus sur les morphismes d'espaces, section
\ref{spaces-more-morphisms-section-thickenings}. Soit
$$
0 \to \mathcal{K} \to \mathcal{F}_T \to \mathcal{Q} \to 0
$$
une suite définissant un élément $x$ de $Q_{\mathcal{F}/X/B}(T)$. Comme
$\mathcal{F}_{T'}$ est plat sur $T'$, nous avons une suite exacte courte
$$
0 \to f_T^*\mathcal{J} \otimes_{\mathcal{O}_{X_T}} \mathcal{F}_T
\xrightarrow{i} \mathcal{F}_{T'} \xrightarrow{\pi} \mathcal{F}_T \to 0
$$
et nous avons
$f_T^*\mathcal{J} \otimes_{\mathcal{O}_{X_T}} \mathcal{F}_T =
\mathcal{I} \otimes_{\mathcal{O}_{X_T}} \mathcal{F}_T$ ; voir Théorie des
déformations, lemme \ref{defos-lemma-deform-module-ringed-topoi}. Pour un
$\mathcal{O}_{X_T}$-module $\mathcal{G}$, nous emploierons l'abréviation
$
f_T^*\mathcal{J} \otimes_{\mathcal{O}_{X_T}} \mathcal{G} =
\mathcal{G} \otimes_{\mathcal{O}_T} \mathcal{J}
$.
Comme $\mathcal{Q}$ est plat sur $T$, nous obtenons une suite exacte courte
$$
0 \to
\mathcal{K} \otimes_{\mathcal{O}_T} \mathcal{J} \to
\mathcal{F}_T \otimes_{\mathcal{O}_T} \mathcal{J} \to
\mathcal{Q} \otimes_{\mathcal{O}_T} \mathcal{J} \to
\to 0
$$
En combinant ce qui précède, nous obtenons une extension canonique
$$
0 \to \mathcal{Q} \otimes_{\mathcal{O}_T} \mathcal{J} \to
\pi^{-1}(\mathcal{K})/i(\mathcal{K} \otimes_{\mathcal{O}_T} \mathcal{J}) \to
\mathcal{K} \to 0
$$
de $\mathcal{O}_{X_T}$-modules. Elle définit une classe canonique
$$
o_x(T') \in
\Ext^1_{\mathcal{O}_{X_T}}(\mathcal{K},
\mathcal{Q} \otimes_{\mathcal{O}_T} \mathcal{J})
$$
Si $o_x(T')$ est nul, nous obtenons une scission de la suite exacte courte qui
le définit, autrement dit un $\mathcal{O}_{X_{T'}}$-sous-module
$\mathcal{K}' \subset \pi^{-1}(\mathcal{K})$ s'insérant dans une suite exacte
courte
$0 \to \mathcal{K} \otimes_{\mathcal{O}_T} \mathcal{J} \to
\mathcal{K}' \to \mathcal{K} \to 0$. Il résulte alors du lemme cité ci-dessus
que $\mathcal{Q}' = \mathcal{F}_{T'}/\mathcal{K}'$ est un relèvement de $x$
en un élément de $Q_{\mathcal{F}/X/B}(T')$. Réciproquement, le lecteur voit
que l'existence d'un relèvement implique que $o_x(T')$ est nul. De plus, si
$x \in Q_{\mathcal{F}/X/B}^{fp}(T)$, alors automatiquement
$x' \in Q_{\mathcal{F}/X/B}^{fp}(T')$, d'après Théorie des déformations,
lemme \ref{defos-lemma-deform-fp-module-ringed-topoi}. Si nous avons un jour
besoin de cette remarque, nous en ferons un lemme, formulerons précisément le
résultat et en donnerons une démonstration détaillée (en fait, tout ce qui
précède vaut dans le cadre de topos annelés arbitraires).
\end{remark}

\begin{remark}[Déformations des quotients]
\label{quot-remark-q-defos}
Dans la situation \ref{quot-situation-q}, {\bf supposons} que $\mathcal{F}$ soit
plat sur $B$. Nous poursuivons la discussion de la remarque
\ref{quot-remark-q-obs}. Supposons $o_x(T') = 0$. Nous affirmons alors que
l'ensemble des relèvements $x' \in Q_{\mathcal{F}/X/B}(T')$ est un espace
principal homogène sous le groupe
$$
\Hom_{\mathcal{O}_{X_T}}(\mathcal{K},
\mathcal{Q} \otimes_{\mathcal{O}_T} \mathcal{J})
$$
En effet, étant donné un quotient
$\mathcal{F}_{T'} \to \mathcal{Q}'$ plat sur $T'$ qui relève le quotient
$\mathcal{Q}$, nous obtenons un diagramme commutatif à lignes et colonnes
exactes
$$
\xymatrix{
& 0 \ar[d] & 0 \ar[d] & 0 \ar[d] \\
0 \ar[r] &
\mathcal{K} \otimes \mathcal{J} \ar[r] \ar[d] &
\mathcal{F}_T \otimes \mathcal{J} \ar[r] \ar[d] &
\mathcal{Q} \otimes \mathcal{J} \ar[r] \ar[d] &
0 \\
0 \ar[r] &
\mathcal{K}' \ar[r] \ar[d] &
\mathcal{F}_{T'} \ar[r] \ar[d] &
\mathcal{Q}' \ar[r] \ar[d] &
0 \\
0 \ar[r] &
\mathcal{K} \ar[d] \ar[r] &
\mathcal{F}_T \ar[d] \ar[r] &
\mathcal{Q} \ar[d] \ar[r] &
0 \\
& 0 & 0 & 0
}
$$
(pour le voir, utiliser les observations faites dans la remarque précédente).
Étant donnée une application
$\varphi : \mathcal{K} \to \mathcal{Q} \otimes \mathcal{J}$, nous pouvons
considérer le sous-faisceau
$\mathcal{K}'_\varphi \subset \mathcal{F}_{T'}$ formé des sections locales
$s$ dont l'image dans $\mathcal{F}_T$ est une section locale $k$ de
$\mathcal{K}$ et dont l'image dans $\mathcal{Q}'$ est la section locale
$\varphi(k)$ de $\mathcal{Q} \otimes \mathcal{J}$. Posons alors
$\mathcal{Q}'_\varphi = \mathcal{F}_{T'}/\mathcal{K}'_\varphi$.
Réciproquement, tout second relèvement de $x$ correspond à l'un des quotients
ainsi construits. Si nous avons un jour besoin de cette remarque, nous en
ferons un lemme, formulerons précisément le résultat et en donnerons une
démonstration détaillée (en fait, tout ce qui précède vaut dans le cadre de
topos annelés arbitraires).
\end{remark}










\section{Le foncteur de Quot}
\label{quot-section-quot}

\noindent
Dans cette section, nous démontrons que le foncteur de Quot est un espace
algébrique.

\begin{situation}
\label{quot-situation-quot}
Soit $S$ un schéma. Soit $f : X \to B$ un morphisme d'espaces algébriques sur
$S$. Supposons que $f$ soit de présentation finie. Soit $\mathcal{F}$ un
$\mathcal{O}_X$-module quasi-cohérent. Pour tout schéma $T$ sur $B$, nous
désignerons par $X_T$ le changement de base de $X$ à $T$ et par
$\mathcal{F}_T$ l'image inverse de $\mathcal{F}$ par le morphisme de
projection $X_T = X \times_S T \to X$. Pour un tel $T$, posons
$$
\Quotfunctor_{\mathcal{F}/X/B}(T) =
\left\{
\begin{matrix}
\text{quotients }\mathcal{F}_T \to \mathcal{Q}\text{ où }
\mathcal{Q}\text{ est un module quasi-cohérent sur }\\
\mathcal{O}_{X_T}\text{, de présentation finie et plat sur }T\\
\text{à support propre sur }T
\end{matrix}
\right\}
$$
D'après Catégories dérivées des espaces, lemme
\ref{spaces-perfect-lemma-base-change-module-support-proper-over-base}, c'est
un sous-foncteur du foncteur $Q^{fp}_{\mathcal{F}/X/B}$ étudié dans la section
\ref{quot-section-functor-quotients}. Nous obtenons donc un foncteur
\begin{equation}
\label{quot-equation-quot}
\Quotfunctor_{\mathcal{F}/X/B} : (\Sch/B)^{opp} \longrightarrow \textit{Ensembles}
\end{equation}
C'est le {\it foncteur de Quot} associé à $\mathcal{F}/X/B$.
\end{situation}

\noindent
Dans la situation \ref{quot-situation-quot}, nous regardons parfois
$\Quotfunctor_{\mathcal{F}/X/B}$ comme un foncteur
$(\Sch/S)^{opp} \to \textit{Ensembles}$ muni d'un morphisme
$\Quotfunctor_{\mathcal{F}/X/B} \to B$. En effet, si $T$ est un schéma sur
$S$, un élément de $\Quotfunctor_{\mathcal{F}/X/B}(T)$ est un couple
$(h, \mathcal{Q})$, où $h$ est un morphisme $h : T \to B$ et où $Q$ est un
quotient $\mathcal{F}_T \to \mathcal{Q}$ de présentation finie et plat sur
$T$, sur $X_T = X \times_{B, h} T$, à support propre sur $T$. En particulier,
lorsque nous disons que $\Quotfunctor_{\mathcal{F}/X/B}$ est un espace
algébrique, nous entendons que le foncteur correspondant
$(\Sch/S)^{opp} \to \textit{Ensembles}$ est un espace algébrique.

\begin{lemma}
\label{quot-lemma-quot-sheaf}
Dans la situation \ref{quot-situation-quot}. Le foncteur
$\Quotfunctor_{\mathcal{F}/X/B}$ satisfait à la propriété de faisceau pour la
topologie fpqc.
\end{lemma}

\begin{proof}
Dans le lemme \ref{quot-lemma-q-sheaf}, nous avons vu que le foncteur
$\text{Q}^{fp}_{\mathcal{F}/X/S}$ est un faisceau. Rappelons que, pour un
schéma $T$ sur $S$, le sous-ensemble
$\Quotfunctor_{\mathcal{F}/X/S}(T) \subset \text{Q}_{\mathcal{F}/X/S}(T)$
sélectionne les quotients dont le support est propre sur $T$. Cela définit un
sous-faisceau d'après Descente sur les espaces, lemme
\ref{spaces-descent-lemma-descending-property-proper}, combiné à Morphismes
d'espaces, lemme \ref{spaces-morphisms-lemma-flat-pullback-support}, qui montre
que la formation du support schématique commute au changement de base plat.
\end{proof}

\noindent
Vérification de cohérence : $\Quotfunctor_{\mathcal{F}/X/B}$ joue le même rôle
parmi les espaces algébriques sur $S$.

\begin{lemma}
\label{quot-lemma-extend-quot-to-spaces}
Dans la situation \ref{quot-situation-quot}. Soit $T$ un espace algébrique sur $S$.
Nous avons
$$
\Mor_{\Sh((\Sch/S)_{fppf})}(T,  \Quotfunctor_{\mathcal{F}/X/B}) =
\left\{
\begin{matrix}
(h, \mathcal{F}_T \to \mathcal{Q}) \text{ où }
h : T \to B \text{ et où}\\
\mathcal{Q}\text{ est de présentation finie et}\\
\text{plate sur }T\text{ à support propre sur }T
\end{matrix}
\right\}
$$
où $\mathcal{F}_T$ désigne l'image inverse de $\mathcal{F}$ sur l'espace
algébrique $X \times_{B, h} T$.
\end{lemma}

\begin{proof}
Observons que les membres de gauche et de droite de l'égalité sont des
sous-ensembles des membres de gauche et de droite de la seconde égalité du
lemme \ref{quot-lemma-extend-q-to-spaces}. Pour voir que ces sous-ensembles se
correspondent par l'identification donnée dans la démonstration de ce lemme, il
suffit de montrer ceci : étant donnés $h : T \to B$, un morphisme \'etale
surjectif $U \to T$ et un $\mathcal{O}_{X_T}$-module quasi-cohérent de type
fini $\mathcal{Q}$, les conditions suivantes sont équivalentes :
\begin{enumerate}
\item le support schématique de $\mathcal{Q}$ est propre sur $T$, et
\item le support schématique de $(X_U \to X_T)^*\mathcal{Q}$ est propre sur
$U$.
\end{enumerate}
Cela résulte de Descente sur les espaces, lemme
\ref{spaces-descent-lemma-descending-property-proper}, combiné à Morphismes
d'espaces, lemme \ref{spaces-morphisms-lemma-flat-pullback-support}, qui montre
que la formation du support schématique commute au changement de base plat.
\end{proof}

\begin{proposition}
\label{quot-proposition-quot}
Soit $S$ un schéma. Soit $f : X \to B$ un morphisme d'espaces algébriques sur
$S$. Soit $\mathcal{F}$ un faisceau quasi-cohérent sur $X$. Si $f$ est de
présentation finie et séparé, alors
$\Quotfunctor_{\mathcal{F}/X/B}$ est un espace algébrique. Si $\mathcal{F}$
est de présentation finie, alors
$\Quotfunctor_{\mathcal{F}/X/B} \to B$ est localement de présentation finie.
\end{proposition}

\begin{proof}
D'après le lemme \ref{quot-lemma-quot-sheaf},
$\Quotfunctor_{\mathcal{F}/X/B}$ est un faisceau pour la topologie fppf. Soit
$\textit{Quot}_{\mathcal{F}/X/B}$ le champ en groupoïdes correspondant à
$\Quotfunctor_{\mathcal{F}/X/S}$ ; voir Champs algébriques, section
\ref{algebraic-section-split}. D'après Champs algébriques, proposition
\ref{algebraic-proposition-algebraic-stack-no-automorphisms}, il suffit de
montrer que $\textit{Quot}_{\mathcal{F}/X/B}$ est un champ algébrique.
Considérons le $1$-morphisme de champs en groupoïdes
$$
\textit{Quot}_{\mathcal{F}/X/S}
\longrightarrow
\Cohstack_{X/B}
$$
sur $(\Sch/S)_{fppf}$, qui associe au quotient
$\mathcal{F}_T \to \mathcal{Q}$ le module $\mathcal{Q}$. D'après le théorème
\ref{quot-theorem-coherent-algebraic-general}, nous savons que
$\Cohstack_{X/B}$ est un champ algébrique. D'après Champs algébriques, lemme
\ref{algebraic-lemma-representable-morphism-to-algebraic}, il suffit de montrer
que ce $1$-morphisme est représentable par des espaces algébriques.

\medskip\noindent
Soit $T$ un schéma sur $S$ et soit l'objet $(h, \mathcal{G})$ de
$\Cohstack_{X/B}$ sur $T$ qui correspond à un $1$-morphisme
$\xi : (\Sch/T)_{fppf} \to \Cohstack_{X/B}$. Le $2$-produit fibré
$$
\mathcal{Z} =
(\Sch/T)_{fppf}
\times_{\xi, \Cohstack_{X/B}}
\textit{Quot}_{\mathcal{F}/X/S}
$$
est un champ en sétoïdes ; voir Champs, lemme
\ref{stacks-lemma-2-fibre-product-gives-stack-in-setoids}. Le faisceau
d'ensembles correspondant (c'est-à-dire le foncteur ; voir Champs, lemmes
\ref{stacks-lemma-2-fibre-product-gives-stack-in-setoids} et
\ref{stacks-lemma-when-stack-in-sets}) associe à un schéma $T'/T$ l'ensemble
des surjections $u : \mathcal{F}_{T'} \to \mathcal{G}_{T'}$ de modules
quasi-cohérents sur $X_{T'}$. Nous voyons donc que $\mathcal{Z}$ est
représentable par un sous-espace ouvert (d'après Platitude sur les espaces,
lemme \ref{spaces-flat-lemma-F-surj-open}) de l'espace algébrique
$\mathit{Hom}(\mathcal{F}_T, \mathcal{G})$ de la proposition
\ref{quot-proposition-hom}.
\end{proof}

\begin{remark}[Quot au moyen des axiomes d'Artin]
\label{quot-remark-quot-via-artins-axioms}
Soit $S$ un schéma noethérien dont tous les anneaux locaux sont des anneaux G.
Soit $X$ un espace algébrique sur $S$ dont le morphisme structural
$f : X \to S$ est de présentation finie et séparé. Soit $\mathcal{F}$ un
faisceau quasi-cohérent de présentation finie sur $X$, plat sur $S$. Nous
esquissons dans cette remarque comment employer les axiomes d'Artin pour
démontrer que $\Quotfunctor_{\mathcal{F}/X/S}$ est un espace algébrique
localement de présentation finie sur $S$, sans utiliser l'algébricité du champ
des faisceaux cohérents comme dans la démonstration de la proposition
\ref{quot-proposition-quot}.

\medskip\noindent
Vérifions les conditions énumérées dans Axiomes d'Artin, proposition
\ref{artin-proposition-spaces-diagonal-representable}. On peut voir comme suit
la représentabilité de la diagonale de
$\Quotfunctor_{\mathcal{F}/X/S}$ : supposons donnés deux quotients
$\mathcal{F}_T \to \mathcal{Q}_i$, $i = 1, 2$. Désignons par
$\mathcal{K}_1$ le noyau du premier. Il faut alors montrer que le lieu de $T$
au-dessus duquel $u : \mathcal{K}_1 \to \mathcal{Q}_2$ devient nul est
représentable. Cela résulte par exemple de Platitude sur les espaces, lemme
\ref{spaces-flat-lemma-F-zero-closed-proper}, ou d'une discussion antérieure
du faisceau $\mathit{Hom}$ dans ce chapitre. Les axiomes [0] (faisceau), [1]
(limites), [2] (Rim--Schlessinger) résultent des lemmes
\ref{quot-lemma-quot-sheaf}, \ref{quot-lemma-q-limit-preserving} et
\ref{quot-lemma-q-RS-star} (moyennant un travail supplémentaire pour traiter la
condition de propreté). L'axiome [3] (dimension finie des espaces tangents)
résulte de la description des déformations infinitésimales dans la remarque
\ref{quot-remark-q-defos} et de la finitude de la cohomologie des faisceaux
cohérents sur les espaces algébriques propres sur des corps (Cohomologie des
espaces, lemme \ref{spaces-cohomology-lemma-proper-pushforward-coherent}).
L'axiome [4] (effectivité des objets formels) résulte du théorème d'existence
de Grothendieck (Plus sur les morphismes d'espaces, théorème
\ref{spaces-more-morphisms-theorem-grothendieck-existence}). Comme toujours,
l'axiome [5] (ouverture de la versalité) est le plus délicat à vérifier. On
peut par exemple utiliser la théorie de l'obstruction décrite dans la remarque
\ref{quot-remark-q-obs} et la description des déformations de la remarque
\ref{quot-remark-q-defos}, en appliquant le critère d'Axiomes d'Artin, lemme
\ref{artin-lemma-get-openness-obstruction-theory}. Comparer avec la seconde
démonstration du lemme \ref{quot-lemma-coherent-defo-thy}.
\end{remark}









\section{Le foncteur de Hilbert}
\label{quot-section-hilb}

\noindent
Dans cette section, nous démontrons que le foncteur Hilb est un espace
algébrique.

\begin{situation}
\label{quot-situation-hilb}
Soit $S$ un schéma. Soit $f : X \to B$ un morphisme d'espaces algébriques sur
$S$. Supposons que $f$ soit de présentation finie. Pour tout schéma $T$ sur
$B$, nous désignerons par $X_T$ le changement de base de $X$ à $T$. Pour un
tel $T$, posons
$$
\Hilbfunctor_{X/B}(T) =
\left\{
\begin{matrix}
\text{sous-espaces fermés }Z \subset X_T\text{ tels que }Z \to T\\
\text{soit de présentation finie, plat et propre}
\end{matrix}
\right\}
$$
Comme le changement de base conserve les propriétés requises (Espaces, lemme
\ref{spaces-lemma-base-change-immersions}, et Morphismes d'espaces, lemmes
\ref{spaces-morphisms-lemma-base-change-finite-presentation},
\ref{spaces-morphisms-lemma-base-change-flat} et
\ref{spaces-morphisms-lemma-base-change-proper}), nous obtenons un foncteur
\begin{equation}
\label{quot-equation-hilb}
\Hilbfunctor_{X/B} : (\Sch/B)^{opp} \longrightarrow \textit{Ensembles}
\end{equation}
C'est le {\it foncteur de Hilbert} associé à $X/B$.
\end{situation}

\noindent
Dans la situation \ref{quot-situation-hilb}, nous regardons parfois
$\Hilbfunctor_{X/B}$ comme un foncteur
$(\Sch/S)^{opp} \to \textit{Ensembles}$ muni d'un morphisme
$\Hilbfunctor_{X/S} \to B$. En effet, si $T$ est un schéma sur $S$, un
élément de $\Hilbfunctor_{X/B}(T)$ est un couple $(h, Z)$ où $h$ est un
morphisme $h : T \to B$ et où $Z \subset X_T = X \times_{B, h} T$ est un
sous-schéma fermé, plat, propre et de présentation finie sur $T$. En
particulier, lorsque nous disons que $\Hilbfunctor_{X/B}$ est un espace
algébrique, nous entendons que le foncteur correspondant
$(\Sch/S)^{opp} \to \textit{Ensembles}$ est un espace algébrique.

\medskip\noindent
Bien entendu, le foncteur de Hilbert n'est qu'un cas particulier du foncteur
de Quot.

\begin{lemma}
\label{quot-lemma-hilb-is-quot}
Dans la situation \ref{quot-situation-hilb}, nous avons
$\Hilbfunctor_{X/B} = \Quotfunctor_{\mathcal{O}_X/X/B}$.
\end{lemma}

\begin{proof}
Soit $T$ un schéma sur $B$. Étant donné un élément
$Z \in \Hilbfunctor_{X/B}(T)$, nous pouvons considérer le quotient
$\mathcal{O}_{X_T} \to i_*\mathcal{O}_Z$, où $i : Z \to X_T$ est le
morphisme d'inclusion. Remarquons que $i_*\mathcal{O}_Z$ est quasi-cohérent.
Comme $Z \to T$ et $X_T \to T$ sont de présentation finie, nous voyons que
$i$ est de présentation finie (Morphismes d'espaces, lemme
\ref{spaces-morphisms-lemma-finite-presentation-permanence}) ; ainsi,
$i_*\mathcal{O}_Z$ est un $\mathcal{O}_{X_T}$-module de présentation finie
(Descente sur les espaces, lemme
\ref{spaces-descent-lemma-finite-finitely-presented-module}). Comme
$Z \to T$ est propre, nous voyons que $i_*\mathcal{O}_Z$ est à support propre
sur $T$ (au sens défini dans Catégories dérivées des espaces, section
\ref{spaces-perfect-section-proper-over-base}). Comme $\mathcal{O}_Z$ est plat
sur $T$ et que $i$ est affine, nous voyons que $i_*\mathcal{O}_Z$ est plat sur
$T$ (nous omettons un petit argument). Par conséquent,
$\mathcal{O}_{X_T} \to i_*\mathcal{O}_Z$ est un élément de
$\Quotfunctor_{\mathcal{O}_X/X/B}(T)$.

\medskip\noindent
Réciproquement, étant donné un élément
$\mathcal{O}_{X_T} \to \mathcal{Q}$ de
$\Quotfunctor_{\mathcal{O}_X/X/B}(T)$, nous pouvons considérer l'immersion
fermée $i : Z \to X_T$ correspondant au faisceau d'idéaux quasi-cohérent
$\mathcal{I} = \Ker(\mathcal{O}_{X_T} \to \mathcal{Q})$ (Morphismes
d'espaces, lemme \ref{spaces-morphisms-lemma-closed-immersion-ideals}). Par
construction de $Z$, nous avons $\mathcal{Q} = i_*\mathcal{O}_Z$. Nous pouvons
alors lire les arguments précédents à rebours pour voir que $Z$ définit un
élément de $\Hilbfunctor_{X/B}(T)$. Par exemple, $\mathcal{I}$ est
quasi-cohérent de type fini (Modules sur les sites, lemme
\ref{sites-modules-lemma-kernel-surjection-finite-onto-finite-presentation}) ;
donc $i : Z \to X_T$ est de présentation finie (Morphismes d'espaces, lemme
\ref{spaces-morphisms-lemma-closed-immersion-finite-presentation}) ; donc
$Z \to T$ est de présentation finie (Morphismes d'espaces, lemme
\ref{spaces-morphisms-lemma-composition-finite-presentation}). La propreté de
$Z \to T$ résulte de la discussion dans Catégories dérivées des espaces,
section \ref{spaces-perfect-section-proper-over-base}. La platitude de
$Z \to T$ résulte de la platitude de $\mathcal{Q}$ sur $T$.

\medskip\noindent
Nous omettons la vérification (immédiate) que les deux constructions ci-dessus
sont inverses l'une de l'autre.
\end{proof}

\noindent
Vérification de cohérence : le faisceau $\Hilbfunctor_{X/B}$ joue le même rôle
parmi les espaces algébriques sur $S$.

\begin{lemma}
\label{quot-lemma-extend-hilb-to-spaces}
Dans la situation \ref{quot-situation-hilb}. Soit $T$ un espace algébrique sur $S$.
Nous avons
$$
\Mor_{\Sh((\Sch/S)_{fppf})}(T, \Hilbfunctor_{X/B}) =
\left\{
\begin{matrix}
(h, Z)\text{ où }h : T \to B,\ Z \subset X_T \\
\text{de présentation finie, plat, propre sur }T
\end{matrix}
\right\}
$$
où $X_T = X \times_{B, h} T$.
\end{lemma}

\begin{proof}
D'après le lemme \ref{quot-lemma-hilb-is-quot}, nous avons
$\Hilbfunctor_{X/B} = \Quotfunctor_{\mathcal{O}_X/X/B}$. Nous pouvons donc
appliquer le lemme \ref{quot-lemma-extend-quot-to-spaces} pour voir que le membre de
gauche est en bijection avec l'ensemble des surjections
$\mathcal{O}_{X_T} \to \mathcal{Q}$ qui sont de présentation finie, plates
sur $T$ et à support propre sur $T$. En raisonnant exactement comme dans la
démonstration du lemme \ref{quot-lemma-hilb-is-quot}, nous voyons que ces quotients
correspondent exactement aux immersions fermées $Z \to X_T$ telles que
$Z \to T$ soit propre, plat et de présentation finie.
\end{proof}

\begin{proposition}
\label{quot-proposition-hilb}
Soit $S$ un schéma. Soit $f : X \to B$ un morphisme d'espaces algébriques sur
$S$. Si $f$ est de présentation finie et séparé, alors
$\Hilbfunctor_{X/B}$ est un espace algébrique localement de présentation finie
sur $B$.
\end{proposition}

\begin{proof}
Conséquence immédiate du lemme \ref{quot-lemma-hilb-is-quot} et de la proposition
\ref{quot-proposition-quot}.
\end{proof}






\section{Le champ de Picard}
\label{quot-section-picard-stack}

\noindent
Le champ de Picard d'un morphisme d'espaces algébriques a été introduit dans
Exemples de champs, section \ref{examples-stacks-section-picard-stack}. Nous
déduirons du lemme suivant que c'est un sous-champ ouvert du champ des faisceaux
cohérents (dans les cas favorables).

\begin{lemma}
\label{quot-lemma-picard-stack-open-in-coh}
Soit $S$ un schéma. Soit $f : X \to B$ un morphisme d'espaces algébriques sur
$S$ qui est plat, de présentation finie et propre. L'application naturelle
$$
\Picardstack_{X/B} \longrightarrow \Cohstack_{X/B}
$$
est représentable par des immersions ouvertes.
\end{lemma}

\begin{proof}
Observons que l'application envoie simplement un triplet
$(T, g, \mathcal{L})$ comme dans Exemples de champs, section
\ref{examples-stacks-section-picard-stack}, sur le même triplet
$(T, g, \mathcal{L})$, mais considéré maintenant comme un triplet du genre
décrit dans la situation \ref{quot-situation-coherent}. Cela fonctionne parce que
le $\mathcal{O}_{X_T}$-module inversible $\mathcal{L}$ est certainement un
$\mathcal{O}_{X_T}$-module de présentation finie, qu'il est plat sur $T$ parce
que $X_T \to T$ est plat, et que son support est propre sur $T$ puisque
$X_T \to T$ est propre (Morphismes d'espaces, lemmes
\ref{spaces-morphisms-lemma-base-change-flat} et
\ref{spaces-morphisms-lemma-base-change-proper}). L'énoncé a donc un sens.

\medskip\noindent
Cela dit, il est clair que le contenu du lemme est le suivant : étant donné un
objet $(T, g, \mathcal{F})$ de $\Cohstack_{X/B}$, il existe un sous-schéma
ouvert $U \subset T$ tel que, pour un morphisme de schémas $T' \to T$, les
conditions suivantes soient équivalentes :
\begin{enumerate}
\item[(a)] $T' \to T$ se factorise par $U$,
\item[(b)] l'image inverse $\mathcal{F}_{T'}$ de $\mathcal{F}$ par
$X_{T'} \to X_T$ est inversible.
\end{enumerate}
Soit $W \subset |X_T|$ l'ensemble des points $x \in |X_T|$ tels que
$\mathcal{F}$ soit localement libre au voisinage de $x$. D'après Plus sur les
morphismes d'espaces, lemme
\ref{spaces-more-morphisms-lemma-finite-free-open}.
$W$ est ouvert et la formation de $W$ commute à tout changement de base. Il
est clair que, si $T' \to T$ satisfait à (b), alors
$|X_{T'}| \to |X_T|$ est à valeurs dans $W$. Nous pouvons donc, pour
construire $U$, remplacer $T$ par l'ouvert
$T \setminus f_T(|X_T| \setminus W)$. Après cette opération, nous sommes dans
la situation où $\mathcal{F}$ est localement libre de type fini. Dans ce cas,
nous obtenons une décomposition en réunion disjointe
$X_T = X_0 \amalg X_1 \amalg X_2 \amalg \ldots$ de sous-espaces ouverts et
fermés telle que la restriction de $\mathcal{F}$ à $X_i$ soit localement libre
de rang $i$. Alors il est clair que
$$
U = T \setminus f_T(|X_0| \cup |X_2| \cup |X_3| \cup \ldots )
$$
convient. (Remarquons que si nous supposons $T$ quasi-compact, alors $X_T$ est
quasi-compact ; seul un nombre fini de $X_i$ sont donc non vides et $U$ est
bien ouvert.)
\end{proof}

\begin{proposition}
\label{quot-proposition-pic}
Soit $S$ un schéma. Soit $f : X \to B$ un morphisme d'espaces algébriques sur
$S$. Si $f$ est plat, de présentation finie et propre, alors
$\Picardstack_{X/B}$ est un champ algébrique.
\end{proposition}

\begin{proof}
Conséquence immédiate du lemme \ref{quot-lemma-picard-stack-open-in-coh}, du lemme
de Champs algébriques
\ref{algebraic-lemma-representable-morphism-to-algebraic}, et soit du théorème
\ref{quot-theorem-coherent-algebraic}, soit du théorème
\ref{quot-theorem-coherent-algebraic-general}
\end{proof}

\section{Le foncteur de Picard}
\label{quot-section-picard-functor}

\noindent
Dans cette section, nous revenons sur le foncteur de Picard étudié dans
Schémas de Picard des courbes, section \ref{pic-section-picard-functor}.
La discussion sera plus générale, car nous voulons étudier le foncteur de
Picard d'un morphisme d'espaces algébriques comme dans la section consacrée au
champ de Picard, voir la section \ref{quot-section-picard-stack}.

\medskip\noindent
Soit $S$ un schéma et soit $X$ un espace algébrique sur $S$.
Un faisceau inversible sur $X$ est un $\mathcal{O}_X$-module inversible
sur $X_\etale$, voir
Modules sur les sites, définition \ref{sites-modules-definition-invertible-sheaf}.
Le groupe des classes d'isomorphisme de modules inversibles est noté
$\Pic(X)$, voir
Modules sur les sites, définition \ref{sites-modules-definition-pic}.
Étant donné un morphisme $f : X \to Y$ d'espaces algébriques sur $S$,
l'image inverse définit un homomorphisme de groupes $\Pic(Y) \to \Pic(X)$.
La correspondance
$X \leadsto \Pic(X)$ est un foncteur contravariant de la catégorie
des schémas vers la catégorie des groupes abéliens. Ce foncteur n'est pas
représentable, mais il se trouve qu'une variante relative de cette construction
l'est parfois.

\begin{situation}
\label{quot-situation-pic}
Soit $S$ un schéma.
Soit $f : X \to B$ un morphisme d'espaces algébriques sur $S$.
Nous définissons
$$
\Picardfunctor_{X/B} : (\Sch/B)^{opp} \longrightarrow \textit{Ensembles}
$$
comme le faisceau fppf associé au foncteur qui, à un schéma $T$
sur $B$, associe le groupe $\Pic(X_T)$.
\end{situation}

\noindent
Dans la situation \ref{quot-situation-pic}, nous considérons parfois
$\Picardfunctor_{X/B}$ comme un foncteur $(\Sch/S)^{opp} \to \textit{Ensembles}$
muni d'un morphisme $\Picardfunctor_{X/B} \to B$. De ce point de vue,
nous définissons $\Picardfunctor_{X/B}$ comme le faisceau fppf associé au
foncteur
$$
T/S \longmapsto \{(h, \mathcal{L}) \mid
h : T \to B,\ \mathcal{L} \in \Pic(X \times_{B, h} T)\}
$$
En particulier, lorsque nous disons que $\Picardfunctor_{X/B}$ est un espace
algébrique, nous entendons que le foncteur correspondant
$(\Sch/S)^{opp} \to \textit{Ensembles}$ est un espace algébrique.

\medskip\noindent
Une remarque souvent utilisée est que, si $T$ est un schéma sur $B$, alors
$\Picardfunctor_{X_T/T}$ est la restriction de $\Picardfunctor_{X/B}$ à
$(\Sch/T)_{fppf}$.

\begin{lemma}
\label{quot-lemma-pic-over-pic}
Dans la situation \ref{quot-situation-pic},
le foncteur $\Picardfunctor_{X/B}$ est le faisceau associé au
foncteur $T \mapsto \Ob(\Picardstack_{X/B, T})/\cong$.
\end{lemma}

\begin{proof}
Puisque la catégorie fibre $\Picardstack_{X/B, T}$ du champ de Picard
$\Picardstack_{X/B}$ au-dessus de $T$ est la catégorie des faisceaux inversibles
sur $X_T$ (voir la section \ref{quot-section-picard-stack} et
Exemples de champs, section \ref{examples-stacks-section-picard-stack}),
cela découle immédiatement des définitions.
\end{proof}

\noindent
Il n'est en fait pas trivial de déterminer la valeur de $\Picardfunctor_{X/B}$
sur les schémas $T$ au-dessus de $B$. Voici un lemme qui aide à accomplir cette
tâche.

\begin{lemma}
\label{quot-lemma-flat-geometrically-connected-fibres}
Dans la situation \ref{quot-situation-pic}.
Si $\mathcal{O}_T \to f_{T, *}\mathcal{O}_{X_T}$ est un isomorphisme
pour tout schéma $T$ sur $B$, alors
$$
0 \to \Pic(T) \to \Pic(X_T) \to \Picardfunctor_{X/B}(T)
$$
est une suite exacte pour tout $T$.
\end{lemma}

\begin{proof}
Nous pouvons remplacer $B$ par $T$ et $X$ par $X_T$, et supposer que $B = T$
afin d'alléger les notations. Soit $\mathcal{N}$ un
$\mathcal{O}_B$-module inversible. Si $f^*\mathcal{N} \cong \mathcal{O}_X$, alors
l'hypothèse donne $f_*f^*\mathcal{N} \cong f_*\mathcal{O}_X \cong \mathcal{O}_B$.
Comme $\mathcal{N}$ est localement trivial, l'application canonique
$\mathcal{N} \to f_*f^*\mathcal{N}$ est localement un isomorphisme (car
$\mathcal{O}_B \to f_*f^*\mathcal{O}_B$ est un isomorphisme par hypothèse).
Nous en concluons que
$\mathcal{N} \to f_*f^*\mathcal{N} \to \mathcal{O}_B$ est un isomorphisme,
et voyons que $\mathcal{N}$ est trivial. Cela prouve que la première flèche
est injective.

\medskip\noindent
Soit $\mathcal{L}$ un $\mathcal{O}_X$-module inversible appartenant au
noyau de $\Pic(X) \to \Picardfunctor_{X/B}(B)$. Il existe alors un
recouvrement fppf $\{B_i \to B\}$ tel que l'image inverse de $\mathcal{L}$
soit le faisceau inversible trivial sur $X_{B_i}$. Choisissons une section
trivialisante $s_i$. Alors $\text{pr}_0^*s_i$ et $\text{pr}_1^*s_j$ sont toutes
deux des sections trivialisantes de $\mathcal{L}$ sur
$X_{B_i \times_B B_j}$ et diffèrent donc par une unité multiplicative
$$
f_{ij} \in
\Gamma(X_{S_i \times_B B_j}, \mathcal{O}_{X_{B_i \times_B B_j}}^*) =
\Gamma(B_i \times_B B_j, \mathcal{O}_{B_i \times_N B_j}^*)
$$
(l'égalité résultant de notre hypothèse sur l'image directe des faisceaux
structuraux). Bien entendu, ces éléments satisfont à la condition de cocycle
sur $B_i \times_B B_j \times_B B_k$ ; ils définissent donc une donnée de
descente sur des faisceaux inversibles relativement au recouvrement fppf
$\{B_i \to B\}$. D'après Descente, proposition
\ref{descent-proposition-fpqc-descent-quasi-coherent}, il existe un
$\mathcal{O}_B$-module inversible $\mathcal{N}$ muni de trivialisations sur les
$B_i$ dont la donnée de descente associée est $\{f_{ij}\}$. (La proposition
s'applique parce que $B$ est un schéma, par le remplacement effectué au début
de la démonstration.) Alors $f^*\mathcal{N} \cong \mathcal{L}$, car le
foncteur des données de descente vers les modules est pleinement fidèle.
\end{proof}

\begin{lemma}
\label{quot-lemma-flat-geometrically-connected-fibres-with-section}
Dans la situation \ref{quot-situation-pic}, soit $\sigma : B \to X$ une section.
Supposons que $\mathcal{O}_T \to f_{T, *}\mathcal{O}_{X_T}$ soit un isomorphisme
pour tout $T$ sur $B$. Alors
$$
0 \to \Pic(T) \to \Pic(X_T) \to \Picardfunctor_{X/B}(T) \to 0
$$
est une suite exacte scindée, une rétraction étant donnée par
$\sigma_T^* : \Pic(X_T) \to \Pic(T)$.
\end{lemma}

\begin{proof}
Notons $K(T) = \Ker(\sigma_T^* : \Pic(X_T) \to \Pic(T))$.
Comme $\sigma$ est une section de $f$, nous voyons que $\Pic(X_T)$ est la somme
directe de $\Pic(T)$ et de $K(T)$.
Ainsi, d'après le lemme \ref{quot-lemma-flat-geometrically-connected-fibres},
nous avons $K(T) \subset \Picardfunctor_{X/B}(T)$ pour tout $T$. De plus,
il ressort clairement de la construction que $\Picardfunctor_{X/B}$ est le
faisceau associé au préfaisceau $K$. Pour achever la démonstration, il suffit
de montrer que $K$ satisfait à la condition de faisceau pour les recouvrements
fppf, ce que nous faisons dans le paragraphe suivant.

\medskip\noindent
Soit $\{T_i \to T\}$ un recouvrement fppf. Soient $\mathcal{L}_i$ des éléments
de $K(T_i)$ qui ont les mêmes images dans $K(T_i \times_T T_j)$ pour tous
$i$ et $j$. Choisissons, pour tout $i$, un isomorphisme
$\alpha_i : \mathcal{O}_{T_i} \to \sigma_{T_i}^*\mathcal{L}_i$.
Choisissons un isomorphisme
$$
\varphi_{ij} :
\mathcal{L}_i|_{X_{T_i \times_T T_j}}
\longrightarrow
\mathcal{L}_j|_{X_{T_i \times_T T_j}}
$$
Si l'application
$$
\alpha_j|_{T_i \times_T T_j} \circ
\sigma_{T_i \times_T T_j}^*\varphi_{ij} \circ
\alpha_i|_{T_i \times_T T_j} :
\mathcal{O}_{T_i \times_T T_j} \to \mathcal{O}_{T_i \times_T T_j}
$$
n'est pas la multiplication par $1$, mais par un certain $u_{ij}$, nous pouvons
multiplier $\varphi_{ij}$ par $u_{ij}^{-1}$ pour y remédier. Cela fait,
considérons l'endomorphisme
$$
\varphi_{ki}|_{X_{T_i \times_T T_j \times_T T_k}} \circ
\varphi_{jk}|_{X_{T_i \times_T T_j \times_T T_k}} \circ
\varphi_{ij}|_{X_{T_i \times_T T_j \times_T T_k}}
\quad\text{sur}\quad
\mathcal{L}_i|_{X_{T_i \times_T T_j \times_T T_k}}
$$
qui est donné par la multiplication par une certaine section $f_{ijk}$ du
faisceau structural de $X_{T_i \times_T T_j \times_T T_k}$. Par notre choix
de $\varphi_{ij}$, nous voyons que l'image inverse de cette application par
$\sigma$ est la multiplication par $1$. Notre hypothèse sur les fonctions sur
$X$ donne alors $f_{ijk} = 1$. Nous obtenons ainsi une donnée de descente pour
le recouvrement fppf
$\{X_{T_i} \to X\}$. D'après Descente sur les espaces, proposition
\ref{spaces-descent-proposition-fpqc-descent-quasi-coherent}, il existe un
$\mathcal{O}_{X_T}$-module inversible $\mathcal{L}$ et un isomorphisme
$\alpha : \mathcal{O}_T \to \sigma_T^*\mathcal{L}$ dont les images inverses
sur $X_{T_i}$ redonnent $(\mathcal{L}_i, \alpha_i)$ (nous omettons un petit
détail). Ainsi, $\mathcal{L}$ définit un objet de $K(T)$, comme voulu.
\end{proof}

\noindent
Dans la situation \ref{quot-situation-pic}, soit $\sigma : B \to X$ une section.
Nous notons $\Picardstack_{X/B, \sigma}$ la catégorie définie comme suit :
\begin{enumerate}
\item Un objet est un quadruplet $(T, h, \mathcal{L}, \alpha)$, où
$(T, h, \mathcal{L})$ est un objet de $\Picardstack_{X/B}$ au-dessus de $T$ et
$\alpha : \mathcal{O}_T \to \sigma_T^*\mathcal{L}$ est un isomorphisme.
\item Un morphisme $(g, \varphi) : (T, h, \mathcal{L}, \alpha)
\to (T', h', \mathcal{L}', \alpha')$ est donné par un morphisme de schémas
$g : T \to T'$ tel que $h = h' \circ g$ et un isomorphisme
$\varphi : (g')^*\mathcal{L}' \to \mathcal{L}$ tel que
$\sigma_T^*\varphi \circ g^*\alpha' = \alpha$.
Ici, $g' : X_{T'} \to X_T$ est le changement de base de $g$.
\end{enumerate}
Il existe un foncteur d'oubli naturel fidèle
$$
\Picardstack_{X/B, \sigma} \longrightarrow
\Picardstack_{X/B}
$$
Nous considérons ainsi $\Picardstack_{X/B, \sigma}$ comme une catégorie
sur $(\Sch/S)_{fppf}$.

\begin{lemma}
\label{quot-lemma-pic-with-section-stack}
Dans la situation \ref{quot-situation-pic}, soit $\sigma : B \to X$ une section.
Alors $\Picardstack_{X/B, \sigma}$, défini ci-dessus, est un champ en
groupoïdes sur $(\Sch/S)_{fppf}$.
\end{lemma}

\begin{proof}
Nous savons déjà que $\Picardstack_{X/B}$ est un champ en groupoïdes
sur $(\Sch/S)_{fppf}$, d'après
Exemples de champs, lemme \ref{examples-stacks-lemma-picard-stack}.
Montrons la descente des objets pour $\Picardstack_{X/B, \sigma}$.
Soit $\{T_i \to T\}$ un recouvrement fppf, soit
$\xi_i = (T_i, h_i, \mathcal{L}_i, \alpha_i)$ un objet de
$\Picardstack_{X/B, \sigma}$ au-dessus de $T_i$, et soit
$\varphi_{ij} : \text{pr}_0^*\xi_i \to \text{pr}_1^*\xi_j$
une donnée de descente. En appliquant le résultat pour $\Picardstack_{X/B}$,
nous voyons que nous pouvons supposer donné un objet $(T, h, \mathcal{L})$
de $\Picardstack_{X/B}$ au-dessus de $T$ dont l'image inverse est $\xi_i$
pour tout $i$. Nous obtenons alors
$$
\alpha_i : \mathcal{O}_{T_i} \to \sigma_{T_i}^*\mathcal{L}_i =
(T_i \to T)^*\sigma_T^*\mathcal{L}
$$
Puisque les applications $\varphi_{ij}$ sont compatibles aux $\alpha_i$,
nous voyons que les images inverses de $\alpha_i$ et $\alpha_j$ sont la même
application sur $T_i \times_T T_j$. Par descente des faisceaux quasi-cohérents
(Descente, proposition \ref{descent-proposition-fpqc-descent-quasi-coherent},
nous voyons que les $\alpha_i$ sont les restrictions d'une unique application
$\alpha : \mathcal{O}_T \to \sigma_T^*\mathcal{L}$, comme voulu.
Nous omettons la démonstration de la descente des morphismes.
\end{proof}

\begin{lemma}
\label{quot-lemma-compare-pic-with-section}
Dans la situation \ref{quot-situation-pic}, soit $\sigma : B \to X$ une section.
Le morphisme $\Picardstack_{X/B, \sigma} \to \Picardstack_{X/B}$
est représentable, surjectif et lisse.
\end{lemma}

\begin{proof}
Soit $T$ un schéma et soit $(\Sch/T)_{fppf} \to \Picardstack_{X/B}$
donné par l'objet $\xi = (T, h, \mathcal{L})$ de $\Picardstack_{X/B}$
au-dessus de $T$. Nous devons montrer que
$$
(\Sch/T)_{fppf} \times_{\xi, \Picardstack_{X/B}} \Picardstack_{X/B, \sigma}
$$
est représentable par un schéma $V$ et que le morphisme correspondant
$V \to T$ est surjectif et lisse. Voir
Champs algébriques, sections \ref{algebraic-section-representable-morphism},
\ref{algebraic-section-morphisms-representable-by-algebraic-spaces} et
\ref{algebraic-section-representable-properties}.
Le foncteur d'oubli $\Picardstack_{X/B, \sigma} \to \Picardstack_{X/B}$
est fidèle sur les catégories fibres et, pour $T'/T$, l'ensemble des classes
d'isomorphisme est l'ensemble des isomorphismes
$$
\alpha' : \mathcal{O}_{T'} \longrightarrow (T' \to T)^*\sigma_T^*\mathcal{L}
$$
Voir Champs algébriques, lemme
\ref{algebraic-lemma-criterion-map-representable-spaces-fibred-in-groupoids}.
Nous savons que ce foncteur est représentable par un schéma affine $U$ de
présentation finie sur $T$, d'après la proposition \ref{quot-proposition-isom}
(appliquée à $\text{id} : T \to T$, à $\mathcal{O}_T$ et à
$\sigma^*\mathcal{L}$). En travaillant localement pour la topologie de Zariski
sur $T$, nous pouvons supposer que $\sigma_T^*\mathcal{L}$ est isomorphe à
$\mathcal{O}_T$ ; nous voyons alors que notre foncteur est représentable par
$\mathbf{G}_m \times T$ sur $T$. Ainsi,
$U \to T$ est, localement pour la topologie de Zariski sur $T$, de la forme
de la projection $\mathbf{G}_m \times T \to T$, laquelle est bien lisse et
surjective.
\end{proof}

\begin{lemma}
\label{quot-lemma-flat-geometrically-connected-fibres-with-section-functor-stack}
Dans la situation \ref{quot-situation-pic}, soit $\sigma : B \to X$ une section.
Si $\mathcal{O}_T \to f_{T, *}\mathcal{O}_{X_T}$ est un isomorphisme
pour tout $T$ sur $B$, alors
$\Picardstack_{X/B, \sigma} \to (\Sch/S)_{fppf}$
est fibré en sétoïdes et l'ensemble de ses classes d'isomorphisme au-dessus de
$T$ est donné par
$$
\coprod\nolimits_{h : T \to B}
\Ker(\sigma_T^* : \Pic(X \times_{B, h} T) \to \Pic(T))
$$
\end{lemma}

\begin{proof}
Si $\xi = (T, h, \mathcal{L}, \alpha)$ est un objet de
$\Picardstack_{X/B, \sigma}$ au-dessus de $T$, alors un automorphisme
$\varphi$ de $\xi$ est donné par la multiplication par une section globale
inversible $u$ du faisceau structural de $X_T$ telle qu'en outre
$\sigma_T^*u = 1$. Alors $u = 1$ par notre hypothèse selon laquelle
$\mathcal{O}_T \to f_{T, *}\mathcal{O}_{X_T}$ est un isomorphisme.
Ainsi, $\Picardstack_{X/B, \sigma}$ est fibré en sétoïdes sur
$(\Sch/S)_{fppf}$. Étant donnés $T$ et $h : T \to B$, l'ensemble des classes
d'isomorphisme de couples $(\mathcal{L}, \alpha)$ est le même que l'ensemble
des classes d'isomorphisme des $\mathcal{L}$ tels que
$\sigma_T^*\mathcal{L} \cong \mathcal{O}_T$ (sans spécification de
l'isomorphisme). Cela est clair, car deux choix de $\alpha$ diffèrent par une
unité globale sur $T$, ce qui revient au même qu'une unité globale sur $X_T$.
\end{proof}

\begin{proposition}
\label{quot-proposition-pic-functor}
Soit $S$ un schéma. Soit $f : X \to B$ un morphisme d'espaces algébriques
sur $S$. Supposons que
\begin{enumerate}
\item $f$ soit plat, de présentation finie et propre, et
\item $\mathcal{O}_T \to f_{T, *}\mathcal{O}_{X_T}$ soit un isomorphisme
pour tout schéma $T$ sur $B$.
\end{enumerate}
Alors $\Picardfunctor_{X/B}$ est un espace algébrique.
\end{proposition}

\noindent
Dans la situation de la proposition, le champ algébrique
$\Picardstack_{X/B}$ est une gerbe au-dessus de l'espace algébrique
$\Picardfunctor_{X/B}$. Après développement de la théorie générale des gerbes,
cela fournit une démonstration plus courte de la proposition (mais faisant
appel à une théorie plus générale).

\begin{proof}
Il existe un morphisme surjectif, plat et de présentation finie
$B' \to B$ d'espaces algébriques tel que le changement de base
$X' = X \times_B B'$ sur $B'$ possède une section : nous pouvons en effet
prendre $B' = X$.
Observons que $\Picardfunctor_{X'/B'} = B' \times_B \Picardfunctor_{X/B}$.
Ainsi, $\Picardfunctor_{X'/B'} \to \Picardfunctor_{X/B}$ est représentable
par des espaces algébriques, surjectif, plat et de présentation finie.
Par conséquent, si nous pouvons montrer que $\Picardfunctor_{X'/B'}$ est un
espace algébrique, il s'ensuit que $\Picardfunctor_{X/B}$ est un espace
algébrique d'après Amorçage, théorème
\ref{bootstrap-theorem-final-bootstrap}.
Nous nous ramenons ainsi au cas décrit dans le paragraphe suivant.

\medskip\noindent
En plus des hypothèses de la proposition, supposons donnée une section
$\sigma : B \to X$. D'après la proposition \ref{quot-proposition-pic},
$\Picardstack_{X/B}$ est un champ algébrique.
D'après le lemme \ref{quot-lemma-compare-pic-with-section} et
Champs algébriques, lemme
\ref{algebraic-lemma-representable-morphism-to-algebraic},
$\Picardstack_{X/B, \sigma}$ est un champ algébrique.
D'après le lemme
\ref{quot-lemma-flat-geometrically-connected-fibres-with-section-functor-stack}
et Champs algébriques, lemme
\ref{algebraic-lemma-characterize-representable-by-space},
nous voyons que $T \mapsto \Ker(\sigma_T^* : \Pic(X_T) \to \Pic(T))$
est un espace algébrique.
D'après le lemme \ref{quot-lemma-flat-geometrically-connected-fibres-with-section},
ce foncteur est le même que $\Picardfunctor_{X/B}$.
\end{proof}

\begin{lemma}
\label{quot-lemma-diagonal-pic}
Sous les hypothèses et avec les notations de la proposition
\ref{quot-proposition-pic-functor}.
Alors la diagonale
$\Picardfunctor_{X/B} \to \Picardfunctor_{X/B} \times_B \Picardfunctor_{X/B}$
est représentable par des immersions. Autrement dit,
$\Picardfunctor_{X/B} \to B$ est localement séparé.
\end{lemma}

\begin{proof}
Soit $T$ un schéma sur $B$ et soient $s, t \in \Picardfunctor_{X/B}(T)$.
Nous voulons montrer qu'il existe un sous-schéma localement fermé
$Z \subset T$ tel que $s|_Z = t|_Z$ et tel qu'un morphisme $T' \to T$ se
factorise par $Z$ si et seulement si $s|_{T'} = t|_{T'}$.

\medskip\noindent
Nous commençons par ramener le problème général au cas où $s$ et $t$
proviennent de modules inversibles sur $X_T$. Nous suggérons au lecteur de
sauter cette étape. Choisissons un recouvrement fppf
$\{T_i \to T\}_{i \in I}$ tel que $s|_{T_i}$ et $t|_{T_i}$ proviennent de
$\Pic(X_{T_i})$ pour tout $i$.
Supposons que nous puissions démontrer le résultat pour tous les couples
$s|_{T_i}, t|_{T_i}$. Nous obtenons alors des sous-schémas localement fermés
$Z_i \subset T_i$ possédant la propriété universelle voulue.
Il s'ensuit que $Z_i$ et $Z_j$ ont la même image inverse schématique dans
$T_i \times_T T_j$.
Cela détermine une donnée de descente sur $Z_i/T_i$.
Comme $Z_i \to T_i$ est localement quasi-fini, il résulte de
Compléments sur les morphismes, lemme
\ref{more-morphisms-lemma-separated-locally-quasi-finite-morphisms-fppf-descend}
que nous obtenons un morphisme localement quasi-fini $Z \to T$ dont le
changement de base redonne $Z_i \to T_i$. Alors $Z \to T$ est une immersion
d'après Descente, lemme
\ref{descent-lemma-descending-fppf-property-immersion}.
Enfin, puisque $\Picardfunctor_{X/B}$ est un faisceau fppf, nous concluons que
$s|_Z = t|_Z$ et que $Z$ satisfait à la propriété universelle mentionnée
ci-dessus.

\medskip\noindent
Supposons que $s$ et $t$ proviennent de modules inversibles
$\mathcal{V}$, $\mathcal{W}$ sur $X_T$.
Posons $\mathcal{L} = \mathcal{V} \otimes \mathcal{W}^{\otimes -1}$
Nous cherchons un sous-schéma localement fermé $Z$ de $T$ tel que
$T' \to T$ se factorise par $Z$ si et seulement si
$\mathcal{L}_{X_{T'}}$ est l'image inverse d'un faisceau inversible sur $T'$,
voir le lemme \ref{quot-lemma-flat-geometrically-connected-fibres}.
L'existence de $Z$ résulte donc de
Compléments sur les morphismes d'espaces, lemme
\ref{spaces-more-morphisms-lemma-diagonal-picard-flat-proper}.
\end{proof}

\section{Morphismes relatifs}
\label{quot-section-relative-morphisms}

\noindent
Nous poursuivons la discussion de Critères de représentabilité, section
\ref{criteria-section-relative-morphisms}.
Dans cette section, partant d'un schéma $S$ et de morphismes d'espaces
algébriques $Z \to B$ et $X \to B$ sur $S$, nous avons construit un foncteur
$$
\mathit{Mor}_B(Z, X) : (\Sch/B)^{opp} \longrightarrow \textit{Ensembles}, \quad
T \longmapsto \{f : Z_T \to X_T\}
$$
Nous considérons parfois $\mathit{Mor}_B(Z, X)$ comme un foncteur
$(\Sch/S)^{opp} \to \textit{Ensembles}$ muni d'un morphisme
$\mathit{Mor}_B(Z, X) \to B$.
En effet, si $T$ est un schéma sur $S$, alors un élément de
$\mathit{Mor}_B(Z, X)(T)$ est un couple $(f, h)$, où $h$ est un morphisme
$h : T \to B$ et où
$f : Z \times_{B, h} T \to X \times_{B, h} T$
est un morphisme d'espaces algébriques sur $T$. En particulier, lorsque nous
disons que $\mathit{Mor}_B(Z, X)$ est un espace algébrique, nous entendons que
le foncteur correspondant $(\Sch/S)^{opp} \to \textit{Ensembles}$ est un espace
algébrique.

\begin{lemma}
\label{quot-lemma-Mor-into-Hilb}
Soit $S$ un schéma. Considérons des morphismes d'espaces algébriques
$Z \to B$ et $X \to B$ sur $S$.
Si $X \to B$ est séparé et si $Z \to B$ est de présentation finie,
plat et propre, alors il existe une transformation injective naturelle de
foncteurs
$$
\mathit{Mor}_B(Z, X) \longrightarrow \Hilbfunctor_{Z \times_B X/B}
$$
qui envoie un morphisme $f : Z_T \to X_T$ sur son graphe.
\end{lemma}

\begin{proof}
Étant donné un schéma $T$ sur $B$ et un morphisme $f_T : Z_T \to X_T$
sur $T$, le graphe de $f$ est le morphisme
$\Gamma_f = (\text{id}, f) : Z_T \to Z_T \times_T X_T = (Z \times_B X)_T$.
Rappelons que les propriétés d'être séparé, plat, propre ou de présentation
finie sont des propriétés des morphismes d'espaces algébriques qui sont stables
par changement de base (Morphismes d'espaces, lemmes
\ref{spaces-morphisms-lemma-base-change-separated},
\ref{spaces-morphisms-lemma-base-change-flat},
\ref{spaces-morphisms-lemma-base-change-proper} et
\ref{spaces-morphisms-lemma-base-change-finite-presentation}).
Ainsi, $\Gamma_f$ est une immersion fermée d'après
Morphismes d'espaces, lemme \ref{spaces-morphisms-lemma-semi-diagonal}.
De plus, $\Gamma_f(Z_T)$ est plat, propre et de présentation finie sur $T$.
Donc $\Gamma_f(Z_T)$ définit un élément de
$\Hilbfunctor_{Z \times_B X/B}(T)$.
Pour montrer que la transformation est injective, il suffit de montrer que deux
morphismes ayant le même graphe sont égaux. Cela est vrai parce que, si
$Y \subset (Z \times_B X)_T$ est le graphe d'un morphisme $f$, nous pouvons
retrouver $f$ en composant l'inverse de
$\text{pr}_1|_Y : Y \to Z_T$ avec $\text{pr}_2|_Y$.
\end{proof}

\begin{lemma}
\label{quot-lemma-Mor-into-Hilb-open}
Hypothèses et notations comme dans le lemme \ref{quot-lemma-Mor-into-Hilb}.
La transformation
$\mathit{Mor}_B(Z, X) \longrightarrow \Hilbfunctor_{Z \times_B X/B}$
est représentable par des immersions ouvertes.
\end{lemma}

\begin{proof}
Soit $T$ un schéma sur $B$ et soit $Y \subset (Z \times_B X)_T$
un élément de $\Hilbfunctor_{Z \times_B X/B}(T)$. Nous voyons alors que
$Y$ est le graphe d'un morphisme $Z_T \to X_T$ sur $T$ si et seulement si
$k = \text{pr}_1|_Y : Y \to Z_T$ est un isomorphisme. D'après
Compléments sur les morphismes d'espaces, lemme
\ref{spaces-more-morphisms-lemma-where-isomorphism}, il existe un sous-schéma
ouvert $V \subset T$ tel que, pour tout morphisme de schémas $T' \to T$,
$k_{T'} : Y_{T'} \to Z_{T'}$ soit un isomorphisme si et seulement si
$T' \to T$ se factorise par $V$.
Cela prouve le lemme.
\end{proof}

\begin{proposition}
\label{quot-proposition-Mor}
Soit $S$ un schéma. Soient $Z \to B$ et $X \to B$ des morphismes d'espaces
algébriques sur $S$. Supposons que $X \to B$ soit de présentation finie et
séparé, et que $Z \to B$ soit de présentation finie, plat et propre. Alors
$\mathit{Mor}_B(Z, X)$ est un espace algébrique localement de présentation
finie sur $B$.
\end{proposition}

\begin{proof}
Conséquence immédiate du lemme \ref{quot-lemma-Mor-into-Hilb-open}
et de la proposition \ref{quot-proposition-hilb}.
\end{proof}







\section{Le champ des espaces algébriques}
\label{quot-section-stack-of-spaces}

\noindent
Cette section poursuit la discussion commencée dans
Exemples de champs, sections
\ref{examples-stacks-section-stack-of-spaces},
\ref{examples-stacks-section-stack-of-finite-type-spaces} et
\ref{examples-stacks-section-stack-in-groupoids-of-finite-type-spaces}.
En travaillant sur $\mathbf{Z}$, la discussion qui s'y trouve montre que nous
avons un champ en groupoïdes
$$
p'_{ft} : \Spacesstack'_{ft} \longrightarrow \Sch_{fppf}
$$
qui paramètre les familles (non plates) d'espaces algébriques de type fini.
Plus précisément, un objet\footnote{Nous effectuons toujours un remplacement
comme dans Exemples de champs, lemme
\ref{examples-stacks-lemma-stack-ft-spaces}.}
de $\Spacesstack'_{ft}$ est un morphisme de type fini $X \to S$
d'un espace algébrique $X$ vers un schéma $S$, et un morphisme
$(X' \to S') \to (X \to S)$ est donné par un couple $(f, g)$, où
$f : X' \to X$ est un morphisme d'espaces algébriques et
$g : S' \to S$ est un morphisme de schémas, qui s'inscrivent dans un diagramme
commutatif
$$
\xymatrix{
X' \ar[d] \ar[r]_f & X \ar[d] \\
S' \ar[r]^g & S
}
$$
induisant un isomorphisme $X' \to S' \times_S X$ ; autrement dit, le
diagramme est cartésien dans la catégorie des espaces algébriques.
Le foncteur $p'_{ft}$ envoie $(X \to S)$ sur $S$ et $(f, g)$ sur $g$.
Nous définissons une sous-catégorie pleine
$$
\Spacesstack'_{fp, flat, proper} \subset
\Spacesstack'_{ft}
$$
formée des objets $X \to S$ de $\Spacesstack'_{ft}$ tels que
$X \to S$ soit de présentation finie, plat et propre.
Nous notons
$$
p'_{fp, flat, proper} :
\Spacesstack'_{fp, flat, proper}
\longrightarrow
\Sch_{fppf}
$$
la restriction du foncteur $p'_{ft}$ à la sous-catégorie indiquée.
Nous rappelons d'abord les résultats déjà obtenus dans les références
ci-dessus, puis nous commençons à y ajouter d'autres résultats.

\begin{lemma}
\label{quot-lemma-spaces-fibred-in-groupoids}
La catégorie $\Spacesstack'_{ft}$ est fibrée en groupoïdes
sur $\Sch_{fppf}$. Il en est de même de
$\Spacesstack'_{fp, flat, proper}$.
\end{lemma}

\begin{proof}
Nous l'avons vu dans
Exemples de champs, section
\ref{examples-stacks-section-stack-in-groupoids-of-finite-type-spaces}
pour $\Spacesstack'_{ft}$, et cela entraîne aisément le résultat dans l'autre
cas. Cependant, démontrons-le aussi directement en vérifiant les conditions
(1) et (2) de Catégories, définition
\ref{categories-definition-fibred-groupoids}.

\medskip\noindent
Condition (1). Soit $X \to S$ un objet de $\Spacesstack'_{ft}$ et soit
$S' \to S$ un morphisme de schémas. Posons alors $X' = S' \times_S X$.
Notons que $X' \to S'$ est de type fini d'après
Morphismes d'espaces, lemme
\ref{spaces-morphisms-lemma-base-change-finite-type},
ce qui donne un morphisme $(X' \to S') \to (X \to S)$
au-dessus de $S' \to S$.
On raisonne de même dans l'autre cas en utilisant
Morphismes d'espaces, lemmes
\ref{spaces-morphisms-lemma-base-change-finite-presentation},
\ref{spaces-morphisms-lemma-base-change-flat} et
\ref{spaces-morphisms-lemma-base-change-proper}.

\medskip\noindent
Condition (2). Considérons des morphismes
$(f, g) : (X' \to S') \to (X \to S)$ et
$(a, b) : (Y \to T) \to (X \to S)$ de $\Spacesstack'_{ft}$.
Étant donné un morphisme $h : T \to S'$ tel que $g \circ h = b$, nous devons
montrer qu'il existe un unique morphisme
$(k, h) : (Y \to T) \to (X' \to S')$ de $\Spacesstack'_{ft}$ tel que
$(f, g) \circ (k, h) = (a, b)$.
Cela résulte clairement de l'égalité $X' = S' \times_S X$.
Le même raisonnement vaut donc pour toute sous-catégorie pleine de
$\Spacesstack'_{ft}$ qui satisfait à (1).
\end{proof}

\begin{lemma}
\label{quot-lemma-spaces-diagonal}
La diagonale
$$
\Delta : \Spacesstack'_{fp, flat, proper} \longrightarrow
\Spacesstack'_{fp, flat, proper} \times \Spacesstack'_{fp, flat, proper}
$$
est représentable par des espaces algébriques.
\end{lemma}

\begin{proof}
Nous utiliserons le critère (2) de
Champs algébriques, lemme \ref{algebraic-lemma-representable-diagonal}.
Soit $S$ un schéma et soient $X$ et $Y$ des espaces algébriques de présentation
finie sur $S$, plats sur $S$ et propres sur $S$. Nous devons montrer que le
foncteur
$$
\mathit{Isom}_S(X, Y) : (\Sch/S)_{fppf} \longrightarrow \textit{Ensembles}, \quad
T \longmapsto \{f : X_T \to Y_T \text{ est un isomorphisme}\}
$$
est un espace algébrique. Un argument élémentaire montre que
$\mathit{Isom}_S(X, Y)$ s'insère dans un produit fibré
$$
\xymatrix{
\mathit{Isom}_S(X, Y) \ar[r] \ar[d] & S \ar[d]_{(\text{id}, \text{id})} \\
\mathit{Mor}_S(X, Y) \times \mathit{Mor}_S(Y, X) \ar[r] &
\mathit{Mor}_S(X, X) \times \mathit{Mor}_S(Y, Y)
}
$$
La flèche du bas envoie $(\varphi, \psi)$ sur
$(\psi \circ \varphi, \varphi \circ \psi)$.
D'après la proposition \ref{quot-proposition-Mor}, les foncteurs de la ligne
inférieure sont des espaces algébriques sur $S$.
Le résultat découle donc du fait que la catégorie des espaces algébriques sur
$S$ admet les produits fibrés.
\end{proof}

\begin{lemma}
\label{quot-lemma-spaces-stack}
La catégorie $\Spacesstack'_{ft}$ est un champ en groupoïdes
sur $\Sch_{fppf}$. Il en est de même de
$\Spacesstack'_{fp, flat, proper}$.
\end{lemma}

\begin{proof}
La raison pour laquelle ce lemme est vrai tient dans le slogan suivant : toute
donnée de descente fppf pour les espaces algébriques est effective, voir
Amorçage, section \ref{bootstrap-section-applications}.
Plus précisément, le lemme pour $\Spacesstack'_{ft}$ découle de
Exemples de champs, lemme
\ref{examples-stacks-lemma-stack-of-finite-type-spaces},
comme nous l'avons vu dans Exemples de champs, section
\ref{examples-stacks-section-stack-in-groupoids-of-finite-type-spaces}.
Cependant, rappelons la démonstration. Nous devons vérifier les conditions
(1), (2) et (3) de Champs, définition
\ref{stacks-definition-stack-in-groupoids}.

\medskip\noindent
Nous avons vu la propriété (1) au lemme
\ref{quot-lemma-spaces-fibred-in-groupoids}.

\medskip\noindent
La propriété (2) résulte du lemme \ref{quot-lemma-spaces-diagonal} dans le cas de
$\Spacesstack'_{fp, flat, proper}$.
Dans le cas de $\Spacesstack'_{ft}$, elle découle de
Exemples de champs, lemme
\ref{examples-stacks-lemma-pre-stack-of-spaces}
(et c'est réellement la référence « correcte »).

\medskip\noindent
La condition (3) pour $\Spacesstack'_{ft}$ se vérifie comme suit. Supposons
donnés
\begin{enumerate}
\item un recouvrement fppf $\{U_i \to U\}_{i \in I}$ dans $\Sch_{fppf}$,
\item pour tout $i \in I$, un espace algébrique $X_i$ de type fini sur
$U_i$, et
\item pour tous $i, j \in I$, un isomorphisme
$\varphi_{ij} : X_i \times_U U_j \to U_i \times_U X_j$ d'espaces algébriques
sur $U_i \times_U U_j$, satisfaisant à la condition de cocycle sur
$U_i \times_U U_j \times_U U_k$.
\end{enumerate}
Nous devons montrer qu'il existe un espace algébrique $X$ de type fini sur $U$
et des isomorphismes $X_{U_i} \cong X_i$ sur $U_i$ redonnant les
isomorphismes $\varphi_{ij}$. Cela résulte de
Amorçage, lemme \ref{bootstrap-lemma-descend-algebraic-space}, partie (2).
D'après Descente sur les espaces, lemme
\ref{spaces-descent-lemma-descending-property-finite-type},
$X \to U$ est de type fini.
Dans le cas de $\Spacesstack'_{fp, flat, proper}$, on utilise en outre
Descente sur les espaces, lemmes
\ref{spaces-descent-lemma-descending-property-finite-presentation},
\ref{spaces-descent-lemma-descending-property-flat} et
\ref{spaces-descent-lemma-descending-property-proper}
à la dernière étape.
\end{proof}

\noindent
Vérification de cohérence : les champs $\Spacesstack'_{ft}$ et
$\Spacesstack'_{fp, flat, proper}$ jouent le même rôle parmi les espaces
algébriques.

\begin{lemma}
\label{quot-lemma-extend-spaces-to-spaces}
Soit $T$ un espace algébrique sur $\mathbf{Z}$. Soit $\mathcal{S}_T$
le champ algébrique correspondant (Champs algébriques, sections
\ref{algebraic-section-split},
\ref{algebraic-section-representable-by-algebraic-spaces} et
\ref{algebraic-section-stacks-spaces}).
Nous avons une équivalence de catégories
$$
\left\{
\begin{matrix}
\text{morphismes d'espaces algébriques }\\
X \to T\text{ de type fini}
\end{matrix}
\right\}
\longrightarrow
\Mor_{\textit{Cat}/\Sch_{fppf}}(\mathcal{S}_T, \Spacesstack'_{ft})
$$
et une équivalence de catégories
$$
\left\{
\begin{matrix}
\text{morphismes d'espaces algébriques }X \to T\\
\text{de présentation finie, plats et propres}
\end{matrix}
\right\}
\longrightarrow
\Mor_{\textit{Cat}/\Sch_{fppf}}(\mathcal{S}_T,
\Spacesstack'_{fp, flat, proper})
$$
\end{lemma}

\begin{proof}
Nous allons déduire ce lemme du fait qu'il vaut pour les schémas
(essentiellement par construction des champs) et du fait que les données de
descente fppf pour les espaces algébriques sur des espaces algébriques sont
effectives. Nous encourageons vivement le lecteur à omettre la démonstration.

\medskip\noindent
La construction de gauche à droite dans chacune des flèches est immédiate :
étant donné $X \to T$ de type fini, le foncteur
$\mathcal{S}_T \to  \Spacesstack'_{ft}$ associe à $U/T$ le changement de base
$X_U \to U$. Nous allons expliquer comment construire un quasi-inverse.

\medskip\noindent
Si $T$ est un schéma, il existe un quasi-inverse d'après le lemme de
$2$-Yoneda, voir Catégories, lemme
\ref{categories-lemma-yoneda-2category}.
Soit $p : U \to T$ un morphisme surjectif \'etale, où $U$ est un schéma.
Soit $R = U \times_T U$, de projections $s, t : R \to U$.
Observons que nous obtenons des morphismes
$$
\xymatrix{
\mathcal{S}_{U \times_T U \times_T U} \ar@<2ex>[r] \ar[r] \ar@<-2ex>[r] &
\mathcal{S}_R \ar@<1ex>[r] \ar@<-1ex>[r] &
\mathcal{S}_U \ar[r] &
\mathcal{S}_T
}
$$
qui satisfont strictement à diverses compatibilités.

\medskip\noindent
Soit $G : \mathcal{S}_T \to \Spacesstack'_{ft}$ un foncteur sur
$\Sch_{fppf}$. La restriction de $G$ à $\mathcal{S}_U$ par l'application
ci-dessus correspond, par le lemme de $2$-Yoneda, à un morphisme de type fini
$X_U \to U$ d'espaces algébriques. Puisque $p \circ s = p \circ t$, nous
voyons que $R \times_{s, U} X_U$ et $R \times_{t, U} X_U$ correspondent tous
deux à la restriction de $G$ à $\mathcal{S}_R$. Nous obtenons donc un
isomorphisme canonique
$\varphi : X_U \times_{U, t} R \to R \times_{s, U} X_U$ sur $R$.
Cet isomorphisme satisfait à la condition de cocycle en vertu des diverses
compatibilités du diagramme ci-dessus.
Nous avons ainsi une donnée de descente qui est effective d'après
Amorçage, lemme \ref{bootstrap-lemma-descend-algebraic-space}, partie (2).
Autrement dit, nous obtenons un objet $X \to T$ de la catégorie du membre de
droite. Nous omettons de vérifier que la construction $G \leadsto X$ est
fonctorielle et qu'elle est quasi-inverse de l'autre construction.
Dans le cas de $\Spacesstack'_{fp, flat, proper}$, on utilise en outre
Descente sur les espaces, lemmes
\ref{spaces-descent-lemma-descending-property-finite-presentation},
\ref{spaces-descent-lemma-descending-property-flat} et
\ref{spaces-descent-lemma-descending-property-proper}
à la dernière étape pour voir que $X \to T$ est de présentation finie,
plat et propre.
\end{proof}

\begin{remark}
\label{quot-remark-spaces-base-change}
Soit $B$ un espace algébrique sur $\Spec(\mathbf{Z})$.
Soit $B\textit{-Espaces}'_{ft}$ la catégorie formée des couples
$(X \to S, h : S \to B)$, où $X \to S$ est un objet de
$\Spacesstack'_{ft}$ et $h : S \to B$ est un morphisme.
Un morphisme $(X' \to S', h') \to (X \to S, h)$ dans
$B\textit{-Espaces}'_{ft}$ est un morphisme $(f, g)$ dans
$\Spacesstack'_{ft}$ tel que $h \circ g = h'$.
Dans cette situation, le diagramme
$$
\xymatrix{
B\textit{-Espaces}'_{ft} \ar[r] \ar[d] & \Spacesstack'_{ft} \ar[d] \\
(\Sch/B)_{fppf} \ar[r] & \Sch_{fppf}
}
$$
est un carré $2$-cartésien. Cette remarque triviale sera parfois utile
pour déduire des résultats du cas absolu $\Spacesstack'_{ft}$ au cas des
familles sur un espace algébrique de base donné.
Bien entendu, une construction analogue vaut pour
$B\textit{-Espaces}'_{fp, flat, proper}$
\end{remark}

\begin{lemma}
\label{quot-lemma-spaces-limits}
Le champ
$p'_{fp, flat, proper} :
\Spacesstack'_{fp, flat, proper} \to \Sch_{fppf}$ commute aux limites
(Axiomes d'Artin, définition \ref{artin-definition-limit-preserving}).
\end{lemma}

\begin{proof}
Soit $T = \lim T_i$ la limite d'un système projectif filtrant de schémas
affines. D'après Limites d'espaces, lemme
\ref{spaces-limits-lemma-descend-finite-presentation},
la catégorie des espaces algébriques de présentation finie sur $T$ est la
limite inductive des catégories d'espaces algébriques de présentation finie
sur les $T_i$.
Pour achever la démonstration, utilisons le fait que la platitude et la
propreté descendent à travers la limite, voir
Limites d'espaces, lemmes
\ref{spaces-limits-lemma-descend-flat} et
\ref{spaces-limits-lemma-eventually-proper}.
\end{proof}

\begin{lemma}
\label{quot-lemma-spaces-RS-star}
Soit
$$
\xymatrix{
T \ar[r] \ar[d] & T' \ar[d] \\
S \ar[r] & S'
}
$$
une somme amalgamée dans la catégorie des schémas, où
$T \to T'$ est un épaississement et $T \to S$ est affine, voir
Compléments sur les morphismes, lemme
\ref{more-morphisms-lemma-pushout-along-thickening}.
Alors le foncteur sur les catégories fibres
$$
\begin{matrix}
\Spacesstack'_{fp, flat, proper, S'} \\
\downarrow \\
\Spacesstack'_{fp, flat, proper, S}
\times_{\Spacesstack'_{fp, flat, proper, T}}
\Spacesstack'_{fp, flat, proper, T'}
\end{matrix}
$$
est une équivalence.
\end{lemma}

\begin{proof}
Le foncteur est une équivalence si l'on retranche « propre » de la liste des
conditions et que l'on remplace « de présentation finie » par
« localement de présentation finie », voir Sommes amalgamées d'espaces, lemme
\ref{spaces-pushouts-lemma-equivalence-categories-spaces-pushout-flat}.
Il suffit donc de montrer qu'étant donné un morphisme $X' \to S'$ d'un espace
algébrique vers $S'$ qui est plat et localement de présentation finie,
$X' \to S'$ est propre si et seulement si
$S \times_{S'} X' \to S$ et $T' \times_{S'} X' \to T'$ sont propres.
Une implication résulte du fait que la propreté est préservée par changement de
base (Morphismes d'espaces, lemme
\ref{spaces-morphisms-lemma-base-change-proper}) et l'autre, du fait que la
propreté de $S \times_{S'} X' \to S$ entraîne celle de $X' \to S'$, d'après
Compléments sur les morphismes d'espaces, lemme
\ref{spaces-more-morphisms-lemma-thicken-property-morphisms-cartesian}.
\end{proof}

\begin{lemma}
\label{quot-lemma-spaces-tangent-space}
Soit $k$ un corps et soit $x = (X \to \Spec(k))$ un objet de
$\mathcal{X} = \Spacesstack'_{fp, flat, proper}$ au-dessus de $\Spec(k)$.
\begin{enumerate}
\item Si $k$ est de type fini sur $\mathbf{Z}$, alors les espaces vectoriels
$T\mathcal{F}_{\mathcal{X}, k, x}$ et
$\text{Inf}(\mathcal{F}_{\mathcal{X}, k, x})$
(voir Axiomes d'Artin, section \ref{artin-section-tangent-spaces})
sont de dimension finie, et
\item en général, les espaces vectoriels $T_x(k)$ et $\text{Inf}_x(k)$
(voir Axiomes d'Artin, section \ref{artin-section-inf})
sont de dimension finie.
\end{enumerate}
\end{lemma}

\begin{proof}
La discussion d'Axiomes d'Artin, section \ref{artin-section-tangent-spaces},
ne s'applique qu'aux corps de type fini sur le schéma de base
$\Spec(\mathbf{Z})$.
Notre champ satisfait à (RS*) d'après le lemme \ref{quot-lemma-spaces-RS-star},
et nous pouvons appliquer
Axiomes d'Artin, lemme \ref{artin-lemma-properties-lift-RS-star},
pour obtenir les espaces vectoriels $T_x(k)$ et $\text{Inf}_x(k)$ mentionnés en
(2). De plus, dans le cas de type fini, ces espaces coïncident avec ceux
mentionnés en (1), d'après Axiomes d'Artin, remarque
\ref{artin-remark-compare-deformation-spaces}.
Cela étant réglé, nous pouvons commencer la démonstration.
Observons que l'épaississement du premier ordre
$\Spec(k) \to \Spec(k[\epsilon]) = \Spec(k[k])$
a pour module conormal $k$. Ainsi, les formules d'après
Théorie des déformations, lemme \ref{defos-lemma-deform-spaces},
qui décrivent les déformations infinitésimales de $X$ et les automorphismes
infinitésimaux de $X$, deviennent
$$
T_x(k) = \Ext^1_{\mathcal{O}_X}(\NL_{X/k}, \mathcal{O}_X)
\quad\text{et}\quad
\text{Inf}_x(k) = \Ext^0_{\mathcal{O}_X}(\NL_{X/k}, \mathcal{O}_X)
$$
D'après Compléments sur les morphismes d'espaces, lemme
\ref{spaces-more-morphisms-lemma-netherlander-fp}, et le fait que $X$ est
noethérien, nous voyons que les faisceaux de cohomologie de $\NL_{X/k}$ sont
cohérents et nuls sauf aux degrés $0$ et $-1$.
D'après Catégories dérivées d'espaces, lemme
\ref{spaces-perfect-lemma-ext-finite}, les groupes $\Ext$ affichés sont des
espaces vectoriels de dimension finie sur $k$, ce qui achève la démonstration.
\end{proof}

\noindent
Prenons garde au fait que l'ouverture de la versalité (démontrée dans le lemme
suivant) est un peu étrange, car notre champ ne satisfait pas à l'effectivité
formelle, voir Exemples, section
\ref{examples-section-proper-spaces-not-algebraic}.
Plus tard, nous appliquerons l'ouverture de la versalité à des sous-champs
convenables de $\Spacesstack'_{fp, flat, proper}$ qui satisfont à l'effectivité
formelle, afin de conclure que ces champs sont algébriques.

\begin{lemma}
\label{quot-lemma-spaces-defo-thy}
Le champ en groupoïdes
$\mathcal{X} = \Spacesstack'_{fp, flat, proper}$ satisfait à l'ouverture de la
versalité sur $\Spec(\mathbf{Z})$.
De même, après changement de base (remarque \ref{quot-remark-spaces-base-change}),
l'ouverture de la versalité vaut sur tout schéma de base noethérien $S$.
\end{lemma}

\begin{proof}
Pour la démonstration « usuelle » de ce fait, nous renvoyons à la discussion de
la remarque qui suit cette démonstration. Nous allons le démontrer à l'aide de
Axiomes d'Artin, lemme \ref{artin-lemma-SGE-implies-openness-versality}.
Nous avons déjà vu que la diagonale de $\mathcal{X}$ est représentable par des
espaces algébriques, qu'il satisfait à (RS*) et qu'il commute aux
limites ; voir les lemmes \ref{quot-lemma-spaces-diagonal},
\ref{quot-lemma-spaces-RS-star} et
\ref{quot-lemma-spaces-limits}.
Il nous suffit donc de voir que $\mathcal{X}$ satisfait à l'effectivité
formelle forte formulée dans
Axiomes d'Artin, lemme \ref{artin-lemma-SGE-implies-openness-versality}.

\medskip\noindent
Soit $(R_n)$ un système projectif d'anneaux tel que $R_n \to R_m$ soit
surjectif à noyau de carré nul pour tous $n \geq m$. Soit
$X_n \to \Spec(R_n)$ un morphisme de présentation finie, plat et propre, où
$X_n$ est un espace algébrique, et soit $X_{n + 1} \to X_n$ un morphisme
sur $\Spec(R_{n + 1})$ induisant un isomorphisme
$X_n = X_{n + 1} \times_{\Spec(R_{n + 1})} \Spec(R_n)$.
Nous devons trouver un morphisme plat, propre et de présentation finie
$X \to \Spec(\lim R_n)$ dont la source est un espace algébrique, tel que
$X_n$ soit le changement de base de $X$ pour tout $n$.

\medskip\noindent
Soit $I_n = \Ker(R_n \to R_1)$. Nous pouvons considérer
$(X_1 \subset X_n) \to (\Spec(R_1) \subset \Spec(R_n))$
comme un morphisme d'épaississements du premier ordre. (Nous prions le lecteur
de parcourir une partie du contenu consacré aux épaississements d'espaces
algébriques dans Compléments sur les morphismes d'espaces, section
\ref{spaces-more-morphisms-section-thickenings}, avant de poursuivre.)
Le faisceau structural de $X_n$ est une extension
$$
0 \to \mathcal{O}_{X_1} \otimes_{R_1} I_n \to
\mathcal{O}_{X_n} \to \mathcal{O}_{X_1} \to 0
$$
au-dessus de $0 \to I_n \to R_n \to R_1$, voir
Compléments sur les morphismes d'espaces, lemme
\ref{spaces-more-morphisms-lemma-deform}.
Considérons l'extension
$$
0 \to \lim \mathcal{O}_{X_1} \otimes_{R_1} I_n \to
\lim \mathcal{O}_{X_n} \to \mathcal{O}_{X_1} \to 0
$$
au-dessus de $0 \to \lim I_n \to \lim R_n \to R_1 \to 0$.
La suite affichée est exacte, car le $R^1\lim$ du système de noyaux est nul
d'après Catégories dérivées d'espaces, lemme
\ref{spaces-perfect-lemma-Rlim-quasi-coherent}.
Observons que l'application
$$
\mathcal{O}_{X_1} \otimes_{R_1} \lim I_n \longrightarrow
\lim \mathcal{O}_{X_1} \otimes_{R_1} I_n
$$
induit un isomorphisme après application du foncteur $DQ_X$, voir
Catégories dérivées d'espaces, lemme
\ref{spaces-perfect-lemma-pullback-and-limits}.
Nous obtenons donc une unique extension
$$
0 \to \mathcal{O}_{X_1} \otimes_{R_1} \lim I_n \to
\mathcal{O}' \to \mathcal{O}_{X_1} \to 0
$$
au-dessus de $0 \to \lim I_n \to \lim R_n \to R_1 \to 0$,
par l'équivalence de catégories de
Théorie des déformations, lemme
\ref{defos-lemma-thickening-over-thickening-space-quasi-coherent}.
Le faisceau $\mathcal{O}'$ détermine un épaississement du premier ordre
d'espaces algébriques $X_1 \subset X$ au-dessus de
$\Spec(R_1) \subset \Spec(\lim R_n)$, d'après
Compléments sur les morphismes d'espaces, lemme
\ref{spaces-more-morphisms-lemma-first-order-thickening}.
Observons que $X \to \Spec(\lim R_n)$ est plat d'après le lemme déjà utilisé
de Compléments sur les morphismes d'espaces,
\ref{spaces-more-morphisms-lemma-deform}.
D'après Compléments sur les morphismes d'espaces, lemme
\ref{spaces-more-morphisms-lemma-deform-property},
$X \to \Spec(\lim R_n)$ est propre et de présentation finie.
Cela achève la démonstration.
\end{proof}

\begin{remark}
\label{quot-remark-spaces-defo-thy}
Le lemme \ref{quot-lemma-spaces-defo-thy} peut aussi se démontrer soit à l'aide de
Axiomes d'Artin, lemme \ref{artin-lemma-dual-openness}
(comme dans la première démonstration du
lemme \ref{quot-lemma-coherent-defo-thy}), soit à l'aide d'une théorie des
obstructions comme dans Axiomes d'Artin, lemme
\ref{artin-lemma-get-openness-obstruction-theory}
(comme dans la seconde démonstration du
lemme \ref{quot-lemma-coherent-defo-thy}).
Dans les deux cas, on utilise la théorie des déformations et des obstructions
développée dans Cotangent, section
\ref{cotangent-section-deformations-ringed-topoi}, pour traduire les propriétés
nécessaires des déformations et des obstructions en groupes $\Ext$ auxquels
Catégories dérivées d'espaces, lemme
\ref{spaces-perfect-lemma-compute-ext}, peut être appliqué.
La seconde méthode (qui utilise une théorie des obstructions et donc le
complexe cotangent tout entier) est peut-être la méthode « standard » employée
dans la plupart des références.
\end{remark}

\section{Le champ des schémas propres polarisés}
\label{quot-section-polarized}

\noindent
Pour étudier le champ des schémas propres polarisés, il suffit de travailler
sur $\mathbf{Z}$, puisque nous pourrons ensuite effectuer l'image inverse sur
tout schéma ou espace algébrique voulu (voir la remarque
\ref{quot-remark-polarized-base-change}).

\begin{situation}
\label{quot-situation-polarized}
Nous définissons une catégorie $\Polarizedstack$ comme suit. Ses objets sont
les couples $(X \to S, \mathcal{L})$ où
\begin{enumerate}
\item $X \to S$ est un morphisme de schémas propre, plat et de présentation
finie, et
\item $\mathcal{L}$ est un $\mathcal{O}_X$-module inversible relativement
ample sur $X/S$
(Morphismes, définition \ref{morphisms-definition-relatively-ample}).
\end{enumerate}
Un morphisme $(X' \to S', \mathcal{L}') \to (X \to S, \mathcal{L})$
entre objets est donné par un triplet $(f, g, \varphi)$, où
$f : X' \to X$ et $g : S' \to S$ sont des morphismes de schémas qui
s'inscrivent dans un diagramme commutatif
$$
\xymatrix{
X' \ar[d] \ar[r]_f & X \ar[d] \\
S' \ar[r]^g & S
}
$$
induisant un isomorphisme $X' \to S' \times_S X$ ; autrement dit, le
diagramme est cartésien, et
$\varphi : f^*\mathcal{L} \to \mathcal{L}'$ est un isomorphisme.
La composition est définie de manière évidente (voir
Exemples de champs, sections
\ref{examples-stacks-section-stack-of-spaces} et
\ref{examples-stacks-section-stack-of-quasi-coherent-sheaves}).
Le foncteur d'oubli
$$
p : \Polarizedstack \longrightarrow \Sch_{fppf},\quad
(X \to S, \mathcal{L}) \longmapsto S
$$
est celui par lequel nous considérons $\Polarizedstack$ comme une catégorie
sur $\Sch_{fppf}$
(voir la section \ref{quot-section-conventions} pour les notations).
\end{situation}

\noindent
Dans la section précédente, nous avons déjà accompli une partie substantielle
du travail sur le champ $\Spacesstack'_{fp, flat, proper}$ des espaces
algébriques de présentation finie, plats et propres. Pour utiliser ces
résultats, nous considérons le foncteur d'oubli
\begin{equation}
\label{quot-equation-over-proper-spaces}
\Polarizedstack \longrightarrow
\Spacesstack'_{fp, flat, proper},\quad
(X \to S, \mathcal{L}) \longmapsto (X \to S)
\end{equation}
Ce foncteur sera un outil utile dans la suite.
Observons que, si $(X \to S)$ appartient à l'image essentielle de
(\ref{quot-equation-over-proper-spaces}), alors $X$ et $S$ sont des schémas.

\begin{lemma}
\label{quot-lemma-polarized-fibred-in-groupoids}
La catégorie $\Polarizedstack$ est fibrée en groupoïdes sur
$\Spacesstack'_{fp, flat, proper}$.
La catégorie $\Polarizedstack$ est fibrée en groupoïdes sur $\Sch_{fppf}$.
\end{lemma}

\begin{proof}
Nous vérifions les conditions (1) et (2) de
Catégories, définition \ref{categories-definition-fibred-groupoids}.

\medskip\noindent
Condition (1). Soit $(X \to S, \mathcal{L})$ un objet de
$\Polarizedstack$ et soit $(X' \to S') \to (X \to S)$ un morphisme de
$\Spacesstack'_{fp, flat, proper}$. Soit alors $\mathcal{L}'$ l'image inverse
de $\mathcal{L}$ sur $X'$.
Observons que $X, S, S'$ sont des schémas ; par conséquent, $X'$ est aussi un
schéma (en tant que produit fibré de schémas). Alors $\mathcal{L}'$ est ample
sur $X'/S'$ d'après
Morphismes, lemme \ref{morphisms-lemma-ample-base-change}.
Nous obtenons ainsi un morphisme
$(X' \to S', \mathcal{L}') \to (X \to S, \mathcal{L})$
au-dessus de $(X' \to S') \to (X \to S)$.

\medskip\noindent
Condition (2). Considérons des morphismes
$(f, g, \varphi) : (X' \to S', \mathcal{L}') \to (X \to S, \mathcal{L})$ et
$(a, b, \psi) : (Y \to T, \mathcal{N}) \to (X \to S, \mathcal{L})$
de $\Polarizedstack$. Étant donné un morphisme
$(k, h) : (Y \to T) \to (X' \to S')$ de
$\Spacesstack'_{fp, flat, proper}$ tel que
$(f, g) \circ (k, h) = (a, b)$, nous devons montrer qu'il existe un unique
morphisme
$(k, h, \chi) : (Y \to T, \mathcal{N}) \to (X' \to S', \mathcal{L}')$
de $\Polarizedstack$ tel que
$(f, g, \varphi) \circ (k, h, \chi) = (a, b, \psi)$.
Nous pouvons simplement prendre
$$
\chi = \psi \circ (k^*\varphi)^{-1}
$$
Cela prouve la condition (2). Une composée de foncteurs définissant des
catégories fibrées définit une catégorie fibrée, voir
Catégories, lemme \ref{categories-lemma-fibred-over-fibred}.
Nous voyons ainsi que $\Polarizedstack$ est fibrée en groupoïdes sur
$\Sch_{fppf}$ (à strictement parler, nous devrions vérifier que les catégories
fibres sont des groupoïdes et appliquer
Catégories, lemme \ref{categories-lemma-fibred-groupoids}).
\end{proof}

\begin{lemma}
\label{quot-lemma-polarized-stack}
La catégorie $\Polarizedstack$ est un champ en groupoïdes sur
$\Spacesstack'_{fp, flat, proper}$ (muni de la topologie induite, voir
Champs, définition \ref{stacks-definition-topology-inherited}).
La catégorie $\Polarizedstack$ est un champ en groupoïdes sur $\Sch_{fppf}$.
\end{lemma}

\begin{proof}
Nous démontrons que $\Polarizedstack$ est un champ en groupoïdes sur
$\Spacesstack'_{fp, flat, proper}$ en vérifiant les conditions (1), (2) et (3)
de Champs, définition \ref{stacks-definition-stack-in-groupoids}.
Nous avons déjà vu (1) au lemme
\ref{quot-lemma-polarized-fibred-in-groupoids}.

\medskip\noindent
Un recouvrement de $\Spacesstack'_{fp, flat, proper}$ s'obtient comme suit :
soit $X \to S$ un objet de $\Spacesstack'_{fp, flat, proper}$.
Supposons que $\{S_i \to S\}_{i \in I}$ soit un recouvrement de
$\Sch_{fppf}$. Posons $X_i = S_i \times_S X$. Alors
$\{(X_i \to S_i) \to (X \to S)\}_{i \in I}$ est un recouvrement de
$\Spacesstack'_{fp, flat, proper}$, et tout recouvrement de
$\Spacesstack'_{fp, flat, proper}$ est isomorphe à l'un de ceux-ci.
Posons $S_{ij} = S_i \times_S S_j$ et $X_{ij} = S_{ij} \times_S X$, de sorte
que $(X_{ij} \to S_{ij}) =
(X_i \to S_i) \times_{(X \to S)} (X_j \to S_j)$.
Supposons ensuite que $\mathcal{L}, \mathcal{N}$ soient des faisceaux
inversibles amples sur $X/S$, de sorte que
$(X \to S, \mathcal{L})$ et $(X \to S, \mathcal{N})$ soient deux objets de
$\Polarizedstack$ au-dessus de l'objet $(X \to S)$.
Pour vérifier la descente des morphismes, supposons donnés des morphismes
$(\text{id}, \text{id}, \varphi_i)$ de
$(X_i \to S_i, \mathcal{L}|_{X_i})$ vers
$(X_i \to S_i, \mathcal{N}|_{X_i})$
dont les changements de base en morphismes de
$(X_{ij} \to S_{ij}, \mathcal{L}|_{X_{ij}})$ vers
$(X_{ij} \to S_{ij}, \mathcal{N}|_{X_{ij}})$
coïncident. Alors
$\varphi_i : \mathcal{L}|_{X_i} \to \mathcal{N}|_{X_i}$
sont des isomorphismes de modules inversibles sur $X_i$ tels que les
restrictions de $\varphi_i$ et $\varphi_j$ à $X_{ij}$ soient les mêmes
isomorphismes.
Par descente des faisceaux quasi-cohérents
(Descente sur les espaces, proposition
\ref{spaces-descent-proposition-fpqc-descent-quasi-coherent}),
nous obtenons un unique isomorphisme $\varphi : \mathcal{L} \to \mathcal{N}$
dont la restriction à $X_i$ redonne $\varphi_i$.

\medskip\noindent
La descente des objets se démontre exactement de la même manière.
Supposons en effet que
$\{(X_i \to S_i) \to (X \to S)\}_{i \in I}$ soit un recouvrement de
$\Spacesstack'_{fp, flat, proper}$ comme ci-dessus.
Supposons donnés des objets $(X_i \to S_i, \mathcal{L}_i)$ de
$\Polarizedstack$ au-dessus de $(X_i \to S_i)$ et une donnée de descente
$$
(\text{id}, \text{id}, \varphi_{ij}) :
(X_{ij} \to S_{ij}, \mathcal{L}_i|_{X_{ij}})
\to
(X_{ij} \to S_{ij}, \mathcal{L}_j|_{X_{ij}})
$$
satisfaisant à la condition de cocycle évidente sur
$(X_{ijk} \to S_{ijk})$ pour tout triplet d'indices.
Alors, par descente des faisceaux quasi-cohérents
(Descente sur les espaces, proposition
\ref{spaces-descent-proposition-fpqc-descent-quasi-coherent}),
nous obtenons un unique $\mathcal{O}_X$-module inversible $\mathcal{L}$ et des
isomorphismes $\mathcal{L}|_{X_i} \to \mathcal{L}_i$ redonnant la donnée de
descente $\varphi_{ij}$.
Pour montrer que $(X \to S, \mathcal{L})$ est un objet de
$\Polarizedstack$, il nous reste à prouver que $\mathcal{L}$ est ample.
Cela résulte de Descente sur les espaces, lemme
\ref{spaces-descent-lemma-descending-property-ample}.

\medskip\noindent
Puisque nous avons déjà vu que $\Spacesstack'_{fp, flat, proper}$ est un champ
en groupoïdes sur $\Sch_{fppf}$ (lemme \ref{quot-lemma-spaces-stack}), il s'ensuit
maintenant formellement que $\Polarizedstack$ est un champ en groupoïdes sur
$\Sch_{fppf}$. Voir Champs, lemme \ref{stacks-lemma-stack-over-stack}.
\end{proof}

\noindent
Vérification de cohérence : le champ $\Polarizedstack$ joue le même rôle parmi
les espaces algébriques.

\begin{lemma}
\label{quot-lemma-extend-polarized-to-spaces}
Soit $T$ un espace algébrique sur $\mathbf{Z}$. Soit $\mathcal{S}_T$
le champ algébrique correspondant (Champs algébriques, sections
\ref{algebraic-section-split},
\ref{algebraic-section-representable-by-algebraic-spaces} et
\ref{algebraic-section-stacks-spaces}).
Nous avons une équivalence de catégories
$$
\left\{
\begin{matrix}
(X \to T, \mathcal{L})\text{ où }X \to T\text{ est un morphisme}\\
\text{d'espaces algébriques, propre, plat et de}\\
\text{présentation finie, et }\mathcal{L}\text{ ample sur }X/T
\end{matrix}
\right\}
\longrightarrow
\Mor_{\textit{Cat}/\Sch_{fppf}}(\mathcal{S}_T, \Polarizedstack)
$$
\end{lemma}

\begin{proof}
Omis. Indications : raisonner exactement comme dans la démonstration du
lemme \ref{quot-lemma-extend-spaces-to-spaces} et utiliser
Descente sur les espaces, proposition
\ref{spaces-descent-proposition-fpqc-descent-quasi-coherent}, pour faire
descendre le faisceau inversible dans la construction du foncteur
quasi-inverse. La propriété d'amplitude relative descend d'après
Descente sur les espaces, lemme
\ref{spaces-descent-lemma-descending-property-ample}.
\end{proof}

\begin{remark}
\label{quot-remark-polarized-base-change}
Soit $B$ un espace algébrique sur $\Spec(\mathbf{Z})$.
Soit $B\textit{-Polarized}$ la catégorie formée des triplets
$(X \to S, \mathcal{L}, h : S \to B)$, où
$(X \to S, \mathcal{L})$ est un objet de $\Polarizedstack$ et
$h : S \to B$ est un morphisme.
Un morphisme
$(X' \to S', \mathcal{L}', h') \to (X \to S, \mathcal{L}, h)$ dans
$B\textit{-Polarized}$ est un morphisme $(f, g, \varphi)$ dans
$\Polarizedstack$ tel que $h \circ g = h'$.
Dans cette situation, le diagramme
$$
\xymatrix{
B\textit{-Polarized} \ar[r] \ar[d] & \Polarizedstack \ar[d] \\
(\Sch/B)_{fppf} \ar[r] & \Sch_{fppf}
}
$$
est un carré $2$-cartésien. Cette remarque triviale sera parfois utile
pour déduire des résultats du cas absolu $\Polarizedstack$ au cas des familles
sur un espace algébrique de base donné.
\end{remark}

\begin{lemma}
\label{quot-lemma-polarized-to-spaces-algebraic}
Le foncteur (\ref{quot-equation-over-proper-spaces}) définit un $1$-morphisme
$$
\Polarizedstack \to \Spacesstack'_{fp, flat, proper}
$$
de champs en groupoïdes sur $\Sch_{fppf}$ qui est algébrique au sens de
Critères de représentabilité, définition
\ref{criteria-definition-algebraic}.
\end{lemma}

\begin{proof}
D'après les lemmes \ref{quot-lemma-spaces-stack} et
\ref{quot-lemma-polarized-stack}, l'énoncé a un sens. Pour le démontrer,
choisissons un schéma $S$ et un objet
$\xi = (X \to S)$ de $\Spacesstack'_{fp, flat, proper}$ au-dessus de $S$.
Nous devons montrer que
$$
\mathcal{X} = (\Sch/S)_{fppf} \times_{\xi, \Spacesstack'_{fp, flat, proper}}
\Polarizedstack
$$
est un champ algébrique sur $S$. Observons qu'un objet de $\mathcal{X}$ est
donné par un couple $(T/S, \mathcal{L})$, où $T$ est un schéma sur $S$ et où
$\mathcal{L}$ est un $\mathcal{O}_{X_T}$-module inversible qui est ample sur
$X_T/T$. Les morphismes sont définis de manière évidente.
En particulier, nous voyons immédiatement que nous avons une inclusion
$$
\mathcal{X} \subset \Picardstack_{X/S}
$$
de catégories sur $(\Sch/S)_{fppf}$, induisant l'égalité sur les ensembles de
morphismes. Puisque $\Picardstack_{X/S}$ est un champ algébrique d'après la
proposition \ref{quot-proposition-pic}, il suffit de montrer que l'inclusion
ci-dessus est représentable par des immersions ouvertes. C'est exactement le
contenu de Descente sur les espaces, lemme
\ref{spaces-descent-lemma-ample-in-neighbourhood}.
\end{proof}

\begin{lemma}
\label{quot-lemma-polarized-diagonal}
La diagonale
$$
\Delta : \Polarizedstack \longrightarrow
\Polarizedstack \times \Polarizedstack
$$
est représentable par des espaces algébriques.
\end{lemma}

\begin{proof}
C'est une conséquence formelle des lemmes
\ref{quot-lemma-polarized-to-spaces-algebraic} et
\ref{quot-lemma-spaces-diagonal}.
Voir Critères de représentabilité, lemme
\ref{criteria-lemma-diagonals-and-algebraic-morphisms}.
\end{proof}

\begin{lemma}
\label{quot-lemma-polarized-limits}
Le champ en groupoïdes $\Polarizedstack$ commute aux limites
(Axiomes d'Artin, définition \ref{artin-definition-limit-preserving}).
\end{lemma}

\begin{proof}
Soit $I$ un ensemble filtrant et soit $(A_i, \varphi_{ii'})$ un système
d'anneaux indexé par $I$. Posons $S = \Spec(A)$ et $S_i = \Spec(A_i)$.
Nous devons montrer que, sur les catégories fibres, nous avons
$$
\Polarizedstack_S = \colim \Polarizedstack_{S_i}
$$
Nous savons que la catégorie des schémas de présentation finie sur $S$ est la
limite inductive des catégories des schémas de présentation finie sur les
$S_i$, voir
Limites, lemme \ref{limits-lemma-descend-finite-presentation}.
De plus, étant donné $X_i \to S_i$ de présentation finie, de limite
$X \to S$, la catégorie des $\mathcal{O}_X$-modules inversibles
$\mathcal{L}$ est la limite inductive des catégories des
$\mathcal{O}_{X_i}$-modules inversibles $\mathcal{L}_i$, voir
Limites, lemmes \ref{limits-lemma-descend-modules-finite-presentation} et
\ref{limits-lemma-descend-invertible-modules}.
Si $X \to S$ est propre et plat, alors, pour $i$ assez grand, le morphisme
$X_i \to S_i$ est également propre et plat, voir
Limites, lemmes \ref{limits-lemma-eventually-proper} et
\ref{limits-lemma-descend-flat-finite-presentation}.
Enfin, si $\mathcal{L}$ est ample sur $X$, alors $\mathcal{L}_i$ est ample sur
$X_i$ pour $i$ assez grand, voir
Limites, lemme \ref{limits-lemma-limit-ample}.
La réunion de ces faits achève la démonstration.
\end{proof}

\begin{lemma}
\label{quot-lemma-polarized-RS-star}
Dans la situation \ref{quot-situation-coherent}. Soit
$$
\xymatrix{
T \ar[r] \ar[d] & T' \ar[d] \\
S \ar[r] & S'
}
$$
une somme amalgamée dans la catégorie des schémas, où
$T \to T'$ est un épaississement et $T \to S$ est affine, voir
Compléments sur les morphismes, lemme
\ref{more-morphisms-lemma-pushout-along-thickening}.
Alors le foncteur sur les catégories fibres
$$
\Polarizedstack_{S'}
\longrightarrow
\Polarizedstack_S \times_{\Polarizedstack_T} \Polarizedstack_{T'}
$$
est une équivalence.
\end{lemma}

\begin{proof}
D'après Compléments sur les morphismes, lemme
\ref{more-morphisms-lemma-equivalence-categories-schemes-over-pushout-flat},
il existe une équivalence
$$
\textit{flat-lfp}_{S'}
\longrightarrow
\textit{flat-lfp}_S \times_{\textit{flat-lfp}_T} \textit{flat-lfp}_{T'}
$$
où $\textit{flat-lfp}_S$ désigne la catégorie des schémas plats et localement
de présentation finie sur $S$.
Soit $X'/S'$ dans le membre de gauche correspondant au triplet
$(X/S, Y'/T', \varphi)$ du membre de droite.
Posons $Y = T \times_{T'} Y'$, qui est isomorphe à $T \times_S X$ par
$\varphi$. Alors Compléments sur les morphismes, lemme
\ref{more-morphisms-lemma-scheme-over-pushout-flat-modules}, montre que nous
avons une équivalence
$$
\textit{QCoh-flat}_{X'/S'}
\longrightarrow
\textit{QCoh-flat}_{X/S}
\times_{\textit{QCoh-flat}_{Y/T}} \textit{QCoh-flat}_{Y'/T'}
$$
où $\textit{QCoh-flat}_{X/S}$ désigne la catégorie des
$\mathcal{O}_X$-modules quasi-cohérents plats sur $S$.
Puisque $X \to S$, $Y \to T$, $X' \to S'$ et $Y' \to T'$ sont plats, cela
s'applique en particulier aux modules inversibles et donne une équivalence de
catégories
$$
\textit{Pic}(X')
\longrightarrow
\textit{Pic}(X) \times_{\textit{Pic}(Y)} \textit{Pic}(Y')
$$
où $\textit{Pic}(X)$ désigne la catégorie des $\mathcal{O}_X$-modules
inversibles. Il y a ici un petit point à vérifier : il faut montrer que, si un
objet $\mathcal{F}'$ de $\textit{QCoh-flat}_{X'/S'}$ a pour images inverses des
modules inversibles sur $X$ et $Y'$, alors $\mathcal{F}'$ est un
$\mathcal{O}_{X'}$-module inversible.
Il résulte du lemme cité que $\mathcal{F}'$ est un
$\mathcal{O}_{X'}$-module de présentation finie.
D'après Compléments sur les morphismes, lemme
\ref{more-morphisms-lemma-flat-and-free-at-point-fibre}, il suffit de vérifier
que la restriction de $\mathcal{F}'$ aux fibres de $X' \to S'$ est inversible.
Mais les fibres de $X' \to S'$ sont les mêmes que celles de $X \to S$ ; ces
restrictions sont donc inversibles.

\medskip\noindent
Compte tenu de ce qui précède, nous obtenons une équivalence de catégories si
nous supprimons l'hypothèse (pour la catégorie des objets au-dessus de $S$) que
$X \to S$ soit propre et l'hypothèse que $\mathcal{L}$ soit ample.
Il est maintenant clair que, si $X' \to S'$ est propre, alors
$X \to S$ et $Y' \to T'$ sont propres (Morphismes, lemme
\ref{morphisms-lemma-base-change-proper}).
Réciproquement, si $X \to S$ et $Y' \to T'$ sont propres, alors
$X' \to S'$ est propre d'après
Compléments sur les morphismes, lemme
\ref{more-morphisms-lemma-thicken-property-morphisms-cartesian}.
De même, si $\mathcal{L}'$ est ample sur $X'/S'$, alors
$\mathcal{L}'|_X$ est ample sur $X/S$ et
$\mathcal{L}'|_{Y'}$ est ample sur $Y'/T'$
(Morphismes, lemme \ref{morphisms-lemma-ample-base-change}).
Enfin, si $\mathcal{L}'|_X$ est ample sur $X/S$ et si
$\mathcal{L}'|_{Y'}$ est ample sur $Y'/T'$, alors
$\mathcal{L}'$ est ample sur $X'/S'$ d'après
Compléments sur les morphismes, lemme
\ref{more-morphisms-lemma-thicken-property-relatively-ample}.
\end{proof}

\begin{lemma}
\label{quot-lemma-polarized-tangent-space}
Soit $k$ un corps et soit $x = (X \to \Spec(k), \mathcal{L})$
un objet de $\mathcal{X} = \Polarizedstack$ au-dessus de $\Spec(k)$.
\begin{enumerate}
\item Si $k$ est de type fini sur $\mathbf{Z}$, alors les espaces vectoriels
$T\mathcal{F}_{\mathcal{X}, k, x}$ et
$\text{Inf}(\mathcal{F}_{\mathcal{X}, k, x})$
(voir Axiomes d'Artin, section \ref{artin-section-tangent-spaces})
sont de dimension finie, et
\item en général, les espaces vectoriels $T_x(k)$ et $\text{Inf}_x(k)$
(voir Axiomes d'Artin, section \ref{artin-section-inf})
sont de dimension finie.
\end{enumerate}
\end{lemma}

\begin{proof}
La discussion d'Axiomes d'Artin, section \ref{artin-section-tangent-spaces},
ne s'applique qu'aux corps de type fini sur le schéma de base
$\Spec(\mathbf{Z})$. Notre champ satisfait à (RS*) d'après le
lemme \ref{quot-lemma-polarized-RS-star}, et nous pouvons appliquer
Axiomes d'Artin, lemme \ref{artin-lemma-properties-lift-RS-star},
pour obtenir les espaces vectoriels $T_x(k)$ et $\text{Inf}_x(k)$ mentionnés en
(2). De plus, dans le cas de type fini, ces espaces coïncident avec ceux
mentionnés dans la partie (1), d'après Axiomes d'Artin, remarque
\ref{artin-remark-compare-deformation-spaces}.
Cela étant réglé, nous pouvons commencer la démonstration.

\medskip\noindent
Une démonstration consiste à employer un argument analogue à celui de la
démonstration du lemme \ref{quot-lemma-spaces-tangent-space} ; cela nous obligerait
à développer une théorie des déformations pour les couples formés d'un schéma
et d'un module quasi-cohérent.
Une autre démonstration utiliserait le résultat du
lemme \ref{quot-lemma-spaces-tangent-space}, l'algébricité de
$\Polarizedstack \to \Spacesstack'_{fp, flat, proper}$ et un calcul de l'espace
de déformations d'un module inversible. Nous allons toutefois plutôt traduire
la question en une question de déformation d'algèbres graduées sur $k$, puis en
déduire le résultat.

\medskip\noindent
Soit $\mathcal{C}_k$ la catégorie des $k$-algèbres locales artiniennes $A$ de
corps résiduel $k$. Notre objet $x$ de $\mathcal{X}$ au-dessus de $k$ donne une
catégorie de prédéformations
$p : \mathcal{F} \to \mathcal{C}_k$, voir
Axiomes d'Artin, section \ref{artin-section-predeformation-categories}.
Ainsi, $\mathcal{F}(A)$ est la catégorie des triplets
$(X_A, \mathcal{L}_A, \alpha)$, où $(X_A, \mathcal{L}_A)$ est un objet de
$\Polarizedstack$ au-dessus de $A$ et où $\alpha$ est un isomorphisme
$(X_A, \mathcal{L}_A) \times_{\Spec(A)} \Spec(k) \cong (X, \mathcal{L})$.
D'autre part, soit $q : \mathcal{G} \to \mathcal{C}_k$ la catégorie cofibrée
en groupoïdes définie dans
Problèmes de déformations, exemple
\ref{examples-defos-example-graded-algebras}.
Choisissons $d_0 \gg 0$ (nous verrons plus bas quelle grandeur suffit).
Soit $P$ la $k$-algèbre graduée
$$
P = k \oplus \bigoplus\nolimits_{d \geq d_0} H^0(X, \mathcal{L}^{\otimes d})
$$
Alors $y = (k, P)$ est un objet de $\mathcal{G}(k)$.
Soit $\mathcal{G}_y$ la catégorie de prédéformations de
Théorie formelle des déformations, remarque
\ref{formal-defos-remark-localize-cofibered-groupoid}.
Étant donné $(X_A, \mathcal{F}_A, \alpha)$ comme ci-dessus, posons
$$
Q = A \oplus \bigoplus\nolimits_{d \geq d_0} H^0(X_A, \mathcal{L}_A^{\otimes d})
$$
L'isomorphisme $\alpha$ induit une application $\beta : Q \to P$.
Par la théorie des déformations des schémas projectifs
(Compléments sur les morphismes, lemme
\ref{more-morphisms-lemma-deform-projective}), nous obtenons un
$1$-morphisme
$$
\mathcal{F} \longrightarrow \mathcal{G}_y,\quad
(X_A, \mathcal{F}_A, \alpha) \longmapsto (Q, \beta : Q \to P)
$$
de catégories cofibrées en groupoïdes sur $\mathcal{C}_k$.
En fait, ce foncteur est une équivalence, avec pour quasi-inverse celui donné
par $Q \mapsto \underline{\text{Proj}}_A(Q)$.
En effet, le schéma $X_A = \underline{\text{Proj}}_A(Q)$ est plat sur $A$
d'après Diviseurs, lemme \ref{divisors-lemma-relative-proj-flat}.
Posons $\mathcal{L}_A = \mathcal{O}_{X_A}(1)$ ; celui-ci est plat sur $A$
d'après le même lemme. Nous tirons de $\beta$ un isomorphisme
$(X_A, \mathcal{L}_A) \times_{\Spec(A)} \Spec(k) = (X, \mathcal{L})$.
Nous pouvons alors déduire toutes les propriétés voulues du couple
$(X_A, \mathcal{L}_A)$ des propriétés correspondantes de
$(X, \mathcal{L})$ à l'aide des techniques de
Compléments sur les morphismes, sections
\ref{more-morphisms-section-morphisms-thickenings} et
\ref{more-morphisms-section-deform}.
Nous omettons certains détails.

\medskip\noindent
En conclusion, nous voyons que
$T\mathcal{F} = T\mathcal{G}_y = T_y\mathcal{G}$ et
$\text{Inf}(\mathcal{F}) = \text{Inf}_y(\mathcal{G})$.
Ces espaces vectoriels sont de dimension finie d'après
Problèmes de déformations, lemme
\ref{examples-defos-lemma-graded-algebras-TI}, ce qui achève la démonstration.
\end{proof}

\begin{lemma}[Effectivité formelle forte pour les schémas polarisés]
\label{quot-lemma-polarized-strong-effectiveness}
\begin{slogan}
Le théorème d'algébrisation de Grothendieck reste valable dans le cadre non
noethérien si l'on suppose la platitude et la présentation finie.
\end{slogan}
Soit $(R_n)$ un système projectif d'anneaux à applications de transition
surjectives dont les noyaux sont localement nilpotents. Posons
$R = \lim R_n$.
Posons $S_n = \Spec(R_n)$ et $S = \Spec(R)$. Considérons un diagramme
commutatif
$$
\xymatrix{
X_1 \ar[r]_{i_1} \ar[d] & X_2 \ar[r]_{i_2} \ar[d] & X_3 \ar[r] \ar[d] &
\ldots \\
S_1 \ar[r] & S_2 \ar[r] & S_3 \ar[r] & \ldots
}
$$
de schémas dont les carrés sont cartésiens. Supposons donnés
$(\mathcal{L}_n, \varphi_n)$, où chaque $\mathcal{L}_n$ est un faisceau
inversible sur $X_n$ et où
$\varphi_n : i_n^*\mathcal{L}_{n + 1} \to \mathcal{L}_n$ est un isomorphisme.
Si
\begin{enumerate}
\item $X_n \to S_n$ est propre, plat et de présentation finie, et
\item $\mathcal{L}_1$ est ample sur $X_1$,
\end{enumerate}
alors il existe un morphisme de schémas $X \to S$ propre, plat et de
présentation finie, un $\mathcal{O}_X$-module inversible ample
$\mathcal{L}$, ainsi que des isomorphismes
$X_n \cong X \times_S S_n$ et
$\mathcal{L}_n \cong \mathcal{L}|_{X_n}$ compatibles aux morphismes $i_n$ et
$\varphi_n$.
\end{lemma}

\begin{proof}
Choisissons $d_0$ pour $X_1 \to S_1$ et $\mathcal{L}_1$ comme dans
Compléments sur les morphismes, lemme
\ref{more-morphisms-lemma-deform-projective}.
Pour tout $n \geq 1$, posons
$$
A_n = R_n \oplus
\bigoplus\nolimits_{d \geq d_0} H^0(X_n, \mathcal{L}_n^{\otimes d})
$$
D'après le lemme, chaque $A_n$ est une $R_n$-algèbre graduée de présentation
finie dont les composantes homogènes $(A_n)_d$ sont des $R_n$-modules
projectifs finis, telle que $X_n = \text{Proj}(A_n)$ et
$\mathcal{L}_n = \mathcal{O}_{\text{Proj}(A_n)}(1)$.
Le lemme garantit aussi que les applications
$$
A_1 \leftarrow A_2 \leftarrow A_3 \leftarrow \ldots
$$
induisent des isomorphismes $A_n = A_m \otimes_{R_m} R_n$ pour $n \leq m$.
Nous posons
$$
B = \bigoplus\nolimits_{d \geq 0} B_d
\quad\text{avec}\quad
B_d = \lim_n (A_n)_d
$$
D'après Compléments d'algèbre, lemme
\ref{more-algebra-lemma-lim-finite-projective-gives-finite-projective},
$B_d$ est un $R$-module projectif fini et $B \otimes_R R_n = A_n$.
Ainsi, le schéma
$$
X = \text{Proj}(B)
\quad\text{et}\quad
\mathcal{L} = \mathcal{O}_X(1)
$$
est plat sur $S$, et $\mathcal{L}$ est un $\mathcal{O}_X$-module
quasi-cohérent plat sur $S$, voir
Diviseurs, lemme \ref{divisors-lemma-relative-proj-flat}.
Comme la formation de Proj commute au changement de base
(Constructions, lemme \ref{constructions-lemma-base-change-map-proj}),
nous obtenons des isomorphismes canoniques
$$
X \times_S S_n = X_n
\quad\text{et}\quad
\mathcal{L}|_{X_n} \cong \mathcal{L}_n
$$
compatibles aux applications de transition du système.
Nous pouvons donc considérer $X_1 \subset X$ comme un sous-schéma fermé.
Nous montrerons plus bas que $B$ est de présentation finie sur $R$.
D'après Diviseurs, lemmes \ref{divisors-lemma-relative-proj-proper} et
\ref{divisors-lemma-relative-proj-finite-presentation},
cela implique que $X \to S$ est de présentation finie et propre, et que
$\mathcal{L} = \mathcal{O}_X(1)$ est de présentation finie en tant que
$\mathcal{O}_X$-module.
Puisque la restriction de $\mathcal{L}$ au changement de base
$X_1 \to S_1$ est inversible, il résulte de
Compléments sur les morphismes, lemme
\ref{more-morphisms-lemma-finite-free-open}, que $\mathcal{L}$ est inversible
sur un voisinage ouvert de $X_1$ dans $X$.
Comme $X \to S$ est fermé et comme $\Ker(R \to R_1)$ est contenu dans le
radical de Jacobson
(Compléments d'algèbre, lemme \ref{more-algebra-lemma-limit-henselian}),
nous voyons que tout voisinage ouvert de $X_1$ dans $X$ est égal à $X$.
Ainsi, $\mathcal{L}$ est inversible. Enfin, l'ensemble des points de $S$ où
$\mathcal{L}$ est ample sur la fibre est ouvert dans $S$
(Compléments sur les morphismes, lemme
\ref{more-morphisms-lemma-ample-in-neighbourhood})
et contient $S_1$ ; il est donc égal à $S$. Ainsi, $X \to S$ et
$\mathcal{L}$ possèdent toutes les propriétés requises dans l'énoncé du lemme.

\medskip\noindent
Démontrons l'assertion ci-dessus.
Choisissons une présentation
$A_1 = R_1[X_1, \ldots, X_s]/(F_1, \ldots, F_t)$ où les $X_i$ sont des
variables de degrés $d_i$ et les $F_j$ des polynômes homogènes en les $X_i$ de
degré $e_j$.
Nous pouvons alors choisir une application
$$
\Psi : R[X_1, \ldots, X_s] \longrightarrow B
$$
relevant l'application $R_1[X_1, \ldots, X_s] \to A_1$.
Puisque chaque $B_d$ est projectif fini sur $R$, nous concluons du lemme de
Nakayama (Algèbre, lemme \ref{algebra-lemma-NAK}, en utilisant de nouveau le
fait que $\Ker(R \to R_1)$ est contenu dans le radical de Jacobson de $R$) que
$\Psi$ est surjective. Comme $- \otimes_R R_1$ est exact à droite, nous pouvons
trouver $G_1, \ldots, G_t \in \Ker(\Psi)$ dont les images sont
$F_1, \ldots, F_t$ dans $R_1[X_1, \ldots, X_s]$.
Observons que $\Ker(\Psi)_d$ est un $R$-module projectif fini pour tout
$d \geq 0$, en tant que noyau de la surjection
$R[X_1, \ldots, X_s]_d \to B_d$ de $R$-modules projectifs finis.
Nous concluons encore une fois du lemme de Nakayama que $\Ker(\Psi)$ est
engendré par $G_1, \ldots, G_t$.
\end{proof}

\begin{lemma}
\label{quot-lemma-polarized-existence}
Considérons le champ $\Polarizedstack$ sur le schéma de base
$\Spec(\mathbf{Z})$. Alors tout objet formel est effectif.
\end{lemma}

\begin{proof}
Pour les définitions des notions du lemme, voir
Axiomes d'Artin, section \ref{artin-section-formal-objects}.
Il ressort des définitions que le lemme résulte immédiatement du lemme plus
général \ref{quot-lemma-polarized-strong-effectiveness}.
\end{proof}

\begin{lemma}
\label{quot-lemma-polarized-defo-thy}
Le champ en groupoïdes $\Polarizedstack$ satisfait à l'ouverture de la
versalité sur $\Spec(\mathbf{Z})$.
De même, après changement de base (remarque
\ref{quot-remark-polarized-base-change}), l'ouverture de la versalité vaut sur tout
schéma de base noethérien $S$.
\end{lemma}

\begin{proof}
Cela résulte de
Axiomes d'Artin, lemme \ref{artin-lemma-SGE-implies-openness-versality},
et des lemmes \ref{quot-lemma-polarized-diagonal},
\ref{quot-lemma-polarized-RS-star},
\ref{quot-lemma-polarized-limits} et
\ref{quot-lemma-polarized-strong-effectiveness}.
Pour la démonstration « usuelle » de ce fait, voir la discussion de la remarque
qui suit cette démonstration.
\end{proof}

\begin{remark}
\label{quot-remark-polarized-defo-thy}
Le lemme \ref{quot-lemma-polarized-defo-thy} peut également se démontrer à l'aide
d'une théorie des obstructions comme dans
Axiomes d'Artin, lemme \ref{artin-lemma-get-openness-obstruction-theory}
(comme dans la seconde démonstration du
lemme \ref{quot-lemma-coherent-defo-thy}).
Pour cela, il faut généraliser la théorie des déformations et des obstructions
développée dans
Cotangent, section \ref{cotangent-section-deformations-ringed-topoi}, au cas des
couples formés d'espaces algébriques et de modules quasi-cohérents.
Une autre possibilité consiste à utiliser le fait que le $1$-morphisme
$\Polarizedstack \to \Spacesstack'_{fp, flat, proper}$ est algébrique
(lemme \ref{quot-lemma-polarized-to-spaces-algebraic}) et le fait que nous
connaissons l'ouverture de la versalité pour le but
(lemme \ref{quot-lemma-spaces-defo-thy} et remarque
\ref{quot-remark-spaces-defo-thy}).
\end{remark}

\begin{theorem}[Algébricité du champ des schémas polarisés]
\label{quot-theorem-polarized-algebraic}
Le champ $\Polarizedstack$ (situation \ref{quot-situation-polarized})
est algébrique. En fait, pour tout espace algébrique $B$, le champ
$B\textit{-Polarized}$ (remarque \ref{quot-remark-polarized-base-change})
est algébrique.
\end{theorem}

\begin{proof}
Le cas absolu résulte de
Axiomes d'Artin, lemme \ref{artin-lemma-diagonal-representable},
et des lemmes \ref{quot-lemma-polarized-diagonal},
\ref{quot-lemma-polarized-RS-star},
\ref{quot-lemma-polarized-limits},
\ref{quot-lemma-polarized-existence} et
\ref{quot-lemma-polarized-defo-thy}.
Le cas sur $B$ en résulte, compte tenu de la description de
$B\textit{-Polarized}$ comme $2$-produit fibré dans la remarque
\ref{quot-remark-polarized-base-change} et du fait que les champs algébriques
admettent les $2$-produits fibrés, voir
Champs algébriques, lemme \ref{algebraic-lemma-2-fibre-product}.
\end{proof}

\section{Le champ des courbes}
\label{quot-section-curves}

\noindent
Dans cette section, nous prouvons que le champ des courbes est algébrique.
Pour une discussion plus approfondie des espaces de modules de courbes, nous
renvoyons le lecteur à Modules de courbes, section
\ref{moduli-curves-section-introduction}.

\medskip\noindent
Dans le Projet Champs, une courbe est une variété de dimension $1$.
Cependant, lorsque nous parlons de familles de courbes, nous autorisons souvent
les fibres à être réductibles et/ou non réduites. Dans cette section, le champ
des courbes « paramétrera les schémas propres de dimension $\leq 1$ ».
Il se trouve toutefois que, pour obtenir la bonne notion de famille, il nous
faut autoriser l'espace total de notre famille à être un espace algébrique.
Cela conduit à la définition suivante.

\begin{situation}
\label{quot-situation-curves}
Nous définissons une catégorie $\Curvesstack$ comme suit :
\begin{enumerate}
\item Les objets sont les {\it familles de courbes}. Plus précisément, un
objet est un morphisme $f : X \to S$ où la base $S$ est un schéma,
l'{\it espace total} $X$ est un espace algébrique, et où
$f$ est plat, propre, de présentation finie et de dimension relative
$\leq 1$ (Morphismes d'espaces, définition
\ref{spaces-morphisms-definition-relative-dimension}).
\item Un morphisme $(X' \to S') \to (X \to S)$ entre objets est donné par un
couple $(f, g)$, où $f : X' \to X$ est un morphisme d'espaces algébriques et
$g : S' \to S$ est un morphisme de schémas, qui s'inscrivent dans un diagramme
commutatif
$$
\xymatrix{
X' \ar[d] \ar[r]_f & X \ar[d] \\
S' \ar[r]^g & S
}
$$
induisant un isomorphisme $X' \to S' \times_S X$ ; autrement dit, le
diagramme est cartésien.
\end{enumerate}
Le foncteur d'oubli
$$
p : \Curvesstack \longrightarrow \Sch_{fppf},\quad
(X \to S) \longmapsto S
$$
est celui par lequel nous considérons $\Curvesstack$ comme une catégorie sur
$\Sch_{fppf}$
(voir la section \ref{quot-section-conventions} pour les notations).
\end{situation}

\noindent
Il résulte de Espaces sur les corps, lemme
\ref{spaces-over-fields-lemma-codim-1-point-in-schematic-locus},
et, plus généralement, de
Compléments sur les morphismes d'espaces, lemme
\ref{spaces-more-morphisms-lemma-projective-over-complete},
que, si $S$ est le spectre d'un corps, d'un anneau local artinien ou d'un
anneau local noethérien complet, alors, pour toute famille de courbes
$X \to S$, l'espace total $X$ est un schéma.
En revanche, il existe des familles de courbes sur $\mathbf{A}^1_k$ dont
l'espace total n'est pas un schéma, voir
Exemples, section \ref{examples-section-family-of-curves}.

\medskip\noindent
Il est clair que
\begin{equation}
\label{quot-equation-curves-over-proper-spaces}
\Curvesstack \subset \Spacesstack'_{fp, flat, proper}
\end{equation}
et qu'un objet $X \to S$ de $\Spacesstack'_{fp, flat, proper}$ appartient à
$\Curvesstack$ si et seulement si $X \to S$ est de dimension relative
$\leq 1$. Nous utiliserons ce fait pour vérifier les axiomes d'Artin pour
$\Curvesstack$.

\begin{lemma}
\label{quot-lemma-curves-fibred-in-groupoids}
La catégorie $\Curvesstack$ est fibrée en groupoïdes sur $\Sch_{fppf}$.
\end{lemma}

\begin{proof}
En utilisant le plongement (\ref{quot-equation-curves-over-proper-spaces}),
la description de son image et le fait correspondant pour
$\Spacesstack'_{fp, flat, proper}$
(lemme \ref{quot-lemma-spaces-fibred-in-groupoids}), nous nous ramenons à l'énoncé
suivant : étant donné un morphisme
$$
\xymatrix{
X' \ar[r] \ar[d] & X \ar[d] \\
S' \ar[r] & S
}
$$
dans $\Spacesstack'_{fp, flat, proper}$ (rappelons que cela implique en
particulier que le diagramme est cartésien), si $X \to S$ est de dimension
relative $\leq 1$, alors $X' \to S'$ est de dimension relative $\leq 1$.
Cela résulte de Morphismes d'espaces, lemme
\ref{spaces-morphisms-lemma-dimension-fibre-after-base-change}.
\end{proof}

\begin{lemma}
\label{quot-lemma-curves-stack}
La catégorie $\Curvesstack$ est un champ en groupoïdes sur $\Sch_{fppf}$.
\end{lemma}

\begin{proof}
En utilisant le plongement (\ref{quot-equation-curves-over-proper-spaces}),
la description de son image et le fait correspondant pour
$\Spacesstack'_{fp, flat, proper}$ (lemme \ref{quot-lemma-spaces-stack}),
nous nous ramenons à l'énoncé suivant : étant donné un objet $X \to S$ de
$\Spacesstack'_{fp, flat, proper}$ et un recouvrement fppf
$\{S_i \to S\}_{i \in I}$, les assertions suivantes sont équivalentes :
\begin{enumerate}
\item $X \to S$ est de dimension relative $\leq 1$, et
\item pour tout $i$, le changement de base $X_i \to S_i$ est de dimension
relative $\leq 1$.
\end{enumerate}
Cela résulte de Morphismes d'espaces, lemme
\ref{spaces-morphisms-lemma-dimension-fibre-after-base-change}.
\end{proof}

\begin{lemma}
\label{quot-lemma-curves-diagonal}
La diagonale
$$
\Delta : \Curvesstack \longrightarrow \Curvesstack \times \Curvesstack
$$
est représentable par des espaces algébriques.
\end{lemma}

\begin{proof}
Cela résulte immédiatement du plongement pleinement fidèle
(\ref{quot-equation-curves-over-proper-spaces}) et du fait correspondant pour
$\Spacesstack'_{fp, flat, proper}$
(lemme \ref{quot-lemma-spaces-diagonal}).
\end{proof}

\begin{remark}
\label{quot-remark-curves-base-change}
Soit $B$ un espace algébrique sur $\Spec(\mathbf{Z})$.
Soit $B\text{-}\Curvesstack$ la catégorie formée des couples
$(X \to S, h : S \to B)$, où $X \to S$ est un objet de $\Curvesstack$ et
$h : S \to B$ est un morphisme.
Un morphisme $(X' \to S', h') \to (X \to S, h)$ dans
$B\text{-}\Curvesstack$ est un morphisme $(f, g)$ dans $\Curvesstack$ tel que
$h \circ g = h'$.
Dans cette situation, le diagramme
$$
\xymatrix{
B\text{-}\Curvesstack \ar[r] \ar[d] & \Curvesstack \ar[d] \\
(\Sch/B)_{fppf} \ar[r] & \Sch_{fppf}
}
$$
est un carré $2$-cartésien. Cette remarque triviale sera parfois utile
pour déduire des résultats du cas absolu $\Curvesstack$ au cas des familles de
courbes sur un espace algébrique de base donné.
\end{remark}

\begin{lemma}
\label{quot-lemma-curves-limits}
Le champ $\Curvesstack \to \Sch_{fppf}$ commute aux limites
(Axiomes d'Artin, définition \ref{artin-definition-limit-preserving}).
\end{lemma}

\begin{proof}
En utilisant le plongement (\ref{quot-equation-curves-over-proper-spaces}),
la description de son image et le fait correspondant pour
$\Spacesstack'_{fp, flat, proper}$ (lemme \ref{quot-lemma-spaces-limits}),
nous nous ramenons à l'énoncé suivant :
soit $T = \lim T_i$ la limite d'un système projectif filtrant de schémas
affines. Soit $i \in I$ et soit $X_i \to T_i$ un objet de
$\Spacesstack'_{fp, flat, proper}$ au-dessus de $T_i$.
Supposons que $T \times_{T_i} X_i \to T$ soit de dimension relative
$\leq 1$. Alors, pour un certain $i' \geq i$, le morphisme
$T_{i'} \times_{T_i} X_i \to T_i$ est de dimension relative $\leq 1$.
Cela résulte de Limites d'espaces, lemme
\ref{spaces-limits-lemma-eventually-relative-dimension}.
\end{proof}

\begin{lemma}
\label{quot-lemma-curves-RS-star}
Soit
$$
\xymatrix{
T \ar[r] \ar[d] & T' \ar[d] \\
S \ar[r] & S'
}
$$
une somme amalgamée dans la catégorie des schémas, où
$T \to T'$ est un épaississement et $T \to S$ est affine, voir
Compléments sur les morphismes, lemme
\ref{more-morphisms-lemma-pushout-along-thickening}.
Alors le foncteur sur les catégories fibres
$$
\Curvesstack_{S'}
\longrightarrow
\Curvesstack_S
\times_{\Curvesstack_T}
\Curvesstack_{T'}
$$
est une équivalence.
\end{lemma}

\begin{proof}
En utilisant le plongement (\ref{quot-equation-curves-over-proper-spaces}),
la description de son image et le fait correspondant pour
$\Spacesstack'_{fp, flat, proper}$ (lemme \ref{quot-lemma-spaces-RS-star}),
nous nous ramenons à l'énoncé suivant : étant donné un morphisme
$X' \to S'$ d'un espace algébrique vers $S'$ qui est de présentation finie,
plat et propre, alors $X' \to S'$ est de dimension relative $\leq 1$ si et
seulement si $S \times_{S'} X' \to S$ et
$T' \times_{S'} X' \to T'$ sont de dimension relative $\leq 1$.
Une implication résulte du fait que la propriété d'être de dimension relative
$\leq 1$ est préservée par changement de base
(Morphismes d'espaces, lemme
\ref{spaces-morphisms-lemma-dimension-fibre-after-base-change}).
L'autre résulte du fait que la propriété d'être de dimension relative
$\leq 1$ se vérifie sur les fibres et du fait que les fibres de $X' \to S'$
(aux points du schéma $S'$) sont les mêmes que celles de
$S \times_{S'} X' \to S$, puisque $S \to S'$ est un épaississement d'après
Compléments sur les morphismes, lemme
\ref{more-morphisms-lemma-pushout-along-thickening}.
\end{proof}

\begin{lemma}
\label{quot-lemma-curves-tangent-space}
Soit $k$ un corps et soit $x = (X \to \Spec(k))$ un objet de
$\mathcal{X} = \Curvesstack$ au-dessus de $\Spec(k)$.
\begin{enumerate}
\item Si $k$ est de type fini sur $\mathbf{Z}$, alors les espaces vectoriels
$T\mathcal{F}_{\mathcal{X}, k, x}$ et
$\text{Inf}(\mathcal{F}_{\mathcal{X}, k, x})$
(voir Axiomes d'Artin, section \ref{artin-section-tangent-spaces})
sont de dimension finie, et
\item en général, les espaces vectoriels $T_x(k)$ et $\text{Inf}_x(k)$
(voir Axiomes d'Artin, section \ref{artin-section-inf})
sont de dimension finie.
\end{enumerate}
\end{lemma}

\begin{proof}
Cela résulte immédiatement du plongement pleinement fidèle
(\ref{quot-equation-curves-over-proper-spaces}) et du fait correspondant pour
$\Spacesstack'_{fp, flat, proper}$
(lemme \ref{quot-lemma-spaces-tangent-space}).
\end{proof}

\begin{lemma}
\label{quot-lemma-curves-existence}
Considérons le champ $\Curvesstack$ sur le schéma de base
$\Spec(\mathbf{Z})$. Alors tout objet formel est effectif.
\end{lemma}

\begin{proof}
Pour les définitions des notions du lemme, voir
Axiomes d'Artin, section \ref{artin-section-formal-objects}.
Soit $(A, \mathfrak m, \kappa)$ un anneau local noethérien complet.
Soit $(X_n \to \Spec(A/\mathfrak m^n))$ un objet formel de
$\Curvesstack$ au-dessus de $A$.
D'après Compléments sur les morphismes d'espaces, lemme
\ref{spaces-more-morphisms-lemma-formal-algebraic-space-proper-reldim-1},
il existe un morphisme projectif $X \to \Spec(A)$ et un système compatible
d'isomorphismes
$X \times_{\Spec(A)} \Spec(A/\mathfrak m^n) \cong X_n$.
D'après Compléments sur les morphismes, lemme
\ref{more-morphisms-lemma-check-flatness-on-infinitesimal-nbhds},
$X \to \Spec(A)$ est plat. D'après Compléments sur les morphismes, lemme
\ref{more-morphisms-lemma-dimension-fibres-proper-flat},
$X \to \Spec(A)$ est de dimension relative $\leq 1$.
Cela prouve le lemme.
\end{proof}

\begin{lemma}
\label{quot-lemma-curves-defo-thy}
Le champ en groupoïdes $\mathcal{X} = \Curvesstack$ satisfait à l'ouverture
de la versalité sur $\Spec(\mathbf{Z})$.
De même, après changement de base (remarque \ref{quot-remark-curves-base-change}),
l'ouverture de la versalité vaut sur tout schéma de base noethérien $S$.
\end{lemma}

\begin{proof}
Cela résulte immédiatement du plongement pleinement fidèle
(\ref{quot-equation-curves-over-proper-spaces}) et du fait correspondant pour
$\Spacesstack'_{fp, flat, proper}$
(lemme \ref{quot-lemma-spaces-defo-thy}).
\end{proof}

\begin{theorem}[Algébricité du champ des courbes]
\label{quot-theorem-curves-algebraic}
\begin{reference}
Voir \cite[proposition 3.3, page 8]{dJHS} et
\cite[appendice B par Jack Hall, théorème B.1]{Smyth}.
\end{reference}
Le champ $\Curvesstack$ (situation \ref{quot-situation-curves})
est algébrique. En fait, pour tout espace algébrique $B$, le champ
$B\text{-}\Curvesstack$ (remarque \ref{quot-remark-curves-base-change})
est algébrique.
\end{theorem}

\begin{proof}
Le cas absolu résulte de
Axiomes d'Artin, lemme \ref{artin-lemma-diagonal-representable},
et des lemmes \ref{quot-lemma-curves-diagonal},
\ref{quot-lemma-curves-RS-star},
\ref{quot-lemma-curves-limits},
\ref{quot-lemma-curves-existence} et
\ref{quot-lemma-curves-defo-thy}.
Le cas sur $B$ en résulte, compte tenu de la description de
$B\text{-}\Curvesstack$ comme $2$-produit fibré dans la remarque
\ref{quot-remark-curves-base-change} et du fait que les champs algébriques admettent
les $2$-produits fibrés, voir
Champs algébriques, lemme \ref{algebraic-lemma-2-fibre-product}.
\end{proof}

\begin{lemma}
\label{quot-lemma-curves-open-and-closed-in-spaces}
Le $1$-morphisme (\ref{quot-equation-curves-over-proper-spaces})
$$
\Curvesstack \longrightarrow \Spacesstack'_{fp, flat, proper}
$$
est représentable par des immersions ouvertes et fermées.
\end{lemma}

\begin{proof}
Puisque (\ref{quot-equation-curves-over-proper-spaces}) est un plongement pleinement
fidèle de catégories, il suffit de montrer ce qui suit : étant donné un objet
$X \to S$ de $\Spacesstack'_{fp, flat, proper}$, il existe un sous-schéma
ouvert et fermé $U \subset S$ tel qu'un morphisme $S' \to S$ se factorise par
$U$ si et seulement si le changement de base $X' \to S'$ de $X \to S$ est de
dimension relative $\leq 1$.
Cela résulte immédiatement de
Compléments sur les morphismes d'espaces, lemme
\ref{spaces-more-morphisms-lemma-dimension-fibres-proper-flat}.
\end{proof}

\begin{remark}
\label{quot-remark-alternative-approach-curves}
Considérons le $2$-produit fibré
$$
\xymatrix{
\Curvesstack \times_{\Spacesstack'_{fp, flat, proper}}
\Polarizedstack \ar[r] \ar[d] &
\Polarizedstack \ar[d] \\
\Curvesstack \ar[r] &
\Spacesstack'_{fp, flat, proper}
}
$$
Ce produit fibré paramètre les courbes polarisées, c'est-à-dire les familles de
courbes munies d'un faisceau inversible relativement ample.
Il se trouve que la flèche verticale de gauche
$$
\textit{PolarizedCurves} \longrightarrow \Curvesstack
$$
est algébrique, lisse et surjective. En effet, ce $1$-morphisme est algébrique
(en tant que changement de base de la flèche du
lemme \ref{quot-lemma-polarized-to-spaces-algebraic}), tout point appartient à son
image, et il n'y a pas d'obstruction à déformer les faisceaux inversibles sur
les courbes (voir la démonstration du lemme \ref{quot-lemma-curves-existence}).
Cela donne une autre approche de l'algébricité de $\Curvesstack$.
En effet, d'après le lemme \ref{quot-lemma-curves-open-and-closed-in-spaces},
$\textit{PolarizedCurves}$ est un sous-champ ouvert et fermé du champ
algébrique $\Polarizedstack$, et tout champ en groupoïdes qui est le but d'un
morphisme algébrique lisse provenant d'un champ algébrique est un champ
algébrique.
\end{remark}

\section{Espaces de modules de complexes sur un morphisme propre}
\label{quot-section-moduli-complexes}

\noindent
Le titre et le contenu de cette section sont tirés de
\cite{lieblich-complexes}. Soit $S$ un schéma et soit
$f : X \to B$ un morphisme propre, plat et de présentation finie d'espaces
algébriques. Nous allons démontrer qu'il existe un champ algébrique
$$
\Complexesstack_{X/B}
$$
qui paramètre les « familles » d'objets de $D^b_{\textit{Coh}}$ des fibres
dont les auto-Ext négatifs s'annulent. Plus précisément, une famille est donnée
par un objet relativement parfait de la catégorie dérivée de l'espace total ;
cette notion quelque peu technique est étudiée dans
Compléments sur les morphismes d'espaces, section
\ref{spaces-more-morphisms-section-relatively-perfect}.

\medskip\noindent
Déjà lorsque $X$ est un espace algébrique propre sur un corps $k$, nous
obtenons un champ algébrique très intéressant. En effet, il existe un
plongement
$$
\Cohstack_{X/k} \longrightarrow \Complexesstack_{X/k}
$$
car, pour tout $\mathcal{O}$-module $\mathcal{F}$ (sur tout topos annelé),
nous avons $\Ext^i_\mathcal{O}(\mathcal{F}, \mathcal{F}) = 0$ pour $i < 0$.
Bien que cela montre assurément que notre champ est non vide, la véritable
motivation pour étudier $\Complexesstack_{X/k}$ est qu'il existe souvent des
objets de la catégorie dérivée $D^b_{\textit{Coh}}(\mathcal{O}_X)$ dont les
auto-Ext négatifs s'annulent et dont les faisceaux de cohomologie sont non nuls
en plus d'un degré.
Par exemple, $X$ pourrait être équivalent au sens dérivé à un autre espace
algébrique propre $Y$ sur $k$, c'est-à-dire que nous aurions une équivalence
$k$-linéaire
$$
F : D^b_{\textit{Coh}}(\mathcal{O}_Y)
\longrightarrow
D^b_{\textit{Coh}}(\mathcal{O}_X)
$$
Il existe des cas où cela se produit sans que $F$ soit donné par un
isomorphisme entre $X$ et $Y$ ; par exemple, dans le cas d'une variété
abélienne et de sa duale. Dans cette situation, $F$ induit un isomorphisme de
champs algébriques
$$
\Complexesstack_{Y/k}
\longrightarrow
\Complexesstack_{X/k}
$$
(insérer ici une référence future) et, en particulier, le champ des faisceaux
cohérents sur $Y$ s'envoie dans le champ des complexes sur $X$. En renversant
ce point de vue, si nous comprenons assez bien la géométrie de
$\Complexesstack_{X/k}$, nous pouvons tenter de l'utiliser pour étudier tous
les $Y$ qui lui sont équivalents au sens dérivé.

\begin{lemma}
\label{quot-lemma-complexes-open-neg-exts-vanishing}
Soit $S$ un schéma.
Soit $f : X \to Y$ un morphisme d'espaces algébriques sur $S$.
Supposons que $f$ soit propre, plat et de présentation finie.
Soient $K, E \in D(\mathcal{O}_X)$. Supposons que $K$ soit pseudo-cohérent et
que $E$ soit $Y$-parfait (Compléments sur les morphismes d'espaces, définition
\ref{spaces-more-morphisms-definition-relatively-perfect}).
Pour un corps $k$ et un morphisme $y : \Spec(k) \to Y$, notons $K_y$, $E_y$
les images inverses sur la fibre $X_y$.
\begin{enumerate}
\item Il existe un ouvert $W \subset Y$ caractérisé par la propriété
$$
y \in |W|
\Leftrightarrow
\Ext^i_{\mathcal{O}_{X_y}}(K_y, E_y) = 0
\text{ pour }i < 0.
$$
\item Pour tout morphisme $V \to Y$ se factorisant par $W$, nous avons
$$
\Ext^i_{\mathcal{O}_{X_V}}(K_V, E_V) = 0
\quad\text{pour}\quad i < 0
$$
où $X_V$ est le changement de base de $X$, et $K_V$ et $E_V$ sont les images
inverses dérivées de $K$ et $E$ sur $X_V$.
\item Le foncteur $V \mapsto \Hom_{\mathcal{O}_{X_V}}(K_V, E_V)$
est un faisceau sur $(\textit{Espaces}/W)_{fppf}$ représentable par un espace
algébrique affine et de présentation finie sur $W$.
\end{enumerate}
\end{lemma}

\begin{proof}
Pour tout morphisme $V \to Y$, le complexe $K_V$ est pseudo-cohérent
(Cohomologie sur les sites, lemme
\ref{sites-cohomology-lemma-pseudo-coherent-pullback})
et $E_V$ est $V$-parfait (Compléments sur les morphismes d'espaces, lemme
\ref{spaces-more-morphisms-lemma-base-change-relatively-perfect}).
Observons aussi qu'étant donnés $y : \Spec(k) \to Y$, une extension de corps
$k'/k$ et le morphisme induit $y' : \Spec(k') \to Y$, nous avons
$$
\Ext^i_{\mathcal{O}_{X_{y'}}}(K_{y'}, E_{y'}) =
\Ext^i_{\mathcal{O}_{X_y}}(K_y, E_y) \otimes_k k'
$$
d'après Catégories dérivées de schémas, lemme
\ref{perfect-lemma-affine-morphism-and-hom-out-of-perfect}.
Ainsi, l'annulation en (1) est bien une propriété du point induit
$y \in |Y|$. Nous utiliserons ces deux observations sans plus les mentionner
dans la démonstration.

\medskip\noindent
Supposons d'abord que $Y$ soit un schéma affine. Nous pouvons alors appliquer
Compléments sur les morphismes d'espaces, lemme
\ref{spaces-more-morphisms-lemma-compute-ext-rel-perfect},
et trouver un $L \in D(\mathcal{O}_Y)$ pseudo-cohérent qui « calcule
universellement » $Rf_*R\SheafHom(K, E)$ au sens décrit dans ce lemme.
En développant les définitions, nous obtenons, pour un point $y \in Y$,
l'égalité
$$
\Ext^i_{\kappa(y)}(L \otimes_{\mathcal{O}_Y}^\mathbf{L} \kappa(y),
\kappa(y)) = \Ext^i_{\mathcal{O}_{X_y}}(K_y, E_y)
$$
Nous en concluons que
$$
H^i(L \otimes_{\mathcal{O}_Y}^\mathbf{L} \kappa(y)) = 0
\text{ pour } i > 0 \Leftrightarrow
\Ext^i_{\mathcal{O}_{X_y}}(K_y, E_y) = 0 \text{ pour }i < 0.
$$
D'après Catégories dérivées de schémas, lemme
\ref{perfect-lemma-jump-loci}, l'ensemble $W$ des $y \in Y$ où cela se produit
définit un ouvert de $Y$.
Cet ouvert $W$ satisfait alors à la condition de (1) pour tous les morphismes
provenant de spectres de corps, par l'« universalité » de $L$.

\medskip\noindent
Revenons au cas où $Y$ est un espace algébrique général.
Choisissons un recouvrement \'etale $\{V_i \to Y\}$ par des schémas affines
$V_i$. Nous voyons alors que l'image inverse du sous-ensemble $W \subset |Y|$
est le sous-ensemble correspondant $W_i \subset |V_i|$ pour
$X_{V_i}$, $K_{V_i}$, $E_{V_i}$.
D'après le paragraphe précédent, $W_i$ est ouvert, donc $W$ est ouvert.
Cela prouve (1) en général. En outre, les parties (2) et (3) sont entièrement
formulées en termes de la catégorie $\textit{Espaces}/W$ et des restrictions
$X_W$, $K_W$, $E_W$. Nous sommes donc ramenés au cas $W = Y$.

\medskip\noindent
Supposons $W = Y$. Nous affirmons que, pour tout espace algébrique $V$ sur
$Y$, le complexe $Rf_{V, *}R\SheafHom(K_V, E_V)$ a ses faisceaux de
cohomologie nuls en degrés $< 0$. Cela prouvera (2), car
$$
\Ext^i_{\mathcal{O}_{X_V}}(K_V, E_V) =
H^i(X_V, R\SheafHom(K_V, E_V)) = 
H^i(V, Rf_{V, *}R\SheafHom(K_V, E_V))
$$
d'après Cohomologie sur les sites, lemmes
\ref{sites-cohomology-lemma-section-RHom-over-U} et
\ref{sites-cohomology-lemma-Leray-unbounded}, et l'annulation des faisceaux de
cohomologie entraîne que le groupe de cohomologie $H^i$ est nul pour $i < 0$
d'après Catégories dérivées, lemme
\ref{derived-lemma-negative-vanishing}.

\medskip\noindent
Pour prouver l'affirmation, nous pouvons travailler localement pour la
topologie \'etale sur $V$. En particulier, nous pouvons supposer que $Y$ est
affine et que $W = Y$.
Soit $L \in D(\mathcal{O}_Y)$ comme dans le deuxième paragraphe de la
démonstration. Pour un espace algébrique $V$ sur $Y$, notons $L_V$ l'image
inverse dérivée de $L$ sur $V$. (Une propriété importante que nous utiliserons
est que $L$ « fonctionne » pour tous les espaces algébriques $V$ sur $Y$, et
pas seulement pour les $V$ affines.)
Comme $W = Y$, nous avons $H^i(L) = 0$ pour $i > 0$
(utiliser Compléments d'algèbre, lemme
\ref{more-algebra-lemma-lift-pseudo-coherent-from-residue-field},
pour passer des fibres aux germes). Donc $H^i(L_V) = 0$ pour $i > 0$.
La propriété qui définit $L$ est que
$$
Rf_{V, *}R\SheafHom(K_V, E_V) = R\SheafHom(L_V, \mathcal{O}_V)
$$
Puisque $L_V$ est concentré en degrés $\leq 0$, nous concluons que
$R\SheafHom(L_V, \mathcal{O}_V)$ est concentré en degrés $\geq 0$, ce qui
prouve l'affirmation. Cela achève la démonstration de (2).

\medskip\noindent
Supposons $W = Y$, sans faire d'hypothèse sur l'espace algébrique $Y$.
Comme nous avons (2), il résulte de
Espaces simpliciaux, lemme
\ref{spaces-simplicial-lemma-fppf-neg-ext-zero-hom}, que le foncteur $F$ donné
par $F(V) = \Hom_{\mathcal{O}_{X_V}}(K_V, E_V)$ est un
faisceau\footnote{Pour vérifier la propriété de faisceau relativement à un
recouvrement $\{V_i \to V\}_{i \in I}$, considérons d'abord
l'hyperrecouvrement fppf de {\v C}ech $a : V_\bullet \to V$, où
$V_n = \coprod_{i_0 \ldots i_n} V_{i_0} \times_V \ldots \times_V V_{i_n}$,
puis posons $U_\bullet = V_\bullet \times_{a, V} X_V$. Alors
$U_\bullet \to X_V$ est un hyperrecouvrement fppf auquel nous pouvons
appliquer Espaces simpliciaux, lemme
\ref{spaces-simplicial-lemma-fppf-neg-ext-zero-hom}.}
sur $(\textit{Espaces}/Y)_{fppf}$. Ainsi, pour prouver que $F$ est un espace
algébrique et que $F \to Y$ est affine et de présentation finie, nous pouvons
travailler localement pour la topologie \'etale sur $Y$ ; voir
Amorçage, lemme \ref{bootstrap-lemma-locally-algebraic-space-finite-type}, et
Morphismes d'espaces, lemmes \ref{spaces-morphisms-lemma-affine-local} et
\ref{spaces-morphisms-lemma-finite-presentation-local}.
Nous concluons qu'il suffit de prouver que $F$ est un espace algébrique affine
de présentation finie sur $Y$ lorsque $Y$ est un schéma affine. Dans ce cas,
nous revenons à notre complexe pseudo-cohérent $L \in D(\mathcal{O}_Y)$.
Comme $H^i(L) = 0$ pour $i > 0$, nous pouvons représenter $L$ par un complexe
de la forme
$$
\ldots \to \mathcal{O}_Y^{\oplus m_1} \to \mathcal{O}_Y^{\oplus m_0}
\to 0 \to \ldots
$$
dont le dernier terme est en degré $0$, voir Compléments d'algèbre, lemme
\ref{more-algebra-lemma-pseudo-coherent}. En combinant les deux formules
affichées plus haut dans la démonstration, nous trouvons que
$$
F(V) =
\Ker(
\Hom_V(\mathcal{O}_V^{\oplus m_0}, \mathcal{O}_V)
\to 
\Hom_V(\mathcal{O}_V^{\oplus m_1}, \mathcal{O}_V)
)
$$
Autrement dit, il existe un diagramme cartésien
$$
\xymatrix{
F \ar[d] \ar[r] & Y \ar[d]^0 \\
\mathbf{A}_Y^{m_0} \ar[r] & \mathbf{A}_Y^{m_1}
}
$$
qui prouve ce que nous voulions.
\end{proof}

\begin{lemma}
\label{quot-lemma-complexes}
Soit $S$ un schéma. Soit $f : X \to Y$ un morphisme d'espaces algébriques
sur $S$. Supposons que $f$ soit propre, plat et de présentation finie.
Soit $E \in D(\mathcal{O}_X)$.
Supposons que
\begin{enumerate}
\item $E$ soit $S$-parfait (Compléments sur les morphismes d'espaces,
définition \ref{spaces-more-morphisms-definition-relatively-perfect}), et
\item pour tout point $s \in S$, nous ayons
$$
\Ext^i_{\mathcal{O}_{X_s}}(E_s, E_s) = 0
\quad\text{pour}\quad i < 0
$$
où $E_s$ est l'image inverse sur la fibre $X_s$.
\end{enumerate}
Alors
\begin{enumerate}
\item[(a)] (1) et (2) sont préservées par tout changement de base $V \to Y$,
\item[(b)] $\Ext^i_{\mathcal{O}_{X_V}}(E_V, E_V) = 0$ pour $i < 0$ et tout
$V$ sur $Y$,
\item[(c)] $V \mapsto \Hom_{\mathcal{O}_{X_V}}(E_V, E_V)$ est représentable
par un espace algébrique affine et de présentation finie sur $Y$.
\end{enumerate}
Ici, $X_V$ est le changement de base de $X$ et $E_V$ l'image inverse dérivée
de $E$ sur $X_V$.
\end{lemma}

\begin{proof}
Conséquence immédiate du lemme
\ref{quot-lemma-complexes-open-neg-exts-vanishing}.
\end{proof}

\begin{situation}
\label{quot-situation-complexes}
Soit $S$ un schéma. Soit $f : X \to B$ un morphisme d'espaces algébriques
sur $S$. Supposons que $f$ soit propre, plat et de présentation finie.
Nous notons $\Complexesstack_{X/B}$ la catégorie dont les objets sont les
triplets $(T, g, E)$ où
\begin{enumerate}
\item $T$ est un schéma sur $S$,
\item $g : T \to B$ est un morphisme sur $S$, et, en posant
$X_T = T \times_{g, B} X$,
\item $E$ est un objet de $D(\mathcal{O}_{X_T})$ satisfaisant aux conditions
(1) et (2) du lemme \ref{quot-lemma-complexes}.
\end{enumerate}
Un morphisme $(T, g, E) \to (T', g', E')$ est donné par un couple
$(h, \varphi)$ où
\begin{enumerate}
\item $h : T \to T'$ est un morphisme de schémas sur $B$
(c'est-à-dire $g' \circ h = g$), et
\item $\varphi : L(h')^*E' \to E$ est un isomorphisme de
$D(\mathcal{O}_{X_T})$, où $h' : X_T \to X_{T'}$ est le changement de base de
$h$.
\end{enumerate}
\end{situation}

\noindent
Ainsi, $\Complexesstack_{X/B}$ est une catégorie et la règle
$$
p : \Complexesstack_{X/B} \longrightarrow (\Sch/S)_{fppf},
\quad
(T, g, E) \longmapsto T
$$
est un foncteur. Pour un schéma $T$ sur $S$, nous notons
$\Complexesstack_{X/B, T}$ la catégorie fibre de $p$ au-dessus de $T$.
Ces catégories fibres sont des groupoïdes.

\begin{lemma}
\label{quot-lemma-complexes-fibred-in-groupoids}
Dans la situation \ref{quot-situation-complexes}, le foncteur
$p : \Complexesstack_{X/B} \longrightarrow (\Sch/S)_{fppf}$
est fibré en groupoïdes.
\end{lemma}

\begin{proof}
Nous montrons que $p$ est fibré en groupoïdes en vérifiant les conditions (1)
et (2) de Catégories, définition
\ref{categories-definition-fibred-groupoids}.
Étant donnés un objet $(T', g', E')$ de $\Complexesstack_{X/B}$ et un morphisme
$h : T \to T'$ de schémas sur $S$, nous pouvons poser $g = h \circ g'$ et
$E = L(h')^*E'$, où $h' : X_T \to X_{T'}$ est le changement de base de $h$.
Il est alors clair que nous obtenons un morphisme
$(T, g, E) \to (T', g', E')$ de $\Complexesstack_{X/B}$ au-dessus de $h$.
Cela prouve (1).
Pour (2), supposons donnés des morphismes
$$
(h_1, \varphi_1) : (T_1, g_1, E_1) \to (T, g, E)
\quad\text{et}\quad
(h_2, \varphi_2) : (T_2, g_2, E_2) \to (T, g, E)
$$
de $\Complexesstack_{X/B}$ et un morphisme $h : T_1 \to T_2$ tel que
$h_2 \circ h = h_1$. Nous pouvons alors prendre pour $\varphi$ la composée
$$
L(h')^*E_2
\xrightarrow{L(h')^*\varphi_2^{-1}}
L(h')^*L(h_2)^*E = L(h_1)^*E
\xrightarrow{\varphi_1}
E_1
$$
et obtenir ainsi le morphisme
$(h, \varphi) : (T_1, g_1, E_1) \to (T_2, g_2, E_2)$
qui atteste la condition (2).
\end{proof}

\begin{lemma}
\label{quot-lemma-complexes-diagonal}
Dans la situation \ref{quot-situation-complexes}. Notons
$\mathcal{X} = \Complexesstack_{X/B}$. Alors
$\Delta : \mathcal{X} \to \mathcal{X} \times \mathcal{X}$ est représentable
par des espaces algébriques.
\end{lemma}

\begin{proof}
Considérons deux objets $x = (T, g, E)$ et $y = (T, g', E')$ de
$\mathcal{X}$ au-dessus d'un schéma $T$. Nous devons montrer que
$\mathit{Isom}_\mathcal{X}(x, y)$ est un espace algébrique sur $T$, voir
Champs algébriques, lemme \ref{algebraic-lemma-representable-diagonal}.
Si, pour $h : T' \to T$, les restrictions $x|_{T'}$ et $y|_{T'}$ sont
isomorphes dans la catégorie fibre $\mathcal{X}_{T'}$, alors
$g \circ h = g' \circ h$.
Il existe donc une transformation de préfaisceaux
$$
\mathit{Isom}_\mathcal{X}(x, y) \longrightarrow \text{Égalisateur}(g, g')
$$
Puisque la diagonale de $B$ est représentable (par des schémas), cet égalisateur
est un schéma. Nous pouvons donc remplacer $T$ par cet égalisateur et $E$, $E'$
par leurs images inverses. Nous pouvons ainsi supposer $g = g'$.

\medskip\noindent
Supposons $g = g'$. Après avoir remplacé $B$ par $T$ et $X$ par $X_T$, nous
arrivons au problème suivant. Étant donnés $E, E' \in D(\mathcal{O}_X)$
satisfaisant aux conditions (1), (2) du lemme \ref{quot-lemma-complexes}, nous devons
montrer que $\mathit{Isom}(E, E')$ est un espace algébrique.
Ici, $\mathit{Isom}(E, E')$ est le foncteur
$$
(\Sch/B)^{opp} \to \textit{Ensembles},\quad
T \mapsto \{\varphi : E_T \to E'_T
\text{ isomorphisme dans }D(\mathcal{O}_{X_T})\}
$$
où $E_T$ et $E'_T$ sont les images inverses dérivées de $E$ et $E'$ sur
$X_T$. Soient maintenant $W \subset B$, resp.\ $W' \subset B$, les sous-espaces
ouverts de $B$ associés à $E, E'$, resp.\ à $E', E$, par le
lemme \ref{quot-lemma-complexes-open-neg-exts-vanishing}.
Il est clair que, s'il existe un isomorphisme $E_T \to E'_T$ comme dans la
définition de $\mathit{Isom}(E, E')$, alors $T \to B$ se factorise à la fois
par $W$ et par $W'$ (parce que nous avons la condition (1) pour $E$ et $E'$, et
que nous aurons évidemment $E_t \cong E'_t$, donc aucune application non nulle
$E_t[i] \to E_t$ ou $E'_t[i] \to E_t$ sur la fibre $X_t$ pour $i > 0$).
Nous pouvons donc remplacer $B$ par l'ouvert $W \cap W'$.
Dans ce cas, le foncteur $H = \SheafHom(E, E')$
$$
(\Sch/B)^{opp} \to \textit{Ensembles},\quad T \mapsto
\Hom_{\mathcal{O}_{X_T}}(E_T, E'_T)
$$
est un espace algébrique affine et de présentation finie sur $B$ d'après le
lemme \ref{quot-lemma-complexes-open-neg-exts-vanishing}.
Il en est de même de
$H' = \SheafHom(E', E)$,
$I = \SheafHom(E, E)$ et
$I' = \SheafHom(E', E')$.
Nous pouvons donc reprendre l'argument de la démonstration de la
proposition \ref{quot-proposition-isom} pour voir que
$$
\mathit{Isom}(E, E') = (H' \times_B H) \times_{c, I \times_B I', \sigma} B
$$
pour certains morphismes $c$ et $\sigma$. Ainsi,
$\mathit{Isom}(E, E')$ est un espace algébrique.
\end{proof}

\begin{lemma}
\label{quot-lemma-complexes-stack}
Dans la situation \ref{quot-situation-complexes}, le foncteur
$p : \Complexesstack_{X/B} \longrightarrow (\Sch/S)_{fppf}$
est un champ en groupoïdes.
\end{lemma}

\begin{proof}
Pour prouver que $\Complexesstack_{X/B}$ est un champ en groupoïdes, nous
devons montrer que les préfaisceaux $\mathit{Isom}$ sont des faisceaux et que
les données de descente sont effectives. L'assertion sur $\mathit{Isom}$
résulte du lemme \ref{quot-lemma-complexes-diagonal}, voir
Champs algébriques, lemme \ref{algebraic-lemma-representable-diagonal}.
Démontrons l'assertion sur les données de descente.

\medskip\noindent
Supposons que $\{a_i : T_i \to T\}$ soit un recouvrement fppf de schémas sur
$S$. Soit $(\xi_i, \varphi_{ij})$ une donnée de descente pour
$\{T_i \to T\}$ à valeurs dans $\Complexesstack_{X/B}$.
Pour chaque $i$, nous pouvons écrire $\xi_i = (T_i, g_i, E_i)$.
Notons $\text{pr}_0 : T_i \times_T T_j \to T_i$ et
$\text{pr}_1 : T_i \times_T T_j \to T_j$ les projections.
La condition
$\xi_i|_{T_i \times_T T_j} \cong \xi_j|_{T_i \times_T T_j}$ implique en
particulier que $g_i \circ \text{pr}_0 = g_j \circ \text{pr}_1$.
Il existe donc un unique morphisme $g : T \to B$ tel que
$g_i = g \circ a_i$, voir
Descente sur les espaces, lemme
\ref{spaces-descent-lemma-fpqc-universal-effective-epimorphisms}.
Notons $X_T = T \times_{g, B} X$. Posons
$X_i = X_{T_i} = T_i \times_{g_i, B} X = T_i \times_{a_i, T} X_T$ et
$$
X_{ij} = X_{T_i} \times_{X_T} X_{T_j} = X_i \times_{X_T} X_j
$$
avec les projections $\text{pr}_i$ et $\text{pr}_j$ vers $X_i$ et $X_j$.
Observons que l'image inverse de $(T_i, g_i, E_i)$ par
$\text{pr}_0 : T_i \times_T T_j \to T_i$ est donnée par
$(T_i \times_T T_j, g_i \circ \text{pr}_0, L\text{pr}_i^*E_i)$.
Une donnée de descente pour $\{T_i \to T\}$ dans $\Complexesstack_{X/B}$ est
donc donnée par les objets $(T_i, g \circ a_i, E_i)$ et, pour chaque couple
$i, j$, par un isomorphisme dans $D\mathcal{O}_{X_{ij}})$
$$
\varphi_{ij} :
L\text{pr}_i^*E_i \longrightarrow L\text{pr}_j^*E_j
$$
satisfaisant à la condition de cocycle sur l'image inverse de $X$ sur
$T_i \times_T T_j \times_T T_k$.
En utilisant l'annulation des Ext négatifs fournie par (b) du
lemme \ref{quot-lemma-complexes}, nous pouvons appliquer
Espaces simpliciaux, lemme
\ref{spaces-simplicial-lemma-fppf-glue-neg-ext-zero}, pour obtenir la
descente\footnote{Pour le vérifier, considérons d'abord l'hyperrecouvrement
fppf de {\v C}ech $a : T_\bullet \to T$, où
$T_n = \coprod_{i_0 \ldots i_n} T_{i_0} \times_T \ldots \times_T T_{i_n}$,
puis posons $U_\bullet = T_\bullet \times_{a, T} X_T$. Alors
$U_\bullet \to X_T$ est un hyperrecouvrement fppf auquel nous pouvons
appliquer Espaces simpliciaux, lemme
\ref{spaces-simplicial-lemma-fppf-glue-neg-ext-zero}.}
pour ces complexes. Autrement dit, nous trouvons qu'il existe un objet $E$ de
$D_\QCoh(\mathcal{O}_{X_T})$ dont la restriction à $X_{T_i}$ est $E_i$,
compatiblement aux $\varphi_{ij}$. Rappelons qu'être $T$-parfait signifie être
pseudo-cohérent et avoir une dimension de Tor localement finie sur
$f^{-1}\mathcal{O}_T$.
Ainsi, $E$ est $T$-parfait par application de
Compléments sur les morphismes d'espaces, lemmes
\ref{spaces-more-morphisms-lemma-pseudo-coherent-descends-fpqc} et
\ref{spaces-more-morphisms-lemma-tor-amplitude-descends-fppf}.
Enfin, nous devons vérifier pour $E$ la condition (2) du
lemme \ref{quot-lemma-complexes}.
Cela résulte immédiatement de la description de l'ouvert $W$ dans le
lemme \ref{quot-lemma-complexes-open-neg-exts-vanishing} et du fait que (2) vaut
pour $E_i$ sur $X_{T_i}/T_i$.
\end{proof}

\begin{remark}
\label{quot-remark-complexes-base-change}
Dans la situation \ref{quot-situation-complexes}, la règle
$(T, g, E) \mapsto (T, g)$ définit un $1$-morphisme
$$
\Complexesstack_{X/B} \longrightarrow \mathcal{S}_B
$$
de champs en groupoïdes
(voir le lemme \ref{quot-lemma-complexes-stack},
Champs algébriques, section \ref{algebraic-section-split}, et
Exemples de champs, section
\ref{examples-stacks-section-stack-associated-to-sheaf}).
Soit $B' \to B$ un morphisme d'espaces algébriques sur $S$.
Soit $\mathcal{S}_{B'} \to \mathcal{S}_B$ le $1$-morphisme associé de champs
fibrés en ensembles. Posons $X' = X \times_B B'$.
Nous obtenons un champ en groupoïdes
$\Complexesstack_{X'/B'} \to (\Sch/S)_{fppf}$
associé au changement de base $f' : X' \to B'$. Dans cette situation, le
diagramme
$$
\vcenter{
\xymatrix{
\Complexesstack_{X'/B'} \ar[r] \ar[d] &
\Complexesstack_{X/B} \ar[d] \\
\mathcal{S}_{B'} \ar[r] & \mathcal{S}_B
}
}
\quad
\begin{matrix}
\text{ou dans} \\
\text{une autre} \\
\text{notation}
\end{matrix}
\quad
\vcenter{
\xymatrix{
\Complexesstack_{X'/B'} \ar[r] \ar[d] &
\Complexesstack_{X/B} \ar[d] \\
\Sch/B' \ar[r] & \Sch/B
}
}
$$
est un carré $2$-cartésien. Cette remarque triviale sera parfois utile
pour changer l'espace algébrique de base.
\end{remark}

\begin{lemma}
\label{quot-lemma-complexes-limits}
Dans la situation \ref{quot-situation-complexes}, supposons que $B \to S$ soit
localement de présentation finie. Alors
$p : \Complexesstack_{X/B} \to (\Sch/S)_{fppf}$ commute aux limites
(Axiomes d'Artin, définition \ref{artin-definition-limit-preserving}).
\end{lemma}

\begin{proof}
Notons $B(T)$ la catégorie discrète dont les objets sont les $S$-morphismes
$T \to B$. Soit $T = \lim T_i$ une limite filtrante de schémas affines sur
$S$. En associant à un objet $(T, h, E)$ de $\Complexesstack_{X/B, T}$ l'objet
$h$ de $B(T)$, nous obtenons un diagramme commutatif de catégories fibres
$$
\xymatrix{
\colim \Complexesstack_{X/B, T_i} \ar[r] \ar[d] &
\Complexesstack_{X/B, T} \ar[d] \\
\colim B(T_i) \ar[r] & B(T)
}
$$
Nous devons montrer que la flèche horizontale supérieure est une équivalence.
Puisque nous avons supposé que $B$ est localement de présentation finie sur
$S$, il résulte de
Limites d'espaces, remarque \ref{spaces-limits-remark-limit-preserving},
que la flèche horizontale inférieure est une équivalence. Cela signifie que
nous pouvons supposer que $T = \lim T_i$ est une limite filtrante de schémas
affines sur $B$. Notons $g_i : T_i \to B$ et $g : T \to B$ les morphismes
correspondants. Posons $X_i = T_i \times_{g_i, B} X$ et
$X_T = T \times_{g, B} X$.
Observons que $X_T = \colim X_i$.
D'après Compléments sur les morphismes d'espaces, lemme
\ref{spaces-more-morphisms-lemma-descend-relatively-perfect},
la catégorie des objets $T$-parfaits de $D(\mathcal{O}_{X_T})$ est la limite
inductive des catégories des objets $T_i$-parfaits de
$D(\mathcal{O}_{X_{T_i}})$.
Il nous reste donc seulement à démontrer qu'étant donné un objet $T_i$-parfait
$E_i$ de $D(\mathcal{O}_{X_{T_i}})$ tel que l'image inverse dérivée $E$ de
$E_i$ sur $X_T$ satisfasse à la condition (2) du lemme
\ref{quot-lemma-complexes}, alors, quitte à augmenter $i$, $E_i$ satisfait à la
condition (2) du lemme \ref{quot-lemma-complexes}.
Soit $W \subset |T_i|$ l'ouvert construit dans le
lemme \ref{quot-lemma-complexes-open-neg-exts-vanishing} pour $E_i$ et $E_i$.
Par l'hypothèse sur $E$, nous trouvons que $T \to T_i$ se factorise par $T$.
Il existe donc un $i' \geq i$ tel que $T_{i'} \to T_i$ se factorise par $W$,
voir Limites, lemme
\ref{limits-lemma-limit-contained-in-constructible}
Alors $i'$ convient par construction de $W$.
\end{proof}

\begin{lemma}
\label{quot-lemma-complexes-RS-star}
Dans la situation \ref{quot-situation-complexes}. Soit
$$
\xymatrix{
Z \ar[r] \ar[d] & Z' \ar[d] \\
Y \ar[r] & Y'
}
$$
une somme amalgamée dans la catégorie des schémas sur $S$, où
$Z \to Z'$ est un épaississement d'ordre fini et $Z \to Y$ est affine, voir
Compléments sur les morphismes, lemme
\ref{more-morphisms-lemma-pushout-along-thickening}.
Alors le foncteur sur les catégories fibres
$$
\Complexesstack_{X/B, Y'}
\longrightarrow
\Complexesstack_{X/B, Y}
\times_{\Complexesstack_{X/B, Z}}
\Complexesstack_{X/B, Z'}
$$
est une équivalence.
\end{lemma}

\begin{proof}
Observons que l'application correspondante
$$
B(Y') \longrightarrow B(Y) \times_{B(Z)} B(Z')
$$
est une bijection, voir Sommes amalgamées d'espaces, lemme
\ref{spaces-pushouts-lemma-pushout-along-thickening-schemes}.
Ainsi, en utilisant le diagramme commutatif
$$
\xymatrix{
\Complexesstack_{X/B, Y'} \ar[r] \ar[d] &
\Complexesstack_{X/B, Y}
\times_{\Complexesstack_{X/B, Z}}
\Complexesstack_{X/B, Z'}
\ar[d] \\
B(Y') \ar[r] & B(Y) \times_{B(Z)} B(Z')
}
$$
nous voyons que nous pouvons supposer que $Y'$ est un schéma sur $B'$.
D'après la remarque \ref{quot-remark-complexes-base-change}, nous pouvons remplacer
$B$ par $Y'$ et $X$ par $X \times_B Y'$.
Nous pouvons donc supposer $B = Y'$.

\medskip\noindent
Supposons $B = Y'$. Nous prouvons d'abord la pleine fidélité de notre foncteur.
Pour cela, soient $\xi_1, \xi_2$ deux objets de $\Complexesstack_{X/B}$
au-dessus de $Y'$. Nous devons alors montrer que
$$
\mathit{Isom}(\xi_1, \xi_2)(Y') \longrightarrow
\mathit{Isom}(\xi_1, \xi_2)(Y)
\times_{\mathit{Isom}(\xi_1, \xi_2)(Z)}
\mathit{Isom}(\xi_1, \xi_2)(Z')
$$
est bijective. Or nous savons déjà que $\mathit{Isom}(\xi_1, \xi_2)$ est un
espace algébrique sur $B = Y'$. Cette bijectivité résulte donc de
Axiomes d'Artin, lemme \ref{artin-lemma-pushout} (ou du lemme déjà mentionné de
Sommes amalgamées d'espaces,
\ref{spaces-pushouts-lemma-pushout-along-thickening-schemes}).

\medskip\noindent
Surjectivité essentielle. Soit $(E_Y, E_{Z'}, \alpha)$ un triplet, où
$E_Y \in D(\mathcal{O}_Y)$ et $E_{Z'} \in D(\mathcal{O}_{X_{Z'}})$ sont des
objets tels que $(Y, Y \to B, E_Y)$ soit un objet de
$\Complexesstack_{X/B}$ au-dessus de $Y$, que
$(Z', Z' \to B, E_{Z'})$ soit un objet de $\Complexesstack_{X/B}$ au-dessus de
$Z'$, et que
$\alpha : L(X_Z \to X_Y)^*E_Y \to L(X_Z \to X_{Z'})^*E_{Z'}$
soit un isomorphisme dans $D(\mathcal{O}_{Z'})$.
C'est-à-dire que
$$
((Y, Y \to B, E_Y), (Z', Z' \to B, E_{Z'}), \alpha)
$$
est un objet du but de la flèche de notre lemme.
Observons que le diagramme
$$
\xymatrix{
X_Z \ar[r] \ar[d] & X_{Z'} \ar[d] \\
X_Y \ar[r] & X_{Y'}
}
$$
est une somme amalgamée, avec $X_Z \to X_Y$ affine et
$X_Z \to X_{Z'}$ un épaississement
(voir Sommes amalgamées d'espaces, lemme
\ref{spaces-pushouts-lemma-equivalence-categories-spaces-pushout-flat}).
Par conséquent, d'après Sommes amalgamées d'espaces, lemme
\ref{spaces-pushouts-lemma-pushout-along-thickening-derived},
nous trouvons un objet $E_{Y'} \in D(\mathcal{O}_{X_{Y'}})$ ainsi que des
isomorphismes
$L(X_Y \to X_{Y'})^*E_{Y'} \to E_Y$ et
$L(X_{Z'} \to X_{Y'})^*E_{Y'} \to E_Z$
compatibles à $\alpha$. Il est clair que, si nous montrons que $E_{Y'}$ est
$Y'$-parfait, alors nous aurons terminé, car la propriété (2) du
lemme \ref{quot-lemma-complexes} est une propriété portant sur les points (et $Y$
et $Y'$ ont les mêmes points).
Cela résulte de Compléments sur les morphismes d'espaces, lemme
\ref{spaces-more-morphisms-lemma-thickening-relatively-perfect}.
\end{proof}

\begin{lemma}
\label{quot-lemma-complexes-tangent-space}
Dans la situation \ref{quot-situation-complexes}, supposons que $S$ soit un schéma
localement noethérien et que $B \to S$ soit localement de présentation finie.
Soit $k$ un corps de type fini sur $S$ et soit
$x_0 = (\Spec(k), g_0, E_0)$ un objet de
$\mathcal{X} = \Complexesstack_{X/B}$ au-dessus de $k$.
Alors les espaces $T\mathcal{F}_{\mathcal{X}, k, x_0}$ et
$\text{Inf}(\mathcal{F}_{\mathcal{X}, k, x_0})$
(Axiomes d'Artin, section \ref{artin-section-tangent-spaces})
sont de dimension finie.
\end{lemma}

\begin{proof}
Observons que, d'après le lemme \ref{quot-lemma-complexes-RS-star}, notre champ en
groupoïdes $\mathcal{X}$ satisfait à la propriété (RS*) définie dans
Axiomes d'Artin, section \ref{artin-section-RS-star}.
En particulier, $\mathcal{X}$ satisfait à (RS).
Par conséquent, toutes les catégories de prédéformations associées sont des
catégories de déformations
(Axiomes d'Artin, lemme \ref{artin-lemma-deformation-category}),
et l'énoncé a un sens.

\medskip\noindent
Dans ce paragraphe, nous montrons que nous pouvons nous ramener au cas
$B = \Spec(k)$. Posons $X_0 = \Spec(k) \times_{g_0, B} X$ et notons
$\mathcal{X}_0 = \Complexesstack_{X_0/k}$. Dans la remarque
\ref{quot-remark-complexes-base-change}, nous avons vu que $\mathcal{X}_0$ est le
$2$-produit fibré de $\mathcal{X}$ et $\Spec(k)$ sur $B$ comme catégories
fibrées en groupoïdes sur $(\Sch/S)_{fppf}$. Ainsi, d'après
Axiomes d'Artin, lemme \ref{artin-lemma-fibre-product-tangent-spaces},
nous nous ramenons à démontrer que $B$, $\Spec(k)$ et $\mathcal{X}_0$ ont des
espaces tangents et des espaces d'automorphismes infinitésimaux de dimension
finie. Les espaces tangents de $B$ et de $\Spec(k)$ sont de dimension finie
d'après Axiomes d'Artin, lemme \ref{artin-lemma-finite-dimension}, et leurs
$\text{Inf}$ sont bien sûr nuls.
Il suffit donc de traiter $\mathcal{X}_0$.

\medskip\noindent
Soit $k[\epsilon]$ l'algèbre des nombres duaux sur $k$.
Soit $\Spec(k[\epsilon]) \to B$ la composée de $g_0 : \Spec(k) \to B$ et du
morphisme $\Spec(k[\epsilon]) \to \Spec(k)$ provenant de l'inclusion
$k \to k[\epsilon]$.
Posons $X_0 = \Spec(k) \times_B X$ et
$X_\epsilon = \Spec(k[\epsilon]) \times_B X$.
Observons que $X_\epsilon$ est un épaississement du premier ordre de $X_0$,
plat sur l'épaississement du premier ordre
$\Spec(k) \to \Spec(k[\epsilon])$.
Observons que $X_0$ et $X_\epsilon$ donnent lieu à des petits topos \'etales
canoniquement équivalents, voir
Compléments sur les morphismes d'espaces, section
\ref{spaces-more-morphisms-section-thickenings}.
D'après Compléments sur les morphismes d'espaces, lemme
\ref{spaces-more-morphisms-lemma-thickening-relatively-perfect},
$T\mathcal{F}_{\mathcal{X}_0, k, x_0}$ est l'ensemble des classes
d'isomorphisme de relèvements de $E_0$ à $X_\epsilon$ au sens de
Théorie des déformations, lemme
\ref{defos-lemma-first-order-defos-complex}.
Nous concluons que
$$
T\mathcal{F}_{\mathcal{X}_0, k, x_0} =
\Ext^1_{\mathcal{O}_{X_0}}(E_0, E_0)
$$
Nous avons utilisé ici l'identification $\epsilon k[\epsilon] \cong k$ de
$k[\epsilon]$-modules. En utilisant encore une fois
Théorie des déformations, lemme
\ref{defos-lemma-first-order-defos-complex}, nous voyons qu'il existe une
surjection
$$
\text{Inf}(\mathcal{F}_{\mathcal{X}, k, x_0})
\leftarrow
\Ext^0_{\mathcal{O}_{X_0}}(E_0, E_0)
$$
d'espaces vectoriels sur $k$. Comme $E_0$ est pseudo-cohérent, il appartient à
$D^-_{\textit{Coh}}(\mathcal{O}_{X_0})$ d'après
Catégories dérivées d'espaces, lemme
\ref{spaces-perfect-lemma-identify-pseudo-coherent-noetherian}.
Puisque $E_0$ est localement de dimension de Tor finie et que $X_0$ est
quasi-compact, nous avons
$E_0 \in D^b_{\textit{Coh}}(\mathcal{O}_{X_0})$.
Les $\Ext$ ci-dessus sont donc des espaces vectoriels de dimension finie sur
$k$ d'après Catégories dérivées d'espaces, lemme
\ref{spaces-perfect-lemma-ext-finite}.
\end{proof}

\begin{lemma}
\label{quot-lemma-complexes-strong-effectiveness}
Dans la situation \ref{quot-situation-complexes}, supposons que $B = S$ soit
localement noethérien. Alors l'effectivité formelle forte au sens de
Axiomes d'Artin, remarque \ref{artin-remark-strong-effectiveness},
vaut pour $p : \Complexesstack_{X/S} \to (\Sch/S)_{fppf}$.
\end{lemma}

\begin{proof}
Soit $(R_n)$ un système projectif de $S$-algèbres à applications de transition
surjectives dont les noyaux sont localement nilpotents. Posons
$R = \lim R_n$. Soit $(\xi_n)$ un système d'objets de
$\Complexesstack_{X/B}$ au-dessus des $(\Spec(R_n))$. Nous devons montrer que
$(\xi_n)$ est effectif, c'est-à-dire qu'il existe un objet $\xi$ de
$\Complexesstack_{X/B}$ au-dessus de $\Spec(R)$.

\medskip\noindent
Écrivons $X_R = \Spec(R) \times_S X$ et
$X_n = \Spec(R_n) \times_S X$.
Bien entendu, $X_n$ est le changement de base de $X_R$ par $R \to R_n$.
Puisque $S = B$, nous voyons que $\xi_n$ correspond simplement à un objet
$R_n$-parfait $E_n \in D(\mathcal{O}_{X_n})$ satisfaisant à la condition (2)
du lemme \ref{quot-lemma-complexes}. En particulier, $E_n$ est pseudo-cohérent.
Les isomorphismes
$\xi_{n + 1}|_{\Spec(R_n)} \cong \xi_n$ correspondent à des isomorphismes
$L(X_n \to X_{n + 1})^*E_{n + 1} \to E_n$.
Par conséquent, d'après Platitude sur les espaces, théorème
\ref{spaces-flat-theorem-existence-derived}, nous trouvons un objet
pseudo-cohérent $E$ de $D(\mathcal{O}_{X_R})$ dont $E_n$ est l'image inverse
dérivée de $E$ pour tout $n$, compatiblement aux isomorphismes de transition.

\medskip\noindent
Observons que $(R, \Ker(R \to R_1))$ est une paire hensélienne, voir
Compléments d'algèbre, lemme \ref{more-algebra-lemma-limit-henselian}.
En particulier, $\Ker(R \to R_1)$ est contenu dans le radical de Jacobson de
$R$. Nous pouvons alors appliquer
Compléments sur les morphismes d'espaces, lemme
\ref{spaces-more-morphisms-lemma-henselian-relatively-perfect},
pour voir que $E$ est $R$-parfait.

\medskip\noindent
Enfin, nous devons vérifier la condition (2) du lemme \ref{quot-lemma-complexes}.
D'après le lemme \ref{quot-lemma-complexes-open-neg-exts-vanishing}, l'ensemble des
points $t$ de $\Spec(R)$ où les auto-Ext négatifs de $E_t$ s'annulent est un
ouvert. Comme cette condition est vraie dans $V(\Ker(R \to R_1))$ et que
$\Ker(R \to R_1)$ est contenu dans le radical de Jacobson de $R$, nous
concluons qu'elle vaut en tout point.
\end{proof}

\begin{theorem}[Algébricité des espaces de modules de complexes sur un morphisme propre]
\label{quot-theorem-complexes-algebraic}
\begin{reference}
\cite{lieblich-complexes}
\end{reference}
Soit $S$ un schéma. Soit $f : X \to B$ un morphisme d'espaces algébriques
sur $S$. Supposons que $f$ soit propre, plat et de présentation finie.
Alors $\Complexesstack_{X/B}$ est un champ algébrique sur $S$.
\end{theorem}

\begin{proof}
Posons $\mathcal{X} = \Complexesstack_{X/B}$. Nous avons vu que $\mathcal{X}$
est un champ en groupoïdes sur $(\Sch/S)_{fppf}$ dont la diagonale est
représentable par des espaces algébriques
(lemmes \ref{quot-lemma-complexes-stack} et \ref{quot-lemma-complexes-diagonal}).
Il suffit donc de trouver un schéma $W$ et un morphisme surjectif et lisse
$W \to \mathcal{X}$.

\medskip\noindent
Soit $B'$ un schéma et soit $B' \to B$ un morphisme surjectif \'etale.
Posons $X' = B' \times_B X$ et notons $f' : X' \to B'$ la projection.
Alors $\mathcal{X}' = \Complexesstack_{X'/B'}$ est égal au $2$-produit fibré
de $\mathcal{X}$ et de la catégorie fibrée en ensembles associée à $B'$ sur la
catégorie fibrée en ensembles associée à $B$
(remarque \ref{quot-remark-complexes-base-change}). D'après le contenu de
Champs algébriques, section
\ref{algebraic-section-representable-properties}, le morphisme
$\mathcal{X}' \to \mathcal{X}$ est surjectif et \'etale.
Il suffit donc de prouver le résultat pour $\mathcal{X}'$.
Autrement dit, nous pouvons supposer que $B$ est un schéma.

\medskip\noindent
Supposons que $B$ soit un schéma. Dans ce cas, nous pouvons remplacer $S$ par
$B$, voir Champs algébriques, section
\ref{algebraic-section-change-base-scheme}.
Nous pouvons donc supposer $S = B$.

\medskip\noindent
Supposons $S = B$. Choisissons un recouvrement ouvert affine
$S = \bigcup U_i$. Notons $\mathcal{X}_i$ la restriction de $\mathcal{X}$ à
$(\Sch/U_i)_{fppf}$. Si nous pouvons trouver des schémas $W_i$ sur $U_i$ et
des morphismes lisses surjectifs $W_i \to \mathcal{X}_i$, alors nous posons
$W = \coprod W_i$ et obtenons un morphisme lisse surjectif
$W \to \mathcal{X}$. Nous pouvons donc supposer que $S = B$ est affine.

\medskip\noindent
Supposons $S = B$ affine, disons $S = \Spec(\Lambda)$.
Écrivons $\Lambda = \colim \Lambda_i$ comme une limite inductive filtrante où
chaque $\Lambda_i$ est de type fini sur $\mathbf{Z}$. Pour un certain $i$,
nous pouvons trouver un morphisme d'espaces algébriques
$X_i \to \Spec(\Lambda_i)$ propre, plat et de présentation finie, dont le
changement de base à $\Lambda$ est $X$. Voir Limites d'espaces, lemmes
\ref{spaces-limits-lemma-descend-finite-presentation},
\ref{spaces-limits-lemma-descend-flat} et
\ref{spaces-limits-lemma-eventually-proper}.
Si nous montrons que $\Complexesstack_{X_i/\Spec(\Lambda_i)}$ est un champ
algébrique, alors il résulte du changement de base
(remarque \ref{quot-remark-complexes-base-change} et
Champs algébriques, section \ref{algebraic-section-change-base-scheme})
que $\mathcal{X}$ est un champ algébrique.
Nous pouvons donc supposer que $\Lambda$ est une $\mathbf{Z}$-algèbre de type
fini.

\medskip\noindent
Supposons que $S = B = \Spec(\Lambda)$ soit affine de type fini sur
$\mathbf{Z}$. Dans ce cas, nous allons vérifier les conditions (1), (2), (3),
(4) et (5) de
Axiomes d'Artin, lemme \ref{artin-lemma-diagonal-representable}, pour conclure
que $\mathcal{X}$ est un champ algébrique.
Notons que $\Lambda$ est un G-anneau, voir
Compléments d'algèbre, proposition
\ref{more-algebra-proposition-ubiquity-G-ring}.
Ainsi, tous les anneaux locaux de $S$ sont des G-anneaux. Donc (5) vaut.
Pour vérifier (2), nous devons vérifier les axiomes [-1], [0], [1], [2] et [3]
de Axiomes d'Artin, section \ref{artin-section-axioms}.
Nous omettons la vérification de [-1], et les axiomes
[0], [1], [2], [3] correspondent respectivement aux
lemmes \ref{quot-lemma-complexes-stack},
\ref{quot-lemma-complexes-limits},
\ref{quot-lemma-complexes-RS-star},
\ref{quot-lemma-complexes-tangent-space}.
La condition (3) résulte du lemme
\ref{quot-lemma-complexes-strong-effectiveness}.
La condition (1) est le lemme \ref{quot-lemma-complexes-diagonal}.

\medskip\noindent
Il reste à montrer la condition (4), qui est l'ouverture de la versalité.
Pour cela, nous allons utiliser
Axiomes d'Artin, lemme \ref{artin-lemma-SGE-implies-openness-versality}.
Nous avons déjà vu que la diagonale de $\mathcal{X}$ est représentable par des
espaces algébriques, qu'il satisfait à (RS*) et qu'il commute aux
limites (voir les lemmes utilisés ci-dessus).
Il nous suffit donc de voir que $\mathcal{X}$ satisfait à l'effectivité
formelle forte formulée dans
Axiomes d'Artin, lemme \ref{artin-lemma-SGE-implies-openness-versality}.
Cela résulte du lemme \ref{quot-lemma-complexes-strong-effectiveness}, et la
démonstration est achevée.
\end{proof}















% OK, start here.
%

\chapter{Propriétés des champs algébriques}



\phantomsection
\label{stacks-properties-section-phantom}


\section{Introduction}
\label{stacks-properties-section-introduction}

\noindent
Voir
Champs algébriques, section \ref{algebraic-section-introduction},
pour une brève introduction aux champs algébriques, et lire
une partie de ce chapitre pour connaître les fondements que nous adoptons.
Dans ce chapitre-là, nous nous attachons à distinguer soigneusement
les schémas, les espaces algébriques et les champs algébriques ; à partir
du présent chapitre, nous emploierons l'abus de langage usuel selon lequel
toutes ces notions s'emploient indifféremment.

\medskip\noindent
L'objectif de ce chapitre est d'introduire quelques notions et
propriétés fondamentales des champs algébriques. Une référence
essentielle pour le cas des champs algébriques quasi-séparés à diagonale
représentable est \cite{LM-B}.



\section{Conventions et abus de langage}
\label{stacks-properties-section-conventions}

\noindent
Fixons un grand site fppf $\Sch_{fppf}$.
Tous les schémas appartiennent à $\Sch_{fppf}$.
Et tous les anneaux $A$ considérés ont la propriété que
$\Spec(A)$ est (isomorphe) à un objet de ce grand site.

\medskip\noindent
Fixons également un schéma de base $S$, qui, d'après les conventions ci-dessus, est un élément
de $\Sch_{fppf}$. Le lecteur qui ne s'intéresse qu'au
cas absolu peut prendre $S = \Spec(\mathbf{Z})$.

\medskip\noindent
Voici nos conventions concernant les champs algébriques :
\begin{enumerate}
\item Lorsque nous dirons {\it champ algébrique}, nous entendrons un champ
algébrique au-dessus de $S$, c'est-à-dire une catégorie fibrée en groupoïdes
$p : \mathcal{X} \to (\Sch/S)_{fppf}$
qui satisfait aux conditions de
Champs algébriques, Définition \ref{algebraic-definition-algebraic-stack}.
\item Nous dirons que $f : \mathcal{X} \to \mathcal{Y}$ est un
{\it morphisme de champs algébriques} pour désigner un $1$-morphisme
de champs algébriques au-dessus de $S$, c'est-à-dire un $1$-morphisme de catégories fibrées
en groupoïdes au-dessus de $(\Sch/S)_{fppf}$, voir
Champs algébriques,
Définition \ref{algebraic-definition-morphism-algebraic-stacks}.
\item Un {\it $2$-morphisme} $\alpha : f \to g$ désignera
un $2$-morphisme dans la $2$-catégorie des champs algébriques au-dessus de
$S$, voir
Champs algébriques,
Définition \ref{algebraic-definition-morphism-algebraic-stacks}.
\item Étant donnés des morphismes $\mathcal{X} \to \mathcal{Z}$
et $\mathcal{Y} \to \mathcal{Z}$ de champs algébriques,
nous appelons abusivement le $2$-produit fibré
$\mathcal{X} \times_\mathcal{Z} \mathcal{Y}$ le {\it produit fibré}.
\item Nous noterons $\mathcal{X} \times_S \mathcal{Y}$ le
produit des champs algébriques $\mathcal{X}$, $\mathcal{Y}$.
\item Nous abuserons souvent de la notation et dirons que deux champs algébriques
$\mathcal{X}$ et $\mathcal{Y}$ sont {\it isomorphes} s'ils sont
équivalents dans cette $2$-catégorie.
\end{enumerate}
Voici nos conventions concernant les espaces algébriques.
\begin{enumerate}
\item Si nous disons que $X$ est un {\it espace algébrique}, nous entendons que
$X$ est un espace algébrique au-dessus de $S$, c'est-à-dire que $X$ est un préfaisceau sur
$(\Sch/S)_{fppf}$ qui satisfait aux conditions de
Espaces, Définition \ref{spaces-definition-algebraic-space}.
\item Un {\it morphisme d'espaces algébriques} $f :X \to Y$ est un morphisme
d'espaces algébriques au-dessus de $S$ tel que défini dans
Espaces, Définition \ref{spaces-definition-morphism-algebraic-spaces}.
\item Nous ne distinguerons {\bf pas} un espace algébrique $X$
du champ algébrique $\mathcal{S}_X \to (\Sch/S)_{fppf}$
auquel il donne lieu, voir
Champs algébriques, Lemme \ref{algebraic-lemma-representable-algebraic}.
\item En particulier, un {\it morphisme} $f : X \to \mathcal{Y}$ de $X$
vers un champ algébrique $\mathcal{Y}$ désigne un morphisme
$f : \mathcal{S}_X \to \mathcal{Y}$ de champs algébriques.
Il en va de même pour les morphismes $\mathcal{Y} \to X$.
\item De plus, étant donné un champ algébrique $\mathcal{X}$, nous dirons que
{\it $\mathcal{X}$ est un espace algébrique} pour indiquer que $\mathcal{X}$
est représentable par un espace algébrique, voir
Champs algébriques,
Définition \ref{algebraic-definition-representable-by-algebraic-space}.
\item Nous emploierons la convention de notation suivante : si nous
désignons un champ algébrique par une majuscule romaine
(telle que $X, Y, Z, A, B, \ldots$), son
champ d'inertie sera trivial, et ce sera donc un espace algébrique, voir
Champs algébriques,
Proposition \ref{algebraic-proposition-algebraic-stack-no-automorphisms}.
\end{enumerate}
Voici nos conventions concernant les schémas.
\begin{enumerate}
\item Si nous disons que $X$ est un {\it schéma}, nous entendons que
$X$ est un schéma au-dessus de $S$, c'est-à-dire que $X$ est un objet de $(\Sch/S)_{fppf}$.
\item Par {\it morphisme de schémas}, nous entendons un morphisme de schémas au-dessus de $S$.
\item Nous ne distinguerons {\bf pas} un schéma $X$ du
champ algébrique $\mathcal{S}_X \to (\Sch/S)_{fppf}$ auquel il donne lieu,
voir
Champs algébriques, Lemme \ref{algebraic-lemma-representable-algebraic}.
\item En particulier, un {\it morphisme} $f : X \to \mathcal{Y}$ d'un
schéma $X$ vers un champ algébrique $\mathcal{Y}$ désigne un morphisme
$f : \mathcal{S}_X \to \mathcal{Y}$ de champs algébriques.
Il en va de même pour les morphismes $\mathcal{Y} \to X$.
\item De plus, étant donné un champ algébrique $\mathcal{X}$, nous dirons que
{\it $\mathcal{X}$ est un schéma} pour indiquer que $\mathcal{X}$
est représentable, voir
Champs algébriques, section \ref{algebraic-section-representable}.
\end{enumerate}
Voici nos conventions concernant les morphismes de champs algébriques :
\begin{enumerate}
\item Un morphisme $f : \mathcal{X} \to \mathcal{Y}$
de champs algébriques est {\it représentable}, ou
{\it représentable par des schémas}, si, pour tout schéma
$T$ et tout morphisme $T \to \mathcal{Y}$, le produit fibré
$T \times_\mathcal{Y} \mathcal{X}$ est un schéma.
Voir
Champs algébriques, section \ref{algebraic-section-representable-morphism}.
\item Un morphisme $f : \mathcal{X} \to \mathcal{Y}$
de champs algébriques est {\it représentable par des espaces algébriques}
si, pour tout schéma $T$ et tout morphisme $T \to \mathcal{Y}$, le produit fibré
$T \times_\mathcal{Y} \mathcal{X}$ est un espace algébrique.
Voir Champs algébriques,
Définition \ref{algebraic-definition-representable-by-algebraic-spaces}.
Dans ce cas, $Z \times_\mathcal{Y} \mathcal{X}$ est un espace
algébrique dès que $Z \to \mathcal{Y}$ est un morphisme dont la source est
un espace algébrique, voir
Champs algébriques,
Lemme \ref{algebraic-lemma-base-change-by-space-representable-by-space}.
\item Nous pourrons abuser du langage et dire qu'un diagramme de champs algébriques
commute si le diagramme est $2$-commutatif dans la $2$-catégorie des
champs algébriques.
\end{enumerate}
Remarquons que tout morphisme $X \to \mathcal{Y}$ d'un espace algébrique
vers un champ algébrique est représentable par des espaces algébriques, voir
Champs algébriques, Lemme \ref{algebraic-lemma-representable-diagonal}.
Nous utiliserons ce résultat fondamental sans plus le rappeler.




\section{Propriétés des morphismes représentables par des espaces algébriques}
\label{stacks-properties-section-properties-morphisms}

\noindent
Nous étudierons les propriétés des morphismes (quelconques) de champs algébriques
dans un chapitre qui leur est propre. Pour les morphismes représentables par des espaces algébriques, nous savons ce que
signifie être surjectif, lisse, ou \'etale, etc. Cela s'applique en particulier
aux morphismes $X \to \mathcal{Y}$ des espaces algébriques vers les champs algébriques.
Dans cette section, nous rappelons le fonctionnement de cette notion, nous dressons la liste des propriétés
auxquelles elle s'applique et nous démontrons quelques lemmes faciles.

\medskip\noindent
Notre premier lemme affirme qu'un morphisme est représentable par des espaces algébriques
s'il le devient après un changement de base par un morphisme plat,
localement de présentation finie et surjectif.

\begin{lemma}
\label{stacks-properties-lemma-check-representable-covering}
Soit $f : \mathcal{X} \to \mathcal{Y}$ un morphisme de champs algébriques.
Soit $W$ un espace algébrique et soit $W \to \mathcal{Y}$ un morphisme surjectif,
localement de présentation finie et plat. Les assertions suivantes sont équivalentes :
\begin{enumerate}
\item $f$ est représentable par des espaces algébriques, et
\item $W \times_\mathcal{Y} \mathcal{X}$ est un espace algébrique.
\end{enumerate}
\end{lemma}

\begin{proof}
L'implication (1) $\Rightarrow$ (2) est
Champs algébriques,
Lemme \ref{algebraic-lemma-base-change-by-space-representable-by-space}.
Réciproquement, soit $W \to \mathcal{Y}$ comme en (2). Pour démontrer (1), il
suffit de montrer que $f$ est fidèle sur les catégories fibres, voir
Champs algébriques,
Lemme \ref{algebraic-lemma-characterize-representable-by-algebraic-spaces}.
L'hypothèse (2) implique en particulier que
$W \times_\mathcal{Y} \mathcal{X} \to W$ est fidèle.
La fidélité de $f$ résulte donc de
Champs, Lemme \ref{stacks-lemma-faithful-descent}.
\end{proof}

\noindent
Soit $P$ une propriété des morphismes d'espaces algébriques qui est
locale pour la topologie fppf sur le but et préservée par tout changement de base.
Soit $f : \mathcal{X} \to \mathcal{Y}$ un morphisme de champs algébriques
représentable par des espaces algébriques. Nous disons alors que
{\it $f$ possède la propriété $P$} si et seulement si, pour tout schéma $T$
et tout morphisme $T \to \mathcal{Y}$, le morphisme d'espaces algébriques
$T \times_\mathcal{Y} \mathcal{X} \to T$ possède la propriété $P$, voir
Champs algébriques,
Définition \ref{algebraic-definition-relative-representable-property}.

\medskip\noindent
Il se trouve que si $f : \mathcal{X} \to \mathcal{Y}$ est représentable
par des espaces algébriques et possède la propriété $P$, alors, pour tout morphisme de champs algébriques
$\mathcal{Y}' \to \mathcal{Y}$, le changement de base
$\mathcal{Y}' \times_\mathcal{Y} \mathcal{X} \to \mathcal{Y}'$
possède la propriété $P$, voir
Champs algébriques,
Lemmes \ref{algebraic-lemma-base-change-representable-by-spaces} et
\ref{algebraic-lemma-base-change-representable-transformations-property}.
Si la propriété $P$ est préservée par composition, cela vaut
également pour les morphismes de champs algébriques représentables par
des espaces algébriques, voir
Champs algébriques,
Lemmes \ref{algebraic-lemma-composition-representable-by-spaces} et
\ref{algebraic-lemma-composition-representable-transformations-property}.
En outre, dans ce cas, les produits
$\mathcal{X}_1 \times \mathcal{X}_2 \to \mathcal{Y}_1 \times \mathcal{Y}_2$
de morphismes représentables par des espaces algébriques et possédant la propriété $\mathcal{P}$
possèdent la propriété $\mathcal{P}$, voir
Champs algébriques, Lemme
\ref{algebraic-lemma-product-representable-transformations-property}.

\medskip\noindent
Enfin, si nous avons deux propriétés $P, P'$ de morphismes d'espaces algébriques
qui sont locales pour la topologie fppf sur le but et préservées par tout changement de base,
et si $P(f) \Rightarrow P'(f)$ pour tout morphisme $f$, alors la même
implication vaut pour les propriétés correspondantes des morphismes de champs algébriques
représentables par des espaces algébriques, voir
Champs algébriques, Lemme
\ref{algebraic-lemma-representable-transformations-property-implication}.
Nous l'utiliserons {\it sans autre mention} ci-dessous et dans les
chapitres suivants.

\medskip\noindent
La discussion ci-dessus s'applique à chacune des propriétés suivantes des
morphismes d'espaces algébriques :
\begin{enumerate}
\item quasi-compact, voir
Morphismes d'espaces,
Lemme \ref{spaces-morphisms-lemma-base-change-quasi-compact}
et
Descente sur les espaces,
Lemme \ref{spaces-descent-lemma-descending-property-quasi-compact},
\item quasi-séparé, voir
Morphismes d'espaces,
Lemme \ref{spaces-morphisms-lemma-base-change-separated}
et
Descente sur les espaces,
Lemme \ref{spaces-descent-lemma-descending-property-quasi-separated},
\item universellement fermé, voir
Morphismes d'espaces,
Lemme \ref{spaces-morphisms-lemma-base-change-universally-closed}
et
Descente sur les espaces,
Lemme \ref{spaces-descent-lemma-descending-property-universally-closed},
\item universellement ouvert, voir
Morphismes d'espaces,
Lemme \ref{spaces-morphisms-lemma-base-change-universally-open}
et
Descente sur les espaces,
Lemme \ref{spaces-descent-lemma-descending-property-universally-open},
\item universellement submersif, voir
Morphismes d'espaces,
Lemme \ref{spaces-morphisms-lemma-base-change-universally-submersive}
et
Descente sur les espaces,
Lemme \ref{spaces-descent-lemma-descending-property-universally-submersive},
\item un homéomorphisme universel, voir
Morphismes d'espaces,
Lemme \ref{spaces-morphisms-lemma-base-change-universal-homeomorphism}
et
Descente sur les espaces,
Lemme \ref{spaces-descent-lemma-descending-property-universal-homeomorphism},
\item surjectif, voir
Morphismes d'espaces,
Lemme \ref{spaces-morphisms-lemma-base-change-surjective}
et
Descente sur les espaces,
Lemme \ref{spaces-descent-lemma-descending-property-surjective},
\item universellement injectif, voir
Morphismes d'espaces,
Lemme \ref{spaces-morphisms-lemma-base-change-universally-injective}
et
Descente sur les espaces,
Lemme \ref{spaces-descent-lemma-descending-property-universally-injective},
\item localement de type fini, voir
Morphismes d'espaces,
Lemme \ref{spaces-morphisms-lemma-base-change-finite-type}
et
Descente sur les espaces,
Lemme \ref{spaces-descent-lemma-descending-property-locally-finite-type},
\item localement de présentation finie, voir
Morphismes d'espaces,
Lemme \ref{spaces-morphisms-lemma-base-change-finite-presentation}
et
Descente sur les espaces, Lemme
\ref{spaces-descent-lemma-descending-property-locally-finite-presentation},
\item de type fini, voir
Morphismes d'espaces,
Lemme \ref{spaces-morphisms-lemma-base-change-finite-type}
et
Descente sur les espaces,
Lemme \ref{spaces-descent-lemma-descending-property-finite-type},
\item de présentation finie, voir
Morphismes d'espaces,
Lemme \ref{spaces-morphisms-lemma-base-change-finite-presentation}
et
Descente sur les espaces, Lemme
\ref{spaces-descent-lemma-descending-property-finite-presentation},
\item plat, voir
Morphismes d'espaces,
Lemme \ref{spaces-morphisms-lemma-base-change-flat}
et
Descente sur les espaces,
Lemme \ref{spaces-descent-lemma-descending-property-flat},
\item une immersion ouverte, voir
Morphismes d'espaces,
section \ref{spaces-morphisms-section-immersions}
et
Descente sur les espaces,
Lemme \ref{spaces-descent-lemma-descending-property-open-immersion},
\item un isomorphisme, voir
Descente sur les espaces,
Lemme \ref{spaces-descent-lemma-descending-property-isomorphism},
\item affine, voir
Morphismes d'espaces,
Lemme \ref{spaces-morphisms-lemma-base-change-affine}
et
Descente sur les espaces,
Lemme \ref{spaces-descent-lemma-descending-property-affine},
\item une immersion fermée, voir
Morphismes d'espaces, section \ref{spaces-morphisms-section-immersions}
et
Descente sur les espaces,
Lemme \ref{spaces-descent-lemma-descending-property-closed-immersion},
\item séparé, voir
Morphismes d'espaces,
Lemme \ref{spaces-morphisms-lemma-base-change-separated}
et
Descente sur les espaces,
Lemme \ref{spaces-descent-lemma-descending-property-separated},
\item propre, voir
Morphismes d'espaces,
Lemme \ref{spaces-morphisms-lemma-base-change-proper}
et
Descente sur les espaces,
Lemme \ref{spaces-descent-lemma-descending-property-proper},
\item quasi-affine, voir
Morphismes d'espaces,
Lemme \ref{spaces-morphisms-lemma-base-change-quasi-affine}
et
Descente sur les espaces,
Lemme \ref{spaces-descent-lemma-descending-property-quasi-affine},
\item entier, voir
Morphismes d'espaces,
Lemme \ref{spaces-morphisms-lemma-base-change-integral}
et
Descente sur les espaces,
Lemme \ref{spaces-descent-lemma-descending-property-integral},
\item fini, voir
Morphismes d'espaces,
Lemme \ref{spaces-morphisms-lemma-base-change-integral}
et
Descente sur les espaces,
Lemme \ref{spaces-descent-lemma-descending-property-finite},
\item quasi-fini (localement), voir
Morphismes d'espaces,
Lemme \ref{spaces-morphisms-lemma-base-change-quasi-finite}
et
Descente sur les espaces,
Lemme \ref{spaces-descent-lemma-descending-property-quasi-finite},
\item syntomique, voir
Morphismes d'espaces,
Lemme \ref{spaces-morphisms-lemma-base-change-syntomic}
et
Descente sur les espaces,
Lemme \ref{spaces-descent-lemma-descending-property-syntomic},
\item lisse, voir
Morphismes d'espaces,
Lemme \ref{spaces-morphisms-lemma-base-change-smooth}
et
Descente sur les espaces,
Lemme \ref{spaces-descent-lemma-descending-property-smooth},
\item non ramifié, voir
Morphismes d'espaces,
Lemme \ref{spaces-morphisms-lemma-base-change-unramified}
et
Descente sur les espaces,
Lemme \ref{spaces-descent-lemma-descending-property-unramified},
\item \'etale, voir
Morphismes d'espaces,
Lemme \ref{spaces-morphisms-lemma-base-change-etale}
et
Descente sur les espaces,
Lemme \ref{spaces-descent-lemma-descending-property-etale},
\item fini localement libre, voir
Morphismes d'espaces,
Lemme \ref{spaces-morphisms-lemma-base-change-finite-locally-free}
et
Descente sur les espaces,
Lemme \ref{spaces-descent-lemma-descending-property-finite-locally-free},
\item un monomorphisme, voir
Morphismes d'espaces,
Lemme \ref{spaces-morphisms-lemma-base-change-monomorphism}
et
Descente sur les espaces,
Lemme \ref{spaces-descent-lemma-descending-property-monomorphism},
\item une immersion, voir
Morphismes d'espaces, section \ref{spaces-morphisms-section-immersions}
et
Descente sur les espaces,
Lemme \ref{spaces-descent-lemma-descending-fppf-property-immersion},
\item localement séparé, voir
Morphismes d'espaces,
Lemme \ref{spaces-morphisms-lemma-base-change-separated}
et
Descente sur les espaces,
Lemme \ref{spaces-descent-lemma-descending-fppf-property-locally-separated},
\end{enumerate}

\begin{lemma}
\label{stacks-properties-lemma-property-spaces-too}
Soit $P$ une propriété des morphismes d'espaces algébriques comme ci-dessus.
Soit $f : \mathcal{X} \to \mathcal{Y}$ un morphisme de champs algébriques
représentable par des espaces algébriques. Les assertions suivantes sont équivalentes :
\begin{enumerate}
\item $f$ possède $P$,
\item pour tout espace algébrique $Z$ et tout morphisme $Z \to \mathcal{Y}$,
le morphisme $Z \times_\mathcal{Y} \mathcal{X} \to Z$ possède $P$.
\end{enumerate}
\end{lemma}

\begin{proof}
L'implication (2) $\Rightarrow$ (1) est immédiate. Supposons (1).
Soit $Z \to \mathcal{Y}$ comme en (2). Choisissons un schéma $U$ et un
morphisme \'etale surjectif $U \to Z$. Par hypothèse, le morphisme
$U \times_\mathcal{Y} \mathcal{X} \to U$ possède $P$. Mais le diagramme
$$
\xymatrix{
U \times_\mathcal{Y} \mathcal{X} \ar[d] \ar[r] &
Z \times_\mathcal{Y} \mathcal{X} \ar[d] \\
U \ar[r] & Z
}
$$
est cartésien; la flèche verticale de droite possède donc $P$, puisque
$\{U \to Z\}$ est un recouvrement fppf.
\end{proof}

\noindent
Le lemme suivant nous dit qu'il suffit de vérifier $P$
après un changement de base par un morphisme surjectif, plat et localement
de présentation finie.

\begin{lemma}
\label{stacks-properties-lemma-check-property-covering}
Soit $P$ une propriété des morphismes d'espaces algébriques comme ci-dessus.
Soit $f : \mathcal{X} \to \mathcal{Y}$ un morphisme de champs algébriques
représentable par des espaces algébriques.
Soit $W$ un espace algébrique et soit $W \to \mathcal{Y}$ un morphisme surjectif,
localement de présentation finie et plat.
Posons $V = W \times_\mathcal{Y} \mathcal{X}$. Alors
$$
(f\text{ possède }P) \Leftrightarrow (\text{la projection }V \to W\text{ possède }P).
$$
\end{lemma}

\begin{proof}
L'implication de gauche à droite résulte du
Lemme \ref{stacks-properties-lemma-property-spaces-too}.
Supposons que $V \to W$ possède $P$. Soit $T$ un schéma et soit
$T \to \mathcal{Y}$ un morphisme. Considérons le diagramme commutatif
$$
\xymatrix{
T \times_\mathcal{Y} \mathcal{X} \ar[d] &
T \times_\mathcal{Y} V \ar[d] \ar[l] \ar[r] &
V \ar[d] \\
T & T \times_\mathcal{Y} W \ar[l] \ar[r] & W
}
$$
d'espaces algébriques. Les carrés sont cartésiens.
Le morphisme inférieur de gauche est un morphisme surjectif et plat, localement de
présentation finie; $\{T \times_\mathcal{Y} V \to T\}$ est donc un
recouvrement fppf. Par conséquent, le fait que la flèche verticale de droite possède la propriété
$P$ implique que la flèche verticale de gauche possède la propriété $P$.
\end{proof}

\begin{lemma}
\label{stacks-properties-lemma-check-property-weak-covering}
Soit $P$ une propriété des morphismes d'espaces algébriques comme ci-dessus.
Soit $f : \mathcal{X} \to \mathcal{Y}$ un morphisme de champs algébriques
représentable par des espaces algébriques.
Soit $\mathcal{Z} \to \mathcal{Y}$ un morphisme de champs algébriques qui
est représentable par des espaces algébriques, surjectif, plat et
localement de présentation finie.
Posons $\mathcal{W} = \mathcal{Z} \times_\mathcal{Y} \mathcal{X}$. Alors
$$
(f\text{ possède }P) \Leftrightarrow
(\text{la projection }\mathcal{W} \to \mathcal{Z}\text{ possède }P).
$$
\end{lemma}

\begin{proof}
Choisissons un espace algébrique $W$ et un morphisme
$W \to \mathcal{Z}$ qui soit surjectif, plat et localement de présentation
finie. D'après la discussion ci-dessus, le composé
$W \to \mathcal{Y}$ est aussi surjectif, plat et localement de présentation
finie. Notons
$V = W \times_\mathcal{Z} \mathcal{W} = V \times_\mathcal{Y} \mathcal{X}$.
Par le
Lemme \ref{stacks-properties-lemma-check-property-covering},
nous voyons que $f$ possède $\mathcal{P}$ si et seulement si $V \to W$ la possède,
et que $\mathcal{W} \to \mathcal{Z}$ possède $\mathcal{P}$ si et seulement
si $V \to W$ la possède. Le lemme en résulte.
\end{proof}

\begin{lemma}
\label{stacks-properties-lemma-check-property-after-precomposing}
Soit $P$ une propriété des morphismes d'espaces algébriques comme ci-dessus.
Soit $\tau \in \{\etale, smooth, syntomic, fppf\}$.
Soient $\mathcal{X} \to \mathcal{Y}$ et $\mathcal{Y} \to \mathcal{Z}$
des morphismes de champs algébriques représentables par des espaces algébriques.
Supposons que
\begin{enumerate}
\item $\mathcal{X} \to \mathcal{Y}$ soit surjectif et
\'etale, lisse, syntomique, ou plat et localement de présentation finie,
\item le composé possède $P$, et
\item $P$ soit locale sur la source pour la topologie $\tau$.
\end{enumerate}
Alors $\mathcal{Y} \to \mathcal{Z}$ possède la propriété $P$.
\end{lemma}

\begin{proof}
Soit $Z$ un schéma et soit $Z \to \mathcal{Z}$ un morphisme.
Posons $X = \mathcal{X} \times_\mathcal{Z} Z$,
$Y = \mathcal{Y} \times_\mathcal{Z} Z$. D'après (1), $\{X \to Y\}$
est un recouvrement pour la topologie $\tau$ d'espaces algébriques, et, d'après (2), $X \to Z$ possède la propriété
$P$. D'après (3), il en résulte que $Y \to Z$ possède la propriété $P$,
ce qui conclut la démonstration.
\end{proof}

\begin{lemma}
\label{stacks-properties-lemma-representable-in-terms-presentations}
Soit $g : \mathcal{X}' \to \mathcal{X}$ un morphisme de champs algébriques
qui soit représentable par des espaces algébriques. Soit $[U/R] \to \mathcal{X}$
une présentation. Posons $U' = U \times_\mathcal{X} \mathcal{X}'$,
et $R' = R \times_\mathcal{X} \mathcal{X}'$.
Il existe alors un groupoïde en espaces algébriques de la forme
$(U', R', s', t', c')$, une présentation $[U'/R'] \to \mathcal{X}'$,
et le diagramme
$$
\xymatrix{
[U'/R'] \ar[d]_{[\text{pr}]} \ar[r] & \mathcal{X}' \ar[d]^g \\
[U/R] \ar[r] & \mathcal{X}
}
$$
est $2$-commutatif, où le morphisme $[\text{pr}]$ provient d'un
morphisme de groupoïdes
$\text{pr} : (U', R', s', t', c') \to (U, R, s, t, c)$.
\end{lemma}

\begin{proof}
Puisque $U \to \mathcal{Y}$ est surjectif et lisse, voir
Champs algébriques,
Lemme \ref{algebraic-lemma-smooth-quotient-smooth-presentation},
le changement de base $U' \to \mathcal{X}'$ est aussi surjectif et lisse.
Par suite, d'après
Champs algébriques,
Lemme \ref{algebraic-lemma-stack-presentation},
il suffit de montrer que $R' = U' \times_{\mathcal{X}'} U'$ afin
d'obtenir un groupoïde lisse $(U', R', s', t', c')$ et une présentation
$[U'/R'] \to \mathcal{X}'$.
Comme $R = V \times_\mathcal{Y} V$ (voir
Groupoïdes en espaces,
Lemme \ref{spaces-groupoids-lemma-quotient-stack-2-cartesian}),
cela résulte de
$$
R' =
U \times_\mathcal{X} U \times_\mathcal{X} \mathcal{X}' =
(U \times_\mathcal{X} \mathcal{X}')
\times_{\mathcal{X}'}
(U \times_\mathcal{X} \mathcal{X}')
$$
voir
Catégories, Lemmes \ref{categories-lemma-associativity-2-fibre-product} et
\ref{categories-lemma-2-fibre-product-erase-factor}.
Il est clair que les morphismes de projection $U' \to U$ et $R' \to R$ donnent le
morphisme de groupoïdes désiré
$\text{pr} : (U', R', s', t', c') \to (U, R, s, t, c)$.
On obtient donc le morphisme $[\text{pr}]$ des champs quotients grâce à
Groupoïdes en espaces,
Lemme \ref{spaces-groupoids-lemma-quotient-stack-functorial}.

\medskip\noindent
Il nous reste à montrer que le diagramme est $2$-commutatif.
Il est clair que le diagramme
$$
\xymatrix{
U' \ar[d]_{\text{pr}_U} \ar[r]_{f'} & \mathcal{X}' \ar[d]^g \\
U \ar[r]^f & \mathcal{X}
}
$$
est $2$-commutatif, où $\text{pr}_U : U' \to U$ est la projection.
Il existe une $2$-flèche canonique
$\tau : f \circ t \to f \circ s$ dans $\Mor(R, \mathcal{X})$
provenant de $R = U \times_\mathcal{X} U$, $t = \text{pr}_0$, et
$s = \text{pr}_1$. En utilisant l'isomorphisme
$R' \to U' \times_{\mathcal{X}'} U'$, nous obtenons de même un isomorphisme
$\tau' : f' \circ t' \to f' \circ s'$. Remarquons que
$g \circ f' \circ t' = f \circ t \circ \text{pr}_R$ et
$g \circ f' \circ s' = f \circ s \circ \text{pr}_R$, où
$\text{pr}_R : R' \to R$ est la projection. Il est donc légitime de se demander
si
\begin{equation}
\label{stacks-properties-equation-verify}
\tau \star \text{id}_{\text{pr}_R} = \text{id}_g \star \tau'.
\end{equation}
Nous avançons maintenant deux assertions : (1) si l'équation (\ref{stacks-properties-equation-verify}) est satisfaite,
alors le diagramme est $2$-commutatif, et (2) l'équation (\ref{stacks-properties-equation-verify}) est satisfaite.
Nous omettons la démonstration de ces deux assertions. Indications : la partie (1) résulte de la
construction de $f = f_{can}$ et $f' = f'_{can}$ dans
Champs algébriques, Lemme \ref{algebraic-lemma-map-space-into-stack}.
La partie (2) résulte d'un examen attentif des définitions.
\end{proof}

\begin{remark}
\label{stacks-properties-remark-representable-over}
Soit $\mathcal{Y}$ un champ algébrique. Considérons la
$2$-catégorie suivante :
\begin{enumerate}
\item un objet est un morphisme $f : \mathcal{X} \to \mathcal{Y}$
représentable par des espaces algébriques,
\item un $1$-morphisme
$(g, \beta) :
(f_1 : \mathcal{X}_1 \to \mathcal{Y})
\to
(f_2 : \mathcal{X}_2 \to \mathcal{Y})$
consiste en un morphisme $g : \mathcal{X}_1 \to \mathcal{X}_2$ et un
$2$-morphisme $\beta : f_1 \to f_2 \circ g$, et
\item un $2$-morphisme entre
$(g, \beta), (g', \beta') :
(f_1 : \mathcal{X}_1 \to \mathcal{Y})
\to
(f_2 : \mathcal{X}_2 \to \mathcal{Y})$
est un $2$-morphisme $\alpha : g \to g'$ tel que
$(\text{id}_{f_2} \star \alpha) \circ \beta = \beta'$.
\end{enumerate}
Notons cette $2$-catégorie $\textit{Espaces}/\mathcal{Y}$, par
analogie avec la notation de
Topologies sur les espaces, section \ref{spaces-topologies-section-procedure}.
Nous affirmons que, dans cette $2$-catégorie, les catégories de morphismes
$$
\Mor_{\textit{Espaces}/\mathcal{Y}}(
(f_1 : \mathcal{X}_1 \to \mathcal{Y}),
(f_2 : \mathcal{X}_2 \to \mathcal{Y}))
$$
sont toutes des setoïdes. En effet, un $2$-morphisme $\alpha$ est une règle qui, à tout
objet $x_1$ de $\mathcal{X}_1$, associe un isomorphisme
$\alpha_{x_1} : g(x_1) \longrightarrow g'(x_1)$
dans la catégorie fibre correspondante de $\mathcal{X}_2$, de telle sorte que le diagramme
$$
\xymatrix{
& f_2(x_1) \ar[ld]_{\beta_{x_1}} \ar[rd]^{\beta'_{x_1}} \\
f_2(g(x_1)) \ar[rr]^{f_2(\alpha_{x_1})} & &
f_2(g'(x_1))
}
$$
commute. Mais, puisque $f_2$ est fidèle (voir
Champs algébriques,
Lemme \ref{algebraic-lemma-characterize-representable-by-algebraic-spaces}),
cela signifie que, si $\alpha_{x_1}$ existe, il est unique ! Autrement dit, la
$2$-catégorie $\textit{Espaces}/\mathcal{Y}$ est très proche d'être
une catégorie. Plus précisément, si nous remplaçons les $1$-morphismes par les classes
d'isomorphisme de $1$-morphismes, nous obtenons une catégorie. Nous effectuerons souvent
ce remplacement sans autre mention.
\end{remark}




\section{Points des champs algébriques}
\label{stacks-properties-section-points}

\noindent
Soit $\mathcal{X}$ un champ algébrique. Soient $K, L$ deux corps,
et soient $p : \Spec(K) \to \mathcal{X}$ et
$q : \Spec(L) \to \mathcal{X}$ des morphismes.
Nous disons que $p$ et $q$ sont {\it équivalents} s'il existe un
corps $\Omega$ et un diagramme $2$-commutatif
$$
\xymatrix{
\Spec(\Omega) \ar[r] \ar[d] &
\Spec(L) \ar[d]^q \\
\Spec(K) \ar[r]^p &
\mathcal{X}.
}
$$

\begin{lemma}
\label{stacks-properties-lemma-equivalence}
La notion ci-dessus définit bien une relation d'équivalence sur les
morphismes des spectres de corps vers le champ algébrique $\mathcal{X}$.
\end{lemma}

\begin{proof}
Il est clair que la relation est réflexive et symétrique.
Il reste donc à montrer qu'elle est transitive. Cela revient
à ceci : étant donné un diagramme
$$
\xymatrix{
\Spec(\Omega) \ar[r]_b \ar[d]_a &
\Spec(L) \ar[d]^q & \Spec(\Omega') \ar[l]^{b'} \ar[d]^{a'} \\
\Spec(K) \ar[r]^p &
\mathcal{X} &
\Spec(K') \ar[l]_{p'}
}
$$
dont les deux carrés sont $2$-commutatifs, nous devons montrer que $p$ est équivalent à
$p'$. Par le lemme de $2$-Yoneda (voir
Champs algébriques, section \ref{algebraic-section-2-yoneda}),
les morphismes $p$, $p'$ et $q$ sont donnés par des objets
$x$, $x'$ et $y$ dans les catégories fibres de $\mathcal{X}$ au-dessus de
$\Spec(K)$, $\Spec(K')$ et $\Spec(L)$. La
$2$-commutativité des carrés signifie qu'il existe des isomorphismes
$\alpha : a^*x \to b^*y$ et $\alpha' : (a')^*x' \to (b')^*y$
dans les catégories fibres
de $\mathcal{X}$ au-dessus de $\Spec(\Omega)$ et $\Spec(\Omega')$.
Choisissons un corps $\Omega''$ et des plongements
$\Omega \to \Omega''$ et $\Omega' \to \Omega''$ coïncidant sur $L$.
Nous pouvons alors prolonger le diagramme ci-dessus en
$$
\xymatrix{
& \Spec(\Omega'') \ar[ld]_c \ar[d]^{q'} \ar[rd]^{c'} \\
\Spec(\Omega) \ar[r]_b \ar[d]_a &
\Spec(L) \ar[d]^q & \Spec(\Omega') \ar[l]^{b'} \ar[d]^{a'} \\
\Spec(K) \ar[r]^p &
\mathcal{X} &
\Spec(K') \ar[l]_{p'}
}
$$
avec des triangles commutatifs, et
$$
(q')^*(\alpha')^{-1} \circ (q')^*\alpha :
(a \circ c)^*x
\longrightarrow
(a' \circ c')^*x'
$$
est un isomorphisme dans la catégorie fibre au-dessus de $\Spec(\Omega'')$.
Ainsi, $p$ est équivalent à $p'$, comme souhaité.
\end{proof}

\begin{definition}
\label{stacks-properties-definition-points}
Soit $\mathcal{X}$ un champ algébrique.
Un {\it point} de $\mathcal{X}$ est une classe d'équivalence de morphismes
de spectres de corps vers $\mathcal{X}$.
L'ensemble des points de $\mathcal{X}$ est noté $|\mathcal{X}|$.
\end{definition}

\noindent
Cela concorde avec notre définition des points des espaces algébriques, voir
Propriétés des espaces, Définition \ref{spaces-properties-definition-points}.
De plus, pour un schéma, nous retrouvons la notion usuelle de point, voir
Propriétés des espaces, Lemme \ref{spaces-properties-lemma-scheme-points}.
Si $f : \mathcal{X} \to \mathcal{Y}$ est un morphisme de champs algébriques,
il existe une application induite $|f| : |\mathcal{X}| \to |\mathcal{Y}|$ qui
envoie un représentant $x : \Spec(K) \to \mathcal{X}$ sur le
représentant $f \circ x : \Spec(K) \to \mathcal{Y}$. Celle-ci est
bien définie : en effet, lorsqu'ils sont $2$-isomorphes, des $1$-morphismes restent $2$-isomorphes après
pré- ou postcomposition par un $1$-morphisme, car on peut les pré- ou postcomposer horizontalement
par l'identité du $1$-morphisme donné. Cela vaut
dans toute $(2, 1)$-catégorie (stricte). Si
$$
\xymatrix{
\mathcal{X} \ar[d] \ar[r] & \mathcal{Y} \ar[d] \\
\mathcal{W} \ar[r] & \mathcal{Z}
}
$$
est un diagramme $2$-commutatif de champs algébriques, alors le diagramme
d'ensembles
$$
\xymatrix{
|\mathcal{X}| \ar[d] \ar[r] & |\mathcal{Y}| \ar[d] \\
|\mathcal{W}| \ar[r] & |\mathcal{Z}|
}
$$
est commutatif.
En particulier, si $\mathcal{X} \to \mathcal{Y}$ est une équivalence,
alors $|\mathcal{X}| \to |\mathcal{Y}|$ est une bijection.

\begin{lemma}
\label{stacks-properties-lemma-points-cartesian}
Soit
$$
\xymatrix{
\mathcal{Z} \times_\mathcal{Y} \mathcal{X} \ar[r] \ar[d] &
\mathcal{X} \ar[d] \\
\mathcal{Z} \ar[r] & \mathcal{Y}
}
$$
un produit fibré de champs algébriques. Alors l'application d'ensembles
de points
$$
|\mathcal{Z} \times_\mathcal{Y} \mathcal{X}|
\longrightarrow
|\mathcal{Z}| \times_{|\mathcal{Y}|} |\mathcal{X}|
$$
est surjective.
\end{lemma}

\begin{proof}
En effet, supposons donnés des corps $K$, $L$ et des morphismes
$\Spec(K) \to \mathcal{X}$, $\Spec(L) \to \mathcal{Z}$.
L'hypothèse qu'ils coïncident comme éléments de $|\mathcal{Y}|$ signifie alors
qu'il existe une extension commune $M/K$ et $M/L$
telle que
$\Spec(M) \to \Spec(K) \to \mathcal{X} \to \mathcal{Y}$ et
$\Spec(M) \to \Spec(L) \to \mathcal{Z} \to \mathcal{Y}$
soient $2$-isomorphes. Et c'est exactement la condition qui donne un
morphisme $\Spec(M) \to \mathcal{Z} \times_\mathcal{Y} \mathcal{X}$.
\end{proof}

\begin{lemma}
\label{stacks-properties-lemma-characterize-surjective}
Soit $f : \mathcal{X} \to \mathcal{Y}$ un morphisme de champs algébriques
représentable par des espaces algébriques. Les assertions suivantes sont équivalentes :
\begin{enumerate}
\item $|f| : |\mathcal{X}| \to |\mathcal{Y}|$ est surjective, et
\item $f$ est surjectif
(au sens de la section \ref{stacks-properties-section-properties-morphisms}).
\end{enumerate}
\end{lemma}

\begin{proof}
Supposons (1). Soit $T \to \mathcal{Y}$ un morphisme dont la source est un schéma.
Pour démontrer (2), nous devons montrer que le morphisme d'espaces algébriques
$T \times_\mathcal{Y} \mathcal{X} \to T$ est surjectif. D'après
Morphismes d'espaces, Définition \ref{spaces-morphisms-definition-surjective},
cela signifie que nous devons montrer que
$|T \times_\mathcal{Y} \mathcal{X}| \to |T|$ est surjective.
En appliquant le
Lemme \ref{stacks-properties-lemma-points-cartesian},
nous voyons que cela résulte de (1).

\medskip\noindent
Réciproquement, supposons (2). Soit $y : \Spec(K) \to \mathcal{Y}$ un
morphisme du spectre d'un corps vers $\mathcal{Y}$. Par hypothèse, le
morphisme
$\Spec(K) \times_{y, \mathcal{Y}} \mathcal{X} \to \Spec(K)$
d'espaces algébriques est surjectif. D'après
Morphismes d'espaces, Définition \ref{spaces-morphisms-definition-surjective},
cela signifie qu'il existe une extension de corps
$K'/K$ et un morphisme
$\Spec(K') \to \Spec(K) \times_{y, \mathcal{Y}} \mathcal{X}$
tels que le carré de gauche du diagramme
$$
\xymatrix{
\Spec(K') \ar[r] \ar[d] &
\Spec(K) \times_{y, \mathcal{Y}} \mathcal{X} \ar[d] \ar[r] &
\mathcal{X} \ar[d]
\\
\Spec(K) \ar@{=}[r] &
\Spec(K) \ar[r]^-y &
\mathcal{Y}
}
$$
soit commutatif. Cela montre que $|X| \to |\mathcal{Y}|$ est surjective.
\end{proof}

\noindent
Voici un lemme qui explique comment calculer l'ensemble des points à partir
d'une présentation.

\begin{lemma}
\label{stacks-properties-lemma-points-presentation}
Soit $\mathcal{X}$ un champ algébrique.
Soit $\mathcal{X} = [U/R]$ une présentation de $\mathcal{X}$, voir
Champs algébriques, Définition \ref{algebraic-definition-presentation}.
Alors l'image de $|R| \to |U| \times |U|$ est une relation d'équivalence
et $|\mathcal{X}|$ est le quotient de $|U|$ par cette relation d'équivalence.
\end{lemma}

\begin{proof}
L'hypothèse signifie que nous avons un groupoïde lisse $(U, R, s, t, c)$
en espaces algébriques, et une équivalence $f : [U/R] \to \mathcal{X}$.
Nous pouvons supposer que $\mathcal{X} = [U/R]$.
Le morphisme induit $p : U \to \mathcal{X}$ est lisse et surjectif, voir
Champs algébriques,
Lemme \ref{algebraic-lemma-smooth-quotient-smooth-presentation}.
Par conséquent, $|U| \to |\mathcal{X}|$ est surjective d'après le
Lemme \ref{stacks-properties-lemma-characterize-surjective}.
Remarquons que $R = U \times_\mathcal{X} U$, voir
Groupoïdes en espaces,
Lemme \ref{spaces-groupoids-lemma-quotient-stack-2-cartesian}.
Le
Lemme \ref{stacks-properties-lemma-points-cartesian}
implique donc que l'application
$$
|R| \longrightarrow |U| \times_{|\mathcal{X}|} |U|
$$
est surjective. L'image de $|R| \to |U| \times |U|$ est donc
exactement l'ensemble des couples $(u_1, u_2) \in |U| \times |U|$
tels que $u_1$ et $u_2$ aient la même image dans $|\mathcal{X}|$.
En combinant ces deux assertions, nous obtenons le résultat du lemme.
\end{proof}

\begin{remark}
\label{stacks-properties-remark-more-general-presentation}
Le résultat du
Lemme \ref{stacks-properties-lemma-points-presentation}
se généralise comme suit.
Soit $\mathcal{X}$ un champ algébrique.
Soit $U$ un espace algébrique et soit $f : U \to \mathcal{X}$ un morphisme surjectif
(ce qui a un sens d'après la
section \ref{stacks-properties-section-properties-morphisms}).
Soit $R = U \times_\mathcal{X} U$, soit $(U, R, s, t, c)$ le groupoïde
en espaces algébriques et soit $f_{can} : [U/R] \to \mathcal{X}$ le
morphisme canonique construit dans
Champs algébriques, Lemme \ref{algebraic-lemma-map-space-into-stack}.
Alors l'image de $|R| \to |U| \times |U|$ est une relation d'équivalence
et $|\mathcal{X}| = |U|/|R|$. La démonstration du
Lemme \ref{stacks-properties-lemma-points-presentation}
s'applique sans changement. (Bien entendu, en général, $[U/R]$ n'est pas un champ
algébrique et, en général, $f_{can}$ n'est pas un isomorphisme.)
\end{remark}

\begin{lemma}
\label{stacks-properties-lemma-topology-points}
Il existe une unique topologie sur les ensembles de points
des champs algébriques qui possède les propriétés suivantes :
\begin{enumerate}
\item pour tout morphisme de champs algébriques $\mathcal{X} \to \mathcal{Y}$,
l'application $|\mathcal{X}| \to |\mathcal{Y}|$ est continue, et
\item pour tout morphisme $U \to \mathcal{X}$ qui est plat et localement
de présentation finie, où $U$ est un espace algébrique,
l'application d'espaces topologiques $|U| \to |\mathcal{X}|$ est continue et ouverte.
\end{enumerate}
\end{lemma}

\begin{proof}
Choisissons un morphisme $p : U \to \mathcal{X}$ qui soit
surjectif, plat et localement de présentation finie,
où $U$ est un espace algébrique. Un tel morphisme existe par définition d'un champ
algébrique, puisqu'un morphisme lisse est plat et localement de présentation finie
(voir
Morphismes d'espaces,
Lemmes \ref{spaces-morphisms-lemma-smooth-locally-finite-presentation} et
\ref{spaces-morphisms-lemma-smooth-flat}).
Nous définissons une topologie sur $|\mathcal{X}|$ par la règle suivante :
$W \subset |\mathcal{X}|$ est ouvert si et seulement si $|p|^{-1}(W)$ est ouvert
dans $|U|$. Pour montrer que cette définition est indépendante du choix de $p$, soit
$p' : U' \to \mathcal{X}$ un autre morphisme surjectif, plat et
localement de présentation finie d'un espace algébrique vers
$\mathcal{X}$. Posons $U'' = U \times_\mathcal{X} U'$,
de sorte que nous avons un diagramme $2$-commutatif
$$
\xymatrix{
U'' \ar[r] \ar[d] & U' \ar[d] \\
U \ar[r] & \mathcal{X}
}
$$
Puisque $U \to \mathcal{X}$ et $U' \to \mathcal{X}$ sont surjectifs, plats et
localement de présentation finie, nous voyons que $U'' \to U'$ et $U'' \to U$
sont surjectifs, plats et localement de présentation finie, voir
Lemme \ref{stacks-properties-lemma-property-spaces-too}.
Par conséquent, les applications $|U''| \to |U'|$ et $|U''| \to |U|$ sont continues, ouvertes
et surjectives, voir
Morphismes d'espaces,
Définition \ref{spaces-morphisms-definition-surjective} et
Lemme \ref{spaces-morphisms-lemma-fppf-open}.
Cela implique clairement que notre définition est indépendante du choix
de $p : U \to \mathcal{X}$.

\medskip\noindent
Soit $f : \mathcal{X} \to \mathcal{Y}$ un morphisme de champs algébriques.
D'après
Champs algébriques, Lemme \ref{algebraic-lemma-lift-morphism-presentations},
nous pouvons trouver un diagramme $2$-commutatif
$$
\xymatrix{
U \ar[d]_x \ar[r]_a & V \ar[d]^y \\
\mathcal{X} \ar[r]^f & \mathcal{Y}
}
$$
dont les flèches verticales sont lisses et surjectives.
Considérons le diagramme commutatif associé
$$
\xymatrix{
|U| \ar[d]_{|x|} \ar[r]_{|a|} & |V| \ar[d]^{|y|} \\
|\mathcal{X}| \ar[r]^{|f|} & |\mathcal{Y}|
}
$$
d'ensembles. Si $W \subset |\mathcal{Y}|$ est ouvert, alors, par la définition
ci-dessus, cela signifie exactement que $|y|^{-1}(W)$ est ouvert dans $|V|$. Puisque
$|a|$ est continue, nous en concluons que
$|a|^{-1}|y|^{-1}(W) = |x|^{-1}|f|^{-1}(W)$ est ouvert dans $|W|$, ce qui signifie,
par définition, que $|f|^{-1}(W)$ est ouvert dans $|\mathcal{X}|$.
Ainsi, $|f|$ est continue.

\medskip\noindent
Enfin, nous devons montrer que si $U$ est un espace algébrique et si
$U \to \mathcal{X}$ est plat et localement de présentation finie, alors
$|U| \to |\mathcal{X}|$ est ouverte. Soit $V \to \mathcal{X}$ un morphisme surjectif,
plat et localement de présentation finie, où $V$ est un espace algébrique.
Considérons le diagramme commutatif
$$
\xymatrix{
|U \times_\mathcal{X} V| \ar[r]_e \ar[rd]_f &
|U| \times_{|\mathcal{X}|} |V| \ar[d]_c \ar[r]_d &
|V| \ar[d]^b \\
& |U| \ar[r]^a & |\mathcal{X}|
}
$$
Le morphisme $U \times_\mathcal{X} V \to U$ est surjectif, c'est-à-dire que
$f : |U \times_\mathcal{X} V| \to |U|$ est surjective.
La flèche horizontale supérieure de gauche est surjective, voir
Lemme \ref{stacks-properties-lemma-points-cartesian}.
Le morphisme $U \times_\mathcal{X} V \to V$ est plat et localement de présentation
finie; $d \circ e : |U \times_\mathcal{X} V| \to |V|$ est donc ouverte,
voir
Morphismes d'espaces, Lemme \ref{spaces-morphisms-lemma-fppf-open}.
Choisissons un ouvert $W \subset |U|$. Les propriétés ci-dessus impliquent que
$b^{-1}(a(W)) = (d \circ e)(f^{-1}(W))$ est ouvert, ce qui, par construction, signifie
que $a(W)$ est ouvert, comme souhaité.
\end{proof}

\begin{definition}
\label{stacks-properties-definition-topological-space}
Soit $\mathcal{X}$ un champ algébrique.
L'{\it espace topologique} sous-jacent à $\mathcal{X}$ est l'ensemble des points
$|\mathcal{X}|$ muni de la topologie construite dans le
Lemme \ref{stacks-properties-lemma-topology-points}.
\end{definition}

\noindent
Cette définition n'entre pas en conflit avec la topologie déjà existante
sur $|\mathcal{X}|$ lorsque $\mathcal{X}$ est un espace algébrique.

\begin{lemma}
\label{stacks-properties-lemma-space-locally-quasi-compact}
Soit $\mathcal{X}$ un champ algébrique.
Tout point de $|\mathcal{X}|$ possède un système fondamental de
voisinages ouverts quasi-compacts.
En particulier, $|\mathcal{X}|$ est localement quasi-compact au sens de
Topologie, Définition \ref{topology-definition-locally-quasi-compact}.
\end{lemma}

\begin{proof}
Cela résulte formellement du fait qu'il existe un schéma
$U$ et une application surjective, ouverte et continue $U \to |\mathcal{X}|$
d'espaces topologiques. En effet, si $U \to \mathcal{X}$ est surjectif et
lisse, alors le
Lemme \ref{stacks-properties-lemma-topology-points}
garantit que $|U| \to |\mathcal{X}|$ est continue, surjective et ouverte.
\end{proof}











\section{Morphismes surjectifs}
\label{stacks-properties-section-surjective}

\noindent
Soit $f : \mathcal{X} \to \mathcal{Y}$ un morphisme de champs algébriques
représentable par des espaces algébriques. Dans la
section \ref{stacks-properties-section-properties-morphisms},
nous avons déjà défini ce que signifie pour $f$ d'être surjectif. Dans le
Lemme \ref{stacks-properties-lemma-characterize-surjective},
nous avons vu que cela équivaut à demander que
$|f| : |\mathcal{X}| \to |\mathcal{Y}|$ soit surjective.
Cela ouvre la voie à la définition suivante.

\begin{definition}
\label{stacks-properties-definition-surjective}
Soit $f : \mathcal{X} \to \mathcal{Y}$ un morphisme de champs algébriques.
Nous disons que $f$ est {\it surjectif} si l'application
$|f| : |\mathcal{X}| \to |\mathcal{Y}|$ des espaces topologiques associés
est surjective.
\end{definition}

\noindent
Voici quelques lemmes.

\begin{lemma}
\label{stacks-properties-lemma-composition-surjective}
Le composé de morphismes surjectifs est surjectif.
\end{lemma}

\begin{proof}
Omis.
\end{proof}

\begin{lemma}
\label{stacks-properties-lemma-base-change-surjective}
Tout changement de base d'un morphisme surjectif est surjectif.
\end{lemma}

\begin{proof}
Omis. Indication : utiliser le
Lemme \ref{stacks-properties-lemma-points-cartesian}.
\end{proof}

\begin{lemma}
\label{stacks-properties-lemma-descent-surjective}
Soit $f : \mathcal{X} \to \mathcal{Y}$ un morphisme de champs algébriques.
Soit $\mathcal{Y}' \to \mathcal{Y}$ un morphisme surjectif de champs
algébriques. Si le changement de base $f' : \mathcal{Y}' \times_\mathcal{Y} \mathcal{X}
\to \mathcal{Y}'$ de $f$ est surjectif, alors $f$ est surjectif.
\end{lemma}

\begin{proof}
Cela résulte immédiatement du
Lemme \ref{stacks-properties-lemma-points-cartesian}.
\end{proof}

\begin{lemma}
\label{stacks-properties-lemma-surjective-permanence}
Soient $\mathcal{X} \to \mathcal{Y} \to \mathcal{Z}$ des morphismes de
champs algébriques. Si $\mathcal{X} \to \mathcal{Z}$ est surjectif,
alors $\mathcal{Y} \to \mathcal{Z}$ l'est aussi.
\end{lemma}

\begin{proof}
Immédiat.
\end{proof}












\section{Champs algébriques quasi-compacts}
\label{stacks-properties-section-quasi-compact}

\noindent
La définition suivante équivaut à la définition pour les espaces
algébriques d'après
Propriétés des espaces, Lemme \ref{spaces-properties-lemma-quasi-compact-space}.

\begin{definition}
\label{stacks-properties-definition-quasi-compact}
Soit $\mathcal{X}$ un champ algébrique.
Nous disons que $\mathcal{X}$ est {\it quasi-compact}
si et seulement si $|\mathcal{X}|$ est quasi-compact.
\end{definition}

\begin{lemma}
\label{stacks-properties-lemma-quasi-compact-stack}
Soit $\mathcal{X}$ un champ algébrique.
Les assertions suivantes sont équivalentes :
\begin{enumerate}
\item $\mathcal{X}$ est quasi-compact,
\item il existe un morphisme lisse surjectif $U \to \mathcal{X}$,
où $U$ est un schéma affine,
\item il existe un morphisme lisse surjectif $U \to \mathcal{X}$,
où $U$ est un schéma quasi-compact,
\item il existe un morphisme lisse surjectif $U \to \mathcal{X}$,
où $U$ est un espace algébrique quasi-compact, et
\item il existe un morphisme surjectif $\mathcal{U} \to \mathcal{X}$
de champs algébriques tel que $\mathcal{U}$ soit quasi-compact.
\end{enumerate}
\end{lemma}

\begin{proof}
Nous utiliserons le
Lemme \ref{stacks-properties-lemma-characterize-surjective}.
Supposons que $\mathcal{U}$ et $\mathcal{U} \to \mathcal{X}$ soient comme en (5).
Puisque $|\mathcal{U}| \to |\mathcal{X}|$ est surjective et
continue, nous en concluons que $|\mathcal{X}|$ est quasi-compact.
Ainsi, (5) implique (1). Les implications
(2) $\Rightarrow$ (3) $\Rightarrow$ (4) $\Rightarrow$ (5)
sont immédiates. Supposons (1), c'est-à-dire que $\mathcal{X}$ est quasi-compact, c'est-à-dire que
$|\mathcal{X}|$ est quasi-compact. Choisissons un schéma $U$ et un morphisme lisse
surjectif $U \to \mathcal{X}$. Puisque $|U| \to |\mathcal{X}|$
est ouverte, nous voyons qu'il existe un ouvert quasi-compact $U' \subset U$
tel que $|U'| \to |X|$ soit surjective (et toujours lisse).
Choisissons un recouvrement ouvert affine fini $U' = U_1 \cup \ldots \cup U_n$.
Alors $U_1 \amalg \ldots \amalg U_n \to \mathcal{X}$
est un morphisme lisse surjectif dont la source est un
schéma affine (Schémas, Lemme \ref{schemes-lemma-disjoint-union-affines}).
Ainsi, (2) est satisfaite.
\end{proof}

\begin{lemma}
\label{stacks-properties-lemma-finite-disjoint-quasi-compact}
Une réunion disjointe finie de champs algébriques quasi-compacts est
un champ algébrique quasi-compact.
\end{lemma}

\begin{proof}
Cela résulte clairement du fait topologique correspondant.
\end{proof}


\section{Propriétés des champs algébriques définies par des propriétés des schémas}
\label{stacks-properties-section-types-properties}

\noindent
Toute propriété des schémas locale pour la topologie lisse donne naissance à une
propriété correspondante des champs algébriques grâce au lemme suivant. Remarquons qu'une
propriété des schémas qui est locale pour la topologie lisse est aussi locale pour la topologie \'etale,
puisque tout recouvrement \'etale est aussi un recouvrement lisse. Ainsi, pour une propriété $P$
des schémas locale pour la topologie lisse, nous savons ce que signifie dire qu'un
espace algébrique possède $P$, voir
Propriétés des espaces, section \ref{spaces-properties-section-types-properties}.

\begin{lemma}
\label{stacks-properties-lemma-type-property}
Soit $\mathcal{P}$ une propriété des schémas qui est locale pour la topologie
lisse, voir
Descente, Définition \ref{descent-definition-property-local}.
Soit $\mathcal{X}$ un champ algébrique. Les assertions suivantes sont équivalentes :
\begin{enumerate}
\item il existe un schéma $U$ et un morphisme lisse surjectif
$U \to \mathcal{X}$ tels que le schéma $U$ possède la propriété $\mathcal{P}$,
\item pour tout schéma $U$ et tout morphisme lisse $U \to \mathcal{X}$,
le schéma $U$ possède la propriété $\mathcal{P}$,
\item il existe un espace algébrique $U$ et un morphisme lisse surjectif
$U \to \mathcal{X}$ tels que l'espace algébrique $U$ possède la propriété $\mathcal{P}$, et
\item pour tout espace algébrique $U$ et tout morphisme lisse
$U \to \mathcal{X}$, l'espace algébrique $U$ possède la propriété $\mathcal{P}$.
\end{enumerate}
Si $\mathcal{X}$ est un schéma $U$, cela équivaut à $\mathcal{P}(U)$.
Si $\mathcal{X}$ est un espace algébrique $X$, cela équivaut à ce que
$X$ possède la propriété $\mathcal{P}$.
\end{lemma}

\begin{proof}
Soit $U \to \mathcal{X}$ un morphisme surjectif et lisse, où $U$ est un espace algébrique.
Soit $V \to \mathcal{X}$ un morphisme lisse, où $V$ est un espace algébrique.
Choisissons des schémas $U'$ et $V'$ et des morphismes \'etales surjectifs
$U' \to U$ et $V' \to V$. Enfin, choisissons un schéma $W$ et un
morphisme \'etale surjectif $W \to V' \times_\mathcal{X} U'$.
Alors $W \to V'$ et $W \to U'$ sont des morphismes lisses de schémas,
car ce sont des composés de morphismes \'etales et lisses d'espaces algébriques, voir
Morphismes d'espaces, Lemmes \ref{spaces-morphisms-lemma-etale-smooth} et
\ref{spaces-morphisms-lemma-composition-smooth}.
De plus, $W \to V'$ est surjectif puisque $U' \to \mathcal{X}$ est surjectif.
Nous avons donc
$$
\mathcal{P}(U) \Leftrightarrow
\mathcal{P}(U') \Rightarrow
\mathcal{P}(W) \Rightarrow
\mathcal{P}(V') \Leftrightarrow \mathcal{P}(V)
$$
où les équivalences résultent de la définition de la propriété $\mathcal{P}$ pour les
espaces algébriques, et les deux implications résultent de
Descente, Définition \ref{descent-definition-property-local}.
Cela démontre (3) $\Rightarrow$ (4).

\medskip\noindent
Les implications (2) $\Rightarrow$ (1), (1) $\Rightarrow$ (3),
et (4) $\Rightarrow$ (2) sont immédiates.
\end{proof}

\begin{definition}
\label{stacks-properties-definition-type-property}
Soit $\mathcal{X}$ un champ algébrique.
Soit $\mathcal{P}$ une propriété des schémas qui est
locale pour la topologie lisse.
Nous disons que $\mathcal{X}$ {\it possède la propriété $\mathcal{P}$}
si l'une quelconque des conditions équivalentes du
Lemme \ref{stacks-properties-lemma-type-property}
est satisfaite.
\end{definition}

\begin{remark}
\label{stacks-properties-remark-list-properties-local-smooth-topology}
Voici une liste de propriétés qui sont locales pour la topologie lisse
(gardons à l'esprit que les topologies fpqc, fppf et syntomique sont
plus fines que la topologie lisse) :
\begin{enumerate}
\item localement noethérien, voir
Descente, Lemme \ref{descent-lemma-Noetherian-local-fppf},
\item de Jacobson, voir
Descente, Lemme \ref{descent-lemma-Jacobson-local-fppf},
\item localement noethérien et $(S_k)$, voir
Descente, Lemme \ref{descent-lemma-Sk-local-syntomic},
\item Cohen-Macaulay, voir
Descente, Lemme \ref{descent-lemma-CM-local-syntomic},
\item réduit, voir
Descente, Lemme \ref{descent-lemma-reduced-local-smooth},
\item normal, voir
Descente, Lemme \ref{descent-lemma-normal-local-smooth},
\item localement noethérien et $(R_k)$, voir
Descente, Lemme \ref{descent-lemma-Rk-local-smooth},
\item régulier, voir
Descente, Lemme \ref{descent-lemma-regular-local-smooth},
\item de Nagata, voir
Descente, Lemme \ref{descent-lemma-Nagata-local-smooth}.
\end{enumerate}
\end{remark}

\noindent
Toute propriété des germes de schémas locale pour la topologie lisse donne naissance à une
propriété correspondante des champs algébriques. Remarquons qu'une propriété des germes locale pour la topologie lisse
est aussi locale pour la topologie \'etale. Ainsi, pour une propriété $P$ des germes de
schémas locale pour la topologie lisse, nous savons ce que signifie dire qu'un espace algébrique $X$ possède
la propriété $P$ en $x \in |X|$, voir
Propriétés des espaces, section \ref{stacks-properties-section-types-properties}.

\begin{lemma}
\label{stacks-properties-lemma-local-source-target-at-point}
Soit $\mathcal{X}$ un champ algébrique.
Soit $x \in |\mathcal{X}|$ un point de $\mathcal{X}$.
Soit $\mathcal{P}$ une propriété des germes de schémas locale pour la topologie lisse, voir
Descente, Définition \ref{descent-definition-local-at-point}.
Les assertions suivantes sont équivalentes :
\begin{enumerate}
\item pour tout morphisme lisse $U \to \mathcal{X}$, où $U$ est un schéma,
et tout $u \in U$ tel que $a(u) = x$, on a $\mathcal{P}(U, u)$,
\item il existe un morphisme lisse $U \to \mathcal{X}$, où $U$ est un schéma,
et un $u \in U$ tel que $a(u) = x$, pour lesquels on a $\mathcal{P}(U, u)$,
\item pour tout morphisme lisse $U \to \mathcal{X}$, où $U$ est un espace algébrique,
et tout $u \in |U|$ tel que $a(u) = x$, l'espace algébrique $U$ possède la propriété
$\mathcal{P}$ en $u$, et
\item il existe un morphisme lisse $U \to \mathcal{X}$, où $U$ est un
espace algébrique, et un $u \in |U|$ tel que $a(u) = x$, tels que l'espace algébrique
$U$ possède la propriété $\mathcal{P}$ en $u$.
\end{enumerate}
Si $\mathcal{X}$ est représentable, cela équivaut à
$\mathcal{P}(\mathcal{X}, x)$. Si $\mathcal{X}$ est un espace algébrique, alors
cela équivaut à ce que $\mathcal{X}$ possède la propriété $\mathcal{P}$ en $x$.
\end{lemma}

\begin{proof}
Soient $a : U \to \mathcal{X}$ et $u \in |U|$ comme en (3). Soit
$b : V \to \mathcal{X}$ un autre morphisme lisse, où $V$ est un espace
algébrique, et soit $v \in |V|$ tel que $b(v) = x$ également.
Choisissons un schéma $U'$, un morphisme \'etale $U' \to U$ et $u' \in U'$
d'image $u$. Choisissons un schéma $V'$, un morphisme \'etale $V' \to V$
et $v' \in V'$ d'image $v$. D'après le
Lemme \ref{stacks-properties-lemma-points-cartesian},
il existe un point $\overline{w} \in |V' \times_\mathcal{X} U'|$
d'images $u'$ et $v'$. Choisissons un schéma $W$ et un morphisme \'etale surjectif
$W \to V' \times_\mathcal{X} U'$. Nous pouvons choisir un
$w \in |W|$ d'image $\overline{w}$ (voir
Propriétés des espaces,
Lemme \ref{spaces-properties-lemma-characterize-surjective}).
Alors $W \to V'$ et $W \to U'$ sont des morphismes lisses de schémas,
car ce sont des composés de morphismes \'etales et lisses d'espaces algébriques, voir
Morphismes d'espaces, Lemmes \ref{spaces-morphisms-lemma-etale-smooth} et
\ref{spaces-morphisms-lemma-composition-smooth}.
Par conséquent,
$$
\mathcal{P}(U, u)
\Leftrightarrow
\mathcal{P}(U', u')
\Leftrightarrow
\mathcal{P}(W, w)
\Leftrightarrow
\mathcal{P}(V', v')
\Leftrightarrow
\mathcal{P}(V, v)
$$
Les deux équivalences extérieures résultent de
Propriétés des espaces,
Définition \ref{spaces-properties-definition-property-at-point},
et les deux autres de la définition d'une propriété locale pour la topologie lisse
des germes de schémas. Cela démontre (4) $\Rightarrow$ (3).

\medskip\noindent
Les implications (1) $\Rightarrow$ (2), (2) $\Rightarrow$ (4),
et (3) $\Rightarrow$ (1) sont immédiates.
\end{proof}

\begin{definition}
\label{stacks-properties-definition-property-at-point}
Soit $\mathcal{P}$ une propriété des germes de schémas qui est
locale pour la topologie lisse. Soit $\mathcal{X}$ un champ algébrique.
Soit $x \in |\mathcal{X}|$.
Nous disons que $\mathcal{X}$ {\it possède la propriété $\mathcal{P}$ en $x$}
si l'une quelconque des conditions équivalentes du
Lemme \ref{stacks-properties-lemma-local-source-target-at-point}
est satisfaite.
\end{definition}










\section{Monomorphismes de champs algébriques}
\label{stacks-properties-section-monomorphisms}

\noindent
Nous définissons comme suit un monomorphisme de champs algébriques.
Nous verrons dans le
Lemme \ref{stacks-properties-lemma-monomorphism}
que cela est compatible avec la notion $2$-catégorique correspondante.

\begin{definition}
\label{stacks-properties-definition-monomorphism}
Soit $f : \mathcal{X} \to \mathcal{Y}$ un morphisme de champs algébriques.
Nous disons que $f$ est un {\it monomorphisme}
s'il est représentable par des espaces algébriques et est un monomorphisme au sens de la
section \ref{stacks-properties-section-properties-morphisms}.
\end{definition}

\noindent
Commençons par quelques lemmes élémentaires.

\begin{lemma}
\label{stacks-properties-lemma-base-change-monomorphism}
Soit $\mathcal{X} \to \mathcal{Y}$ un morphisme de champs algébriques.
Soit $\mathcal{Z} \to \mathcal{Y}$ un monomorphisme.
Alors $\mathcal{Z} \times_\mathcal{Y} \mathcal{X} \to \mathcal{X}$
est un monomorphisme.
\end{lemma}

\begin{proof}
Cela résulte de la discussion générale de la
section \ref{stacks-properties-section-properties-morphisms}.
\end{proof}

\begin{lemma}
\label{stacks-properties-lemma-composition-monomorphism}
Les composés de monomorphismes de champs algébriques sont des monomorphismes.
\end{lemma}

\begin{proof}
Cela résulte de la discussion générale de la
section \ref{stacks-properties-section-properties-morphisms}
et de
Morphismes d'espaces,
Lemme \ref{spaces-morphisms-lemma-composition-monomorphism}.
\end{proof}

\begin{lemma}
\label{stacks-properties-lemma-monomorphism}
Soit $f : \mathcal{X} \to \mathcal{Y}$ un morphisme de champs algébriques.
Les assertions suivantes sont équivalentes :
\begin{enumerate}
\item $f$ est un monomorphisme,
\item $f$ est pleinement fidèle,
\item la diagonale
$\Delta_f : \mathcal{X} \to \mathcal{X} \times_\mathcal{Y} \mathcal{X}$
est une équivalence, et
\item il existe un espace algébrique $W$ et un morphisme surjectif et plat
$W \to \mathcal{Y}$, localement de présentation finie, tels que
$V = \mathcal{X} \times_\mathcal{Y} W$ soit un espace algébrique et que le
morphisme $V \to W$ soit un monomorphisme d'espaces algébriques.
\end{enumerate}
\end{lemma}

\begin{proof}
L'équivalence de (1) et (4) résulte de la discussion générale de la
section \ref{stacks-properties-section-properties-morphisms},
et en particulier des
Lemmes \ref{stacks-properties-lemma-check-representable-covering} et
\ref{stacks-properties-lemma-check-property-covering}.

\medskip\noindent
L'équivalence de (2) et (3) est
Catégories, Lemme \ref{categories-lemma-fully-faithful-diagonal-equivalence}.

\medskip\noindent
Supposons satisfaites les conditions équivalentes (2) et (3). Alors $f$ est représentable
par des espaces algébriques d'après
Champs algébriques,
Lemme \ref{algebraic-lemma-characterize-representable-by-algebraic-spaces}.
En outre, le lemme de $2$-Yoneda, combiné à la pleine fidélité,
implique que, pour tout schéma $T$, le foncteur
$$
\Mor(T, \mathcal{X})
\longrightarrow
\Mor(T, \mathcal{Y})
$$
est pleinement fidèle. Étant donné un morphisme $y : T \to \mathcal{Y}$, il existe donc,
à un unique $2$-isomorphisme près, au plus un morphisme $x : T \to \mathcal{X}$,
tel que $y \cong f \circ x$. En particulier, étant donné un morphisme de schémas
$h : T' \to T$, il existe au plus un relèvement
$\tilde h : T' \to T \times_\mathcal{Y} \mathcal{X}$ de $h$.
Ainsi, $T \times_\mathcal{Y} \mathcal{X} \to T$ est un monomorphisme
d'espaces algébriques, ce qui démontre (1).

\medskip\noindent
Enfin, supposons (1). Alors, pour tout schéma $T$ et tout morphisme
$y : T \to \mathcal{Y}$, le produit fibré $T \times_\mathcal{Y} \mathcal{X}$
est un espace algébrique, et $T \times_\mathcal{Y} \mathcal{X} \to T$
est un monomorphisme. Il existe donc exactement un couple, à isomorphisme unique près,
$(x, \alpha)$, où $x : T \to \mathcal{X}$ est un morphisme
et $\alpha : f \circ x \to y$ est un $2$-morphisme. En appliquant
le lemme de $2$-Yoneda, cela signifie exactement que $f$ est pleinement fidèle, c'est-à-dire
que (2) est satisfaite.
\end{proof}

\begin{lemma}
\label{stacks-properties-lemma-monomorphism-injective-points}
\begin{slogan}
Les monomorphismes de champs sont injectifs sur les points.
\end{slogan}
Un monomorphisme de champs algébriques induit une application injective entre les
ensembles de points.
\end{lemma}

\begin{proof}
Soit $f : \mathcal{X} \to \mathcal{Y}$ un monomorphisme de champs algébriques.
Supposons que $x_i : \Spec(K_i) \to \mathcal{X}$ soient des morphismes tels que
$f \circ x_1$ et $f \circ x_2$ définissent le même élément de $|\mathcal{Y}|$.
En appliquant la définition, nous trouvons une extension commune $\Omega$ munie des
morphismes correspondants $c_i : \Spec(\Omega) \to \Spec(K_i)$ et d'un
$2$-isomorphisme $\beta : f \circ x_1 \circ c_1 \to f \circ x_1 \circ c_2$.
Comme $f$ est pleinement fidèle, voir
Lemme \ref{stacks-properties-lemma-monomorphism},
nous pouvons relever $\beta$ en un isomorphisme
$\alpha : x_1 \circ c_1 \to x_1 \circ c_2$.
Ainsi, $x_1$ et $x_2$ définissent le même point de $|\mathcal{X}|$,
comme souhaité.
\end{proof}

\begin{lemma}
\label{stacks-properties-lemma-monomorphism-diagonal}
Soient $\mathcal{X} \to \mathcal{X}' \to \mathcal{Y}$ des morphismes
de champs algébriques. Si $\mathcal{X} \to \mathcal{X}'$ est un monomorphisme,
alors le diagramme canonique
$$
\xymatrix{
\mathcal{X} \ar[r] \ar[d] &
\mathcal{X} \times_\mathcal{Y} \mathcal{X} \ar[d] \\
\mathcal{X}' \ar[r] &
\mathcal{X}' \times_\mathcal{Y} \mathcal{X}'
}
$$
est un carré cartésien.
\end{lemma}

\begin{proof}
Nous avons $\mathcal{X} = \mathcal{X} \times_{\mathcal{X}'} \mathcal{X}$
d'après le Lemme \ref{stacks-properties-lemma-monomorphism}. Le résultat en découle alors en appliquant
Catégories, Lemme \ref{categories-lemma-fibre-product-after-map}.
\end{proof}







\section{Immersions de champs algébriques}
\label{stacks-properties-section-immersions}

\noindent
Les immersions de champs algébriques sont définies comme suit.

\begin{definition}
\label{stacks-properties-definition-immersion}
Immersions.
\begin{enumerate}
\item Un morphisme de champs algébriques est appelé une {\it immersion ouverte}
s'il est représentable et est une immersion ouverte
au sens de la
section \ref{stacks-properties-section-properties-morphisms}.
\item Un morphisme de champs algébriques est appelé une {\it immersion fermée}
s'il est représentable et est une immersion fermée
au sens de la
section \ref{stacks-properties-section-properties-morphisms}.
\item Un morphisme de champs algébriques est appelé une {\it immersion}
s'il est représentable et est une immersion
au sens de la
section \ref{stacks-properties-section-properties-morphisms}.
\end{enumerate}
\end{definition}

\noindent
Ce n'est pas pour nous la manière la plus commode de concevoir les immersions.
Il nous est un peu plus commode de considérer une immersion
comme un morphisme de champs algébriques représentable par des espaces
algébriques et qui est une immersion au sens de la
section \ref{stacks-properties-section-properties-morphisms}.
Il en va de même pour les immersions fermées et ouvertes.
Comme cela équivaut clairement à la notion qui vient d'être définie, nous
utiliserons cette caractérisation sans autre mention.
Nous démontrons quelques lemmes simples sur cette notion.

\begin{lemma}
\label{stacks-properties-lemma-base-change-immersion}
Soit $\mathcal{X} \to \mathcal{Y}$ un morphisme de champs algébriques.
Soit $\mathcal{Z} \to \mathcal{Y}$ une immersion
(fermée, resp.\ ouverte).
Alors $\mathcal{Z} \times_\mathcal{Y} \mathcal{X} \to \mathcal{X}$
est une immersion (fermée, resp.\ ouverte).
\end{lemma}

\begin{proof}
Cela résulte de la discussion générale de la
section \ref{stacks-properties-section-properties-morphisms}.
\end{proof}

\begin{lemma}
\label{stacks-properties-lemma-composition-immersion}
Les composés d'immersions de champs algébriques sont des immersions.
Il en va de même des immersions fermées et des immersions ouvertes.
\end{lemma}

\begin{proof}
Cela résulte de la discussion générale de la
section \ref{stacks-properties-section-properties-morphisms}
et de
Espaces, Lemme \ref{spaces-lemma-composition-immersions}.
\end{proof}

\begin{lemma}
\label{stacks-properties-lemma-check-immersion-covering}
Soit $f : \mathcal{X} \to \mathcal{Y}$ un morphisme de champs algébriques.
Soit $W$ un espace algébrique et soit $W \to \mathcal{Y}$ un morphisme surjectif
et plat, localement de présentation finie. Les assertions suivantes
sont équivalentes :
\begin{enumerate}
\item $f$ est une immersion (ouverte, resp.\ fermée), et
\item $V = W \times_\mathcal{Y} \mathcal{X}$ est un espace algébrique, et
$V \to W$ est une immersion (ouverte, resp.\ fermée).
\end{enumerate}
\end{lemma}

\begin{proof}
Cela résulte de la discussion générale de la
section \ref{stacks-properties-section-properties-morphisms}
et en particulier des
Lemmes \ref{stacks-properties-lemma-check-representable-covering} et
\ref{stacks-properties-lemma-check-property-covering}.
\end{proof}

\begin{lemma}
\label{stacks-properties-lemma-immersion-monomorphism}
Une immersion est un monomorphisme.
\end{lemma}

\begin{proof}
Voir
Morphismes d'espaces,
Lemme \ref{spaces-morphisms-lemma-immersions-monomorphisms}.
\end{proof}

\begin{lemma}
\label{stacks-properties-lemma-subspace-induced-topology}
Si $f : \mathcal{X} \to \mathcal{Y}$ est une immersion, alors
$|f| : |\mathcal{X}| \to |\mathcal{Y}|$ est un homéomorphisme
sur une partie localement fermée. Si $f$ est une immersion fermée, resp.\ ouverte,
alors $|f|$ est fermée, resp.\ ouverte.
\end{lemma}

\begin{proof}
Omis.
\end{proof}

\noindent
Les deux lemmes suivants expliquent comment considérer les
immersions à l'aide de présentations.

\begin{lemma}
\label{stacks-properties-lemma-immersion-into-presentation}
Soit $(U, R, s, t, c)$ un groupoïde lisse en espaces algébriques.
Soit $i : \mathcal{Z} \to [U/R]$ une immersion.
Il existe alors un sous-espace localement fermé $R$-invariant
$Z \subset U$ et une présentation $[Z/R_Z] \to \mathcal{Z}$,
où $R_Z$ est la restriction de $R$ à $Z$, tels que
$$
\xymatrix{
[Z/R_Z] \ar[dr] \ar[rr] & & \mathcal{Z} \ar[ld]^i \\
& [U/R]
}
$$
soit $2$-commutatif. Si $i$ est une immersion fermée (resp.\ ouverte),
alors $Z$ est un sous-espace fermé (resp.\ ouvert) de $U$.
\end{lemma}

\begin{proof}
D'après le
Lemme \ref{stacks-properties-lemma-representable-in-terms-presentations},
nous obtenons un diagramme commutatif
$$
\xymatrix{
[U'/R'] \ar[dr] \ar[rr] & & \mathcal{Z} \ar[ld] \\
& [U/R]
}
$$
où $U' = \mathcal{Z} \times_{[U/R]} U$ et
$R' = \mathcal{Z} \times_{[U/R]} R$.
Puisque $\mathcal{Z} \to [U/R]$ est une immersion, nous voyons que
$U' \to U$ est une immersion d'espaces algébriques. Soit $Z \subset U$
le sous-espace localement fermé tel que $U' \to U$ se factorise par
$Z$ et induise un isomorphisme $U' \to Z$.
Il ressort clairement de la construction de $R'$ que
$R' = U' \times_{U, t} R = R \times_{s, U} U'$.
Cela implique que $Z \cong U'$ est $R$-invariant et que l'image de
$R' \to R$ identifie $R'$ à la restriction
$R_Z = s^{-1}(Z) = t^{-1}(Z)$ de $R$ à $Z$. Le lemme en résulte.
\end{proof}

\begin{lemma}
\label{stacks-properties-lemma-immersion-presentation}
Soit $(U, R, s, t, c)$ un groupoïde lisse en espaces algébriques.
Soit $\mathcal{X} = [U/R]$ le champ algébrique associé, voir
Champs algébriques,
Théorème \ref{algebraic-theorem-smooth-groupoid-gives-algebraic-stack}.
Soit $Z \subset U$ un sous-espace localement fermé $R$-invariant. Alors
$$
[Z/R_Z] \longrightarrow [U/R]
$$
est une immersion de champs algébriques, où $R_Z$ est la restriction
de $R$ à $Z$. Si $Z \subset U$ est ouvert (resp.\ fermé), alors le morphisme
est une immersion ouverte (resp.\ fermée) de champs algébriques.
\end{lemma}

\begin{proof}
Rappelons que, d'après
Groupoïdes en espaces,
Définition \ref{spaces-groupoids-definition-invariant-open}
(voir aussi la discussion qui suit la définition),
on a $R_Z = s^{-1}(Z) = t^{-1}(Z)$ comme sous-espaces localement fermés
de $R$. Les deux morphismes $R_Z \to Z$ sont donc lisses comme changements de base
de $s$ et $t$. Par conséquent,
$(Z, R_Z, s|_{R_Z}, t|_{R_Z}, c|_{R_Z \times_{s, Z, t} R_Z})$ est
un groupoïde lisse en espaces algébriques, et l'on voit que
$[Z/R_Z]$ est un champ algébrique, voir
Champs algébriques,
Théorème \ref{algebraic-theorem-smooth-groupoid-gives-algebraic-stack}.
Les hypothèses de
Groupoïdes en espaces,
Lemme \ref{spaces-groupoids-lemma-criterion-fibre-product}
sont toutes satisfaites; il en résulte que l'on a un carré $2$-cartésien
$$
\xymatrix{
Z \ar[d] \ar[r] & [Z/R_Z] \ar[d] \\
U \ar[r] & [U/R]
}
$$
Il résulte de ceci et du
Lemme \ref{stacks-properties-lemma-check-representable-covering}
que $[Z/R_Z] \to [U/R]$ est représentable par des espaces algébriques;
alors, il résulte du
Lemme \ref{stacks-properties-lemma-check-property-covering}
que la flèche verticale de droite est une immersion (resp.\ une immersion fermée,
resp.\ une immersion ouverte) si et seulement si la flèche verticale de gauche l'est.
\end{proof}

\noindent
Nous pouvons définir les sous-champs ouverts, fermés et localement fermés comme suit.

\begin{definition}
\label{stacks-properties-definition-substacks}
Soit $\mathcal{X}$ un champ algébrique.
\begin{enumerate}
\item Un {\it sous-champ ouvert} de $\mathcal{X}$ est une sous-catégorie strictement pleine
$\mathcal{X}' \subset \mathcal{X}$ telle que $\mathcal{X}'$ soit un champ
algébrique et que $\mathcal{X}' \to \mathcal{X}$ soit une immersion ouverte.
\item Un {\it sous-champ fermé} de $\mathcal{X}$ est une sous-catégorie strictement pleine
$\mathcal{X}' \subset \mathcal{X}$ telle que $\mathcal{X}'$ soit un champ
algébrique et que $\mathcal{X}' \to \mathcal{X}$ soit une immersion fermée.
\item Un {\it sous-champ localement fermé} de $\mathcal{X}$ est une sous-catégorie
strictement pleine $\mathcal{X}' \subset \mathcal{X}$ telle que $\mathcal{X}'$
soit un champ algébrique et que $\mathcal{X}' \to \mathcal{X}$ soit une immersion.
\end{enumerate}
\end{definition}

\noindent
Cette définition doit être employée avec précaution. En effet, si
$f : \mathcal{X} \to \mathcal{Y}$ est une équivalence de champs algébriques
et si $\mathcal{X}' \subset \mathcal{X}$ est un sous-champ ouvert, la sous-catégorie
$f(\mathcal{X}')$ n'est pas nécessairement un sous-champ ouvert de $\mathcal{Y}$.
Le problème vient de ce qu'elle peut ne pas être une sous-catégorie {\it strictement}
pleine; mais c'est aussi le seul problème.
Voici un énoncé formel.

\begin{lemma}
\label{stacks-properties-lemma-substack-image}
Pour toute immersion $i : \mathcal{Z} \to \mathcal{X}$, il existe un
unique sous-champ localement fermé $\mathcal{X}' \subset \mathcal{X}$
tel que $i$ se factorise en la composée d'une
équivalence $i' : \mathcal{Z} \to \mathcal{X}'$
et du morphisme d'inclusion $\mathcal{X}' \to \mathcal{X}$.
Si $i$ est une immersion fermée (resp.\ ouverte), alors $\mathcal{X}'$
est un sous-champ fermé (resp.\ ouvert) de $\mathcal{X}$.
\end{lemma}

\begin{proof}
Omis.
\end{proof}

\begin{lemma}
\label{stacks-properties-lemma-substacks-presentation}
Soit $[U/R] \to \mathcal{X}$ une présentation d'un champ algébrique.
Il existe une bijection canonique
$$
\text{sous-champs localement fermés }\mathcal{Z}\text{ de }\mathcal{X}
\longrightarrow
R\text{-invariants parmi les sous-espaces localement fermés }Z\text{ de }U
$$
qui envoie $\mathcal{Z}$ sur $U \times_\mathcal{X} \mathcal{Z}$.
De plus, un morphisme de champs algébriques $f : \mathcal{Y} \to \mathcal{X}$
se factorise par $\mathcal{Z}$ si et seulement si
$\mathcal{Y} \times_\mathcal{X} U \to U$ se factorise par $Z$.
Il en va de même des sous-champs fermés et des sous-champs ouverts.
\end{lemma}

\begin{proof}
D'après les Lemmes \ref{stacks-properties-lemma-immersion-into-presentation} et
\ref{stacks-properties-lemma-immersion-presentation},
l'application est une bijection.
Si $\mathcal{Y} \to \mathcal{X}$ se factorise par $\mathcal{Z}$,
le changement de base $\mathcal{Y} \times_\mathcal{X} U \to U$
se factorise bien sûr par $Z$. Réciproquement, supposons que $\mathcal{Y} \to \mathcal{X}$
soit un morphisme tel que $\mathcal{Y} \times_\mathcal{X} U \to U$
se factorise par $Z$. Nous allons montrer que, pour tout schéma $T$ et tout
morphisme $T \to \mathcal{Y}$,
donné par un objet $y$ de la catégorie fibre de $\mathcal{Y}$ au-dessus de $T$,
l'objet $y$ appartient en fait à la catégorie fibre de $\mathcal{Z}$ au-dessus de $T$.
En effet, le produit fibré $T \times_\mathcal{X} U$ est un espace
algébrique et $T \times_\mathcal{X} U \to T$ est un morphisme lisse surjectif.
Il existe donc un recouvrement fppf $\{T_i \to T\}$ tel que
$T_i \to T$ se factorise par $T \times_\mathcal{X} U \to T$ pour tout $i$.
Alors $T_i \to \mathcal{X}$ se factorise par $\mathcal{Y} \times_\mathcal{X} U$,
et donc par $Z \subset U$. Ainsi, $y|_{T_i}$ est un objet de
$\mathcal{Z}$ (car $Z$ est le produit fibré de $U$ et de
$\mathcal{Z}$ au-dessus de $\mathcal{X}$).
Puisque $\mathcal{Z}$ est un sous-champ strictement plein, nous concluons
que $y$ est un objet de $\mathcal{Z}$, comme voulu.
\end{proof}

\begin{lemma}
\label{stacks-properties-lemma-open-substacks}
Soit $\mathcal{X}$ un champ algébrique. La règle
$\mathcal{U} \mapsto |\mathcal{U}|$ définit une bijection qui préserve les inclusions
entre les sous-champs ouverts de $\mathcal{X}$ et les parties ouvertes
de $|\mathcal{X}|$.
\end{lemma}

\begin{proof}
Choisissons une présentation $[U/R] \to \mathcal{X}$, voir
Champs algébriques, Lemme \ref{algebraic-lemma-stack-presentation}.
D'après le
Lemme \ref{stacks-properties-lemma-substacks-presentation},
les sous-champs ouverts correspondent aux sous-schémas ouverts $R$-invariants
de $U$. D'autre part, les
Lemmes \ref{stacks-properties-lemma-points-presentation} et \ref{stacks-properties-lemma-topology-points}
assurent que ceux-ci correspondent bijectivement aux parties ouvertes de $|\mathcal{X}|$.
\end{proof}

\begin{lemma}
\label{stacks-properties-lemma-open-image-substack}
Soit $\mathcal X$ un champ algébrique. Soit $U$ un espace algébrique et
$U \to \mathcal X$ un morphisme lisse surjectif. Pour une immersion ouverte
$V \hookrightarrow U$, il existe un champ algébrique $\mathcal Y$, une
immersion ouverte $\mathcal Y \to \mathcal X$ et un morphisme lisse
surjectif $V \to \mathcal Y$.
\end{lemma}

\begin{proof}
Nous définissons une catégorie fibrée en groupoïdes $\mathcal Y$ en prenant pour catégorie
fibre $\mathcal{Y}_T$ au-dessus d'un objet $T$ de $(\Sch/S)_{fppf}$
la sous-catégorie pleine de $\mathcal{X}_T$ formée de tous les
$y \in \Ob(\mathcal{X}_T)$ tels que le morphisme de projection
$V \times_{\mathcal X, y} T \to T$ soit surjectif. Or, pour tout morphisme
$x : T \to \mathcal X$, le produit $2$-fibré
$T \times_{x, \mathcal X} \mathcal Y$ a pour catégorie fibre au-dessus de $T'$ celle formée
des triplets $(f : T' \to T, y \in \mathcal{X}_{T'}, f^*x \simeq y)$ tels que
$V \times_{\mathcal X, y} T' \to T'$ soit surjectif.
Remarquons que $T \times_{x, \mathcal X} \mathcal Y$ est fibré en setoïdes puisque
$\mathcal Y \to \mathcal X$ est fidèle (voir
Champs, Lemme \ref{stacks-lemma-2-fibre-product-gives-stack-in-setoids}).
L'isomorphisme $f^*x \simeq y$ donne alors le diagramme
$$
\xymatrix{
V \times_{\mathcal X, y} T' \ar[d] \ar[r] &
V \times_{\mathcal X, x} T \ar[r] \ar[d] &
V \ar[d] \\
T' \ar[r]^f &
T \ar[r]^x &
\mathcal X
}
$$
où les deux carrés sont cartésiens. Le morphisme
$V \times_{\mathcal X, x} T \to T$ est lisse par changement de base, et donc ouvert.
Soit $T_0 \subset T$ son image. Les carrés cartésiens montrent que
$V \times_{\mathcal X, y} T' \to T'$ est surjectif si et seulement si $f$ se factorise
par $T_0$. Par conséquent, $T \times_{x, \mathcal X} \mathcal Y$ est représentable par
$T_0$, de sorte que l'inclusion $\mathcal Y \to \mathcal X$ est une immersion ouverte.
D'après
Champs algébriques, Lemme \ref{algebraic-lemma-open-fibred-category-is-algebraic},
nous concluons que $\mathcal{Y}$ est un champ algébrique.
Enfin, si nous désignons le morphisme $V \to \mathcal X$ par $g$, alors
$V \times_{\mathcal X} V \to V$ est surjectif (la diagonale fournit une
section). Ainsi, $g$ appartient à l'image de $\mathcal{Y}_V \to \mathcal{X}_V$, c'est-à-dire
que l'on obtient un morphisme $g' : V \to \mathcal{Y}$ s'inscrivant dans le
diagramme commutatif
$$
\xymatrix{
V \ar[r] \ar[d]^{g'} & U \ar[d] \\
\mathcal{Y} \ar[r] & \mathcal{X}
}
$$
Puisque $V \times_{g, \mathcal X} \mathcal Y \to V$ est un
monomorphisme, c'est en fait un isomorphisme car $(1, g')$ en définit une section.
Par conséquent, $g' : V \to \mathcal Y$ est un morphisme lisse, puisqu'il est le
changement de base du morphisme lisse $g : V \to \mathcal{X}$.
Il est surjectif par notre construction de $\mathcal{Y}$, ce qui achève
la démonstration du lemme.
\end{proof}

\begin{lemma}
\label{stacks-properties-lemma-union-open-substacks}
Soit $\mathcal X$ un champ algébrique et soient $\mathcal{X}_i \subset \mathcal X$
une famille de sous-champs ouverts indexée par $i \in I$. Il existe alors un
sous-champ ouvert, que nous notons
$\bigcup_{i\in I} \mathcal{X}_i \subset \mathcal X$, tel que
les $\mathcal{X}_i$ soient des sous-champs ouverts qui le recouvrent.
\end{lemma}

\begin{proof}
Nous définissons une sous-catégorie fibrée
$\mathcal{X}' = \bigcup_{i \in I} \mathcal{X}_i$
en prenant pour catégorie fibre au-dessus d'un objet $T$ de $(\Sch/S)_{fppf}$
la sous-catégorie pleine de $\mathcal{X}_T$ formée de tous les
$x \in \Ob(\mathcal{X}_T)$ tels que le morphisme
$\coprod_{i \in I} (\mathcal{X}_i \times_{\mathcal X} T) \to T$
soit surjectif. Soit $x_i \in \Ob((\mathcal{X}_i)_T)$.
Alors $(x_i, 1)$ définit une section de
$\mathcal{X}_i \times_{\mathcal X} T \to T$, et donc un isomorphisme. Ainsi,
$\mathcal{X}_i \subset \mathcal{X}'$ est une sous-catégorie pleine.
Soit maintenant $x \in \Ob(\mathcal{X}_T)$. Alors
$\mathcal{X}_i \times_{\mathcal X} T$ est représentable
par un sous-schéma ouvert $T_i \subset T$. Le produit $2$-fibré
$\mathcal{X}' \times_{\mathcal X} T$ a pour fibre au-dessus de $T'$ les
$(y \in \mathcal{X}_{T'}, f : T' \to T, f^*x \simeq y)$ tels que
$\coprod (\mathcal{X}_i \times_{\mathcal X, y} T') \to T'$ soit surjectif.
L'isomorphisme $f^*x \simeq y$ induit un isomorphisme
$\mathcal{X}_i \times_{\mathcal X, y} T' \simeq T_i \times_T T'$.
Alors les $T_i \times_T T'$ recouvrent $T'$ si et seulement si $f$ se factorise par
$\bigcup T_i$. Nous avons donc un diagramme
$$
\xymatrix{
T_i \ar[r] \ar[d] &
\bigcup T_i \ar[r] \ar[d] &
T \ar[d] \\
\mathcal{X}_i \ar[r] &
\mathcal{X}' \ar[r] &
\mathcal{X}
}
$$
dont les deux carrés sont cartésiens. D'après
Champs algébriques, Lemme \ref{algebraic-lemma-open-fibred-category-is-algebraic},
nous concluons que $\mathcal{X'} \subset \mathcal{X}$ est algébrique et constitue un
sous-champ ouvert. Il ressort aussi des carrés cartésiens ci-dessus que le
morphisme $\coprod_{i \in I} \mathcal{X}_i \to \mathcal{X}'$ est surjectif, ce
qui achève la démonstration du lemme.
\end{proof}

\begin{lemma}
\label{stacks-properties-lemma-quasi-compact-finite-subcover}
Soit $\mathcal X$ un champ algébrique et soit $\mathcal X' \subset \mathcal X$
un sous-champ ouvert quasi-compact. Supposons donnée une famille de sous-
champs ouverts $\mathcal{X}_i \subset \mathcal X$ indexée par $i \in I$ telle
que $\mathcal{X}' \subset \bigcup_{i \in I} \mathcal{X}_i$,
où la réunion est définie comme au Lemme \ref{stacks-properties-lemma-union-open-substacks}.
Il existe alors une partie finie $I' \subset I$ telle que
$\mathcal{X}' \subset \bigcup_{i \in I'} \mathcal{X}_i$.
\end{lemma}

\begin{proof}
Puisque $\mathcal X$ est algébrique, il existe un schéma $U$ muni d'un morphisme
lisse surjectif $U \to \mathcal X$. Soit $U_i \subset U$ le sous-schéma ouvert
représentant $\mathcal{X}_i \times_{\mathcal X} U$, et $U' \subset U$ le
sous-schéma ouvert représentant $\mathcal{X}' \times_{\mathcal X} U$. Par
hypothèse, $U'\subset \bigcup_{i\in I} U_i$. D'après la démonstration du
Lemme \ref{stacks-properties-lemma-quasi-compact-stack},
il existe un ouvert quasi-compact $V \subset U'$ tel que $V \to \mathcal{X}'$
soit un morphisme lisse surjectif. Il existe donc une partie finie
$I' \subset I$ telle que $V \subset \bigcup_{i \in I'}
U_i$. Nous affirmons que $\mathcal{X}' \subset \bigcup_{i \in I'} \mathcal{X}_i$.
Prenons $x \in \Ob(\mathcal{X}'_T)$ pour
$T \in \Ob((\Sch/S)_{fppf})$.
Comme $\mathcal{X}' \to \mathcal{X}$ est un monomorphisme, nous avons les carrés
cartésiens
$$
\xymatrix{
V \times_\mathcal{X} T \ar[r] \ar[d] &
T \ar[d]^x \ar@{=}[r] &
T \ar[d]^x \\
V \ar[r] &
\mathcal{X}' \ar[r] &
\mathcal X
}
$$
Par changement de base, $V \times_{\mathcal X} T \to T$ est surjectif. Par conséquent,
$\bigcup_{i \in I'} U_i \times_{\mathcal X} T \to T$ est lui aussi surjectif.
Soit $T_i \subset T$ le sous-schéma ouvert représentant
$\mathcal{X}_i \times_{\mathcal X} T$.
Un argument formel donne le carré cartésien
$$
\xymatrix{
U_i \times_{\mathcal{X}_i} T_i \ar[r] \ar[d] &
U \times_{\mathcal X} T \ar[d] \\
T_i \ar[r] & T
}
$$
où les flèches verticales sont surjectives par changement de base. Puisque
$U_i \times_{\mathcal{X}_i} T_i \simeq U_i \times_{\mathcal X} T$,
nous obtenons $\bigcup_{i \in I'} T_i = T$. Ainsi,
$x$ est un objet de $(\bigcup_{i\in I'} \mathcal{X}_i)_T$ par
définition de la réunion. Observons que l'inclusion
$\mathcal{X}' \subset \bigcup_{i \in I'} \mathcal{X}_i$
est automatiquement un sous-champ ouvert.
\end{proof}

\begin{lemma}
\label{stacks-properties-lemma-zariski-open-cover-stack-is-space}
Soit $\mathcal X$ un champ algébrique.
Soit $\mathcal{X}_i$, $i \in I$, une famille de sous-champs ouverts de $\mathcal{X}$.
Supposons que
\begin{enumerate}
\item $\mathcal{X} = \bigcup_{i \in I} \mathcal{X}_i$, et que
\item chaque $\mathcal{X}_i$ soit un espace algébrique.
\end{enumerate}
Alors $\mathcal{X}$ est un espace algébrique.
\end{lemma}

\begin{proof}
Appliquons
Champs, Lemme \ref{stacks-lemma-stack-in-setoids-descent}
au morphisme $\coprod_{i \in I} \mathcal{X}_i \to \mathcal{X}$
et au morphisme $\text{id} : \mathcal{X} \to \mathcal{X}$ pour
voir que $\mathcal{X}$ est un champ en setoïdes.
Ainsi, $\mathcal{X}$ est un espace algébrique, voir
Champs algébriques,
Proposition \ref{algebraic-proposition-algebraic-stack-no-automorphisms}.
\end{proof}

\begin{lemma}
\label{stacks-properties-lemma-zariski-open-cover-stack-is-scheme}
Soit $\mathcal X$ un champ algébrique.
Soit $\mathcal{X}_i$, $i \in I$, une famille de sous-champs ouverts de $\mathcal{X}$.
Supposons que
\begin{enumerate}
\item $\mathcal{X} = \bigcup_{i \in I} \mathcal{X}_i$, et que
\item chaque $\mathcal{X}_i$ soit un schéma.
\end{enumerate}
Alors $\mathcal{X}$ est un schéma.
\end{lemma}

\begin{proof}
D'après le
Lemme \ref{stacks-properties-lemma-zariski-open-cover-stack-is-space},
on voit que $\mathcal{X}$ est un espace algébrique. Comme tout espace
algébrique possède un plus grand sous-espace ouvert qui soit un schéma, voir
Propriétés des espaces, Lemme \ref{spaces-properties-lemma-subscheme},
on voit que $\mathcal{X}$ est un schéma.
\end{proof}


\noindent
Le lemme suivant est l'analogue de
Compléments sur les groupoïdes, Lemme \ref{more-groupoids-lemma-local-source}.

\begin{lemma}
\label{stacks-properties-lemma-local-source}
Soient $\mathcal{P}, \mathcal{Q}, \mathcal{R}$ des propriétés de morphismes
d'espaces algébriques. Supposons que
\begin{enumerate}
\item $\mathcal{P}, \mathcal{Q}, \mathcal{R}$ soient locales pour la topologie fppf sur le but
et stables par tout changement de base,
\item $\text{lisse} \Rightarrow \mathcal{R}$,
\item pour tout morphisme $f : X \to Y$ qui possède $\mathcal{Q}$, il existe un
plus grand sous-espace ouvert $W(\mathcal{P}, f) \subset X$ tel que
$f|_{W(\mathcal{P}, f)}$ possède $\mathcal{P}$, et
\item pour tout morphisme $f : X \to Y$ qui possède $\mathcal{Q}$,
et tout morphisme $Y' \to Y$ qui possède $\mathcal{R}$, on ait
$Y' \times_Y W(\mathcal{P}, f) = W(\mathcal{P}, f')$, où
$f' : X_{Y'} \to Y'$ est le changement de base de $f$.
\end{enumerate}
Soit $f : \mathcal{X} \to \mathcal{Y}$ un morphisme de champs algébriques
représentable par des espaces algébriques. Supposons que $f$ possède $\mathcal{Q}$. Alors
\begin{enumerate}
\item[(A)] il existe un plus grand sous-champ ouvert
$\mathcal{X}' \subset \mathcal{X}$ tel que $f|_{\mathcal{X}'}$ possède
$\mathcal{P}$, et
\item[(B)] si $\mathcal{Z} \to \mathcal{Y}$ est un morphisme de champs
algébriques représentable par des espaces algébriques qui possède $\mathcal{R}$,
alors $\mathcal{Z} \times_\mathcal{Y} \mathcal{X}'$ est le plus grand sous-champ
ouvert de $\mathcal{Z} \times_\mathcal{Y} \mathcal{X}$ au-dessus duquel
le changement de base $\text{id}_\mathcal{Z} \times f$ possède la propriété $\mathcal{P}$.
\end{enumerate}
\end{lemma}

\begin{proof}
Choisissons un schéma $V$ et un morphisme lisse surjectif $V \to \mathcal{Y}$.
Posons $U = V \times_\mathcal{Y} \mathcal{X}$ et soit $f' : U \to V$ le
changement de base de $f$. Le morphisme d'espaces algébriques $f' : U \to V$ possède
la propriété $\mathcal{Q}$. L'hypothèse (3) fournit donc l'ouvert $W(\mathcal{P}, f') \subset U$.
Remarquons que
$U \times_\mathcal{X} U =
(V \times_\mathcal{Y} V) \times_\mathcal{Y} \mathcal{X}$,
si bien que le morphisme $f'' : U \times_\mathcal{X} U \to V \times_\mathcal{Y} V$
est le changement de base de $f$ par l'une ou l'autre projection
$V \times_\mathcal{Y} V \to V$. Par notre choix de $V$, ces projections
sont lisses, et possèdent donc la propriété $\mathcal{R}$ d'après (2). Ainsi, (4) montre
que les images inverses de $W(\mathcal{P}, f')$ par les deux projections
$\text{pr}_i : U \times_\mathcal{X} U \to U$ coïncident. Autrement dit,
$W(\mathcal{P}, f')$ est un sous-espace $R$-invariant de $U$ (où
$R = U \times_\mathcal{X} U$). Soit $\mathcal{X}'$ le sous-champ ouvert de
$\mathcal{X}$ qui correspond à $W(\mathcal{P}, f)$ par le
Lemme \ref{stacks-properties-lemma-immersion-into-presentation}.
Par construction, $W(\mathcal{P}, f') = \mathcal{X}' \times_\mathcal{Y} V$,
donc $f|_{\mathcal{X}'}$ possède la propriété $\mathcal{P}$ d'après le
Lemme \ref{stacks-properties-lemma-check-property-covering}.
De plus, $\mathcal{X}'$ est le plus grand sous-champ ouvert
tel que $f|_{\mathcal{X}'}$ possède $\mathcal{P}$, puisque la même maximalité vaut
pour $W(\mathcal{P}, f)$. Cela prouve (A).

\medskip\noindent
Enfin, si $\mathcal{Z} \to \mathcal{Y}$ est un
morphisme de champs algébriques représentable par des espaces algébriques qui possède
$\mathcal{R}$, posons $T = V \times_\mathcal{Y} \mathcal{Z}$; on voit alors que
$T \to V$ est un morphisme d'espaces algébriques possédant
la propriété $\mathcal{R}$. Soit $f'_T : T \times_V U \to T$ le changement de base
de $f'$. De nouveau, (4) montre que $W(\mathcal{P}, f'_T)$
est l'image inverse de $W(\mathcal{P}, f)$ dans $T \times_V U$. Cela implique
(B); nous omettons certains détails.
\end{proof}

\begin{remark}
\label{stacks-properties-remark-local-source-warning}
Avertissement :
Le Lemme \ref{stacks-properties-lemma-local-source}
doit être utilisé avec précaution. Par exemple, il s'applique à
$\mathcal{P}=$``plat'', $\mathcal{Q}=$``vide'',
et $\mathcal{R}=$``plat et localement de présentation finie''. Mais, étant donné un
morphisme d'espaces algébriques $f : X \to Y$, le plus grand sous-
espace ouvert $W \subset X$ tel que $f|_W$ soit plat n'est {\it pas} l'ensemble des points
où $f$ est plat !
\end{remark}

\begin{remark}
\label{stacks-properties-remark-local-source-apply}
Malgré l'avertissement de la
Remarque \ref{stacks-properties-remark-local-source-warning},
il existe des cas où le
Lemme \ref{stacks-properties-lemma-local-source}
peut être utilisé sans ambiguïté.
Nous en donnons une liste. Dans chaque cas, nous omettons la vérification des
hypothèses (1) et (2), et donnons des références qui impliquent
(3) et (4). Voici la liste :
\begin{enumerate}
\item
\label{stacks-properties-item-rel-dim-leq-d}
$\mathcal{Q} = $``localement de type fini'', $\mathcal{R} = \emptyset$,
et $\mathcal{P} =$``dimension relative $\leq d$''.
Voir
Morphismes d'espaces,
Définition \ref{spaces-morphisms-definition-relative-dimension}
et
Morphismes d'espaces, Lemmes
\ref{spaces-morphisms-lemma-openness-bounded-dimension-fibres} et
\ref{spaces-morphisms-lemma-dimension-fibre-after-base-change}.
\item
\label{stacks-properties-item-loc-quasi-finite}
$\mathcal{Q} =$``localement de type fini'', $\mathcal{R} = \emptyset$,
et $\mathcal{P} =$``localement quasi-fini''.
Il s'agit du cas $d = 0$ du point précédent, voir
Morphismes d'espaces, Lemme
\ref{spaces-morphisms-lemma-locally-quasi-finite-rel-dimension-0}.
D'autre part, les propriétés (3) et (4) sont explicitées dans
Morphismes d'espaces, Lemme
\ref{spaces-morphisms-lemma-locally-finite-type-quasi-finite-part}.
\item
\label{stacks-properties-item-unramified}
$\mathcal{Q} = $``localement de type fini'', $\mathcal{R} = \emptyset$,
et $\mathcal{P} =$``non ramifié''. C'est le
Morphismes d'espaces, Lemme \ref{spaces-morphisms-lemma-where-unramified}.
\item
\label{stacks-properties-item-flat}
$\mathcal{Q} =$``localement de présentation finie'',
$\mathcal{R} =$``plat et localement de présentation finie'', et
$\mathcal{P} =$``plat''. Voir
Compléments sur les morphismes d'espaces, Théorème
\ref{spaces-more-morphisms-theorem-openness-flatness} et
Lemme \ref{spaces-more-morphisms-lemma-flat-locus-base-change}.
Remarquons qu'ici $W(\mathcal{P}, f)$ est toujours exactement l'ensemble des points
où le morphisme $f$ est plat, car nous ne considérons cet ouvert que
lorsque $f$ possède $\mathcal{Q}$ (voir loc. cit.).
\item
\label{stacks-properties-item-etale}
$\mathcal{Q} =$``localement de présentation finie'',
$\mathcal{R} =$``plat et localement de présentation finie'', et
$\mathcal{P}=$``\'etale''. Cela résulte de la combinaison de
(\ref{stacks-properties-item-unramified}) et (\ref{stacks-properties-item-flat}), puisqu'un morphisme non ramifié
qui est plat et localement de présentation finie est \'etale, voir
Morphismes d'espaces,
Lemme \ref{spaces-morphisms-lemma-unramified-flat-lfp-etale}.
\item Ajouter ici d'autres cas selon les besoins (comparer avec la liste plus longue de
Compléments sur les groupoïdes, Remarque \ref{more-groupoids-remark-local-source-apply}).
\end{enumerate}
\end{remark}











\section{Champs algébriques réduits}
\label{stacks-properties-section-reduced}

\noindent
Nous avons déjà défini les champs algébriques réduits dans la
section \ref{stacks-properties-section-types-properties}.

\begin{lemma}
\label{stacks-properties-lemma-reduced-closed-substack}
Soit $\mathcal{X}$ un champ algébrique.
Soit $T \subset |\mathcal{X}|$ une partie fermée.
Il existe un unique sous-champ fermé $\mathcal{Z} \subset \mathcal{X}$
possédant les propriétés suivantes :
(a) on a $|\mathcal{Z}| = T$, et (b) $\mathcal{Z}$ est réduit.
\end{lemma}

\begin{proof}
Soit $U \to \mathcal{X}$ un morphisme lisse surjectif, où $U$ est un
espace algébrique. Posons $R = U \times_\mathcal{X} U$; on a alors une
présentation $[U/R] \to \mathcal{X}$, voir
Champs algébriques, Lemme \ref{algebraic-lemma-stack-presentation}.
Comme d'habitude, nous notons $s, t : R \to U$ les deux morphismes de projection lisses.
D'après le Lemme \ref{stacks-properties-lemma-points-presentation},
on voit que $T$ correspond à une partie fermée $T' \subset |U|$ telle
que $|s|^{-1}(T') = |t|^{-1}(T')$.
Soit $Z \subset U$ l'espace algébrique obtenu en munissant
$T'$ de la structure réduite induite, voir
Propriétés des espaces,
Définition \ref{spaces-properties-definition-reduced-induced-space}.
Les produits fibrés
$Z \times_{U, t} R$ et $R \times_{s, U} Z$ sont des sous-espaces fermés
de $R$
(Espaces, Lemme \ref{spaces-lemma-base-change-immersions}).
Les projections $Z \times_{U, t} R \to Z$ et
$R \times_{s, U} Z \to Z$ sont lisses d'après
Morphismes d'espaces, Lemme \ref{spaces-morphisms-lemma-base-change-smooth}.
Ainsi, puisque $Z$ est réduit, il s'ensuit que
$Z \times_{U, t} R$ et $R \times_{s, U} Z$ sont réduits, voir
Remarque \ref{stacks-properties-remark-list-properties-local-smooth-topology}.
Comme
$$
|Z \times_{U, t} R| = |t|^{-1}(T') = |s|^{-1}(T') = R \times_{s, U} Z
$$
l'unicité de
Propriétés des espaces,
Lemme \ref{spaces-properties-lemma-reduced-closed-subspace}
entraîne que $Z \times_{U, t} R = R \times_{s, U} Z$.
Ainsi, $Z$ est un sous-espace fermé $R$-invariant de $U$.
Par la correspondance du
Lemme \ref{stacks-properties-lemma-substacks-presentation},
nous obtenons un sous-champ fermé $\mathcal{Z} \subset \mathcal{X}$
tel que $Z = \mathcal{Z} \times_\mathcal{X} U$. Alors
$[Z/R_Z] \to \mathcal{Z}$ est une présentation
(Lemme \ref{stacks-properties-lemma-immersion-into-presentation}).
Enfin, $|\mathcal{Z}| = |Z|/|R_Z| = |T'|/\sim$ est la
partie fermée donnée $T$. Nous omettons la démonstration de l'unicité.
\end{proof}

\begin{lemma}
\label{stacks-properties-lemma-reduced-stack-determined-by-points}
Soit $\mathcal{X}$ un champ algébrique.
Si $\mathcal{X}' \subset \mathcal{X}$
est un sous-champ fermé, si $\mathcal{X}$ est réduit et si
$|\mathcal{X}'| = |\mathcal{X}|$, alors $\mathcal{X}' = \mathcal{X}$.
\end{lemma}

\begin{proof}
Choisissons une présentation $[U/R] \to \mathcal{X}$ avec $U$ un schéma.
Comme $\mathcal{X}$ est réduit, $U$ est réduit (par définition
des champs algébriques réduits). D'après le
Lemme \ref{stacks-properties-lemma-substacks-presentation},
$\mathcal{X}'$ correspond à un sous-schéma fermé $R$-invariant $Z \subset U$.
Mais $|Z| \subset |U|$ est l'image inverse de $|\mathcal{X}'|$, et
donc $|Z| = |U|$. Ainsi, $Z$ est un sous-schéma fermé de $U$ dont les ensembles
de points sous-jacents coïncident. D'après
Schémas, Lemme \ref{schemes-lemma-map-into-reduction},
l'application $\text{id}_U : U \to U$ se factorise par $Z \to U$, et donc
$Z = U$, c'est-à-dire $\mathcal{X}' = \mathcal{X}$.
\end{proof}

\begin{lemma}
\label{stacks-properties-lemma-map-into-reduction}
Soient $\mathcal{X}$, $\mathcal{Y}$ des champs algébriques.
Soit $\mathcal{Z} \subset \mathcal{X}$ un sous-champ fermé.
Supposons $\mathcal{Y}$ réduit.
Un morphisme $f : \mathcal{Y} \to \mathcal{X}$ se factorise par
$\mathcal{Z}$ si et seulement si
$f(|\mathcal{Y}|) \subset |\mathcal{Z}|$.
\end{lemma}

\begin{proof}
Supposons $f(|\mathcal{Y}|) \subset |\mathcal{Z}|$. Considérons
$\mathcal{Y} \times_\mathcal{X} \mathcal{Z} \to \mathcal{Y}$.
Il existe une équivalence
$\mathcal{Y} \times_\mathcal{X} \mathcal{Z} \to \mathcal{Y}'$
où $\mathcal{Y}'$ est un sous-champ fermé de $\mathcal{Y}$, voir les
Lemmes \ref{stacks-properties-lemma-base-change-immersion} et
\ref{stacks-properties-lemma-substack-image}.
À l'aide des
Lemmes \ref{stacks-properties-lemma-points-cartesian},
\ref{stacks-properties-lemma-monomorphism-injective-points} et
\ref{stacks-properties-lemma-immersion-monomorphism},
on voit que $|\mathcal{Y}'| = |\mathcal{Y}|$. Le lemme est donc ramené au
Lemme \ref{stacks-properties-lemma-reduced-stack-determined-by-points}.
\end{proof}

\begin{definition}
\label{stacks-properties-definition-reduced-induced-stack}
Soit $\mathcal{X}$ un champ algébrique.
Soit $Z \subset |\mathcal{X}|$ une partie fermée.
Une {\it structure de champ algébrique sur $Z$} est donnée par un sous-champ fermé
$\mathcal{Z}$ de $\mathcal{X}$ tel que $|\mathcal{Z}|$ soit égal à $Z$.
La {\it structure de champ algébrique réduite induite}
sur $Z$ est celle construite au
Lemme \ref{stacks-properties-lemma-reduced-closed-substack}.
La {\it réduction $\mathcal{X}_{red}$ de $\mathcal{X}$}
est la structure de champ algébrique réduite induite sur $|\mathcal{X}|$.
\end{definition}


\noindent
En fait, cela permet de définir la structure de champ algébrique réduite
induite sur une partie localement fermée.

\begin{remark}
\label{stacks-properties-remark-stack-structure-locally-closed-subset}
Soit $X$ un champ algébrique.
Soit $T \subset |\mathcal{X}|$ une partie localement fermée.
Soit $\partial T$ la frontière de $T$ dans
l'espace topologique $|\mathcal{X}|$. En formule,
$$
\partial T = \overline{T} \setminus T.
$$
Soit $\mathcal{U} \subset \mathcal{X}$ le sous-champ ouvert de $X$ tel que
$|\mathcal{U}| = |\mathcal{X}| \setminus \partial T$, voir
Lemme \ref{stacks-properties-lemma-open-substacks}.
Soit $\mathcal{Z}$ le sous-champ fermé réduit de $\mathcal{U}$ tel que
$|\mathcal{Z}| = T$, obtenu en prenant la structure réduite induite
de sous-espace fermé, voir
Définition \ref{stacks-properties-definition-reduced-induced-stack}.
Par construction, $\mathcal{Z} \to \mathcal{U}$ est une immersion fermée de
champs algébriques et $\mathcal{U} \to \mathcal{X}$ est une immersion ouverte;
ainsi, $\mathcal{Z} \to \mathcal{X}$ est une immersion de champs algébriques d'après le
Lemme \ref{stacks-properties-lemma-composition-immersion}.
Remarquons que $\mathcal{Z}$ est un champ algébrique réduit et que
$|\mathcal{Z}| = T$ comme parties de $|X|$. On dit parfois que
$\mathcal{Z}$ est la {\it structure de sous-champ réduite induite} sur $T$.
\end{remark}













\section{Gerbes résiduelles}
\label{stacks-properties-section-residual-gerbe}

\noindent
Dans le projet Champs, nous voudrions définir la {\it gerbe résiduelle}
d'un champ algébrique $\mathcal{X}$ en un point $x \in |\mathcal{X}|$
comme un monomorphisme de champs algébriques
$m_x : \mathcal{Z}_x \to \mathcal{X}$ où $\mathcal{Z}_x$ est un champ
algébrique réduit ayant un unique point, envoyé par $m_x$ sur $x$.
Il s'avère que cette notion soulève de nombreuses difficultés; l'existence n'est pas
claire en général, pas plus que l'unicité.
Nous résolvons le problème d'unicité en imposant une condition légèrement plus forte
aux champs algébriques $\mathcal{Z}_x$.
Nous examinons cela plus en détail au moyen de quelques lemmes simples
sur les champs algébriques réduits ayant un unique point.

\begin{lemma}
\label{stacks-properties-lemma-flat-cover-by-field}
Soit $\mathcal{Z}$ un champ algébrique. Soit $k$ un corps et soit
$\Spec(k) \to \mathcal{Z}$ un morphisme plat surjectif. Alors tout
morphisme $\Spec(k') \to \mathcal{Z}$, où $k'$ est un corps, est
plat et surjectif.
\end{lemma}

\begin{proof}
Considérons le carré fibré
$$
\xymatrix{
T \ar[d] \ar[r] & \Spec(k) \ar[d] \\
\Spec(k') \ar[r] & \mathcal{Z}
}
$$
Remarquons que $T \to \Spec(k')$ est plat et surjectif; ainsi, $T$
n'est pas vide. D'autre part, $T \to \Spec(k)$ est plat puisque
$k$ est un corps. Il s'ensuit que $T \to \mathcal{Z}$ est plat et surjectif.
D'après
Morphismes d'espaces, Lemme \ref{spaces-morphisms-lemma-flat-permanence}
(via la discussion de la
section \ref{stacks-properties-section-properties-morphisms}),
on en déduit que $\Spec(k') \to \mathcal{Z}$ est plat. Il est clairement
surjectif puisque, par hypothèse, $|\mathcal{Z}|$ est réduit à un point.
\end{proof}

\begin{lemma}
\label{stacks-properties-lemma-unique-point}
Soit $\mathcal{Z}$ un champ algébrique. Les assertions suivantes sont équivalentes :
\begin{enumerate}
\item $\mathcal{Z}$ est réduit et $|\mathcal{Z}|$ est réduit à un point,
\item il existe un morphisme plat surjectif $\Spec(k) \to \mathcal{Z}$,
où $k$ est un corps, et
\item il existe un morphisme localement de type fini, surjectif et plat
$\Spec(k) \to \mathcal{Z}$, où $k$ est un corps.
\end{enumerate}
\end{lemma}

\begin{proof}
Supposons (1). Soit $W$ un schéma et
soit $W \to \mathcal{Z}$ un morphisme lisse surjectif. Alors $W$ est
un schéma réduit. Soit $\eta \in W$ un point générique d'une composante
irréductible de $W$. Comme $W$ est réduit, on a
$\mathcal{O}_{W, \eta} = \kappa(\eta)$. Il s'ensuit que le morphisme canonique
$\eta = \Spec(\kappa(\eta)) \to W$ est plat. On voit que le
composé $\eta \to \mathcal{Z}$ est plat (voir
Morphismes d'espaces, Lemme \ref{spaces-morphisms-lemma-composition-flat}).
Il est aussi surjectif puisque $|\mathcal{Z}|$ est réduit à un point. Autrement dit,
(2) est vérifiée.

\medskip\noindent
Supposons (2). Soit $W$ un schéma et
soit $W \to \mathcal{Z}$ un morphisme lisse surjectif. Choisissons un corps
$k$ et un morphisme plat surjectif $\Spec(k) \to \mathcal{Z}$.
Alors $W \times_\mathcal{Z} \Spec(k)$ est un espace algébrique lisse
sur $k$, donc régulier (voir
Espaces sur les corps, Lemme \ref{spaces-over-fields-lemma-smooth-regular})
et en particulier réduit. Comme $W \times_\mathcal{Z} \Spec(k) \to W$
est plat et surjectif, nous concluons que $W$ est réduit
(Descente sur les espaces, Lemme \ref{spaces-descent-lemma-descend-reduced}).
Autrement dit, (1) est vérifiée.

\medskip\noindent
Il est clair que (3) implique (2). Enfin, supposons (2). Prenons un schéma affine
non vide $W$ et un morphisme lisse $W \to \mathcal{Z}$. Prenons un point fermé
$w \in W$ et posons $k = \kappa(w)$. Le composé
$$
\Spec(k) \xrightarrow{w} W \longrightarrow \mathcal{Z}
$$
est localement de type fini d'après
Morphismes d'espaces, Lemmes
\ref{spaces-morphisms-lemma-composition-finite-type} et
\ref{spaces-morphisms-lemma-smooth-locally-finite-type}.
Il est aussi plat et surjectif d'après le
Lemme \ref{stacks-properties-lemma-flat-cover-by-field}.
Ainsi, (3) est vérifiée.
\end{proof}

\noindent
Le lemme suivant distingue une classe légèrement meilleure de champs algébriques
réduits à un point que celle du lemme précédent.

\begin{lemma}
\label{stacks-properties-lemma-unique-point-better}
Soit $\mathcal{Z}$ un champ algébrique. Les assertions suivantes sont équivalentes :
\begin{enumerate}
\item $\mathcal{Z}$ est réduit, localement noethérien, et $|\mathcal{Z}|$
est réduit à un point, et
\item il existe un morphisme localement de présentation finie, surjectif et plat
$\Spec(k) \to \mathcal{Z}$, où $k$ est un corps.
\end{enumerate}
\end{lemma}

\begin{proof}
Supposons (2). D'après le
Lemme \ref{stacks-properties-lemma-unique-point},
on voit que $\mathcal{Z}$ est réduit et que $|\mathcal{Z}|$ est réduit à un point.
Soit $W$ un schéma et soit $W \to \mathcal{Z}$ un morphisme lisse
surjectif. Choisissons un corps $k$ et un morphisme localement de présentation finie,
surjectif et plat $\Spec(k) \to \mathcal{Z}$.
Alors $W \times_\mathcal{Z} \Spec(k)$ est un espace algébrique
lisse sur $k$, donc localement noethérien (voir
Morphismes d'espaces, Lemme
\ref{spaces-morphisms-lemma-locally-finite-type-locally-noetherian}).
Comme $W \times_\mathcal{Z} \Spec(k) \to W$
est plat, surjectif et localement de présentation finie, on voit
que $\{W \times_\mathcal{Z} \Spec(k) \to W\}$ est un recouvrement fppf,
et nous concluons que $W$ est localement noethérien
(Descente sur les espaces, Lemme
\ref{spaces-descent-lemma-descend-locally-Noetherian}).
Autrement dit, (1) est vérifiée.

\medskip\noindent
Supposons (1). Prenons un schéma affine non vide $W$ et un morphisme lisse
$W \to \mathcal{Z}$. Prenons un point fermé $w \in W$ et posons
$k = \kappa(w)$. Comme $W$ est localement noethérien, le morphisme
$w : \Spec(k) \to W$ est de présentation finie, voir
Morphismes, Lemme \ref{morphisms-lemma-closed-immersion-finite-presentation}.
Le composé
$$
\Spec(k) \xrightarrow{w} W \longrightarrow \mathcal{Z}
$$
est donc localement de présentation finie d'après
Morphismes d'espaces, Lemmes
\ref{spaces-morphisms-lemma-composition-finite-presentation} et
\ref{spaces-morphisms-lemma-smooth-locally-finite-presentation}.
Il est aussi plat et surjectif d'après le
Lemme \ref{stacks-properties-lemma-flat-cover-by-field}.
Ainsi, (2) est vérifiée.
\end{proof}

\begin{lemma}
\label{stacks-properties-lemma-monomorphism-into-point}
Soit $\mathcal{Z}' \to \mathcal{Z}$ un monomorphisme de champs algébriques.
Supposons qu'il existe un corps $k$ et un morphisme localement de présentation finie,
surjectif et plat $\Spec(k) \to \mathcal{Z}$. Alors soit $\mathcal{Z}'$
est vide, soit $\mathcal{Z}' \to \mathcal{Z}$ est une équivalence.
\end{lemma}

\begin{proof}
Nous pouvons supposer que $\mathcal{Z}'$ est non vide. Dans ce cas, le
produit fibré $T = \mathcal{Z}' \times_\mathcal{Z} \Spec(k)$
est non vide, voir
Lemme \ref{stacks-properties-lemma-points-cartesian}.
Or $T$ est un espace algébrique et la projection $T \to \Spec(k)$
est un monomorphisme. Ainsi, $T = \Spec(k)$, voir
Morphismes d'espaces, Lemme
\ref{spaces-morphisms-lemma-monomorphism-toward-field}.
Nous concluons que $\Spec(k) \to \mathcal{Z}$ se factorise par
$\mathcal{Z}'$. Supposons que le morphisme $z : \Spec(k) \to \mathcal{Z}$
soit donné par l'objet $\xi$ au-dessus de $\Spec(k)$. Nous venons de voir que
$\xi$ est isomorphe à un objet $\xi'$ de $\mathcal{Z}'$ au-dessus de
$\Spec(k)$. Puisque $z$
est surjectif, plat et localement de présentation finie, on voit que
tout objet de $\mathcal{Z}$ au-dessus d'un schéma quelconque est localement pour fppf isomorphe
à un tiré en arrière de $\xi$, et donc aussi à un tiré en arrière de $\xi'$. Par descente des
objets pour les champs en groupoïdes, cela implique que
$\mathcal{Z}' \to \mathcal{Z}$ est essentiellement surjectif (et aussi
pleinement fidèle, voir
Lemme \ref{stacks-properties-lemma-monomorphism}).
D'où le résultat.
\end{proof}

\begin{lemma}
\label{stacks-properties-lemma-improve-unique-point}
Soit $\mathcal{Z}$ un champ algébrique. Supposons que $\mathcal{Z}$ satisfasse
aux conditions équivalentes du
Lemme \ref{stacks-properties-lemma-unique-point}.
Il existe alors une unique sous-catégorie strictement pleine
$\mathcal{Z}' \subset \mathcal{Z}$ telle que
$\mathcal{Z}'$ soit un champ algébrique satisfaisant aux conditions équivalentes du
Lemme \ref{stacks-properties-lemma-unique-point-better}.
Le morphisme d'inclusion $\mathcal{Z}' \to \mathcal{Z}$ est un monomorphisme
de champs algébriques.
\end{lemma}

\begin{proof}
La dernière assertion résulte immédiatement de la première et du
Lemme \ref{stacks-properties-lemma-monomorphism}.
Choisissons un corps $k$ et un morphisme $\Spec(k) \to \mathcal{Z}$
qui soit surjectif, plat et localement de type fini.
Posons $U = \Spec(k)$ et $R = U \times_\mathcal{Z} U$.
Les projections $s, t : R \to U$ sont localement de type fini.
Comme $U$ est le spectre d'un corps, il s'ensuit que
$s, t$ sont plates et localement de présentation finie (d'après
Morphismes d'espaces, Lemme
\ref{spaces-morphisms-lemma-noetherian-finite-type-finite-presentation}).
On voit que $\mathcal{Z}' = [U/R]$ est un champ algébrique d'après
Critères de représentabilité,
Théorème \ref{criteria-theorem-flat-groupoid-gives-algebraic-stack}.
D'après
Champs algébriques, Lemme \ref{algebraic-lemma-map-space-into-stack},
nous obtenons un morphisme canonique
$$
f : \mathcal{Z}' \longrightarrow \mathcal{Z}
$$
qui est pleinement fidèle. Ce morphisme est donc représentable par des
espaces algébriques, voir
Champs algébriques, Lemme
\ref{algebraic-lemma-characterize-representable-by-algebraic-spaces},
et c'est un monomorphisme, voir
Lemme \ref{stacks-properties-lemma-monomorphism}.
D'après
Critères de représentabilité,
Lemme \ref{criteria-lemma-flat-quotient-flat-presentation},
le morphisme $U \to \mathcal{Z}'$ est surjectif, plat et localement de
présentation finie. Ainsi, $\mathcal{Z}'$ est un champ algébrique qui satisfait
aux conditions équivalentes du
Lemme \ref{stacks-properties-lemma-unique-point-better}.
D'après
Champs algébriques, Lemme \ref{algebraic-lemma-equivalent},
nous pouvons remplacer $\mathcal{Z}'$ par son image essentielle dans $\mathcal{Z}$.
Nous avons donc prouvé toutes les assertions du lemme, à l'exception de
l'unicité de $\mathcal{Z}' \subset \mathcal{Z}$. Supposons que
$\mathcal{Z}'' \subset \mathcal{Z}$ soit un second champ algébrique de ce type.
Alors les projections
$$
\mathcal{Z}'
\longleftarrow
\mathcal{Z}' \times_\mathcal{Z} \mathcal{Z}''
\longrightarrow
\mathcal{Z}''
$$
sont des monomorphismes. Le champ algébrique du milieu est non vide d'après le
Lemme \ref{stacks-properties-lemma-points-cartesian}.
Les deux projections sont donc des isomorphismes d'après le
Lemme \ref{stacks-properties-lemma-monomorphism-into-point},
d'où le résultat.
\end{proof}


\begin{example}
\label{stacks-properties-example-distinct}
Voici un exemple où le morphisme construit dans le
Lemme \ref{stacks-properties-lemma-improve-unique-point}
n'est pas un isomorphisme. Cet exemple montre qu'il est nécessaire d'imposer
aux gerbes résiduelles d'être localement noethériennes dans la
Définition \ref{stacks-properties-definition-residual-gerbe}.
En fait, l'exemple est même un espace algébrique !
Soit $\text{Gal}(\overline{\mathbf{Q}}/\mathbf{Q})$ le groupe de
Galois absolu de $\mathbf{Q}$ muni de la topologie profinie.
Soit
$$
U =
\Spec(\overline{\mathbf{Q}})
\times_{\Spec(\mathbf{Q})}
\Spec(\overline{\mathbf{Q}}) =
\text{Gal}(\overline{\mathbf{Q}}/\mathbf{Q}) \times
\Spec(\overline{\mathbf{Q}})
$$
(nous omettons une explication précise du sens du dernier signe égal).
Notons $G$ le groupe de Galois absolu
$\text{Gal}(\overline{\mathbf{Q}}/\mathbf{Q})$ muni de la topologie
discrète, considéré comme un schéma en groupes constant sur
$\Spec(\overline{\mathbf{Q}})$, voir
Groupoïdes, Exemple \ref{groupoids-example-constant-group}.
Alors $G$ agit librement et transitivement sur $U$.
Soit $X = U/G$, voir
Espaces, Définition \ref{spaces-definition-quotient}.
Alors $X$ est un espace algébrique réduit non noethérien ayant exactement un point.
De plus, $X$ possède un point (localement) de type fini :
$$
x :
\Spec(\overline{\mathbf{Q}}) \longrightarrow
U \longrightarrow X
$$
En effet, tout point de $U$ est en fait fermé !
Comme $X$ est un espace algébrique sur $\overline{\mathbf{Q}}$,
il s'ensuit que $x$ est un monomorphisme. Ainsi, $x$ est le morphisme
construit dans le
Lemme \ref{stacks-properties-lemma-improve-unique-point},
mais $x$ n'est pas un isomorphisme. En fait,
$\Spec(\overline{\mathbf{Q}}) \to X$ est la gerbe résiduelle de $X$ en $x$.
\end{example}

\noindent
Nous verrons plus loin que, sous des hypothèses modérées sur le champ algébrique
$\mathcal{X}$, les conditions équivalentes du lemme suivant sont satisfaites
pour tout point $x \in |\mathcal{X}|$ (voir Morphismes de champs, section
\ref{stacks-morphisms-section-existence-residual-gerbes}).

\begin{lemma}
\label{stacks-properties-lemma-residual-gerbe}
Soit $\mathcal{X}$ un champ algébrique. Soit $x \in |\mathcal{X}|$ un point.
Les assertions suivantes sont équivalentes :
\begin{enumerate}
\item il existe un champ algébrique $\mathcal{Z}$ et un monomorphisme
$\mathcal{Z} \to \mathcal{X}$ tels que $|\mathcal{Z}|$ soit réduit à un point
et que l'image de $|\mathcal{Z}|$ dans $|\mathcal{X}|$ soit $x$,
\item il existe un champ algébrique réduit $\mathcal{Z}$ et un monomorphisme
$\mathcal{Z} \to \mathcal{X}$ tels que $|\mathcal{Z}|$ soit réduit à un point
et que l'image de $|\mathcal{Z}|$ dans $|\mathcal{X}|$ soit $x$,
\item il existe un champ algébrique $\mathcal{Z}$, un monomorphisme
$f : \mathcal{Z} \to \mathcal{X}$ et un morphisme plat surjectif
$z : \Spec(k) \to \mathcal{Z}$, où $k$ est un corps tel que
$x = f(z)$.
\end{enumerate}
De plus, si ces conditions sont satisfaites, il existe une unique
sous-catégorie strictement pleine $\mathcal{Z}_x \subset \mathcal{X}$
telle que $\mathcal{Z}_x$ soit un champ algébrique réduit et localement
noethérien, et que $|\mathcal{Z}_x|$ soit réduit à un point envoyé sur $x$
par l'application $|\mathcal{Z}_x| \to |\mathcal{X}|$.
\end{lemma}

\begin{proof}
Si $\mathcal{Z} \to \mathcal{X}$ est comme en (1), alors
$\mathcal{Z}_{red} \to \mathcal{X}$ est comme en (2). (Voir la
section \ref{stacks-properties-section-reduced}
pour la notion de réduction d'un champ algébrique.)
Ainsi, (1) implique (2).
Il est immédiat que (2) implique (1).
L'équivalence de (2) et (3) résulte immédiatement du
Lemme \ref{stacks-properties-lemma-unique-point}.

\medskip\noindent
À ce stade, nous avons établi l'équivalence de (1) -- (3).
Choisissons un monomorphisme $f : \mathcal{Z} \to \mathcal{X}$ comme en (2).
Remarquons qu'il en résulte que $f$ est pleinement fidèle, voir
Lemme \ref{stacks-properties-lemma-monomorphism}.
Notons $\mathcal{Z}' \subset \mathcal{X}$ l'image essentielle du foncteur
$f$. Alors $f : \mathcal{Z} \to \mathcal{Z}'$ est une équivalence, et donc
$\mathcal{Z}'$ est un champ algébrique, voir
Champs algébriques, Lemme \ref{algebraic-lemma-equivalent}.
Appliquons le
Lemme \ref{stacks-properties-lemma-improve-unique-point}
pour obtenir une sous-catégorie strictement pleine $\mathcal{Z}_x \subset \mathcal{Z}'$
comme dans l'énoncé du lemme.
Cela prouve toutes les assertions du lemme, sauf l'unicité.

\medskip\noindent
Pour prouver l'unicité, supposons que
$\mathcal{Z}_x \subset \mathcal{X}$
et
$\mathcal{Z}'_x \subset \mathcal{X}$
soient deux sous-catégories strictement pleines comme dans l'énoncé du lemme.
Alors les projections
$$
\mathcal{Z}'_x
\longleftarrow
\mathcal{Z}'_x \times_\mathcal{X} \mathcal{Z}_x
\longrightarrow
\mathcal{Z}_x
$$
sont des monomorphismes. Le champ algébrique du milieu est non vide d'après le
Lemme \ref{stacks-properties-lemma-points-cartesian}.
Les deux projections sont donc des isomorphismes d'après le
Lemme \ref{stacks-properties-lemma-monomorphism-into-point},
d'où le résultat.
\end{proof}

\noindent
Ces explications nous permettent maintenant de poser la définition suivante.

\begin{definition}
\label{stacks-properties-definition-residual-gerbe}
Soit $\mathcal{X}$ un champ algébrique. Soit $x \in |\mathcal{X}|$.
\begin{enumerate}
\item On dit que la {\it gerbe résiduelle de $\mathcal{X}$ en $x$ existe}
si les conditions équivalentes (1), (2) et (3) du
Lemme \ref{stacks-properties-lemma-residual-gerbe}
sont satisfaites.
\item Si la gerbe résiduelle de $\mathcal{X}$ en $x$ existe, la
{\it gerbe résiduelle de $\mathcal{X}$ en $x$}\footnote{Cela entre en conflit avec
\cite{LM-B} dans l'esprit, mais non en fait. Plus précisément, au chapitre 11, ils associent
à tout point de tout champ algébrique quasi-séparé une gerbe (pas nécessairement
algébrique), qu'ils appellent la gerbe résiduelle. Nous verrons dans
Morphismes de champs, Lemme
\ref{stacks-morphisms-lemma-every-point-residual-gerbe}
que, sur un champ algébrique quasi-séparé, tout point
possède une gerbe résiduelle en notre sens, qui est alors équivalente à la leur. Pour
plus de renseignements sur ce sujet, voir
\cite[Annexe B]{rydh_etale_devissage}.}
est la sous-catégorie strictement pleine
$\mathcal{Z}_x \subset \mathcal{X}$ construite au
Lemme \ref{stacks-properties-lemma-residual-gerbe}.
\end{enumerate}
\end{definition}

\noindent
En particulier, nous savons que $\mathcal{Z}_x$ (si elle existe) est un
champ algébrique localement noethérien et réduit, et qu'il existe un
corps et un morphisme surjectif, plat et localement de présentation finie
$$
\Spec(k) \longrightarrow \mathcal{Z}_x.
$$
Nous verrons dans
Morphismes de champs, Lemme
\ref{stacks-morphisms-lemma-noetherian-singleton-stack-gerbe}
que $\mathcal{Z}_x$ est une gerbe. L'existence des gerbes résiduelles est
étudiée dans Morphismes de champs, section
\ref{stacks-morphisms-section-existence-residual-gerbes}.

\begin{example}
\label{stacks-properties-example-residual-gerbe}
Soit $X$ un schéma et soit $x \in X$ un point. Alors le monomorphisme
$x \to X$ est la gerbe résiduelle de $X$ en $x$, où, comme d'habitude, nous identifions
$x$ au schéma $x = \Spec(\kappa(x))$. Si $X$ est un espace algébrique
et $x \in |X|$, alors la gerbe résiduelle en $x$ (appelée
l'espace résiduel) existe toujours, voir
Espaces décents, section \ref{decent-spaces-section-singleton}.
\end{example}

\noindent
La gerbe résiduelle, si elle existe, est un champ algébrique régulier
d'après le lemme suivant.

\begin{lemma}
\label{stacks-properties-lemma-residual-gerbe-regular}
Un champ algébrique réduit et localement noethérien $\mathcal{Z}$ tel que
$|\mathcal{Z}|$ soit réduit à un point est régulier.
\end{lemma}

\begin{proof}
Soit $W \to \mathcal{Z}$ un morphisme lisse surjectif, où $W$ est un schéma.
Soit $k$ un corps et soit $\Spec(k) \to \mathcal{Z}$ un morphisme surjectif,
plat et localement de présentation finie (voir
Lemme \ref{stacks-properties-lemma-unique-point-better}).
L'espace algébrique $T = W \times_\mathcal{Z} \Spec(k)$ est
lisse sur $k$, donc en particulier régulier, voir
Espaces sur les corps, Lemme \ref{spaces-over-fields-lemma-smooth-regular}.
Comme $T \to W$ est localement de présentation finie, plat et surjectif, il
s'ensuit que $W$ est régulier, voir
Descente sur les espaces, Lemme \ref{spaces-descent-lemma-descend-regular}.
Par définition, cela signifie que $\mathcal{Z}$ est régulier.
\end{proof}

\begin{lemma}
\label{stacks-properties-lemma-residual-gerbe-points}
Soit $\mathcal{X}$ un champ algébrique. Soit $x \in |\mathcal{X}|$.
Supposons que la gerbe résiduelle $\mathcal{Z}_x$ de $\mathcal{X}$ existe.
Soit $f : \Spec(K) \to \mathcal{X}$ un morphisme, où $K$ est un corps,
dans la classe d'équivalence de $x$. Alors $f$ se factorise par le morphisme
d'inclusion $\mathcal{Z}_x \to \mathcal{X}$.
\end{lemma}

\begin{proof}
Choisissons un corps $k$ et un morphisme surjectif, plat et localement de présentation finie
$\Spec(k) \to \mathcal{Z}_x$. Posons
$T = \Spec(K) \times_\mathcal{X} \mathcal{Z}_x$. D'après le
Lemme \ref{stacks-properties-lemma-points-cartesian},
on voit que $T$ est non vide. Comme $\mathcal{Z}_x \to \mathcal{X}$
est un monomorphisme, on voit que $T \to \Spec(K)$ est un monomorphisme.
Ainsi, d'après
Morphismes d'espaces, Lemme
\ref{spaces-morphisms-lemma-monomorphism-toward-field},
on voit que $T = \Spec(K)$, ce qui prouve le lemme.
\end{proof}

\begin{lemma}
\label{stacks-properties-lemma-residual-gerbe-unique}
Soit $\mathcal{X}$ un champ algébrique. Soit $x \in |\mathcal{X}|$.
Soit $\mathcal{Z}$ un champ algébrique satisfaisant aux conditions équivalentes du
Lemme \ref{stacks-properties-lemma-unique-point-better},
et soit $\mathcal{Z} \to \mathcal{X}$ un monomorphisme tel que l'image
de $|\mathcal{Z}| \to |\mathcal{X}|$ soit $x$. Alors la gerbe résiduelle
$\mathcal{Z}_x$ de $\mathcal{X}$ en $x$ existe, et
$\mathcal{Z} \to \mathcal{X}$ se factorise sous la forme
$\mathcal{Z} \to \mathcal{Z}_x \to \mathcal{X}$, où la première flèche
est une équivalence.
\end{lemma}

\begin{proof}
Soit $\mathcal{Z}_x \subset \mathcal{X}$ la sous-catégorie pleine correspondant
à l'image essentielle du foncteur $\mathcal{Z} \to \mathcal{X}$.
Alors $\mathcal{Z} \to \mathcal{Z}_x$ est une équivalence, donc
$\mathcal{Z}_x$ est un champ algébrique, voir
Champs algébriques, Lemme \ref{algebraic-lemma-equivalent}.
Comme $\mathcal{Z}_x$ hérite de toutes les propriétés de $\mathcal{Z}$ par
cette équivalence, il résulte clairement de l'unicité du
Lemme \ref{stacks-properties-lemma-residual-gerbe}
que $\mathcal{Z}_x$ est la gerbe résiduelle de $\mathcal{X}$ en $x$.
\end{proof}

\begin{lemma}
\label{stacks-properties-lemma-residual-gerbe-functorial}
Soit $f : \mathcal{X} \to \mathcal{Y}$ un morphisme de champs algébriques.
Soit $x \in |\mathcal{X}|$, d'image $y \in |\mathcal{Y}|$.
Si les gerbes résiduelles $\mathcal{Z}_x \subset \mathcal{X}$
et $\mathcal{Z}_y \subset \mathcal{Y}$ de $x$ et de $y$ existent,
alors $f$ induit un diagramme commutatif
$$
\xymatrix{
\mathcal{X} \ar[d]_f & \mathcal{Z}_x \ar[l] \ar[d] \\
\mathcal{Y} & \mathcal{Z}_y \ar[l]
}
$$
\end{lemma}

\begin{proof}
Choisissons un corps $k$ et un morphisme surjectif, plat et localement de présentation finie
$\Spec(k) \to \mathcal{Z}_x$. Le morphisme
$\Spec(k) \to \mathcal{Y}$ se factorise par $\mathcal{Z}_y$ d'après le
Lemme \ref{stacks-properties-lemma-residual-gerbe-points}.
Ainsi, $\mathcal{Z}_x \times_\mathcal{Y} \mathcal{Z}_y$
est un sous-champ non vide de $\mathcal{Z}_x$,
donc égal à $\mathcal{Z}_x$ d'après le Lemme \ref{stacks-properties-lemma-monomorphism-into-point}.
\end{proof}

\begin{lemma}
\label{stacks-properties-lemma-residual-gerbe-isomorphic}
Soit $f : \mathcal{X} \to \mathcal{Y}$ un morphisme de champs algébriques.
Soit $x \in |\mathcal{X}|$ dont l'image est $y \in |\mathcal{Y}|$.
Supposons que les gerbes résiduelles $\mathcal{Z}_x \subset \mathcal{X}$
et $\mathcal{Z}_y \subset \mathcal{Y}$ de $x$ et de $y$ existent
et qu'il existe un morphisme $\Spec(k) \to \mathcal{X}$
dans la classe d'équivalence de $x$ tel que
$$
\Spec(k) \times_\mathcal{X} \Spec(k)
\longrightarrow
\Spec(k) \times_\mathcal{Y} \Spec(k)
$$
soit un isomorphisme. Alors $\mathcal{Z}_x \to \mathcal{Z}_y$
est un isomorphisme.
\end{lemma}

\begin{proof}
Soit $k'/k$ une extension de corps. Alors
$$
\Spec(k') \times_\mathcal{X} \Spec(k')
\longrightarrow
\Spec(k') \times_\mathcal{Y} \Spec(k')
$$
est le changement de base du morphisme du lemme par le
morphisme fidèlement plat $\Spec(k' \otimes k') \to \Spec(k \otimes k)$.
Ainsi, la propriété décrite dans le lemme ne dépend pas du
choix du morphisme $\Spec(k) \to \mathcal{X}$ dans la
classe d'équivalence de $x$. Ainsi, nous pouvons supposer que
$\Spec(k) \to \mathcal{Z}_x$ est surjectif, plat et
localement de présentation finie. Dans cette situation, nous avons
$$
\mathcal{Z}_x = [\Spec(k)/R]
$$
où $R = \Spec(k) \times_\mathcal{X} \Spec(k)$. Voir la
démonstration du Lemme \ref{stacks-properties-lemma-improve-unique-point}.
Comme on a aussi $R = \Spec(k) \times_\mathcal{Y} \Spec(k)$,
nous concluons que le morphisme $\mathcal{Z}_x \to \mathcal{Z}_y$
du Lemme \ref{stacks-properties-lemma-residual-gerbe-functorial}
est pleinement fidèle d'après le
Lemme sur les champs algébriques \ref{algebraic-lemma-map-space-into-stack}.
On conclut, par exemple, à l'aide du Lemme \ref{stacks-properties-lemma-residual-gerbe-unique}.
\end{proof}








\section{Dimension d'un champ}
\label{stacks-properties-section-dimension}

\noindent
Nous pouvons définir la dimension d'un champ algébrique $\mathcal{X}$
en un point $x$ à l'aide de la notion de dimension d'un espace algébrique
en un point (Propriétés des espaces, Définition
\ref{spaces-properties-definition-dimension-at-point}).
Dans le lemme suivant, le résultat peut être $\infty$, soit parce que
$\mathcal{X}$ n'est pas quasi-compact, soit parce que nous rencontrons le
phénomène décrit dans Exemples, section
\ref{examples-section-Noetherian-infinite-dimension}.

\begin{lemma}
\label{stacks-properties-lemma-dimension-at-point-well-defined}
Soit $\mathcal{X}$ un champ algébrique localement noethérien sur un schéma $S$.
Soit $x \in |\mathcal{X}|$ un point de $\mathcal{X}$.
Soit $[U/R] \to \mathcal{X}$ une présentation
(Champs algébriques, Définition \ref{algebraic-definition-presentation})
où $U$ est un schéma. Soit $u \in U$ un point dont l'image est $x$.
Soit $e : U \to R$ le morphisme « identité » et soit $s : R \to U$ le
morphisme « source », qui est un morphisme lisse d'espaces algébriques. Soit $R_u$
la fibre de $s : R \to U$ au-dessus de $u$. L'élément
$$
\dim_x(\mathcal{X}) = \dim_u(U) - \dim_{e(u)}(R_u) \in \mathbf{Z} \cup \infty
$$
est indépendant du choix de la présentation et du point $u$ au-dessus de $x$.
\end{lemma}

\begin{proof}
Puisque $R \to U$ est lisse, le schéma $R_u$ est lisse sur $\kappa(u)$
et est donc de dimension finie. D'autre part, le schéma $U$
est localement noethérien, mais cela ne garantit pas que
$\dim_u(U)$ soit finie. Ainsi, la différence est un élément de
$\mathbf{Z} \cup \{\infty\}$.

\medskip\noindent
Soient $[U'/R'] \to \mathcal{X}$ et $u' \in U'$ une seconde présentation
où $U'$ est un schéma et où l'image de $u'$ est $x$. Considérons l'espace
algébrique $P = U \times_\mathcal{X} U'$. D'après le
Lemme \ref{stacks-properties-lemma-points-cartesian}, il existe $p \in |P|$ dont les images sont
$u$ et $u'$. Comme $P \to U$ et $P \to U'$ sont lisses, on voit que
$\dim_p(P) = \dim_u(U) + \dim_p(P_u)$ et
$\dim_p(P) = \dim_{u'}(U') + \dim_p(P_{u'})$; voir le
Lemme de Morphismes d'espaces
\ref{spaces-morphisms-lemma-smoothness-dimension-spaces}.
Notons que
$$
R'_{u'} = \Spec(\kappa(u')) \times_\mathcal{X} U'
\quad\text{et}\quad
P_u = \Spec(\kappa(u)) \times_\mathcal{X} U'
$$
Représentons $p \in |P|$ par un morphisme $\Spec(\Omega) \to P$.
Comme $p$ s'envoie à la fois sur $u$ et sur $u'$, il induit un $2$-morphisme
entre les composées
$\Spec(\Omega) \to \Spec(\kappa(u')) \to \mathcal{X}$ et
$\Spec(\Omega) \to \Spec(\kappa(u)) \to \mathcal{X}$,
qui définit à son tour un isomorphisme
$$
\Spec(\Omega) \times_{\Spec(\kappa(u'))} R'_{u'}
\cong
\Spec(\Omega) \times_{\Spec(\kappa(u))} P_u
$$
d'espaces algébriques sur $\Spec(\Omega)$ envoyant le point $\Omega$-rationnel
$(1, e'(u'))$ sur $(1, p)$ (certains détails sont omis). On en déduit que
$$
\dim_{e'(u')}(R'_{u'}) = \dim_p(P_u)
$$
d'après le Lemme de Morphismes d'espaces
\ref{spaces-morphisms-lemma-dimension-fibre-after-base-change}.
Par symétrie, on a
$\dim_{e(u)}(R_u) = \dim_p(P_{u'})$.
En réunissant toutes ces égalités, on obtient l'indépendance des choix.
\end{proof}

\noindent
Nous pouvons utiliser le lemme précédent pour poser la définition suivante.

\begin{definition}
\label{stacks-properties-definition-dimension-at-point}
Soit $\mathcal{X}$ un champ algébrique localement noethérien sur un schéma $S$.
Soit $x \in |\mathcal{X}|$ un point de $\mathcal{X}$.
Soit $[U/R] \to \mathcal{X}$ une présentation
(Champs algébriques, Définition \ref{algebraic-definition-presentation})
où $U$ est un schéma,
et soit $u \in U$ un point dont l'image est $x$.
Nous définissons la {\it dimension de $\mathcal{X}$ en $x$} comme
l'élément $\dim_x(\mathcal{X}) \in \mathbf{Z} \cup \infty$
tel que 
$$
\dim_x(\mathcal{X}) = \dim_u(U)-\dim_{e(u)}(R_u).
$$
Les notations sont celles du Lemme \ref{stacks-properties-lemma-dimension-at-point-well-defined}.
\end{definition}

\noindent
La dimension d'un champ en un point coïncide avec la notion usuelle
lorsque $\mathcal{X}$ est un schéma (Topologie, Définition
\ref{topology-definition-Krull}),
ou, plus généralement, lorsque $\mathcal{X}$ est un espace algébrique localement noethérien
(Propriétés des espaces, Définition
\ref{spaces-properties-definition-dimension-at-point}).

\begin{definition}
\label{stacks-properties-definition-dimension}
Soit $S$ un schéma. Soit $\mathcal{X}$
un champ algébrique localement noethérien sur $S$.
La {\it dimension} $\dim(\mathcal{X})$ de $\mathcal{X}$ est définie par
$$
\dim(\mathcal{X}) = \sup\nolimits_{x \in |\mathcal{X}|} \dim_x(\mathcal{X})
$$
\end{definition}

\noindent
Cette définition de la dimension coïncide avec la notion usuelle
si $\mathcal{X}$ est un schéma
(Propriétés, Lemme \ref{properties-lemma-dimension})
ou un espace algébrique (Propriétés des espaces, Définition
\ref{spaces-properties-definition-dimension}).

\begin{remark}
\label{stacks-properties-remark-dimension-empty-stack}
Si $\mathcal{X}$ est un champ non vide de type fini sur un corps,
alors $\dim(\mathcal{X})$ est un entier. Pour un champ algébrique
localement noethérien quelconque $\mathcal{X}$,
$\dim(\mathcal{X})$ appartient à $Z\cup \{\pm \infty\}$,
et $\dim(\mathcal{X}) = -\infty$ si et seulement si $\mathcal{X}$
est vide.
\end{remark}

\begin{example}
\label{stacks-properties-example-dimension-quotient-stack}
Soit $X$ un schéma de type fini sur un corps $k$, et soit $G$
un schéma en groupes de type fini sur $k$ qui opère sur $X$.
Alors la dimension du champ quotient $[X/G]$ est égale à
$\dim(X)-\dim(G)$. En particulier, la dimension du champ classifiant
$BG=[\Spec(k)/G]$ est $-\dim(G)$. Ainsi, la dimension d'un champ algébrique
peut être un entier négatif, contrairement à ce qui se produit pour les schémas
ou les espaces algébriques.
\end{example}






\section{Irréductibilité locale}
\label{stacks-properties-section-irreducible-local-ring}

\noindent
Nous avons défini le nombre de branches géométriques d'un schéma en un point
dans Propriétés, section \ref{properties-section-unibranch},
et celui d'un espace algébrique en un point dans Propriétés des espaces, section
\ref{spaces-properties-section-irreducible-local-ring}.
Soit $n \in \mathbf{N}$. Pour un anneau local $A$, posons
$$
P_n(A) = \text{le nombre de branches géométriques de }A\text{ vaut }n
$$
Pour un homomorphisme lisse d'anneaux $A \to B$ et un idéal premier $\mathfrak q$
de $B$ situé au-dessus de $\mathfrak p$ de $A$, on a
$$
P_n(A_\mathfrak p) \Leftrightarrow P_n(B_\mathfrak q)
$$
d'après le Lemme de Compléments d'algèbre
\ref{more-algebra-lemma-invariance-number-branches-smooth}.
Comme dans Propriétés des espaces, Remarque
\ref{spaces-properties-remark-list-properties-local-ring-local-etale-topology},
nous pouvons utiliser $P_n$ pour définir une propriété locale pour la topologie \'etale $\mathcal{P}_n$
des germes $(U, u)$ de schémas en posant
$\mathcal{P}_n(U, u) = P_n(\mathcal{O}_{U, u})$.
La propriété correspondante $\mathcal{P}_n$
d'un espace algébrique $X$ en un point $x$
(voir Propriétés des espaces, Définition
\ref{spaces-properties-definition-property-at-point})
est précisément la propriété
``le nombre de branches géométriques de $X$ en $x$ est $n$''; voir
Propriétés des espaces, Définition
\ref{spaces-properties-definition-number-of-geometric-branches}.
De plus, la propriété $\mathcal{P}_n$ est locale pour la topologie lisse; voir
Descente, Définition \ref{descent-definition-local-at-point}.
Cela résulte soit de l'équivalence affichée ci-dessus,
soit du Lemme de Compléments sur les morphismes
\ref{more-morphisms-lemma-number-of-branches-and-smooth}.
Ainsi, la Définition \ref{stacks-properties-definition-property-at-point}
s'applique et nous obtenons une notion pour les champs algébriques en un point.

\begin{definition}
\label{stacks-properties-definition-number-of-geometric-branches}
Soit $\mathcal{X}$ un champ algébrique. Soit $x \in |\mathcal{X}|$.
\begin{enumerate}
\item Le {\it nombre de branches géométriques de $\mathcal{X}$ en $x$}
vaut soit $n \in \mathbf{N}$ si les conditions équivalentes du
Lemme \ref{stacks-properties-lemma-local-source-target-at-point} sont satisfaites pour
$\mathcal{P}_n$ définie ci-dessus, soit $\infty$.
\item Nous disons que $\mathcal{X}$ est {\it géométriquement unibranche en $x$}
si le nombre de branches géométriques de $\mathcal{X}$ en $x$ vaut $1$.
\end{enumerate}
\end{definition}






\section{Conditions de finitude et points}
\label{stacks-properties-section-points-conditions}

\noindent
Cette section est l'analogue de
Espaces décents, section \ref{decent-spaces-section-points-monomorphisms}
pour les points des champs algébriques.

\begin{lemma}
\label{stacks-properties-lemma-UR-quasi-compact-above-x}
Soit $\mathcal{X}$ un champ algébrique. Soit $x \in |\mathcal{X}|$
un point. Les conditions suivantes sont équivalentes :
\begin{enumerate}
\item un morphisme $\Spec(k) \to \mathcal{X}$ dans la classe d'équivalence
de $x$ est quasi-compact, et
\item tout morphisme $\Spec(k) \to \mathcal{X}$ dans la classe d'équivalence
de $x$ est quasi-compact.
\end{enumerate}
\end{lemma}

\begin{proof}
Soit $\Spec(k) \to \mathcal{X}$ dans la classe d'équivalence
de $x$. Soit $k'/k$ une extension de corps.
Nous devons alors montrer que $\Spec(k) \to \mathcal{X}$ est
quasi-compact si et seulement si $\Spec(k') \to \mathcal{X}$
est quasi-compact. Cela résulte du
Lemme de Morphismes d'espaces
\ref{spaces-morphisms-lemma-surjection-from-quasi-compact}
et du principe donné par le Lemme de Champs algébriques
\ref{algebraic-lemma-representable-transformations-property-implication}.
\end{proof}

\noindent
On dit parfois qu'un point $x \in |\mathcal{X}|$ satisfaisant
aux conditions équivalentes du Lemme \ref{stacks-properties-lemma-UR-quasi-compact-above-x}
est un ``{\it point quasi-compact}''.













% OK, commencer ici.
%

\chapter{Morphismes de champs algébriques}



\phantomsection
\label{stacks-morphisms-section-phantom}


\section{Introduction}
\label{stacks-morphisms-section-introduction}

\noindent
Dans ce chapitre, nous introduisons certains types de morphismes de champs
algébriques. Une référence pour le cas des champs algébriques quasi-séparés à
diagonale représentable est \cite{LM-B}.

\medskip\noindent
Le but est d'étendre à des morphismes de champs algébriques la définition de
chacun des types de morphismes d'espaces algébriques. Chaque cas présente de
légères différences, et il semble préférable de les traiter séparément.

\medskip\noindent
Pour les morphismes de champs algébriques représentables par des espaces
algébriques, nous avons déjà défini un grand nombre de types de morphismes, voir
Propriétés des champs,
section \ref{stacks-properties-section-properties-morphisms}.
Dans chaque cas correspondant du présent chapitre, nous devons vérifier que la
définition générale est compatible avec celle qui y a été donnée.




\section{Conventions et abus de langage}
\label{stacks-morphisms-section-conventions}

\noindent
Nous continuons d'employer les conventions et les abus de langage introduits
dans Propriétés des champs, section
\ref{stacks-properties-section-conventions}.




\section{Propriétés des diagonales}
\label{stacks-morphisms-section-diagonals}

\noindent
La diagonale d'un champ algébrique est étroitement liée aux faisceaux
$\mathit{Isom}$, voir
Champs algébriques, Lemme \ref{algebraic-lemma-representable-diagonal}.
Par la seconde propriété qui définit un champ algébrique, ces faisceaux
$\mathit{Isom}$ sont toujours des espaces algébriques.

\begin{lemma}
\label{stacks-morphisms-lemma-isom-locally-finite-type}
Soit $\mathcal{X}$ un champ algébrique.
Soit $T$ un schéma et soient $x,y$ des objets de la catégorie fibre de
$\mathcal{X}$ au-dessus de $T$. Alors le morphisme
$\mathit{Isom}_\mathcal{X}(x,y) \to T$ est localement de type fini.
\end{lemma}

\begin{proof}
Par Champs algébriques, Lemme
\ref{algebraic-lemma-stack-presentation}, nous pouvons supposer que
$\mathcal{X}=[U/R]$ pour un groupoïde lisse en espaces algébriques.
Par Descente sur les espaces,
Lemme \ref{spaces-descent-lemma-descending-property-locally-finite-type},
il suffit de vérifier la propriété localement pour la topologie fppf sur $T$.
Nous pouvons donc supposer que $x,y$ proviennent de morphismes
$x',y':T\to U$. Par Groupoïdes en espaces,
Lemme \ref{spaces-groupoids-lemma-quotient-stack-morphisms}, nous avons alors
$\mathit{Isom}_\mathcal{X}(x,y)=T\times_{(y',x'),U\times_S U}R$.
Il suffit donc de prouver que $R\to U\times_S U$ est localement de type fini.
Cela résulte du fait que le composé
$s:R\to U\times_S U\to U$ est lisse (donc localement de type fini, voir
Morphismes d'espaces, Lemmes
\ref{spaces-morphisms-lemma-smooth-locally-finite-presentation} et
\ref{spaces-morphisms-lemma-finite-presentation-finite-type})
et Morphismes d'espaces, Lemme
\ref{spaces-morphisms-lemma-permanence-finite-type}.
\end{proof}

\begin{lemma}
\label{stacks-morphisms-lemma-isom-pseudo-torsor-aut}
Soit $\mathcal{X}$ un champ algébrique.
Soit $T$ un schéma et soient $x,y$ des objets de la catégorie fibre de
$\mathcal{X}$ au-dessus de $T$. Alors
\begin{enumerate}
\item $\mathit{Isom}_\mathcal{X}(y,y)$ est un espace algébrique en groupes
sur $T$, et
\item $\mathit{Isom}_\mathcal{X}(x,y)$ est un pseudo-torseur sous
$\mathit{Isom}_\mathcal{X}(y,y)$ sur $T$.
\end{enumerate}
\end{lemma}

\begin{proof}
Voir Groupoïdes en espaces,
Définitions \ref{spaces-groupoids-definition-group-space} et
\ref{spaces-groupoids-definition-pseudo-torsor}.
Le lemme résulte immédiatement du fait que $\mathcal{X}$ est un champ en
groupoïdes.
\end{proof}

\noindent
Soit $f:\mathcal{X}\to\mathcal{Y}$ un morphisme de champs algébriques.
La {\it diagonale de $f$} est le morphisme
$$
\Delta_f :
\mathcal{X}
\longrightarrow
\mathcal{X} \times_\mathcal{Y} \mathcal{X}
$$
Voici deux propriétés que possède tout morphisme diagonal.

\begin{lemma}
\label{stacks-morphisms-lemma-properties-diagonal}
\begin{slogan}
Les diagonales des morphismes de champs algébriques sont représentables par
des espaces algébriques et localement de type fini.
\end{slogan}
Soit $f:\mathcal{X}\to\mathcal{Y}$ un morphisme de champs algébriques.
Alors
\begin{enumerate}
\item $\Delta_f$ est représentable par des espaces algébriques, et
\item $\Delta_f$ est localement de type fini.
\end{enumerate}
\end{lemma}

\begin{proof}
Soit $T$ un schéma et soit
$a:T\to\mathcal{X}\times_\mathcal{Y}\mathcal{X}$ un morphisme. Par la
définition du produit fibré et le lemme de $2$-Yoneda, le morphisme $a$ est
donné par un triplet $a=(x,x',\alpha)$, où $x,x'$ sont des objets de
$\mathcal{X}$ au-dessus de $T$ et
$\alpha:f(x)\to f(x')$ est un morphisme dans la catégorie fibre de
$\mathcal{Y}$ au-dessus de $T$. Par la définition d'un champ algébrique, les
faisceaux $\mathit{Isom}_\mathcal{X}(x,x')$ et
$\mathit{Isom}_\mathcal{Y}(f(x),f(x'))$ sont des espaces algébriques sur
$T$. Dans ce langage, $\alpha$ définit une section du morphisme
$\mathit{Isom}_\mathcal{Y}(f(x),f(x'))\to T$. Un point à valeurs dans $T'$ de
$\mathcal{X}\times_{\mathcal{X}\times_\mathcal{Y}\mathcal{X},a}T$, où
$T'\to T$ est un schéma au-dessus de $T$, est la même chose qu'un
isomorphisme $x|_{T'}\to x'|_{T'}$ dont l'image par
$f$ est $\alpha|_{T'}$. Nous voyons donc que
\begin{equation}
\label{stacks-morphisms-equation-diagonal}
\vcenter{
\xymatrix{
\mathcal{X} \times_{\mathcal{X} \times_\mathcal{Y} \mathcal{X}, a} T
\ar[d] \ar[r] &
\mathit{Isom}_\mathcal{X}(x, x') \ar[d] \\
T\ar[r]^-\alpha &
\mathit{Isom}_\mathcal{Y}(f(x), f(x'))
}
}
\end{equation}
est un carré cartésien de faisceaux sur $T$. En particulier,
$\mathcal{X} \times_{\mathcal{X} \times_\mathcal{Y} \mathcal{X}, a} T$
est un espace algébrique, ce qui démontre l'assertion (1) du lemme.

\medskip\noindent
Pour démontrer la seconde assertion, nous devons montrer que la flèche
verticale de gauche du Diagramme (\ref{stacks-morphisms-equation-diagonal}) est localement de
type fini. Par le Lemme \ref{stacks-morphisms-lemma-isom-locally-finite-type}, l'espace
algébrique $\mathit{Isom}_\mathcal{X}(x,x')$ est localement de type fini sur
$T$. La flèche verticale de droite du Diagramme
(\ref{stacks-morphisms-equation-diagonal}) est donc localement de type fini, voir
Morphismes d'espaces, Lemme
\ref{spaces-morphisms-lemma-permanence-finite-type}. Nous concluons par
Morphismes d'espaces, Lemme
\ref{spaces-morphisms-lemma-base-change-finite-type}.
\end{proof}

\begin{lemma}
\label{stacks-morphisms-lemma-properties-diagonal-representable}
Soit $f:\mathcal{X}\to\mathcal{Y}$ un morphisme de champs algébriques
représentable par des espaces algébriques. Alors
\begin{enumerate}
\item $\Delta_f$ est représentable (par des schémas),
\item $\Delta_f$ est localement de type fini,
\item $\Delta_f$ est un monomorphisme,
\item $\Delta_f$ est séparé, et
\item $\Delta_f$ est localement quasi-fini.
\end{enumerate}
\end{lemma}

\begin{proof}
Nous avons déjà vu au Lemme \ref{stacks-morphisms-lemma-properties-diagonal} que $\Delta_f$
est représentable par des espaces algébriques. Les assertions (2) -- (5) ont
donc un sens, voir Propriétés des champs,
section \ref{stacks-properties-section-properties-morphisms}. Le Lemme
\ref{stacks-morphisms-lemma-properties-diagonal} garantit aussi (2). Soit
$T\to\mathcal{X}\times_\mathcal{Y}\mathcal{X}$ un morphisme, et considérons
le Diagramme (\ref{stacks-morphisms-equation-diagonal}). Par Champs algébriques, Lemme
\ref{algebraic-lemma-criterion-map-representable-spaces-fibred-in-groupoids},
la flèche verticale de droite est injective comme application de faisceaux,
c'est-à-dire qu'elle est un monomorphisme d'espaces algébriques. Par suite, le
morphisme
$T\times_{\mathcal{X}\times_\mathcal{Y}\mathcal{X}}\mathcal{X}\to T$
est lui aussi un monomorphisme. Ainsi, (3) est satisfaite. Nous savons déjà que
$T\times_{\mathcal{X}\times_\mathcal{Y}\mathcal{X}}\mathcal{X}\to T$ est
localement de type fini. Morphismes d'espaces, Lemme
\ref{spaces-morphisms-lemma-monomorphism-loc-finite-type-loc-quasi-finite},
permet donc de conclure que
$T\times_{\mathcal{X}\times_\mathcal{Y}\mathcal{X}}\mathcal{X}\to T$
est localement quasi-fini et séparé.
Cela démontre (4) et (5). Enfin, Morphismes d'espaces, Proposition
\ref{spaces-morphisms-proposition-locally-quasi-finite-separated-over-scheme},
implique que
$T\times_{\mathcal{X}\times_\mathcal{Y}\mathcal{X}}\mathcal{X}$ est un
schéma, ce qui démontre (1).
\end{proof}

\begin{lemma}
\label{stacks-morphisms-lemma-representable-separated-diagonal-closed}
Soit $f:\mathcal{X}\to\mathcal{Y}$ un morphisme de champs algébriques
représentable par des espaces algébriques. Les assertions suivantes sont
équivalentes :
\begin{enumerate}
\item $f$ est séparé ;
\item $\Delta_f$ est une immersion fermée ;
\item $\Delta_f$ est propre ;
\item $\Delta_f$ est universellement fermé.
\end{enumerate}
\end{lemma}

\begin{proof}
Les énoncés « $f$ est séparé », « $\Delta_f$ est une immersion fermée »,
« $\Delta_f$ est universellement fermé » et « $\Delta_f$ est propre » font
référence aux notions définies dans Propriétés des champs,
section \ref{stacks-properties-section-properties-morphisms}. Choisissons un
schéma $V$ et un morphisme lisse surjectif $V\to\mathcal{Y}$. Posons
$U=\mathcal{X}\times_\mathcal{Y}V$, qui est par hypothèse un espace
algébrique ; le morphisme $U\to\mathcal{X}$ est lisse et surjectif. Par
Catégories, Lemme \ref{categories-lemma-base-change-diagonal}, et Propriétés
des champs, Lemme \ref{stacks-properties-lemma-check-property-covering}, pour
toute propriété $P$ considérée dans ce dernier lemme, $\Delta_f$ possède $P$
si et seulement si $\Delta_{U/V}:U\to U\times_VU$ possède $P$. L'équivalence
de (2), (3) et (4) résulte donc de Morphismes d'espaces, Lemme
\ref{spaces-morphisms-lemma-separated-diagonal-proper}, appliqué à $U\to V$.
De plus, si (1) est satisfaite, $U\to V$ est séparé, donc
$\Delta_{U/V}$ est une immersion fermée : (2) est satisfaite. Supposons enfin
(2). Soient $T$ un schéma et $a:T\to\mathcal{Y}$ un morphisme. Posons
$T'=\mathcal{X}\times_\mathcal{Y}T$. Pour démontrer (1), il faut montrer que
le morphisme d'espaces algébriques $T'\to T$ est séparé. En appliquant de
nouveau Catégories, Lemme \ref{categories-lemma-base-change-diagonal}, nous
voyons que $\Delta_{T'/T}$ est le changement de base de $\Delta_f$.
L'hypothèse (2) implique donc que $\Delta_{T'/T}$ est une immersion fermée,
et ainsi que $T'\to T$ est séparé, comme voulu.
\end{proof}

\begin{lemma}
\label{stacks-morphisms-lemma-representable-quasi-separated-diagonal-quasi-compact}
Soit $f:\mathcal{X}\to\mathcal{Y}$ un morphisme de champs algébriques
représentable par des espaces algébriques. Les assertions suivantes sont
équivalentes :
\begin{enumerate}
\item $f$ est quasi-séparé ;
\item $\Delta_f$ est quasi-compact ;
\item $\Delta_f$ est de type fini.
\end{enumerate}
\end{lemma}

\begin{proof}
Les énoncés « $f$ est quasi-séparé », « $\Delta_f$ est quasi-compact » et
« $\Delta_f$ est de type fini » font référence aux notions définies dans
Propriétés des champs,
section \ref{stacks-properties-section-properties-morphisms}. Remarquons que
(2) et (3) sont équivalentes, puisque $\Delta_f$ est localement de type fini
par le Lemme \ref{stacks-morphisms-lemma-properties-diagonal-representable} (et Champs
algébriques, Lemme
\ref{algebraic-lemma-representable-transformations-property-implication}).
Choisissons un schéma $V$ et un morphisme lisse surjectif
$V\to\mathcal{Y}$. Posons $U=\mathcal{X}\times_\mathcal{Y}V$, qui est par
hypothèse un espace algébrique ; le morphisme $U\to\mathcal{X}$ est lisse et
surjectif. Par Catégories, Lemme
\ref{categories-lemma-base-change-diagonal}, et Propriétés des champs, Lemme
\ref{stacks-properties-lemma-check-property-covering}, $\Delta_f$ est
quasi-compact si et seulement si
$\Delta_{U/V}:U\to U\times_VU$ est quasi-compact. Si (1) est satisfaite,
$U\to V$ est quasi-séparé, donc $\Delta_{U/V}$ est quasi-compact : (2) est
satisfaite. Supposons (2). Soient $T$ un schéma et
$a:T\to\mathcal{Y}$ un morphisme. Posons
$T'=\mathcal{X}\times_\mathcal{Y}T$. Pour démontrer (1), il faut montrer que
le morphisme d'espaces algébriques $T'\to T$ est quasi-séparé. En appliquant
de nouveau Catégories, Lemme \ref{categories-lemma-base-change-diagonal}, nous
voyons que $\Delta_{T'/T}$ est le changement de base de $\Delta_f$.
L'hypothèse (2) implique donc que $\Delta_{T'/T}$ est quasi-compact, et ainsi
que $T'\to T$ est quasi-séparé, comme voulu.
\end{proof}

\begin{lemma}
\label{stacks-morphisms-lemma-representable-locally-separated-diagonal-immersion}
Soit $f:\mathcal{X}\to\mathcal{Y}$ un morphisme de champs algébriques
représentable par des espaces algébriques. Les assertions suivantes sont
équivalentes :
\begin{enumerate}
\item $f$ est localement séparé ;
\item $\Delta_f$ est une immersion.
\end{enumerate}
\end{lemma}

\begin{proof}
Les énoncés « $f$ est localement séparé » et « $\Delta_f$ est une immersion »
font référence aux notions définies dans Propriétés des champs,
section \ref{stacks-properties-section-properties-morphisms}.
Démonstration omise. Indication : raisonner comme dans les démonstrations des
Lemmes \ref{stacks-morphisms-lemma-representable-separated-diagonal-closed} et
\ref{stacks-morphisms-lemma-representable-quasi-separated-diagonal-quasi-compact}.
\end{proof}






\section{Axiomes de séparation}
\label{stacks-morphisms-section-separated}

\noindent
Soit $\mathcal{X}=[U/R]$ une présentation d'un champ algébrique. Les
propriétés de la diagonale de $\mathcal{X}$ sur $S$ sont alors celles du
morphisme $j:R\to U\times_SU$. Par exemple, si
$\mathcal{X}=[S/G]$ pour un groupe lisse $G$ en espaces algébriques sur $S$,
alors $j$ est le morphisme structural $G\to S$. La diagonale elle-même n'est
donc pas automatiquement séparée (contrairement à ce qui se produit pour les
schémas et les espaces algébriques). Dire que $[S/G]$ est quasi-séparé sur
$S$ devrait certainement impliquer que $G\to S$ est quasi-compact ; nous
hésitons toutefois à dire que $[S/G]$ est quasi-séparé sur $S$ sans exiger aussi que
$G\to S$ soit quasi-séparé. Autrement dit, imposer seulement que le morphisme
diagonal soit quasi-compact ne correspond pas vraiment à l'intuition que nous
avons d'un « champ algébrique quasi-séparé » : il faut aussi exiger que la
diagonale elle-même soit quasi-séparée.

\medskip\noindent
Qu'en est-il des « champs algébriques séparés » ? Nous avons vu dans
Morphismes d'espaces,
Lemme \ref{spaces-morphisms-lemma-separated-diagonal-proper}, qu'un espace
algébrique est séparé si et seulement si sa diagonale est propre. C'est aussi
la condition habituellement employée pour définir les champs algébriques
séparés. Dans l'exemple $[S/G]\to S$ ci-dessus, cela signifie que $G\to S$
est un schéma en groupes propre. Ainsi, les champs algébriques de la forme
$[\Spec(k)/E]$ sont propres sur $k$ lorsque $E$ est une courbe elliptique sur
$k$ (insérer ici une référence ultérieure). Dans certaines situations, il peut
être plus naturel de supposer la diagonale finie.

\begin{definition}
\label{stacks-morphisms-definition-separated}
Soit $f:\mathcal{X}\to\mathcal{Y}$ un morphisme de champs algébriques.
\begin{enumerate}
\item Nous disons que $f$ est {\it DM} si $\Delta_f$ est non
ramifié\footnote{Les lettres DM signifient Deligne--Mumford. Si $f$ est DM,
alors, pour tout schéma $T$ et tout morphisme $T\to\mathcal{Y}$, le produit fibré
$\mathcal{X}_T = \mathcal{X} \times_\mathcal{Y} T$
est un champ algébrique sur $T$ dont la diagonale est non ramifiée,
c'est-à-dire que $\mathcal{X}_T$ est DM. Il s'ensuit que $\mathcal{X}_T$ est
un champ de Deligne--Mumford, voir le Théorème \ref{stacks-morphisms-theorem-DM}. Autrement dit,
un morphisme DM est un morphisme dont les « fibres » sont des champs de
Deligne--Mumford. Cela motive, espérons-le, au moins la terminologie.}.
\item Nous disons que $f$ est {\it quasi-DM} si $\Delta_f$ est localement
quasi-fini\footnote{Si $f$ est quasi-DM, alors les « fibres »
$\mathcal{X}_T$ de $\mathcal{X}\to\mathcal{Y}$ sont quasi-DM. Un champ
algébrique $\mathcal{X}$ est quasi-DM si et seulement s'il existe un schéma
$U$ et un morphisme plat surjectif $U\to\mathcal{X}$, de présentation finie
et localement quasi-fini, voir le Théorème \ref{stacks-morphisms-theorem-quasi-DM}. Notons la
ressemblance avec la propriété d'être de Deligne--Mumford, définie par
l'existence d'un recouvrement étale par un schéma.}.
\item Nous disons que $f$ est {\it séparé} si $\Delta_f$ est propre.
\item Nous disons que $f$ est {\it quasi-séparé} si $\Delta_f$ est
quasi-compact et quasi-séparé.
\end{enumerate}
\end{definition}

\noindent
Dans cette définition, nous utilisons le fait que $\Delta_f$ est représentable
par des espaces algébriques ainsi que Propriétés des champs,
section \ref{stacks-properties-section-properties-morphisms}, afin de donner
un sens aux conditions imposées à $\Delta_f$. Notons que ces définitions sont
compatibles avec les notions déjà existantes lorsque $f$ est représentable par
des espaces algébriques, voir les Lemmes
\ref{stacks-morphisms-lemma-representable-quasi-separated-diagonal-quasi-compact} et
\ref{stacks-morphisms-lemma-representable-separated-diagonal-closed}. On peut caractériser ces
conditions d'une manière intéressante en considérant les diagonales d'ordre
supérieur, voir le Lemme \ref{stacks-morphisms-lemma-definition-separated}.

\begin{definition}
\label{stacks-morphisms-definition-absolute-separated}
Soit $\mathcal{X}$ un champ algébrique sur le schéma de base $S$. Notons
$p:\mathcal{X}\to S$ le morphisme structural.
\begin{enumerate}
\item Nous disons que $\mathcal{X}$ est {\it DM sur $S$} si
$p:\mathcal{X}\to S$ est DM.
\item Nous disons que $\mathcal{X}$ est {\it quasi-DM sur $S$} si
$p:\mathcal{X}\to S$ est quasi-DM.
\item Nous disons que $\mathcal{X}$ est {\it séparé sur $S$} si
$p:\mathcal{X}\to S$ est séparé.
\item Nous disons que $\mathcal{X}$ est {\it quasi-séparé sur $S$} si
$p:\mathcal{X}\to S$ est quasi-séparé.
\item Nous disons que $\mathcal{X}$ est {\it DM} si $\mathcal{X}$ est
DM\footnote{Le Théorème \ref{stacks-morphisms-theorem-DM} montre que cela équivaut à dire que
$\mathcal{X}$ est un champ de Deligne--Mumford.} sur $\Spec(\mathbf{Z})$.
\item Nous disons que $\mathcal{X}$ est {\it quasi-DM} si $\mathcal{X}$ est
quasi-DM sur $\Spec(\mathbf{Z})$.
\item Nous disons que $\mathcal{X}$ est {\it séparé} si $\mathcal{X}$ est
séparé sur $\Spec(\mathbf{Z})$.
\item Nous disons que $\mathcal{X}$ est {\it quasi-séparé} si
$\mathcal{X}$ est quasi-séparé sur $\Spec(\mathbf{Z})$.
\end{enumerate}
Dans les quatre dernières définitions, nous considérons $\mathcal{X}$ comme un
champ algébrique sur $\Spec(\mathbf{Z})$ au moyen de Champs algébriques,
Définition \ref{algebraic-definition-viewed-as}.
\end{definition}

\noindent
Dans chaque cas, nous disposons donc d'une notion absolue et d'une notion
relative au schéma de base fixé (que notre abus de notation introduit dans
Propriétés des champs, section \ref{stacks-properties-section-conventions},
conduit généralement à ne pas mentionner). Nous verrons que
(1) $\Leftrightarrow$ (5) et (2) $\Leftrightarrow$ (6) au Lemme
\ref{stacks-morphisms-lemma-separated-implies-morphism-separated}. Nous consacrons quelque
temps à démontrer des résultats usuels sur ces notions.

\begin{lemma}
\label{stacks-morphisms-lemma-trivial-implications}
Soit $f : \mathcal{X} \to \mathcal{Y}$ un morphisme de champs algébriques.
\begin{enumerate}
\item Si $f$ est séparé, alors $f$ est quasi-séparé.
\item Si $f$ est DM, alors $f$ est quasi-DM.
\item Si $f$ est représentable par des espaces algébriques, alors $f$ est DM.
\end{enumerate}
\end{lemma}

\begin{proof}
Pour (1), remarquons qu'un morphisme propre d'espaces algébriques est
quasi-compact et quasi-séparé, voir Morphismes d'espaces, Définition
\ref{spaces-morphisms-definition-proper}. Pour (2), remarquons qu'un
morphisme non ramifié d'espaces algébriques est localement quasi-fini, voir
Morphismes d'espaces, Lemme
\ref{spaces-morphisms-lemma-unramified-quasi-finite}.
Enfin, (3) résulte du Lemme \ref{stacks-morphisms-lemma-properties-diagonal-representable}.
\end{proof}

\begin{lemma}
\label{stacks-morphisms-lemma-base-change-separated}
Tous les axiomes de séparation énumérés dans la Définition
\ref{stacks-morphisms-definition-separated} sont stables par changement de base.
\end{lemma}

\begin{proof}
Soient $f : \mathcal{X} \to \mathcal{Y}$ et
$\mathcal{Y}' \to \mathcal{Y}$ des morphismes de champs algébriques.
Soit $f' : \mathcal{Y}' \times_\mathcal{Y} \mathcal{X} \to \mathcal{Y}'$
le changement de base de $f$ par $\mathcal{Y}' \to \mathcal{Y}$.
Alors $\Delta_{f'}$ est le changement de base de $\Delta_f$ par le morphisme
$\mathcal{X}' \times_{\mathcal{Y}'} \mathcal{X}' \to
\mathcal{X} \times_\mathcal{Y} \mathcal{X}$, voir Catégories, Lemme
\ref{categories-lemma-base-change-diagonal}. D'après les résultats de
Propriétés des champs,
section \ref{stacks-properties-section-properties-morphisms}, chacune des
propriétés de la diagonale employées dans la Définition
\ref{stacks-morphisms-definition-separated} est stable par changement de base. Le lemme en
résulte.
\end{proof}

\begin{lemma}
\label{stacks-morphisms-lemma-check-separated-covering}
Soit $f : \mathcal{X} \to \mathcal{Y}$ un morphisme de champs algébriques.
Soit $W \to \mathcal{Y}$ un morphisme surjectif, plat et localement de
présentation finie, où $W$ est un espace algébrique. Si le changement de base
$W \times_\mathcal{Y} \mathcal{X} \to W$ possède l'une des propriétés de
séparation de la Définition \ref{stacks-morphisms-definition-separated}, alors $f$ la possède
également.
\end{lemma}

\begin{proof}
Notons $g : W \times_\mathcal{Y} \mathcal{X} \to W$ le changement de base.
Alors $\Delta_g$ est le changement de base de $\Delta_f$ par le morphisme
$q : W \times_\mathcal{Y} (\mathcal{X} \times_\mathcal{Y} \mathcal{X})
\to \mathcal{X} \times_\mathcal{Y} \mathcal{X}$. Comme $q$ est le changement
de base de $W \to \mathcal{Y}$, nous voyons que $q$ est représentable par des espaces
algébriques, surjectif, plat et localement de présentation finie. Le résultat
découle donc de Propriétés des champs, Lemme
\ref{stacks-properties-lemma-check-property-weak-covering}.
\end{proof}

\begin{lemma}
\label{stacks-morphisms-lemma-change-of-base-separated}
Soit $S$ un schéma. La propriété d'être quasi-DM sur $S$, quasi-séparé sur
$S$ ou séparé sur $S$ (voir la Définition
\ref{stacks-morphisms-definition-absolute-separated}) est stable par changement du schéma de
base, voir Champs algébriques, Définition
\ref{algebraic-definition-change-of-base}.
\end{lemma}

\begin{proof}
Cela résulte immédiatement du Lemme \ref{stacks-morphisms-lemma-base-change-separated}.
\end{proof}

\begin{lemma}
\label{stacks-morphisms-lemma-fibre-product-after-map}
Soient $f : \mathcal{X} \to \mathcal{Z}$, $g : \mathcal{Y} \to \mathcal{Z}$
et $\mathcal{Z} \to \mathcal{T}$ des morphismes de champs algébriques.
Considérons le morphisme induit
$i : \mathcal{X} \times_\mathcal{Z} \mathcal{Y} \to
\mathcal{X} \times_\mathcal{T} \mathcal{Y}$.
Alors
\begin{enumerate}
\item $i$ est représentable par des espaces algébriques et localement de type
fini ;
\item si $\Delta_{\mathcal{Z}/\mathcal{T}}$ est quasi-séparé, alors $i$ est
quasi-séparé ;
\item si $\Delta_{\mathcal{Z}/\mathcal{T}}$ est séparé, alors $i$ est séparé ;
\item si $\mathcal{Z} \to \mathcal{T}$ est DM, alors $i$ est non ramifié ;
\item si $\mathcal{Z} \to \mathcal{T}$ est quasi-DM, alors $i$ est localement
quasi-fini ;
\item si $\mathcal{Z} \to \mathcal{T}$ est séparé, alors $i$ est propre ;
\item si $\mathcal{Z} \to \mathcal{T}$ est quasi-séparé, alors $i$ est
quasi-compact et quasi-séparé.
\end{enumerate}
\end{lemma}

\begin{proof}
Le diagramme suivant
$$
\xymatrix{
\mathcal{X} \times_\mathcal{Z} \mathcal{Y} \ar[r]_i \ar[d] &
\mathcal{X} \times_\mathcal{T} \mathcal{Y} \ar[d] \\
\mathcal{Z} \ar[r]^-{\Delta_{\mathcal{Z}/\mathcal{T}}} \ar[r] &
\mathcal{Z} \times_\mathcal{T} \mathcal{Z}
}
$$
est un diagramme de $2$-produit fibré, voir Catégories, Lemme
\ref{categories-lemma-fibre-product-after-map}. Par conséquent, $i$ est le
changement de base du morphisme diagonal
$\Delta_{\mathcal{Z}/\mathcal{T}}$. Le lemme résulte donc du Lemme
\ref{stacks-morphisms-lemma-properties-diagonal} et des résultats de Propriétés des champs,
section \ref{stacks-properties-section-properties-morphisms}.
\end{proof}

\begin{lemma}
\label{stacks-morphisms-lemma-semi-diagonal}
Soit $\mathcal{T}$ un champ algébrique. Soit
$g : \mathcal{X} \to \mathcal{Y}$ un morphisme de champs algébriques sur
$\mathcal{T}$. Considérons le graphe
$i : \mathcal{X} \to \mathcal{X} \times_\mathcal{T} \mathcal{Y}$ de $g$. Alors
\begin{enumerate}
\item $i$ est représentable par des espaces algébriques et localement de type
fini ;
\item si $\mathcal{Y} \to \mathcal{T}$ est DM, alors $i$ est non ramifié ;
\item si $\mathcal{Y} \to \mathcal{T}$ est quasi-DM, alors $i$ est localement
quasi-fini ;
\item si $\mathcal{Y} \to \mathcal{T}$ est séparé, alors $i$ est propre ;
\item si $\mathcal{Y} \to \mathcal{T}$ est quasi-séparé, alors $i$ est
quasi-compact et quasi-séparé.
\end{enumerate}
\end{lemma}

\begin{proof}
Il s'agit du cas particulier du Lemme \ref{stacks-morphisms-lemma-fibre-product-after-map}
appliqué au morphisme
$\mathcal{X} = \mathcal{X} \times_\mathcal{Y} \mathcal{Y} \to
\mathcal{X} \times_\mathcal{T} \mathcal{Y}$.
\end{proof}

\begin{lemma}
\label{stacks-morphisms-lemma-section-immersion}
Soit $f : \mathcal{X} \to \mathcal{T}$ un morphisme de champs algébriques.
Soit $s : \mathcal{T} \to \mathcal{X}$ un morphisme tel que
$f \circ s$ soit $2$-isomorphe à $\text{id}_\mathcal{T}$. Alors
\begin{enumerate}
\item $s$ est représentable par des espaces algébriques et localement de type
fini ;
\item si $f$ est DM, alors $s$ est non ramifié ;
\item si $f$ est quasi-DM, alors $s$ est localement quasi-fini ;
\item si $f$ est séparé, alors $s$ est propre ;
\item si $f$ est quasi-séparé, alors $s$ est quasi-compact et quasi-séparé.
\end{enumerate}
\end{lemma}

\begin{proof}
Il s'agit du cas particulier du Lemme \ref{stacks-morphisms-lemma-semi-diagonal} appliqué à
$g = s$ et $\mathcal{Y} = \mathcal{T}$ ; dans ce cas,
$i : \mathcal{T} \to \mathcal{T} \times_\mathcal{T} \mathcal{X}$
est $2$-isomorphe à $s$.
\end{proof}

\begin{lemma}
\label{stacks-morphisms-lemma-composition-separated}
Tous les axiomes de séparation énumérés dans la Définition
\ref{stacks-morphisms-definition-separated} sont stables par composition des morphismes.
\end{lemma}

\begin{proof}
Soient $f : \mathcal{X} \to \mathcal{Y}$ et
$g : \mathcal{Y} \to \mathcal{Z}$ des morphismes de champs algébriques
auxquels s'applique l'axiome considéré. La diagonale
$\Delta_{\mathcal{X}/\mathcal{Z}}$ est le composé
$$
\mathcal{X} \longrightarrow
\mathcal{X} \times_\mathcal{Y} \mathcal{X} \longrightarrow
\mathcal{X} \times_\mathcal{Z} \mathcal{X}.
$$
Notre axiome de séparation est défini en imposant à la diagonale une certaine
propriété $\mathcal{P}$. D'après le Lemme
\ref{stacks-morphisms-lemma-fibre-product-after-map} ci-dessus, la seconde flèche possède elle
aussi cette propriété. Le lemme en résulte, puisque le composé de morphismes
représentables par des espaces algébriques et possédant la propriété
$\mathcal{P}$ possède encore la propriété $\mathcal{P}$, voir la discussion
générale dans Propriétés des champs,
section \ref{stacks-properties-section-properties-morphisms}
et Morphismes d'espaces, Lemmes
\ref{spaces-morphisms-lemma-composition-unramified},
\ref{spaces-morphisms-lemma-composition-quasi-finite},
\ref{spaces-morphisms-lemma-composition-proper},
\ref{spaces-morphisms-lemma-composition-quasi-compact} et
\ref{spaces-morphisms-lemma-composition-separated}.
\end{proof}

\begin{lemma}
\label{stacks-morphisms-lemma-separated-over-separated}
Soit $f : \mathcal{X} \to \mathcal{Y}$ un morphisme de champs algébriques
sur le schéma de base $S$.
\begin{enumerate}
\item Si $\mathcal{Y}$ est DM sur $S$ et $f$ est DM, alors $\mathcal{X}$ est
DM sur $S$.
\item Si $\mathcal{Y}$ est quasi-DM sur $S$ et $f$ est quasi-DM, alors
$\mathcal{X}$ est quasi-DM sur $S$.
\item Si $\mathcal{Y}$ est séparé sur $S$ et $f$ est séparé, alors
$\mathcal{X}$ est séparé sur $S$.
\item Si $\mathcal{Y}$ est quasi-séparé sur $S$ et $f$ est quasi-séparé,
alors $\mathcal{X}$ est quasi-séparé sur $S$.
\item Si $\mathcal{Y}$ est DM et $f$ est DM, alors $\mathcal{X}$ est DM.
\item Si $\mathcal{Y}$ est quasi-DM et $f$ est quasi-DM, alors $\mathcal{X}$
est quasi-DM.
\item Si $\mathcal{Y}$ est séparé et $f$ est séparé, alors $\mathcal{X}$ est
séparé.
\item Si $\mathcal{Y}$ est quasi-séparé et $f$ est quasi-séparé, alors
$\mathcal{X}$ est quasi-séparé.
\end{enumerate}
\end{lemma}

\begin{proof}
Les assertions (1), (2), (3) et (4) résultent immédiatement du Lemme
\ref{stacks-morphisms-lemma-composition-separated} et de la Définition
\ref{stacks-morphisms-definition-absolute-separated}. Pour (5), (6), (7) et (8), considérons
$\mathcal{X}$ et $\mathcal{Y}$ comme des champs algébriques sur
$\Spec(\mathbf{Z})$ et appliquons le Lemme
\ref{stacks-morphisms-lemma-composition-separated}. Détails omis.
\end{proof}

\noindent
Le lemme suivant diffère quelque peu de son analogue pour les espaces
algébriques. Pour les comparer, voir Morphismes d'espaces, Lemme
\ref{spaces-morphisms-lemma-compose-after-separated}.

\begin{lemma}
\label{stacks-morphisms-lemma-compose-after-separated}
Soient $f : \mathcal{X} \to \mathcal{Y}$ et
$g : \mathcal{Y} \to \mathcal{Z}$ des morphismes de champs algébriques.
\begin{enumerate}
\item Si $g \circ f$ est DM, alors $f$ l'est aussi.
\item Si $g \circ f$ est quasi-DM, alors $f$ l'est aussi.
\item Si $g \circ f$ est séparé et $\Delta_g$ est séparé, alors $f$ est
séparé.
\item Si $g \circ f$ est quasi-séparé et si $\Delta_g$ est quasi-séparé,
alors $f$ est quasi-séparé.
\end{enumerate}
\end{lemma}

\begin{proof}
Considérons la factorisation
$$
\mathcal{X} \to
\mathcal{X} \times_\mathcal{Y} \mathcal{X} \to
\mathcal{X} \times_\mathcal{Z} \mathcal{X}
$$
du morphisme diagonal de $g \circ f$. Les deux morphismes sont représentables
par des espaces algébriques, voir les Lemmes
\ref{stacks-morphisms-lemma-properties-diagonal} et \ref{stacks-morphisms-lemma-fibre-product-after-map}. Ainsi,
pour tout schéma $T$ et tout morphisme
$T \to \mathcal{X} \times_\mathcal{Y} \mathcal{X}$, nous obtenons des
morphismes d'espaces algébriques
$$
A = \mathcal{X} \times_{(\mathcal{X} \times_\mathcal{Z} \mathcal{X})} T
\longrightarrow
B = (\mathcal{X} \times_\mathcal{Y} \mathcal{X})
\times_{(\mathcal{X} \times_\mathcal{Z} \mathcal{X})} T
\longrightarrow
T
$$
Si $g \circ f$ est DM (resp.\ quasi-DM), alors le composé $A \to T$ est non
ramifié (resp.\ localement quasi-fini). Par conséquent, $A \to B$ est non
ramifié (resp.\ localement quasi-fini) d'après Morphismes d'espaces, Lemme
\ref{spaces-morphisms-lemma-permanence-unramified}
(resp.\ Morphismes d'espaces, Lemme
\ref{spaces-morphisms-lemma-permanence-quasi-finite}). Considérons maintenant
le diagramme
$$
\xymatrix{
A \times_B T \ar[r] \ar[d] & T \ar[d] \\
A \ar[r] \ar[d] & B \ar[r] \ar[d] & T \ar[d] \\
\mathcal{X} \ar[r] &
\mathcal{X} \times_\mathcal{Y} \mathcal{X} \ar[r] &
\mathcal{X} \times_\mathcal{Z} \mathcal{X}
}
$$
dont tous les carrés sont $2$-cartésiens et où la flèche $T \to B$ provient
de la flèche $T \to \mathcal{X} \times_\mathcal{Y} \mathcal{X}$ de départ.
La flèche $A \times_B T \to T$ est non ramifiée (resp.\ localement
quasi-finie) d'après Morphismes d'espaces, Lemme
\ref{spaces-morphisms-lemma-base-change-unramified}
(resp.\ Morphismes d'espaces, Lemme
\ref{spaces-morphisms-lemma-base-change-quasi-finite}). Comme
$A \times_B T$ est égal à
$\mathcal{X} \times_{\mathcal{X} \times_\mathcal{Y} \mathcal{X}} T$, nous
concluons que $f$ est DM (resp.\ quasi-DM). Cela démontre (1) et (2).

\medskip\noindent
Démonstration de (4). Supposons que $g \circ f$ soit quasi-séparé et que
$\Delta_g$ soit quasi-séparé. Considérons la factorisation
$$
\mathcal{X} \to
\mathcal{X} \times_\mathcal{Y} \mathcal{X} \to
\mathcal{X} \times_\mathcal{Z} \mathcal{X}
$$
du morphisme diagonal de $g \circ f$. Les deux morphismes sont représentables
par des espaces algébriques, et le second est quasi-séparé, voir les Lemmes
\ref{stacks-morphisms-lemma-properties-diagonal} et \ref{stacks-morphisms-lemma-fibre-product-after-map}. Ainsi,
pour tout schéma $T$ et tout morphisme
$T \to \mathcal{X} \times_\mathcal{Y} \mathcal{X}$, nous obtenons des
morphismes d'espaces algébriques
$$
A = \mathcal{X} \times_{(\mathcal{X} \times_\mathcal{Z} \mathcal{X})} T
\longrightarrow
B = (\mathcal{X} \times_\mathcal{Y} \mathcal{X})
\times_{(\mathcal{X} \times_\mathcal{Z} \mathcal{X})} T
\longrightarrow
T
$$
tel que $B \to T$ soit quasi-séparé. Le composé $A \to T$ est quasi-compact
et quasi-séparé, puisque nous avons supposé que $g \circ f$ est quasi-séparé.
Par conséquent, $A \to B$ est quasi-séparé d'après Morphismes d'espaces,
Lemme \ref{spaces-morphisms-lemma-compose-after-separated}. De plus,
$A \to B$ est quasi-compact d'après Morphismes d'espaces, Lemme
\ref{spaces-morphisms-lemma-quasi-compact-permanence}. Ainsi, $f$ est
quasi-séparé.

\medskip\noindent
Démonstration de (3). Supposons que $g \circ f$ soit séparé et que $\Delta_g$
soit séparé. Considérons la factorisation
$$
\mathcal{X} \to
\mathcal{X} \times_\mathcal{Y} \mathcal{X} \to
\mathcal{X} \times_\mathcal{Z} \mathcal{X}
$$
du morphisme diagonal de $g \circ f$. Les deux morphismes sont représentables
par des espaces algébriques, et le second est séparé, voir les Lemmes
\ref{stacks-morphisms-lemma-properties-diagonal} et \ref{stacks-morphisms-lemma-fibre-product-after-map}. Ainsi,
pour tout schéma $T$ et tout morphisme
$T \to \mathcal{X} \times_\mathcal{Y} \mathcal{X}$, nous obtenons des
morphismes d'espaces algébriques
$$
A = \mathcal{X} \times_{(\mathcal{X} \times_\mathcal{Z} \mathcal{X})} T
\longrightarrow
B = (\mathcal{X} \times_\mathcal{Y} \mathcal{X})
\times_{(\mathcal{X} \times_\mathcal{Z} \mathcal{X})} T
\longrightarrow
T
$$
tel que $B \to T$ soit séparé. Le composé $A \to T$ est propre, puisque nous
avons supposé que $g \circ f$ est séparé. Par conséquent, $A \to B$ est propre
d'après Morphismes d'espaces, Lemme
\ref{spaces-morphisms-lemma-universally-closed-permanence}, ce qui signifie
que $f$ est séparé.
\end{proof}

\begin{lemma}
\label{stacks-morphisms-lemma-separated-implies-morphism-separated}
Soit $\mathcal{X}$ un champ algébrique sur le schéma de base $S$.
\begin{enumerate}
\item
$\mathcal{X}$ est DM $\Leftrightarrow$
$\mathcal{X}$ est DM sur $S$.
\item
$\mathcal{X}$ est quasi-DM $\Leftrightarrow$
$\mathcal{X}$ est quasi-DM sur $S$.
\item Si $\mathcal{X}$ est séparé, alors $\mathcal{X}$ est séparé sur $S$.
\item Si $\mathcal{X}$ est quasi-séparé, alors $\mathcal{X}$ est quasi-séparé
sur $S$.
\end{enumerate}
Soit $f : \mathcal{X} \to \mathcal{Y}$ un morphisme de champs algébriques
sur le schéma de base $S$.
\begin{enumerate}
\item[(5)] Si $\mathcal{X}$ est DM sur $S$, alors $f$ est DM.
\item[(6)] Si $\mathcal{X}$ est quasi-DM sur $S$, alors $f$ est quasi-DM.
\item[(7)] Si $\mathcal{X}$ est séparé sur $S$ et si
$\Delta_{\mathcal{Y}/S}$ est séparé, alors $f$ est séparé.
\item[(8)] Si $\mathcal{X}$ est quasi-séparé sur $S$ et si
$\Delta_{\mathcal{Y}/S}$ est quasi-séparé, alors $f$ est quasi-séparé.
\end{enumerate}
\end{lemma}

\begin{proof}
Les assertions (5), (6), (7) et (8) résultent immédiatement du Lemme
\ref{stacks-morphisms-lemma-compose-after-separated} et de Espaces, Définition
\ref{spaces-definition-separated}. Pour démontrer (3) et (4), considérons
$\mathcal{X}$ et $\mathcal{Y}$ comme des champs algébriques sur
$\Spec(\mathbf{Z})$ et appliquons le Lemme
\ref{stacks-morphisms-lemma-compose-after-separated}. De même, pour démontrer (1) et (2),
considérons $\mathcal{X}$ comme un champ algébrique sur
$\Spec(\mathbf{Z})$ et considérons les morphismes
$$
\mathcal{X} \longrightarrow
\mathcal{X} \times_S \mathcal{X} \longrightarrow
\mathcal{X} \times_{\Spec(\mathbf{Z})} \mathcal{X}
$$
Les deux flèches sont représentables par des espaces algébriques. La seconde
est non ramifiée et localement quasi-finie, car elle est le changement de base
de l'immersion $\Delta_{S/\mathbf{Z}}$. Le composé est donc non ramifié
(resp.\ localement quasi-fini) si et seulement si la première flèche est non
ramifiée (resp.\ localement quasi-finie), voir Morphismes d'espaces, Lemmes
\ref{spaces-morphisms-lemma-composition-unramified} et
\ref{spaces-morphisms-lemma-permanence-unramified}
(resp.\ Morphismes d'espaces, Lemmes
\ref{spaces-morphisms-lemma-composition-quasi-finite} et
\ref{spaces-morphisms-lemma-permanence-quasi-finite}).
\end{proof}

\begin{lemma}
\label{stacks-morphisms-lemma-properties-covering-imply-diagonal}
Soit $\mathcal{X}$ un champ algébrique.
Soit $W$ un espace algébrique et soit $f : W \to \mathcal{X}$
un morphisme surjectif, plat et localement de présentation finie.
\begin{enumerate}
\item Si $f$ est non ramifié (c'est-à-dire étale, ou encore si
$\mathcal{X}$ est de Deligne--Mumford), alors $\mathcal{X}$ est DM.
\item Si $f$ est localement quasi-fini, alors $\mathcal{X}$ est quasi-DM.
\end{enumerate}
\end{lemma}

\begin{proof}
Remarquons que, si $f$ est non ramifié, alors il est étale d'après
Morphismes d'espaces, Lemme
\ref{spaces-morphisms-lemma-unramified-flat-lfp-etale}.
Cela explique la parenthèse de (1).
Supposons que $f$ soit non ramifié (resp.\ localement quasi-fini). Il faut
montrer que
$\Delta_\mathcal{X} : \mathcal{X} \to \mathcal{X} \times \mathcal{X}$
est non ramifié (resp.\ localement quasi-fini). Remarquons que
$W \times W \to \mathcal{X} \times \mathcal{X}$ est lui aussi surjectif,
plat et localement de présentation finie. Il suffit donc de montrer que
$$
W \times_{\mathcal{X} \times \mathcal{X}, \Delta_\mathcal{X}} \mathcal{X}
=
W \times_\mathcal{X} W
\longrightarrow
W \times W
$$
est non ramifié (resp.\ localement quasi-fini), voir Propriétés des champs,
Lemme
\ref{stacks-properties-lemma-check-property-covering}.
Par hypothèse, le morphisme $\text{pr}_i : W \times_\mathcal{X} W \to W$
est non ramifié (resp.\ localement quasi-fini). La flèche affichée est donc
non ramifiée (resp.\ localement quasi-finie) d'après Morphismes d'espaces,
Lemme
\ref{spaces-morphisms-lemma-permanence-unramified}
(resp.\ Morphismes d'espaces, Lemme
\ref{spaces-morphisms-lemma-permanence-quasi-finite}).
\end{proof}

\begin{lemma}
\label{stacks-morphisms-lemma-monomorphism-separated}
Un monomorphisme de champs algébriques est séparé et DM.
Il en va de même des immersions de champs algébriques.
\end{lemma}

\begin{proof}
Si $f : \mathcal{X} \to \mathcal{Y}$ est un monomorphisme de champs
algébriques, alors $\Delta_f$ est un isomorphisme, voir Propriétés des
champs, Lemme \ref{stacks-properties-lemma-monomorphism}.
Comme un isomorphisme d'espaces algébriques est propre et non ramifié,
$f$ est séparé et DM. La seconde assertion résulte de la première, puisqu'une
immersion est un monomorphisme, voir Propriétés des champs, Lemme
\ref{stacks-properties-lemma-immersion-monomorphism}.
\end{proof}

\begin{lemma}
\label{stacks-morphisms-lemma-separation-properties-residual-gerbe}
Soit $\mathcal{X}$ un champ algébrique et soit $x \in |\mathcal{X}|$.
Supposons que la gerbe résiduelle $\mathcal{Z}_x$ de $\mathcal{X}$ en $x$
existe. Si $\mathcal{X}$ est DM, resp.\ quasi-DM, resp.\ séparé, resp.\
quasi-séparé, alors $\mathcal{Z}_x$ l'est aussi.
\end{lemma}

\begin{proof}
En effet, $\mathcal{Z}_x \to \mathcal{X}$ est un monomorphisme, donc est
DM et séparé d'après le Lemme \ref{stacks-morphisms-lemma-monomorphism-separated}.
On conclut en appliquant le Lemme \ref{stacks-morphisms-lemma-separated-over-separated}.
\end{proof}


















\section{Champs d'inertie}
\label{stacks-morphisms-section-inertia}

\noindent
Le champ d'inertie (relatif) d'un champ en groupoïdes est défini dans
Champs, section \ref{stacks-section-the-inertia-stack}.
La construction proprement dite, dans le cadre des catégories fibrées, ainsi
que certaines de ses propriétés se trouvent dans Catégories, section
\ref{categories-section-inertia}.

\begin{lemma}
\label{stacks-morphisms-lemma-inertia}
Soit $\mathcal{X}$ un champ algébrique. Alors le champ d'inertie
$\mathcal{I}_\mathcal{X}$ est lui aussi un champ algébrique. Le morphisme
$$
\mathcal{I}_\mathcal{X} \longrightarrow \mathcal{X}
$$
est représentable par des espaces algébriques et localement de type fini.
Plus généralement, soit $f : \mathcal{X} \to \mathcal{Y}$ un morphisme de
champs algébriques. Alors l'inertie relative
$\mathcal{I}_{\mathcal{X}/\mathcal{Y}}$ est un champ algébrique et le
morphisme
$$
\mathcal{I}_{\mathcal{X}/\mathcal{Y}} \longrightarrow \mathcal{X}
$$
est représentable par des espaces algébriques et localement de type fini.
\end{lemma}

\begin{proof}
D'après Catégories, Lemme
\ref{categories-lemma-inertia-fibred-category}, il existe des équivalences
$$
\mathcal{I}_\mathcal{X} \to
\mathcal{X} \times_{\Delta, \mathcal{X} \times_S \mathcal{X}, \Delta}
\mathcal{X}
\quad\text{et}\quad
\mathcal{I}_{\mathcal{X}/\mathcal{Y}} \to
\mathcal{X}
\times_{\Delta, \mathcal{X} \times_\mathcal{Y} \mathcal{X}, \Delta}
\mathcal{X}
$$
qui montrent que les champs d'inertie sont des champs algébriques.
Soit $T \to \mathcal{X}$ le morphisme donné par l'objet $x$ de la catégorie
fibre de $\mathcal{X}$ au-dessus de $T$. Nous obtenons alors un carré de
$2$-produit fibré
$$
\xymatrix{
\mathit{Isom}_\mathcal{X}(x, x) \ar[d] \ar[r] &
\mathcal{I}_\mathcal{X} \ar[d] \\
T \ar[r]^x & \mathcal{X}
}
$$
Cela résulte immédiatement de la définition de $\mathcal{I}_\mathcal{X}$.
Puisque $\mathit{Isom}_\mathcal{X}(x, x)$ est toujours un espace algébrique
localement de type fini sur $T$ (voir le Lemme
\ref{stacks-morphisms-lemma-isom-locally-finite-type}), nous concluons que
$\mathcal{I}_\mathcal{X} \to \mathcal{X}$ est représentable par des espaces
algébriques et localement de type fini. Enfin, pour l'inertie relative, nous
obtenons
$$
\vcenter{
\xymatrix{
\mathit{Isom}_\mathcal{X}(x, x) \ar[d] &
K \ar[l] \ar[d] \ar[r] &
\mathcal{I}_{\mathcal{X}/\mathcal{Y}} \ar[d] \\
\mathit{Isom}_\mathcal{Y}(f(x), f(x)) &
T \ar[l]_-e \ar[r]^x & \mathcal{X}
}
}
$$
où les deux carrés sont des $2$-produits fibrés. Cela résulte de Catégories,
Lemme \ref{categories-lemma-relative-inertia-as-fibre-product}.
La flèche verticale de gauche est un morphisme d'espaces algébriques
localement de type fini sur $T$; elle est donc localement de type fini, voir
Morphismes d'espaces, Lemme
\ref{spaces-morphisms-lemma-permanence-finite-type}.
Ainsi, $K$ est un espace algébrique et $K \to T$ est localement de type fini.
Ceci démontre l'assertion relative à l'inertie relative.
\end{proof}

\begin{remark}
\label{stacks-morphisms-remark-inertia-is-group-in-spaces}
Soit $f : \mathcal{X} \to \mathcal{Y}$ un morphisme de champs algébriques.
Dans Propriétés des champs, Remarque
\ref{stacks-properties-remark-representable-over}, nous avons vu que la
$2$-catégorie des morphismes $\mathcal{Z} \to \mathcal{X}$ représentables
par des espaces algébriques, ayant $\mathcal{X}$ pour but, est une catégorie.
Dans cette catégorie, le champ d'inertie de $\mathcal{X}/\mathcal{Y}$ est un
{\it objet en groupes}. Rappelons qu'un objet de
$\mathcal{I}_{\mathcal{X}/\mathcal{Y}}$
est simplement un couple $(x, \alpha)$, où $x$ est un objet de $\mathcal{X}$
et $\alpha$ un automorphisme de $x$, dans la catégorie fibre de
$\mathcal{X}$ à laquelle appartient $x$, tel que
$f(\alpha) = \text{id}$. La composition
$$
c :
\mathcal{I}_{\mathcal{X}/\mathcal{Y}}
\times_\mathcal{X} \mathcal{I}_{\mathcal{X}/\mathcal{Y}}
\longrightarrow
\mathcal{I}_{\mathcal{X}/\mathcal{Y}}
$$
est donnée sur les objets par la règle
$$
((x, \alpha), (x', \alpha'), \beta) \mapsto
(x, \alpha \circ \beta^{-1} \circ \alpha' \circ \beta)
$$
qui a un sens puisque $\beta : x \to x'$ est un isomorphisme dans la
catégorie fibre, par notre définition des produits fibrés. L'élément neutre
$e : \mathcal{X} \to \mathcal{I}_{\mathcal{X}/\mathcal{Y}}$ est donné par le
foncteur $x \mapsto (x, \text{id}_x)$. Nous omettons la vérification des
axiomes d'un objet en groupes.
\end{remark}

\noindent
Soit $f : \mathcal{X} \to \mathcal{Y}$ un morphisme de champs algébriques
et soit $\mathcal{I}_{\mathcal{X}/\mathcal{Y}}$ son champ d'inertie.
Soit $T$ un schéma et soit $x$ un objet de $\mathcal{X}$ au-dessus de $T$.
Posons $y = f(x)$. Nous avons vu, dans la démonstration du Lemme
\ref{stacks-morphisms-lemma-inertia}, que, pour tout schéma $T$ et tout objet $x$ de
$\mathcal{X}$ au-dessus de $T$, il existe une suite exacte de faisceaux de
groupes
\begin{equation}
\label{stacks-morphisms-equation-exact-sequence-isom}
0 \to
\mathit{Isom}_{\mathcal{X}/\mathcal{Y}}(x, x) \to
\mathit{Isom}_\mathcal{X}(x, x) \to
\mathit{Isom}_\mathcal{Y}(y, y)
\end{equation}
Les structures de groupe sur les deuxième et troisième termes sont celles
définies au Lemme \ref{stacks-morphisms-lemma-isom-pseudo-torsor-aut}, et la suite donne un
sens au premier terme. Il existe en outre un carré cartésien canonique
$$
\xymatrix{
\mathit{Isom}_{\mathcal{X}/\mathcal{Y}}(x, x) \ar[d] \ar[r] &
\mathcal{I}_{\mathcal{X}/\mathcal{Y}} \ar[d] \\
T \ar[r]^x & \mathcal{X}
}
$$
En fait, la structure de groupe sur
$\mathcal{I}_{\mathcal{X}/\mathcal{Y}}$ considérée dans la Remarque
\ref{stacks-morphisms-remark-inertia-is-group-in-spaces} induit la structure de groupe sur
$\mathit{Isom}_{\mathcal{X}/\mathcal{Y}}(x, x)$. Cela nous permet de définir
le faisceau
$\mathit{Isom}_{\mathcal{X}/\mathcal{Y}}(x, x)$
également pour les morphismes d'espaces algébriques vers $\mathcal{X}$.
Nous formalisons ceci dans la définition suivante.

\begin{definition}
\label{stacks-morphisms-definition-isom}
Soit $f : \mathcal{X} \to \mathcal{Y}$ un morphisme de champs algébriques.
Soit $Z$ un espace algébrique.
\begin{enumerate}
\item Soit $x : Z \to \mathcal{X}$ un morphisme. Nous posons
$$
\mathit{Isom}_{\mathcal{X}/\mathcal{Y}}(x, x) =
Z \times_{x, \mathcal{X}} \mathcal{I}_{\mathcal{X}/\mathcal{Y}}
$$
Nous le munissons d'une structure d'espace algébrique en groupes sur $Z$ par
changement de base de la loi de composition considérée dans la Remarque
\ref{stacks-morphisms-remark-inertia-is-group-in-spaces}. Nous appellerons parfois
$\mathit{Isom}_{\mathcal{X}/\mathcal{Y}}(x, x)$ le {\it faisceau relatif
d'automorphismes de $x$}.
\item Soient $x_1, x_2 : Z \to \mathcal{X}$ des morphismes. Posons
$y_i = f \circ x_i$. Soit $\alpha : y_1 \to y_2$ un $2$-morphisme.
Alors $\alpha$ détermine un morphisme
$\Delta^\alpha : Z \to Z \times_{y_1, \mathcal{Y}, y_2} Z$, et nous posons
$$
\mathit{Isom}_{\mathcal{X}/\mathcal{Y}}^\alpha(x_1, x_2) =
(Z \times_{x_1, \mathcal{X}, x_2} Z)
\times_{Z \times_{y_1, \mathcal{Y}, y_2} Z, \Delta^\alpha} Z.
$$
Nous appellerons parfois
$\mathit{Isom}_{\mathcal{X}/\mathcal{Y}}^\alpha(x_1, x_2)$
le {\it faisceau relatif des isomorphismes de $x_1$ vers $x_2$}.
\end{enumerate}
Si $\mathcal{Y} = \Spec(\mathbf{Z})$, ou plus généralement si
$\mathcal{Y}$ est un espace algébrique, nous employons les notations
$\mathit{Isom}_\mathcal{X}(x, x)$ et $\mathit{Isom}_\mathcal{X}(x_1, x_2)$
et nous parlons du {\it faisceau d'automorphismes de $x$} et du
{\it faisceau des isomorphismes de $x_1$ vers $x_2$}.
\end{definition}

\begin{lemma}
\label{stacks-morphisms-lemma-isom-pseudo-torsor-aut-over-space}
Soit $f : \mathcal{X} \to \mathcal{Y}$ un morphisme de champs algébriques.
Soit $Z$ un espace algébrique et soient
$x_i : Z \to \mathcal{X}$, $i = 1, 2$, des morphismes. Alors
\begin{enumerate}
\item $\mathit{Isom}_{\mathcal{X}/\mathcal{Y}}(x_2, x_2)$
est un espace algébrique en groupes sur $Z$,
\item il existe une suite exacte de groupes
$$
0 \to \mathit{Isom}_{\mathcal{X}/\mathcal{Y}}(x_2, x_2)
\to \mathit{Isom}_\mathcal{X}(x_2, x_2)
\to \mathit{Isom}_\mathcal{Y}(f \circ x_2, f \circ x_2)
$$
\item il existe un morphisme d'espaces algébriques
$
\mathit{Isom}_\mathcal{X}(x_1, x_2)
\to \mathit{Isom}_\mathcal{Y}(f \circ x_1, f \circ x_2)
$
telle que, pour tout $2$-morphisme
$\alpha : f \circ x_1 \to f \circ x_2$, nous obtenions un diagramme
cartésien
$$
\xymatrix{
\mathit{Isom}_{\mathcal{X}/\mathcal{Y}}^\alpha(x_1, x_2) \ar[d] \ar[r] &
Z \ar[d]^\alpha \\
\mathit{Isom}_\mathcal{X}(x_1, x_2) \ar[r] &
\mathit{Isom}_\mathcal{Y}(f \circ x_1, f \circ x_2)
}
$$
\item pour tout $2$-morphisme $\alpha : f \circ x_1 \to f \circ x_2$,
l'espace algébrique
$\mathit{Isom}_{\mathcal{X}/\mathcal{Y}}^\alpha(x_1, x_2)$ est un
pseudo-torseur sous
$\mathit{Isom}_{\mathcal{X}/\mathcal{Y}}(x_2, x_2)$ sur $Z$.
\end{enumerate}
\end{lemma}

\begin{proof}
L'assertion (1) résulte de la Définition \ref{stacks-morphisms-definition-isom}.
L'assertion (2) résulte de la suite exacte
(\ref{stacks-morphisms-equation-exact-sequence-isom}) localement pour la topologie étale sur
$Z$. L'assertion (3) se voit en explicitant les définitions. Localement sur
$Z$ pour la topologie étale, l'assertion (4) se réduit à l'assertion (2) du
Lemme \ref{stacks-morphisms-lemma-isom-pseudo-torsor-aut}.
\end{proof}

\begin{lemma}
\label{stacks-morphisms-lemma-cartesian-square-inertia}
Soient $\pi : \mathcal{X} \to \mathcal{Y}$ et
$f : \mathcal{Y}' \to \mathcal{Y}$ des morphismes de champs algébriques.
Posons $\mathcal{X}' = \mathcal{X} \times_\mathcal{Y} \mathcal{Y}'$.
Alors les deux carrés du diagramme
$$
\xymatrix{
\mathcal{I}_{\mathcal{X}'/\mathcal{Y}'} \ar[r]
\ar[d]_{
\text{Catégories, Équation}\ (\ref{categories-equation-functorial})
} &
\mathcal{X}' \ar[r]_{\pi'} \ar[d] & \mathcal{Y}' \ar[d]^f \\
\mathcal{I}_{\mathcal{X}/\mathcal{Y}} \ar[r] &
\mathcal{X} \ar[r]^\pi & \mathcal{Y}
}
$$
sont des carrés de produit fibré.
\end{lemma}

\begin{proof}
Le champ d'inertie $\mathcal{I}_{\mathcal{X}'/\mathcal{Y}'}$ est défini
comme la catégorie des couples $(x', \alpha')$, où $x'$ est un objet de
$\mathcal{X}'$ et $\alpha'$ un automorphisme de $x'$ tel que
$\pi'(\alpha') = \text{id}$, voir Catégories, section
\ref{categories-section-inertia}. Supposons que $x'$ soit au-dessus du
schéma $U$ et s'envoie sur l'objet $x$ de $\mathcal{X}$. Par la construction
du $2$-produit fibré dans Catégories, Lemme
\ref{categories-lemma-2-product-categories-over-C}, nous voyons que
$x' = (U, x, y', \beta)$, où $y'$ est un objet de $\mathcal{Y}'$ au-dessus
de $U$ et $\beta$ un isomorphisme
$\beta : \pi(x) \to f(y')$ dans la catégorie fibre de $\mathcal{Y}$
au-dessus de $U$. Par la construction même du $2$-produit fibré,
l'automorphisme $\alpha'$ est un couple $(\alpha, \gamma)$, où $\alpha$ est
un automorphisme de $x$ au-dessus de $U$ et $\gamma$ un automorphisme de
$y'$ au-dessus de $U$, tels que $\alpha$ et $\gamma$ soient compatibles par
$\beta$. La condition $\pi'(\alpha') = \text{id}$ signifie que
$\gamma = \text{id}$; la compatibilité de $\alpha, \beta, \gamma$ équivaut
alors exactement à la condition $\pi(\alpha) = \text{id}$, c'est-à-dire que
$(x, \alpha)$ est un objet de $\mathcal{I}_{\mathcal{X}/\mathcal{Y}}$.
Nous voyons ainsi que le carré de gauche est un carré de produit fibré
(certains détails sont omis).
\end{proof}

\begin{lemma}
\label{stacks-morphisms-lemma-monomorphism-cartesian-square-inertia}
Soit $f : \mathcal{X} \to \mathcal{Y}$ un monomorphisme de champs
algébriques. Alors le diagramme
$$
\xymatrix{
\mathcal{I}_\mathcal{X} \ar[r] \ar[d] &
\mathcal{X} \ar[d] \\
\mathcal{I}_\mathcal{Y} \ar[r] &
\mathcal{Y}
}
$$
est un carré de produit fibré.
\end{lemma}

\begin{proof}
Cela résulte immédiatement du fait que $f$ est pleinement fidèle (voir
Propriétés des champs, Lemme
\ref{stacks-properties-lemma-monomorphism}) et de la définition de l'inertie
dans Catégories, section \ref{categories-section-inertia}. En effet, un objet
de $\mathcal{I}_\mathcal{X}$ au-dessus d'un schéma $T$ est la même chose
qu'un couple $(x, \alpha)$ constitué d'un objet $x$ de $\mathcal{X}$
au-dessus de $T$ et d'un morphisme $\alpha : x \to x$ dans la catégorie
fibre de $\mathcal{X}$ au-dessus de $T$. Comme $f$ est pleinement fidèle,
se donner $\alpha$ revient à se donner un morphisme
$\beta : f(x) \to f(x)$ dans la catégorie fibre de $\mathcal{Y}$ au-dessus
de $T$. Nous pouvons donc considérer les objets de
$\mathcal{I}_\mathcal{X}$ au-dessus de $T$ comme les triplets
$((y, \beta), x, \gamma)$, où $y$ est un objet de $\mathcal{Y}$ au-dessus
de $T$, où $\beta : y \to y$ appartient à $\mathcal{Y}_T$, et où
$\gamma : y \to f(x)$ est un isomorphisme au-dessus de $T$, c'est-à-dire
comme les objets de $\mathcal{I}_\mathcal{Y} \times_\mathcal{Y}
\mathcal{X}$ au-dessus de $T$.
\end{proof}

\begin{lemma}
\label{stacks-morphisms-lemma-presentation-inertia}
Soit $\mathcal{X}$ un champ algébrique. Soit
$[U/R] \to \mathcal{X}$ une présentation. Soit $G/U$ l'espace algébrique
en groupes stabilisateur associé au groupoïde $(U, R, s, t, c)$. Alors
$$
\xymatrix{
G \ar[d] \ar[r] & U \ar[d] \\
\mathcal{I}_\mathcal{X} \ar[r] & \mathcal{X}
}
$$
est un diagramme de produit fibré.
\end{lemma}

\begin{proof}
C'est immédiat d'après Groupoïdes en espaces, Lemme
\ref{spaces-groupoids-lemma-2-cartesian-inertia}.
\end{proof}









\section{Diagonales d'ordre supérieur}
\label{stacks-morphisms-section-higher-diagonals}

\noindent
Soit $f : \mathcal{X} \to \mathcal{Y}$ un morphisme de champs algébriques.
Dans cette situation, il est naturel de considérer non seulement la diagonale
$$
\Delta_f : \mathcal{X} \to \mathcal{X} \times_\mathcal{Y} \mathcal{X}
$$
mais aussi la diagonale de la diagonale, c'est-à-dire le morphisme
$$
\Delta_{\Delta_f} :
\mathcal{X}
\longrightarrow
\mathcal{X} \times_{(\mathcal{X} \times_\mathcal{Y} \mathcal{X})} \mathcal{X}
$$
C'est pourquoi nous employons parfois la terminologie suivante. Nous appelons
$\Delta_{f, 0} = f$ la {\it diagonale d'ordre zéro},
$\Delta_{f, 1} = \Delta_f$ la {\it première diagonale} et
$\Delta_{f, 2} = \Delta_{\Delta_f}$ la {\it seconde diagonale}.
Remarquons que $\Delta_{f, 1}$ est représentable par des espaces algébriques
et localement de type fini, voir le Lemme \ref{stacks-morphisms-lemma-properties-diagonal}.
Par conséquent, $\Delta_{f, 2}$ est représentable, est un monomorphisme, et
est localement de type fini, séparé et localement quasi-fini, voir le Lemme
\ref{stacks-morphisms-lemma-properties-diagonal-representable}.

\medskip\noindent
Nous pouvons décrire la seconde diagonale à l'aide du champ d'inertie relatif.
En effet, le produit fibré
$\mathcal{X}
\times_{(\mathcal{X} \times_\mathcal{Y} \mathcal{X})} \mathcal{X}$
est équivalent au champ d'inertie relatif
$\mathcal{I}_{\mathcal{X}/\mathcal{Y}}$ d'après Catégories, Lemme
\ref{categories-lemma-inertia-fibred-category}. De plus, par cette
identification, la seconde diagonale devient la {\it section unité}
$$
\Delta_{f, 2} = e : \mathcal{X} \to \mathcal{I}_{\mathcal{X}/\mathcal{Y}}
$$
du champ d'inertie relatif. Par analogie avec ce qui se passe pour les
groupoïdes en espaces algébriques (Groupoïdes en espaces, Lemme
\ref{spaces-groupoids-lemma-diagonal}), nous avons les équivalences suivantes.

\begin{lemma}
\label{stacks-morphisms-lemma-diagonal-diagonal}
Soit $f : \mathcal{X} \to \mathcal{Y}$ un morphisme de champs algébriques.
\begin{enumerate}
\item
Les conditions suivantes sont équivalentes :
\begin{enumerate}
\item $\mathcal{I}_{\mathcal{X}/\mathcal{Y}} \to \mathcal{X}$
est séparé,
\item $\Delta_{f, 1} = \Delta_f :
\mathcal{X} \to \mathcal{X} \times_\mathcal{Y} \mathcal{X}$
est séparé, et
\item $\Delta_{f, 2} = e :
\mathcal{X} \to \mathcal{I}_{\mathcal{X}/\mathcal{Y}}$
est une immersion fermée.
\end{enumerate}
\item
Les conditions suivantes sont équivalentes :
\begin{enumerate}
\item $\mathcal{I}_{\mathcal{X}/\mathcal{Y}} \to \mathcal{X}$
est quasi-séparé,
\item $\Delta_{f, 1} = \Delta_f :
\mathcal{X} \to \mathcal{X} \times_\mathcal{Y} \mathcal{X}$
est quasi-séparé, et
\item $\Delta_{f, 2} = e :
\mathcal{X} \to \mathcal{I}_{\mathcal{X}/\mathcal{Y}}$
est quasi-compact.
\end{enumerate}
\item
Les conditions suivantes sont équivalentes :
\begin{enumerate}
\item $\mathcal{I}_{\mathcal{X}/\mathcal{Y}} \to \mathcal{X}$
est localement séparé,
\item $\Delta_{f, 1} = \Delta_f :
\mathcal{X} \to \mathcal{X} \times_\mathcal{Y} \mathcal{X}$
est localement séparé, et
\item $\Delta_{f, 2} = e :
\mathcal{X} \to \mathcal{I}_{\mathcal{X}/\mathcal{Y}}$
est une immersion.
\end{enumerate}
\item
Les conditions suivantes sont équivalentes :
\begin{enumerate}
\item $\mathcal{I}_{\mathcal{X}/\mathcal{Y}} \to \mathcal{X}$
est non ramifié,
\item $f$ est DM.
\end{enumerate}
\item
Les conditions suivantes sont équivalentes :
\begin{enumerate}
\item $\mathcal{I}_{\mathcal{X}/\mathcal{Y}} \to \mathcal{X}$
est localement quasi-fini,
\item $f$ est quasi-DM.
\end{enumerate}
\end{enumerate}
\end{lemma}

\begin{proof}
Démonstration de (1), (2) et (3).
Choisissons un espace algébrique $U$ et un morphisme lisse surjectif
$U \to \mathcal{X}$. Alors
$G = U \times_\mathcal{X} \mathcal{I}_{\mathcal{X}/\mathcal{Y}}$
est un espace algébrique sur $U$ (Lemme \ref{stacks-morphisms-lemma-inertia}). En fait,
$G$ est un espace algébrique en groupes sur $U$, par la loi de groupe sur
l'inertie relative construite dans la Remarque
\ref{stacks-morphisms-remark-inertia-is-group-in-spaces}. De plus,
$G \to \mathcal{I}_{\mathcal{X}/\mathcal{Y}}$ est surjectif et lisse comme
changement de base de $U \to \mathcal{X}$. Enfin, le changement de base de
$e : \mathcal{X} \to \mathcal{I}_{\mathcal{X}/\mathcal{Y}}$
par $G \to \mathcal{I}_{\mathcal{X}/\mathcal{Y}}$ est la section unité
$U \to G$ de $G/U$. L'équivalence de (a) et (c) résulte donc de Groupoïdes
en espaces, Lemme \ref{spaces-groupoids-lemma-group-scheme-separated}.
Puisque $\Delta_{f, 2}$ est la diagonale de $\Delta_f$, nous avons
(b) $\Leftrightarrow$ (c) par définition.

\medskip\noindent
Démonstration de (4) et (5). Rappelons que (4)(b) signifie que $\Delta_f$
est non ramifié, tandis que (5)(b) signifie que $\Delta_f$ est localement
quasi-fini. Choisissons un schéma $Z$ et un morphisme
$a : Z \to \mathcal{X} \times_\mathcal{Y} \mathcal{X}$.
Alors $a = (x_1, x_2, \alpha)$, où $x_i : Z \to \mathcal{X}$ et où
$\alpha : f \circ x_1  \to f \circ x_2$ est un $2$-morphisme. Rappelons que
$$
\vcenter{
\xymatrix{
\mathit{Isom}_{\mathcal{X}/\mathcal{Y}}^\alpha(x_1, x_2)
\ar[d] \ar[r] &
Z \ar[d] \\
\mathcal{X} \ar[r]^{\Delta_f} &
\mathcal{X} \times_\mathcal{Y} \mathcal{X}
}
}
\quad\text{et}\quad
\vcenter{
\xymatrix{
\mathit{Isom}_{\mathcal{X}/\mathcal{Y}}(x_2, x_2)
\ar[d] \ar[r] &
Z \ar[d]^{x_2} \\
\mathcal{I}_{\mathcal{X}/\mathcal{Y}} \ar[r] &
\mathcal{X}
}
}
$$
sont des carrés cartésiens. D'après le Lemme
\ref{stacks-morphisms-lemma-isom-pseudo-torsor-aut-over-space}, l'espace algébrique
$\mathit{Isom}_{\mathcal{X}/\mathcal{Y}}^\alpha(x_1, x_2)$ est un
pseudo-torseur sous
$\mathit{Isom}_{\mathcal{X}/\mathcal{Y}}(x_2, x_2)$ sur $Z$. Les
équivalences de (4) et (5) résultent donc de Groupoïdes en espaces, Lemme
\ref{spaces-groupoids-lemma-pseudo-torsor-implications}.
\end{proof}

\begin{lemma}
\label{stacks-morphisms-lemma-second-diagonal}
Soit $f : \mathcal{X} \to \mathcal{Y}$ un morphisme de champs algébriques.
Les conditions suivantes sont équivalentes :
\begin{enumerate}
\item le morphisme $f$ est représentable par des espaces algébriques,
\item la seconde diagonale de $f$ est un isomorphisme,
\item le champ en groupes $ \mathcal{I}_{\mathcal{X}/\mathcal{Y}}$
est trivial sur $\mathcal X$, et
\item pour tout schéma $T$ et tout morphisme $x : T \to \mathcal{X}$,
le noyau de $\mathit{Isom}_\mathcal{X}(x, x) \to
\mathit{Isom}_\mathcal{Y}(f(x), f(x))$ est trivial.
\end{enumerate}
\end{lemma}

\begin{proof}
Montrons d'abord l'équivalence de (1) et (2). En effet, $f$ est représentable
par des espaces algébriques si et seulement si $f$ est fidèle, voir Champs
algébriques, Lemme
\ref{algebraic-lemma-characterize-representable-by-algebraic-spaces}.
D'autre part, $f$ est fidèle si et seulement si, pour tout objet $x$ de
$\mathcal{X}$ au-dessus d'un schéma $T$, le foncteur $f$ induit une injection
$\mathit{Isom}_\mathcal{X}(x, x) \to
\mathit{Isom}_\mathcal{Y}(f(x), f(x))$,
ce qui a lieu si et seulement si le noyau $K$ est trivial, ou encore si et
seulement si $e : T \to K$ est un isomorphisme pour tout
$x : T \to \mathcal{X}$. Puisque
$K = T \times_{x, \mathcal{X}} \mathcal{I}_{\mathcal{X}/\mathcal{Y}}$,
comme nous l'avons vu ci-dessus, cela démontre l'équivalence de (1) et (2).
Pour démontrer l'équivalence de (2) et (3), il suffit, d'après ce qui précède,
de remarquer qu'un champ en groupes est trivial si et seulement si sa section
unité est un isomorphisme. Enfin, l'équivalence de (3) et (4) résulte des
définitions : dans la démonstration du Lemme \ref{stacks-morphisms-lemma-inertia}, nous avons
vu que le noyau de (4) correspond au produit fibré
$T \times_{x, \mathcal{X}} \mathcal{I}_{\mathcal{X}/\mathcal{Y}}$ sur $T$.
\end{proof}

\noindent
Ce lemme fournit la hiérarchie suivante pour les morphismes de champs
algébriques.

\begin{lemma}
\label{stacks-morphisms-lemma-hierarchy}
Un morphisme $f : \mathcal{X} \to \mathcal{Y}$ de champs algébriques est
\begin{enumerate}
\item un monomorphisme si et seulement si $\Delta_{f, 1}$ est un
isomorphisme, et
\item représentable par des espaces algébriques si et seulement si
$\Delta_{f, 1}$ est un monomorphisme.
\end{enumerate}
De plus, la seconde diagonale $\Delta_{f, 2}$ est toujours un monomorphisme.
\end{lemma}

\begin{proof}
Rappelons, d'après Propriétés des champs, Lemme
\ref{stacks-properties-lemma-monomorphism}, qu'un morphisme de champs
algébriques est un monomorphisme si et seulement si sa diagonale est un
isomorphisme de champs. Le Lemme \ref{stacks-morphisms-lemma-second-diagonal} peut donc se
reformuler en disant qu'un morphisme est représentable par des espaces
algébriques si sa diagonale est un monomorphisme. En particulier, il montre
que la condition (3) du Lemme
\ref{stacks-morphisms-lemma-properties-diagonal-representable} est en fait une condition
nécessaire et suffisante : un morphisme de champs algébriques est
représentable par des espaces algébriques si et seulement si sa diagonale
est un monomorphisme.
\end{proof}

\begin{lemma}
\label{stacks-morphisms-lemma-first-diagonal-separated-second-diagonal-closed}
Soit $f : \mathcal{X} \to \mathcal{Y}$ un morphisme de champs algébriques.
Alors
\begin{enumerate}
\item $\Delta_{f, 1}$ est séparé $\Leftrightarrow$
$\Delta_{f, 2}$ est une immersion fermée $\Leftrightarrow$
$\Delta_{f, 2}$ est propre $\Leftrightarrow$
$\Delta_{f, 2}$ est universellement fermé,
\item $\Delta_{f, 1}$ est quasi-séparé $\Leftrightarrow$
$\Delta_{f, 2}$ est de type fini $\Leftrightarrow$
$\Delta_{f, 2}$ est quasi-compact, et
\item $\Delta_{f, 1}$ est localement séparé $\Leftrightarrow$
$\Delta_{f, 2}$ est une immersion.
\end{enumerate}
\end{lemma}

\begin{proof}
Cela résulte des Lemmes \ref{stacks-morphisms-lemma-representable-separated-diagonal-closed},
\ref{stacks-morphisms-lemma-representable-quasi-separated-diagonal-quasi-compact} et
\ref{stacks-morphisms-lemma-representable-locally-separated-diagonal-immersion}
appliqués à $\Delta_{f, 1}$.
\end{proof}

\noindent
Le lemme suivant est assez élégant et suggère peut-être une généralisation de
ces conditions aux champs algébriques supérieurs.

\begin{lemma}
\label{stacks-morphisms-lemma-definition-separated}
Soit $f : \mathcal{X} \to \mathcal{Y}$ un morphisme de champs algébriques.
Alors
\begin{enumerate}
\item $f$ est séparé si et seulement si $\Delta_{f, 1}$ et
$\Delta_{f, 2}$ sont universellement fermés, et
\item $f$ est quasi-séparé si et seulement si $\Delta_{f, 1}$ et
$\Delta_{f, 2}$ sont quasi-compacts.
\item $f$ est quasi-DM si et seulement si $\Delta_{f, 1}$ et
$\Delta_{f, 2}$ sont localement quasi-finis.
\item $f$ est DM si et seulement si $\Delta_{f, 1}$ et $\Delta_{f, 2}$
sont non ramifiés.
\end{enumerate}
\end{lemma}

\begin{proof}
Démonstration de (1). Supposons $\Delta_{f, 2}$ et $\Delta_{f, 1}$
universellement fermés. Alors $\Delta_{f, 1}$ est séparé et universellement
fermé d'après le Lemme
\ref{stacks-morphisms-lemma-first-diagonal-separated-second-diagonal-closed}. D'après
Morphismes d'espaces, Lemme
\ref{spaces-morphisms-lemma-universally-closed-quasi-compact}, et Champs
algébriques, Lemme
\ref{algebraic-lemma-representable-transformations-property-implication},
nous voyons que $\Delta_{f, 1}$ est quasi-compact. Il est donc quasi-compact,
séparé, universellement fermé et localement de type fini (d'après le Lemme
\ref{stacks-morphisms-lemma-properties-diagonal}), donc propre. Cela démontre l'implication
``$\Leftarrow$'' de (1). La démonstration de l'implication réciproque est
omise.

\medskip\noindent
Démonstration de (2). Cela résulte immédiatement du Lemme
\ref{stacks-morphisms-lemma-first-diagonal-separated-second-diagonal-closed}.

\medskip\noindent
Démonstration de (3). Cela résulte du fait que $\Delta_{f, 2}$ est toujours
localement quasi-fini, d'après le Lemme
\ref{stacks-morphisms-lemma-properties-diagonal-representable} appliqué à
$\Delta_f = \Delta_{f, 1}$.

\medskip\noindent
Démonstration de (4). Cela résulte du fait que $\Delta_{f, 2}$ est toujours
non ramifié : le Lemme \ref{stacks-morphisms-lemma-properties-diagonal-representable},
appliqué à $\Delta_f = \Delta_{f, 1}$, montre que $\Delta_{f, 2}$ est
localement de type fini et est un monomorphisme. Voir Compléments sur les
morphismes d'espaces, Lemme
\ref{spaces-more-morphisms-lemma-universally-injective-unramified}.
\end{proof}

\begin{lemma}
\label{stacks-morphisms-lemma-separated-implies-isom}
Soit $f : \mathcal{X} \to \mathcal{Y}$ un morphisme séparé
(resp.\ quasi-séparé, resp.\ quasi-DM, resp.\ DM) de champs algébriques.
Alors
\begin{enumerate}
\item pour des espaces algébriques $T_i$, $i = 1, 2$, et des morphismes
$x_i : T_i \to \mathcal{X}$, en posant $y_i = f \circ x_i$, le morphisme
$$
T_1 \times_{x_1, \mathcal{X}, x_2} T_2 \longrightarrow
T_1 \times_{y_1, \mathcal{Y}, y_2} T_2
$$
est propre (resp.\ quasi-compact et quasi-séparé,
resp.\ localement quasi-fini, resp.\ non ramifié),
\item pour un espace algébrique $T$ et des morphismes
$x_i : T \to \mathcal{X}$, $i = 1, 2$, en posant
$y_i = f \circ x_i$, le morphisme
$$
\mathit{Isom}_\mathcal{X}(x_1, x_2) \longrightarrow
\mathit{Isom}_\mathcal{Y}(y_1, y_2)
$$
est propre (resp.\ quasi-compact et quasi-séparé,
resp.\ localement quasi-fini, resp.\ non ramifié).
\end{enumerate}
\end{lemma}

\begin{proof}
Démonstration de (1). Observons que le diagramme
$$
\xymatrix{
T_1 \times_{x_1, \mathcal{X}, x_2} T_2 \ar[d] \ar[r] &
T_1 \times_{y_1, \mathcal{Y}, y_2} T_2 \ar[d] \\
\mathcal{X} \ar[r] & \mathcal{X} \times_\mathcal{Y} \mathcal{X}
}
$$
est cartésien. L'assertion résulte donc du fait que $f$ est séparé
(resp.\ quasi-séparé, resp.\ quasi-DM, resp.\ DM) si et seulement si sa
diagonale est propre (resp.\ quasi-compacte et quasi-séparée,
resp.\ localement quasi-finie, resp.\ non ramifiée).

\medskip\noindent
Démonstration de (2). Cela résulte de l'égalité
$$
\mathit{Isom}_\mathcal{X}(x_1, x_2) =
(T \times_{x_1, \mathcal{X}, x_2} T) \times_{T \times T, \Delta_T} T
$$
et le morphisme de (2) est donc un changement de base du morphisme de (1).
\end{proof}








\section{Morphismes quasi-compacts}
\label{stacks-morphisms-section-quasi-compact}

\noindent
Soit $f$ un morphisme de champs algébriques représentable par des espaces
algébriques. Dans Propriétés des champs, section
\ref{stacks-properties-section-properties-morphisms}
nous avons défini ce que signifie, pour $f$, être quasi-compact.
En voici une autre caractérisation.

\begin{lemma}
\label{stacks-morphisms-lemma-characterize-representable-quasi-compact}
Soit $f : \mathcal{X} \to \mathcal{Y}$ un morphisme de champs algébriques
représentable par des espaces algébriques. Les conditions suivantes sont
équivalentes :
\begin{enumerate}
\item $f$ est quasi-compact (au sens de Propriétés des champs, section
\ref{stacks-properties-section-properties-morphisms}), et
\item pour tout champ algébrique quasi-compact $\mathcal{Z}$ et tout
morphisme $\mathcal{Z} \to \mathcal{Y}$, le champ algébrique
$\mathcal{Z} \times_\mathcal{Y} \mathcal{X}$ est quasi-compact.
\end{enumerate}
\end{lemma}

\begin{proof}
Supposons (1), et soit $\mathcal{Z} \to \mathcal{Y}$ un morphisme de champs
algébriques, avec $\mathcal{Z}$ quasi-compact. D'après Propriétés des champs,
Lemme \ref{stacks-properties-lemma-quasi-compact-stack}, il existe un schéma
quasi-compact $U$ et un morphisme lisse surjectif
$U \to \mathcal{Z}$. Puisque $f$ est représentable par des espaces
algébriques et quasi-compact, la définition montre que
$U \times_\mathcal{Y} \mathcal{X}$ est un espace algébrique et que
$U \times_\mathcal{Y} \mathcal{X} \to U$ est quasi-compact.
Ainsi, $U \times_\mathcal{Y} \mathcal{X}$ est un espace algébrique
quasi-compact. Le morphisme
$U \times_\mathcal{Y} \mathcal{X} \to
\mathcal{Z} \times_\mathcal{Y} \mathcal{X}$
est lisse et surjectif (comme changement de base du morphisme lisse et
surjectif $U \to \mathcal{Z}$). Par conséquent,
$\mathcal{Z} \times_\mathcal{Y} \mathcal{X}$ est quasi-compact par une
nouvelle application de Propriétés des champs, Lemme
\ref{stacks-properties-lemma-quasi-compact-stack}.

\medskip\noindent
Supposons (2). Soit $Z \to \mathcal{Y}$ un morphisme, où $Z$ est un schéma.
Il faut montrer que le morphisme d'espaces algébriques
$p : Z \times_\mathcal{Y} \mathcal{X} \to Z$ est quasi-compact.
Soit $U \subset Z$ un ouvert affine. Alors
$p^{-1}(U) = U \times_\mathcal{Y} \mathcal{X}$,
et l'espace algébrique $U \times_\mathcal{Y} \mathcal{X}$ est quasi-compact
par l'hypothèse (2). Par conséquent, $p$ est quasi-compact, voir Morphismes
d'espaces, Lemme \ref{spaces-morphisms-lemma-quasi-compact-local}.
\end{proof}

\noindent
Ceci motive la définition suivante.

\begin{definition}
\label{stacks-morphisms-definition-quasi-compact}
Soit $f : \mathcal{X} \to \mathcal{Y}$ un morphisme de champs algébriques.
Nous disons que $f$ est {\it quasi-compact} si, pour tout champ algébrique
quasi-compact $\mathcal{Z}$ et tout morphisme
$\mathcal{Z} \to \mathcal{Y}$, le produit fibré
$\mathcal{Z} \times_\mathcal{Y} \mathcal{X}$ est quasi-compact.
\end{definition}

\noindent
D'après le Lemme \ref{stacks-morphisms-lemma-characterize-representable-quasi-compact}
ci-dessus, cette définition coïncide avec la notion déjà existante pour les
morphismes de champs algébriques représentables par des espaces algébriques.
En particulier, cette notion coïncide avec celles déjà définies pour les
morphismes entre champs algébriques et schémas.

\begin{lemma}
\label{stacks-morphisms-lemma-base-change-quasi-compact}
Le changement de base d'un morphisme quasi-compact de champs algébriques par
un morphisme quelconque de champs algébriques est quasi-compact.
\end{lemma}

\begin{proof}
Démonstration omise.
\end{proof}

\begin{lemma}
\label{stacks-morphisms-lemma-composition-quasi-compact}
Le composé de deux morphismes quasi-compacts de champs algébriques est
quasi-compact.
\end{lemma}

\begin{proof}
Démonstration omise.
\end{proof}

\begin{lemma}
\label{stacks-morphisms-lemma-closed-immersion-quasi-compact}
Une immersion fermée de champs algébriques est quasi-compacte.
\end{lemma}

\begin{proof}
Cela résulte du fait que les immersions sont toujours représentables et de
l'énoncé correspondant pour les immersions fermées d'espaces algébriques.
\end{proof}

\begin{lemma}
\label{stacks-morphisms-lemma-surjection-from-quasi-compact}
Soit
$$
\xymatrix{
\mathcal{X} \ar[rr]_f \ar[rd]_p & &
\mathcal{Y} \ar[dl]^q \\
& \mathcal{Z}
}
$$
un diagramme $2$-commutatif de morphismes de champs algébriques.
Si $f$ est surjectif et $p$ quasi-compact, alors $q$ est quasi-compact.
\end{lemma}

\begin{proof}
Soit $\mathcal{T}$ un champ algébrique quasi-compact et soit
$\mathcal{T} \to \mathcal{Z}$ un morphisme. D'après
Propriétés des champs,
Lemme \ref{stacks-properties-lemma-base-change-surjective},
le morphisme
$\mathcal{T} \times_\mathcal{Z} \mathcal{X} \to
\mathcal{T} \times_\mathcal{Z} \mathcal{Y}$
est surjectif et, par hypothèse,
$\mathcal{T} \times_\mathcal{Z} \mathcal{X}$ est quasi-compact. Il s'ensuit
que $\mathcal{T} \times_\mathcal{Z} \mathcal{Y}$ est quasi-compact d'après
Propriétés des champs, Lemme
\ref{stacks-properties-lemma-quasi-compact-stack}.
\end{proof}

\begin{lemma}
\label{stacks-morphisms-lemma-quasi-compact-permanence}
Soient $f : \mathcal{X} \to \mathcal{Y}$ et
$g : \mathcal{Y} \to \mathcal{Z}$ des morphismes de champs algébriques.
Si $g \circ f$ est quasi-compact et si $g$ est quasi-séparé,
alors $f$ est quasi-compact.
\end{lemma}

\begin{proof}
En effet, $f$ est le composé
$(1, f) : \mathcal{X} \to \mathcal{X} \times_\mathcal{Z} \mathcal{Y} \to
\mathcal{Y}$. Le premier morphisme est quasi-compact d'après le
Lemme \ref{stacks-morphisms-lemma-section-immersion}, car c'est une section du morphisme
quasi-séparé
$\mathcal{X} \times_\mathcal{Z} \mathcal{Y} \to \mathcal{X}$
(un changement de base de $g$; voir le
Lemme \ref{stacks-morphisms-lemma-base-change-separated}). Le second morphisme est
quasi-compact puisqu'il est le changement de base de $f$; voir le
Lemme \ref{stacks-morphisms-lemma-base-change-quasi-compact}. Enfin, les composés de
morphismes quasi-compacts sont quasi-compacts; voir le
Lemme \ref{stacks-morphisms-lemma-composition-quasi-compact}.
\end{proof}

\begin{lemma}
\label{stacks-morphisms-lemma-quasi-compact-quasi-separated-permanence}
Soit $f : \mathcal{X} \to \mathcal{Y}$ un morphisme de champs algébriques.
\begin{enumerate}
\item Si $\mathcal{X}$ est quasi-compact et si $\mathcal{Y}$ est
quasi-séparé, alors $f$ est quasi-compact.
\item Si $\mathcal{X}$ est quasi-compact et quasi-séparé et si $\mathcal{Y}$
est quasi-séparé, alors $f$ est quasi-compact et quasi-séparé.
\item Un produit fibré de champs algébriques quasi-compacts et quasi-séparés
est quasi-compact et quasi-séparé.
\end{enumerate}
\end{lemma}

\begin{proof}
L'assertion (1) résulte du
Lemme \ref{stacks-morphisms-lemma-quasi-compact-permanence}. L'assertion (2) résulte de (1)
et du Lemme \ref{stacks-morphisms-lemma-compose-after-separated}. Pour (3), soient
$\mathcal{X} \to \mathcal{Y}$ et $\mathcal{Z} \to \mathcal{Y}$ des
morphismes de champs algébriques quasi-compacts et quasi-séparés. Alors
$\mathcal{X} \times_\mathcal{Y} \mathcal{Z} \to \mathcal{Z}$ est
quasi-compact et quasi-séparé, car il s'obtient par changement de base à
partir de $\mathcal{X} \to \mathcal{Y}$, en vertu de (2) et des Lemmes
\ref{stacks-morphisms-lemma-base-change-quasi-compact} et
\ref{stacks-morphisms-lemma-base-change-separated}. Par conséquent,
$\mathcal{X} \times_\mathcal{Y} \mathcal{Z}$ est quasi-compact et
quasi-séparé en tant que champ algébrique quasi-compact et quasi-séparé sur
$\mathcal{Z}$; voir les Lemmes \ref{stacks-morphisms-lemma-separated-over-separated} et
\ref{stacks-morphisms-lemma-composition-quasi-compact}.
\end{proof}

\begin{lemma}
\label{stacks-morphisms-lemma-reach-points-scheme-theoretic-image}
Soit $f : \mathcal{X} \to \mathcal{Y}$ un morphisme quasi-compact de
champs algébriques. Soit $y \in |\mathcal{Y}|$ un point de l'adhérence de
l'image de $|f|$. Il existe un anneau de valuation $A$, de corps des
fractions $K$, et un diagramme commutatif
$$
\xymatrix{
\Spec(K) \ar[r] \ar[d] & \mathcal{X} \ar[d] \\
\Spec(A) \ar[r] & \mathcal{Y}
}
$$
tels que le point fermé de $\Spec(A)$ s'envoie sur $y$.
\end{lemma}

\begin{proof}
Choisissons un schéma affine $V$, un point $v \in V$ et un morphisme lisse
$V \to \mathcal{Y}$ qui envoie $v$ sur $y$. Considérons le diagramme de
changement de base
$$
\xymatrix{
V \times_\mathcal{Y} \mathcal{X} \ar[r] \ar[d]_g & \mathcal{X} \ar[d]^f \\
V \ar[r] & \mathcal{Y}
}
$$
Rappelons que $|V \times_\mathcal{Y} \mathcal{X}| \to
|V| \times_{|\mathcal{Y}|} |\mathcal{X}|$ est surjectif
(Propriétés des champs, Lemme
\ref{stacks-properties-lemma-points-cartesian}). Comme
$|V| \to |\mathcal{Y}|$ est ouvert (Propriétés des champs, Lemme
\ref{stacks-properties-lemma-topology-points}), nous en déduisons que $v$
appartient à l'adhérence de l'image de $|g|$. Il suffit donc de démontrer le
lemme pour le morphisme quasi-compact $g$
(Lemme \ref{stacks-morphisms-lemma-base-change-quasi-compact}), ce que nous faisons au
paragraphe suivant.

\medskip\noindent
Supposons que $\mathcal{Y} = Y$ soit un schéma affine. Alors $\mathcal{X}$
est quasi-compact, puisque $f$ est quasi-compact
(Définition \ref{stacks-morphisms-definition-quasi-compact}). Choisissons un schéma affine
$W$ et un morphisme lisse surjectif $W \to \mathcal{X}$. L'image de $|f|$
est alors l'image de $W \to Y$. D'après Morphismes, Lemme
\ref{morphisms-lemma-reach-points-scheme-theoretic-image}, nous pouvons
choisir un diagramme
$$
\xymatrix{
\Spec(K) \ar[r] \ar[d] & W \ar[d] \ar[r] & \mathcal{X} \ar[d] \\
\Spec(A) \ar[r] & Y \ar[r] & Y
}
$$
tel que le point fermé de $\Spec(A)$ s'envoie sur $y$.
En composant avec $W \to \mathcal{X}$, nous obtenons le diagramme voulu.
\end{proof}

\begin{lemma}
\label{stacks-morphisms-lemma-check-quasi-compact-covering}
Soit $f : \mathcal{X} \to \mathcal{Y}$ un morphisme de champs algébriques.
Soit $W \to \mathcal{Y}$ un morphisme surjectif, plat et localement de
présentation finie, où $W$ est un espace algébrique. Si le changement de base
$W \times_\mathcal{Y} \mathcal{X} \to W$ est quasi-compact, alors $f$ est
quasi-compact.
\end{lemma}

\begin{proof}
Supposons que $W \times_\mathcal{Y} \mathcal{X} \to W$ soit quasi-compact.
Soit $\mathcal{Z} \to \mathcal{Y}$ un morphisme, où $\mathcal{Z}$ est un
champ algébrique quasi-compact. Choisissons un schéma $U$ et un morphisme
lisse surjectif $U \to W \times_\mathcal{Y} \mathcal{Z}$. Puisque
$U \to \mathcal{Z}$ est plat, surjectif et localement de présentation finie,
et puisque $\mathcal{Z}$ est quasi-compact, il existe un sous-schéma ouvert
quasi-compact $U' \subset U$ tel que $U' \to \mathcal{Z}$ soit surjectif.
Alors
$U' \times_\mathcal{Y} \mathcal{X} =
U' \times_W (W \times_\mathcal{Y} \mathcal{X})$
est quasi-compact par hypothèse et se surjecte sur
$\mathcal{Z} \times_\mathcal{Y} \mathcal{X}$. Par conséquent,
$\mathcal{Z} \times_\mathcal{Y} \mathcal{X}$ est quasi-compact, comme voulu.
\end{proof}











\section{Champs algébriques noethériens}
\label{stacks-morphisms-section-noetherian}

\noindent
Nous avons déjà défini les champs algébriques localement noethériens dans
Propriétés des champs, section
\ref{stacks-properties-section-types-properties}.

\begin{definition}
\label{stacks-morphisms-definition-noetherian}
Soit $\mathcal{X}$ un champ algébrique. Nous disons que $\mathcal{X}$ est
{\it noethérien} si $\mathcal{X}$ est quasi-compact, quasi-séparé et
localement noethérien.
\end{definition}

\noindent
Remarquons qu'un champ algébrique noethérien $\mathcal{X}$ n'est pas seulement
quasi-compact et localement noethérien, mais aussi quasi-séparé. Dans le
langage de la section \ref{stacks-morphisms-section-higher-diagonals}, si l'on note
$p : \mathcal{X} \to \Spec(\mathbf{Z})$ le morphisme structural
``absolu'' (c'est-à-dire le morphisme structural de $\mathcal{X}$ considéré
comme champ algébrique sur $\mathbf{Z}$), alors
$$
\mathcal{X}\text{ noethérien}
\Leftrightarrow
\mathcal{X}\text{ localement noethérien et }
\Delta_{p, 0}, \Delta_{p, 1}, \Delta_{p, 2}
\text{ quasi-compacts}.
$$
Cela impliquera plus loin qu'un champ algébrique de type fini sur un champ
algébrique noethérien n'est pas nécessairement noethérien.

\begin{lemma}
\label{stacks-morphisms-lemma-locally-closed-in-noetherian}
Soit $j : \mathcal{X} \to \mathcal{Y}$ une immersion de champs algébriques.
\begin{enumerate}
\item Si $\mathcal{Y}$ est localement noethérien, alors
$\mathcal{X}$ est localement noethérien et $j$ est quasi-compact.
\item Si $\mathcal{Y}$ est noethérien, alors $\mathcal{X}$ est noethérien.
\end{enumerate}
\end{lemma}

\begin{proof}
Choisissons un schéma $V$ et un morphisme lisse surjectif
$V \to \mathcal{Y}$. Alors $U = \mathcal{X} \times_\mathcal{Y} V$ est un
schéma et $U \to V$ est une immersion; voir Propriétés des champs,
Définition \ref{stacks-properties-definition-immersion}. Rappelons que
$\mathcal{Y}$ est localement noethérien si et seulement si $V$ l'est. Dans ce
cas, $U$ est lui aussi localement noethérien (Morphismes, Lemmes
\ref{morphisms-lemma-immersion-locally-finite-type} et
\ref{morphisms-lemma-finite-type-noetherian}), et $U \to V$ est quasi-compact
(Propriétés, Lemme \ref{properties-lemma-immersion-into-noetherian}). Cela
montre que $j$ est quasi-compact
(Lemme \ref{stacks-morphisms-lemma-check-quasi-compact-covering}) et que $\mathcal{X}$ est
localement noethérien. Enfin, si $\mathcal{Y}$ est noethérien, ce qui précède
montre que $\mathcal{X}$ est quasi-compact et localement noethérien. Pour
achever la démonstration, observons que $j$ est séparé; par conséquent,
$\mathcal{X}$ est quasi-séparé puisque $\mathcal{Y}$ l'est, d'après le
Lemme \ref{stacks-morphisms-lemma-separated-over-separated}.
\end{proof}

\begin{lemma}
\label{stacks-morphisms-lemma-Noetherian-topology}
Soit $\mathcal{X}$ un champ algébrique.
\begin{enumerate}
\item Si $\mathcal{X}$ est localement noethérien, alors $|\mathcal{X}|$ est
un espace topologique localement noethérien.
\item Si $\mathcal{X}$ est quasi-compact et localement noethérien, alors
$|\mathcal{X}|$ est un espace topologique noethérien.
\end{enumerate}
\end{lemma}

\begin{proof}
Supposons $\mathcal{X}$ localement noethérien. Choisissons un schéma $U$ et un
morphisme lisse surjectif $U \to \mathcal{X}$. Comme $\mathcal{X}$ est
localement noethérien, $U$ l'est aussi. D'après Propriétés, Lemme
\ref{properties-lemma-Noetherian-topology}, cela signifie que $|U|$ est un
espace topologique localement noethérien. Puisque
$|U| \to |\mathcal{X}|$ est ouvert et surjectif, nous en déduisons que
$|\mathcal{X}|$ est localement noethérien d'après Topologie, Lemme
\ref{topology-lemma-image-Noetherian}. Cela démontre (1). Si $\mathcal{X}$ est
quasi-compact et localement noethérien, alors $|\mathcal{X}|$ est
quasi-compact et localement noethérien. Par conséquent, $|\mathcal{X}|$ est
noethérien d'après Topologie, Lemme
\ref{topology-lemma-quasi-compact-locally-Noetherian-Noetherian}.
\end{proof}

\begin{lemma}
\label{stacks-morphisms-lemma-locally-Noetherian-quasi-sober}
Soit $\mathcal{X}$ un champ algébrique localement noethérien.
Alors $|\mathcal{X}|$ est quasi-sobre (Topologie, Définition
\ref{topology-definition-generic-point}).
\end{lemma}

\begin{proof}
Il faut démontrer que toute partie fermée irréductible
$T \subset |\mathcal{X}|$ possède un point générique. Choisissons un schéma
affine $U$ et un morphisme lisse $f : U \to \mathcal{X}$ tels que
$f^{-1}(T) \subset |U|$ soit non vide. Puisque $U$ est noethérien, la partie
fermée $f^{-1}(T)$ n'a qu'un nombre fini de composantes irréductibles
(Topologie, Lemme \ref{topology-lemma-Noetherian}). Écrivons
$f^{-1}(T) = Z_1 \cup \ldots \cup Z_n$ pour sa décomposition en composantes
irréductibles. Comme $f$ est ouvert, l'image de
$f|_{f^{-1}(T)} : f^{-1}(T) \to T$ contient une partie ouverte non vide de
$T$. Puisque $T$ est irréductible, cela signifie que $f(f^{-1}(T))$ est
dense. L'irréductibilité de $T$ entraîne alors que $f(Z_i)$ est dense pour un
certain $i$. Si $\xi_i \in Z_i$ est le point générique, $f(\xi_i)$ est donc
un point générique de $T$.
\end{proof}





\section{Morphismes affines}
\label{stacks-morphisms-section-affine}

\noindent
Les morphismes affines de champs algébriques se définissent comme suit.

\begin{definition}
\label{stacks-morphisms-definition-affine}
Un morphisme de champs algébriques est dit {\it affine} s'il est représentable
et affine au sens de Propriétés des champs, section
\ref{stacks-properties-section-properties-morphisms}.
\end{definition}

\noindent
Il nous est un peu plus commode de considérer un morphisme affine de champs
algébriques comme un morphisme de champs algébriques représentable par des
espaces algébriques et affine au sens de Propriétés des champs, section
\ref{stacks-properties-section-properties-morphisms}.
(Rappelons que, par défaut, ``représentable'' signifie dans le Champs Project
représentable par des schémas.) Puisque cela équivaut manifestement à la notion
qui vient d'être définie, nous emploierons cette caractérisation sans autre
mention. Démontrons quelques lemmes simples sur cette notion.

\begin{lemma}
\label{stacks-morphisms-lemma-base-change-affine}
Soit $\mathcal{X} \to \mathcal{Y}$ un morphisme de champs algébriques.
Soit $\mathcal{Z} \to \mathcal{Y}$ un morphisme affine de champs algébriques.
Alors $\mathcal{Z} \times_\mathcal{Y} \mathcal{X} \to \mathcal{X}$ est un
morphisme affine de champs algébriques.
\end{lemma}

\begin{proof}
Cela résulte de la discussion de Propriétés des champs, section
\ref{stacks-properties-section-properties-morphisms}.
\end{proof}

\begin{lemma}
\label{stacks-morphisms-lemma-composition-affine}
Les composés de morphismes affines de champs algébriques sont affines.
\end{lemma}

\begin{proof}
Cela résulte de la discussion de Propriétés des champs, section
\ref{stacks-properties-section-properties-morphisms}, et de Morphismes
d'espaces, Lemme \ref{spaces-morphisms-lemma-composition-affine}.
\end{proof}

\begin{lemma}
\label{stacks-morphisms-lemma-affine-permanence}
Soit
$$
\xymatrix{
\mathcal{X} \ar[rr]_f \ar[rd]_a & & \mathcal{Y} \ar[dl]^b \\
& \mathcal{Z}
}
$$
un diagramme commutatif de morphismes de champs algébriques.
Si $a$ et $\Delta_b$ sont affines, alors $f$ est affine.
\end{lemma}

\begin{proof}
Le changement de base
$\text{pr}_2 : \mathcal{X} \times_\mathcal{Z} \mathcal{Y} \to \mathcal{Y}$
de $a$ est affine d'après le Lemme \ref{stacks-morphisms-lemma-base-change-affine}. Le morphisme
$(1, f) : \mathcal{X} \to \mathcal{X} \times_\mathcal{Z} \mathcal{Y}$
est le changement de base de
$\Delta_b : \mathcal{Y} \to \mathcal{Y} \times_\mathcal{Z} \mathcal{Y}$
par le morphisme $\mathcal{X} \times_\mathcal{Z} \mathcal{Y} \to
\mathcal{Y} \times_\mathcal{Z} \mathcal{Y}$ (voir Catégories, section
\ref{categories-section-2-fibre-products}). Il est donc affine d'après le
Lemme \ref{stacks-morphisms-lemma-base-change-affine}. Le composé
$f = \text{pr}_2 \circ (1, f)$ de morphismes affines est affine d'après le
Lemme \ref{stacks-morphisms-lemma-composition-affine}, ce qui achève la démonstration.
\end{proof}





\section{Morphismes intégraux et finis}
\label{stacks-morphisms-section-integral}

\noindent
Les morphismes intégraux et finis de champs algébriques se définissent comme
suit.

\begin{definition}
\label{stacks-morphisms-definition-integral}
Soit $f : \mathcal{X} \to \mathcal{Y}$ un morphisme de champs algébriques.
\begin{enumerate}
\item Nous disons que $f$ est {\it intégral} si $f$ est représentable et
intégral au sens de Propriétés des champs, section
\ref{stacks-properties-section-properties-morphisms}.
\item Nous disons que $f$ est {\it fini} si $f$ est représentable et fini au
sens de Propriétés des champs, section
\ref{stacks-properties-section-properties-morphisms}.
\end{enumerate}
\end{definition}

\noindent
Il nous est un peu plus commode de considérer un morphisme intégral, resp.\
fini, de champs algébriques comme un morphisme de champs algébriques
représentable par des espaces algébriques et intégral, resp.\ fini, au sens de
Propriétés des champs, section
\ref{stacks-properties-section-properties-morphisms}.
(Rappelons que, par défaut, ``représentable'' signifie dans le Champs Project
représentable par des schémas.) Puisque cela équivaut manifestement à la notion
qui vient d'être définie, nous emploierons cette caractérisation sans autre
mention. Démontrons quelques lemmes simples sur cette notion.

\begin{lemma}
\label{stacks-morphisms-lemma-base-change-integral}
Soit $\mathcal{X} \to \mathcal{Y}$ un morphisme de champs algébriques.
Soit $\mathcal{Z} \to \mathcal{Y}$ un morphisme intégral (ou fini) de champs
algébriques. Alors
$\mathcal{Z} \times_\mathcal{Y} \mathcal{X} \to \mathcal{X}$
est un morphisme intégral (ou fini) de champs algébriques.
\end{lemma}

\begin{proof}
Cela résulte de la discussion de Propriétés des champs, section
\ref{stacks-properties-section-properties-morphisms}.
\end{proof}

\begin{lemma}
\label{stacks-morphisms-lemma-composition-integral}
Les composés de morphismes intégraux, resp.\ finis, de champs algébriques sont
intégraux, resp.\ finis.
\end{lemma}

\begin{proof}
Cela résulte de la discussion de Propriétés des champs, section
\ref{stacks-properties-section-properties-morphisms}, et de Morphismes
d'espaces, Lemme \ref{spaces-morphisms-lemma-composition-integral}.
\end{proof}




\section{Morphismes ouverts}
\label{stacks-morphisms-section-open}

\noindent
Soit $f$ un morphisme de champs algébriques qui est représentable par
des espaces algébriques. Dans Propriétés des champs, section
\ref{stacks-properties-section-properties-morphisms}
nous avons défini ce que signifie pour $f$ d'être universellement ouvert.
En voici une autre caractérisation.

\begin{lemma}
\label{stacks-morphisms-lemma-characterize-representable-universally-open}
Soit $f : \mathcal{X} \to \mathcal{Y}$ un morphisme de
champs algébriques qui est représentable par des espaces algébriques.
Les assertions suivantes sont équivalentes :
\begin{enumerate}
\item $f$ est universellement ouvert (au sens de Propriétés des champs,
section \ref{stacks-properties-section-properties-morphisms}) ;
\item pour tout morphisme de champs algébriques
$\mathcal{Z} \to \mathcal{Y}$, l'application d'espaces topologiques
$|\mathcal{Z} \times_\mathcal{Y} \mathcal{X}| \to |\mathcal{Z}|$ est ouverte.
\end{enumerate}
\end{lemma}

\begin{proof}
Supposons (1), et soit $\mathcal{Z} \to \mathcal{Y}$ comme en (2).
Choisissons un schéma $V$ et un morphisme lisse surjectif
$V \to \mathcal{Z}$. Par hypothèse, le morphisme d'espaces algébriques
$V \times_\mathcal{Y} \mathcal{X} \to V$ est universellement ouvert ;
en particulier, l'application
$|V \times_\mathcal{Y} \mathcal{X}| \to |V|$ est ouverte. D'après
Propriétés des champs, section \ref{stacks-properties-section-points},
dans le diagramme commutatif
$$
\xymatrix{
|V \times_\mathcal{Y} \mathcal{X}| \ar[r] \ar[d] &
|\mathcal{Z} \times_\mathcal{Y} \mathcal{X}| \ar[d] \\
|V| \ar[r] & |\mathcal{Z}|
}
$$
les flèches horizontales sont ouvertes et surjectives ; de plus,
$$
|V \times_\mathcal{Y} \mathcal{X}| \longrightarrow
|V| \times_{|\mathcal{Z}|} |\mathcal{Z} \times_\mathcal{Y} \mathcal{X}|
$$
est surjective. Puisque la flèche verticale de gauche est ouverte,
il s'ensuit que celle de droite est ouverte. Cela démontre (2).
L'implication (2) $\Rightarrow$ (1) résulte des définitions.
\end{proof}

\noindent
Nous pouvons donc adopter la définition naturelle suivante.

\begin{definition}
\label{stacks-morphisms-definition-open}
Soit $f : \mathcal{X} \to \mathcal{Y}$ un morphisme de champs algébriques.
\begin{enumerate}
\item Nous disons que $f$ est {\it ouvert} si l'application d'espaces
topologiques $|\mathcal{X}| \to |\mathcal{Y}|$ est ouverte.
\item Nous disons que $f$ est {\it universellement ouvert} si, pour tout
morphisme de champs algébriques $\mathcal{Z} \to \mathcal{Y}$,
l'application d'espaces topologiques
$$
|\mathcal{Z} \times_\mathcal{Y} \mathcal{X}| \to |\mathcal{Z}|
$$
est ouverte, c'est-à-dire si le changement de base
$\mathcal{Z} \times_\mathcal{Y} \mathcal{X} \to \mathcal{Z}$ est ouvert.
\end{enumerate}
\end{definition}

\begin{lemma}
\label{stacks-morphisms-lemma-base-change-universally-open}
Le changement de base d'un morphisme universellement ouvert de champs
algébriques par un morphisme quelconque de champs algébriques est
universellement ouvert.
\end{lemma}

\begin{proof}
C'est immédiat par définition.
\end{proof}

\begin{lemma}
\label{stacks-morphisms-lemma-composition-universally-open}
Le composé de deux morphismes (universellement) ouverts de champs
algébriques est (universellement) ouvert.
\end{lemma}

\begin{proof}
Démonstration omise.
\end{proof}







\section{Morphismes submersifs}
\label{stacks-morphisms-section-submersive}

\noindent
Soit $f$ un morphisme de champs algébriques qui est représentable par
des espaces algébriques. Dans Propriétés des champs, section
\ref{stacks-properties-section-properties-morphisms}
nous avons défini ce que signifie pour $f$ d'être universellement submersif.
En voici une autre caractérisation.

\begin{lemma}
\label{stacks-morphisms-lemma-characterize-representable-universally-submersive}
Soit $f : \mathcal{X} \to \mathcal{Y}$ un morphisme de
champs algébriques qui est représentable par des espaces algébriques.
Les assertions suivantes sont équivalentes :
\begin{enumerate}
\item $f$ est universellement submersif (au sens de Propriétés des champs,
section \ref{stacks-properties-section-properties-morphisms}) ;
\item pour tout morphisme de champs algébriques
$\mathcal{Z} \to \mathcal{Y}$, l'application d'espaces topologiques
$|\mathcal{Z} \times_\mathcal{Y} \mathcal{X}| \to |\mathcal{Z}|$
est submersive.
\end{enumerate}
\end{lemma}

\begin{proof}
Supposons (1), et soit $\mathcal{Z} \to \mathcal{Y}$ comme en (2).
Choisissons un schéma $V$ et un morphisme lisse surjectif
$V \to \mathcal{Z}$. Par hypothèse, le morphisme d'espaces algébriques
$V \times_\mathcal{Y} \mathcal{X} \to V$ est universellement submersif ;
en particulier, l'application
$|V \times_\mathcal{Y} \mathcal{X}| \to |V|$ est submersive. D'après
Propriétés des champs, section \ref{stacks-properties-section-points},
dans le diagramme commutatif
$$
\xymatrix{
|V \times_\mathcal{Y} \mathcal{X}| \ar[r] \ar[d] &
|\mathcal{Z} \times_\mathcal{Y} \mathcal{X}| \ar[d] \\
|V| \ar[r] & |\mathcal{Z}|
}
$$
les flèches horizontales sont ouvertes et surjectives ; de plus,
$$
|V \times_\mathcal{Y} \mathcal{X}| \longrightarrow
|V| \times_{|\mathcal{Z}|} |\mathcal{Z} \times_\mathcal{Y} \mathcal{X}|
$$
est surjective. Puisque la flèche verticale de gauche est submersive,
il s'ensuit que celle de droite est submersive. Cela démontre (2).
L'implication (2) $\Rightarrow$ (1) résulte des définitions.
\end{proof}

\noindent
Nous pouvons donc adopter la définition naturelle suivante.

\begin{definition}
\label{stacks-morphisms-definition-submersive}
Soit $f : \mathcal{X} \to \mathcal{Y}$ un morphisme de champs algébriques.
\begin{enumerate}
\item Nous disons que $f$ est {\it submersif}\footnote{Cette notion est
très différente de celle d'une submersion entre variétés différentielles.}
si l'application continue $|\mathcal{X}| \to |\mathcal{Y}|$ est submersive ;
voir Topologie, Définition \ref{topology-definition-submersive}.
\item Nous disons que $f$ est {\it universellement submersif} si, pour tout
morphisme de champs algébriques $\mathcal{Y}' \to \mathcal{Y}$,
le changement de base
$\mathcal{Y}' \times_\mathcal{Y} \mathcal{X} \to \mathcal{Y}'$
est submersif.
\end{enumerate}
\end{definition}

\noindent
Notons qu'un morphisme submersif est en particulier surjectif.

\begin{lemma}
\label{stacks-morphisms-lemma-base-change-universally-submersive}
Le changement de base d'un morphisme universellement submersif de champs
algébriques par un morphisme quelconque de champs algébriques est
universellement submersif.
\end{lemma}

\begin{proof}
C'est immédiat par définition.
\end{proof}

\begin{lemma}
\label{stacks-morphisms-lemma-composition-universally-submersive}
Le composé de deux morphismes (universellement) submersifs de champs
algébriques est (universellement) submersif.
\end{lemma}

\begin{proof}
Démonstration omise.
\end{proof}












\section{Morphismes universellement fermés}
\label{stacks-morphisms-section-universally-closed}

\noindent
Soit $f$ un morphisme de champs algébriques qui est représentable par
des espaces algébriques. Dans Propriétés des champs, section
\ref{stacks-properties-section-properties-morphisms}
nous avons défini ce que signifie pour $f$ d'être universellement fermé.
En voici une autre caractérisation.

\begin{lemma}
\label{stacks-morphisms-lemma-characterize-representable-universally-closed}
Soit $f : \mathcal{X} \to \mathcal{Y}$ un morphisme de
champs algébriques qui est représentable par des espaces algébriques.
Les assertions suivantes sont équivalentes :
\begin{enumerate}
\item $f$ est universellement fermé (au sens de Propriétés des champs,
section \ref{stacks-properties-section-properties-morphisms}) ;
\item pour tout morphisme de champs algébriques
$\mathcal{Z} \to \mathcal{Y}$, l'application d'espaces topologiques
$|\mathcal{Z} \times_\mathcal{Y} \mathcal{X}| \to |\mathcal{Z}|$ est fermée.
\end{enumerate}
\end{lemma}

\begin{proof}
Supposons (1), et soit $\mathcal{Z} \to \mathcal{Y}$ comme en (2).
Choisissons un schéma $V$ et un morphisme lisse surjectif
$V \to \mathcal{Z}$. Par hypothèse, le morphisme d'espaces algébriques
$V \times_\mathcal{Y} \mathcal{X} \to V$ est universellement fermé ;
en particulier, l'application
$|V \times_\mathcal{Y} \mathcal{X}| \to |V|$ est fermée. D'après
Propriétés des champs, section \ref{stacks-properties-section-points},
dans le diagramme commutatif
$$
\xymatrix{
|V \times_\mathcal{Y} \mathcal{X}| \ar[r] \ar[d] &
|\mathcal{Z} \times_\mathcal{Y} \mathcal{X}| \ar[d] \\
|V| \ar[r] & |\mathcal{Z}|
}
$$
les flèches horizontales sont ouvertes et surjectives ; de plus,
$$
|V \times_\mathcal{Y} \mathcal{X}| \longrightarrow
|V| \times_{|\mathcal{Z}|} |\mathcal{Z} \times_\mathcal{Y} \mathcal{X}|
$$
est surjective. Puisque la flèche verticale de gauche est fermée,
il s'ensuit que celle de droite est fermée. Cela démontre (2).
L'implication (2) $\Rightarrow$ (1) résulte des définitions.
\end{proof}

\noindent
Nous pouvons donc adopter la définition naturelle suivante.

\begin{definition}
\label{stacks-morphisms-definition-closed}
Soit $f : \mathcal{X} \to \mathcal{Y}$ un morphisme de champs algébriques.
\begin{enumerate}
\item Nous disons que $f$ est {\it fermé} si l'application d'espaces
topologiques $|\mathcal{X}| \to |\mathcal{Y}|$ est fermée.
\item Nous disons que $f$ est {\it universellement fermé} si, pour tout
morphisme de champs algébriques $\mathcal{Z} \to \mathcal{Y}$,
l'application d'espaces topologiques
$$
|\mathcal{Z} \times_\mathcal{Y} \mathcal{X}| \to |\mathcal{Z}|
$$
est fermée, c'est-à-dire si le changement de base
$\mathcal{Z} \times_\mathcal{Y} \mathcal{X} \to \mathcal{Z}$ est fermé.
\end{enumerate}
\end{definition}

\begin{lemma}
\label{stacks-morphisms-lemma-base-change-universally-closed}
Le changement de base d'un morphisme universellement fermé de champs
algébriques par un morphisme quelconque de champs algébriques est
universellement fermé.
\end{lemma}

\begin{proof}
C'est immédiat par définition.
\end{proof}

\begin{lemma}
\label{stacks-morphisms-lemma-composition-universally-closed}
Le composé de deux morphismes (universellement) fermés de champs
algébriques est (universellement) fermé.
\end{lemma}

\begin{proof}
Démonstration omise.
\end{proof}

\begin{lemma}
\label{stacks-morphisms-lemma-universally-closed-local}
Soit $f : \mathcal{X} \to \mathcal{Y}$ un morphisme de champs algébriques.
Les assertions suivantes sont équivalentes :
\begin{enumerate}
\item $f$ est universellement fermé ;
\item pour tout schéma $Z$ et tout morphisme $Z \to \mathcal{Y}$,
la projection $|Z \times_\mathcal{Y} \mathcal{X}| \to |Z|$
est fermée ;
\item pour tout schéma affine $Z$ et tout morphisme $Z \to \mathcal{Y}$,
la projection $|Z \times_\mathcal{Y} \mathcal{X}| \to |Z|$
est fermée ;
\item il existe un espace algébrique $V$ et un morphisme lisse surjectif
$V \to \mathcal{Y}$ tels que $V \times_\mathcal{Y} \mathcal{X} \to V$
soit un morphisme universellement fermé de champs algébriques.
\end{enumerate}
\end{lemma}

\begin{proof}
Nous omettons de démontrer que (1) entraîne (2) et que (2) entraîne (3).

\medskip\noindent
Supposons (3). Choisissons un morphisme lisse surjectif
$V \to \mathcal{Y}$. Nous allons montrer que
$V \times_\mathcal{Y} \mathcal{X} \to V$ est un morphisme
universellement fermé de champs algébriques. Soit
$\mathcal{Z} \to V$ un morphisme d'un champ algébrique vers $V$.
Soit $W \to \mathcal{Z}$ un morphisme lisse surjectif, où
$W = \coprod W_i$ est une réunion disjointe de schémas affines.
Nous avons alors le diagramme commutatif suivant :
$$
\xymatrix{
\coprod_i |W_i \times_\mathcal{Y} \mathcal{X}| \ar@{=}[r] \ar[d] &
|W \times_\mathcal{Y} \mathcal{X}| \ar[r] \ar[d] &
|\mathcal{Z} \times_\mathcal{Y} \mathcal{X}| \ar[d] \ar@{=}[r] &
|\mathcal{Z} \times_V (V \times_\mathcal{Y} \mathcal{X})| \ar[ld] \\
\coprod |W_i| \ar@{=}[r] &
|W| \ar[r] &
|\mathcal{Z}|
}
$$
Il faut montrer que la flèche sud-est est fermée. Les flèches
horizontales du milieu sont surjectives et ouvertes
(Propriétés des champs, Lemme
\ref{stacks-properties-lemma-topology-points}).
En vertu de (3) et du fait que les $W_i$ sont affines, les flèches
verticales de gauche sont fermées. Il s'ensuit que la flèche verticale
de droite est fermée.

\medskip\noindent
Supposons (4). Nous allons montrer que $f$ est universellement fermé.
Soit $\mathcal{Z} \to \mathcal{Y}$ un morphisme de champs algébriques.
Considérons le diagramme
$$
\xymatrix{
|(V \times_\mathcal{Y} \mathcal{Z})
\times_V (V \times_\mathcal{Y} \mathcal{X})| \ar@{=}[r] \ar[rd] &
|V \times_\mathcal{Y} \mathcal{X}| \ar[r] \ar[d] &
|\mathcal{Z} \times_\mathcal{Y} \mathcal{X}| \ar[d] \\
 &
|V \times_\mathcal{Y} \mathcal{Z}| \ar[r] &
|\mathcal{Z}|
}
$$
La flèche sud-ouest est fermée par hypothèse. Les flèches horizontales
sont surjectives et ouvertes, car les morphismes de champs algébriques
correspondants sont lisses et surjectifs (voir la référence ci-dessus).
Il s'ensuit que la flèche verticale de droite est fermée.
\end{proof}










\section{Morphismes universellement injectifs}
\label{stacks-morphisms-section-universally-injective}

\noindent
Soit $f$ un morphisme de champs algébriques qui est représentable par
des espaces algébriques. Dans Propriétés des champs, section
\ref{stacks-properties-section-properties-morphisms}
nous avons défini ce que signifie pour $f$ d'être universellement injectif
(ou radiciel, dans la terminologie classique des morphismes représentables).
En voici une autre caractérisation.

\begin{lemma}
\label{stacks-morphisms-lemma-characterize-representable-universally-injective}
Soit $f : \mathcal{X} \to \mathcal{Y}$ un morphisme de
champs algébriques qui est représentable par des espaces algébriques.
Les assertions suivantes sont équivalentes :
\begin{enumerate}
\item $f$ est universellement injectif (au sens de Propriétés des champs,
section \ref{stacks-properties-section-properties-morphisms}) ;
\item pour tout morphisme de champs algébriques
$\mathcal{Z} \to \mathcal{Y}$, l'application
$|\mathcal{Z} \times_\mathcal{Y} \mathcal{X}| \to |\mathcal{Z}|$
est injective.
\end{enumerate}
\end{lemma}

\begin{proof}
Supposons (1), et soit $\mathcal{Z} \to \mathcal{Y}$ comme en (2).
Choisissons un schéma $V$ et un morphisme lisse surjectif
$V \to \mathcal{Z}$. Par hypothèse, le morphisme d'espaces algébriques
$V \times_\mathcal{Y} \mathcal{X} \to V$ est universellement injectif ;
en particulier, l'application
$|V \times_\mathcal{Y} \mathcal{X}| \to |V|$ est injective. D'après
Propriétés des champs, section \ref{stacks-properties-section-points},
dans le diagramme commutatif
$$
\xymatrix{
|V \times_\mathcal{Y} \mathcal{X}| \ar[r] \ar[d] &
|\mathcal{Z} \times_\mathcal{Y} \mathcal{X}| \ar[d] \\
|V| \ar[r] & |\mathcal{Z}|
}
$$
les flèches horizontales sont ouvertes et surjectives ; de plus,
$$
|V \times_\mathcal{Y} \mathcal{X}| \longrightarrow
|V| \times_{|\mathcal{Z}|} |\mathcal{Z} \times_\mathcal{Y} \mathcal{X}|
$$
est surjective. Puisque la flèche verticale de gauche est injective,
il s'ensuit que celle de droite est injective. Cela démontre (2).
L'implication (2) $\Rightarrow$ (1) résulte des définitions.
\end{proof}

\noindent
Nous pouvons donc adopter la définition naturelle suivante.

\begin{definition}
\label{stacks-morphisms-definition-universally-injective}
Soit $f : \mathcal{X} \to \mathcal{Y}$ un morphisme de champs algébriques.
Nous disons que $f$ est {\it universellement injectif} si, pour tout
morphisme de champs algébriques $\mathcal{Z} \to \mathcal{Y}$,
l'application
$$
|\mathcal{Z} \times_\mathcal{Y} \mathcal{X}| \to |\mathcal{Z}|
$$
est injective.
\end{definition}

\begin{lemma}
\label{stacks-morphisms-lemma-base-change-universally-injective}
Le changement de base d'un morphisme universellement injectif de champs
algébriques par un morphisme quelconque de champs algébriques est
universellement injectif.
\end{lemma}

\begin{proof}
C'est immédiat par définition.
\end{proof}

\begin{lemma}
\label{stacks-morphisms-lemma-composition-universally-injective}
Le composé de deux morphismes universellement injectifs de champs
algébriques est universellement injectif.
\end{lemma}

\begin{proof}
Démonstration omise.
\end{proof}

\begin{lemma}
\label{stacks-morphisms-lemma-universally-injective}
Soit $f : \mathcal{X} \to \mathcal{Y}$ un morphisme de champs algébriques.
Les assertions suivantes sont équivalentes :
\begin{enumerate}
\item $f$ est universellement injectif ;
\item $\Delta : \mathcal{X} \to \mathcal{X} \times_\mathcal{Y} \mathcal{X}$
est surjective ;
\item pour tout corps algébriquement clos, tous
$x_1, x_2 : \Spec(k) \to \mathcal{X}$ et toute $2$-flèche
$\beta : f \circ x_1 \to f \circ x_2$, il existe une
$2$-flèche $\alpha : x_1 \to x_2$ telle que
$\beta = \text{id}_f \star \alpha$.
\end{enumerate}
\end{lemma}

\begin{proof}
(1) $\Rightarrow$ (2). Si $f$ est universellement injectif, la première
projection
$|\mathcal{X} \times_\mathcal{Y} \mathcal{X}| \to |\mathcal{X}|$
est injective, ce qui implique que $|\Delta|$ est surjective.

\medskip\noindent
(2) $\Rightarrow$ (1). Supposons $\Delta$ surjective. Tout changement
de base de $\Delta$ est alors surjectif (voir Propriétés des champs, section
\ref{stacks-properties-section-surjective}).
Comme la diagonale d'un changement de base de $f$ est un changement
de base de $\Delta$, il suffit de montrer que
$|\mathcal{X}| \to |\mathcal{Y}|$ est injective.
Dans le cas contraire, Propriétés des champs, Lemme
\ref{stacks-properties-lemma-points-cartesian}
montre que la première projection
$|\mathcal{X} \times_\mathcal{Y} \mathcal{X}| \to |\mathcal{X}|$
n'est pas injective. Cela signifie évidemment que $|\Delta|$ n'est
pas surjective.

\medskip\noindent
(3) $\Rightarrow$ (2). Soit
$t \in |\mathcal{X} \times_\mathcal{Y} \mathcal{X}|$.
Nous pouvons représenter $t$ par un morphisme
$t : \Spec(k) \to \mathcal{X} \times_\mathcal{Y} \mathcal{X}$
où $k$ est un corps algébriquement clos.
D'après notre construction des $2$-produits fibrés, nous pouvons
représenter $t$ par $(x_1, x_2, \beta)$, où
$x_1, x_2 : \Spec(k) \to \mathcal{X}$ et
$\beta : f \circ x_1 \to f \circ x_2$ est un $2$-morphisme.
L'assertion (3) implique alors qu'il existe un $2$-morphisme
$\alpha : x_1 \to x_2$ dont l'image est $\beta$.
Cela signifie exactement que $\Delta(x_1) = (x_1, x_1, \text{id})$
est isomorphe à $t$. L'assertion (2) est donc satisfaite.

\medskip\noindent
(2) $\Rightarrow$ (3). Soient
$x_1, x_2 : \Spec(k) \to \mathcal{X}$ des morphismes, où $k$ est un
corps algébriquement clos. Soit
$\beta : f \circ x_1 \to f \circ x_2$ un $2$-morphisme.
Comme au paragraphe précédent, nous obtenons un morphisme
$t = (x_1, x_2, \beta) :
\Spec(k) \to \mathcal{X} \times_\mathcal{Y} \mathcal{X}$.
D'après le Lemme \ref{stacks-morphisms-lemma-properties-diagonal},
$$
T = \mathcal{X}
\times_{\Delta, \mathcal{X} \times_\mathcal{Y} \mathcal{X}, t} \Spec(k)
$$
est un espace algébrique localement de type fini sur $\Spec(k)$.
La condition (2) implique que $T$ est non vide. Comme $k$ est
algébriquement clos, il existe alors un point à valeurs dans $k$
appartenant à $T$.
En explicitant les définitions, cela signifie qu'il existe un morphisme
$\alpha : x_1 \to x_2$ dans $\Mor(\Spec(k), \mathcal{X})$
tel que $\beta = \text{id}_f \star \alpha$.
\end{proof}

\begin{lemma}
\label{stacks-morphisms-lemma-universally-injective-point}
Soit $f : \mathcal{X} \to \mathcal{Y}$ un morphisme universellement
injectif de champs algébriques. Soit
$y : \Spec(k) \to \mathcal{Y}$ un morphisme, où $k$ est un corps
algébriquement clos. Si $y$ appartient à l'image de
$|\mathcal{X}| \to |\mathcal{Y}|$, il existe un morphisme
$x : \Spec(k) \to \mathcal{X}$ tel que $y = f \circ x$.
\end{lemma}

\begin{proof}
Remarquons d'abord que ce lemme n'est pas trivial : l'hypothèse selon
laquelle $y$ appartient à l'image de $|f|$ signifie seulement que l'on
peut relever $y$ en un morphisme vers $\mathcal{X}$ après avoir
éventuellement remplacé $k$ par un corps d'extension. Pour démontrer le
lemme, nous pouvons effectuer le changement de base de $f$ par $y$.
Nous pouvons donc supposer que $\mathcal{X}$ est un champ algébrique
non vide muni d'un morphisme universellement injectif
$\mathcal{X} \to \Spec(k)$, et nous voulons trouver, à valeurs dans $k$,
un point de $\mathcal{X}$. Nous pouvons remplacer $\mathcal{X}$ par sa réduction.
Choisissons un corps $k'$ et un morphisme plat surjectif et localement
de type fini $\Spec(k') \to \mathcal{X}$ ; voir Propriétés des champs,
Lemme \ref{stacks-properties-lemma-unique-point}.
Comme $\mathcal{X} \to \Spec(k)$ est universellement injectif,
$$
\Spec(k') \times_\mathcal{X} \Spec(k') \to \Spec(k' \otimes_k k')
$$
est surjectif, car c'est le changement de base du morphisme surjectif
$\Delta : \mathcal{X} \to \mathcal{X} \times_{\Spec(k)} \mathcal{X}$
(Lemme \ref{stacks-morphisms-lemma-universally-injective}).
Comme $k$ est algébriquement clos, $k' \otimes_k k'$ est un domaine
(Algèbre, Lemme
\ref{algebra-lemma-geometrically-integral-any-integral-base-change}).
Soit $\xi \in \Spec(k') \times_\mathcal{X} \Spec(k')$
un point dont l'image est le point générique de $\Spec(k' \otimes_k k')$.
Soit $U$ la composante connexe de
$\Spec(k') \times_\mathcal{X} \Spec(k')$ contenant $\xi$, munie de
sa structure réduite induite de sous-schéma fermé. Les deux projections
$U \to \Spec(k')$ sont alors localement de type fini, puisque c'était
le cas des projections
$\Spec(k') \times_\mathcal{X} \Spec(k') \to \Spec(k')$
qui sont des changements de base du morphisme
$\Spec(k') \to \mathcal{X}$. En appliquant Variétés, Proposition
\ref{varieties-proposition-unique-base-field}, nous voyons que les
clôtures intégrales des deux images de $k'$ dans
$\Gamma(U, \mathcal{O}_U)$ sont égales. Dans $\kappa(\xi)$, cela signifie
que tout élément de la forme $\lambda \otimes 1$ est algébriquement
dépendant du sous-corps
$$
1 \otimes k' \subset 
(\text{corps des fractions de }k' \otimes_k k') \subset
\kappa(\xi).
$$
Comme $k$ est algébriquement clos, cela n'est possible que si
$k' = k$, ce qui achève la démonstration.
\end{proof}

\begin{lemma}
\label{stacks-morphisms-lemma-universally-injective-local}
Soit $f : \mathcal{X} \to \mathcal{Y}$ un morphisme de champs algébriques.
Les assertions suivantes sont équivalentes :
\begin{enumerate}
\item $f$ est universellement injectif ;
\item pour tout schéma affine $Z$ et tout morphisme
$Z \to \mathcal{Y}$, le morphisme
$Z \times_\mathcal{Y} \mathcal{X} \to Z$
est universellement injectif ;
\item compléter ici.
\end{enumerate}
\end{lemma}

\begin{proof}
L'implication (1) $\Rightarrow$ (2) est immédiate.
Supposons (2). Nous allons montrer que
$\Delta_f : \mathcal{X} \to \mathcal{X} \times_\mathcal{Y} \mathcal{X}$
est surjective, ce qui implique (1) d'après le
Lemme \ref{stacks-morphisms-lemma-universally-injective}.
Considérons un schéma affine $V$ et un morphisme lisse
$V \to \mathcal{Y}$. Puisque
$g : V \times_\mathcal{Y} \mathcal{X} \to V$
est universellement injectif d'après (2), sa diagonale
$\Delta_g$ est surjective.
Or $\Delta_g$ est le changement de base de $\Delta_f$
par le morphisme lisse $V \to \mathcal{Y}$.
Comme la famille de ces morphismes $V \to \mathcal{Y}$ est
conjointement surjective, nous concluons que $\Delta_f$ est surjective.
\end{proof}

\begin{lemma}
\label{stacks-morphisms-lemma-check-universally-injective-covering}
Soit $f : \mathcal{X} \to \mathcal{Y}$ un morphisme de champs algébriques.
Soit $W \to \mathcal{Y}$ un morphisme plat surjectif et localement
de présentation finie, où $W$ est un espace algébrique. Si le changement
de base $W \times_\mathcal{Y} \mathcal{X} \to W$ est universellement
injectif, alors $f$ est universellement injectif.
\end{lemma}

\begin{proof}
Observons que la diagonale $\Delta_g$ du morphisme
$g : W \times_\mathcal{Y} \mathcal{X} \to W$
est le changement de base de $\Delta_f$ par $W \to \mathcal{Y}$.
Ainsi, si $\Delta_g$ est surjective, alors $\Delta_f$ l'est aussi
d'après Propriétés des champs, Lemme
\ref{stacks-properties-lemma-check-property-covering}.
Le lemme résulte donc de la caractérisation (2) du
Lemme \ref{stacks-morphisms-lemma-universally-injective}.
\end{proof}








\section{Homéomorphismes universels}
\label{stacks-morphisms-section-universal-homeomorphisms}

\noindent
Soit $f$ un morphisme de champs algébriques qui est représentable par
des espaces algébriques. Dans Propriétés des champs, section
\ref{stacks-properties-section-properties-morphisms}
nous avons défini ce que signifie pour $f$ d'être un homéomorphisme universel.
En voici une autre caractérisation.

\begin{lemma}
\label{stacks-morphisms-lemma-characterize-representable-universal-homeomorphism}
Soit $f : \mathcal{X} \to \mathcal{Y}$ un morphisme de
champs algébriques qui est représentable par des espaces algébriques.
Les assertions suivantes sont équivalentes :
\begin{enumerate}
\item $f$ est un homéomorphisme universel (Propriétés des champs,
section \ref{stacks-properties-section-properties-morphisms}) ;
\item pour tout morphisme de champs algébriques
$\mathcal{Z} \to \mathcal{Y}$, l'application d'espaces topologiques
$|\mathcal{Z} \times_\mathcal{Y} \mathcal{X}| \to |\mathcal{Z}|$ est
un homéomorphisme.
\end{enumerate}
\end{lemma}

\begin{proof}
Supposons (1), et soit $\mathcal{Z} \to \mathcal{Y}$ comme en (2).
Choisissons un schéma $V$ et un morphisme lisse surjectif
$V \to \mathcal{Z}$. Par hypothèse, le morphisme d'espaces algébriques
$V \times_\mathcal{Y} \mathcal{X} \to V$ est un homéomorphisme universel ;
en particulier, l'application
$|V \times_\mathcal{Y} \mathcal{X}| \to |V|$ est un homéomorphisme.
D'après Propriétés des champs, section
\ref{stacks-properties-section-points}, dans le diagramme commutatif
$$
\xymatrix{
|V \times_\mathcal{Y} \mathcal{X}| \ar[r] \ar[d] &
|\mathcal{Z} \times_\mathcal{Y} \mathcal{X}| \ar[d] \\
|V| \ar[r] & |\mathcal{Z}|
}
$$
les flèches horizontales sont ouvertes et surjectives ; de plus,
$$
|V \times_\mathcal{Y} \mathcal{X}| \longrightarrow
|V| \times_{|\mathcal{Z}|} |\mathcal{Z} \times_\mathcal{Y} \mathcal{X}|
$$
est surjective. Puisque la flèche verticale de gauche est un
homéomorphisme, il s'ensuit que celle de droite est un homéomorphisme.
Cela démontre (2). L'implication (2) $\Rightarrow$ (1) résulte
des définitions.
\end{proof}

\noindent
Nous pouvons donc adopter la définition naturelle suivante.

\begin{definition}
\label{stacks-morphisms-definition-universal-homeomorphism}
Soit $f : \mathcal{X} \to \mathcal{Y}$ un morphisme de champs algébriques.
Nous disons que $f$ est un {\it homéomorphisme universel} si, pour tout
morphisme de champs algébriques $\mathcal{Z} \to \mathcal{Y}$,
l'application d'espaces topologiques
$$
|\mathcal{Z} \times_\mathcal{Y} \mathcal{X}| \to |\mathcal{Z}|
$$
est un homéomorphisme.
\end{definition}

\begin{lemma}
\label{stacks-morphisms-lemma-base-change-universal-homeomorphism}
Le changement de base d'un homéomorphisme universel de champs algébriques
par un morphisme quelconque de champs algébriques est un homéomorphisme
universel.
\end{lemma}

\begin{proof}
C'est immédiat par définition.
\end{proof}

\begin{lemma}
\label{stacks-morphisms-lemma-composition-universal-homeomorphism}
Le composé de deux homéomorphismes universels de champs algébriques
est un homéomorphisme universel.
\end{lemma}

\begin{proof}
Démonstration omise.
\end{proof}

\begin{lemma}
\label{stacks-morphisms-lemma-check-universal-homeomorphism-covering}
Soit $f : \mathcal{X} \to \mathcal{Y}$ un morphisme de champs algébriques.
Soit $W \to \mathcal{Y}$ un morphisme plat surjectif et localement
de présentation finie, où $W$ est un espace algébrique. Si le changement
de base $W \times_\mathcal{Y} \mathcal{X} \to W$ est un homéomorphisme
universel, alors $f$ est un homéomorphisme universel.
\end{lemma}

\begin{proof}
Supposons que $g : W \times_\mathcal{Y} \mathcal{X} \to W$ soit un
homéomorphisme universel. Alors $g$ est universellement injectif ;
par conséquent, $f$ est universellement injectif d'après le
Lemme \ref{stacks-morphisms-lemma-check-universally-injective-covering}.
D'autre part, soit $\mathcal{Z} \to \mathcal{Y}$ un morphisme,
où $\mathcal{Z}$ est un champ algébrique. Choisissons un schéma $U$
et un morphisme lisse surjectif
$U \to W \times_\mathcal{Y} \mathcal{Z}$. Considérons le diagramme
$$
\xymatrix{
W \times_\mathcal{Y} \mathcal{X} \ar[d]^g &
U \times_\mathcal{Y} \mathcal{X} \ar[d] \ar[l] \ar[r] &
\mathcal{Z} \times_\mathcal{Y} \mathcal{X} \ar[d] \\
W &
U \ar[l] \ar[r] &
\mathcal{Z}
}
$$
Par l'hypothèse sur $g$, la flèche verticale médiane induit un
homéomorphisme sur les espaces topologiques. Les morphismes
$U \to \mathcal{Z}$ et
$U \times_\mathcal{Y} \mathcal{X} \to
\mathcal{Z} \times_\mathcal{Y} \mathcal{X}$
sont plats, surjectifs et localement de présentation finie ; ils
induisent donc des applications ouvertes sur les espaces topologiques.
Nous concluons que
$|\mathcal{Z} \times_\mathcal{Y} \mathcal{X}| \to |\mathcal{Z}|$
est ouverte. La surjectivité est facile à démontrer ; nous omettons
la démonstration.
\end{proof}










\section{Propriétés de morphismes locales pour la topologie lisse
sur la source et la cible}
\label{stacks-morphisms-section-local-source-target}

\noindent
Étant donnée une propriété des morphismes d'espaces algébriques qui est
{\it locale pour la topologie lisse sur la source et la cible}, voir
Descente sur les espaces,
Définition \ref{spaces-descent-definition-local-source-target},
nous pouvons nous en servir pour définir une propriété correspondante
des morphismes de champs algébriques, en imposant l'une ou l'autre des
conditions équivalentes du lemme suivant.

\begin{lemma}
\label{stacks-morphisms-lemma-local-source-target}
Soit $\mathcal{P}$ une propriété des morphismes d'espaces algébriques
qui est locale pour la topologie lisse sur la source et la cible.
Soit $f : \mathcal{X} \to \mathcal{Y}$ un morphisme de champs algébriques.
Considérons des diagrammes commutatifs
$$
\xymatrix{
U \ar[d]_a \ar[r]_h & V \ar[d]^b \\
\mathcal{X} \ar[r]^f & \mathcal{Y}
}
$$
où $U$ et $V$ sont des espaces algébriques et les flèches verticales
sont lisses. Les assertions suivantes sont équivalentes :
\begin{enumerate}
\item pour tout diagramme comme ci-dessus tel que, de plus,
$U \to \mathcal{X} \times_\mathcal{Y} V$ soit lisse,
le morphisme $h$ possède la propriété $\mathcal{P}$ ;
\item pour un certain diagramme comme ci-dessus où
$a : U \to \mathcal{X}$ est surjectif, le morphisme $h$
possède la propriété $\mathcal{P}$.
\end{enumerate}
Si $\mathcal{X}$ et $\mathcal{Y}$ sont représentables par des espaces
algébriques, ces assertions équivalent encore à ce que $f$, considéré
comme morphisme d'espaces algébriques, possède la propriété $\mathcal{P}$.
Si, de plus, $\mathcal{P}$ est stable par tout changement de base et
locale pour la topologie fppf sur la base, alors, pour les morphismes $f$
représentables par des espaces algébriques, elles équivalent encore à ce
que $f$ possède la propriété $\mathcal{P}$ au sens de Propriétés des champs,
section \ref{stacks-properties-section-properties-morphisms}.
\end{lemma}

\begin{proof}
Démontrons l'implication (1) $\Rightarrow$ (2). Choisissons un espace
algébrique $V$ et un morphisme lisse surjectif $V \to \mathcal{Y}$.
Choisissons un espace algébrique $U$ et un morphisme lisse surjectif
$U \to \mathcal{X} \times_\mathcal{Y} V$. Le morphisme
$U \to \mathcal{X}$ est lui aussi lisse et surjectif, car il est le
composé du changement de base
$\mathcal{X} \times_\mathcal{Y} V \to \mathcal{X}$ et du morphisme choisi
$U \to \mathcal{X} \times_\mathcal{Y} V$. Nous obtenons donc un diagramme
comme en (1). Si (1) est satisfaite, $h : U \to V$ possède la propriété
$\mathcal{P}$ ; puisque $U \to \mathcal{X}$ est surjectif, (2) est satisfaite.

\medskip\noindent
Réciproquement, supposons (2) satisfaite et soient $U, V, a, b, h$
comme en (2). Soit ensuite $U', V', a', b', h'$ un diagramme quelconque
comme en (1). Considérons
$$
\xymatrix{
U \ar[d] \ar[r]_h & V \ar[d] \\
\mathcal{X} \ar[r]^f & \mathcal{Y}
}
\quad\quad
\xymatrix{
U' \ar[d] \ar[r]_{h'} & V' \ar[d] \\
\mathcal{X} \ar[r]^f & \mathcal{Y}
}
$$
Pour montrer que (2) implique (1), il faut prouver que $h'$ possède
$\mathcal{P}$. Considérons à cette fin le diagramme commutatif
$$
\xymatrix{
U \ar[dd]^h &
U \times_\mathcal{X} U' \ar[d] \ar[l] \ar@/^6ex/[dd]^{(h, h')} \ar[r] &
U' \ar[dd]^{h'} \\
& U \times_\mathcal{Y} V' \ar[lu] \ar[d] & \\
V &
V \times_\mathcal{Y} V' \ar[l] \ar[r] &
V'
}
$$
d'espaces algébriques. Les flèches horizontales sont lisses, car ce sont
des changements de base des morphismes lisses
$V \to \mathcal{Y}$, $V' \to \mathcal{Y}$, $U \to \mathcal{X}$ et
$U' \to \mathcal{X}$. D'autre part,
$$
\xymatrix{
U \times_\mathcal{X} U' \ar[d] \ar[r] & U' \ar[d] \\
U \times_\mathcal{Y} V' \ar[r] & \mathcal{X} \times_\mathcal{Y} V'
}
$$
est cartésien ; la flèche verticale de gauche est donc lisse puisque
$U', V', a', b', h'$ sont comme en (1).
Comme $\mathcal{P}$ est locale pour la topologie lisse sur la cible
d'après Descente sur les espaces, Lemme
\ref{spaces-descent-lemma-local-source-target-implies}, partie (2),
le changement de base
$U \times_\mathcal{Y} V' \to V \times_\mathcal{Y} V'$
possède $\mathcal{P}$. Comme $\mathcal{P}$ est locale pour la topologie
lisse sur la source d'après Descente sur les espaces, Lemme
\ref{spaces-descent-lemma-local-source-target-implies}, partie (1),
nous pouvons précomposer par le morphisme lisse
$U \times_\mathcal{X} U' \to U \times_\mathcal{Y} V'$ et conclure que
$(h, h')$ possède $\mathcal{P}$.
Comme $V \times_\mathcal{Y} V' \to V'$ est lisse, Descente sur les espaces,
Lemme \ref{spaces-descent-lemma-local-source-target-implies}, partie (3),
montre que $U \times_\mathcal{X} U' \to V'$ possède $\mathcal{P}$.
Enfin, puisque $U \times_\mathcal{X} U' \to U'$ est lisse et surjectif
et que $\mathcal{P}$ est locale pour la topologie lisse sur la source
(même lemme), nous concluons que $h'$ possède $\mathcal{P}$. Cela achève
la démonstration de l'équivalence de (1) et (2).

\medskip\noindent
Si $\mathcal{X}$ et $\mathcal{Y}$ sont représentables, Descente sur les
espaces, Lemme \ref{spaces-descent-lemma-local-source-target-characterize},
s'applique et montre que (1) et (2) équivalent à ce que $f$ possède
$\mathcal{P}$.

\medskip\noindent
Enfin, supposons $f$ représentable, soient $U, V, a, b, h$ comme dans
la partie (2) du lemme, et supposons $\mathcal{P}$ stable par tout
changement de base. Il faut montrer que, pour tout schéma $Z$ et tout
morphisme $Z \to \mathcal{Y}$, le changement de base
$Z \times_\mathcal{Y} \mathcal{X} \to Z$
possède la propriété $\mathcal{P}$. Considérons le diagramme
$$
\xymatrix{
Z \times_\mathcal{Y} U \ar[d] \ar[r] &
Z \times_\mathcal{Y} V \ar[d] \\
Z \times_\mathcal{Y} \mathcal{X} \ar[r] &
Z
}
$$
La flèche horizontale supérieure est un changement de base de $h$ et
possède donc la propriété $\mathcal{P}$. La flèche verticale de gauche
est lisse et surjective, et celle de droite est lisse. Descente sur les
espaces, Lemme \ref{spaces-descent-lemma-local-source-target-characterize},
s'applique donc et montre que
$Z \times_\mathcal{Y} \mathcal{X} \to Z$ possède la propriété
$\mathcal{P}$.
\end{proof}

\begin{definition}
\label{stacks-morphisms-definition-P}
Soit $\mathcal{P}$ une propriété des morphismes d'espaces algébriques
qui est locale pour la topologie lisse sur la source et la cible.
Nous disons qu'un morphisme $f : \mathcal{X} \to \mathcal{Y}$ de champs
algébriques {\it possède la propriété $\mathcal{P}$} si les conditions
équivalentes du Lemme \ref{stacks-morphisms-lemma-local-source-target} sont satisfaites.
\end{definition}

\begin{remark}
\label{stacks-morphisms-remark-composition}
Soit $\mathcal{P}$ une propriété des morphismes d'espaces algébriques
qui est locale pour la topologie lisse sur la source et la cible et
stable par composition. La propriété des morphismes de champs algébriques
définie dans la Définition \ref{stacks-morphisms-definition-P} est alors stable par
composition. En effet, soient $f : \mathcal{X} \to \mathcal{Y}$ et
$g : \mathcal{Y} \to \mathcal{Z}$ des morphismes de champs algébriques
possédant la propriété $\mathcal{P}$. Choisissons un espace algébrique
$W$ et un morphisme lisse surjectif $W \to \mathcal{Z}$. Choisissons
un espace algébrique $V$ et un morphisme lisse surjectif
$V \to \mathcal{Y} \times_\mathcal{Z} W$. Enfin, choisissons un espace
algébrique $U$ et un morphisme lisse surjectif
$U \to \mathcal{X} \times_\mathcal{Y} V$. Par définition, les morphismes
$V \to W$ et $U \to V$ possèdent la propriété $\mathcal{P}$.
Ainsi, $U \to W$ possède la propriété $\mathcal{P}$, puisque nous avons
supposé $\mathcal{P}$ stable par composition. La définition montre alors
de nouveau que
$g \circ f : \mathcal{X} \to \mathcal{Z}$ possède la propriété
$\mathcal{P}$.
\end{remark}

\begin{remark}
\label{stacks-morphisms-remark-base-change}
Soit $\mathcal{P}$ une propriété des morphismes d'espaces algébriques
qui est locale pour la topologie lisse sur la source et la cible et
stable par changement de base. La propriété des morphismes de champs
algébriques définie dans la Définition \ref{stacks-morphisms-definition-P} est alors stable
par changement de base. En effet, soient
$f : \mathcal{X} \to \mathcal{Y}$ et
$g : \mathcal{Y}' \to \mathcal{Y}$ des morphismes de champs algébriques,
et supposons que $f$ possède la propriété $\mathcal{P}$. Choisissons
un espace algébrique $V$ et un morphisme lisse surjectif
$V \to \mathcal{Y}$. Choisissons un espace algébrique $U$ et un morphisme
lisse surjectif $U \to \mathcal{X} \times_\mathcal{Y} V$. Enfin,
choisissons un espace algébrique $V'$ et un morphisme lisse surjectif
$V' \to \mathcal{Y}' \times_\mathcal{Y} V$. Par définition, le morphisme
$U \to V$ possède la propriété $\mathcal{P}$. Ainsi,
$V' \times_V U \to V'$ possède la propriété $\mathcal{P}$, puisque nous
avons supposé $\mathcal{P}$ stable par changement de base. En considérant
le diagramme
$$
\xymatrix{
V' \times_V U \ar[r] \ar[d] &
\mathcal{Y}' \times_\mathcal{Y} \mathcal{X} \ar[r] \ar[d] &
\mathcal{X} \ar[d] \\
V' \ar[r] & \mathcal{Y}' \ar[r] & \mathcal{Y}
}
$$
nous voyons que la flèche horizontale supérieure gauche est lisse et
surjective ; la définition montre donc que la projection
$\mathcal{Y}' \times_\mathcal{Y} \mathcal{X} \to \mathcal{Y}'$
possède la propriété $\mathcal{P}$.
\end{remark}

\begin{remark}
\label{stacks-morphisms-remark-implication}
Soient $\mathcal{P}, \mathcal{P}'$ des propriétés des morphismes d'espaces
algébriques qui sont locales pour la topologie lisse sur la source et
la cible. Supposons que $\mathcal{P} \Rightarrow \mathcal{P}'$ pour les
morphismes d'espaces algébriques. Alors
$\mathcal{P} \Rightarrow \mathcal{P}'$ également pour les propriétés des
morphismes de champs algébriques définies dans la
Définition \ref{stacks-morphisms-definition-P} à l'aide de $\mathcal{P}$ et
$\mathcal{P}'$. Cela résulte immédiatement de la définition.
\end{remark}









\section{Morphismes de type fini}
\label{stacks-morphisms-section-finite-type}

\noindent
La propriété « localement de type fini » des morphismes d'espaces
algébriques est locale pour la topologie lisse sur la source et la cible ;
voir Descente sur les espaces, Remarque
\ref{spaces-descent-remark-list-local-source-target}.
Elle est également stable par changement de base et locale pour la
topologie fpqc sur la cible ; voir Morphismes d'espaces, Lemme
\ref{spaces-morphisms-lemma-base-change-finite-type}, et Descente sur les
espaces, Lemme
\ref{spaces-descent-lemma-descending-property-locally-finite-type}.
Par conséquent, d'après le Lemme \ref{stacks-morphisms-lemma-local-source-target}
ci-dessus, nous pouvons définir comme suit ce que signifie pour un morphisme
de champs algébriques d'être localement de type fini. Cette définition
coïncide, lorsque le morphisme est représentable par des espaces algébriques,
avec la notion déjà définie dans Propriétés des champs, section
\ref{stacks-properties-section-properties-morphisms}.

\begin{definition}
\label{stacks-morphisms-definition-locally-finite-type}
Soit $f : \mathcal{X} \to \mathcal{Y}$ un morphisme de champs algébriques.
\begin{enumerate}
\item Nous disons que $f$ est {\it localement de type fini} si les
conditions équivalentes du Lemme \ref{stacks-morphisms-lemma-local-source-target}
sont satisfaites avec
$\mathcal{P} = \text{localement de type fini}$.
\item Nous disons que $f$ est {\it de type fini} s'il est localement
de type fini et quasi-compact.
\end{enumerate}
\end{definition}

\begin{lemma}
\label{stacks-morphisms-lemma-composition-finite-type}
Le composé de morphismes de type fini est de type fini.
Il en va de même pour les morphismes localement de type fini.
\end{lemma}

\begin{proof}
Le résultat découle de la Remarque \ref{stacks-morphisms-remark-composition} et du
Lemme \ref{spaces-morphisms-lemma-composition-finite-type} de Morphismes d'espaces.
\end{proof}

\begin{lemma}
\label{stacks-morphisms-lemma-base-change-finite-type}
Un changement de base d'un morphisme de type fini est de type fini.
Il en va de même pour les morphismes localement de type fini.
\end{lemma}

\begin{proof}
Le résultat découle de la Remarque \ref{stacks-morphisms-remark-base-change} et du
Lemme \ref{spaces-morphisms-lemma-base-change-finite-type} de Morphismes d'espaces.
\end{proof}

\begin{lemma}
\label{stacks-morphisms-lemma-immersion-locally-finite-type}
Une immersion est localement de type fini.
\end{lemma}

\begin{proof}
Le résultat découle de la Remarque \ref{stacks-morphisms-remark-implication} et du
Lemme \ref{spaces-morphisms-lemma-immersion-locally-finite-type} de Morphismes d'espaces.
\end{proof}

\begin{lemma}
\label{stacks-morphisms-lemma-locally-finite-type-locally-noetherian}
Soit $f : \mathcal{X} \to \mathcal{Y}$ un morphisme de champs algébriques.
Si $f$ est localement de type fini et si $\mathcal{Y}$ est localement
noethérien, alors $\mathcal{X}$ est localement noethérien.
\end{lemma}

\begin{proof}
Soit
$$
\xymatrix{
U \ar[d] \ar[r] & V \ar[d] \\
\mathcal{X} \ar[r] & \mathcal{Y}
}
$$
un diagramme commutatif où $U$ et $V$ sont des schémas,
$V \to \mathcal{Y}$ est lisse et surjectif, et
$U \to V \times_\mathcal{Y} \mathcal{X}$ est lisse et surjectif.
Alors $U \to V$ est localement de type fini. Si $\mathcal{Y}$ est
localement noethérien, $V$ l'est aussi. D'après Morphismes, Lemme
\ref{morphisms-lemma-finite-type-noetherian}, $U$ est localement
noethérien, ce qui signifie que $\mathcal{X}$ est localement noethérien.
\end{proof}

\noindent
Les deux lemmes suivants seront améliorés plus loin (après l'étude des
morphismes de champs algébriques localement de présentation finie).

\begin{lemma}
\label{stacks-morphisms-lemma-check-finite-type-covering}
Soit $f : \mathcal{X} \to \mathcal{Y}$ un morphisme de champs algébriques.
Soit $W \to \mathcal{Y}$ un morphisme plat surjectif et localement de
présentation finie, où $W$ est un espace algébrique. Si le changement
de base $W \times_\mathcal{Y} \mathcal{X} \to W$ est localement de
type fini, alors $f$ est localement de type fini.
\end{lemma}

\begin{proof}
Choisissons un espace algébrique $V$ et un morphisme lisse surjectif
$V \to \mathcal{Y}$. Choisissons un espace algébrique $U$ et un morphisme
lisse surjectif $U \to V \times_\mathcal{Y} \mathcal{X}$.
Il faut montrer que $U \to V$ est localement de type fini.
Effectuons maintenant partout le changement de base par
$W \to \mathcal{Y}$ : posons
$U' = W \times_\mathcal{Y} U$,
$V' = W \times_\mathcal{Y} V$,
$\mathcal{X}' = W \times_\mathcal{Y} \mathcal{X}$,
et $\mathcal{Y}' = W \times_\mathcal{Y} \mathcal{Y} = W$.
Par changement de base,
$U' \to V' \times_{\mathcal{Y}'} \mathcal{X}'$ est encore lisse.
La définition des morphismes de champs algébriques localement de type
fini et l'hypothèse selon laquelle
$\mathcal{X}' \to \mathcal{Y}'$ est localement de type fini montrent donc
que $U' \to V'$ est localement de type fini. Puisque $V' \to V$ est
plat, surjectif et localement de présentation finie, comme changement
de base de $W \to \mathcal{Y}$, Descente sur les espaces, Lemme
\ref{spaces-descent-lemma-descending-property-locally-finite-type},
montre que $U \to V$ est localement de type fini, ce qui conclut.
\end{proof}

\begin{lemma}
\label{stacks-morphisms-lemma-check-finite-type-precompose}
Soient $\mathcal{X} \to \mathcal{Y} \to \mathcal{Z}$ des morphismes de
champs algébriques. Supposons $\mathcal{X} \to \mathcal{Z}$ localement
de type fini et $\mathcal{X} \to \mathcal{Y}$ représentable par des
espaces algébriques, plat, surjectif et localement de présentation finie.
Alors $\mathcal{Y} \to \mathcal{Z}$ est localement de type fini.
\end{lemma}

\begin{proof}
Choisissons un espace algébrique $W$ et un morphisme lisse surjectif
$W \to \mathcal{Z}$. Choisissons un espace algébrique $V$ et un morphisme
lisse surjectif $V \to W \times_\mathcal{Z} \mathcal{Y}$. Posons
$U = V \times_\mathcal{Y} \mathcal{X}$, qui est un espace algébrique.
Le morphisme $U \to V$ est plat, surjectif et localement de présentation
finie, tandis que $U \to W$ est localement de type fini.
Le lemme se ramène donc au cas des morphismes d'espaces algébriques.
Ce cas est Descente sur les espaces, Lemme
\ref{spaces-descent-lemma-locally-finite-type-fppf-local-source}.
\end{proof}

\begin{lemma}
\label{stacks-morphisms-lemma-finite-type-permanence}
Soient $f : \mathcal{X} \to \mathcal{Y}$ et
$g : \mathcal{Y} \to \mathcal{Z}$ des morphismes de champs algébriques.
Si $g \circ f : \mathcal{X} \to \mathcal{Z}$ est localement de type fini,
alors $f : \mathcal{X} \to \mathcal{Y}$ est localement de type fini.
\end{lemma}

\begin{proof}
On peut trouver un diagramme
$$
\xymatrix{
U \ar[r] \ar[d] & V \ar[r] \ar[d] & W \ar[d] \\
\mathcal{X} \ar[r] & \mathcal{Y} \ar[r] & \mathcal{Z}
}
$$
où $U$, $V$ et $W$ sont des schémas, la flèche verticale
$W \to \mathcal{Z}$ est lisse et surjective, la flèche
$V \to \mathcal{Y} \times_\mathcal{Z} W$ est lisse et surjective,
et la flèche $U \to \mathcal{X} \times_\mathcal{Y} V$ est lisse et
surjective. Alors $U \to \mathcal{X} \times_\mathcal{Z} W$ est lui aussi
lisse et surjectif (comme composé d'un morphisme lisse surjectif et du
changement de base d'un tel morphisme). Par définition, $U \to W$ est
localement de type fini. D'après Morphismes, Lemme
\ref{morphisms-lemma-permanence-finite-type}, $U \to V$ est donc
localement de type fini, ce qui signifie à son tour, par définition,
que $\mathcal{X} \to \mathcal{Y}$ est localement de type fini.
\end{proof}










\section{Points de type fini}
\label{stacks-morphisms-section-points-finite-type}

\noindent
Soit $\mathcal{X}$ un champ algébrique.
Un point de type fini $x \in |\mathcal{X}|$ est un point qui peut être représenté
par un morphisme $\Spec(k) \to \mathcal{X}$ localement de type fini.
Les points de type fini remplacent avantageusement les points fermés pour les
espaces algébriques et les champs algébriques. Par exemple, il y en a toujours
``assez''.

\begin{lemma}
\label{stacks-morphisms-lemma-point-finite-type}
Soit $\mathcal{X}$ un champ algébrique.
Soit $x \in |\mathcal{X}|$. Les conditions suivantes sont équivalentes :
\begin{enumerate}
\item Il existe un morphisme $\Spec(k) \to \mathcal{X}$
localement de type fini qui représente $x$.
\item Il existe un schéma $U$, un point fermé $u \in U$ et un morphisme
lisse $\varphi : U \to \mathcal{X}$ tels que $\varphi(u) = x$.
\end{enumerate}
\end{lemma}

\begin{proof}
Soient $u \in U$ et $U \to \mathcal{X}$ comme en (2). Alors
$\Spec(\kappa(u)) \to U$ est de type fini, et $U \to \mathcal{X}$ est
représentable et localement de type fini (par Morphismes d'espaces,
Lemmes \ref{spaces-morphisms-lemma-etale-locally-finite-presentation} et
\ref{spaces-morphisms-lemma-finite-presentation-finite-type}).
La condition (1) résulte donc du
Lemme \ref{stacks-morphisms-lemma-composition-finite-type}.

\medskip\noindent
Réciproquement, supposons que $\Spec(k) \to \mathcal{X}$ soit localement de
type fini et représente $x$. Soit $U \to \mathcal{X}$ un morphisme lisse
surjectif, où $U$ est un schéma. Par hypothèse,
$U \times_\mathcal{X} \Spec(k) \to U$ est un morphisme d'espaces algébriques
localement de type fini. Choisissons un point de type fini $v$ de
$U \times_\mathcal{X} \Spec(k)$ (il en existe au moins un, voir Morphismes
d'espaces, Lemme \ref{spaces-morphisms-lemma-identify-finite-type-points}).
Par Morphismes d'espaces,
Lemme \ref{spaces-morphisms-lemma-finite-type-points-morphism},
l'image $u \in U$ de $v$ est un point de type fini de $U$.
Par Morphismes, Lemme \ref{morphisms-lemma-identify-finite-type-points},
quitte à restreindre $U$, nous pouvons donc supposer que $u$ est un point
fermé de $U$ ; autrement dit, (2) est satisfaite.
\end{proof}

\begin{definition}
\label{stacks-morphisms-definition-finite-type-point}
Soit $\mathcal{X}$ un champ algébrique. On dit qu'un point
$x \in |\mathcal{X}|$ est un {\it point de type fini}\footnote{Il s'agit
d'un léger abus de langage, puisqu'il serait peut-être plus exact de dire
``point localement de type fini''.} s'il satisfait aux conditions équivalentes
du Lemme \ref{stacks-morphisms-lemma-point-finite-type}. On note
$\mathcal{X}_{\text{ft-pts}}$ l'ensemble des points de type fini de
$\mathcal{X}$.
\end{definition}

\noindent
On peut décrire comme suit l'ensemble des points de type fini.

\begin{lemma}
\label{stacks-morphisms-lemma-identify-finite-type-points}
Soit $\mathcal{X}$ un champ algébrique. On a
$$
\mathcal{X}_{\text{ft-pts}} =
\bigcup\nolimits_{\varphi : U \to \mathcal{X}\text{ lisse}} |\varphi|(U_0)
$$
où $U_0$ est l'ensemble des points fermés de $U$.
Ici, on peut faire parcourir à $U$ soit tous les schémas lisses sur
$\mathcal{X}$, soit tous les schémas affines lisses sur $\mathcal{X}$.
\end{lemma}

\begin{proof}
C'est une conséquence immédiate du
Lemme \ref{stacks-morphisms-lemma-point-finite-type}.
\end{proof}

\begin{lemma}
\label{stacks-morphisms-lemma-finite-type-points-morphism}
Soit $f : \mathcal{X} \to \mathcal{Y}$ un morphisme de champs algébriques.
Si $f$ est localement de type fini, alors
$f(\mathcal{X}_{\text{ft-pts}}) \subset \mathcal{Y}_{\text{ft-pts}}$.
\end{lemma}

\begin{proof}
Soit $x \in \mathcal{X}_{\text{ft-pts}}$. Représentons $x$ par un morphisme
localement de type fini $x : \Spec(k) \to \mathcal{X}$. Alors
$f \circ x$ est localement de type fini par le
Lemme \ref{stacks-morphisms-lemma-composition-finite-type}.
Par conséquent, $f(x) \in \mathcal{Y}_{\text{ft-pts}}$.
\end{proof}

\begin{lemma}
\label{stacks-morphisms-lemma-finite-type-points-surjective-morphism}
Soit $f : \mathcal{X} \to \mathcal{Y}$ un morphisme de champs algébriques.
Si $f$ est localement de type fini et surjectif, alors
$f(\mathcal{X}_{\text{ft-pts}}) = \mathcal{Y}_{\text{ft-pts}}$.
\end{lemma}

\begin{proof}
On a $f(\mathcal{X}_{\text{ft-pts}}) \subset
\mathcal{Y}_{\text{ft-pts}}$ par le
Lemme \ref{stacks-morphisms-lemma-finite-type-points-morphism}.
Soit $y \in |\mathcal{Y}|$ un point de type fini. Représentons $y$ par un
morphisme $\Spec(k) \to \mathcal{Y}$ localement de type fini.
Puisque $f$ est surjectif, le champ algébrique
$\mathcal{X}_k = \Spec(k) \times_\mathcal{Y} \mathcal{X}$ est non vide ; il
possède donc un point de type fini $x \in |\mathcal{X}_k|$ par le
Lemme \ref{stacks-morphisms-lemma-identify-finite-type-points}.
Or $\mathcal{X}_k \to \mathcal{X}$ est localement de type fini, car il se
déduit de $\Spec(k) \to \mathcal{Y}$ par changement de base
(Lemme \ref{stacks-morphisms-lemma-base-change-finite-type}).
L'image de $x$ dans $\mathcal{X}$ est donc un point de type fini par le
Lemme \ref{stacks-morphisms-lemma-finite-type-points-morphism}, et elle s'envoie sur $y$ par
construction.
\end{proof}

\begin{lemma}
\label{stacks-morphisms-lemma-enough-finite-type-points}
Soit $\mathcal{X}$ un champ algébrique.
Pour toute partie localement fermée $T \subset |\mathcal{X}|$, on a
$$
T \not = \emptyset
\Rightarrow
T \cap \mathcal{X}_{\text{ft-pts}} \not = \emptyset.
$$
En particulier, pour toute partie fermée $T \subset |\mathcal{X}|$,
$T \cap \mathcal{X}_{\text{ft-pts}}$ est dense dans $T$.
\end{lemma}

\begin{proof}
Soit $i : \mathcal{Z} \to \mathcal{X}$ la structure réduite de sous-champ
induite sur $T$, voir Propriétés des champs,
Remarque \ref{stacks-properties-remark-stack-structure-locally-closed-subset}.
Une immersion est localement de type fini, voir le
Lemme \ref{stacks-morphisms-lemma-immersion-locally-finite-type}.
Par le Lemme \ref{stacks-morphisms-lemma-finite-type-points-morphism}, on a donc
$\mathcal{Z}_{\text{ft-pts}} \subset \mathcal{X}_{\text{ft-pts}} \cap T$.
Enfin, tout schéma affine non vide $U$ muni d'un morphisme lisse vers
$\mathcal{Z}$ possède au moins un point fermé ; ainsi $\mathcal{Z}$ possède
au moins un point de type fini par le
Lemme \ref{stacks-morphisms-lemma-identify-finite-type-points}.
Le lemme en résulte.
\end{proof}

\noindent
Voici une autre caractérisation, plus technique, d'un point de type fini d'un
champ algébrique. Elle montre notamment que la gerbe résiduelle de
$\mathcal{X}$ en $x$ existe dès que $x$ est un point de type fini !

\begin{lemma}
\label{stacks-morphisms-lemma-point-finite-type-monomorphism}
Soit $\mathcal{X}$ un champ algébrique.
Soit $x \in |\mathcal{X}|$. Les conditions suivantes sont équivalentes :
\begin{enumerate}
\item $x$ est un point de type fini,
\item il existe un champ algébrique $\mathcal{Z}$ dont l'espace topologique
sous-jacent $|\mathcal{Z}|$ est réduit à un point, ainsi qu'un morphisme
localement de type fini $f : \mathcal{Z} \to \mathcal{X}$ tel que
$\{x\} = |f|(|\mathcal{Z}|)$, et
\item la gerbe résiduelle $\mathcal{Z}_x$ de $\mathcal{X}$ en $x$ existe et
le morphisme d'inclusion $\mathcal{Z}_x \to \mathcal{X}$ est localement de
type fini.
\end{enumerate}
\end{lemma}

\begin{proof}
(Tous les morphismes qui interviennent dans ce paragraphe sont représentables
par des espaces algébriques ; les conventions et les résultats de
Propriétés des champs,
section \ref{stacks-properties-section-properties-morphisms},
sont donc applicables.)
Supposons que $x$ soit un point de type fini. Choisissons un schéma affine
$U$, un point fermé $u \in U$ et un morphisme lisse
$\varphi : U \to \mathcal{X}$ tels que $\varphi(u) = x$, voir le
Lemme \ref{stacks-morphisms-lemma-identify-finite-type-points}.
Comme d'habitude, posons $u = \Spec(\kappa(u))$. Posons
$R = u \times_\mathcal{X} u$ ; on obtient ainsi un groupoïde en espaces
algébriques $(u, R, s, t, c)$, voir Champs algébriques,
Lemme \ref{algebraic-lemma-map-space-into-stack}.
Chacun des morphismes de projection $R \to u$ est le composé
$$
R = u \times_\mathcal{X} u \to
u \times_\mathcal{X} U \to
u \times_\mathcal{X} \mathcal{X} = u
$$
où la première flèche est de type fini (c'est un changement de base de
l'immersion fermée de schémas $u \to U$) et la seconde est lisse (c'est un
changement de base du morphisme lisse $U \to \mathcal{X}$). Par conséquent,
$s, t : R \to u$ sont localement de type fini (comme composés, voir
Morphismes d'espaces,
Lemme \ref{spaces-morphisms-lemma-composition-finite-type}).
Comme $u$ est le spectre d'un corps, les morphismes $s, t$ sont plats et localement de
présentation finie (par Morphismes d'espaces, Lemme
\ref{spaces-morphisms-lemma-noetherian-finite-type-finite-presentation}).
Le champ $\mathcal{Z} = [u/R]$ est donc algébrique par Critères de
représentabilité,
Théorème \ref{criteria-theorem-flat-groupoid-gives-algebraic-stack}.
Par Champs algébriques,
Lemme \ref{algebraic-lemma-map-space-into-stack},
on obtient un morphisme canonique
$$
f : \mathcal{Z} \longrightarrow \mathcal{X}
$$
qui est pleinement fidèle. Ce morphisme est donc représentable par des espaces
algébriques, voir Champs algébriques, Lemme
\ref{algebraic-lemma-characterize-representable-by-algebraic-spaces},
et c'est un monomorphisme, voir Propriétés des champs,
Lemme \ref{stacks-properties-lemma-monomorphism}.
Il s'ensuit que la gerbe résiduelle
$\mathcal{Z}_x \subset \mathcal{X}$ de $\mathcal{X}$ en $x$ existe et que
$f$ induit une équivalence $\mathcal{Z} \to \mathcal{Z}_x$, voir
Propriétés des champs,
Lemme \ref{stacks-properties-lemma-residual-gerbe-unique}.
Par construction, le diagramme
$$
\xymatrix{
u \ar[d] \ar[r] & U \ar[d] \\
\mathcal{Z} \ar[r]^f & \mathcal{X}
}
$$
est commutatif. Par Critères de représentabilité,
Lemme \ref{criteria-lemma-flat-quotient-flat-presentation},
la flèche verticale de gauche est surjective, plate et localement de
présentation finie. Considérons
$$
\xymatrix{
u \times_\mathcal{X} U \ar[d] \ar[r] &
\mathcal{Z} \times_\mathcal{X} U \ar[r] \ar[d] & U \ar[d] \\
u \ar[r] & \mathcal{Z} \ar[r]^f & \mathcal{X}
}
$$
Comme $u \to \mathcal{X}$ est localement de type fini, son changement de base
$u \times_\mathcal{X} U \to U$ est localement de type fini. De plus,
$u \times_\mathcal{X} U \to \mathcal{Z} \times_\mathcal{X} U$ est
surjectif, plat et localement de présentation finie, car il se déduit de
$u \to \mathcal{Z}$ par changement de base. Ainsi,
$\{u \times_\mathcal{X} U \to \mathcal{Z} \times_\mathcal{X} U\}$
est un recouvrement fppf d'espaces algébriques ; on en conclut que
$\mathcal{Z} \times_\mathcal{X} U \to U$ est localement de type fini par
Descente sur les espaces, Lemme
\ref{spaces-descent-lemma-locally-finite-presentation-fppf-local-source}.
Par définition, cela signifie que $f$ est localement de type fini (en effet,
la flèche verticale $\mathcal{Z} \times_\mathcal{X} U \to \mathcal{Z}$ est
lisse, comme changement de base de $U \to \mathcal{X}$, et surjective puisque
$\mathcal{Z}$ n'a qu'un point). Comme
$\mathcal{Z} = \mathcal{Z}_x$, la condition (3) est satisfaite.

\medskip\noindent
Il est clair que (3) implique (2).
Si (2) est satisfaite, alors $x$ est un point de type fini de $\mathcal{X}$
par le Lemme \ref{stacks-morphisms-lemma-finite-type-points-morphism} et le
Lemme \ref{stacks-morphisms-lemma-enough-finite-type-points}, qui montre que
$\mathcal{Z}_{\text{ft-pts}}$ est non vide, autrement dit que l'unique point
de $\mathcal{Z}$ est un point de type fini de $\mathcal{Z}$.
\end{proof}







\section{Groupes d'automorphismes}
\label{stacks-morphisms-section-automorphism-groups}

\noindent
Soit $\mathcal{X}$ un champ algébrique. Soit $x \in |\mathcal{X}|$
correspondant à $x : \Spec(k) \to \mathcal{X}$. Dans cette situation, nous
dirons souvent : ``soit $G_x/k$ l'espace algébrique en groupes des
automorphismes de $x$''. Cela signifie simplement que
$$
G_x = \mathit{Isom}_\mathcal{X}(x, x) =
\Spec(k) \times_\mathcal{X} \mathcal{I}_\mathcal{X}
$$
est l'espace algébrique en groupes des automorphismes de $x$. C'est un espace
algébrique en groupes sur $\Spec(k)$. Si $k'/k$ est une extension de corps,
l'espace algébrique en groupes des automorphismes du morphisme induit
$x' : \Spec(k') \to \mathcal{X}$ se déduit de $G_x$ par changement de base
à $\Spec(k')$.

\begin{lemma}
\label{stacks-morphisms-lemma-automorphism-group-scheme}
Dans la situation ci-dessus, $G_x$ est un schéma si l'une des conditions
suivantes est satisfaite :
\begin{enumerate}
\item $\Delta : \mathcal{X} \to \mathcal{X} \times \mathcal{X}$
est quasi-séparée,
\item $\Delta : \mathcal{X} \to \mathcal{X} \times \mathcal{X}$
est localement séparée,
\item $\mathcal{X}$ est quasi-DM,
\item $\mathcal{I}_\mathcal{X} \to \mathcal{X}$
est quasi-séparé,
\item $\mathcal{I}_\mathcal{X} \to \mathcal{X}$
est localement séparé, ou
\item $\mathcal{I}_\mathcal{X} \to \mathcal{X}$
est localement quasi-fini.
\end{enumerate}
\end{lemma}

\begin{proof}
Le Lemme \ref{stacks-morphisms-lemma-diagonal-diagonal} montre que
(1) $\Rightarrow$ (4), (2) $\Rightarrow$ (5) et (3) $\Rightarrow$ (6).
Dans le cas (4), $G_x$ est un espace algébrique quasi-séparé ; dans le cas
(5), $G_x$ est un espace algébrique localement séparé.
Dans les deux cas, $G_x$ est un espace algébrique décent
(Espaces décents,
section \ref{decent-spaces-section-reasonable-decent} et
Lemme \ref{decent-spaces-lemma-locally-separated-decent}).
Ensuite, $G_x$ est séparé par Compléments sur les groupoïdes en espaces,
Lemme \ref{spaces-more-groupoids-lemma-group-scheme-over-field-separated} ;
on en conclut que $G_x$ est un schéma par Compléments sur les groupoïdes en
espaces, Proposition
\ref{spaces-more-groupoids-proposition-group-space-scheme-over-field}.
Dans le cas (6), $G_x \to \Spec(k)$ est localement quasi-fini ; $G_x$ est donc
un schéma par Espaces sur des corps,
Lemme \ref{spaces-over-fields-lemma-locally-quasi-finite-over-field}.
\end{proof}

\begin{lemma}
\label{stacks-morphisms-lemma-property-automorphism-groups}
Soit $\mathcal{X}$ un champ algébrique et soit $x \in |\mathcal{X}|$ un point.
Soit $P$ une propriété des espaces algébriques sur des corps qui soit
invariante par extension du corps de base ; par exemple,
$P(X/k) = X \to \Spec(k)\text{ est fini}$.
Les conditions suivantes sont équivalentes :
\begin{enumerate}
\item pour un certain morphisme $x : \Spec(k) \to \mathcal{X}$ représentant
la classe de $x$, l'espace algébrique en groupes d'automorphismes $G_x/k$
vérifie $P$, et
\item pour tout morphisme $x : \Spec(k) \to \mathcal{X}$ représentant
la classe de $x$, l'espace algébrique en groupes d'automorphismes $G_x/k$
vérifie $P$.
\end{enumerate}
\end{lemma}

\begin{proof}
Démonstration omise.
\end{proof}

\begin{remark}
\label{stacks-morphisms-remark-property-automorphism-groups}
Soit $P$ une propriété des espaces algébriques sur des corps qui soit
invariante par extension du corps de base. Étant donnés un champ algébrique
$\mathcal{X}$ et $x \in |\mathcal{X}|$, on dit que le groupe d'automorphismes
de $\mathcal{X}$ en $x$ vérifie $P$ si les conditions équivalentes du
Lemme \ref{stacks-morphisms-lemma-property-automorphism-groups} sont satisfaites.
Par exemple, on dit que {\it le groupe d'automorphismes de $\mathcal{X}$
en $x$ est fini} si $G_x \to \Spec(k)$ est fini pour tout représentant
$x : \Spec(k) \to \mathcal{X}$ de $x$.
On procède de même pour les propriétés lisse, propre, etc.
(Il s'agit manifestement d'un abus de langage, mais nous pensons qu'il
n'entraînera ni confusion ni imprécision.)
\end{remark}

\begin{lemma}
\label{stacks-morphisms-lemma-iso-automorphism-groups}
Soit $f : \mathcal{X} \to \mathcal{Y}$ un morphisme de champs algébriques.
Soit $x \in |\mathcal{X}|$ un point. Les conditions suivantes sont équivalentes :
\begin{enumerate}
\item pour un certain morphisme $x : \Spec(k) \to \mathcal{X}$ représentant
la classe de $x$, si l'on pose $y = f \circ x$, l'application
$G_x \to G_y$ entre les espaces algébriques en groupes d'automorphismes est
un isomorphisme, et
\item pour tout morphisme $x : \Spec(k) \to \mathcal{X}$ représentant
la classe de $x$, si l'on pose $y = f \circ x$, l'application
$G_x \to G_y$ entre les espaces algébriques en groupes d'automorphismes est
un isomorphisme.
\end{enumerate}
\end{lemma}

\begin{proof}
Cela résulte du fait qu'être un isomorphisme est une propriété locale pour la
topologie fpqc sur la cible, voir Descente sur les espaces, Lemme
\ref{spaces-descent-lemma-descending-property-isomorphism}.
En effet, supposons que $k'/k$ soit une extension de corps, notons
$x' : \Spec(k') \to \mathcal{X}$ le composé et posons $y' = f \circ x'$.
Alors le morphisme $G_{x'} \to G_{y'}$ se déduit de $G_x \to G_y$ par le
changement de base $\Spec(k') \to \Spec(k)$.
Par conséquent, $G_x \to G_y$ est un isomorphisme si et seulement si
$G_{x'} \to G_{y'}$ en est un. La propriété vaut donc pour toute la classe
d'équivalence dès qu'elle vaut pour un représentant.
\end{proof}

\begin{remark}
\label{stacks-morphisms-remark-identify-automorphism-groups}
Soit $f : \mathcal{X} \to \mathcal{Y}$ un morphisme de champs algébriques,
et soit $x \in |\mathcal{X}|$ un point. Pour exprimer que $f$ et $x$
satisfont aux conditions équivalentes du
Lemme \ref{stacks-morphisms-lemma-iso-automorphism-groups}, la littérature dit que
{\it $f$ préserve les stabilisateurs en $x$} ou que
{\it $f$ reflète les points fixes en $x$}.
Nous préférons dire que {\it $f$ induit un isomorphisme entre les groupes
d'automorphismes en $x$ et en $f(x)$}.
\end{remark}





\section{Présentations et propriétés des champs algébriques}
\label{stacks-morphisms-section-presentations-properties}

\noindent
Soit $(U, R, s, t, c)$ un groupoïde en espaces algébriques.
Si $s, t : R \to U$ sont plats et localement de présentation finie,
alors le champ quotient $[U/R]$ est un champ algébrique, voir
Critères de représentabilité, Théorème
\ref{criteria-theorem-flat-groupoid-gives-algebraic-stack}.
Dans cette section, nous étudions les conséquences des propriétés de
$(U, R, s, t, c)$ pour le champ algébrique $[U/R]$.

\begin{lemma}
\label{stacks-morphisms-lemma-properties-diagonal-from-presentation}
Soit $(U, R, s, t, c)$ un groupoïde en espaces algébriques tel que
$s, t : R \to U$ soient plats et localement de présentation finie.
Considérons le champ algébrique $\mathcal{X} = [U/R]$ (voir ci-dessus).
\begin{enumerate}
\item Si $R \to U \times U$ est séparé, alors
$\Delta_\mathcal{X}$ est séparée.
\item Si $U$ et $R$ sont séparés, alors $\Delta_\mathcal{X}$ est séparée.
\item Si $R \to U \times U$ est localement quasi-fini, alors $\mathcal{X}$
est quasi-DM.
\item Si $s, t : R \to U$ sont localement quasi-finis, alors
$\mathcal{X}$ est quasi-DM.
\item Si $R \to U \times U$ est propre, alors $\mathcal{X}$ est séparé.
\item Si $s, t : R \to U$ sont propres et si $U$ est séparé, alors
$\mathcal{X}$ est séparé.
\item Ajouter d'autres cas ici.
\end{enumerate}
\end{lemma}

\begin{proof}
Remarquons que le morphisme $U \to \mathcal{X}$ est surjectif, plat et
localement de présentation finie par Critères de représentabilité, Lemme
\ref{criteria-lemma-flat-quotient-flat-presentation}.
Il en est donc de même de
$U \times U \to \mathcal{X} \times \mathcal{X}$. On a le diagramme cartésien
$$
\xymatrix{
R  = U \times_\mathcal{X} U \ar[r] \ar[d] & U \times U \ar[d] \\
\mathcal{X} \ar[r] & \mathcal{X} \times \mathcal{X}
}
$$
(voir Groupoïdes en espaces, Lemme
\ref{spaces-groupoids-lemma-quotient-stack-2-cartesian}).
Ainsi, $\Delta_\mathcal{X}$ possède l'une des propriétés énumérées dans
Propriétés des champs, section
\ref{stacks-properties-section-properties-morphisms}
si et seulement si le morphisme $R \to U \times U$ la possède, voir
Propriétés des champs, Lemme
\ref{stacks-properties-lemma-check-property-covering}.
Cela démontre (1), (3) et (5).
La condition de (2) implique que $R \to U \times U$ est séparé ; (2) résulte
donc de (1). La condition de (4) implique celle de (3) ; (4) résulte donc de
(3). La condition de (6) implique celle de (5) par Morphismes d'espaces, Lemme
\ref{spaces-morphisms-lemma-universally-closed-permanence} ; (6) résulte donc
de (5).
\end{proof}

\begin{lemma}
\label{stacks-morphisms-lemma-points-presentation}
Soit $(U, R, s, t, c)$ un groupoïde en espaces algébriques tel que
$s, t : R \to U$ soient plats et localement de présentation finie.
Considérons le champ algébrique $\mathcal{X} = [U/R]$ (voir ci-dessus).
Alors l'image de $|R| \to |U| \times |U|$ est une relation d'équivalence et
$|\mathcal{X}|$ est le quotient de $|U|$ par cette relation d'équivalence.
\end{lemma}

\begin{proof}
Le morphisme induit $p : U \to \mathcal{X}$ est surjectif, plat et localement
de présentation finie, voir Critères de représentabilité, Lemme
\ref{criteria-lemma-flat-quotient-flat-presentation}.
Ainsi, $|U| \to |\mathcal{X}|$ est surjectif par Propriétés des champs, Lemme
\ref{stacks-properties-lemma-characterize-surjective}.
Notons que $R = U \times_\mathcal{X} U$, voir Groupoïdes en espaces,
Lemme \ref{spaces-groupoids-lemma-quotient-stack-2-cartesian}.
Par conséquent, Propriétés des champs,
Lemme \ref{stacks-properties-lemma-points-cartesian}, montre que l'application
$$
|R| \longrightarrow |U| \times_{|\mathcal{X}|} |U|
$$
est surjective. L'image de $|R| \to |U| \times |U|$ est donc exactement
l'ensemble des couples $(u_1, u_2) \in |U| \times |U|$ tels que $u_1$ et
$u_2$ aient la même image dans $|\mathcal{X}|$.
En combinant ces deux assertions, on obtient le résultat du lemme.
\end{proof}






\section{Présentations particulières des champs algébriques}
\label{stacks-morphisms-section-presentations}

\noindent
Dans cette section, nous démontrons deux théorèmes importants.
Le premier caractérise les champs quasi-DM $\mathcal{X}$ comme les champs de
la forme $\mathcal{X} = [U/R]$ pour lesquels $s, t : R \to U$ sont localement
quasi-finis (ainsi que plats et localement de présentation finie).
Le second affirme que les champs algébriques DM sont de Deligne--Mumford.

\medskip\noindent
Le lemme suivant donne un critère pour qu'une ``tranche'' d'une présentation
reste plate sur le champ algébrique.

\begin{lemma}
\label{stacks-morphisms-lemma-slice}
Soit $\mathcal{X}$ un champ algébrique.
Considérons un diagramme cartésien
$$
\xymatrix{
U \ar[d] & F \ar[l]^p \ar[d] \\
\mathcal{X} & \Spec(k) \ar[l]
}
$$
où $U$ est un espace algébrique, $k$ un corps, et où $U \to \mathcal{X}$
est plat et localement de présentation finie. Soient
$f_1, \ldots, f_r \in \Gamma(U, \mathcal{O}_U)$
et $z \in |F|$ tels que les images de $f_1, \ldots, f_r$ forment une suite
régulière dans l'anneau local $\mathcal{O}_{F, \overline{z}}$.
Alors, après avoir remplacé $U$ par un sous-espace ouvert contenant $p(z)$,
le morphisme
$$
V(f_1, \ldots, f_r) \longrightarrow \mathcal{X}
$$
est plat et localement de présentation finie.
\end{lemma}

\begin{proof}
Choisissons un schéma $W$ et un morphisme lisse surjectif
$W \to \mathcal{X}$. Choisissons une extension de corps $k'/k$ et un
morphisme $w : \Spec(k') \to W$ tels que
$\Spec(k') \to W \to \mathcal{X}$ soit $2$-isomorphe à
$\Spec(k') \to \Spec(k) \to \mathcal{X}$. Cela est possible puisque
$W \to \mathcal{X}$ est surjectif. Considérons le diagramme commutatif
$$
\xymatrix{
U \ar[d] &
U \times_\mathcal{X} W \ar[l]^-{\text{pr}_0} \ar[d] &
F' \ar[l]^-{p'} \ar[d] \\
\mathcal{X} &
W \ar[l] &
\Spec(k') \ar[l]
}
$$
dont les deux carrés sont cartésiens. D'après le choix de $w$, on a
$F' = F \times_{\Spec(k)} \Spec(k')$. Ainsi $F' \to F$ est surjectif, et
l'on peut choisir un point $z' \in |F'|$ qui s'envoie sur $z$.
Comme $F' \to F$ est plat, l'homomorphisme
$\mathcal{O}_{F, \overline{z}} \to \mathcal{O}_{F', \overline{z}'}$ est
plat, voir Morphismes d'espaces,
Lemme \ref{spaces-morphisms-lemma-flat-at-point-etale-local-rings}.
Les images de $f_1, \ldots, f_r$ forment donc une suite régulière dans
$\mathcal{O}_{F', \overline{z}'}$, voir Algèbre,
Lemme \ref{algebra-lemma-flat-increases-depth}.
Notons que $U \times_\mathcal{X} W \to W$ est un morphisme d'espaces
algébriques plat et localement de présentation finie. Par Compléments sur les
morphismes d'espaces, Lemme \ref{spaces-more-morphisms-lemma-slice}, il existe
donc un sous-espace ouvert $U'$ de $U \times_\mathcal{X} W$, contenant
$p'(z')$, tel que l'intersection
$U' \cap (V(f_1, \ldots, f_r) \times_\mathcal{X} W)$ soit plate et localement
de présentation finie sur $W$. Notons que $\text{pr}_0(U')$ est un
sous-espace ouvert de $U$ contenant $p(z)$, car $\text{pr}_0$ est lisse, donc
ouvert. Or
$U' \cap (V(f_1, \ldots, f_r) \times_\mathcal{X} W) \to \mathcal{X}$
est plat et localement de présentation finie comme composé
$$
U' \cap (V(f_1, \ldots, f_r) \times_\mathcal{X} W) \to W \to \mathcal{X}.
$$
Par conséquent, Propriétés des champs,
Lemme \ref{stacks-properties-lemma-check-property-after-precomposing},
montre que $\text{pr}_0(U') \cap V(f_1, \ldots, f_r) \to \mathcal{X}$ est
plat et localement de présentation finie, comme souhaité.
\end{proof}

\begin{lemma}
\label{stacks-morphisms-lemma-quasi-finite-at-point}
Soit $\mathcal{X}$ un champ algébrique. Considérons un diagramme cartésien
$$
\xymatrix{
U \ar[d] & F \ar[l]^p \ar[d] \\
\mathcal{X} & \Spec(k) \ar[l]
}
$$
où $U$ est un espace algébrique, $k$ un corps, et où $U \to \mathcal{X}$ est
localement de type fini. Soit $z \in |F|$ tel que $\dim_z(F) = 0$.
Alors, après avoir remplacé $U$ par un sous-espace ouvert contenant $p(z)$,
le morphisme
$$
U \longrightarrow \mathcal{X}
$$
est localement quasi-fini.
\end{lemma}

\begin{proof}
Comme $f : U \to \mathcal{X}$ est localement de type fini, il existe un ouvert
maximal $W(f) \subset U$ tel que la restriction
$f|_{W(f)} : W(f) \to \mathcal{X}$ soit localement quasi-finie, voir
Propriétés des champs, Remarque
\ref{stacks-properties-remark-local-source-apply}
(\ref{stacks-properties-item-loc-quasi-finite}).
Il suffit donc de montrer que $p(z)$ appartient à $W(f)$.
De plus, la remarque citée ci-dessus montre aussi que la formation de $W(f)$
commute à tout changement de base par un morphisme représentable par des
espaces algébriques. Il suffit donc de montrer que le morphisme
$F \to \Spec(k)$ est localement quasi-fini en $z$. Cela résulte immédiatement
de Morphismes d'espaces,
Lemme \ref{spaces-morphisms-lemma-locally-quasi-finite-rel-dimension-0}.
\end{proof}

\noindent
Un champ quasi-DM admet un ``recouvrement'' localement quasi-fini par un schéma.

\begin{theorem}
\label{stacks-morphisms-theorem-quasi-DM}
Soit $\mathcal{X}$ un champ algébrique. Les conditions suivantes sont
équivalentes :
\begin{enumerate}
\item $\mathcal{X}$ est quasi-DM, et
\item il existe un schéma $W$ et un morphisme $W \to \mathcal{X}$ surjectif,
plat, localement de présentation finie et localement quasi-fini.
\end{enumerate}
\end{theorem}

\begin{proof}
L'implication (2) $\Rightarrow$ (1) est le contenu du
Lemme \ref{stacks-morphisms-lemma-properties-covering-imply-diagonal}.
Supposons (1).
Soit $x \in |\mathcal{X}|$ un point de type fini. Nous allons construire un
schéma sur $\mathcal{X}$ qui ``convient'' dans un voisinage de $x$. À la fin
de la démonstration, nous prendrons la réunion disjointe de tous ces schémas.

\medskip\noindent
Soient $U$ un schéma affine, $U \to \mathcal{X}$ un morphisme lisse et
$u \in U$ un point fermé qui s'envoie sur $x$, voir le
Lemme \ref{stacks-morphisms-lemma-point-finite-type}.
Notons, comme d'habitude, $u = \Spec(\kappa(u))$. Considérons le diagramme
commutatif suivant
$$
\xymatrix{
u \ar[d] & R \ar[l] \ar[d] \\
U \ar[d] & F \ar[d] \ar[l]^p \\
\mathcal{X} & u \ar[l]
}
$$
dont les deux carrés sont des carrés de produit fibré ; en particulier,
$R = u \times_\mathcal{X} u$. Dans la démonstration du
Lemme \ref{stacks-morphisms-lemma-point-finite-type-monomorphism}, nous avons vu que
$(u, R, s, t, c)$ est un groupoïde en espaces algébriques dont les morphismes
$s, t$ sont localement de type fini. Soit $G \to u$ l'espace algébrique en groupes
stabilisateur (voir Groupoïdes en espaces, Définition
\ref{spaces-groupoids-definition-stabilizer-groupoid}).
Notons que
$$
G = R \times_{(u \times u)} u =
(u \times_\mathcal{X} u) \times_{(u \times u)} u =
\mathcal{X} \times_{\mathcal{X} \times \mathcal{X}} u.
$$
Comme $\mathcal{X}$ est quasi-DM, $G$ est localement quasi-fini sur $u$. Par
Compléments sur les groupoïdes en espaces, Lemme
\ref{spaces-more-groupoids-lemma-groupoid-on-field-dimension-equal-stabilizer}
on a $\dim(R) = 0$.

\medskip\noindent
Soit $e : u \to R$ l'identité du groupoïde. Le composé
$u \to R \to u$ est le morphisme identité de $u$.
Notons que $R \subset F$ est un sous-espace fermé, puisque $u \subset U$ est
un sous-schéma fermé. On peut donc aussi voir $e$ comme un point de $F$ ;
notons-le $z$. Considérons les homomorphismes d'anneaux locaux étales
$$
\mathcal{O}_{U, u}
\xrightarrow{p^\sharp}
\mathcal{O}_{F, \overline{z}}
\longrightarrow
\mathcal{O}_{R, \overline{e}}
$$
L'anneau $\mathcal{O}_{R, \overline{e}}$ est de dimension $0$ d'après le
premier paragraphe. D'autre part, le noyau de la seconde flèche est
$p^\sharp(\mathfrak m_u)\mathcal{O}_{F, \overline{z}}$, puisque $R$ est défini
dans $F$ par $\mathfrak m_u$. Ainsi,
$$
\mathfrak m_{\overline{z}} =
\sqrt{p^\sharp(\mathfrak m_u)\mathcal{O}_{F, \overline{z}}}
$$
D'autre part, puisque le morphisme $U \to \mathcal{X}$ est lisse,
$F \to u$ est un morphisme lisse d'espaces algébriques.
Cela signifie que $F$ est un espace algébrique régulier
(Espaces sur des corps, Lemme \ref{spaces-over-fields-lemma-smooth-regular}).
Par conséquent, $\mathcal{O}_{F, \overline{z}}$ est un anneau local régulier
(Propriétés des espaces, Lemme \ref{spaces-properties-lemma-regular}).
Rappelons qu'un anneau local régulier est de Cohen--Macaulay
(Algèbre, Lemme \ref{algebra-lemma-regular-ring-CM}).
Posons $d = \dim(\mathcal{O}_{F, \overline{z}})$. Par Algèbre,
Lemme \ref{algebra-lemma-find-sequence-image-regular}, on peut trouver
$f_1, \ldots, f_d \in \mathcal{O}_{U, u}$ dont les images
$p^\sharp(f_1), \ldots, p^\sharp(f_d)$ forment une suite régulière dans
$\mathcal{O}_{F, \overline{z}}$. Par le Lemme \ref{stacks-morphisms-lemma-slice}, quitte à
restreindre $U$, nous pouvons supposer que
$Z = V(f_1, \ldots, f_d) \to \mathcal{X}$ est plat et localement de
présentation finie. Par construction,
$F_Z = Z \times_\mathcal{X} u$ est un sous-espace fermé de
$F = U \times_\mathcal{X} u$, le point $z$ appartient à ce sous-espace fermé,
et
$$
\dim(\mathcal{O}_{F_Z, \overline{z}}) = 0.
$$
Par Morphismes d'espaces,
Lemme \ref{spaces-morphisms-lemma-dimension-fibre-at-a-point},
il s'ensuit que $\dim_z(F_Z) = 0$, car le degré de transcendance de $z$
relativement à $u$ est nul. Le Lemme \ref{stacks-morphisms-lemma-quasi-finite-at-point} montre
donc que, quitte à restreindre encore $U$, le morphisme
$Z \to \mathcal{X}$ est localement quasi-fini.

\medskip\noindent
Nous concluons que, pour tout point de type fini $x$ de $\mathcal{X}$, il
existe un morphisme $f_x : Z_x \to \mathcal{X}$ localement quasi-fini, plat
et localement de présentation finie, tel que $x$ appartienne à l'image de
$|f_x|$. Posons $W = \coprod_x Z_x$ et $f = \coprod f_x$. Alors $f$ est plat,
localement de présentation finie et localement quasi-fini. En particulier,
l'image de $|f|$ est ouverte, voir Propriétés des champs,
Lemme \ref{stacks-properties-lemma-topology-points}.
Par construction, l'image contient tous les points de type fini de
$\mathcal{X}$ ; le morphisme $f$ est donc surjectif par le
Lemme \ref{stacks-morphisms-lemma-enough-finite-type-points} (et Propriétés des champs,
Lemme \ref{stacks-properties-lemma-characterize-surjective}).
\end{proof}

\begin{lemma}
\label{stacks-morphisms-lemma-DM-residual-gerbe}
Soit $\mathcal{Z}$ un champ algébrique DM, localement noethérien et réduit,
tel que $|\mathcal{Z}|$ soit réduit à un point. Alors il existe un corps $k$
et un morphisme étale surjectif $\Spec(k) \to \mathcal{Z}$.
\end{lemma}

\begin{proof}
Par Propriétés des champs,
Lemme \ref{stacks-properties-lemma-unique-point-better},
il existe un corps $k$ et un morphisme $\Spec(k) \to \mathcal{Z}$ surjectif,
plat et localement de présentation finie. Posons $U = \Spec(k)$ et
$R = U \times_\mathcal{Z} U$ ; on obtient ainsi un groupoïde en espaces
algébriques $(U, R, s, t, c)$, voir Champs algébriques, Lemme
\ref{algebraic-lemma-criterion-map-representable-spaces-fibred-in-groupoids}.
Notons que, par Champs algébriques,
Remarque \ref{algebraic-remark-flat-fp-presentation}, on a une équivalence
$$
f_{can} : [U/R] \longrightarrow \mathcal{Z}
$$
Les projections $s, t : R \to U$ sont localement de présentation finie.
Comme $\mathcal{Z}$ est DM, l'espace algébrique en groupes stabilisateur
$$
G = U \times_{U \times U} R = U \times_{U \times U} (U \times_\mathcal{Z} U) =
U \times_{\mathcal{Z} \times \mathcal{Z}, \Delta_\mathcal{Z}} \mathcal{Z}
$$
est non ramifié sur $U$. En particulier, $\dim(G) = 0$ et, par Compléments
sur les groupoïdes en espaces, Lemme
\ref{spaces-more-groupoids-lemma-groupoid-on-field-dimension-equal-stabilizer}
on a $\dim(R) = 0$. Il s'ensuit que $R$ est un schéma, voir Espaces sur des
corps, Lemme
\ref{spaces-over-fields-lemma-locally-finite-type-dim-zero}.
Par Variétés, Lemme \ref{varieties-lemma-algebraic-scheme-dim-0},
$R$ (ainsi que $G$) est réunion disjointe de spectres d'anneaux locaux
artiniens, finis sur $k$ aussi bien par $s$ que par $t$. Soit
$P = \Spec(A) \subset R$ le sous-schéma ouvert et fermé dont le point
sous-jacent est l'identité $e$ du schéma en groupoïdes $(U, R, s, t, c)$.
Comme $s \circ e = t \circ e = \text{id}_{\Spec(k)}$, l'anneau $A$ est local
artinien et son corps résiduel s'identifie à $k$ aussi bien par
$s^\sharp : k \to A$ que par $t^\sharp : k \to A$.
Notons que $s, t : \Spec(A) \to \Spec(k)$ sont finis (par le lemme cité
ci-dessus). Puisque $G \to \Spec(k)$ est non ramifié, on a
$$
G \cap P = P \times_{U \times U} U = \Spec(A \otimes_{k \otimes k} k)
$$
qui est non ramifié sur $k$. D'autre part,
$A \otimes_{k \otimes k} k$ est local, comme quotient de $A$, et admet une
surjection sur $k$. On en conclut que $A \otimes_{k \otimes k} k = k$. Il s'ensuit que
$P \to U \times U$ est universellement injectif (car $P$ n'a qu'un point,
de corps résiduel $k$), non ramifié (d'après le calcul ci-dessus de la fibre
sur l'unique point image) et de type fini (parce que $s, t$ le sont) ; c'est
donc un monomorphisme (voir Morphismes étales, Lemme
\ref{etale-lemma-universally-injective-unramified}).
Ainsi, $s|_P, t|_P : P \to U$ définissent une relation d'équivalence finie
plate. La Proposition \ref{groupoids-proposition-finite-flat-equivalence} de
Groupoïdes montre alors que $U/P$ existe et est un schéma $\overline{U}$.
De plus, $U \to \overline{U}$ est fini localement libre et
$P = U \times_{\overline{U}} U$.
En fait, $\overline{U} = \Spec(k_0)$, où $k_0 \subset k$ est l'anneau des
fonctions invariantes par $R$. Comme $k$ est un corps, la définition donnée
dans Groupoïdes, Équation (\ref{groupoids-equation-invariants}), montre que
$k_0$ est un corps.

\medskip\noindent
Nous affirmons que
\begin{equation}
\label{stacks-morphisms-equation-etale-covering}
\Spec(k_0) = \overline{U} = U/P \to [U/R] = \mathcal{Z}
\end{equation}
est le morphisme étale surjectif cherché. Il résulte de Propriétés des champs,
Lemme \ref{stacks-properties-lemma-flat-cover-by-field}, que ce morphisme est
surjectif. Il suffit donc de montrer que (\ref{stacks-morphisms-equation-etale-covering}) est
étale\footnote{Nous invitons le lecteur à trouver sa propre démonstration de
ce fait. L'argument a en réalité beaucoup en commun avec l'argument final de
la démonstration d'Amorçage, Théorème
\ref{bootstrap-theorem-final-bootstrap} ; il devrait donc probablement être
isolé quelque part dans un lemme qui lui soit propre.}.
Au lieu de démontrer directement l'étalité, appliquons d'abord Amorçage,
Lemme \ref{bootstrap-lemma-divide-subgroupoid}. Il existe alors un schéma en
groupoïdes
$(\overline{U}, \overline{R}, \overline{s}, \overline{t}, \overline{c})$
tel que $(U, R, s, t, c)$ soit la restriction de
$(\overline{U}, \overline{R}, \overline{s}, \overline{t}, \overline{c})$
par le morphisme quotient $U \to \overline{U}$.
(Nous avons vérifié ci-dessus toutes les hypothèses du lemme, hormis que
$j : R \to U \times U$ est séparé et localement quasi-fini ; cela résulte du
fait que $R$ est un schéma séparé localement quasi-fini sur $k$.) Comme
$U \to \overline{U}$ est fini localement libre,
$[U/R] \to [\overline{U}/\overline{R}]$ est une équivalence, voir Groupoïdes
en espaces,
Lemme \ref{spaces-groupoids-lemma-quotient-stack-restrict-equivalence}.

\medskip\noindent
Notons que $s, t$ sont les changements de base des morphismes
$\overline{s}, \overline{t}$ par $U \to \overline{U}$.
Comme $\{U \to \overline{U}\}$ est un recouvrement fppf, on en conclut que
$\overline{s}, \overline{t}$ sont plats, localement de présentation finie et
localement quasi-finis, voir Descente, Lemmes
\ref{descent-lemma-descending-property-flat},
\ref{descent-lemma-descending-property-locally-finite-presentation} et
\ref{descent-lemma-descending-property-quasi-finite}.
Considérons le diagramme commutatif
$$
\xymatrix{
U \times_{\overline{U}} U \ar@{=}[r] \ar[rd] & P \ar[r] \ar[d] & R \ar[d] \\
& \overline{U} \ar[r]^{\overline{e}} & \overline{R}
}
$$
Il est général, pour les restrictions, que les quatre sommets extérieurs
forment un diagramme cartésien. L'égalité montre que le carré intérieur est
cartésien. Comme $P$ est ouvert dans $R$, on en conclut que $\overline{e}$
est une immersion ouverte par Descente,
Lemme \ref{descent-lemma-descending-property-open-immersion}.

\medskip\noindent
Bien entendu, si $\overline{e}$ est une immersion ouverte et si
$\overline{s}, \overline{t}$ sont plats et localement de présentation finie,
alors les morphismes $\overline{t}, \overline{s}$ sont étales.
On peut le voir, par exemple, en appliquant Compléments sur les groupoïdes,
Lemme \ref{more-groupoids-lemma-sheaf-differentials}, qui montre que
$\Omega_{\overline{R}/\overline{U}} = 0$ implique que
$\overline{s}, \overline{t} : \overline{R} \to \overline{U}$ sont non
ramifiés (voir Morphismes,
Lemme \ref{morphisms-lemma-unramified-omega-zero}), ce qui implique à son tour
que $\overline{s}, \overline{t}$ sont étales (voir Morphismes,
Lemme \ref{morphisms-lemma-flat-unramified-etale}).
Ainsi, $\mathcal{Z} = [\overline{U}/\overline{R}]$ est une présentation étale
du champ algébrique $\mathcal{Z}$, et l'on conclut que
$\overline{U} \to \mathcal{Z}$ est étale par Propriétés des champs, Lemme
\ref{stacks-properties-lemma-check-property-covering}.
\end{proof}

\begin{lemma}
\label{stacks-morphisms-lemma-etale-at-point}
Soit $\mathcal{X}$ un champ algébrique. Considérons un diagramme cartésien
$$
\xymatrix{
U \ar[d] & F \ar[l]^p \ar[d] \\
\mathcal{X} & \Spec(k) \ar[l]
}
$$
où $U$ est un espace algébrique, $k$ un corps, et où $U \to \mathcal{X}$ est
plat et localement de présentation finie. Soit $z \in |F|$ tel que
$F \to \Spec(k)$ soit non ramifié en $z$. Alors, après avoir remplacé $U$ par
un sous-espace ouvert contenant $p(z)$, le morphisme
$$
U \longrightarrow \mathcal{X}
$$
est étale.
\end{lemma}

\begin{proof}
Comme $f : U \to \mathcal{X}$ est plat et localement de présentation finie,
il existe un ouvert maximal $W(f) \subset U$ tel que la restriction
$f|_{W(f)} : W(f) \to \mathcal{X}$ soit étale, voir Propriétés des champs,
Remarque
\ref{stacks-properties-remark-local-source-apply}
(\ref{stacks-properties-item-etale}).
Il suffit donc de montrer que $p(z)$ appartient à $W(f)$.
De plus, la remarque citée ci-dessus montre que la formation de $W(f)$ commute
à tout changement de base par un morphisme représentable par des espaces
algébriques. Il suffit donc de montrer que $F \to \Spec(k)$ est étale en $z$.
Ce morphisme est plat et localement de présentation finie, comme changement
de base de $U \to \mathcal{X}$, et, puisque $F \to \Spec(k)$ est non ramifié
en $z$ par hypothèse,
le résultat découle donc de Morphismes d'espaces,
Lemme \ref{spaces-morphisms-lemma-unramified-flat-lfp-etale}.
\end{proof}

\noindent
Un champ DM est un champ de Deligne--Mumford.

\begin{theorem}
\label{stacks-morphisms-theorem-DM}
Soit $\mathcal{X}$ un champ algébrique. Les conditions suivantes sont
équivalentes :
\begin{enumerate}
\item $\mathcal{X}$ est DM,
\item $\mathcal{X}$ est de Deligne--Mumford, et
\item il existe un schéma $W$ et un morphisme étale surjectif
$W \to \mathcal{X}$.
\end{enumerate}
\end{theorem}

\begin{proof}
Rappelons que (3) est la définition de (2), voir Champs algébriques,
Définition \ref{algebraic-definition-deligne-mumford}.
L'implication (3) $\Rightarrow$ (1) est le contenu du
Lemme \ref{stacks-morphisms-lemma-properties-covering-imply-diagonal}.
Supposons (1). Soit $x \in |\mathcal{X}|$ un point de type fini.
Nous allons construire un schéma sur $\mathcal{X}$ qui ``convient'' dans un
voisinage de $x$. À la fin de la démonstration, nous prendrons la réunion
disjointe de tous ces schémas.

\medskip\noindent
Par le Lemme \ref{stacks-morphisms-lemma-point-finite-type-monomorphism}, la gerbe résiduelle
$\mathcal{Z}_x$ de $\mathcal{X}$ en $x$ existe et
$\mathcal{Z}_x \to \mathcal{X}$ est localement de type fini. Par le
Lemme \ref{stacks-morphisms-lemma-separation-properties-residual-gerbe}, le champ algébrique
$\mathcal{Z}_x$ est DM. Par le Lemme \ref{stacks-morphisms-lemma-DM-residual-gerbe}, il existe
un corps $k$ et un morphisme étale surjectif
$z : \Spec(k) \to \mathcal{Z}_x$.
En particulier, le composé $x : \Spec(k) \to \mathcal{X}$ est localement de
type fini (par Morphismes d'espaces, Lemmes
\ref{spaces-morphisms-lemma-composition-finite-type} et
\ref{spaces-morphisms-lemma-etale-locally-finite-type}).

\medskip\noindent
Choisissons un schéma $U$ et un morphisme lisse $U \to \mathcal{X}$ tels que
$x$ appartienne à l'image de $|U| \to |\mathcal{X}|$.
Considérons le carré de produit fibré suivant
$$
\xymatrix{
U \ar[d] & F  \ar[l] \ar[d] \\
\mathcal{X} & \Spec(k) \ar[l]_-x
}
$$
autrement dit $F = U \times_{\mathcal{X}, x} \Spec(k)$. Par Propriétés des
champs, Lemme \ref{stacks-properties-lemma-points-cartesian}, l'espace $F$
est non vide. Comme $\mathcal{Z}_x \to \mathcal{X}$ est un monomorphisme,
on a
$$
\Spec(k) \times_{z, \mathcal{Z}_x, z} \Spec(k)
=
\Spec(k) \times_{x, \mathcal{X}, x} \Spec(k)
$$
dont les projections vers $\Spec(k)$ sont étales par construction de $z$.
Comme
$$
F \times_U F =
(\Spec(k) \times_\mathcal{X} \Spec(k))
\times_{\Spec(k)} F
$$
les projections $F \times_U F \to F$ sont elles aussi étales.
Il s'ensuit que $\Delta_{F/U} : F \to F \times_U F$ est étale (voir
Morphismes d'espaces, Lemme \ref{spaces-morphisms-lemma-etale-permanence}).
Par Morphismes d'espaces, Lemme
\ref{spaces-morphisms-lemma-etale-universally-injective-open}
cela implique que $\Delta_{F/U}$ est une immersion ouverte ; enfin, par
Morphismes d'espaces, Lemme
\ref{spaces-morphisms-lemma-diagonal-unramified-morphism}
il s'ensuit que $F \to U$ est non ramifié.

\medskip\noindent
Choisissons un schéma affine non vide $V$ et un morphisme étale $V \to F$.
(On pourrait éviter ce choix en travaillant directement avec $F$, mais il
semble ainsi plus facile d'expliquer ce qui se passe.) Représentons la
situation par le diagramme
$$
\xymatrix{
U \ar[d] & F  \ar[l] \ar[d] & V \ar[l] \ar[ld] \\
\mathcal{X} & \Spec(k) \ar[l]_-x
}
$$
Alors $V \to \Spec(k)$ est un morphisme lisse de schémas et $V \to U$ un
morphisme non ramifié de schémas (voir Morphismes d'espaces, Lemmes
\ref{spaces-morphisms-lemma-composition-smooth} et
\ref{spaces-morphisms-lemma-composition-unramified}).
Choisissons un point fermé $v \in V$ tel que l'extension
$k \subset \kappa(v)$ soit finie séparable, voir Variétés,
Lemme \ref{varieties-lemma-smooth-separable-closed-points-dense}.
Soit $u \in U$ le point image. L'anneau local $\mathcal{O}_{V, v}$ est
régulier (voir Variétés, Lemme \ref{varieties-lemma-smooth-regular}), et
l'homomorphisme d'anneaux locaux
$$
\varphi : \mathcal{O}_{U, u} \longrightarrow \mathcal{O}_{V, v}
$$
induit par le morphisme $V \to U$ vérifie
$\varphi(\mathfrak m_u)\mathcal{O}_{V, v} = \mathfrak m_v$,
voir Morphismes, Lemme \ref{morphisms-lemma-unramified-at-point}.
On peut donc trouver $f_1, \ldots, f_d \in \mathcal{O}_{U, u}$ tels que les
images $\varphi(f_1), \ldots, \varphi(f_d)$ forment une base de
$\mathfrak m_v/\mathfrak m_v^2$ sur $\kappa(v)$.
Comme $\mathcal{O}_{V, v}$ est un anneau local régulier, il s'ensuit que
$\varphi(f_1), \ldots, \varphi(f_d)$ forment une suite régulière dans
$\mathcal{O}_{V, v}$ (voir Algèbre,
Lemme \ref{algebra-lemma-regular-ring-CM}).
Après avoir remplacé $U$ par un voisinage ouvert de $u$, on peut supposer que
$f_1, \ldots, f_d \in \Gamma(U, \mathcal{O}_U)$. En remplaçant encore $U$
par un voisinage ouvert éventuellement plus petit de $u$, on peut supposer que
$V(f_1, \ldots, f_d) \to \mathcal{X}$ est plat et localement de présentation
finie, voir le Lemme \ref{stacks-morphisms-lemma-slice}.
Par construction,
$$
V(f_1, \ldots, f_d) \times_\mathcal{X} \Spec(k)
\longleftarrow
V(f_1, \ldots, f_d) \times_U V
$$
est étale, et $V(f_1, \ldots, f_d) \times_U V$ est le sous-schéma fermé
$T \subset V$ défini par $f_1|_V, \ldots, f_d|_V$.
Par construction, $v \in T$ et
$$
\mathcal{O}_{T, v} =
\mathcal{O}_{V, v}/(\varphi(f_1), \ldots, \varphi(f_d)) = \kappa(v)
$$
qui est une extension finie séparable de $k$. Il s'ensuit que
$T \to \Spec(k)$ est non ramifié en $v$, voir Morphismes,
Lemme \ref{morphisms-lemma-unramified-at-point}.
Par définition d'un morphisme non ramifié d'espaces algébriques, cela signifie
que $V(f_1, \ldots, f_d) \times_\mathcal{X} \Spec(k) \to \Spec(k)$ est non
ramifié au point image de $v$ dans
$V(f_1, \ldots, f_d) \times_\mathcal{X} \Spec(k)$.
En appliquant le Lemme \ref{stacks-morphisms-lemma-etale-at-point}, on voit que, quitte à
restreindre encore $U$ à un voisinage ouvert de $u$, le morphisme
$V(f_1, \ldots, f_d) \to \mathcal{X}$ est étale.

\medskip\noindent
Nous concluons que, pour tout point de type fini $x$ de $\mathcal{X}$, il
existe un morphisme étale $f_x : W_x \to \mathcal{X}$ tel que $x$ appartienne
à l'image de $|f_x|$. Posons $W = \coprod_x W_x$ et $f = \coprod f_x$.
Alors $f$ est étale. En particulier, l'image de $|f|$ est ouverte, voir
Propriétés des champs, Lemme \ref{stacks-properties-lemma-topology-points}.
Par construction, l'image contient tous les points de type fini de
$\mathcal{X}$ ; le morphisme $f$ est donc surjectif par le
Lemme \ref{stacks-morphisms-lemma-enough-finite-type-points} (et Propriétés des champs,
Lemme \ref{stacks-properties-lemma-characterize-surjective}).
\end{proof}

\noindent
Voici un corollaire utile qui montre que les ``fibres'' d'un morphisme DM
de champs algébriques sont de Deligne--Mumford.

\begin{lemma}
\label{stacks-morphisms-lemma-DM}
Soit $f : \mathcal{X} \to \mathcal{Y}$ un morphisme DM de champs
algébriques. Alors
\begin{enumerate}
\item pour tout champ algébrique DM $\mathcal{Z}$ et tout morphisme
$\mathcal{Z} \to \mathcal{Y}$, il existe un schéma et un morphisme
étale surjectif
$U \to \mathcal{X} \times_\mathcal{Y} \mathcal{Z}$.
\item pour tout espace algébrique $Z$ et tout morphisme
$Z \to \mathcal{Y}$, il existe un schéma et un morphisme
étale surjectif
$U \to \mathcal{X} \times_\mathcal{Y} Z$.
\end{enumerate}
\end{lemma}

\begin{proof}
Preuve de (1). Comme $f$ est DM, son changement de base
$\mathcal{X} \times_\mathcal{Y} \mathcal{Z} \to \mathcal{Z}$ est DM
par le Lemme \ref{stacks-morphisms-lemma-base-change-separated}.
Puisque $\mathcal{Z}$ est DM, il en résulte que
$\mathcal{X} \times_\mathcal{Y} \mathcal{Z}$ est DM
par le Lemme \ref{stacks-morphisms-lemma-separated-over-separated}. Il existe donc un
schéma $U$ et un morphisme étale surjectif
$U \to \mathcal{X} \times_\mathcal{Y} \mathcal{Z}$, voir le
Théorème \ref{stacks-morphisms-theorem-DM}.
La partie (2) est un cas particulier de (1), puisqu'un espace algébrique
(considéré comme un champ algébrique) est DM par le
Lemme \ref{stacks-morphisms-lemma-trivial-implications}.
\end{proof}






\section{Le lieu de Deligne--Mumford}
\label{stacks-morphisms-section-DM}

\noindent
Tout champ algébrique admet un plus grand sous-champ ouvert qui soit un champ
de Deligne--Mumford ; cela est plus ou moins clair, mais nous en donnons tout
de même la démonstration ci-dessous. Bien entendu, ce sous-champ peut être
vide, par exemple si $X = [\Spec(\mathbf{Z})/\mathbf{G}_{m, \mathbf{Z}}]$.
Nous caractériserons ci-dessous les points du lieu DM.

\begin{lemma}
\label{stacks-morphisms-lemma-open-DM-locus}
Soit $\mathcal{X}$ un champ algébrique. Il existe des sous-champs ouverts
$$
\mathcal{X}'' \subset \mathcal{X}' \subset \mathcal{X}
$$
tels que $\mathcal{X}''$ soit DM et $\mathcal{X}'$ quasi-DM, et ce sont les
plus grands sous-champs ouverts possédant respectivement ces propriétés.
\end{lemma}

\begin{proof}
Tout ce que nous affirmons ici est que, si
$\mathcal{U} \subset \mathcal{X}$ et $\mathcal{V} \subset \mathcal{X}$
sont des sous-champs ouverts qui sont DM, alors le sous-champ ouvert
$\mathcal{W} \subset \mathcal{X}$ défini par
$|\mathcal{W}| = |\mathcal{U}| \cup |\mathcal{V}|$ est lui aussi DM.
(Il en va de même pour quasi-DM.) Même s'il y a là un léger artifice,
utilisons le Théorème \ref{stacks-morphisms-theorem-DM} pour le démontrer.
D'après ce théorème, on peut choisir des schémas $U$ et $V$ et des morphismes
étales surjectifs $U \to \mathcal{U}$ et $V \to \mathcal{V}$.
Alors, bien entendu, $U \amalg V \to \mathcal{W}$ est étale et surjectif.
Le cas quasi-DM se démontre exactement de la même manière en utilisant le
Théorème \ref{stacks-morphisms-theorem-quasi-DM}.
\end{proof}

\begin{lemma}
\label{stacks-morphisms-lemma-points-DM-locus}
Soit $\mathcal{X}$ un champ algébrique. Soit $x \in |\mathcal{X}|$ le point
correspondant à $x : \Spec(k) \to \mathcal{X}$. Soit $G_x/k$ l'espace
algébrique en groupes des automorphismes de $x$. Alors
\begin{enumerate}
\item $x$ appartient au lieu DM de $\mathcal{X}$ si et seulement si
$G_x \to \Spec(k)$ est non ramifié, et
\item $x$ appartient au lieu quasi-DM de $\mathcal{X}$ si et seulement si
$G_x \to \Spec(k)$ est localement quasi-fini.
\end{enumerate}
\end{lemma}

\begin{proof}
Preuve de (2). Choisissons un schéma $U$ et un morphisme lisse surjectif
$U \to \mathcal{X}$. Considérons le produit fibré
$$
\xymatrix{
G \ar[r] \ar[d] & \mathcal{I}_\mathcal{X} \ar[d] \\
U \ar[r] & \mathcal{X}
}
$$
Rappelons que $G$ est l'espace algébrique en groupes des automorphismes de
$U \to \mathcal{X}$. Par Groupoïdes en espaces, Lemme
\ref{spaces-groupoids-lemma-open-over-which-unramified-or-lqf}
il existe un plus grand sous-schéma ouvert $U' \subset U$ tel que
$G_{U'} \to U'$ soit localement quasi-fini.
En outre, la formation de $U'$ commute à tout changement de base.
En particulier, les deux images réciproques de $U'$ dans
$R = U \times_\mathcal{X} U$ sont le même sous-espace ouvert de $R$
(les deux morphismes $R \to \mathcal{X}$ sont en effet isomorphes, et leurs
espaces algébriques en groupes d'automorphismes sont donc isomorphes).
Par conséquent, $U'$ est l'image réciproque d'un sous-champ ouvert
$\mathcal{X}' \subset \mathcal{X}$ par Propriétés des champs, Lemme
\ref{stacks-properties-lemma-substacks-presentation}
et l'on a un diagramme cartésien
$$
\xymatrix{
G_{U'} \ar[r] \ar[d] & \mathcal{I}_{\mathcal{X}'} \ar[d] \\
U' \ar[r] & \mathcal{X}'
}
$$
Ainsi, le morphisme $\mathcal{I}_{\mathcal{X}'} \to \mathcal{X}'$ est
localement quasi-fini, et le Lemme \ref{stacks-morphisms-lemma-diagonal-diagonal}, partie (5),
montre que $\mathcal{X}'$ est quasi-DM. D'autre part, si
$\mathcal{W} \subset \mathcal{X}$ est un sous-champ ouvert quasi-DM,
l'image réciproque $W \subset U$ de $\mathcal{W}$ est contenue dans $U'$
par construction de $U'$, puisque
$\mathcal{I}_\mathcal{W} =
\mathcal{W} \times_\mathcal{X} \mathcal{I}_\mathcal{X}$
est localement quasi-fini sur $\mathcal{W}$.
Ainsi, $\mathcal{X}'$ est le lieu quasi-DM.
Enfin, choisissons une extension de corps $K/k$ et un diagramme
$2$-commutatif
$$
\xymatrix{
\Spec(K) \ar[r] \ar[d] & \Spec(k) \ar[d]^x \\
U \ar[r] & \mathcal{X}
}
$$
On obtient alors un isomorphisme
$G_x \times_{\Spec(k)} \Spec(K) \cong G \times_U \Spec(K)$
entre espaces algébriques en groupes sur $K$. Par conséquent, $G_x$ est
localement quasi-fini sur $k$ si et seulement si l'image de
$\Spec(K) \to U$ est contenue dans $U'$ (utiliser la commutation de la
formation de $U'$ aux changements de base et Groupoïdes en espaces,
Lemme
\ref{spaces-groupoids-lemma-open-over-which-unramified-or-lqf}
appliqué à $\Spec(K) \to \Spec(k)$ et à $G_x$).
Ceci achève la preuve de (2). Celle de (1) est exactement la même.
\end{proof}









\section{Morphismes localement quasi-finis}
\label{stacks-morphisms-section-locally-quasi-finite}

\noindent
La propriété ``localement quasi-fini'' des morphismes d'espaces algébriques
n'est pas locale pour la topologie lisse sur la source et la cible ; nous ne
pouvons donc pas utiliser les résultats de la
section \ref{stacks-morphisms-section-local-source-target} pour définir les morphismes
localement quasi-finis de champs algébriques. Nous savons déjà ce que signifie
être localement quasi-fini pour un morphisme de champs algébriques
représentable par des espaces algébriques, voir Propriétés des champs, section
\ref{stacks-properties-section-properties-morphisms}.
Pour trouver une condition convenant aux morphismes généraux, faisons
l'observation suivante.

\begin{lemma}
\label{stacks-morphisms-lemma-representable-by-spaces-quasi-finite}
Soit $f : \mathcal{X} \to \mathcal{Y}$ un morphisme de champs algébriques.
Supposons que $f$ soit représentable par des espaces algébriques.
Les conditions suivantes sont équivalentes :
\begin{enumerate}
\item $f$ est localement quasi-fini (au sens de Propriétés des champs,
section \ref{stacks-properties-section-properties-morphisms}), et
\item $f$ est localement de type fini et, pour tout morphisme
$\Spec(k) \to \mathcal{Y}$, où $k$ est un corps, l'espace
$|\Spec(k) \times_\mathcal{Y} \mathcal{X}|$ est discret.
\end{enumerate}
\end{lemma}

\begin{proof}
Supposons (1). Le morphisme d'espaces algébriques
$\mathcal{X}_k \to \Spec(k)$ est alors localement quasi-fini, car il est un
changement de base de $f$. Ainsi, $|\mathcal{X}_k|$ est discret par
Morphismes d'espaces, Lemme
\ref{spaces-morphisms-lemma-locally-quasi-finite}.
Réciproquement, supposons (2). Choisissons un morphisme lisse surjectif
$V \to \mathcal{Y}$, où $V$ est un schéma. Il suffit de montrer que le
morphisme d'espaces algébriques $V \times_\mathcal{Y} \mathcal{X} \to V$
est localement quasi-fini, voir Propriétés des champs, Lemme
\ref{stacks-properties-lemma-check-property-covering}.
Par hypothèse, le morphisme $V \times_\mathcal{Y} \mathcal{X} \to V$ est
localement de type fini. Pour tout morphisme $\Spec(k) \to V$, où $k$ est
un corps,
$$
\Spec(k) \times_V (V \times_\mathcal{Y} \mathcal{X}) =
\Spec(k) \times_\mathcal{Y} \mathcal{X}
$$
a par hypothèse un espace de points discret. On en conclut que
$V \times_\mathcal{Y} \mathcal{X} \to V$ est localement quasi-fini par
Morphismes d'espaces, Lemme
\ref{spaces-morphisms-lemma-locally-quasi-finite}.
\end{proof}

\noindent
Un morphisme de champs algébriques représentable par des espaces algébriques
est quasi-DM, voir le Lemme \ref{stacks-morphisms-lemma-trivial-implications}.
Avec le lemme précédent, cela montre que la définition suivante est compatible
avec la notion déjà existante dans le cas des morphismes représentables par
des espaces algébriques.

\begin{definition}
\label{stacks-morphisms-definition-locally-quasi-finite}
Soit $f : \mathcal{X} \to \mathcal{Y}$ un morphisme de champs algébriques.
On dit que $f$ est {\it localement quasi-fini} si $f$ est quasi-DM, localement
de type fini et si, pour tout morphisme $\Spec(k) \to \mathcal{Y}$, où
$k$ est un corps, l'espace $|\mathcal{X}_k|$ est discret.
\end{definition}

\noindent
La condition que $f$ soit quasi-DM est naturelle. Par exemple, soit $k$ un
corps et considérons le morphisme
$\pi : [\Spec(k)/\mathbf{G}_m] \to \Spec(k)$
dont les fibres sont réduites à un point et qui est localement de type fini.
Comme nous le verrons plus loin, ce morphisme est lisse de dimension relative
$-1$, tandis que nous voulons que les morphismes localement quasi-finis aient
la dimension relative $0$. Notons aussi que la section
$\Spec(k) \to [\Spec(k)/\mathbf{G}_m]$ n'a pas de fibres discrètes et n'est
donc pas localement quasi-finie. Nous souhaitons en outre la propriété de
permanence suivante : si $f : \mathcal{X} \to \mathcal{X}'$ est un morphisme
de champs algébriques localement quasi-fini sur le champ algébrique
$\mathcal{Y}$, alors $f$ est localement quasi-fini (un énoncé légèrement plus
fort est en fait vrai, voir le
Lemme \ref{stacks-morphisms-lemma-locally-quasi-finite-permanence}).

\medskip\noindent
La définition précédente est également justifiée par le
Lemme \ref{stacks-morphisms-lemma-characterize-locally-quasi-finite} ci-dessous, qui caractérise
la propriété d'être localement quasi-fini par l'existence de
``présentations'' ou de ``recouvrements'' convenables de
$\mathcal{X}$ et de $\mathcal{Y}$.

\begin{lemma}
\label{stacks-morphisms-lemma-base-change-locally-quasi-finite}
Tout changement de base d'un morphisme localement quasi-fini est localement
quasi-fini.
\end{lemma}

\begin{proof}
Nous l'avons vu pour les morphismes quasi-DM au
Lemme \ref{stacks-morphisms-lemma-base-change-separated} et pour les morphismes localement de
type fini au Lemme \ref{stacks-morphisms-lemma-base-change-finite-type}.
La condition sur les fibres se conserve manifestement par changement de base.
\end{proof}

\begin{lemma}
\label{stacks-morphisms-lemma-locally-quasi-finite-over-field}
Soit $\mathcal{X} \to \Spec(k)$ un morphisme localement quasi-fini, où
$\mathcal{X}$ est un champ algébrique et $k$ un corps. Soit
$f : V \to \mathcal{X}$ un morphisme localement quasi-fini, où $V$ est un
schéma. Alors $V \to \Spec(k)$ est localement quasi-fini.
\end{lemma}

\begin{proof}
Le Lemme \ref{stacks-morphisms-lemma-composition-finite-type} montre que
$V \to \Spec(k)$ est localement de type fini.
Supposons, afin d'obtenir une contradiction, que $V \to \Spec(k)$ ne soit
pas localement quasi-fini. Il existe alors une spécialisation non triviale
$v \leadsto v'$ de points de $V$, voir Morphismes, Lemme
\ref{morphisms-lemma-quasi-finite-at-point-characterize}.
En particulier,
$\text{trdeg}_k(\kappa(v)) > \text{trdeg}_k(\kappa(v'))$, voir Morphismes,
Lemme \ref{morphisms-lemma-dimension-fibre-specialization}.
Comme $|\mathcal{X}|$ est discret, on a $|f|(v) = |f|(v')$.
Considérons $R = V \times_\mathcal{X} V$. Alors $R$ est un espace algébrique
et les projections $s, t : R \to V$ sont localement quasi-finies comme
changements de base de $V \to \mathcal{X}$ (ce morphisme est représentable
par des espaces algébriques, si bien que cela résulte de la discussion de
Propriétés des champs, section
\ref{stacks-properties-section-properties-morphisms}).
Par Propriétés des champs,
Lemme \ref{stacks-properties-lemma-points-cartesian}, il existe
$r \in |R|$ tel que $s(r) = v$ et $t(r) = v'$.
Par Morphismes d'espaces,
Lemme \ref{spaces-morphisms-lemma-compare-tr-deg}, le degré de transcendance
de $v/k$ est égal à celui de $r/k$, lui-même égal à celui de $v'/k$.
Cette contradiction démontre le lemme.
\end{proof}

\begin{lemma}
\label{stacks-morphisms-lemma-composition-locally-quasi-finite}
Un composé de morphismes localement quasi-finis est localement quasi-fini.
\end{lemma}

\begin{proof}
Nous l'avons vu pour les morphismes quasi-DM au
Lemme \ref{stacks-morphisms-lemma-composition-separated} et pour les morphismes localement de
type fini au Lemme \ref{stacks-morphisms-lemma-composition-finite-type}.
Supposons que $\mathcal{X} \to \mathcal{Y}$ et
$\mathcal{Y} \to \mathcal{Z}$ soient localement quasi-finis. Soit $k$ un
corps et soit $\Spec(k) \to \mathcal{Z}$ un morphisme.
Il suffit de montrer que $|\mathcal{X}_k|$ est discret. Par le
Lemme \ref{stacks-morphisms-lemma-base-change-locally-quasi-finite}, les morphismes
$\mathcal{X}_k \to \mathcal{Y}_k$ et
$\mathcal{Y}_k \to \Spec(k)$ sont localement quasi-finis.
En particulier, $\mathcal{Y}_k$ est un champ algébrique quasi-DM, voir le
Lemme \ref{stacks-morphisms-lemma-separated-implies-morphism-separated}.
Par le Théorème \ref{stacks-morphisms-theorem-quasi-DM}, il existe un schéma $V$ et un
morphisme surjectif, plat, localement de présentation finie et localement
quasi-fini $V \to \mathcal{Y}_k$. Le
Lemme \ref{stacks-morphisms-lemma-locally-quasi-finite-over-field} montre que $V$ est
localement quasi-fini sur $k$ ; en particulier, $|V|$ est discret. Le morphisme
$V \times_{\mathcal{Y}_k} \mathcal{X}_k \to \mathcal{X}_k$ est surjectif,
plat et localement de présentation finie ; l'application
$|V \times_{\mathcal{Y}_k} \mathcal{X}_k| \to |\mathcal{X}_k|$ est donc
surjective et ouverte. Il suffit ainsi de montrer que
$|V \times_{\mathcal{Y}_k} \mathcal{X}_k|$ est discret.
Notons que $V$ est réunion disjointe de spectres de $k$-algèbres locales
artiniennes $A_i$, de corps résiduels $k_i$, voir Variétés,
Lemme \ref{varieties-lemma-algebraic-scheme-dim-0}.
Il suffit donc de montrer que chacun des espaces
$$
|\Spec(A_i) \times_{\mathcal{Y}_k} \mathcal{X}_k| =
|\Spec(k_i) \times_{\mathcal{Y}_k} \mathcal{X}_k| =
|\Spec(k_i) \times_\mathcal{Y} \mathcal{X}|
$$
est discret, ce qui résulte de l'hypothèse que
$\mathcal{X} \to \mathcal{Y}$ est localement quasi-fini.
\end{proof}

\noindent
Avant de caractériser les morphismes localement quasi-finis au moyen de
recouvrements, faisons-le pour les morphismes quasi-DM.

\begin{lemma}
\label{stacks-morphisms-lemma-characterize-quasi-DM}
Soit $f : \mathcal{X} \to \mathcal{Y}$ un morphisme de champs algébriques.
Les conditions suivantes sont équivalentes :
\begin{enumerate}
\item $f$ est quasi-DM,
\item pour tout morphisme $V \to \mathcal{Y}$, où $V$ est un espace
algébrique, il existe un morphisme surjectif, plat, localement de présentation
finie et localement quasi-fini
$U \to \mathcal{X} \times_\mathcal{Y} V$, où $U$ est un espace algébrique,
et
\item il existe des espaces algébriques $U$, $V$, un morphisme
$V \to \mathcal{Y}$ surjectif, plat et localement de présentation finie,
et un morphisme $U \to \mathcal{X} \times_\mathcal{Y} V$ surjectif, plat,
localement de présentation finie et localement quasi-fini.
\end{enumerate}
\end{lemma}

\begin{proof}
L'implication (2) $\Rightarrow$ (3) est immédiate.

\medskip\noindent
Supposons (1) et soit $V \to \mathcal{Y}$ comme en (2). Alors
$\mathcal{X} \times_\mathcal{Y} V \to V$ est quasi-DM, voir le
Lemme \ref{stacks-morphisms-lemma-base-change-separated}.
Par le Lemme \ref{stacks-morphisms-lemma-trivial-implications}, l'espace algébrique $V$ est DM,
donc quasi-DM. Ainsi, $\mathcal{X} \times_\mathcal{Y} V$ est quasi-DM par le
Lemme \ref{stacks-morphisms-lemma-separated-over-separated}.
On peut donc appliquer le Théorème \ref{stacks-morphisms-theorem-quasi-DM} pour obtenir le
morphisme $U \to \mathcal{X} \times_\mathcal{Y} V$ de (2).

\medskip\noindent
Supposons (3). Soient $V \to \mathcal{Y}$ et
$U \to \mathcal{X} \times_\mathcal{Y} V$ comme en (3).
Pour montrer que $f$ est quasi-DM, il suffit de montrer que
$\mathcal{X} \times_\mathcal{Y} V \to V$ est quasi-DM, voir le
Lemme \ref{stacks-morphisms-lemma-check-separated-covering}.
Le Lemme \ref{stacks-morphisms-lemma-properties-covering-imply-diagonal} montre que
$\mathcal{X} \times_\mathcal{Y} V$ est quasi-DM.
Par conséquent, $\mathcal{X} \times_\mathcal{Y} V \to V$ est quasi-DM par
le Lemme \ref{stacks-morphisms-lemma-separated-implies-morphism-separated}, et (1) est
vérifiée. Ceci achève la démonstration du lemme.
\end{proof}

\begin{lemma}
\label{stacks-morphisms-lemma-characterize-locally-quasi-finite}
Soit $f : \mathcal{X} \to \mathcal{Y}$ un morphisme de champs algébriques.
Les conditions suivantes sont équivalentes :
\begin{enumerate}
\item $f$ est localement quasi-fini,
\item $f$ est quasi-DM et, pour tout morphisme $V \to \mathcal{Y}$, où
$V$ est un espace algébrique, et tout morphisme localement quasi-fini
$U \to \mathcal{X} \times_\mathcal{Y} V$, où $U$ est un espace algébrique,
le morphisme $U \to V$ est localement quasi-fini,
\item pour tout morphisme $V \to \mathcal{Y}$ issu d'un espace algébrique
$V$, il existe un morphisme surjectif, plat, localement de présentation finie
et localement quasi-fini $U \to \mathcal{X} \times_\mathcal{Y} V$, où
$U$ est un espace algébrique, tel que $U \to V$ soit localement quasi-fini,
\item il existe des espaces algébriques $U$, $V$, un morphisme surjectif,
plat et localement de présentation finie $V \to \mathcal{Y}$, et un morphisme
$U \to \mathcal{X} \times_\mathcal{Y} V$ surjectif, plat, localement de
présentation finie et localement quasi-fini, tels que $U \to V$ soit
localement quasi-fini.
\end{enumerate}
\end{lemma}

\begin{proof}
Supposons (1). Alors $f$ est quasi-DM par hypothèse. Soient
$V \to \mathcal{Y}$ et $U \to \mathcal{X} \times_\mathcal{Y} V$
comme en (2). Par le Lemme \ref{stacks-morphisms-lemma-composition-locally-quasi-finite}, le
composé $U \to \mathcal{X} \times_\mathcal{Y} V \to V$ est localement
quasi-fini. Ainsi, (1) implique (2).

\medskip\noindent
Supposons (2). Soit $V \to \mathcal{Y}$ comme en (3). Par le
Lemme \ref{stacks-morphisms-lemma-characterize-quasi-DM}, on peut trouver un espace algébrique
$U$ et un morphisme surjectif, plat, localement de présentation finie et
localement quasi-fini $U \to \mathcal{X} \times_\mathcal{Y} V$.
Par (2), le composé $U \to V$ est localement quasi-fini. Ainsi, (2)
implique (3).

\medskip\noindent
Il est immédiat que (3) implique (4).

\medskip\noindent
Supposons (4). Nous allons établir (1), ce qui achèvera la démonstration.
Le Lemme \ref{stacks-morphisms-lemma-characterize-quasi-DM} montre que $f$ est quasi-DM.
Pour montrer que $f$ est localement de type fini, il suffit de montrer que
$g : \mathcal{X} \times_\mathcal{Y} V \to V$ est localement de type fini,
voir le Lemme \ref{stacks-morphisms-lemma-check-finite-type-covering}.
Il suffit ensuite de vérifier que le précomposé de $g$ avec
$h : U \to \mathcal{X} \times_\mathcal{Y} V$ est localement de type fini,
voir le Lemme \ref{stacks-morphisms-lemma-check-finite-type-precompose}.
Cela résulte de l'hypothèse que $g \circ h : U \to V$ est localement
quasi-fini ; ainsi, $f$ est localement de type fini.
Enfin, soit $k$ un corps et soit $\Spec(k) \to \mathcal{Y}$ un morphisme.
Alors $V \times_\mathcal{Y} \Spec(k)$ est un espace algébrique non vide,
localement de présentation finie sur $k$. Il existe donc une extension finie
$k'/k$ et un morphisme $\Spec(k') \to V$ tels que le diagramme
$$
\xymatrix{
\Spec(k') \ar[r] \ar[d] & V \ar[d] \\
\Spec(k) \ar[r] & \mathcal{Y}
}
$$
soit commutatif (détails omis). Alors
$\mathcal{X}_{k'} \to \mathcal{X}_k$ est représentable (par des schémas),
surjectif et fini localement libre. En particulier,
$|\mathcal{X}_{k'}| \to |\mathcal{X}_k|$ est surjectif et ouvert.
Il suffit donc de montrer que $|\mathcal{X}_{k'}|$ est discret. Puisque
$$
U \times_V \Spec(k') =
U \times_{\mathcal{X} \times_\mathcal{Y} V} \mathcal{X}_{k'}
$$
le morphisme $U \times_V \Spec(k') \to \mathcal{X}_{k'}$ est surjectif,
plat et localement de présentation finie (comme changement de base de
$U \to \mathcal{X} \times_\mathcal{Y} V$). L'application
$|U \times_V \Spec(k')| \to |\mathcal{X}_{k'}|$ est donc surjective et
ouverte. Il suffit ainsi de montrer que $|U \times_V \Spec(k')|$ est discret.
Cela résulte du fait que $U \to V$ est localement quasi-fini (soit par notre
définition ci-dessus, soit par la définition originelle pour les morphismes
d'espaces algébriques, au moyen de Morphismes d'espaces,
Lemme \ref{spaces-morphisms-lemma-locally-quasi-finite}).
\end{proof}

\begin{lemma}
\label{stacks-morphisms-lemma-locally-quasi-finite-permanence}
Soient $\mathcal{X} \to \mathcal{Y} \to \mathcal{Z}$ des morphismes de
champs algébriques. Supposons que $\mathcal{X} \to \mathcal{Z}$ soit
localement quasi-fini et que $\mathcal{Y} \to \mathcal{Z}$ soit quasi-DM.
Alors $\mathcal{X} \to \mathcal{Y}$ est localement quasi-fini.
\end{lemma}

\begin{proof}
Écrivons $\mathcal{X} \to \mathcal{Y}$ comme le composé
$$
\mathcal{X} \longrightarrow
\mathcal{X} \times_\mathcal{Z} \mathcal{Y} \longrightarrow
\mathcal{Y}
$$
La seconde flèche est localement quasi-finie comme changement de base de
$\mathcal{X} \to \mathcal{Z}$, voir le
Lemme \ref{stacks-morphisms-lemma-base-change-locally-quasi-finite}.
La première est localement quasi-finie par le
Lemme \ref{stacks-morphisms-lemma-semi-diagonal}, puisque
$\mathcal{Y} \to \mathcal{Z}$ est quasi-DM.
Ainsi, $\mathcal{X} \to \mathcal{Y}$ est localement quasi-fini par le
Lemme \ref{stacks-morphisms-lemma-composition-locally-quasi-finite}.
\end{proof}

















\section{Morphismes quasi-finis}
\label{stacks-morphisms-section-quasi-finite}

\noindent
Nous avons défini les morphismes ``localement quasi-finis'' de champs
algébriques à la section \ref{stacks-morphisms-section-locally-quasi-finite}, et les morphismes
``quasi-compacts'' de champs algébriques à la
section \ref{stacks-morphisms-section-quasi-compact}. Comme un morphisme d'espaces algébriques
est, par définition, quasi-fini si et seulement s'il est à la fois localement
quasi-fini et quasi-compact (Morphismes d'espaces, Définition
\ref{spaces-morphisms-definition-locally-quasi-finite}),
nous pouvons définir comme suit ce que signifie être quasi-fini pour un
morphisme de champs algébriques. Cette définition coïncide avec la notion déjà
définie dans Propriétés des champs,
section \ref{stacks-properties-section-properties-morphisms}, lorsque le
morphisme est représentable par des espaces algébriques.

\begin{definition}
\label{stacks-morphisms-definition-quasi-finite}
\begin{reference}
\cite{rydh_approx}
\end{reference}
Soit $f : \mathcal{X} \to \mathcal{Y}$ un morphisme de champs algébriques.
On dit que $f$ est {\it quasi-fini} si $f$ est localement quasi-fini
(Définition \ref{stacks-morphisms-definition-locally-quasi-finite})
et quasi-compact (Définition \ref{stacks-morphisms-definition-quasi-compact}).
\end{definition}

\begin{lemma}
\label{stacks-morphisms-lemma-composition-quasi-finite}
Le composé de morphismes quasi-finis est quasi-fini.
\end{lemma}

\begin{proof}
Il suffit de combiner les
Lemmes \ref{stacks-morphisms-lemma-composition-locally-quasi-finite} et
\ref{stacks-morphisms-lemma-composition-quasi-compact}.
\end{proof}

\begin{lemma}
\label{stacks-morphisms-lemma-base-change-quasi-finite}
Tout changement de base d'un morphisme quasi-fini est quasi-fini.
\end{lemma}

\begin{proof}
Il suffit de combiner les
Lemmes \ref{stacks-morphisms-lemma-base-change-locally-quasi-finite} et
\ref{stacks-morphisms-lemma-base-change-quasi-compact}.
\end{proof}

\begin{lemma}
\label{stacks-morphisms-lemma-quasi-finite-permanence}
Soient $f : \mathcal{X} \to \mathcal{Y}$ et
$g : \mathcal{Y} \to \mathcal{Z}$ des morphismes de champs algébriques.
Si $g \circ f$ est quasi-fini et si $g$ est quasi-séparé et quasi-DM,
alors $f$ est quasi-fini.
\end{lemma}

\begin{proof}
Il suffit de combiner les
Lemmes \ref{stacks-morphisms-lemma-locally-quasi-finite-permanence} et
\ref{stacks-morphisms-lemma-quasi-compact-permanence}.
\end{proof}












\section{Morphismes plats}
\label{stacks-morphisms-section-flat}

\noindent
La propriété ``être plat'' des morphismes d'espaces algébriques est locale
pour la topologie lisse sur la source et la cible, voir Descente sur les
espaces, Remarque \ref{spaces-descent-remark-list-local-source-target}.
Elle est aussi stable par changement de base et locale pour la topologie fpqc
sur la cible, voir Morphismes d'espaces,
Lemme \ref{spaces-morphisms-lemma-base-change-flat}, et Descente sur les
espaces, Lemme
\ref{spaces-descent-lemma-descending-property-flat}.
Par le Lemme \ref{stacks-morphisms-lemma-local-source-target} ci-dessus, nous pouvons donc
définir comme suit ce que signifie être plat pour un morphisme de champs
algébriques. Cette définition coïncide avec la notion déjà définie dans
Propriétés des champs,
section \ref{stacks-properties-section-properties-morphisms}, lorsque le
morphisme est représentable par des espaces algébriques.

\begin{definition}
\label{stacks-morphisms-definition-flat}
Soit $f : \mathcal{X} \to \mathcal{Y}$ un morphisme de champs algébriques.
On dit que $f$ est {\it plat} si les conditions équivalentes du
Lemme \ref{stacks-morphisms-lemma-local-source-target} sont satisfaites avec
$\mathcal{P} = \text{plat}$.
\end{definition}

\begin{lemma}
\label{stacks-morphisms-lemma-composition-flat}
Le composé de morphismes plats est plat.
\end{lemma}

\begin{proof}
Il suffit de combiner la Remarque \ref{stacks-morphisms-remark-composition} avec Morphismes
d'espaces, Lemme
\ref{spaces-morphisms-lemma-composition-flat}.
\end{proof}

\begin{lemma}
\label{stacks-morphisms-lemma-base-change-flat}
Tout changement de base d'un morphisme plat est plat.
\end{lemma}

\begin{proof}
Il suffit de combiner la Remarque \ref{stacks-morphisms-remark-base-change} avec Morphismes
d'espaces, Lemme
\ref{spaces-morphisms-lemma-base-change-flat}.
\end{proof}

\begin{lemma}
\label{stacks-morphisms-lemma-descent-flat}
Soit $f : \mathcal{X} \to \mathcal{Y}$ un morphisme de champs algébriques.
Soit $\mathcal{Z} \to \mathcal{Y}$ un morphisme plat surjectif de champs
algébriques. Si le changement de base
$\mathcal{Z} \times_\mathcal{Y} \mathcal{X} \to \mathcal{Z}$
est plat, alors $f$ est plat.
\end{lemma}

\begin{proof}
Choisissons un espace algébrique $W$ et un morphisme lisse surjectif
$W \to \mathcal{Z}$. Alors $W \to \mathcal{Z}$ est surjectif et plat
(Morphismes d'espaces, Lemme \ref{spaces-morphisms-lemma-smooth-flat}) ;
par conséquent, $W \to \mathcal{Y}$ est surjectif et plat (par Propriétés des
champs, Lemme \ref{stacks-properties-lemma-composition-surjective}, et le
Lemme \ref{stacks-morphisms-lemma-composition-flat}).
Puisque le changement de base de
$\mathcal{Z} \times_\mathcal{Y} \mathcal{X} \to \mathcal{Z}$
par $W \to \mathcal{Z}$ est un morphisme plat
(Lemme \ref{stacks-morphisms-lemma-base-change-flat}), nous pouvons remplacer
$\mathcal{Z}$ par $W$.

\medskip\noindent
Choisissons un espace algébrique $V$ et un morphisme lisse surjectif
$V \to \mathcal{Y}$. Choisissons un espace algébrique $U$ et un morphisme
lisse surjectif $U \to V \times_\mathcal{Y} \mathcal{X}$.
Il faut montrer que $U \to V$ est plat. Effectuons maintenant partout le
changement de base par $W \to \mathcal{Y}$ : posons
$U' = W \times_\mathcal{Y} U$,
$V' = W \times_\mathcal{Y} V$,
$\mathcal{X}' = W \times_\mathcal{Y} \mathcal{X}$,
et $\mathcal{Y}' = W \times_\mathcal{Y} \mathcal{Y} = W$.
Par changement de base,
$U' \to V' \times_{\mathcal{Y}'} \mathcal{X}'$ est encore lisse.
Notre définition des morphismes plats de champs algébriques et l'hypothèse que
$\mathcal{X}' \to \mathcal{Y}'$ est plat montrent donc que
$U' \to V'$ est plat. Enfin, $V' \to V$ est surjectif comme changement de
base de $W \to \mathcal{Y}$ ; ainsi, $U \to V$ est plat par Morphismes
d'espaces, Lemme
\ref{spaces-morphisms-lemma-base-change-module-flat} (2)
et la démonstration est achevée.
\end{proof}

\begin{lemma}
\label{stacks-morphisms-lemma-flat-permanence}
Soient $\mathcal{X} \to \mathcal{Y} \to \mathcal{Z}$ des morphismes de
champs algébriques. Si $\mathcal{X} \to \mathcal{Z}$ est plat et si
$\mathcal{X} \to \mathcal{Y}$ est surjectif et plat, alors
$\mathcal{Y} \to \mathcal{Z}$ est plat.
\end{lemma}

\begin{proof}
Choisissons un espace algébrique $W$ et un morphisme lisse surjectif
$W \to \mathcal{Z}$. Choisissons un espace algébrique $V$ et un morphisme
lisse surjectif $V \to W \times_\mathcal{Z} \mathcal{Y}$. Choisissons un
espace algébrique $U$ et un morphisme lisse surjectif
$U \to V \times_\mathcal{Y} \mathcal{X}$.
Nous savons que $U \to V$ est plat et que $U \to W$ est plat.
De plus, puisque $\mathcal{X} \to \mathcal{Y}$ est surjectif,
$U \to V$ est surjectif (comme composé de morphismes surjectifs).
Le lemme se ramène donc au cas des morphismes d'espaces algébriques, qui est
Morphismes d'espaces, Lemme
\ref{spaces-morphisms-lemma-flat-permanence}.
\end{proof}

\begin{lemma}
\label{stacks-morphisms-lemma-lift-valuation-ring-through-flat-morphism}
Soit $f : \mathcal{X} \to \mathcal{Y}$ un morphisme plat de champs
algébriques. Soit $\Spec(A) \to \mathcal{Y}$ un morphisme, où $A$ est un
anneau de valuation. Si le point fermé de $\Spec(A)$ s'envoie sur un point de
$|\mathcal{Y}|$ appartenant à l'image de
$|\mathcal{X}| \to |\mathcal{Y}|$, alors il existe un diagramme commutatif
$$
\xymatrix{
\Spec(A') \ar[r] \ar[d] & \mathcal{X} \ar[d] \\
\Spec(A) \ar[r] & \mathcal{Y}
}
$$
où $A \to A'$ est une extension d'anneaux de valuation
(Compléments d'algèbre, Définition
\ref{more-algebra-definition-extension-valuation-rings}).
\end{lemma}

\begin{proof}
Le changement de base $\mathcal{X}_A \to \Spec(A)$ est plat
(Lemme \ref{stacks-morphisms-lemma-base-change-flat}) et le point fermé de
$\Spec(A)$ appartient à l'image de $|\mathcal{X}_A| \to |\Spec(A)|$
(Propriétés des champs,
Lemme \ref{stacks-properties-lemma-points-cartesian}).
Nous pouvons donc supposer que $\mathcal{Y} = \Spec(A)$. Soit
$U \to \mathcal{X}$ un morphisme lisse surjectif, où $U$ est un schéma.
On peut alors appliquer Morphismes d'espaces, Lemme
\ref{spaces-morphisms-lemma-lift-valuation-ring-through-flat-morphism}
au morphisme $U \to \Spec(A)$ pour conclure.
\end{proof}







\section{Plat en un point}
\label{stacks-morphisms-section-flat-at-point}

\noindent
Il nous reste à développer le formalisme général nécessaire pour exprimer
qu'un morphisme de champs algébriques possède une propriété donnée en un
point. Pour l'instant, le lemme suivant suffit.

\begin{lemma}
\label{stacks-morphisms-lemma-flat-at-point}
Soit $f : \mathcal{X} \to \mathcal{Y}$ un morphisme de champs algébriques.
Soit $x \in |\mathcal{X}|$. Considérons des diagrammes commutatifs
$$
\vcenter{
\xymatrix{
U \ar[d]_a \ar[r]_h & V \ar[d]^b \\
\mathcal{X} \ar[r]^f & \mathcal{Y}
}
}
\quad\text{avec des points}
\vcenter{
\xymatrix{
u \in |U| \ar[d] \\
x \in |\mathcal{X}|
}
}
$$
où $U$ et $V$ sont des espaces algébriques, $b$ est plat, et
$(a, h) : U \to \mathcal{X} \times_\mathcal{Y} V$
est plat. Les conditions suivantes sont équivalentes :
\begin{enumerate}
\item $h$ est plat en $u$ pour l'un des diagrammes ci-dessus ;
\item $h$ est plat en $u$ pour tout diagramme comme ci-dessus.
\end{enumerate}
\end{lemma}

\begin{proof}
Supposons donné un second diagramme $U', V', u', a', b', h'$ satisfaisant
aux conditions du lemme. Nous pouvons alors considérer
$$
\xymatrix{
U \ar[d] & U \times_\mathcal{X} U' \ar[l] \ar[d] \ar[r] & U' \ar[d] \\
V & V \times_\mathcal{Y} V' \ar[l] \ar[r] & V'
}
$$
D'après Propriétés des champs, Lemme
\ref{stacks-properties-lemma-points-cartesian}, il existe un point
$u'' \in |U \times_\mathcal{X} U'|$ qui s'envoie sur $u$ et $u'$.
Si $h$ est plat en $u$, alors le changement de base
$U \times_V (V \times_\mathcal{Y} V') \to V \times_\mathcal{Y} V'$
est plat en tout point au-dessus de $u$ ; voir Morphismes d'espaces, Lemme
\ref{spaces-morphisms-lemma-base-change-module-flat}. D'autre part, le
morphisme
$$
U \times_\mathcal{X} U' \to
U \times_\mathcal{X} (\mathcal{X} \times_\mathcal{Y} V') =
U \times_\mathcal{Y} V' =
U \times_V (V \times_\mathcal{Y} V')
$$
est plat, car c'est un changement de base de $(a', h')$ ; voir le
Lemme \ref{stacks-morphisms-lemma-base-change-flat}. En composant et en utilisant Morphismes
d'espaces, Lemme \ref{spaces-morphisms-lemma-composition-module-flat}, nous
concluons que $U \times_\mathcal{X} U' \to V \times_\mathcal{Y} V'$ est plat
en $u''$. En composant ensuite avec le morphisme plat
$V \times_\mathcal{Y} V' \to V'$, nous voyons que
$U \times_\mathcal{X} U' \to V'$ est plat en $u''$. Enfin, puisque
$U \times_\mathcal{X} U' \to U'$ est plat en $u''$ et que $u''$ s'envoie
sur $u'$, nous concluons que $U' \to V'$ est plat en $u'$ d'après Morphismes
d'espaces, Lemme \ref{spaces-morphisms-lemma-flat-permanence}.
\end{proof}

\begin{definition}
\label{stacks-morphisms-definition-flat-at-point}
Soit $f : \mathcal{X} \to \mathcal{Y}$ un morphisme de champs algébriques.
Soit $x \in |\mathcal{X}|$. Nous disons que $f$ est {\it plat en $x$} si les
conditions équivalentes du Lemme \ref{stacks-morphisms-lemma-flat-at-point} sont satisfaites.
\end{definition}





\section{Morphismes de présentation finie}
\label{stacks-morphisms-section-finite-presentation}

\noindent
La propriété ``être localement de présentation finie'' des morphismes
d'espaces algébriques est locale pour la topologie lisse sur la source et la
cible, voir Descente sur les espaces, Remarque
\ref{spaces-descent-remark-list-local-source-target}. Elle est aussi stable
par changement de base et locale pour la topologie fpqc sur la cible, voir
Morphismes d'espaces, Lemme
\ref{spaces-morphisms-lemma-base-change-finite-presentation}, et Descente sur
les espaces, Lemme
\ref{spaces-descent-lemma-descending-property-locally-finite-presentation}.
Par le Lemme \ref{stacks-morphisms-lemma-local-source-target} ci-dessus, nous pouvons donc
définir comme suit ce que signifie être localement de présentation finie pour
un morphisme de champs algébriques. Cette définition coïncide avec la notion
déjà définie dans Propriétés des champs,
section \ref{stacks-properties-section-properties-morphisms}, lorsque le
morphisme est représentable par des espaces algébriques.

\begin{definition}
\label{stacks-morphisms-definition-locally-finite-presentation}
Soit $f : \mathcal{X} \to \mathcal{Y}$ un morphisme de champs algébriques.
\begin{enumerate}
\item On dit que $f$ est {\it localement de présentation finie} si les
conditions équivalentes du Lemme \ref{stacks-morphisms-lemma-local-source-target} sont
satisfaites avec
$\mathcal{P} = \text{localement de présentation finie}$.
\item On dit que $f$ est {\it de présentation finie} s'il est localement de
présentation finie, quasi-compact et quasi-séparé.
\end{enumerate}
\end{definition}

\noindent
Notons qu'un morphisme de présentation finie n'est {\bf pas} simplement un
morphisme quasi-compact qui est localement de présentation finie.

\begin{lemma}
\label{stacks-morphisms-lemma-composition-finite-presentation}
Le composé de morphismes de présentation finie est de présentation finie.
Il en va de même pour les morphismes localement de présentation finie.
\end{lemma}

\begin{proof}
Il suffit de combiner la Remarque \ref{stacks-morphisms-remark-composition} avec Morphismes
d'espaces, Lemme
\ref{spaces-morphisms-lemma-composition-finite-presentation}.
\end{proof}

\begin{lemma}
\label{stacks-morphisms-lemma-base-change-finite-presentation}
Un changement de base d'un morphisme de présentation finie est de présentation
finie. Il en va de même pour les morphismes localement de présentation finie.
\end{lemma}

\begin{proof}
Il suffit de combiner la Remarque \ref{stacks-morphisms-remark-base-change} avec Morphismes
d'espaces, Lemme
\ref{spaces-morphisms-lemma-base-change-finite-presentation}.
\end{proof}

\begin{lemma}
\label{stacks-morphisms-lemma-finite-presentation-finite-type}
Un morphisme localement de présentation finie est localement de type fini.
Un morphisme de présentation finie est de type fini.
\end{lemma}

\begin{proof}
Il suffit de combiner la Remarque \ref{stacks-morphisms-remark-implication} avec Morphismes
d'espaces, Lemme
\ref{spaces-morphisms-lemma-finite-presentation-finite-type}.
\end{proof}

\begin{lemma}
\label{stacks-morphisms-lemma-noetherian-finite-type-finite-presentation}
Soit $f : \mathcal{X} \to \mathcal{Y}$ un morphisme de champs algébriques.
\begin{enumerate}
\item Si $\mathcal{Y}$ est localement noethérien et $f$ localement de type
fini, alors $f$ est localement de présentation finie.
\item Si $\mathcal{Y}$ est localement noethérien et si $f$ est de type fini
et quasi-séparé, alors $f$ est de présentation finie.
\end{enumerate}
\end{lemma}

\begin{proof}
Supposons $f : \mathcal{X} \to \mathcal{Y}$ localement de type fini et
$\mathcal{Y}$ localement noethérien. Il existe donc un diagramme comme dans le
Lemme \ref{stacks-morphisms-lemma-local-source-target}, avec $h$ localement de type fini et la
flèche verticale $a$ surjective. D'après Morphismes d'espaces, Lemme
\ref{spaces-morphisms-lemma-noetherian-finite-type-finite-presentation}
$h$ est localement de présentation finie. Ainsi,
$\mathcal{X} \to \mathcal{Y}$ est localement de présentation finie par
définition. Cela prouve (1). Si $f$ est de type fini et quasi-séparé, il est
aussi quasi-compact et quasi-séparé, et (2) en résulte immédiatement.
\end{proof}

\begin{lemma}
\label{stacks-morphisms-lemma-finite-presentation-permanence}
Soient $f : \mathcal{X} \to \mathcal{Y}$ et
$g : \mathcal{Y} \to \mathcal{Z}$ des morphismes de champs algébriques.
Si $g \circ f$ est localement de présentation finie et si $g$ est localement
de type fini, alors $f$ est localement de présentation finie.
\end{lemma}

\begin{proof}
Choisissons un espace algébrique $W$ et un morphisme lisse surjectif
$W \to \mathcal{Z}$. Choisissons un espace algébrique $V$ et un morphisme
lisse surjectif $V \to \mathcal{Y} \times_\mathcal{Z} W$. Choisissons un
espace algébrique $U$ et un morphisme lisse surjectif
$U \to \mathcal{X} \times_\mathcal{Y} V$. Le lemme résulte de l'application
de Morphismes d'espaces, Lemme
\ref{spaces-morphisms-lemma-finite-presentation-permanence}
aux morphismes $U \to V \to W$.
\end{proof}

\begin{lemma}
\label{stacks-morphisms-lemma-diagonal-morphism-finite-type}
Soit $f : \mathcal{X} \to \mathcal{Y}$ un morphisme de champs algébriques,
de diagonale
$\Delta : \mathcal{X} \to \mathcal{X} \times_\mathcal{Y} \mathcal{X}$.
Si $f$ est localement de type fini, alors $\Delta$ est localement de
présentation finie. Si $f$ est quasi-séparé et localement de type fini, alors
$\Delta$ est de présentation finie.
\end{lemma}

\begin{proof}
Notons que $\Delta$ est un morphisme au-dessus de $\mathcal{X}$ (par la
seconde projection). Si $f$ est localement de type fini, alors
$\mathcal{X}$ est de présentation finie sur $\mathcal{X}$ et
$\text{pr}_2 : \mathcal{X} \times_\mathcal{Y} \mathcal{X} \to \mathcal{X}$
est localement de type fini d'après le
Lemme \ref{stacks-morphisms-lemma-base-change-finite-type}. La première assertion résulte donc
du Lemme \ref{stacks-morphisms-lemma-finite-presentation-permanence}. La seconde résulte de la
première et des définitions (car, par définition, dire que $f$ est
quasi-séparé signifie que $\Delta_f$ est quasi-compact et quasi-séparé).
\end{proof}

\begin{lemma}
\label{stacks-morphisms-lemma-open-immersion-locally-finite-presentation}
Une immersion ouverte est localement de présentation finie.
\end{lemma}

\begin{proof}
Compte tenu de Propriétés des champs, Définition
\ref{stacks-properties-definition-immersion}, cela résulte de Morphismes
d'espaces, Lemme
\ref{spaces-morphisms-lemma-open-immersion-locally-finite-presentation}.
\end{proof}

\begin{lemma}
\label{stacks-morphisms-lemma-check-property-after-fppf-base-change}
Soit $P$ une propriété de morphismes d'espaces algébriques qui est locale pour
la topologie fppf sur la cible et stable par changement de base arbitraire.
Soit $f : \mathcal{X} \to \mathcal{Y}$ un morphisme de champs algébriques
représentable par des espaces algébriques. Soit
$\mathcal{Z} \to \mathcal{Y}$ un morphisme de champs algébriques surjectif,
plat et localement de présentation finie. Posons
$\mathcal{W} = \mathcal{Z} \times_\mathcal{Y} \mathcal{X}$. Alors
$$
(f\text{ possède }P) \Leftrightarrow
(\text{la projection }\mathcal{W} \to \mathcal{Z}\text{ possède }P).
$$
Pour la signification de cet énoncé, voir Propriétés des champs, section
\ref{stacks-properties-section-properties-morphisms}.
\end{lemma}

\begin{proof}
Choisissons un espace algébrique $W$ et un morphisme
$W \to \mathcal{Z}$ surjectif, plat et localement de présentation finie.
D'après Propriétés des champs, Lemme
\ref{stacks-properties-lemma-composition-surjective}
et les Lemmes \ref{stacks-morphisms-lemma-composition-flat} et
\ref{stacks-morphisms-lemma-composition-finite-presentation}, le composé
$W \to \mathcal{Y}$ est encore surjectif, plat et localement de présentation
finie. Posons
$V = W \times_\mathcal{Z} \mathcal{W} = W \times_\mathcal{Y} \mathcal{X}$.
D'après Propriétés des champs, Lemme
\ref{stacks-properties-lemma-check-property-covering}
$f$ possède $P$ si et seulement si $V \to W$ la possède, et
$\mathcal{W} \to \mathcal{Z}$ possède $P$ si et seulement si $V \to W$
la possède. Le lemme en résulte.
\end{proof}

\begin{lemma}
\label{stacks-morphisms-lemma-descent-property}
Soit $\mathcal{P}$ une propriété de morphismes d'espaces algébriques qui est
locale pour la topologie lisse sur la source et la cible, et locale pour la
topologie fppf sur la cible. Soit
$f : \mathcal{X} \to \mathcal{Y}$ un morphisme de champs algébriques. Soit
$\mathcal{Z} \to \mathcal{Y}$ un morphisme de champs algébriques surjectif,
plat et localement de présentation finie. Si le changement de base
$\mathcal{Z} \times_\mathcal{Y} \mathcal{X} \to \mathcal{Z}$
possède $\mathcal{P}$, alors $f$ possède $\mathcal{P}$.
\end{lemma}

\begin{proof}
Supposons que
$\mathcal{Z} \times_\mathcal{Y} \mathcal{X} \to \mathcal{Z}$ possède
$\mathcal{P}$. Choisissons un espace algébrique $W$ et un morphisme lisse
surjectif $W \to \mathcal{Z}$. Observons que
$W \times_\mathcal{Z} \mathcal{Z} \times_\mathcal{Y} \mathcal{X} =
W \times_\mathcal{Y} \mathcal{X}$. Par la définition même de ce que signifie,
pour $\mathcal{Z} \times_\mathcal{Y} \mathcal{X} \to \mathcal{Z}$, posséder
$\mathcal{P}$ (voir la
Définition \ref{stacks-morphisms-definition-P} et le Lemme \ref{stacks-morphisms-lemma-local-source-target}),
le morphisme $W \times_\mathcal{Y} \mathcal{X} \to W$ possède donc
$\mathcal{P}$. D'autre part, $W \to \mathcal{Z}$ est surjectif, plat et
localement de présentation finie (Morphismes d'espaces, Lemmes
\ref{spaces-morphisms-lemma-smooth-flat} et
\ref{spaces-morphisms-lemma-smooth-locally-finite-presentation}) ; par
conséquent, $W \to \mathcal{Y}$ est surjectif, plat et localement de
présentation finie (d'après Propriétés des champs, Lemme
\ref{stacks-properties-lemma-composition-surjective}
et les Lemmes \ref{stacks-morphisms-lemma-composition-flat} et
\ref{stacks-morphisms-lemma-composition-finite-presentation}). Nous pouvons donc remplacer
$\mathcal{Z}$ par $W$.

\medskip\noindent
Choisissons un espace algébrique $V$ et un morphisme lisse surjectif
$V \to \mathcal{Y}$. Choisissons un espace algébrique $U$ et un morphisme
lisse surjectif $U \to V \times_\mathcal{Y} \mathcal{X}$. Il s'agit de
montrer que $U \to V$ possède $\mathcal{P}$. Effectuons partout le changement
de base par $W \to \mathcal{Y}$ et posons
$U' = W \times_\mathcal{Y} U$,
$V' = W \times_\mathcal{Y} V$,
$\mathcal{X}' = W \times_\mathcal{Y} \mathcal{X}$,
et $\mathcal{Y}' = W \times_\mathcal{Y} \mathcal{Y} = W$. Le morphisme
$U' \to V' \times_{\mathcal{Y}'} \mathcal{X}'$ demeure lisse par changement
de base. Le Lemme \ref{stacks-morphisms-lemma-local-source-target}, utilisé dans la définition
de ce que signifie pour $\mathcal{X}' \to \mathcal{Y}' = W$ posséder
$\mathcal{P}$, montre donc que $U' \to V'$ possède
$\mathcal{P}$. Enfin, $V' \to V$ est surjectif, plat et localement de
présentation finie, car c'est un changement de base de
$W \to \mathcal{Y}$ ; il s'ensuit que $U \to V$ possède $\mathcal{P}$,
puisque $\mathcal{P}$ est locale pour la topologie fppf sur la cible.
\end{proof}

\begin{lemma}
\label{stacks-morphisms-lemma-descent-finite-presentation}
Soit $f : \mathcal{X} \to \mathcal{Y}$ un morphisme de champs algébriques.
Soit $\mathcal{Z} \to \mathcal{Y}$ un morphisme de champs algébriques
surjectif, plat et localement de présentation finie. Si le changement de base
$\mathcal{Z} \times_\mathcal{Y} \mathcal{X} \to \mathcal{Z}$
est localement de présentation finie, alors $f$ est localement de présentation
finie.
\end{lemma}

\begin{proof}
La propriété ``être localement de présentation finie'' satisfait aux
conditions du Lemme \ref{stacks-morphisms-lemma-descent-property}. Nous avons vu dans
l'introduction de cette section qu'elle est locale pour la topologie lisse sur
la source et la cible ; sa localité pour la topologie fppf sur la cible est
Descente sur les espaces, Lemme
\ref{spaces-descent-lemma-descending-property-locally-finite-presentation}.
\end{proof}

\begin{lemma}
\label{stacks-morphisms-lemma-flat-finite-presentation-permanence}
Soient $\mathcal{X} \to \mathcal{Y} \to \mathcal{Z}$ des morphismes de
champs algébriques. Si $\mathcal{X} \to \mathcal{Z}$ est localement de
présentation finie et si $\mathcal{X} \to \mathcal{Y}$ est surjectif, plat et
localement de présentation finie, alors $\mathcal{Y} \to \mathcal{Z}$ est
localement de présentation finie.
\end{lemma}

\begin{proof}
Choisissons un espace algébrique $W$ et un morphisme lisse surjectif
$W \to \mathcal{Z}$. Choisissons un espace algébrique $V$ et un morphisme
lisse surjectif $V \to W \times_\mathcal{Z} \mathcal{Y}$. Choisissons un
espace algébrique $U$ et un morphisme lisse surjectif
$U \to V \times_\mathcal{Y} \mathcal{X}$. Nous savons que $U \to V$ est plat
et localement de présentation finie, et que $U \to W$ est localement de
présentation finie. De plus, puisque $\mathcal{X} \to \mathcal{Y}$ est
surjectif, $U \to V$ est surjectif (comme composé de morphismes surjectifs).
Le lemme se ramène donc au cas des morphismes d'espaces algébriques, lequel est
Descente sur les espaces, Lemme
\ref{spaces-descent-lemma-locally-finite-presentation-fppf-local-source}.
\end{proof}

\begin{lemma}
\label{stacks-morphisms-lemma-surjective-flat-locally-finite-presentation}
Soit $f : \mathcal{X} \to \mathcal{Y}$ un morphisme de champs algébriques
surjectif, plat et localement de présentation finie. Pour tout schéma $U$ et
tout objet $y$ de $\mathcal{Y}$ au-dessus de $U$, il existe alors un
recouvrement fppf $\{U_i \to U\}$ et des objets $x_i$ de $\mathcal{X}$
au-dessus de $U_i$ tels que $f(x_i) \cong y|_{U_i}$ dans
$\mathcal{Y}_{U_i}$.
\end{lemma}

\begin{proof}
Nous pouvons considérer $y$ comme un morphisme $U \to \mathcal{Y}$. D'après
Propriétés des champs, Lemme
\ref{stacks-properties-lemma-base-change-surjective}
et les Lemmes \ref{stacks-morphisms-lemma-base-change-finite-presentation} et
\ref{stacks-morphisms-lemma-base-change-flat}, le morphisme
$\mathcal{X} \times_\mathcal{Y} U \to U$ est surjectif, plat et localement de
présentation finie. Soient $V$ un schéma et
$V \to \mathcal{X} \times_\mathcal{Y} U$ un morphisme lisse surjectif.
Le morphisme $V \to \mathcal{X} \times_\mathcal{Y} U$ est aussi surjectif,
plat et localement de présentation finie (voir
Morphismes d'espaces, Lemmes
\ref{spaces-morphisms-lemma-smooth-flat} et
\ref{spaces-morphisms-lemma-smooth-locally-finite-presentation}). Ainsi,
$V \to U$ est lui aussi surjectif, plat et localement de présentation finie ;
voir Propriétés des champs, Lemme
\ref{stacks-properties-lemma-composition-surjective}
et les Lemmes \ref{stacks-morphisms-lemma-composition-finite-presentation} et
\ref{stacks-morphisms-lemma-composition-flat}. Par conséquent, $\{V \to U\}$ est le
recouvrement fppf recherché et $x : V \to \mathcal{X}$ est l'objet recherché.
\end{proof}

\begin{lemma}
\label{stacks-morphisms-lemma-surjective-family-flat-locally-finite-presentation}
Soit $f_j : \mathcal{X}_j \to \mathcal{X}$, $j \in J$, une famille de
morphismes de champs algébriques, plats et localement de présentation finie,
qui sont conjointement surjectifs, c'est-à-dire
$|\mathcal{X}| = \bigcup |f_j|(|\mathcal{X}_j|)$.
Alors, pour tout schéma $U$ et tout objet $x$ de $\mathcal{X}$ au-dessus de
$U$, il existe un recouvrement fppf $\{U_i \to U\}_{i \in I}$, une application
$a : I \to J$ et des objets $x_i$ de $\mathcal{X}_{a(i)}$ au-dessus de $U_i$
tels que $f_{a(i)}(x_i) \cong x|_{U_i}$ dans $\mathcal{X}_{U_i}$.
\end{lemma}

\begin{proof}
Appliquons le Lemme
\ref{stacks-morphisms-lemma-surjective-flat-locally-finite-presentation} au morphisme
$\coprod_{j \in J} \mathcal{X}_j \to \mathcal{X}$. (Il subsiste ici un léger
problème ensembliste, dû à notre formalisme, que nous ignorons.) Enfin, un
morphisme $x_i : U_i \to \coprod_{j \in J} \mathcal{X}_j$ équivaut à la
donnée d'une décomposition en union disjointe
$U_i = \coprod U_{i, j}$ et de morphismes
$U_{i, j} \to \mathcal{X}_j$. Le recouvrement fppf
$\{U_{i, j} \to U\}$ et les morphismes
$U_{i, j} \to \mathcal{X}_j$ conviennent donc.
\end{proof}

\begin{lemma}
\label{stacks-morphisms-lemma-fppf-open}
Soit $f : \mathcal{X} \to \mathcal{Y}$ plat et localement de présentation
finie. Alors $|f| : |\mathcal{X}| \to |\mathcal{Y}|$ est ouverte.
\end{lemma}

\begin{proof}
Choisissons un schéma $V$ et un morphisme lisse surjectif
$V \to \mathcal{Y}$. Choisissons un schéma $U$ et un morphisme lisse
surjectif $U \to V \times_\mathcal{Y} \mathcal{X}$. Par hypothèse, le
morphisme de schémas $U \to V$ est plat et localement de présentation finie.
Ainsi, $U \to V$ est ouvert d'après Morphismes, Lemme
\ref{morphisms-lemma-fppf-open}. Par construction de la topologie sur
$|\mathcal{Y}|$, l'application $|V| \to |\mathcal{Y}|$ est ouverte.
L'application $|U| \to |\mathcal{X}|$ est surjective. Le résultat découle
alors de ces faits par topologie élémentaire.
\end{proof}

\begin{lemma}
\label{stacks-morphisms-lemma-descent-quasi-compact}
Soit $f : \mathcal{X} \to \mathcal{Y}$ un morphisme de champs algébriques.
Soit $\mathcal{Z} \to \mathcal{Y}$ un morphisme de champs algébriques
surjectif, plat et localement de présentation finie. Si le changement de base
$\mathcal{Z} \times_\mathcal{Y} \mathcal{X} \to \mathcal{Z}$
est quasi-compact, alors $f$ est quasi-compact.
\end{lemma}

\begin{proof}
Il faut montrer que, pour tout $\mathcal{Y}' \to \mathcal{Y}$ avec
$\mathcal{Y}'$ quasi-compact, le champ
$\mathcal{Y}' \times_\mathcal{Y} \mathcal{X}$ est quasi-compact. Posons
$\mathcal{Z}' = \mathcal{Z} \times_\mathcal{Y} \mathcal{Y}'$. Alors
$|\mathcal{Z}'| \to |\mathcal{Y}'|$ est ouverte ; voir le
Lemme \ref{stacks-morphisms-lemma-fppf-open}. Il existe donc un sous-champ ouvert quasi-compact
$\mathcal{W} \subset \mathcal{Z}'$ dont l'image contient
$\mathcal{Y}'$. Puisque
$\mathcal{Z} \times_\mathcal{Y} \mathcal{X} \to \mathcal{Z}$
est quasi-compact, nous savons que
$$
\mathcal{W} \times_\mathcal{Z} \mathcal{Z} \times_\mathcal{Y} \mathcal{X} =
\mathcal{W} \times_\mathcal{Y} \mathcal{X}
$$
est quasi-compact. Or l'application
$\mathcal{W} \times_\mathcal{Y} \mathcal{X} \to
\mathcal{Y}' \times_\mathcal{Y} \mathcal{X}$
est surjective, ce qui conclut. Quelques détails sont omis.
\end{proof}

\begin{lemma}
\label{stacks-morphisms-lemma-check-separated-on-ui-cover}
Soient $f : \mathcal{X} \to \mathcal{Y}$ et
$g : \mathcal{Y} \to \mathcal{Z}$ des morphismes composables de champs
algébriques, de composé
$h = g \circ f : \mathcal{X} \to \mathcal{Z}$.
Si $f$ est surjectif, plat, localement de présentation finie et
universellement injectif, et si $h$ est séparé, alors $g$ est séparé.
\end{lemma}

\begin{proof}
Considérons le diagramme
$$
\xymatrix{
\mathcal{X} \ar[r]_\Delta \ar[rd] &
\mathcal{X} \times_\mathcal{Y} \mathcal{X} \ar[r] \ar[d] &
\mathcal{X} \times_\mathcal{Z} \mathcal{X} \ar[d] \\
& \mathcal{Y} \ar[r] & \mathcal{Y} \times_\mathcal{Z} \mathcal{Y}
}
$$
Le carré est cartésien. Il faut montrer que la flèche horizontale inférieure
est propre. Nous savons déjà qu'elle est représentable par des espaces
algébriques et localement de type fini
(Lemme \ref{stacks-morphisms-lemma-properties-diagonal}). Puisque la flèche verticale de droite
est surjective, plate et localement de présentation finie, il suffit de
montrer que la flèche horizontale supérieure de droite est propre
(Lemme \ref{stacks-morphisms-lemma-check-property-after-fppf-base-change}). Comme $h$ est
séparé, le composé des flèches horizontales supérieures est propre.

\medskip\noindent
Puisque $f$ est universellement injectif, $\Delta$ est surjective
(Lemme \ref{stacks-morphisms-lemma-universally-injective}). Comme le composé de $\Delta$ avec
la projection
$\mathcal{X} \times_\mathcal{Y} \mathcal{X} \to \mathcal{X}$
est l'identité, $\Delta$ est universellement fermé. D'après Morphismes
d'espaces, Lemme
\ref{spaces-morphisms-lemma-image-universally-closed-separated}
nous concluons que
$\mathcal{X} \times_\mathcal{Y} \mathcal{X} \to
\mathcal{X} \times_\mathcal{Z} \mathcal{X}$ est séparé, puisque
$\mathcal{X} \to \mathcal{X} \times_\mathcal{Z} \mathcal{X}$ est séparé.
Nous utilisons ici le fait que les implications entre propriétés de
morphismes d'espaces algébriques se transportent aux mêmes implications entre
propriétés de morphismes de champs algébriques représentables par des espaces
algébriques ; ceci est expliqué dans Propriétés des champs, section
\ref{stacks-properties-section-properties-morphisms}.
Enfin, le même principe et Morphismes d'espaces, Lemme
\ref{spaces-morphisms-lemma-image-proper-is-proper}, permettent de conclure que
$\mathcal{X} \times_\mathcal{Y} \mathcal{X} \to
\mathcal{X} \times_\mathcal{Z} \mathcal{X}$ est propre.
\end{proof}














\section{Gerbes}
\label{stacks-morphisms-section-gerbes}

\noindent
Une classe importante de champs algébriques est formée des champs de la forme
$[B/G]$, où $B$ est un espace algébrique et $G$ un espace algébrique en groupes
sur $B$, plat et localement de présentation finie (agissant trivialement sur
$B$) ; voir Critères de représentabilité, Lemme
\ref{criteria-lemma-BG-algebraic}. Un champ algébrique est une gerbe lorsqu'il
est localement de cette forme pour la topologie fppf ; voir le
Lemme \ref{stacks-morphisms-lemma-local-structure-gerbe}. Dans cette section, nous examinons
brièvement cette notion et la notion relative correspondante.

\begin{definition}
\label{stacks-morphisms-definition-gerbe}
Soit $f : \mathcal{X} \to \mathcal{Y}$ un morphisme de champs algébriques.
On dit que $\mathcal{X}$ est une {\it gerbe sur} $\mathcal{Y}$ si
$\mathcal{X}$ est une gerbe sur $\mathcal{Y}$ comme champs en
groupoïdes sur $(\Sch/S)_{fppf}$ ; voir Champs, Définition
\ref{stacks-definition-gerbe-over-stack-in-groupoids}. On dit qu'un champ
algébrique $\mathcal{X}$ est une {\it gerbe} s'il existe un morphisme
$\mathcal{X} \to X$, où $X$ est un espace algébrique, qui fasse de
$\mathcal{X}$ une gerbe sur $X$.
\end{definition}

\noindent
La condition que $\mathcal{X}$ soit une gerbe sur $\mathcal{Y}$ se
définit uniquement à partir de la topologie et de la théorie des catégories
sous-jacentes aux champs algébriques considérés ; mais, comme nous le verrons,
elle a des conséquences géométriques. Elle implique par exemple que
$\mathcal{X} \to \mathcal{Y}$ est surjectif, plat et localement de présentation
finie ; voir le Lemme \ref{stacks-morphisms-lemma-gerbe-fppf}. La notion absolue est plus
délicate à interpréter, car il n'est pas évident à première vue que $X$ soit
bien déterminé. Il l'est néanmoins.

\begin{lemma}
\label{stacks-morphisms-lemma-gerbe-over-iso-classes}
Soit $\mathcal{X}$ un champ algébrique. Si $\mathcal{X}$ est une gerbe, alors
le faisceau associé au préfaisceau
$$
(\Sch/S)_{fppf}^{opp} \to \textit{Ens}, \quad
U \mapsto \Ob(\mathcal{X}_U)/\!\!\cong
$$
est un espace algébrique, et $\mathcal{X}$ est une gerbe sur celui-ci.
\end{lemma}

\begin{proof}
(Dans cette démonstration, l'abus de langage introduit à la
section \ref{stacks-morphisms-section-conventions} se révèle particulièrement utile.)
Choisissons un morphisme $\pi : \mathcal{X} \to X$, où $X$ est un espace
algébrique, qui fasse de $\mathcal{X}$ une gerbe sur $X$. Il suffit de
montrer que $X$ est le faisceau associé au préfaisceau $\mathcal{F}$ affiché
dans le lemme. Il existe manifestement une application
$c : \mathcal{F} \to X$. Nous utiliserons les propriétés (2)(a) et (2)(b) de
Champs, Lemme \ref{stacks-lemma-when-gerbe}, pour montrer que l'application
$c^\# : \mathcal{F}^\# \to X$ est surjective et injective, donc un
isomorphisme ; voir Sites, Lemme \ref{sites-lemma-mono-epi-sheaves}.

Surjectivité : soit $T$ un schéma et $f : T \to X$. D'après la propriété
(2)(a), il existe un recouvrement fppf $\{h_i : T_i \to T\}$ et des morphismes
$x_i : T_i \to \mathcal{X}$ tels que $f \circ h_i$ corresponde à
$\pi \circ x_i$. Ainsi, $f|_{T_i}$ appartient à l'image de $c$.

Injectivité : soient $T$ un schéma et $x, y : T \to \mathcal{X}$ des
morphismes tels que $c \circ x = c \circ y$. D'après (2)(b), il existe un
recouvrement $\{T_i \to T\}$ et des morphismes
$x|_{T_i} \to y|_{T_i}$ dans la catégorie fibre $\mathcal{X}_{T_i}$.
Les restrictions $x|_{T_i}, y|_{T_i}$ sont donc égales dans
$\mathcal{F}(T_i)$. Cela prouve, comme voulu, que $x, y$ définissent la même
section de $\mathcal{F}^\#$ au-dessus de $T$.
\end{proof}

\begin{lemma}
\label{stacks-morphisms-lemma-base-change-gerbe}
Soit
$$
\xymatrix{
\mathcal{X}' \ar[r] \ar[d] & \mathcal{X} \ar[d] \\
\mathcal{Y}' \ar[r] & \mathcal{Y}
}
$$
un produit fibré de champs algébriques. Si $\mathcal{X}$ est une gerbe
sur $\mathcal{Y}$, alors $\mathcal{X}'$ est une gerbe sur
$\mathcal{Y}'$.
\end{lemma}

\begin{proof}
Cela résulte immédiatement des définitions et de Champs, Lemme
\ref{stacks-lemma-base-change-gerbe}.
\end{proof}

\begin{lemma}
\label{stacks-morphisms-lemma-composition-gerbe}
Soient $\mathcal{X} \to \mathcal{Y}$ et
$\mathcal{Y} \to \mathcal{Z}$ des morphismes de champs algébriques. Si
$\mathcal{X}$ est une gerbe sur $\mathcal{Y}$ et $\mathcal{Y}$ une
gerbe sur $\mathcal{Z}$, alors $\mathcal{X}$ est une gerbe sur
de $\mathcal{Z}$.
\end{lemma}

\begin{proof}
Cela résulte immédiatement de Champs, Lemme
\ref{stacks-lemma-composition-gerbe}.
\end{proof}

\begin{lemma}
\label{stacks-morphisms-lemma-gerbe-descent}
Soit
$$
\xymatrix{
\mathcal{X}' \ar[r] \ar[d] & \mathcal{X} \ar[d] \\
\mathcal{Y}' \ar[r] & \mathcal{Y}
}
$$
un produit fibré de champs algébriques. Si
$\mathcal{Y}' \to \mathcal{Y}$ est surjectif, plat et localement de
présentation finie, et si $\mathcal{X}'$ est une gerbe sur
$\mathcal{Y}'$, alors $\mathcal{X}$ est une gerbe sur $\mathcal{Y}$.
\end{lemma}

\begin{proof}
Cela résulte immédiatement du
Lemme \ref{stacks-morphisms-lemma-surjective-flat-locally-finite-presentation} et de Champs,
Lemme \ref{stacks-lemma-gerbe-descent}.
\end{proof}

\begin{lemma}
\label{stacks-morphisms-lemma-gerbe-with-section}
Soit $\pi : \mathcal{X} \to U$ un morphisme d'un champ algébrique vers un
espace algébrique, et soit $x : U \to \mathcal{X}$ une section de $\pi$.
Posons $G = \mathit{Isom}_\mathcal{X}(x, x)$ ; voir la
Définition \ref{stacks-morphisms-definition-isom}. Si $\mathcal{X}$ est une gerbe sur
$U$, alors :
\begin{enumerate}
\item il existe une équivalence canonique de champs en groupoïdes
$$
x_{can} : [U/G] \longrightarrow \mathcal{X}.
$$
où $[U/G]$ est le champ quotient pour l'action triviale de $G$ sur $U$ ;
\item $G \to U$ est plat et localement de présentation finie ;
\item $U \to \mathcal{X}$ est surjectif, plat et localement de présentation
finie.
\end{enumerate}
\end{lemma}

\begin{proof}
Posons $R = U \times_{x, \mathcal{X}, x} U$. Le morphisme
$R \to U \times U$ se factorise par la diagonale
$\Delta_U : U \to U \times U$, puisqu'il se factorise par
$U \times_U U = U$. Ainsi $R = G$, car
\begin{align*}
G & = \mathit{Isom}_\mathcal{X}(x, x) \\
& = U \times_{x, \mathcal{X}} \mathcal{I}_\mathcal{X} \\
& = U \times_{x, \mathcal{X}}
(\mathcal{X}
\times_{\Delta, \mathcal{X} \times_S \mathcal{X}, \Delta}
\mathcal{X}) \\
& = (U \times_{x, \mathcal{X}, x} U) \times_{U \times U, \Delta_U} U \\
& = R \times_{U \times U, \Delta_U} U \\
& = R
\end{align*}
où, pour la quatrième égalité, on utilise Catégories, Lemme
\ref{categories-lemma-diagonal-2}.
Soient $t, s : R \to U$ les projections. La loi de composition
$c : R \times_{s, U, t} R \to R$ construite sur $R$ dans Champs algébriques,
Lemme \ref{algebraic-lemma-map-space-into-stack}, coïncide avec la loi de
groupe sur $G$ (démonstration omise). Ainsi, Champs algébriques, Lemme
\ref{algebraic-lemma-map-space-into-stack}, fournit un $1$-morphisme canonique
pleinement fidèle
$$
x_{can} : [U/G] \longrightarrow \mathcal{X}
$$
de champs en groupoïdes sur $(\Sch/S)_{fppf}$. Pour voir que c'est une
équivalence, il suffit de montrer qu'il est essentiellement surjectif. Il
suffit pour cela de montrer que tout objet de $\mathcal{X}$ au-dessus d'un
schéma $T$ provient, localement pour la topologie fppf, de $x$ par un morphisme
$T \to U$ ; voir Champs, Lemme
\ref{stacks-lemma-characterize-essentially-surjective-when-ff}. Or cela résulte
de la condition selon laquelle $\pi$ fait de $\mathcal{X}$ une gerbe sur
$U$ ; voir la propriété (2)(a) de Champs, Lemme
\ref{stacks-lemma-when-gerbe}.

\medskip\noindent
D'après Critères de représentabilité, Lemme
\ref{criteria-lemma-BG-algebraic}, le morphisme $G \to U$ est plat et
localement de présentation finie. Enfin, $U \to \mathcal{X}$ est surjectif,
plat et localement de présentation finie d'après Critères de représentabilité,
Lemme
\ref{criteria-lemma-flat-quotient-flat-presentation}.
\end{proof}

\begin{lemma}
\label{stacks-morphisms-lemma-local-structure-gerbe}
Soit $\pi : \mathcal{X} \to \mathcal{Y}$ un morphisme de champs algébriques.
Les conditions suivantes sont équivalentes :
\begin{enumerate}
\item $\mathcal{X}$ est une gerbe sur $\mathcal{Y}$ ;
\item il existe un espace algébrique $U$, un espace algébrique en groupes $G$
sur $U$, plat et localement de présentation finie, et un morphisme
$U \to \mathcal{Y}$ surjectif, plat et localement de présentation finie, tels
que $\mathcal{X} \times_\mathcal{Y} U \cong [U/G]$ au-dessus de $U$.
\end{enumerate}
\end{lemma}

\begin{proof}
Supposons (2). D'après le Lemme \ref{stacks-morphisms-lemma-gerbe-descent}, il suffit, pour
prouver (1), de montrer que $[U/G]$ est une gerbe sur $U$. Cela
résulte immédiatement de Groupoïdes en espaces, Lemme
\ref{spaces-groupoids-lemma-group-quotient-gerbe}.

\medskip\noindent
Supposons (1). Tout changement de base de $\pi$ est une gerbe ; voir le
Lemme \ref{stacks-morphisms-lemma-base-change-gerbe}. Choisissons d'abord un schéma $V$ et un
morphisme lisse surjectif $V \to \mathcal{Y}$. Nous pouvons ainsi supposer que
$\pi : \mathcal{X} \to V$ est une gerbe sur un schéma. Il existe alors
un recouvrement fppf $\{V_i \to V\}$ tel que la catégorie fibre
$\mathcal{X}_{V_i}$ soit non vide ; voir Champs, Lemme
\ref{stacks-lemma-when-gerbe} (2)(a). Notons que
$U = \coprod V_i \to V$ est surjectif, plat et localement de présentation
finie. Nous pouvons donc remplacer $V$ par $U$ et supposer que
$\pi : \mathcal{X} \to U$ est une gerbe sur le schéma $U$, et qu'il
existe un objet $x$ de $\mathcal{X}$ au-dessus de $U$. D'après le
Lemme \ref{stacks-morphisms-lemma-gerbe-with-section}, on a alors
$\mathcal{X} = [U/G]$ au-dessus de $U$, pour un certain espace algébrique en
groupes $G$ sur $U$, plat et localement de présentation finie.
\end{proof}

\begin{lemma}
\label{stacks-morphisms-lemma-gerbe-fppf}
Soit $\pi : \mathcal{X} \to \mathcal{Y}$ un morphisme de champs algébriques.
Si $\mathcal{X}$ est une gerbe sur $\mathcal{Y}$, alors $\pi$ est
surjectif, plat et localement de présentation finie.
\end{lemma}

\begin{proof}
D'après Propriétés des champs, Lemme
\ref{stacks-properties-lemma-descent-surjective}
et les Lemmes \ref{stacks-morphisms-lemma-descent-flat} et
\ref{stacks-morphisms-lemma-descent-finite-presentation}, il suffit de démontrer le lemme après
avoir remplacé $\pi$ par son changement de base par un morphisme
$\mathcal{Y}' \to \mathcal{Y}$ surjectif, plat et localement de présentation
finie. D'après le Lemme \ref{stacks-morphisms-lemma-local-structure-gerbe}, nous pouvons
supposer que $\mathcal{Y} = U$ est un espace algébrique et que
$\mathcal{X} = [U/G]$ au-dessus de $U$. Le morphisme $U \to [U/G]$ est alors
surjectif, plat et localement de présentation finie ; voir le
Lemme \ref{stacks-morphisms-lemma-gerbe-with-section}. Il en résulte que $\pi$ est surjectif,
plat et localement de présentation finie d'après Propriétés des champs, Lemme
\ref{stacks-properties-lemma-surjective-permanence}, et les
Lemmes \ref{stacks-morphisms-lemma-flat-permanence} et
\ref{stacks-morphisms-lemma-flat-finite-presentation-permanence}.
\end{proof}

\begin{proposition}
\label{stacks-morphisms-proposition-when-gerbe}
Soit $\mathcal{X}$ un champ algébrique. Les conditions suivantes sont
équivalentes :
\begin{enumerate}
\item $\mathcal{X}$ est une gerbe ;
\item $\mathcal{I}_\mathcal{X} \to \mathcal{X}$ est plat et localement de
présentation finie.
\end{enumerate}
\end{proposition}

\begin{proof}
Supposons (1). Choisissons un morphisme $\mathcal{X} \to X$ vers un espace
algébrique $X$ qui fasse de $\mathcal{X}$ une gerbe sur $X$. Soit
$X' \to X$ un morphisme surjectif, plat et localement de présentation finie,
et posons $\mathcal{X}' = X' \times_X \mathcal{X}$. Le champ $\mathcal{X}'$
est une gerbe sur $X'$ d'après le
Lemme \ref{stacks-morphisms-lemma-base-change-gerbe}. Les deux carrés du diagramme
$$
\xymatrix{
\mathcal{I}_{\mathcal{X}'} \ar[r] \ar[d] &
\mathcal{X}' \ar[r] \ar[d] & X' \ar[d] \\
\mathcal{I}_\mathcal{X} \ar[r] &
\mathcal{X} \ar[r] & X
}
$$
sont cartésiens ; voir le Lemme \ref{stacks-morphisms-lemma-cartesian-square-inertia}. Pour
montrer que $\mathcal{I}_\mathcal{X} \to \mathcal{X}$ est plat et localement
de présentation finie, il suffit donc de le vérifier après un tel changement
de base, d'après les Lemmes \ref{stacks-morphisms-lemma-descent-flat} et
\ref{stacks-morphisms-lemma-descent-finite-presentation}. Le
Lemme \ref{stacks-morphisms-lemma-local-structure-gerbe} permet alors de supposer que
$\mathcal{X} = [U/G]$. D'après le Lemme \ref{stacks-morphisms-lemma-gerbe-with-section}, $G$
est plat et localement de présentation finie sur $U$, et
$x : U \to [U/G]$ est surjectif, plat et localement de présentation finie.
De plus, le changement de base de $\mathcal{I}_\mathcal{X}$ par $x$ est $G$ ;
une nouvelle descente, c'est-à-dire les Lemmes \ref{stacks-morphisms-lemma-descent-flat} et
\ref{stacks-morphisms-lemma-descent-finite-presentation}, montre donc (2).

\medskip\noindent
Réciproquement, supposons (2). Choisissons une présentation lisse
$\mathcal{X} = [U/R]$ ; voir Champs algébriques, section
\ref{algebraic-section-stack-to-presentation}. Notons $G \to U$ l'espace
algébrique en groupes stabilisateur du groupoïde
$(U, R, s, t, c, e, i)$ ; voir Groupoïdes en espaces, Définition
\ref{spaces-groupoids-definition-stabilizer-groupoid}.
D'après le Lemme \ref{stacks-morphisms-lemma-presentation-inertia}, $G \to U$ est plat et
localement de présentation finie, comme changement de base de
$\mathcal{I}_\mathcal{X} \to \mathcal{X}$ ; voir les
Lemmes \ref{stacks-morphisms-lemma-base-change-flat} et
\ref{stacks-morphisms-lemma-base-change-finite-presentation}. Considérons l'action
$$
a : G \times_{U, t} R \to R, \quad (g, r) \mapsto c(g, r)
$$
de $G$ sur $R$. Pour tout schéma $T$, cette action est libre sur les points à
valeurs dans $T$, puisque $R$ est un groupoïde. Ainsi $R' = R/G$ est un espace
algébrique, et le morphisme quotient $\pi : R \to R'$ est surjectif, plat et
localement de présentation finie d'après Amorçage, Lemme
\ref{bootstrap-lemma-quotient-free-action}. Les projections
$s, t : R \to U$ sont $G$-invariantes ; il existe donc des morphismes
$s' , t' : R' \to U$ tels que $s = s' \circ \pi$ et
$t = t' \circ \pi$. Puisque $s, t : R \to U$ sont plats et localement de
présentation finie, $s', t'$ le sont aussi ; voir Morphismes d'espaces, Lemme
\ref{spaces-morphisms-lemma-flat-permanence}, et Descente sur les espaces,
Lemme
\ref{spaces-descent-lemma-locally-finite-presentation-fppf-local-source}.
Considérons le morphisme
$$
j' = (t', s') : R' \longrightarrow U \times U.
$$
Nous affirmons que c'est un monomorphisme. Soient en effet $T$ un schéma et
$a, b : T \to R'$ des morphismes ayant la même image dans $U \times U$. Par
définition du quotient $R' = R/G$, il existe un recouvrement fppf
$\{h_j : T_j \to T\}$ tel que
$a \circ h_j = \pi \circ a_j$ et $b \circ h_j = \pi \circ b_j$, pour des
morphismes $a_j, b_j : T_j \to R$. Puisque $a_j, b_j$ ont la même image dans
$U \times U$, l'élément $g_j = c(a_j, i(b_j))$ est un point à valeurs dans
$T_j$ de $G$ tel que $c(g_j, b_j) = a_j$. Autrement dit, $a_j$ et $b_j$ ont la
même image dans $R'$, ce qui prouve l'affirmation. Comme
$j : R \to U \times U$ est une prérelation d'équivalence (voir Groupoïdes en
espaces, Lemme
\ref{spaces-groupoids-lemma-groupoid-pre-equivalence})
et que $R \to R'$ est surjectif (comme morphisme de faisceaux), le morphisme
$j' : R' \to U \times U$ est une relation d'équivalence. Par conséquent,
Amorçage, Théorème \ref{bootstrap-theorem-final-bootstrap}, montre que
$X = U/R'$ est un espace algébrique. Enfin, nous affirmons que le morphisme
$$
\mathcal{X} = [U/R] \longrightarrow X = U/R'
$$
fait de $\mathcal{X}$ une gerbe sur $X$. Cela résulte de Groupoïdes
en espaces, Lemme \ref{spaces-groupoids-lemma-when-gerbe}, puisque
$R \to R'$ est surjectif, plat et localement de présentation finie (au besoin,
utiliser Amorçage, Lemme
\ref{bootstrap-lemma-surjective-flat-locally-finite-presentation}, pour voir
que cela entraîne l'hypothèse requise).
\end{proof}

\begin{lemma}
\label{stacks-morphisms-lemma-gerbe-isom-fppf}
Soit $f : \mathcal{X} \to \mathcal{Y}$ un morphisme de champs algébriques
qui fait de $\mathcal{X}$ une gerbe sur $\mathcal{Y}$. Alors :
\begin{enumerate}
\item $\mathcal{I}_{\mathcal{X}/\mathcal{Y}} \to \mathcal{X}$
est plat et localement de présentation finie ;
\item $\mathcal{X} \to \mathcal{X} \times_\mathcal{Y} \mathcal{X}$
est surjectif, plat et localement de présentation finie ;
\item étant donnés des espaces algébriques $T_i$, $i = 1, 2$, et des
morphismes $x_i : T_i \to \mathcal{X}$, avec $y_i = f \circ x_i$, le
morphisme
$$
T_1 \times_{x_1, \mathcal{X}, x_2} T_2 \longrightarrow
T_1 \times_{y_1, \mathcal{Y}, y_2} T_2
$$
est surjectif, plat et localement de présentation finie ;
\item étant donnés un espace algébrique $T$ et des morphismes
$x_i : T \to \mathcal{X}$, $i = 1, 2$, avec $y_i = f \circ x_i$, le morphisme
$$
\mathit{Isom}_\mathcal{X}(x_1, x_2) \longrightarrow
\mathit{Isom}_\mathcal{Y}(y_1, y_2)
$$
est surjectif, plat et localement de présentation finie.
\end{enumerate}
\end{lemma}

\begin{proof}
Démonstration de (1). Choisissons un schéma $Y$ et un morphisme lisse
surjectif $Y \to \mathcal{Y}$. Posons
$\mathcal{X}' = \mathcal{X} \times_\mathcal{Y} Y$. Le
Lemme \ref{stacks-morphisms-lemma-cartesian-square-inertia} fournit les carrés cartésiens
$$
\xymatrix{
\mathcal{I}_{\mathcal{X}'} \ar[r] \ar[d] &
\mathcal{X}' \ar[r] \ar[d] & Y \ar[d] \\
\mathcal{I}_{\mathcal{X}/\mathcal{Y}} \ar[r] &
\mathcal{X} \ar[r] & \mathcal{Y}
}
$$
D'après les Lemmes \ref{stacks-morphisms-lemma-descent-flat} et
\ref{stacks-morphisms-lemma-descent-finite-presentation}, il suffit de montrer que
$\mathcal{I}_{\mathcal{X}'} \to \mathcal{X}'$ est plat et localement de
présentation finie. Cela résulte de la
Proposition \ref{stacks-morphisms-proposition-when-gerbe} (car $\mathcal{X}'$ est une gerbe
au-dessus de $Y$ d'après le Lemme \ref{stacks-morphisms-lemma-base-change-gerbe}).

\medskip\noindent
Démonstration de (2). Avec les notations ci-dessus, nous pouvons supposer que
$\mathcal{X}' = [Y/G]$ pour un espace algébrique en groupes $G$ sur $Y$, plat
et localement de présentation finie ; voir le
Lemme \ref{stacks-morphisms-lemma-local-structure-gerbe}. Le changement de base du morphisme
$\Delta : \mathcal{X} \to \mathcal{X} \times_\mathcal{Y} \mathcal{X}$
au-dessus de $\mathcal{Y}$ par $Y \to \mathcal{Y}$ est le morphisme
$\Delta' : \mathcal{X}' \to \mathcal{X}' \times_Y \mathcal{X}'$.
Il suffit donc de montrer que $\Delta'$ est surjectif, plat et localement de
présentation finie (voir les Lemmes \ref{stacks-morphisms-lemma-descent-flat} et
\ref{stacks-morphisms-lemma-descent-finite-presentation}). Autrement dit, il faut montrer que
$$
[Y/G] \longrightarrow [Y/G \times_Y G]
$$
est surjectif, plat et localement de présentation finie. En effet, son
changement de base par le morphisme surjectif, plat et localement de
présentation finie $Y \to [Y/G \times_Y G]$ est le morphisme $G \to Y$.

\medskip\noindent
Démonstration de (3). Observons que le diagramme
$$
\xymatrix{
T_1 \times_{x_1, \mathcal{X}, x_2} T_2 \ar[d] \ar[r] &
T_1 \times_{y_1, \mathcal{Y}, y_2} T_2 \ar[d] \\
\mathcal{X} \ar[r] & \mathcal{X} \times_\mathcal{Y} \mathcal{X}
}
$$
est cartésien. Ainsi (3) résulte de (2).

\medskip\noindent
Démonstration de (4). Cela résulte de l'égalité
$$
\mathit{Isom}_\mathcal{X}(x_1, x_2) =
(T \times_{x_1, \mathcal{X}, x_2} T) \times_{T \times T, \Delta_T} T
$$
car le morphisme de (4) est ainsi un changement de base du morphisme de (3).
\end{proof}

\begin{proposition}
\label{stacks-morphisms-proposition-when-gerbe-over}
Soit $f : \mathcal{X} \to \mathcal{Y}$ un morphisme de champs algébriques.
Les conditions suivantes sont équivalentes :
\begin{enumerate}
\item $\mathcal{X}$ est une gerbe sur $\mathcal{Y}$ ;
\item $f : \mathcal{X} \to \mathcal{Y}$ et
$\Delta : \mathcal{X} \to \mathcal{X} \times_\mathcal{Y} \mathcal{X}$
sont surjectifs, plats et localement de présentation finie.
\end{enumerate}
\end{proposition}

\begin{proof}
L'implication (1) $\Rightarrow$ (2) résulte des
Lemmes \ref{stacks-morphisms-lemma-gerbe-fppf} et \ref{stacks-morphisms-lemma-gerbe-isom-fppf}.

\medskip\noindent
Supposons (2). Il suffit de prouver (1) après avoir changé $f$ de base par un
morphisme $\mathcal{Y}' \to \mathcal{Y}$ surjectif, plat et localement de
présentation finie ; voir le Lemme \ref{stacks-morphisms-lemma-gerbe-descent} (le changement
de base de la diagonale de $f$ est la diagonale du changement de base). Nous
pouvons donc supposer que $\mathcal{Y}$ est un schéma $Y$. Dans ce cas,
$\mathcal{I}_\mathcal{X} \to \mathcal{X}$ est un changement de base de
$\Delta$, et la Proposition \ref{stacks-morphisms-proposition-when-gerbe} montre que
$\mathcal{X}$ est une gerbe. Il reste à montrer que $\mathcal{X}$ est une
gerbe sur $Y$. Soit $\mathcal{X} \to X$ le morphisme du
Lemme \ref{stacks-morphisms-lemma-gerbe-over-iso-classes} qui fait de $\mathcal{X}$ une gerbe
au-dessus de l'espace algébrique $X$ classifiant les classes d'isomorphisme
d'objets de $\mathcal{X}$. Le morphisme $f : \mathcal{X} \to Y$ se factorise
manifestement en $\mathcal{X} \to X \to Y$. Comme $f$ est surjectif, plat et
localement de présentation finie, $X \to Y$ est surjectif comme morphisme de
faisceaux fppf (utiliser par exemple le
Lemme \ref{stacks-morphisms-lemma-surjective-flat-locally-finite-presentation}). D'autre part,
$X \to Y$ est aussi injectif : pour tout schéma $T$ et tous points à valeurs
dans $T$, $x_1, x_2$, de $X$ ayant même image dans $Y$, on peut
d'abord, localement pour la topologie fppf sur $T$, relever $x_1, x_2$ en des
objets $\xi_1, \xi_2$ de $\mathcal{X}$ au-dessus de $T$, puis déduire que
$\xi_1$ et $\xi_2$ sont localement isomorphes pour la topologie fppf, puisque
$\Delta : \mathcal{X} \to \mathcal{X} \times_Y \mathcal{X}$
est surjectif, plat et localement de présentation finie. On a donc
$x_1 = x_2$ par construction de $X$. Ainsi $X = Y$, ce qui achève la
démonstration.
\end{proof}

\noindent
Nous disposons maintenant d'un formalisme suffisant pour montrer que les
gerbes résiduelles, lorsqu'elles existent, sont des gerbes.

\begin{lemma}
\label{stacks-morphisms-lemma-noetherian-singleton-stack-gerbe}
Soit $\mathcal{Z}$ un champ algébrique réduit et localement noethérien tel que
$|\mathcal{Z}|$ soit un singleton. Alors $\mathcal{Z}$ est une gerbe sur
un espace algébrique $Z$ réduit et localement noethérien, dont $|Z|$ est un
singleton.
\end{lemma}

\begin{proof}
D'après Propriétés des champs, Lemme
\ref{stacks-properties-lemma-unique-point-better}, il existe un morphisme
surjectif, plat et localement de présentation finie
$\Spec(k) \to \mathcal{Z}$, où $k$ est un corps. Le morphisme
$\mathcal{I}_\mathcal{Z} \times_\mathcal{Z} \Spec(k) \to \Spec(k)$ est alors
représentable par des espaces algébriques et localement de type fini (comme
changement de base de $\mathcal{I}_\mathcal{Z} \to \mathcal{Z}$ ; voir les
Lemmes \ref{stacks-morphisms-lemma-inertia} et \ref{stacks-morphisms-lemma-base-change-finite-type}). Il est donc
localement de présentation finie ; voir Morphismes d'espaces, Lemme
\ref{spaces-morphisms-lemma-noetherian-finite-type-finite-presentation}.
Il est aussi plat, puisque $k$ est un corps. Les
Lemmes \ref{stacks-morphisms-lemma-descent-flat} et
\ref{stacks-morphisms-lemma-descent-finite-presentation} montrent donc que
$\mathcal{I}_\mathcal{Z} \to \mathcal{Z}$ est plat et localement de
présentation finie. La Proposition \ref{stacks-morphisms-proposition-when-gerbe} implique que
$\mathcal{Z}$ est une gerbe. Soit $\pi : \mathcal{Z} \to Z$ un morphisme vers
un espace algébrique tel que $\mathcal{Z}$ soit une gerbe sur $Z$.
Le morphisme $\pi$ est surjectif, plat et localement de présentation finie
d'après le Lemme \ref{stacks-morphisms-lemma-gerbe-fppf}. Par composition,
$\Spec(k) \to Z$ est donc surjectif, plat et localement de présentation finie ;
voir Propriétés des champs, Lemme
\ref{stacks-properties-lemma-composition-surjective}, et les
Lemmes \ref{stacks-morphisms-lemma-composition-flat} et
\ref{stacks-morphisms-lemma-composition-finite-presentation}. Enfin, Propriétés des champs,
Lemme \ref{stacks-properties-lemma-unique-point-better}, montre que $|Z|$ est
un singleton et que $Z$ est localement noethérien et réduit.
\end{proof}

\begin{lemma}
\label{stacks-morphisms-lemma-gerbe-bijection-points}
Soit $f : \mathcal{X} \to \mathcal{Y}$ un morphisme de champs algébriques.
Si $\mathcal{X}$ est une gerbe sur $\mathcal{Y}$, alors $f$ est un
homéomorphisme universel.
\end{lemma}

\begin{proof}
D'après le Lemme \ref{stacks-morphisms-lemma-base-change-gerbe}, l'hypothèse sur $f$ est stable
par changement de base. Il suffit donc de montrer que l'application
$|\mathcal{X}| \to |\mathcal{Y}|$ est un homéomorphisme d'espaces
topologiques. Soient $k$ un corps et $y$ un objet de $\mathcal{Y}$ au-dessus de
$\Spec(k)$. D'après la propriété (2)(a) de Champs, Lemme
\ref{stacks-lemma-when-gerbe}, il existe un recouvrement fppf
$\{T_i \to \Spec(k)\}$ et des objets $x_i$ de $\mathcal{X}$ au-dessus de
$T_i$, avec $f(x_i) \cong y|_{T_i}$. Choisissons $i$ tel que
$T_i \not = \emptyset$, puis un morphisme $\Spec(K) \to T_i$ pour un certain
corps $K$. Alors $k \subset K$ et $x_i|_K$ est un objet de $\mathcal{X}$
au-dessus de $y|_K$. L'application
$|\mathcal{X}| \to |\mathcal{Y}|$ est donc surjective. L'application
$|\mathcal{X}| \to |\mathcal{Y}|$ est aussi injective. En effet, si
$x, x'$ sont des objets de $\mathcal{X}$ au-dessus de
$\Spec(k)$ dont les images $f(x), f(x')$ deviennent isomorphes (après une
extension) dans $\mathcal{Y}$, la propriété (2)(b) de Champs, Lemme
\ref{stacks-lemma-when-gerbe}, garantit l'existence d'une extension de $k$ sur
laquelle $x$ et $x'$ deviennent isomorphes (détails omis). Ainsi,
$|\mathcal{X}| \to |\mathcal{Y}|$ est continue et bijective ; il suffit de
montrer qu'elle est ouverte. Cela résulte des
Lemmes \ref{stacks-morphisms-lemma-gerbe-fppf} et \ref{stacks-morphisms-lemma-fppf-open}.
\end{proof}

\begin{lemma}
\label{stacks-morphisms-lemma-gerbe-diagonal-quasi-compact}
Soit $f : \mathcal{X} \to \mathcal{Y}$ un morphisme de champs algébriques tel
que $\mathcal{X}$ soit une gerbe sur $\mathcal{Y}$. Si
$\Delta_\mathcal{X}$ est quasi-compact, alors $\Delta_\mathcal{Y}$ l'est aussi.
\end{lemma}

\begin{proof}
Considérons le diagramme
$$
\xymatrix{
\mathcal{X} \ar[r] &
\mathcal{X} \times_\mathcal{Y} \mathcal{X} \ar[r] \ar[d] &
\mathcal{X} \times \mathcal{X} \ar[d] \\
&
\mathcal{Y} \ar[r] &
\mathcal{Y} \times \mathcal{Y}
}
$$
La Proposition \ref{stacks-morphisms-proposition-when-gerbe-over} montre que la flèche
supérieure de gauche est surjective. Puisque le composé des flèches
horizontales supérieures est quasi-compact, le
Lemme \ref{stacks-morphisms-lemma-surjection-from-quasi-compact} montre que la flèche supérieure
de droite est quasi-compact. Le carré est cartésien et la flèche verticale de
droite est surjective, plate et localement de présentation finie. On conclut
donc par le Lemme \ref{stacks-morphisms-lemma-descent-quasi-compact}.
\end{proof}

\noindent
Le lemme suivant affirme que les gerbes résiduelles existent en tout point
d'un champ algébrique qui est une gerbe.

\begin{lemma}
\label{stacks-morphisms-lemma-gerbe-residual-gerbe-exists}
Soit $\mathcal{X}$ un champ algébrique. Si $\mathcal{X}$ est une gerbe, alors,
pour tout $x \in |\mathcal{X}|$, la gerbe résiduelle de $\mathcal{X}$ en $x$
existe.
\end{lemma}

\begin{proof}
Soit $\pi : \mathcal{X} \to X$ un morphisme de $\mathcal{X}$ vers un espace
algébrique $X$ qui fasse de $\mathcal{X}$ une gerbe sur $X$. Soit
$Z_x \to X$ l'espace résiduel de $X$ en $x$ ; voir Espaces décents, Définition
\ref{decent-spaces-definition-residual-space}. Posons
$\mathcal{Z} = \mathcal{X} \times_X Z_x$. D'après le
Lemme \ref{stacks-morphisms-lemma-base-change-gerbe}, le champ algébrique $\mathcal{Z}$ est une
gerbe sur $Z_x$. Ainsi $|\mathcal{Z}| = |Z_x|$ est un singleton
(Lemme \ref{stacks-morphisms-lemma-gerbe-bijection-points}). Comme
$\mathcal{Z} \to Z_x$ est localement de présentation finie, en tant que
changement de base de $\pi$ (voir les Lemmes \ref{stacks-morphisms-lemma-gerbe-fppf} et
\ref{stacks-morphisms-lemma-base-change-finite-presentation}), le champ $\mathcal{Z}$ est
localement noethérien ; voir le
Lemme \ref{stacks-morphisms-lemma-locally-finite-type-locally-noetherian}. La gerbe résiduelle
$\mathcal{Z}_x$ de $\mathcal{X}$ en $x$ existe donc et vaut
$\mathcal{Z}_x = \mathcal{Z}_{red}$, la réduction du champ algébrique
$\mathcal{Z}$. En effet, nous avons vu que $|\mathcal{Z}_{red}|$ est un
singleton qui s'envoie sur $x \in |\mathcal{X}|$ ; ce champ est réduit par
construction et localement noethérien (la réduction d'un champ algébrique
localement noethérien est localement noethérienne ; détails omis).
\end{proof}









\section{Stratification par des gerbes}
\label{stacks-morphisms-section-stratify}

\noindent
Le but de cette section est de montrer que de nombreux champs algébriques
$\mathcal{X}$ possèdent une « stratification » par des sous-champs
localement fermés $\mathcal{X}_i \subset \mathcal{X}$ tels que chaque
$\mathcal{X}_i$ soit une gerbe. Cela montre qu'en un certain sens les gerbes
sont les briques élémentaires dont est construit tout champ algébrique.
Notons que, par stratification, nous entendons ici uniquement que
$$
|\mathcal{X}| = \bigcup\nolimits_i |\mathcal{X}_i|
$$
est une stratification de l'espace topologique associé à $\mathcal{X}$,
et rien de plus (dans cette section). On peut donc, sans inconvénient,
remplacer $\mathcal{X}$ par sa réduction (voir Propriétés des champs,
section \ref{stacks-properties-section-reduced}) pour étudier cette
stratification.

\medskip\noindent
La proposition suivante affirme qu'il existe (presque toujours) un
sous-champ ouvert dense de la réduction de $\mathcal{X}$.

\begin{proposition}
\label{stacks-morphisms-proposition-open-stratum}
Soit $\mathcal{X}$ un champ algébrique réduit tel que
$\mathcal{I}_\mathcal{X} \to \mathcal{X}$ soit quasi-compact.
Il existe alors un sous-champ ouvert dense
$\mathcal{U} \subset \mathcal{X}$ qui est une gerbe.
\end{proposition}

\begin{proof}
D'après la Proposition \ref{stacks-morphisms-proposition-when-gerbe}, il suffit de trouver un
sous-champ ouvert dense $\mathcal{U}$ tel que
$\mathcal{I}_\mathcal{U} \to \mathcal{U}$ soit plat et localement de
présentation finie. Notons que
$\mathcal{I}_\mathcal{U} =
\mathcal{I}_\mathcal{X} \times_\mathcal{X} \mathcal{U}$ ; voir le
Lemme \ref{stacks-morphisms-lemma-cartesian-square-inertia}.

\medskip\noindent
Choisissons une présentation $\mathcal{X} = [U/R]$. Soit $G \to U$
l'espace algébrique en groupes stabilisateur du groupoïde $R$. D'après le
Lemme \ref{stacks-morphisms-lemma-presentation-inertia}, $G \to U$ est le changement de base
de $\mathcal{I}_\mathcal{X} \to \mathcal{X}$ ; il est donc quasi-compact
(par hypothèse) et localement de type fini (d'après le
Lemme \ref{stacks-morphisms-lemma-inertia}). Soit $W \subset U$ le plus grand sous-schéma
ouvert (éventuellement vide) tel que la restriction $G_W \to W$ soit plate
et localement de présentation finie (nous omettons la preuve de l'existence de
$W$ ; indication : utiliser le caractère local des propriétés). D'après
Morphismes d'espaces, Proposition
\ref{spaces-morphisms-proposition-generic-flatness-reduced}
$W \subset U$ est dense. Notons que $W \subset U$ est $R$-invariant d'après
Compléments sur les groupoïdes en espaces, Lemme
\ref{spaces-more-groupoids-lemma-property-G-invariant}.
Ainsi, $W$ correspond à un sous-champ ouvert
$\mathcal{U} \subset \mathcal{X}$ d'après Propriétés des champs, Lemme
\ref{stacks-properties-lemma-substacks-presentation}.
Comme l'application $|U| \to |\mathcal{X}|$ est ouverte et que
$|W| \subset |U|$ est dense, $\mathcal{U}$ est dense dans $\mathcal{X}$.
Enfin, le morphisme $\mathcal{I}_\mathcal{U} \to \mathcal{U}$ est plat et
localement de présentation finie, car son changement de base par le morphisme
lisse surjectif $W \to \mathcal{U}$ est le morphisme $G_W \to W$, qui est
plat et localement de présentation finie par construction. Voir les
Lemmes \ref{stacks-morphisms-lemma-descent-flat} et
\ref{stacks-morphisms-lemma-descent-finite-presentation}.
\end{proof}

\noindent
La proposition précédente implique immédiatement que tout point d'un champ
algébrique à inertie quasi-compacte possède une gerbe résiduelle, comme nous
le montrerons dans le Lemme \ref{stacks-morphisms-lemma-every-point-residual-gerbe}.
Il n'existe cependant pas toujours de stratification finie par des gerbes.
En voici un exemple.

\begin{example}
\label{stacks-morphisms-example-infinite-stratification}
Soit $k$ un corps. Prenons $U = \Spec(k[x_0, x_1, x_2, \ldots])$ et faisons
agir $\mathbf{G}_m$ par
$t(x_0, x_1, x_2, \ldots) = (tx_0, t^p x_1, t^{p^2} x_2, \ldots)$
où $p$ est un nombre premier. Posons
$\mathcal{X} = [U/\mathbf{G}_m]$. C'est un champ algébrique. Il existe une
stratification de $\mathcal{X}$ dont les strates sont les suivantes :
\begin{enumerate}
\item $\mathcal{X}_0$ est le lieu où $x_0$ est non nul,
\item $\mathcal{X}_1$ est le lieu où $x_0$ est nul mais $x_1$ est non nul,
\item $\mathcal{X}_2$ est le lieu où $x_0, x_1$ sont nuls mais $x_2$ est non
nul,
\item et ainsi de suite, et
\item $\mathcal{X}_{\infty}$ est le lieu où tous les $x_i$ sont nuls.
\end{enumerate}
Chaque strate est une gerbe sur un schéma, de groupe $\mu_{p^i}$ pour
$\mathcal{X}_i$ et $\mathbf{G}_m$ pour $\mathcal{X}_{\infty}$.
Les strates sont des sous-champs localement fermés réduits. Il n'existe pas
de stratification plus grossière ayant les mêmes propriétés.
\end{example}

\noindent
Néanmoins, par induction transfinie, la
Proposition \ref{stacks-morphisms-proposition-open-stratum} permet de construire des
stratifications par des gerbes, éventuellement infinies... !

\begin{lemma}
\label{stacks-morphisms-lemma-every-point-in-a-stratum}
Soit $\mathcal{X}$ un champ algébrique tel que
$\mathcal{I}_\mathcal{X} \to \mathcal{X}$ soit quasi-compact.
Il existe alors un ensemble d'indices bien ordonné $I$ et, pour tout
$i \in I$, un sous-champ localement fermé réduit
$\mathcal{U}_i \subset \mathcal{X}$ tels que
\begin{enumerate}
\item chaque $\mathcal{U}_i$ est une gerbe,
\item on a $|\mathcal{X}| = \bigcup_{i \in I} |\mathcal{U}_i|$,
\item $T_i = |\mathcal{X}| \setminus \bigcup_{i' < i} |\mathcal{U}_{i'}|$
est fermé dans $|\mathcal{X}|$ pour tout $i \in I$, et
\item $|\mathcal{U}_i|$ est ouvert dans $T_i$.
\end{enumerate}
On peut de plus s'arranger pour que, soit (a)
$|\mathcal{U}_i| \subset T_i$ soit dense, soit (b) $\mathcal{U}_i$ soit
quasi-compact. Dans le cas (a), si l'on choisit $\mathcal{U}_i$ aussi grand
que possible (voir la démonstration pour les détails), la stratification est
canonique.
\end{lemma}

\begin{proof}
Soit $T \subset |\mathcal{X}|$ une partie fermée non vide. Pour tout $T$ de
cette sorte, nous allons trouver (resp.\ choisir) un sous-champ localement
fermé réduit $\mathcal{U}(T) \subset \mathcal{X}$ tel que
$|\mathcal{U}(T)| \subset T$ soit une partie ouverte dense
(resp.\ une partie ouverte non vide et quasi-compacte) de cette dernière. En effet,
d'après Propriétés des champs, Lemme
\ref{stacks-properties-lemma-reduced-closed-substack}
il existe un unique sous-champ fermé réduit
$\mathcal{X}' \subset \mathcal{X}$ tel que $T = |\mathcal{X}'|$.
Notons que $\mathcal{I}_{\mathcal{X}'} =
\mathcal{I}_\mathcal{X} \times_\mathcal{X} \mathcal{X}'$ d'après le
Lemme \ref{stacks-morphisms-lemma-monomorphism-cartesian-square-inertia}.
Ainsi, $\mathcal{I}_{\mathcal{X}'} \to \mathcal{X}'$ est quasi-compact en
tant que changement de base ; voir le
Lemme \ref{stacks-morphisms-lemma-base-change-quasi-compact}. La
Proposition \ref{stacks-morphisms-proposition-open-stratum} implique donc qu'il existe un
sous-champ ouvert dense maximal $\mathcal{U} \subset \mathcal{X}'$ qui est
une gerbe (voir la démonstration de cette proposition). Dans le cas (a),
posons $\mathcal{U}(T) = \mathcal{U}$ (ce choix est canonique) ; dans le cas
(b), choisissons simplement un ouvert non vide quasi-compact
$\mathcal{U}(T) \subset \mathcal{U}$ ; voir Propriétés des champs, Lemme
\ref{stacks-properties-lemma-space-locally-quasi-compact}
(l'axiome du choix permet d'effectuer ces choix simultanément pour tous les
$T$).

\medskip\noindent
Par récurrence transfinie, construisons pour tout ordinal $\alpha$ une partie
fermée $T_\alpha \subset |\mathcal{X}|$. Pour $\alpha = 0$, posons
$T_0 = |\mathcal{X}|$. Une fois $T_\alpha$ construit, posons
$$
T_{\alpha + 1} = T_\alpha \setminus |\mathcal{U}(T_\alpha)|.
$$
Si $\beta$ est un ordinal limite, posons
$$
T_\beta = \bigcap\nolimits_{\alpha < \beta} T_\alpha.
$$
Nous affirmons que $T_\alpha = \emptyset$ pour tout $\alpha$ assez grand.
Supposons en effet que $T_\alpha \not = \emptyset$ pour tout $\alpha$.
On obtiendrait alors une application injective de la classe des ordinaux dans
l'ensemble des parties de $|\mathcal{X}|$, ce qui est contradictoire.

\medskip\noindent
L'affirmation implique le lemme. Posons en effet
$$
I = \{\alpha \mid \mathcal{U}(T_\alpha) \not = \emptyset \}.
$$
D'après l'affirmation, c'est un ensemble bien ordonné. Pour
$i = \alpha \in I$, posons $\mathcal{U}_i = \mathcal{U}(T_\alpha)$.
Ainsi, $\mathcal{U}_i$ est un sous-champ localement fermé réduit et une
gerbe, ce qui donne (1). Par construction, $T_i = T_\alpha$ si
$i = \alpha \in I$, d'où (3). De même, (4) ainsi que (a) ou (b) résultent
de notre choix de $\mathcal{U}(T)$. Enfin, pour établir (2), soit
$x \in |\mathcal{X}|$. Il existe un plus petit ordinal $\beta$ tel que
$x \not \in T_\beta$ (puisque les ordinaux sont bien ordonnés). Cet ordinal
$\beta$ est nécessairement un successeur d'après la définition de $T_\beta$
pour les ordinaux limites. Ainsi, $\beta = \alpha + 1$ et
$x \in |\mathcal{U}(T_\alpha)|$, ce qui achève la preuve.
\end{proof}

\begin{remark}
\label{stacks-morphisms-remark-order-type}
On peut s'interroger sur le type d'ordre des stratifications canoniques
obtenues par la construction de type (a) du
Lemme \ref{stacks-morphisms-lemma-every-point-in-a-stratum}. Il est naturel de conjecturer que
l'ensemble bien ordonné $I$ est de {\it cardinal} au plus $\aleph_0$.
Nous ignorons si cela est vrai ou faux. Si vous le savez, veuillez écrire à
\href{mailto:stacks.project@gmail.com}{stacks.project@gmail.com}.
\end{remark}







\section{L'espace topologique d'un champ algébrique}
\label{stacks-morphisms-section-topology}

\noindent
Dans cette section, nous appliquons les résultats précédents à l'espace
topologique $|\mathcal{X}|$ associé à un champ algébrique.

\begin{lemma}
\label{stacks-morphisms-lemma-spectral-qc-diagonal-qc}
Soit $\mathcal{X}$ un champ algébrique quasi-compact dont la diagonale
$\Delta$ est quasi-compacte. Alors $|\mathcal{X}|$ est un espace spectral.
\end{lemma}

\begin{proof}
Choisissons un schéma affine $U$ et un morphisme lisse surjectif
$U \to \mathcal{X}$ ; voir Propriétés des champs, Lemme
\ref{stacks-properties-lemma-quasi-compact-stack}. L'application
$|U| \to |\mathcal{X}|$ est alors continue, ouverte et surjective ; voir
Propriétés des champs, Lemme
\ref{stacks-properties-lemma-topology-points}. Les ouverts quasi-compacts de
$|\mathcal{X}|$ forment donc une base de la topologie. Pour deux ouverts
quasi-compacts $W_1, W_2 \subset |\mathcal{X}|$, on peut choisir des ouverts
quasi-compacts $V_1, V_2$ de $U$ dont les images sont respectivement $W_1$ et
$W_2$. Comme $\Delta$ est quasi-compacte,
$$
V_1 \times_\mathcal{X} V_2 =
(V_1 \times V_2) \times_{\mathcal{X} \times \mathcal{X}, \Delta} \mathcal{X}
$$
est quasi-compact. L'image de $|V_1 \times_\mathcal{X} V_2|$ dans
$|\mathcal{X}|$ est $W_1 \cap W_2$ d'après Propriétés des champs, Lemme
\ref{stacks-properties-lemma-points-cartesian}. Ainsi,
$W_1 \cap W_2$ est quasi-compact. Pour conclure, il suffit de montrer que
$|\mathcal{X}|$ est sobre ; voir Topologie, Définition
\ref{topology-definition-spectral-space}.

\medskip\noindent
Soit $T \subset |\mathcal{X}|$ une partie fermée irréductible. Il faut
montrer que $T$ possède un unique point générique. Soit
$\mathcal{Z} \subset \mathcal{X}$ le sous-champ fermé muni de la structure
réduite induite correspondant à $T$ ; voir Propriétés des champs, Définition
\ref{stacks-properties-definition-reduced-induced-stack}.
Comme $\mathcal{Z} \to \mathcal{X}$ est une immersion fermée,
$\Delta_\mathcal{Z}$ est quasi-compacte : montrons d'abord que
$\mathcal{Z} \to \mathcal{X} \times \mathcal{X}$ est quasi-compact comme
composé de $\mathcal{Z} \to \mathcal{X}$ et de $\Delta$, puis écrivons
$\mathcal{Z} \to \mathcal{X} \times \mathcal{X}$ comme le composé de
$\Delta_\mathcal{Z}$ et de
$\mathcal{Z} \times \mathcal{Z} \to \mathcal{X} \times \mathcal{X}$, et
appliquons le Lemme \ref{stacks-morphisms-lemma-quasi-compact-permanence} ainsi que le fait que
$\mathcal{Z} \times \mathcal{Z} \to \mathcal{X} \times \mathcal{X}$ est
séparé. Nous sommes ainsi ramenés au cas traité dans le paragraphe suivant.

\medskip\noindent
Supposons $\mathcal{X}$ quasi-compact, $\Delta$ quasi-compacte,
$\mathcal{X}$ réduit et $|\mathcal{X}|$ irréductible. Il faut montrer que
$|\mathcal{X}|$ possède un unique point générique. Comme
$\mathcal{I}_\mathcal{X} \to \mathcal{X}$ est un changement de base de
$\Delta$, le morphisme
$\mathcal{I}_\mathcal{X} \to \mathcal{X}$ est quasi-compact
(Lemme \ref{stacks-morphisms-lemma-base-change-quasi-compact}). D'après la
Proposition \ref{stacks-morphisms-proposition-open-stratum}, il existe donc un sous-champ
ouvert dense $\mathcal{U} \subset \mathcal{X}$ qui est une gerbe. Autrement
dit, $|\mathcal{U}| \subset |\mathcal{X}|$ est ouvert dense. Nous pouvons
donc supposer que $\mathcal{X}$ est une gerbe. Soit
$\mathcal{X} \to X$ un morphisme faisant de $\mathcal{X}$ une gerbe sur
l'espace algébrique $X$.
D'après le Lemme \ref{stacks-morphisms-lemma-gerbe-bijection-points}, on a alors
$|\mathcal{X}| \cong |X|$. En particulier, $X$ est quasi-compact. D'après le
Lemme \ref{stacks-morphisms-lemma-gerbe-diagonal-quasi-compact}, la diagonale de $X$ est
quasi-compacte, c'est-à-dire que $X$ est un espace algébrique quasi-séparé.
Alors $|X|$ est spectral d'après Propriétés des espaces, Lemme
\ref{spaces-properties-lemma-quasi-compact-quasi-separated-spectral}
et le résultat en découle.
\end{proof}

\begin{lemma}
\label{stacks-morphisms-lemma-spectral-qcqs}
Soit $\mathcal{X}$ un champ algébrique quasi-compact et quasi-séparé.
Alors $|\mathcal{X}|$ est un espace spectral.
\end{lemma}

\begin{proof}
C'est un cas particulier du Lemme \ref{stacks-morphisms-lemma-spectral-qc-diagonal-qc}.
\end{proof}

\begin{lemma}
\label{stacks-morphisms-lemma-sober-qs}
Soit $\mathcal{X}$ un champ algébrique dont la diagonale est quasi-compacte
(par exemple, si $\mathcal{X}$ est quasi-séparé). Il existe alors un
recouvrement ouvert $|\mathcal{X}| = \bigcup U_i$ tel que chaque $U_i$ soit
spectral. En particulier, $|\mathcal{X}|$ est un espace sobre.
\end{lemma}

\begin{proof}
Conséquence immédiate du Lemme \ref{stacks-morphisms-lemma-spectral-qc-diagonal-qc}.
\end{proof}






\section{Existence des gerbes résiduelles}
\label{stacks-morphisms-section-existence-residual-gerbes}

\noindent
La gerbe résiduelle d'un point d'un champ algébrique est définie dans
Propriétés des champs, Définition
\ref{stacks-properties-definition-residual-gerbe}.
Nous avons déjà montré que les gerbes résiduelles existent aux points
de type fini (Lemme \ref{stacks-morphisms-lemma-point-finite-type-monomorphism})
et en tout point d'une gerbe (Lemme \ref{stacks-morphisms-lemma-gerbe-residual-gerbe-exists}).
Dans cette section, nous montrons que les gerbes résiduelles existent sur
de nombreux champs algébriques. Commençons par l'application annoncée de la
Proposition \ref{stacks-morphisms-proposition-open-stratum}.

\begin{lemma}
\label{stacks-morphisms-lemma-every-point-residual-gerbe}
Soit $\mathcal{X}$ un champ algébrique tel que
$\mathcal{I}_\mathcal{X} \to \mathcal{X}$ soit quasi-compact.
Alors la gerbe résiduelle de $\mathcal{X}$ en $x$ existe pour tout
$x \in |\mathcal{X}|$.
\end{lemma}

\begin{proof}
Soit $T = \overline{\{x\}} \subset |\mathcal{X}|$ l'adhérence de $x$.
D'après Propriétés des champs, Lemme
\ref{stacks-properties-lemma-reduced-closed-substack}
il existe un sous-champ fermé réduit $\mathcal{X}' \subset \mathcal{X}$
tel que $T = |\mathcal{X}'|$. Notons que
$\mathcal{I}_{\mathcal{X}'} =
\mathcal{I}_\mathcal{X} \times_\mathcal{X} \mathcal{X}'$ d'après le
Lemme \ref{stacks-morphisms-lemma-monomorphism-cartesian-square-inertia}.
Ainsi $\mathcal{I}_{\mathcal{X}'} \to \mathcal{X}'$ est quasi-compact
comme changement de base, voir le
Lemme \ref{stacks-morphisms-lemma-base-change-quasi-compact}.
La Proposition \ref{stacks-morphisms-proposition-open-stratum} implique donc qu'il existe
un sous-champ ouvert dense
$\mathcal{U} \subset \mathcal{X}'$
qui est une gerbe. Notons que $x \in |\mathcal{U}|$, car
$\{x\} \subset T$ est lui aussi dense. Il existe donc une gerbe résiduelle
$\mathcal{Z}_x \subset \mathcal{U}$ de $\mathcal{U}$ en $x$, d'après le
Lemme \ref{stacks-morphisms-lemma-gerbe-residual-gerbe-exists}.
Il résulte immédiatement des définitions que $\mathcal{Z}_x \to \mathcal{X}$
est une gerbe résiduelle de $\mathcal{X}$ en $x$.
\end{proof}

\noindent
Si le champ est quasi-DM, les gerbes résiduelles existent également.
En particulier, elles existent toujours sur les champs de Deligne--Mumford.

\begin{lemma}
\label{stacks-morphisms-lemma-every-point-residual-gerbe-quasi-DM}
Soit $\mathcal{X}$ un champ algébrique quasi-DM.
Alors la gerbe résiduelle de $\mathcal{X}$ en $x$ existe pour tout
$x \in |\mathcal{X}|$.
\end{lemma}

\begin{proof}
Choisissons un schéma $U$ et un morphisme $U \to \mathcal{X}$ surjectif,
plat, localement de présentation finie et localement quasi-fini, voir le
Théorème \ref{stacks-morphisms-theorem-quasi-DM}.
Posons $R = U \times_\mathcal{X} U$. Les projections $s, t : R \to U$
sont surjectives, plates, localement de présentation finie et localement
quasi-finies, car ce sont des changements de base du morphisme
$U \to \mathcal{X}$. Il existe un morphisme canonique
$[U/R] \to \mathcal{X}$ (voir Champs algébriques, Lemme
\ref{algebraic-lemma-map-space-into-stack}), qui est une équivalence puisque
$U \to \mathcal{X}$ est surjectif, plat et localement de présentation finie,
voir Champs algébriques, Remarque
\ref{algebraic-remark-flat-fp-presentation}.
Nous pouvons donc supposer que $\mathcal{X} = [U/R]$, où
$(U, R, s, t, c)$ est un groupoïde en espaces algébriques tel que
$s, t : R \to U$ soient surjectifs, plats, localement de présentation finie
et localement quasi-finis. Posons
$$
U' = \coprod\nolimits_{u \in U\text{ se projetant sur }x} \Spec(\kappa(u)).
$$
Le morphisme canonique $U' \to U$ est un monomorphisme. Posons
$$
R' = U' \times_\mathcal{X} U' =
R \times_{(U \times U)} (U' \times U')
$$
Comme $U' \to U$ est un monomorphisme, chacune des deux projections
$s', t' : R' \to U'$ se factorise en un monomorphisme suivi d'un morphisme
localement quasi-fini. Puisque $U'$ est une somme disjointe de spectres
de corps, le Lemme
\ref{spaces-over-fields-lemma-mono-towards-locally-quasi-finite-over-field}
d'Espaces sur des corps montre que les morphismes
$s', t' : R' \to U'$ sont localement quasi-finis. Comme $U'$ est une
somme disjointe de spectres de corps, les morphismes $s', t'$ sont aussi
plats. Enfin, comme $s', t'$ sont localement quasi-finis, $s', t'$ sont
localement de type fini et, par suite, $s', t'$ sont localement de
présentation finie (en effet,
$U'$ est une somme disjointe de spectres de corps et, en particulier,
est localement noethérien, si bien que le Lemme
\ref{spaces-morphisms-lemma-noetherian-finite-type-finite-presentation}
de Morphismes d'espaces s'applique). Ainsi $\mathcal{Z} = [U'/R']$ est un
champ algébrique d'après Critères de représentabilité, Théorème
\ref{criteria-theorem-flat-groupoid-gives-algebraic-stack}.
Comme $R'$ est la restriction de $R$ par $U' \to U$, le morphisme
$\mathcal{Z} \to \mathcal{X}$ est un monomorphisme d'après Groupoïdes en
espaces, Lemme
\ref{spaces-groupoids-lemma-quotient-stack-restrict}
et Propriétés des champs, Lemme
\ref{stacks-properties-lemma-monomorphism}.
Puisque $\mathcal{Z} \to \mathcal{X}$ est un monomorphisme, l'application
$|\mathcal{Z}| \to |\mathcal{X}|$ est injective, voir Propriétés des champs,
Lemme
\ref{stacks-properties-lemma-monomorphism-injective-points}.
D'après Propriétés des champs, Lemme
\ref{stacks-properties-lemma-points-cartesian}, l'application
$$
|U'| = |\mathcal{Z} \times_\mathcal{X} U'|
\longrightarrow
|\mathcal{Z}| \times_{|\mathcal{X}|} |U'|
$$
est surjective. Par le choix de $U'$, il en résulte que l'image de
$|\mathcal{Z}| \to |\mathcal{X}|$ est $\{x\}$.
Ainsi $|\mathcal{Z}|$ est réduit à un point.
Enfin, $U'$ est par construction localement noethérien et réduit ; autrement
dit, $\mathcal{Z}$ est réduit et localement noethérien. L'image essentielle
de $\mathcal{Z} \to \mathcal{X}$ est donc la gerbe résiduelle de
$\mathcal{X}$ en $x$, voir Propriétés des champs, Lemme
\ref{stacks-properties-lemma-residual-gerbe-unique}.
\end{proof}

\begin{lemma}
\label{stacks-morphisms-lemma-every-point-residual-gerbe-locally-Noetherian}
Soit $\mathcal{X}$ un champ algébrique localement noethérien.
Alors la gerbe résiduelle de $\mathcal{X}$ en $x$ existe pour tout
$x \in |\mathcal{X}|$.
\end{lemma}

\begin{proof}
Choisissons un schéma affine $U$ et un morphisme lisse
$U \to \mathcal{X}$ tels que $x$ appartienne à l'image de l'application
continue ouverte $|U| \to |\mathcal{X}|$. Remplaçons $\mathcal{X}$ par le
sous-champ ouvert correspondant à l'image de $|U| \to |\mathcal{X}|$, voir
Propriétés des champs, Lemme \ref{stacks-properties-lemma-open-substacks}.
Nous pouvons ainsi supposer que $\mathcal{X} = [U/R]$ pour un groupoïde lisse
$(U, R, s, t, c)$ en espaces algébriques, où $U$ est un schéma affine
noethérien, voir Champs algébriques, section
\ref{algebraic-section-stack-to-presentation}.

\medskip\noindent
Soit $E \subset |U|$ l'image réciproque de
$\{x\} \subset |\mathcal{X}|$. Bien sûr, $E \not = \emptyset$.
Puisque $|U|$ est un espace topologique noethérien, nous pouvons choisir
$u \in E$ tel que $\overline{\{u\}} \cap E = \{u\}$.
Comme d'habitude, nous identifions $u = \Spec(\kappa(u))$ au spectre de son
corps résiduel. Écrivons
$$
F = u \times_{U, t} R = u \times_\mathcal{X} U
\quad\text{et}\quad
R' = (u \times u) \times_{(U \times U), (t, s)} R = u \times_\mathcal{X} u
$$
Munissons en outre $Z = \overline{\{u\}} \subset U$ de la structure réduite
de sous-schéma induite. Notons $p : F \to U$ le morphisme induit par la
seconde projection (on utilise $s : R \to U$ dans la première description
de $F$ comme produit fibré). Alors $E$ est l'image ensembliste de $p$.
Le morphisme $R' \to F$ est un monomorphisme qui se factorise par l'image
réciproque $p^{-1}(Z)$ de $Z$. Cette image réciproque
$p^{-1}(Z) \subset F$ est un sous-schéma fermé ; l'image de la restriction
$p|_{p^{-1}(Z)} : p^{-1}(Z) \to Z$ est, ensemblistement, contenue dans
$\{u\} \subset Z$ par le choix soigneux de $u \in E$ fait plus haut.
Puisque $u = \lim W$, où la limite porte sur les
sous-schémas ouverts affines non vides du schéma réduit irréductible $Z$,
le morphisme $p|_{p^{-1}(Z)} : p^{-1}(Z) \to Z$ se factorise par
$u \to Z$. Cela implique clairement que $R' = p^{-1}(Z)$. En particulier,
le morphisme $t' : R' \to u$ est localement de présentation finie : c'est la
composée de l'immersion fermée $p^{-1}(Z) \to F$ d'espaces algébriques
localement noethériens et du morphisme lisse $\text{pr}_1 : F \to u$ ; voir
Morphismes d'espaces, Lemmes
\ref{spaces-morphisms-lemma-locally-finite-type-locally-noetherian},
\ref{spaces-morphisms-lemma-closed-immersion-finite-presentation} et
\ref{spaces-morphisms-lemma-composition-finite-presentation}.
La restriction $(u, R', s', t', c')$ de $(U, R, s, t, c)$ par
$u \to U$ est donc un groupoïde en espaces algébriques dont $s'$ et $t'$
sont plats et localement de présentation finie. Ainsi
$\mathcal{Z} = [u/R']$ est un champ algébrique d'après Critères de
représentabilité, Théorème
\ref{criteria-theorem-flat-groupoid-gives-algebraic-stack}.
Comme $R'$ est la restriction de $R$ par $u \to U$, le morphisme
$\mathcal{Z} \to \mathcal{X}$ est un monomorphisme d'après Groupoïdes en
espaces, Lemme
\ref{spaces-groupoids-lemma-quotient-stack-restrict}
et Propriétés des champs, Lemme
\ref{stacks-properties-lemma-monomorphism}.
Le champ $\mathcal{Z}$ est alors isomorphe à la gerbe résiduelle d'après les
résultats de Propriétés des champs, section
\ref{stacks-properties-section-residual-gerbe}.
\end{proof}













\section{Structure locale \'etale}
\label{stacks-morphisms-section-etale-local}

\noindent
Dans cette section, nous commen\c{c}ons \`a examiner ce que l'on peut dire
de la structure locale \'etale d'un champ alg\'ebrique.

\begin{lemma}
\label{stacks-morphisms-lemma-quotient-etale}
Soit $Y$ un espace alg\'ebrique.
Soit $(U, R, s, t, c)$ un groupo\"ide en espaces alg\'ebriques sur $Y$.
Supposons que $U \to Y$ soit plat et localement de pr\'esentation finie
et que $R \to U \times_Y U$ soit une immersion ouverte.
Alors $X = [U/R] = U/R$ est un espace alg\'ebrique et $X \to Y$
est \'etale.
\end{lemma}

\begin{proof}
Le champ quotient $[U/R]$ est un champ alg\'ebrique d'apr\`es
Crit\`eres de repr\'esentabilit\'e, th\'eor\`eme
\ref{criteria-theorem-flat-groupoid-gives-algebraic-stack}.
D'autre part, puisque $R \to U \times U$ est un monomorphisme,
ce champ est un espace alg\'ebrique (en vertu de notre abus de langage et d'apr\`es
Champs alg\'ebriques, proposition
\ref{algebraic-proposition-algebraic-stack-no-automorphisms})
et il est bien s\^ur \'egal \`a l'espace alg\'ebrique $U/R$
(d'Amor\c{c}age, th\'eor\`eme \ref{bootstrap-theorem-final-bootstrap}).
Comme $U \to X$ est surjectif, plat et localement de pr\'esentation finie
(Amor\c{c}age, lemme \ref{bootstrap-lemma-covering-quotient}),
nous en d\'eduisons que $X \to Y$ est plat et localement de pr\'esentation finie d'apr\`es
Morphismes d'espaces, lemme \ref{spaces-morphisms-lemma-flat-permanence}
et
Descente sur les espaces, lemme
\ref{spaces-descent-lemma-flat-finitely-presented-permanence}.
Enfin, consid\'erons le diagramme cart\'esien
$$
\xymatrix{
R \ar[d] \ar[r] & U \times_Y U \ar[d] \\
X \ar[r] & X \times_Y X
}
$$
Puisque la fl\`eche verticale de droite est surjective, plate et
localement de pr\'esentation finie (petit d\'etail omis), nous
voyons que $X \to X \times_Y X$ est une immersion ouverte, car la fl\`eche
horizontale sup\'erieure poss\`ede cette propri\'et\'e par hypoth\`ese (utiliser
Propri\'et\'es des champs, lemme
\ref{stacks-properties-lemma-check-property-covering}).
Ainsi $X \to Y$ est non ramifi\'e d'apr\`es
Morphismes d'espaces, lemme
\ref{spaces-morphisms-lemma-diagonal-unramified-morphism}.
Nous concluons \`a l'aide de
Morphismes d'espaces, lemme
\ref{spaces-morphisms-lemma-unramified-flat-lfp-etale}.
\end{proof}

\begin{lemma}
\label{stacks-morphisms-lemma-quasi-splitting-etale}
Soit $S$ un sch\'ema.
Soit $(U, R, s, t, c)$ un groupo\"ide en espaces alg\'ebriques sur $S$.
Supposons que $s, t$ soient plats et localement de pr\'esentation finie.
Soit $P \subset R$ un sous-espace ouvert tel que
$(U, P, s|_P, t|_P, c|_{P \times_{s, U, t} P})$ soit un
groupo\"ide en espaces alg\'ebriques sur $S$. Alors
$$
[U/P] \longrightarrow [U/R]
$$
est un morphisme de champs alg\'ebriques qui
est repr\'esentable par des espaces alg\'ebriques, surjectif et \'etale.
\end{lemma}

\begin{proof}
Puisque $P \subset R$ est ouvert, $s|_P$ et $t|_P$
sont plats et localement de pr\'esentation finie.
Ainsi $[U/R]$ et $[U/P]$ sont des champs alg\'ebriques d'apr\`es
Crit\`eres de repr\'esentabilit\'e, th\'eor\`eme
\ref{criteria-theorem-flat-groupoid-gives-algebraic-stack}.
Pour voir que le morphisme est repr\'esentable par des espaces alg\'ebriques,
il suffit de montrer que $[U/P] \to [U/R]$ est fid\`ele sur
chaque cat\'egorie fibre ; voir
Champs alg\'ebriques, lemme
\ref{algebraic-lemma-characterize-representable-by-algebraic-spaces}.
Cela r\'esulte imm\'ediatement du fait que $P \to R$ est un monomorphisme
et de la description explicite des champs quotients donn\'ee dans
Groupo\"ides en espaces, lemme
\ref{spaces-groupoids-lemma-quotient-stack-objects}.
Ceci \'etabli, nous savons ce que signifie le fait que
$[U/P] \to [U/R]$ soit surjectif et \'etale, d'apr\`es
Champs alg\'ebriques, d\'efinition
\ref{algebraic-definition-relative-representable-property}.
La surjectivit\'e r\'esulte par exemple de
Crit\`eres de repr\'esentabilit\'e,
lemme \ref{criteria-lemma-representable-by-spaces-surjective}
et de la description des objets des champs quotients
(voir le lemme cit\'e ci-dessus) sur les spectres de corps.
Il reste \`a prouver que notre morphisme est \'etale.

\medskip\noindent
Pour cela, il suffit de montrer que $U \times_{[U/R]} [U/P] \to U$
est \'etale ; voir Propri\'et\'es des champs, lemme
\ref{stacks-properties-lemma-check-property-covering}.
D'apr\`es Groupo\"ides en espaces, lemme
\ref{spaces-groupoids-lemma-cartesian-square-of-morphism},
le produit fibr\'e s'identifie \`a $[R/P \times_{s, U, t} R]$,
le morphisme vers $U$ \'etant induit par $s : R \to U$.
Les applications $s', t' : P \times_{s, U, t} R \to R$ sont d\'efinies par
$s' : (p, r) \mapsto r$ et $t' : (p, r) \mapsto c(p, r)$. Puisque
$P \subset R$ est ouvert, nous en d\'eduisons que
$(t', s') : P \times_{s, U, t} R \to R \times_{s, U, s} R$
est une immersion ouverte.
Nous pouvons donc appliquer le lemme \ref{stacks-morphisms-lemma-quotient-etale}
pour conclure.
\end{proof}

\begin{lemma}
\label{stacks-morphisms-lemma-etale-local-quasi-DM}
Soit $\mathcal{X}$ un champ alg\'ebrique. Supposons que $\mathcal{X}$ soit
quasi-DM \`a diagonale s\'epar\'ee (de mani\`ere \'equivalente,
$\mathcal{I}_\mathcal{X} \to \mathcal{X}$ est localement quasi-fini et
s\'epar\'e). Soit $x \in |\mathcal{X}|$. Il existe alors un
morphisme de champs alg\'ebriques
$$
\mathcal{U} \longrightarrow \mathcal{X}
$$
ayant les propri\'et\'es suivantes :
\begin{enumerate}
\item il existe un point $u \in |\mathcal{U}|$ dont l'image est $x$,
\item $\mathcal{U} \to \mathcal{X}$ est repr\'esentable par des espaces alg\'ebriques et
\'etale,
\item $\mathcal{U} = [U/R]$, o\`u $(U, R, s, t, c)$ est un sch\'ema
en groupo\"ides, avec $U$ et $R$ affines, et $s, t$ finis, plats et
localement de pr\'esentation finie.
\end{enumerate}
\end{lemma}

\begin{proof}
(L'assertion entre parenth\`eses r\'esulte des \'equivalences du
lemme \ref{stacks-morphisms-lemma-diagonal-diagonal}).
Choisissons un sch\'ema affine $U$ et un morphisme $U \to \mathcal{X}$
plat, localement de pr\'esentation finie et localement quasi-fini, tel que $x$
soit l'image d'un point $u \in U$. Cela est possible d'apr\`es
le th\'eor\`eme \ref{stacks-morphisms-theorem-quasi-DM} et l'hypoth\`ese que $\mathcal{X}$
est quasi-DM. Soit $(U, R, s, t, c)$ le groupo\"ide en espaces alg\'ebriques
obtenu en posant $R = U \times_\mathcal{X} U$ ; voir
Champs alg\'ebriques, lemme \ref{algebraic-lemma-map-space-into-stack}.
Soit $\mathcal{X}' \subset \mathcal{X}$ le sous-champ ouvert correspondant
\`a l'image ouverte de $|U| \to |\mathcal{X}|$
(Propri\'et\'es des champs, lemmes
\ref{stacks-properties-lemma-topology-points} et
\ref{stacks-properties-lemma-open-substacks}).
Il est clair que nous pouvons remplacer $\mathcal{X}$ par le sous-champ ouvert $\mathcal{X}'$.
Nous pouvons donc supposer que $U \to \mathcal{X}$ est surjectif ; alors
Champs alg\'ebriques, remarque \ref{algebraic-remark-flat-fp-presentation}
donne $\mathcal{X} = [U/R]$.
Remarquons que $s, t : R \to U$ sont plats, localement de pr\'esentation finie
et localement quasi-finis.
Comme $R = U \times U \times_{(\mathcal{X} \times \mathcal{X})} \mathcal{X}$
et que la diagonale de $\mathcal{X}$ est s\'epar\'ee, nous voyons que
$R$ est s\'epar\'e. Par cons\'equent, $s, t : R \to U$ sont s\'epar\'es. Il s'ensuit
que $R$ est un sch\'ema d'apr\`es
Morphismes d'espaces, proposition
\ref{spaces-morphisms-proposition-locally-quasi-finite-separated-over-scheme},
appliqu\'ee \`a $s : R \to U$.

\medskip\noindent
Nous avons v\'erifi\'e ci-dessus que toutes les hypoth\`eses de
Compl\'ements sur les groupo\"ides en espaces, lemme
\ref{spaces-more-groupoids-lemma-quasi-splitting-affine-scheme}
sont satisfaites pour $(U, R, s, t, c)$ et $u$.
Nous pouvons donc trouver un voisinage \'etale \'el\'ementaire
$(U', u') \to (U, u)$ tel que la restriction $R'$ de $R$ \`a $U'$
soit quasi-scind\'ee au-dessus de $u'$. Remarquons que $R' = U' \times_\mathcal{X} U'$
(petit d\'etail omis ; indication : transitivit\'e des produits fibr\'es).
En rempla\c{c}ant $(U, R, s, t, c)$ par $(U', R', s', t', c')$ et en r\'etr\'ecissant
$\mathcal{X}$ comme ci-dessus, nous pouvons supposer que $(U, R, s, t, c)$ admet
un quasi-scindage au-dessus de $u$ (le point $u$ ne joue d\'esormais plus
aucun r\^ole, comme le montre la note de bas de page de
Compl\'ements sur les groupo\"ides en espaces, d\'efinition
\ref{spaces-more-groupoids-definition-split-at-point}).
Soit $P \subset R$ un quasi-scindage de $R$ au-dessus de $u$.
Le lemme \ref{stacks-morphisms-lemma-quasi-splitting-etale}
montre que
$$
\mathcal{U} = [U/P] \longrightarrow [U/R] = \mathcal{X}
$$
poss\`ede toutes les propri\'et\'es voulues.
\end{proof}

\begin{lemma}
\label{stacks-morphisms-lemma-etale-local-quasi-DM-at-x}
Soit $\mathcal{X}$ un champ alg\'ebrique. Supposons que $\mathcal{X}$ soit
quasi-DM \`a diagonale s\'epar\'ee (de mani\`ere \'equivalente,
$\mathcal{I}_\mathcal{X} \to \mathcal{X}$ est localement quasi-fini et
s\'epar\'e). Soit $x \in |\mathcal{X}|$. Supposons que le
groupe d'automorphismes de $\mathcal{X}$ en $x$ soit fini
(remarque \ref{stacks-morphisms-remark-property-automorphism-groups}).
Il existe alors un morphisme de champs alg\'ebriques
$$
g : \mathcal{U} \longrightarrow \mathcal{X}
$$
ayant les propri\'et\'es suivantes :
\begin{enumerate}
\item il existe un point $u \in |\mathcal{U}|$ dont l'image est $x$, et
$g$ induit un isomorphisme entre les groupes d'automorphismes en $u$ et en $x$
(remarque \ref{stacks-morphisms-remark-identify-automorphism-groups}),
\item $\mathcal{U} \to \mathcal{X}$ est repr\'esentable par des espaces alg\'ebriques et
\'etale,
\item $\mathcal{U} = [U/R]$, o\`u $(U, R, s, t, c)$ est un sch\'ema
en groupo\"ides, avec $U$ et $R$ affines, et $s, t$ finis, plats et
localement de pr\'esentation finie.
\end{enumerate}
\end{lemma}

\begin{proof}
Remarquons que $G_x$ est un sch\'ema en groupes d'apr\`es
le lemme \ref{stacks-morphisms-lemma-automorphism-group-scheme}.
Le d\'ebut de cette preuve est {\bf exactement} le m\^eme que celui de la preuve
du lemme \ref{stacks-morphisms-lemma-etale-local-quasi-DM}.
Nous pouvons donc supposer que $\mathcal{X} = [U/R]$, o\`u $(U, R, s, t, c)$
et $u \in U$, d'image $x$, satisfont \`a toutes les hypoth\`eses de
Compl\'ements sur les groupo\"ides en espaces, lemme
\ref{spaces-more-groupoids-lemma-quasi-splitting-affine-scheme}.
Notre hypoth\`ese sur $G_x$ implique que $G_u$ est fini sur $u$.
Toutes les hypoth\`eses de
Compl\'ements sur les groupo\"ides en espaces, lemme
\ref{spaces-more-groupoids-lemma-splitting-affine-scheme}
sont donc satisfaites.
Nous pouvons par cons\'equent trouver un voisinage \'etale \'el\'ementaire
$(U', u') \to (U, u)$ tel que la restriction $R'$ de $R$ \`a $U'$
soit scind\'ee au-dessus de $u'$. Remarquons que $R' = U' \times_\mathcal{X} U'$
(petit d\'etail omis ; indication : transitivit\'e des produits fibr\'es).
En rempla\c{c}ant $(U, R, s, t, c)$ par $(U', R', s', t', c')$ et en r\'etr\'ecissant
$\mathcal{X}$ comme ci-dessus, nous pouvons supposer que $(U, R, s, t, c)$ admet
un scindage au-dessus de $u$. Soit $P \subset R$ un scindage de $R$ au-dessus de $u$.
Le lemme \ref{stacks-morphisms-lemma-quasi-splitting-etale} montre que
$$
\mathcal{U} = [U/P] \longrightarrow [U/R] = \mathcal{X}
$$
est repr\'esentable par des espaces alg\'ebriques et \'etale. Par construction,
$G_u$ est contenu dans $P$ ; ce morphisme induit donc l'isomorphisme voulu
du groupe d'automorphismes en $u$ sur celui du point image.
\end{proof}

\begin{lemma}
\label{stacks-morphisms-lemma-etale-local-quasi-DM-at-x-inertia}
Soit $\mathcal{X}$ un champ alg\'ebrique. Supposons que $\mathcal{X}$ soit
quasi-DM \`a diagonale s\'epar\'ee (de mani\`ere \'equivalente,
$\mathcal{I}_\mathcal{X} \to \mathcal{X}$ est localement quasi-fini et
s\'epar\'e). Soit $x \in |\mathcal{X}|$. Supposons que $x$ puisse \^etre repr\'esent\'e
par un morphisme quasi-compact $\Spec(k) \to \mathcal{X}$.
Il existe alors un morphisme de champs alg\'ebriques
$$
g : \mathcal{U} \longrightarrow \mathcal{X}
$$
ayant les propri\'et\'es suivantes :
\begin{enumerate}
\item il existe un point $u \in |\mathcal{U}|$ dont l'image est $x$, et
$g$ induit un isomorphisme entre les gerbes r\'esiduelles en $u$ et en $x$,
\item $\mathcal{U} \to \mathcal{X}$ est repr\'esentable par des espaces alg\'ebriques et
\'etale,
\item $\mathcal{U} = [U/R]$, o\`u $(U, R, s, t, c)$ est un sch\'ema
en groupo\"ides, avec $U$ et $R$ affines, et $s, t$ finis, plats et
localement de pr\'esentation finie.
\end{enumerate}
\end{lemma}

\begin{proof}
Le d\'ebut de cette preuve est {\bf exactement} le m\^eme que celui de la preuve
du lemme \ref{stacks-morphisms-lemma-etale-local-quasi-DM}.
Nous pouvons donc supposer que $\mathcal{X} = [U/R]$, o\`u $(U, R, s, t, c)$
et $u \in U$, d'image $x$, satisfont \`a toutes les hypoth\`eses de
Compl\'ements sur les groupo\"ides en espaces, lemme
\ref{spaces-more-groupoids-lemma-quasi-splitting-affine-scheme}.
Remarquons que $u = \Spec(\kappa(u)) \to \mathcal{X}$ est quasi-compact ; voir
Propri\'et\'es des champs, lemme
\ref{stacks-properties-lemma-UR-quasi-compact-above-x}.
Consid\'erons le diagramme cart\'esien
$$
\xymatrix{
F \ar[d] \ar[r] & U \ar[d] \\
u \ar[r]^u & \mathcal{X}
}
$$
Puisque $U$ est un sch\'ema affine et que $F \to U$ est quasi-compact,
$F$ est quasi-compact. Puisque $U \to \mathcal{X}$
est localement quasi-fini, $F \to u$ est
localement quasi-fini. Ainsi $F \to u$ est quasi-fini,
et $F$ est un sch\'ema affine dont l'espace topologique
sous-jacent est fini et discret (Espaces sur des corps, lemme
\ref{spaces-over-fields-lemma-locally-quasi-finite-over-field}).
Remarquons que nous avons un monomorphisme $u \times_\mathcal{X} u \to F$.
En particulier, l'ensemble $\{r \in R : s(r) = u, t(r) = u\}$,
qui est l'image de $|u \times_\mathcal{X} u| \to |R|$,
est fini. Nous en concluons que toutes les hypoth\`eses de
Compl\'ements sur les groupo\"ides en espaces, lemme
\ref{spaces-more-groupoids-lemma-strong-splitting-affine-scheme}
sont satisfaites.

\medskip\noindent
Nous pouvons donc trouver un voisinage \'etale \'el\'ementaire
$(U', u') \to (U, u)$ tel que la restriction $R'$ de $R$ \`a $U'$
soit fortement scind\'ee au-dessus de $u'$. Remarquons que $R' = U' \times_\mathcal{X} U'$
(petit d\'etail omis ; indication : transitivit\'e des produits fibr\'es).
En rempla\c{c}ant $(U, R, s, t, c)$ par $(U', R', s', t', c')$ et en r\'etr\'ecissant
$\mathcal{X}$ comme ci-dessus, nous pouvons supposer que $(U, R, s, t, c)$ admet
un scindage fort au-dessus de $u$. Soit $P \subset R$ un scindage fort
de $R$ au-dessus de $u$. Le lemme \ref{stacks-morphisms-lemma-quasi-splitting-etale} montre que
$$
\mathcal{U} = [U/P] \longrightarrow [U/R] = \mathcal{X}
$$
est repr\'esentable par des espaces alg\'ebriques et \'etale. Puisque $P \subset R$
est ouvert et contient $\{r \in R : s(r) = u, t(r) = u\}$ par construction,
nous voyons que
$u \times_\mathcal{U} u \to u \times_\mathcal{X} u$ est un isomorphisme.
L'assertion sur les gerbes r\'esiduelles r\'esulte alors de
Propri\'et\'es des champs, lemme
\ref{stacks-properties-lemma-residual-gerbe-isomorphic}
(nous remarquons que les gerbes r\'esiduelles en question existent d'apr\`es
le lemme \ref{stacks-morphisms-lemma-every-point-residual-gerbe-quasi-DM}).
\end{proof}



\section{Morphismes lisses}
\label{stacks-morphisms-section-smooth}

\noindent
La propriété ``être lisse'' des morphismes d'espaces
algébriques est locale pour la topologie lisse sur la source et la cible, voir
Descente sur les espaces, Remarque \ref{spaces-descent-remark-list-local-source-target}.
Elle est également stable par changement de base et locale pour la topologie fpqc sur la cible, voir
Morphismes d'espaces,
Lemme \ref{spaces-morphisms-lemma-base-change-smooth}
et
Descente sur les espaces, Lemme
\ref{spaces-descent-lemma-descending-property-smooth}.
Par conséquent, d'après le
Lemme \ref{stacks-morphisms-lemma-local-source-target}
ci-dessus, nous pouvons définir comme suit ce que signifie être lisse pour un
morphisme de champs algébriques, et cette définition coïncide avec la notion
déjà définie dans
Propriétés des champs,
section \ref{stacks-properties-section-properties-morphisms}
lorsque le morphisme est représentable par des espaces algébriques.

\begin{definition}
\label{stacks-morphisms-definition-smooth}
Soit $f : \mathcal{X} \to \mathcal{Y}$ un morphisme de champs algébriques.
Nous disons que $f$ est {\it lisse} si les conditions équivalentes du
Lemme \ref{stacks-morphisms-lemma-local-source-target}
sont satisfaites avec $\mathcal{P} = \text{lisse}$.
\end{definition}

\begin{lemma}
\label{stacks-morphisms-lemma-composition-smooth}
Le composé de morphismes lisses est lisse.
\end{lemma}

\begin{proof}
Le résultat découle de
la Remarque \ref{stacks-morphisms-remark-composition}
et du
Lemme \ref{spaces-morphisms-lemma-composition-smooth}
de Morphismes d'espaces.
\end{proof}

\begin{lemma}
\label{stacks-morphisms-lemma-base-change-smooth}
Un changement de base d'un morphisme lisse est lisse.
\end{lemma}

\begin{proof}
Le résultat découle de
la Remarque \ref{stacks-morphisms-remark-base-change}
et du
Lemme \ref{spaces-morphisms-lemma-base-change-smooth}
de Morphismes d'espaces.
\end{proof}

\begin{lemma}
\label{stacks-morphisms-lemma-descent-smooth}
Soit $f : \mathcal{X} \to \mathcal{Y}$ un morphisme de champs algébriques.
Soit $\mathcal{Z} \to \mathcal{Y}$ un morphisme de champs algébriques surjectif,
plat et localement de présentation finie. Si le changement de base
$\mathcal{Z} \times_\mathcal{Y} \mathcal{X} \to \mathcal{Z}$
est lisse, alors $f$ est lisse.
\end{lemma}

\begin{proof}
La propriété ``être lisse''
satisfait aux conditions du Lemme \ref{stacks-morphisms-lemma-descent-property}.
Nous avons vu dans l'introduction de cette section qu'elle est locale pour la
topologie lisse sur la source et la cible ; le caractère local pour la topologie fppf sur la cible résulte de
Descente sur les espaces, Lemme
\ref{spaces-descent-lemma-descending-property-smooth}.
\end{proof}

\begin{lemma}
\label{stacks-morphisms-lemma-smooth-locally-finite-presentation}
Un morphisme lisse de champs algébriques est localement de présentation finie.
\end{lemma}

\begin{proof}
Démonstration omise.
\end{proof}

\begin{lemma}
\label{stacks-morphisms-lemma-where-smooth}
Soit $f : \mathcal{X} \to \mathcal{Y}$ un morphisme de champs algébriques.
Il existe un plus grand sous-champ ouvert $\mathcal{U} \subset \mathcal{X}$
tel que $f|_\mathcal{U} : \mathcal{U} \to \mathcal{Y}$ soit lisse.
En outre, la formation de cet ouvert commute aux opérations suivantes :
\begin{enumerate}
\item la précomposition par des morphismes lisses,
\item le changement de base par des morphismes plats et localement de
présentation finie,
\item le changement de base par des morphismes plats, pourvu que $f$ soit localement de
présentation finie.
\end{enumerate}
\end{lemma}

\begin{proof}
Choisissons un diagramme commutatif
$$
\xymatrix{
U \ar[d]_a \ar[r]_h & V \ar[d]^b \\
\mathcal{X} \ar[r]^f & \mathcal{Y}
}
$$
où $U$ et $V$ sont des espaces algébriques, les flèches verticales sont lisses
et $a : U \to \mathcal{X}$ est surjectif. Il existe un ouvert maximal
$U' \subset U$ tel que $h_{U'} : U' \to V$ soit lisse, voir
Morphismes d'espaces, Lemme \ref{spaces-morphisms-lemma-where-smooth}.
Soit $\mathcal{U} \subset \mathcal{X}$ le sous-champ ouvert
correspondant à l'image de $|U'| \to |\mathcal{X}|$
(Propriétés des champs, Lemmes \ref{stacks-properties-lemma-topology-points} et
\ref{stacks-properties-lemma-open-substacks}).
D'après l'équivalence du Lemme \ref{stacks-morphisms-lemma-local-source-target},
nous voyons que $f|_\mathcal{U} : \mathcal{U} \to \mathcal{Y}$ est lisse
et que $\mathcal{U}$ est le plus grand sous-champ ouvert possédant cette
propriété.

\medskip\noindent
L'assertion (1) résulte du fait que la propriété d'être lisse
est locale pour la topologie lisse sur la source (cette propriété a même servi à définir les
morphismes lisses de champs algébriques).
Les assertions (2) et (3) résultent du cas des espaces algébriques, voir
Morphismes d'espaces, Lemme \ref{spaces-morphisms-lemma-where-smooth}.
\end{proof}

\begin{lemma}
\label{stacks-morphisms-lemma-smooth-quotient-stack}
Soit $X \to Y$ un morphisme lisse d'espaces algébriques.
Soit $G$ un espace algébrique en groupes sur $Y$, plat
et localement de présentation finie sur $Y$. Soit donnée une action de $G$ sur $X$ au-dessus de $Y$.
Alors le champ quotient $[X/G]$ est lisse sur $Y$.
\end{lemma}

\noindent
Cela reste vrai même si $G$ n'est pas lisse sur $Y$ !

\begin{proof}
Le champ quotient $[X/G]$ est un champ algébrique d'après
Critères de représentabilité, Théorème
\ref{criteria-theorem-flat-groupoid-gives-algebraic-stack}.
La lissité de $[X/G]$ sur $Y$ résulte du fait que la lissité
descend par les recouvrements fppf :
Choisissons un morphisme lisse surjectif $U \to [X/G]$, où $U$ est un schéma.
La lissité de $[X/G]$ sur $Y$ équivaut à celle de $U$ sur $Y$.
Notons que $U \times_{[X/G]} X$ est lisse sur $X$, donc lisse
sur $Y$ (car un composé de morphismes lisses est lisse).
D'autre part, $U \times_{[X/G]} X \to U$ est localement de
présentation finie, plat et surjectif (car c'est
le changement de base de $X \to [X/G]$, qui possède ces propriétés,
par exemple d'après Critères de représentabilité, Lemme
\ref{criteria-lemma-flat-quotient-flat-presentation}).
Nous pouvons donc appliquer Descente sur les espaces,
Lemme \ref{spaces-descent-lemma-smooth-permanence}.
\end{proof}

\begin{lemma}
\label{stacks-morphisms-lemma-gerbe-smooth}
Soit $\pi : \mathcal{X} \to \mathcal{Y}$ un morphisme de champs algébriques.
Si $\mathcal{X}$ est une gerbe sur $\mathcal{Y}$, alors $\pi$ est surjectif
et lisse.
\end{lemma}

\begin{proof}
Nous avons établi la surjectivité au Lemme \ref{stacks-morphisms-lemma-gerbe-fppf}.
D'après le Lemme \ref{stacks-morphisms-lemma-descent-smooth},
il suffit de démontrer le lemme après avoir remplacé $\pi$ par son changement de base
par un morphisme surjectif, plat et localement de présentation finie
$\mathcal{Y}' \to \mathcal{Y}$. D'après le
Lemme \ref{stacks-morphisms-lemma-local-structure-gerbe},
nous pouvons supposer que $\mathcal{Y} = U$ est un espace algébrique et que
$\mathcal{X} = [U/G]$ au-dessus de $U$, avec $G \to U$ plat et
localement de présentation finie.
Le résultat découle alors du Lemme \ref{stacks-morphisms-lemma-smooth-quotient-stack}.
\end{proof}




\section[Morphismes DM : localité étale-lisse]{Types de morphismes dont la propriété est de nature locale pour la topologie \'etale sur la source et pour la topologie lisse sur le but}
\label{stacks-morphisms-section-etale-smooth-local-source-target}

\noindent
Soit une propriété de morphismes d'espaces algébriques qui est
{\it de nature locale pour la topologie \'etale sur la source et pour la topologie lisse sur le but}, voir
Descente sur les espaces,
Définition \ref{spaces-descent-definition-etale-smooth-local-source-target}.
Nous pouvons l'utiliser pour définir une propriété correspondante
des morphismes DM de champs algébriques, en imposant l'une ou l'autre
des conditions équivalentes du lemme ci-dessous.

\begin{lemma}
\label{stacks-morphisms-lemma-etale-smooth-local-source-target}
Soit $\mathcal{P}$ une propriété de morphismes d'espaces algébriques
qui est de nature locale pour la topologie \'etale sur la source et pour la topologie lisse sur le but.
Soit $f : \mathcal{X} \to \mathcal{Y}$ un morphisme DM de champs algébriques.
Considérons des diagrammes commutatifs
$$
\xymatrix{
U \ar[d]_a \ar[r]_h & V \ar[d]^b \\
\mathcal{X} \ar[r]^f & \mathcal{Y}
}
$$
où $U$ et $V$ sont des espaces algébriques, $V \to \mathcal{Y}$ est lisse,
et $U \to \mathcal{X} \times_\mathcal{Y} V$ est \'etale.
Les conditions suivantes sont équivalentes :
\begin{enumerate}
\item pour tout diagramme comme ci-dessus, le morphisme $h$ possède la propriété $\mathcal{P}$ ;
\item pour un diagramme comme ci-dessus dans lequel $a : U \to \mathcal{X}$ est surjectif,
le morphisme $h$ possède la propriété $\mathcal{P}$.
\end{enumerate}
Si $\mathcal{X}$ et $\mathcal{Y}$ sont représentables par des espaces algébriques,
cela équivaut également au fait que $f$ (considéré comme morphisme d'espaces algébriques)
possède la propriété $\mathcal{P}$. Si $\mathcal{P}$ est en outre stable par
tout changement de base et de nature locale pour la topologie fppf sur la base, alors, pour les morphismes $f$
qui sont représentables par des espaces algébriques,
cela équivaut également au fait que $f$ possède la propriété $\mathcal{P}$ au sens
de
Propriétés des champs,
section \ref{stacks-properties-section-properties-morphisms}.
\end{lemma}

\begin{proof}
Montrons l'implication (1) $\Rightarrow$ (2). Choisissons un espace
algébrique $V$ et un morphisme lisse surjectif $V \to \mathcal{Y}$.
Comme $f$ est DM, il existe un schéma $U$ et un morphisme \'etale surjectif
$U \to V \times_\mathcal{Y} \mathcal{X}$, voir
Lemme \ref{stacks-morphisms-lemma-DM}. Nous disposons ainsi des données géométriques requises en (2).
Notons que $U \to \mathcal{X}$ est également lisse et surjectif, puisqu'il est le
composé du changement de base
$\mathcal{X} \times_\mathcal{Y} V \to \mathcal{X}$ et du morphisme choisi
$U \to \mathcal{X} \times_\mathcal{Y} V$. Nous obtenons donc un
diagramme comme en (1). Par conséquent, si (1) est vérifiée, $h : U \to V$ possède la propriété
$\mathcal{P}$, ce qui signifie que (2) est vérifiée puisque $U \to \mathcal{X}$ est surjectif.

\medskip\noindent
Réciproquement, supposons que (2) soit vérifiée et soient $U, V, a, b, h$ comme en (2).
Soient ensuite $U', V', a', b', h'$ les données d'un diagramme quelconque comme en (1).
Considérons
$$
\xymatrix{
U \ar[d] \ar[r]_h & V \ar[d] \\
\mathcal{X} \ar[r]^f & \mathcal{Y}
}
\quad\quad
\xymatrix{
U' \ar[d] \ar[r]_{h'} & V' \ar[d] \\
\mathcal{X} \ar[r]^f & \mathcal{Y}
}
$$
Pour montrer que (2) implique (1), nous devons prouver que $h'$ possède $\mathcal{P}$.
À cette fin, considérons le diagramme commutatif
$$
\xymatrix{
U \ar[d]^h &
U \times_\mathcal{X} U' \ar[l] \ar[d]^{(h, h')} \ar[r] &
U' \ar[d]^{h'} \\
V &
V \times_\mathcal{Y} V' \ar[l] \ar[r] &
V'
}
$$
d'espaces algébriques. Notons que les flèches horizontales sont
lisses comme changements de base des morphismes lisses
$V \to \mathcal{Y}$, $V' \to \mathcal{Y}$, $U \to \mathcal{X}$ et
$U' \to \mathcal{X}$. Notons que les carrés
$$
\xymatrix{
U \ar[d] & U \times_\mathcal{X} U' \ar[l] \ar[d] &
U \times_\mathcal{X} U' \ar[d] \ar[r] & U' \ar[d] \\
V \times_\mathcal{Y} \mathcal{X} &
V \times_\mathcal{Y} U' \ar[l] &
U \times_\mathcal{Y} V' \ar[r] &
\mathcal{X} \times_\mathcal{Y} V'
}
$$
sont cartésiens ; les flèches verticales sont donc \'etales par nos hypothèses sur
$U', V', a', b', h'$ et $U, V, a, b, h$.
Puisque $\mathcal{P}$ est de nature locale pour la topologie lisse sur le but d'après
Descente sur les espaces, Lemme
\ref{spaces-descent-lemma-etale-smooth-local-source-target-implies}, partie (2),
nous voyons que le changement de base
$t : U \times_\mathcal{Y} V' \to V \times_\mathcal{Y} V'$ de $h$
possède $\mathcal{P}$. Puisque $\mathcal{P}$ est de nature locale pour la topologie \'etale sur la source d'après
Descente sur les espaces, Lemme
\ref{spaces-descent-lemma-etale-smooth-local-source-target-implies}, partie (1),
et que $s : U \times_\mathcal{X} U' \to U \times_\mathcal{Y} V'$ est \'etale,
nous voyons que le morphisme $(h, h') = t \circ s$ possède $\mathcal{P}$.
Considérons le diagramme
$$
\xymatrix{
U \times_\mathcal{X} U' \ar[r]_{(h, h')} \ar[d] &
V \times_\mathcal{Y} V' \ar[d] \\
U' \ar[r]^{h'} & V'
}
$$
La flèche verticale gauche est surjective, la flèche verticale droite est lisse et
le morphisme induit
$$
U \times_\mathcal{X} U'
\longrightarrow
U' \times_{V'} (V \times_\mathcal{Y} V') = V \times_\mathcal{Y} U'
$$
est \'etale, comme on l'a vu ci-dessus. Par conséquent, d'après
Descente sur les espaces, Définition
\ref{spaces-descent-definition-etale-smooth-local-source-target}, partie (3),
nous concluons que $h'$ possède $\mathcal{P}$. Ceci achève la démonstration de
l'équivalence de (1) et (2).

\medskip\noindent
Si $\mathcal{X}$ et $\mathcal{Y}$ sont représentables, alors
Descente sur les espaces,
Lemme \ref{spaces-descent-lemma-etale-smooth-local-source-target-characterize},
s'applique et montre que les conditions (1) et (2) équivalent au fait que $f$ possède
$\mathcal{P}$.

\medskip\noindent
Enfin, supposons $f$ représentable, soient $U, V, a, b, h$
comme dans la partie (2) du lemme, et supposons que $\mathcal{P}$ est stable par
changement de base arbitraire. Nous devons montrer que pour tout schéma
$Z$ et tout morphisme $Z \to \mathcal{Y}$, le changement de base
$Z \times_\mathcal{Y} \mathcal{X} \to Z$
possède la propriété $\mathcal{P}$. Considérons le diagramme
$$
\xymatrix{
Z \times_\mathcal{Y} U \ar[d] \ar[r] &
Z \times_\mathcal{Y} V \ar[d] \\
Z \times_\mathcal{Y} \mathcal{X} \ar[r] &
Z
}
$$
Notons que la flèche horizontale supérieure est un changement de base de $h$ et
possède donc la propriété $\mathcal{P}$. La flèche verticale gauche est
surjective, et le morphisme induit
$$
Z \times_\mathcal{Y} U \longrightarrow
(Z \times_\mathcal{Y} \mathcal{X}) \times_Z (Z \times_\mathcal{Y} V)
$$
est \'etale, tandis que la flèche verticale droite est lisse. Ainsi,
Descente sur les espaces,
Lemme \ref{spaces-descent-lemma-etale-smooth-local-source-target-characterize}
s'applique et montre que $Z \times_\mathcal{Y} \mathcal{X} \to Z$
possède la propriété $\mathcal{P}$.
\end{proof}

\begin{definition}
\label{stacks-morphisms-definition-etale-smooth-P}
Soit $\mathcal{P}$ une propriété de morphismes d'espaces algébriques
qui est de nature locale pour la topologie \'etale sur la source et pour la topologie lisse sur le but.
On dit qu'un morphisme DM $f : \mathcal{X} \to \mathcal{Y}$ de champs algébriques
{\it possède la propriété $\mathcal{P}$} si les conditions équivalentes du
Lemme \ref{stacks-morphisms-lemma-etale-smooth-local-source-target}
sont vérifiées.
\end{definition}

\begin{remark}
\label{stacks-morphisms-remark-etale-smooth-composition}
Soit $\mathcal{P}$ une propriété de morphismes d'espaces algébriques
qui est de nature locale pour la topologie \'etale sur la source et pour la topologie lisse sur le but, et
stable par composition. Alors la propriété des morphismes DM de champs algébriques
définie dans la Définition \ref{stacks-morphisms-definition-etale-smooth-P}
est stable par composition. En effet, soient $f : \mathcal{X} \to \mathcal{Y}$
et $g : \mathcal{Y} \to \mathcal{Z}$ des morphismes DM de champs algébriques
possédant la propriété $\mathcal{P}$. D'après le Lemme \ref{stacks-morphisms-lemma-composition-separated},
le composé $g \circ f$ est DM. Choisissons un espace algébrique $W$ et un
morphisme lisse surjectif $W \to \mathcal{Z}$. Choisissons un espace algébrique
$V$ et un morphisme \'etale surjectif $V \to \mathcal{Y} \times_\mathcal{Z} W$
(Lemme \ref{stacks-morphisms-lemma-DM}).
Choisissons un espace algébrique $U$ et un morphisme \'etale surjectif
$U \to \mathcal{X} \times_\mathcal{Y} V$. Alors les morphismes
$V \to W$ et $U \to V$ possèdent la propriété $\mathcal{P}$ par définition.
Par suite, $U \to W$ possède la propriété $\mathcal{P}$, puisque nous avons supposé que
$\mathcal{P}$ est stable par composition. Ainsi, par définition encore,
nous voyons que $g \circ f : \mathcal{X} \to \mathcal{Z}$ possède la
propriété $\mathcal{P}$.
\end{remark}

\begin{remark}
\label{stacks-morphisms-remark-etale-smooth-base-change}
Soit $\mathcal{P}$ une propriété de morphismes d'espaces algébriques
qui est de nature locale pour la topologie \'etale sur la source et pour la topologie lisse sur le but, et
stable par changement de base. Alors la propriété des
morphismes DM de champs algébriques définie dans la
Définition \ref{stacks-morphisms-definition-etale-smooth-P}
est stable par changement de base arbitraire.
En effet, soit $f : \mathcal{X} \to \mathcal{Y}$
un morphisme DM de champs algébriques,
soit $g : \mathcal{Y}' \to \mathcal{Y}$ un morphisme de champs algébriques,
et supposons que $f$ possède la propriété $\mathcal{P}$.
Alors le changement de base
$\mathcal{Y}' \times_\mathcal{Y} \mathcal{X} \to \mathcal{Y}'$
est un morphisme DM d'après le Lemme \ref{stacks-morphisms-lemma-base-change-separated}.
Choisissons un espace algébrique $V$
et un morphisme lisse surjectif $V \to \mathcal{Y}$. Choisissons un espace
algébrique $U$ et un morphisme \'etale surjectif
$U \to \mathcal{X} \times_\mathcal{Y} V$ (Lemme \ref{stacks-morphisms-lemma-DM}).
Enfin, choisissons un espace algébrique
$V'$ et un morphisme lisse surjectif
$V' \to \mathcal{Y}' \times_\mathcal{Y} V$. Alors le morphisme
$U \to V$ possède la propriété $\mathcal{P}$ par définition.
Par suite, $V' \times_V U \to V'$ possède la propriété $\mathcal{P}$, puisque nous avons supposé que
$\mathcal{P}$ est stable par changement de base. Considérons le diagramme
$$
\xymatrix{
V' \times_V U \ar[r] \ar[d] &
\mathcal{Y}' \times_\mathcal{Y} \mathcal{X} \ar[r] \ar[d] &
\mathcal{X} \ar[d] \\
V' \ar[r] & \mathcal{Y}' \ar[r] & \mathcal{Y}
}
$$
nous voyons que la flèche horizontale supérieure de gauche est surjective et que
$$
V' \times_V U \to V' \times_{\mathcal{Y}'}
(\mathcal{Y}' \times_\mathcal{Y} \mathcal{X}) =
V' \times_V (\mathcal{X} \times_\mathcal{Y} V)
$$
est \'etale comme changement de base de $U \to \mathcal{X} \times_\mathcal{Y} V$ ;
d'où, par définition, la projection
$\mathcal{Y}' \times_\mathcal{Y} \mathcal{X} \to \mathcal{Y}'$ possède la
propriété $\mathcal{P}$.
\end{remark}

\begin{remark}
\label{stacks-morphisms-remark-etale-smooth-implication}
Soient $\mathcal{P}, \mathcal{P}'$ des propriétés de morphismes d'espaces algébriques
qui sont de nature locale pour la topologie \'etale sur la source et pour la topologie lisse sur le but.
Supposons que l'on ait $\mathcal{P} \Rightarrow \mathcal{P}'$ pour les morphismes
d'espaces algébriques. Alors on a également $\mathcal{P} \Rightarrow \mathcal{P}'$
pour les propriétés des morphismes de champs algébriques définies dans la
Définition \ref{stacks-morphisms-definition-etale-smooth-P}
à partir de $\mathcal{P}$ et de $\mathcal{P}'$. Cela résulte immédiatement de la définition.
\end{remark}



\section{Morphismes \'etales}
\label{stacks-morphisms-section-etale}

\noindent
Un morphisme \'etale de champs alg\'ebriques ne saurait \^etre d\'efini comme un
morphisme lisse de dimension relative $0$. En effet, le morphisme
$$
[\mathbf{A}^1_k/\mathbf{G}_{m, k}] \longrightarrow \Spec(k)
$$
est lisse de dimension relative $0$, quelle que soit l'action choisie
du sch\'ema en groupes $\mathbf{G}_{m, k}$ sur $\mathbf{A}^1_k$.
Cela ne correspond pas \`a l'id\'ee usuelle selon laquelle les morphismes \'etales
devraient identifier les espaces tangents. L'exemple ci-dessus n'est pas quasi-fini,
mais le morphisme
$$
\mathcal{X} = [\Spec(k)/\mu_{p, k}] \longrightarrow \Spec(k)
$$
est lisse et quasi-fini (section \ref{stacks-morphisms-section-locally-quasi-finite}).
Toutefois, si la caract\'eristique de $k$ est $p > 0$, on constate
que le morphisme repr\'esentable $\Spec(k) \to \mathcal{X}$
n'est pas \'etale, puisque son changement de base
$\mu_{p, k} = \Spec(k) \times_\mathcal{X} \Spec(k) \to \Spec(k)$
est un morphisme d'un sch\'ema non r\'eduit vers le spectre d'un corps.
Ainsi, si l'on d\'efinissait un morphisme \'etale comme un morphisme lisse et localement quasi-fini,
alors l'analogue du lemme de Morphismes d'espaces
\ref{spaces-morphisms-lemma-etale-permanence} serait faux.

\medskip\noindent
Notre d\'emarche consistera plut\^ot \`a partir des exigences suivantes :
``\^etre \'etale'' devra \^etre une propri\'et\'e stable par changement de base et, si
$\mathcal{X} \to X$ est un morphisme \'etale d'un champ alg\'ebrique
vers un sch\'ema, alors $\mathcal{X}$ devra \^etre de Deligne--Mumford.
Autrement dit, nous exigerons qu'un morphisme \'etale soit DM et
nous utiliserons les r\'esultats de la
section \ref{stacks-morphisms-section-etale-smooth-local-source-target} pour d\'efinir
les morphismes \'etales de champs alg\'ebriques.

\medskip\noindent
Dans le lemme \ref{stacks-morphisms-lemma-characterize-etale}, nous caract\'eriserons
les morphismes \'etales de champs alg\'ebriques comme ceux qui sont
(a) localement de pr\'esentation finie,
(b) plats et (c) \`a diagonale \'etale.

\medskip\noindent
La propri\'et\'e ``\^etre \'etale'' des morphismes d'espaces alg\'ebriques
est de nature locale pour la topologie \'etale sur la source et pour la topologie lisse
sur le but; voir Descente sur les espaces, remarque
\ref{spaces-descent-remark-list-etale-smooth-local-source-target}.
Elle est aussi stable par changement de base et de nature locale pour la topologie fpqc sur le but; voir
Morphismes d'espaces,
lemme \ref{spaces-morphisms-lemma-base-change-etale}
et
Descente sur les espaces,
lemme \ref{spaces-descent-lemma-descending-property-etale}.
Par cons\'equent, d'apr\`es
le lemme \ref{stacks-morphisms-lemma-etale-smooth-local-source-target}
ci-dessus, nous pouvons d\'efinir comme suit ce que signifie, pour un morphisme
de champs alg\'ebriques, \^etre \'etale; cette d\'efinition co\"incide avec la notion
d\'ej\`a d\'efinie dans Propri\'et\'es des champs, section
\ref{stacks-properties-section-properties-morphisms},
lorsque le morphisme est repr\'esentable par des espaces alg\'ebriques,
car un tel morphisme est automatiquement DM d'apr\`es
le lemme \ref{stacks-morphisms-lemma-trivial-implications}.

\begin{definition}
\label{stacks-morphisms-definition-etale}
Soit $f : \mathcal{X} \to \mathcal{Y}$ un morphisme de champs alg\'ebriques.
On dit que $f$ est {\it \'etale} si $f$ est DM et si les conditions \'equivalentes du
lemme \ref{stacks-morphisms-lemma-etale-smooth-local-source-target}
sont v\'erifi\'ees avec $\mathcal{P} = \etale$.
\end{definition}

\noindent
Nous utiliserons sans autre mention le fait que cette d\'efinition co\"incide avec la notion
d\'ej\`a existante de morphisme \'etale lorsque $f$ est repr\'esentable
par des espaces alg\'ebriques ou lorsque $\mathcal{X}$ et $\mathcal{Y}$ sont
repr\'esentables par des espaces alg\'ebriques.

\begin{lemma}
\label{stacks-morphisms-lemma-composition-etale}
Le compos\'e de deux morphismes \'etales est \'etale.
\end{lemma}

\begin{proof}
Cela résulte de la remarque \ref{stacks-morphisms-remark-etale-smooth-composition}
et du lemme \ref{spaces-morphisms-lemma-composition-etale} de
Morphismes d'espaces.
\end{proof}

\begin{lemma}
\label{stacks-morphisms-lemma-base-change-etale}
Tout changement de base d'un morphisme \'etale est \'etale.
\end{lemma}

\begin{proof}
Cela résulte de la remarque \ref{stacks-morphisms-remark-etale-smooth-base-change}
et du lemme \ref{spaces-morphisms-lemma-base-change-etale} de
Morphismes d'espaces.
\end{proof}

\begin{lemma}
\label{stacks-morphisms-lemma-open-immersion-etale}
Toute immersion ouverte est \'etale.
\end{lemma}

\begin{proof}
Soit $j : \mathcal{U} \to \mathcal{X}$ une immersion ouverte de
champs alg\'ebriques. Puisque $j$ est repr\'esentable, il est DM d'apr\`es
le lemme \ref{stacks-morphisms-lemma-trivial-implications}. D'autre part,
si $X \to \mathcal{X}$ est un morphisme lisse et surjectif,
o\`u $X$ est un sch\'ema, alors $U = \mathcal{U} \times_\mathcal{X} X$
est un sous-sch\'ema ouvert de $X$. Par cons\'equent, $U \to X$ est \'etale
(Morphismes, lemme \ref{morphisms-lemma-open-immersion-etale})
et la d\'efinition montre que $j$ est \'etale.
\end{proof}

\begin{lemma}
\label{stacks-morphisms-lemma-etale}
Soit $f : \mathcal{X} \to \mathcal{Y}$ un morphisme de champs alg\'ebriques.
Les conditions suivantes sont \'equivalentes :
\begin{enumerate}
\item $f$ est \'etale,
\item $f$ est DM et, pour tout morphisme $V \to \mathcal{Y}$
o\`u $V$ est un espace alg\'ebrique, et tout morphisme \'etale
$U \to V \times_\mathcal{Y} \mathcal{X}$ o\`u $U$ est un espace alg\'ebrique,
le morphisme $U \to V$ est \'etale,
\item il existe un morphisme surjectif, localement de pr\'esentation finie et plat
$W \to \mathcal{Y}$, o\`u $W$ est un espace alg\'ebrique, et un
morphisme \'etale surjectif $T \to W \times_\mathcal{Y} \mathcal{X}$
o\`u $T$ est un espace alg\'ebrique, tel que le morphisme $T \to W$ soit \'etale.
\end{enumerate}
\end{lemma}

\begin{proof}
Supposons (1). Alors $f$ est DM et, comme la propri\'et\'e d'\^etre \'etale est stable
par changement de base, la condition (2) est v\'erifi\'ee.

\medskip\noindent
Supposons (2). Choisissons un sch\'ema $V$ et un morphisme \'etale surjectif
$V \to \mathcal{Y}$. Choisissons un sch\'ema $U$ et un morphisme \'etale surjectif
$U \to V \times_\mathcal{Y} \mathcal{X}$ (lemme \ref{stacks-morphisms-lemma-DM}).
La condition (3) est donc v\'erifi\'ee.

\medskip\noindent
Supposons que $W \to \mathcal{Y}$ et $T \to W \times_\mathcal{Y} \mathcal{X}$
soient comme dans (3). V\'erifions d'abord que $f$ est DM. En effet, il suffit de v\'erifier
que $W \times_\mathcal{Y} \mathcal{X} \to W$ est DM; voir
le lemme \ref{stacks-morphisms-lemma-check-separated-covering}.
D'apr\`es le lemme \ref{stacks-morphisms-lemma-compose-after-separated},
il suffit de v\'erifier que $W \times_\mathcal{Y} \mathcal{X}$ est de Deligne--Mumford.
Cela r\'esulte de l'existence de $T \to W \times_\mathcal{Y} \mathcal{X}$
par le sens facile du th\'eor\`eme \ref{stacks-morphisms-theorem-DM}.

\medskip\noindent
Supposons que $f$ soit DM et que $W \to \mathcal{Y}$ et
$T \to W \times_\mathcal{Y} \mathcal{X}$ soient comme dans (3).
Soit $V$ un espace alg\'ebrique, soit $V \to \mathcal{Y}$ un morphisme lisse surjectif,
soit $U$ un espace alg\'ebrique, et soit
$U \to V \times_\mathcal{Y} \mathcal{X}$ un morphisme \'etale surjectif
(lemme \ref{stacks-morphisms-lemma-DM}). Il faut v\'erifier que $U \to V$ est \'etale.
Il suffit de montrer que $U \times_\mathcal{Y} W \to V \times_\mathcal{Y} W$
est \'etale d'apr\`es Descente sur les espaces, lemme
\ref{spaces-descent-lemma-descending-property-etale}.
Nous pouvons remplacer $\mathcal{X}, \mathcal{Y}, W, T, U, V$ respectivement par
$\mathcal{X} \times_\mathcal{Y} W, W, W, T, U \times_\mathcal{Y} W,
V \times_\mathcal{Y} W$ (un petit d\'etail est omis).
Nous pouvons donc supposer que $Y = \mathcal{Y}$ est un espace alg\'ebrique, qu'il existe
un espace alg\'ebrique $T$ et un morphisme \'etale surjectif
$T \to \mathcal{X}$ tel que $T \to Y$ soit \'etale, et que $U$ et $V$
soient comme ci-dessus. Dans ce cas, on sait que
$$
U \to V\text{ est \'etale}
\Leftrightarrow
\mathcal{X} \to Y\text{ est \'etale}
\Leftrightarrow
T \to Y\text{ est \'etale}
$$
par l'\'equivalence des propri\'et\'es (1) et (2) du
lemme \ref{stacks-morphisms-lemma-etale-smooth-local-source-target}
et la d\'efinition \ref{stacks-morphisms-definition-etale}.
Cela ach\`eve la d\'emonstration.
\end{proof}

\begin{lemma}
\label{stacks-morphisms-lemma-etale-permanence}
Soient $\mathcal{X}, \mathcal{Y}$ des champs alg\'ebriques \'etales sur
un champ alg\'ebrique $\mathcal{Z}$. Tout morphisme
$\mathcal{X} \to \mathcal{Y}$ au-dessus de $\mathcal{Z}$ est \'etale.
\end{lemma}

\begin{proof}
Le morphisme $\mathcal{X} \to \mathcal{Y}$ est DM d'apr\`es
le lemme \ref{stacks-morphisms-lemma-compose-after-separated}.
Soit $W \to \mathcal{Z}$ un morphisme lisse surjectif
dont la source est un espace alg\'ebrique. Soit
$V \to \mathcal{Y} \times_\mathcal{Z} W$ un morphisme \'etale surjectif
dont la source est un espace alg\'ebrique
(lemme \ref{stacks-morphisms-lemma-DM}).
Soit $U \to \mathcal{X} \times_\mathcal{Y} V$ un morphisme \'etale surjectif
dont la source est un espace alg\'ebrique
(lemme \ref{stacks-morphisms-lemma-DM}).
Alors
$$
U \longrightarrow \mathcal{X} \times_\mathcal{Z} W
$$
est \'etale surjectif en tant que compos\'e de
$U \to \mathcal{X} \times_\mathcal{Y} V$
et du changement de base de $V \to \mathcal{Y} \times_\mathcal{Z} W$
par $\mathcal{X} \times_\mathcal{Z} W \to \mathcal{Y} \times_\mathcal{Z} W$.
Il suffit donc de montrer que $U \to V$ est \'etale.
Or $U \to W$ et $V \to W$ sont \'etales, puisque
$\mathcal{X} \to \mathcal{Z}$ et $\mathcal{Y} \to \mathcal{Z}$
le sont; la conclusion r\'esulte donc de la version du lemme pour les
espaces alg\'ebriques, \`a savoir
Morphismes d'espaces, lemme \ref{spaces-morphisms-lemma-etale-permanence}.
\end{proof}








\section{Morphismes non ramifi\'es}
\label{stacks-morphisms-section-unramified}

\noindent
Pour justifier notre choix de la d\'efinition des morphismes
non ramifi\'es, nous renvoyons le lecteur \`a la discussion
de la section sur les morphismes \'etales, section \ref{stacks-morphisms-section-etale}.

\medskip\noindent
Dans le Lemme \ref{stacks-morphisms-lemma-characterize-unramified}, nous caract\'eriserons
les morphismes non ramifi\'es de champs alg\'ebriques comme les morphismes
qui sont localement de type fini et dont la diagonale est \'etale.

\medskip\noindent
La propri\'et\'e ``non ramifi\'e'' des morphismes d'espaces alg\'ebriques est de
nature locale pour la topologie \'etale sur la source et pour la topologie lisse sur le but, voir
Descente sur les espaces, Remarque
\ref{spaces-descent-remark-list-etale-smooth-local-source-target}.
Elle est aussi stable par changement de base et locale sur la cible pour la
topologie fpqc, voir Morphismes d'espaces,
Lemme \ref{spaces-morphisms-lemma-base-change-unramified}
et
Descente sur les espaces,
Lemme \ref{spaces-descent-lemma-descending-property-unramified}.
Par cons\'equent, d'apr\`es le
Lemme \ref{stacks-morphisms-lemma-etale-smooth-local-source-target}
ci-dessus, nous pouvons d\'efinir comme suit ce que signifie, pour un morphisme
de champs alg\'ebriques, d'\^etre non ramifi\'e; cette d\'efinition co\"incide avec
la notion d\'ej\`a d\'efinie dans Propri\'et\'es des champs, section
\ref{stacks-properties-section-properties-morphisms},
lorsque le morphisme est repr\'esentable par des espaces alg\'ebriques,
car un tel morphisme est automatiquement DM d'apr\`es le
Lemme \ref{stacks-morphisms-lemma-trivial-implications}.

\begin{definition}
\label{stacks-morphisms-definition-unramified}
Soit $f : \mathcal{X} \to \mathcal{Y}$ un morphisme de champs alg\'ebriques.
On dit que $f$ est {\it non ramifi\'e} si $f$ est DM et si les conditions
\'equivalentes du Lemme \ref{stacks-morphisms-lemma-etale-smooth-local-source-target}
sont satisfaites pour $\mathcal{P} =$``non ramifi\'e''.
\end{definition}

\noindent
Nous utiliserons sans autre mention le fait que cette d\'efinition co\"incide
avec la notion d\'ej\`a existante de morphisme non ramifi\'e lorsque $f$ est
repr\'esentable par des espaces alg\'ebriques, ou lorsque $\mathcal{X}$ et
$\mathcal{Y}$ sont repr\'esentables par des espaces alg\'ebriques.

\begin{lemma}
\label{stacks-morphisms-lemma-composition-unramified}
Le compos\'e de morphismes non ramifi\'es est non ramifi\'e.
\end{lemma}

\begin{proof}
Cela r\'esulte de la Remarque \ref{stacks-morphisms-remark-etale-smooth-composition}
et de
Morphismes d'espaces, Lemme \ref{spaces-morphisms-lemma-composition-unramified}.
\end{proof}

\begin{lemma}
\label{stacks-morphisms-lemma-base-change-unramified}
Tout morphisme d\'eduit par changement de base d'un morphisme non ramifi\'e est non ramifi\'e.
\end{lemma}

\begin{proof}
Cela r\'esulte de
la Remarque \ref{stacks-morphisms-remark-etale-smooth-base-change}
et de
Morphismes d'espaces, Lemme \ref{spaces-morphisms-lemma-base-change-unramified}.
\end{proof}

\begin{lemma}
\label{stacks-morphisms-lemma-etale-unramified}
Un morphisme \'etale est non ramifi\'e.
\end{lemma}

\begin{proof}
Cela r\'esulte de la Remarque \ref{stacks-morphisms-remark-etale-smooth-implication}
et de Morphismes d'espaces, Lemme
\ref{spaces-morphisms-lemma-etale-unramified}.
\end{proof}

\begin{lemma}
\label{stacks-morphisms-lemma-immersion-unramified}
Une immersion est non ramifi\'ee.
\end{lemma}

\begin{proof}
Soit $j : \mathcal{Z} \to \mathcal{X}$ une immersion de
champs alg\'ebriques. Comme $j$ est repr\'esentable, il est DM d'apr\`es le
Lemme \ref{stacks-morphisms-lemma-trivial-implications}. D'autre part,
si $X \to \mathcal{X}$ est un morphisme lisse et surjectif,
o\`u $X$ est un sch\'ema, alors $Z = \mathcal{Z} \times_\mathcal{X} X$
est un sous-sch\'ema localement ferm\'e de $X$. Par cons\'equent, $Z \to X$
est non ramifi\'e (Morphismes, Lemmes \ref{morphisms-lemma-open-immersion-unramified} et
\ref{morphisms-lemma-closed-immersion-unramified}),
et la d\'efinition montre que $j$ est non ramifi\'e.
\end{proof}

\begin{lemma}
\label{stacks-morphisms-lemma-unramified}
Soit $f : \mathcal{X} \to \mathcal{Y}$ un morphisme de champs alg\'ebriques.
Les conditions suivantes sont \'equivalentes :
\begin{enumerate}
\item $f$ est non ramifi\'e,
\item $f$ est DM et, pour tout morphisme $V \to \mathcal{Y}$
o\`u $V$ est un espace alg\'ebrique, et tout morphisme \'etale
$U \to V \times_\mathcal{Y} \mathcal{X}$ o\`u $U$ est un espace alg\'ebrique,
le morphisme $U \to V$ est non ramifi\'e,
\item il existe un morphisme $W \to \mathcal{Y}$ surjectif,
localement de pr\'esentation finie et plat, o\`u $W$ est un espace alg\'ebrique,
et un morphisme $T \to W \times_\mathcal{Y} \mathcal{X}$ surjectif et \'etale,
o\`u $T$ est un espace alg\'ebrique, tel que le morphisme $T \to W$ soit non ramifi\'e.
\end{enumerate}
\end{lemma}

\begin{proof}
Supposons (1). Alors $f$ est DM et, puisque la propri\'et\'e d'\^etre non
ramifi\'e est stable par changement de base, la condition (2) est satisfaite.

\medskip\noindent
Supposons (2). Choisissons un sch\'ema $V$ et un morphisme \'etale surjectif
$V \to \mathcal{Y}$. Choisissons un sch\'ema $U$ et un morphisme \'etale surjectif
$U \to V \times_\mathcal{Y} \mathcal{X}$ (Lemme \ref{stacks-morphisms-lemma-DM}).
La condition (3) est donc satisfaite.

\medskip\noindent
Supposons que $W \to \mathcal{Y}$ et $T \to W \times_\mathcal{Y} \mathcal{X}$
soient comme en (3). Montrons d'abord que $f$ est DM. Il suffit en effet de
v\'erifier que $W \times_\mathcal{Y} \mathcal{X} \to W$ est DM, voir le
Lemme \ref{stacks-morphisms-lemma-check-separated-covering}.
D'apr\`es le Lemme \ref{stacks-morphisms-lemma-compose-after-separated},
il suffit de v\'erifier que $W \times_\mathcal{Y} \mathcal{X}$ est un champ
de Deligne--Mumford. Cela r\'esulte de l'existence de
$T \to W \times_\mathcal{Y} \mathcal{X}$ par le sens facile du
Th\'eor\`eme \ref{stacks-morphisms-theorem-DM}.

\medskip\noindent
Supposons que $f$ soit DM et que $W \to \mathcal{Y}$ et
$T \to W \times_\mathcal{Y} \mathcal{X}$ soient comme en (3).
Soit $V$ un espace alg\'ebrique, soit $V \to \mathcal{Y}$ un morphisme
lisse surjectif, soit $U$ un espace alg\'ebrique, et soit
$U \to V \times_\mathcal{Y} \mathcal{X}$ un morphisme surjectif et \'etale
(Lemme \ref{stacks-morphisms-lemma-DM}). Nous devons v\'erifier que $U \to V$ est non ramifi\'e.
Il suffit de montrer que $U \times_\mathcal{Y} W \to V \times_\mathcal{Y} W$
est non ramifi\'e d'apr\`es Descente sur les espaces, Lemme
\ref{spaces-descent-lemma-descending-property-unramified}.
Nous pouvons remplacer $\mathcal{X}, \mathcal{Y}, W, T, U, V$ par
$\mathcal{X} \times_\mathcal{Y} W, W, W, T, U \times_\mathcal{Y} W,
V \times_\mathcal{Y} W$ (un d\'etail mineur est omis).
Nous pouvons donc supposer que $Y = \mathcal{Y}$ est un espace alg\'ebrique,
qu'il existe un espace alg\'ebrique $T$ et un morphisme \'etale surjectif
$T \to \mathcal{X}$ tel que $T \to Y$ soit non ramifi\'e, et que $U$ et $V$
soient comme ci-dessus. Dans ce cas, nous savons que
$$
U \to V\text{ est non ramifi\'e}
\Leftrightarrow
\mathcal{X} \to Y\text{ est non ramifi\'e}
\Leftrightarrow
T \to Y\text{ est non ramifi\'e}
$$
par l'\'equivalence des propri\'et\'es (1) et (2) du
Lemme \ref{stacks-morphisms-lemma-etale-smooth-local-source-target}
et par la D\'efinition \ref{stacks-morphisms-definition-unramified}.
Cela ach\`eve la d\'emonstration.
\end{proof}

\begin{lemma}
\label{stacks-morphisms-lemma-unramified-quasi-finite}
Un morphisme non ramifi\'e de champs alg\'ebriques est localement quasi-fini.
\end{lemma}

\begin{proof}
Cela r\'esulte du
Lemme \ref{stacks-morphisms-lemma-unramified} (qui caract\'erise les morphismes non ramifi\'es),
du Lemme \ref{stacks-morphisms-lemma-characterize-locally-quasi-finite} (qui caract\'erise
les morphismes localement quasi-finis), et de
Morphismes d'espaces, Lemme \ref{spaces-morphisms-lemma-unramified-quasi-finite}
(qui est l'\'enonc\'e correspondant pour les espaces alg\'ebriques).
\end{proof}

\begin{lemma}
\label{stacks-morphisms-lemma-permanence-unramified}
Soient $\mathcal{X} \to \mathcal{Y} \to \mathcal{Z}$ des
morphismes de champs alg\'ebriques.
Si $\mathcal{X} \to \mathcal{Z}$ est non ramifi\'e et si
$\mathcal{Y} \to \mathcal{Z}$ est DM, alors
$\mathcal{X} \to \mathcal{Y}$ est non ramifi\'e.
\end{lemma}

\begin{proof}
Supposons que $\mathcal{X} \to \mathcal{Z}$ soit non ramifi\'e. D'apr\`es le
Lemme \ref{stacks-morphisms-lemma-compose-after-separated}, le morphisme
$\mathcal{X} \to \mathcal{Y}$ est DM. Choisissons un diagramme commutatif
$$
\xymatrix{
U \ar[d] \ar[r] & V \ar[d] \ar[r] & W \ar[d] \\
\mathcal{X} \ar[r] & \mathcal{Y} \ar[r] & \mathcal{Z}
}
$$
o\`u $U, V, W$ sont des espaces alg\'ebriques,
o\`u $W \to \mathcal{Z}$ est lisse et surjectif,
$V \to \mathcal{Y} \times_\mathcal{Z} W$ est \'etale et surjectif, et
$U \to \mathcal{X} \times_\mathcal{Y} V$ est \'etale et surjectif
(voir le Lemme \ref{stacks-morphisms-lemma-DM}). Alors
$U \to \mathcal{X} \times_\mathcal{Z} W$ est lui aussi surjectif et \'etale.
Nous savons donc que $U \to W$ est non ramifi\'e, et il nous faut montrer que
$U \to V$ est non ramifi\'e. Cela r\'esulte de
Morphismes d'espaces, Lemme \ref{spaces-morphisms-lemma-permanence-unramified}.
\end{proof}

\begin{lemma}
\label{stacks-morphisms-lemma-characterize-unramified}
Soit $f : \mathcal{X} \to \mathcal{Y}$ un morphisme de champs alg\'ebriques.
Les conditions suivantes sont \'equivalentes :
\begin{enumerate}
\item $f$ est non ramifi\'e, et
\item $f$ est localement de type fini et sa diagonale est \'etale.
\end{enumerate}
\end{lemma}

\begin{proof}
Supposons que $f$ soit non ramifi\'e. Alors $f$ est DM; nous pouvons donc choisir
des espaces alg\'ebriques $U$, $V$, un morphisme lisse surjectif
$V \to \mathcal{Y}$ et un morphisme \'etale surjectif
$U \to \mathcal{X} \times_\mathcal{Y} V$ (Lemme \ref{stacks-morphisms-lemma-DM}).
Comme $f$ est non ramifi\'e, le morphisme induit $U \to V$ est non ramifi\'e.
Ainsi, $U \to V$ est localement de type fini
(Morphismes d'espaces, Lemme
\ref{spaces-morphisms-lemma-unramified-locally-finite-type}),
et nous en concluons que $f$ est localement de type fini. La diagonale
$\Delta : \mathcal{X} \to \mathcal{X} \times_\mathcal{Y} \mathcal{X}$
est un morphisme de champs alg\'ebriques au-dessus de $\mathcal{Y}$.
Le changement de base de $\Delta$ par le morphisme lisse surjectif
$V \to \mathcal{Y}$ est la diagonale du changement de base de
$f$, c'est-\`a-dire de $\mathcal{X}_V = \mathcal{X} \times_\mathcal{Y} V \to V$.
Autrement dit, le diagramme
$$
\xymatrix{
\mathcal{X}_V \ar[r] \ar[d] &  \mathcal{X}_V \times_V \mathcal{X}_V \ar[d] \\
\mathcal{X} \ar[r] & \mathcal{X} \times_\mathcal{Y} \mathcal{X}
}
$$
est cart\'esien. Comme la fl\`eche verticale de droite est lisse et surjective,
il suffit de montrer que la fl\`eche horizontale sup\'erieure est \'etale, d'apr\`es
Propri\'et\'es des champs, Lemme
\ref{stacks-properties-lemma-check-property-weak-covering}.
Consid\'erons le diagramme commutatif
$$
\xymatrix{
U \ar[d] \ar[r] & U \times_V U \ar[d] \\
\mathcal{X}_V \ar[r] & \mathcal{X}_V \times_V \mathcal{X}_V
}
$$
Toutes les fl\`eches sont repr\'esentables par des espaces alg\'ebriques,
les fl\`eches verticales sont \'etales, celle de gauche est surjective, et
la fl\`eche horizontale sup\'erieure est une immersion ouverte d'apr\`es
Morphismes d'espaces, Lemme
\ref{spaces-morphisms-lemma-diagonal-unramified-morphism}.
Cela donne le r\'esultat voulu : tout d'abord,
$U \to \mathcal{X}_V \times_V \mathcal{X}_V$ est \'etale
comme compos\'e de morphismes \'etales; nous pouvons ensuite appliquer
Propri\'et\'es des champs, Lemme
\ref{stacks-properties-lemma-check-property-after-precomposing}
pour voir que $\mathcal{X}_V \to \mathcal{X}_V \times_V \mathcal{X}_V$
est \'etale, puisque la propri\'et\'e d'\^etre \'etale (pour les morphismes
d'espaces alg\'ebriques) est locale sur la source pour la topologie \'etale
(Descente sur les espaces, Lemme \ref{spaces-descent-lemma-etale-etale-local-source}).

\medskip\noindent
Supposons que $f$ soit localement de type fini et que sa diagonale soit \'etale.
Alors $f$ est DM par d\'efinition (puisque les morphismes \'etales d'espaces
alg\'ebriques sont non ramifi\'es). Comme ci-dessus, cela signifie que nous
pouvons choisir des espaces alg\'ebriques $U$, $V$, un morphisme lisse surjectif
$V \to \mathcal{Y}$ et un morphisme \'etale surjectif
$U \to \mathcal{X} \times_\mathcal{Y} V$ (Lemme \ref{stacks-morphisms-lemma-DM}).
Pour achever la d\'emonstration, nous devons montrer que $U \to V$ est non ramifi\'e.
Nous savons d\'ej\`a que $U \to V$ est localement de type fini.
En raisonnant comme ci-dessus, nous obtenons un diagramme commutatif
$$
\xymatrix{
U \ar[d] \ar[r] & U \times_V U \ar[d] \\
\mathcal{X}_V \ar[r] & \mathcal{X}_V \times_V \mathcal{X}_V
}
$$
o\`u toutes les fl\`eches sont repr\'esentables par des espaces alg\'ebriques,
les fl\`eches verticales sont \'etales, et la fl\`eche horizontale inf\'erieure
est \'etale comme changement de base de $\Delta$.
Il s'ensuit que $U \to U \times_V U$ est \'etale,
par exemple d'apr\`es le Lemme \ref{stacks-morphisms-lemma-etale-permanence}\footnote{On peut
aussi le d\'eduire directement, et assez facilement, de
Morphismes d'espaces, Lemme \ref{spaces-morphisms-lemma-etale-permanence}.}.
Ainsi, $U \to U \times_V U$ est un monomorphisme \'etale,
donc une immersion ouverte (Morphismes d'espaces, Lemme
\ref{spaces-morphisms-lemma-etale-universally-injective-open}).
Le morphisme $U \to V$ est alors non ramifi\'e d'apr\`es
Morphismes d'espaces, Lemme
\ref{spaces-morphisms-lemma-diagonal-unramified-morphism}.
\end{proof}

\begin{lemma}
\label{stacks-morphisms-lemma-characterize-etale}
Soit $f : \mathcal{X} \to \mathcal{Y}$ un morphisme de champs alg\'ebriques.
Les conditions suivantes sont \'equivalentes :
\begin{enumerate}
\item $f$ est \'etale, et
\item $f$ est localement de pr\'esentation finie, plat et non ramifi\'e,
\item $f$ est localement de pr\'esentation finie, plat, et sa diagonale
est \'etale.
\end{enumerate}
\end{lemma}

\begin{proof}
L'\'equivalence de (2) et (3) r\'esulte imm\'ediatement du
Lemme \ref{stacks-morphisms-lemma-characterize-unramified}. Ainsi, dans chaque cas,
le morphisme $f$ est DM.
Nous pouvons donc choisir des espaces alg\'ebriques $U$, $V$,
un morphisme lisse surjectif
$V \to \mathcal{Y}$ et un morphisme \'etale surjectif
$U \to \mathcal{X} \times_\mathcal{Y} V$ (Lemme \ref{stacks-morphisms-lemma-DM}).
Pour achever la d\'emonstration, nous devons montrer que
$U \to V$ est \'etale si et seulement s'il est localement de pr\'esentation
finie, plat et non ramifi\'e.
Cela r\'esulte de Morphismes d'espaces, Lemme
\ref{spaces-morphisms-lemma-unramified-flat-lfp-etale}
(et, plus imm\'ediatement, de Morphismes d'espaces, Lemmes
\ref{spaces-morphisms-lemma-etale-unramified},
\ref{spaces-morphisms-lemma-etale-locally-finite-presentation} et
\ref{spaces-morphisms-lemma-etale-flat}).
\end{proof}






\section{Morphismes propres}
\label{stacks-morphisms-section-proper}

\noindent
La notion de morphisme propre joue un rôle important en géométrie algébrique.
Voici la définition d'un morphisme propre de champs algébriques.

\begin{definition}
\label{stacks-morphisms-definition-proper}
Soit $f : \mathcal{X} \to \mathcal{Y}$
un morphisme de champs algébriques.
On dit que $f$ est {\it propre} si $f$ est séparé, de type fini et
universellement fermé.
\end{definition}

\noindent
Cette définition est compatible avec la notion déjà existante de morphisme propre
d'espaces algébriques : un morphisme d'espaces algébriques est propre si et seulement si
il est séparé, de type fini et universellement fermé
(Morphismes d'espaces, Définition \ref{spaces-morphisms-definition-proper}),
et nous avons déjà vérifié la compatibilité de ces notions dans le
Lemme \ref{stacks-morphisms-lemma-representable-separated-diagonal-closed},
la section \ref{stacks-morphisms-section-finite-type} et le
Lemme \ref{stacks-morphisms-lemma-characterize-representable-universally-closed}.
De même, si $f : \mathcal{X} \to \mathcal{Y}$ est un
morphisme de champs algébriques représentable par des espaces algébriques,
nous avons défini ce que signifie pour $f$ d'être propre dans
Propriétés des champs, section
\ref{stacks-properties-section-properties-morphisms}.
Cependant, la discussion de cette section montre que cela
équivaut à exiger que $f$ soit séparé, de type fini et
universellement fermé, et les mêmes références que ci-dessus établissent la
compatibilité.

\begin{lemma}
\label{stacks-morphisms-lemma-base-change-proper}
Tout changement de base d'un morphisme propre est propre.
\end{lemma}

\begin{proof}
Voir les
Lemmes \ref{stacks-morphisms-lemma-base-change-separated},
\ref{stacks-morphisms-lemma-base-change-finite-type} et
\ref{stacks-morphisms-lemma-base-change-universally-closed}.
\end{proof}

\begin{lemma}
\label{stacks-morphisms-lemma-composition-proper}
Le composé de morphismes propres est propre.
\end{lemma}

\begin{proof}
Voir les
Lemmes \ref{stacks-morphisms-lemma-composition-separated},
\ref{stacks-morphisms-lemma-composition-finite-type} et
\ref{stacks-morphisms-lemma-composition-universally-closed}.
\end{proof}

\begin{lemma}
\label{stacks-morphisms-lemma-closed-immersion-proper}
Une immersion fermée de champs algébriques est un morphisme propre de
champs algébriques.
\end{lemma}

\begin{proof}
Une immersion fermée est, par définition, représentable
(Propriétés des champs, Définition
\ref{stacks-properties-definition-immersion}).
L'assertion résulte donc de la discussion de
Propriétés des champs, section
\ref{stacks-properties-section-properties-morphisms}
et du résultat correspondant pour les morphismes d'espaces algébriques ; voir
Morphismes d'espaces, Lemme
\ref{spaces-morphisms-lemma-closed-immersion-proper}.
\end{proof}

\begin{lemma}
\label{stacks-morphisms-lemma-universally-closed-permanence}
Considérons un diagramme commutatif
$$
\xymatrix{
\mathcal{X} \ar[rr] \ar[rd] & &
\mathcal{Y} \ar[ld] \\
& \mathcal{Z} &
}
$$
de champs algébriques.
\begin{enumerate}
\item Si $\mathcal{X} \to \mathcal{Z}$ est universellement fermé et
$\mathcal{Y} \to \mathcal{Z}$ est séparé,
alors le morphisme $\mathcal{X} \to \mathcal{Y}$ est universellement fermé.
En particulier, l'image de $|\mathcal{X}|$ dans $|\mathcal{Y}|$ est fermée.
\item Si $\mathcal{X} \to \mathcal{Z}$ est propre et
$\mathcal{Y} \to \mathcal{Z}$ est séparé, alors
le morphisme $\mathcal{X} \to \mathcal{Y}$ est propre.
\end{enumerate}
\end{lemma}

\begin{proof}
Supposons que $\mathcal{X} \to \mathcal{Z}$ soit universellement fermé et que
$\mathcal{Y} \to \mathcal{Z}$ soit séparé.
Factorisons le morphisme sous la forme
$\mathcal{X} \to \mathcal{X} \times_\mathcal{Z} \mathcal{Y} \to \mathcal{Y}$.
Le premier morphisme est propre (Lemme \ref{stacks-morphisms-lemma-semi-diagonal}),
donc universellement fermé.
La projection $\mathcal{X} \times_\mathcal{Z} \mathcal{Y} \to \mathcal{Y}$
est déduite par changement de base d'un morphisme universellement fermé, donc
elle est universellement fermée ; voir le
Lemme \ref{stacks-morphisms-lemma-base-change-universally-closed}.
Ainsi, $\mathcal{X} \to \mathcal{Y}$ est universellement fermé comme composé
de morphismes universellement fermés ; voir le
Lemme \ref{stacks-morphisms-lemma-composition-universally-closed}.
Cela prouve (1). Pour en déduire (2), combinons (1) avec les
Lemmes \ref{stacks-morphisms-lemma-compose-after-separated},
\ref{stacks-morphisms-lemma-quasi-compact-permanence} et
\ref{stacks-morphisms-lemma-finite-type-permanence}.
\end{proof}

\begin{lemma}
\label{stacks-morphisms-lemma-image-proper-is-proper}
Soit $\mathcal{Z}$ un champ algébrique.
Soit $f : \mathcal{X} \to \mathcal{Y}$
un morphisme de champs algébriques au-dessus de $\mathcal{Z}$.
Si $\mathcal{X}$ est universellement fermé au-dessus de $\mathcal{Z}$
et si $f$ est surjectif, alors $\mathcal{Y}$
est universellement fermé au-dessus de $\mathcal{Z}$.
En particulier, si de plus $\mathcal{Y}$ est
séparé et de type fini au-dessus de $\mathcal{Z}$,
alors $\mathcal{Y}$ est propre au-dessus de $\mathcal{Z}$.
\end{lemma}

\begin{proof}
Supposons que $\mathcal{X}$ soit universellement fermé et que $f$ soit surjectif.
Notons $p : \mathcal{X} \to \mathcal{Z}$ et
$q : \mathcal{Y} \to \mathcal{Z}$ les morphismes structuraux.
Soit $\mathcal{Z}' \to \mathcal{Z}$ un morphisme de champs algébriques.
Le changement de base $f' : \mathcal{X}' \to \mathcal{Y}'$
de $f$ par $\mathcal{Z}' \to \mathcal{Z}$ est surjectif
(Propriétés des champs, Lemme
\ref{stacks-properties-lemma-base-change-surjective}) et le changement
de base $p' : \mathcal{X}' \to \mathcal{Z}'$ de $p$ est fermé.
Notons $q' : \mathcal{Y}' \to \mathcal{Z}'$ le changement de base de $q$.
Si $T \subset |\mathcal{Y}'|$ est fermé, alors
$(f')^{-1}(T) \subset |\mathcal{X}'|$ est fermé, donc
$p'((f')^{-1}(T)) = q'(T)$ est fermé. Ainsi, $q'$ est fermé.
\end{proof}






\section{Image schématique}
\label{stacks-morphisms-section-scheme-theoretic-image}

\noindent
Voici la définition.

\begin{definition}
\label{stacks-morphisms-definition-scheme-theoretic-image}
Soit $f : \mathcal{X} \to \mathcal{Y}$ un morphisme de champs algébriques.
L'{\it image schématique} de $f$ est le plus petit sous-champ fermé
$\mathcal{Z} \subset \mathcal{Y}$ par lequel $f$
se factorise\footnote{Nous verrons au
lemme \ref{stacks-morphisms-lemma-scheme-theoretic-image-existence}
que l'image schématique existe toujours.}.
\end{definition}

\noindent
Nous désignons souvent encore par $f : \mathcal{X} \to \mathcal{Z}$ la factorisation de $f$.
Si le morphisme $f$ n'est pas quasi-compact, alors, en général, la
construction de l'image schématique ne commute pas à la
restriction aux sous-champs ouverts de $\mathcal{Y}$. Toutefois, si $f$ est
quasi-compact, l'image schématique commute au changement de base plat
(lemme \ref{stacks-morphisms-lemma-existence-plus-flat-base-change}).

\begin{lemma}
\label{stacks-morphisms-lemma-cover-upstairs}
Soit $f : \mathcal{X} \to \mathcal{Y}$ un morphisme de champs algébriques.
Soit $g : \mathcal{W} \to \mathcal{X}$ un morphisme de champs algébriques
surjectif, plat et localement de présentation finie.
Alors l'image schématique de $f$ existe si et seulement si celle de
$f \circ g$ existe et, dans ce cas, ces images schématiques
sont les mêmes.
\end{lemma}

\begin{proof}
Supposons que $\mathcal{Z} \subset \mathcal{Y}$
soit un sous-champ fermé et que $f \circ g$ se factorise par $\mathcal{Z}$.
Pour démontrer le lemme, il suffit de montrer
que $f$ se factorise par $\mathcal{Z}$.
Considérons un schéma $T$ et un morphisme $T \to \mathcal{X}$
donné par un objet $x$ de la catégorie fibre de $\mathcal{X}$ au-dessus de $T$.
Nous allons montrer que $f(x)$ appartient en fait à la catégorie fibre de $\mathcal{Z}$
au-dessus de $T$. En effet, la projection $T \times_\mathcal{X} \mathcal{W} \to T$
est un morphisme surjectif, plat et localement de présentation finie.
Il existe donc un recouvrement fppf $\{T_i \to T\}$ tel que
$T_i \to T$ se factorise par $T \times_\mathcal{X} \mathcal{W} \to T$
pour tout $i$. Alors $T_i \to \mathcal{X}$ se factorise par $\mathcal{W}$
et donc $T_i \to \mathcal{Y}$ se factorise par $\mathcal{Z}$.
Ainsi $f(x)|_{T_i}$ est un objet de $\mathcal{Z}$.
Puisque $\mathcal{Z}$ est un sous-champ strictement plein, nous en concluons
que $f(x)$ est un objet de $\mathcal{Z}$, comme voulu.
\end{proof}

\begin{lemma}
\label{stacks-morphisms-lemma-scheme-theoretic-image-existence}
Soit $f : \mathcal{Y} \to \mathcal{X}$ un morphisme de champs algébriques.
Alors l'image schématique de $f$ existe.
\end{lemma}

\begin{proof}
Choisissons un schéma $V$ et un morphisme lisse surjectif $V \to \mathcal{Y}$.
D'après le lemme \ref{stacks-morphisms-lemma-cover-upstairs}, nous pouvons remplacer $\mathcal{Y}$ par $V$.
Il suffit donc de montrer que, si $X \to \mathcal{X}$ est un morphisme d'un
schéma vers un champ algébrique, alors l'image schématique existe.
Choisissons un schéma $U$ et un morphisme lisse surjectif $U \to \mathcal{X}$.
Posons $R = U \times_\mathcal{X} U$.
Nous avons $\mathcal{X} = [U/R]$ d'après
Champs algébriques, lemme \ref{algebraic-lemma-stack-presentation}.
D'après Propriétés des champs, lemme
\ref{stacks-properties-lemma-substacks-presentation},
les sous-champs fermés $\mathcal{Z}$ de $\mathcal{X}$
sont en correspondance $1$ à $1$ avec les sous-schémas fermés invariants
par $R$, $Z \subset U$.
Soit $Z_1 \subset U$ l'image schématique de
$X \times_\mathcal{X} U \to U$.
Remarquons que $X \to \mathcal{X}$ se factorise par $\mathcal{Z}$
si et seulement si $X \times_\mathcal{X} U \to U$ se factorise par
le sous-schéma fermé correspondant, invariant par $R$, à savoir $Z$
(détails omis ; indication : ceci résulte de ce que
$X \times_\mathcal{X} U \to X$ est lisse et surjectif).
Il nous faut donc montrer qu'il existe un plus petit sous-schéma fermé
invariant par $R$, $Z \subset U$, et contenant $Z_1$.

\medskip\noindent
Soit $\mathcal{I}_1 \subset \mathcal{O}_U$ le faisceau quasi-cohérent
d'idéaux correspondant au sous-schéma fermé $Z_1 \subset U$.
Soient $Z_\alpha$, $\alpha \in A$, tous les sous-schémas fermés invariants
par $R$ de $U$ et contenant $Z_1$.
Pour $\alpha \in A$, soit $\mathcal{I}_\alpha \subset \mathcal{O}_U$
le faisceau quasi-cohérent d'idéaux correspondant au sous-schéma fermé
$Z_\alpha \subset U$. L'inclusion $Z_1 \subset Z_\alpha$
signifie que $\mathcal{I}_\alpha \subset \mathcal{I}_1$.
L'invariance par $R$ de $Z_\alpha$ signifie que
$$
s^{-1}\mathcal{I}_\alpha \cdot \mathcal{O}_R =
t^{-1}\mathcal{I}_\alpha \cdot \mathcal{O}_R
$$
comme faisceaux quasi-cohérents d'idéaux sur l'espace algébrique $R$.
Considérons l'image
$$
\mathcal{I} =
\Im\left(
\bigoplus\nolimits_{\alpha \in A} \mathcal{I}_\alpha \to \mathcal{I}_1
\right) =
\Im\left(
\bigoplus\nolimits_{\alpha \in A} \mathcal{I}_\alpha \to \mathcal{O}_U
\right)
$$
Comme les sommes directes de faisceaux quasi-cohérents sont quasi-cohérentes
et que les images de morphismes entre faisceaux quasi-cohérents sont
quasi-cohérentes, nous voyons que $\mathcal{I}$ est quasi-cohérent.
Comme le foncteur image inverse est exact et commute aux sommes directes, nous obtenons
$$
s^{-1}\mathcal{I} \cdot \mathcal{O}_R =
t^{-1}\mathcal{I} \cdot \mathcal{O}_R
$$
Par conséquent, $\mathcal{I}$ définit un sous-schéma fermé invariant par $R$,
$Z \subset U$ qui est contenu dans tout $Z_\alpha$ et contient
$Z_1$, comme voulu.
\end{proof}

\begin{lemma}
\label{stacks-morphisms-lemma-factor-factor}
Soit
$$
\xymatrix{
\mathcal{X}_1 \ar[d] \ar[r]_{f_1} & \mathcal{Y}_1 \ar[d] \\
\mathcal{X}_2 \ar[r]^{f_2} & \mathcal{Y}_2
}
$$
un diagramme commutatif de champs algébriques.
Soient $\mathcal{Z}_i \subset \mathcal{Y}_i$, $i = 1, 2$, les
images schématiques de $f_i$. Alors le morphisme
$\mathcal{Y}_1 \to \mathcal{Y}_2$ induit un morphisme
$\mathcal{Z}_1 \to \mathcal{Z}_2$ et un
diagramme commutatif
$$
\xymatrix{
\mathcal{X}_1 \ar[r] \ar[d] &
\mathcal{Z}_1 \ar[d] \ar[r] &
\mathcal{Y}_1 \ar[d] \\
\mathcal{X}_2 \ar[r] &
\mathcal{Z}_2 \ar[r] &
\mathcal{Y}_2
}
$$
\end{lemma}

\begin{proof}
L'image réciproque schématique de $\mathcal{Z}_2$ dans $\mathcal{Y}_1$
est un sous-champ fermé de $\mathcal{Y}_1$ par lequel
$f_1$ se factorise. Par conséquent, $\mathcal{Z}_1$ y est contenu.
Ceci démontre le lemme.
\end{proof}

\begin{lemma}
\label{stacks-morphisms-lemma-existence-plus-flat-base-change}
Soit $f : \mathcal{X} \to \mathcal{Y}$ un morphisme quasi-compact
de champs algébriques. Alors la formation de l'image schématique
commute au changement de base plat.
\end{lemma}

\begin{proof}
Soit $\mathcal{Y}' \to \mathcal{Y}$ un morphisme plat de champs algébriques.
Choisissons un schéma $V$ et un morphisme lisse surjectif $V \to \mathcal{Y}$.
Choisissons un schéma $V'$ et un
morphisme lisse surjectif $V' \to \mathcal{Y}' \times_\mathcal{Y} V$.
Nous supposerons que $V = \coprod_{i \in I} V_i$ est une somme
disjointe de schémas affines et que
$V' = \coprod_{i \in I} \coprod_{j \in J_i} V_{i, j}$
est une somme disjointe de schémas affines, chaque $V_{i, j}$ s'envoyant dans $V_i$.
Soient
\begin{enumerate}
\item $\mathcal{Z} \subset \mathcal{Y}$ l'image schématique de $f$,
\item $\mathcal{Z}' \subset \mathcal{Y}'$ l'image schématique
du changement de base de $f$ par $\mathcal{Y}' \to \mathcal{Y}$,
\item $Z \subset V$ l'image schématique
du changement de base de $f$ par $V \to \mathcal{Y}$,
\item $Z' \subset V'$ l'image schématique
du changement de base de $f$ par $V' \to \mathcal{Y}$.
\end{enumerate}
Si nous pouvons montrer que
(a) $Z = V \times_\mathcal{Y} \mathcal{Z}$,
(b) $Z' = V' \times_{\mathcal{Y}'} \mathcal{Z}'$, et
(c) $Z' = V' \times_V Z$,
alors le lemme en résulte : l'inclusion
$\mathcal{Z}' \to \mathcal{Z} \times_\mathcal{Y} \mathcal{Y}'$
(lemme \ref{stacks-morphisms-lemma-factor-factor})
est nécessairement un isomorphisme, car son changement de base par le morphisme
lisse surjectif $V' \to \mathcal{Y}'$ en est un.

\medskip\noindent
Démonstration de (a). Posons $R = V \times_\mathcal{Y} V$.
D'après Propriétés des champs, lemme
\ref{stacks-properties-lemma-substacks-presentation},
l'application $\mathcal{Z} \mapsto \mathcal{Z} \times_\mathcal{Y} V$
définit une correspondance bijective entre les sous-champs fermés
de $\mathcal{Y}$ et les sous-espaces fermés invariants par $R$ dans $V$.
De plus, $f : \mathcal{X} \to \mathcal{Y}$ se factorise par $\mathcal{Z}$
si et seulement si le changement de base
$g : \mathcal{X} \times_\mathcal{Y} V \to V$ se factorise par
$\mathcal{Z} \times_\mathcal{Y} V$.
Nous affirmons que l'image schématique $Z \subset V$ de $g$
est invariante par $R$. Cette assertion implique (a) d'après ce qui précède.

\medskip\noindent
Pour chaque $i$, le morphisme $\mathcal{X} \times_\mathcal{Y} V_i \to V_i$
est quasi-compact ; par conséquent, $\mathcal{X} \times_\mathcal{Y} V_i$ est
quasi-compact. Nous pouvons donc choisir un schéma affine $W_i$ et un
morphisme lisse surjectif $W_i \to \mathcal{X} \times_\mathcal{Y} V_i$.
Remarquons que $W = \coprod W_i$ est un schéma muni d'un
morphisme lisse surjectif $W \to \mathcal{X} \times_\mathcal{Y} V$
tel que le composé $W \to V$ défini par $g$ soit quasi-compact.
Soit $Z \to V$ l'image schématique de $W \to V$, voir
Morphismes, section
\ref{morphisms-section-scheme-theoretic-image}, et
Morphismes d'espaces, section
\ref{spaces-morphisms-section-scheme-theoretic-image}.
Il résulte du lemme \ref{stacks-morphisms-lemma-cover-upstairs}
que $Z \subset V$ est l'image schématique de $g$.
Pour montrer que $Z$ est invariant par $R$, nous affirmons que les deux
$$
\text{pr}_0^{-1}(Z), \text{pr}_1^{-1}(Z) \subset R = V \times_\mathcal{Y} V
$$
sont l'image schématique de $\mathcal{X} \times_\mathcal{Y} R \to R$.
En effet, Morphismes d'espaces, lemme
\ref{spaces-morphisms-lemma-flat-base-change-scheme-theoretic-image},
montre d'abord que $\text{pr}_0^{-1}(Z)$ est l'image schématique
du composé
$$
W \times_{V, \text{pr}_0} R = W \times_\mathcal{Y} V \to
\mathcal{X} \times_\mathcal{Y} R \to R
$$
Comme la première flèche est lisse et surjective, nous voyons que
$\text{pr}_0^{-1}(Z)$ est l'image schématique de
$\mathcal{X} \times_\mathcal{Y} R \to R$. Le même argument montre que
$\text{pr}_1^{-1}(Z)$ l'est aussi. Ainsi $Z$ est invariant par $R$.

\medskip\noindent
L'assertion (b) se démontre exactement de la même manière que (a).

\medskip\noindent
Démonstration de (c). Soit $Z_i \subset V_i$ l'image schématique
de $\mathcal{X} \times_\mathcal{Y} V_i \to V_i$ et soit
$Z_{i, j} \subset V_{i, j}$ l'image schématique de
$\mathcal{X} \times_\mathcal{Y} V_{i, j} \to V_{i, j}$.
Il suffit clairement de montrer que l'image réciproque de $Z_i$
dans $V_{i, j}$ est $Z_{i, j}$. Nous avons vu ci-dessus que
$Z_i$ est l'image schématique de $W_i \to V_i$
et, de même, que $Z_{i, j}$ est l'image
schématique de $W_i \times_{V_i} V_{i, j} \to V_{i, j}$.
L'égalité résulte donc du cas des schémas
(Morphismes, lemme
\ref{morphisms-lemma-flat-base-change-scheme-theoretic-image})
et du fait que $V_{i, j} \to V_i$ est plat.
\end{proof}

\begin{lemma}
\label{stacks-morphisms-lemma-topology-scheme-theoretic-image}
Soit $f : \mathcal{X} \to \mathcal{Y}$ un morphisme quasi-compact
de champs algébriques. Soit $\mathcal{Z} \subset \mathcal{Y}$
l'image schématique de $f$. Alors $|\mathcal{Z}|$
est l'adhérence de l'image de $|f|$.
\end{lemma}

\begin{proof}
Soit $z \in |\mathcal{Z}|$ un point.
Choisissons un schéma affine $V$, un point $v \in V$ et un morphisme lisse
$V \to \mathcal{Y}$ envoyant $v$ sur $z$.
Alors $\mathcal{X} \times_\mathcal{Y} V$ est un champ algébrique quasi-compact.
Nous pouvons donc trouver un schéma affine $W$ et un morphisme lisse
surjectif $W \to \mathcal{X} \times_\mathcal{Y} V$.
D'après le lemme \ref{stacks-morphisms-lemma-existence-plus-flat-base-change},
l'image schématique de
$\mathcal{X} \times_\mathcal{Y} V \to V$ est
$Z = \mathcal{Z} \times_\mathcal{Y} V$.
Ainsi, l'image réciproque de $|\mathcal{Z}|$ dans $|V|$ est $|Z|$ d'après
Propriétés des champs, lemme \ref{stacks-properties-lemma-points-cartesian}.
D'après le lemme \ref{stacks-morphisms-lemma-cover-upstairs}, $Z$ est
l'image schématique de $W \to V$.
D'après Morphismes d'espaces, lemme
\ref{spaces-morphisms-lemma-quasi-compact-scheme-theoretic-image},
nous voyons que l'image de $|W| \to |Z|$ est dense.
Par conséquent, l'image de $|\mathcal{X} \times_\mathcal{Y} V| \to |Z|$
est dense. Remarquons que $v \in |Z|$.
Puisque $|V| \to |\mathcal{Y}|$ est ouvert, un argument topologique
montre que $z$ appartient à l'adhérence de l'image de $|f|$, comme voulu.
\end{proof}

\begin{lemma}
\label{stacks-morphisms-lemma-scheme-theoretic-image-of-partial-section}
Soit $f : \mathcal{X} \to \mathcal{Y}$ un morphisme de champs algébriques
représentable par des espaces algébriques et séparé.
Soit $\mathcal{V} \subset \mathcal{Y}$ un sous-champ ouvert tel que
$\mathcal{V} \to \mathcal{Y}$ soit quasi-compact.
Soit $s : \mathcal{V} \to \mathcal{X}$ un morphisme tel que
$f \circ s = \text{id}_\mathcal{V}$.
Soit $\mathcal{Y}'$ l'image schématique de $s$.
Alors $\mathcal{Y}' \to \mathcal{Y}$ est un isomorphisme au-dessus de $\mathcal{V}$.
\end{lemma}

\begin{proof}
D'après le lemme \ref{stacks-morphisms-lemma-quasi-compact-permanence},
le morphisme $s : \mathcal{V} \to \mathcal{X}$ est quasi-compact.
Par conséquent, la construction de l'image schématique $\mathcal{Y}'$
de $s$ commute au changement de base plat d'après
le lemme \ref{stacks-morphisms-lemma-existence-plus-flat-base-change}.
Ainsi, pour démontrer le lemme,
nous pouvons supposer que $\mathcal{Y}$ est représentable par un espace algébrique,
et nous sommes ramenés au cas des espaces algébriques, qui est
Morphismes d'espaces, lemme
\ref{spaces-morphisms-lemma-scheme-theoretic-image-of-partial-section}.
\end{proof}






\section{Crit\`eres valuatifs}
\label{stacks-morphisms-section-valuative}

\noindent
Il faut prendre garde lors de l'\'enonc\'e du crit\`ere valuatif. En effet, dans
la formulation, nous devons parler de diagrammes commutatifs, mais nous
travaillons dans une $2$-cat\'egorie et nous devons nous assurer que les $2$-morphismes
se composent eux aussi correctement !

\begin{definition}
\label{stacks-morphisms-definition-fill-in-diagram}
Soit $f : \mathcal{X} \to \mathcal{Y}$ un morphisme de champs alg\'ebriques.
Consid\'erons le diagramme $2$-commutatif en fl\`eches pleines
\begin{equation}
\label{stacks-morphisms-equation-diagram}
\vcenter{
\xymatrix{
\Spec(K) \ar[r]_-x \ar[d]_j & \mathcal{X} \ar[d]^f \\
\Spec(A) \ar[r]^-y \ar@{..>}[ru] & \mathcal{Y}
}
}
\end{equation}
o\`u $A$ est un anneau de valuation et $K$ son corps des fractions. Soit
$$
\gamma : y \circ j \longrightarrow f \circ x
$$
un $2$-morphisme t\'emoignant de la $2$-commutativit\'e du diagramme.
(Notations comme dans Cat\'egories, sections \ref{categories-section-formal-cat-cat}
et \ref{categories-section-2-categories}.)
\'Etant donn\'es (\ref{stacks-morphisms-equation-diagram}) et $\gamma$,
une {\it fl\`eche pointill\'ee} est un triplet $(a, \alpha, \beta)$ constitu\'e d'un
morphisme $a : \Spec(A) \to \mathcal{X}$ et de $2$-fl\`eches
$\alpha : a \circ j \to x$, $\beta : y \to f \circ a$
telles que
$\gamma = (\text{id}_f \star \alpha) \circ (\beta \star \text{id}_j)$,
autrement dit telles que
$$
\xymatrix{
& f \circ a \circ j \ar[rd]^{\text{id}_f \star \alpha} \\
y \circ j \ar[ru]^{\beta \star \text{id}_j} \ar[rr]^\gamma & &
f \circ x
}
$$
soit commutatif. Un {\it morphisme de fl\`eches pointill\'ees}
$(a, \alpha, \beta) \to (a', \alpha', \beta')$ est une
$2$-fl\`eche $\theta : a \to a'$ telle que
$\alpha = \alpha' \circ (\theta \star \text{id}_j)$ et
$\beta' = (\text{id}_f \star \theta) \circ \beta$.
\end{definition}

\noindent
La d\'efinition pr\'ec\'edente est un cas particulier de Cat\'egories, 
D\'efinition \ref{categories-definition-dotted-arrows}.
La cat\'egorie des fl\`eches pointill\'ees d\'epend en g\'en\'eral de $\gamma$.
Si $\mathcal{Y}$ est repr\'esentable par un espace alg\'ebrique
(ou si les groupes d'automorphismes des objets au-dessus des corps sont triviaux),
il existe bien entendu au plus un $\gamma$, et il n'est pas n\'ecessaire
de v\'erifier la commutativit\'e du triangle. Plus g\'en\'eralement, nous
disposons du Lemme \ref{stacks-morphisms-lemma-cat-dotted-arrows-independent}.
La commutativit\'e du triangle est importante dans la d\'emonstration
de la compatibilit\'e au changement de base, voir la
d\'emonstration du Lemme \ref{stacks-morphisms-lemma-cat-dotted-arrows-base-change}.

\begin{lemma}
\label{stacks-morphisms-lemma-cat-dotted-arrows}
Dans la situation de la D\'efinition \ref{stacks-morphisms-definition-fill-in-diagram},
la cat\'egorie des fl\`eches pointill\'ees est un groupo\"\i de. Si $\Delta_f$
est s\'epar\'e, alors c'est un s\'eto\"\i de.
\end{lemma}

\begin{proof}
Comme les $2$-fl\`eches sont inversibles, il est clair que la cat\'egorie des
fl\`eches pointill\'ees est un groupo\"\i de. Pour une fl\`eche pointill\'ee $(a, \alpha, \beta)$,
un automorphisme de $(a, \alpha, \beta)$ est un $2$-morphisme
$\theta : a \to a$ qui v\'erifie deux conditions. La premi\`ere condition
$\beta = (\text{id}_f \star \theta) \circ \beta$ exprime que
$\theta$ d\'efinit un morphisme
$(a, \theta) : \Spec(A) \to \mathcal{I}_{\mathcal{X}/\mathcal{Y}}$.
La seconde condition
$\alpha = \alpha \circ (\theta \star \text{id}_j)$
implique que la restriction de $(a, \theta)$ \`a $\Spec(K)$
est l'identit\'e. On a le diagramme
$$
\xymatrix{
\mathcal{I}_{\mathcal{X}/\mathcal{Y}} \ar[d] & &
\Spec(K) \ar[d]^j \ar[ll]_{(a \circ j, \text{id})} \\
\mathcal{X} & & \Spec(A) \ar[ll]_a \ar[llu]_{(a, \theta)}
}
$$
Autrement dit, si $G \to \Spec(A)$ est l'espace alg\'ebrique en groupes
obtenu par changement de base de l'inertie relative suivant $a$, alors
$\theta$ d\'efinit un point $\theta \in G(A)$ dont l'image
dans $G(K)$ est l'\'el\'ement neutre. Bien entendu, si la section unit\'e $e : \Spec(A) \to G$
est une immersion ferm\'ee, cela ne peut se produire que si
$\theta$ est l'identit\'e.
Le Lemme \ref{stacks-morphisms-lemma-diagonal-diagonal}
donne le r\'esultat voulu.
\end{proof}

\begin{lemma}
\label{stacks-morphisms-lemma-cat-dotted-arrows-independent}
Dans la D\'efinition \ref{stacks-morphisms-definition-fill-in-diagram}, supposons que
$\mathcal{I}_\mathcal{Y} \to \mathcal{Y}$ soit propre
(par exemple si $\mathcal{Y}$ est s\'epar\'e, ou si $\mathcal{Y}$
est s\'epar\'e au-dessus d'un espace alg\'ebrique). Alors la cat\'egorie des fl\`eches pointill\'ees
est ind\'ependante, \`a une \'equivalence non canonique pr\`es, du choix de $\gamma$,
et l'existence d'une fl\`eche pointill\'ee
(pour un certain choix, donc, de mani\`ere \`equivalente, pour tout $\gamma$)
\'equivaut \`a l'existence d'un diagramme
$$
\xymatrix{
\Spec(K) \ar[r]_-x \ar[d]_j & \mathcal{X} \ar[d]^f \\
\Spec(A) \ar[r]^-y \ar[ru]_a & \mathcal{Y}
}
$$
dont les triangles sont $2$-commutatifs
(sans v\'erifier que les $2$-morphismes se composent correctement).
\end{lemma}

\begin{proof}
Soient $\gamma, \gamma' : y \circ j \longrightarrow f \circ x$
deux $2$-morphismes. Alors $\gamma^{-1} \circ \gamma'$
est un automorphisme de $y$ au-dessus de $\Spec(K)$.
Ainsi, si $\mathit{Isom}_\mathcal{Y}(y, y) \to \Spec(A)$
est propre, le crit\`ere valuatif de propret\'e
(Morphismes d'espaces, Lemme \ref{spaces-morphisms-lemma-characterize-proper})
permet de trouver $\delta : y \to y$ dont la restriction \`a
$\Spec(K)$ est $\gamma^{-1} \circ \gamma'$.
Nous pouvons alors utiliser $\delta$ pour d\'efinir une \`equivalence
entre la cat\'egorie des fl\`eches pointill\'ees pour $\gamma$
et la cat\'egorie des fl\`eches pointill\'ees pour $\gamma'$, en
envoyant $(a, \alpha, \beta)$ sur $(a, \alpha, \beta \circ \delta)$.
La derni\`ere assertion est claire.
\end{proof}

\begin{lemma}
\label{stacks-morphisms-lemma-cat-dotted-arrows-base-change}
Supposons donn\'e un diagramme $2$-commutatif
$$
\xymatrix{
\Spec(K) \ar[r]_-{x'} \ar[d]_j &
\mathcal{X}' \ar[d]^p \ar[r]_q &
\mathcal{X} \ar[d]^f \\
\Spec(A) \ar[r]^-{y'} &
\mathcal{Y}' \ar[r]^g &
\mathcal{Y}
}
$$
dont le carr\'e de droite est $2$-cart\'esien. Choisissons une $2$-fl\`eche
$\gamma' : y' \circ j \to p \circ x'$. Posons
$x = q \circ x'$, $y = g \circ y'$ et soit
$\gamma : y \circ j \to f \circ x$ le compos\'e de
$\gamma'$ et de la $2$-fl\`eche implicite dans la $2$-commutativit\'e
du carr\'e de droite. Alors la cat\'egorie des fl\`eches pointill\'ees
pour le carr\'e de gauche et $\gamma'$ est \`equivalente \`a la cat\'egorie des fl\`eches
pointill\'ees pour le rectangle ext\'erieur et $\gamma$.
\end{lemma}

\begin{proof}
(Nous ne savons pas d\'emontrer l'analogue de ce lemme si, au lieu
de la cat\'egorie des fl\`eches pointill\'ees, nous consid\'erons l'ensemble des classes
d'isomorphie des morphismes donnant deux triangles
$2$-commutatifs comme dans le Lemme \ref{stacks-morphisms-lemma-cat-dotted-arrows-independent};
il se peut fort bien, en fait, que cet analogue soit faux.)
Premi\`ere d\'emonstration : ce lemme est un cas particulier de Cat\'egories, 
Lemme \ref{categories-lemma-cat-dotted-arrows-base-change}.
Seconde d\'emonstration : nous pouvons remplacer
$\mathcal{X}'$ par le $2$-produit fibr\'e
$\mathcal{Y}' \times_\mathcal{Y} \mathcal{X}$
tel qu'il est d\'ecrit dans Cat\'egories, Lemme
\ref{categories-lemma-2-product-categories-over-C}.
L'objet $x'$ devient alors le triplet $(y' \circ j, x, \gamma)$.
On passe alors d'une fl\`eche pointill\'ee $(a, \alpha, \beta)$ pour le
rectangle ext\'erieur \`a une fl\`eche pointill\'ee $(a', \alpha', \beta')$
pour le carr\'e de gauche en prenant $a' = (y', a, \beta)$,
$\alpha' = (\text{id}_{y' \circ j}, \alpha)$ et
$\beta' = \text{id}_{y'}$. D\'etails omis.
\end{proof}

\begin{lemma}
\label{stacks-morphisms-lemma-cat-dotted-arrows-composition}
Supposons donn\'e un diagramme $2$-commutatif
$$
\xymatrix{
\Spec(K) \ar[r]_-x \ar[dd]_j & \mathcal{X} \ar[d]^f \\
& \mathcal{Y} \ar[d]^g \\
\Spec(A) \ar[r]^-z & \mathcal{Z}
}
$$
Choisissons une $2$-fl\`eche $\gamma : z \circ j \to g \circ f \circ x$.
Soit $\mathcal{C}$ la cat\'egorie des fl\`eches pointill\'ees pour
le rectangle ext\'erieur et $\gamma$. Soit $\mathcal{C}'$ la
cat\'egorie des fl\`eches pointill\'ees pour le carr\'e
$$
\xymatrix{
\Spec(K) \ar[r]_-{f \circ x} \ar[d]_j & \mathcal{Y} \ar[d]^g \\
\Spec(A) \ar[r]^-z & \mathcal{Z}
}
$$
et $\gamma$. Alors $\mathcal{C}$ est \`equivalente \`a une cat\'egorie $\mathcal{C}''$ 
qui poss\`ede la propri\'et\'e suivante : il existe 
un foncteur $\mathcal{C}'' \to \mathcal{C}'$
qui fait de $\mathcal{C}''$ une cat\'egorie fibr\'ee en groupo\"\i des sur
$\mathcal{C}'$ et dont les cat\'egories-fibres sont des cat\'egories de fl\`eches pointill\'ees
pour certains carr\'es de la forme
$$
\xymatrix{
\Spec(K) \ar[r]_-x \ar[d]_j & \mathcal{X} \ar[d]^f \\
\Spec(A) \ar[r]^-y & \mathcal{Y}
}
$$
et certains choix de $y \circ j \to f \circ x$.
\end{lemma}

\begin{proof}
Ce lemme est un cas particulier de Cat\'egories, 
Lemme \ref{categories-lemma-cat-dotted-arrows-composition}.
\end{proof}

\begin{definition}
\label{stacks-morphisms-definition-uniqueness}
Soit $f : \mathcal{X} \to \mathcal{Y}$ un morphisme de champs alg\'ebriques.
Nous disons que $f$ v\'erifie la {\it partie d'unicit\'e du crit\`ere valuatif}
si, pour tout diagramme (\ref{stacks-morphisms-equation-diagram}) et tout $\gamma$
comme dans la D\'efinition \ref{stacks-morphisms-definition-fill-in-diagram},
la cat\'egorie des fl\`eches pointill\'ees est soit vide, soit
un s\'eto\"\i de ayant exactement une classe d'isomorphie.
\end{definition}

\begin{lemma}
\label{stacks-morphisms-lemma-base-change-uniqueness}
Le changement de base, par un morphisme quelconque de champs alg\'ebriques,
d'un morphisme de champs alg\'ebriques qui v\'erifie la
partie d'unicit\'e du crit\`ere valuatif est un morphisme de champs alg\'ebriques
qui v\'erifie la partie d'unicit\'e du crit\`ere valuatif.
\end{lemma}

\begin{proof}
Cela r\'esulte du Lemme \ref{stacks-morphisms-lemma-cat-dotted-arrows-base-change}
et de la d\'efinition.
\end{proof}

\begin{lemma}
\label{stacks-morphisms-lemma-composition-uniqueness}
Le compos\'e de morphismes de champs alg\'ebriques qui v\'erifient la
partie d'unicit\'e du crit\`ere valuatif est encore un
morphisme de champs alg\'ebriques qui v\'erifie la
partie d'unicit\'e du crit\`ere valuatif.
\end{lemma}

\begin{proof}
Cela r\'esulte du Lemme \ref{stacks-morphisms-lemma-cat-dotted-arrows-composition}
et de la d\'efinition.
\end{proof}

\begin{lemma}
\label{stacks-morphisms-lemma-uniqueness-representable}
Soit $f : \mathcal{X} \to \mathcal{Y}$ un morphisme de champs alg\'ebriques
qui est repr\'esentable par des espaces alg\'ebriques. Alors les conditions suivantes sont \`equivalentes :
\begin{enumerate}
\item $f$ v\'erifie la partie d'unicit\'e du crit\`ere valuatif ;
\item pour tout sch\'ema $T$ et tout morphisme $T \to \mathcal{Y}$,
le morphisme $\mathcal{X} \times_\mathcal{Y} T \to T$ v\'erifie
la partie d'unicit\'e du crit\`ere valuatif en tant que morphisme
d'espaces alg\'ebriques.
\end{enumerate}
\end{lemma}

\begin{proof}
Cela r\'esulte du Lemme \ref{stacks-morphisms-lemma-cat-dotted-arrows-base-change}
et de la d\'efinition.
\end{proof}

\begin{definition}
\label{stacks-morphisms-definition-existence}
Soit $f : \mathcal{X} \to \mathcal{Y}$ un morphisme de champs alg\'ebriques.
Nous disons que $f$ v\'erifie la {\it partie d'existence du crit\`ere valuatif}
si, pour tout diagramme (\ref{stacks-morphisms-equation-diagram}) et tout $\gamma$
comme dans la D\'efinition \ref{stacks-morphisms-definition-fill-in-diagram},
il existe une extension de corps $K'/K$, un anneau de valuation $A' \subset K'$
dominant $A$, tels que la cat\'egorie des fl\`eches pointill\'ees pour le
rectangle ext\'erieur du diagramme
$$
\xymatrix{
\Spec(K') \ar[r] \ar@/^2em/[rr]_{x'} \ar[d]_{j'} &
\Spec(K) \ar[d]_j \ar[r]_-x &
\mathcal{X} \ar[d]^f \\
\Spec(A') \ar[r] \ar@/_2em/[rr]^{y'} &
\Spec(A) \ar[r]^-y &
\mathcal{Y}
}
$$
et la $2$-fl\`eche induite $\gamma' : y' \circ j' \to f \circ x'$ soit non vide.
\end{definition}

\noindent
Nous avons d\'ej\`a vu, dans le cas des morphismes d'espaces alg\'ebriques,
qu'il est n\'ecessaire d'autoriser des extensions des corps des fractions
pour obtenir la bonne notion de crit\`ere valuatif.
Voir Morphismes d'espaces, Exemple
\ref{spaces-morphisms-example-finite-separable-needed}.
Toutefois, pour les morphismes entre espaces alg\'ebriques s\'epar\'es, une telle
extension n'est pas n\'ecessaire
(Morphismes d'espaces, Lemme \ref{spaces-morphisms-lemma-usual-enough}).
Cependant, pour les morphismes entre champs alg\'ebriques, une extension peut
\^etre n\'ecessaire m\^eme si $\mathcal{X}$ et $\mathcal{Y}$ sont tous deux s\'epar\'es.
Consid\'erons par exemple le morphisme de champs alg\'ebriques
$$
[\Spec(\mathbf{C})/G] \to \Spec(\mathbf{C})
$$
au-dessus du sch\'ema de base $\Spec(\mathbf{C})$,
o\`u $G$ est un groupe d'ordre $2$. La source et le but sont tous deux des
champs alg\'ebriques s\'epar\'es, et le morphisme est propre. Il
v\'erifie donc les parties d'unicit\'e et d'existence du crit\`ere valuatif
(voir le Lemme \ref{stacks-morphisms-lemma-criterion-proper}).
Mais, d'autre part, il existe un diagramme
$$
\xymatrix{
\Spec(K) \ar[r] \ar[d] & [\Spec(\mathbf{C})/G] \ar[d] \\
\Spec(A) \ar[r] & \Spec(\mathbf{C})
}
$$
pour lequel aucune fl\`eche pointill\'ee n'existe lorsque $A = \mathbf{C}[[t]]$ et
$K = \mathbf{C}((t))$. En effet, la fl\`eche horizontale sup\'erieure est donn\'ee
par un torseur sous $G$ sur le spectre de $K = \mathbf{C}((t))$. Comme tout torseur sous $G$
sur le spectre de l'anneau local strictement hens\'elien $A = \mathbf{C}[[t]]$ est trivial, nous voyons
que, si une fl\`eche pointill\'ee existe toujours, alors tout torseur sous $G$ sur
$K$ est trivial. Ce n'est pas le cas, car $G = \{+1, -1\}$
et, d'apr\`es la th\'eorie de Kummer, les torseurs sous $G$ sur $K$ sont class\'es par
$K^*/(K^*)^2$, qui n'est pas trivial.

\begin{lemma}
\label{stacks-morphisms-lemma-base-change-existence}
Le changement de base, par un morphisme quelconque de champs alg\'ebriques,
d'un morphisme de champs alg\'ebriques qui v\'erifie la
partie d'existence du crit\`ere valuatif est un morphisme de champs alg\'ebriques
qui v\'erifie la partie d'existence du crit\`ere valuatif.
\end{lemma}

\begin{proof}
Cela r\'esulte du Lemme \ref{stacks-morphisms-lemma-cat-dotted-arrows-base-change}
et de la d\'efinition.
\end{proof}

\begin{lemma}
\label{stacks-morphisms-lemma-composition-existence}
Le compos\'e de morphismes de champs alg\'ebriques qui v\'erifient la
partie d'existence du crit\`ere valuatif est encore un
morphisme de champs alg\'ebriques qui v\'erifie la
partie d'existence du crit\`ere valuatif.
\end{lemma}

\begin{proof}
Cela r\'esulte du Lemme \ref{stacks-morphisms-lemma-cat-dotted-arrows-composition}
et de la d\'efinition.
\end{proof}

\begin{lemma}
\label{stacks-morphisms-lemma-existence-representable}
Soit $f : \mathcal{X} \to \mathcal{Y}$ un morphisme de champs alg\'ebriques
qui est repr\'esentable par des espaces alg\'ebriques. Alors les conditions suivantes sont \`equivalentes :
\begin{enumerate}
\item $f$ v\'erifie la partie d'existence du crit\`ere valuatif ;
\item pour tout sch\'ema $T$ et tout morphisme $T \to \mathcal{Y}$,
le morphisme $\mathcal{X} \times_\mathcal{Y} T \to T$ v\'erifie
la partie d'existence du crit\`ere valuatif en tant que morphisme
d'espaces alg\'ebriques.
\end{enumerate}
\end{lemma}

\begin{proof}
Cela r\'esulte du Lemme \ref{stacks-morphisms-lemma-cat-dotted-arrows-base-change}
et de la d\'efinition.
\end{proof}

\begin{lemma}
\label{stacks-morphisms-lemma-closed-immersion-valuative-criteria}
Une immersion ferm\'ee de champs alg\'ebriques v\'erifie \`a la fois
les parties d'existence et d'unicit\'e du crit\`ere valuatif.
\end{lemma}

\begin{proof}
D\'emonstration omise. Indication : se ramener au cas d'une immersion ferm\'ee de
sch\'emas \`a l'aide des Lemmes \ref{stacks-morphisms-lemma-uniqueness-representable} et
\ref{stacks-morphisms-lemma-existence-representable}.
\end{proof}




\section{Critère valuatif pour la seconde diagonale}
\label{stacks-morphisms-section-valuative-second}

\noindent
La réciproque a déjà été démontrée dans le
Lemme \ref{stacks-morphisms-lemma-cat-dotted-arrows}.
Le critère proprement dit est le suivant.

\begin{lemma}
\label{stacks-morphisms-lemma-setoids-and-diagonal}
Soit $f : \mathcal{X} \to \mathcal{Y}$ un morphisme de champs algébriques.
Si le morphisme $\Delta_f$ est quasi-séparé et si, pour tout diagramme
(\ref{stacks-morphisms-equation-diagram}) et tout choix de $\gamma$ comme dans la
Définition \ref{stacks-morphisms-definition-fill-in-diagram},
la catégorie des flèches pointillées
est un setoïde, alors $\Delta_f$ est séparé.
\end{lemma}

\begin{proof}
Nous allons donner une démonstration détaillée, mais nous recommandons vivement au
lecteur de chercher sa propre démonstration en s'inspirant de l'argument
donné dans la preuve du Lemme \ref{stacks-morphisms-lemma-cat-dotted-arrows}.

\medskip\noindent
Supposons que $\Delta_f$ soit quasi-séparé et que, pour tout diagramme
(\ref{stacks-morphisms-equation-diagram}) et tout choix de $\gamma$ comme dans la
Définition \ref{stacks-morphisms-definition-fill-in-diagram},
la catégorie des flèches pointillées soit un setoïde.
D'après le Lemme \ref{stacks-morphisms-lemma-diagonal-diagonal}, il suffit de montrer que
$e : \mathcal{X} \to \mathcal{I}_{\mathcal{X}/\mathcal{Y}}$
est une immersion fermée. D'après le
Lemme \ref{stacks-morphisms-lemma-first-diagonal-separated-second-diagonal-closed},
il suffit même de montrer que $e = \Delta_{f, 2}$ est
universellement fermé.
L'un ou l'autre de ces lemmes montre aussi que $e = \Delta_{f, 2}$ est quasi-compact
en vertu de notre hypothèse de quasi-séparation de $\Delta_f$.

\medskip\noindent
Dans ce paragraphe, nous allons montrer que $e$ vérifie la partie
d'existence du critère valuatif. Considérons le diagramme $2$-commutatif en flèches pleines
$$
\xymatrix{
\Spec(K) \ar[r]_x \ar[d]_j & \mathcal{X} \ar[d]^e \\
\Spec(A) \ar[r]^{(a, \theta)} & \mathcal{I}_{\mathcal{X}/\mathcal{Y}}
}
$$
et soit $\alpha : (a, \theta) \circ j \to e \circ x$ un $2$-morphisme quelconque
témoignant de la $2$-commutativité du diagramme (nous employons $\alpha$ à la place
de la lettre $\gamma$ utilisée dans la Définition \ref{stacks-morphisms-definition-fill-in-diagram}).
Notons que $f \circ \theta = \text{id}$; nous utiliserons ce fait ci-dessous.
On a $e \circ x = (x, \text{id}_x)$ et
$(a, \theta) \circ j = (a \circ j, \theta \star \text{id}_j)$.
Ainsi, $\alpha$ est une $2$-flèche $\alpha : a \circ j \to x$
compatible avec $\theta \star \text{id}_j$ et $\text{id}_x$.
Posons $y = f \circ x$ et $\beta = \text{id}_{f \circ a}$.
En reprenant en sens inverse les arguments donnés dans la preuve du
Lemme \ref{stacks-morphisms-lemma-cat-dotted-arrows},
on voit que $\theta$ est un automorphisme de la
flèche pointillée $(a, \alpha, \beta)$ pour laquelle
$$
\gamma : y \circ j \to f \circ x
\quad\text{égal à}\quad
\text{id}_f \star \alpha : f \circ a \circ j \to f \circ x
$$
D'autre part, $\text{id}_a$ est aussi un automorphisme; on en conclut
$\theta = \text{id}_a$ grâce à l'hypothèse sur $f$.
On peut alors prendre pour flèche pointillée du diagramme affiché ci-dessus
le morphisme $a : \Spec(A) \to \mathcal{X}$ muni des deux $2$-morphismes suivants :
le premier, $(a, \text{id}_a) \circ j \to (x, \text{id}_x)$, est donné par $\alpha$,
et le second, $(a, \theta) \to e \circ a$, par $\text{id}_a$.

\medskip\noindent
D'après le Lemme \ref{stacks-morphisms-lemma-base-change-existence}, tout changement de base de $e$
vérifie la partie d'existence du critère valuatif.
Comme $e$ est représentable par des espaces algébriques, il suffit de
montrer que $e$ est universellement fermé après changement de base
par un morphisme $I \to \mathcal{I}_{\mathcal{X}/\mathcal{Y}}$
surjectif et lisse, où $I$ est un espace algébrique
(voir Propriétés des champs, section
\ref{stacks-properties-section-properties-morphisms}).
Ce changement de base $e' : X' \to I'$ est un
morphisme quasi-compact d'espaces algébriques qui
vérifie la partie d'existence du critère valuatif;
il est donc universellement fermé d'après
Morphismes d'espaces, Lemme
\ref{spaces-morphisms-lemma-quasi-compact-existence-universally-closed}.
\end{proof}






\section{Critère valuatif pour la diagonale}
\label{stacks-morphisms-section-valuative-diagonal}

\noindent
Le résultat est le Lemme \ref{stacks-morphisms-lemma-uniqueness-and-diagonal}.
Nous commençons par énoncer et démontrer un lemme auxiliaire formel.

\begin{lemma}
\label{stacks-morphisms-lemma-helper-diagonal}
Soit $f : \mathcal{X} \to \mathcal{Y}$ un morphisme de champs algébriques.
Considérons le diagramme $2$-commutatif en flèches pleines
$$
\xymatrix{
\Spec(K) \ar[rr]_-x \ar[d]_j & &
\mathcal{X} \ar[d]^{\Delta_f} \\
\Spec(A) \ar[rr]^{(a_1, a_2, \varphi)} \ar@{..>}[rru] & &
\mathcal{X} \times_\mathcal{Y} \mathcal{X}
}
$$
où $A$ est un anneau de valuation de corps des fractions $K$. Soit
$\gamma : (a_1, a_2, \varphi) \circ j \longrightarrow \Delta_f \circ x$
un $2$-morphisme témoignant de la $2$-commutativité du diagramme.
Alors :
\begin{enumerate}
\item En écrivant $\gamma = (\alpha_1, \alpha_2)$ avec
$\alpha_i : a_i \circ j \to x$, on obtient deux flèches pointillées
$(a_1, \alpha_1, \text{id})$ et $(a_2, \alpha_2, \varphi)$ dans
le diagramme
$$
\xymatrix{
\Spec(K) \ar[r]_-x \ar[d]_j & \mathcal{X} \ar[d]^f \\
\Spec(A) \ar[r]^-{f \circ a_1} \ar@{..>}[ru] & \mathcal{Y}
}
$$
\item La catégorie des flèches pointillées associée au diagramme initial
et à $\gamma$ est un setoïde; l'ensemble des classes d'isomorphisme
de ses objets s'identifie à l'ensemble des morphismes
$(a_1, \alpha_1, \text{id}) \to (a_2, \alpha_2, \varphi)$ de
la catégorie des flèches pointillées.
\end{enumerate}
\end{lemma}

\begin{proof}
Comme $\Delta_f$ est représentable par des espaces algébriques (la diagonale
de $\Delta_f$ est donc séparée), le Lemme
\ref{stacks-morphisms-lemma-cat-dotted-arrows} montre que la catégorie des flèches pointillées
du premier diagramme commutatif du lemme est un setoïde.
Toutes les autres assertions du lemme résultent de calculs
de diagrammes dans une $2$-catégorie, que nous omettons.
\end{proof}

\begin{lemma}
\label{stacks-morphisms-lemma-uniqueness-and-diagonal}
Soit $f : \mathcal{X} \to \mathcal{Y}$ un morphisme de champs algébriques.
Supposons que $f$ soit quasi-séparé.
Si $f$ vérifie la partie d'unicité du critère valuatif,
alors $f$ est séparé.
\end{lemma}

\begin{proof}
L'hypothèse sur $f$ signifie que $\Delta_f$ est quasi-compact
et quasi-séparé (Définition \ref{stacks-morphisms-definition-separated}).
Il nous faut montrer que $\Delta_f$ est propre.
Le Lemme \ref{stacks-morphisms-lemma-setoids-and-diagonal} montre que $\Delta_f$
est séparé. D'après le Lemme \ref{stacks-morphisms-lemma-properties-diagonal},
on sait que $\Delta_f$ est localement de type fini.
Pour achever la preuve, il reste à montrer que
$\Delta_f$ est universellement fermé. Un argument formel
(voir le Lemme \ref{stacks-morphisms-lemma-helper-diagonal}) montre que
la partie d'unicité du critère valuatif entraîne
l'existence d'une flèche pointillée dans tout diagramme en flèches pleines de la forme suivante :
$$
\xymatrix{
\Spec(K) \ar[d] \ar[r] & \mathcal{X} \ar[d]^{\Delta_f} \\
\Spec(A) \ar[r] \ar@{..>}[ru] & \mathcal{X} \times_\mathcal{Y} \mathcal{X}
}
$$
Cette propriété étant préservée par tout changement de base,
tout changement de base de $\Delta_f$ par un espace algébrique
muni d'un morphisme vers $\mathcal{X} \times_\mathcal{Y} \mathcal{X}$
vérifie la partie d'existence du critère valuatif et
est donc universellement fermé d'après le critère valuatif
pour les morphismes d'espaces algébriques; voir
Morphismes d'espaces, Lemme
\ref{spaces-morphisms-lemma-quasi-compact-existence-universally-closed}.
\end{proof}

\noindent
Voici une réciproque.

\begin{lemma}
\label{stacks-morphisms-lemma-converse-uniqueness-and-diagonal}
Soit $f : \mathcal{X} \to \mathcal{Y}$ un morphisme de champs algébriques.
Si $f$ est séparé, alors $f$ vérifie la
partie d'unicité du critère valuatif.
\end{lemma}

\begin{proof}
Puisque $f$ est séparé, toutes les catégories de flèches pointillées
sont des setoïdes d'après le Lemme \ref{stacks-morphisms-lemma-cat-dotted-arrows}.
Considérons un diagramme
$$
\xymatrix{
\Spec(K) \ar[r]_-x \ar[d]_j & \mathcal{X} \ar[d]^f \\
\Spec(A) \ar[r]^-y \ar@{..>}[ru] & \mathcal{Y}
}
$$
et un $2$-morphisme $\gamma : y \circ j \to f \circ x$ comme dans la
Définition \ref{stacks-morphisms-definition-fill-in-diagram}. Considérons deux
objets $(a, \alpha, \beta)$ et $(a', \beta', \alpha')$
de la catégorie des flèches pointillées. Pour achever la preuve, il
faut montrer que ces objets sont isomorphes. L'isomorphisme
$$
f \circ a \xrightarrow{\beta^{-1}} y \xrightarrow{\beta'} f \circ a'
$$
signifie que $(a, a', \beta' \circ \beta^{-1})$ définit un morphisme
$\Spec(A) \to \mathcal{X} \times_\mathcal{Y} \mathcal{X}$.
D'autre part, $\alpha$ et $\alpha'$ définissent
une $2$-flèche
$$
(a, a', \beta' \circ \beta^{-1}) \circ j =
(a \circ j, a' \circ j,
(\beta' \star \text{id}_j) \circ (\beta \star \text{id}_j)^{-1})
\xrightarrow{(\alpha, \alpha')} (x, x, \text{id}) = \Delta_f \circ x
$$
On utilise ici que $(a, \alpha, \beta)$ et $(a', \alpha', \beta')$
sont toutes deux des flèches pointillées associées à $\gamma$.
On obtient un diagramme commutatif
$$
\xymatrix{
\Spec(K) \ar[d]_j \ar[rr]_x & & \mathcal{X} \ar[d]^{\Delta_f} \\
\Spec(A) \ar[rr]^{(a, a', \beta' \circ \beta^{-1})} & &
\mathcal{X} \times_\mathcal{Y} \mathcal{X}
}
$$
dont la $2$-commutativité est donnée par $(\alpha, \alpha')$. Or
$\Delta_f$ est représentable par des espaces algébriques
(Lemme \ref{stacks-morphisms-lemma-properties-diagonal})
et propre puisque $f$ est séparé. Le
Lemme \ref{stacks-morphisms-lemma-existence-representable}
et le critère valuatif de propreté pour les espaces algébriques
(Morphismes d'espaces, Lemme \ref{spaces-morphisms-lemma-characterize-proper})
montrent donc qu'il existe une flèche pointillée.
En explicitant la construction, on voit que cela signifie que
$(a, \alpha, \beta)$ et $(a', \alpha', \beta')$
sont isomorphes dans la catégorie des flèches pointillées, comme souhaité.
\end{proof}





 
\section{Critère valuatif pour les morphismes universellement fermés}
\sectionmark{Critère valuatif de fermeture universelle}
\label{stacks-morphisms-section-valutive-criterion}

\noindent
Voici un énoncé.

\begin{lemma}
\label{stacks-morphisms-lemma-quasi-compact-existence-universally-closed}
Soit $f : \mathcal{X} \to \mathcal{Y}$ un morphisme de champs algébriques.
Supposons que
\begin{enumerate}
\item $f$ soit quasi-compact, et que
\item $f$ vérifie la partie d'existence du critère valuatif.
\end{enumerate}
Alors $f$ est universellement fermé.
\end{lemma}

\begin{proof}
D'après les Lemmes \ref{stacks-morphisms-lemma-base-change-quasi-compact}
et \ref{stacks-morphisms-lemma-base-change-existence},
les propriétés (1) et (2) sont préservées par
tout changement de base. D'après le Lemme \ref{stacks-morphisms-lemma-universally-closed-local},
il suffit de montrer que $|T \times_\mathcal{Y} \mathcal{X}| \to |T|$
est fermée chaque fois que $T$ est un schéma affine muni d'un morphisme vers $\mathcal{Y}$.
Il suffit donc de montrer ceci : si $f : \mathcal{X} \to Y$ est un
morphisme quasi-compact d'un champ algébrique vers un schéma affine
qui vérifie la partie d'existence du critère valuatif,
alors $|f|$ est fermée. Soit $T \subset |\mathcal{X}|$ une partie fermée.
Pour conclure la démonstration, il faut montrer que $f(T)$ est fermé.

\medskip\noindent
Soit $\mathcal{Z} \subset \mathcal{X}$ le sous-champ algébrique réduit
induit par $T$ (Propriétés des champs,
Définition \ref{stacks-properties-definition-reduced-induced-stack}).
Alors $i : \mathcal{Z} \to \mathcal{X}$ est une immersion fermée
et il faut montrer que l'image de $|\mathcal{Z}| \to |Y|$
est fermée. Comme une immersion fermée est quasi-compacte
(Lemme \ref{stacks-morphisms-lemma-closed-immersion-quasi-compact})
et vérifie la partie d'existence du critère valuatif
(Lemme \ref{stacks-morphisms-lemma-closed-immersion-valuative-criteria}),
comme le composé de morphismes quasi-compacts est quasi-compact
(Lemme \ref{stacks-morphisms-lemma-composition-quasi-compact})
et comme la composition préserve la propriété de vérifier
la partie d'existence du critère valuatif
(Lemme \ref{stacks-morphisms-lemma-composition-existence}),
nous en concluons qu'il suffit de montrer ceci : si $f : \mathcal{X} \to Y$
est un morphisme quasi-compact d'un champ algébrique vers un schéma affine
qui vérifie la partie d'existence du critère valuatif,
alors $|f|(|\mathcal{X}|)$ est fermé.

\medskip\noindent
Puisque $\mathcal{X}$ est quasi-compact (car quasi-compact au-dessus du
schéma affine $Y$), nous pouvons choisir un schéma affine $U$ et un morphisme
lisse et surjectif $U \to \mathcal{X}$
(Propriétés des champs, Lemme
\ref{stacks-properties-lemma-quasi-compact-stack}).
Supposons que $y \in Y$ appartienne à l'adhérence de l'image de $U \to Y$
(autrement dit, à l'adhérence de l'image de $|f|$).
Alors, d'après
Morphismes, Lemme \ref{morphisms-lemma-reach-points-scheme-theoretic-image},
il existe un anneau de valuation $A$ de corps des fractions $K$
et un diagramme commutatif
$$
\xymatrix{
\Spec(K) \ar[r] \ar[d] & U \ar[d] \\
\Spec(A) \ar[r] & Y
}
$$
tel que le point fermé de $\Spec(A)$ ait pour image $y$. Par hypothèse,
il existe une extension $K'/K$, un anneau de valuation $A' \subset K'$
dominant $A$, ainsi que la flèche pointillée du diagramme suivant :
$$
\xymatrix{
\Spec(K') \ar[r] \ar[d] &
\Spec(K) \ar[r] \ar[d] &
U \ar[d] \ar[r] &
\mathcal{X} \ar[d]^f \\
\Spec(A') \ar[r] \ar@{..>}[rrru] &
\Spec(A) \ar[r] &
Y  \ar@{=}[r] &
Y
}
$$
Ainsi, $y$ appartient à l'image de $|f|$, ce qui conclut.
\end{proof}

\noindent
Voici une réciproque.

\begin{lemma}
\label{stacks-morphisms-lemma-converse-existence-universally-closed}
Soit $f : \mathcal{X} \to \mathcal{Y}$ un morphisme de champs algébriques.
Supposons que
\begin{enumerate}
\item $f$ soit quasi-séparé, et que
\item $f$ soit universellement fermé.
\end{enumerate}
Alors $f$ vérifie la partie d'existence du critère valuatif.
\end{lemma}

\begin{proof}
Considérons un diagramme en flèches pleines
$$
\xymatrix{
\Spec(K) \ar[r]_-x \ar[d]_j & \mathcal{X} \ar[d]^f \\
\Spec(A) \ar[r]^-y \ar@{..>}[ru] & \mathcal{Y}
}
$$
où $A$ est un anneau de valuation de corps des fractions $K$
et où $\gamma : y \circ j \longrightarrow f \circ x$ est comme dans la
Définition \ref{stacks-morphisms-definition-fill-in-diagram}. D'après le
Lemme \ref{stacks-morphisms-lemma-cat-dotted-arrows-base-change},
pour trouver une flèche pointillée (après avoir éventuellement remplacé
$K$ par une extension et $A$ par un anneau de valuation le dominant),
nous pouvons remplacer $\mathcal{Y}$ par $\Spec(A)$ et $\mathcal{X}$
par $\Spec(A) \times_\mathcal{Y} \mathcal{X}$. Bien entendu,
nous utilisons ici le fait que les propriétés d'être
quasi-séparé et universellement fermé sont préservées par changement de base.
Nous sommes ainsi ramenés au cas traité au paragraphe suivant.

\medskip\noindent
Considérons un diagramme en flèches pleines
$$
\xymatrix{
\Spec(K) \ar[r]_-x \ar[d]_j & \mathcal{X} \ar[d]^f \\
\Spec(A) \ar@{=}[r] \ar@{..>}[ru] & \Spec(A)
}
$$
où $A$ est un anneau de valuation de corps des fractions $K$, comme dans la
Définition \ref{stacks-morphisms-definition-fill-in-diagram}.
D'après le Lemme \ref{stacks-morphisms-lemma-quasi-compact-permanence} et puisque
$f$ est quasi-séparé, nous savons que
le morphisme $x$ est quasi-compact.
Comme $f$ est universellement fermé, nous savons en particulier
que $|f|(\overline{\{x\}})$ est fermé dans $\Spec(A)$.
Puisque cette image contient le point générique de $\Spec(A)$,
il existe un point $x' \in |\mathcal{X}|$ dans l'adhérence
de $x$ qui s'envoie sur le point fermé de $\Spec(A)$.
D'après le Lemme \ref{stacks-morphisms-lemma-reach-points-scheme-theoretic-image},
il existe un diagramme commutatif
$$
\xymatrix{
\Spec(K') \ar[r] \ar[d] & \Spec(K) \ar[d] \\
\Spec(A') \ar[r] & \mathcal{X}
}
$$
tel que le point fermé de $\Spec(A')$ s'envoie sur $x' \in |\mathcal{X}|$.
Il s'ensuit que $\Spec(A') \to \Spec(A)$ envoie le point fermé
sur le point fermé, c'est-à-dire que $A'$ domine $A$, ce qui achève la démonstration.
\end{proof}





\section{Critère valuatif de propreté}
\label{stacks-morphisms-section-valutive-criterion-properness}

\noindent
Voici l'énoncé.

\begin{lemma}
\label{stacks-morphisms-lemma-criterion-proper}
Soit $f : \mathcal{X} \to \mathcal{Y}$ un morphisme de champs algébriques.
Supposons que $f$ soit de type fini et quasi-séparé.
Alors les conditions suivantes sont équivalentes :
\begin{enumerate}
\item $f$ est propre, et
\item $f$ vérifie à la fois les parties d'unicité et d'existence
du critère valuatif.
\end{enumerate}
\end{lemma}

\begin{proof}
Un morphisme propre est exactement un morphisme séparé, de type fini et
universellement fermé. Le lemme résulte donc des Lemmes
\ref{stacks-morphisms-lemma-uniqueness-and-diagonal},
\ref{stacks-morphisms-lemma-converse-uniqueness-and-diagonal},
\ref{stacks-morphisms-lemma-quasi-compact-existence-universally-closed} et
\ref{stacks-morphisms-lemma-converse-existence-universally-closed}.
\end{proof}





 
 
\section{Morphismes localement d'intersection compl\`ete}
\label{stacks-morphisms-section-lci}

\noindent
La propri\'et\'e ``\^etre un morphisme localement d'intersection compl\`ete''
des morphismes d'espaces alg\'ebriques est
locale pour la topologie lisse sur la source et sur le but; voir
Descente sur les espaces, Lemme \ref{spaces-descent-lemma-smooth-local-source-target}
et
Compl\'ements sur les morphismes d'espaces, Lemmes
\ref{spaces-more-morphisms-lemma-descending-property-lci} et
\ref{spaces-more-morphisms-lemma-lci-syntomic-local-source}.
D'apr\`es le Lemme \ref{stacks-morphisms-lemma-local-source-target} ci-dessus, nous pouvons d\'efinir
comme suit la notion de morphisme localement d'intersection compl\`ete
pour les champs alg\'ebriques; elle co\"incide avec la
notion d\'ej\`a d\'efinie dans
Compl\'ements sur les morphismes d'espaces,
section \ref{spaces-more-morphisms-section-lci},
lorsque la source et le but sont tous deux des espaces alg\'ebriques.

\begin{definition}
\label{stacks-morphisms-definition-lci}
Soit $f : \mathcal{X} \to \mathcal{Y}$ un morphisme de champs alg\'ebriques.
On dit que $f$ est un {\it morphisme localement d'intersection compl\`ete} ou
{\it de Koszul} si les conditions \`equivalentes du
Lemme \ref{stacks-morphisms-lemma-local-source-target}
sont satisfaites pour $\mathcal{P} = \text{localement d'intersection compl\`ete}$.
\end{definition}

\begin{lemma}
\label{stacks-morphisms-lemma-composition-lci}
Le compos\'e de morphismes localement d'intersection compl\`ete est
un morphisme localement d'intersection compl\`ete.
\end{lemma}

\begin{proof}
Le résultat découle de
la Remarque \ref{stacks-morphisms-remark-composition}
et du
Lemme \ref{spaces-more-morphisms-lemma-composition-lci}
de Compl\'ements sur les morphismes d'espaces.
\end{proof}

\begin{lemma}
\label{stacks-morphisms-lemma-flat-base-change-lci}
Un changement de base plat d'un morphisme localement d'intersection compl\`ete est
un morphisme localement d'intersection compl\`ete.
\end{lemma}

\begin{proof}
D\'emonstration omise. Indication : raisonner exactement comme dans la Remarque \ref{stacks-morphisms-remark-base-change}
(mais seulement lorsque $\mathcal{Y}' \to \mathcal{Y}$ est plat) en utilisant
Compl\'ements sur les morphismes d'espaces, Lemme
\ref{spaces-more-morphisms-lemma-flat-base-change-lci}.
\end{proof}

\begin{lemma}
\label{stacks-morphisms-lemma-lci-permanence}
Soit
$$
\xymatrix{
\mathcal{X} \ar[rr]_f \ar[rd] & & \mathcal{Y} \ar[ld] \\
& \mathcal{Z}
}
$$
un diagramme commutatif de morphismes de champs alg\'ebriques.
Supposons que $\mathcal{Y} \to \mathcal{Z}$ soit lisse et que
$\mathcal{X} \to \mathcal{Z}$ soit un morphisme localement d'intersection compl\`ete.
Alors $f : \mathcal{X} \to \mathcal{Y}$ est un morphisme
localement d'intersection compl\`ete.
\end{lemma}

\begin{proof}
Choisissons un sch\'ema $W$ et un morphisme lisse surjectif $W \to \mathcal{Z}$.
Choisissons un sch\'ema $V$ et un morphisme lisse surjectif
$V \to W \times_\mathcal{Z} \mathcal{Y}$.
Choisissons un sch\'ema $U$ et un morphisme lisse surjectif
$U \to V \times_\mathcal{Y} \mathcal{X}$.
Alors $U \to W$ est un morphisme localement d'intersection compl\`ete de sch\'emas et
$V \to W$ est un morphisme lisse de sch\'emas. Par le r\'esultat pour les sch\'emas
(Compl\'ements sur les morphismes, Lemme \ref{more-morphisms-lemma-lci-permanence}),
nous concluons que $U \to V$ est un morphisme localement d'intersection compl\`ete.
Par d\'efinition, cela signifie que $f$ est un morphisme localement d'intersection compl\`ete.
\end{proof}





\section{Morphismes préservant les stabilisateurs}
\label{stacks-morphisms-section-stabilizer-preserving}

\noindent
Dans la littérature, on dit qu'un morphisme $f : \mathcal{X} \to \mathcal{Y}$ de
champs algébriques {\it préserve les stabilisateurs} ou
{\it reflète les points fixes} si le morphisme induit
$\mathcal{I}_\mathcal{X} \to
\mathcal{X} \times_\mathcal{Y} \mathcal{I}_\mathcal{Y}$
est un isomorphisme. Un tel morphisme induit un isomorphisme
entre les groupes d'automorphismes (remarque \ref{stacks-morphisms-remark-identify-automorphism-groups})
en tout point de $\mathcal{X}$.
Dans cette section, nous démontrons quelques lemmes élémentaires sur cette notion.

\begin{lemma}
\label{stacks-morphisms-lemma-stabilizer-preserving}
Soit $f : \mathcal{X} \to \mathcal{Y}$ un morphisme de champs algébriques.
Si $\mathcal{I}_\mathcal{X} \to
\mathcal{X} \times_\mathcal{Y} \mathcal{I}_\mathcal{Y}$ est un isomorphisme,
alors $f$ est représentable par des espaces algébriques.
\end{lemma}

\begin{proof}
Cela résulte immédiatement du lemme \ref{stacks-morphisms-lemma-second-diagonal}.
\end{proof}

\begin{remark}
\label{stacks-morphisms-remark-get-property-auts-from-diagonal}
Soit $f : \mathcal{X} \to \mathcal{Y}$ un morphisme de champs algébriques.
Soit $U \to \mathcal{X}$ un morphisme dont la source est un espace algébrique.
Soit $G \to H$ l'image réciproque du morphisme
$\mathcal{I}_\mathcal{X} \to
\mathcal{X} \times_\mathcal{Y} \mathcal{I}_\mathcal{Y}$
sur $U$. Si $\Delta_f$ est non ramifié, \'etale, etc.,
il en va de même de $G \to H$. En effet,
$$
\xymatrix{
U \times_\mathcal{X} U \ar[r] \ar[d] & \mathcal{X} \ar[d]^{\Delta_f} \\
U \times_\mathcal{Y} U \ar[r] & \mathcal{X} \times_\mathcal{Y} \mathcal{X}
}
$$
est cartésien et le morphisme $G \to H$ se déduit de la
flèche verticale de gauche par le changement de base défini par la diagonale $U \to U \times U$.
Comparer avec la démonstration du lemme \ref{stacks-morphisms-lemma-separated-implies-isom}.
\end{remark}

\begin{lemma}
\label{stacks-morphisms-lemma-aut-iso-unramified}
Soit $f : \mathcal{X} \to \mathcal{Y}$ un morphisme non ramifié
de champs algébriques. Les conditions suivantes sont équivalentes :
\begin{enumerate}
\item $\mathcal{I}_\mathcal{X} \to
\mathcal{X} \times_\mathcal{Y} \mathcal{I}_\mathcal{Y}$
est un isomorphisme, et
\item $f$ induit un isomorphisme entre les groupes d'automorphismes en $x$ et en $f(x)$
(remarque \ref{stacks-morphisms-remark-identify-automorphism-groups}) pour tout
$x \in |\mathcal{X}|$.
\end{enumerate}
\end{lemma}

\begin{proof}
Choisissons un schéma $U$ et un morphisme lisse surjectif $U \to \mathcal{X}$.
Notons $G \to H$ l'image réciproque du morphisme
$\mathcal{I}_\mathcal{X} \to
\mathcal{X} \times_\mathcal{Y} \mathcal{I}_\mathcal{Y}$
sur $U$. D'après la remarque \ref{stacks-morphisms-remark-get-property-auts-from-diagonal} et le
lemme \ref{stacks-morphisms-lemma-characterize-unramified}, le morphisme $G \to H$ est \'etale.
La condition (1) équivaut à dire que
$G \to H$ est un isomorphisme (cela résulte par exemple de l'application du
lemme de Propriétés des champs
\ref{stacks-properties-lemma-check-property-covering}).
La condition (2) équivaut à dire que
pour tout $u \in U$, le morphisme de fibres $G_u \to H_u$
est un isomorphisme. Ainsi, l'implication (1) $\Rightarrow$ (2) est triviale.
Si (2) est satisfaite, alors $G \to H$ est un morphisme d'espaces algébriques surjectif, universellement injectif,
et \'etale. Un tel morphisme est un isomorphisme d'après le
lemme de Morphismes d'espaces
\ref{spaces-morphisms-lemma-etale-universally-injective-open}.
\end{proof}

\begin{lemma}
\label{stacks-morphisms-lemma-stabilizer-preserving-unramified}
\begin{reference}
\cite[Proposition 3.5]{rydh_quotients} et
\cite[Proposition 2.5]{alper_quotient}
\end{reference}
Soit $f : \mathcal{X} \to \mathcal{Y}$ un morphisme de champs algébriques.
Supposons que
\begin{enumerate}
\item $f$ soit représentable par des espaces algébriques et non ramifié, et
\item $\mathcal{I}_\mathcal{Y} \to \mathcal{Y}$ soit propre.
\end{enumerate}
Alors l'ensemble des $x \in |\mathcal{X}|$ tels que $f$ induise un
isomorphisme entre les groupes d'automorphismes en $x$ et en $f(x)$
(remarque \ref{stacks-morphisms-remark-identify-automorphism-groups}) est ouvert.
Notons $\mathcal{U} \subset \mathcal{X}$ le sous-champ ouvert correspondant ;
alors le morphisme
$\mathcal{I}_\mathcal{U} \to
\mathcal{U} \times_\mathcal{Y} \mathcal{I}_\mathcal{Y}$
est un isomorphisme.
\end{lemma}

\begin{proof}
Choisissons un schéma $U$ et un morphisme lisse surjectif $U \to \mathcal{X}$.
Notons $G \to H$ l'image réciproque du morphisme
$\mathcal{I}_\mathcal{X} \to
\mathcal{X} \times_\mathcal{Y} \mathcal{I}_\mathcal{Y}$
sur $U$. D'après la remarque \ref{stacks-morphisms-remark-get-property-auts-from-diagonal} et le
lemme \ref{stacks-morphisms-lemma-characterize-unramified}, le morphisme
$G \to H$ est \'etale. Comme $f$ est représentable par des espaces algébriques,
nous voyons que $G \to H$ est un monomorphisme. Ainsi, $G \to H$ est une immersion
ouverte, voir
le lemme de Morphismes d'espaces
\ref{spaces-morphisms-lemma-etale-universally-injective-open}.
Par hypothèse, $H \to U$ est propre.

\medskip\noindent
Ces préparatifs étant faits, démontrons le lemme.
L'image réciproque du sous-ensemble de $|\mathcal{X}|$ considéré dans le lemme est
manifestement l'ensemble des $u \in U$ tels que $G_u \to H_u$ soit un isomorphisme
(puisque, après tout, $G_u$ est un sous-groupe ouvert de l'espace algébrique en groupes $H_u$).
Cet ensemble est ouvert, car son complément est
l'image du fermé $|H| \setminus |G|$, et l'application $|H| \to |U|$
est fermée. D'après le
lemme de Propriétés des champs \ref{stacks-properties-lemma-open-substacks},
nous pouvons considérer le sous-champ ouvert correspondant $\mathcal{U}$ de $\mathcal{X}$.
La dernière assertion du lemme découle de l'application du
lemme \ref{stacks-morphisms-lemma-aut-iso-unramified}
à $\mathcal{U} \to \mathcal{Y}$.
\end{proof}

\begin{lemma}
\label{stacks-morphisms-lemma-base-change-stabilizer-preserving}
Soit
$$
\xymatrix{
\mathcal{X}' \ar[r] \ar[d]_{f'} & \mathcal{X} \ar[d]^f \\
\mathcal{Y}' \ar[r] & \mathcal{Y}
}
$$
un diagramme cartésien de champs algébriques.
\begin{enumerate}
\item Soit $x' \in |\mathcal{X}'|$, d'image $x \in |\mathcal{X}|$.
Si $f$ induit un isomorphisme entre les groupes d'automorphismes en
$x$ et en $f(x)$ (remarque \ref{stacks-morphisms-remark-identify-automorphism-groups}), alors
$f'$ induit un isomorphisme entre les groupes d'automorphismes en $x'$ et en $f(x')$.
\item Si $\mathcal{I}_\mathcal{X} \to
\mathcal{X} \times_\mathcal{Y} \mathcal{I}_\mathcal{Y}$ est un isomorphisme,
alors $\mathcal{I}_{\mathcal{X}'} \to
\mathcal{X}' \times_{\mathcal{Y}'} \mathcal{I}_{\mathcal{Y}'}$
est un isomorphisme.
\end{enumerate}
\end{lemma}

\begin{proof}
Démonstration omise.
\end{proof}

\begin{lemma}
\label{stacks-morphisms-lemma-stabilizer-preserving-points-cartesian}
Soit
$$
\xymatrix{
\mathcal{X}' \ar[r] \ar[d]_{f'} & \mathcal{X} \ar[d]^f \\
\mathcal{Y}' \ar[r]^g & \mathcal{Y}
}
$$
un diagramme cartésien de champs algébriques. Si $f$ induit en tout point un isomorphisme
entre les groupes d'automorphismes
(remarque \ref{stacks-morphisms-remark-identify-automorphism-groups}),
alors l'application
$$
\Mor(\Spec(k), \mathcal{X}')
\longrightarrow
\Mor(\Spec(k), \mathcal{Y}') \times \Mor(\Spec(k), \mathcal{X})
$$
est injective sur les classes d'isomorphisme pour tout corps $k$.
\end{lemma}

\begin{proof}
Il s'agit de montrer qu'étant donné $(y', x)$, il existe au plus un $x'$
qui ait cette image.
D'après notre construction des $2$-produits fibrés, un morphisme
$x'$ est donné par un triplet $(x, y', \alpha)$,
où $\alpha : g \circ y' \to f \circ x$ est un $2$-morphisme.
Supposons maintenant donné un second triplet $(x, y', \beta)$ de cette sorte.
Alors $\alpha$ et $\beta$ diffèrent par un point à valeurs dans $k$, noté
$\epsilon$, de l'espace algébrique en groupes d'automorphismes $G_{f(x)}$.
Puisque, par hypothèse, $f$ induit un isomorphisme $G_x \to G_{f(x)}$,
nous pouvons relever $\epsilon$ en un point à valeurs dans $k$, noté
$\gamma$, de $G_x$. Alors $(\gamma, \text{id}) : (x, y', \alpha) \to
(x, y', \beta)$ est l'isomorphisme recherché.
\end{proof}

\begin{lemma}
\label{stacks-morphisms-lemma-etale-iso}
Soit $f : \mathcal{X} \to \mathcal{Y}$ un morphisme de champs algébriques.
Supposons que $f$ soit \'etale, que $f$ induise un isomorphisme
entre les groupes d'automorphismes en tout point
(remarque \ref{stacks-morphisms-remark-identify-automorphism-groups}),
et que, pour tout corps algébriquement clos $k$, le foncteur
$$
f : \Mor(\Spec(k), \mathcal{X}) \longrightarrow \Mor(\Spec(k), \mathcal{Y})
$$
soit une équivalence. Alors $f$ est un isomorphisme.
\end{lemma}

\begin{proof}
D'après le lemme \ref{stacks-morphisms-lemma-universally-injective}, $f$ est
universellement injectif. En combinant les
lemmes \ref{stacks-morphisms-lemma-stabilizer-preserving} et
\ref{stacks-morphisms-lemma-aut-iso-unramified},
nous voyons que $f$ est représentable par des espaces algébriques.
Par conséquent, $f$ est une immersion ouverte d'après le lemme de Morphismes d'espaces
\ref{spaces-morphisms-lemma-etale-universally-injective-open}.
Pour conclure, remarquons que la condition de l'énoncé garantit également
que $f$ est surjectif.
\end{proof}






 
\section{Normalisation}
\label{stacks-morphisms-section-normalization}

\noindent
Cette section est l'analogue de
Morphismes d'espaces, section \ref{spaces-morphisms-section-normalization}.

\begin{lemma}
\label{stacks-morphisms-lemma-prepare-normalization}
Soit $\mathcal{X}$ un champ algébrique. Les conditions suivantes sont équivalentes :
\begin{enumerate}
\item il existe un morphisme lisse surjectif $U \to \mathcal{X}$, où $U$
est un schéma, tel que tout ouvert quasi-compact de $U$ ait
un nombre fini de composantes irréductibles,
\item pour tout schéma $U$ et tout morphisme lisse
$U \to \mathcal{X}$, tout ouvert quasi-compact de $U$ a un nombre fini de
composantes irréductibles,
\item pour tout espace algébrique $Y$ et tout morphisme lisse $Y \to \mathcal{X}$,
l'espace $Y$ satisfait aux conditions équivalentes de
Morphismes d'espaces, Lemme \ref{spaces-morphisms-lemma-prepare-normalization},
et
\item pour tout champ algébrique quasi-compact $\mathcal{Y}$ lisse sur
$\mathcal{X}$, l'espace $|\mathcal{Y}|$ a un nombre fini de composantes
irréductibles.
\end{enumerate}
\end{lemma}

\begin{proof}
L'équivalence de (1), (2) et (3) résulte de
Descente, Lemme \ref{descent-lemma-locally-finite-nr-irred-local-fppf},
Propriétés des champs, Lemme \ref{stacks-properties-lemma-type-property}, et de
Morphismes d'espaces, Lemme \ref{spaces-morphisms-lemma-prepare-normalization}.
Ces références montrent aussi clairement que la condition (4)
entraîne la condition (1). Réciproquement, supposons que les conditions équivalentes
(1), (2) et (3) soient satisfaites, et soit $\mathcal{Y} \to \mathcal{X}$ un
morphisme lisse de champs algébriques, avec $\mathcal{Y}$
quasi-compact. Nous pouvons alors choisir un schéma affine $V$ et un morphisme lisse
surjectif $V \to \mathcal{Y}$ d'après Propriétés des champs, Lemme
\ref{stacks-properties-lemma-quasi-compact-stack}.
Comme $V$ a un nombre fini de composantes irréductibles d'après (2), et que
l'application $|V| \to |\mathcal{Y}|$ est surjective et continue, on en conclut que
$|\mathcal{Y}|$ a un nombre fini de composantes irréductibles d'après Topologie, Lemme
\ref{topology-lemma-surjective-continuous-irreducible-components}.
\end{proof}

\begin{lemma}
\label{stacks-morphisms-lemma-normalization}
Soit $\mathcal{X}$ un champ algébrique satisfaisant aux
conditions équivalentes du Lemme \ref{stacks-morphisms-lemma-prepare-normalization}.
Alors il existe un morphisme intégral de champs algébriques
$$
\mathcal{X}^\nu \longrightarrow \mathcal{X}
$$
tel que, pour tout schéma $U$ et tout morphisme lisse $U \to \mathcal{X}$,
le produit fibré $\mathcal{X}^\nu \times_\mathcal{X} U$
soit le normalisé de $U$.
\end{lemma}

\begin{proof}
Soit $U \to \mathcal{X}$ un morphisme lisse surjectif, où $U$ est un schéma.
Posons $R = U \times_\mathcal{X} U$. Rappelons que l'on obtient un groupoïde lisse
$(U, R, s, t, c)$ en espaces algébriques et une présentation
$\mathcal{X} = [U/R]$ de $\mathcal{X}$ ; voir Champs algébriques,
Lemmes \ref{algebraic-lemma-map-space-into-stack} et
\ref{algebraic-lemma-stack-presentation}, ainsi que la
Définition \ref{algebraic-definition-presentation}.
L'hypothèse sur $\mathcal{X}$ signifie que le normalisé
$U^\nu$ de $U$ est défini ; voir Morphismes, Définition
\ref{morphisms-definition-normalization}.
D'après Morphismes d'espaces, Lemme
\ref{spaces-morphisms-lemma-normalization},
la formation du normalisé commute au changement de base par un morphisme lisse
d'espaces algébriques. Ainsi, le normalisé $R^\nu$ de $R$
est isomorphe tant à $R \times_{s, U} U^\nu$ qu'à $U^\nu \times_{U, t} R$.
Nous obtenons donc deux morphismes lisses
$s^\nu : R^\nu \to U^\nu$ et $t^\nu : R^\nu \to U^\nu$
d'espaces algébriques. Un calcul formel sur les produits fibrés montre que
$R^\nu \times_{s^\nu, U^\nu, t^\nu} R^\nu$ est le normalisé
de $R \times_{s, U, t} R$. Par conséquent, le morphisme lisse
$c : R \times_{s, U, t} R \to R$ se prolonge lui aussi en $c^\nu$.
De même, l'inverse $i : R \to R$ (qui est un isomorphisme) induit
un isomorphisme $i^\nu : R^\nu \to R^\nu$. Enfin, l'identité
$e : U \to R$ admet un relèvement $e^\nu : U^\nu \to R^\nu$, par exemple
parce que $e$ est une section de $s$ et que $R^\nu = R \times_{U, s} U^\nu$.
Nous affirmons que $(U^\nu, R^\nu, s^\nu, t^\nu, c^\nu)$ est
un groupoïde lisse en espaces algébriques. Pour le voir, il faut vérifier
les axiomes (1), (2)(a), (2)(b), (3)(a) et (3)(b) de
Groupoïdes, section \ref{groupoids-section-groupoids}
pour $(U^\nu, R^\nu, s^\nu, t^\nu, c^\nu, e^\nu, i^\nu)$.
Pour (1), par exemple, il faut vérifier que les deux morphismes
$a, b : R^\nu \times_{s^\nu, U^\nu, t^\nu} R^\nu
\times_{s^\nu, U^\nu, t^\nu} R^\nu \to R^\nu$ que l'on obtient coïncident.
Cela résulte du fait que les deux morphismes correspondants
$R \times_{s, U, t} R \times_{s, U, t} R \to R$
coïncident, et que les morphismes $a$ et $b$ sont les prolongements uniques
de ce morphisme aux normalisés correspondants. Les autres
axiomes se traitent de même.

\medskip\noindent
Considérons le champ algébrique $\mathcal{X}^\nu = [U^\nu/R^\nu]$
(Champs algébriques, Théorème
\ref{algebraic-theorem-smooth-groupoid-gives-algebraic-stack}).
Puisque nous disposons d'un morphisme
$(U^\nu, R^\nu, s^\nu, t^\nu, c^\nu) \to (U, R, s, t, c)$
de groupoïdes en espaces algébriques, nous obtenons un morphisme
$\nu : \mathcal{X}^\nu \to \mathcal{X}$ de champs algébriques.
Comme $R^\nu = R \times_{s, U} U^\nu$, on voit que
$U^\nu = \mathcal{X}^\nu \times_\mathcal{X} U$ d'après
Groupoïdes en espaces, Lemme
\ref{spaces-groupoids-lemma-criterion-fibre-product}.
En particulier, puisque $U^\nu \to U$ est intégral, $\nu$ est intégral.
Nous omettons de vérifier que la propriété de changement de base
énoncée dans le lemme est satisfaite pour tout morphisme lisse d'un schéma vers
$\mathcal{X}$.
\end{proof}

\noindent
Ceci conduit à la définition suivante.

\begin{definition}
\label{stacks-morphisms-definition-normalization}
Soit $\mathcal{X}$ un champ algébrique satisfaisant aux
conditions équivalentes du Lemme \ref{stacks-morphisms-lemma-prepare-normalization}.
Nous définissons la {\it normalisation} de $\mathcal{X}$ comme le morphisme
$$
\nu : \mathcal{X}^\nu \longrightarrow \mathcal{X}
$$
construit dans le Lemme \ref{stacks-morphisms-lemma-normalization}.
\end{definition}
 
  
   
    
     
      




\section{Points et spécialisations}
\label{stacks-morphisms-section-specializations}

\noindent
Cette section est l'analogue de la section
\ref{decent-spaces-section-specializations} d'Espaces décents.

\begin{lemma}
\label{stacks-morphisms-lemma-specialization}
Soit $\mathcal{X}$ un champ algébrique. Soit $f : U \to \mathcal{X}$
un morphisme lisse, où $U$ est un espace algébrique. Soit
$x' \leadsto x$ une spécialisation de points de $|\mathcal{X}|$.
Soit $u \in |U|$ tel que $f(u) = x$. Si $(\mathcal{X}, x')$ vérifie les
conditions équivalentes de Propriétés des champs,
Lemme \ref{stacks-properties-lemma-UR-quasi-compact-above-x},
alors il existe une spécialisation
$u' \leadsto u$ dans $|U|$ telle que $f(u') = x'$.
\end{lemma}

\begin{proof}
Choisissons un morphisme \'etale $(U_1, u_1) \to (U, u)$, où $U_1$
est un schéma affine. Remplaçons alors $U$ par $U_1$.
Nous pouvons donc supposer que $U$ est un schéma affine.
Considérons l'espace algébrique $R = U \times_\mathcal{X} U$
muni des projections lisses $t, s : R \to U$.
Choisissons un point $w \in U$ d'image $x'$ ; cela est possible puisque
l'application $f : |U| \to |\mathcal{X}|$ est ouverte.
D'après notre hypothèse sur $x'$, la fibre
$F' = t^{-1}(w) = R \times_{t, U} w$
de $t : R \to U$ au-dessus de $w$ est un espace algébrique quasi-compact.
Choisissons un schéma affine $T$ et un morphisme \'etale surjectif
$T \to F'$. Le fait que $x' \leadsto x$ signifie que $u$ appartient à
l'adhérence de l'image du morphisme
$$
T \to F' \to R \xrightarrow{s} U
$$
En effet, cette image est la fibre de $|U| \to |\mathcal{X}|$
au-dessus de $x'$ ; si $u \in V \subset |U|$, pour un ouvert disjoint
de cette fibre, alors $f(V)$ est un voisinage ouvert de $x$
qui ne contient pas $x'$ ; contradiction.
Ainsi, d'après Morphismes, Lemme
\ref{morphisms-lemma-reach-points-scheme-theoretic-image},
il existe $u' \in |U|$, dans la fibre
de $|U| \to |\mathcal{X}|$ au-dessus de $x'$, qui se spécialise en $u$.
\end{proof}


\section{Champs algébriques décents}
\label{stacks-morphisms-section-reasonable-decent}

\noindent
Cette section est l'analogue de
Espaces décents, section \ref{decent-spaces-section-reasonable-decent}.
En particulier, la définition suivante est compatible avec
la notion d'espace algébrique décent qui y est définie.

\begin{definition}
\label{stacks-morphisms-definition-decent}
Soit $\mathcal{X}$ un champ algébrique. Nous disons que $\mathcal{X}$ est
{\it décent} si, pour tout $x \in |\mathcal{X}|$, les conditions équivalentes
de Propriétés des champs, Lemme
\ref{stacks-properties-lemma-UR-quasi-compact-above-x} sont satisfaites.
\end{definition}

\noindent
Certains reformuleraient cette définition en disant que
tout point de $\mathcal{X}$ est quasi-compact.
Une condition un peu plus forte consisterait à demander que tout morphisme
$\Spec(k) \to \mathcal{X}$ appartenant à la classe d'équivalence de $x$ soit
quasi-séparé aussi bien que quasi-compact.

\begin{lemma}
\label{stacks-morphisms-lemma-quasi-separated-decent}
Un champ algébrique quasi-séparé $\mathcal{X}$ est décent.
Plus généralement, si $\Delta : \mathcal{X} \to \mathcal{X} \times \mathcal{X}$
est quasi-compact, alors $\mathcal{X}$ est décent.
\end{lemma}

\begin{proof}
En effet, si $\mathcal{X}$ est quasi-séparé, alors
tout morphisme $f : T \to \mathcal{X}$ dont la source est un schéma
quasi-compact $T$ est quasi-compact, voir
Lemme \ref{stacks-morphisms-lemma-quasi-compact-permanence}.
Si l'on sait que $\Delta$ est quasi-compact, on utilise la description
$$
T \times_{f, \mathcal{X}, f'} T' =
(T \times T')
\times_{(f, f'), \mathcal{X} \times \mathcal{X}, \Delta}
\mathcal{X}
$$
pour le démontrer. Nous omettons les détails.
\end{proof}

\begin{lemma}
\label{stacks-morphisms-lemma-representable-decent}
Soit $f : \mathcal{X} \to \mathcal{Y}$ un morphisme
de champs algébriques. Supposons que $\mathcal{Y}$ soit décent
et que $f$ soit représentable (par des schémas), ou que $f$ soit représentable
par des espaces algébriques et quasi-séparé.
Alors $\mathcal{X}$ est décent.
\end{lemma}

\begin{proof}
Soit $x \in |\mathcal{X}|$, d'image $y \in |\mathcal{Y}|$.
Choisissons un morphisme $y : \Spec(k) \to \mathcal{Y}$
dans la classe d'équivalence définissant $y$.
Posons $\mathcal{X}_y = \Spec(k) \times_{y, \mathcal{Y}} \mathcal{X}$.
Choisissons un point $x' \in |\mathcal{X}_y|$ d'image $x$, voir
Propriétés des champs, Lemme \ref{stacks-properties-lemma-points-cartesian}.
Choisissons un morphisme $x' : \Spec(k') \to \mathcal{X}_y$ dans la classe
d'équivalence de $x'$. Considérons le diagramme
$$
\xymatrix{
\Spec(k') \ar[r]_{x'} &
\mathcal{X}_y \ar[r] \ar[d] &
\mathcal{X} \ar[d] \\
&
\Spec(k) \ar[r]^y &
\mathcal{Y}
}
$$
Le morphisme $y$ est quasi-compact puisque $\mathcal{Y}$ est décent.
Par conséquent, $\mathcal{X}_y \to \mathcal{X}$ est quasi-compact, car il
s'obtient par changement de base
(Lemme \ref{stacks-morphisms-lemma-base-change-quasi-compact}).
Pour conclure, il suffit donc de montrer que $x'$ est quasi-compact
(Lemme \ref{stacks-morphisms-lemma-composition-quasi-compact}).
Si $f$ est représentable, alors $\mathcal{X}_y$ est un schéma
et $x'$ est quasi-compact.
Si $f$ est représentable par des espaces algébriques et quasi-séparé,
alors $\mathcal{X}_y$ est un espace algébrique quasi-séparé
et donc décent (Espaces décents, Lemme
\ref{decent-spaces-lemma-properties-trivial-implications}).
\end{proof}

\begin{lemma}
\label{stacks-morphisms-lemma-qc-decent}
Soit $f : \mathcal{X} \to \mathcal{Y}$ un morphisme de champs algébriques.
Si $f$ est quasi-compact et surjectif et si $\mathcal{X}$ est décent,
alors $\mathcal{Y}$ est décent.
\end{lemma}

\begin{proof}
Soit $x : \Spec(k) \to \mathcal{X}$ un morphisme, où $k$ est un corps,
et notons $y = f \circ x$. Puisque $f$ est surjectif, tout point
de $|\mathcal{Y}|$ s'obtient de cette manière, voir Propriétés des champs,
Lemme \ref{stacks-properties-lemma-characterize-surjective}.
Soient $T$ un schéma affine et $T \to \mathcal{Y}$ un morphisme.
Il suffit de montrer que
$T \times_{\mathcal{Y}, y} \Spec(k)$ est quasi-compact, voir
Lemme \ref{stacks-morphisms-lemma-check-quasi-compact-covering}.
Nous avons
$$
(T \times_{\mathcal{Y}} \mathcal{X}) \times_{\mathcal{X}, x} \Spec(k) =
T \times_{\mathcal{Y}, y} \Spec(k)
$$
Le morphisme $T \times_{\mathcal{Y}} \mathcal{X} \to T$ est quasi-compact ;
ainsi, $T \times_\mathcal{Y} \mathcal{X}$ est quasi-compact. Comme $x$ est
un morphisme quasi-compact puisque $\mathcal{X}$ est décent,
le produit fibré affiché est quasi-compact.
\end{proof}

\begin{lemma}
\label{stacks-morphisms-lemma-gerbe-decent}
Soit $f : \mathcal{X} \to \mathcal{Y}$ un morphisme de champs algébriques.
Si $\mathcal{X}$ est une gerbe sur $\mathcal{Y}$ et si $\mathcal{X}$ est décent,
alors $\mathcal{Y}$ est décent.
\end{lemma}

\begin{proof}
Supposons que $\mathcal{X}$ soit une gerbe sur $\mathcal{Y}$ et que
$\mathcal{X}$ soit décent. Remarquons que $f$ est un homéomorphisme universel
d'après le Lemme \ref{stacks-morphisms-lemma-gerbe-bijection-points}.
Le lemme résulte donc du Lemme \ref{stacks-morphisms-lemma-qc-decent}.
\end{proof}


\section{Points des champs alg\'ebriques d\'ecents}
\label{stacks-morphisms-section-points-decent}

\noindent
Cette section est l'analogue de la section
\ref{decent-spaces-section-points} d'Espaces d\'ecents.
Nous ne savons pas si l'espace topologique associ\'e
\`a un champ alg\'ebrique d\'ecent est toujours sobre ; voir
la Proposition \ref{stacks-morphisms-proposition-decent-sober} pour un r\'esultat
un peu plus faible.

\begin{lemma}
\label{stacks-morphisms-lemma-kolmogorov}
Soit $\mathcal{X}$ un champ alg\'ebrique d\'ecent.
Alors $|\mathcal{X}|$ est un espace de Kolmogoroff (voir
Topologie, D\'efinition \ref{topology-definition-generic-point}).
\end{lemma}

\begin{proof}
Soient $x_1, x_2 \in |\mathcal{X}|$ tels que $x_1 \leadsto x_2$ et
$x_2 \leadsto x_1$. Nous devons montrer que $x_1 = x_2$.
Soit $\mathcal{Z} \subset \mathcal{X}$ le sous-champ ferm\'e
r\'eduit tel que $|\mathcal{Z}|$ soit \`egal \`a
$\overline{\{x_1\}} = \overline{\{x_2\}}$.
D'apr\`es le Lemme \ref{stacks-morphisms-lemma-representable-decent},
le champ $\mathcal{Z}$ est d\'ecent.
En rempla\c{c}ant $\mathcal{X}$ par $\mathcal{Z}$, nous nous ramenons
au cas examin\'e dans le paragraphe suivant.

\medskip\noindent
Supposons $|\mathcal{X}|$ irr\'eductible, ayant $x_1$ et $x_2$
pour points g\'en\'eriques. Choisissons un sch\'ema affine $U$, des points $u_1, u_2 \in U$
et un morphisme lisse $f : U \to \mathcal{X}$ tels que
$f(u_i) = x_i$. Il existe alors un troisi\`eme point $u_3 \in U$
qui est le point g\'en\'erique d'une composante irr\'eductible de $U$
et dont l'image $x_3 \in |\mathcal{X}|$ est {\it elle aussi} un
point g\'en\'erique de $|\mathcal{X}|$. En effet, il suffit de
choisir pour $u_3$ un point g\'en\'erique quelconque d'une composante irr\'eductible
passant par $u_1$ (ou par $u_2$, si l'on pr\'ef\`ere).
Dans le paragraphe suivant, nous montrerons que $x_1 = x_3$ et $x_2 = x_3$,
ce qui donnera le r\'esultat voulu.

\medskip\noindent
Par sym\'etrie, il suffit de prouver que $x_1 = x_3$. Puisque $x_1$
est un point g\'en\'erique de $|\mathcal{X}|$, on a une sp\'ecialisation
$x_1 \leadsto x_3$.
D'apr\`es le Lemme \ref{stacks-morphisms-lemma-specialization}, il existe une sp\'ecialisation
$u'_1 \leadsto u_3$ dans $U$ (!) dont l'image est $x_1 \leadsto x_3$.
Or $u_3$ est le point g\'en\'erique d'une composante irr\'eductible ;
on a donc $u'_1 = u_3$, comme voulu.
\end{proof}

\begin{lemma}
\label{stacks-morphisms-lemma-decent-locally-noetherian-sober}
Soit $\mathcal{X}$ un champ alg\'ebrique d\'ecent et localement noeth\'erien.
Alors $|\mathcal{X}|$ est un espace topologique sobre et localement noeth\'erien.
\end{lemma}

\begin{proof}
D'apr\`es le Lemme \ref{stacks-morphisms-lemma-Noetherian-topology}, l'espace topologique $|\mathcal{X}|$
est localement noeth\'erien. D'apr\`es le Lemme \ref{stacks-morphisms-lemma-kolmogorov}, l'espace topologique
$|\mathcal{X}|$ est un espace de Kolmogoroff. D'apr\`es le
Lemme \ref{stacks-morphisms-lemma-locally-Noetherian-quasi-sober}, l'espace topologique
$|\mathcal{X}|$ est quasi-sobre. Cela ach\`eve la d\'emonstration ; voir
Topologie, D\'efinition \ref{topology-definition-generic-point}.
\end{proof}

\begin{proposition}
\label{stacks-morphisms-proposition-decent-sober}
Soit $\mathcal{X}$ un champ alg\'ebrique d\'ecent tel que
$\mathcal{I}_\mathcal{X} \to \mathcal{X}$ soit quasi-compact.
Alors $|\mathcal{X}|$ est sobre.
\end{proposition}

\begin{proof}
D'apr\`es le Lemme \ref{stacks-morphisms-lemma-kolmogorov}, $|\mathcal{X}|$ est un espace de Kolmogoroff
(en fait, nous allons le red\'emontrer).
Soit $T \subset |\mathcal{X}|$ une partie ferm\'ee irr\'eductible.
Nous devons montrer que $T$ poss\`ede un point g\'en\'erique.
Soit $\mathcal{Z} \subset \mathcal{X}$ le sous-champ ferm\'e induit
correspondant \`a $T$, muni de sa structure r\'eduite ; voir Propri\'et\'es des champs, D\'efinition
\ref{stacks-properties-definition-reduced-induced-stack}.
Puisque $\mathcal{Z} \to \mathcal{X}$ est une immersion ferm\'ee,
$\mathcal{Z}$ est un champ alg\'ebrique d\'ecent ; voir le
Lemme \ref{stacks-morphisms-lemma-representable-decent}.
De plus, le morphisme $\mathcal{I}_\mathcal{Z} \to \mathcal{Z}$
est le changement de base de $\mathcal{I}_\mathcal{X} \to \mathcal{X}$
(Lemme \ref{stacks-morphisms-lemma-monomorphism-cartesian-square-inertia}).
Par cons\'equent, $\mathcal{I}_\mathcal{Z} \to \mathcal{Z}$ est quasi-compact
(Lemme \ref{stacks-morphisms-lemma-base-change-quasi-compact}).
Nous nous ramenons ainsi au cas examin\'e dans le paragraphe suivant.

\medskip\noindent
Supposons que $\mathcal{X}$ soit d\'ecent, que
$\mathcal{I}_\mathcal{X} \to \mathcal{X}$ soit quasi-compact,
que $\mathcal{X}$ soit r\'eduit et que $|\mathcal{X}|$ soit irr\'eductible.
Nous devons montrer que $|\mathcal{X}|$ poss\`ede un point g\'en\'erique.
D'apr\`es la Proposition \ref{stacks-morphisms-proposition-open-stratum},
il existe un sous-champ ouvert dense $\mathcal{U} \subset \mathcal{X}$
qui est une gerbe. Autrement dit, $|\mathcal{U}| \subset |\mathcal{X}|$
est ouvert et dense.
Nous pouvons donc supposer en outre que $\mathcal{X}$ est une gerbe.
Supposons que $\mathcal{X} \to X$ munisse $\mathcal{X}$ d'une structure de gerbe sur
l'espace alg\'ebrique $X$. Alors $|\mathcal{X}| \cong |X|$ d'apr\`es le
Lemme \ref{stacks-morphisms-lemma-gerbe-bijection-points}.
En particulier, $X$ est quasi-compact et $|X|$ est irr\'eductible.
De plus, le Lemme \ref{stacks-morphisms-lemma-gerbe-decent} montre
que $X$ est un espace alg\'ebrique d\'ecent.
Alors $|\mathcal{X}| = |X|$ est sobre d'apr\`es
Espaces d\'ecents, Proposition \ref{decent-spaces-proposition-reasonable-sober},
et poss\`ede donc un point g\'en\'erique (unique).
\end{proof}




%
\section{Champs alg\'ebriques int\`egres}
\label{stacks-morphisms-section-integral-stacks}

\noindent
Cette section est l'analogue de Espaces sur des corps, section
\ref{spaces-over-fields-section-integral-spaces}.
Les consid\'erations de cette section et le r\'esultat de la
Proposition \ref{stacks-morphisms-proposition-decent-sober} nous am\`enent \`a d\'efinir
comme suit un champ alg\'ebrique int\`egre (et cette d\'efinition est compatible
avec les d\'efinitions d\'ej\`a existantes des sch\'emas int\`egres et
des espaces alg\'ebriques int\`egres).

\begin{definition}
\label{stacks-morphisms-definition-integral-algebraic-stack}
On dit qu'un champ alg\'ebrique $\mathcal{X}$ est
{\it int\`egre} s'il est r\'eduit et d\'ecent, si
$\mathcal{I}_\mathcal{X} \to \mathcal{X}$ est quasi-compact,
et si $|\mathcal{X}|$ est irr\'eductible.
\end{definition}

\noindent
Remarquons que si $\mathcal{X}$ est quasi-s\'epar\'e, il suffit,
pour qu'il soit int\`egre, que $\mathcal{X}$ soit r\'eduit et que
$|\mathcal{X}|$ soit irr\'eductible; voir le Lemme \ref{stacks-morphisms-lemma-totaro-case}.

\begin{lemma}
\label{stacks-morphisms-lemma-structure-integral-stack}
Soit $\mathcal{X}$ un champ alg\'ebrique int\`egre. Alors
\begin{enumerate}
\item $|\mathcal{X}|$ est sobre et irr\'eductible, et poss\`ede un unique point g\'en\'erique,
\item il existe un sous-champ ouvert $\mathcal{U} \subset \mathcal{X}$
qui est une gerbe sur un sch\'ema int\`egre $U$.
\end{enumerate}
\end{lemma}

\begin{proof}
La Proposition \ref{stacks-morphisms-proposition-decent-sober} montre que $|\mathcal{X}|$
est sobre. Il est bien entendu aussi irr\'eductible et poss\`ede donc un unique
point g\'en\'erique $x$ (par d\'efinition de la sobri\'et\'e).
La Proposition \ref{stacks-morphisms-proposition-open-stratum} montre
qu'il existe un sous-champ ouvert dense $\mathcal{U} \subset \mathcal{X}$
qui est une gerbe sur un espace alg\'ebrique $U$.
Alors $U$ est un espace alg\'ebrique d\'ecent d'apr\`es le
Lemme \ref{stacks-morphisms-lemma-gerbe-decent} (et puisque $\mathcal{U}$
est d\'ecent d'apr\`es le Lemme \ref{stacks-morphisms-lemma-representable-decent}).
Comme $|U| = |\mathcal{U}|$, on voit que $|U|$ est irr\'eductible.
Enfin, puisque $\mathcal{U}$ est r\'eduit, le morphisme
$\mathcal{U} \to U$ se factorise par $U_{red}$; voir
Propri\'et\'es des champs, Lemme \ref{stacks-properties-lemma-map-into-reduction}.
Or $\mathcal{U} \to U$ est plat, localement de pr\'esentation finie
et surjectif (Lemme \ref{stacks-morphisms-lemma-gerbe-fppf}); il s'ensuit que
$U = U_{red}$, c'est-\`a-dire que $U$ est r\'eduit (un petit d\'etail est omis).
Ainsi, $U$ est un espace alg\'ebrique int\`egre; voir
Espaces sur des corps, D\'efinition
\ref{spaces-over-fields-definition-integral-algebraic-space}.
Enfin, nous pouvons remplacer $U$ (et, corr\'elativement, $\mathcal{U}$)
par un sous-espace ouvert et supposer que $U$ est un sch\'ema int\`egre; voir
la discussion de Espaces sur des corps, section
\ref{spaces-over-fields-section-integral-spaces}.
\end{proof}

\begin{lemma}
\label{stacks-morphisms-lemma-totaro-case}
Soit $\mathcal{X}$ un champ alg\'ebrique r\'eduit et quasi-s\'epar\'e
dont l'espace topologique sous-jacent $|\mathcal{X}|$ est irr\'eductible.
Alors $\mathcal{X}$ est int\`egre.
\end{lemma}

\begin{proof}
Si $\mathcal{X}$ est quasi-s\'epar\'e, alors $\mathcal{X}$ est d\'ecent
d'apr\`es le Lemme \ref{stacks-morphisms-lemma-quasi-separated-decent}.
Si $\mathcal{X}$ est quasi-s\'epar\'e, alors
$\Delta : \mathcal{X} \to \mathcal{X} \times \mathcal{X}$ est quasi-compact;
par cons\'equent, $\mathcal{I}_\mathcal{X} \to \mathcal{X}$ est quasi-compact
comme changement de base de $\Delta$ par $\Delta$; voir le
Lemme \ref{stacks-morphisms-lemma-base-change-quasi-compact}.
On voit ainsi que toutes les hypoth\`eses de la
D\'efinition \ref{stacks-morphisms-definition-integral-algebraic-stack} sont satisfaites
(et aussi que l'on peut remplacer ``quasi-s\'epar\'e'' par
``$\Delta_\mathcal{X}$ est quasi-compact'').
\end{proof}

\begin{lemma}
\label{stacks-morphisms-lemma-decent-irreducible-closed}
Soit $\mathcal{X}$ un champ alg\'ebrique d\'ecent tel que
$\mathcal{I}_\mathcal{X} \to \mathcal{X}$ soit quasi-compact.
Il existe des bijections canoniques entre les ensembles suivants :
\begin{enumerate}
\item l'ensemble des points de $\mathcal{X}$, c'est-\`a-dire $|\mathcal{X}|$,
\item l'ensemble des parties ferm\'ees irr\'eductibles de $|\mathcal{X}|$,
\item l'ensemble des sous-champs ferm\'es int\`egres de $\mathcal{X}$.
\end{enumerate}
La bijection de (1) vers (2) envoie $x$ sur $\overline{\{x\}}$.
La bijection de (3) vers (2) envoie $\mathcal{Z}$ sur $|\mathcal{Z}|$.
\end{lemma}

\begin{proof}
La premi\`ere application d\'efinit une bijection entre (1) et (2), puisque
$|\mathcal{X}|$ est sobre d'apr\`es la Proposition \ref{stacks-morphisms-proposition-decent-sober}.
\'Etant donn\'ee une partie ferm\'ee irr\'eductible $T \subset |\mathcal{X}|$,
il existe un unique sous-champ ferm\'e r\'eduit $\mathcal{Z} \subset \mathcal{X}$
tel que $|\mathcal{Z}| = T$ : $\mathcal{Z}$ est pr\'ecis\'ement
la structure r\'eduite de sous-champ induite sur $T$; voir
Propri\'et\'es des champs, D\'efinition
\ref{stacks-properties-definition-reduced-induced-stack}.
Alors $\mathcal{Z}$ est un champ alg\'ebrique int\`egre : il est
d\'ecent (Lemme \ref{stacks-morphisms-lemma-representable-decent}), le
morphisme $\mathcal{I}_\mathcal{Z} \to \mathcal{Z}$ est quasi-compact
(comme changement de base de $\mathcal{I}_\mathcal{X} \to \mathcal{X}$; voir
le Lemme \ref{stacks-morphisms-lemma-monomorphism-cartesian-square-inertia}),
$\mathcal{Z}$ est r\'eduit, et $|\mathcal{Z}|$ est irr\'eductible.
\end{proof}






\section{Gerbes r\'esiduelles}
\label{stacks-morphisms-section-residual-gerbe}

\noindent
Cette section prolonge Propri\'et\'es des champs, section
\ref{stacks-properties-section-residual-gerbe}.

\begin{lemma}
\label{stacks-morphisms-lemma-gerbe-residual-gerbe}
Soit $\pi : \mathcal{X} \to \mathcal{Y}$ un morphisme de champs alg\'ebriques.
Soit $x \in |\mathcal{X}|$, d'image $y \in |\mathcal{Y}|$.
Supposons que la gerbe r\'esiduelle $\mathcal{Z}_y \subset \mathcal{Y}$
de $\mathcal{Y}$ en $y$ existe et que $\mathcal{X}$ soit une gerbe
sur $\mathcal{Y}$. Alors
$\mathcal{Z}_x = \mathcal{Z}_y \times_\mathcal{Y} \mathcal{X}$ est
la gerbe r\'esiduelle de $\mathcal{X}$ en $x$.
\end{lemma}

\begin{proof}
Le morphisme $\mathcal{Z}_x \to \mathcal{X}$ est un monomorphisme, car
il se d\'eduit par changement de base du monomorphisme $\mathcal{Z}_y \to \mathcal{Y}$.
Le morphisme $\pi$ est un hom\'eomorphisme universel d'apr\`es le
Lemme \ref{stacks-morphisms-lemma-gerbe-bijection-points}; par cons\'equent, $|\mathcal{Z}_x| = \{x\}$.
Enfin, le morphisme $\mathcal{Z}_x \to \mathcal{Z}_y$ est lisse,
car il se d\'eduit par changement de base du morphisme lisse $\pi$, voir le
Lemme \ref{stacks-morphisms-lemma-gerbe-smooth}.
Puisque $\mathcal{Z}_y$ est r\'eduit et localement noeth\'erien, il en
est donc de m\^eme de $\mathcal{Z}_x$ (d\'etails omis). Ainsi, $\mathcal{Z}_x$ est
la gerbe r\'esiduelle de $\mathcal{X}$ en $x$ d'apr\`es
Propri\'et\'es des champs, D\'efinition
\ref{stacks-properties-definition-residual-gerbe}.
\end{proof}

\begin{lemma}
\label{stacks-morphisms-lemma-factor-through-residual-space}
Soit $f : \mathcal{Y} \to \mathcal{X}$ un morphisme de champs alg\'ebriques.
Soit $x \in |\mathcal{X}|$ un point. Supposons que
\begin{enumerate}
\item $\mathcal{X}$ soit d\'ecent ou localement noeth\'erien (ou les deux),
\item $\mathcal{I}_\mathcal{X} \to \mathcal{X}$ soit quasi-compact,
\item $|f|(|\mathcal{Y}|)$ soit contenu dans $\{x\} \subset |\mathcal{X}|$, et
\item $\mathcal{Y}$ soit r\'eduit.
\end{enumerate}
Alors $f$ se factorise par la gerbe r\'esiduelle $\mathcal{Z}_x$
de $\mathcal{X}$ en $x$ (dont l'existence est assur\'ee par le
Lemme \ref{stacks-morphisms-lemma-every-point-residual-gerbe} ou le
Lemme \ref{stacks-morphisms-lemma-every-point-residual-gerbe-locally-Noetherian}).
\end{lemma}

\begin{proof}
Soit $T = \overline{\{x\}} \subset |\mathcal{X}|$ l'adh\'erence de $x$.
D'apr\`es Propri\'et\'es des champs, Lemme
\ref{stacks-properties-lemma-reduced-closed-substack},
il existe un sous-champ ferm\'e r\'eduit $\mathcal{X}' \subset \mathcal{X}$
tel que $T = |\mathcal{X}'|$. D'apr\`es Propri\'et\'es des champs, Lemme
\ref{stacks-properties-lemma-map-into-reduction}, le
morphisme $f$ se factorise par $\mathcal{X}'$. Si $\mathcal{X}$ est d\'ecent,
alors, d'apr\`es le Lemme \ref{stacks-morphisms-lemma-representable-decent}, le champ
$\mathcal{X}'$ est d\'ecent. Si $\mathcal{X}$ est localement noeth\'erien,
alors $\mathcal{X}'$ est localement noeth\'erien (d\'etails omis).
Le morphisme $\mathcal{I}_{\mathcal{X}'} \to \mathcal{X}'$
se d\'eduit par changement de base de $\mathcal{I}_\mathcal{X} \to \mathcal{X}$ d'apr\`es le
Lemme \ref{stacks-morphisms-lemma-monomorphism-cartesian-square-inertia};
le Lemme \ref{stacks-morphisms-lemma-base-change-quasi-compact} montre donc que
$\mathcal{I}_{\mathcal{X}'} \to \mathcal{X}'$ est quasi-compact.
On est ainsi ramen\'e au cas examin\'e dans le paragraphe suivant.

\medskip\noindent
Supposons que $\mathcal{X}$ soit r\'eduit et que $x \in |\mathcal{X}|$
soit un point g\'en\'erique. La Proposition \ref{stacks-morphisms-proposition-open-stratum}
entra\^ine l'existence d'un sous-champ ouvert dense
$\mathcal{U} \subset \mathcal{X}$ qui est une gerbe.
Notons que $x \in |\mathcal{U}|$.
En reprenant les arguments ci-dessus, on se ram\`ene au cas examin\'e
dans le paragraphe suivant.

\medskip\noindent
Supposons que $\mathcal{X} \to X$ soit une gerbe sur l'espace alg\'ebrique $X$.
Si $\mathcal{X}$ est d\'ecent, alors les
Lemmes \ref{stacks-morphisms-lemma-gerbe-bijection-points} et \ref{stacks-morphisms-lemma-qc-decent}
montrent que l'espace $X$ est d\'ecent. Si $\mathcal{X}$ est localement noeth\'erien,
alors $X$ est localement noeth\'erien par descente fppf (d\'etails omis).
Le r\'esultat correspondant vaut donc pour $X$, voir Espaces d\'ecents, Lemme
\ref{decent-spaces-lemma-factor-through-residual-space-decent} ou
\ref{decent-spaces-lemma-factor-through-residual-space-Noetherian}
(un petit d\'etail est omis).
En appliquant le Lemme \ref{stacks-morphisms-lemma-gerbe-residual-gerbe}, on conclut que le r\'esultat
vaut aussi pour $\mathcal{X}$.
\end{proof}

\begin{remark}
\label{stacks-morphisms-lemma-factor-through-residual-space-Noetherian}
Nous ignorons si le Lemme \ref{stacks-morphisms-lemma-factor-through-residual-space}
reste vrai lorsqu'on suppose seulement que $\mathcal{X}$ est localement noeth\'erien,
c'est-\`a-dire lorsqu'on supprime l'hypoth\`ese de quasi-compacit\'e de l'inertie. Dans ce cas,
si $x$ est un point ferm\'e, le r\'esultat est certainement vrai, comme le montre le
lemme beaucoup plus simple qui suit.
\end{remark}

\begin{lemma}
\label{stacks-morphisms-lemma-residual-space-closed}
Soit $\mathcal{X}$ un champ alg\'ebrique localement noeth\'erien.
Soit $x \in |\mathcal{X}|$, dont la gerbe r\'esiduelle est
$\mathcal{Z}_x \subset \mathcal{X}$
(Lemme \ref{stacks-morphisms-lemma-every-point-residual-gerbe-locally-Noetherian}).
Alors $x$ est un point ferm\'e de $|\mathcal{X}|$ si et seulement si
le morphisme $\mathcal{Z}_x \to \mathcal{X}$ est une immersion ferm\'ee.
\end{lemma}

\begin{proof}
Si $\mathcal{Z}_x \to \mathcal{X}$ est une immersion ferm\'ee, alors
$x$ est un point ferm\'e de $|\mathcal{X}|$, voir par exemple le
Lemme \ref{stacks-morphisms-lemma-closed-immersion-proper}.
R\'eciproquement, supposons que $x$ soit un point ferm\'e de $|\mathcal{X}|$.
Soit $\mathcal{Z} \subset \mathcal{X}$
le sous-champ ferm\'e r\'eduit tel que $|\mathcal{Z}| = \{x\}$
(Propri\'et\'es des champs,
Lemme \ref{stacks-properties-lemma-reduced-closed-substack}).
Alors $\mathcal{Z}$ est un champ alg\'ebrique localement noeth\'erien
d'apr\`es les Lemmes \ref{stacks-morphisms-lemma-immersion-locally-finite-type} et
\ref{stacks-morphisms-lemma-locally-finite-type-locally-noetherian}.
Comme $\mathcal{Z}$ est aussi r\'eduit et que $|\mathcal{Z}| = \{x\}$,
il s'ensuit que $\mathcal{Z} = \mathcal{Z}_x$ est, par d\'efinition, la
gerbe r\'esiduelle.
\end{proof}














\bibliography{my}
\bibliographystyle{amsalpha}
\end{document}
